% Compile with XeLaTeX (or LuaLaTeX), twice.
\documentclass[11pt,a4paper]{article}
\usepackage[margin=25mm,headheight=14pt]{geometry}
\usepackage{fontspec}
\setmainfont{Latin Modern Roman}
\setmonofont[Scale=.88]{Latin Modern Mono}
\usepackage{amsmath,amssymb,amsthm,mathtools}
\usepackage{unicode-math}
\setmathfont{Latin Modern Math}
\usepackage{microtype,booktabs,longtable,array,enumitem,xcolor,fancyhdr,xurl}
\usepackage[breaklinks,colorlinks,linkcolor=blue!45!black,urlcolor=blue!45!black,citecolor=blue!45!black]{hyperref}
\hypersetup{pdftitle={Analytic and arithmetic reductions toward Dirichlet GRH},pdfauthor={Independent mathematical reconstruction}}
\setlength{\parskip}{3pt}
\setlength{\emergencystretch}{2em}
\setlist{topsep=3pt,itemsep=3pt,parsep=0pt}
\pagestyle{fancy}
\fancyhf{}
\fancyhead[L]{\small Reductions toward Dirichlet GRH}
\fancyhead[R]{\small Working manuscript}
\fancyfoot[C]{\thepage}
\newtheorem{theorem}{Theorem}[section]
\newtheorem{proposition}[theorem]{Proposition}
\newtheorem{lemma}[theorem]{Lemma}
\newtheorem{corollary}[theorem]{Corollary}
\newtheorem{example}[theorem]{Example}
\theoremstyle{remark}
\newtheorem{remark}[theorem]{Remark}
\newcommand{\Z}{\mathbb Z}
\newcommand{\Q}{\mathbb Q}
\newcommand{\R}{\mathbb R}
\newcommand{\C}{\mathbb C}
\newcommand{\OO}{\mathcal O}
\DeclareMathOperator{\Rea}{Re}
\DeclareMathOperator{\rank}{rank}
\DeclareMathOperator{\ord}{ord}
\DeclareMathOperator{\im}{im}
\DeclareMathOperator{\length}{length}
\DeclareMathOperator{\Sha}{Sha}
\title{Analytic and arithmetic reductions\\toward Dirichlet GRH}
\author{A conventional mathematical reconstruction: working draft}
\date{21 November 2026}
\begin{document}
\maketitle
\begin{abstract}
We develop components of a conventional reconstruction of the announced
proof of RH and Dirichlet GRH. A complete one-sided Mellin--Laplace
argument reduces the conclusion to upper bounds for fixed positive
combinations of smoothed prime sums on arbitrarily large multiplicative
intervals. A uniform quasi-Riemann hypothesis gives a complete
simultaneous selection of the required finite family of quadratic
twists by weighted second moments. We reconstruct the cyclotomic
common-interval argument, retaining finite torsion, and prove its
coarse cohomological input by an independent algebraic $K$-theory
argument with the classical comparison theorems stated explicitly.
We verify its
compatibility with the selected primes. Integral polynomial
Weyl duality is proved directly by Gaussian elimination and divided
row operators. We also prove the quotient-metric divisor formula,
its finite-group dual exponent, homology torsion bounds from an
explicit arithmetic complex, and local Gaussian, stationary-phase,
convolution and unfolding formulas used by the theta estimates.
The arithmetic complex, including its exactness and orbit count,
is constructed explicitly. The full local orbital germ and
positive-majorant induction are also proved, with the specified
finite cokernels and self-dual measures retained.
A complete comparison with units and ideal classes, cyclotomic
enlargement, explicit finite-field homology functionals, and integral
arithmetic cycles now proves the coarse strict Selmer estimate by the
original cohomological route as well.
We compute sharp elementary-divisor and singular-value bounds for
the actual coefficient maps. A direct packing and integration argument
strengthens the global ideal sum, with its necessary logarithmic factor;
the exact period-lattice covolumes and the global Poisson formula
are proved with their hypotheses retained. We prove integral spanning
by initial minors, the compact branching law, and the exact integral
branching maps with their adjoint constants. These give the actual
lowering maps for every auxiliary character. Adelic class coordinates,
the extra cosets, and the initial class-period bound are reconstructed.
The archimedean coefficient identities and separated interpolation
sets are proved. The terminal radial density is positive, and its
finite orthogonal decomposition is constructed pointwise.
The exact hook product proves the required equality of the integral
coefficient lattices at the selected primes. Primitive unitary local
densities, including the dyadic case, and the source mass ratio are
verified from the explicitly stated classical global mass formula.
We identify the source's formal degree by a matrix-ball integral
and the classical discrete-series theorems, and prove its column
dimensions. The ordinary CM theta laws, including finite-combination
norms and exact torsion phases, are reconstructed with a common unit
period. The higher-dimensional moment measures and both matrix supports
are constructed, with the exact selected-prime denominator. The exact
neighboring-lattice incidence identity, its theta normalization, and
the weighted Hecke commutation law are proved.
The actual polynomial kernel has the required mixed-derivative
normalization; these results give the selected-prime column bound
on its prescribed initial lattice.
All four global differential lifts and the
regular spectral inverse are constructed with their exact radial
actions.
We prove the terminal quantum determinant, all terminal commutation
relations, and the exact eighth-power theta normalization.
A separate classical special-value argument proves the entire
opposite-parity Selmer estimate, including primes dividing the
field degree, with a stronger bound. A direct residual class-group
argument proves the exceptional trivial-character case in the
complementary parity. Restriction to fixed cyclotomic fields gives
the fixed-conductor case with its additive common-interval bound;
a bounded-degree argument proves the exponent-one estimate for
the bounded-order characters. An integral Fourier comparison transfers
the remaining character bounds to the required whole-field eigenspace,
with an explicit polynomial error at primes dividing the field degree.
The actual adjoint-vanishing implication follows from Thorne's classical
theorem, with its coefficient, polarization and local conditions checked.
One finite coefficient field contains all
theta jets required in each fixed construction. The lowest compact-type
integral supplies coefficient decay and the distant orbital sum.
The automorphic projection kernel, rare cusp-volume estimate, and
integrated nearby-unipotent bound are proved. An explicit Sobolev
estimate and iteration of zero constant terms control the deepest
cusp and complete the ambient multiplicity comparison.
The full-group majorant,
geometric integration formula, regularized Jacquet pairing and its
adjointness are proved with the actual derivative and support conditions.
A direct noncommutative calculation constructs the actual recursive
radial frame and proves its positive quadratic form on compactly
supported vectors, with every spectral domain specified.
The four group operators give rapid Fourier decay for the actual
Jacquet transforms. Compact lower-rank families establish pointwise
smooth marginal densities, local regularity of the generated slices,
and terminal compactness. A common positive Hermitian complement and
Schur's lemma prove degenerate off-diagonal vanishing under absolute
convergence. Exact Gram negative moments improve both orbital volume
leaves, and a sharper dyadic argument supplies the final summation.
A direct lattice-slope comparison and
several numerical bounds are sharpened. The quadratic Frobenius
determinant argument supplies real-zero and auxiliary-prime estimates;
exact prime removal gives a conditional improvement of the M\"obius
exponent with its required shorter mean square stated explicitly.

\textbf{This version does not prove GRH.} The central uniform arithmetic
construction remains incomplete: the auxiliary-level-prime theta
bounds, the noncompact integrated theta estimates, the remaining
cohomological components, and their final assembly are not proved in
this manuscript.
The conditional uses of the quasi-Riemann theorems retain their
hypotheses; no fixed zero-free region is promoted to the critical line.
All new deductions presented as results are proved below, with their
classical background inputs identified.
\end{abstract}
\setcounter{tocdepth}{1}
\hypersetup{bookmarksdepth=2}
\begingroup
\small
\setlength{\parskip}{0pt}
\tableofcontents
\endgroup
\newpage
\section{Purpose, notation, and the assertion still to be proved}
\label{sec:purpose}

This manuscript begins a conventional reconstruction of the announced proof
of the Riemann hypothesis and its extension to Dirichlet $L$-functions
\cite{alpo-announcement}.
It is a mathematical working paper, not a completed proof of those
hypotheses. The results below have proofs in ordinary mathematical
language. Where a result from the supplied papers is used as an input,
that input is stated explicitly. In particular, the arithmetic estimate
\eqref{eq:remaining-estimate} is not proved in this version. Calling it
an auxiliary lemma would not make the reconstruction complete.

The aim is to separate the analytic conclusion, which can be proved quite
concisely, from the arithmetic construction that must produce the required
estimate. This also identifies precisely what can be gained from the
earlier work on zero-free regions and from the quadratic determinant
arguments. The latter arguments are reconstructed below because their
divisibility and nonvanishing mechanisms are useful in their own right;
their conclusions do not supply \eqref{eq:remaining-estimate}.

We write $s=\sigma+it$, use natural logarithms, and put
$y_+=\max(y,0)$. The notation $A\ll B$ means $|A|\le CB$ for a
positive constant; subscripts specify permitted dependencies.
A Dirichlet character modulo $q$ is a homomorphism
$(\Z/q\Z)^\times\to\C^\times$, extended periodically to $\Z$ by
zero on integers not coprime to $q$. It is primitive if it is not induced
from a proper divisor of $q$. Its conductor is the least modulus from
which it is induced. The character of conductor one is denoted by $1$.
For $\sigma>1$,
\[
 L(s,\chi)=\sum_{n\ge1}\frac{\chi(n)}{n^s}
 =\prod_p(1-\chi(p)p^{-s})^{-1},\qquad L(s,1)=\zeta(s).
\]
The von Mangoldt function is $\Lambda(p^j)=\log p$ for primes $p$
and integers $j\ge1$, and $\Lambda(n)=0$ otherwise.

We use the classical analytic continuation and functional equations of
these functions, and their nonvanishing on $\Re s=1$ away from the
pole of $\zeta$. In particular $L(1,\chi)\ne0$ for nonprincipal
$\chi$. For primitive $\chi$ of conductor $q>1$, with
$\chi(-1)=(-1)^a$, the completed function
\[
 \Lambda_\chi(s)=(q/\pi)^{(s+a)/2}
       \Gamma((s+a)/2)L(s,\chi)
\]
is entire and satisfies
$\Lambda_\chi(s)=w_\chi\Lambda_{\bar\chi}(1-s)$ with
$|w_\chi|=1$. Also
$\zeta(s)=(s-1)^{-1}+\gamma+O(s-1)$ at $s=1$.
These classical facts are background inputs; the arguments below do not
claim to reprove the analytic theory of $L$-functions from its
foundations. References and normalizations are given in
\cite{dlmf-dirichlet,dlmf-zeta,dlmf-zeros}.

For $x\ge1$ define the smoothed sums
\begin{align}
 S_\chi(x)&=\sum_{n\le x}\Lambda(n)\chi(n)(x/n-1),\label{eq:Schi}\\
 G_\chi(x)&=S_\chi(x)+x\frac{L'(1,\chi)}{L(1,\chi)}
                  &&(\chi\ne1),\label{eq:Gchi}\\
 G_1(x)&=\sum_{n\le x}\Lambda(n)(x/n-1)
          -x\log x+(1+\gamma)x.\label{eq:Gone}
\end{align}
Only primitive characters will be used in these definitions. In
particular, \eqref{eq:Gone} refers to conductor one, not to the
principal character at an arbitrary larger modulus.

Fix a primitive character $\chi$ modulo $q$, let $o$ be its order, and
let $\mathcal C_\chi$ be the set of distinct primitive characters
inducing the powers $\chi^j$, $0\le j<o$. Put
\begin{equation}\label{eq:Hchi}
 H_\chi(x)=G_1(x)+\frac1{3o+1}
       \sum_{\substack{\psi\in\mathcal C_\chi\\\psi\ne1}}
               \Re G_\psi(x).
\end{equation}
The set and its weights are fixed once $\chi$ is fixed; they do not
change with $x$ or with an error exponent.

Here is a precise sufficient arithmetic assertion. For each primitive
$\chi$, there is a fixed $b>1$ such that, for every $\varepsilon>0$,
there are arbitrarily large $Y$ for which
\begin{equation}\label{eq:remaining-estimate}
             H_\chi(x)\le x^{1/2+\varepsilon}
                    \qquad(Y\le x\le Y^b).
\end{equation}
The proof below shows that this assertion implies Dirichlet GRH.
Its proof from the proposed automorphic and arithmetic constructions
remains the principal unfinished part of the present manuscript.
The one-sided inequality and the quantifier ``arbitrarily large'' are
intentional: an absolute-value bound at every sufficiently large $x$
is unnecessary.

\section{Smoothed sums and their Mellin transforms}
\label{sec:mellin}

\begin{lemma}\label{lem:mellin}
For $\Re s>1$ one has
\begin{align}
 \int_1^\infty G_\chi(x)x^{-s-1}\,dx
  &=-\frac{L'(s,\chi)/L(s,\chi)
                  -sL'(1,\chi)/L(1,\chi)}{s(s-1)}
                       &&(\chi\ne1),\label{eq:mellin-chi}\\
 \int_1^\infty G_1(x)x^{-s-1}\,dx
  &=-\frac{\zeta'(s)/\zeta(s)}{s(s-1)}
      -\frac1{(s-1)^2}+\frac{1+\gamma}{s-1}.
                       \label{eq:mellin-zeta}
\end{align}
The expressions on the right extend meromorphically to $\Re s>0$,
are holomorphic at $s=1$, and have, at a zero $\rho$ of the
corresponding $L$-function in $0<\Re\rho<1$, residue
\begin{equation}\label{eq:residue}
                    -\frac{m_\chi(\rho)}{\rho(\rho-1)},
\end{equation}
where $m_\chi(\rho)$ is the positive integer multiplicity of that zero.
All their poles in $\Re s>0$ are of this form and are simple.
\end{lemma}

\begin{proof}
For $n\ge1$ and $\sigma>1$, direct integration gives
\[
 \int_1^\infty(x/n-1)_+x^{-s-1}\,dx
 =\frac{n^{-s}}{s(s-1)}.
\]
The same calculation with $\sigma$ in place of $s$ bounds the integral
of the absolute value by $n^{-\sigma}/(\sigma(\sigma-1))$.
Consequently the exchange of sum and integral is justified by
$\sum\Lambda(n)n^{-\sigma}<\infty$. Logarithmic differentiation of
the absolutely convergent Euler product gives
\[
       \sum_{n\ge1}\Lambda(n)\chi(n)n^{-s}
                         =-\frac{L'(s,\chi)}{L(s,\chi)}.
\]
Finally,
$\int_1^\infty x^{-s}\,dx=(s-1)^{-1}$ and
$\int_1^\infty x^{-s}\log x\,dx=(s-1)^{-2}$.
These identities prove both formulas, including absolute convergence
of the integrals for $\sigma>1$.

In \eqref{eq:mellin-chi} the numerator vanishes at $s=1$; its
apparent singularity there is removable. For \eqref{eq:mellin-zeta},
use $\zeta'/\zeta=-1/(s-1)+\gamma+O(s-1)$.
The coefficients of both $(s-1)^{-2}$ and $(s-1)^{-1}$ cancel.
If $L(s,\chi)=(s-\rho)^m h(s)$ with $h$ holomorphic and
$h(\rho)\ne0$, then $L'/L=m/(s-\rho)+h'/h$.
Since $\rho(\rho-1)\ne0$, this proves \eqref{eq:residue}.
There are no other singularities in the indicated half-plane, by
the classical analytic properties stated above. The zeros of a
nonzero holomorphic function are isolated, so the pole set is locally
finite.
\end{proof}


\begin{lemma}\label{lem:positive-family}
Let $\chi_1,\ldots,\chi_r$ be primitive characters, possibly including
the conductor-one character, and let $c_j>0$ be real numbers.
Set $H(x)=\sum_jc_j\Re G_{\chi_j}(x)$ and $g(u)=H(e^u)$.
Then $|g(u)|\ll e^{2u}$ for $u\ge0$. Its Laplace transform extends
meromorphically to $\Re s>0$. At every zero $\rho$ of one of the
$L(s,\chi_j)$ in that half-plane the transform has a genuine simple
pole, with residue
\begin{equation}\label{eq:family-residue}
 -\frac{1}{\rho(\rho-1)}
    \sum_{j=1}^r\frac{c_j}{2}
                 \bigl(m_{\chi_j}(\rho)+m_{\bar\chi_j}(\rho)\bigr).
\end{equation}
Here a multiplicity is zero if the corresponding function does not
vanish. The residue at a real pole is positive.
\end{lemma}

\begin{proof}
The bound $\Lambda(n)\le\log n$ gives
$|S_\chi(x)|\le x\sum_{2\le n\le x}(\log n)/n
\ll x(1+\log^2x)\ll x^2$. The correction terms in
\eqref{eq:Gchi}--\eqref{eq:Gone} obey the same bound for each fixed
character. Thus $g$ is continuous and has the stated growth.

Conjugating the absolutely convergent Dirichlet series, and then
continuing analytically, gives
$L(s,\bar\chi)=\overline{L(\bar s,\chi)}$ and
$G_{\bar\chi}(x)=\overline{G_\chi(x)}$.
Hence $\Re G_\chi=(G_\chi+G_{\bar\chi})/2$.
The substitution $x=e^u$ identifies each Mellin integral in
Lemma~\ref{lem:mellin} with the corresponding Laplace integral.
Formula~\eqref{eq:family-residue} follows by addition. Its integer
multiplicities are nonnegative, and at least one is positive at a
zero under consideration. Thus the residue cannot cancel.
For real $0<\rho<1$, the factor $-1/(\rho(\rho-1))$ is positive.
\end{proof}

\section{A one-sided criterion on multiplicative intervals}
\label{sec:one-sided}

Only elementary complex analysis is needed in the following argument.
The proof is included because replacing a one-sided estimate by an
unjustified absolute-value estimate would invalidate the reduction.


\begin{proposition}\label{prop:landau-independent}
Let $g:[0,\infty)\to\R$ be measurable and satisfy
$|g(u)|\le C_0e^{2u}$. Suppose that its Laplace transform
\[
 \Phi(s)=\int_0^\infty g(u)e^{-su}\,du\qquad(\Rea s>2)
\]
extends meromorphically to $\Rea s>0$, with only isolated simple poles,
and that every real pole has positive real residue.
Fix $0<\theta<1$ and $\delta>0$. Suppose that for every $\eta>0$
there are arbitrarily large $U$ and a nonnegative function $R_U\ge g$
on $[U,(1+\delta)U]$ such that $R_U(u)e^{-\theta u}$ is integrable there and
\[
 \int_U^{(1+\delta)U}R_U(u)e^{-\theta u}\,du\le e^{\eta U}.
\]
Then $\Phi$ has no pole with real part greater than $\theta$.
\end{proposition}

\begin{proof}
We give the tail estimate and the pole argument explicitly.

Put $a=1+\delta/2$ and $b=1+\delta$. Both numbers
\[
 q_- = \frac{\exp(1-1/a)}{a},\qquad
 q_+ = \frac ba\exp(1-b/a)
\]
are strictly less than $1$. For the first, use
$\log a>1-1/a$ for $a>1$; for the second, use
$\log r<r-1$ for $r=b/a>1$. Choose $\lambda>0$ so large that
\[
 \lambda a\ge1,\qquad
 q:=\max(q_-,q_+)<\frac{\lambda}{\lambda+2},
 \qquad \sigma:=\lambda+2>2.
\]
For $U\ge1$, set $k=\lceil\lambda aU\rceil$, so that $k\ge U$.
Elementary exponential tilting and
$\int_0^\infty u^ke^{-tu}\,du=k!t^{-k-1}$ give
\begin{align*}
 \int_0^U\frac{u^ke^{-\lambda u}}{k!}\,du
 &\le \frac{e^{\lambda(a-1)U}}{(\lambda a)^{k+1}}
 \le \frac1{\lambda a}\left(\frac{q_-}{\lambda}\right)^k,\\
 \int_{bU}^\infty\frac{u^ke^{-\lambda u}}{k!}\,du
 &\le \frac{e^{-\lambda(b-a)U}}{(\lambda a/b)^{k+1}}
 \le \frac{b\exp((b-a)/a)}{\lambda a}
          \left(\frac{q_+}{\lambda}\right)^k.
\end{align*}
For the first inequality, write the integrand as
$e^{\lambda(a-1)u}u^ke^{-\lambda au}/k!$ and bound the first factor
by its value at $U$. For the second, write it as
$e^{-\lambda(1-a/b)u}u^ke^{-(\lambda a/b)u}/k!$ and bound the first
factor by its value at $bU$. The final bounds use
$U\le k/(\lambda a)$ and $U\ge(k-1)/(\lambda a)$, respectively.
Thus, for a constant $C_1$ depending only on the chosen $\delta,\lambda$,
\begin{equation}\label{eq:landau-tail-independent}
 \int_{[0,\infty)\setminus[U,bU]}
      \frac{u^ke^{-(\sigma-2)u}}{k!}\,du
 \le C_1\sigma^{-k}.
\end{equation}

For $\Rea s>2$, differentiation under the integral is justified by
the exponential bound on $g$. Write
\[
 A_k(s)=\frac{(-1)^k}{k!}\Phi^{(k)}(s)
       =\int_0^\infty\frac{g(u)u^ke^{-su}}{k!}\,du.
\]
The triangle inequality gives, for every real $\gamma$,
\[
 |A_k(\sigma+i\gamma)|+\Rea A_k(\sigma)
 \le\int_0^\infty (|g(u)|+g(u))
               \frac{u^ke^{-\sigma u}}{k!}\,du.
\]
On $[U,bU]$, use $|g|+g\le2R_U$ and
\[
 \frac{u^ke^{-(\sigma-\theta)u}}{k!}
 \le(\sigma-\theta)^{-k},
\]
which follows from the power series for the exponential.
Off this interval, use $|g|+g\le2C_0e^{2u}$ and
\eqref{eq:landau-tail-independent}. Since $\theta>0$, there is a constant
$C_2>0$, independent of $U,k,\eta,\gamma$, such that, for every
$\eta>0$ and arbitrarily large suitable $U$,
\begin{equation}\label{eq:landau-coeff-independent}
 |A_k(\sigma+i\gamma)|+\Rea A_k(\sigma)
 \le C_2\frac{e^{\eta U}}{(\sigma-\theta)^k}
 \quad\text{for all }\gamma\in\R,
 \qquad k\ge U.
\end{equation}

Suppose, for a contradiction, that a pole $\rho_0$ has
$\Rea\rho_0>\theta$. All poles have real part at most $2$, since the
original Laplace integral is holomorphic on $\Rea s>2$.
Let $Z$ be the pole set. Every compact subset of $\Rea s>0$ meets $Z$
in a finite set. If the closed disk with centre $\sigma$ and radius
$\sigma-\theta$ contains a pole, let $D$ be the distance to its nearest
pole; otherwise put $D=\sigma-\theta$. In either case
\[
 0<D\le\sigma-\theta,
 \qquad B(\sigma,D)\cap Z=\varnothing.
\]

First suppose there is a pole $\rho_1$ with
$\sigma-\Rea\rho_1<D$. There is then a real $\gamma$ for which
$z=\sigma+i\gamma$ has a unique nearest pole $\rho$ at distance $d<D$.
Here is the elementary justification. Start with
$\gamma=\Im\rho_1$ and vary it in a sufficiently small interval so
that $\rho_1$ remains at distance less than $D$. Only finitely many
poles can enter the disks of radius $D$ under this variation. For two
distinct such poles, equality of their distances from
$\sigma+i\gamma$ occurs at at most one value of $\gamma$: squared
distances have a difference affine in $\gamma$, and if their imaginary
parts coincide their different real parts, both less than $\sigma$,
give unequal distances for every $\gamma$. Avoid these finitely many
exceptional values. Choose $r$ with $d<r<D$ such that $B(z,r)$ contains
only $\rho$, and choose $r_1$ with $d<r_1<r$.

If $c\ne0$ is the residue at $\rho$, subtracting $c/(s-\rho)$ makes
$\Phi$ holomorphic on $B(z,r)$ after filling the removable singularity.
Cauchy's estimates consequently give constants $M,M'>0$ such that
\[
 \left|A_k(z)-\frac{c}{(z-\rho)^{k+1}}\right|\le Mr_1^{-k},
 \qquad |A_k(\sigma)|\le M'r_1^{-k}
 \quad(k\ge0).
\]
The second estimate uses the pole-free disk $B(\sigma,D)$.
Combining these inequalities with~\eqref{eq:landau-coeff-independent}
gives
\begin{equation}\label{eq:landau-final-independent}
 \frac{|c|}{d}\,d^{-k}
 \le (M+M')r_1^{-k}
       +C_2e^{\eta U}(\sigma-\theta)^{-k}.
\end{equation}

It remains to consider the case in which
$D\le\sigma-\Rea\rho$ for every pole $\rho$.
The pole-free closed-disk alternative is impossible because of
$\rho_0$. A nearest pole $\rho$ must satisfy
\[
 D\le\sigma-\Rea\rho\le|\sigma-\rho|=D.
\]
Hence it is the real pole $\beta=\sigma-D$, and it is the unique nearest
pole. Moreover $\beta>\theta$, since
$D\le\sigma-\Rea\rho_0<\sigma-\theta$.
Choose $D<r_1<r\le\sigma-\theta$ so that $B(\sigma,r)$ contains no
other pole. Its positive residue $c$ gives
\[
 \Rea A_k(\sigma)\ge cD^{-k-1}-Mr_1^{-k}
\]
by the same Cauchy estimate. Equation~\eqref{eq:landau-coeff-independent},
with $\gamma=0$, gives~\eqref{eq:landau-final-independent} with
$d=D$, $|c|=c$, and $M'=0$.

In either case $0<d<r_1\le\sigma-\theta$. Fix
$0<\eta<\log(r_1/d)$. Since $k\ge U$, multiplying
\eqref{eq:landau-final-independent} by $d^k$ shows
\[
 0<\frac{|c|}{d}
 \le (M+M'+C_2)
      \exp\bigl((\eta-\log(r_1/d))U\bigr).
\]
The right-hand side tends to zero along the arbitrarily large admissible
$U$, which is a contradiction.
\end{proof}


\begin{theorem}\label{thm:interval-grh}
Suppose that \eqref{eq:remaining-estimate} holds for every primitive
Dirichlet character, with $b>1$ fixed for each character. Then every
zero of every Dirichlet $L$-function in $0<\Re s<1$ has real part
$1/2$.
\end{theorem}

\begin{proof}
Fix $\chi$ and put $g(u)=H_\chi(e^u)$. Lemma~\ref{lem:positive-family}
supplies a meromorphic Laplace transform, with positive residues at
all real poles. Given $\eta>0$, choose
$0<\varepsilon<\eta/(2b)$. On any sufficiently large interval
$[Y,Y^b]$ in \eqref{eq:remaining-estimate}, set $U=\log Y$ and
$R_U(u)=e^{(1/2+\varepsilon)u}$. Then
\[
 \int_U^{bU}R_U(u)e^{-u/2}\,du
    =\frac{e^{b\varepsilon U}-e^{\varepsilon U}}{\varepsilon}
    \le \varepsilon^{-1}e^{b\varepsilon U}
    \le e^{\eta U}
\]
for all sufficiently large such $U$. Proposition~\ref{prop:landau-independent},
with $\theta=1/2$ and $\delta=b-1$, excludes every pole with
real part greater than $1/2$. Since $\chi\in\mathcal C_\chi$ and its
coefficient is positive, Lemma~\ref{lem:positive-family} excludes
every zero of $L(s,\chi)$ in that half-plane. The same argument
includes $\zeta$.

For a primitive character, a zero with $0<\Re s<1/2$ would give a
zero of $L(1-s,\bar\chi)$ with real part greater than $1/2$, by the
functional equation. The gamma factors are finite and nonzero
throughout the open critical strip. Thus every primitive zero there
has real part $1/2$.

Finally, if a character $\xi$ modulo $N$ is induced by $\xi^*$ of
conductor $f$, comparison of Euler factors, followed by analytic
continuation, gives
\[
 L(s,\xi)=L(s,\xi^*)
       \prod_{\substack{p\mid N\\p\nmid f}}
                    (1-\xi^*(p)p^{-s}).
\]
When $\Re s>0$, each factor is nonzero since
$|\xi^*(p)p^{-s}|=p^{-\Re s}<1$. Hence the result extends to all
characters. The usual Riemann hypothesis follows by taking conductor
one and using the classical location of the trivial zeta zeros.
\end{proof}

\begin{remark}[Essential quantifiers and signs]
The characters and positive coefficients in $H_\chi$ must be fixed
while $Y$ tends to infinity. One cannot replace them at each $Y$ by
an unrelated family and apply the argument to a single Laplace
transform. Similarly, $b>1$ must be fixed. The constants and the
starting point of the admissible intervals may depend on
$\varepsilon$ and on $\chi$.

Two counterexamples delimit the analytic theorem. If
$g(u)=-e^{\beta u}$ with $\beta>\theta$, the majorant $R=0$ satisfies
every one-sided estimate but the transform has residue $-1$ at
$\beta$; thus positivity at real poles cannot be omitted.
If $g(u)=e^{\beta u}\cos u$, there are arbitrarily late negative
intervals of fixed length, although its transform has poles at
$\beta\pm i$. Such intervals cannot replace $[U,bU]$ with fixed
$b>1$. These are counterexamples to weaker criteria, not to the
criterion just proved.
\end{remark}

\section{The numerical balance required of the arithmetic construction}
\label{sec:numerical}

The following scalar implication records why a number of stages of
order $\sqrt{x}$ is appropriate. It is useful independently of the
particular construction of the parameters. All its inequalities are
assumptions; proving that the required arithmetic objects supply
them is a separate task.



\begin{proposition}\label{prop:size}
Let $x,\varepsilon,C,n,n_T,n_2,m,b,k,A,E,G$ be real numbers satisfying
\[
 x\ge3,\qquad \varepsilon>0,\qquad C\ge0,\qquad n>0,\qquad b>0,
 \qquad 0\le A\le2xb,
\]
\[
 n_T\ge\frac n2x^{-\varepsilon/8},\qquad
 x^{-\varepsilon}\le\frac{n_2}{n}\le1,\qquad
 x^{1/2-\varepsilon}\le m\le x^{1/2+\varepsilon/2},
\]
\[
 xm\le k\le n^2x^{1/2-\varepsilon},\qquad
 x^{3\varepsilon/8}\ge26.
\]
Suppose also that
\begin{align*}
E\le x^{C\varepsilon}\biggl[
&n_TA\left(1+\frac{m^2}{k}\right)\\
&+m\left(2nmbx^\varepsilon+
 \frac{n_2kbx^\varepsilon}{n^2x^{-\varepsilon}}\right)\biggr]
\end{align*}
and
\[
 n_Tmb\frac{n_2}{2n}\bigl(G-4x^{1/2+\varepsilon}\bigr)
       \le E+\frac{2n_Tmb}{x}.
\]
Then $G\le x^{1/2+(C+3)\varepsilon}$.
\end{proposition}
\begin{proof}
Set $Q=n_Tmb>0$, $d=n_2/n>0$, and $Y=x^{1/2+13\varepsilon/8}$. Dividing the expression in square brackets by $Q$ gives exactly
\[
 \frac{A}{mb}+\frac{Am}{kb}
 +\frac{2nmx^\varepsilon}{n_T}
 +\frac{n_2kx^{2\varepsilon}}{n^2n_T}.
\]
The respective bounds are
\begin{align*}
 \frac{A}{mb}&\le\frac{2x}{m}\le2x^{1/2+\varepsilon},\\
 \frac{Am}{kb}&\le\frac{2xm}{k}\le2,\\
 \frac{2nmx^\varepsilon}{n_T}&\le4mx^{9\varepsilon/8}
                  \le4x^{1/2+13\varepsilon/8},\\
 \frac{n_2kx^{2\varepsilon}}{n^2n_T}
       &\le\frac{2k}{n^2}x^{17\varepsilon/8}
                  \le2x^{1/2+9\varepsilon/8}.
\end{align*}
Each displayed power is at most $Y$, and $Y\ge1$. Hence $E/Q\le10x^{C\varepsilon}Y$. The second assumed inequality implies
\[
 d\bigl(G-4x^{1/2+\varepsilon}\bigr)
       \le 2E/Q+4/x
       \le22x^{C\varepsilon}Y,
\]
because $4/x\le4/3\le2x^{C\varepsilon}Y$. Dividing by $d$ and using $d^{-1}\le x^\varepsilon$ gives
\begin{align*}
 G&\le4x^{1/2+\varepsilon}
           +22x^{1/2+(C+21/8)\varepsilon}\\
  &\le26x^{1/2+(C+21/8)\varepsilon}
   \le x^{1/2+(C+3)\varepsilon}.
\end{align*}
The last step is precisely $26\le x^{3\varepsilon/8}$.
\end{proof}


\begin{corollary}\label{cor:scalar-construction}
Fix $b_0>1$ and a real function $H$ of the form in
Lemma~\ref{lem:positive-family}. Suppose there is $C\ge0$ such that
for every sufficiently small $\varepsilon>0$ and arbitrarily large
$Y$, the hypotheses of Proposition~\ref{prop:size} can be satisfied
at every $x\in[Y,Y^{b_0}]$, with $G=H(x)$ and with that same $C$.
Then none of the $L$-functions occurring in $H$ has a zero with
real part greater than $1/2$.
\end{corollary}

\begin{proof}
For any $\alpha>0$, choose $\varepsilon>0$ in the permitted range
with $(C+3)\varepsilon<\alpha$. Proposition~\ref{prop:size} gives
$H(x)\le x^{1/2+\alpha}$ on the corresponding arbitrarily large
intervals. The proof of Theorem~\ref{thm:interval-grh} applies to
this fixed positive character combination.
\end{proof}

This reduction does not permit $C$ to grow like
$\varepsilon^{-1}$: the exponent $(C+3)\varepsilon$ must tend to
zero. Nor does it permit the main height parameter $A$ to be replaced
by a bound for a different arithmetic function. The two leading
contributions after division by $n_Tmb$ are of sizes $x/m$ and $m$.
For positive $m$, $x/m+m\ge2\sqrt{x}$, with equality at
$m=\sqrt{x}$. This elementary optimization explains the target
scale; it does not construct the theta lifts or their integral
lattices.

In the submitted construction the analogue of $A$ is a height or
weight parameter, compatible with weights $a_t$ satisfying
$|a_t/b-x|\le1$. A bound for a M\"obius sum of size
$x^{7/8+\varepsilon}$ is therefore not a bound
$A\ll b x^{7/8+\varepsilon}$. Making that substitution would be
an error. Likewise, estimates at exponent $7/8$ do not provide the
$x^{1/2+\varepsilon}$ bounds required for the three specially
chosen quadratic-character sums. Those character choices and
estimates remain part of the arithmetic problem.



% Mathematical contribution. No result in this file assumes the conclusion GRH.
% Dependencies: amsmath, amssymb, amsthm; theorem/proposition/lemma/corollary/remark environments.
\section{What the quasi-Riemann hypothesis can contribute}
\label{sec:qrh-contribution}

The two supplied quasi-Riemann-hypothesis papers have different endpoints.\cite{qrh-september,qrh-october}
The paper of 30 September 2026 states zero exclusion in
$\operatorname{Re}s>7/8$ for every finite-order Hecke character of
$\mathbb Q(\sqrt{-3})$, and hence for every Dirichlet character.
The shorter paper of 5 October states the corresponding $11/12$ result.
The latter reduces its conclusion to a particularly concrete mean-square
estimate (its Proposition~3.1). The reductions below can be checked
independently of the substantial cubic-theta calculation that supplies
that estimate. We do not treat that calculation as an elementary lemma
or silently include it among the hypotheses already proved here.

For use below, write $\mathcal Z(\vartheta)$ for the following precise
assertion, where $1/2\leq\vartheta<1$:
\[
 L(s,\chi)\ne0\quad(\operatorname{Re}s>\vartheta)
\]
for every Dirichlet character $\chi$, with the pole of a principal
character at $1$ allowed. A statement about each character with its
own unspecified boundary is weaker than $\mathcal Z(\vartheta)$.

\subsection{A direct Mellin-transform criterion}

For a number field $K$, a finite-order Hecke character may be
viewed as a character of a ray class group: it is a homomorphism from
the group of fractional ideals prime to a fixed integral modulus
$\mathfrak f$ to the roots of unity, trivial on the principal ideals
$(\alpha)$ with $\alpha\equiv1\pmod{\mathfrak f}$ and with
$\alpha$ positive at the real places prescribed by the modulus.
Let $\nu$ be such a character, and let $S$ be a finite set of prime
ideals containing the primes dividing its conductor. The notation
$(\mathfrak a,S)=1$ means that no member of $S$ divides
$\mathfrak a$. Let $\mu_K$ be the ideal M\"obius function: it is zero on ideals
divisible by a prime square and equals $(-1)^r$ on a product of $r$
distinct primes. The norm of an integral ideal $\mathfrak a$ is
$N\mathfrak a=|\mathcal O_K/\mathfrak a|$. Put
\[
 L^S(s,\nu)=\prod_{\mathfrak p\notin S}
       (1-\nu(\mathfrak p)N\mathfrak p^{-s})^{-1}
       \qquad(\operatorname{Re}s>1).
\]
Assume that the number of ideals of norm at most $t$ is $O_K(t)$,
and that $L^S$ has its usual meromorphic continuation to
$\operatorname{Re}s>0$. These standard analytic inputs will be kept
explicit in this subsection.

\begin{proposition}[Compactly supported smoothing]
\label{prop:qrh-mellin}
Fix $0<\alpha<1$. Suppose that, for every
$W\in C_c^\infty(0,\infty)$ and every $\epsilon>0$,
\[
 A_W(X):=\sum_{(\mathfrak a,S)=1}
     \mu_K(\mathfrak a)\nu(\mathfrak a)
     W(N\mathfrak a/X)
       =O_{K,\nu,S,W,\epsilon}(X^{\alpha+\epsilon}).
\]
Then $L(s,\nu)$ has no zero in $\operatorname{Re}s>\alpha$.
No uniformity in $\nu$ of the implied constant is required.
\end{proposition}
\begin{proof}
If $\rho$ were such a zero, choose a nonzero nonnegative function
$\phi\in C_c^\infty(1,2)$ and set
$W(y)=y^{-\rho}\phi(y)$, where $y^{-\rho}=\exp(-\rho\log y)$.
The Mellin transform
\[
 \widehat W(s)=\int_0^\infty W(y)y^s\,\frac{dy}{y}
\]
is entire, and $\widehat W(\rho)=\int\phi(y)\,dy/y>0$.
On each compact subset of $\operatorname{Re}s>\alpha$, choose
$\epsilon$ smaller than the distance of that compact set from the
line $\operatorname{Re}s=\alpha$. The assumed bound, including the
same bound after multiplying the integrand by any fixed power of
$\log X$, proves local uniform convergence and holomorphy of
\[
 M_W(s)=\int_0^\infty A_W(X)X^{-s}\,\frac{dX}{X}.
\]
There is no issue at zero: compact support makes $A_W(X)=0$ for all
sufficiently small $X$. For $\operatorname{Re}s>1$, absolute
convergence permits termwise integration and the substitution
$y=N\mathfrak a/X$ gives
\[
 M_W(s)=\widehat W(s)
       \sum_{(\mathfrak a,S)=1}
        \frac{\mu_K(\mathfrak a)\nu(\mathfrak a)}{N\mathfrak a^s}
       =\frac{\widehat W(s)}{L^S(s,\nu)}.
\]
Consequently the meromorphic functions satisfy
$L^S(s,\nu)M_W(s)=\widehat W(s)$ throughout the connected
half-plane $\operatorname{Re}s>\alpha$, by the identity theorem.
At a zero $\rho$, both $L^S$ and $M_W$ are holomorphic, so evaluation
gives $0=\widehat W(\rho)>0$. Finally, the finitely many factors
$1-\nu(\mathfrak p)N\mathfrak p^{-s}$ have no zeros in
$\operatorname{Re}s>0$, because their second term has absolute
value strictly less than one. Thus deleting them does not change
the zeros under consideration.
\end{proof}

There is an even simpler criterion when an exponentially smoothed
sum can be estimated directly. Define
\[
 F_\chi(X)=\sum_{n\ge1}\mu(n)\chi(n)e^{-n/X}.
\]
If $F_\chi(X)=O_{\chi,\epsilon}(X^{\alpha+\epsilon})$ for every
$\epsilon>0$, then
\[
 \int_0^\infty F_\chi(X)X^{-s}\,\frac{dX}{X}
       =\frac{\Gamma(s)}{L(s,\chi)}
       \qquad(\operatorname{Re}s>1)
\]
extends holomorphically to $\operatorname{Re}s>\alpha$.
At zero, the absolute sum is bounded by
$e^{-1/X}/(1-e^{-1/X})$; at infinity use the assumed estimate.
The preceding proof therefore applies with the single factor
$\Gamma(s)$, which has no zeros. Thus it suffices to estimate one
fixed smooth function; a collection of auxiliary functions is not
inherent in the analytic implication. This is a simplification of
the criterion, not a proof of the required cancellation estimate.

\subsection{The exact cost of amplification by powers of primes}

The following lemma isolates the numerical gain and the error that
must be paid. It applies to ideals as well as to positive integers.
We state it for $K=\mathbb Q(\sqrt{-3})$, where every ideal prime
to $3$ has a unique generator congruent to $1$ modulo $3$.
Write $\mathcal O=\mathbb Z[\omega]$, with
$\omega=(-1+\sqrt{-3})/2$. The six units give the six invertible
classes modulo $3\mathcal O$, so the asserted uniqueness follows
from the fact that $\mathcal O$ is a Euclidean domain.

\begin{lemma}[Power amplification]
\label{lem:qrh-amplification}
Let $k\ge2$, and let $\chi_{\mathfrak a}(u)$ be multiplicative
characters in $u$, of order dividing $k$ away from
$\mathfrak a$, extended by zero when $(u,\mathfrak a)\ne1$.
For a fixed bounded smooth function $W$ supported in a compact
subinterval of $(0,\infty)$, put
\[
 A_u(D)=\sum_{(\mathfrak a,S)=1}
 \mu_K(\mathfrak a)\nu(\mathfrak a)
 \chi_{\mathfrak a}(u)W(N\mathfrak a/D).
\]
Suppose $2\le Y=H^{1/k}\le D$, and that
\[
 \sum_{0<Nu\le H}|A_u(D)|^2\le B(D,H).
\]
If there are $\gg Y/\log Y$ prime ideals outside $S$ with norms
in $(Y/2,Y]$, then
\[
 |A_1(D)|^2\ll
       B(D,H)H^{-1/k}\log(2H)+D^2H^{-2/k}.
\]
The constants can depend on $K,S,W$ and on the constant in the
prime count, but not on $D$ or $H$.
\end{lemma}
\begin{proof}
Let $p$ be the primary generator of one of the indicated prime
ideals. The convention at nonunits gives exactly
\[
 \chi_{\mathfrak a}(p^k)=1_{p\nmid\mathfrak a}.
\]
Thus $A_1(D)-A_{p^k}(D)$ is a sum over ideals divisible by $p$.
There are $O_W(D/Np)$ such ideals, because the remaining ideal
has norm at most a fixed multiple of $D/Np$; the condition
$Y\le D$ absorbs the constant in ideal counting. Hence
\[
 |A_1(D)-A_{p^k}(D)|\ll_W D/Y.
\]
Distinct primary primes have distinct $k$th powers. If there
are $J\gg Y/\log Y$ of them, the inequality
$|z+w|^2\le2|z|^2+2|w|^2$ gives
\[
 J|A_1(D)|^2\le
 2\sum_p|A_{p^k}(D)|^2+O_W(JD^2/Y^2).
\]
Their norms are at most $H$, and the first sum is bounded by
$B(D,H)$. Divide by $J$.
\end{proof}

Suppose more generally that
$B(D,H)\ll_\epsilon D^{a+\epsilon}H^b$ and $H=D^h$.
The exponent furnished by this argument is
\begin{equation}\label{eq:qrh-general-exponent}
 \max\left\{\frac a2+\frac h2\left(b-\frac1k\right),
                    1-\frac hk\right\},
\end{equation}
up to an arbitrary positive loss. In the October paper,
$k=6$, $a=b=1$, and the proved range used for the mean square
is $H=D^{1+\eta}$, $0<\eta\le1/10$. Therefore the first
entry is $11/12+5\eta/12$, the second is
$5/6-\eta/6$, and letting $\eta$ tend to zero gives $11/12$.
This verifies the final deduction from Proposition~3.1 without
any claim about its difficult proof.

Even if the bound $B(D,H)\ll_\epsilon D^{1+\epsilon}H$
were known for every $H=D^h$, the best exponent obtainable
from \emph{this estimate together with this absolute bound for
prime deletion} would be $k/(k+1)$. Indeed the two entries in
\eqref{eq:qrh-general-exponent} are equal at
$h=k/(k+1)$, and one increases while the other decreases.
For $k=6$ this gives $6/7$, still greater than $1/2$.
This is a limitation of these two inequalities, not an
impossibility theorem for character amplification.

One can retain more arithmetic than the absolute deletion
bound. For $p\notin S$, multiplicativity of the M\"obius function
gives the exact identity
\begin{equation}\label{eq:qrh-prime-removal}
 A_1(D)=A_{p^k}(D)-\nu(p)A_{p^k}(D/Np).
\end{equation}
To prove it, split the defining sum according to $p\mid\mathfrak a$;
write $\mathfrak a=p\mathfrak b$ in the latter part, use
$\mu_K(p\mathfrak b)=-\mu_K(\mathfrak b)$ when $p\nmid\mathfrak b$,
and use zero otherwise. An a priori bound for $A_1$ can be used
in this identity to control prime deletion uniformly. Section~\ref{sec:qrh-bootstrap}
proves that improvement and states the additional mean-square estimate
it would require. A bound whose constant is allowed to depend on the
varying prime cannot be inserted into the amplification without a
separate uniformity argument.

For completeness, the positivity step alone cannot force
an exponent below $11/12$ in the range $H=D$. Consider abstract
complex numbers with $J$ distinguished entries, all equal to
$D/\sqrt J$, and all other entries zero. Their sum of squares
is exactly $D^2$, while the common distinguished value has
size $D^{11/12}$ when $J=D^{1/6}$.
This is an exact example for the numerical inequalities, not
an example consisting of M\"obius sums and not a counterexample
to either submitted theorem.

\subsection{Passing from the Eisenstein field to Dirichlet characters}

\begin{proposition}[Quadratic base change at the analytic endpoint]
\label{prop:qrh-basechange}
Let $0<\alpha<1$. If every finite-order Hecke $L$-function of
$K=\mathbb Q(\sqrt{-3})$ has no zero in
$\operatorname{Re}s>\alpha$, with principal poles allowed,
then the same conclusion holds for every Dirichlet $L$-function.
\end{proposition}
\begin{proof}
Let $\epsilon$ be the nontrivial character modulo $3$.
For a Dirichlet character $\chi$, pull back its ideal character
along the norm from $K$ to $\mathbb Q$. At primes where all
characters are unramified, the Euler factors give
\[
 L_K(s,\chi\circ N)=L(s,\chi)L(s,\chi\epsilon).
\]
If a rational prime splits, its two prime ideals have norm $p$
and $\epsilon(p)=1$, so both sides have factor
$(1-\chi(p)p^{-s})^{-2}$. If it is inert, its prime ideal
has norm $p^2$ and $\epsilon(p)=-1$, so both sides have
factor $(1-\chi(p)^2p^{-2s})^{-1}$.
At the finitely many excluded primes insert the usual Euler
factors; they are nonzero for $\operatorname{Re}s>0$.
Equality in the half-plane of absolute convergence extends
meromorphically.

Away from $1$, a zero of either Dirichlet factor would give
a zero of the product, because neither factor has a pole
there. At $1$, a cancellation could occur only when one of
the factors is principal, so the other is induced by
$\epsilon$. But
\[
 L(1,\epsilon)=\sum_{j=0}^\infty
      \left(\frac1{3j+1}-\frac1{3j+2}\right)>0.
\]
The series converges and every displayed summand is positive.
Imprimitive Euler factors at $1$ are also nonzero. Thus a
pole cannot conceal a zero, completing the proof.
\end{proof}

\subsection{A useful uniform consequence: logarithmic derivatives}

We now give a consequence of $\mathcal Z(\vartheta)$ that
really does shorten one of the auxiliary-prime arguments in
the submitted proof. The hypothesis remains explicit: the
following proof is independent of how $\mathcal Z(\vartheta)$
is obtained.

\begin{lemma}[A logarithm on a disk]
\label{lem:qrh-log-disk}
Let $f$ be holomorphic on $|z-z_0|<R$, and suppose
$\operatorname{Re}f\le M$ there. For $0<r<R$,
\[
 \sup_{|z-z_0|\le r}|f(z)-f(z_0)|
 \le \frac{2r}{R-r}\bigl(M-\operatorname{Re}f(z_0)\bigr).
\]
\end{lemma}
\begin{proof}
Put $g=f-f(z_0)$ and $A=M-\operatorname{Re}f(z_0)\ge0$.
If $A=0$, the maximum principle for the real part makes
$g=0$. If $A>0$, the function $g/(2A-g)$ maps the disk
into the closed unit disk, since $\operatorname{Re}g\le A$,
and vanishes at $z_0$. Schwarz's lemma after rescaling gives
$|g/(2A-g)|\le |z-z_0|/R$. Rearrangement gives the stated
bound. Applying the argument first on smaller closed disks
and taking a limit avoids assumptions at the boundary.
\end{proof}

\begin{proposition}[Uniform logarithmic derivative]
\label{prop:qrh-log-derivative}
Assume $\mathcal Z(\vartheta)$. For every fixed
$\sigma>\vartheta$, there is a constant $C_{\vartheta,\sigma}$
such that, for every nonprincipal Dirichlet character modulo
$q$ and all real $t$,
\[
 \left|\frac{L'}{L}(\sigma+it,\chi)\right|
 \le C_{\vartheta,\sigma}\log\bigl(q(|t|+3)\bigr).
\]
The same estimate holds at every point in any fixed closed
vertical strip whose real parts lie in
$(\vartheta,\infty)$. In particular
$|L'(1,\chi)/L(1,\chi)|\ll_\vartheta\log(q+3)$.
For a principal character, subtract its simple-pole term
before making the corresponding assertion.
\end{proposition}
\begin{proof}
It is enough first to consider primitive nonprincipal
characters and $\vartheta<\sigma\le2$. The partial sums
$A(u)=\sum_{n\le u}\chi(n)$ have absolute value at most $q$:
a complete period sums to zero. Partial summation therefore
gives, for $\operatorname{Re}s>0$,
\[
 L(s,\chi)=s\int_1^\infty A(u)u^{-s-1}\,du,
 \qquad |L(s,\chi)|\le\frac{q|s|}{\operatorname{Re}s}.
\]
This also proves holomorphy there by locally uniform
convergence. Choose radii
\[
 2-\sigma<r<R<2-\vartheta
\]
depending only on $\vartheta,\sigma$. On the disk of radius
$R$ centered at $2+it$, the zero-free hypothesis gives a
holomorphic logarithm $f=\log L$. Choose its value at the
center by the absolutely convergent Euler-product logarithm;
that value is bounded absolutely, uniformly in $q,t$.
The preceding bound gives
$\operatorname{Re}f\le\log(q(|t|+3))+C_{\vartheta,\sigma}$
on the larger disk. Lemma~\ref{lem:qrh-log-disk} bounds $f$
on the disk of radius $r$ by the same quantity.
A further fixed smaller circle around $\sigma+it$ lies
inside that disk. Cauchy's derivative formula bounds
$f'(\sigma+it)=L'/L(\sigma+it,\chi)$ as asserted.
The argument is uniform as $\sigma$ varies in a compact
subinterval of $(\vartheta,2]$. For larger real parts the
Euler product gives the bound directly.

If $\chi$ is induced from $\chi^*$, then
\[
 \frac{L'}L(s,\chi)=\frac{L'}L(s,\chi^*)+
 \sum_{\substack{p\mid q\\p\nmid q^*}}
       \frac{\chi^*(p)\log p}{p^s-\chi^*(p)}.
\]
For $\operatorname{Re}s\ge\sigma>0$, the absolute value of
the sum is at most
$(2^\sigma-1)^{-1}\sum_{p\mid q}\log p
\le(2^\sigma-1)^{-1}\log q$.

For the primitive principal character use
$F(s)=(s-1)\zeta(s)$ instead of $L(s,\chi)$. It is
holomorphic and zero-free in the relevant disks, including
at $1$. The elementary continuation formula
\[
 \zeta(s)=\frac{s}{s-1}
          -s\int_1^\infty\{u\}u^{-s-1}\,du
\]
bounds $F$ polynomially in $|t|+3$ on those disks.
At their centers a logarithm of $F$ is
$\log(1+it)+\log\zeta(2+it)$, of size
$O(\log(|t|+3))$. The same argument bounds $F'/F$;
finally $\zeta'/\zeta=F'/F-1/(s-1)$.
\end{proof}

\subsection{Primes in progressions without an exceptional-zero case}

Let $\Lambda(n)$ equal $\log p$ when $n=p^j$ is a prime
power, and zero otherwise. Choose a fixed nonnegative
$\Phi\in C_c^\infty(1,2)$ with $\int\Phi(y)\,dy>0$.

\begin{proposition}[A smooth prime-counting estimate]
\label{prop:qrh-smooth-primes}
Assume $\mathcal Z(\vartheta)$ and fix
$\vartheta<\sigma<1$. Uniformly in $q\ge1$, $(a,q)=1$
and $X\ge2$,
\[
 \sum_{n\equiv a\ (q)}\Lambda(n)\Phi(n/X)
 =\frac{X}{\varphi(q)}\int_0^\infty\Phi(y)\,dy
   +O_{\vartheta,\sigma,\Phi}
      \bigl(X^\sigma\log(q+3)\bigr).
\]
Consequently, for every $B>1/(1-\vartheta)$ there is
$C_B$ such that every reduced class modulo $q$ contains
a prime in $(X,2X)$ whenever $X\ge C_B(q+1)^B$.
\end{proposition}
\begin{proof}
For $c>1$, absolute convergence of the Euler product and
Mellin inversion give
\[
 \sum_n\Lambda(n)\chi(n)\Phi(n/X)
 =\frac1{2\pi i}\int_{(c)}
    -\frac{L'}L(s,\chi)\widehat\Phi(s)X^s\,ds.
\]
Mellin inversion here is ordinary Fourier inversion after
the substitution $y=e^u$. Repeated integration by parts in
$u$ gives, for every integer $A\ge0$,
$|\widehat\Phi(v+it)|\ll_{A,\Phi}(1+|t|)^{-A}$,
uniformly for $v$ in a fixed compact interval.
Proposition~\ref{prop:qrh-log-derivative} therefore permits
moving the line from $c$ to $\sigma$. On the horizontal
segments the integral tends to zero; the vertical integral
is absolutely bounded by
$O(X^\sigma\log(q+3))$, since
$\log(q(|t|+3))\le\log(q+3)+\log(|t|+3)$.
The only pole crossed is at $1$ for a principal character,
with residue $X\widehat\Phi(1)$.
Character orthogonality gives the progression formula.
There are $\varphi(q)$ terms in that average and the factor
$1/\varphi(q)$ cancels their number, so the error does not
acquire an extra factor of $\varphi(q)$.

Terms with $n=p^j$, $j\ge2$, contribute
$O_\Phi(X^{1/2}(\log X)^2)$: for each
$2\le j\le\log(2X)/\log2$ there are at most
$(2X)^{1/j}\le(2X)^{1/2}$ possibilities for $p$, and
$\log p\le\log(2X)$. Since $\sigma>1/2$, this is $O_\sigma(X^\sigma)$,
although it is useful to retain its smaller exponent when
comparing it with the main term.
Choose $\sigma$ with $\vartheta<\sigma<1-1/B$.
For $X\ge C_B(q+1)^B$, the ratios of both errors to
$X/\varphi(q)$ tend to zero uniformly as $C_B$ tends
to infinity: use $\varphi(q)\le q$,
$qX^{\sigma-1}\log(q+3)\ll
 C_B^{\sigma-1}(q+1)^{1-B(1-\sigma)}\log(q+3)$,
and the analogous bound for prime powers.
The weighted sum over actual primes is consequently
positive, proving existence.
\end{proof}

\begin{corollary}[The progression input in the submitted proof]
\label{cor:qrh-ProgAt}
Assume $\mathcal Z(\vartheta)$. There is a constant $C=C_\vartheta\ge1$ such that, for every positive integer $M$,
every integer $a$ coprime to $M$, and every
$X\ge\max(3,M^C)$, there is a prime $p$ satisfying
\[
 X<p\le X^2,\qquad p\equiv a\pmod M,
 \qquad
 \left|\frac{L'(1,\chi_p)}{L(1,\chi_p)}\right|
      \le C\log p,
\]
where $\chi_p(n)=(n/p)$ is the quadratic character modulo
$p$. For the fixed boundary $\vartheta=7/8$, the constant is absolute.
This supplies the ordinary progression input used in the
submitted argument.
\end{corollary}
\begin{proof}
For $M\ge2$, enlarge $C$ so that $M^C$ dominates
$C_B(M+1)^B$ in Proposition~\ref{prop:qrh-smooth-primes}.
The prime then lies in $(X,2X)\subset(X,X^2]$.
Proposition~\ref{prop:qrh-log-derivative} at $s=1$ gives
the logarithmic derivative bound for every such prime,
not merely for a selected subfamily. For $M=1$,
Bertrand's postulate supplies a prime in $(X,2X)$ for
$X\ge3$, and the same bound applies. One may also cover
bounded $X$ by finitely many primes and increase $C$.
\end{proof}

\begin{corollary}[The two auxiliary quadratic fields]
\label{cor:qrh-auxiliary-pair}
Assume $\mathcal Z(\vartheta)$ and fix $\epsilon>0$.
For all sufficiently large $V$ there are distinct primes
$q,d$ such that
\[
 q\equiv3\pmod8,\quad d\equiv7\pmod{16},\quad q\mid d+1,
 \quad p\mid d+1\quad
  \text{for every odd prime }p\le V^{\epsilon/5},
\]
\[
 \tfrac12V^{\epsilon/5}<\log q\le V^{\epsilon/5},\qquad
 \tfrac12V^{\epsilon/5}(\log V)^2<\log d
                         \le V^{\epsilon/5}(\log V)^2,
\]
and
\[
 \left|\frac{L'(1,\chi_d)}{L(1,\chi_d)}\right|
       \le V^\epsilon.
\]
In particular $(-d/p)=1$ for every indicated odd prime $p$.
Since $d\equiv3\pmod4$, quadratic reciprocity identifies
$\chi_d(n)=(n/d)$ with the primitive quadratic character
$\bigl(\frac{-d}{n}\bigr)$ of the field
$\mathbb Q(\sqrt{-d})$. For odd $n$ the reciprocity signs
cancel against $(-1/n)$; the supplementary law at $2$ gives
the same equality for even $n$, and multiplicativity completes
the verification.
\end{corollary}
\begin{proof}
Put $Y=V^{\epsilon/5}$. Apply
Proposition~\ref{prop:qrh-smooth-primes} to the class $3$
modulo $8$ at $X=e^{3Y/4}$. The resulting prime $q$
has the required logarithmic bounds and satisfies $q>Y$
for large $V$.
Let $P=\prod_{3\le p\le Y}p$ and $M=16qP$.
The Chinese remainder theorem gives a reduced class
$a\pmod M$ with $a\equiv7\pmod{16}$ and
$a\equiv-1\pmod{qP}$.
We only need the elementary bound
$\log P\le Y\log Y$: there are at most $Y$ factors,
each at most $Y$. Thus $\log M=O_\epsilon(Y\log V)$.
Take $X=e^{3Y(\log V)^2/4}$. For every fixed $B$,
$X\ge C_B(M+1)^B$ once $V$ is sufficiently large.
A prime $d\in(X,2X)$ in the class $a$ has all the
required size and congruence conditions. It is distinct
from $q$ because $d$ is much larger. Moreover
$-d\equiv1\pmod p$ proves $(-d/p)=1$.
Finally Proposition~\ref{prop:qrh-log-derivative} bounds
the logarithmic derivative by
$O_\vartheta(\log d)=O_\vartheta(Y(\log V)^2)$,
which is less than $Y^5=V^\epsilon$ eventually.
\end{proof}

These corollaries remove an exceptional-zero selection
argument from the ordinary Dirichlet part of the proof.\cite{alpo-prime-supply}
They do not replace the distinct assertion about primes
represented by
\[
 (1+Mu)^2+d(Mv)^2.
\]
That assertion is a prime theorem for principal ray classes
in the varying imaginary quadratic field
$\mathbb Q(\sqrt{-d})$. A Hecke character of this field
need not be a norm pullback of a Dirichlet character.
Neither $\mathcal Z(\vartheta)$ nor the supplied theorem
for the one fixed field $\mathbb Q(\sqrt{-3})$ supplies
that varying-field theorem. Retaining this distinction is
essential for a valid simplification.

\subsection{Why these estimates do not supply the remaining arithmetic argument}

In the final scalar estimate of the submitted argument,\cite{alpo-final-scalar}
$A\le2xb$ is an upper bound for an auxiliary height or
weight parameter. It is not a bound for a M\"obius sum.
After division by the common positive factor, its contribution
includes $A/(mb)$, while another contribution is proportional
to $m$. Thus the elementary balance is $x/m+m$, whose
minimum is $2\sqrt{x}$. A proof of
$\mathcal Z(7/8)$ does not change $A$ into
$O(bx^{7/8})$.

The RH subcase separately requires square-root bounds for
three selected quadratic characters. The full Dirichlet-GRH
construction requires them for a finite family of twists of
fixed characters by the selected quadratic characters, uniformly
on the chosen interval of scales. A uniform ordinary bound
$O(x^{7/8+\epsilon})$ is weaker and cannot be substituted
directly. However, selection by averaging can use the fixed
zero-free half-plane in a different way; the continuation below
gives a conditional proof of that selection, including the fixed
twists. The estimates for the theta lifts, the local divisibility
and operator norms, and the Galois cohomology still remain
separate arithmetic inputs.

\subsection{Two checks on proposed analytic shortcuts}

The catalogue is useful for locating exact reductions and
for recording unsuccessful estimates.\cite{qrh-catalogue} Its statements that
a component is conditional, source-relative, or only checked
as an algebraic inequality have mathematical force: they
cannot be dropped when using the component. The following
obstruction, corresponding to its entry G280, admits a
short complete proof.

\begin{proposition}[Finite Dirichlet polynomials and absolute convergence]
\label{prop:qrh-absolute-obstruction}
Let $\chi$ be a Dirichlet character and let
$B(s)=\sum_{d\le D}b_d d^{-s}$ be a finite Dirichlet
polynomial. Expand
$1-L(s,\chi)B(s)=\sum_{n\ge1}r_n n^{-s}$ for
$\operatorname{Re}s>1$. For real $\sigma\le1$, the
inequality $\sum_n|r_n|n^{-\sigma}<1$ is impossible.
\end{proposition}
\begin{proof}
If $b_1=0$, then $r_1=1$. If $b_1\ne0$, every prime
$p>D$ coprime to the modulus has
$r_p=-b_1\chi(p)$, so $|r_p|=|b_1|$.
The series $\sum_p1/p$ diverges. To recall its elementary
proof, convergence would bound
$\prod_{p\le y}(1-p^{-1})^{-1}$ uniformly in $y$;
expanding that product gives at least
$\sum_{n\le y}1/n$, a contradiction. Removing finitely
many primes does not alter divergence. Since
$p^{-\sigma}\ge p^{-1}$ for $\sigma\le1$, the
absolute residual series is infinite.
\end{proof}

A different technical issue is whether a Mellin multiplier
could vanish at the very zero one wants to exclude.
Catalogue entry G276 suggests the following useful general
fact, which can be proved without asymptotic integration.

\begin{proposition}[A finite family of Mellin transforms]
\label{prop:qrh-mellin-family}
Let $a>0$ and let $P(y)=\sum_{j=0}^d c_jy^j$ be a real
polynomial of degree $d\ge1$. For
$\operatorname{Re}s>0$ put
\[
 F(s)=\int_0^\infty y^{s-1}e^{-ay-iP(y)}\,dy.
\]
The $d$ numbers $F(s),F(s+1),\ldots,F(s+d-1)$ never
vanish simultaneously.
\end{proposition}
\begin{proof}
Exponential decay and local integrability at zero justify
holomorphy and differentiation under the integral on every
compact subset of $\operatorname{Re}s>0$.
Integration of the derivative of
$y^s e^{-ay-iP(y)}$ gives
\[
 sF(s)=aF(s+1)+i\sum_{j=1}^d j c_jF(s+j).
\]
For $d\ge2$, the coefficient $idc_d$ is nonzero.
If the first $d$ consecutive values vanished, this relation,
then its translates by positive integers, would make every
$F(s+n)$ vanish. For $d=1$ the same conclusion uses
$a+ic_1\ne0$.

Now the finite complex measure
$y^{s-1}e^{-ay-iP(y)}\,dy$ has Laplace transform
\[
 H(z)=\int_0^\infty y^{s-1}e^{-ay-iP(y)-zy}\,dy,
       \qquad\operatorname{Re}z>-a.
\]
All its derivatives at zero would vanish. Holomorphy
would give $H\equiv0$. Push the measure forward under
$u=e^{-y}$ to $[0,1]$. Its moments are $H(n)=0$ for
all integers $n\ge0$. Polynomial approximation of
continuous functions on $[0,1]$ and finite total variation
then show that this measure is zero. Its original density
is nonzero everywhere on $(0,\infty)$, a contradiction.
\end{proof}

This last proposition removes a possible cancellation by
a fixed Mellin multiplier. It supplies no bound for the
arithmetic sums multiplying that factor. In particular,
it cannot turn a bound with exponent $7/8$ into one with
exponent $1/2$.

\subsection{An exact cancellation in the double Poisson argument}

The reindexing in Section 7.3 of the October paper contains a useful
finite identity. Its strength comes from performing the cancellation
before replacing any summand by its absolute value. The following
weighted form is slightly more general and has a short proof.

\begin{proposition}[A squarefree factorization identity]
\label{prop:qrh-factorization}
Let $f$ be squarefree in a unique factorization domain, or let
$f$ be a squarefree integral ideal of a number field. Let $k\ne0$
be an element, or a nonzero integral ideal, respectively.
Assign an arbitrary complex number $c(p)$ to each prime dividing
$f$, and extend $c$ multiplicatively to its squarefree divisors.
Then
\[
 \sum_{Cew=f}\mu(e)\mu(w)c(e)
    1_{e\mid k}\,1_{(w,k)=1}
           \sum_{d\mid(C,k)}c(d)
       =1_{f\mid k}.
\]
Factorizations count associates only once in the element formulation.
Since $f$ is squarefree, the three factors $C,e,w$ are automatically
pairwise coprime.
\end{proposition}
\begin{proof}
Both sides factor over the primes dividing $f$. Fix such a prime
$p$, and put $I=1_{p\mid k}$. Placing $p$ in $C$ contributes
$1+c(p)I$, placing it in $e$ contributes $-c(p)I$, and placing
it in $w$ contributes $-(1-I)$. Their sum is $I$. Multiplication
over $p\mid f$ gives the result. The empty product handles $f=1$.
\end{proof}

For $c(p)=1$ this is the cancellation used after the two Poisson
summations: it forces the new auxiliary divisor to divide the new
frequency. It is valid at shared primes, including the terms where
that frequency is divisible by a prime in the original exclusion.
Dropping those terms would change the identity. The weighted form
shows that bounded prime weights can sometimes be retained through
this particular step; a theorem about the complete surrounding
mean square would still have to pay all other terms and the
resulting uniform constants.


\section{A conditional improvement using exact prime removal}
\label{sec:qrh-bootstrap}

The absolute deletion error in Lemma~\ref{lem:qrh-amplification}
loses information about the M\"obius function. Retaining that
information gives a better conditional implication. The new
mean-square estimate required by the implication is stated explicitly;
it is not established by the supplied estimates in their present ranges.

Use the notation of Lemma~\ref{lem:qrh-amplification}, fix
$k\ge2$, and suppose that the finite-order twists and their zero
extensions satisfy
\[
 A_1(T)=A_{p^k}(T)-\nu(p)A_{p^k}(T/Np)
       \qquad(p\notin S).
\]
For a fixed compactly supported $W$, all these sums vanish for
sufficiently small positive $T$. This point makes the following
iteration exact, without an infinite unestimated remainder.

\begin{lemma}[Uniform control of the deleted prime]
\label{lem:qrh-better-deletion}
Suppose $\beta>0$ and
$A_1(T)\ll_\epsilon T^{\beta+\epsilon}$ for every
$\epsilon>0$, with fixed arithmetic data and fixed $W$.
Then, uniformly in primes $p\notin S$ and all $T>0$,
\[
 A_{p^k}(T)=\sum_{j\ge0}\nu(p)^j A_1(T/(Np)^j),
 \qquad
 |A_{p^k}(T)|\ll_\epsilon T^{\beta+\epsilon}.
\]
Consequently, for $Y/2<Np\le Y$,
\[
 |A_1(D)-A_{p^k}(D)|
       \ll_\epsilon(D/Y)^{\beta+\epsilon}.
\]
\end{lemma}
\begin{proof}
Solve the displayed prime-removal identity for $A_{p^k}(T)$
and iterate. For each $T>0$, the sum $A_{p^k}(T/(Np)^J)$
vanishes once $J$ is large. Thus the expansion is a finite identity.
Choose $v>0$ containing the upper endpoint of the support
of $W$. The assumed estimate can be made valid for every
$T>0$ by increasing its constant: the sum is zero for
$T<1/v$, and is bounded on every fixed compact range of
positive $T$. Hence
\[
 |A_{p^k}(T)|\le
 C_\epsilon T^{\beta+\epsilon}
       \sum_{j\ge0}(Np)^{-j(\beta+\epsilon)}
 \le\frac{C_\epsilon T^{\beta+\epsilon}}
                {1-2^{-\beta-\epsilon}}.
\]
The denominator is positive and independent of $p$.
Apply this bound at $T=D/Np$, and use $Np>Y/2$,
to obtain the last assertion.
\end{proof}

\begin{proposition}[A conditional exponent improvement]
\label{prop:qrh-exponent-improvement}
Let $1/2<\beta\le1$ and $0<h<k$. Suppose that, for every
$\epsilon>0$, the following two estimates hold for the same
fixed $W$ and arithmetic data:
\[
 A_1(D)\ll_\epsilon D^{\beta+\epsilon},\qquad
 \sum_{0<Nu\le D^h}|A_u(D)|^2
       \ll_{h,\epsilon}D^{1+h+\epsilon}.
\]
Suppose also that the prime-ideal count used in
Lemma~\ref{lem:qrh-amplification} is available. Then
\[
 A_1(D)\ll_{h,\epsilon}D^{\beta'(h)+\epsilon},
 \qquad
 \beta'(h)=\max\left\{
    \frac12+\frac{k-1}{2k}h,
    \beta\left(1-\frac hk\right)\right\}.
\]
In particular, if the mean-square estimate is available at
\[
 h=h_*(\beta):=\frac{k(2\beta-1)}{k-1+2\beta},
\]
then
\begin{equation}\label{eq:qrh-bootstrap-map}
 \beta'=T_k(\beta):=\frac{k\beta}{k-1+2\beta}<\beta.
\end{equation}
All implied constants may depend on $\beta,h,W$ and on the
fixed character, but must be independent of the varying prime
$p$ and of $D$.
\end{proposition}
\begin{proof}
Set $H=D^h$ and $Y=H^{1/k}$. The proof of
Lemma~\ref{lem:qrh-amplification}, now using
Lemma~\ref{lem:qrh-better-deletion}, gives
\[
 |A_1(D)|^2\ll
 D^{1+h-h/k+\epsilon_1}\log(2D)
       +(D/Y)^{2(\beta+\epsilon_2)}.
\]
Choose the preliminary positive losses $\epsilon_1,\epsilon_2$
small enough in terms of any requested final loss, absorb
$\log(2D)$ into a further arbitrarily small power, and take
square roots. The first resulting exponent increases with
$h$ and the second decreases. They coincide at $h_*(\beta)$;
substitution gives \eqref{eq:qrh-bootstrap-map}. Since
$\beta>1/2$, one has $0<h_*(\beta)<k$ and
$k-1+2\beta>k$, so the new exponent is strictly smaller.
\end{proof}

For sixth-power characters this would give
\[
 \beta=\frac78,\quad h_* =\frac23,
       \quad T_6(\beta)=\frac79;
 \qquad
 \beta=\frac{11}{12},\quad h_* =\frac{30}{41},
       \quad T_6(\beta)=\frac{33}{41}.
\]
Thus the first useful extension of the known $7/8$ result
would concern a mean square over $Nu\le D^{2/3}$.
That is substantially shorter than the available range
$Nu\le D^{1+\eta}$.

If the hypothesis at $h_*(\beta)$ could be deduced from the
current exponent bound for every $\beta>1/2$, iteration
would approach the critical line. More precisely, if
$\beta_{n+1}=T_k(\beta_n)$, then
\[
 \frac1{\beta_n}
 =2-\left(\frac{k-1}{k}\right)^n
               \left(2-\frac1{\beta_0}\right),
 \qquad \beta_n\downarrow\frac12.
\]
This follows by taking reciprocals in
\eqref{eq:qrh-bootstrap-map} and solving the affine
recurrence. For any fixed zero to the right of $1/2$,
only finitely many iterations would be needed to obtain
a contradictory bound in Proposition~\ref{prop:qrh-mellin}.
Thus constants may deteriorate with the number of
iterations; no uniform limit of those constants is being
assumed. For the conclusion about every Hecke character, these
implications must hold for every fixed character and every
compactly supported smooth $W$, with a common exponent;
their implied constants may depend on those fixed data.
The missing assertion is the implication from the
current exponent to the shorter mean square.

\subsection{What a fixed lower limit on the range permits}

Suppose the mean-square estimate is available only at
exponents $h\ge h_0$, where $0<h_0<k$. The best consequence of the last
proposition is obtained at $h=\max(h_0,h_*(\beta))$
whenever this lies in the permitted range. In particular,
a strict improvement over $\beta$ requires
\[
 h_0<\frac{k(2\beta-1)}{k-1}.
\]
Indeed the first entry in $\beta'(h)$ must be smaller than
$\beta$, while the second is smaller for every $h>0$.
A fixed $h_0$ leaves the lower bound
$1/2+(k-1)h_0/(2k)$ on the exponent obtainable by this
particular argument. With $k=6$ and $h_0=1$, that bound
is $11/12$. A previously known $7/8$ estimate therefore
does not improve when inserted into the existing
$h>1$ mean-square range.

One must also avoid a circular reformulation. If the bound
\[
 \sum_{0<Nu\le D^h}|A_u(D)|^2\ll_{h,\epsilon}
                         D^{1+h+\epsilon}
\]
were simply assumed for every $h>0$, its single term
$u=1$ would already give
$A_1(D)\ll_\epsilon D^{1/2+\epsilon}$ by taking $h$
sufficiently small. Amplification would not be needed.
Even at $h=h_*(\beta)$, the single term gives a strict,
though weaker, improvement, because
$h_*(\beta)<2\beta-1$. Thus proving this shorter
mean-square estimate is itself substantive new cancellation;
it is not a formal consequence of changing parameters.

\subsection{The precise failure of the existing iteration in this range}

The October proof transforms the original scale $D$ and
range $H$ by Poisson summation. In the contribution whose
common-divisor parameter is $1$, its new parameters have
sizes
\[
 X=D,\qquad F=1,\qquad \mathcal H\asymp D^2/H,
          \qquad\Sigma=XF=D.
\]
The mean-square theorem in its Proposition 5.1 requires
$\mathcal H\le\Sigma D^{-\kappa}$ for a fixed
$\kappa>0$. This requires $H$ to be greater than $D$
by a fixed power. If $H=D^h$ with $h<1$, the ratio
$\mathcal H/\Sigma\asymp D^{1-h}$ instead tends to
infinity. Fixed scale constants do not change this
conclusion.

This is also visible directly in the cube-divisor step.
The splitting threshold in its equation (5.6) satisfies
\[
 H_c^3=\min(X,X^2/\mathcal H^2)<1
\]
for large $D$. The unit divisor $b=1$ is therefore in
the remaining sum, and its new length is $L_b=X$.
An attempted reuse of the displayed supremum for that
remaining sum would include the original mean square itself.
There is no decrease of the parameters. In fact, the lemma
stating that bound already assumes $\mathcal H\le\Sigma$;
we are diagnosing why its proof does not extend to this
case, not applying the lemma outside its hypotheses. The next transfer proposition requires
$\mathcal H\le\Sigma$ and is unavailable at the desired
$\Sigma=X$. Thus the existing finite induction does not
prove the shorter-range estimate.

In the September paper, the unmarked case of the principal
inverse moment requires $r<m-c$ when the norm length is
$D=Z^r$ and the outer range is $H=Z^m$. Its canonical
estimate also requires, even when the additional excluded
ideal is the unit ideal,
$1\le Z^{F_0-M-c}$ with $c>0$. These conditions enforce
the same strict direction of inequality between the
corresponding ranges. The longer-length estimate later
obtained for sixth-power-free elements does not remedy this
problem: the amplifying elements $p^6$ have valuation six
and are excluded from that index set.

These are exact limitations of the statements and
recurrences presently available. They do not prove that a
new argument cannot extend the mean square. They locate
where such an argument is needed and prevent a purported
iteration to the critical line from concealing that step.

\subsection{A counterexample to enlarging the coefficient class}

The shorter mean square also cannot be obtained by applying an
identical estimate for arbitrary bounded squarefree coefficients.
Here is an explicit obstruction. Fix $0<h<k/(k-1)$, put
$H=D^h$, and take the fixed column coefficient $c_{\mathfrak a}$
to be $1$ when $\mathfrak a$ is prime outside $S$, and
zero otherwise. This coefficient sequence is independent
of $D$ and of the outer index. Choose a fixed nonnegative
$W\in C_c^\infty((1/3,3/2))$ which is at least one
on $[1/2,1]$. For every prime ideal $p\notin S$ of norm in
$[H^{1/k}/2,H^{1/k}]$, one has $Np<D/3$ when $D$ is
large, because $h/k<1$. Thus every prime column on which
$W(N\mathfrak a/D)\ne0$ is different from $p$. Thus the zero-extended character
is exactly $1$ on every column at the argument $p^k$.
The prime-ideal count gives
\[
 \left|\sum_{\mathfrak a}\!c_{\mathfrak a}
      \chi_{\mathfrak a}(p^k)W(N\mathfrak a/D)\right|
       \gg\frac D{\log D}.
\]
Their primary generators have distinct $k$th powers:
equality of the powers would give equality of the prime
ideals by unique factorization, and then equality of their
primary generators. Multiplicativity of the norm gives
$N(p^k)=(Np)^k\le H$. There are consequently
$\gg H^{1/k}/\log D$ distinct such arguments.
Consequently the complete mean square is at least
\[
       \gg\frac{D^2H^{1/k}}{(\log D)^3}.
\]
The ratio of this lower bound to $DH$ grows like
$D^{1-h(1-1/k)}/(\log D)^3$, which tends to infinity.
It even contradicts an estimate $D^{1+\epsilon}H$ for
sufficiently small fixed $\epsilon>0$.
This is an actual counterexample for the enlarged
coefficient class. It is not a counterexample for the
M\"obius coefficients in the supplied papers. It shows
why their special arithmetic identities cannot be replaced
by a statement about arbitrary bounded coefficients.
If a fixed finite-order twist $\nu$ is retained explicitly,
replace $c_{\mathfrak a}$ on prime ideals by
$\overline{\nu(\mathfrak a)}$; the same argument applies.
The counterexample concerns the complete outer set. It does
not assert anything about a restricted set deliberately
excluding $k$th powers. Fixed exclusions remove only
finitely many primes and hence do not change the argument.

The later catalogue entries R107 and R109 retain the same distinction.
Their amplified inverse-moment exponent, with base $D$, inverse
length exponent $1$, no additional prime factors, and outer
range exponent $h<1$, is $(h+5)/6$ after normalizing the inverse
sum by $D^{-1/2}$. Thus their unnormalized bound has exponent
$(11+h)/6$, which exceeds the desired exponent $1+h$ by
$5(1-h)/6$. Their index set is also restricted to
sixth-power-free elements. The more elaborate cofactor bookkeeping
in those entries consequently does not provide the missing
shorter-range mean square.\cite{qrh-catalogue}

There is another useful numerical check. The September marked
inverse moment, with inverse length $D$ and additional prime factors
of total length $D^z$, requires outer range exponent
\[
 h>1+2z,\qquad 3h>2+8z.
\]
If one tries to use those prime factors merely as positive
amplifiers at arguments $p^6$, their product has size
$D^{z/2}$ up to logarithms. The same positivity calculation
then gives the first exponent
\[
 \frac12+\frac{5h}{12}-\frac z2.
\]
The first required inequality already makes this greater than
$11/12+z/3$. Taking $z>0$ therefore does not improve the
unmarked $11/12$ consequence by this particular use of the
marked estimate. The actual $7/8$ proof uses a different
compensated expression and a more detailed estimate; it is
not the direct positive-amplification argument just tested.


% Standalone mathematical sections for inclusion in the reconstruction manuscript.
% Requires amsmath, amssymb, amsthm; theorem, lemma, proposition, corollary,
% and remark environments. All labels and citation keys have the prefix sz-.
\section{An elementary determinant estimate for quadratic characters}
\label{sz-section}

This section isolates the part of the suggested auxiliary literature that can
be proved with particularly little machinery. The argument originates in the
OpenAI proof of the exclusion of Landau--Siegel zeros and its subsequent
simplifications, including Kerr's earlier interpolation argument
\cite{sz-kerr}. The rectangular form below follows the use of derivatives
in Green's argument, as explained by Tao, and the reduction modulo the largest
inert prime used by Pescador and van Dijk
\cite{sz-openai,sz-tao,sz-pescador,sz-vandijk}.
We give all of the determinant argument, with an explicit finite inequality
and a slightly sharper bound for the determinant. No result about automorphic
representations is used in this section.

For a nonsquare integer $d$, set $A=\mathbb Z[\sqrt d]$ and write
$\overline{u+v\sqrt d}=u-v\sqrt d$. An odd prime $p$ will be called inert
here if the Legendre symbol $(d/p)$ is $-1$. This definition makes sense even
when $A$ is not the maximal order. It implies
\begin{equation}\label{sz-frobenius}
 A/pA\simeq\mathbb F_p[T]/(T^2-d)\simeq\mathbb F_{p^2},
 \qquad \alpha^p\equiv\bar\alpha\pmod {pA}\quad(\alpha\in A).
\end{equation}
Indeed, $T^2-d$ is irreducible modulo $p$, and Euler's criterion gives
$(\sqrt d)^p=-\sqrt d$ in that field.

\begin{lemma}[Interpolation in characteristic $p$]\label{sz-interpolation}
Let $k$ be a field of characteristic $p$, let $1\le s\le p$, and let
$a_1,\ldots,a_p$ be distinct elements of $k$. The map
\[
 F\longmapsto\bigl((\partial_z^jF)(a_i,a_i^p)\bigr)_{
          1\le i\le p,\ 0\le j<s}
\]
from the polynomials with $\deg_z F<p$ and $\deg_w F<s$ to $k^{ps}$
is an isomorphism.
\end{lemma}
\begin{proof}
The source and target have dimension $ps$, so injectivity suffices. Put
$G(z)=F(z,z^p)$. The monomials $z^rw^t$ become the distinct monomials
$z^{r+pt}$, for $0\le r<p$ and $0\le t<s$. Thus $F\ne0$ implies
$G\ne0$ and $\deg G<ps$. Since $(z^p)'=0$, the chain rule gives
$G^{(j)}(a_i)=(\partial_z^jF)(a_i,a_i^p)$ for every $j$. Vanishing of the
displayed evaluations forces $G$ to have a zero of order at least $s$ at
each $a_i$: the coefficients of $(z-a_i)^j$, for $j<s$, are
$G^{(j)}(a_i)/j!$, and these factorials are invertible because $j<p$.
The distinct factors $(z-a_i)^s$ then give $\deg G\ge ps$, a contradiction.
\end{proof}

\begin{proposition}[A finite determinant inequality]\label{sz-finite}
Let $R$ be an odd inert prime, let $1\le S\le R$, and let
$\alpha_1,\ldots,\alpha_R\in A$ have distinct reductions modulo $RA$.
Suppose that both complex embeddings of every $\alpha_i$ have modulus
at most $H$, where $H\ge1$. Put $N=RS$ and, for each odd inert prime $p$, put
\begin{align}
 E_p&=\sum_{r=0}^{R-1}\sum_{s=0}^{S-1}
        \min\{\lfloor r/p\rfloor,S-1-s\}\label{sz-E}\\
    &=\sum_{j=1}^{S-1}(S-j)\max\{R-jp,0\}.\nonumber
\end{align}
Then
\begin{equation}\label{sz-finite-ineq}
 2\sum_{\substack{p\ \mathrm{odd}\\ (d/p)=-1}} E_p\log p
 \le N\log N+N(R-1)\log H+N(S-1)\log R.
\end{equation}
The sum is finite, since $E_p=0$ for $p\ge R$.
\end{proposition}
\begin{proof}
Form the $N\times N$ matrix
\[
 B_{(r,s),(i,j)}=(\partial_z^j z^rw^s)(\alpha_i,\bar\alpha_i),
 \quad 0\le r<R,\quad 0\le s<S,\quad1\le i\le R,\quad0\le j<S.
\]
Its determinant $\Delta$ belongs to $A$. Reducing modulo $RA$,
\eqref{sz-frobenius} and Lemma~\ref{sz-interpolation} show that $B$ is
invertible. In particular $\Delta\ne0$. Consequently
$M=|\Delta\bar\Delta|$ is a positive rational integer.

Fix an odd inert prime $p$ and set
$h_{r,s}=\min\{\lfloor r/p\rfloor,S-1-s\}$. Replace the row
corresponding to $z^rw^s$ by the evaluation row of
\[
 z^{r-ph_{r,s}}(z^p-w)^{h_{r,s}}w^s.
\]
The leading term is $z^rw^s$. Every other term has smaller $z$-degree
and still has $z$-degree below $R$ and $w$-degree below $S$.
Ordering the rows by their $z$-degree shows that the change of basis is
triangular over $\mathbb Z$ with diagonal entries $1$, so it preserves
$\Delta$.

Every ordinary $z$-derivative of $z^p-w$, evaluated at
$(\alpha_i,\bar\alpha_i)$, lies in $pA$. For derivative order $0$ this
is \eqref{sz-frobenius}; for orders $1\le j\le p$ the coefficient
$p(p-1)\cdots(p-j+1)$ is divisible by $p$; for higher orders the derivative
is zero. In the product rule each of the $h_{r,s}$ factors therefore
contributes a factor $p$, even if some factors are differentiated.
The transformed row lies in $p^{h_{r,s}}A^N$. Factoring these powers
from the rows gives $\Delta\in p^{E_p}A$, and taking its norm gives
$p^{2E_p}\mid M$. Distinct primes may be multiplied, whence
\begin{equation}\label{sz-norm-lower}
 \log M\ge2\sum_p E_p\log p.
\end{equation}
The second formula for $E_p$ counts, for each $j\ge1$, the pairs
$(r,s)$ with $r\ge jp$ and $s\le S-1-j$.

We bound the determinant separately under each complex embedding.
An entry with $j\le r$ has modulus at most
$R^jH^{r+s-j}$, and an entry with $j>r$ is zero. Divide row $(r,s)$
by $H^{r+s}$ and column $(i,j)$ by $(R/H)^j$. Every resulting entry
has modulus at most $1$, so Hadamard's inequality bounds its determinant
by $N^{N/2}$. Moreover,
\[
 \sum_{r,s}(r+s)-\sum_{i,j}j=\frac{N(R-1)}2,
 \qquad \sum_{i,j}j=\frac{N(S-1)}2.
\]
Thus
\[
 |\Delta|\le N^{N/2}H^{N(R-1)/2}R^{N(S-1)/2}
\]
under either embedding. Multiplying the bounds for the two embeddings,
taking logarithms, and using \eqref{sz-norm-lower} proves the result.
\end{proof}

\begin{remark}[A number-field version of the finite inequality]
Proposition~\ref{sz-finite} only needs the following data: a number field
$K$ of degree $e$, an order $\mathcal A\subset K$, an automorphism
$\tau$ of $\mathcal A$, and elements $\alpha_1,\ldots,\alpha_R$.
For every prime $p$ to be used in the lower bound, assume
$\alpha_i^p-\tau(\alpha_i)\in p\mathcal A$ for every $i$.
Require the same Frobenius congruence at the auxiliary prime $R$,
and assume also that the images of the
$\alpha_i$ in some residue field of characteristic $R$ are distinct.
If all embeddings of the $\alpha_i$ and $\tau(\alpha_i)$ have modulus
at most $H$, the identical inequality \eqref{sz-finite-ineq} holds, with
the sum restricted to these primes. To prove this, replace
$\bar\alpha_i$ by $\tau(\alpha_i)$ in the matrix. Reduction to the
chosen residue field again makes the determinant nonzero. The row changes
give $\Delta\in p^{E_p}\mathcal A$, so
$p^{eE_p}$ divides the nonzero integer $N_{K/\mathbb Q}(\Delta)$.
The archimedean calculation gives the $e$th power of the bound for one
embedding. Multiplying the logarithmic comparison by $2/e$ gives exactly
\eqref{sz-finite-ineq}. No unramifiedness assertion is implicit in this
version; the stated congruences and residue-field hypothesis are required.
\end{remark}

\begin{remark}[Ordinary derivatives are essential]
One cannot replace every $\partial_z^j$ by the divided derivative
$\partial_z^j/j!$ in the divisibility proof. At order $p$ the divided
derivative of $z^p-w$ is $1$, which is not divisible by $p$.
The restriction $S\le R$ concerns the auxiliary prime proving nonvanishing;
the other inert primes can be much smaller than $S$. Thus invertibility
of factorials at $R$ does not justify dividing by them at all those primes.
\end{remark}

Here are the elementary prime estimates needed to pass from the finite
inequality to a uniform statement. All implied constants below are absolute
and effective.

\begin{lemma}\label{sz-mertens}
For $x\ge2$ and $y\ge2$,
\[
 \vartheta(x):=\sum_{p\le x}\log p\ll x,
 \qquad \sum_{p\le x}\frac{\log p}{p}=\log x+O(1),
 \qquad \sum_{p\ge y}\frac{\log p}{p^2}\ll y^{-1}.
\]
\end{lemma}
\begin{proof}
Every prime between $n$ and $2n$ divides $\binom{2n}{n}$, so
$\vartheta(2n)-\vartheta(n)\le2n\log2$. Summation on dyadic intervals
gives $\vartheta(x)\ll x$. Hence, writing $\Lambda$ for the von Mangoldt
function,
\[
 \psi(x):=\sum_{m\le x}\Lambda(m)
 =\sum_{k\ge1}\vartheta(x^{1/k})
 \ll x+\sqrt x\log(2x)\ll x.
\]
For an integer $n\ge2$, prime factorization of $n!$ gives
\[
 \log(n!)=\sum_{m\le n}\Lambda(m)\lfloor n/m\rfloor
 =n\sum_{m\le n}\frac{\Lambda(m)}m+O(\psi(n)).
\]
Integral comparison gives $\log(n!)=n\log n+O(n)$, so
$\sum_{m\le n}\Lambda(m)/m=\log n+O(1)$. The terms with $m=p^k$,
$k\ge2$, have bounded total, because
$\sum_{p,k\ge2}(\log p)/p^k\le\sum_{n\ge2}(\log n)/(n(n-1))<\infty$.
This proves the second estimate, also for real $x$ by taking its integer
part. The last estimate follows by partial summation from
$\vartheta(t)\ll t$.
\end{proof}

\begin{theorem}[Uniform bound for inert primes]\label{sz-inert}
For every nonsquare integer $d$ and every $X\ge2$,
\begin{equation}\label{sz-inert-bound}
 I_d(X):=\sum_{\substack{p\le X\ p\ {
m odd}\ (d/p)=-1}}
            \frac{\log p}{p}
 \le\frac12\log X+\frac14\log|d|+\log\log(3X)+O(1).
\end{equation}
The implied constant is independent of both $d$ and $X$.
\end{theorem}
\begin{proof}
If there is no odd inert prime below $X$, the assertion is immediate.
Otherwise let $R$ be the largest such prime. Then $I_d(X)=I_d(R)$.
If $R<e^3$, the result is immediate with an absolute constant. If
$R\le\sqrt{|d|}$, Lemma~\ref{sz-mertens} gives
\[
 I_d(R)\le\log R+O(1)
 \le\tfrac12\log R+\tfrac14\log|d|+O(1).
\]
We may therefore assume $R\ge e^3$ and $R>\sqrt{|d|}$.

Set $L=\log R$, $S=\lceil R/L\rceil$, and $N=RS$. Then $1\le S<R$.
Consider the quadratic integers $u+v\sqrt d$ with
\[
 0\le u<\sqrt R\,|d|^{1/4},
 \qquad0\le v<\sqrt R\,|d|^{-1/4},\qquad u,v\in\mathbb Z.
\]
There are at least $R$ such integers because the two side lengths have
product $R$. Each coordinate lies in $[0,R)$, so all reductions are distinct
in $A/RA$. Select any $R$ of them. Both embeddings have modulus below
$H=2\sqrt R\,|d|^{1/4}<2R$. Apply Proposition~\ref{sz-finite},
and let $M$ be the positive determinant norm constructed in its proof.

For $a\in\mathbb Z_{\ge0}$ one has
$\sum_{s<S}\min\{a,S-1-s\}\ge Sa-a^2$.
For $a<S$ this follows by summing the arithmetic progression; for $a\ge S$
the right side is nonpositive. Consequently
\begin{align*}
 E_p&\ge S\sum_{r<R}\lfloor r/p\rfloor-
              \sum_{r<R}\lfloor r/p\rfloor^2\\
 &\ge\frac{SR(R-1)}{2p}-SR-\frac{R^3}{3p^2}.
\end{align*}
Sum this inequality only over inert primes with $L\le p\le R$.
By Lemma~\ref{sz-mertens}, and because $S\ge R/L$,
\[
 2\sum_pE_p\log p
 \ge NR\sum_{\substack{L\le p\le R\\(d/p)=-1}}
           \frac{\log p}{p}-O(NR).
\]
Indeed, replacing $R-1$ by $R$ costs $O(N\log R)$, the floor errors
cost $O(N\vartheta(R))$, and the quadratic error is
$O(R^3/L)=O(NR)$.
On the other hand, \eqref{sz-finite-ineq} gives
\[
 \frac{\log M}{NR}
 \le\frac{\log N}{R}+\left(1-\frac1R\right)\log H
       +\frac{S-1}{R}\log R
 \le\log H+O(1).
\]
Here $(S-1)\log R/R\le1$ and $N\le R^2$.
Restore the primes below $L$, whose contribution is at most
$\log L+O(1)$ by Lemma~\ref{sz-mertens}. The result is
\[
 I_d(R)\le\log H+\log L+O(1)
 =\tfrac12\log R+\tfrac14\log|d|+\log\log R+O(1).
\]
Finally $R\le X$ proves \eqref{sz-inert-bound}.
\end{proof}

The parameter choice can be checked directly instead of guessed. The same
calculation with any integer $1\le S\le R/2$, and cutoff $R/S$ for the
inert primes, gives
\[
 I_d(R)\le\log H+\frac SR\log R+\log(R/S)+O(1).
\]
Writing $u=S/R$, the varying part of this upper bound is
$u\log R-\log u$. Its derivative is $\log R-1/u$, so its minimum occurs
at $u=1/\log R$. This recovers the choice above. Increasing the number of
derivatives arbitrarily does not improve this estimate: it increases the
term $u\log R$. This calculation describes a limitation of this particular
determinant bound, and is not a claim of optimality among all possible
arithmetic constructions.

\subsection{The analytic consequence and its exact scope}

A nonprincipal real primitive Dirichlet character is quadratic:
there is a fundamental discriminant $D$ with $|D|=q$ for which
$\chi(n)=(D/n)$, the Kronecker symbol. A fundamental discriminant is
the discriminant of a quadratic number field. In particular, there is
a squarefree $d\ne1$, $|d|\le q$, such that $\chi(p)=(d/p)$ at every
odd prime. Thus Theorem~\ref{sz-inert} applies to the inert primes
of each real primitive character.

The following use of complex analysis is explicit. We use the classical
analytic continuation, functional equation, and Hadamard product for a
primitive Dirichlet $L$-function. If $a=(1-\chi(-1))/2$, these give, for
real $\sigma>1$ and real primitive nonprincipal $\chi$,
\begin{equation}\label{sz-hadamard}
 \sum_\rho\Re\frac1{\sigma-\rho}
 =\frac12\log(q/\pi)
   +\frac12\frac{\Gamma'}{\Gamma}\left(\frac{\sigma+a}{2}\right)
   +\frac{L'}{L}(\sigma,\chi).
\end{equation}
The zeros here are the nontrivial zeros, with multiplicity; all summands
are nonnegative and their sum converges. The same identity holds at
$\sigma=1$ when $L(1,\chi)\ne0$. These are classical analytic inputs,
not consequences of the determinant calculation. They are also the inputs
used in the last step of \cite{sz-openai,sz-pescador,sz-vandijk}.

\begin{corollary}\label{sz-zero-gap}
Every real zero $\beta\in(0,1)$ of a real primitive nonprincipal character
of conductor $q$ satisfies
\[
 (1-\beta)\log q\ge1-O\left(\frac{\log\log(3q)}{\log q}\right).
\]
In particular there is an effective absolute constant $c>0$ such that
$(1-\beta)\log q\ge c$ for all these characters and zeros.
\end{corollary}
\begin{proof}
For $0<\eta\le1$, partial summation in Theorem~\ref{sz-inert}, including
the harmless prime $2$, gives
\begin{equation}\label{sz-laplace}
 \sum_{\chi(p)=-1}\frac{\log p}{p^{1+\eta}}
 \le\frac1{2\eta}+\frac14\log q+O(\log(2/\eta)).
\end{equation}
For completeness, the sum is
$\eta\int_1^\infty S_-(t)t^{-1-\eta}\,dt$, where
$S_-(t)=\sum_{p\le t,\chi(p)=-1}(\log p)/p$.
The integrals of $1$ and $\log t$ against this probability measure are
$1$ and $1/\eta$. With $u=\log t$, concavity of the logarithm bounds
the integral of $\log(\log3+u)$ by $\log(\log3+1/\eta)$.

The Euler products, with all prime powers of exponent at least two
contributing $O(1)$, give
\begin{align*}
 \frac{L'}{L}(1+\eta,\chi)
 &=2\sum_{\chi(p)=-1}\frac{\log p}{p^{1+\eta}}
     -\sum_p\frac{\log p}{p^{1+\eta}}
     +\sum_{p\mid q}\frac{\log p}{p^{1+\eta}}+O(1)\\
 &\le\frac12\log q+O(\log(2/\eta)+\log\log(3q)).
\end{align*}
Here $\sum_p(\log p)/p^{1+\eta}=1/\eta+O(1)$ follows by partial
summation from Lemma~\ref{sz-mertens}. Also
$\sum_{p\mid q}(\log p)/p\ll\log\log(3q)$: split at $p=\log q$,
use Lemma~\ref{sz-mertens} below that point, and use
$\sum_{p\mid q}\log p\le\log q$ above it.

Retain only $\rho=\beta$ in \eqref{sz-hadamard}. The gamma term is
bounded for $1\le\sigma\le2$, uniformly in the parity, and hence
\[
 \frac1{1-\beta+\eta}
 \le\log q+O(\log(2/\eta)+\log\log(3q)).
\]
For large $q$, take $\eta=(\log q)^{-2}$ and invert the positive
right side. This gives the asserted asymptotic inequality and a positive
uniform constant for all sufficiently large $q$.

To see that the extension to the finitely many smaller conductors is
effective, use the classical theorem $L(1,\chi)\ne0$. Periodicity gives
$|\sum_{u<n\le v}\chi(n)|\le q/2$ and hence an effective tail bound
$q/(2N)$ for $\sum_{n>N}\chi(n)/n$. Computing partial sums therefore
eventually certifies a positive lower bound for $|L(1,\chi)|$.
Summation by parts also gives
$|L'(s,\chi)|\le 2q/e$ for $1/2\le s\le1$: the total variation of
$(\log t)t^{-s}$ on $[1,\infty)$ is $2/(es)$, and the partial sums
are bounded by $q/2$. The mean value inequality supplies a positive
zero-free interval at $1$. Taking the minimum over the finite list and
also over $\tfrac12\log q$ for zeros below $1/2$ gives an effective $c$.
\end{proof}

The conclusion is a logarithmic gap from $1$, not a half-plane with boundary
$1/2$. Neither a nonreal zero with real part $3/4$ nor a real zero bounded
away from $1$ contradicts the inequalities just proved. For example, the
bound obtained from \eqref{sz-hadamard} at a hypothetical real zero
$\beta=3/4$ compares a number tending to $4$ with a quantity of size
$\log q$. It cannot exclude such a zero as $q$ grows. This observation
identifies the exact limitation of this method; it does not assert that
any Dirichlet $L$-function has such a zero.

\section{Pairing two character tails before estimating them}
\label{sz-tails-section}

There is another useful simplification in Guo's paper \cite{sz-guo}.
It avoids a spurious factor $(1-\beta)^{-1}$ in a hyperbola estimate.
The next version applies to any bounded sequence with bounded interval sums
and records the best constant supplied by this particular pairing.

\begin{proposition}\label{sz-paired-tails}
Let $c:\mathbb N\to\mathbb C$ satisfy $|c(n)|\le1$ and
\[
 \left|\sum_{u<n\le v}c(n)\right|\le B\qquad(0\le u\le v).
\]
For real $s>0$, set $L_c(s)=\sum_{n\ge1}c(n)n^{-s}$, which converges
by summation by parts. Suppose $0<\beta<1$ and $L_c(\beta)=0$.
Put $\varepsilon=1-\beta$ and $a(n)=\sum_{d\mid n}c(d)$.
Then, for every $t\ge1$,
\begin{align}\label{sz-paired-main}
 \sum_{n\le t}\frac{a(n)}{n^\beta}
 &=\frac{L_c(1)}{\varepsilon}t^\varepsilon+E(t),\\
 |E(t)|&\le
 \left[B\bigl(1+2\beta^{\beta/(1-\beta)}\bigr)
                  +1+\tfrac12\log t\right]t^{1/2-\beta}.\nonumber
\end{align}
For $\beta\ge1/2$, the bracket can be replaced by
$2B+1+\tfrac12\log t$.
\end{proposition}
\begin{proof}
Write $A_u(v)=\sum_{u<n\le v}c(n)$. Summation by parts gives
\[
 R_s(u):=\sum_{n>u}\frac{c(n)}{n^s}
      =s\int_u^\infty A_u(v)v^{-s-1}\,dv,
 \qquad |R_s(u)|\le Bu^{-s}.
\]
Set $u=\sqrt t$. Splitting the pairs $dm\le t$ according to $d\le u$
or $d>u$ gives
\[
 \sum_{n\le t}\frac{a(n)}{n^\beta}
 =\sum_{d\le u}d^{-\beta}\sum_{m\le t/d}\frac{c(m)}{m^\beta}
  +\sum_{m\le u}\frac{c(m)}{m^\beta}
            \sum_{u<d\le t/m}d^{-\beta}.
\]
The first sum has modulus at most $Bt^{1/2-\beta}$, because its inner
sum is $-R_\beta(t/d)$. For real $v\ge u\ge1$, integral summation gives
\[
 \sum_{u<n\le v}n^{-\beta}
 =\frac{v^\varepsilon-u^\varepsilon}{\varepsilon}+r(u,v),
 \qquad |r(u,v)|\le u^{-\beta}.
\]
The error contributed to the second sum is therefore at most
\[
 u^{-\beta}\sum_{m\le u}m^{-\beta}
 \le u^{\varepsilon-\beta}(1+\log u).
\]
Its two main terms are
\[
 \frac{L_c(1)}\varepsilon t^\varepsilon
 -\frac{t^\varepsilon}\varepsilon R_1(u)
 +\frac{u^\varepsilon}\varepsilon R_\beta(u).
\]
Pair the last two terms before taking absolute values. Their sum is
\[
 u^\varepsilon\sum_{m>u}c(m)g(m),\qquad
 g(v)=\frac{v^{-\beta}-u^\varepsilon v^{-1}}\varepsilon.
\]
The function $g$ is nonnegative, vanishes at $u$ and at infinity, and
increases and then decreases, with maximum at $v=u\beta^{-1/\varepsilon}$.
A direct substitution gives
$\max g=u^{-\beta}\beta^{\beta/\varepsilon}$.
Consequently its total variation is twice this maximum. A final summation
by parts bounds the paired terms by
\[
 2B\beta^{\beta/\varepsilon}u^{\varepsilon-\beta}.
\]
Adding the three errors gives \eqref{sz-paired-main}. Finally
$\beta^{\beta/(1-\beta)}$ decreases from $1/2$ to $e^{-1}$ on
$[1/2,1)$, since the derivative of its logarithm is
$(\log\beta+1-\beta)/(1-\beta)^2\le0$.
\end{proof}

For a nonprincipal real character, one may take $B=q/2$ by deleting complete
periods from an interval and comparing the remaining block with its
complement. In that case $a=1*\chi$ is nonnegative: at a prime power,
\[
 a(p^k)=
 \begin{cases}
 k+1,&\chi(p)=1,\\
 1,&\chi(p)=0,\\
 (1+(-1)^k)/2,&\chi(p)=-1.
 \end{cases}
\]
This is why the following comparison applies to real characters.

\begin{lemma}[A weighted positive-sum comparison]\label{sz-positive}
Suppose that $\chi$ is a nonprincipal real character,
$1/2\le\beta<1$, and $L(\beta,\chi)=0$. Put $\varepsilon=1-\beta$,
\[
 H_1(y)=\sum_{n\le y}\frac{(1*\chi)(n)}n,\quad
 B_0(t)=(q+1+\tfrac12\log t)t^{-1/2},\quad
 K=\log x\log y+\tfrac12(\log x)^2.
\]
If $1\le y\le x$ and $\varepsilon\log(xy)\le1$, then
\begin{align*}
 \sum_{y<p\le x}\frac{(1+\chi(p))\log p}{p}
 \le\varepsilon K+\frac{1}{H_1(y)}
 \bigl[B_0(y)(\varepsilon K+\log y)
       +B_0(xy)\log(xy)+2B_0(y)+2y^{-1/2}\bigr].
\end{align*}
\end{lemma}
\begin{proof}
Write $a=1*\chi$ and $N(t)=\sum_{n\le t}a(n)(\log n)/n$.
The map $(n,p)\mapsto np$, with $n\le y<p\le x$, is injective, since
equality $np=n'p'$ with distinct primes would force $p\mid n'$, although
$n'<p$. Multiplicativity and nonnegativity give
\[
 H_1(y)\sum_{y<p\le x}\frac{a(p)\log p}{p}\le N(xy)-N(y).
\]
Use Proposition~\ref{sz-paired-tails} with $c=\chi$ and put
$D(t)=t^{-\varepsilon}E(t)$; thus $|D(t)|\le B_0(t)$.
Partial summation, applied to $(\log t)t^{-\varepsilon}$, gives exactly
\begin{align*}
 N(xy)-N(y)&=L(1,\chi)K+R,\\
 R&=D(xy)\log(xy)-D(y)\log y
      +\int_y^{xy}(\varepsilon\log t-1)D(t)\frac{dt}{t}.
\end{align*}
Because $0\le\varepsilon\log t\le1$ on the interval,
\[
 |R|\le B_0(xy)\log(xy)+B_0(y)\log y
             +\int_y^\infty B_0(t)\frac{dt}{t}.
\]
The last integral is $2B_0(y)+2y^{-1/2}$. Also, positivity gives
\[
 \frac{L(1,\chi)}\varepsilon
 =y^{-\varepsilon}\left(\sum_{n\le y}\frac{a(n)}{n^\beta}-E(y)\right)
 \le H_1(y)+B_0(y).
\]
Substitute this inequality and divide by $H_1(y)\ge a(1)=1$.
\end{proof}

The pairing is useful whenever a proof needs a uniform error as a real zero
approaches $1$. It does not produce high multiplicity for an automorphic
theta lift, nor does positivity of $1*\chi$ extend to arbitrary complex
characters. Those are distinct issues in a proof of the full generalized
Riemann hypothesis.

\section{What these results simplify in the auxiliary choices}
\label{sz-integration}

The original construction chooses auxiliary primes of size
\[
 \log q\asymp x^{\epsilon/5},\qquad
 \log d\asymp x^{\epsilon/5}(\log x)^2,
\]
and imposes, among other requirements, a bound of the form
$|L'(1,\chi_{-d})/L(1,\chi_{-d})|\le x^\epsilon$.
The following deduction explains exactly how uniform exclusion of real
zeros helps with that condition.

\begin{proposition}\label{sz-log-derivative}
Assume the classical zero-free region for primitive Dirichlet $L$-functions,
with its possible simple real zero, and the uniform real-zero gap of
Corollary~\ref{sz-zero-gap}. Then, for every real primitive nonprincipal
character of conductor $q\ge3$,
\[
 \left|\frac{L'}{L}(1,\chi)\right|\ll\log q
\]
with an effective absolute implied constant.
\end{proposition}
\begin{proof}
The two assumptions imply an absolute $c_0>0$ such that every nontrivial
zero $\rho=b+i\gamma$ with $|\gamma|\le1$ satisfies
$1-b\ge c_0/\log(3q)$. Put $h=1/\log q$ and $\sigma_0=1+h$.
For $a=1-b>0$,
\[
 \frac{\Re(1-\rho)^{-1}}{\Re(\sigma_0-\rho)^{-1}}
 =\frac{a}{a+h}\frac{(a+h)^2+\gamma^2}{a^2+\gamma^2}.
\]
If $|\gamma|\le1$, this is at most $(a+h)/a\ll1$ by the gap.
If $|\gamma|>1$, use $0<a<1$ and $0<h<1$ to bound the ratio by $4$.
All terms are nonnegative, so summing gives
\[
 \sum_\rho\Re(1-\rho)^{-1}
 \ll\sum_\rho\Re(\sigma_0-\rho)^{-1}.
\]
By \eqref{sz-hadamard}, the sum on the right is at most
$\tfrac12\log q+|L'/L(\sigma_0,\chi)|+O(1)$.
The Euler product and $|\chi(n)|\le1$ give
\[
 |L'/L(\sigma_0,\chi)|
 \le\sum_{n\ge1}\Lambda(n)n^{-\sigma_0}
 =\frac1h+O(1)\ll\log q,
\]
where the equality also follows from the proof of
Lemma~\ref{sz-mertens}. The Hadamard formula at $1$ now bounds
$L'/L(1,\chi)$ above by $O(\log q)$ and below by
$-\tfrac12\log q-O(1)$. This derivative is real, proving the claim.
\end{proof}

Consequently, for every fixed $\epsilon>0$, the auxiliary derivative bound
holds for every chosen prime $d$ satisfying
$\log d\le x^{\epsilon/5}(\log x)^2$, once $x$ is sufficiently large.
Indeed,
\[
 \frac{x^{\epsilon/5}(\log x)^2}{x^\epsilon}
 =x^{-4\epsilon/5}(\log x)^2\longrightarrow0.
\]
No exceptional-character avoidance is needed solely for this derivative
condition. This is a proved simplification of one requirement, rather than
a replacement of all the auxiliary-prime arguments.

There is a second, stronger simplification if the all-character
quasi-Riemann theorem is used. The source calls a primitive real character
of conductor at most $Q$ good when it has no zero with
\[
 \Re s>399/400,\qquad |\Im s|\le Q.
\]
A zero-free half-plane $\Re s>7/8$ makes this condition automatic,
uniformly in $Q$. The corresponding exclusion of bad characters can then
be dropped, after retaining the conductor bounds in each application.
Uniform exclusion of Landau--Siegel zeros by itself does not imply this
fixed-width half-plane condition.

Two substantial restrictions remain. First, the source also uses ray-class
$L$-functions over a varying imaginary quadratic field. A theorem stated
only for Dirichlet characters over $\mathbb Q$ does not immediately remove
an exceptional-zero term for those Hecke $L$-functions; a separate descent
or zero-free theorem with the required conductor uniformity would have to
be supplied. Second, the construction needs square-root bounds for selected
smoothed character sums, including a fixed finite family of character
twists for GRH. A bound with exponent $7/8$ does not directly establish
those bounds for every fixed character. A separate second-moment argument
over the auxiliary primes does establish the required simultaneous
selection under a uniform QRH hypothesis; that argument is developed in
the continuation. Neither the determinant theorem nor the character-tail
pairing above proves the estimates for a growing chain of theta lifts
which supply the missing square-root exponent in the proposed GRH argument.


% Optional mathematical section for the reconstruction manuscript.
% Requires amsmath, amssymb, amsthm and the standard theorem environments.
% Every assertion proved here is independent of the unverified construction
% of the particular automorphic lattices in the proposed GRH proof.
\section{A shorter slope comparison for a family of integral maps}
\label{cs-section}

The final arithmetic comparison in the proposed argument\cite{alpo-central-comparison} has two distinct
parts. One constructs integral maps, controls their heights and local
divisibility, and proves that their complex realizations are nearly isotropic.
The other turns these facts into an inequality between slopes. We give a
short proof of the second part. Its hypotheses are concrete; establishing
them for the required automorphic family remains a separate problem.

\subsection{Slopes and the height of a map}

Let $K$ be a number field of degree $d$. A Hermitian $\mathcal O_K$-lattice
$E$ is a finite projective $\mathcal O_K$-module with a positive definite
Hermitian norm at each complex embedding, compatible with conjugation.
Its rank is its dimension over $K$. Give its real Minkowski space the
metric in which a complex place contributes twice the real part of the
Hermitian form. Write $\operatorname{covol}(E)$ for the resulting real
covolume and put
\[
 \widehat{\deg}(E)
 =\frac1d\left(-\log\operatorname{covol}(E)
                  +\frac{\operatorname{rk}E}{2}\log|\Delta_K|\right),
 \qquad \mu(E)=\frac{\widehat{\deg}(E)}{\operatorname{rk}E}.
\]
Submodules are given the restricted metrics and quotients the quotient
metrics. A submodule is called saturated if its quotient is torsion free.
For a nonzero integral linear map $f:E\to F$, its height is
\[
 h(f)=\frac1d\sum_{v\mid\infty}n_v\log\|f_v\|_{\rm op},
 \qquad n_v=[K_v:\mathbb R].
\]
We allow the zero map and assign it height $-\infty$.

We recall the slope facts that will be used. Every nonzero $E$ has a
filtration by saturated submodules
\[
 0=E_0\subset E_1\subset\cdots\subset E_a=E
\]
whose successive quotients are semistable and have strictly decreasing
slopes. Semistable means that no nonzero saturated submodule has larger
slope than the whole lattice. List the successive slopes with their
rank multiplicities as
$\lambda_1(E)\ge\cdots\ge\lambda_r(E)$.
Then their mean is $\mu(E)$; write
$\mu_{\max}(E)=\lambda_1(E)$ and
$\mu_{\min}(E)=\lambda_r(E)$. For a real $t$, let $E_{\ge t}$ denote the
step consisting of all successive quotients with slope at least $t$.

For completeness, the construction uses only covolumes. In each fixed
rank, the degrees of saturated submodules are bounded above and attain
their maximum: their underlying primitive integral sublattices have
nonzero exterior products in a fixed discrete lattice, and there are
only finitely many such products of bounded norm. Choose a submodule
of maximal slope and, among those, maximal rank. It is unique. Indeed,
for two saturated submodules $U,V$ the determinant formula for an exact
sequence and the decrease of quotient norms give
\[
 \widehat{\deg}U+\widehat{\deg}V
 \le\widehat{\deg}(U\cap V)
       +\widehat{\deg}\bigl((U+V)^{\rm sat}\bigr).
\]
If $U,V$ have maximal slope, both terms on the right are at most that
slope times their respective ranks; the ranks add to
$\operatorname{rk}U+\operatorname{rk}V$. Thus their saturated sum also has
maximal slope. Maximality of the chosen rank proves uniqueness.
Repeat in the quotient. The slopes strictly decrease, since otherwise
the inverse image of the next maximal submodule contradicts the choice
at the preceding step. The process terminates after at most $r$ steps.
Degree additivity proves the assertion about the mean. It also shows
that every nonzero torsion-free quotient, with its quotient metrics,
has slope at least $\mu_{\min}(E)$ and every
nonzero submodule has slope at most $\mu_{\max}(E)$: intersect a submodule,
or filter a quotient, by the displayed filtration, and use semistability
on its successive images. Saturation and quotient norms can only improve
the needed inequalities.

\begin{lemma}[Slopes under an integral map]\label{cs-functoriality}
If $f:E\to F$ is nonzero, then
\[
 \mu_{\min}(E)\le\mu_{\max}(F)+h(f).
\]
If $h(f)\le0$, then $f(E_{\ge t})\subset F_{\ge t}$ for every real $t$.
\end{lemma}
\begin{proof}
Let $Q=E/\ker f$ and let $J$ be the saturation of $f(E)$ in $F$.
Put $q=\operatorname{rk}Q=\operatorname{rk}J$. On $q$th exterior powers
the operator norm is at most $\|f_v\|_{\rm op}^{q}$.
The image lattice has finite index in $J$; its index contributes a
nonnegative logarithm to its covolume. Multiplying the determinant
inequalities at the infinite places therefore gives
\[
 \widehat{\deg}(Q)\le\widehat{\deg}(J)+q h(f).
\]
Divide by $q$ and use the preceding submodule and quotient properties.
For the second assertion, compose the restriction to $E_{\ge t}$ with
$F\to F/F_{\ge t}$. Restriction and orthogonal quotient do not increase
operator norms. If this composite were nonzero, the first assertion
would give $t\le\mu_{\min}(E_{\ge t})
\le\mu_{\max}(F/F_{\ge t})<t$, a contradiction.
\end{proof}

\subsection{The matrix estimate}

For a complex matrix $A$, put
$\|A\|_{\mathrm{HS}}^2=\operatorname{tr}(A^*A)$.
An orthogonal projection always refers to the ordinary Euclidean inner
product, independently of the arithmetic metrics above.

\begin{lemma}[Images of subspaces]\label{cs-dimension}
Let $A_i:\mathbb C^{d_i}\to\mathbb C^D$ be a finite family of matrices,
where $D\ge1$ and $R=\sum_i d_i>0$.
Set $Q=\sum_i A_iA_i^*$ and suppose that
\begin{align}
 \sum_i\|A_i^*A_i-I\|_{\mathrm{HS}}^2&\le eR,
       &0\le e\le\tfrac14,\label{cs-diagonal}\\
 \operatorname{tr}(Q^2)&\le
       \frac{1+\eta}{D}(\operatorname{tr}Q)^2,
       &\eta\ge0.\label{cs-isotropic}
\end{align}
If $V_i\subset\mathbb C^{d_i}$ and all $A_i(V_i)$ lie in a subspace
$W\subset\mathbb C^D$, put $t=\sum_i\dim V_i$ and $r=\dim W$. Then
\begin{equation}\label{cs-rank-bound}
 \frac rD\ge
 \frac{(t/R-\sqrt e)_+^2}{(1+\eta)(1+\sqrt e)^2},
 \qquad z_+=\max\{z,0\}.
\end{equation}
\end{lemma}
\begin{proof}
Let $P_i$ project onto $V_i$ and put $C=\sum_i A_iP_iA_i^*$.
Then $0\le C\le Q$ and $\operatorname{rank}C\le r$. Cauchy--Schwarz gives
\[
 |\operatorname{tr}Q-R|\le\sqrt{R}
     \left(\sum_i\|A_i^*A_i-I\|_{\mathrm{HS}}^2\right)^{1/2}
 \le\sqrt e\,R
\]
and, using $\sum_i\|P_i\|_{\mathrm{HS}}^2=t$,
\[
 \operatorname{tr}C
 =t+\sum_i\operatorname{tr}\bigl(P_i(A_i^*A_i-I)\bigr)
 \ge t-\sqrt{eRt}\ge t-\sqrt e\,R.
\]
The eigenvalues of $C$ are nonnegative. Thus
$(\operatorname{tr}C)^2\le r\operatorname{tr}(C^2)$.
Moreover
$\operatorname{tr}(Q^2-C^2)
=\operatorname{tr}((Q-C)(Q+C))\ge0$, since the trace of a product
of two positive semidefinite matrices is nonnegative. Combine these
facts with \eqref{cs-isotropic} and
$\operatorname{tr}Q\le(1+\sqrt e)R$. If $t/R\le\sqrt e$, the desired
inequality is simply $r\ge0$; otherwise division gives
\eqref{cs-rank-bound}.
\end{proof}

\begin{proposition}[Slope comparison]\label{cs-slope}
Let $I$ be a finite nonempty index set. Let $E_i$ for $i\in I$, and $F$,
be Hermitian $\mathcal O_K$-lattices with positive ranks
$d_i,n$, put $R=\sum_{i\in I}d_i$, and let $f_i:E_i\to F$ be integral linear maps with
$h(f_i)\le h_i\in\mathbb R$. Suppose that
\begin{equation}\label{cs-numerical}
 \frac1R\sum_i d_i\bigl(\mu(E_i)-h_i\bigr)\ge s,
 \quad \mu_{\max}(E_i)-h_i\le s+u,
 \quad \mu_{\min}(F)\ge s-H,
 \qquad u,H\ge0.
\end{equation}
Fix one complex embedding of $K$. Suppose that $D\ge n$,
$n\ge(1-\delta)D$, and $0\le\delta<1$. Suppose that matrices
$A_i:\mathbb C^{d_i}\to\mathbb C^D$ satisfy
\eqref{cs-diagonal}--\eqref{cs-isotropic}, and that their first $n$ rows
represent the maps $f_i$ in some complex bases of $E_i$ and a single
common basis of $F$. Positive scalar factors depending on $i$ are allowed.
Then
\begin{equation}\label{cs-slope-bound}
 s-\mu(F)\le\frac{2u\log(1+H/u)+(\eta+4\sqrt e)H}{1-\delta}.
\end{equation}
At $u=0$ the expression $u\log(1+H/u)$ means its continuous value $0$.
\end{proposition}
\begin{proof}
Multiply every norm on $E_i$ by $\exp(h_i)$. All of its slopes decrease
by $h_i$, and the height of $f_i$ decreases by $h_i$. Thus the maps now
have nonpositive height. By Lemma~\ref{cs-functoriality}, the source
subspaces with slopes at least $s-z$ map into $F_{\ge s-z}$.

For $z>0$, let $t(z)$ be the sum of the ranks of these source subspaces.
Every adjusted source slope is at most $s+u$, and their weighted mean is
at least $s$. The total of their nonnegative deficits from $s+u$ is at
most $uR$. Each slope smaller than $s-z$ has deficit larger than $u+z$.
Consequently
\[
 a(z):=t(z)/R\ge\frac z{u+z}.
\]
Write $r(z)=\operatorname{rk}F_{\ge s-z}$. The images in $\mathbb C^D$
lie in the direct sum of this $r(z)$-dimensional target subspace and the
last $D-n$ coordinate directions. Lemma~\ref{cs-dimension} gives
\[
 1-\frac{r(z)}n\le\frac1{1-\delta}
 \left(1-\frac{(a(z)-\sqrt e)_+^2}
 {(1+\eta)(1+\sqrt e)^2}\right).
\]
For $0\le a\le1$, $q\ge0$, and $\eta\ge0$ one has
\[
 1-\frac{(a-q)_+^2}{(1+\eta)(1+q)^2}
 \le\eta+4q+2(1-a).
\]
To see this, put $b=(a-q)_+/(1+q)$. Then $0\le b\le a\le1$,
$a-b\le2q$, and
$1-b^2/(1+\eta)\le\eta+(1-a^2)+(a^2-b^2)
\le\eta+2(1-a)+4q$.
It follows that
\[
 1-r(z)/n\le\frac{\eta+4\sqrt e+2u/(u+z)}{1-\delta}.
\]
Every target slope is at least $s-H$. If these slopes are $\nu_j$,
integration of the last inequality therefore yields
\begin{align*}
 s-\mu(F)&\le\frac1n\sum_j(s-\nu_j)_+
   =\int_0^H\left(1-\frac{r(z)}n\right)\,dz\\
 &\le\frac{(\eta+4\sqrt e)H+2u\log(1+H/u)}{1-\delta}.
\end{align*}
When $u=0$, the source dimension is $R$ for every $z>0$ and the same
proof directly gives the stated limiting formula.
\end{proof}

\begin{lemma}[Pairwise estimates suffice]\label{cs-pairwise}
Let $A_i:\mathbb C^{d_i}\to\mathbb C^D$, $D\ge1$, be a finite family,
$R=\sum_i d_i>0$, and $Q=\sum_iA_iA_i^*$.
Suppose $D_*>0$, $\eta_0\ge-1$, $\eta_1\ge0$, $0\le e\le1/4$, and
\[
 \|A_i^*A_j\|_{\mathrm{HS}}^2
       \le(1+\eta_0)d_id_j/D_*\quad(i\ne j),\qquad
 \|A_i^*A_i-I\|_{\mathrm{HS}}^2\le e d_i,
 \qquad D\le(1+\eta_1)D_*.
\]
Then \eqref{cs-diagonal} holds and
\[
 \frac{D\operatorname{tr}(Q^2)}{(\operatorname{tr}Q)^2}
 \le
 \frac{(1+\eta_1)
   \bigl(1+\eta_0+(D_*/R)(1+\sqrt e)^2\bigr)}{(1-\sqrt e)^2}.
\]
Thus any nonnegative $\eta$ at least the right side minus $1$ is
permitted in Proposition~\ref{cs-slope}.
\end{lemma}
\begin{proof}
Cyclicity of trace gives
$\operatorname{tr}(Q^2)=\sum_{i,j}\|A_i^*A_j\|_{\mathrm{HS}}^2$.
The off-diagonal terms sum to at most $(1+\eta_0)R^2/D_*$.
The triangle inequality gives
$\|A_i^*A_i\|_{\mathrm{HS}}\le(1+\sqrt e)\sqrt{d_i}$.
Therefore
\[
 \operatorname{tr}(Q^2)
 \le(1+\eta_0)R^2/D_*+(1+\sqrt e)^2R.
\]
The proof of Lemma~\ref{cs-dimension} gives
$\operatorname{tr}Q\ge(1-\sqrt e)R>0$. Division proves the claim.
\end{proof}

\begin{lemma}[Restriction to subspaces]\label{cs-compression}
Let $B_i:\mathbb C^{r_i}\to\mathbb C^D$, $D\ge1$, be a finite family satisfying the pairwise and
diagonal inequalities of Lemma~\ref{cs-pairwise}, with parameters
$D_*,\eta_0,\eta_1,e_0$. Let
$V_i:\mathbb C^{d_i}\to\mathbb C^{r_i}$ satisfy $V_i^*V_i=I$,
and put $A_i=B_iV_i$, $R_f=\sum_i r_i>0$, $R=\sum_i d_i$.
Suppose $R_f-R\le\theta R_f$, where $0\le\theta<1$, and set
\[
 q=\sqrt{e_0},\qquad
 \kappa=1-\theta-q>0.
\]
Then $R>0$, and the compressed matrices satisfy
\[
 \sum_i\|A_i^*A_i-I\|_{\mathrm{HS}}^2
       \le\frac{e_0}{1-\theta}R,
\]
and, for $Q=\sum_iA_iA_i^*$,
\[
 \frac{D\operatorname{tr}(Q^2)}{(\operatorname{tr}Q)^2}
 \le\frac{(1+\eta_1)
 [1+\eta_0+(D_*/R_f)(1+q)^2]}{\kappa^2}.
\]
After omitting zero-dimensional domains, these are the matrix hypotheses
of Proposition~\ref{cs-slope} whenever
$e_0/(1-\theta)\le1/4$, with the displayed bound determining $1+\eta$.
\end{lemma}
\begin{proof}
Orthogonal compression decreases the Hilbert--Schmidt norm, so
\[
 \|A_i^*A_i-I\|_{\mathrm{HS}}
 =\|V_i^*(B_i^*B_i-I)V_i\|_{\mathrm{HS}}
 \le\|B_i^*B_i-I\|_{\mathrm{HS}}.
\]
This follows also by extending the columns of $V_i$ to an orthonormal
basis and summing the squares of matrix entries. Summing and using
$R\ge(1-\theta)R_f>0$ proves the first assertion.

Put $Q_f=\sum_iB_iB_i^*$ and let $P_i=V_iV_i^*$.
Then $0\le Q\le Q_f$, and $\sum_i\operatorname{rank}P_i=R$.
Cauchy--Schwarz gives
\begin{align*}
 \operatorname{tr}Q
 &=R+\sum_i\operatorname{tr}\bigl(P_i(B_i^*B_i-I)\bigr)\\
 &\ge R-\sqrt{e_0R_fR}\ge(1-\theta-q)R_f.
\end{align*}
Thus $\operatorname{tr}Q\ge\kappa R_f>0$.
As in Lemma~\ref{cs-dimension},
$\operatorname{tr}(Q^2)\le\operatorname{tr}(Q_f^2)$.
Use the upper bound for this last trace from the proof of
Lemma~\ref{cs-pairwise} and divide by $\kappa^2R_f^2$.
\end{proof}

In the source application, $D_*/R_f\le\sigma$,
$\eta_0,\eta_1,\sqrt{e_0}\le\sigma/2$, $0<\sigma\le1/20$,
and $0\le\theta\le1/4$; a negative $\eta_1$ can harmlessly be replaced
by $0$. These imply $\kappa>0$ and $e_0/(1-\theta)<1/4$.
In fact the displayed trace-ratio bound is at most $1+7\sigma+4\theta$.
Indeed, its numerator is at most $1+9\sigma/4$, and
$\kappa\ge1-t$ for $t=\theta+\sigma/2\le11/40$.
For $0\le t\le1/3$ the inequality $(1-t)^{-2}\le1+4t$ follows by
multiplication and $2-7t+4t^2\ge0$. Hence the trace ratio is at most
$(1+9\sigma/4)(1+4\theta+2\sigma)
\le1+7\sigma+4\theta$, using $\theta\le1/4$ and $\sigma\le1/20$.
Thus the passage from the pairwise estimates
on the full matrices to the required estimate on their restrictions
is justified by ordinary matrix inequalities. The quantitative
pairwise estimates themselves are still an input about the theta maps.

% Optional refinement of the independent comparison in central_step.tex.
% This subsection uses the notation and labels introduced there.
\subsection{A refinement using the complementary subspaces}
\label{qrh-sharper-slope-section}

The same hypotheses give a second bound that can be useful when the
source slope spread is much smaller than the permitted target deficit.
The simple estimate in Proposition~\ref{cs-slope} remains available;
one may use whichever of the two bounds is smaller.

\begin{proposition}[Refined slope comparison]
\label{qrh-sharper-slope}
Keep the lattice, map, rank, and numerical hypotheses of
Proposition~\ref{cs-slope}. In its matrix hypotheses, assume
$0\le e<1$ in place of $e\le1/4$, and retain
\eqref{cs-isotropic} with $\eta\ge0$.
Put $q=\sqrt e$ and, for $u,H\ge0$, define
\[
 \mathcal A=
 \frac{u\log(1+H/u)
       +2q\sqrt u\bigl(\sqrt{u+H}-\sqrt u\bigr)}{1-q}.
\]
At $u=0$ this expression means its continuous value $0$;
at $H=0$ it also equals $0$. Then
\begin{equation}\label{qrh-sharper-slope-bound}
 s-\mu(F)\le
 \frac{\mathcal A+\eta H/2+
       \sqrt{\mathcal A\eta H+(\eta H)^2/4}}{1-\delta}.
\end{equation}
If $u=0$, the stronger bound
\[
 s-\mu(F)\le\frac{\eta H}{(1+\eta)(1-\delta)}
\]
holds. In particular, when $e=\eta=0$,
\begin{equation}\label{qrh-exact-slope-bound}
 s-\mu(F)\le\frac{u\log(1+H/u)}{1-\delta}.
\end{equation}
All assertions include $H=0$ and $u=0$ with the stated conventions.
\end{proposition}

\begin{proof}
As in Proposition~\ref{cs-slope}, rescale the source norms by
$\exp(h_i)$ so that all map heights are nonpositive. For $z>0$,
let $V_i(z)$ be the complex realization, in the chosen source basis,
of the subspace of adjusted slopes at least $s-z$. Set
\[
 t(z)=\sum_i\dim V_i(z),\qquad c(z)=1-t(z)/R.
\]
The source slope assumptions give
\begin{equation}\label{qrh-source-complement}
 0\le c(z)\le\frac{u}{u+z}.
\end{equation}
When $u=0$, every adjusted source slope equals $s$, so $c(z)=0$.
By Lemma~\ref{cs-functoriality}, the images of $V_i(z)$ under $A_i$
lie in the subspace $W(z)$ obtained by adjoining the last $D-n$
coordinate directions to $F_{\ge s-z}$. Write
\[
 r(z)=\operatorname{rk}F_{\ge s-z},\qquad
 x(z)=\frac{D-\dim W(z)}D=\frac{n-r(z)}D.
\]

Fix $z$. Let $P_i$ be the orthogonal projection onto $V_i(z)$,
let $P$ be the orthogonal projection onto $W(z)^\perp$, and put
\[
 Q=\sum_i A_iA_i^*,\qquad
 C=\sum_i A_iP_iA_i^*,\qquad T=\operatorname{tr}Q.
\]
Thus $0\le C\le Q$ and $PC=0$. Since $0\le P\le I$, one has
$\operatorname{tr}(PQ)\le\operatorname{tr}(Q-C)$.
Using Cauchy--Schwarz for the sum of the Hilbert--Schmidt inner
products gives
\begin{align*}
 \operatorname{tr}(Q-C)
 &=R-t+\sum_i\operatorname{tr}
       \bigl((I-P_i)(A_i^*A_i-I)\bigr)\\
 &\le R-t+\sqrt{(R-t)eR}
   =R(c+q\sqrt c).
\end{align*}
The diagonal hypothesis also gives $T\ge(1-q)R>0$.
Consequently
\begin{equation}\label{qrh-complement-trace}
 \frac{\operatorname{tr}(PQ)}T
       \le d(z):=\frac{c(z)+q\sqrt{c(z)}}{1-q}.
\end{equation}

The isotropy hypothesis is exactly the variance bound
\[
 \left\|Q-\frac TD I\right\|_{\mathrm{HS}}^2
       \le\frac{\eta T^2}D.
\]
Since $\operatorname{tr}P=Dx$ and
$\|P-xI\|_{\mathrm{HS}}^2=Dx(1-x)$, another application of
Cauchy--Schwarz yields
\[
 \frac{\operatorname{tr}(PQ)}T
 \ge x-\sqrt{\eta x(1-x)}.
\]
Combining this with \eqref{qrh-complement-trace}, and using
$0\le x\le1$, gives
\begin{equation}\label{qrh-complement-dimension}
 x(z)\le d(z)+\sqrt{\eta x(z)}.
\end{equation}

Integrate over $0<z<H$. By \eqref{qrh-source-complement},
\[
 \int_0^H d(z)\,dz\le\mathcal A,
\]
because the integrals of $u/(u+z)$ and
$\sqrt{u/(u+z)}$ are respectively
$u\log(1+H/u)$ and
$2\sqrt u(\sqrt{u+H}-\sqrt u)$.
If $X=\int_0^H x(z)\,dz$, Cauchy--Schwarz in
\eqref{qrh-complement-dimension} therefore gives
\[
 X\le\mathcal A+\sqrt{\eta HX}.
\]
Solving the quadratic inequality for $\sqrt X$ proves
\[
 X\le\mathcal A+\frac{\eta H}{2}
              +\sqrt{\mathcal A\eta H+\frac{(\eta H)^2}{4}}.
\]
Every target slope is at least $s-H$. Hence, exactly as in the
proof of Proposition~\ref{cs-slope},
\[
 s-\mu(F)\le\int_0^H\frac{n-r(z)}n\,dz
       =\frac Dn X\le\frac X{1-\delta},
\]
which proves \eqref{qrh-sharper-slope-bound}.

If $u=0$, then $c(z)=0$ and
$\operatorname{tr}(PQ)=0$. The sharper variance inequality gives
$x\le\sqrt{\eta x(1-x)}$. For $x>0$, squaring and dividing by
$x$ gives $x\le\eta/(1+\eta)$; the same bound holds when $x=0$.
Integrating proves the asserted endpoint improvement. If $e=\eta=0$,
then $\mathcal A=u\log(1+H/u)$ and the general bound is
\eqref{qrh-exact-slope-bound}. When $H=0$, the target slope
assumption already gives $\mu(F)\ge s$, so all displayed bounds
hold without an integration interval.
\end{proof}

The special case $e=\eta=0$ means that the matrices are isometries
on their source spaces and that $Q=(R/D)I$. It gives coefficient
$1$ in front of $u\log(1+H/u)$, in place of the coefficient $2$
in the simpler comparison. For nonzero diagonal error,
\[
 \mathcal A\le
 \frac{u\log(1+H/u)+2\sqrt e\,\sqrt{uH}}{1-\sqrt e},
\]
so this form can reduce the effect of the diagonal error when
$u$ is much smaller than $H$. The mixed term involving $\eta$
is explicit in \eqref{qrh-sharper-slope-bound}; no assertion of
uniform improvement over every parameter choice is needed.
This refinement concerns only the comparison of slopes. It does
not establish the integral automorphic maps, their height bounds,
or their matrix estimates.


\subsection{The arithmetic inequality still required}

Here is a precise reduction of the remaining numerical comparison.
It avoids auxiliary parameters that disappear by cancellation.

\begin{lemma}[Combination of the local estimates]\label{cs-combination}
Let $I$ be a nonempty finite set, with positive weights $d_i$ and
weighted mean denoted by a bar. Let
$\nu_i\le\nu_i^+$, $h_i$, $M_i$, $M_i^+$, $D_i$, $E_i$, and $\ell_i$
be real numbers. The scalar parameters
$c,F,G,A,\mathcal M,K,M_-,M_+,E_+,V,L,L_0$ are also real.
Suppose
\begin{align*}
 c+h_i-F&\le D_i,&
 \nu_i&\ge M_i+G+F-E_i-\ell_i,\\
 \nu_i^+&\le M_i^++G+F-E_i+E_+,&
 -M_-&\le M_i,\qquad M_i^+\le M_+,\\
 D_i+E_i&\le\mathcal M+K,&
 (-E_i-h_i)-(-E_j-h_j)&\le V,\\
 A-G&\le L,&\overline\ell&\le L_0.
\end{align*}
Put $s=\overline{\nu-h}$ and
$U=\max_i(\nu_i^+-h_i)-s$. Suppose $H\ge0$ and
\[
 s-c\le\frac{2U\log(1+H/U)+(\eta+4\sqrt e)H}{1-\delta},
 \qquad 0\le\delta<1,\quad\eta,e\ge0.
\]
Then, with $B=M_++M_-+V+E_++L_0$, one has $0\le U\le B$ and
\[
 A-\mathcal M\le L+K+M_-+L_0+
 \frac{2B\log(1+H/B)+(\eta+4\sqrt e)H}{1-\delta}.
\]
If $H\le Xb$ with $X\ge0$, $b>0$, the last numerator is at most
$2(B+b)\log(1+X)+(\eta+4\sqrt e)Xb$.
\end{lemma}
\begin{proof}
The first, second, fifth and lower-slope inequalities give
$s\ge G+c-\mathcal M-K-M_--\overline\ell$.
Comparing the upper bound for $\nu_i^+-h_i$ with the weighted lower
bound for $\nu_j-h_j$ gives
$U\le M_++M_-+V+E_++\overline\ell\le B$.
Also $U\ge0$, since $\nu_i\le\nu_i^+$.
The function $z\mapsto z\log(1+H/z)$ is increasing for $z>0$:
its derivative is
$\log(1+H/z)-H/(H+z)\ge0$, by integrating
$1/(1+t)\ge1/(1+H/z)$ from $0$ to $H/z$.
The asserted inequality now follows by addition. Finally, concavity of
the logarithm, with weights $B/(B+b)$ and $b/(B+b)$, gives
\[
 B\log(1+H/B)\le(B+b)\log(1+H/(B+b))
 \le(B+b)\log(1+X).
\]
The limiting cases $B=0$ or $H=0$ follow by continuity.
\end{proof}

For comparison with the proposed arithmetic construction, let $m\ge1$ be
its integer number of lifts, let $\mathcal T$ be the finite nonempty
index set of interpolation orders $a_t>0$, and let $T_i>0$ be its
auxiliary primes, $g>0$ its product of two local factors, and $w_{j\rho}>0$
its weights, where $f\ge1$ is an integer,
$1\le j\le m$, and $1\le\rho\le f$. Write $\mathcal R^\vee$
for the target Hermitian lattice, $h_i$ for the map-height bounds,
$F\in\mathbb R$ for the local correction, and $v>0$ for the
archimedean normalization. Set
\[
 S=\sum_t a_t,\qquad
 \log\Omega=\gamma+\log(\pi/v)-\ell,
\]
where $\gamma$ is Euler's constant and $\ell$ is the relevant real
logarithmic derivative of a quadratic $L$-function at $1$. The required
divisor inequality has the concrete form
\begin{equation}\label{cs-required-divisor}
 \mu(\mathcal R^\vee)+h_i-F\le
 \sum_t\{m(a_t\log a_t-a_t-a_t\log\Omega)-a_t\log T_i\}
       +\mathcal E_D.
\end{equation}
The slopes here and below may all be multiplied by the same fixed
normalizing factor; the preceding homogeneous inequalities are unchanged.
Let $\mathcal L,B_0,N_c\ge0$ denote the error parameters whose
arithmetic interpretation and simultaneous bounds must be supplied by
that construction. The proposed bound for the error is
\[
 \mathcal E_D=C|\mathcal T|a\mathcal L
             +Cm(B_0+N_c)x^{2\varepsilon},
 \qquad a=xb,\quad b=\exp(x^\varepsilon),\quad
 x^{1/2-\varepsilon}\le m\le x^{1/2+\varepsilon/2}.
\]
The multiplicative constant $C$ may depend on the family and on
$\varepsilon$, while the exponent $2$ must be independent of
$\varepsilon$. The assertion is required for every word and all
sufficiently large parameters in the selected family.
In particular, a multiplicative constant depending on $\varepsilon$
can be absorbed by increasing that threshold. An exponent coefficient
growing like $1/\varepsilon$ cannot be treated in the same way.

The remaining cancellation is elementary. Suppose
\begin{align*}
 E_i&\le mS\log g+\mathcal E_E,\\
 \log T_i&\ge
   \sum_{j=1}^m\frac1f\sum_{\rho=1}^f\log(w_{j\rho}v/\pi)
       +m\log g-\mathcal E_C.
\end{align*}
If $D_i$ denotes the right side of \eqref{cs-required-divisor}, then
\[
 D_i+E_i\le\mathcal M+S\mathcal E_C+\mathcal E_D+\mathcal E_E,
\]
where
\[
 \mathcal M=\sum_{j=1}^m\sum_t
 a_t\left(\log a_t-1-\gamma+\ell
                   -\frac1f\sum_{\rho=1}^f\log w_{j\rho}\right).
\]
This follows by direct substitution: both $mS\log g$ and
$mS\log(v/\pi)$ cancel. Together with comparison of the source slopes,
the mean loss bound, and Proposition~\ref{cs-slope},
Lemma~\ref{cs-combination} proves the final scalar estimate. It removes
the term $\log|I|$ and replaces the coefficient $4$ of
$U\log(1+H/U)$ in the displayed source comparison by $2$.
These improvements simplify constants and exposition. They do not
prove \eqref{cs-required-divisor}, nor the matrix estimates for the
automorphic maps, nor by themselves the required bounds for the selected
twisted-character sums.  The latter involve a fixed finite family of
base characters; only the RH subcase reduces to three quadratic
characters.  A separate argument by averaging over the auxiliary primes
can establish this selection under a suitable uniform QRH hypothesis.

The source of \eqref{cs-required-divisor} is more than the
nonvanishing of one integer determinant. At each finite place it needs
pairing divisibility for every vector in a prescribed lattice, and a
bound for a rank-normalized sum of valuations of bordered minors.
At every infinite place it needs upper bounds for operator norms of a
chain of rectangular matrices, together with lower bounds for the
eigenvalues in the quotient metrics. These are separate assertions
about the constructed matrices. The elementary determinant of the
preceding quadratic argument does not establish them.

\begin{remark}[A counterexample to a possible simplification]\label{cs-counterexample}
Nonvanishing modulo every prime does not imply the matrix estimate
\eqref{cs-isotropic} with a small $\eta$. For any integer $M\ge1$, take
the columns of the unimodular matrix
\[
 \begin{pmatrix}1&M\\0&1\end{pmatrix}
\]
and normalize each column to length $1$. Thus
$A_1=(1,0)^t$ and $A_2=(M,1)^t/\sqrt{M^2+1}$.
Both diagonal errors vanish, but
\[
 \operatorname{tr}Q=2,\qquad
 \operatorname{tr}(Q^2)=2+\frac{2M^2}{M^2+1},\qquad
 \frac{2\operatorname{tr}(Q^2)}{(\operatorname{tr}Q)^2}-1
       =\frac{M^2}{M^2+1}.
\]
The smallest permissible $\eta$ tends to $1$. This is a counterexample
to replacing the theta estimates by determinant nonvanishing; it is
not a counterexample to the theta estimates asserted for the specific
automorphic construction.
\end{remark}


\section{Simultaneous selection of the auxiliary quadratic characters}
\label{qs-section}

The three quadratic characters used in the RH subcase do not describe
the entire requirement for Dirichlet GRH. For a fixed modulus one needs
a finite family of its characters, each twisted by either of two new
quadratic characters or by their product. The same two primes must work
throughout an interval of scales. This section proves that selection
from the following explicit analytic hypothesis:
\begin{equation}\label{qs-qrh}
 L(s,\chi)\ne0\qquad(\operatorname{Re}s>11/12)
\end{equation}
for every primitive Dirichlet character, with the pole of $\zeta$ at
$1$ allowed. This is the endpoint stated in the supplied October
quasi-Riemann-hypothesis paper; the September endpoint is stronger.
\cite{qrh-october,qrh-september}
The conclusion below remains conditional on \eqref{qs-qrh} in this
reconstruction. Its proof does not assume the Riemann hypothesis for
any of the selected characters.

For a nonprincipal Dirichlet character $\chi$, write
\[
 F_\chi(s)=-\frac{L'}{L}(s,\chi),\qquad
 G_\chi(y)=\sum_{n\le y}\Lambda(n)\chi(n)(y/n-1)-yF_\chi(1).
\]
Thus the second term in $G_\chi$ is
$+yL'(1,\chi)/L(1,\chi)$. For an odd prime $p$, let
$\chi_p(n)=(n/p)$ be the Legendre character, extended by zero at
multiples of $p$. In a product of these characters, a factor with
exponent zero is omitted; it is not a principal character at an
unnecessarily enlarged modulus. Each resulting primitive character
is considered at its own conductor.

\begin{theorem}[Simultaneous quadratic selection]\label{qs-selection}
Assume \eqref{qs-qrh}. Fix an integer $M\ge1$, real numbers
$0<\varepsilon<1/52$ and $A\ge1$. For all sufficiently large $Y$,
with threshold depending only on $M,\varepsilon,A$, put
\[
 z=Y^{\varepsilon/5},\qquad z'=z(\log Y)^2.
\]
There are distinct primes $q<d$, both coprime to $M$, such that
\begin{gather*}
 e^{z/2}<q<e^z,\quad q\equiv3\pmod8,\qquad
 e^{z'/2}<d<e^{z'},\quad d\equiv7\pmod{16},\\
 q\mid d+1,\qquad
 \ell\mid d+1\quad\hbox{for every odd prime }\ell\le z.
\end{gather*}
For every primitive character $\xi$ whose conductor divides $8M$,
and all $i,j\in\{0,1\}$ with $i+j>0$, the character
$\chi=\xi\chi_q^i\chi_d^j$ is primitive and nonprincipal, and
simultaneously for every $Y\le y\le Y^A$,
\begin{equation}\label{qs-bounds}
 |G_\chi(y)|\le y^{1/2+\varepsilon/2},\qquad
 \left|\frac{L'}{L}(1,\chi)\right|\le4\log y.
\end{equation}
Every such $L$-function is also zero-free in
$\operatorname{Re}s>399/400$, at every height.
\end{theorem}

The finite family of base characters is fixed before $Y$ varies.
The primes $q,d$, and consequently their quadratic twists, may change
with $Y$. This order of quantifiers is essential: the theorem gives
no square-root estimate for an arbitrary fixed character at every
large scale.

An imprimitive base character modulo $M$ is used through its primitive
inducer, whose conductor divides $M$. This supplies the primitive
character determined by its values away from $2Mqd$.
Primitive and imprimitive centered sums are not silently identified.
If a nonprincipal $\chi$ modulo $Q$ is induced from $\chi^*$ of
conductor $f$, direct subtraction of the finite Euler factors gives
the exact correction
\[
 G_\chi(y)-G_{\chi^*}(y)=
 \sum_{\substack{p\mid Q\\p\nmid f}}
   \sum_{r\ge1}(\log p)\chi^*(p)^r\min(1,y/p^r).
\]
Indeed, subtract the missing prime powers in the finite defining sum
and add their convergent geometric series in $-L'/L(1)$; the two
terms combine into the displayed minimum. The theorem is stated
at the primitive conductors, so these corrections are not needed
in its selection argument.

\subsection{Uniform analytic estimates with a smooth cutoff}

We record the conductor dependence used below. Under
\eqref{qs-qrh}, put $\sigma=23/24$. Uniformly for nonprincipal
characters modulo $Q$,
\begin{equation}\label{qs-logderivative}
 |F_\chi(v+it)|\ll\log\bigl((Q+3)(|t|+3)\bigr)
       \qquad(\sigma\le v\le2).
\end{equation}
Here and below the constants in analytic estimates are absolute once
the displayed lines and smooth functions have been fixed. For clarity,
this estimate does not mean that the value of $L$ is bounded below by
an absolute constant. The proof uses a holomorphic logarithm. For a
primitive nonprincipal character, its partial sums have modulus at
most $Q$, so partial summation gives
$|L(s,\chi)|\le Q|s|/\operatorname{Re}s$ in the right half-plane.
On a disk centered at $2+it$ with radius strictly between
$2-\sigma$ and $2-11/12$, take the logarithm whose value at the
center is supplied by the absolutely convergent Euler product.
That center value is bounded absolutely. The disk estimate of
Lemma~\ref{lem:qrh-log-disk}, followed by Cauchy's derivative
formula on a smaller disk, gives \eqref{qs-logderivative}.
The chosen radii are fixed, so no conductor enters its constant.

If $\chi$ is induced from $\chi^*$, its additional contribution is
\[
 \frac{L'}L(s,\chi)-\frac{L'}L(s,\chi^*)
 =\sum_{\substack{p\mid Q\\p\nmid\operatorname{cond}\chi^*}}
       \frac{\chi^*(p)\log p}{p^s-\chi^*(p)}.
\]
On the displayed strip this is at most
$(2^\sigma-1)^{-1}\log Q$ in modulus. The additional Euler
factors have their zeros on $\operatorname{Re}s=0$, so they do
not change the zero-free hypothesis in use. For a principal
character one instead bounds $F_\chi(s)-1/(s-1)$ by the same
quantity, using $(s-1)\zeta(s)$ in the disk argument.

Fix a smooth function $w:[0,\infty)\to[0,1]$ with
$w(u)=1$ for $u\le1$ and $w(u)=0$ for $u\ge2$.

\begin{lemma}[A finite smooth approximation]\label{qs-cutoff}
Under \eqref{qs-qrh}, for every nonprincipal character modulo $Q$
and $Z\ge2$,
\[
 F_\chi(1)=\sum_{n\le2Z}\frac{\Lambda(n)\chi(n)}n w(n/Z)
       +O\bigl(Z^{-1/24}\log(Q+3)\bigr).
\]
Consequently, uniformly for $1\le y\le Z$,
\begin{equation}\label{qs-smooth-identity}
 G_\chi(y)=-\sum_{n\le2Z}\Lambda(n)\chi(n)
                    \min(1,y/n)w(n/Z)
       +O\bigl(yZ^{-1/24}\log(Q+3)\bigr).
\end{equation}
There is also an absolute constant $C$ such that
\begin{equation}\label{qs-loglog}
 |F_\chi(1)|\le192\log\log(Q+3)+C.
\end{equation}
These statements exclude principal characters: $F_\chi(1)$ would
then have a pole.
\end{lemma}
\begin{proof}
For $\operatorname{Re}s>0$, the Mellin transform of $w$ is
\[
 \widehat w(s)=\int_0^\infty w(u)u^{s-1}\,du
       =-\frac1s\int_1^2 w'(u)u^s\,du.
\]
The second expression continues it meromorphically, with only a
simple pole at $0$, of residue $1$. Repeated integration by parts
after $u=e^v$ proves rapid decay on every fixed vertical strip away
from this pole. Mellin inversion and absolute convergence on
$\operatorname{Re}s=1$ give
\[
 \sum_n\frac{\Lambda(n)\chi(n)}n w(n/Z)
 =\frac1{2\pi i}\int_{(1)}F_\chi(1+s)\widehat w(s)Z^s\,ds.
\]
Move the contour to $\operatorname{Re}s=-1/24$.
Equation~\eqref{qs-logderivative} and rapid decay justify the shift
and bound the new integral by $O(Z^{-1/24}\log(Q+3))$.
The only crossed pole is that of $\widehat w$ at $0$, with residue
$F_\chi(1)$. This proves the first assertion with the stated sign.
When $n\le y\le Z$, $w(n/Z)=1$. Substitution in the definition
of $G_\chi$ therefore gives \eqref{qs-smooth-identity} exactly,
apart from the displayed analytic error.

We shall use the elementary bound
\begin{equation}\label{qs-chebyshev}
 \psi(t):=\sum_{n\le t}\Lambda(n)\le4t\qquad(t\ge1).
\end{equation}
Indeed, every prime power in $(m,2m]$ contributes its logarithm
to $\binom{2m}{m}$, and all other contributions to its prime
valuations are nonnegative. Thus
$\psi(2m)-\psi(m)\le2m\log2$.
Sum this for powers of two and use the smallest power of two
at least $t$ to obtain \eqref{qs-chebyshev}.
Partial summation now gives
\[
 \sum_{n\le T}\frac{\Lambda(n)}n\le4(1+\log T)
       \qquad(T\ge1).
\]
Apply the first assertion with $Z=(\log(Q+3))^{48}$.
The error is bounded absolutely, and the polynomial has norm at
most $4(1+\log(2Z))$. This is
$192\log\log(Q+3)+O(1)$, proving \eqref{qs-loglog}.
\end{proof}

\subsection{Quadratic orthogonality in a progression}

Choose a nonnegative $\Phi\in C_c^\infty(1/2,1)$ with
$I_\Phi=\int\Phi>0$. For $(a,m)=1$ and $U\ge2$, give a prime
$p\equiv a\pmod m$ weight $(\log p)\Phi(p/U)$. Let
\[
 W(U;m,a)=\sum_{p\equiv a\ (m)}(\log p)\Phi(p/U).
\]
Whenever $W>0$, expectation below means the weighted average with
these weights divided by $W$.

\begin{lemma}[Weighted correlations]\label{qs-correlations}
Assume \eqref{qs-qrh}. Uniformly in $m,a,U$,
\[
 W(U;m,a)=\frac{I_\Phi U}{\varphi(m)}
       +O\bigl(U^{23/24}\log(m+3)+U^{1/2}(\log(2U))^2\bigr).
\]
For odd squarefree $s>1$ coprime to $m$,
\[
 \left|\sum_{p\equiv a\ (m)}(\log p)\Phi(p/U)(p/s)\right|
 \ll U^{23/24}\log(ms+3)+U^{1/2}(\log(2U))^2.
\]
In particular, if $W\ge I_\Phi U/(2\varphi(m))$, then
for every such $s\le N^2$,
\begin{equation}\label{qs-rho}
 |\mathbb E(p/s)|\ll\varphi(m)
 \left(U^{-1/24}\log(mN^2+3)
             +U^{-1/2}(\log(2U))^2\right).
\end{equation}
\end{lemma}
\begin{proof}
For a Dirichlet character $\alpha$ modulo $L$, Mellin inversion gives
\[
 \sum_n\Lambda(n)\alpha(n)\Phi(n/U)
 =\frac1{2\pi i}\int_{(2)}F_\alpha(s)
                          \widehat\Phi(s)U^s\,ds.
\]
Moving the line to $23/24$ gives an error
$O(U^{23/24}\log(L+3))$ by \eqref{qs-logderivative}.
The pole at $1$ contributes $I_\Phi U$ precisely when $\alpha$
is principal. This argument includes imprimitive characters and
their local Euler factors, as described above.

Expand the progression condition with the characters modulo $m$.
For the first formula exactly one character is principal. For the
second, the products of these characters with $(\cdot/s)$ are all
nonprincipal modulo $ms$. To verify the latter assertion, the Chinese
remainder theorem permits an integer equal to $1$ modulo $m$ on
which the nontrivial character $(\cdot/s)$ has value $-1$.
The factor $1/\varphi(m)$ in the character average cancels the
number of error terms. Finally remove all prime powers of exponent
at least two. Their absolute total is
$O(U^{1/2}(\log(2U))^2)$, independently of the progression or
twist. Division by the lower bound for $W$ proves
\eqref{qs-rho}.
\end{proof}

The next lemma is elementary and independent of
\eqref{qs-qrh}. It makes explicit why forcing splitting at a short
initial set of primes does not destroy the second moment.

\begin{lemma}[A finite second moment]\label{qs-second-moment}
Let $\mathcal P$ be a finite probability space of primes $p>N$
in one reduced class $a\pmod m$, where $8\mid m$.
Let $2\le t\le N$, $1\le z\le t^{1/5}$, and suppose every prime
divisor of $m$ at most $N$ is at most $z$. Suppose that
\[
 |\mathbb E(p/s)|\le\rho
\]
for every odd squarefree $s\le N^2$, $s>1$, $(s,m)=1$.
Then, for every complex sequence $|\kappa(n)|\le1$,
\begin{equation}\label{qs-moment}
 \mathbb E\left|\sum_{n\le N}\Lambda(n)\min(1,t/n)
                         \kappa(n)(n/p)\right|^2
 \le C\left(t\log(2t)+\rho t^2\log^2(eN)\right)
\end{equation}
with an absolute constant $C$.
The same conclusion holds if the correlation hypothesis omits
at most one odd squarefree integer $s_0>1$.
\end{lemma}
\begin{proof}
Put $c_n=\Lambda(n)\min(1,t/n)$. Only prime powers occur.
For a fixed prime $\ell$, let
\[
 b_\ell=\sum_{j\ge1}c_{\ell^j},
\]
where the sum can be truncated at $\ell^j\le N$; all estimates
below remain valid for the infinite sum. If $\ell\le t$, a
geometric tail gives $b_\ell\le\log t+2\log\ell\le3\log t$.
If $\ell>t$, it gives
$b_\ell\le2t\log\ell/\ell$.
It follows that
\begin{equation}\label{qs-prime-diagonal}
 \sum_\ell b_\ell^2\ll t\log(2t),\qquad
 \max_\ell b_\ell\ll\log(2t).
\end{equation}
Here are the elementary details of the sum. From
\eqref{qs-chebyshev}, splitting the primes at $\sqrt t$ gives
$\pi(t)\le10t/\log t$ for $t\ge2$.
This bounds the terms $\ell\le t$. For the others, partial
summation with $\vartheta(u)=\sum_{\ell\le u}\log\ell\le4u$
gives
\[
 \sum_{\ell>t}\frac{(\log\ell)^2}{\ell^2}
 \le4\int_t^\infty\frac{2\log u-1}{u^2}\,du
 \ll\frac{\log(2t)}t.
\]
The claimed bounds follow, with a harmless adjustment on $2\le t<3$.

All even powers contribute at most
\[
 \sum_{r\ge1}\Lambda(r)\min(1,t/r^2)\le12\sqrt t
\]
by \eqref{qs-chebyshev} and partial summation at $\sqrt t$.
The total contribution from primes dividing $m$ is at most
$3z\log t$, since those contributing primes are at most $z$.
Because $z\le t^{1/5}$ and
$\log t\le(10/3)t^{3/10}$, these terms together have total
absolute value
\begin{equation}\label{qs-forced-part}
 R\ll\sqrt t.
\end{equation}

Quadratic reciprocity shows that $(\ell/p)$ is constant on the
given progression when $\ell\mid m$: the residue modulo $4$
is fixed, the residue modulo each odd such $\ell$ is fixed, and
the value at $2$ is fixed by the residue modulo $8$.
Group the sum in \eqref{qs-moment} into these terms and the even
powers, followed by odd powers of primes $\ell\nmid m$.
It becomes
\[
 A_0+\sum_{\ell\nmid m}A_\ell(\ell/p),
       \qquad |A_0|\le R,\quad |A_\ell|\le b_\ell.
\]
The diagonal quadratic factors equal $1$, since $p>N$.
The other factors are, up to a constant sign fixed by the
progression, $(p/\ell)$ or $(p/(\ell\ell'))$ for distinct
primes outside $m$. Expanding the square and using the correlation
bound therefore gives
\[
 \mathbb E|\cdots|^2\le R^2+\sum_\ell b_\ell^2
                 +\rho\left(\sum_{n\le N}c_n\right)^2.
\]
Partial summation gives
$\sum_{n\le N}c_n\le4t(1+\log(N/t))\le4t\log(eN)$.
Together with \eqref{qs-prime-diagonal} and
\eqref{qs-forced-part}, this proves the bound.

If there is one exceptional $s_0$, it affects these correlations
only when $s_0$ is one prime or the product of two distinct primes.
In the first case the extra term is at most
$2R\max b_\ell$; in the second it is at most
$2(\max b_\ell)^2$. Both are $O(t\log(2t))$.
Thus one exceptional correlation does not change the conclusion.
No multiplicativity of $\kappa$ has been used.
\end{proof}

\subsection{Uniformity on an interval of scales}

\begin{lemma}[A maximal second moment]\label{qs-maximal}
For the same finite probability space, put
\[
 S_p(t;\kappa)=\sum_{n\le N}\Lambda(n)\min(1,t/n)
                                      \kappa(n)(n/p).
\]
Suppose that, for every $Y\le t\le X$ and every
$|\kappa(n)|\le1$,
$\mathbb E|S_p(t;\kappa)|^2\le C t\log(2t)$, where
$2\le Y\le X\le N$. For each fixed such sequence and each
$\tau>0$, the probability that
\[
 |S_p(t;\kappa)|>\tfrac12t^{1/2+\tau}
       \quad\hbox{for some }t\in[Y,X]
\]
is at most
\begin{equation}\label{qs-bad-probability}
 16C Y^{-2\tau}
 \left\{\frac{\log(4Y)}{1-2^{-2\tau}}
       +\frac{(\log2)2^{-2\tau}}{(1-2^{-2\tau})^2}\right\}.
\end{equation}
For any fixed finite set of sequences, multiply this bound by
their number to control failure for at least one sequence.
\end{lemma}
\begin{proof}
The function $S_p(t;\kappa)$ is continuous and piecewise linear.
At all but finitely many points,
\[
 D_p(t):=S'_p(t;\kappa)
       =\sum_{t<n\le N}\frac{\Lambda(n)}n\kappa(n)(n/p).
\]
Its second moment is at most $C\log(2t)/t$: multiply by $t$
and apply the hypothesis to
$\kappa(n)1_{n>t}$. The uniformity in bounded sequences is what
permits this step.

On a block $[T,W]$ with $T\le W\le2T$, the fundamental theorem
of calculus and Cauchy--Schwarz give
\[
 \sup_{T\le t\le W}|S_p(t;\kappa)|^2
 \le2|S_p(T;\kappa)|^2+2T\int_T^W|D_p(v)|^2\,dv.
\]
Taking expectations yields a bound $4CT\log(4T)$.
Markov's inequality therefore bounds failure in that block by
$16CT^{-2\tau}\log(4T)$. Cover $[Y,X]$ by the blocks starting
at $T=2^jY$ and sum this bound. The two geometric sums are
$\sum r^j=(1-r)^{-1}$ and
$\sum jr^j=r(1-r)^{-2}$, with $r=2^{-2\tau}$.
This proves \eqref{qs-bad-probability}; the finite-family assertion
is the union bound. There is no discretization loss depending on
the length $X/Y$.
\end{proof}

\subsection{Proof of simultaneous selection}

\begin{proof}[Proof of Theorem~\ref{qs-selection}]
Let $\mathcal F$ be the finite set of primitive characters whose
conductors divide $8M$, and put
\[
 X=Y^A,\qquad Z=Y^{48(A+1)},\qquad N=\lfloor2Z\rfloor.
\]
Its cardinality depends on $M$ alone. We use the finite sums with
the bounded sequence $\kappa(n)=\xi(n)w(n/Z)$, and later also
with $\kappa(n)=\xi(n)\chi_q(n)w(n/Z)$.

First take $U_1=e^z$, $m=8$, $a=3$, and use the probability
measure preceding Lemma~\ref{qs-correlations}. Its total weight
is positive for all sufficiently large $Y$. Every prime in its
support is greater than $U_1/2$, hence greater than $N^2$ and
$8M$. For every fixed $L>0$, the right side of
\eqref{qs-rho} is at most $Y^{-L}$ for all sufficiently large
$Y$: $z=Y^{\varepsilon/5}$ grows faster than every multiple of
$\log Y$, whereas $N$ is a fixed power of $Y$.
This proves the correlation estimate uniformly for all the
integers needed in Lemma~\ref{qs-second-moment}.

For $Y\le t\le X$, the only prime divisor of $8$ is harmless
in \eqref{qs-forced-part}. More generally $z\le t^{1/5}$,
because $\varepsilon<1$. Choose $L=A+2$. The error term in
\eqref{qs-moment} is absorbed into $Ct\log(2t)$, uniformly
over this interval and over all bounded sequences. Apply
Lemma~\ref{qs-maximal} with $\tau=\varepsilon/2$ to all
$\xi\in\mathcal F$. Their total failure probability tends
to zero. We can consequently choose a prime $q$ in the support
such that all their finite sums have modulus at most
$\tfrac12t^{1/2+\varepsilon/2}$ throughout $[Y,X]$.
This prime satisfies $q\equiv3\pmod8$ and
$e^{z/2}<q<e^z$ once $Y$ is large.

Now put
\[
 P_z=\prod_{\substack{3\le\ell\le z\\\ell\ {
m prime}}}\ell,
       \qquad m=16qP_z,
\]
and choose by the Chinese remainder theorem the reduced residue
class $a$ given by
$a\equiv7\pmod{16}$ and $a\equiv-1\pmod{qP_z}$.
Take $U_2=e^{z'}$. The bound \eqref{qs-chebyshev} gives
\[
 \log m\le\log16+5z,
\]
uniformly in the previously chosen $q$.
The modulus satisfies the local condition in
Lemma~\ref{qs-second-moment}: the prime $q$ is greater than
$N$, and every other odd prime divisor is at most $z$.

Again the total prime weight is positive, and for every fixed
$L>0$ the correlation bound is at most $Y^{-L}$, uniformly in
$q$. Indeed its logarithm is bounded above, apart from
$O_{A,M,\varepsilon}(\log Y+\log z')$, by
\[
 5z-z'/24=5z-z(\log Y)^2/24,
\]
which is eventually smaller than $-L\log Y$. The prime-power
error has the still smaller exponent $5z-z'/2$.
All primes in the support exceed $N^2$ and $q$.

Apply the same second-moment and maximal estimates to the at
most $2|\mathcal F|$ sequences
$\xi(n)\chi_q(n)^i w(n/Z)$, $\xi\in\mathcal F$, $i=0,1$.
Their absolute values are at most one; the constants in the
second moment do not depend on the already chosen $q$.
The total failure probability again tends to zero. Choose $d$
outside this exceptional set. It has all the stated congruences
and sizes, and $q<d$. Both primes exceed $8M$.

The product of primitive characters of relatively prime
conductors is primitive: by the Chinese remainder theorem its
restriction to each conductor factor is the original primitive
character, so a character inducing the product must contain
each of those factors in its conductor. Applying this to
$\xi$, $\chi_q$, and $\chi_d$ shows that the conductor of each
selected twist is
$f q^i d^j$, where $f=\operatorname{cond}\xi\mid8M$.
It is nonprincipal, since at least one new prime occurs to the
first power in its conductor.

For all selected twists this conductor $Q$ satisfies
\[
 \log(Q+3)\ll_M z(\log Y)^2.
\]
Equation~\eqref{qs-smooth-identity} has error, uniformly on
$[Y,X]$, at most
\[
 C_M Y^A Y^{-2(A+1)}Y^{\varepsilon/5}(\log Y)^2=o(1).
\]
For large $Y$ this is at most
$\tfrac12Y^{1/2+\varepsilon/2}$.
The choices of $q,d$ therefore give the first bound in
\eqref{qs-bounds} for all twists simultaneously.

For the derivative bound, the same conductor estimate implies
\[
 \log\log(Q+3)
 \le\frac{\varepsilon}{5}\log Y+2\log\log Y+O_M(1).
\]
Equation~\eqref{qs-loglog} now yields
\[
 |F_\chi(1)|
 \le\frac{192\varepsilon}{5}\log Y
                 +384\log\log Y+O_M(1)
 \le4\log Y\le4\log y
\]
for sufficiently large $Y$. Here
$192\varepsilon/5<192/260<4$, leaving a fixed positive
margin. The zero-free assertion follows directly from
\eqref{qs-qrh}, since $399/400>11/12$.
\end{proof}

For the RH subcase take the principal base character of conductor
$1$. The three twists are precisely $\chi_q,\chi_d,\chi_q\chi_d$.
Since $q,d\equiv3\pmod4$, quadratic reciprocity identifies
$\chi_q(n)$ and $\chi_d(n)$ with the Kronecker characters
$(-q/n)$ and $(-d/n)$, including the supplementary rule at $2$.
For each odd prime $\ell\mid d+1$, $-d\equiv1\pmod\ell$,
so $\ell$ splits in $\mathbb Q(\sqrt{-d})$. The same holds for
$\ell=q$. At $2$, the congruence $d\equiv7\pmod{16}$ makes
$X^2-X+(d+1)/4$ have the two distinct roots $0,1$ modulo $2$.
Thus the required splitting statements are also supplied.

The argument is uniform if $z'$ is allowed to vary between
$z(\log Y)^2$ and $A^2z(\log Y)^2$: its lower endpoint
controls correlations, and its upper endpoint controls
$\log(Q+3)$. No constant depends on the growing primes or on
the intermediate scale $y$.

\subsection{The larger-scale character estimate}

The same analytic input supplies the separate small relative error
needed on the much larger scales. For nonprincipal $\chi$ modulo
$Q$, Mellin inversion of $(u-1)_+$ gives
\[
 \sum_{n\le y}\Lambda(n)\chi(n)(y/n-1)
 =\frac1{2\pi i}\int_{(2)}
             F_\chi(s)\frac{y^s}{s(s-1)}\,ds.
\]
Moving the line to $23/24$ crosses the simple pole with residue
$yF_\chi(1)$. Subtracting it gives
\begin{equation}\label{qs-large-scale}
 |G_\chi(y)|\ll y^{23/24}\log(Q+3)\qquad(y\ge1).
\end{equation}
The vertical integral is absolutely convergent by
\eqref{qs-logderivative}. For the primitive principal character
one uses instead
\[
 G_1(y)=\sum_{n\le y}\Lambda(n)(y/n-1)
                   -y(\log y-1-\gamma).
\]
Since $-\zeta'/\zeta(s)=1/(s-1)-\gamma+O(s-1)$, the residue
of the double pole is exactly the term subtracted here, and the
same argument gives $|G_1(y)|\ll y^{23/24}$.
This treats the principal case separately rather than assigning
it a nonexistent logarithmic derivative at $1$.

For example, if $\log(Q+3)\ll x^{\varepsilon/5}(\log x)^2$
and $y\ge\tfrac13\exp(x^{\varepsilon/3})$, then
\eqref{qs-large-scale} gives, for all sufficiently large $x$,
\[
 |G_\chi(y)|\le y\exp(-x^{\varepsilon/8}).
\]
This follows because $y^{-1/24}$ decays faster than the displayed
polynomial conductor factor and than $\exp(-x^{\varepsilon/8})$.
Thus a fixed QRH gap is sufficient for this larger-scale input
as well as for the simultaneous selection, though it is
insufficient for a direct square-root estimate for every fixed
character.

\subsection{A necessary restriction and the remaining scope}

The restriction on prime divisors of the progression modulus in
Lemma~\ref{qs-second-moment} has real content. If it were removed,
take a large integer $t$, put
$m=16\prod_{\substack{3\le\ell\le t\\\ell\ {\rm prime}}}\ell$,
and choose the reduced class
$a\equiv7\pmod{16}$, $a\equiv-1\pmod\ell$ for every such
odd prime. Quadratic reciprocity and the supplementary rule at
$2$ give $(n/p)=1$ for every $n\le t$ and every prime $p>t$
in this class. With $N=t$ and $\kappa=1$, the sum is therefore
$\psi(t)$ for every $p$. This is $\gg t$: indeed
\[
 \psi(2r)\ge\log\binom{2r}{r}
       \ge2r\log2-\log(2r+1),
\]
and monotonicity handles intervening real $t$.
The second moment is consequently $\gg t^2$, contradicting
$O(t\log t)$. Under \eqref{qs-qrh} the progression contains
primes for all sufficiently large averaging scales, so the
example is nonempty. This is a counterexample to dropping a
hypothesis, not to the actual selected-prime construction.

The finite second-moment and maximal lemmas have been proved here
without QRH. QRH is used to establish the weighted prime
correlations uniformly in the growing progression modulus, to
control the smooth logarithmic-derivative remainder, and to make
all required zero-free boxes automatic. The allowance of one
exceptional correlation also explains why the ordinary
exceptional-character formulation can be handled in the original
second moment; proving its full analytic input without QRH is a
separate route.

This section completes the selected-character part conditional on
the stated QRH theorem. It does not construct or estimate the
automorphic interpolation maps, their local minors, or their
nearly isotropic complex realizations. Nor does it supply primes
in varying quadratic ray classes. Those parts of the proposed
GRH proof remain separate obligations.


\section{Cyclotomic towers and a common interval for the cohomological bounds}
\label{sec:ct-intervals}

We reconstruct the part of the cohomological argument\cite{alpo-uniform-source} that treats the
prime $2$ and a fixed finite set of odd primes. The algebraic estimates
below allow finite torsion in the modules. The only deep classical
input in the uniform generator bound is the following precise form of
the Ferrero--Washington theorem: if $K/\Q$ is finite abelian, the
$p$-ranks of the ideal class groups of the fields
$K(\mu_{p^n})$ are bounded as $n$ varies. Here the $p$-rank of a finite
abelian group $A$ means $\dim_{\mathbb F_p} A/pA$.
This is the usual consequence of the vanishing of the Iwasawa
$\mu$-invariant and the structure theorem for the unramified Iwasawa
module; we use it as a classical input\cite{ct-fw,ct-sharifi}.
The final interval theorem is first stated with an explicit coarse
cohomological hypothesis. Theorem~\ref{kt-fourteen} below proves
that hypothesis, using precisely stated classical $K$-theory and
duality theorems.

\subsection{The local conditions and cyclotomic descent}

Fix a prime $p$ and a number field $K$ containing $\mu_p$ if $p$ is
odd, and $\mu_4$ if $p=2$. Put
$K_\infty=K(\mu_{p^\infty})$ and
$\Gamma=\operatorname{Gal}(K_\infty/K)$. The cyclotomic character
$\varepsilon_p$ identifies $\Gamma$ with $1+p^c\Z_p$, where
$c\ge1$, and $c\ge2$ for $p=2$. Equivalently, $c$ is the greatest
integer for which $\mu_{p^c}\subset K$. Choose a topological generator
$\gamma$ and put $u=\varepsilon_p(\gamma)$; thus
$v_p(u-1)=c$. Let $K_m$ be the subfield with
$[K_m:K]=p^m$.

For $e\ge1$, let $A_e=\Q_p/\Z_p(-e)$: its underlying group is
$\Q_p/\Z_p$ and $g$ acts by multiplication by
$\varepsilon_p(g)^{-e}$. Define
\[
 \mathcal S_p(F,e)=
 \ker\left(H^1(G_{F,S},A_e)\longrightarrow
                  \prod_{v\mid p}H^1(F_v,A_e)\right),
 \qquad S=\{v:v\mid p\}\cup\{v:v\mid\infty\}.
\]
Here $G_{F,S}$ is the Galois group of the maximal extension unramified
outside $S$, and $H^1$ is continuous group cohomology. In terms of the
absolute Galois group, this imposes zero restriction to inertia away
from $p$ and zero restriction to the \emph{full decomposition group}
at each prime above $p$. Define
\[
 \lambda_p(F,e)=\length_{\Z_p}\mathcal S_p(F,e)
       \quad\hbox{in }\mathbb N\cup\{\infty\}.
\]
When it is finite, $\#\mathcal S_p(F,e)=p^{\lambda_p(F,e)}$.

Let $D_p(K_\infty)$ be the discrete group of continuous homomorphisms
$G_{K_\infty}\to\Q_p/\Z_p$ that vanish on every finite inertia
group and on every decomposition group above $p$. Its Pontryagin
dual is the compact $\Z_p$-module
\[
 X_p(K_\infty)=
 \operatorname{Hom}_{\Z_p}(D_p(K_\infty),\Q_p/\Z_p).
\]
The topology on this dual is the topology of pointwise convergence
on finite subgroups. The action on $D_p$ is
$(\gamma\varphi)(h)=\varphi(\gamma^{-1}h\gamma)$; the action on
$X_p$ is the contragredient action. We write $T$ for the action of
$\gamma$ on $X_p$.

\begin{proposition}[Descent with the full local conditions]
\label{ct-descent}
For every $e\ge1$ and $m\ge0$, there is a canonical isomorphism
\[
 \mathcal S_p(K_m,e)^\vee\simeq
 X_p(K_\infty)\big/
       \bigl((u^eT)^{p^m}-1\bigr)X_p(K_\infty).
\]
In particular, the two sides have the same $\Z_p$-length, with
infinite length permitted.
\end{proposition}

\begin{proof}
The group $G_{K_\infty}$ acts trivially on $A_e$.
For a procyclic pro-$p$ group generated by $\tau$, continuous
cohomology on discrete $p$-primary modules is computed in degrees
$0$ and $1$ by $B\xrightarrow{\tau-1}B$, and vanishes above
degree $1$. Indeed, its completed group ring is $\Z_p[[t]]$,
with $t=\tau-1$, and
$0\to\Z_p[[t]]\xrightarrow{t}\Z_p[[t]]\to\Z_p\to0$
is the continuous free resolution of the trivial module.
For $B=A_e$, multiplication by
$u^{-ep^m}-1$ is surjective on $\Q_p/\Z_p$, since this is a
nonzero scalar and that group is divisible. The scalar is nonzero
because $u$ has infinite order and $e\ne0$. Inflation--restriction
therefore gives
\[
 H^1(K_m,A_e)\simeq
 H^1(K_\infty,A_e)^{\Gamma^{p^m}}.
\]

At a prime not above $p$, the cyclotomic extension is unramified,
so the inertia group is unchanged on passing to $K_\infty$.
At a prime above $p$, the image of the decomposition group in
$\Gamma$ is open and nontrivial: a finite extension of $\Q_p$
contains only finitely many $p$-power roots of unity. The same
divisibility argument makes the local inflation kernel zero.
Thus a global class has zero localization there if and only if its
restriction to the decomposition group over $K_\infty$ is zero.
The displayed global isomorphism consequently identifies the strict
Selmer group with the invariants of $D_p(K_\infty)(-e)$.

Pontryagin duality is exact here, since $\Q_p/\Z_p$ is injective
as a $\Z_p$-module. The dual of invariants is the corresponding
group of coinvariants. Dualizing the twist $-e$ replaces it by $e$;
the action on the dual is therefore $u^eT$. This proves the formula.
It also proves the asserted equality of lengths, by dualizing finite
composition series; either side has finite length precisely when
the other does.
\end{proof}

\subsection{A uniform bound as the quadratic fields vary}

We first give the elementary group argument that controls the cost
of a quadratic extension. If $V$ is profinite and $S\subset V$ is
stable under conjugation, write
\[
 C(V,S)=\{\varphi:V\to\mathbb F_2:
       \varphi\text{ is a continuous homomorphism},\ 
       \varphi|_S=0\}.
\]

\begin{lemma}[A quadratic extension]
\label{ct-index-two}
Let $U$ be an open subgroup of index $2$ in $V$, and suppose
$S\subset U$ is stable under conjugation by $V$. If
$\dim C(V,S)=r<\infty$, then
\[
                 \dim C(U,S)\le 2r.
\]
\end{lemma}

\begin{proof}
Choose $\sigma\in V\setminus U$. Its action $a$ on
$W=C(U,S)$ is an involution, because $\sigma^2\in U$ and inner
automorphisms act trivially on homomorphisms to an abelian group.
In characteristic $2$, $(a-1)^2=0$, so
$\operatorname{im}(a-1)\subset W^a$. It remains to show
$\dim W^a\le r$.

For $\varphi\in W^a$, the only obstruction to extending
$\varphi$ to $V$, with value $0$ at $\sigma$, is
$\varphi(\sigma^2)$. To check this directly, write each element
of $V$ uniquely as $h\sigma^i$ with $h\in U$ and $i\in\{0,1\}$,
and set $\widetilde\varphi(h\sigma^i)=\varphi(h)$.
The multiplication law gives a homomorphism exactly when
$\varphi$ is invariant and $\varphi(\sigma^2)=0$.
Continuity follows since $U$ is open. It still vanishes on $S$.
The restriction map from $C(V,S)$ to $W^a$ has a one-dimensional
kernel, the character of $V/U$, and its image has codimension at
most one in $W^a$. Hence $\dim W^a\le(r-1)+1=r$.
Finally, $W/W^a$ injects into $W^a$ via $a-1$, proving both
finiteness of $W$ and the asserted bound.
\end{proof}

\begin{lemma}[Descent of unramified characters]
\label{ct-character-descent}
Let $K_\infty/K$ be a cyclotomic $\Z_p$-extension.
Every continuous homomorphism
$G_{K_\infty}\to\mathbb F_p$ unramified at all finite primes
is the restriction of such a homomorphism on $G_{K_n}$ for some
$n$. A finite family can be descended to a common level.
\end{lemma}

\begin{proof}
Put $G=G_K$ and $H=G_{K_\infty}$. Choose an open normal subgroup
$N$ of $G$ such that $N\cap H$ lies in the kernel of the given
homomorphism $\varphi$. The rule $hn\mapsto\varphi(h)$ extends
it to $HN$: the ambiguity and the commutators $[n,h]$ lie in
$N\cap H$, where $\varphi$ vanishes. Since the decreasing compact subgroups $G_{K_n}$
have intersection $H$, some $G_{K_n}$ lies in $HN$.

At primes away from $p$, inertia lies in $H$, so the extension
already vanishes on inertia. The union of the conjugates, inside
$G_{K_n}$, of its finitely many inertia groups above $p$ is compact:
it is a finite union of continuous images of products of compact
groups. If the extension failed to vanish on that union at every
deeper level, fix this compact union $C$ at the initial level and
consider the decreasing compact sets
$C\cap G_{K_m}\cap\widetilde\varphi^{-1}(\mathbb F_p\setminus\{0\})$.
If all were nonempty, they
would have nonempty intersection. That would give an element of
$H$ in an inertia group on which $\varphi$ is nonzero, a
contradiction. A deeper level is therefore unramified also at $p$.
Taking the largest of finitely many levels gives the last assertion.
\end{proof}

\begin{proposition}[Uniform generator bound]
\label{ct-uniform-generators}
Let $M/\Q$ be finite abelian, of degree $d_M$ and conductor $f$.
There is an integer $g$, depending only on $M$, with the following
property. If $q,d$ are primes coprime to $2f$ satisfying
\[
                    q\equiv3\pmod8,
       \qquad       d\equiv7\pmod {16},
\]
and $L=M(i,\sqrt{-q},\sqrt{-d})$, then
$X_2(L_\infty)$ is generated over $\Z_2$ by at most
\[
                         R=4(g+6d_M)
\]
elements. This bound is independent of $q$ and $d$.
For any fixed abelian field $K$ containing the requisite roots of
unity, $X_p(K_\infty)$ is finitely generated over $\Z_p$ as well.
\end{proposition}

\begin{proof}
Apply Ferrero--Washington to the fixed field $M(i)$. By global
class field theory, the dimension of its unramified
$\mathbb F_2$-valued characters at each finite cyclotomic level
equals the $2$-rank of the ideal class group. Choose a uniform
bound $g$. Lemma~\ref{ct-character-descent}, applied to any finite
family, bounds the space of unramified characters of
$M(\mu_{2^\infty})$ by $g$.

The number of primes above an odd prime $a$ in
$\Q(\mu_{2^\infty})$ is the index of the closure of $\langle a\rangle$
in $\Z_2^\times$. For $a\equiv3\pmod8$ this index is $2$;
for $a\equiv7\pmod {16}$ it is $4$. Here is a direct verification.
If $a\equiv3\pmod4$ and $s=v_2(a^2-1)$, then
\[
 v_2(a^{2^k}-1)=s+k-1\quad(k\ge1).
\]
This follows by induction, since each subsequent factor
$a^{2^k}+1$ has valuation $1$. For $n\ge s$, the order of $a$
modulo $2^n$ is $2^{n-s+1}$ and the index is $2^{s-2}$.
The two congruences give $s=3$ and $s=4$, respectively.
Adjoining $M$ increases a number of primes by at most $d_M$.
There are consequently at most $2d_M$ primes above $q$ and
$4d_M$ above $d$ in $M(\mu_{2^\infty})$.

Allow ramification at those primes. The space of characters can
increase in dimension by at most one for each such prime. Indeed,
wild inertia has pro-odd order and every homomorphism from its
tame quotient to $\mathbb F_2$ is determined by one value.
Evaluation on one tame generator at each prime has the unramified
character space as its kernel. Conjugate primes give the same
condition. Thus the dimension after allowing $q,d$ is at most
$g+6d_M$.

Adjoin $\sqrt{-q}$ and then $\sqrt{-d}$. Each extension has
degree at most $2$ and is unramified outside the allowed primes:
the quadratic discriminants are $-q$ and $-d$.
Apply Lemma~\ref{ct-index-two} at every nontrivial step, taking
$S$ to be the union of the inertia groups at the other primes.
The dimension is multiplied by at most $4$. Imposing
unramifiedness again at $q,d$, and complete splitting above $2$,
can only decrease it. Therefore
\[
             \dim_{\mathbb F_2}D_2(L_\infty)[2]\le R.
\]

Duality identifies the dual of this last group with
$X_2(L_\infty)/2X_2(L_\infty)$. A compact pro-$p$ abelian group
$X$ with $\dim X/pX\le R$ is generated by $R$ elements as a
$\Z_p$-module. To see this without a freeness assumption, choose
lifts of a basis and let $Y$ be the image of $\Z_p^R$.
It is compact, hence closed. The quotient $X/Y$ satisfies
$X/Y=p(X/Y)$, so every finite $p$-group quotient is zero.
A profinite group is separated by its finite quotients, giving
$X/Y=0$. This proves the uniform assertion. The same descent
and duality argument, with Ferrero--Washington for a fixed $K$,
proves the last assertion.
\end{proof}

\begin{remark}
The congruences control the numbers of primes in the infinite
cyclotomic tower. Replacing them by mere oddness would remove this
particular uniform bound: for $a\equiv1\pmod {2^k}$ the closure
of $\langle a\rangle$ has index at least $2^{k-1}$ in
$\Z_2^\times$. This observation is not a counterexample to the
proposition, whose congruences are essential and are satisfied in
the submitted construction. The additional requirement $q\mid d+1$
is not needed for this proposition.
\end{remark}

\subsection{A determinant bound that retains finite torsion}

Write $v_p(0)=\infty$, and normalize $|p|_p=p^{-1}$.
If $X$ is generated by $R\ge1$ elements over $\Z_p$, and
$T$ is an endomorphism whose reduction on $X/pX$ is unipotent,
put
\[
 j=1+\min\{k\ge0:p^k\ge R\},\qquad
 b(e,m)=\length_{\Z_p}X/((u^eT)^{p^m}-1)X,
\]
where $v_p(u-1)=c\ge1$, with $c\ge2$ for $p=2$.

\begin{lemma}[Growth in a cyclotomic layer]
\label{ct-growth}
For all $e,m\ge0$,
\[
                    b(e,m)\le b(e,j)+R\max(m-j,0).
\]
\end{lemma}

\begin{proof}
Set $g=u^eT$. Its reduction modulo $p$ is unipotent on a vector
space of dimension at most $R$, so $(g-1)^R X\subset pX$.
If $p^k\ge R$, the binomial theorem in characteristic $p$ gives
$(g^{p^k}-1)X\subset pX$. Raising once more to the $p$th power
gives $(g^{p^{k+1}}-1)X\subset p^2X$: in the binomial expansion,
the first term is divisible by $p^2$, the intermediate terms have
an additional factor $p$, and the last has image in $p^pX$.
This proves the containment at level $j$, and at every larger level.

For any endomorphism $h$ with $(h-1)X\subset p^2X$, put
$Q=1+h+\cdots+h^{p-1}$. Then
$(Q-p)X\subset p^2X$, and hence
$QX\subset pX$ and $pX=QX+p^2X$.
Nakayama's lemma applied to the finitely generated module
$pX/QX$ gives $QX=pX$. As $Q$ commutes with $h-1$,
\[
                    (h^p-1)X=p(h-1)X.
\]
If $A=(h-1)X$, the exact sequence
$0\to A/pA\to X/pA\to X/A\to0$ shows that the cokernel
length increases by at most $R$: $A$ is an image of $X$, so
$\dim A/pA\le R$. The inequality remains true if a length is
infinite. Iteration gives the assertion for $m\ge j$.
For $m<j$, the factorization of $g^{p^j}-1$ by
$g^{p^m}-1$ gives a surjection of the cokernel at level $j$
onto the cokernel at level $m$.
\end{proof}

\begin{proposition}[At most $R$ exceptional weights]
\label{ct-centres}
There are $s_1,\ldots,s_R\in\Z_p$, independent of $e$ and $m$,
such that
\[
 b(e,m)\le
       \sum_{i=1}^R\bigl(c+\max(m,j)+v_p(e-s_i)\bigr).
\]
Consequently, for every integer $Z\ge1$ there is a set
$E_Z\subset\{1,\ldots,Z\}$ of cardinality at most $R$ such that
\[
 b(e,m)\le R\bigl(c+\max(m,j)+\lfloor\log_p Z\rfloor\bigr)
       \quad(1\le e\le Z,\ e\notin E_Z).
\]
In particular every length in this last inequality is finite.
\end{proposition}

\begin{proof}
Choose a surjection $\pi:\Z_p^R\to X$ and lift $T^{p^j}$ to a
matrix $B\in M_R(\Z_p)$, so that $\pi B=T^{p^j}\pi$.
For $a\in\Z_p$, the cokernel of $aB-1$ surjects onto that of
$aT^{p^j}-1$. Smith normal form gives
\[
 \length X/(aT^{p^j}-1)X\le v_p\det(aB-1).
\]
The convention $v_p(0)=\infty$ handles a singular matrix.
No torsion subgroup of $X$ has been discarded.

In a finite splitting field over $\Q_p$, let $z_i$ be the
eigenvalues of $B$, counted with multiplicity. The elementary
lifting-the-exponent formula gives
\[
 |u^{p^j a}-u^{p^j b}|_p
          =p^{-(c+j)}|a-b|_p\qquad(a,b\in\Z_p).
\]
For completeness, for integral $a-b$ prime to $p$ the first term
of the binomial expansion has strictly smallest valuation; raising
to the $p$th power increases that valuation by exactly one.
The condition $c\ge2$ when $p=2$ ensures strictness there.
Approximation by integers gives the displayed identity on $\Z_p$,
and defines the continuous powers used in it.

For each $z_i$ choose $s_i\in\Z_p$ that minimizes
$|z_i-u^{-p^j s}|_p$. Such a point exists by compactness.
The ultrametric inequality and the minimizing property give
\[
 |z_i-u^{-p^j e}|_p
       \ge |u^{-p^j e}-u^{-p^j s_i}|_p
       =p^{-(c+j)}|e-s_i|_p.
\]
Since $|u|_p=1$, the left side equals
$|u^{ep^j}z_i-1|_p$. Multiply these inequalities and use
$\det(u^{ep^j}B-1)=\prod_i(u^{ep^j}z_i-1)$.
This proves the bound at level $j$; Lemma~\ref{ct-growth}
proves it at every level.

Put $k=\lfloor\log_p Z\rfloor$. Each ball
$\{e\in\{1,\ldots,Z\}:|e-s_i|_p<p^{-k}\}$ contains at most
one integer. Two distinct integers in it would have difference
divisible by $p^{k+1}$, although its ordinary absolute value is
less than $p^{k+1}$. The union of these balls is the required
$E_Z$. Outside it, each $v_p(e-s_i)\le k$.
\end{proof}

For the module in Proposition~\ref{ct-descent}, the unipotence
hypothesis is automatic once it is finitely generated: the action
on the finite space $X_p/pX_p$ is continuous and factors through
a finite quotient of $\Gamma\simeq\Z_p$, hence through a
$p$-group. If $T^{p^a}=1$ on that space, then
$(T-1)^{p^a}=0$. The continuity follows either from the compact
dual construction or directly from the finite orbit of each
continuous finite-valued character. Thus all hypotheses in the
two preceding results hold for the actual cyclotomic modules.

\begin{example}[Why two tempting replacements fail]
Let $X=\Z/p^a\Z$, $a\ge1$, $T=1$, and $u=1+p^c$, with
the same restrictions on $c$. Then
\[
       b(e,m)=\min\{a,c+m+v_p(e)\}\quad(e\ge1).
\]
The equality follows from the lifting-the-exponent formula.
Passing to $X$ modulo its torsion would give the zero module and
would incorrectly make this positive length vanish. Also, the
strict inequality in the exceptional balls is necessary:
with $s=1$, $k\ge1$, and $Z=1+p^k<p^{k+1}$, both $1$ and $1+p^k$
have distance at most $p^{-k}$ from $s$. These examples refute
the proposed simplifications, not the actual results above or the
qualified source statements.
\end{example}

\subsection{Avoiding all the exceptional weights on one interval}

\begin{lemma}[A common multiplicative interval]
\label{ct-common-interval}
Let $E$ be a finite set of positive integers with $\#E\le R$,
let $C_t\ge\eta>0$, let $b>1$, and put
\[
             J=b^{R+1}(C_t/\eta)^R.
\]
For every $X>1$, there is $Y$ satisfying
\[
                 X^{1/J}\le Y,\qquad Y^b\le X,
\]
such that, simultaneously for every $x\in[Y,Y^b]$ and $e\in E$,
\[
                     e\le x^{C_t}\ \Longrightarrow\ e\le x^\eta.
\]
\end{lemma}

\begin{proof}
We prove the corresponding real assertion. Given $U>0$, $B\ge1$,
and a set $A\subset\R$ of at most $R$ points, there is
$y\ge U/[b(bB)^R]$ with $by\le U$ such that
$a\le v\Longrightarrow Ba\le v$ for all $a\in A$ and
$v\in[y,by]$.
Induct on $R$. If no point lies in $(U/(bB),U]$, take $y=U/b$:
every relevant point satisfies $Ba\le U/b$.
Otherwise choose a point $a$ in that interval, remove it, and
apply the induction hypothesis with $U'=U/(bB)$.
The resulting interval lies below $a$, so the removed point
imposes no condition; the bound on $y$ is exactly the required one.
For the empty set the choice $y=U/b$ starts the induction.
Now set $U=\log X$, $B=C_t/\eta$,
$A=\{\log(e)/C_t:e\in E\}$, and $Y=e^y$.
Taking logarithms gives the stated implications and endpoints.
\end{proof}

For $b=2$, this gives $J=2^{R+1}(C_t/\eta)^R$, improving the
factor $2^{R+2}$ in the corresponding source argument. More
significantly, the exceptional sets for all coefficient primes are
united \emph{before} the lemma is applied. Intersecting intervals
chosen separately would not prove that a common interval exists.

\subsection{The estimate at the exceptional weights}

We state the coarse estimate used at small weights as an explicit
hypothesis. There is an absolute $C_K\ge1$ such that, for every
totally complex number field $F$ containing $\mu_p$ (and $\mu_4$
when $p=2$) and every $e\ge1$, the group
$H^2(G_{F,S},\Z_p(e+1))$ is finite and
\begin{equation}
 \label{ct-coarse-hypothesis}
 \log\#H^2(G_{F,S},\Z_p(e+1))
 \le \exp\bigl((2e[F:\Q])^{C_K}\log(2|D_F|)\bigr).
\end{equation}
Here $D_F$ is the absolute field discriminant. Poitou--Tate
duality identifies the dual of $\mathcal S_p(F,e)$ with the
subgroup of this $H^2$ group whose localizations vanish. Thus
\begin{equation}
 \label{ct-small-weight}
 \lambda_p(F,e)\log p
 \le \log\#H^2(G_{F,S},\Z_p(e+1)).
\end{equation}
We use this precise consequence of duality; it is a statement
about a subgroup and includes finiteness. Theorem~\ref{kt-fourteen}
proves \eqref{ct-coarse-hypothesis} with $C_K=14$, using the
stated classical comparison theorems. Finite generation alone
would not give this order estimate.

\begin{proposition}[The common interval under the coarse estimate]
\label{ct-conditional-interval}
Assume \eqref{ct-coarse-hypothesis}. Fix $M$, a finite set $P$ of
odd primes, and for each $p\in P$ a fixed abelian field $K_p$
containing $\mu_p$. Fix $C_t\ge1$, $A>0$, and $0<\epsilon<1$.
Choose generator bounds $R_2,R_p\ge1$ as in
Proposition~\ref{ct-uniform-generators}; set
\[
 R_*=R_2+\sum_{p\in P}R_p,\quad
 \eta=\frac{\epsilon}{2C_K},\quad
 J=2^{R_*+1}(C_t/\eta)^{R_*}.
\]
For each coefficient prime use the fixed layer
$j_p=1+\min\{k:p^k\ge R_p\}$.
For all sufficiently large $X$, uniformly for the primes $q,d$
in Proposition~\ref{ct-uniform-generators} satisfying
\[
                     \log(qd)\le X^{\epsilon/(4J)},
\]
there is $Y$ with $X^{1/J}\le Y$ and $Y^2\le X$ such that,
for all $x\in[Y,Y^2]$, all integers $0\le h\le A(1+\log x)$,
and all integers $1\le t\le x^{C_t}$,
\begin{equation}
 \label{ct-summed-bound}
 (\log2)\sum_{1\le e<t}\lambda_2(E_h,e)
   +\sum_{p\in P}(\log p)
          \sum_{1\le e<t}\lambda_p((K_p)_{j_p},e)
        \le \exp(x^{4\epsilon/5}),
\end{equation}
where
$E_h=M\Q(\mu_{2^{h+2}})^+(\sqrt{-q},\sqrt{-d})$.
All lengths in this inequality are finite. The threshold for $X$
may depend on the fixed data and on $\epsilon$, but not on $q,d$.
\end{proposition}

\begin{proof}
Put $L=M(i,\sqrt{-q},\sqrt{-d})$. Its degree is at most
$8d_M$, so the largest $c$ with $\mu_{2^c}\subset L$ is bounded
in terms of $M$: $2^{c-1}\le[L:\Q]\le8d_M$.
Its fixed layer $L_{j_2}$ has bounded degree, and
\begin{equation}
 \label{ct-discriminant}
        \log(2|D_{L_{j_2}}|)\ll_M 1+\log(qd),
\end{equation}
where the constant may also depend on the fixed $j_2$.
Here is a direct verification of this last uniformity. Choose a
fixed integer $c_0$ bounding $c$, and put
$B=M(i,\mu_{2^{c_0+j_2}})$. The field $L_{j_2}$ is a subfield
of $B(\sqrt{-q},\sqrt{-d})$. For a quadratic extension
$F(\sqrt a)/F$ with integral $a$, the relative discriminant
divides $(4a)$: the order $\OO_F[\sqrt a]$ has that
discriminant and passage to the maximal order divides it by the
square of an index. If the extension is trivial the same upper
bound holds. Applying the discriminant tower formula twice gives
\[
 |D_{B(\sqrt{-q},\sqrt{-d})}|
 \le |D_B|^4(4q)^{2[B:\Q]}(4d)^{2[B:\Q]}.
\]
The discriminant tower formula also bounds the discriminant of
any subfield by that of this field. This proves
\eqref{ct-discriminant}. The odd-prime fixed layers have fixed
degree and discriminant.

Apply Proposition~\ref{ct-centres} with
$Z=\lfloor X^{C_t}\rfloor$ to all these modules. Unite their
exceptional sets; the union has at most $R_*$ elements. Lemma
\ref{ct-common-interval} gives the required $[Y,Y^2]$. For any
$x$ in it, every exceptional weight under consideration is at
most $x^\eta$. For a nonexceptional weight the length is
$O(1+\log X+h)=O(1+\log x)$, since
$\log X\le J\log x$. The constants are uniform in $q,d$.

For an exceptional weight, the growth lemma reduces to the fixed
layer, at a cost $O(1+\log x)$. At that layer, use
\eqref{ct-small-weight} and \eqref{ct-coarse-hypothesis}.
As $e\le x^\eta$, the degree is bounded, and
$X^{1/J}\le x$, equation \eqref{ct-discriminant} gives
\[
 (2e[F:\Q])^{C_K}\log(2|D_F|)
       \le C x^{\epsilon/2}x^{\epsilon/4}
       =C x^{3\epsilon/4}.
\]
Thus each exceptional length, multiplied by $\log p$, is at most
$C(1+\log x)+\exp(Cx^{3\epsilon/4})$.

To pass from the $2$-power layers of $L$ to $E_h$, observe that
$E_h(i)=L(\mu_{2^{h+2}})=L_m$ for some $m\le h+2$ and
$[E_h(i):E_h]\le2$. Restriction carries the strict local
conditions to the same strict conditions. Its kernel is a
subgroup of
$H^1(\operatorname{Gal}(E_h(i)/E_h),A_e^{G_{E_h(i)}})$,
which has order at most $2$. Indeed, this $H^1$ is killed by $2$
and is a subquotient of a subgroup of $\Q_2/\Z_2$, so its
order is at most $2$. Hence
\[
                    \lambda_2(E_h,e)\le1+\lambda_2(L_m,e).
\]

In each prime-wise sum use that prime's own exceptional set.
There are at most $R_*$ exceptional terms across these sums,
counted with multiplicity, and at most
$O(x^{C_t})$ other terms. The left side of
\eqref{ct-summed-bound} is therefore bounded by
\[
 Cx^{C_t}(1+\log x)+C\exp(Cx^{3\epsilon/4}).
\]
For fixed $C,C_t,\epsilon>0$ this is at most
$\exp(x^{4\epsilon/5})$ once $x$ is sufficiently large:
$3/4<4/5$ and every fixed multiple of $\log x$ is
$o(x^{4\epsilon/5})$. Increasing the threshold for $X$ makes
this hold throughout the chosen interval. The determinant bounds
and the coarse estimate separately give finiteness in the two
cases, so no infinite length has been silently converted to zero.
\end{proof}

This proves the common-interval mechanism and its uniformity in
the auxiliary quadratic fields. The initially displayed coarse
hypothesis is supplied by Theorem~\ref{kt-fourteen}.
The full cohomological estimate needed
for GRH additionally controls all the other primes and all relevant
character components. Those other estimates remain to be reconstructed. In
particular, this section does not replace the automorphy and
deformation arguments used for those components.


\section{Homology torsion and the remaining cohomological estimate}
\label{ht-section}

The coarse cohomological estimate used in
Section~\ref{sec:ct-intervals} has two different components: a bound for
torsion in arithmetic group homology, and a comparison with the
strict Selmer group.  We give complete proofs of the elementary
homological bounds and of the final local contribution to that
comparison. The geometric construction is supplied in
Section~\ref{ac-section}; the intervening global Selmer count
remains to be reconstructed. These distinctions matter: the submitted
archive contains both an alternative comparison through algebraic
$K$-theory and the class-counting route used in its final assembly.
An interface for the former does not prove the latter.

\subsection{An integral boundary matrix}

For an abelian group $M$, write $M_{\mathrm{tors}}$ for its torsion
subgroup.  Give $\mathbb Z^a$ its usual coordinate basis and the
norm $\|x\|_1=\sum_i|x_i|$.

\begin{lemma}\label{ht-matrix}
Let $L\subset\mathbb Z^a$ be generated by vectors of $\ell^1$ norm
at most $B$, where $B\ge1$.  The generating family may be infinite.
Then
\[
             \#(\mathbb Z^a/L)_{\mathrm{tors}}\le B^a.
\]
\end{lemma}
\begin{proof}
Put $r=\dim_{\mathbb Q}(L\otimes\mathbb Q)$ and
$L^{\mathrm{sat}}=(L\otimes\mathbb Q)\cap\mathbb Z^a$.
The torsion group in question is $L^{\mathrm{sat}}/L$.
If $r=0$, both lattices are zero.  Otherwise choose $r$ independent
members $v_1,\ldots,v_r$ of the given generating family and set
$L_0=\sum_i\mathbb Zv_i$.  Choose $r$ coordinate rows on which their
matrix has nonzero determinant.  Projection onto those coordinates
is injective on $L\otimes\mathbb Q$, so
\[
 [L^{\mathrm{sat}}:L]\le[L^{\mathrm{sat}}:L_0]
 \le\left|\det(v_1,\ldots,v_r)_{\mathrm{selected\ rows}}\right|
 \le\prod_i\|v_i\|_2\le B^r\le B^a.
\]
The middle inequality follows by placing the projected lattice
$L^{\mathrm{sat}}$ between the projected $L_0$ and $\mathbb Z^r$.
The determinant bound is Hadamard's inequality.
\end{proof}

\subsection{Finite stabilizers and an acyclic permutation complex}

For a group $\Gamma$, $H_j(\Gamma,\mathbb Z)$ denotes group homology
with trivial integral coefficients.  Its homogeneous bar complex
has in degree $j$ the free abelian group on tuples
$(g_0,\ldots,g_j)$, modulo simultaneous left translation, and
boundary the alternating sum of the tuples obtained by deleting
one coordinate.

\begin{lemma}\label{ht-finite-group}
If $S$ is a finite group of order $s$ and $j\ge1$, then
\[
 sH_j(S,\mathbb Z)=0,\qquad
 \#H_j(S,\mathbb Z)\le s^{s^j}.
\]
\end{lemma}
\begin{proof}
In the homogeneous complex define
\[
 A_j[g_0,\ldots,g_j]=\sum_{h\in S}[h,g_0,\ldots,g_j].
\]
The sum makes this well defined on the simultaneous-translation
quotient.  Deleting the first coordinate and then cancelling the
other terms gives
$\partial A_j+A_{j-1}\partial=s\,\mathrm{id}$ in positive
degrees.  Thus $s$ kills homology.  The degree-$j$ chain group
has rank $s^j$, since a tuple has a unique translate with first
coordinate $1$.  Its subgroup of cycles is free of rank at most
$s^j$.  Hence its homology quotient has at most $s^j$ generators.
A group killed by $s$ with that many generators has order at
most $s^{s^j}$.  This includes $s=1$, when the positive homology
is zero.
\end{proof}

An \emph{augmented acyclic permutation complex} for $\Gamma$ will
mean an exact complex of $\mathbb Z[\Gamma]$-modules
\[
 \cdots\longrightarrow \mathbb Z[X_2]\longrightarrow
 \mathbb Z[X_1]\longrightarrow\mathbb Z[X_0]
 \longrightarrow\mathbb Z\longrightarrow0,
\]
where each $X_i$ is a $\Gamma$-set, $\mathbb Z[X_i]$ means finite
integral linear combinations of its elements, and the augmentation
sends every member of $X_0$ to $1$.  No freeness of the
$\Gamma$-action is assumed.

\begin{proposition}\label{ht-permutation}
Fix $m\ge1$.  Suppose such a complex has at most $A$ orbits on
$X_i$ for $0\le i\le m$, all their stabilizers have order at most
$s$, and every column of the map
$\mathbb Z[X_{m+1}]\to\mathbb Z[X_m]$ has $\ell^1$ norm at most
$B$.  Here $A,s,B$ are positive integers.  Then
\begin{equation}\label{ht-torsion-bound}
 \#H_m(\Gamma,\mathbb Z)_{\mathrm{tors}}
 \le B^A s^{A\sum_{j=1}^m s^j}
 \le B^A s^{mAs^m}.
\end{equation}
In particular the torsion subgroup is finite.
\end{proposition}
\begin{proof}
Let $P_\bullet\to\mathbb Z$ be the free right bar
$\mathbb Z[\Gamma]$-resolution and form the first-quadrant double
complex $P_q\otimes_{\mathbb Z[\Gamma]}\mathbb Z[X_p]$.
Computing first in the $p$ direction gives the bar complex for
$H_*(\Gamma,\mathbb Z)$: $P_q$ is free, so tensoring it preserves
the exact augmented complex.  Computing first in the $q$ direction
gives a spectral sequence
\[
 E^1_{p,q}
 =\bigoplus_{\sigma\in\Gamma\backslash X_p}
           H_q(\Gamma_\sigma,\mathbb Z)
 \ \Longrightarrow\ H_{p+q}(\Gamma,\mathbb Z).
\]
For completeness, this identification uses only
$\mathbb Z[\Gamma/\Gamma_\sigma]
=\mathbb Z[\Gamma]\otimes_{\mathbb Z[\Gamma_\sigma]}\mathbb Z$.
The restricted complex $P_\bullet$ is still a free resolution
over $\mathbb Z[\Gamma_\sigma]$, so the corresponding tensor
complex computes the displayed stabilizer homology.  The two
filtrations of the total complex are bounded in every total
degree; thus the spectral sequence gives a finite filtration
in the degree under consideration.

For $q\ge1$ and $p+q=m$, Lemma~\ref{ht-finite-group} bounds the
order of $E^1_{p,q}$ by $s^{As^q}$.  Every later term in that
position is a subquotient, so has no greater order.
The row $q=0$ is the complex of coinvariants
$\mathbb Z[\Gamma\backslash X_p]$.  Passing to orbits cannot
increase a column's $\ell^1$ norm.  Its $m$th homology embeds
in the cokernel of the incoming boundary.  Therefore
Lemma~\ref{ht-matrix} bounds the torsion of $E^2_{m,0}$
by $B^A$.

For $r\ge2$ the homological differential has bidegree
$(-r,r-1)$.  Consequently no differential can enter
$(m,0)$ on such a page.  The successive terms there are
subgroups, rather than arbitrary quotients, of $E^2_{m,0}$.
Their torsion orders remain at most $B^A$.
Finally, in an exact sequence $0\to U\to V\to W\to0$
with finite torsion subgroups, the map
$V_{\mathrm{tors}}\to W_{\mathrm{tors}}$ has kernel contained
in $U_{\mathrm{tors}}$, so
$\#V_{\mathrm{tors}}\le\#U_{\mathrm{tors}}\#W_{\mathrm{tors}}$.
Apply this successively to the filtration of
$H_m(\Gamma,\mathbb Z)$.  Multiplying the bounds proves the
first inequality; $s^j\le s^m$ proves the second.
\end{proof}

For the ordered simplices in the arithmetic construction,
the boundary of an $(m+1)$-simplex has $m+2$ terms, with signs.
Repeated faces only decrease the $\ell^1$ norm.  Thus $B=m+2$.
Using ordered simplices is also what makes the modules permutation
modules with trivial stabilizer coefficients; an action on
unoriented simplices cannot silently be substituted if a stabilizer
reverses orientation.

\begin{example}\label{ht-stabilizer-example}
Take $\Gamma=\mathbb Z/\ell^a\mathbb Z$ and the augmented complex
with $X_0$ a single fixed point and all higher $X_i$ empty.
It is exact, and there is just one orbit, but
$H_1(\Gamma,\mathbb Z)=\mathbb Z/\ell^a\mathbb Z$.
Thus an orbit count alone gives no torsion bound.
The finite-stabilizer factor in
\eqref{ht-torsion-bound} cannot be omitted.
\end{example}

\subsection{The actual stabilizer bound}

\begin{lemma}\label{ht-congruence-three}
Let $R$ be a ring whose additive group is free of finite rank $d$
over $\mathbb Z$.  Reduction modulo $3$ is injective on every
finite subgroup of $R^\times$.  Such a subgroup has order
at most $3^d$.
\end{lemma}
\begin{proof}
It is enough to show that $1+3R$ contains no nonidentity unit
of finite order.  Otherwise a suitable power has prime order
$\ell$ and is still congruent to $1$ modulo $3$.
Write this power as $1+3^a b$, where $a\ge1$ and
$b\notin3R$.  Such a largest $a$ exists because $R$ is free
of finite rank.
If $\ell\ne3$, expand $(1+3^a b)^\ell=1$, divide the resulting
equality by $3^a$ in the torsion-free additive group, and reduce
modulo $3$.  This gives $\ell b=0$ in $R/3R$, a contradiction.
If $\ell=3$, division by $3^{a+1}$ instead gives
\[
                    b+3^a b^2+3^{2a-1}b^3=0,
\]
again contradicting $b\notin3R$ after reduction modulo $3$.
The binomial expansions are valid in a noncommutative ring
because only $1$ and the single element $b$ occur.
Two elements of a finite subgroup with the same reduction
have quotient in $1+3R$, proving injectivity.  The target
ring has $3^d$ elements.
\end{proof}

Let $F$ be totally imaginary of degree $n$, let
$\Gamma=\operatorname{GL}_N(\mathcal O_F)$, and embed it
diagonally at the complex places.  A positive definite
Hermitian form at every complex place has compact unitary
stabilizer.  The matrices over $\mathcal O_F$ form a closed
discrete subset of the product of complex matrix spaces,
since $\mathcal O_F$ is a lattice in its Minkowski embedding.
The arithmetic stabilizer is therefore finite.  The
stabilizer of an ordered tuple of such forms is a subgroup
of the stabilizer of its first entry.  Applying
Lemma~\ref{ht-congruence-three} to
$R=M_N(\mathcal O_F)$, whose integral rank is $nN^2$, gives
\begin{equation}\label{ht-arithmetic-stabilizer}
                         s\le3^{nN^2}.
\end{equation}
This proves the numerical stabilizer bound used in the
source, once positivity of its vertices has been supplied.

\subsection{The geometric input and its numerical consequence}

Put $D=|\operatorname{disc}F|$ and $d=nN$.
The geometric input is the existence of the above exact
permutation complex, formed from ordered tuples of vertices
of the Hermitian reduction polyhedron, with
\begin{equation}\label{ht-geometric-input}
 \#(\Gamma\backslash X_p)
       \le (2Dd)^{12(p+1)d^3N}.
\end{equation}
Its vertices must be positive definite, so that
\eqref{ht-arithmetic-stabilizer} applies.
The submitted files construct this complex, establish its
acyclicity, and obtain its orbit bound by integral matrix
reduction\cite{alpo-uniform-source}. Section~\ref{ac-section}
gives complete ordinary proofs of that construction and its bound.
We first prove the numerical consequence, so its dependence on
the geometric estimate is explicit.

\begin{proposition}\label{ht-arithmetic-numerical}
Assume this geometric input for $F$ and $N$.
If $E\ge1$, $1\le N\le4E+6$, and $0\le m\le2E+1$, then
\begin{equation}\label{ht-power-thirteen}
 \log\#H_m(\operatorname{GL}_N(\mathcal O_F),\mathbb Z)_{\mathrm{tors}}
       \le \exp\bigl((2En)^{13}\log(2D)\bigr).
\end{equation}
The exponent $13$ suffices here in place of the exponent $24$
used in the corresponding numerical estimate of the source.
This improvement uses $n\ge2$, which follows from $F$ being
totally imaginary.
\end{proposition}
\begin{proof}
For $m=0$ the homology is $\mathbb Z$, so the assertion is
immediate.  Suppose $m\ge1$.  Proposition~\ref{ht-permutation},
with
\[
 A\le\exp\bigl(12(m+1)d^3N\log(2Dd)\bigr),\quad
 s=3^{dN},\quad B=m+2,
\]
gives
\[
 \log\#H_m(\Gamma,\mathbb Z)_{\mathrm{tors}}
 \le A\{\log(m+2)+m s^m\log s\}.
\]
Since $\log s\le s$, $\log(m+2)\le m+1$, and
$2m+1\le e^{2m}$, the expression in braces is at most
$s^{m+1}e^{2m}$.  Since $\log3<2$, the last display is bounded by
\[
 \exp\{12(m+1)d^3N\log(2Dd)+2(m+1)dN+2m\}.
\]
Write $x=2En\ge4$ and $L=\log(2D)\ge\log2>1/2$.
The stated ranges give
\[
 m+1\le x,\qquad m\le x,\qquad
 d\le5x,\qquad N\le\tfrac52x.
\]
For $x\ge4$, $\log(5x)\le x$: it holds at $4$
because $e^4>20$, and the derivative of $x-\log(5x)$
is positive there and afterwards.  Consequently
$\log(2Dd)\le L+x\le(2x+1)L\le\tfrac94xL$.
The exponent in the preceding display is therefore at most
\[
 8438x^6L+25x^3+2x
       <8439x^6L\le x^{13}L.
\]
Indeed the first product has coefficient
$12\cdot5^3\cdot(5/2)\cdot(9/4)=8437.5$.
Also
$25x^3+2x\le(50/4^3+4/4^5)x^6L<x^6L$,
and $x^7\ge4^7=16384>8439$.  This proves the assertion.
\end{proof}

\subsection{The local terms in global duality}

Let $p$ be a rational prime and $F$ a totally imaginary number
field of degree $n$.  The negative twist
$A_e=\mathbb Q_p/\mathbb Z_p(-e)$, $e\ge1$, is as in
Section~\ref{sec:ct-intervals}.  For a place $v\mid p$, let
$U_v\subset\mathbb Z_p^\times$ be the image of the local
cyclotomic character and let $I_v$ be its index.

\begin{lemma}\label{ht-local-zero}
One has $I_v\le[F_v:\mathbb Q_p]\le n$ and
\begin{equation}\label{ht-local-bound}
 \length_{\mathbb Z_p}H^0(F_v,A_e)
       \le2+v_p(I_v)+v_p(e)
       \le2+\lfloor\log_p n\rfloor+v_p(e).
\end{equation}
\end{lemma}
\begin{proof}
The cyclotomic extension of $\mathbb Q_p$ has Galois group
$\mathbb Z_p^\times$.  Its intersection with $F_v$ is finite,
and restriction identifies the index $I_v$ with the degree
of that intersection.  This proves the first assertion.
Every element of $\mathbb Z_p^\times$ raised to $I_v$
belongs to $U_v$.  In particular
$a=(1+p^2)^{I_v}$ belongs to $U_v$.  An invariant in $A_e$
is killed by $a^{-e}-1$, or equivalently by $a^e-1$.
The elementary lifting-the-exponent formula, also valid at
$p=2$ because $1+p^2\equiv1\pmod4$, gives
\[
             v_p((1+p^2)^{I_ve}-1)=2+v_p(I_v)+v_p(e).
\]
The subgroup of $\mathbb Q_p/\mathbb Z_p$ killed by an element
of valuation $b$ is cyclic of order $p^b$.  This proves
\eqref{ht-local-bound}.  Finally $p^{v_p(I_v)}\le I_v\le n$.
\end{proof}

Let $\mathcal S_p(F,e)$ be the strict Selmer group defined in
Section~\ref{sec:ct-intervals}; it has zero localization at every
$v\mid p$.  The form of Poitou--Tate and local Tate duality
used here identifies
\[
 \ker\!\left(H^2(G_{F,S},\mathbb Z_p(e+1))
              \longrightarrow
              \bigoplus_{v\mid p}H^2(F_v,\mathbb Z_p(e+1))\right)
          \simeq \mathcal S_p(F,e)^\vee
\]
and $H^2(F_v,\mathbb Z_p(e+1))\simeq H^0(F_v,A_e)^\vee$.
These are classical duality inputs, with the same twists
and local conditions as in the earlier section; no proof of
the general duality theorems is claimed here.
Lemma~\ref{ht-local-zero} now proves the following consequence.

\begin{proposition}\label{ht-h2-from-selmer}
If $\mathcal S_p(F,e)$ is finite, so is
$H^2(G_{F,S},\mathbb Z_p(e+1))$, and
\begin{equation}\label{ht-h2-additive}
 \log\#H^2(G_{F,S},\mathbb Z_p(e+1))
 \le\log\#\mathcal S_p(F,e)
       +n\{2+\lfloor\log_p n\rfloor+v_p(e)\}\log p.
\end{equation}
If $F$ contains $\mu_p$ and, for some $C\ge1$,
\[
 \log\#\mathcal S_p(F,e)
       \le\exp\bigl((2en)^C\log(2D)\bigr),
\]
then the coarse $H^2$ estimate of Section~\ref{sec:ct-intervals}
holds with $C_K=C+2$.
\end{proposition}
\begin{proof}
Global duality embeds the quotient by the displayed kernel
into the direct sum of local groups.  They are finite by
Lemma~\ref{ht-local-zero}; there are at most $n$ places above
$p$.  Taking lengths proves \eqref{ht-h2-additive}, and in
particular proves finiteness rather than presupposing it.

Since $\mu_p\subset F$, one has $p-1\le n$.
Thus the local contribution is at most
\[
 n\{2\log(n+1)+\log n+\log e\}.
\]
Put $x=2en\ge4$ and $L=\log(2D)>1/2$.  The last expression
is at most $2x^2$, since each logarithm is at most $x$ and
$n\le x/2$.  Moreover $2x^2\le e^{x^2L}$ for $x\ge4$:
the inequality holds at $4$ with $L=1/2$, and
$x^2/2-\log(2x^2)$ is increasing afterwards.
The total in \eqref{ht-h2-additive} is therefore at most
\[
 e^{x^CL}+e^{x^2L}
 \le2e^{x^{\max(C,2)}L}
 \le e^{x^{\max(C,2)+1}L}
 \le e^{x^{C+2}L}.
\]
The penultimate inequality follows because
$(x-1)x^{\max(C,2)}L>\log2$.
\end{proof}

The last proposition clarifies the remaining task.
The source's actual class-counting argument must bound the
strict Selmer group uniformly in the field and the weight.
It uses the homology bound, finite-field homology classes,
relations among their diagonal translates, a class-field
detection step, and an odd-degree rational rank bound.
The new homology argument proves the part that follows from
the exact arithmetic complex and its orbit count, and
Proposition~\ref{ht-h2-from-selmer} proves the final passage
from a strict Selmer bound to $H^2$, with the sharper local
term $v_p(e)$ in place of $e$. Section~\ref{ac-section}
establishes the geometric input \eqref{ht-geometric-input}.
The intervening global class count remains an explicit proof
obligation in our reconstruction, not an alleged missing
assumption of the submitted formal proof.


\section{The arithmetic complex and its orbit bound}
\label{ac-section}

This section proves the geometric input in
\eqref{ht-geometric-input}.  In particular, the homology estimate in
Proposition~\ref{ht-arithmetic-numerical} no longer requires that input
as an additional hypothesis.  The construction is the ordered nerve
of the normal cones of the Hermitian reduction polyhedron used in
the submitted source\cite{alpo-uniform-source}.  We give the exact
arithmetic definitions and prove its exactness, the stabilizer bound,
and its orbit count.  This does not prove the separate global Selmer
class count discussed at the end of Section~\ref{ht-section}.

\subsection{The lattice and the Hermitian reduction polyhedron}

Let $F$ be a totally imaginary number field, let
$n=[F:\mathbb Q]=2s$, and let
$D=|\operatorname{disc}F|$.  Fix one embedding
$\sigma_v:F\longrightarrow\mathbb C$ from each conjugate pair.
Let $N\ge1$, and put
\[
 \Lambda=\mathcal O_F^N,\qquad
 d=nN,\qquad
 \Gamma=\operatorname{GL}_N(\mathcal O_F),\qquad
 \mathcal H=\prod_{v=1}^s\operatorname{Herm}_N(\mathbb C).
\]
Here $\operatorname{Herm}_N(\mathbb C)$ is the real vector space of
Hermitian $N$ by $N$ matrices.  Its dimension is $N^2$, and
$\dim_{\mathbb R}\mathcal H=sN^2\le dN$.  For $H=(H_v)\in\mathcal H$
and $y\in\Lambda$, define
\[
 q_H(y)=\sum_{v=1}^s \sigma_v(y)^*H_v\sigma_v(y).
\]
Every summand is real.  The star denotes conjugate transpose;
no involution of the number field is assumed.
Write
\[
 \begin{aligned}
 \mathcal P&=\{H:H_v\text{ is positive semidefinite for every }v\},\\
 \mathcal R&=\{H:q_H(y)\ge1\text{ for every }0\ne y\in\Lambda\}.
 \end{aligned}
 \tag{AC1}
\]
The group acts by
\[
 (gH)_v=\sigma_v(g)^{-*}H_v\sigma_v(g)^{-1}.
 \tag{AC2}
\]
Consequently $q_{gH}(y)=q_H(g^{-1}y)$, and the action preserves
$\mathcal R$ and $\mathcal P$.

We use the elementary facts that $\mathcal O_F$ is a free abelian
group of rank $n$, that the field norm of a nonzero algebraic integer
is a nonzero integer, and that the discriminant of an integral basis
is nonzero.  The volume normalization needed below follows directly
from these facts.  If $\omega_1,\ldots,\omega_n$ is an integral basis,
the determinant of the matrix of all $n$ embeddings has absolute
value $\sqrt D$.  Replacing each conjugate pair of rows by its real
and imaginary parts divides its absolute determinant by $2$.
Thus the image of $\mathcal O_F$ in $\mathbb C^s=\mathbb R^n$
is a full lattice of covolume $2^{-s}\sqrt D$.

\begin{lemma}\label{ac-positive-gram}
Let $G$ be a real symmetric $r$ by $r$ matrix, with $r\ge1$.
If
\[
                 z^{\mathsf T}Gz\ge1
       \quad(0\ne z\in\mathbb Z^r),
\]
then $G$ is positive definite and $\det G\ge r^{-r}$.
\end{lemma}
\begin{proof}
First $G$ is positive semidefinite.  Otherwise continuity supplies
a rational vector on which its quadratic form is negative, and
clearing denominators contradicts the hypothesis.

We recall and prove the elementary lattice estimate used here.
If a measurable set $A\subset\mathbb R^r$ has volume greater than
the covolume of a lattice $L$, then two distinct points of $A$
differ by a lattice vector.  Indeed, let $P$ be a half-open
fundamental parallelepiped.  Integration over $P$ of
$\sum_{\ell\in L}\mathbf1_A(x+\ell)$ gives $\operatorname{vol}(A)$.
If no two such points existed, this integer-valued sum would be at
most one everywhere, contradicting the volume inequality.
Applying this assertion to one half of an open centrally symmetric
convex set proves that such a set of volume greater than
$2^r\operatorname{covol}(L)$ contains a nonzero lattice point:
the difference of the two points belongs to the original set.

Suppose first that $G$ is positive definite, and choose a real
invertible matrix $T$ with $G=T^{\mathsf T}T$.  Apply the preceding
argument to $T\mathbb Z^r$ and the cube $(-R,R)^r$, where
$R>|\det T|^{1/r}$.  For some $0\ne z\in\mathbb Z^r$,
\[
                  1\le\|Tz\|^2\le rR^2.
\]
Letting $R$ decrease to $|\det T|^{1/r}$ gives
$1\le r(\det G)^{1/r}$.  For a merely positive semidefinite $G$,
apply this bound to $G+\varepsilon I$ and let
$\varepsilon\downarrow0$.  Its determinant is at least $r^{-r}$
in the limit.  Hence $\det G>0$, and positive semidefiniteness
then implies positive definiteness.
\end{proof}

\begin{proposition}\label{ac-polyhedron}
The set $\mathcal R$ is nonempty, closed and convex, and every
$H\in\mathcal R$ is positive definite at every complex place.
The defining inequalities are locally finite on $\mathcal R$:
for every $H\in\mathcal R$, there are a neighborhood $U$ of $H$
and a finite set $T\subset\Lambda\setminus\{0\}$ such that
\[
          q_K(y)>1\qquad(K\in U,\ y\notin T).
          \tag{AC3}
\]
\end{proposition}
\begin{proof}
Closedness and convexity follow because each $q_H(y)$ is real
linear in $H$.  The tuple of identity matrices belongs to
$\mathcal R$.  Indeed, a nonzero coordinate $a\in\mathcal O_F$
of $y$ has $|\sigma_v(a)|\ge1$ for at least one $v$;
otherwise $|\operatorname{Norm}_{F/\mathbb Q}(a)|<1$.

Choose the integral basis above in every coordinate of $\Lambda$.
Its embedding identifies $\mathbb Z^d$ with a full lattice in
$\prod_v\mathbb C^N$, and identifies the associated real vector
spaces.  The real quadratic form $q_H$ therefore has a symmetric
Gram matrix $G_H$ in these coordinates.  For $H\in\mathcal R$,
Lemma~\ref{ac-positive-gram} makes $G_H$ positive definite.
Since the embedding is a real vector-space isomorphism, a vector
supported at any one complex place shows that $H_v$ is positive
definite at that place.

Finally choose $c>0$ and a sufficiently small neighborhood $U$
such that
\[
 q_K(z)\ge c\sum_v\|z_v\|^2
       \quad(K\in U,\ z\in\prod_v\mathbb C^N).
\]
The lattice has only finitely many points in the ball
$\sum_v\|\sigma_v(y)\|^2\le c^{-1}$.  Taking these points for $T$
proves (AC3).
\end{proof}

For $H\in\mathcal R$ let
\[
                  M(H)=\{y\in\Lambda:q_H(y)=1\}.
\]
This is a finite set, and does not contain zero.
Call $H$ a vertex if the linear functionals
$K\mapsto q_K(y)$, for $y\in M(H)$, span $\mathcal H^*$.
Equivalently,
\[
 q_K(y)=0\ (y\in M(H))\quad\Longrightarrow\quad K=0.
 \tag{AC4}
\]
This is the actual vertex definition in the source.  It agrees
with the usual definition as an extreme point of $\mathcal R$.
To see this, if (AC4) holds and $H$ is the midpoint of
$H_1,H_2\in\mathcal R$, all its active inequalities are active
at both endpoints, so $H_1-H_2=0$.
Conversely, if a nonzero $K$ annihilates all active inequalities,
(AC3) and the finitely many remaining strict inequalities show
that $H\pm\varepsilon K\in\mathcal R$ for sufficiently small
$\varepsilon>0$.  Thus $H$ is not extreme.
We will repeatedly use this last perturbation argument.

\subsection{Normal cones and an exact ordered complex}

Let $V$ be the set of vertices of $\mathcal R$.  Define
\[
 \mathcal U=\{\ell\in\mathcal H^*:
            \ell(K)>0\text{ for all }0\ne K\in\mathcal P\},
 \qquad
 C_H=\{\ell:\ell(H)\le\ell(K)\text{ for all }K\in\mathcal R\}.
 \tag{AC5}
\]
The cone $\mathcal U$ is convex and nonempty: the sum of the
matrix traces belongs to it.  Each $C_H$ is closed and convex.

\begin{lemma}\label{ac-normal-cover}
The sets $C_H$, $H\in V$, cover $\mathcal U$.
For every finite nonempty set $E\subset\mathcal U$, only finitely
many of these cones meet $\operatorname{conv}(E)$.
\end{lemma}
\begin{proof}
Fix $\ell\in\mathcal U$.  On the intersection of $\mathcal P$
with the unit sphere, $\ell$ has a strictly positive minimum.
It follows that $\ell(K)\ge c\|K\|$ on $\mathcal P$ for some
$c>0$.  Therefore every sublevel
\[
          \{K\in\mathcal R:\ell(K)\le a\}
\]
is compact.  As $\mathcal R$ is nonempty, $\ell$ attains its
minimum there, and its minimizing face is nonempty and compact.
A nonempty compact convex set in a finite-dimensional real
space has an extreme point: one can successively maximize the
coordinates of a fixed basis, restricting to the maximizing
face after each step, until a singleton remains.
An extreme point of this minimizing face is extreme in
$\mathcal R$, because the value of $\ell$ at each endpoint
of any segment with midpoint in the face must also be minimal.
Hence it belongs to $V$, proving the covering assertion.

We need one consequence of (AC3).  A compact subset
$K\subset\mathcal R$ contains only finitely many vertices.
Cover $K$ by finitely many neighborhoods as in (AC3), and let
$T$ be the union of their finite sets of inequalities.
Every active set $M(H)$ for $H\in K$ is contained in $T$.
For a vertex, this set determines $H$ uniquely by the
equations $q_H(y)=1$.  There are only finitely many subsets
of $T$.

Now fix $H_0\in\mathcal R$, and write
$E=\{\ell_1,\ldots,\ell_t\}$.  If $C_H$ meets
$\operatorname{conv}(E)$, some convex combination
$\ell=\sum_i a_i\ell_i$ is minimized at $H$.  Thus
\[
                 \sum_i a_i\{\ell_i(H)-\ell_i(H_0)\}\le0.
\]
For some $i$ with $a_i>0$, $\ell_i(H)\le\ell_i(H_0)$.
The vertex $H$ therefore lies in the finite union of compact
sublevels defined by these $t$ inequalities.  The preceding
paragraph proves the desired finiteness.
\end{proof}

For $p\ge0$, define $X_p$ to be the set of ordered tuples
$(H_0,\ldots,H_p)$ of vertices such that
\[
                   \mathcal U\cap C_{H_0}\cap\cdots
                            \cap C_{H_p}\ne\varnothing.
                   \tag{AC6}
\]
Repetitions are allowed.  Deleting an entry defines a map
$X_p\to X_{p-1}$.  Put
\[
 C_p=\mathbb Z[X_p],\quad
 \partial(H_0,\ldots,H_p)
      =\sum_{i=0}^p(-1)^i(H_0,\ldots,\widehat H_i,\ldots,H_p),
 \quad
 \epsilon(H)=1.
 \tag{AC7}
\]
Here $\mathbb Z[X]$ denotes the free abelian group on $X$,
with finite support.  The familiar cancellation of pairs of
deletions gives $\partial^2=0$.

We prove the needed exactness directly, including the issue
that the cover consists of closed sets.

\begin{lemma}\label{ac-convex-nerve}
Let $K$ be a nonempty compact convex subset of a
finite-dimensional real vector space, covered by finitely many
closed convex sets $A_i$.  Form the ordered tuple complex of
the nerve: a tuple is allowed precisely when
$K\cap\bigcap_j A_{i_j}\ne\varnothing$.  Its augmented integral
chain complex, with repetitions and the boundary (AC7), is exact.
\end{lemma}
\begin{proof}
Replace $A_i$ by $K\cap A_i$ and discard empty sets.
We induct on the number of subsets of the finite index set
whose corresponding intersection is nonempty, counting the
empty subset as well.

If all the $A_i$ have a common point, every tuple of indices
is allowed.  Fix an index $a$.  Prepending $a$ to a tuple
defines a contraction: for $p\ge0$,
\[
       \partial(a,i_0,\ldots,i_p)
          =(i_0,\ldots,i_p)-(a,\partial(i_0,\ldots,i_p)),
\]
and in degree $-1$ send $1$ to $(a)$.  This proves exactness.

Otherwise choose a minimal collection of the $A_i$ with
empty intersection.  Write it as $\{A_i\}\cup\{A_j:j\in J\}$.
Both $A_i$ and $B=\bigcap_{j\in J}A_j$ are nonempty compact
convex sets, and they are disjoint.
They admit a strict separating linear functional $f$:
\[
                         f(A_i)<u<f(B)
                         \tag{AC8}
\]
for some real number $u$.
For completeness, take a pair of nearest points $a\in A_i$,
$b\in B$.  Their distance is positive by compactness.
With $f(x)=\langle b-a,x\rangle$, differentiating the squared
distance along segments in the two convex sets gives
$f(x)\le f(a)<f(b)\le f(y)$ for $x\in A_i$, $y\in B$.
Any number between $f(a)$ and $f(b)$ gives (AC8).

Set $K_-=K\cap\{f\le u\}$, $K_+=K\cap\{f\ge u\}$,
and $K_0=K\cap\{f=u\}$.  They are nonempty compact convex
sets, covered by the corresponding intersections with $A_i$.
All three nerves have strictly fewer faces: $K_-$ and $K_0$
lose the face $J$, and $K_+$ loses the vertex $i$.
Their augmented complexes are therefore exact by induction.

The nerve on $K$ is the union of the nerves on $K_-$ and
$K_+$.  Their intersection is the nerve on $K_0$.
Indeed, witnesses on both sides for the same tuple can be
joined by a segment lying in every set of that tuple, and
this segment meets the separating hyperplane.
These statements remain true for ordered tuples with repetitions.

Here is the chain argument for this union.  Decompose a cycle
$z$ as $z_-+z_+$, using chains in the two subcomplexes.
Then $\partial z_-=-\partial z_+$ is supported in their
intersection.  Exactness there supplies a chain $w$ with
$\partial w=\partial z_-$.  The chains $z_--w$ and $z_++w$
are cycles in their respective subcomplexes, hence boundaries.
Their sum is $z$.  In degree zero the same argument uses
surjectivity of the augmentation of the intersection; in
degree $-1$ surjectivity follows from nonemptiness.
This completes the induction.
\end{proof}

\begin{theorem}\label{ac-exact}
The augmented sequence
\[
       \cdots\longrightarrow\mathbb Z[X_2]
       \longrightarrow\mathbb Z[X_1]
       \longrightarrow\mathbb Z[X_0]
       \overset{\epsilon}{\longrightarrow}\mathbb Z
       \longrightarrow0
\]
is an exact sequence of permutation $\mathbb Z\Gamma$-modules.
Every column of the boundary
$\mathbb Z[X_{p+1}]\to\mathbb Z[X_p]$ has integral
$\ell^1$-norm at most $p+2$.
\end{theorem}
\begin{proof}
The action (AC2), and its contragredient action on linear
functionals, preserve (AC6).  Deletions and augmentation are
equivariant, so the modules are permutation modules.
The boundary estimate counts its $p+2$ signed terms;
coinciding faces can only decrease the $\ell^1$-norm.

A cycle has finite support.  Choose a witness in $\mathcal U$
for each tuple in its support, and choose an additional point
of $\mathcal U$ if necessary.  Their convex hull $K$ is
nonempty, compact, and contained in $\mathcal U$.
By Lemma~\ref{ac-normal-cover}, only finitely many $C_H$
meet $K$, and they cover $K$.  Apply
Lemma~\ref{ac-convex-nerve} to their intersections with $K$.
It supplies a filling using only tuples already allowed
by (AC6).  The same covering lemma ensures that $X_0$
is nonempty, so augmentation is surjective.
\end{proof}

\subsection{Common minimal vectors and determinant bounds}

\begin{lemma}\label{ac-common-minima}
For every tuple $(H_0,\ldots,H_p)\in X_p$, the common minimal
vectors
\[
                 \bigcap_{j=0}^p M(H_j)
\]
span $F^N$ over $F$.
\end{lemma}
\begin{proof}
Let $\ell\in\mathcal U$ witness (AC6), and put
$\overline H=(p+1)^{-1}\sum_jH_j$.
Each $H_j$ minimizes $\ell$ on $\mathcal R$, so
$\overline H$ does too.  Since every $q_{H_j}(y)\ge1$,
\[
                  M(\overline H)=\bigcap_jM(H_j).
\]
Suppose this set does not span $F^N$.  There exists a nonzero
$F$-linear functional $a:F^N\to F$ vanishing on it.
At each place let $K_v$ be the positive semidefinite rank-one
form whose value at $z$ is $|\sigma_v(a)z|^2$.
This defines a nonzero $K\in\mathcal P$ with
$q_K(y)=0$ for every $y\in M(\overline H)$.
The perturbation argument following (AC4) gives
$\overline H-\varepsilon K\in\mathcal R$ for some
$\varepsilon>0$.  But
$\ell(\overline H-\varepsilon K)<\ell(\overline H)$,
a contradiction.
\end{proof}

Consequently we may choose $N$ common minimal vectors as the
columns of an integral matrix $A\in M_N(\mathcal O_F)$ with
$\det A\ne0$.  Such columns are an $F$-basis, and are not
claimed to be an $\mathcal O_F$-basis of $\Lambda$.
Define the integers
\[
       I_0=D^N d^d,\qquad R_0=N^N I_0^2,\qquad
       I(A)=|\operatorname{Norm}_{F/\mathbb Q}(\det A)|.
       \tag{AC9}
\]

\begin{lemma}\label{ac-det}
For every $H\in\mathcal R$,
\[
                  I_0\prod_v\det H_v\ge1.
                  \tag{AC10}
\]
If the columns of $A$ are minimal for $H$, put
$H'_v=\sigma_v(A)^*H_v\sigma_v(A)$.  Then
\[
 \begin{split}
 &1\le I(A)\le I_0,\\
 &H'_v>0,\qquad (H'_v)_{aa}\le1,\qquad
                I_0^{-1}\le\det H'_v\le1
                   \quad\text{for every }v,a.
 \end{split}
 \tag{AC11}
\]
\end{lemma}
\begin{proof}
Let $G_H$ be the real Gram matrix in the integral coordinates
used above.  Its determinant is
\[
       \det G_H=2^{-d}D^N\left(\prod_v\det H_v\right)^2.
       \tag{AC12}
\]
Indeed the squared determinant of the lattice embedding is
$2^{-d}D^N$.  For a positive Hermitian matrix, its quadratic
form on the underlying real space has determinant equal to
the square of its complex determinant: unitary diagonalization
reduces this to the fact that each eigenvalue occurs twice.
This proves (AC12).

Write $P=\prod_v\det H_v>0$.
Lemma~\ref{ac-positive-gram} gives
$P^2\ge 2^dD^{-N}d^{-d}$, and hence
\[
                       (I_0P)^2\ge D^N(2d)^d\ge1.
\]
This proves (AC10); it also shows that our intermediate
determinant estimate is stronger than the one required.

For each column $a$ of $A$, minimality says
$\sum_v(H'_v)_{aa}=1$.  Every term is positive, so
$(H'_v)_{aa}\le1$.  Hadamard's inequality gives
$0<\det H'_v\le1$.  This inequality follows, for example,
by writing $H'_v$ as a Gram matrix: successive orthogonal
projections do not increase the lengths of its Gram vectors.
The determinant transformation formula gives
\[
                  \prod_v\det H'_v=I(A)P.
                  \tag{AC13}
\]
Here each complex embedding is represented once, so
$\prod_v|\sigma_v(\det A)|^2=I(A)$.
The left side of (AC13) is at most one, whence
$I(A)\le P^{-1}\le I_0$.
Each factor on that left side is at least their product,
because all the other factors are at most one.  Since
$I(A)\ge1$, this proves
$\det H'_v\ge I(A)P\ge I_0^{-1}$.
\end{proof}

\subsection{Integral matrix reduction without a principal ideal hypothesis}

\begin{lemma}\label{ac-scaled-inverse}
For $A\in M_N(\mathcal O_F)$ with $\det A\ne0$, the matrix
\[
                           B_A=I(A)A^{-1}
                           \tag{AC14}
\]
belongs to $M_N(\mathcal O_F)$.
For every integer $J\ge1$, there is a set $\mathscr A_J$
of at most $J^{1+dN}$ integral matrices such that every
$A$ with $0<I(A)\le J$ satisfies
$gA\in\mathscr A_J$ for some $g\in\Gamma$.
\end{lemma}
\begin{proof}
If $0\ne\alpha\in\mathcal O_F$, the group
$\mathcal O_F/\alpha\mathcal O_F$ has order
$|\operatorname{Norm}_{F/\mathbb Q}(\alpha)|$.
Indeed multiplication by $\alpha$ on an integral basis has
that determinant, and the index of an integral full-rank
matrix is the absolute value of its determinant, as follows
from integer elementary row and column reduction.
Lagrange's theorem in this finite additive group therefore
gives
\[
          |\operatorname{Norm}_{F/\mathbb Q}(\alpha)|
                      \in\alpha\mathcal O_F.
\]
For $\alpha=\det A$, the adjugate formula now proves (AC14).
In particular $AB_A=B_AA=I(A)I_N$.

Fix a positive integer $I$, and consider only matrices with
$I(A)=I$.  Assign to $A$ the residue of $B_A$ in
$M_N(\mathcal O_F/I\mathcal O_F)$.
If $A,A'$ have the same residue, write
$B_A-B_{A'}=IX$ with $X\in M_N(\mathcal O_F)$.
Then
\[
       A'A^{-1}=I_N+A'X,\qquad
       AA'^{-1}=I_N-AX.
       \tag{AC15}
\]
Both matrices on the right are integral and are inverse
to each other.  Thus $g=A'A^{-1}$ belongs to $\Gamma$
and satisfies $gA=A'$.

There are exactly $I^{nN^2}=I^{dN}$ possible residues,
because the additive rank of $\mathcal O_F$ is $n$.
Choose one representative matrix for each residue that
actually occurs.  Equal residues imply equal left
$\Gamma$-orbits by (AC15), so these representatives meet
every such orbit.  They need not be distinct orbit
representatives, which only improves the upper bound.
Summing over $1\le I\le J$ gives at most
\[
                         \sum_{I=1}^J I^{dN}
                             \le J^{1+dN}
\]
matrices.  No ideal in this argument is required to be
principal except the explicitly displayed ideal
$\alpha\mathcal O_F$.
\end{proof}

\subsection{Counting vertices with fixed minimal columns}

For an integral matrix $A$ with nonzero determinant, define
\[
 \begin{split}
 T(A)=\{y\in\Lambda:\;&q_H(y)=1
       \text{ for some }H\in\mathcal R\\
       &\text{for which every column of }A\text{ is minimal}\},
 \end{split}
\]
and let $V(A)$ be the set of vertices for which every column
of $A$ is minimal.  These sets are allowed to be empty.

\begin{lemma}\label{ac-cramer}
Let $G$ be a positive definite Hermitian $N$ by $N$ matrix
such that $G_{aa}\le1$ for every $a$.  If $c^*Gc\le1$, then
\[
                          (\det G)|c_a|\le N^N
                          \quad(1\le a\le N).
                          \tag{AC16}
\]
\end{lemma}
\begin{proof}
Cauchy--Schwarz for the inner product defined by $G$ gives
$|G_{ij}|\le1$ and $|(Gc)_i|\le1$.
Cramer's rule says that $(\det G)c_a$ is the determinant
of the matrix obtained by replacing column $a$ of $G$
by $Gc$.  Every entry of this matrix has absolute value
at most one.  Its determinant is a sum of $N!$ signed
products, so its absolute value is at most
$N!\le N^N$.
\end{proof}

\begin{lemma}\label{ac-short-count}
One has
\[
             |T(A)|\le(2R_0+1)^d,\qquad
             |V(A)|\le\bigl((2R_0+1)^d+1\bigr)^{dN}.
             \tag{AC17}
\]
In particular both sets are finite.
\end{lemma}
\begin{proof}
Suppose $y\in T(A)$, and choose the form $H$ in its definition.
Write $I=I(A)$, $B=B_A$, and $z=By$.
Since $Az=Iy$, and $H'_v=\sigma_v(A)^*H_v\sigma_v(A)$,
\[
              \sum_v \sigma_v(z)^*H'_v\sigma_v(z)=I^2.
\]
Every term is nonnegative.  For each place apply
Lemma~\ref{ac-cramer} to $G=H'_v$ and
$c=I^{-1}\sigma_v(z)$.  Using (AC11) gives
\[
       |\sigma_v(z_a)|
          \le I N^N I_0
          \le N^N I_0^2=R_0
                  \qquad(v,a).
                  \tag{AC18}
\]

For any integer $R\ge0$, there are at most $(2R+1)^n$
algebraic integers $a\in\mathcal O_F$ satisfying
$|\sigma_v(a)|\le R$ for all $v$.
To prove this, reduce modulo $(2R+1)\mathcal O_F$.
If two such integers $a,b$ have the same residue, write
$a-b=(2R+1)c$ with $c\in\mathcal O_F$.  Then
\[
                    |\sigma_v(c)|\le\frac{2R}{2R+1}<1
\]
at every place.  The norm of $c$ is an integer of absolute
value less than one, so $c=0$.  Reduction is injective on
the set being counted, and the quotient has
$(2R+1)^n$ elements.  Applying this with $R=R_0$ in each
of the $N$ coordinates bounds the possible $z$ by
$(2R_0+1)^d$.  The map $y\mapsto B_Ay$ is injective
because $B_A$ is invertible over $F$.  This proves the
first assertion in (AC17).

If $H\in V(A)$, its active functionals span
$\mathcal H^*$ by definition.  Choose at most
$\dim_{\mathbb R}\mathcal H\le dN$ of its minimal vectors
whose functionals still span this dual space.
All of them belong to $T(A)$, and the equations
$q_K(y)=1$ for this chosen set uniquely determine $K=H$.
Choose one such determining subset for each vertex.
Two vertices assigned the same subset must be equal.
A set of $t$ elements has at most $(t+1)^r$ subsets of
size at most $r$: order the elements, list each subset
in order, and pad its list to length $r$ with an extra
symbol.  Taking $t=|T(A)|$ and $r=dN$ proves the
second assertion.
\end{proof}

\begin{theorem}\label{ac-orbit}
For every $p\ge0$, the ordered tuple set $X_p$ has finitely
many $\Gamma$-orbits, and
\[
 \begin{split}
 \#(\Gamma\backslash X_p)
   &\le I_0^{\,1+dN}
       \bigl((2R_0+1)^d+1\bigr)^{dN(p+1)}\\
   &\le (2Dd)^{12(p+1)d^3N}.
 \end{split}
 \tag{AC19}
\]
\end{theorem}
\begin{proof}
For a tuple in $X_p$, choose the matrix $A$ of common
minimal columns furnished by Lemma~\ref{ac-common-minima}.
By Lemma~\ref{ac-det}, $I(A)\le I_0$.
Choose $g\in\Gamma$ such that $gA$ belongs to the set
$\mathscr A_{I_0}$ in Lemma~\ref{ac-scaled-inverse}.
Because $q_{gH_j}(gy)=q_{H_j}(y)$, every vertex of the
translated tuple belongs to $V(gA)$.
Every orbit thus meets
\[
              \bigcup_{A'\in\mathscr A_{I_0}}V(A')^{p+1}.
\]
Only the tuples satisfying (AC6) need be included, so
counting all displayed tuples is an upper bound.
Lemma~\ref{ac-short-count} proves the first inequality.

Here are the numerical estimates for the second inequality.
Put $b=2Dd\ge2$.  Since $1\le N\le d$,
\[
 \begin{aligned}
 I_0=D^N d^d&\le b^{2d},&
 N^N&\le b^d,&
 R_0&\le b^{5d},\\
 2R_0+1&\le b^{5d+2},&
 (2R_0+1)^d+1&\le b^{(5d+2)d+1}.
 \end{aligned}
\]
For the fourth inequality use $R_0\ge1$ and
$2R_0+1\le3R_0\le b^2R_0$; for the fifth use $2\le b$.
The exponent resulting from the first line of (AC19)
is therefore at most
\[
       2d+2d^2N+(5d^3N+2d^2N+dN)(p+1).
\]
Each of its five terms is bounded respectively by
$2$, $2$, $5$, $2$, and $1$ times $(p+1)d^3N$.
Their sum is at most $12(p+1)d^3N$.
\end{proof}

\subsection{Stabilizers, hypotheses, and the resulting homology estimate}

\begin{proposition}\label{ac-stabilizer}
Every stabilizer of a tuple in $X_p$ is finite and has
order at most $3^{nN^2}=3^{dN}$.
\end{proposition}
\begin{proof}
All vertices are positive definite by
Proposition~\ref{ac-polyhedron}.  At the complex places,
the stabilizer of one such tuple is contained in the
product of the unitary stabilizers of its first entry.
This is a compact subset of the product of complex
matrix spaces.  Matrices over $\mathcal O_F$ form a
closed discrete subset of that real vector space,
so their intersection with this compact set is finite.

For the quantitative assertion apply
Lemma~\ref{ht-congruence-three} to
$M_N(\mathcal O_F)$, whose additive rank is $nN^2$.
That lemma proves that reduction modulo $3$ is injective
on every finite subgroup of its unit group.  The residue
ring has $3^{nN^2}$ elements.
\end{proof}

Theorem~\ref{ac-exact}, Theorem~\ref{ac-orbit}, and
Proposition~\ref{ac-stabilizer} supply all the hypotheses
of Proposition~\ref{ht-permutation} for this arithmetic
group.  In degree $m\ge1$ one may take
\[
       A=(2Dd)^{12(m+1)d^3N},\qquad
       s=3^{dN},\qquad B=m+2.
\]
In particular the geometric input of
Proposition~\ref{ht-arithmetic-numerical} is now proved.
For $E\ge1$, $1\le N\le4E+6$, and $0\le m\le2E+1$,
its established numerical argument gives the unconditional
inequality
\[
 \log\#H_m(\operatorname{GL}_N(\mathcal O_F),\mathbb Z)_{\rm tors}
       \le\exp\bigl((2En)^{13}\log(2D)\bigr).
 \tag{AC20}
\]
No class-number estimate is imported in the proof of
(AC20).  The remaining passage from this homology estimate
to the strict Selmer group in the submitted proof is
a different assertion.

\paragraph{Boundary cases and normalization checks.}
The source also allows $N=0$ in some intermediate statements.
In that case $\Lambda=0$, $\mathcal H=0$, $\mathcal R$
has one point, and $\Gamma$ is trivial.  The ordered
complex has a single generator in every degree and
is contracted by prepending its unique vertex.
Every orbit and stabilizer has cardinality one.
With $0^0=1$, the unsimplified first bound in (AC19)
also equals one.  The simplified bound involving $d$
is only asserted above for $N\ge1$.

The normalization in (AC12) can be checked already at
$F=\mathbb Q(i)$, $N=1$.  For $H=(t)$ and the integral
basis $1,i$, one has $G_H=tI_2$, $D=4$, and
$2^2\det G_H=4t^2=D(\det H)^2$.
Removing the factor $2^{-d}$ would give a false identity.
This check concerns the convention of using one embedding
per conjugate pair; the source includes the correct factor.

Positive definiteness is also essential to the stabilizer
argument.  On $\mathcal O_F^2$, the semidefinite form
$|x_1|^2$ at every complex place is preserved by all
matrices
\[
                      \begin{pmatrix}1&0\\a&1\end{pmatrix},
                         \qquad a\in\mathcal O_F.
\]
Its arithmetic stabilizer is infinite.  It is not a
counterexample to the proposition: the integral vector
$(0,1)$ has value zero, so this form is excluded from
$\mathcal R$ by (AC1).  Proposition~\ref{ac-polyhedron}
proves exactly the positivity needed to exclude this
possibility.

Finally, the construction uses the ordered tuples (AC6),
with repetitions and the unsigned action on their
generators.  Replacing them by oriented simplices without
accounting for orientation characters would change the
permutation-module argument.  The proof above makes no
such replacement.  No counterexample satisfying the actual
hypotheses was found.




\section{The strict Selmer comparison and its numerical bound}
\label{sc-section}

The arithmetic complex in Section~\ref{ac-section} proves the
homology torsion estimate used by the submitted class-counting
argument.  It does not by itself bound a strict Selmer group.
We now state the comparison assertions precisely,
prove one elementary improvement in their numerical factor,
and give the resulting conditional exponent.
This section isolates the comparisons before proving the numerical
consequence. Sections~\ref{sl-section}--\ref{sp-section} then prove
the class-number estimate and assertions~\ref{sc-count}--\ref{sc-layer}.
Theorem~\ref{sw-few} later proves the finite-field homology
comparison in assertion~\ref{sc-few}, with an improved scalar exponent.

\subsection{The precise arithmetic objects}

Let $\ell$ be prime and let $E\ge1$.  For a number field $L$
put $n_L=[L:\mathbb Q]$, $D_L=|\operatorname{disc}L|$,
$r_2(L)=\#\{\text{complex places of }L\}$, and
\[
 t_L(N,m)=\#H_m(\operatorname{GL}_N(\mathcal O_L),
                       \mathbb Z)_{\rm tors},
 \qquad
 g_L=2n_L+\lfloor\log_2 h_L\rfloor+2,
 \tag{SC1}
\]
where $h_L$ is the ideal class number.
For the totally imaginary fields occurring below, the relevant
homology torsion groups are finite by Section~\ref{ac-section}.

Let $\mathcal S_\ell(L,E)$ denote the strict Selmer group with
coefficients $\mathbb Q_\ell/\mathbb Z_\ell(-E)$: it consists
of the classes in $H^1(G_L,\mathbb Q_\ell/\mathbb Z_\ell(-E))$
which are unramified outside $\ell$ and have zero restriction
to every decomposition group above $\ell$.
Equivalently it is the kernel of localization above $\ell$
in $H^1(G_{L,S},\mathbb Q_\ell/\mathbb Z_\ell(-E))$,
where $S$ is the set of primes above $\ell$.
Write
\[
       s_\ell(L,E)=
       \length_{\mathbb Z_\ell}\mathcal S_\ell(L,E)
             \in\mathbb N\cup\{\infty\}.
       \tag{SC2}
\]
When this length is finite, the group has order
$\ell^{s_\ell(L,E)}$.
These are the negative twists used by the source.

Consider a finite cyclic Galois extension $K/L$ with generator
$\gamma$, and integers $\nu\ge c\ge1$, with $c\ge2$ if $\ell=2$.
Suppose that $K$ contains a primitive $\ell^\nu$-th root
$\zeta$ satisfying
\[
          \gamma(\zeta)=\zeta^{\,1+\ell^c},
          \qquad [K:L]=\ell^{\nu-c}.
          \tag{SC3}
\]
Put $q=\ell^\nu$, $\kappa=1+\ell^c$, and $a=\kappa^E$.
For a fractional divisor $\mathfrak d=\sum_P d_P P$
on $K$, call $\alpha\in K^\times$ admissible for
$\mathfrak d$ if
\[
                v_P(\alpha)=d_{\gamma P}-a d_P
                       \quad\text{for every finite prime }P.
                \tag{SC4}
\]
Define the explicitly weighted product
\[
       N_{\gamma,a}(\alpha)
          =\prod_{j=0}^{[K:L]-1}\gamma^j(\alpha)^{a^j}.
          \tag{SC5}
\]
This is not the ordinary field norm.

Let $T_{\gamma,a,q}$ be the set of functions
$\varphi:\{P\text{ finite in }K\}\to\mathbb Z/q\mathbb Z$
such that
\[
 \sum_P v_P(\alpha)\varphi(P)=0\quad(\alpha\in K^\times),
 \qquad
 \varphi(\gamma^{-1}P)=a\varphi(P)\quad(P).
 \tag{SC6}
\]
The sum is finite because a principal divisor has finite
support.  The first condition identifies these functions
with homomorphisms out of the ideal class group; the second
is the displayed equivariance condition.  In particular
the inverse on $\gamma$ is part of the statement.

For $b,B\in\mathbb N$, write $\mathcal A(b,B)$ for the
following assertion about these fixed data.  Every finite
set $S\subset K^\times$ has cardinality at most $B$
provided that each $\alpha\in S$ is admissible for some
fractional divisor and that, for all $\alpha,\alpha'\in S$,
\[
 \begin{split}
 N_{\gamma,a}(\alpha)^{\ell^b}
     &=N_{\gamma,a}(\alpha')^{\ell^b}\delta^q
       \text{ for some }\delta\in K^\times\\
 &\hspace{35mm}\Longrightarrow\quad\alpha=\alpha'.
 \end{split}
 \tag{SC7}
\]
Thus this assertion bounds a set of distinct classes
represented by the weighted products, after the specified
power is taken.  It does not assert a bound on the infinite
set of all admissible elements.

\subsection{The four comparison assertions}

The following are the exact mathematical assertions used in this
route of the submitted proof\cite{alpo-uniform-source}.

\begin{enumerate}
\item\label{sc-few}
For totally imaginary $L$ and data (SC3), assume
\[
                 2E+1<\ell^c,\qquad 2E+1<r_2(L).
                 \tag{SC8}
\]
Then $\mathcal A(b,B)$ holds with
\[
 \begin{split}
 b&=2E+(2E+1)(c+2E+1),\\
 B&=q^{r_2(L)}t_L(4E+6,2E+1)t_L(4E+6,2E).
 \end{split}
 \tag{SC9}
\]
In the source this assertion is obtained from constructed
finite-field homology functionals, integral homology
classes with prescribed reductions, an odd-degree rational
rank estimate, and detection of global powers by residue
fields.  Those constructions and comparisons are not
replaced here by the homology torsion bound.

\item\label{sc-count}
For totally imaginary $L$, any data (SC3), and any $b,B$
for which $\mathcal A(b,B)$ holds, one has
\[
 |T_{\gamma,a,q}|\,q^{r_2(L)}
       \le B\,\ell^{(b+v_\ell(E))g_L+c+v_\ell(E)}.
 \tag{SC10}
\]
This is the separate comparison with class groups and
units.  It does not require the two inequalities (SC8).

\item\label{sc-sha}
Suppose $L$ is totally imaginary and $c\ge1$, with $c\ge2$
at $\ell=2$, is the largest integer such that
$\mu_{\ell^c}\subset L$.  If there is a single integer $B$
such that
\[
                         |T_{\gamma,\kappa^E,\ell^\nu}|
                              \le B
                              \tag{SC11}
\]
for every $\nu$ and every datum (SC3) over $L$, then
$s_\ell(L,E)$ is finite and
\[
                              \ell^{s_\ell(L,E)}\le B.
                              \tag{SC12}
\]
Uniformity in the entire cyclotomic tower is essential.

\item\label{sc-layer}
Suppose $L$ contains $\mu_\ell$ if $\ell$ is odd and
$\mu_4$ if $\ell=2$.  There is a totally imaginary
cyclotomic extension $L'/L$ of degree $d$, and a largest
integer $c$ with $\mu_{\ell^c}\subset L'$, such that
\[
 \begin{gathered}
 c\ge1,\quad \ell=2\Longrightarrow c\ge2,\qquad
 2E+1<\ell^c,\quad 2E+1<r_2(L'),\\
 n_{L'}=dn_L,\quad
 d\le(n_L+1)(4E+4),\quad
 \ell^c\le2n_{L'},\\
 D_{L'}\le D_L^d(\ell^c)^{n_{L'}},\qquad
 s_\ell(L,E)\le s_\ell(L',E).
 \end{gathered}
 \tag{SC13}
\]
The numerical statement alone does not prove the last
inequality; that is the Selmer comparison in this assertion.
\end{enumerate}

For the numerical conclusion below we also isolate the
classical class-number bound used by the source:
\[
 \log h_L\le\tfrac12\log D_L+
             n_L\log(1+\tfrac12\log D_L).
 \tag{SC14}
\]
Its source proof invokes the existence of a small ideal
in every ideal class and a count of ideals of bounded norm.
Those inputs are not proved in this section.
We will only need the weaker consequence
\[
                        g_L\le2n_L+4n_L\log(2D_L)+2.
                        \tag{SC15}
\]
Indeed $\log(1+x)\le x$ gives
$\log h_L\le n_L\log D_L$, since $n_L\ge1$.
The inequality $\log2>1/2$ then gives
$\lfloor\log_2h_L\rfloor\le2n_L\log D_L
\le4n_L\log(2D_L)$.

\subsection{An exact valuation in the class construction}

The large value of $b$ in (SC9) comes partly from bounding
a product of differences of powers.  That product admits
a simple exact evaluation of its valuation.

\begin{lemma}\label{sc-vandermonde}
Let $\ell$ be prime, $c\ge1$, and $c\ge2$ if $\ell=2$.
Put $\kappa=1+\ell^c$.  For $E\ge1$ define
\[
             P_E=\prod_{\substack{0\le j\le2E\\j\ne E}}
                       (\kappa^E-\kappa^j)\in\mathbb Z.
\]
Then $P_E\ne0$ and
\[
                   v_\ell(P_E)=2Ec+2v_\ell(E!).
                   \tag{SC16}
\]
Consequently, in any $\mathbb Z/\ell^\nu\mathbb Z$-module,
\[
           \ell^{2E}P_E z=0
             \quad\Longrightarrow\quad
           \ell^{\,2E(c+1)+2v_\ell(E!)}z=0.
           \tag{SC17}
\]
\end{lemma}
\begin{proof}
We first establish the valuation identity
\[
                   v_\ell(\kappa^m-1)=c+v_\ell(m)
                              \quad(m\ge1).
                   \tag{SC18}
\]
If $t$ has valuation $h\ge1$ (and $h\ge2$ for $\ell=2$),
and $u$ is a positive integer prime to $\ell$, the binomial
expansion of $(1+t)^u-1$ has first term $ut$ of valuation
$h$ and all later terms of valuation at least $2h>h$.
If instead the exponent is $\ell$, its first term
$\ell t$ has valuation $h+1$.  All terms with
$2\le j<\ell$ have valuation at least $1+2h>h+1$,
because $\ell\mid\binom{\ell}{j}$.
The last term $t^\ell$ also has strictly greater valuation:
for odd $\ell$, $\ell h\ge h+2$; for $\ell=2$ this uses
$h\ge2$.  Thus raising $1+t$ to the $\ell$-th power
increases the valuation of its difference from one by
exactly one.  Write $m=\ell^r u$ with $\ell\nmid u$,
first apply the prime-to-$\ell$ calculation, and then
apply the second calculation $r$ times.  This proves
(SC18).

The powers of $\kappa>1$ are distinct, so $P_E\ne0$.
For $j\ne E$ the unit factor
$\kappa^{\min(j,E)}$ and (SC18) give
\[
             v_\ell(\kappa^E-\kappa^j)
                          =c+v_\ell(|E-j|).
\]
The distances $|E-j|$ run through $1,\ldots,E$ twice.
Adding the valuations proves (SC16).
Write $P_E=\ell^\delta u$ with
$\delta=2Ec+2v_\ell(E!)$ and $\ell\nmid u$.
Multiplication by $u$ is invertible modulo $\ell^\nu$.
Cancelling it proves (SC17), including the case
$\delta\ge\nu$.
\end{proof}

The source uses precisely $P_E$ at this stage, with the
index in its definition shifted by $e=E-1$.
Its much larger valuation allowance is
$(2E+1)(c+2E+1)$.  Formula (SC17) supplies the smaller
power
\[
                         b_{\rm new}=2E(c+1)+2v_\ell(E!)
                         \tag{SC19}
\]
in the same finite-field detection calculation.
This proves the required cancellation of the scalar.
It does not prove the other constructions entering
assertion~\ref{sc-few}, nor does assertion~\ref{sc-few}
with its original exponent alone imply its stronger
version with $b_{\rm new}$.  Theorem~\ref{sw-few} carries
out those constructions with this scalar estimate and proves
the strengthened version directly.
For example, at $\ell=2$, $c=2$, $E=1$, the original
exponent is $17$, whereas (SC19) is $6$.

The restriction at $2$ cannot be dropped.
For $\ell=2$, $c=1$, and $m=2$ one has
$v_2(3^2-1)=3$, while $c+v_2(m)=2$.
The source includes $c\ge2$ in this case, so this is
a counterexample to the weakened statement only.

\subsection{The conditional exponent}

\begin{proposition}\label{sc-exponent}
Assume assertions~\ref{sc-few}--\ref{sc-layer} and the
bound (SC15).  Suppose that, for some real $C$ and every
totally imaginary number field $F$,
\[
 \log t_F(4E+6,m)
      \le\exp\bigl((2E n_F)^C\log(2D_F)\bigr)
       \qquad(0\le m\le2E+1).
 \tag{SC20}
\]
If $L$ contains the roots of unity required in
assertion~\ref{sc-layer}, then $s_\ell(L,E)$ is finite and
\[
 s_\ell(L,E)\log\ell
   \le\exp\bigl((2E n_L)^{\,5(\max(C,9)+2)+11}
                          \log(2D_L)\bigr).
 \tag{SC21}
\]
In particular the proved homology exponent $C=13$ gives
the conditional strict Selmer exponent $86$.
\end{proposition}
\begin{proof}
Choose $L',c,d$ from (SC13), and write
$n'=n_{L'}$, $D'=D_{L'}$, $r'=r_2(L')$, and
\[
 T_1=t_{L'}(4E+6,2E+1),\quad
 T_2=t_{L'}(4E+6,2E),\quad
 b=2E+(2E+1)(c+2E+1),\quad
 v=v_\ell(E).
\]
Apply assertions~\ref{sc-few} and \ref{sc-count} to
each datum (SC3) over $L'$.  The positive factor $q^{r'}$
cancels, leaving the bound, uniform in $\nu$,
\[
                |T_{\gamma,\kappa^E,\ell^\nu}|
                         \le T_1T_2\ell^X,
              \qquad X=(b+v)g_{L'}+c+v.
 \tag{SC22}
\]
Assertion~\ref{sc-sha} gives a finite length $s'=s_\ell(L',E)$
and $\ell^{s'}\le T_1T_2\ell^X$.  Hence
\[
           s'\log\ell\le\log T_1+\log T_2+X\log\ell.
           \tag{SC23}
\]

Put $u=2En'\ge2$ and $Y=\log(2D')>1/2$.
The inequalities in (SC13) give
\[
 \log\ell\le u,\quad c\le u,\quad v\le u,\quad
 b\le u+2u^2,\quad
 g_{L'}\le u+2uY+2.
\]
For the bound on $b$, use $2E\le u$,
$2E+1\le n'$, $c\le2n'$, and $2n'\le u$.
It follows that
$b+v\le3u^2$, $g_{L'}\le6uY$, and $c+v\le4uY$.
Consequently
\[
                  X\log\ell\le19u^4Y\le u^9Y,
                  \tag{SC24}
\]
since $u^5\ge32>19$.
By (SC20), (SC23) is at most
$2\exp(u^CY)+u^9Y$.

Let $A=u^{\max(C,9)}Y\ge1$.  Then
\[
 2\exp(u^CY)+u^9Y
       \le2e^A+A\le3e^A\le e^{4A}
       \le\exp\bigl(u^{\max(C,9)+2}Y\bigr).
 \tag{SC25}
\]
Here $A\le e^A$, $3\le e^{3A}$, and $u^2\ge4$
justify the three last comparisons.

It remains to express the bound in the original field.
Put $x=2E n_L\ge2$, $Y_0=\log(2D_L)>1/2$,
$w=\ell^c$, and $C_1=\max(C,9)+2\ge0$.
The extension bounds imply
\[
 d\le16E n_L=8x,\quad
 n'\le16E n_L^2\le4x^2,\quad
 u=2En'\le8x^2\le x^5,\quad w\le8x^2.
\]
Using the discriminant inequality in (SC13) yields
\[
 \begin{split}
 Y&\le\log2+d\log D_L+n'\log w\\
  &\le(1+8x)Y_0+32x^4
   \le73x^4Y_0
   \le x^{11}Y_0,
 \end{split}
\]
where $\log w\le w$, $Y_0>1/2$, and $x^7\ge128>73$
were used.  Thus $u^{C_1}Y\le x^{5C_1+11}Y_0$.
Combining this with (SC25) and
$s_\ell(L,E)\le s'$ proves (SC21).
The value for $C=13$ is $5(15)+11=86$.
\end{proof}

The original source inserts $C=24$ and obtains $141$
by this same exponent propagation.  The improvement to
$86$ follows once the comparison assertions are supplied.
The following sections prove the class-number input and the
other comparisons, and Theorem~\ref{sw-few} proves
assertion~\ref{sc-few}; Corollary~\ref{sw-coarse-consequence}
records the resulting bound.
It is not a proof of their conclusion from the geometric
complex alone.



\section{Class numbers and the cyclotomic enlargement}
\label{sl-section}

We prove two of the inputs isolated in Section~\ref{sc-section}:
the class-number estimate and the cyclotomic enlargement
assertion~\ref{sc-layer}.  The latter includes the inequality for
strict Selmer lengths.  It is a statement about this particular
cyclotomic tower and nonzero twist; it is not an assertion about
restriction in an arbitrary extension of prime-power degree.
The finite-field homology construction and the comparison with
units in assertions~\ref{sc-few} and \ref{sc-count} are separate.

\subsection{An elementary bound for the class number}

For a number field $F$, write $n=[F:\mathbb Q]=r_1+2r_2$,
$D=|\operatorname{disc}F|$, and let $N\mathfrak a$ be the
absolute norm of a nonzero integral ideal.  Use ordinary Lebesgue
measure on $\mathbb R^{r_1}\times\mathbb C^{r_2}$, with one
embedding from each conjugate pair.  The image of $\mathfrak a$
has covolume
\[
                 2^{-r_2}\sqrt D\,N\mathfrak a .
                 \tag{SL1}
\]
Indeed, for $\mathcal O_F$ the matrix of all embeddings has
squared determinant equal to the discriminant.  Replacing each
pair of complex rows by its real and imaginary parts changes
the absolute determinant by $1/2$.  Passing to an ideal
multiplies the covolume by its index $N\mathfrak a$.

\begin{lemma}\label{sl-small-ideal}
Every ideal class of $F$ contains a nonzero integral ideal of norm
at most
\[
          \left(\frac4\pi\right)^{r_2}
                 \frac{n!}{n^n}\sqrt D
             \ \le\ \sqrt D .
             \tag{SL2}
\]
\end{lemma}
\begin{proof}
The elementary lattice-volume argument in the proof of
Lemma~\ref{ac-positive-gram} shows that an open, convex,
centrally symmetric body of volume greater than $2^n$ times a
lattice's covolume contains a nonzero lattice point.
Apply it to the image of an integral ideal $\mathfrak a$ and
the body
\[
 B_T=\left\{(x,z):
       \sum_{i=1}^{r_1}|x_i|+
                  2\sum_{j=1}^{r_2}|z_j|<T\right\}.
\]
Integration in polar coordinates gives
\[
       \operatorname{vol}(B_T)
             =2^{r_1}\left(\frac\pi2\right)^{r_2}
                         \frac{T^n}{n!}.
                         \tag{SL3}
\]
For clarity, put $u_i=|x_i|$ and $v_j=2|z_j|$.
The measures become $2\,du_i$ and $(\pi/2)v_j\,dv_j$.
The integral of $\prod_jv_j$ over
$u_i,v_j\ge0$, $\sum u_i+\sum v_j<T$, is $T^n/n!$;
successive integration, using
$\int_0^T t^a(T-t)^b\,dt=a!b!T^{a+b+1}/(a+b+1)!$,
proves this formula.

Thus, whenever
$T^n>(4/\pi)^{r_2}n!\sqrt D\,N\mathfrak a$, there is
$0\ne\alpha\in\mathfrak a$ in $B_T$.  The arithmetic--geometric
mean inequality, applied to the $n$ nonnegative numbers
consisting of each $|x_i|$ and two copies of each $|z_j|$, gives
\[
        |N_{F/\mathbb Q}\alpha|
              \le(T/n)^n .
\]
Let $T$ decrease to the indicated threshold.  A bounded region
contains only finitely many lattice points, so one nonzero
$\alpha$ attains the resulting non-strict bound.
The ideal $(\alpha)\mathfrak a^{-1}$ is integral and has norm
$|N_{F/\mathbb Q}\alpha|/N\mathfrak a$.  Taking $\mathfrak a$
in the inverse of the desired class proves the first inequality.

For the second, induction gives $n!/n^n\le2^{1-n}$:
the ratio of the left-hand quantities at $n+1$ and $n$ is
$(n/(n+1))^n\le1/2$, by the binomial inequality
$(1+1/n)^n\ge2$.  Since $4/\pi<2$ and
$r_2\le n-1$ for $n\ge2$, the desired product is at most
$2^{1-n+r_2}\le1$.  For $n=1$ it is exactly one.
\end{proof}

\begin{lemma}\label{sl-ideal-count}
For every real $X\ge1$ the number of nonzero integral ideals
of norm at most $X$ is at most
\[
                    X(1+\log X)^{n-1}.
                    \tag{SL4}
\]
\end{lemma}
\begin{proof}
There are at most $n$ prime ideals above any rational prime $p$,
because their ramification indices and residue degrees satisfy
$\sum_{\mathfrak p\mid p}e_{\mathfrak p}f_{\mathfrak p}=n$.
For each $p$, assign distinct labels in $\{1,\ldots,n\}$ to
these prime ideals.  If
$\mathfrak a=\prod_{\mathfrak p}\mathfrak p^{a_{\mathfrak p}}$,
define integers $b_1,\ldots,b_n\ge1$ by placing the factor
$p^{a_{\mathfrak p}}$ in the coordinate with the label of
$\mathfrak p$.  Unique factorization of integers and ideals
shows that $\mathfrak a\mapsto(b_1,\ldots,b_n)$ is injective.
Moreover
\[
        \prod_i b_i
          =\prod_{\mathfrak p}p^{a_{\mathfrak p}}
          \le \prod_{\mathfrak p}
                         p^{f_{\mathfrak p}a_{\mathfrak p}}
          =N\mathfrak a.
\]
No assumption that the prime ideals are principal is involved.

Let $A_j(X)$ count $j$-tuples of positive integers with product
at most $X$.  The bound
$A_j(X)\le X(1+\log X)^{j-1}$ follows by induction.
For $j=1$, it is $\lfloor X\rfloor\le X$.  For $j\ge2$,
\[
 \begin{split}
 A_j(X)&=\sum_{a\le X}A_{j-1}(X/a)\\
       &\le X(1+\log X)^{j-2}\sum_{a\le X}\frac1a
        \le X(1+\log X)^{j-1}.
 \end{split}
\]
The last inequality follows by comparing the terms from $a=2$
onwards with the integral of $1/t$ over $[1,X]$.
The injection proves (SL4).
\end{proof}

\begin{corollary}\label{sl-class-number}
For every number field,
\[
       \log h_F\le\frac12\log D+
             (n-1)\log\left(1+\frac12\log D\right).
             \tag{SL5}
\]
In particular (SC14) and (SC15) hold.
\end{corollary}
\begin{proof}
Choose one ideal of norm at most $\sqrt D$ in each class by
Lemma~\ref{sl-small-ideal}, apply Lemma~\ref{sl-ideal-count}
with $X=\sqrt D$, and take logarithms.
The discriminant is a nonzero integer, so $D\ge1$ and all
quantities are defined.  Replacing $n-1$ by $n$ gives (SC14);
the calculation following that formula gives (SC15).
\end{proof}

\subsection{The cyclotomic degree and discriminant}

Let $p$ be prime and let $L$ contain $\mu_p$ if $p$ is odd,
and $\mu_4$ if $p=2$.  Write $c_0$ for the largest integer
such that $\mu_{p^{c_0}}\subset L$.  Then $c_0\ge1$ and
$c_0\ge2$ at $p=2$.  The cyclotomic character identifies
\[
 \operatorname{Gal}(L(\mu_{p^\infty})/L)
                         =1+p^{c_0}\mathbb Z_p,
 \qquad
 L_m=L(\mu_{p^{c_0+m}}),\quad [L_m:L]=p^m .
 \tag{SL6}
\]
Here the equality denotes the image under that character.
We recall the justification to fix both the degree and the
largest-root convention.

The polynomial $\Phi_{p^a}(1+X)$ is Eisenstein at $p$:
its constant term is $p$, its leading coefficient is one, and
its reduction modulo $p$ is $X^{p^{a-1}(p-1)}$.
For the last assertion use
$\Phi_{p^a}(Y)=\sum_{j=0}^{p-1}Y^{jp^{a-1}}$ and
the binomial identity in characteristic $p$.
Eisenstein's criterion follows directly by reducing a proposed
factorization into monic integral polynomials modulo $p$:
both constant terms would be divisible by $p$, contradicting
the constant term $p$.  The usual clearing of denominators
for primitive polynomials passes from rational to integral
factorizations.  Therefore
$[\mathbb Q(\mu_{p^a}):\mathbb Q]=p^{a-1}(p-1)$.
The automorphisms $\zeta\mapsto\zeta^b$ consequently fill the
whole unit group modulo $p^a$.

For $a\ge1$ (and $a\ge2$ at $2$), the group
$1+p^a\mathbb Z_p$ is topologically generated by $1+p^a$.
Indeed the valuation formula (SC18) shows that this element
has order $p^{b-a}$ modulo $p^b$, equal to the number of
elements in the indicated congruence group.  Every nontrivial
closed subgroup is $1+p^b\mathbb Z_p$ for some $b\ge a$:
choose a nonidentity element with least valuation of its
difference from one; the same valuation argument shows that
its powers fill that congruence group, and minimality puts
every element of the subgroup in it.
The image of the cyclotomic character for $L$ is a closed
subgroup of finite index in the cyclotomic Galois group.
It lies in $1+p\mathbb Z_p$, or in $1+4\mathbb Z_2$.
The largest-root definition forces its index parameter to
be $c_0$, proving (SL6).  It follows as well that the largest
$p$-power roots in $L_m$ have order exactly $p^{c_0+m}$.

\begin{lemma}\label{sl-cyclotomic-discriminant}
If $L'=L(\zeta)$, where $\zeta$ is a root of unity of order
$N\ge1$, and $d=[L':L]$, then
\[
                  D_{L'}\le D_L^d N^{[L':\mathbb Q]} .
                  \tag{SL7}
\]
\end{lemma}
\begin{proof}
Put $R=\mathcal O_L$, $B=\mathcal O_{L'}$ and
$A=R[\zeta]\subset B$.  The minimal polynomial $f$ of
$\zeta$ is monic in $R[X]$ and divides $X^N-1$ there.
The trace dual $A^\#=\{x:\operatorname{Tr}_{L'/L}(xA)
\subset R\}$ equals $f'(\zeta)^{-1}A$.
To see this, use the basis $1,\zeta,\ldots,\zeta^{d-1}$.
Lagrange interpolation on the distinct roots of $f$ gives
\[
 \operatorname{Tr}\bigl(\zeta^j/f'(\zeta)\bigr)
     =0\quad(0\le j<d-1),\qquad
 \operatorname{Tr}\bigl(\zeta^{d-1}/f'(\zeta)\bigr)=1.
\]
For larger $j$ these traces lie in $R$, by reducing $X^j$
modulo $f$.  Thus the pairing matrix for the bases
$\zeta^i$ and $\zeta^j/f'(\zeta)$ has zeros below the
opposite diagonal and ones on that diagonal.  Its determinant
is a unit in $R$.  This proves the trace-dual assertion over
$R$ itself, without assuming that $R$ is principal.

As $B^\#\subset A^\#$, we have
$f'(\zeta)B^\#\subset A\subset B$.
The relative different is
$\mathfrak D_{L'/L}=(B^\#)^{-1}$, so
$f'(\zeta)\in\mathfrak D_{L'/L}$.
Differentiating $X^N-1=f(X)g(X)$ and evaluating at $\zeta$
gives
$N\zeta^{N-1}=f'(\zeta)g(\zeta)$.
Since $\zeta$ is a unit, $N$ belongs to the different.
Hence $(N)\subset\mathfrak D_{L'/L}$ and
$N_{L'/\mathbb Q}\mathfrak D_{L'/L}\le N^{[L':\mathbb Q]}$.
The discriminant identity
$D_{L'}=D_L^d N_{L'/\mathbb Q}\mathfrak D_{L'/L}$
proves (SL7).  This identity follows by the determinant of
the trace pairing in a tower, or equivalently by transitivity
of the trace-dual fractional ideals.
\end{proof}

\subsection{The strict Selmer comparison in a layer}

\begin{proposition}\label{sl-selmer-monotonicity}
For $E\ge1$ and every $m\ge0$,
\[
                    s_p(L,E)\le s_p(L_m,E).
                    \tag{SL8}
\]
\end{proposition}
\begin{proof}
Use the common infinite tower and the compact module $X$
of Proposition~\ref{ct-descent}.  If $u$ is the cyclotomic
value of a generator and $T$ its action on $X$, put
$U=u^E T$.  That proposition identifies the two Pontryagin
duals with $X/(U-1)X$ and $X/(U^{p^m}-1)X$.
The polynomial identity
\[
 U^{p^m}-1=(U-1)(1+U+\cdots+U^{p^m-1})
\]
gives a surjection from the latter quotient to the former.
Dualizing gives an injection of the strict Selmer groups,
and the length inequality follows, with infinite length
allowed.  The descent proposition proves preservation of
the full local conditions; no assertion that ramification
conditions are automatically unchanged under an arbitrary
extension is being used.
\end{proof}

\begin{theorem}\label{sl-enlargement}
Assertion~\ref{sc-layer} is valid.
\end{theorem}
\begin{proof}
Let $n=[L:\mathbb Q]$ and choose the least $m\ge0$ such that
$p^m\ge4E+4$.  Put $d=p^m$, $L'=L_m$ and $c=c_0+m$.
Because $4E+4>1$, we have $m\ge1$, and minimality gives
\[
                    4E+4\le d<p(4E+4).
                    \tag{SL9}
\]
The hypotheses on roots of unity imply that $L$ is totally
imaginary and $n\ge2$.  They also imply $p\le n+1$:
for odd $p$ use $p-1=[\mathbb Q(\mu_p):\mathbb Q]\le n$;
for $p=2$ the inequality is immediate.
Consequently $d\le(n+1)(4E+4)$ and
$[L':\mathbb Q]=dn$.
The field $L'$ is totally imaginary, so
$r_2(L')=dn/2\ge d>2E+1$.
Also $p^c=p^{c_0}d\ge2d>2E+1$.

The largest-root assertion follows from (SL6).
Since $\mathbb Q(\mu_{p^c})\subset L'$,
\[
        p^{c-1}(p-1)\le n'=[L':\mathbb Q],
        \qquad p^c\le\frac{p}{p-1}n'\le2n'.
\]
Lemma~\ref{sl-cyclotomic-discriminant}, with $N=p^c$,
gives $D_{L'}\le D_L^d(p^c)^{n'}$.
Finally Proposition~\ref{sl-selmer-monotonicity} gives the
required length inequality.  These are all the assertions
in (SC13).
\end{proof}

Thus the exponent calculation in Proposition~\ref{sc-exponent}
no longer requires an unproved class-number estimate or
cyclotomic enlargement.  Its remaining assumptions are the
finite-field homology and class-group comparisons in
assertions~\ref{sc-few}--\ref{sc-sha}, to the extent not proved
in subsequent sections.


\section{The comparison with units and ideal classes}
\label{su-section}

This section proves assertion~\ref{sc-count}.  We first give
a proof of the classical Dirichlet unit theorem
\cite[Chapter 5]{milne-ant} from the lattice-volume argument
and ideal count already established.
The comparison below, including the determinant of the action,
the invariant Kummer classes, and all cardinality estimates, is
proved here.  No homology-class construction from
assertion~\ref{sc-few} is assumed except its stated finite-set
bound $\mathcal A(b,B)$.

\subsection{The logarithmic unit lattice}

\begin{theorem}\label{su-unit-theorem}
Let $F$ be a number field of degree $n=r_1+2r_2$.
For one embedding at each infinite place put
\[
 \lambda(u)=(\log|\sigma_vu|)_v,\qquad
 H=\{x:\ \sum_v w_vx_v=0\},
 \quad
 w_v=\begin{cases}1&v\text{ real},\\2&v\text{ complex}.\end{cases}
\]
The image $\lambda(\mathcal O_F^\times)$ is a full lattice
in $H$, and its kernel is the finite cyclic group $\mu(F)$.
Consequently $\mathcal O_F^\times/\mu(F)$ is free abelian
of rank $r_1+r_2-1$.
\end{theorem}
\begin{proof}
A unit has norm of absolute value one, so its logarithm
belongs to $H$.  If these logarithms belong to a bounded
set, all conjugates of the units are bounded.
By (SL1), only finitely many such algebraic integers
exist.  Thus the logarithmic image has finite intersection
with every bounded subset of $H$, and is discrete.
Its kernel is a finite group.  Every element of this
finite group is a root of unity, and every root of unity
lies in the kernel.  A finite subgroup of
$\mathbb C^\times$ is cyclic: if its exponent is $m$,
the equation $X^m=1$ has at most $m$ roots, whereas the
classification of finite abelian groups gives at least
$m$ elements, with equality only for a cyclic group.
Equivalently, it is a subgroup of the cyclic group of
complex $m$-th roots of unity.

It remains to prove that the logarithmic image spans $H$.
Choose $C>1$ so large that
\[
       2^{r_1}\pi^{r_2}C^n
                         >2^n2^{-r_2}\sqrt{D_F}.
\]
For each $x\in H$, apply the lattice-volume lemma to the
product of intervals and disks
\[
    |y_v|<Ce^{x_v}\quad(v\text{ real}),\qquad
    |z_v|<Ce^{x_v}\quad(v\text{ complex}).
\]
Its volume is $2^{r_1}\pi^{r_2}C^n$, independent of $x$,
because $\sum_vw_vx_v=0$.  Hence it contains
$0\ne\alpha\in\mathcal O_F$.
Write $b_v=\log|\sigma_v\alpha|-x_v$.
Each $b_v<\log C$, and
\[
          0\le\log|N_{F/\mathbb Q}\alpha|
                    =\sum_vw_vb_v<n\log C.
\]
The lower inequality holds because the nonzero norm is
an integer.  For each $v$ it follows that
\[
          -\frac{n-w_v}{w_v}\log C
                    \le b_v<\log C .
\]
Thus $\lambda(\alpha)-x$ stays in a fixed bounded box.

The principal ideal $(\alpha)$ has norm at most $C^n$.
By Lemma~\ref{sl-ideal-count}, only finitely many such
ideals occur.  Choose a generator $\beta_i$ for each
principal ideal in that finite list.  If
$(\alpha)=(\beta_i)$, then $u=\alpha/\beta_i$ is a unit,
and
$\lambda(u)-x=\lambda(\alpha)-x-\lambda(\beta_i)$
lies in a bounded set independent of $x$.
Consequently every point of $H$ is a bounded distance
from a logarithm of a unit.  A proper real linear
subspace of $H$ cannot have this property, as points
on a line transverse to it have unbounded distance.
The logarithmic image therefore spans $H$.

Choose a real basis of $H$ from this discrete subgroup.
The half-open parallelepiped of its integral span is
bounded, so it contains only finitely many elements of
the subgroup.  Reduction into that parallelepiped shows
that the integral span has finite index in the subgroup.
The subgroup is thus finitely generated, torsion-free,
and of rank $\dim H=r_1+r_2-1$; its parallelepipeds
give a lattice in all of $H$.  This proves the theorem,
including $\dim H=0$.
\end{proof}

\subsection{The divisor sequence and the finite comparison group}

Fix the data (SC3), with $K/L$ cyclic of degree
$d=p^{\nu-c}$, and put
\[
 q=p^\nu,\qquad a=(1+p^c)^E,\qquad
 v=v_p(E),\qquad r=r_2(L).
\]
Thus $a>1$, $a\equiv1\pmod p$, and $L$ is totally imaginary.
The hypotheses also imply $K=L(\mu_{p^\nu})$:
the restriction of $\gamma$ to the primitive root has order
$p^{\nu-c}=d$ by (SC18), so it has trivial kernel in
$\operatorname{Gal}(K/L)$.
We write $Y=K^\times$ additively for the next group calculations:
$y+y'$ means the product of the corresponding nonzero field
elements, and $my$ means their $m$-th power.
Let $I=\bigoplus_P\mathbb Z P$ be the group of fractional
divisors and let $\partial:Y\to I$ be the principal-divisor
map.  Its kernel is the unit group $U_K=\mathcal O_K^\times$,
and its cokernel is the ideal class group $C_K$.
Write $\mu=\mu(K)$ for the torsion subgroup of $Y$.

Let $g$ act on $Y$ by $\gamma^{-1}$ and on $I$ by
$g(P)=\gamma^{-1}P$.  In coordinates,
$(gD)_P=D_{\gamma P}$, so $\partial$ commutes with $g$.
Set
\[
 \delta=g-a,\qquad
 V=\{y\in Y:\partial y\in\delta I\},\qquad
 W=V/(\delta Y+\mu),\qquad
 C=C_K/\delta C_K .
 \tag{SU1}
\]
Every torsion element has zero divisor, and
$\partial(\delta Y)\subset\delta I$, so this quotient is defined.
Multiplicatively, $y\in V$ means exactly that it is admissible
in (SC4).  Since $g^d=1$ and $a>1$, $\delta$ is injective on
the torsion-free group $I$.  Indeed $\delta x=0$ implies
$x=a^dx$, and $a^d-1\ne0$.

\begin{lemma}\label{su-exact}
There is an exact sequence of finite groups
\[
 0\longrightarrow
 (U_K/\mu)/\delta(U_K/\mu)
 \longrightarrow W
 \longrightarrow\ker(\delta:C_K\to C_K)
 \longrightarrow0.
 \tag{SU2}
\]
Moreover
\[
 |W|=\frac{(a^d-1)^r}{a-1}\,|C|,
 \qquad
 v_p\!\left(\frac{(a^d-1)^r}{a-1}\right)
                    =r(\nu+v)-c-v .
 \tag{SU3}
\]
\end{lemma}
\begin{proof}
For $y\in V$ there is a unique $D\in I$ with
$\delta D=\partial y$.  Send the class of $y$ to the ideal
class $[D]$.  Changing $y$ by $\delta z+t$, $t\in\mu$,
changes $D$ by $\partial z$, so the map is well-defined.
Its image is precisely $\ker\delta$ on $C_K$: an ideal
class killed by $\delta$ has a representative $D$ with
$\delta D$ principal.  Its kernel is represented by units,
because $D=\partial z$ implies $y-\delta z\in U_K$.
Finally
\[
             U_K\cap(\delta Y+\mu)=\delta U_K+\mu:
\]
if a unit is $\delta z+t$, taking divisors and using
injectivity of $\delta$ on $I$ gives $\partial z=0$.
This proves the exact sequence.

Theorem~\ref{su-unit-theorem} identifies $(U_K/\mu)\otimes\mathbb R$
with the hyperplane $H$ of coordinate sum zero in
$\mathbb R^{r_2(K)}$.  All its places are complex, so
the weights in the product formula are equal.
The $d$ automorphisms of $K/L$ act freely on the complex
places over each of the $r$ complex places of $L$.
To verify freeness, fix one complex embedding of $L$.
Its $d$ extensions to $K$ form a free Galois orbit;
none is identified with another by conjugation, since
that would change the chosen embedding of $L$.
Thus the permutation of coordinates is a disjoint union
of $r$ cycles of length $d$.

The characteristic polynomial on each cycle is $X^d-1$.
The quotient of $\mathbb R^{rd}$ by $H$ is the trivial
one-dimensional representation.  Therefore
\[
        |\det(\delta|_H)|=\frac{(a^d-1)^r}{a-1}.
\]
An injective endomorphism of a free integral lattice has
cokernel order equal to the absolute determinant of its
matrix: integer row and column operations reduce the
matrix to Smith normal form.  It follows that the unit
quotient is finite with the displayed order.
The group $C_K$ is finite by Lemma~\ref{sl-small-ideal}.
For an endomorphism of a finite group, the kernel and
cokernel have the same order; both orders are the total
order divided by the image order.  This proves the first
formula (SU3) and finiteness of $W$.

Finally (SC18) gives
$v_p(a-1)=c+v$ and $v_p(a^d-1)=\nu+v$.
Taking valuations proves the second formula.
\end{proof}

The equivariant class-group functions in (SC6) identify with
$\operatorname{Hom}(C,\mathbb Z/q\mathbb Z)$.
Indeed vanishing on principal divisors factors a function
through $C_K$, and vanishing on $(g-a)P$ is precisely
$\varphi(\gamma^{-1}P)=a\varphi(P)$.
For a finite abelian group $A$, write $|A|_p=p^{v_p(|A|)}$.
Decomposition into cyclic groups gives
\[
 |\operatorname{Hom}(A,\mathbb Z/p^\nu\mathbb Z)|
        =\prod_i p^{\min(e_i,\nu)}\le |A|_p,
 \tag{SU4}
\]
where $p^{e_i}$ are the orders of its cyclic $p$-primary
factors.  Consequently (SU3) implies
\[
            |T_{\gamma,a,q}|q^r
                   \le |W|_p\,p^{c+v}.
                   \tag{SU5}
\]
In obtaining this inequality we discarded the additional
nonnegative exponent $rv$ in the unit factor.

\subsection{A uniform bound for the \texorpdfstring{$p$}{p}-rank}

Define
\[
 \begin{split}
 \mathcal K_p(L)
   &=\{x\in L^\times/(L^\times)^p:
                 v_{\mathfrak p}(x)\equiv0\pmod p
                       \text{ if }\mathfrak p\nmid p\},\\
 \mathcal I_p(K)
   &=\{[y]\in Y/(pY+\mu):
            gy-y\in pY+\mu,\quad \partial y\in pI\}.
 \end{split}
 \tag{SU6}
\]
These definitions are independent of representatives.
The group $\mathcal K_p(L)$ is a Kummer group with the
displayed valuation conditions, not the strict Selmer group.

\begin{lemma}\label{su-base-kummer}
If $n=[L:\mathbb Q]$ and $h_L$ is the class number, then
\[
           |\mathcal K_p(L)|
                \le p^{\,2n+\lfloor\log_2h_L\rfloor}.
                \tag{SU7}
\]
\end{lemma}
\begin{proof}
First impose the valuation condition also at primes above
$p$, and call the resulting subgroup $\mathcal K_p^0(L)$.
Sending $\alpha$ to the ideal class of
$\frac1p\operatorname{div}(\alpha)$ maps this group to
$C_L[p]$.  Its kernel consists exactly of the image of
$\mathcal O_L^\times/(\mathcal O_L^\times)^p$:
if the divided divisor is principal, dividing $\alpha$
by the corresponding $p$-th power gives a unit.
By Theorem~\ref{su-unit-theorem} and cyclicity of roots of unity this
kernel has order at most $p^{r_1(L)+r_2(L)}\le p^n$.
If $|C_L[p]|=p^t$, then $2^t\le p^t\le h_L$, whence
$t\le\lfloor\log_2h_L\rfloor$.

Reduction of the remaining valuations modulo $p$ injects
$\mathcal K_p(L)/\mathcal K_p^0(L)$ into one copy of
$\mathbb Z/p\mathbb Z$ for each prime above $p$.
There are at most $n$ such primes.
Multiplying the three bounds proves finiteness and (SU7).
\end{proof}

\begin{lemma}[Cyclic Hilbert theorem 90]\label{su-hilbert90}
If $K/L$ is cyclic with generator $\gamma$ and
$w\in K^\times$ has norm one, then $w=\gamma(z)/z$
for some $z\in K^\times$.
\end{lemma}
\begin{proof}
Put $c_0=1$ and
$c_j=(w\gamma(w)\cdots\gamma^{j-1}(w))^{-1}$ for $j\ge1$.
For $b\in K$ define $z_b=\sum_{j=0}^{d-1}c_j\gamma^j(b)$.
The norm condition gives $c_d=1$, and shifting the summation
index gives $\gamma(z_b)=wz_b$.
For some $b$, $z_b\ne0$, by linear independence of distinct
field automorphisms over $K$.
Here is the usual short proof of that independence.
A nonzero relation with the fewest terms, after rescaling,
has coefficient one at its first automorphism.
Choose $x$ on which a second automorphism differs from the
first, evaluate the relation at $xy$, and subtract the
relation at $y$ multiplied by the first automorphism of $x$.
This is a nonzero relation with fewer terms, a contradiction.
A single automorphism cannot have a nonzero scalar multiple
equal to the zero map.  Taking $z=z_b$ proves the lemma.
\end{proof}

\begin{lemma}\label{su-invariant-kummer}
For the cyclotomic data under consideration,
\[
                |\mathcal I_p(K)|
                     \le p^2|\mathcal K_p(L)|.
                     \tag{SU8}
\]
\end{lemma}
\begin{proof}
Choose a generator $\rho$ of the finite cyclic group
$\mu(K)$.  Its finiteness also follows directly from
(SL1): roots of unity are integral and all their conjugates
have absolute value one, so they belong to a compact set
in a lattice.  A finite subgroup of $\mathbb C^\times$
is cyclic, and any complex embedding proves cyclicity here.

Every class in $\mathcal I_p(K)$ has a representative
$\alpha\in K^\times$ with all valuations divisible by $p$
and with
\[
           \gamma(\alpha)=\alpha w^p\rho^j,
              \qquad 0\le j<p,\quad w\in K^\times .
              \tag{SU9}
\]
The invariance condition with $\gamma^{-1}$ is equivalent
to that with $\gamma$.  Reduction of the exponent of
$\rho$ modulo $p$, absorbing its $p$-th power into $w$,
proves (SU9).
Taking ordinary field norms gives
\[
                     (N_{K/L}w)^p=(N_{K/L}\rho)^{-j}.
\]
There are at most $p$ choices for $j$ and at most $p$
values of $N_{K/L}w$ for each $j$, since a polynomial of
degree $p$ over a field has at most $p$ roots.

Compare two representatives $\alpha,\alpha'$ with the
same $j$ and the same norm of $w,w'$.  Put
$\beta=\alpha/\alpha'$ and $v=w/w'$.
Then $\gamma(\beta)/\beta=v^p$ and $N_{K/L}v=1$.
By Lemma~\ref{su-hilbert90}, $v=\gamma(z)/z$.
It follows that $\beta=xz^p$ with $x\in L^\times$.
The extension $K=L(\mu_{p^\nu})$ is unramified outside $p$:
the proof of Lemma~\ref{sl-cyclotomic-discriminant} puts
$p^\nu$ in its relative different, so that the different
is the unit ideal at every other prime.
Thus for $P\mid\mathfrak p\nmid p$,
$v_P(x)=v_{\mathfrak p}(x)$.
The valuations of $\alpha$ and $\alpha'$ and those of
$z^p$ show that this number is divisible by $p$.
Hence $[x]\in\mathcal K_p(L)$.

For each of the at most $p^2$ pairs
$(j,N_{K/L}w)$ that occurs, fix one representative.
Every class with that pair differs from its representative
by the image of a class in $\mathcal K_p(L)$.
This covers $\mathcal I_p(K)$ by at most $p^2$ sets of
size at most $|\mathcal K_p(L)|$, proving (SU8).
No uniqueness of the pair assigned to a class is needed.
\end{proof}

\begin{lemma}\label{su-rank}
With $g_L$ as in (SC1),
\[
                           |W[p]|\le p^{g_L}.
                           \tag{SU10}
\]
\end{lemma}
\begin{proof}
There is an injection $W[p]\to\mathcal I_p(K)$.
To construct it, choose $u\in V$ representing a class
$w\in W[p]$, and write
\[
                         pu=\delta y+t,\qquad t\in\mu .
                         \tag{SU11}
\]
Such $y,t$ exist by the definition of $W$.
Write $\partial u=\delta D$.
Taking divisors in (SU11) and using injectivity of
$\delta$ on $I$ gives $\partial y=pD$.
Since $p\mid a-1$, equation (SU11) also gives
$gy-y\in pY+\mu$.  Thus $[y]\in\mathcal I_p(K)$.

This class is independent of choices.  If $u$ changes by
$\delta z+t_0$, then subtraction of the two equations
(SU11) shows
$\delta(y'-y-pz)\in\mu$.  Whenever $\delta x\in\mu$,
one has $x\in\mu$: apply the polynomial
$g^{d-1}+ag^{d-2}+\cdots+a^{d-1}$ and use
$g^d=1$ to obtain $(1-a^d)x\in\mu$.
The torsion subgroup of an abelian group is saturated,
meaning that a nonzero multiple of an element is torsion
only if the element is torsion.  Therefore
$y'-y-pz\in\mu$, as required.
If $[y]=0$, write $y=pz+t_1$; then (SU11) gives
$p(u-\delta z)\in\mu$, so $u-\delta z\in\mu$ and $w=0$.
The construction respects addition, hence it is injective.
Combine (SU7), (SU8), and
$g_L=2n+\lfloor\log_2h_L\rfloor+2$.
\end{proof}

\subsection{The weighted product and its kernel}

In the additive notation set
\[
       N_a=\sum_{j=0}^{d-1}a^j g^{d-j},\qquad
       s=p^b,\qquad
       R=\frac{a^d-1}{q}.
       \tag{SU12}
\]
The symbol $R$ denotes an integer in this subsection.
Since $g=\gamma^{-1}$ and $\gamma^d=1$, the element
$N_a y$ is exactly the weighted product in (SC5).
Telescoping gives
\[
                    \delta N_a=N_a\delta=(1-a^d)g=-qRg.
                    \tag{SU13}
\]
Also $v_p(R)=v$, by (SC18).

\begin{lemma}\label{su-weighted-map}
The formula
\[
          f:W\longrightarrow Y/(qY+\mu),\qquad
                     [y]\longmapsto[sN_a y]
                     \tag{SU14}
\]
defines a homomorphism.  Its kernel is killed by $sR$.
If $\mathcal A(b,B)$ holds, then $|\operatorname{im}f|\le B$.
\end{lemma}
\begin{proof}
Equation (SU13) shows that $sN_a(\delta Y)\subset qY$.
An endomorphism of an abelian group preserves its torsion
subgroup, so the formula descends to $W$.
If $f([y])=0$, write $sN_a y=qz+t$ with $t\in\mu$.
Apply $\delta$ and use (SU13):
\[
                      q(-sRg y-\delta z)=\delta t.
\]
The right-hand side is torsion.  Saturation of $\mu$
gives $sRg y\in\delta Y+\mu$.
The automorphism $g$ preserves all groups in (SU1);
applying its inverse shows $sRy\in\delta Y+\mu$.
Thus $sR$ kills the kernel.

The group $W$ is finite by Lemma~\ref{su-exact}.
Choose one representative $y\in V$ for each value of $f$.
These give distinct elements $\alpha\in K^\times$, each
admissible for a divisor.  If two weighted products raised
to $p^b$ differ by a $q$-th power, their images in the
even coarser quotient $Y/(qY+\mu)$ coincide.
They were chosen for distinct values, so their
representatives are equal.  This finite set therefore
satisfies exactly (SC7).  Assertion $\mathcal A(b,B)$
bounds its size by $B$.
We have not replaced equality modulo powers by equality
modulo roots of unity in the hypothesis.
\end{proof}

\begin{lemma}\label{su-finite-group}
Let $W$ be a finite abelian group and $f:W\to A$ a
homomorphism.  Suppose $|W[p]|\le p^g$ and the
$p$-primary part of $\ker f$ is killed by $p^k$.
Then
\[
                         |W|_p\le|\operatorname{im}f|p^{kg}.
                         \tag{SU15}
\]
\end{lemma}
\begin{proof}
Multiplication by $p$ maps $W[p^{j+1}]$ into
$W[p^j]$ with kernel $W[p]$.  Induction gives
$|W[p^k]|\le|W[p]|^k\le p^{kg}$.
The $p$-primary part of $\ker f$ is a subgroup of
$W[p^k]$.  Now use $|W|=|\ker f||\operatorname{im}f|$
and take $p$-primary parts, bounding that of the image
by its full order.
\end{proof}

\begin{theorem}\label{su-count}
Assertion~\ref{sc-count} holds.  More explicitly,
\[
 |T_{\gamma,a,q}|q^{r_2(L)}
       \le B\,p^{(b+v_p(E))g_L+c+v_p(E)}
       \quad\text{if }\mathcal A(b,B)\text{ holds}.
 \tag{SU16}
\]
\end{theorem}
\begin{proof}
Write $R=p^vR_0$ with $p\nmid R_0$.
Lemma~\ref{su-weighted-map} shows that the $p$-primary
kernel of $f$ is killed by $p^{b+v}$, because
multiplication by $R_0$ is invertible on a finite
$p$-group.
Lemmas~\ref{su-rank} and \ref{su-finite-group} give
$|W|_p\le Bp^{(b+v)g_L}$.
Substitution in (SU5) proves the theorem.
\end{proof}

The factor $q^{r_2(L)}$ has thus been derived from the
determinant of the unit action, rather than inserted as
a rank heuristic.  The extra two powers of $p$ in $g_L$
come from the root-of-unity exponent and the norm value
in (SU9).  The argument covers the boundary case
$\nu=c$, when $K=L$ and $d=1$, and it retains the
condition $c\ge2$ at $p=2$ needed in (SC18).


\section{From equivariant class-group characters to strict Selmer classes}
\label{sp-section}

We prove assertion~\ref{sc-sha}, the uniform comparison between
class-group characters in the cyclotomic tower and the strict
Selmer group.  The argument gives all the restriction,
continuity, local-condition and twist calculations explicitly.
Its classical input is Hilbert class field reciprocity:
for a totally imaginary number field $K$, every continuous
finite-image abelian character of $G_K$ which vanishes on all
finite inertia groups factors uniquely through
\[
                  \operatorname{Cl}(K)
                    \simeq\operatorname{Gal}(H_K/K),
                  \tag{SP1}
\]
where $H_K$ is the maximal unramified abelian extension.
The Artin map sends a prime class to its arithmetic Frobenius.
We use this theorem in its usual form; we do not give a proof
of global class field theory\cite[Chapter V, \S3,
Theorems 3.5--3.6 and Example 3.9]{milne-cft}.

Fix a totally imaginary field $L$, a prime $\ell$, and $E\ge1$.
Let $c$ be the largest integer such that
$\mu_{\ell^c}\subset L$, and suppose $c\ge1$ and $c\ge2$
if $\ell=2$.
Let $\chi:G_L\to\mathbb Z_\ell^\times$ be the cyclotomic
character, and let
\[
 A=\mathbb Q_\ell/\mathbb Z_\ell,\qquad
 \rho(g)=\chi(g)^{-E},\qquad
 L_\infty=L(\mu_{\ell^\infty}),\qquad
 G_\infty=\operatorname{Gal}(\overline L/L_\infty).
 \tag{SP2}
\]
Thus $A$ with this action is the negative Tate twist
$\mathbb Q_\ell/\mathbb Z_\ell(-E)$.
All cochains and characters below are continuous, and $A$
has the discrete topology.

The cyclotomic calculation in Section~\ref{sl-section} gives
\[
        \chi(G_L)=1+\ell^c\mathbb Z_\ell,\qquad
        G_\infty=\ker\chi.
        \tag{SP3}
\]
In particular, with $\kappa=1+\ell^c$, the quotient
$G_L/G_\infty$ is topologically generated by an element
with cyclotomic value $\kappa$.
For $\nu\ge c$, put $K_\nu=L(\mu_{\ell^\nu})$.
Then $[K_\nu:L]=\ell^{\nu-c}$, and there is a generator
$\gamma_\nu$ with
$\gamma_\nu(\zeta)=\zeta^\kappa$ for a primitive
$\ell^\nu$-th root $\zeta$.  Thus these fields and generators
satisfy exactly (SC3).

\begin{lemma}\label{sp-restriction}
Restriction induces an injective map
\[
 H^1(G_L,A)\longrightarrow
                \operatorname{Hom}_{\rm cont}(G_\infty,A).
 \tag{SP4}
\]
\end{lemma}
\begin{proof}
The action on $G_\infty$ is trivial, so a cocycle restricts
to a homomorphism and every coboundary restricts to zero.
Suppose a cocycle $z$ restricts to zero.  The identity
\[
                     z(gh)=z(g)+\rho(g)z(h)
                     \tag{SP5}
\]
shows that it is constant on cosets of $G_\infty$.
It descends continuously to the quotient
$\Gamma=G_L/G_\infty$.  Continuity follows from the quotient
topology, since the quotient map is open.

Let $\tau$ be the topological generator with
$\chi(\tau)=\kappa$.  The nonzero scalar
$\kappa^{-E}-1\in\mathbb Z_\ell$ acts surjectively on $A$:
it is a unit times a power of $\ell$, and $A$ is
$\ell$-divisible.  Choose $b\in A$ such that
\[
                        (\kappa^{-E}-1)b=z(\tau).
\]
Subtract the coboundary $g\mapsto(\rho(g)-1)b$.
The resulting cocycle vanishes at $\tau$, hence at every
positive and negative integral power of $\tau$ by (SP5).
Those powers are dense in $\Gamma$, and the cocycle is
continuous with Hausdorff target.  It therefore vanishes
everywhere.  The original cocycle was a coboundary,
proving injectivity.
\end{proof}

The nonzero twist is essential to this proof.
For the trivial action, $H^1(\mathbb Z_\ell,A)$ contains
the nonzero characters
$x\mapsto \ell^{-r}x\bmod\mathbb Z_\ell$.
No vanishing statement for that action is used here.

\begin{lemma}\label{sp-finite-level}
Let $S$ be a finite set of distinct elements of
$\mathcal S_\ell(L,E)$.  For all sufficiently large
$\nu\ge c$, there is an injective map of sets
\[
                    S\longrightarrow
                    T_{\gamma_\nu,\kappa^E,\ell^\nu}.
                    \tag{SP6}
\]
The sufficient value of $\nu$ may depend on $S$.
\end{lemma}
\begin{proof}
For each $x\in S$ choose a continuous cocycle
$z_x:G_L\to A$ representing it.
The image of each $z_x$ is finite, because $G_L$ is compact
and $A$ is discrete.  As $S$ is finite, one power
$\ell^r$ annihilates all values of all these cocycles.

For each place $v\mid\ell$ choose one decomposition group
$D_v\subset G_L$.  The strict local condition says that
there exists $b_{x,v}\in A$ with
\[
                 z_x(h)=(\rho(h)-1)b_{x,v}
                         \qquad(h\in D_v).
                 \tag{SP7}
\]
There are only finitely many pairs $(x,v)$.
Increase $r$ so that $\ell^r$ also annihilates every
$b_{x,v}$.  Choose any $\nu\ge\max(c,r)$ and write
$K=K_\nu$, $G_K=\operatorname{Gal}(\overline L/K)$.

For $h\in G_K$, $\chi(h)\equiv1\pmod{\ell^\nu}$.
Consequently $\rho(h)$ acts trivially on $A[\ell^r]$.
It follows from (SP5) that $z_x|_{G_K}$ is a homomorphism,
with values in $A[\ell^\nu]$.

We verify its local conditions at every prime of $K$.
For arbitrary $g,h\in G_L$, the cocycle identity gives
\[
 z_x(ghg^{-1})
       =\rho(g)z_x(h)+(1-\rho(h))z_x(g).
       \tag{SP8}
\]
Here $\rho(ghg^{-1})=\rho(h)$ because the cyclotomic
character is abelian.  Combining (SP7) and (SP8), on the
conjugate decomposition group $gD_vg^{-1}$ the cocycle
is the coboundary with value
\[
                          \rho(g)b_{x,v}-z_x(g).
\]
This value is still annihilated by $\ell^r$, uniformly
in $g$.  Its coboundary vanishes on
$G_K\cap gD_vg^{-1}$.  These intersections are the
decomposition groups of primes of $K$ above $\ell$,
with the corresponding embeddings into $\overline L$.
Thus $z_x|_{G_K}$ is zero on each such decomposition group.

At a prime of $L$ away from $\ell$, the cyclotomic character
is unramified, so the action on inertia is trivial.
The unramified local condition on $x$ therefore implies
that $z_x$ itself, not just its cohomology class, is zero
on inertia: every coboundary for a trivial action is zero.
This remains true on conjugate inertia groups by (SP8).
Inertia groups at primes of $K$ away from $\ell$ are
subgroups of those groups.  We have proved that the
homomorphism $z_x|_{G_K}$ is unramified everywhere.
The unramifiedness of the cyclotomic character away from
$\ell$ can be seen at each finite level from
$X^{\ell^a}-1$.  Its reduction has $\ell^a$ distinct roots
in the algebraic residue field.  Hensel's lemma lifts all
of them to the maximal unramified local extension.

Apply (SP1).  There is a character
$\theta_x:\operatorname{Cl}(K)\to A[\ell^\nu]$
whose pullback to $G_K$ is $z_x|_{G_K}$.
Use the isomorphism
\[
       \mathbb Z/\ell^\nu\mathbb Z
              \longrightarrow A[\ell^\nu],
       \qquad \overline u\longmapsto
                    u/\ell^\nu\bmod\mathbb Z_\ell
                    \tag{SP9}
\]
to regard its values as residues, and define
$\varphi_x(P)=\theta_x([P])$.
As principal divisors have trivial class,
\[
                       \sum_Pv_P(\alpha)\varphi_x(P)=0.
\]

Choose $g\in G_L$ inducing $\gamma_\nu$ with $\chi(g)=\kappa$,
as permitted by (SP3).
For $h\in G_K$, equation (SP8) simplifies to
\[
                         z_x(ghg^{-1})=\kappa^{-E}z_x(h),
                         \tag{SP10}
\]
because $\rho(h)-1$ annihilates $z_x(g)$.
Conjugation by $g$ takes the arithmetic Frobenius of $P$
to that of $\gamma_\nu P$ in the maximal unramified
abelian quotient.  To see that this quotient is preserved,
conjugation preserves finite abelian unramified extensions
of $K$ and hence their maximal such extension $H_K$;
the Frobenius assertion follows on residue fields.
Thus
\[
                \varphi_x(\gamma_\nu P)
                         =\kappa^{-E}\varphi_x(P),
 \quad\text{or equivalently}\quad
                \varphi_x(\gamma_\nu^{-1}P)
                         =\kappa^E\varphi_x(P).
\]
Together with the principal-divisor relation, these are
exactly (SC6), including its inverse on $\gamma_\nu$.

Finally suppose $\varphi_x=\varphi_y$.
Prime ideal classes generate $\operatorname{Cl}(K)$,
so $\theta_x=\theta_y$, and hence
$z_x|_{G_K}=z_y|_{G_K}$.
Their restrictions to $G_\infty\subset G_K$ are equal.
Lemma~\ref{sp-restriction} implies $x=y$.
This proves the injection (SP6).
\end{proof}

\begin{theorem}\label{sp-sha}
Assume that a fixed integer $B$ satisfies
\[
                   |T_{\gamma,\kappa^E,\ell^\nu}|\le B
\]
for every cyclic datum (SC3) over $L$.
Then $\mathcal S_\ell(L,E)$ is finite and
\[
              \#\mathcal S_\ell(L,E)
                    =\ell^{s_\ell(L,E)}\le B.
                    \tag{SP11}
\]
\end{theorem}
\begin{proof}
For every finite subset $S$ of $\mathcal S_\ell(L,E)$,
Lemma~\ref{sp-finite-level} and the assumed uniform bound
give $|S|\le B$.
If the group had more than $B$ elements, one could choose
a subset of $B+1$ distinct elements, a contradiction.
Thus the group is finite of order at most $B$.

Every element has $\ell$-power order.  A sufficiently
large power $\ell^a$ therefore kills the finite group.
The filtration by $\ell^i\mathcal S_\ell(L,E)$ has finite
$\mathbb F_\ell$-vector-space quotients.  The sum of their
dimensions is its $\mathbb Z_\ell$-length, and multiplying
their cardinalities gives
$\#\mathcal S_\ell(L,E)=\ell^{s_\ell(L,E)}$.
\end{proof}

This proves assertion~\ref{sc-sha}, with the exact local
conditions and uniform quantifiers in its source statement.
The separate unit calculation in Theorem~\ref{su-count}
proves assertion~\ref{sc-count}, and
Theorem~\ref{sl-enlargement} proves assertion~\ref{sc-layer}.
The homological finite-set bound in assertion~\ref{sc-few}
is proved subsequently in Theorem~\ref{sw-few}.
There is no use of generic injectivity of restriction
in a finite extension of degree divisible by $\ell$:
the injective restriction proved above is to the
infinite cyclotomic kernel, and the finite-level maps
are constructed separately for each finite set.

\begin{corollary}\label{sp-conditional-coarse}
Assume the finite-field homology comparison in
assertion~\ref{sc-few}, with its displayed hypotheses,
for every totally imaginary number field and every cyclic
datum to which that assertion applies.
Let $L$ be totally imaginary, let $E\ge1$, and let $\ell$
be prime.  Suppose $\mu_\ell\subset L$ if $\ell$ is odd,
and $\mu_4\subset L$ if $\ell=2$.
Write $n=[L:\mathbb Q]$, $D=|\operatorname{disc}L|$,
and let $S$ be the primes above $\ell$.
Then the strict Selmer group is finite and
\[
 \log\#\mathcal S_\ell(L,E)
       \le\exp\bigl((2En)^{86}\log(2D)\bigr).
 \tag{SP12}
\]
Using Poitou--Tate and local Tate duality in the forms
stated before Proposition~\ref{ht-h2-from-selmer}, one also has
\[
 \log\#H^2(G_{L,S},\mathbb Z_\ell(E+1))
       \le\exp\bigl((2En)^{88}\log(2D)\bigr),
 \tag{SP13}
\]
and this cohomology group is finite.
Thus this modular implication requires only the finite-set
comparison, which Theorem~\ref{sw-few} supplies.
The classical reciprocity and duality theorems used in this
implication remain the explicitly stated inputs.
\end{corollary}
\begin{proof}
Section~\ref{ac-section} proves the exact arithmetic complex
and orbit estimate.  Proposition~\ref{ht-arithmetic-numerical}
therefore gives (SC20) with $C=13$.
Corollary~\ref{sl-class-number} gives the class-number bound
(SC15).  Theorems~\ref{su-count}, \ref{sp-sha}, and
\ref{sl-enlargement} prove the other three comparison
assertions.  Proposition~\ref{sc-exponent} now applies and
gives $5(\max(13,9)+2)+11=86$, proving (SP12).
The assumed roots of unity imply $\mu_\ell\subset L$
also when $\ell=2$.  Proposition~\ref{ht-h2-from-selmer}
applied with $C=86$ gives $C_K=86+2=88$ and (SP13).
\end{proof}



\section{Detection of global powers at degree-one primes}
\label{sd-section}

We now prove the power-detection statement used in
assertion~\ref{sc-few}.  The argument applies to every positive
integer exponent.  Its classical input is the following
qualitative form of the Chebotarev theorem: in a finite
Galois extension of \(\mathbb Q\), every conjugacy class
occurs as the Frobenius class of infinitely many unramified
rational primes.\cite[Theorem 8.31]{milne-ant}
No effective prime bound is needed.

A prime ideal \(P\) of a number field \(F\) has
\emph{absolute residue degree one} if its residue field
is \(\mathbb F_r\), where \(r\) is the rational prime
below \(P\).
For \(\beta\in F^\times\), reduction modulo \(P\) is
defined as a nonzero residue when \(v_P(\beta)=0\).
Equivalently, there are \(a,b\in\mathcal O_F\) with
\(b\notin P\), \(\beta b=a\), and \(a\notin P\).
The residue is then \(\overline a/\overline b\).
This is precisely the residue relation used by the source:
the definition of the local ring
\(\mathcal O_{F,P}\) gives the converse as well.

\begin{lemma}\label{sd-degree-one}
Let \(K/F\) be a finite Galois extension of number fields,
let \(\sigma\in\operatorname{Gal}(K/F)\), and let
\(\mathcal B\) be any finite set of primes of \(F\).
There exists an unramified prime \(P\notin\mathcal B\)
of absolute residue degree one, and a prime of \(K\)
above \(P\), whose arithmetic Frobenius is \(\sigma\).
\end{lemma}
\begin{proof}
Choose a finite Galois extension \(N/\mathbb Q\)
containing \(K\).  The restriction map
\(\operatorname{Gal}(N/F)\to\operatorname{Gal}(K/F)\)
is surjective: an \(F\)-embedding of the subfield \(K\)
extends to an \(F\)-automorphism of the normal
extension \(N/F\).
Choose a lift \(\widetilde\sigma\).

Chebotarev provides infinitely many rational primes
whose Frobenius class in \(\operatorname{Gal}(N/\mathbb Q)\)
is the conjugacy class of \(\widetilde\sigma\).
Omit the finitely many rational primes lying below
\(\mathcal B\).
For any remaining prime \(r\), choose a prime \(Q\) of
\(N\) at which the Frobenius element is exactly
\(\widetilde\sigma\), rather than a different conjugate.
This is possible because changing \(Q\) by an
automorphism conjugates its Frobenius.
Put \(P=Q\cap\mathcal O_F\).

For \(a\in\mathcal O_F\), the element
\(\widetilde\sigma\) fixes \(a\), whereas its action
on the residue field at \(Q\) is \(x\mapsto x^r\).
Every element of \(\mathcal O_F/P\) is therefore
fixed by the \(r\)-power map.  The polynomial
\(X^r-X\) has at most \(r\) roots, and the residue
field contains the \(r\) elements of \(\mathbb F_r\).
Consequently \(\mathcal O_F/P=\mathbb F_r\).
All these primes are unramified in \(N/\mathbb Q\).
Since \(NP=r\), the restriction of
\(\widetilde\sigma\) to \(K\) is also the arithmetic
Frobenius of \(K/F\) at the chosen prime above \(P\).
It is \(\sigma\), as required.
\end{proof}

\begin{lemma}\label{sd-roots-reduce}
Let \(F\) contain all \(q\)-th roots of unity, and let
\(P\nmid q\) be a finite prime.
Reduction modulo \(P\) is injective on \(\mu_q\).
\end{lemma}
\begin{proof}
Every root of unity and its inverse are algebraic integers,
so these reductions are defined and nonzero.
If \(\zeta\in\mu_q\) reduces to \(1\) and \(\zeta\ne1\),
then
\[
             1+\zeta+\cdots+\zeta^{q-1}=0.
\]
Reducing this equality modulo \(P\) gives \(q=0\),
contrary to \(P\nmid q\).
Apply this to the quotient of two roots of unity
with equal reduction.
\end{proof}

\begin{theorem}\label{sd-power}
Let \(F\) be a number field, let \(q\ge1\), and suppose
\(\mu_q\subset F\).
Let \(\beta\in F^\times\) and let \(\mathcal B\) be
a finite set of primes.
Suppose that, at every prime \(P\notin\mathcal B\)
of absolute residue degree one at which \(\beta\)
is a unit, its residue is a \(q\)-th power in
\((\mathcal O_F/P)^\times\).
Then \(\beta\) is a \(q\)-th power in \(F^\times\).
\end{theorem}
\begin{proof}
The case \(q=1\) is immediate.
For \(q>1\), choose \(\alpha\) with
\(\alpha^q=\beta\) and put \(K=F(\alpha)\).
The polynomial is separable and all its roots
are \(\zeta\alpha\), \(\zeta\in\mu_q\subset F\).
Thus \(K/F\) is Galois.

If \(\alpha\notin F\), choose
\(\sigma\in\operatorname{Gal}(K/F)\) with
\(\sigma\alpha\ne\alpha\).
Enlarge \(\mathcal B\) by the primes dividing \(q\)
and by the finite support of the principal divisor
of \(\beta\).
Lemma~\ref{sd-degree-one} gives a prime \(P\) and a
prime \(Q\) of \(K\) above it with Frobenius \(\sigma\).
Let \(r=NP\), a rational prime.
The equality \(\alpha^q=\beta\) shows that \(\alpha\)
is a unit at \(Q\).
By the hypothesis there is \(t\in\mathbb F_r^\times\)
with \(t^q=\overline\beta\).

All \(q\)-th roots of unity reduce to elements of
\(\mathbb F_r\), and they remain distinct by
Lemma~\ref{sd-roots-reduce}.
The \(q\) roots of \(X^q-\overline\beta\) in any
extension of the residue field are consequently the
elements \(\overline\zeta t\), all in \(\mathbb F_r\).
In particular \(\overline\alpha\in\mathbb F_r^\times\).
The Frobenius property gives
\[
             \overline{\sigma\alpha}
                  =\overline\alpha^{\,r}
                  =\overline\alpha.
\]
But \(\sigma\alpha/\alpha\in\mu_q\), so its reduction
is \(1\).  The same injectivity lemma forces
\(\sigma\alpha/\alpha=1\), a contradiction.
Hence \(\alpha\in F^\times\), proving the theorem.
\end{proof}

\begin{corollary}\label{sd-source}
Let \(L\subset F\) be number fields, let
\(\mu_q\subset F\), let \(\beta\in F^\times\),
and let \(0\ne b\in\mathcal O_F\).
Assume that for every prime \(P\nmid b\) such that
the natural map
\[
                 \mathcal O_L\longrightarrow\mathcal O_F/P
\]
is surjective, \(\beta\) has a nonzero residue
which is a \(q\)-th power.
Then \(\beta\in(F^\times)^q\).
In particular this proves the source assertion
\(\mathrm{SDetect}\), with \(q=\ell^\nu\).
\end{corollary}
\begin{proof}
At every absolute degree-one prime \(P\), the residue
field is \(\mathbb F_r\).  The image of \(\mathcal O_L\)
contains the image of \(\mathbb Z\), which already
equals \(\mathbb F_r\); the displayed map is therefore
surjective.  Apply Theorem~\ref{sd-power}, with the
primes dividing \(b\) excluded.
No normality assumption on \(F/L\) is required.
\end{proof}

The roots-of-unity hypothesis is essential to this argument
and to the resulting statement.  A precise counterexample
after deleting it is
\[
                      F=\mathbb Q,\qquad q=8,\qquad\beta=16.
\]
The element \(16\) is not an eighth power in \(\mathbb Q\),
because its \(2\)-adic valuation is \(4\).
Nevertheless it is an eighth power modulo every odd
rational prime \(r\).
Indeed at least one of \(2,-2,-1\) is a square in
\(\mathbb F_r^\times\): the square subgroup has index
two, and their product \(4\) is a square, so they cannot
all belong to the other coset.
If \(t^2=2\) or \(t^2=-2\), then \(t^8=16\).
If \(i^2=-1\), then \((1+i)^2=2i\), and hence
\((1+i)^8=16\).
All these elements are nonzero when \(r\) is odd.
This example does not satisfy \(\mu_8\subset\mathbb Q\),
which is an explicit hypothesis of the source.

The detection step thus has a complete conventional proof,
with Chebotarev as the stated classical input.
It supplies the final separation argument in
\(\mathrm{SFew}\); it does not construct the finite-field
homology classes or the integral cycles whose values must
be separated.


\section{The cyclic bar cycle and its relation}
\label{sh-section}

Two elementary homology inputs to the source's class construction
can be proved directly.  We give the actual cycle, compute the
homology group in which its diagonal images lie, and construct
a relation with a coefficient invertible modulo \(\ell^\nu\).
These are the assertions called \(\mathrm{SZc}\) and
\(\mathrm{SRel}\) in the source.

\subsection{The bar complex and the cyclic cycle}

For a group \(G\) and a commutative coefficient ring \(A\)
with trivial \(G\)-action, let \(C_j(G;A)\) be the free
\(A\)-module on the tuples
\([g_1|\cdots|g_j]\), with \(C_0(G;A)=A\).
For \(j\ge1\), put
\[
\begin{split}
 \partial[g_1|\cdots|g_j]
 ={}&[g_2|\cdots|g_j]\\
 &+\sum_{i=1}^{j-1}(-1)^i
       [g_1|\cdots|g_ig_{i+1}|\cdots|g_j]\\
 &+(-1)^j[g_1|\cdots|g_{j-1}].
\end{split}
\tag{SH1}
\]
In degree one the two empty tuples cancel.
The identity \(\partial^2=0\) follows by pairing the
two orders of deleting or multiplying any pair of
adjacent positions; the two terms have opposite signs.
The homology of this complex is \(H_j(G;A)\).
A homomorphism of groups acts on a tuple entry by entry
and commutes with \(\partial\).

Let \(q\ge2\), write \(C_q=\mathbb Z/q\mathbb Z\)
additively with generator \(t=\overline1\), and put
\(A=\mathbb Z/q\mathbb Z\).
For \(E\ge1\), define the following finite bar chain:
\[
 z_E=\sum_{j_1,\ldots,j_E\in C_q}
              [j_1|t|j_2|t|\cdots|j_E|t]
                   \in C_{2E}(C_q;A).
 \tag{SH2}
\]

\begin{lemma}\label{sh-cycle}
One has \(\partial z_E=0\).
\end{lemma}
\begin{proof}
The first face in (SH1) forgets \(j_1\).  Its sum over
\(j_1\) is therefore \(q\) times the same chain and is
zero in \(A\).
For \(1\le i<E\), pair the face multiplying
\(j_i\) and \(t\), with sign \(-1\), with the face
multiplying \(t\) and \(j_{i+1}\), with sign \(+1\).
After summing over all indices their unsigned values
are equal, since translation by \(t\) permutes \(C_q\)
in either variable.
Finally, the face multiplying \(j_E,t\), with sign
\(-1\), cancels the last face, with sign \(+1\),
by the same translation in \(j_E\).
These are all \(2E+1\) faces, proving the assertion.
\end{proof}

\subsection{The exact homology group}

\begin{proposition}\label{sh-cyclic-square}
For every \(m\ge0\),
\[
       H_m(C_q\times C_q;\mathbb Z/q\mathbb Z)
                    \simeq(\mathbb Z/q\mathbb Z)^{m+1}.
 \tag{SH3}
\]
\end{proposition}
\begin{proof}
Temporarily write \(t\) multiplicatively and put
\(B=A[C_q]\), \(N=1+t+\cdots+t^{q-1}\).
There is an augmented free \(B\)-resolution of the
trivial module \(A\):
\[
 \cdots\xrightarrow{\,t-1\,}B
       \xrightarrow{\,N\,}B
       \xrightarrow{\,t-1\,}B
       \xrightarrow{\,\varepsilon\,}A\longrightarrow0.
 \tag{SH4}
\]
To verify exactness, write \(u=\sum_{i=0}^{q-1}a_it^i\).
The equation \((t-1)u=0\) says that all \(a_i\) are
equal, so its kernel is \(AN=\operatorname{im}N\).
Also \(Nu=(\sum_i a_i)N\).
The kernel of \(N\) is therefore the same as the kernel
of augmentation.  The equations
\(b_{i-1}-b_i=a_i\), with cyclic indices, have a solution
precisely when \(\sum_i a_i=0\); successive partial
sums construct one.  Hence this kernel is
\((t-1)B\), as required.
These calculations hold over \(A\) without division.

The augmented resolution (SH4) is split exact as a
complex of \(A\)-modules.  Indeed \(AN\) is free of
rank one with a coordinate retraction, and the kernel
of augmentation is free of rank \(q-1\), with the
last coefficient determined by the others.
The short exact sequences just described consequently
split over \(A\).

Take the tensor product over \(A\) of two copies of
(SH4) and form the total complex.
It is a free resolution of \(A\) over
\[
           A[C_q]\otimes_A A[C_q]
                        =A[C_q\times C_q].
\]
For clarity, exactness follows by first taking homology
in one of the two nonnegative degrees: tensoring
preserves the explicitly split exact sequences above,
so only its degree-zero augmentation survives; the
other resolution then has homology \(A\) in degree
zero and zero otherwise.
There are exactly \(m+1\) summands
\(A[C_q\times C_q]\) in total degree \(m\).

After tensoring with the trivial module \(A\),
the maps \(t-1\) become zero and the maps \(N\)
become multiplication by \(q=0\) in \(A\).
The resulting complex has zero differentials and
degree-\(m\) term \(A^{m+1}\), proving (SH3).
This resolution computes the same homology as the
bar complex: lift the augmentation successively
between two free resolutions, and lift differences
of the resulting composites to construct chain
homotopies.  Thus their tensor complexes have
canonically isomorphic homology.
\end{proof}

\begin{lemma}\label{sh-primitive-relation}
Let \(A=\mathbb Z/\ell^\nu\mathbb Z\), where \(\ell\)
is prime and \(\nu\ge1\).
For any \(r+1\) vectors \(v_0,\ldots,v_r\in A^r\),
there are \(a_0,\ldots,a_r\in A\), at least one of
them a unit, with
\[
                         \sum_{i=0}^r a_iv_i=0.
\]
\end{lemma}
\begin{proof}
Lift the vectors to the columns of an integer
\(r\) by \(r+1\) matrix.
It has a nonzero rational kernel vector; clear
denominators and divide by the greatest common
divisor of its integer coordinates.
The resulting primitive integer vector is still
in the kernel.
At least one coordinate is not divisible by \(\ell\).
Reduce the relation modulo \(\ell^\nu\); that
coordinate becomes a unit.
The same proof works for \(r=0\), taking the
single coefficient to be \(1\).
\end{proof}

\begin{corollary}\label{sh-diagonal-relation}
Let \(q=\ell^\nu\), \(\nu\ge1\), and \(E\ge1\).
For \(0\le t\le2E+1\), define
\[
       d_t:C_q\longrightarrow C_q\times C_q,
                   \qquad x\longmapsto(x,tx).
\]
There are \(a_t\in\mathbb Z/q\mathbb Z\), an index
\(t_0\) with \(a_{t_0}\) a unit, and a bar chain
\(b\in C_{2E+1}(C_q\times C_q;\mathbb Z/q\mathbb Z)\)
such that
\[
                   \partial b
                      =\sum_{t=0}^{2E+1}a_t(d_t)_*z_E.
 \tag{SH5}
\]
\end{corollary}
\begin{proof}
All the displayed images are cycles by
Lemma~\ref{sh-cycle}.
Their \(2E+2\) homology classes belong to the module
\(A^{2E+1}\) computed in Proposition~\ref{sh-cyclic-square}.
Apply Lemma~\ref{sh-primitive-relation}.
The resulting zero homology class means precisely
that the indicated cycle is a boundary, providing \(b\).
\end{proof}

Under \(E=e+1\), the cycle (SH2) is exactly the
source's \(\mathrm{zc}(q,e)\), and the maps \(d_t\)
are its \(\mathrm{dt}(q,t)\).
Thus (SH5) has \(2e+4\) terms and proves the
actual relation required there.
The unused case \(\nu=0\) in the broad source
interface has the zero coefficient ring; all chains
are zero and the relation is immediate.

The prime-power restriction in
Lemma~\ref{sh-primitive-relation} cannot simply be
removed from its general module statement.
In \(\mathbb Z/6\mathbb Z\), a relation
\(3a+2b=0\) forces \(a\) even and \(b\) divisible
by \(3\), so neither coefficient can be a unit.
The source uses a prime-power modulus, as required.

These proofs establish the actual cyclic cycle and
relation, but do not yet supply the additive
finite-field cohomology functional or the decorated
arithmetic lattices that enter the later class
construction.


\section{Isomorphisms of ideal lattices with prescribed determinant}
\label{si-section}

The arithmetic cycle construction requires an isomorphism
between two direct sums of fractional ideals, with its
determinant specified as an element of the number field.
An equality of ideal classes alone would not specify this
determinant.  We prove the exact statement by elementary
two-dimensional operations.

Let \(R\) be a Dedekind domain with fraction field \(F\).
All fractional ideals below are nonzero.  Thus they are
invertible, and their products have unique factorizations
into integral powers of the nonzero prime ideals of \(R\).
For a fractional ideal \(I\), write \(v_P(I)\) for its
exponent at \(P\), and set
\[
 I^{-1}=\{x\in F:xI\subset R\}.
\]
In particular \(II^{-1}=R\).

\begin{lemma}\label{si-two-elements}
For fractional ideals \(C,D\), there are
\[
 p\in C\setminus\{0\},\qquad q\in D,\qquad
 r\in D^{-1},\qquad s\in C^{-1}
\]
such that \(ps-qr=1\).
\end{lemma}
\begin{proof}
Choose any nonzero \(p\in C\).
The ideal \(A=(p)C^{-1}\) is a nonzero integral ideal.
There are only finitely many prime ideals containing it;
denote their set by \(\mathcal P\).
For each \(P\in\mathcal P\), choose
\(q_P\in D\setminus PD\).
Such an element exists: a finite set of generators of
\(D\) has minimum \(P\)-valuation \(v_P(D)\), whereas
every element of \(PD\) has strictly larger valuation.

The Chinese remainder theorem gives elements
\(e_P\in R\) which are \(1\) modulo \(P\) and
\(0\) modulo every other member of \(\mathcal P\).
Put \(q=\sum_{P\in\mathcal P}e_Pq_P\in D\).
Its class in \(D/PD\) is the nonzero class of \(q_P\).
Consequently the integral ideal \((q)D^{-1}\) is not
contained in any prime of \(\mathcal P\).
If \(\mathcal P\) is empty, take \(q=0\), since
then \(A=R\).

It follows that
\[
                 (p)C^{-1}+(q)D^{-1}=R.
\]
Indeed a proper ideal on the left would lie in a maximal
ideal containing \(A\), and every such ideal belongs
to \(\mathcal P\).
Writing \(1=ps+qw\) with
\(s\in C^{-1}\) and \(w\in D^{-1}\), and setting
\(r=-w\), proves the assertion.
\end{proof}

\begin{lemma}\label{si-move}
For any fractional ideals \(I,J,C\), there exists
\(M\in\operatorname{SL}_2(F)\) such that
\[
                       M(I\oplus J)
                            =CI\oplus C^{-1}J.
\]
\end{lemma}
\begin{proof}
Apply Lemma~\ref{si-two-elements} to \(C\) and
\(D=CIJ^{-1}\), and put
\[
       M=\begin{pmatrix}p&q\\r&s\end{pmatrix},
       \qquad
       M^{-1}=\begin{pmatrix}s&-q\\-r&p\end{pmatrix}.
\]
The determinant is exactly \(1\).
The containments
\[
 pI\subset CI,\quad qJ\subset CI,\quad
 rI\subset C^{-1}J,\quad sJ\subset C^{-1}J
\]
show one inclusion in the asserted equality.
For the inverse, use
\[
 sCI\subset I,\quad qC^{-1}J\subset I,\quad
 rCI\subset J,\quad pC^{-1}J\subset J.
\]
These follow immediately from the four ideal memberships
in Lemma~\ref{si-two-elements} and
\(D^{-1}=C^{-1}I^{-1}J\).
Thus the restriction of \(M\) is an isomorphism of the
specified \(R\)-modules.
\end{proof}

\begin{proposition}\label{si-equal-product}
Let \(I_1,\ldots,I_n\) and \(J_1,\ldots,J_n\) be
fractional ideals with
\[
                         \prod_{i=1}^n I_i
                              =\prod_{i=1}^n J_i.
\]
Then an element of \(\operatorname{SL}_n(F)\) maps
\(\bigoplus_i I_i\) onto \(\bigoplus_i J_i\).
For \(n=0\), use the unique automorphism of the zero
space, whose determinant is \(1\).
\end{proposition}
\begin{proof}
For \(n=1\) the ideals are equal and the identity works.
For \(n\ge2\), designate the last coordinate as an
auxiliary coordinate.
For each \(i<n\), let \(I'_i\) be the ideal currently
in coordinate \(i\), and put \(C=J_i(I'_i)^{-1}\).
Apply Lemma~\ref{si-move} in coordinates \(i,n\),
acting as the identity on the other coordinates.
The new ideal in coordinate \(i\) is \(J_i\),
the product of all coordinate ideals is unchanged,
and the determinant of the operation is \(1\).
After these \(n-1\) operations, the first \(n-1\)
coordinate ideals are the required ones.
Cancellation of their invertible product shows that
the last ideal is \(J_n\) as well.
The product of the matrices has determinant \(1\).
\end{proof}

\begin{theorem}\label{si-prescribed}
Let \(I_i,J_i\), \(1\le i\le n\), be fractional
ideals, and let \(\beta\in F^\times\).
There exists \(g\in\operatorname{GL}_n(F)\) with
\[
          g\Bigl(\bigoplus_i I_i\Bigr)
                   =\bigoplus_i J_i,
          \qquad \det g=\beta,
\]
if and only if
\[
             \prod_iJ_i=(\beta)\prod_iI_i,
\]
provided that \(n>0\).  For \(n=0\), existence is
equivalent to \(\beta=1\).
\end{theorem}
\begin{proof}
For necessity when \(n>0\), take the top exterior
power inside \(\bigwedge^n_F F^n=F\).
The exterior power of the first ideal lattice is
\(\prod_iI_i\), that of the second is
\(\prod_iJ_i\), and the induced map is multiplication
by \(\det g\).

For sufficiency, multiply the first coordinate by
\(\beta\).  This has determinant \(\beta\) and
replaces the product of the source ideals by
\((\beta)\prod_iI_i=\prod_iJ_i\).
Apply Proposition~\ref{si-equal-product} afterwards.
The composite has the required image and determinant.
The statement in dimension zero follows from the
definition of the empty determinant.
\end{proof}

To match the source, specialize to
\(R=\mathcal O_F\) and let
\(a_i,b_i\) be finitely supported integral divisors.
Define
\[
 I(a_i)=\prod_P P^{a_i(P)},\qquad
 I(b_i)=\prod_P P^{b_i(P)}.
\]
Membership in \(I(a_i)\) means
\(v_P(x)\ge a_i(P)\) at every prime, with \(x=0\)
always allowed.  These are exactly the valuation
inequalities defining the source's ideal lattice.
The hypothesis
\[
              v_P(\beta)=\sum_i\bigl(b_i(P)-a_i(P)\bigr)
              \qquad\text{for every }P
\]
is exactly the equality of fractional ideals in
Theorem~\ref{si-prescribed}.
Thus the theorem proves the source assertion
\(\mathrm{SStein}\), including its condition
that the index set be nonempty or \(\beta=1\).
There is no assumption that an individual ideal is
principal.

Both qualifications matter.  In dimension one an
isomorphism of the required kind would force the
ratio ideal to be principal, so arbitrary pairs of
ideals cannot replace the displayed hypothesis.
In dimension zero all divisor sums vanish, so the
divisor hypothesis alone would permit any unit
\(\beta\); nevertheless the determinant remains
\(1\).  For example, \(\beta=-1\) over \(\mathbb Q\)
would contradict the conclusion if the source's
empty-index condition were deleted.
Neither example satisfies all the actual hypotheses.


\section{An integral interpolation of cyclic characters}
\label{se-section}

The next input prescribes one pushforward of an integral
function on a product of two cyclic groups, while making
several other pushforwards zero.  Working over the integers
is necessary: the function will be used as the exponents
of fractional ideals.  No division by a difference of
slopes is permitted in the group ring.

For a finite abelian group \(A\), write
\(\mathbb Z[A]\) for the free abelian group on symbols
\([a]\), with multiplication \([a][b]=[a+b]\).
Its multiplicative identity is \([0]\), which we also
write as \(1\).
A homomorphism \(f:A\to B\) induces
\(f_*:\mathbb Z[A]\to\mathbb Z[B]\) by
\([a]\mapsto[f(a)]\).
Thus if \(\eta=\sum_a\eta(a)[a]\), the coefficient
of \([b]\) in \(f_*\eta\) is
\(\sum_{f(a)=b}\eta(a)\).

\begin{lemma}\label{se-difference}
Let \(w,M\) be nonnegative integers, put \(\kappa=1+w\),
and write
\[
 P_M(T)=\prod_{j=0}^{M-1}(T-\kappa^j)
                    =\sum_{i=0}^M p_iT^i.
\]
Then in \(\mathbb Z[X]\),
\[
                (X^w-1)^M
                    \ \bigm|\ \sum_{i=0}^M p_iX^{\kappa^i}.
\]
\end{lemma}
\begin{proof}
Let \(U\) be the endomorphism of the abelian group
\(\mathbb Z[X]\) given by \(Uf(X)=f(X^\kappa)\),
and put \(D=X^w-1\).
For every polynomial \(y\), we have
\(Uy\equiv y\pmod D\), since \(X^\kappa\equiv X\pmod D\).
Moreover
\[
 UD=D S,\qquad
 S=\sum_{a=0}^{\kappa-1}X^{wa},\qquad
 S\equiv\kappa\pmod D.
\]
These identities also hold when \(w=0\), with
\(D=0\) and \(U\) the identity.
If \(f=D^jy\), then
\[
 (U-\kappa^j)f
   =D^j\bigl((S^j-\kappa^j)Uy
                      +\kappa^j(Uy-y)\bigr)
\]
is divisible by \(D^{j+1}\).
Starting with \(X\) and applying
\((U-1),(U-\kappa),\ldots,(U-\kappa^{M-1})\)
therefore gives a polynomial divisible by \(D^M\).
The operators commute, and \(U^iX=X^{\kappa^i}\).
Their product applied to \(X\) is consequently the
polynomial in the assertion.
The empty product \(M=0\) gives the same conclusion.
\end{proof}

\begin{lemma}\label{se-divisor}
Let \(\ell\) be prime, \(q=\ell^\nu\), and
\(0\le c\le\nu\).
If \(d\) is a nonzero integer with \(|d|<\ell^c\),
then the image of \(d\) in \(\mathbb Z/q\mathbb Z\)
divides the image of \(\ell^c\).
Consequently, in \(\mathbb Z[\mathbb Z/q\mathbb Z]\),
\[
                  [d]-1\ \bigm|\ [\ell^c]-1.
\]
\end{lemma}
\begin{proof}
Write \(d=\ell^a u\) with \(\ell\nmid u\).
The size hypothesis gives \(a<c\).
The residue class of \(u\) is a unit modulo \(q\),
so a solution to \(dv=\ell^c\) is
\(v=\ell^{c-a}u^{-1}\pmod q\).
Choose a nonnegative integer representative \(v_0\).
Then \([\ell^c]=[d]^{v_0}\), and the usual geometric
sum proves divisibility by \([d]-1\).
If \(v_0=0\), the numerator is zero and divisibility
still holds.
\end{proof}

\begin{theorem}\label{se-interpolation}
Let \(\ell\) be prime and let \(c,\nu,M\) be
nonnegative integers satisfying
\[
                 c\le\nu,\qquad M<\ell^c.
\]
Put \(q=\ell^\nu\), \(A=\mathbb Z/q\mathbb Z\),
and \(\kappa=1+\ell^c\).
For \(0\le t\le M\), define
\[
             \chi_t:A^2\longrightarrow A,
             \qquad \chi_t(x,y)=x+ty.
\]
Fix \(t_0\in\{0,\ldots,M\}\), and write
\[
              P_M(T)=\prod_{j=0}^{M-1}(T-\kappa^j)
                       =\sum_{i=0}^Mp_iT^i.
\]
There exists \(\eta\in\mathbb Z[A^2]\) such that
\[
  (\chi_t)_*\eta=0\quad(t\ne t_0),
  \qquad
  (\chi_{t_0})_*\eta
                       =\sum_{i=0}^M p_i[\kappa^i].
\]
\end{theorem}
\begin{proof}
All residue classes in this proof lie in \(A\).
Let
\[
          \Delta=\sum_{i=0}^Mp_i[\kappa^i],
          \qquad D=[\ell^c]-1.
\]
Apply the ring homomorphism
\(\mathbb Z[X]\to\mathbb Z[A]\), \(X\mapsto[1]\),
to Lemma~\ref{se-difference}.
It gives \(D^M\mid\Delta\).
For each \(t\ne t_0\) in the indicated range,
\(0<|t_0-t|\le M<\ell^c\).
Lemma~\ref{se-divisor} therefore gives
\([t_0-t]-1\mid D\).
There are exactly \(M\) such slopes, so
\[
          B:=\prod_{t\ne t_0}([t_0-t]-1)
                      \ \bigm|\ D^M\ \bigm|\ \Delta.
\]
Choose \(Q\in\mathbb Z[A]\) with \(BQ=\Delta\).
This argument multiplies divisibility identities;
it never cancels a zero divisor in the group ring.

Let \(\iota:\mathbb Z[A]\to\mathbb Z[A^2]\)
be induced by \(a\mapsto(a,0)\), and put
\[
 \eta=
 \prod_{t\ne t_0}\bigl([(0,1)]-[(t,0)]\bigr)
 \ \iota\!\left(\left[-\sum_{t\ne t_0}t\right]Q\right).
\]
If \(s\ne t_0\), the factor indexed by \(t=s\)
has zero image under \((\chi_s)_*\), so
\((\chi_s)_*\eta=0\).
For the distinguished slope,
\[
 (\chi_{t_0})_*\bigl([(0,1)]-[(t,0)]\bigr)
       =[t_0]-[t]=[t]([t_0-t]-1).
\]
Also \((\chi_{t_0})_*\iota\) is the identity.
The product of the factors \([t]\) cancels the
explicit unit \([-\sum t]\), leaving \(BQ=\Delta\).
This proves both required pushforward identities.
\end{proof}

This is precisely the source assertion
\(\mathrm{SEta}\).  In the arithmetic cycle
construction, \(M=2E+1\), and the hypothesis is
\(2E+1<\ell^c\).
The proof retains the exact integer coefficients, so
they can be used as positive or negative powers of
fractional ideals.
It also includes \(M=0\), when the product is empty,
and \(c=\nu\), when \([\ell^c]-1=0\).
In the latter case with \(M>0\), the target
\(\Delta\) is zero by the proved divisibility;
there is no illicit division by zero.


\section{Explicit homology functionals over a finite field}
\label{sf-section}

We construct the finite-field functional used to evaluate
the arithmetic cycles.  The construction uses a cyclic
two-cocycle and a sum over projective lines.  In particular,
this step does not require a theorem about the algebraic
\(K\)-groups of finite fields.

We use the bar chains and boundary of
Section~\ref{sh-section}.  For a commutative coefficient
ring \(A\), a degree-\(m\) cochain is an \(A\)-linear
map \(C_m(G;A)\to A\); it is a cocycle if it vanishes
on boundaries.  The group acts trivially on \(A\).

\subsection{A cyclic cocycle with an exact value}

Let \(k\) be a finite field of cardinality \(Q\),
let \(q\ge2\), and let \(\zeta\in k^\times\) have
order \(q\).  Put \(R=\mathbb Z/q\mathbb Z\), and
let \(v:k^\times\to(R,+)\) be a homomorphism.
For \(x\in\mathbb Z\), write
\(\chi_x:C_q\to k^\times\) for
\(\chi_x(j)=\zeta^{xj}\), where \(C_q=\mathbb Z/q\mathbb Z\).
For \(E\ge0\), define
\[
 z_E=\sum_{j_1,\ldots,j_E\in C_q}
             [j_1|1|\cdots|j_E|1]\in C_{2E}(C_q;R).
\]
When \(E=0\), this means the single empty tuple.
For \(E\ge1\), it is the cycle proved in
Lemma~\ref{sh-cycle}.

\begin{lemma}\label{sf-carry}
There is a symmetric integer-valued function
\(u_0:k^\times\times k^\times\to\mathbb Z\) such that
\[
 u_0(b,c)-u_0(ab,c)+u_0(a,bc)-u_0(a,b)=0
\]
and, writing \(u\) for its reduction modulo \(q\),
\[
                \sum_{j\in C_q}u(\zeta^{xj},\zeta^x)=x
                         \qquad\text{in }R.
\]
\end{lemma}
\begin{proof}
The multiplicative group of a finite field is cyclic.
Here is the elementary reason: if \(m\) is the least
common multiple of its element orders, every element is
a root of \(X^m-1\), so its cardinality is at most \(m\).
Lagrange's theorem gives the reverse inequality.
In a finite abelian group, an element of order \(m\)
is obtained by taking, for each prime, a component of
largest prime-power order and multiplying these components.
The component orders are relatively prime.
Thus there is a generator of \(k^\times\).

Put \(N=Q-1=qr\).  We may choose a generator \(g\)
with \(g^r=\zeta\).
Indeed, for an initial generator \(h\), write
\(\zeta=h^{ra}\), where \((a,q)=1\).
Chinese remainders give an integer \(b\equiv a\pmod q\)
which is \(1\) modulo each prime dividing \(N\) but
not \(q\).  Then \((b,N)=1\), and \(g=h^b\) works.

For \(s\in k^\times\), let \(a(s)\in\{0,\ldots,N-1\}\)
be its exponent relative to \(g\), and define
\[
                  u_0(s,t)
                       =\frac{a(s)+a(t)-a(st)}{N}.
\]
The numerator is divisible by \(N\).
Symmetry follows from commutativity, and the displayed
cocycle identity follows by substituting this formula:
every exponent cancels.

Set \(s=\zeta^x\).  Summing over \(j\in C_q\),
the terms \(a(s^j)\) and \(a(s^{j+1})\) cancel,
even when \(s\) has order less than \(q\).  Hence
\[
       \sum_{j\in C_q}u_0(s^j,s)
             =\frac{q\,a(s)}{N}=\frac{a(s)}r.
\]
Since \(g^r=\zeta\), the last integer is the least
nonnegative residue of \(x\) modulo \(q\).
Reducing modulo \(q\) proves the assertion.
\end{proof}

For commuting entries \(h_1,\ldots,h_m\) and an element
\(g\) commuting with them, define the insertion chain
\[
 P_g[h_1|\cdots|h_m]
   =\sum_{i=0}^m(-1)^i
       [h_1|\cdots|h_i|g|h_{i+1}|\cdots|h_m].
\]
This is the source's prism on such chains; its conjugations
disappear because of the stated commutation.

\begin{proposition}\label{sf-cyclic-functional}
For every \(E\ge0\), there is a cocycle
\(t:C_{2E+1}(k^\times;R)\to R\) satisfying
\[
              t(P_g(\chi_x)_*z_E)=x^E v(g)
            \qquad(x\in\mathbb Z,\ g\in k^\times).
\]
\end{proposition}
\begin{proof}
Using Lemma~\ref{sf-carry}, set
\[
 t([h_1|\cdots|h_{2E+1}])
       =\prod_{a=1}^E u(h_{2a-1},h_{2a})\,
                                      v(h_{2E+1}).
\]
This is the cup product \(u^E\smile v\).
For completeness, if \(f\) has degree \(a\) and
\(h\) has degree \(b\), their cup product on a tuple
is \(f\) of its first \(a\) entries times \(h\) of
its last \(b\) entries.
The bar differential gives
\[
       \delta(f\smile h)
          =(\delta f)\smile h+(-1)^a f\smile\delta h.
\]
To check this identity, expand both sides by the faces:
the faces before the dividing position are those of
\(\delta f\), and those after it are those of
\(\delta h\), shifted by \(a\) positions.
The last face of the first group and the first face
of the second group cancel with opposite signs;
the remaining common face multiplies the two entries
at the dividing position.  These are exactly all
the faces of the tuple on the left.
The cocycle identity of Lemma~\ref{sf-carry} and
\(v(ab)=v(a)+v(b)\) therefore imply \(\delta t=0\).

Evaluate \(t\) on an insertion chain of even length.
The terms inserting \(g\) just before and just after
\(h_{2a-1}\) cancel, because \(u\) is symmetric.
Only insertion in the last position remains.  Thus
\[
 t(P_g[h_1|\cdots|h_{2E}])
       =\prod_{a=1}^E u(h_{2a-1},h_{2a})\,v(g).
\]
For \((\chi_x)_*z_E\), the \(E\) summation indices
are independent, so the value is
\[
       \left(\sum_{j\in C_q}u(\zeta^{xj},\zeta^x)\right)^E
                  v(g)=x^Ev(g).
\]
This includes the empty product \(E=0\).
\end{proof}

\subsection{Transfer over projective lines}

We give the chain construction, including the independence
of its choices.  Let a group \(G\) act on a finite set
\(X\) on the right, and let \(M\) be an abelian group.
Suppose
\[
 c(gh,x)=c(g,x)c(h,xg),\qquad c(1,x)=1,
\]
for a function \(c:G\times X\to M\).
For a tuple \([g_1|\cdots|g_m]\), put
\(x_i=xg_1\cdots g_i\), \(x_0=x\), and define
\[
 T_c[g_1|\cdots|g_m]
    =\sum_{x\in X}
       [c(g_1,x_0)|c(g_2,x_1)|\cdots|c(g_m,x_{m-1})].
\]
Extend this map linearly over any commutative ring \(A\).

\begin{lemma}\label{sf-transfer}
The maps \(T_c\) commute with the bar boundary.
If
\[
                c'(g,x)=b(x)^{-1}c(g,x)b(xg)
\]
for a function \(b:X\to M\), then \(T_{c'}\) and
\(T_c\) are chain homotopic.  In particular they give
the same value on cycles after applying a cocycle.
\end{lemma}
\begin{proof}
For an internal face, the two adjacent entries in
the image multiply to \(c(g_ig_{i+1},x_{i-1})\),
by the cocycle identity.  The last face deletes the
same entry on either side.  For the first face,
replace the summation index \(x\) by \(xg_1\).
This is a permutation of \(X\), so the first faces
also agree.  This proves the chain-map assertion.

Here is an explicit homotopy.  On a degree-\(m\) tuple,
write \(c_i=c(g_i,x_{i-1})\),
\(c'_i=c'(g_i,x_{i-1})\), and \(b_i=b(x_i)\).
Set
\[
 K_m[g_1|\cdots|g_m]
 =\sum_{x\in X}\sum_{i=0}^m(-1)^i
          [c_1|\cdots|c_i|b_i|c'_{i+1}|\cdots|c'_m].
\]
Then
\[
                    dK_m+K_{m-1}d=T_{c'}-T_c.
\]
Indeed the faces not adjacent to the inserted entry
are precisely the corresponding terms of
\(-K_{m-1}d\), after the first-face permutation of
\(X\) already described.  The two faces adjacent
to successive inserted entries cancel by
\(c_i b_i=b_{i-1}c'_i\).
The two unpaired end faces are the all-\(c'\) tuple
and minus the all-\(c\) tuple.
This accounts for every face and its sign.
For a cycle \(z\), the identity expresses
\(T_{c'}z-T_cz\) as the boundary \(dK_mz\).
\end{proof}

Let \(V\) now be a finite-dimensional vector space
over \(k\).  Write \(\mathbb P(V)\) for its set
of one-dimensional subspaces; it is empty for \(V=0\).
Choose one nonzero vector in each line.  This defines
a homogeneous function \(\lambda:V\setminus\{0\}\to k^\times\)
by expressing a vector as \(\lambda(v)\) times the
chosen vector in its line.  Thus
\(\lambda(av)=a\lambda(v)\) for \(a\ne0\).

For a representation \(\rho:G\to\operatorname{GL}(V)\),
use the right action \(\ell g=\rho(g)^{-1}\ell\),
and put
\[
                  c_\rho(g,\ell)
                    =\frac{\lambda(v)}{\lambda(\rho(g)^{-1}v)},
                    \qquad 0\ne v\in\ell.
\]
Homogeneity makes this independent of \(v\).
It satisfies the cocycle identity in
Lemma~\ref{sf-transfer}.
Changing \(\lambda\) gives precisely the coboundary
in that lemma, with
\(b(\ell)=\lambda(v)/\lambda'(v)\).

\begin{proposition}\label{sf-additive-transfer}
Let \(A\) be a commutative ring in which \(Q=1\),
and let \(t:C_m(k^\times;A)\to A\) be a cocycle
of positive degree.  For every finite-dimensional
\(k\)-space \(V\), define
\[
                    \theta_V=t\circ T_{c_{\mathrm{id}}}.
\]
These functionals vanish on boundaries.  Their values
on cycles are invariant under linear isomorphisms.
For representations \(\rho\) on \(V\) and \(\sigma\)
on \(W\), and any cycle \(z\in C_m(G;A)\),
\[
 \theta_{V\oplus W}((\rho\oplus\sigma)_*z)
       =\theta_V(\rho_*z)+\theta_W(\sigma_*z).
\]
In dimension one, the scalar representation identifies
\(\theta_k\) with \(t\) on cycles.
\end{proposition}
\begin{proof}
The chain-map property proves the assertion about
boundaries.  Pulling the homogeneous function through
a linear isomorphism gives identical summands under
the resulting bijection of projective lines.
Any other choice gives the same value on cycles by
Lemma~\ref{sf-transfer}.

For the direct sum, choose
\[
 \lambda(v,w)=
 \begin{cases}
     \lambda_V(v),&v\ne0,\\
     \lambda_W(w),&v=0.
 \end{cases}
\]
Send a line represented by \((v,w)\) to \(kv\in\mathbb P(V)\)
when \(v\ne0\), and otherwise to \(kw\in\mathbb P(W)\).
This map is equivariant for \(\rho\oplus\sigma\).
The scalar cocycle is the pullback of the one on
the appropriate summand.
A fiber over \(kv\in\mathbb P(V)\), after fixing
the nonzero representative \(v\), consists uniquely
of the lines \(k(v,w)\) with \(w\in W\).
It has \(Q^{\dim W}\) elements.
A fiber over a line in \(W\) has one element.
Consequently the transfer for the direct sum is
\(Q^{\dim W}T_\rho+T_\sigma=T_\rho+T_\sigma\)
with coefficients in \(A\).
Choice independence proves the claimed formula for
the originally chosen functionals on cycles.

For \(V=k\), there is one projective line.
Choosing \(\lambda(v)=v\) gives
\(c(g)=g\), so its transfer is the identity.
\end{proof}

\subsection{The determinant value}

\begin{lemma}\label{sf-determinant}
For \(g\in\operatorname{GL}(V)\),
\[
                 \prod_{\ell\in\mathbb P(V)}c(g,\ell)
                              =\det(g)\quad\text{in }k^\times.
\]
\end{lemma}
\begin{proof}
Denote the product by \(D(g)\).
The cocycle identity and the permutation of the
projective lines imply \(D(gh)=D(g)D(h)\).
Changing \(\lambda\) does not change the product,
because all factors \(b(\ell)^{-1}b(\ell g)\) cancel.

Choose coordinates and take \(\lambda(v)\) to be
the first nonzero coordinate of \(v\).
For \(g=\operatorname{diag}(a_1,\ldots,a_n)\), a line
whose first nonzero coordinate has index \(i\)
contributes \(a_i\).  There are \(Q^{n-i}\) such
lines, obtained by making that coordinate equal
to \(1\).  Since \(a_i^Q=a_i\) in the finite field,
\[
                      D(g)=\prod_i a_i^{Q^{n-i}}
                            =\prod_i a_i=\det(g).
\]
If \(Q=2\), the target group \(k^\times\) is
trivial and the result already follows.
Otherwise choose \(a\in k\setminus\{0,1\}\).
For \(i\ne j\), let \(T_{ij}(s)=I+sE_{ij}\),
and let \(d\) be diagonal with entry \(a\) at
\(i\) and \(1\) elsewhere.  Then
\[
      dT_{ij}(s)d^{-1}T_{ij}(s)^{-1}=T_{ij}((a-1)s).
\]
Every transvection is therefore a commutator, so
the homomorphism \(D\) is \(1\) on transvections.
Gaussian elimination writes every invertible matrix
as a product of diagonal matrices and transvections:
if a pivot is zero, add a row having a nonzero entry
in that column, then clear the remaining column
entries and continue.  Row scalings supply the
diagonal matrices.
Both \(D\) and determinant agree on these generators,
which proves the lemma.  In dimension zero both
products are \(1\).
\end{proof}

\begin{theorem}\label{sf-functional}
Under the hypotheses of Proposition~\ref{sf-cyclic-functional},
there are functionals
\[
             \theta_V:C_{2E+1}(\operatorname{GL}(V);R)\to R
\]
which vanish on boundaries, are invariant under
linear isomorphisms on cycles, and are additive
under direct sums on cycles.  They satisfy
\[
 \theta_V\bigl(P_g(\chi_{x,V})_*z_E\bigr)
                         =x^E v(\det g),
 \qquad
 \chi_{x,V}(j)=\zeta^{xj}I_V.
\]
\end{theorem}
\begin{proof}
Since \(q\mid Q-1\), we have \(Q=1\) in \(R\).
Apply Proposition~\ref{sf-additive-transfer} to the
cyclic cocycle \(t\).
Scalar matrices act trivially on \(\mathbb P(V)\),
and their scalar cocycle is the scalar itself.
Thus, directly from the tuple formula for transfer,
\[
 T\bigl(P_g(\chi_{x,V})_*z_E\bigr)
     =\sum_{\ell\in\mathbb P(V)}
          P_{c(g,\ell)}(\chi_x)_*z_E.
\]
Indeed, before the inserted \(g\), all projective
lines remain fixed by the preceding scalar entries;
after it, the line may change, but the remaining
scalar cocycle entries do not depend on that line.
Apply Proposition~\ref{sf-cyclic-functional} and then
Lemma~\ref{sf-determinant} to obtain
\[
 x^E\sum_\ell v(c(g,\ell))
       =x^E v\!\left(\prod_\ell c(g,\ell)\right)
       =x^E v(\det g).
\]
If the coefficient modulus is \(q=1\), the same
assertion holds with all functionals zero; no
nontrivial root of unity is then required.
\end{proof}

Taking \(E=e+1\) proves the source assertions
\(\mathrm{SCoc}\), \(\mathrm{STrans}\), and
\(\mathrm{STheta}\), with their stated coefficients
and normalization.  The determinant calculation also
explains why no exceptional claim about
\(\operatorname{GL}_2(\mathbb F_2)\) is needed:
its multiplicative residue group is trivial.

\begin{corollary}\label{sf-primitive-family}
Let \(A\) be a commutative coefficient ring, let
\(\theta\) be any positive-degree additive family
with values in \(A\) as above, let \(B\) be a commutative ring, and let
\(r:B\to k\) be a ring homomorphism.
Composing reduction of matrices with \(\theta\)
gives a family on \(\operatorname{GL}_n(B)\) which
vanishes on boundaries, is invariant on cycles under
\(g\mapsto\operatorname{diag}(g,I)\), and is
multiplied by \(d\) under
\(g\mapsto\operatorname{diag}(g,\ldots,g)\)
with \(d\) copies.
\end{corollary}
\begin{proof}
Reduction is a group homomorphism and therefore
commutes with the bar boundary.
The contribution of a trivial representation on
a positive-degree cycle is zero.
To see this without an acyclicity assumption,
the bar chain group of the trivial group is \(A\)
in every degree; its boundary is the identity in
positive even degrees and zero in odd degrees.
Its positive-degree homology is consequently zero.
The image of the cycle under a trivial representation
is therefore a boundary, annihilated by \(\theta\).
Direct-sum additivity now proves stabilization
invariance, and iteration proves multiplication
by \(d\).  For \(d=0\) the representation is trivial
on the zero space and gives zero.
\end{proof}

This corollary is exactly the source assertion
\(\mathrm{SPhi}\).  The finite-field evaluations and
their stable additivity are now proved.  Constructing
the integral arithmetic cycles with the prescribed
evaluations, and bounding their possible values by
the stable rational homology calculation, remain
separate tasks.


\section{Integral arithmetic cycles with prescribed residue values}
\label{sa-section}

We now construct the arithmetic cycles used in the class
count.  This proves the source assertion \(\mathrm{SClass}\).
The proof retains the integral ideal lattices, the exact
determinants, and the inverse Galois action in the defining
polynomial.  The matrix size is allowed to depend on the
input; the subsequent stable homology argument is where
that dependence must be removed.

\subsection{A cycle in a lattice stabilizer}

For an element \(g\) of a group, put
\(f_g(a)=g^{-1}ag\).  On bar chains define
\[
 P_g[a_1|\cdots|a_m]
  =\sum_{i=0}^m(-1)^i
    [a_1|\cdots|a_i|g|f_g(a_{i+1})|\cdots|f_g(a_m)].
\]
Here and below the coefficient ring is arbitrary and
commutative.

\begin{lemma}\label{sa-prisms}
The following chain identities hold:
\[
       \partial P_g+P_g\partial=(f_g)_*-1,
\]
and
\[
       \partial Q_{g,h}-Q_{g,h}\partial
                 =P_g+P_h(f_g)_*-P_{gh},
\]
where the degree-two homotopy is
\begin{align*}
 Q_{g,h}[a_1|\cdots|a_m]
   =\sum_{0\le i\le j\le m}(-1)^{i+j}
   [&a_1|\cdots|a_i|g|
       f_g(a_{i+1})|\cdots|f_g(a_j)|h|\\
    &f_{gh}(a_{j+1})|\cdots|f_{gh}(a_m)].
\end{align*}
Consequently, on cycles,
\(P_{gh}-P_g-P_h(f_g)_*\) takes values in boundaries.
\end{lemma}
\begin{proof}
Expand the bar boundaries.  In the first identity,
the faces away from the inserted \(g\) give
\(-P_g\partial\).  The adjacent faces cancel using
\(a g=g f_g(a)\).  The two end faces give
\((f_g)_*\) and minus the identity.

For the second identity, the faces away from both
inserted entries give \(Q_{g,h}\partial\).
Faces next to \(g\) or \(h\) cancel between successive
values of \(i\) or \(j\), using
\(a g=g f_g(a)\) and
\(f_g(a)h=h f_{gh}(a)\).
Deleting an initial \(g\) gives \(P_h(f_g)_*\),
and deleting a final \(h\) gives \(P_g\).
When \(i=j\), merging the adjacent \(g,h\) gives
minus \(P_{gh}\): its sign is
\((-1)^{2i}(-1)^{i+1}=-(-1)^i\).
This accounts for all remaining faces, and proves
the displayed identity with its stated sign.
\end{proof}

\begin{lemma}\label{sa-stabilizer-cycle}
Let \(H\subset G\) be a subgroup, let
\(\rho:D\to G\) and \(d_t:C\to D\) be homomorphisms,
and let \(z\in C_m(C;A)\) be a cycle.
Suppose
\[
                 \partial b=\sum_t n_t(d_t)_*z
\]
for a chain \(b\in C_{m+1}(D;A)\).
Let \(g_t\in G\), and fix an index \(t_0\), writing
\(g_0=g_{t_0}\).
Assume that \(g_t\) commutes with
\(\rho(d_t(C))\), that both \(\rho(D)\) and
\(g_0^{-1}\rho(D)g_0\) lie in \(H\), and that
\(g_0^{-1}g_t\in H\) for every \(t\).
Then there is a cycle \(y\in C_{m+1}(H;A)\)
whose image in \(G\) is homologous to
\[
                    \sum_t n_tP_{g_t}(\rho d_t)_*z.
\]
\end{lemma}
\begin{proof}
Regard \(\rho_0=\rho\) and
\(\rho_1=g_0^{-1}\rho g_0\) as homomorphisms into
\(H\), and put \(k_t=g_0^{-1}g_t\in H\).
Commutation gives
\(f_{k_t}\rho_1d_t=\rho_0d_t\).
Set
\[
 y=(\rho_1)_*b-(\rho_0)_*b
                +\sum_t n_tP_{k_t}(\rho_1d_t)_*z.
\]
By the first identity of Lemma~\ref{sa-prisms}, its
boundary is
\[
 ((\rho_1)_*-(\rho_0)_*)\partial b
       +\sum_t n_t((\rho_0d_t)_*-(\rho_1d_t)_*)z=0.
\]
In \(G\), the second identity of that lemma gives
\[
 \sum_t n_tP_{g_t}(\rho d_t)_*z
 \equiv P_{g_0}\rho_*\partial b
                 +\sum_t n_tP_{k_t}(\rho_1d_t)_*z
                    \pmod{\text{boundaries}}.
\]
The first identity, now applied to \(\rho_*b\),
identifies its first term modulo boundaries with
\((\rho_1)_*b-(\rho_0)_*b\).  The result is \(y\).
\end{proof}

\subsection{Integral modules and good reductions}

\begin{lemma}\label{sa-projective}
Every finitely generated torsion-free module over a
Dedekind domain is a direct summand of a finite free
module.
\end{lemma}
\begin{proof}
For a nonzero fractional ideal \(I\), choose finite
families \(u_i\in I\), \(v_i\in I^{-1}\) with
\(\sum_i u_iv_i=1\), using \(II^{-1}=R\).
The maps
\[
 R^r\longrightarrow I,\quad (a_i)\longmapsto\sum_i a_iu_i,
 \qquad
 I\longrightarrow R^r,\quad x\longmapsto(v_ix)_i
\]
have composite the identity on \(I\).
Thus every such ideal is a finite projective module.

For a finitely generated torsion-free module \(M\)
of positive rank, choose a nonzero linear functional
on \(M\otimes_R\operatorname{Frac}(R)\).
Its image on \(M\) is a nonzero fractional ideal
\(I\).  The surjection \(M\to I\) splits by the
preceding construction.  Its kernel is finitely
generated, torsion-free, and has one less rank;
finite generation follows because \(R\) is Noetherian.
Induction expresses \(M\) as a finite direct sum
of fractional ideals.  Taking the direct sum of
their finite free splittings proves the assertion.
Rank zero gives the zero module, since \(M\) is
torsion-free.
\end{proof}

Let \(F/L\) be a finite Galois extension, and put
\(R=\mathcal O_L\), \(S=\mathcal O_F\).
The module \(S\) is finite over \(R\): an integral
basis over \(\mathbb Z\) already generates it over
\(R\).  It is torsion-free, and hence projective
by Lemma~\ref{sa-projective}.

\begin{lemma}\label{sa-split-reduction}
There exists \(0\ne b\in R\) such that, for every
field \(k\) and homomorphism \(r:S\to k\) whose
restriction to \(R\) is surjective and for which
\(r(b)\ne0\), the map
\[
 S\otimes_R k\longrightarrow
              \prod_{\tau\in\operatorname{Gal}(F/L)} k,
 \qquad x\otimes a\longmapsto(a\,r(\tau x))_\tau
\]
is an isomorphism.
We may enlarge \(b\) so that any specified finite
list of elements of \(F^\times\) are units in
\(S[1/b]\), and any specified finite list of
\(L\)-linear automorphisms of \(F^J\) preserve
\(S[1/b]^J\), for a finite index set \(J\).
\end{lemma}
\begin{proof}
Choose an integral \(\vartheta\in S\) with
\(F=L(\vartheta)\).  Such an element can be obtained
without a separability assumption beyond the one
already present: take an integral basis
\(a_1,\ldots,a_n\) of \(S\) over \(\mathbb Z\)
and consider \(\sum_i t^{i-1}a_i\).
For each nonidentity \(\tau\in\operatorname{Gal}(F/L)\),
the equality between this element and its \(\tau\)-image
is a nonzero polynomial equation in the integer \(t\).
Avoiding the finitely many roots for all \(\tau\)
makes the orbit have \([F:L]\) elements, as required.

Let \(f(T)\in R[T]\) be the monic minimal polynomial
of \(\vartheta\).  Its coefficients are integral
because they are symmetric functions of the integral
conjugates, and they lie in \(L\).
Its discriminant \(\delta\in R\) is nonzero.
As \(R[\vartheta]\) spans \(F\) over \(L\),
clearing the denominators of a finite set of generators
of \(S\) gives \(0\ne a\in R\) such that
\(S[1/a]=R[1/a][\vartheta]\).
Take initially \(b=a\delta\).

In a good residue field, the elements
\(r(\tau\vartheta)\) are all distinct: the product
of their squared differences is the nonzero residue
of \(\delta\).
They give all the roots of \(f\) in \(k\).
Chinese remainders applied to these distinct linear
factors give
\(k[T]/(\overline f)\simeq\prod_\tau k\), with
the stated evaluation maps.  This is the required
tensor-product isomorphism.

For the extra conditions, clear denominators of
each specified nonzero element and its inverse.
For an automorphism, clear the denominators of
the images of a finite set of \(R\)-generators of
\(S^J\), both for it and for its inverse.
Multiplying \(b\) by these finitely many nonzero
elements of \(R\) proves all assertions together.
\end{proof}

If a group \(H\) acts on the projective module
\(\mathcal M=S^J\), choose a complement \(P\) and
an isomorphism \(\mathcal M\oplus P\simeq R^N\).
Let \(H\) act trivially on \(P\).
This embeds the representation in \(\operatorname{GL}_N(R)\).
After reduction, \(\mathcal M\otimes_R k\) is the
space in Lemma~\ref{sa-split-reduction}, repeated
over \(J\).
By Corollary~\ref{sf-primitive-family}, any positive-degree
cycle has the same value under an additive family
before and after the trivial complementary summand
is added.

\subsection{The exact arithmetic construction}

Fix a prime \(\ell\), integers \(\nu\ge c\ge1\),
and assume \(c\ge2\) if \(\ell=2\).
Put
\[
 q=\ell^\nu,\qquad \kappa=1+\ell^c,
 \qquad d=\ell^{\nu-c},\qquad A=\mathbb Z/q\mathbb Z.
\]
Suppose \(F/L\) is cyclic of degree \(d\), generated
by \(\gamma\), and that \(F\) contains a primitive
\(q\)-th root \(\zeta\) with
\(\gamma\zeta=\zeta^\kappa\).
The order of \(\kappa\) in \(A^\times\) is exactly
\(d\).  This follows from the valuation identity
\(v_\ell(\kappa^m-1)=c+v_\ell(m)\) proved in
Lemma~\ref{sc-vandermonde}.

Let \(E\ge1\), assume \(M:=2E+1<\ell^c\), and put
\(a=\kappa^E\).
For \(\alpha\in F^\times\), suppose an integral
divisor \(\mathfrak d\) satisfies
\[
       v_P(\alpha)=\mathfrak d(\gamma P)
                                     -a\mathfrak d(P)
                   \qquad\text{for every prime }P.
\]
Define the weighted norm and the integer factor
\[
 \mathcal N(\alpha)=\prod_{j=0}^{d-1}(\gamma^j\alpha)^{a^j},
 \qquad
 c_E=\prod_{\substack{0\le j<M\\j\ne E}}(a-\kappa^j).
\]

\begin{theorem}\label{sa-class}
There are an integer \(N\), a cycle
\(y\in C_{2E+1}(\operatorname{GL}_N(R);A)\),
and \(0\ne b\in R\) with the following property.
Let \(r:S\to k\) be a homomorphism to a finite field,
surjective on \(R\), with \(r(b)\ne0\).
Let \(v:k^\times\to(A,+)\) be a homomorphism, and
let \(\theta\) be an additive family which vanishes
on boundaries.  For \(x\in A\), write
\(\rho_x(j)=\overline\zeta^{\,xj}I_V\) for the scalar
representation of \(C_q\) on \(V\), and suppose
\[
       \theta_V\bigl(P_g(\rho_x)_*z_E\bigr)
                             =x^E v(\det g).
\]
Then \(\mathcal N(\alpha)\) is a unit in \(S[1/b]\),
and
\[
       \theta_{k^N}(\operatorname{red}(y))
                       =c_E\,v\bigl(r(\mathcal N(\alpha))\bigr)
                        \qquad\text{in }A.
\]
Here reduction of a localized unit means its well-defined
nonzero residue.  The family constructed in
Theorem~\ref{sf-functional} has the required properties
whenever \(\overline\zeta\) remains primitive.
\end{theorem}
\begin{proof}
We give the construction in five steps.

\emph{The cyclic relation and ideal exponents.}
Let \(J=A^2\), and let
\(d_t:C_q\to C_q^2\) be \(j\mapsto(j,tj)\),
for \(t=0,\ldots,M\).
Corollary~\ref{sh-diagonal-relation} gives coefficients
\(n_t\in A\), an index \(t_0\) with \(n_{t_0}\)
a unit, and a chain \(b_1\) such that
\[
                 \partial b_1=\sum_{t=0}^M n_t(d_t)_*z_E.
\]
Write
\[
 P(T)=\prod_{j=0}^{M-1}(T-\kappa^j)=\sum_{i=0}^M p_iT^i,
 \qquad Q(T)=\frac{P(T)}{T-a}\in\mathbb Z[T].
\]
Theorem~\ref{se-interpolation} gives integer exponents
\(\eta_j\), \(j\in J\), with zero pushforward under
\(\chi_t(j_1,j_2)=j_1+tj_2\) for \(t\ne t_0\),
and pushforward \(\sum_i p_i[\kappa^i]\) for \(t=t_0\).
Let \(I=\prod_P P^{\mathfrak d(P)}\), and put
\[
        V=F^J,\qquad \mathcal M=S^J,
        \qquad \mathcal M'=\bigoplus_{j\in J}I^{\eta_j}.
\]
On \(V\), let \(C_q^2\) act diagonally by
\[
     \rho(u_1,u_2)(x_j)_j
            =(\zeta^{j_1u_1+j_2u_2}x_j)_j.
\]
Both ideal lattices are preserved by this action.

\emph{Isomorphisms commuting with the characters.}
For \(t\ne t_0\), divide the coordinates into the
blocks \(\chi_t(j)=x\), \(x\in A\).
The product of the target ideals in each block is
\(I^{\sum_{\chi_t(j)=x}\eta_j}=S\).
Theorem~\ref{si-prescribed} gives an \(F\)-linear
isomorphism in each block, of determinant \(1\),
from its standard lattice to its target lattice.
Their direct sum \(g_t\) maps \(\mathcal M\) onto
\(\mathcal M'\) and commutes with \(\rho d_t\).
Empty blocks have determinant \(1\), as required.

For the distinguished slope, write \(\chi=\chi_{t_0}\).
For \(x=\kappa^i\in\langle\kappa\rangle\subset A^\times\),
put \(T_x=\gamma^i\) and \(B(x)=1\).
For all other \(x\), put \(T_x=1\) and \(B(x)=x\).
The equality of the orders of \(\kappa\) and \(\gamma\)
makes \(T_x\) well defined.
The coordinatewise \(L\)-linear isomorphism
\[
             U(x_j)_j=(T_{\chi(j)}^{-1}x_j)_j
\]
preserves \(\mathcal M\) and changes \(\rho d_{t_0}\)
to the diagonal character with weight \(B(\chi(j))\).
The target ideal in its \(j\)-th coordinate has
divisor
\[
             \eta_j\,\mathfrak d(T_{\chi(j)}P).
\]

For a polynomial \(Q(T)=\sum_i q_iT^i\), define
\[
                       \beta=\prod_i(\gamma^{-i}\alpha)^{q_i}.
\]
The divisor hypothesis on \(\alpha\) gives
\begin{align*}
 v_P(\beta)
 &=\sum_i q_i\bigl(\mathfrak d(\gamma^{i+1}P)
                                      -a\mathfrak d(\gamma^iP)\bigr)\\
 &=\sum_i p_i\mathfrak d(\gamma^iP).
\end{align*}
Thus, in the transformed block of weight \(1\),
the product of the target ideals is \((\beta)\).
Indeed this follows by pushing forward \(\eta\)
and using \(T_{\kappa^i}=\gamma^i\), with repeated
residue classes grouped together.
In every other block the product ideal is \(S\):
there \(T_x=1\), and the relevant sum of exponents
is zero.
The weight-one block is nonempty, since
\(j=(1,0)\) has \(\chi(j)=1\).
Theorem~\ref{si-prescribed} therefore gives a block
map \(g'\) with determinant \(\beta\) in this block
and determinant \(1\) in all the others, carrying
\(\mathcal M\) onto \(U\mathcal M'\).
Put \(g_{t_0}=U^{-1}g'U\).
It maps \(\mathcal M\) onto \(\mathcal M'\) and
commutes with \(\rho d_{t_0}\).
It is only required to be \(L\)-linear; the
coordinatewise Galois changes need not be \(F\)-linear.

\emph{The integral cycle.}
Let \(H=\operatorname{Aut}_R(\mathcal M)\).
All \(g_t\) have the same image lattice, so
\(g_{t_0}^{-1}g_t\in H\).
Since \(\rho\) preserves both lattices, both
\(\rho\) and \(g_{t_0}^{-1}\rho g_{t_0}\) take
values in \(H\).
Choose \(b\) as in Lemma~\ref{sa-split-reduction},
large enough that every \(g_t\) preserves
\(\mathcal M[1/b]\) and that \(\alpha\) and
\(\beta\) are units in \(S[1/b]\).
We may include \(q\) as a factor of \(b\).
Apply Lemma~\ref{sa-stabilizer-cycle} inside
\(\operatorname{Aut}_{R[1/b]}(\mathcal M[1/b])\).
It gives a cycle \(y_0\) in \(H\) homologous,
in this localized group, to
\[
                       \sum_t n_tP_{g_t}(\rho d_t)_*z_E.
\]
Divide \(y_0\) by the unit \(n_{t_0}\in A\),
and extend the representation of \(H\) by a trivial
complement as described after Lemma~\ref{sa-split-reduction}.
Its image is the required integral cycle
\(y\in C_{2E+1}(\operatorname{GL}_N(R);A)\).

\emph{Evaluation after reduction.}
Fix a good residue map \(r\).
All the preceding homology relations can be reduced,
because their chains lie in the localized group.
The trivial complementary summand has zero value.
Under the splitting in Lemma~\ref{sa-split-reduction},
the diagonal character on the component indexed by
\(\tau=\gamma^j\) has its weights multiplied by
\(\kappa^j\) modulo \(q\).
The coordinatewise Galois change \(U\) is an
invertible linear change after reduction; explicitly,
it permutes the embedding coordinate by
\(\tau\mapsto\tau T_{\chi(j)}^{-1}\) in coordinate \(j\).
Therefore invariance under isomorphisms and additivity
allow evaluation in the blocks used above.

Each block prism is a cycle: its scalar cyclic
character commutes with its block map, so the first
identity of Lemma~\ref{sa-prisms} and
\(\partial z_E=0\) apply.
For \(t\ne t_0\), every block determinant is \(1\),
and its value is zero.
For \(t=t_0\), only the weight-one block contributes.
Its determinant on the \(\gamma^j\)-component is
\(r(\gamma^j\beta)\).
The normalization assumed for \(\theta\) gives
\[
       \theta\bigl(P_{g_{t_0}}(\rho d_{t_0})_*z_E\bigr)
               =\sum_{j=0}^{d-1}a^j v(r(\gamma^j\beta)).
\]
The chosen division by \(n_{t_0}\) leaves exactly
this value for \(y\).

\emph{The polynomial and the weighted norm.}
On \(S[1/b]^\times\), put \(f(x)=v(r(x))\), and
\[
                         F(x)=\sum_{j=0}^{d-1}a^j f(\gamma^jx).
\]
The base element \(b\) is fixed by the Galois group,
so all these arguments are localized units.
Since \(a^d\equiv1\pmod q\), reindexing the cyclic
sum gives
\(F(\gamma^{-1}x)=aF(x)\) in \(A\).
Also \(F(xy)=F(x)+F(y)\), including negative powers.
Consequently
\[
     F(\beta)=Q(a)F(\alpha)
            =c_E\,f(\mathcal N(\alpha)).
\]
This is precisely the required value, and
\(\mathcal N(\alpha)\) is a localized unit because
every factor defining it is one.
\end{proof}

All the hypotheses and the evaluation constant in
Theorem~\ref{sa-class} agree with \(\mathrm{SClass}\),
after writing \(E=e+1\).
The proof makes no principal-ideal assumption,
does not confuse an \(L\)-linear map with an
\(F\)-linear map, and retains the complementary
module when embedding in a free integral module.
The construction supplies the arithmetic classes.
Proposition~\ref{sb-odd-rank} and Theorem~\ref{sb-upper}
then prove the stable rational homology estimate
\(\mathrm{BOdd}\) and the finite-image bound
\(\mathrm{SUpper}\); Theorem~\ref{sw-few} combines them
with these classes to complete the finite-set comparison.


\section{The odd-degree rank bound and the finite-image estimate}
\label{sb-section}

The remaining homology comparison has a short proof using the
classical rational homology of the algebraic \(K\)-theory space.
We state those classical inputs explicitly.  All subsequent
integral counting arguments are proved below.  This proves the
assertions called \(\mathrm{BOdd}\) and \(\mathrm{SUpper}\) in the
source, with their actual constants; it does not purport to
reproduce the source's longer proof by differential forms.

Let \(L\) be a totally imaginary number field, put
\(R=\mathcal O_L\), and write \(r=r_2(L)\).
For \(N,j\ge0\), retain the notation
\[
 t_L(N,j)=\#H_j(\operatorname{GL}_N(R),\mathbb Z)_{\rm tors}.
\]
The finiteness needed here was proved in
Sections~\ref{ac-section} and~\ref{ht-section}.
All homology has trivial coefficients unless otherwise stated.

\subsection{Stable homology and the repeated direct-sum map}

As in Section~\ref{kt-section}, put
\(X=B\operatorname{GL}(R)^+\), the connected component of the
algebraic \(K\)-theory infinite loop space.  Its multiplication
is induced by direct sum of matrices.  It is homotopy
associative and homotopy commutative, has a unit, and
\[
 H_j(X,\mathbb Z)=H_j(\operatorname{GL}(R),\mathbb Z),
 \qquad \pi_j(X)=K_j(R)\quad(j\ge1).
 \tag{SB1}
\]
We use the following established results.

\begin{enumerate}
\item Integral homological stability identifies
\(H_j(\operatorname{GL}_N(R),\mathbb Z)\) with
\(H_j(X,\mathbb Z)\) for
\(N\ge\max(2j+1,j+2)\).
The stable-range condition for \(R\) was proved in
Section~\ref{kt-section}.
\cite[IV, Corollary 1.14.1]{weibel-kbook}
\item For a path-connected, homotopy associative \(H\)-space,
its rational homology is generated, under its homology
product, by primitive elements.  The primitive elements
in positive degree are precisely the image of its rational
homotopy under the Hurewicz map, and that map is injective.
Here an element \(v\) is primitive when
\(\Delta v=v\otimes1+1\otimes v\), where \(\Delta\)
comes from the diagonal of the space.
\cite[Corollary 4.18, Theorem 5.18, and Appendix]{milnor-moore}
\item For every odd \(j\ge3\),
\(\dim_{\mathbb Q}(K_j(R)\otimes\mathbb Q)=r\).
For \(j=1\), this dimension is \(r-1\): the classical
identity \(K_1(R)=R^\times\) and
Theorem~\ref{su-unit-theorem} give this assertion.
The identity in degree one is the Bass--Milnor--Serre
theorem, not an assertion about arbitrary Dedekind domains.
\cite[III, Example 1.3.6; IV, Theorems 1.17--1.18]{weibel-kbook}
\end{enumerate}

These are classical inputs with their stated hypotheses;
we do not reprove the general stability or rational
homotopy theorems here.  In particular, no quantitative
comparison of integral \(K\)-groups with homology is being
assumed in this section.

For a positive integer \(d\), let
\[
 \psi_{d,N}:\operatorname{GL}_N(R)
       \longrightarrow\operatorname{GL}_{dN}(R),\qquad
       g\longmapsto\operatorname{diag}(g,\ldots,g).
\]
Its stable homology map will be denoted by \(\Psi_d\).
The choice of an ordering of the blocks changes the map
by conjugation and therefore does not change homology.
On \(X\), this is the map which adds a point to itself
\(d\) times under the direct-sum multiplication.

\begin{lemma}\label{sb-weight}
Let \(m\ge1\).  In rational stable homology, the image of
\[
       \prod_{j=2}^{m}(\Psi_d-d^j)
           :H_m(X,\mathbb Q)\longrightarrow H_m(X,\mathbb Q)
 \tag{SB2}
\]
is contained in the space of primitive elements of degree
\(m\).  An empty product, when \(m=1\), means the identity.
\end{lemma}
\begin{proof}
If \(v\) is primitive, its iterated diagonal is the sum
of the \(d\) tensors having \(v\) in one position and
\(1\) elsewhere.  Multiplication therefore gives
\(\Psi_d(v)=dv\).  Homotopy commutativity implies that
the repeated-sum map preserves the homology product:
equivalently, the product and the diagonal in the graded
commutative homology algebra show that
\(\Psi_d(vw)=\Psi_d(v)\Psi_d(w)\).
Consequently a product of \(s\) positive-degree
primitive elements has eigenvalue \(d^s\).

Primitive generation expresses every element in degree
\(m\) as a sum of such products.  Every factor has
degree at least one, so \(1\le s\le m\).
If \(s\ge2\), the factor \(\Psi_d-d^s\) in (SB2)
annihilates the product.  The remaining terms are single
primitive elements of degree \(m\).  This also proves
the statement for \(d=1\); no division by an eigenvalue
difference has been used.
\end{proof}

We now keep track of the finite matrix sizes.  Fix
\(d\ge1\), \(N\ge1\), and an integral homology class
\(x\in H_m(\operatorname{GL}_N(R),\mathbb Z)\).
Let \(N_t=d^tN\), let \(s_t\) be the usual inclusion
from size \(N_t\) to size \(N_{t+1}\), and define
\[
 \begin{split}
 Z_{d,0}(x)&=x,\\
 Z_{d,t+1}(x)&=(\psi_{d,N_t})_*Z_{d,t}(x)
                         -d^{t+2}(s_t)_*Z_{d,t}(x).
 \end{split}
 \tag{SB3}
\]
These are homomorphisms of integral homology groups.
They are exactly the repeated direct-sum differences
used by the source.

\begin{proposition}\label{sb-odd-rank}
If \(m\) is positive and odd and \(N\ge2m+4\), then
\[
 \operatorname{rank}_{\mathbb Z}
 \operatorname{im}\bigl(Z_{d,m-1}:
 H_m(\operatorname{GL}_N(R),\mathbb Z)
 \longrightarrow H_m(\operatorname{GL}_{d^{m-1}N}(R),\mathbb Z)
                         \bigr)\le r .
 \tag{SB4}
\]
The extra hypothesis \(m<r\) in the submitted assertion
is not needed for this rank estimate.
\end{proposition}
\begin{proof}
All sizes in (SB3) are at least \(N\), since \(d\ge1\).
They lie in the integral stable range.  After these
identifications and tensoring with \(\mathbb Q\),
\(Z_{d,m-1}\) becomes (SB2).
Its image is primitive by Lemma~\ref{sb-weight}.
The dimension of that space is \(r\) for odd
\(m\ge3\), and \(r-1\) for \(m=1\), by the classical
inputs above.  Tensoring with \(\mathbb Q\) is exact,
so it computes the rank of the integral image.
This proves (SB4).
\end{proof}

\subsection{Two elementary counting lemmas}

\begin{lemma}\label{sb-finite-rank}
Let \(H\) be an abelian group with finite torsion subgroup,
let \(W\subset H\) have rank at most \(r\), and let
\(f:H\to B\) be a homomorphism to an abelian group
annihilated by a positive integer \(q\).  Every finite
subset of \(f(W)\) has cardinality at most
\[
                     q^r\,\#H_{\rm tors}.
 \tag{SB5}
\]
In particular, \(f(W)\) is finite with that bound.
Finite generation of \(H\) is not required.
\end{lemma}
\begin{proof}
Lift a given finite subset to finitely many elements of
\(W\), and let \(W'\) be the subgroup they generate.
The structure theorem for finitely generated abelian
groups gives
\(W'\simeq\mathbb Z^s\oplus T'\), with
\(s\le r\) and \(T'\subset H_{\rm tors}\).
The image of the free part has at most \(q^s\) elements,
and the image of \(T'\) has at most \(\#H_{\rm tors}\)
elements.  This proves (SB5).
If \(f(W)\) had more elements, choosing a finite subset
of one greater than the bound would contradict (SB5).
\end{proof}

\begin{lemma}\label{sb-coefficients}
For every group \(G\), positive integer \(q\), and
\(m\ge1\), there is a natural exact sequence
\[
 0\longrightarrow H_m(G,\mathbb Z)/qH_m(G,\mathbb Z)
 \longrightarrow H_m(G,\mathbb Z/q)
 \xrightarrow{\ \beta\ }H_{m-1}(G,\mathbb Z)[q]
 \longrightarrow0.
 \tag{SB6}
\]
Here \(H[q]=\{h\in H:qh=0\}\).
\end{lemma}
\begin{proof}
Use the integral bar complex \(C_\bullet\), whose groups
are free abelian, and its quotient \(C_\bullet/qC_\bullet\).
Lift a cycle \(\overline z\) of degree \(m\) to
\(z\in C_m\).  There is \(w\in C_{m-1}\) with
\(\partial z=qw\).  Since \(C_{m-2}\) has no
\(q\)-torsion, \(\partial w=0\); for \(m=1\)
use the ordinary complex with \(\partial:C_0\to0\).
Set \(\beta[\overline z]=[w]\).
Changing the lift by \(qu\) changes \(w\) by
\(\partial u\), and changing the cycle by a boundary
has the same harmless effect.  Thus \(\beta\) is a
well-defined homomorphism.  Its value is killed by
\(q\), because \(qw=\partial z\).

If \([w]=0\), write \(w=\partial u\).  Then
\(z-qu\) is an integral cycle reducing to
\(\overline z\), proving exactness in the middle.
Conversely every integral cycle has zero image under
\(\beta\).  If an integral cycle \(z\) reduces to a
boundary, write \(z=\partial v+qu\).  Applying
\(\partial\) and using the freeness of the chain groups
gives \(\partial u=0\).  Therefore the kernel of
reduction on integral homology is exactly \(qH_m\).
Finally, if \([w]\in H_{m-1}[q]\), choose a cycle
\(w\) and a chain \(z\) with \(qw=\partial z\).
The reduction of \(z\) is a cycle mapping to \([w]\).
This proves surjectivity and every term of (SB6).
\end{proof}

\subsection{The finite-image estimate}

Fix a prime \(\ell\), an integer \(\nu\ge0\), and
\(A=\mathbb Z/\ell^\nu\mathbb Z\).  Let \(B\) be
an \(A\)-module, and suppose for each \(N\ge0\) that
\[
 \Phi_N:C_m(\operatorname{GL}_N(R);A)\longrightarrow B
\]
is \(A\)-linear and vanishes on boundaries.
Assume that, on cycles, it is unchanged by stabilization
and satisfies
\[
     \Phi_{dN}((\psi_{d,N})_*y)=d\Phi_N(y)
                         \qquad(d\ge0).
 \tag{SB7}
\]
These are the precise properties required of the
family; they are proved for the finite-field functionals
in Section~\ref{sf-section}.

\begin{theorem}\label{sb-upper}
Let \(m\ge1\) be odd and \(N_0=2m+4\).
Every finite subset \(S\subset B\) whose elements have
the form
\[
                \ell^{m-1}\Phi_N(y),
       \qquad \partial y=0,
 \tag{SB8}
\]
where \(N\) and \(y\) may depend on the element,
satisfies
\[
 \#S\le (\ell^\nu)^r\,
       t_L(N_0,m)\,t_L(N_0,m-1).
 \tag{SB9}
\]
As in Proposition~\ref{sb-odd-rank}, the extra
source hypothesis \(m<r\) can be omitted here.
\end{theorem}
\begin{proof}
If \(\nu=0\), then \(A\) is the zero ring and every
unital \(A\)-module is zero.  The assertion follows.
Suppose henceforth \(\nu\ge1\), and write
\(q=\ell^\nu\).
There are finitely many elements of \(S\), so choose
one size \(N\ge N_0\) larger than all their given sizes.
Stabilization expresses them all in the form (SB8)
at this common size.

For each size, reduction of integral cycles and then
\(\Phi\) defines a homomorphism
\[
 \phi_N:H_m(\operatorname{GL}_N(R),\mathbb Z)
                        \longrightarrow B.
\]
It is well-defined because integral boundaries reduce
to boundaries.  It inherits stabilization and (SB7).
Applying it to (SB3), with \(d=\ell\), gives
\[
 \phi_{\ell^{m-1}N}(Z_{\ell,m-1}(x))
    =\prod_{j=2}^{m}(\ell-\ell^j)\phi_N(x)
    =\ell^{m-1}u\phi_N(x),
 \quad
 u=\prod_{j=2}^{m}(1-\ell^{j-1}).
 \tag{SB10}
\]
Each factor of \(u\) is congruent to \(1\) modulo
\(\ell\), so \(u\) is invertible modulo \(q\).
Choose \(v\in\mathbb Z\) with \(uv\equiv1\pmod q\).
It follows from (SB10) that every
\(\ell^{m-1}\phi_N(x)\) belongs to
\[
 \phi_{\ell^{m-1}N}
       (\operatorname{im}Z_{\ell,m-1}),
\]
by replacing \(x\) in the right-hand side by \(vx\).
Proposition~\ref{sb-odd-rank} bounds the rank of this
image by \(r\).  Integral stability identifies the
torsion subgroup at size \(\ell^{m-1}N\) with the one
at size \(N_0\).  Lemma~\ref{sb-finite-rank} therefore
bounds the image of all integral cycles, after
multiplication by \(\ell^{m-1}\), by
\(q^r t_L(N_0,m)\).

It remains to account for cycles modulo \(q\) which
need not lift to integral cycles.  Apply (SB6) to
\(G=\operatorname{GL}_N(R)\).
Its right-hand group is a subgroup of the integral
torsion in degree \(m-1\), of order at most
\(t_L(N_0,m-1)\) by stability.
The homomorphism induced by
\(\ell^{m-1}\Phi_N\) maps each coset of the kernel
of \(\beta\) to a translate of the image on that
kernel.  The preceding paragraph bounds that image
by \(q^r t_L(N_0,m)\), and there are at most
\(t_L(N_0,m-1)\) such cosets.  Multiplication of the
two bounds gives (SB9).
\end{proof}

The factor involving degree \(m-1\) is necessary in
this argument: it measures the failure of a cycle
modulo \(q\) to lift to an integral cycle.
Likewise, the repeated direct-sum differences in
(SB3) are essential to isolate the primitive part.
For example, take \(L=\mathbb Q(\zeta_{13})\), so
\(r=6\).  Its stable rational homology in degree three
has dimension \(6+\binom{5}{3}=16\), whereas (SB4)
bounds the rank after the two differences by \(6\).
Indeed the primitive degree-one space has dimension
\(5\), the degree-three space has dimension \(6\),
and the rational homology algebra is the graded
symmetric algebra on its primitives by the same
classical Hopf-algebra theorem.  There are no
degree-two primitives, by Borel's rank theorem.
Thus the extra ten classes are the triple products
in degree one; the factor \(\Psi_d-d^3\) removes them.
This example refutes the unfiltered rank statement,
not the statement used by the submitted proof.

No counterexample satisfying the actual hypotheses of
\(\mathrm{BOdd}\) or \(\mathrm{SUpper}\) was found.
The arguments above prove their mathematical statements
using the explicitly stated classical inputs.


\section{The finite-set bound for the weighted norms}
\label{sw-section}

We can now combine the arithmetic cycles, finite-field
functionals, homology bound, and detection of global
powers.  This proves assertion~\ref{sc-few}, and also
its version with the smaller scalar exponent anticipated
in (SC19).  The classical rational homology theorems used
in Section~\ref{sb-section} and the qualitative
Chebotarev theorem used in Section~\ref{sd-section}
remain the explicit classical inputs to this argument.

Let \(F/L\) be a cyclic extension of number fields with
generator \(\gamma\).  Suppose that \(L\) is totally
imaginary and that
\[
 \nu\ge c\ge1,\qquad c\ge2\text{ if }\ell=2,
 \qquad [F:L]=\ell^{\nu-c},
 \qquad \gamma\zeta=\zeta^{1+\ell^c},
 \tag{SW1}
\]
where \(\ell\) is prime and \(\zeta\in F\) is a
primitive \(q\)-th root of unity, with \(q=\ell^\nu\).
Put
\[
 \kappa=1+\ell^c,\quad E\ge1,\quad
 m=2E+1<\ell^c,\quad a=\kappa^E,\quad A=\mathbb Z/q.
\]
Call \(\alpha\in F^\times\) admissible if there is a
fractional divisor \(\mathfrak d\) on \(F\) such that
\[
             v_P(\alpha)=\mathfrak d(\gamma P)
                                      -a\mathfrak d(P)
                  \quad\text{for all finite primes }P.
\]
As before, define
\[
 \mathcal N(\alpha)=\prod_{j=0}^{[F:L]-1}
                              (\gamma^j\alpha)^{a^j},
 \quad
 c_E=\prod_{\substack{0\le j\le2E\\j\ne E}}
                              (\kappa^E-\kappa^j),
 \quad
 b_*=2E(c+1)+2v_\ell(E!).
 \tag{SW2}
\]

\begin{theorem}\label{sw-few}
Let \(b\ge b_*\) be an integer.  If \(T\subset F^\times\)
is a finite set of admissible elements and
\[
 \left(\frac{\mathcal N(\alpha)}{\mathcal N(\alpha')}
              \right)^{\ell^b}\in(F^\times)^q
          \quad\Longrightarrow\quad\alpha=\alpha'
                       \qquad(\alpha,\alpha'\in T),
 \tag{SW3}
\]
then
\[
 \#T\le q^{r_2(L)}
           t_L(4E+6,2E+1)t_L(4E+6,2E).
 \tag{SW4}
\]
This conclusion does not require \(2E+1<r_2(L)\).
\end{theorem}
\begin{proof}
For each \(\alpha\in T\), choose a divisor witnessing
admissibility.  Theorem~\ref{sa-class} supplies a size
\(N_\alpha\), a cycle
\[
 y_\alpha\in C_m(\operatorname{GL}_{N_\alpha}
                           (\mathcal O_L);A),
\]
and a nonzero \(b_\alpha\in\mathcal O_L\) with its
simultaneous reduction property.  The cycle and
\(b_\alpha\) are chosen before any residue field or
finite-field functional.  Set
\[
                  b_0=\ell\prod_{\alpha\in T}b_\alpha
                         \in\mathcal O_L\setminus\{0\}.
\]
Let \(\mathcal P\) be the primes \(P\) of \(F\)
such that \(P\nmid b_0\) and the natural map
\(\mathcal O_L\to k_P:=\mathcal O_F/P\) is
surjective.

For \(P\in\mathcal P\), reduction is injective on
\(\mu_q\), by Lemma~\ref{sd-roots-reduce}, since
\(P\nmid\ell\).  Thus the reduction \(\overline\zeta\)
has order \(q\), and \(q\mid\#k_P^\times\).
Choose a generator \(g_P\) of the cyclic group
\(k_P^\times\), and define
\[
      v_P:k_P^\times\longrightarrow A,
                       \qquad v_P(g_P^j)=j\pmod q.
\]
The divisibility just proved makes this well-defined.
It is surjective and
\[
                         \ker v_P=(k_P^\times)^q.
 \tag{SW5}
\]
Indeed an exponent maps to zero precisely when it is
divisible by \(q\).

Theorem~\ref{sf-functional} constructs an additive
family \(\theta_P\) with the scalar-prism value
specified by \(v_P\) and \(\overline\zeta\).
Reduce matrices over \(\mathcal O_L\) to \(k_P\)
and apply this family.  Corollary~\ref{sf-primitive-family}
gives linear maps
\[
 \Phi_{P,N}:C_m(\operatorname{GL}_N(\mathcal O_L);A)
                                     \longrightarrow A
\]
which vanish on boundaries, are stable on cycles,
and multiply their value by \(d\) on \(d\) repeated
direct sums.  The defining property of the chosen
arithmetic cycle is
\[
 \Phi_{P,N_\alpha}(y_\alpha)
     =c_E\,v_P\bigl(\overline{\mathcal N(\alpha)}\bigr)
                              \qquad(P\in\mathcal P).
 \tag{SW6}
\]
The residue on the right exists and is nonzero:
\(\mathcal N(\alpha)\) is a unit in
\(\mathcal O_F[1/b_\alpha]\) by
Theorem~\ref{sa-class}.

Take the \(A\)-module
\(B=\prod_{P\in\mathcal P}A\) and let
\(\Phi_N=(\Phi_{P,N})_{P\in\mathcal P}\).
All three properties of the family hold
coordinate by coordinate.  No finiteness assumption
on this product is needed in Theorem~\ref{sb-upper}.
Define
\[
 f:T\longrightarrow B,
      \qquad f(\alpha)=\ell^{2E}\Phi_{N_\alpha}(y_\alpha).
 \tag{SW7}
\]
That theorem bounds \(\#f(T)\) by the right-hand side
of (SW4), since \(m-1=2E\) and \(2m+4=4E+6\).
We prove that \(f\) is injective.

Suppose \(f(\alpha)=f(\alpha')\), and put
\(\beta=\mathcal N(\alpha)/\mathcal N(\alpha')\).
For every \(P\in\mathcal P\), (SW6) gives
\[
                    \ell^{2E}c_E v_P(\overline\beta)=0
                                                \quad\text{in }A.
\]
Lemma~\ref{sc-vandermonde} gives
\(v_\ell(c_E)=2Ec+2v_\ell(E!)\), and proves the
exact cancellation of the factor prime to \(\ell\).
Consequently
\[
 v_P(\overline\beta^{\,\ell^b})
       =\ell^b v_P(\overline\beta)=0
                      \qquad(P\in\mathcal P),
\]
because \(b\ge b_*\).
Equation (SW5) says that the residue of
\(\beta^{\ell^b}\) is a \(q\)-th power at every
such prime.  The element is a unit there, as noted
after (SW6).

Corollary~\ref{sd-source}, applied to
\(\beta^{\ell^b}\) and the nonzero integral element
\(b_0\), now gives
\(\beta^{\ell^b}\in(F^\times)^q\).
The separating hypothesis (SW3) implies
\(\alpha=\alpha'\).  Thus \(f\) is injective, and
the bound for \(\#f(T)\) proves (SW4).
\end{proof}

\begin{corollary}\label{sw-source-few}
Assertion~\ref{sc-few}, with its original exponent
\[
                 b=2E+(2E+1)(c+2E+1),
\]
holds.  Its strengthened version with \(b=b_*\)
also holds under the same hypotheses.
\end{corollary}
\begin{proof}
The valuation of \(E!\) is
\(\sum_{j\ge1}\lfloor E/\ell^j\rfloor\le E\),
by counting the multiples of each prime power and
bounding the resulting geometric series.
Hence
\[
 b_*\le 2E+2Ec+2E
          \le 2E+(2E+1)(c+2E+1).
\]
Apply Theorem~\ref{sw-few} with each of the two
exponents.  The source's additional hypothesis
\(2E+1<r_2(L)\) is permitted but unnecessary here.
\end{proof}

\begin{corollary}\label{sw-coarse-consequence}
With the classical reciprocity and duality results
stated earlier, the bounds (SP12) and (SP13) hold
without assuming assertion~\ref{sc-few} separately.
More explicitly, if \(L\) is totally imaginary,
\(E\ge1\), and \(\ell\) is prime, with
\(\mu_\ell\subset L\) for odd \(\ell\) and
\(\mu_4\subset L\) for \(\ell=2\), then
\[
 \begin{split}
 \log\#\mathcal S_\ell(L,E)
   &\le \exp\bigl((2E[L:\mathbb Q])^{86}
                          \log(2|\operatorname{disc}L|)\bigr),\\
 \log\#H^2(G_{L,S},\mathbb Z_\ell(E+1))
   &\le \exp\bigl((2E[L:\mathbb Q])^{88}
                          \log(2|\operatorname{disc}L|)\bigr),
 \end{split}
\]
where \(S\) is the set of primes above \(\ell\).
Both groups are finite.
\end{corollary}
\begin{proof}
Corollary~\ref{sw-source-few} supplies the remaining
hypothesis of Corollary~\ref{sp-conditional-coarse}.
Its proof combines the already proved homology torsion
estimate, class-number bound, unit comparison,
cyclotomic descent, and strict Selmer comparison.
No additional quantitative comparison is introduced
at this last step.
\end{proof}

The independent algebraic \(K\)-theory route of
Section~\ref{kt-section} gives the stronger coarse
exponent \(14\).  The purpose of the present argument
is different: it completes the conventional proof of
the source's actual finite-set comparison, while
improving its scalar exponent.  This does not by
itself prove the remaining theta and period-map
arguments or the final analytic assertion.


\section{A direct coarse bound from algebraic \texorpdfstring{$K$}{K}-theory}
\label{kt-section}

The coarse cohomology estimate has a shorter route through established
comparison theorems in algebraic \(K\)-theory.  We state these inputs
precisely, including the bound at the prime \(2\), and prove all the
resulting order estimates.  This gives an alternative to
assertion~\ref{sc-few} for the coarse estimate; that assertion itself
is proved independently in Theorem~\ref{sw-few}.

Let \(L\) be a totally imaginary number field, let
\[
 n=[L:\mathbb Q],\qquad D=|\operatorname{disc}L|,\qquad
 h=\#\operatorname{Cl}(L),
\]
and let \(p\) be a prime such that \(\mu_p\subset L\).
For \(p=2\) this last condition is automatic.
Let \(S\) be the finite primes above \(p\), let
\(R=\mathcal O_L[1/p]\), and let \(G_{L,S}\) be the Galois
group of the maximal algebraic extension unramified outside \(S\).
All \(p\)-adic cohomology groups below use continuous cohomology.
For a finite abelian group \(A\), its \(p\)-primary subgroup is
\(A\{p\}=\{a:p^r a=0\text{ for some }r\ge0\}\).

\subsection{The precise classical inputs}

Write \(\operatorname{GL}(\mathcal O_L)\) for the union of the
groups \(\operatorname{GL}_N(\mathcal O_L)\) under the usual
upper-left inclusions.  Quillen's construction gives a connected
infinite loop space
\(X=B\operatorname{GL}(\mathcal O_L)^+\).
It is the base component of the \(K\)-theory infinite
loop space, equivalently the infinite loop space of the
connected cover of its spectrum; taking this cover
preserves all positive homotopy groups.
Here the plus construction preserves integral homology, and
\[
 K_j(\mathcal O_L)=\pi_j(X)\quad(j\ge1),\qquad
 H_j(X,\mathbb Z)=
 H_j(\operatorname{GL}(\mathcal O_L),\mathbb Z).
 \tag{KT1}
\]
An infinite loop space is a pointed space equipped with
compatible deloopings in all positive degrees.  We use this
structure only to apply the Hurewicz theorem with an exponent
bound stated below.

The required results are the following.

\begin{enumerate}
\item Quillen's finite-generation theorem and Borel's rank
theorem imply that \(K_{2E}(\mathcal O_L)\) is finite for
every \(E\ge1\).
\cite[IV, Theorems 6.9 and 1.17--1.18]{weibel-kbook}
\item Homological stability gives
\[
 H_m(\operatorname{GL}_N(\mathcal O_L),\mathbb Z)
       \simeq H_m(\operatorname{GL}(\mathcal O_L),\mathbb Z)
       \quad\text{if }N\ge\max(2m+1,m+2).
 \tag{KT2}
\]
This is the stable-range form of the usual integral
stability theorem.\cite[IV, Corollary 1.14.1]{weibel-kbook}
\item The Quillen--Lichtenbaum comparison, valid also at \(2\)
for totally imaginary \(L\), identifies
\[
 K_{2E}(\mathcal O_L)\{p\}
       \simeq H^2_{\mathrm{\acute et}}
                  (\operatorname{Spec}R,\mathbb Z_p(E+1)),
                    \qquad E\ge1.
 \tag{KT3}
\]
\cite[VI, Theorem 8.2]{weibel-kbook}
For \(p\)-power torsion coefficients, the usual
\'etale--Galois comparison identifies the right-hand
cohomology with that of \(G_{L,S}\), since \(p\) is
invertible on \(\operatorname{Spec}R\).
It also identifies their compatible inverse limits for
\(\mathbb Z_p(E+1)\).
\cite[II, Proposition 2.9]{milne-adt}
\item If \(Y\) is a connected infinite loop space and \(j\ge2\),
the \(p\)-primary part of the kernel of
\[
            \pi_j(Y)\longrightarrow H_j(Y,\mathbb Z)
\]
is killed by \(p^{b_p(j)}\), where one may take
\[
 b_2(j)=\left\lceil\frac{j-1}{2}\right\rceil+3,\qquad
 b_p(j)=\left\lceil\frac{j+2}{2p-2}\right\rceil+1
                  \quad(p>2).
 \tag{KT4}
\]
This is Mathew's exponent theorem with connectivity
parameter \(m=1\).\cite[Theorem 6.3]{mathew-hurewicz}
In particular \(b_p(2E)\le E+3\) for every prime \(p\).
\end{enumerate}

These theorems are used as established inputs; their general
proofs are not reproduced here.  In (KT3), the integral
\(p\)-primary group is the one over \(\mathcal O_L\), not
over the field \(L\).  The comparison includes the
localization at \(p\) displayed on the right.
For the \(p\)-adic comparison, the finite-coefficient
\(H^1\) groups are finite.  One can see this directly:
after passing to \(L(\mu_{p^a})\), the module
\(\mathbb Z/p^a(E+1)\) is isomorphic to \(\mu_{p^a}\).
Kummer theory expresses its \(H^1\) as an extension of
the \(S\)-units modulo \(p^a\)-th powers by a subgroup
of the finite ideal class group.
The \(S\)-unit group is finitely generated: its valuation
map has kernel the ordinary units and image in
\(\mathbb Z^S\), and the ordinary unit theorem was
proved in Theorem~\ref{su-unit-theorem}.
Restriction has finite kernel, since that kernel is
cohomology of a finite group with finite coefficients.
This proves the required finiteness before base change as well.
Inverse systems of finite groups satisfy the
Mittag--Leffler condition, so the possible
\(\varprojlim^1H^1\) term vanishes and the two ordinary
inverse limits agree.

For completeness, the stable-range bound used in (KT2)
does apply to every ring of integers.  The needed ring
condition says that a unimodular triple \((a,b,c)\) can
be shortened to a unimodular pair \((a+cx,b+cy)\).
If \(c=0\), take \(x=y=0\).
Otherwise choose \(x\in\mathcal O_L\) so that \(d=a+cx\ne0\).
There are only finitely many prime ideals dividing \((d)\).
At any such prime dividing \(c\), unimodularity forces
\(b\) to be nonzero modulo that prime.
At each of the other primes, choose \(y\) modulo the prime
so that \(b+cy\ne0\).  The Chinese remainder theorem makes
all these choices simultaneous.
No prime ideal divides both \(d\) and \(b+cy\), proving
the ring condition, or stable range at most \(2\).
Thus the general bound
\(N\ge\max(2m+1,m+\operatorname{sr}(\mathcal O_L))\)
gives (KT2).

\subsection{The number of cyclic \texorpdfstring{$p$}{p}-primary factors}

We also use the classical Brauer--Hasse--Noether theorem
in the following form for the ring \(R\):
\[
 0\longrightarrow\operatorname{Br}(R)
   \longrightarrow\bigoplus_{v\in S}\mathbb Q/\mathbb Z
   \xrightarrow{\ \sum_v\ }\mathbb Q/\mathbb Z
   \longrightarrow0.
 \tag{KT5}
\]
The field is totally imaginary, so there are no real-place
terms.  Here \(\operatorname{Br}(R)=
H^2_{\mathrm{\acute et}}(\operatorname{Spec}R,\mathbb G_m)\);
the maps are the local invariants.
\cite[VI, equation (8.1.1)]{weibel-kbook}
Only the resulting elementary \(p\)-rank bound will be needed.

\begin{lemma}\label{kt-rank}
Put \(M_E=H^2(G_{L,S},\mathbb Z_p(E+1))\).
This is a finite \(p\)-group, and
\[
 \dim_{\mathbb F_p}(M_E/pM_E)
       \le |S|-1+\log_p h
       \le n+\log_p h .
 \tag{KT6}
\]
\end{lemma}
\begin{proof}
Finiteness follows from the first and third classical
inputs.  The coefficient sequence for multiplication by
\(p\) gives an injection
\[
 M_E/pM_E\ \lhook\joinrel\longrightarrow\
       H^2(G_{L,S},\mathbb F_p(E+1)).
\]
Since \(\mu_p\subset L\), all these finite twists have
trivial Galois action.  Choosing a primitive root
identifies the finite coefficient module with \(\mu_p\).
Use the \'etale--Galois comparison to compute the
last group over \(\operatorname{Spec}R\).

The Kummer sequence
\[
 1\longrightarrow\mu_p\longrightarrow\mathbb G_m
       \xrightarrow{\,(\cdot)^p\,}\mathbb G_m
       \longrightarrow1
\]
is exact on the \'etale site: a unit acquires a \(p\)-th
root in a finite \'etale algebra because \(p\) is invertible.
Its cohomology sequence, with
\(H^1(\mathbb G_m)=\operatorname{Pic}(R)\), gives
\[
 0\longrightarrow\operatorname{Pic}(R)/p
  \longrightarrow H^2_{\mathrm{\acute et}}(R,\mu_p)
  \longrightarrow\operatorname{Br}(R)[p]
  \longrightarrow0.
 \tag{KT7}
\]
Every fractional ideal of \(R\) comes from an ideal of
\(\mathcal O_L\); factorization into prime ideals shows
that localization merely kills the classes of the
primes in \(S\).  Thus \(\operatorname{Pic}(R)\) is a
quotient of \(\operatorname{Cl}(L)\), and its quotient
by \(p\) has at most \(h\) elements.

Equation (KT5) identifies \(\operatorname{Br}(R)[p]\)
with the kernel of the sum map
\(\mathbb F_p^{S}\to\mathbb F_p\).
There is at least one prime over \(p\), so its dimension
is exactly \(|S|-1\).  Add the dimensions in (KT7).
Finally \(|S|\le n\), since
\(\sum_{v\mid p}e_vf_v=n\) and each summand is positive.
This proves (KT6).
\end{proof}

\subsection{The complete order comparison}

\begin{proposition}\label{kt-additive}
For every \(E\ge1\),
\[
 \log\#H^2(G_{L,S},\mathbb Z_p(E+1))
 \le
 \log t_L(4E+6,2E)
       +(E+3)\bigl(n\log p+\log h\bigr).
 \tag{KT8}
\]
\end{proposition}
\begin{proof}
Let \(A=K_{2E}(\mathcal O_L)\{p\}\), which is finite.
The ordinary Hurewicz map and (KT1) give
\[
 f:A\longrightarrow
               H_{2E}(\operatorname{GL}(\mathcal O_L),\mathbb Z).
\]
The image of a finite group lies in the torsion subgroup.
For \(m=2E\), the choice \(N=4E+6\) satisfies (KT2),
so
\[
                  |\operatorname{im}f|
                              \le t_L(4E+6,2E).
\]
This is where finiteness of the even \(K\)-group is
needed; it is supplied explicitly above.

By (KT4), \(\ker f\) is killed by \(p^{E+3}\).
Write \(r=\dim_{\mathbb F_p}(A/pA)\).
For any finite abelian \(p\)-group the kernel and
cokernel of multiplication by \(p\) have equal orders,
so \(|A[p]|=p^r\).
For \(j\ge1\), multiplication by \(p\) maps
\(A[p^{j+1}]\) to \(A[p^j]\) and has kernel \(A[p]\).
Induction therefore gives
\[
                 |A[p^j]|\le p^{jr}.
\]
Consequently \(|\ker f|\le p^{(E+3)r}\).
By (KT3) and Lemma~\ref{kt-rank},
\(r\le n+\log_p h\).
The identity \(|A|=|\ker f|\,|\operatorname{im}f|\)
and (KT3) now prove (KT8).
\end{proof}

\begin{theorem}\label{kt-fourteen}
For every totally imaginary \(L\), every prime \(p\)
with \(\mu_p\subset L\), and every \(E\ge1\), one has
\[
 \log\#H^2(G_{L,S},\mathbb Z_p(E+1))
       \le\exp\bigl((2En)^{14}\log(2D)\bigr).
 \tag{KT9}
\]
With the stated global duality input, the strict Selmer
group with coefficient module
\(\mathbb Q_p/\mathbb Z_p(-E)\) is finite and also satisfies
\[
        \log\#\mathcal S_p(L,E)
            \le\exp\bigl((2En)^{14}\log(2D)\bigr).
 \tag{KT10}
\]
\end{theorem}
\begin{proof}
Put \(x=2En\ge4\) and \(Y=\log(2D)\ge\log2>1/2\).
The arithmetic complex and
Proposition~\ref{ht-arithmetic-numerical} give
\[
                 \log t_L(4E+6,2E)\le e^{x^{13}Y}.
\]
Corollary~\ref{sl-class-number}, together with
\(\log(1+t)\le t\) for \(t\ge0\), gives
\[
 \log h\le\tfrac12\log D
        +(n-1)\log(1+\tfrac12\log D)
            \le\tfrac n2\log D\le\tfrac n2Y.
\]
The inclusion \(\mu_p\subset L\) implies \(p-1\le n\).
For odd \(p\), this is the degree of the prime cyclotomic
field, proved by the Eisenstein calculation in
Section~\ref{sl-section}; it also holds for \(p=2\).
Thus \(\log p\le\log(n+1)\le n\).
Since \(E+3\le4E\) and \(n\le x/2\), the extra term in
(KT8) is at most
\[
 4E(n^2+nY/2)=2xn+xY
       \le x^2+xY\le3x^2Y\le x^{13}Y.
\]
Here \(Y>1/2\) and \(x\ge4\) give the penultimate
inequality, and \(x^{11}\ge3\) gives the last.
Therefore (KT8) is at most
\[
 e^{x^{13}Y}+x^{13}Y
       \le2e^{x^{13}Y}\le e^{x^{14}Y}.
\]
For the last step, use
\((x-1)x^{13}Y\ge\log2\).
This proves (KT9), with all small primes included.

The global duality statement before
Proposition~\ref{ht-h2-from-selmer} identifies
\(\mathcal S_p(L,E)^\vee\) with a subgroup of this
finite \(H^2\).  It follows that the Selmer group is
finite and has order no larger than that of \(H^2\),
proving (KT10).
\end{proof}

Theorem~\ref{kt-fourteen} proves the coarse cohomology
input using the explicitly cited classical theorems.
It does not use the source-specific finite-field
homology construction in assertion~\ref{sc-few}.
The independent proof of that construction is given in
Theorem~\ref{sw-few}.  The other arithmetic and analytic
constructions needed for the claimed Riemann hypothesis
remain distinct tasks.


\section{A shorter bound for the opposite-parity Selmer groups}
\label{ms-section}

The entire odd-prime estimate for the opposite-parity summand
admits a classical proof which avoids separating primes dividing
the degree of the field.  Apply the equivariant special-value
theorem over the maximal real subfield: its Galois group then
has order two.  We give the reduction and all numerical estimates.
The classical inputs are Greither--Popescu's cohomological
Coates--Sinnott theorem, Ferrero--Washington's vanishing of the
Iwasawa $\mu$-invariant for abelian number fields, and
Poitou--Tate duality.  These theorems are stated where used;
they are not reproved here.

\subsection{The group and its parity}

Let $E/\mathbb Q$ be an abelian CM extension, $j$ complex
conjugation, $F=E^{\langle j\rangle}$, and $n=[E:\mathbb Q]$.
Choose $N$ with $E\subset\mathbb Q(\zeta_N)$.
Necessarily $N\ge3$, $n$ is even, and $n\le\varphi(N)\le N$.
For an odd prime $p$, let $G_{E,S_p}$ be the Galois group
of the maximal extension unramified away from the places
above $p$ and the infinite places.
The Tate module $\mathbb Z_p(1)$ is the inverse limit of the
$p$-power roots of unity; integer twists are its tensor powers
and their duals.  Set
$\mathbb Q_p/\mathbb Z_p(-e)=
\mathbb Z_p(-e)\otimes_{\mathbb Z_p}\mathbb Q_p/\mathbb Z_p$.
For $e\ge1$ define
\[
 \mathcal S_p(E,e)=
 \ker\left(H^1(G_{E,S_p},\mathbb Q_p/\mathbb Z_p(-e))
       \longrightarrow
       \bigoplus_{v\mid p}H^1(E_v,\mathbb Q_p/\mathbb Z_p(-e))
       \right).
 \tag{MS1}
\]
Cohomology is continuous, and the local maps are restriction.
The action of $j$ conjugates Galois arguments and acts on
the coefficient module.  Since two is a unit in $\mathbb Z_p$,
the two eigenspaces are the images of $(1+j)/2$ and $(1-j)/2$.
We require
\[
 b_p(E,e)=\operatorname{length}_{\mathbb Z_p}
                     \mathcal S_p(E,e)^{\,(-1)^{e+1}}.
 \tag{MS2}
\]
This is the summand called \texttt{ShaMinus} in the source.
It is not always the eigenspace with eigenvalue $-1$.

\begin{theorem}\label{ms-main}
For every integer $t\ge2$ the groups in \textup{(MS2)}, for
$1\le e<t$, are finite and only finitely many are nonzero.
Moreover
\[
 \sum_{\substack{p\ {\rm prime}\\p\ne2}}\log p
       \sum_{e=1}^{t-1} b_p(E,e)
 \le n(t^2+t-2)\log(tN)
 \le \frac98 n t^2\log(tN).
 \tag{MS3}
\]
In particular this supplies the source's estimate
$3nt^2\log(tN)$, including primes dividing $[E:\mathbb Q]$,
without its additional degree-dependent error term.
\end{theorem}

\subsection{The special-value theorem over a group of order two}

For $m\ge2$ put $\epsilon=(-1)^m$ and
$e_\epsilon=(1+\epsilon j)/2$.
Let $S$ consist of the infinite places of $F$ and its
finite places above rational primes dividing $N$.
It contains all places ramified in $E/F$.  Write
\[
 H^i_{S,p}(E,m)=
 H^i_{\mathrm{et}}(\mathcal O_{E,S}[1/p],\mathbb Z_p(m)),
 \qquad
 W_p(E,m)=H^0(E,\mathbb Q_p/\mathbb Z_p(m)).
 \tag{MS4}
\]
The latter group is cyclic of finite order $w_p(E,m)$.
It has $j$-eigenvalue $\epsilon$, because the cyclotomic
character of a complex conjugation is $-1$.

Let $X(E)$ be the group of Dirichlet characters modulo $N$
factoring through $\operatorname{Gal}(E/\mathbb Q)$, and set
\[
 X_\epsilon(E)=\{\chi\in X(E):\chi(-1)=\epsilon\},
 \qquad
 R_m(E,N)=\prod_{\chi\in X_\epsilon(E)} L(1-m,\chi).
 \tag{MS5}
\]
There are $n/2$ factors.  Every character is viewed modulo
$N$, even when its primitive conductor is smaller.
Thus its $L$-function omits the Euler factors at every
prime dividing $N$.

\begin{proposition}[Classical special-value theorem,
specialized]\label{ms-special-value}
For every odd prime $p$,
\[
 \operatorname{length}_{\mathbb Z_p}
                    e_\epsilon H^2_{S,p}(E,m)
       =v_p(w_p(E,m))+v_p(R_m(E,N)).
 \tag{MS6}
\]
The group on the left is finite, and $R_m(E,N)$ is a
nonzero rational number.
\end{proposition}
\begin{proof}
Here is the precise classical input.
For an abelian extension $K/k$ with $k$ totally real,
$p>2$, $m\ge2$, and $S$ containing the infinite and
ramified places, Greither--Popescu
\cite[Theorem 6.11]{greither-popescu} prove
\[
 \begin{split}
 &\operatorname{Ann}_{\mathbb Z_p[G]}
     \bigl(H^1_{\mathrm{et}}(\mathcal O_{K,S}[1/p],
                             \mathbb Z_p(m))_{\mathrm{tors}}\bigr)
                       \Theta_{S,K/k}(1-m)\\
 &\hspace{12mm}=
 e_m(K/k)\operatorname{Fitt}_{\mathbb Z_p[G]}
             H^2_{\mathrm{et}}(\mathcal O_{K,S}[1/p],\mathbb Z_p(m)).
 \end{split}
 \tag{MS7}
\]
The hypothesis is the vanishing of the cyclotomic
$\mu$-invariant for the maximal CM subfield of $K(\mu_p)$.
For $K=E$ that field is abelian over $\mathbb Q$, so
Ferrero--Washington \cite{ct-fw} supplies it.
The idempotent $e_m(K/k)$ is the product, over real places
of $k$, of $(1+(-1)^m j_v)/2$.  For $E/F$ every $j_v$
is $j$, so the product is $e_\epsilon$.
Finiteness of $H^2$ is part of the standard cohomological
setting of this theorem
\cite[\S6.2, preceding Lemma 6.9]{greither-popescu}.

We spell out its specialization.  The maps $j\mapsto1$
and $j\mapsto-1$ give
\[
 \mathbb Z_p[\langle j\rangle]\simeq
                       \mathbb Z_p\times\mathbb Z_p,
 \qquad
 e_\epsilon\mathbb Z_p[\langle j\rangle]\simeq\mathbb Z_p.
 \tag{MS8}
\]
The inverse uses division only by two.  There is no
division by $n$ and no exclusion of $p\mid n$.
The torsion of $H^1_{S,p}(E,m)$ is $W_p(E,m)$
\cite[Lemma 6.9(1)]{greither-popescu}; this also follows
from the exact sequence for
$\mathbb Z_p(m)\subset\mathbb Q_p(m)$, since
$H^0(E,\mathbb Q_p(m))=0$ and $H^1(E,\mathbb Q_p(m))$
is a vector space.
In the $\epsilon$ factor its annihilator is
$(p^{v_p(w_p(E,m))})$.
For a finite $\mathbb Z_p$-module
$M=\bigoplus_i\mathbb Z_p/p^{a_i}$, the diagonal presentation
gives $\operatorname{Fitt}(M)=(p^{\sum_i a_i})$.
The exponent is its length.

The $\epsilon$ component of $\Theta_{S,E/F}(1-m)$ is
$\zeta_{F,S}(1-m)$ for $\epsilon=1$, and
$L_S(1-m,E/F)$ for $\epsilon=-1$.
Artin factorization identifies these with (MS5).
Indeed, induction of the trivial character of
$\langle j\rangle$ to $\operatorname{Gal}(E/\mathbb Q)$
is the sum of the characters with $\chi(j)=1$;
induction of its nontrivial character is the sum with
$\chi(j)=-1$.  Equality of Euler factors proves the
factorization, including the removal of all primes above
the rational divisors of $N$.
Projecting (MS7) to (MS8) proves (MS6).

The character set in (MS5) is invariant under Galois
automorphisms of the field of character values, so the
product is rational.  Every factor is nonzero.
Its primitive character has parity $(-1)^m$; the
functional equation expresses its value at $1-m$ as a
nonzero constant times $L(m,\bar\chi)$.
The absolutely convergent Euler product makes the latter
nonzero.  Each removed Euler factor is
$1-\chi^*(\ell)\ell^{m-1}\ne0$, since $m\ge2$.
\end{proof}

\subsection{Localization and the strict Selmer group}

\begin{lemma}\label{ms-localization}
For $m=e+1\ge2$ and odd $p$,
\[
 b_p(E,e)\le
 \operatorname{length}_{\mathbb Z_p}
                       e_{(-1)^m}H^2_{S,p}(E,m).
 \tag{MS9}
\]
\end{lemma}
\begin{proof}
Apply finite-coefficient Poitou--Tate duality
\cite[Chapter I, Theorem 4.10(a)]{milne-adt}
to $A_a=\mathbb Z/p^a(1-m)$ and its Cartier dual
$T_a=\mathbb Z/p^a(m)$.  It pairs their local restriction
kernels in degrees one and two perfectly.
The pairings respect the inclusions $A_a\longrightarrow A_{a+1}$
and the quotient maps $T_{a+1}\longrightarrow T_a$.
Continuous cohomology with discrete coefficients commutes
with filtered direct limits.  Since there are finitely many
local places, exactness of that limit identifies the direct
limit of the degree-one kernels with (MS1).
All finite-coefficient cohomology groups in question are
finite.  Their inverse systems have the Mittag--Leffler
property (descending sequences of images in a fixed finite
group stabilize); hence their first derived inverse limits
vanish.  Consequently the inverse limit of the degree-two
kernels is the kernel for
$H^2(G_{E,S_p},\mathbb Z_p(m))$.
Pontryagin duality exchanges this inverse limit and the
direct limit just described, giving the required perfect
pairing.  These pairings are equivariant for automorphisms of $E$:
cup products and restriction are natural, and an automorphism
permutes the local invariant maps without changing their sum.
The value group $\mathbb Q_p/\mathbb Z_p$ has trivial
$j$-action.  The eigenspaces of the same sign therefore
pair perfectly; opposite signs pair to zero, since two
is invertible.  Hence $b_p(E,e)$ is the length of a subgroup
of the $(-1)^m$ eigenspace of
$H^2(G_{E,S_p},\mathbb Z_p(m))$.

This is
$H^2_{\mathrm{et}}(\mathcal O_E[1/p],\mathbb Z_p(m))$.
Localization injects it into $H^2_{S,p}(E,m)$.
The potentially preceding term for a removed prime
$v\nmid p$ is
$H^0(\kappa(v),\mathbb Z_p(m-1))$.
Frobenius acts by $(Nv)^{m-1}\ne1$, and multiplication
by $(Nv)^{m-1}-1$ is injective on $\mathbb Z_p$.
That preceding group is zero.
Adding the finitely many primes gives the injection.
It respects $j$, as the set of removed primes is stable.
Projection by $e_{(-1)^m}$ preserves it, proving (MS9)
and finiteness of its left side.
\end{proof}

\subsection{Elementary special-value estimates}

Define the Bernoulli polynomials by
\[
 \frac{z e^{xz}}{e^z-1}
       =\sum_{a\ge0}B_a(x)\frac{z^a}{a!},
 \qquad B_a=B_a(0),
 \tag{MS10}
\]
and put $D_m=m(m+1)!N$ for $m\ge2$.

\begin{lemma}\label{ms-bernoulli}
For every character $\chi$ modulo $N$,
\[
 L(1-m,\chi)=
       -\frac{N^{m-1}}m\sum_{a=1}^N\chi(a)B_m(a/N).
 \tag{MS11}
\]
The number $D_mL(1-m,\chi)$ is an algebraic integer, and
\[
         D_m\le(mN)^m,\qquad |L(1-m,\chi)|\le(mN)^m.
 \tag{MS12}
\]
\end{lemma}
\begin{proof}
For $z>0$, periodicity gives
\[
 \sum_{\nu\ge1}\chi(\nu)e^{-\nu z}
       =\sum_{a=1}^N\chi(a)\frac{e^{-az}}{1-e^{-Nz}}.
\]
Its Mellin integral is $\Gamma(s)L(s,\chi)$ for
$\operatorname{Re}s>1$.
Split that integral at one and subtract sufficiently many
terms of the power series at zero.
The tail at infinity is entire, and a subtracted term
$c z^h$ contributes $c/(s+h)$.
The coefficient of $z^{m-1}$ in the exponential sum is
$(-1)^mN^{m-1}\sum_a\chi(a)B_m(a/N)/m!$, by (MS10).
The residue of $\Gamma(s)$ at $s=1-m$ is
$(-1)^{m-1}/(m-1)!$.
Dividing residues proves (MS11); the same subtraction
supplies the analytic continuation used here.

A denominator bound follows without a theorem about
Bernoulli denominators.  The formal identities
\[
 \frac z{e^z-1}=\sum_{b\ge0}\frac{(-1)^b}{b+1}(e^z-1)^b,
 \qquad
 (e^z-1)^b=b!\sum_{a\ge b}
    \left\{\begin{smallmatrix}a\\b\end{smallmatrix}\right\}
    \frac{z^a}{a!}
\]
give
\[
 B_a=\sum_{b=0}^a\frac{(-1)^b b!}{b+1}
                \left\{\begin{smallmatrix}a\\b\end{smallmatrix}\right\}.
 \tag{MS13}
\]
The braces count partitions into nonempty subsets and
are integers.  For $a\le m$, $(m+1)!$ clears all
denominators in (MS13).
Also $B_m(x)=\sum_{a=0}^m\binom ma B_a x^{m-a}$.
Thus $N^m(m+1)!B_m(b/N)$ is an integer for integer $b$.
Multiplying (MS11) by $D_m$ gives an integral linear
combination of roots of unity and zeros.

Repeated integration by parts in Fourier coefficients,
using $B_a'=aB_{a-1}$ and
$B_a(1)-B_a(0)=0$ for $a\ne1$, gives
\[
 B_m(x)=-\frac{m!}{(2\pi i)^m}
          \sum_{h\ne0}\frac{e^{2\pi ihx}}{h^m},
                  \quad 0\le x\le1,\quad m\ge2.
 \tag{MS14}
\]
The constant coefficient is zero by (MS10).
The series converges absolutely uniformly and has the
same Fourier coefficients as the continuous periodic
function $B_m$, proving equality, including endpoints.
Hence
$|B_m(x)|\le2m!\zeta(m)/(2\pi)^m\le m!$,
using $\zeta(m)\le2$ and $2\pi>4$.
Equation (MS11) gives
$|L(1-m,\chi)|\le N^m(m-1)!\le(mN)^m$.
Finally $m!\le m^{m-1}$ and
$m+1\le3^{m-1}\le N^{m-1}$, so
$m(m+1)!N\le m^m(m+1)N\le(mN)^m$.
\end{proof}

\begin{lemma}\label{ms-roots}
The integer $w_{\mathrm{odd}}(E,m)=\prod_{p\ne2}w_p(E,m)$
is finite and
\[
             \log w_{\mathrm{odd}}(E,m)\le nm\log(mN).
 \tag{MS15}
\]
\end{lemma}
\begin{proof}
If $p^a\mid w_p(E,m)$, the subgroup
$\operatorname{Gal}(E(\zeta_{p^a})/E)$ of
$(\mathbb Z/p^a)^\times$ has exponent dividing $m$.
For odd $p$ it is cyclic, hence has order at most $m$.
Its index is
$[E\cap\mathbb Q(\zeta_{p^a}):\mathbb Q]\le n$.
Therefore $p^{a-1}(p-1)\le nm$ and
$p^a\le3nm/2$.
Only primes $3\le p\le nm+1$ occur.
Since $nm$ is even, there are at most $nm/2$ odd integers
in that interval, whence
\[
 \log w_{\mathrm{odd}}(E,m)
       \le\frac{nm}{2}\log(3nm/2)\le nm\log(mN).
\]
The last inequality uses $n\le N$ and
$3nm/2\le(mN)^2$, since $mN\ge6$.
The prime-power bound also proves finiteness.
\end{proof}

\begin{proof}[Proof of Theorem~\ref{ms-main}]
For fixed $m=e+1$, put $r=n/2$ and $R=R_m(E,N)$.
The number $D_m^rR$ is a nonzero algebraic integer by
Lemma~\ref{ms-bernoulli}, and is rational by
Proposition~\ref{ms-special-value}.  It is therefore a
nonzero rational integer.  So is
$I=w_{\mathrm{odd}}(E,m)D_m^rR$.
Equations (MS6) and (MS9) imply, for odd $p$,
\[
 0\le b_p(E,m-1)\le
       v_p(w_p(E,m))+v_p(R)\le v_p(I).
 \tag{MS16}
\]
The factors $w_q$ for $q\ne p$ are $p$-adic units,
and $D_m^r$ only increases the valuation.
Only primes dividing $I$ can have a nonzero group.
Summing and applying (MS12) and (MS15) yields
\[
 \begin{split}
 \sum_{p\ne2}b_p(E,m-1)\log p
 &\le\log|I|\\
 &\le nm\log(mN)+\frac n2 m\log(mN)+\frac n2 m\log(mN)\\
 &=2nm\log(mN).
 \end{split}
 \tag{MS17}
\]
Using the actual integer $I$ handles possible denominators
of $R$ at two; bounding odd valuations by $\log|R|$ alone
would not be justified.
Sum over $2\le m\le t$ and use
$\log(mN)\le\log(tN)$.
The equality $2\sum_{m=2}^t m=t^2+t-2$ gives the first
inequality in (MS3).
The second is equivalent to $(t-4)^2\ge0$.
\end{proof}

\paragraph{Relation to the submitted proof.}
Taking $N$ to be the conductor proves the required
opposite-parity estimate for every abelian CM field, and
therefore for the composita in the source.
In its notation this supplies \texttt{ThmBcFree} with
the additional constant equal to zero.
This is a separate conventional proof using the named
classical theorems.  It does not certify the source's
individual characterwise annihilator argument.
It does not estimate the complementary eigenspace
$(-1)^e$, whose relevant special value has a zero.


\section{The trivial and quadratic residual characters in the complementary parity}
\label{sz-selmer-section}

The complementary-parity estimate in
Section~\ref{sec:selection-assembly} requires several different
arithmetic arguments. We first prove, without an automorphy input,
the exact vanishing assertion used for its trivial Dirichlet
character. This is the case called $Z$ in the submitted
construction.\cite{alpo-uniform-source} Its proof is shorter than
passing to the full cyclotomic tower: the necessary lifting of
local coboundaries can be done over $\mathbb Q$ itself.

Let $p$ be an odd prime, let $L/\mathbb Q_p$ be finite, and
write $\mathcal O$ for its ring of integers, $\varpi$ for a
uniformizer and $\mathbf k=\mathcal O/\varpi\mathcal O$ for
its residue field. Write $\varepsilon_p:G_{\mathbb Q}\to
\mathbb Z_p^\times$ for the cyclotomic character. For $e\ge1$
put
\[
 A_e=(L/\mathcal O)(-e),\qquad
       g\cdot a=\varepsilon_p(g)^{-e}a,\qquad
 B_e=A_e[\varpi]\simeq\mathbf k(\bar\varepsilon_p^{-e}).
                                                               \tag{SZ1}
\]
These are discrete modules. Fix decomposition groups $D_q$
and inertia groups $I_q$ at the rational primes $q$. Define
\[
 \mathcal T_p(e)=\ker\left[
 H^1(G_{\mathbb Q},A_e)\longrightarrow
 H^1(D_p,A_e)\oplus\prod_{q\ne p}H^1(I_q,A_e)\right].   \tag{SZ2}
\]
This is the strict Selmer group with local condition zero
at $p$ and unramified conditions elsewhere. Adding zero
local conditions at further primes only replaces it by a
subgroup.

\subsection{Lifting a torsion class with its local condition}

We use the following standard form of unramified class-field
reciprocity: continuous homomorphisms from $G_F$ to a finite
abelian group of odd order that vanish on all finite inertia
groups factor through the ordinary ideal class group of $F$.
At a real place the inertia group has order at most two and
therefore has zero image. This is the Hilbert class field
form of the global reciprocity theorem.\cite[Chapter V,
\S3]{milne-cft}

\begin{lemma}\label{sz-torsion-and-local}
Every element of $H^1(G_{\mathbb Q},A_e)$ is killed by a
power of $\varpi$. The image of the cyclotomic character
on $D_p$ equals its image on $G_{\mathbb Q}$, namely
$\mathbb Z_p^\times$. Consequently
\[
                         A_e^{G_{\mathbb Q}}=A_e^{D_p}.
                                                               \tag{SZ3}
\]
Every class in $\mathcal T_p(e)[\varpi]$
is the image of a class $c\in H^1(G_{\mathbb Q},B_e)$
whose restrictions to $D_p$ and all $I_q$, $q\ne p$,
are zero.
\end{lemma}
\begin{proof}
A continuous cocycle from a profinite group to the discrete
module $A_e$ has finite image. Choose a common power of
$\varpi$ killing those finitely many values; it kills the
cohomology class as well.

For $r\ge1$, the polynomial $\Phi_{p^r}(1+X)$ is Eisenstein
at $p$. Its constant term is $p$, and its reduction modulo
$p$ is $X^{p^{r-1}(p-1)}$. Hence
\[
 [\mathbb Q_p(\mu_{p^r}):\mathbb Q_p]=p^{r-1}(p-1).
\]
Its Galois group therefore gives the full group
$(\mathbb Z/p^r\mathbb Z)^\times$ under the cyclotomic
character. Taking inverse limits gives the local assertion;
the global image contains the local one. Formula (SZ3)
follows because the action on $A_e$ factors through that
character.

The exact sequence
\[
 0\longrightarrow B_e\longrightarrow A_e
          \xrightarrow{\,\varpi\,}A_e\longrightarrow0    \tag{SZ4}
\]
shows that a class $x\in H^1(G_{\mathbb Q},A_e)[\varpi]$
lifts to $c\in H^1(G_{\mathbb Q},B_e)$. If $x\in
\mathcal T_p(e)$, the restriction $c|_{D_p}$ belongs to
the image of the connecting map
\[
 A_e^{D_p}/\varpi A_e^{D_p}\longrightarrow H^1(D_p,B_e).
\]
By (SZ3) choose an element of
$A_e^{G_{\mathbb Q}}/\varpi A_e^{G_{\mathbb Q}}$ with
this same image. Subtract its global connecting class from
$c$. Its image in $H^1(G_{\mathbb Q},A_e)$ is unchanged,
and the new $c$ has zero restriction to $D_p$.

For $q\ne p$ the inertia group acts trivially on $A_e$,
since the cyclotomic character is unramified there.
Multiplication by $\varpi$ is surjective on $A_e$, so
(SZ4) shows that
\[
                     H^1(I_q,B_e)\longrightarrow H^1(I_q,A_e)
\]
is injective. The image of $c$ on the right is zero by
(SZ2), so its restriction on the left is zero as well.
The stated lifting conclusion follows.
\end{proof}

\begin{proposition}\label{sz-residual-class}
Let $F_e$ be the fixed field in $\mathbb Q(\mu_p)$ of the
kernel of $\bar\varepsilon_p^{-e}$. There is an injection
\[
       \mathcal T_p(e)[\varpi]\lhook\joinrel\longrightarrow
                  \operatorname{Hom}(\operatorname{Cl}(F_e),B_e),
                                                               \tag{SZ5}
\]
where the target denotes homomorphisms of additive abelian
groups, and $B_e$ is trivial as a $G_{F_e}$-module.
In particular $p\nmid h_{F_e}$ implies $\mathcal T_p(e)=0$.
Moreover
\[
 \dim_{\mathbf k}\mathcal T_p(e)[\varpi]
       \le\dim_{\mathbb F_p}\operatorname{Cl}(F_e)/p.
                                                               \tag{SZ6}
\]
\end{proposition}
\begin{proof}
The final part of Lemma~\ref{sz-torsion-and-local} constructs
a lift with zero restriction at $D_p$ for every $e$.
It is unique: two lifts differ by a global connecting class
from $A_e^{G_{\mathbb Q}}/\varpi A_e^{G_{\mathbb Q}}$.
The local connecting map is injective on
$A_e^{D_p}/\varpi A_e^{D_p}$ by exactness, and (SZ3)
identifies its source with the global one. A connecting
class whose local restriction is zero is therefore zero.
Thus this lift depends $\mathcal O$-linearly on the original
class.

Restrict to $G_{F_e}$. This is injective on
$H^1(G_{\mathbb Q},B_e)$: the inflation kernel is
$H^1(\operatorname{Gal}(F_e/\mathbb Q),B_e)$, which
vanishes by averaging because its order divides $p-1$.
On $G_{F_e}$ the cocycle is a homomorphism. It is
unramified everywhere by the local conditions on its lift,
so class-field reciprocity gives (SZ5). A finite abelian
group with $p$-rank $r$ has a space of homomorphisms to
$\mathbf k$ of $\mathbf k$-dimension $r$, proving (SZ6).
If $p\nmid h_{F_e}$, (SZ5) is zero; the torsion argument
used below then annihilates the full group.
\end{proof}

\begin{theorem}\label{sz-trivial}
If $p-1\mid e$, then $\mathcal T_p(e)=0$.
\end{theorem}
\begin{proof}
The residual action on $B_e$ is trivial. A class furnished
by Lemma~\ref{sz-torsion-and-local} is therefore a continuous
homomorphism $G_{\mathbb Q}\to(B_e,+)$ which is unramified
at every finite prime. Its target is a finite $p$-group.
Class-field reciprocity makes it factor through
$\operatorname{Cl}(\mathbb Q)=0$. Thus
$\mathcal T_p(e)[\varpi]=0$. A nonzero element of a
$\varpi$-primary torsion module has a nonzero multiple
killed by $\varpi$; the first assertion of the lemma
therefore proves $\mathcal T_p(e)=0$.
\end{proof}

\subsection{The quadratic residual character}

\begin{lemma}\label{sz-small-real-class}
If $p\equiv1\pmod4$ is prime and $F=\mathbb Q(\sqrt p)$,
then $h_F<p$.
\end{lemma}
\begin{proof}
The discriminant is $p$. The small-ideal and ideal-count
bounds in Section~\ref{sl-section}, with $[F:\mathbb Q]=2$,
give
\[
                         h_F\le\sqrt p\,(1+\log\sqrt p).
\]
For $u>1$, integrating $1/t<1$ from $1$ to $u$ gives
$\log u<u-1$. Taking $u=\sqrt p$ makes the right side
strictly less than $p$.
\end{proof}

\begin{theorem}\label{sz-quadratic}
Suppose $e$ is even, $p-1\mid2e$, and $p-1\nmid e$.
Then $\mathcal T_p(e)=0$.
\end{theorem}
\begin{proof}
Write $p-1=2a$. The divisibilities say $e=ab$ with $b$
odd. Since $e$ is even, $a$ is even and $p\equiv1\pmod4$;
in particular $p\ge5$. The residual character on $B_e$
is the quadratic character
$\bar\varepsilon_p^{(p-1)/2}$. Its fixed field is
$F=\mathbb Q(\sqrt p)$. One explicit identification is the
quadratic Gauss sum
\[
 \tau=\sum_{u\bmod p}(u/p)\zeta_p^u,\qquad
 \sigma_v\tau=(v/p)\tau,\qquad \tau^2=p.
\]
The last identity follows from
$\tau\bar\tau=p$ and $(-1/p)=1$, as in the quadratic
field calculation of Section~\ref{sec:selection-assembly}.

There is an element of $D_p$ on which this residual
character is $-1$. Its action minus one on $A_e$ is
multiplication by a unit of $\mathcal O$. Hence
$A_e^{D_p}=0$. Given $x\in\mathcal T_p(e)[\varpi]$,
choose its lift $c\in H^1(G_{\mathbb Q},B_e)$ using
(SZ4). The local connecting term at $p$ is now zero,
so $c|_{D_p}=0$. As in the preceding lemma, its
restrictions to every $I_q$, $q\ne p$, are zero.

Restriction to $G_F$ makes $B_e$ a trivial module.
Thus $c|_{G_F}$ is an everywhere unramified homomorphism
from $G_F$ to the additive group of $\mathbf k$.
At primes above $p$ it is zero even on the decomposition
groups; at all other finite primes inertia is already
killed. Class-field reciprocity makes it factor through
$\operatorname{Cl}(F)$. Lemma~\ref{sz-small-real-class}
shows $p\nmid h_F$, so this homomorphism is zero.

The restriction map
$H^1(G_{\mathbb Q},B_e)\to H^1(G_F,B_e)$ is injective.
Indeed its kernel is the image of
$H^1(\operatorname{Gal}(F/\mathbb Q),B_e)$ by the
inflation--restriction sequence, and this latter group
vanishes by averaging over its group of order two; two
is invertible on $B_e$. Thus $c=0$ and $x=0$.
The torsion argument in Theorem~\ref{sz-trivial} finishes
the proof.
\end{proof}

\begin{corollary}\label{sz-exact-case}
For every prime $p\ge5$, every even $e\ge2$ such that
$p-1\mid2e$, and every finite extension $L/\mathbb Q_p$,
one has $\mathcal T_p(e)=0$. The same is true if additional
zero local conditions are imposed at any finite set of
primes. In particular the length for the trivial Dirichlet
character in case $Z$ of the submitted complementary-parity
argument is exactly zero.
\end{corollary}
\begin{proof}
Apply Theorem~\ref{sz-trivial} when $p-1\mid e$ and
Theorem~\ref{sz-quadratic} otherwise. An additional local
condition cuts out a subgroup, so it preserves vanishing.
The source's coefficient action for the trivial character
is exactly $\varepsilon_p^{-e}$, and its level-$N$ group
adds zero conditions at the primes dividing $N$.
Thus its group is a subgroup of (SZ2). Its parity condition
at the trivial character is precisely that $e$ be even.
\end{proof}

This proves one complete arithmetic input, including
arbitrarily ramified coefficient extensions and all the
relevant twists. It does not identify the groups of
nontrivial Dirichlet characters with these trivial-character
groups. The complementary-parity estimate for the other
characters still requires the remaining counting and
restriction arguments, with constants independent of the
selected auxiliary primes.


\section{Restriction for the fixed-conductor characters}
\label{iv-section}

We next prove the arithmetic assertion called case $IV$ in the
submitted construction.\cite{alpo-uniform-source} The main point is
that its character has conductor dividing the original fixed modulus.
Consequently it is trivial over one of the fixed odd-prime fields
already controlled by the common cohomological interval of
Section~\ref{sec:selection-assembly}. We keep the resulting bound
additive in that interval's bound $B$.

\subsection{The character loses its varying conductor}

Fix a primitive Dirichlet character $\chi$ of conductor $M$ and
order $o$, and retain the fields
\[
 E_h=K_\chi R_h(\sqrt{-q},\sqrt{-d}),\qquad
 N_h=\operatorname{lcm}(2^{h+2},M)qd
                                                               \tag{IV1}
\]
from Section~\ref{sec:selection-assembly}. Thus $q,d$ are distinct
odd primes prime to $M$, and $R_h$ is the real subfield of the
$2^{h+2}$-th cyclotomic field. A character $\theta$ of
$\operatorname{Gal}(E_h/\mathbb Q)$ will also denote its
Dirichlet character modulo $N_h$.

For a prime $r$, write $\theta_r$ for the component on
$(\mathbb Z/r^{v_r(N_h)}\mathbb Z)^\times$ under the Chinese
remainder decomposition; it is understood to be the trivial
character when $r\nmid N_h$. Independently, for an odd prime
$p$, write
\[
                    \theta=\kappa\vartheta,
\]
where $\kappa$ has $p$-power order and $\vartheta$ has
order prime to $p$. This factorization is unique in the finite
cyclic group generated by $\theta$. Define
\[
                              \eta=\vartheta/\vartheta_p.
                                                               \tag{IV2}
\]
The denominator in (IV2) is the local component at the rational
prime $p$; it is distinct from the factor of $p$-power order.

\begin{lemma}\label{iv-fixed-conductor}
Suppose $p\ne q,d$ and $\eta$ has odd order. Then
\[
                              \operatorname{cond}\theta\mid M.
                                                               \tag{IV3}
\]
Indeed the local components of $\theta$ at $2,q,d$ are trivial.
\end{lemma}
\begin{proof}
The group of units modulo a power of $2$ is a $2$-group, so
$\theta_2$ has $2$-power order. Its $p$-power order factor is
trivial because $p$ is odd. Removing the local component at $p$
does not affect it. Hence $\eta_2=\theta_2$. Its order divides
both a power of $2$ and the odd order of $\eta$, and is therefore
one.

The inertia groups at $q$ and $d$ in $E_h/\mathbb Q$ have order
at most two. This follows directly from (IV1): the other fields
are unramified there, and the respective added quadratic field
has degree two. Equivalently, the two local Dirichlet components
are trivial or quadratic. Their $p$-power order parts are again
trivial. Since $p\ne q,d$, they are unchanged in (IV2), so their
orders are odd and divide two. They too are trivial.

At every remaining odd prime, the exponent in the modulus $N_h$
is precisely the exponent in $M$. Triviality at $2,q,d$ now
shows that the primitive conductor divides $M$.
\end{proof}

For the selected fixed set of odd primes $P$, the inequalities
$q,d>\max P$ ensure $p\ne q,d$ for every $p\in P$.
Thus the hypothesis of Lemma~\ref{iv-fixed-conductor} is exactly
the odd-order alternative required in case $IV$. No bound for
$q,d$ enters the conclusion.

\subsection{Restriction and the loss in degree one}

For a number field $F$ and $e\ge1$, let
\[
 \mathcal S_p(F,e)=\ker\left[
 H^1(G_F,\mathbb Q_p/\mathbb Z_p(-e))\longrightarrow
 \prod_{v\mid p}H^1(D_v,\mathbb Q_p/\mathbb Z_p(-e))
 \oplus\prod_{v\nmid p}H^1(I_v,\mathbb Q_p/\mathbb Z_p(-e))
 \right].                                                     \tag{IV4}
\]
All primes in (IV4) are finite. Since $p$ is odd, the corresponding
cohomology at a real place vanishes by averaging over its group
of order two. Put
$\lambda_p(F,e)=\operatorname{length}_{\mathbb Z_p}\mathcal S_p(F,e)$,
with value $+\infty$ when this length is infinite.

Let $K/\mathbb Q_p$ be finite, with ring of integers $\mathcal O$
and ramification index $b$. Suppose that the values of
$\theta$ belong to $\mathcal O^\times$. Put
\[
 A(e,\theta)=(K/\mathcal O)(\theta^{-1}\varepsilon_p^{-e}).
                                                               \tag{IV5}
\]
Let $\mathcal T_p(e,\theta)$ be the strict Selmer group defined
as in (SZ2), with this coefficient module, and write
\[
 s_p(e,\theta)=b^{-1}
       \operatorname{length}_{\mathcal O}\mathcal T_p(e,\theta)
                                                               \tag{IV6}
\]
when the length is finite. Additional zero local conditions at
primes dividing $N_h$ only decrease this length.

\begin{lemma}\label{iv-restriction}
Suppose $\operatorname{cond}\theta\mid M$. For an integer
$j\ge0$, set
\[
 F=\mathbb Q(\mu_{p^{j+1}M}),\quad
 a=j+1+v_p(M),\quad n=[F:\mathbb Q].
                                                               \tag{IV7}
\]
If $\mathcal S_p(F,e)$ is finite, then $\mathcal T_p(e,\theta)$
is finite and
\[
             s_p(e,\theta)\le\lambda_p(F,e)+n(a+v_p(e)).
                                                               \tag{IV8}
\]
In particular this statement permits $p\mid n$ and arbitrary
ramification in $K/\mathbb Q_p$.
\end{lemma}
\begin{proof}
The character $\theta$ is trivial on $G_F$. Let
$\Delta=\operatorname{Gal}(F/\mathbb Q)$ and
$A=A(e,\theta)$. Inflation--restriction gives an exact sequence
\[
 0\longrightarrow H^1(\Delta,A^{G_F})
 \longrightarrow H^1(G_{\mathbb Q},A)
 \longrightarrow H^1(G_F,A).
                                                               \tag{IV9}
\]
No division by $|\Delta|$ is used. A class in
$\mathcal T_p(e,\theta)$ restricts into the strict Selmer group
over $F$: a decomposition or inertia group over $F$ is a
subgroup of the corresponding conjugate group over $\mathbb Q$.
The local restriction of a coboundary is still a coboundary.

On $G_F$ the module $A$ is $(K/\mathcal O)(-e)$. Its strict
Selmer group is
\[
                       \mathcal S_p(F,e)\otimes_{\mathbb Z_p}
                                      \mathcal O.              \tag{IV10}
\]
To see this directly, choose a $\mathbb Z_p$-basis of the finite
free module $\mathcal O$. Since
$K/\mathcal O=(\mathbb Q_p/\mathbb Z_p)\otimes\mathcal O$,
the cochain complex, all local cochain complexes, and their
restriction kernels become the direct sums indexed by that
basis. The resulting identification is the canonical
$\mathcal O$-linear map in (IV10). A cyclic summand
$\mathbb Z_p/p^r\mathbb Z_p$ contributes
$\mathcal O/p^r\mathcal O$, of $\mathcal O$-length $br$.
The target of restriction therefore has length
$b\lambda_p(F,e)$.

It remains to bound the left group in (IV9). The cyclotomic
image of $G_F$ is exactly $1+p^a\mathbb Z_p$. Indeed the
Chinese remainder theorem gives
\[
 F\cap\mathbb Q(\mu_{p^\infty})=\mathbb Q(\mu_{p^a}).
\]
For odd $p$ and $a\ge1$,
\[
                   v_p((1+p^a)^e-1)=a+v_p(e).                 \tag{IV11}
\]
For completeness, when $p\nmid e$ the first nonconstant term
of the binomial expansion has valuation $a$ and every other
term has larger valuation. If $v_p(z)\ge1$, the term $pz$
in $(1+z)^p-1$ has valuation $1+v_p(z)$, strictly below
all the other terms, because $p$ is odd. Iterating this
observation proves (IV11) for every $e\ge1$.
The same formula holds with $e$ replaced by $-e$, since
$(1+p^a)^e$ is a unit. Every element of
$1+p^a\mathbb Z_p$ has its $e$-th power congruent to one
modulo $p^{a+v_p(e)}$, while $1+p^a$ attains this valuation.
It follows that
\[
 A^{G_F}=A[p^{a+v_p(e)}],\qquad
 \operatorname{length}_{\mathcal O}A^{G_F}=b(a+v_p(e)).
                                                               \tag{IV12}
\]
Every $1$-cocycle of the finite group $\Delta$ is a function
$\Delta\to A^{G_F}$. Thus its cocycle module is a submodule
of $(A^{G_F})^n$, and taking the quotient by coboundaries
cannot increase length. Consequently
\[
 \operatorname{length}_{\mathcal O}H^1(\Delta,A^{G_F})
                                  \le nb(a+v_p(e)).
\]
Combining this with (IV9)--(IV10) proves finiteness and (IV8).
\end{proof}

\subsection{Summation with fixed constants}

\begin{theorem}\label{iv-summed}
Fix $p,M,j$, and let $F,a,n$ be as in (IV7). If $t\ge2$
and
\[
                 \sum_{1\le e<t}\lambda_p(F,e)\le B,
                                                               \tag{IV13}
\]
then for every character $\theta$ of conductor dividing $M$
all the corresponding strict Selmer groups are finite, and
\[
 \sum_{1\le e<t}s_p(e,\theta)
 \le B+n\left(a+\frac1{p-1}\right)(t-1).
                                                               \tag{IV14}
\]
There is an integer $C$, depending only on $p,M,j$, such that
for every $N\ge1$ the right side is at most $B+(Nt)^C$.
This bound remains true after adding the level-$N$ zero
local conditions.
\end{theorem}
\begin{proof}
Sum (IV8). Counting the multiples of each power of $p$ gives
\[
 \sum_{e=1}^{t-1}v_p(e)
   =\sum_{r\ge1}\left\lfloor\frac{t-1}{p^r}\right\rfloor
   \le\frac{t-1}{p-1}.
\]
This proves (IV14). Choose $C\ge1$ so that
$2^{C-1}\ge n(a+1)$. Since $t\ge2$ and $N\ge1$,
$n(a+1)t\le t^C\le(Nt)^C$.
The last assertion follows from inclusion of Selmer groups.
\end{proof}

\begin{corollary}\label{iv-exact-case}
Fix the original character $\chi$, an odd prime $p$ in the
fixed set $P$, and the fixed cyclotomic layer index $j_p$.
Use the common interval with base field
$\mathbb Q(\mu_{pM})$ at $p$. For every growing field (IV1),
every character $\theta$ in its case $IV$, and every $t$
controlled by that interval, the normalized sum of its
level-$N_h$ Selmer lengths is at most
\[
                              B+(N_ht)^C,
                                                               \tag{IV15}
\]
where $C$ depends only on $\chi,p,j_p$.
In particular $C$ is independent of $h,q,d,t$ and of the
chosen coefficient extension.
\end{corollary}
\begin{proof}
The $j_p$-th layer of the cyclotomic $\mathbb Z_p$-extension
of $\mathbb Q(\mu_{pM})$ is
$\mathbb Q(\mu_{p^{j_p+1}M})$: its degree over the base is
$p^{j_p}$. These are precisely the fixed fields entering
(IV13). The prime selection makes $q,d>\max P$.
Lemma~\ref{iv-fixed-conductor} gives
$\operatorname{cond}\theta\mid M$, and
Theorem~\ref{iv-summed} gives (IV15).
\end{proof}

There are at most $\varphi(M)$ distinct characters of conductor
dividing $M$. Thus, if (IV13) holds with a common $B$ for
$p\in P$, the entire collection just treated contributes at
most
\[
 \varphi(M)\left(\sum_{p\in P}\log p\right)B
  +\varphi(M)\sum_{p\in P}\log p\,
       n_p\left(a_p+\frac1{p-1}\right)(t-1)                 \tag{IV16}
\]
to the sums of character lengths weighted by $\log p$.
Both coefficients in (IV16) are fixed before selecting $q,d$
and the growing $2$-power conductor. This establishes the
required dependence on $B$ for these characters; it does not
replace the separate estimates for the other characters.


\section{The bounded-order characters at exponent one}
\label{sv-section}

The part of case $V$ with small $2$-power conductor admits a
direct use of the coarse Selmer bound. The relevant distinction
is that the degree of the field killing its character is bounded
in terms of the fixed data. The growing conductor therefore
occurs only in the discriminant, and the coarse bound becomes a
fixed power of that conductor. We prove this with the exact
threshold used for case $V_s$ in the submitted source.
\cite{alpo-uniform-source}

Retain $\chi,M,o,E_h,N_h$ from (IV1), and fix an odd prime $p$.
Let $\theta$ be a character of $\operatorname{Gal}(E_h/\mathbb Q)$.
Write $2^c$ for the primitive conductor of its local component
$\theta_2$. Let $\vartheta$ be its part of order prime to $p$,
and put
\[
 \xi=\vartheta/\vartheta_2,\quad
 u=\operatorname{ord}(\xi),\quad
 f=\operatorname{ord}_{\operatorname{lcm}(4,u)}(p),\quad
 \nu=v_2(p^f-1),\quad
 R=16o(pM)^2\,4^\nu.
                                                               \tag{SV1}
\]
The multiplicative order defining $f$ is well defined: $u$ is
prime to $p$ and $p$ is odd. All character orders are finite.

\subsection{The field degree is bounded before the conductor grows}

\begin{lemma}\label{sv-order}
The order of the product of all odd-prime components of
$\theta$ divides $2o$. Consequently $u\mid2o$ and
\[
 \begin{split}
 0<f&\le\varphi(4o),\qquad 2^\nu<p^{\varphi(4o)},\\
 R&\le H_p:=16o(pM)^2p^{2\varphi(4o)}.
 \end{split}                                                   \tag{SV2}
\]
If $2^c<R$, then
\[
                       \operatorname{ord}(\theta)\le2oH_p.
                                                               \tag{SV3}
\]
\end{lemma}
\begin{proof}
At an odd prime the inertia group of $E_h/\mathbb Q$ injects
into the product of the inertia groups of
$K_\chi/\mathbb Q$, $\mathbb Q(\sqrt{-q})/\mathbb Q$, and
$\mathbb Q(\sqrt{-d})/\mathbb Q$. The field $R_h$ is
unramified there. The first group has exponent dividing $o$;
the other two have exponent dividing two. Thus each odd local
character has order dividing $2o$. The order of their product
is the least common multiple of their orders, because their
arguments lie in different Chinese remainder factors. It too
divides $2o$. Taking the part of order prime to $p$ gives
$u\mid2o$.

Set $m=\operatorname{lcm}(4,u)$. Then $m\mid4o$ and
$(m,p)=1$. The order $f$ divides $\varphi(m)$ by Lagrange's
theorem. For positive integers $m\mid n$, one has
$\varphi(m)\mid\varphi(n)$: this follows at once from
$\varphi(r^a)=r^{a-1}(r-1)$ by increasing each prime exponent
and then adding the missing primes. Hence
$f\le\varphi(4o)$. The positive integer $p^f-1$ is divisible
by $2^\nu$, so
\[
                 2^\nu\le p^f-1<p^f\le p^{\varphi(4o)}.
\]
Squaring proves the bound for $R$.

The order of $\theta_2$ divides $\varphi(2^c)$ and hence
divides $2^c$; this also holds for $c=0$. The order of
$\theta$ consequently divides $2o\,2^c$. The hypothesis
$2^c<R\le H_p$ implies (SV3).
\end{proof}

\begin{lemma}\label{sv-field}
Let $F_\theta$ be the fixed field of $\ker\theta$ in
$\mathbb Q(\mu_{N_h})$, and put $F=F_\theta(\mu_p)$.
Under $2^c<R$ this is an abelian, totally imaginary field with
\[
 \begin{split}
 \mu_p&\subset F,\qquad \theta|_{G_F}=1,\qquad
 F\subset\mathbb Q(\mu_{pN_h}),\\
 n=[F:\mathbb Q]&\le n_0:=2opH_p,\qquad
                   |\operatorname{disc}F|\le(pN_h)^n.
 \end{split}                                                   \tag{SV4}
\]
Here $n_0$ depends only on $p,M,o$.
\end{lemma}
\begin{proof}
The first statements follow from the definitions. The image of
$\theta$ is cyclic of order $\operatorname{ord}(\theta)$, so
$[F_\theta:\mathbb Q]=\operatorname{ord}(\theta)$.
Adjoining $\mu_p$ increases this degree by at most $p-1$.
Now use (SV3). The presence of $\mu_p$, with $p$ odd, makes
$F$ totally imaginary.

For completeness, if $F\subset L=\mathbb Q(\mu_m)$ and
$n=[F:\mathbb Q]$, the discriminant tower formula implies
\[
             |\operatorname{disc}F|^{[L:F]}
                                      \le|\operatorname{disc}L|.
\]
Moreover $|\operatorname{disc}L|\le m^{[L:\mathbb Q]}$.
One way to obtain the latter inequality without the explicit
cyclotomic discriminant formula is to write
$X^m-1=\Phi_m(X)Q(X)$. At a primitive root $\zeta$,
\[
                 m\zeta^{m-1}=\Phi'_m(\zeta)Q(\zeta).
\]
Taking absolute norms shows
$|N_{L/\mathbb Q}\Phi'_m(\zeta)|\le m^{[L:\mathbb Q]}$,
because $Q(\zeta)$ is a nonzero algebraic integer. The
field discriminant divides the discriminant of the order
$\mathbb Z[\zeta]$, whose absolute value is this derivative
norm. Combining the two inequalities gives
$|\operatorname{disc}F|\le m^n$. Apply this with $m=pN_h$.
\end{proof}

\subsection{Restriction with a bounded invariant module}

Use the normalized length $s_p(1,\theta)$ of (IV6), with an
arbitrary finite coefficient extension containing the values of
$\theta$. Write $b$ for its ramification index.

\begin{lemma}\label{sv-restriction}
For the field $F$ in (SV4),
\[
                     s_p(1,\theta)\le\lambda_p(F,1)+n^2.
                                                               \tag{SV5}
\]
In particular, finiteness of the group on the right implies
finiteness on the left.
\end{lemma}
\begin{proof}
Since $F$ is finite and contains $\mu_p$, its intersection
with $\mathbb Q(\mu_{p^\infty})$ is
$\mathbb Q(\mu_{p^a})$ for an integer $a\ge1$.
Indeed the Galois group above $\mathbb Q(\mu_p)$ is the
procyclic group $1+p\mathbb Z_p$, whose open subgroups are
$1+p^a\mathbb Z_p$. The degree of this intersection gives
\[
                            (p-1)p^{a-1}\le n,
\]
and in particular $a\le n$. On $G_F$, the coefficient
module is $(K/\mathcal O)(-1)$. The cyclotomic image is
$1+p^a\mathbb Z_p$, so its invariant subgroup is precisely
$(K/\mathcal O)[p^a]$, of length $ba$.

Apply the inflation--restriction sequence (IV9), with this
$F$ and coefficient module. The kernel has length at most
$n\,ba$, since its cocycles form a submodule of the functions
from the group of order $n$ to that invariant subgroup.
The image belongs to the strict Selmer group over $F$,
whose coefficient-extended length is $b\lambda_p(F,1)$
by (IV10). Division by $b$ proves (SV5), since $a\le n$.
The argument again requires no coprimality between $p$ and $n$.
\end{proof}

\begin{theorem}\label{sv-exact-case}
For fixed $p,M,o$, there is an integer $C$ such that, for every
$h,q,d,\theta$ satisfying the preceding conditions and $2^c<R$,
the strict Selmer group at $e=1$ is finite and
\[
                           s_p(1,\theta)\log p\le(pN_h)^C.
                                                               \tag{SV6}
\]
The same conclusion holds with the source's extra zero local
conditions at the primes dividing $N_h$.
The exponent $C$ is independent of $h,q,d$ and of the coefficient
extension. Thus this proves the source's case $V_s$ for every
$p$ in its fixed finite set of odd primes.
\end{theorem}
\begin{proof}
Corollary~\ref{sw-coarse-consequence}, applied at exponent one
to the field in (SV4), proves finiteness and
\[
 \lambda_p(F,1)\log p
 \le\exp\bigl((2n)^{86}\log(2|\operatorname{disc}F|)\bigr).
                                                               \tag{SV7}
\]
Put $x=pN_h\ge3$. By (SV4),
$2|\operatorname{disc}F|\le x^{n+1}$. Hence the right side
of (SV7) is at most $x^A$, where
\[
                            A=(2n_0)^{86}(n_0+1)
                                                               \tag{SV8}
\]
is fixed. By (SV5),
\[
                       s_p(1,\theta)\log p
                                      \le x^A+n_0^2\log p.
\]
Choose an integer $k\ge0$ with $3^k\ge n_0^2\log p$.
Then the latter expression is at most
$2x^{\max(A,k)}\le x^{\max(A,k)+1}$, proving (SV6).
Additional local conditions decrease the group.

The much smaller coarse exponent from
Section~\ref{kt-section} could replace $86$ after the strict
Selmer comparison. Its numerical size is immaterial here:
$n_0$ is already fixed. Applying the same argument to a field
whose degree grew with $h$ would not give a fixed exponent,
and no such step is used.
\end{proof}

\begin{corollary}\label{sv-summed}
For a fixed finite set $P$ of odd primes, there is an integer
$C_P$, depending only on $P,M,o$, such that for every $t\ge2$
\[
 \sum_{p\in P}\ \sum_{\substack{\theta\in
          \widehat{\operatorname{Gal}(E_h/\mathbb Q)}\\2^{c(\theta)}<R_p(\theta)}}
                   s_p(1,\theta)\log p\le(N_ht)^{C_P}.
                                                               \tag{SV9}
\]
Thus the entire case $V_s$ contribution has a fixed polynomial
bound before selecting $h,q,d$.
\end{corollary}
\begin{proof}
For fixed $p$, (SV2) bounds the $2$-part of the primitive
conductor by $H_p$. If $T$ is the largest power of two at most
$H_p$, all these local components factor modulo $T$, so there
are at most $\varphi(T)\le H_p$ of them. At the odd primes
of $M$ there are at most $\varphi(M/2^{v_2(M)})\le M$
choices; at each of $q,d$ there are at most two. Thus there
are at most $4MH_p$ characters in this sum, independently of
$h,q,d$.

Choose $C_p$ from Theorem~\ref{sv-exact-case}. The sum is at
most $A N_h^{C_0}$, where
\[
 C_0=\max_{p\in P}C_p,\qquad
 A=\sum_{p\in P}4MH_p p^{C_p}.
\]
For nonempty $P$, choose $r\ge0$ with $2^r\ge A$.
Then $A N_h^{C_0}\le(N_ht)^{C_0+r}$ for $t\ge2$.
For empty $P$ the sum is zero. All choices involve only the
fixed data.
\end{proof}


\section{The finite character comparison}
\label{fcmp-section}

We now pass from character Selmer groups over $\mathbb Q$ to
the required eigenspace over an abelian CM field. The comparison
is exact at primes not dividing the field degree. At the other
odd primes an integral Fourier transform leaves a bounded
torsion error. In the actual family of fields those primes and
the necessary torsion exponents are fixed before the conductor
grows. A class-field calculation bounds their total contribution
by a fixed multiple of $N^2t\log N$.

\subsection{The local conditions and the parity}

Let $L/\mathbb Q$ be a finite abelian CM extension, let
$\Delta=\operatorname{Gal}(L/\mathbb Q)$, let
$d=[L:\mathbb Q]$, and let $j\in\Delta$ be complex
conjugation. Fix $N\ge3$ such that
$L\subset\mathbb Q(\mu_N)$, so $d\le\varphi(N)\le N$.
For an odd prime $p$ and $e\ge1$, put
$A_e=\mathbb Q_p/\mathbb Z_p(-e)$ and retain
$\mathcal S_p(L,e)$ from (MS1). Its $\Delta$-action
conjugates Galois arguments and acts on the coefficients.
Set
\[
 a_p(L,e)=\operatorname{length}_{\mathbb Z_p}
                  \mathcal S_p(L,e)^{\,j=(-1)^e}.
                                                               \tag{FC1}
\]
Thus (FC1) is exactly the complementary-parity summand.

Choose a finite extension $K/\mathbb Q_p$ containing the
values of every character of $\Delta$, and let $\mathcal O$
be its ring of integers, with ramification index $b$.
For $\theta\in\widehat\Delta$ define
\[
 \begin{split}
 A_{e,\theta}&=(K/\mathcal O)
                         (\theta^{-1}\varepsilon_p^{-e}),\\
 \mathcal V_p(e,\theta)
   &=\ker\left(H^1(\mathbb Q,A_{e,\theta})
              \longrightarrow\prod_{\ell\ {\rm prime}}
                         H^1(\mathbb Q_\ell,A_{e,\theta})\right),\\
 \ell_p(e,\theta)&=b^{-1}
       \operatorname{length}_{\mathcal O}\mathcal V_p(e,\theta).
 \end{split}                                                   \tag{FC2}
\]
We allow infinite length until finiteness is proved.

\begin{lemma}\label{fcmp-local}
For $e\ge1$, the strict Selmer group $\mathcal S_p(L,e)$
is the kernel of restriction to every finite decomposition
group. For a character $\theta$ of conductor dividing $N$,
the group $\mathcal V_p(e,\theta)$ is equivalently defined
by zero restriction to decomposition groups at primes
dividing $pN$, and zero restriction to inertia at all
other finite primes.
\end{lemma}
\begin{proof}
Let $w\nmid p$ be a place of $L$. Inertia acts trivially on
$A_e$, and arithmetic Frobenius acts by $q_w^{-e}$, where
$q_w>1$ is the residue-field cardinality. A class trivial
on inertia comes from a cocycle of the procyclic unramified
quotient. Such a cocycle is determined by its value at
Frobenius, modulo the coboundaries
$(q_w^{-e}-1)A_e$. Multiplication by the nonzero
element $q_w^{-e}-1$ of $\mathbb Z_p$ is surjective on $A_e$.
After subtracting a coboundary the cocycle vanishes at
Frobenius, hence on its dense cyclic subgroup, hence
everywhere by continuity. Thus the unramified local
cohomology group is zero.

The same argument applies to $A_{e,\theta}$ at
$\ell\nmid pN$. Its Frobenius scalar is
$\theta(\ell)^{-1}\ell^{-e}$, which is not one: otherwise
the rational number $\ell^e>1$ would be a root of unity.
Again multiplication by the scalar minus one is
surjective on $K/\mathcal O$. This proves both assertions.
At real places $H^1$ is zero because $p$ is odd and the
decomposition group has order two.
\end{proof}

In particular (FC2) has the additional zero local conditions
used for the character groups in the submitted construction.
The use of $e\ge1$ in Lemma~\ref{fcmp-local} is essential.
At $e=0$ the Frobenius scalar can equal one.

\subsection{Shapiro's lemma and integral Fourier maps}

Put $G=G_{\mathbb Q}$, $H=G_L$, and write $\bar g$ for
the image of $g$ in $\Delta$. Consider
\[
 \mathcal M=\operatorname{Map}(\Delta,K/\mathcal O),\qquad
 (g f)(\delta)=\varepsilon_p(g)^{-e}
                                      f(\bar g^{-1}\delta).
                                                               \tag{FC3}
\]
Shapiro's lemma gives
\[
 H^1(G,\mathcal M)\simeq
 H^1(H,(K/\mathcal O)(-e)).
                                                               \tag{FC4}
\]
For clarity, the full cocycle calculation is as follows.
The coinduced module consists of functions $F:G\to K/\mathcal O$
with $F(hx)=\varepsilon_p(h)^{-e}F(x)$ and action
$(gF)(x)=F(xg)$. The identification with (FC3) is
$F(x)=\varepsilon_p(x)^{-e}f(\bar x^{-1})$.
Choose representatives $s(x)$ for the finite right coset
space $H\backslash G$, with $s(H)=1$, and set
\[
 h(x,g)=s(x)g\,s(xg)^{-1}\in H.
                                                               \tag{FC5}
\]
Given a cocycle $c:H\to K/\mathcal O$, define
$C(g)(s(x))=c(h(x,g))$ and extend by the coinduction rule.
The identity
$h(x,g_1g_2)=h(x,g_1)h(xg_1,g_2)$ verifies the cocycle
identity. In the reverse direction use $c(h)=C(h)(1)$.
For an arbitrary $C$, the difference from the reconstruction
of this $c$ is the coboundary of the function
$U(s(x))=C(s(x))(1)$: the cocycle identity gives
\[
 C(g)(s(x))=c(h(x,g))
       +\varepsilon_p(h(x,g))^{-e}U(s(xg))-U(s(x)).
                                                               \tag{FC6}
\]
The formulas preserve coboundaries and continuity, proving
(FC4).

For a rational place $\ell$, the restriction of the coinduced
module to a decomposition group is the direct sum of the
coinduced modules from the decomposition groups at places
of $L$ above $\ell$. This follows by separating the finitely
many orbits on $G/H$; the formulas (FC5)--(FC6) on each
orbit give the corresponding local assertion. Consequently
(FC4) identifies classes which are zero at every finite
decomposition group on its two sides. By
Lemma~\ref{fcmp-local}, this is the coefficient extension of
$\mathcal S_p(L,e)$ on the right.

There are integral, equivariant maps
\[
 \begin{aligned}
 \mathcal F:\mathcal M&\longrightarrow
                         \bigoplus_{\theta\in\widehat\Delta}A_{e,\theta},
 &(\mathcal F f)_\theta
       &=\sum_{\delta\in\Delta}\theta(\delta)^{-1}f(\delta),\\
 \mathcal G:\bigoplus_{\theta\in\widehat\Delta}A_{e,\theta}
                 &\longrightarrow\mathcal M,
 &(\mathcal G a)(\delta)&=\sum_\theta\theta(\delta)a_\theta.
 \end{aligned}                                                 \tag{FC7}
\]
All coefficients are units of $\mathcal O$. The change of
variable $\delta=\bar g\gamma$ shows that the $\theta$
coordinate transforms by
$\theta(\bar g)^{-1}\varepsilon_p(g)^{-e}$, exactly the
character in (FC2). Character orthogonality gives
\[
                    \mathcal G\mathcal F=d\,\mathrm{id},
             \qquad \mathcal F\mathcal G=d\,\mathrm{id}.
                                                               \tag{FC8}
\]
For example, the sum of a nontrivial character is zero
because translation multiplies the sum by a value different
from one; over $K$ this forces the sum to be zero. The
identities in (FC8) are therefore identities over
$\mathcal O$ and on its quotient $K/\mathcal O$.

The action of $j$ under (FC4) corresponds to
$T_jf(\delta)=f(\delta j)$. To check the coefficient sign,
choose a complex conjugation $c\in G$ above $j$.
For a cocycle $C$ with values in (FC3), put
$c_H(h)=C(h)(1)$. The cocycle identity implies
\[
 C(h)(j)=\varepsilon_p(c)^{-e}c_H(chc^{-1})
                      +(1-\varepsilon_p(h)^{-e})C(c)(j).
                                                               \tag{FC9}
\]
The last term is a coboundary. The first is the usual
action of complex conjugation on $H^1(H,(K/\mathcal O)(-e))$.
On the $\theta$ coordinate of (FC7), $T_j$ acts by
$\theta(j)$. Thus the required characters are precisely
\[
                 X_e(L)=\{\theta\in\widehat\Delta:
                                      \theta(j)=(-1)^e\}.
                                                               \tag{FC10}
\]
The factor $(-1)^e$ has not been lost in changing the
coefficient module or the coinduction convention.

\begin{proposition}\label{fcmp-fourier}
Let $\nu=v_p(d)$.
If $\mathcal V_p(e,\theta)$ is finite for every
$\theta\in X_e(L)$, then $\mathcal S_p(L,e)^{j=(-1)^e}$
is finite and
\[
 a_p(L,e)\le\sum_{\theta\in X_e(L)}\ell_p(e,\theta)
             +\operatorname{length}_{\mathbb Z_p}
                             \mathcal S_p(L,e)[p^\nu].
                                                               \tag{FC11}
\]
When $p\nmid d$, the last term is zero and equality holds
between $a_p(L,e)$ and the character sum.
\end{proposition}
\begin{proof}
Since $p$ is odd, the two projectors $(1\pm T_j)/2$
are integral. Apply (FC7) to the image of the projector
$(1+(-1)^eT_j)/2$, and then to locally trivial cohomology.
The resulting map lands in
$\bigoplus_{\theta\in X_e(L)}\mathcal V_p(e,\theta)$.
By (FC8), its kernel is killed by $d$, equivalently by
$p^\nu$. Under (FC4), it is therefore contained in the
$p^\nu$-torsion of $\mathcal S_p(L,e)\otimes_{\mathbb Z_p}\mathcal O$.
When $\nu>0$, finiteness of this torsion, independently
of finiteness of the whole Selmer group, is proved in
Lemma~\ref{fcmp-torsion} below. Kernel and image are finite,
so the stated eigenspace is finite.

Coefficient extension commutes with cochains, cohomology
and the local kernels because $\mathcal O$ is finite free
over $\mathbb Z_p$. It also commutes with the two projectors
and with $p^\nu$-torsion. For a finite $\mathbb Z_p$-module
$T$,
\[
 \operatorname{length}_{\mathcal O}
                  (T\otimes_{\mathbb Z_p}\mathcal O)
                =b\,\operatorname{length}_{\mathbb Z_p}T.
                                                               \tag{FC12}
\]
This follows first from $\mathcal O/p^a\mathcal O$ and
then from the elementary-divisor decomposition of $T$.
Taking lengths and dividing by $b$ proves (FC11).
If $p\nmid d$, (FC8) makes the two Fourier maps inverse
after division by the unit $d$. They preserve every
local condition and the parity projector, which proves
the equality.
\end{proof}

\begin{remark}\label{fcmp-integral-error}
An exact integral decomposition cannot be substituted for
(FC8) when $p\mid d$. For instance, take $\Delta=\mathbb Z/p\mathbb Z$,
let $K$ contain a primitive root $\zeta_p$, and put
$\pi=\zeta_p-1$. In $\operatorname{Map}(\Delta,K/\mathcal O)$
let $f(0)=1/\pi$, $f(1)=-1/\pi$, and let all other
coordinates be zero. This is a nonzero vector, but its
Fourier coefficient at the character $a\mapsto\zeta_p^{ka}$
is $(1-\zeta_p^{-k})/\pi\in\mathcal O$, hence is zero in
$K/\mathcal O$. The same example occurs in either parity
summand for $\mathbb Z/2\mathbb Z\times\mathbb Z/p\mathbb Z$
by applying the integral projector for the first factor.
This explains the necessity of controlling the kernel;
it is not a counterexample to (FC11).
\end{remark}

\subsection{A torsion bound independent of the twist}

\begin{lemma}\label{fcmp-torsion}
Let $L$ be any number field, let $p$ be odd, let $e\ge1$,
and let $\nu\ge1$. Set
$F=L(\mu_{p^\nu})$, $u=[F:L]$, $n_F=[F:\mathbb Q]$,
and let $s_F$ be the number of places of $F$ above $p$.
Then
\[
 \operatorname{length}_{\mathbb Z_p}\mathcal S_p(L,e)[p^\nu]
       \le \nu u+\nu(n_F+s_F)+\frac{\log h_F}{\log p}
       \le \nu u+2\nu n_F+\frac{\log h_F}{\log p}.
                                                               \tag{FC13}
\]
In particular the torsion on the left is finite. The
bound is independent of $e$.
\end{lemma}
\begin{proof}
Put $B=A_e[p^\nu]$. The coefficient exact sequence
$0\to B\to A_e\xrightarrow{p^\nu}A_e\to0$ shows that
every class in $\mathcal S_p(L,e)[p^\nu]$ lifts to
$H^1(L,B)$. At a place away from $p$, inertia acts
trivially on $A_e$ and multiplication by $p^\nu$ on
$H^0(I_w,A_e)=A_e$ is surjective. Hence
$H^1(I_w,B)\to H^1(I_w,A_e)$ is injective.
Every such lift is consequently unramified away from $p$.
It suffices to bound this larger finite-coefficient
cohomology group; no choice of compatible local
primitives above $p$ is required.

Over $F$, the group $B$ is a trivial cyclic group of
order $p^\nu$. Inflation--restriction has kernel
$H^1(\operatorname{Gal}(F/L),B)$, whose cardinality is at
most $p^{\nu u}$: its cocycles are a subgroup of the
functions from a group of order $u$ to $B$. The image
consists of homomorphisms $G_F\to\mathbb Z/p^\nu\mathbb Z$
unramified away from $p$.

The ray class exact sequence and global reciprocity give
\[
 \#\operatorname{Hom}_{\mathrm{unr}\ {\rm outside}\ p}
                   (G_F,\mathbb Z/p^\nu\mathbb Z)
 \le h_F\prod_{w\mid p}
       \#\operatorname{Hom}_{\mathrm{cont}}
                 (\mathcal O_{F_w}^{\times},\mathbb Z/p^\nu\mathbb Z).
                                                               \tag{FC14}
\]
Here is the precise passage. Every homomorphism on the
left factors through a ray class group with modulus
supported above $p$, by reciprocity. The kernel of the
map from that ray class group to the ideal class group
is a quotient of the product of the corresponding
local unit quotients. The number of homomorphisms from
the class group is at most $h_F$; restriction to that
kernel bounds the remaining choices by the displayed
local factors. The groups of homomorphisms for increasing
moduli form an increasing union. The uniform finite
bound just obtained also bounds this union.
These are the ray class exact sequence and reciprocity
in Milne, Chapter~V, Theorems~1.7 and~3.5--3.6.
\cite{milne-cft}

For a finite extension $K_0/\mathbb Q_p$ of degree $n_0$,
\[
 \#\operatorname{Hom}_{\mathrm{cont}}
        (\mathcal O_{K_0}^{\times},\mathbb Z/p^\nu\mathbb Z)
                         \le p^{\nu(n_0+1)}.                   \tag{FC15}
\]
To justify this bound, the residue-field unit group has
order prime to $p$ and its Teichm\"uller representatives
split it off. The principal units form an abelian pro-$p$
group. On a sufficiently deep principal-unit subgroup,
the convergent logarithm and exponential give inverse
isomorphisms with an ideal of $\mathcal O_{K_0}$: choose
the depth to have $p$-adic valuation greater than
$1/(p-1)$, and the usual series converge and their
formal inverse identities hold. This subgroup is
$\mathbb Z_p^{n_0}$ and has finite index. Thus the
principal units are a finitely generated
$\mathbb Z_p$-module of rank $n_0$. Their torsion is a
finite subgroup of a field's multiplicative group,
and is therefore cyclic. The structure theorem over
$\mathbb Z_p$ gives $\mathbb Z_p^{n_0}$ plus one finite
cyclic $p$-group. Each factor contributes at most
$p^\nu$ characters, proving (FC15).

Now $\sum_{w\mid p}[F_w:\mathbb Q_p]=n_F$.
Equations (FC14)--(FC15) give the bound
$h_Fp^{\nu(n_F+s_F)}$. Combining with the
inflation--restriction kernel and taking $\log_p$
proves the first inequality in (FC13).
The second uses $s_F\le n_F$.
\end{proof}

\subsection{The growing family and its uniform error}

Return to the fields
$L=E_h=K_\chi R_h(\sqrt{-q},\sqrt{-d_0})$
of (IV1), writing $d_0$ for the second auxiliary prime
to distinguish it from the degree $d=[L:\mathbb Q]$.
The fixed primitive character $\chi$ has conductor $M$
and order $o$. The degree of $K_\chi$ is $o$ and that of
$R_h$ is $2^h$; the two quadratic adjunctions show
\[
                            d\mid 2^{h+2}o.
                                                               \tag{FC16}
\]
For each odd prime $p\mid o$, put
\[
 a_p^0=v_p(o),\qquad m_p=\varphi(p^{a_p^0}),\qquad
 C_p=3a_p^0m_p\frac{\log p}{\log2}
       +\frac{m_p^2}{2}
                    \left(1+a_p^0\frac{\log p}{\log2}\right),
 \quad C_o=\sum_{\substack{p\mid o\\p\ne2}}C_p.
                                                               \tag{FC17}
\]
These constants depend only on the original character.
In particular $C_o=0$ if $o$ is a power of two.

\begin{theorem}\label{fcmp-uniform}
Let $P$ be a finite set of odd primes and $t\ge2$ an integer.
Suppose that the groups in (FC2) are finite for every
$p\in P$, $1\le e<t$, and $\theta\in X_e(E_h)$. Then
\[
 \begin{split}
 \sum_{p\in P}\log p\sum_{e=1}^{t-1}a_p(E_h,e)
 &\le
 \sum_{p\in P}\log p\sum_{e=1}^{t-1}
                    \sum_{\theta\in X_e(E_h)}\ell_p(e,\theta)\\
 &\quad+C_o N_h^2(t-1)\log N_h .
 \end{split}                                                   \tag{FC18}
\]
The coefficient of each character length is exactly one,
regardless of the coefficient extension. The error has
a fixed polynomial bound in $N_ht$.
\end{theorem}
\begin{proof}
By (FC16), the Fourier error is zero unless $p\mid o$,
and otherwise $\nu=v_p(d)\le a_p^0$. In (FC13) one may
enlarge $F$ to $E_h(\mu_{p^{a_p^0}})$; the same proof
uses $p^\nu$ coefficients and gives the same formula
with this larger field. Then
\[
 u\le m_p,\quad n_F\le dm_p,\quad
 F\subset\mathbb Q(\mu_{N_hp^{a_p^0}}),\quad
 |\operatorname{disc}F|
                 \le (N_hp^{a_p^0})^{n_F}.                     \tag{FC19}
\]
The discriminant estimate follows from the tower formula
and the cyclotomic derivative-norm argument in the proof
of Lemma~\ref{sv-field}.
Corollary~\ref{sl-class-number} and $\log(1+x)\le x$ give
$\log h_F\le \tfrac12 n_F\log|\operatorname{disc}F|$.
Consequently the error in (FC11), multiplied by $\log p$,
is at most
\[
 a_p^0m_p\log p+2a_p^0dm_p\log p
            +\frac12d^2m_p^2(\log N_h+a_p^0\log p)
 \le C_p d^2\log N_h .
                                                               \tag{FC20}
\]
For the last inequality use $d\ge1$ and $\log N_h\ge\log2$.
Sum (FC20) over the fixed primes dividing $o$ and over
$1\le e<t$, and use $d\le N_h$. Together with (FC11)
this proves (FC18). Since $\log N_h\le N_h$ and $C_o$
is fixed, its error is bounded by $(N_ht)^C$ for one
integer $C$ chosen before $h,q,d_0,t$.
\end{proof}

In the source notation the characters in (FC10) are
exactly the parity-restricted characters of the chosen
abelian field, and (FC2) is the normalized length of
the group with zero local conditions at primes dividing
$pN_h$.\cite{alpo-uniform-source}
Thus (FC18) supplies the whole-field comparison needed
to assemble the complementary-parity estimates.
It does not require that $p$ be prime to the field
degree. In particular a bound $CB$ for the total
fixed-conductor character contribution remains $CB$:
the passage to $E_h$ adds only the polynomial error
in (FC18). The unproved character estimates in the
remaining automorphic cases are still required.


\section{Adjoint vanishing at the automorphic points}
\label{av-section}

A theorem of Thorne gives a substantial simplification of the
complementary-parity argument. It proves the adjoint vanishing
assertion at each of the automorphic points occurring in the
submitted construction. We give the comparison of representations,
polarizations, coefficient fields, and local conditions explicitly.
The theorem is a classical input; its Taylor--Wiles proof is not
reproduced here. The argument below does not assert that every
deformation point is automorphic, and gives no bound for the length
of an integral congruence module.

\subsection{Local conditions and finite coefficient fields}

Let $p$ be a prime, let $F$ be a number field, and let $E$ be a
finite extension of $\mathbb Q_p$. Write $G_F$ for the absolute
Galois group of $F$. If $W$ is a finite-dimensional continuous
$E$-representation of $G_F$, a continuous one-cocycle is a map
$z:G_F\longrightarrow W$ satisfying
\[
                    z(gh)=z(g)+g z(h).                         \tag{AV1}
\]
A coboundary has the form $z(g)=gX-X$ for one $X\in W$.
The quotient of cocycles by coboundaries is $H^1(F,W)$.
For each finite place $v$, fix a decomposition group $G_v$
and its inertia group $I_v$. Put
\[
 H^1_f(F_v,W)=
 \begin{cases}
 \ker\bigl(H^1(G_v,W)\longrightarrow H^1(I_v,W)\bigr),&v\nmid p,\\
 \ker\bigl(H^1(G_v,W)\longrightarrow
       H^1(G_v,W\otimes_{\mathbb Q_p}B_{\mathrm{cris}})\bigr),&v\mid p.
 \end{cases}                                                   \tag{AV2}
\]
Here $B_{\mathrm{cris}}$ is Fontaine's crystalline period ring,
with its usual continuous local Galois action; the tensor product
has the diagonal action. Define
\[
 H^1_f(F,W)=\{\alpha\in H^1(F,W):
                    \alpha|_{G_v}\in H^1_f(F_v,W)\text{ for all finite }v\}.
                                                               \tag{AV3}
\]
Changing a decomposition group by conjugation does not change
this condition. Indeed (AV1) gives
\[
 z(g)=\gamma z(\gamma^{-1}g\gamma)+z(\gamma)-g z(\gamma).
                                                               \tag{AV4}
\]
Thus the cocycles on conjugate local groups differ, after the
natural conjugation isomorphism, by a coboundary. The same
calculation applies after tensoring with $B_{\mathrm{cris}}$.

We shall also use representations and cocycles with values in
$\overline{\mathbb Q}_p$, with its $p$-adic topology. The following
elementary point avoids an implicit assumption on their field
of definition.

\begin{lemma}\label{av-coefficients}
Every compact subgroup of
$\mathrm{GL}_d(\overline{\mathbb Q}_p)$ is contained in
$\mathrm{GL}_d(E)$ for some finite extension
$E/\mathbb Q_p$. Consequently a continuous representation of a
profinite group and one continuous cocycle for it are defined,
in their given bases, over a common finite extension of
$\mathbb Q_p$.
\end{lemma}
\begin{proof}
There are countably many finite subfields of
$\overline{\mathbb Q}_p$ over $\mathbb Q_p$. Here is a useful
justification. If $K=\mathbb Q_p(\alpha)$ is finite, approximate
the coefficients of the monic minimal polynomial of $\alpha$
by rational numbers, keeping the degree fixed. Hensel's lemma
in the complete field $K$ gives a root $\beta\in K$ arbitrarily
close to $\alpha$ of one of the resulting polynomials in
$\mathbb Q[X]$. Take
\[
 |\alpha-\beta|_p<
        \min_{\alpha'\ne\alpha}|\alpha-\alpha'|_p,
                                                               \tag{AV5}
\]
where $\alpha'$ runs over the other $\mathbb Q_p$-conjugates.
Every automorphism fixing $\beta$ preserves the $p$-adic
absolute value, so (AV5) forces it to fix $\alpha$.
Thus $\mathbb Q_p(\alpha)\subseteq\mathbb Q_p(\beta)$,
and the reverse inclusion follows from $\beta\in K$.
The degree-one case is immediate. There are only countably
many such algebraic numbers $\beta$, which proves the assertion.

Enumerate these finite subfields as $E_1,E_2,\ldots$, and let
$H$ be the compact subgroup in question. Each
$H_i=H\cap\mathrm{GL}_d(E_i)$ is closed in $H$, because a
finite extension of $\mathbb Q_p$ is complete and therefore
closed in $\overline{\mathbb Q}_p$. Their union is $H$.
The Baire theorem for the compact metric space $H$ implies
that some $H_i$ has nonempty interior. A subgroup with
nonempty interior is open, so its index in $H$ is finite.
Adjoin to $E_i$ the entries of finitely many coset
representatives. The resulting finite extension contains
every matrix in $H$.

For the last assertion, if $R$ is the representation, the map
\[
       g\longmapsto
       \begin{pmatrix}R(g)&z(g)\\0&1\end{pmatrix}               \tag{AV6}
\]
is a continuous representation by (AV1). Its image is compact.
Apply the first assertion to (AV6).
\end{proof}

The local conditions also descend. Suppose that $R$ and $z$
are defined over $E$ and that a local coboundary equation has
a solution after extending coefficients to
$\overline{\mathbb Q}_p$. Choose an $E$-linear map
$\lambda:\overline{\mathbb Q}_p\longrightarrow E$ with
$\lambda(1)=1$, and apply it to the solution's coordinates.
This gives a solution over $E$. At a place above $p$, apply
$\lambda\otimes1$ to the solution in the algebraic tensor
product with $B_{\mathrm{cris}}$. The coefficient action of
the local Galois group is trivial, so this map commutes with
the diagonal action. This proves the same assertion there.
No topological assertion about $\lambda$ is needed: it is
applied to a single vector giving a coboundary, not to a
variable cocycle. Lemma~\ref{av-coefficients} and this observation
justify the use of (AV3) for the cocycles with algebraic
$p$-adic coefficients in the submitted source.

\subsection{The polarization and its adjoint action}

Let now $F$ be a CM field with maximal totally real subfield
$F^+$. Fix a complex conjugation $c\in G_{F^+}$, so that
$c^2=1$ and $c\notin G_F$. Let $n\ge2$, let
$\varepsilon_p:G_{F^+}\longrightarrow\mathbb Z_p^\times$
be the cyclotomic character, and fix an abstract field
isomorphism $\iota:\overline{\mathbb Q}_p\simeq\mathbb C$.
This isomorphism is not asserted to be continuous.

For a representation
$r:G_F\longrightarrow\mathrm{GL}_n(\overline{\mathbb Q}_p)$,
write $r^c(g)=r(cgc^{-1})$. A matrix $P\in
\mathrm{GL}_n(\overline{\mathbb Q}_p)$ gives the required
polarization when
\[
           r(g)^t P r^c(g)=\varepsilon_p(g)^{1-n}P
                         \qquad(g\in G_F).                    \tag{AV7}
\]
For a symmetric such $P$, the polarized adjoint representation
$\operatorname{Ad}_P(r)$ of $G_{F^+}$ acts on
$M_n(\overline{\mathbb Q}_p)$ by
\[
 \begin{aligned}
 g\cdot X&=r(g)Xr(g)^{-1}&& (g\in G_F),\\
 c\cdot X&=-P^{-1}X^tP.
 \end{aligned}                                                 \tag{AV8}
\]
The square of the second operator is the identity. Transposing
(AV7), or solving it for $r^c(g)$, verifies that this operator
conjugates $\operatorname{Ad}(r(g))$ to
$\operatorname{Ad}(r^c(g))$. These are the defining relations
for the extension from the subgroup $G_F$ of index two.

\begin{theorem}[Classical adjoint vanishing]\label{av-thorne}
Suppose that $\pi$ is a regular algebraic cuspidal automorphic
representation of $\mathrm{GL}_n(\mathbb A_F)$ satisfying
$\pi^c\simeq\pi^\vee$. Let $r_{\pi,\iota}$ be its continuous
semisimple Galois representation, in the Tate normalization.
If
\[
 r_{\pi,\iota}|_{G_{F(\mu_{p^\infty})}}
                         \quad\hbox{is irreducible},            \tag{AV9}
\]
then it has a symmetric polarization as in (AV7), and
\[
 H^1_f\bigl(F^+,\operatorname{Ad}_P(r_{\pi,\iota})\bigr)=0.
                                                               \tag{AV10}
\]
This holds for every prime $p$; no condition on its residual
representation is imposed.
\end{theorem}

The vanishing is Thorne's Theorem~6.1, specialized to the
conjugate-self-dual case. The symmetric pairing and the
minus sign in (AV8) are the conventions in Newton--Thorne,
page~1920; the Frobenius normalization is fixed on page~1923.
\cite{thorne-adjoint-2022,newton-thorne-adjoint-2023}
In the submitted source's convention a polarized pair also
records a character $\chi$ of $F^+$ satisfying
$\chi_v(-1)=(-1)^n$. For the above $\pi$ take
$\chi=\delta_{F/F^+}^n$, where $\delta_{F/F^+}$ is the
quadratic Hecke character of the CM extension. Its pullback
by the idele norm from $F$ is trivial, and its component at
each real place is the sign character. Its Galois character
restricts trivially to $G_F$. This explains both the parity
condition and the multiplier $\varepsilon_p^{1-n}$ in (AV7).
In odd rank one must not replace this auxiliary character by
the trivial character in that particular convention.

\begin{lemma}\label{av-pairing-unique}
If $r$ is irreducible and $P_0,P$ both satisfy (AV7), then
$P$ is a nonzero scalar multiple of $P_0$. Consequently the
operator in the second line of (AV8) is independent of the
choice of polarization.
\end{lemma}
\begin{proof}
Substitution of $P=P_0T$ in (AV7) shows that $T$ commutes
with every $r^c(g)$. An eigenspace of $T$ exists over
$\overline{\mathbb Q}_p$, and is stable under $r^c$.
Irreducibility forces this eigenspace to be the whole space,
so $T$ is scalar. Invertibility makes the scalar nonzero,
and it cancels in $P^{-1}X^tP$.
\end{proof}

Scalar twists cause no difficulty either. For a character
$\alpha$ of $G_F$, conjugation by $\alpha(g)r(g)$ equals
conjugation by $r(g)$. If the polarized extension uses the
same pairing, its second operator in (AV8) is identical too.
In particular scalar twisting changes neither the invariant
subspaces in (AV9) nor this adjoint representation.

\subsection{Comparison with the actual automorphic point}

We record explicitly the elementary uniqueness principle used
when the Galois representation is specified by Frobenius
polynomials. If two continuous semisimple representations
over $\overline{\mathbb Q}_p$ have equal characteristic
polynomials at all but finitely many unramified Frobenius
classes, they are isomorphic. Lemma~\ref{av-coefficients}
puts both over a finite extension, and Chebotarev density
and continuity give equality of their traces at every group
element. To see the last algebraic implication, let $A$ be
the finite-dimensional algebra spanned by the image on the
direct sum of the two representation spaces. The direct
sum is a faithful semisimple $A$-module, so the radical of
$A$ is zero. Over an algebraically closed field, $A$ is a
product of matrix algebras. The trace of each central
primitive idempotent gives the multiplicity of the
corresponding simple module, multiplied by its nonzero
dimension. Equality of traces on group elements implies
equality on their span, hence on these idempotents. The
multiplicities agree, proving the assertion. This is the
usual Brauer--Nesbitt argument, included to specify exactly
what the Frobenius comparison supplies.

Here is the quotient form needed in the source. Let
$q:G_{F^+}\twoheadrightarrow\Gamma$ be a continuous quotient
of profinite groups. Suppose that a quadratic character
$\eta:\Gamma\longrightarrow\{\pm1\}$ and a continuous
character $\epsilon:\Gamma\longrightarrow\mathbb Z_p^\times$
satisfy
\[
 \ker(\eta q)=G_F,\qquad
 \epsilon q=\varepsilon_p.
                                                               \tag{AV11}
\]
Put $\Gamma_0=\ker\eta$, $c_0=q(c)$ and
$\Gamma_1=\ker(\epsilon|_{\Gamma_0})$.

\begin{proposition}\label{av-quotient}
Let $\rho:\Gamma_0\longrightarrow
\mathrm{GL}_n(\overline{\mathbb Q}_p)$ be continuous, and
assume that $\rho|_{\Gamma_1}$ is irreducible. Suppose that
$r=\rho\circ q|_{G_F}$ is attached to a regular algebraic
conjugate-self-dual cuspidal representation as in
Theorem~\ref{av-thorne}, by the Frobenius criterion just
described. Let $P$ be any matrix satisfying the quotient
version of (AV7), and let $R$ be a representation of
$\Gamma$ satisfying
\[
 R(g)X=\rho(g)X\rho(g)^{-1}\quad(g\in\Gamma_0),\qquad
 R(c_0)X=-P^{-1}X^tP.                                          \tag{AV12}
\]
If $z:\Gamma\longrightarrow M_n(\overline{\mathbb Q}_p)$
is a continuous one-cocycle whose pullback along $q$
satisfies (AV2) at every finite place of $F^+$, then there is
one matrix $X$ such that
\[
                          z(g)=R(g)X-X\qquad(g\in\Gamma).
                                                               \tag{AV13}
\]
\end{proposition}
\begin{proof}
Surjectivity of $q$ and the first equality in (AV11) give
$q(G_F)=\Gamma_0$. The second equality then gives the exact
identity
\[
          q\bigl(G_{F(\mu_{p^\infty})}\bigr)=\Gamma_1.
                                                               \tag{AV14}
\]
Indeed, lift an element of $\Gamma_1$ to $G_{F^+}$;
the first equality forces the lift into $G_F$, and the
second forces its cyclotomic character to be one. The
reverse inclusion follows immediately from (AV11).
Thus the assumed irreducibility is precisely (AV9).

The Frobenius criterion identifies $r$ with $r_{\pi,\iota}$.
Pulling $P$ back gives (AV7). By
Lemma~\ref{av-pairing-unique}, it differs by a scalar from
the symmetric pairing of Theorem~\ref{av-thorne}. Hence
$R\circ q$ is its polarized adjoint representation: the
two actions agree on $G_F$ and on $c$, which together
generate $G_{F^+}$. Notice that the proposition did not
need a separate symmetry assumption on $P$.

The representation $R$ is continuous. This follows from
(AV12), continuity of $\rho$, and the two open cosets of
$\Gamma_0$. By Lemma~\ref{av-coefficients} the cocycle
and representation have a common finite coefficient field;
the local descent argument following that lemma applies.
Theorem~\ref{av-thorne} therefore gives
\[
                         z(q(\sigma))=R(q(\sigma))X-X
                         \qquad(\sigma\in G_{F^+}).
\]
Surjectivity of $q$ proves (AV13).
\end{proof}

\begin{corollary}\label{av-source}
The adjoint-vanishing assertion at the automorphic points of
the submitted ordinary deformation construction follows from
the classical theorem above.
\cite{alpo-uniform-source}
\end{corollary}
\begin{proof}
Use $F=L$, $F^+=L^+$, $\Gamma=G_{L^+,S'}$ and the
natural quotient $q$. The source's level data state (AV11)
exactly, and its selected $c_0$ is the image of a genuine
complex conjugation. The irreducibility assumption concerns
exactly the subgroup $\Gamma_1$. The automorphic-point
hypothesis supplies the cuspidal representation, its regular algebraicity,
conjugate self-duality, and the matching semisimple Galois
representation. In the source's Frobenius convention the
polynomial of $r(\sigma^{-1})$, for arithmetic Frobenius
$\sigma$, has Satake roots multiplied by
$q_v^{(n-1)/2}$. This is the same Tate normalization used
in Theorem~\ref{av-thorne}.

The specified polarization and adjoint action are precisely
(AV7) and (AV12). The local condition says that the
restriction of the cocycle is a coboundary on inertia away
from $p$, and becomes one after tensoring with
$B_{\mathrm{cris}}$ above $p$. It initially chooses one
embedding above each place; (AV4) shows independence of
that choice. Its local group consists of pairs of global
and local automorphisms which agree on the chosen embedded
algebraic closure of the number field. Projection to the
local factor identifies it with the Galois group of the
corresponding completion: the local automorphism has a
unique restriction to that algebraic closure, and fixing
the number field is equivalent, by continuity, to fixing
its completion. Thus these are exactly (AV2).
Proposition~\ref{av-quotient}
therefore gives the requested global coboundary conclusion.
All extra ordinary, weight, rank-at-most-four and local
genericity assumptions in the source implication may be
retained, but are unnecessary for this deduction.
\end{proof}

\subsection{The symplectic projection}

For later use we also give the matrix identity which transports
this vanishing to symplectic infinitesimal deformations.
Let $J\in\mathrm{GL}_{2m}(E)$ satisfy $J^t=-J$, and put
\[
 X^*=J^{-1}X^tJ,\qquad
 \mathfrak{sp}(J)=\{X:X^*=-X\},\qquad
                        P_-(X)=\frac{X-X^*}{2}.               \tag{AV15}
\]
The symbol $P_-$ denotes a linear projection and is unrelated
to the polarization matrix $P$ above. The coefficient field
has characteristic zero, including when $p=2$.
Suppose that $r:G\longrightarrow\mathrm{GSp}(J,E)$ is a
representation and $\eta:G\longrightarrow\{\pm1\}$ is
a character. If $A_g(X)=r(g)Xr(g)^{-1}$, then
\[
 R_g(X)=A_g(P_-X)+\eta(g)A_g(X-P_-X)                            \tag{AV16}
\]
defines a representation. To check this, the similitude
identity gives $(A_gX)^*=A_g(X^*)$, so $P_-$ commutes
with every $A_g$. It is idempotent, and (AV16) acts by
$A_g$ on its image and by $\eta(g)A_g$ on its kernel.
The group law follows on these two summands.

When $\eta(c)=-1$, (AV16) gives
\[
 R_c(X)=-r(c)J^{-1}X^tJr(c)^{-1}
       =-(Jr(c)^{-1})^{-1}X^t(Jr(c)^{-1}).                      \tag{AV17}
\]
For $g\in\ker\eta$, its action is the ordinary adjoint
action. Thus (AV17) is exactly the polarized extension
with pairing matrix $Jr(c)^{-1}$, whenever the latter
representation has the automorphic polarization considered
above.

The inclusion $\mathfrak{sp}(J)\hookrightarrow M_{2m}(E)$
and projection $P_-$ are equivariant for these actions.
They commute with restriction to inertia and with the
diagonal action after tensoring with $B_{\mathrm{cris}}$.
They therefore give split maps on the Selmer groups in
(AV3). In particular vanishing for $R$ implies vanishing
for its symplectic summand. Explicitly, if a cocycle
$z(g)\in\mathfrak{sp}(J)$ becomes $R_gX-X$, then
\[
                    z(g)=A_g(P_-X)-P_-X.                       \tag{AV18}
\]
This retains the primitive in the symplectic Lie algebra;
it is stronger than merely finding a primitive in all
matrices. The argument is the ordinary matrix calculation
needed in the source's rank-four rigidity step.

\paragraph{What remains.}
Corollary~\ref{av-source} closes the adjoint-vanishing
implication at the actual automorphic points. To deduce
the remaining complementary-parity estimates, one still
needs the arithmetic construction and automorphy of the
deformations which encode the original rank-one Selmer
classes, and the quantitative integral comparison with
the relevant congruence modules. Rational vanishing does
not bound the length of a torsion module: the modules
$\mathbb Z_p/p^a\mathbb Z_p$ all become zero after
tensoring with $\mathbb Q_p$, while their lengths are
the arbitrary integers $a$. Thus neither the polynomial
exponent nor the coefficient of the common interval
bound $B$ can be inferred from (AV10) alone. These
remaining steps must be checked with their original
uniform quantifiers.


\section{Compatibility of the prime selection and the cohomological interval}
\label{sec:selection-assembly}

This section checks the order of choices in the arithmetic construction.
It uses the selected-character theorem under its stated QRH hypothesis,
and the common-interval theorem under its stated coarse cohomological
hypothesis. Neither hypothesis is strengthened here. In particular, the
common-interval theorem controls the prime $2$ and specified fixed
odd-prime fields; the estimates for all the remaining character components
are identified separately below.

\subsection{The fixed field and the growing cyclotomic extension}

Fix a primitive Dirichlet character $\psi$ of conductor $M\ge1$ and
order $o$. Let $K_\psi$ be the fixed field of $\ker\psi$ in
$\mathbb Q(\mu_M)$. Thus $[K_\psi:\mathbb Q]=o$. Put
\[
 R_h=\mathbb Q(\mu_{2^{h+2}})^+,
 \qquad B_h=K_\psi R_h,
 \qquad E_h=B_h(\sqrt{-q},\sqrt{-d}),
\]
where $q,d$ are distinct primes prime to $2M$, both congruent to $3$
modulo $4$. The superscript $+$ denotes the maximal real subfield.

\begin{lemma}\label{sa-fields}
Write $a=v_2(M)$ and $h_*=\max(a-2,0)$. There is an integer
$1\le c\le o$, depending only on $\psi$, such that, for $h\ge h_*$,
\[
 [B_h:\mathbb Q]=c2^h,\qquad [E_h:\mathbb Q]=4c2^h.
\]
One may take
$c=[K_\psi:K_\psi\cap\bigcup_hR_h]$.
The field $E_h$ is abelian and CM, and is contained in
$\mathbb Q(\mu_{N_h})$, where
\[
 Q_h=\operatorname{lcm}(2^{h+2},M),\qquad N_h=Q_hqd.
\]
In any cyclotomic overfield containing these roots of unity, the
automorphism $\zeta\mapsto\zeta^u$ fixes $E_h$ if and only if
\begin{equation}\label{sa-field-kernel}
 u\equiv\pm1\pmod{2^{h+2}},\quad
 \psi(u)=1,\quad (u/q)=(u/d)=1.
\end{equation}
Here $u$ is prime to the conductor of the overfield.
\end{lemma}

\begin{proof}
Write $M=2^am$ with $m$ odd. The Chinese remainder theorem for
cyclotomic Galois groups shows, for $n\ge a$, that
$\mathbb Q(\mu_M)\cap\mathbb Q(\mu_{2^n})=
\mathbb Q(\mu_{2^a})$; the factors of odd and $2$-power conductor
are independent. Consequently $K_\psi\cap\bigcup_hR_h$ is
already contained in $R_h$ when $h+2\ge a$. The compositum degree
formula and $[R_h:\mathbb Q]=2^h$ give the first assertion.
The integer $c$ is a relative degree and is at most $o$.

For completeness, the quadratic subfield of $\mathbb Q(\mu_q)$
is $\mathbb Q(\sqrt{-q})$ in the present congruence class. Indeed,
for $\zeta=\exp(2\pi i/q)$ and
$\tau=\sum_{u\bmod q}(u/q)\zeta^u$, changing variables gives
$\sigma_v\tau=(v/q)\tau$ and $\overline\tau=(-1/q)\tau$.
Also, writing $u=tv$ in $\tau\overline\tau$, the inner sum over
$v\ne0$ is $q-1$ if $t=1$ and $-1$ otherwise. Hence
$\tau\overline\tau=q$, and $\tau^2=(-1/q)q=-q$.
The same proof applies to $d$. These identities and Galois
correspondence prove the containment and \eqref{sa-field-kernel}.
They also identify any field specified by that fixed-subgroup
condition with $E_h$ up to a $\mathbb Q$-isomorphism.

The field $B_h$ is unramified at $q,d$: at a prime not dividing
its cyclotomic modulus, the polynomial $X^n-1$ has invertible
derivative at every root, so its cyclotomic factors have
discriminant prime to that prime. The three quadratic
subfields of $\mathbb Q(\sqrt{-q},\sqrt{-d})$ have discriminants
$-q,-d,qd$, respectively, so none is contained in $B_h$.
The two fields are therefore linearly disjoint: a nontrivial
intersection with a biquadratic extension would contain a
quadratic subfield. This proves the factor $4$ in the degree.
Finally $E_h$ is abelian and contains the imaginary quadratic
field $\mathbb Q(\sqrt{-q})$. Complex conjugation has order two,
and its fixed field is totally real, so $E_h$ is CM.
\end{proof}

The displayed $Q_h$ is a permissible comparison modulus. If a
construction uses a modulus $Q$ with
$2^{h+2}\mid Q$, $M\mid Q$, and $Q\mid Q_h$, then necessarily
$Q=Q_h$. In particular its odd part is the fixed odd part of $M$.
There is no need to select a new family of base characters modulo
the growing power of $2$.

\begin{lemma}\label{sa-character-families}
Characters of $\operatorname{Gal}(E_h/\mathbb Q)$ trivial on
$\operatorname{Gal}(E_h/K_\psi(\sqrt{-q},\sqrt{-d}))$ are exactly
\[
             \psi^r\chi_q^i\chi_d^j,
 \qquad 0\le r<o,\quad i,j\in\{0,1\}.
\]
They are understood through their primitive Dirichlet inducers,
and $\chi_q(n)=(n/q)$, $\chi_d(n)=(n/d)$.
The family with $(i,j)\ne(0,0)$ has $3o$ distinct members.
Every primitive real character of conductor dividing $Q_hqd$
has the form $\rho\chi_q^i\chi_d^j$, where $\rho$ is primitive
real and $\operatorname{cond}\rho\mid8M$.
\end{lemma}

\begin{proof}
The preceding disjointness gives
$\operatorname{Gal}(K_\psi(\sqrt{-q},\sqrt{-d})/\mathbb Q)
\simeq\operatorname{Gal}(K_\psi/\mathbb Q)\times C_2\times C_2$.
The first factor is cyclic with character group generated by $\psi$.
Taking characters proves the first assertion and distinctness.
If the primitive conductor of $\psi^r$ is $f_r\mid M$, the primitive
conductor after twisting is exactly $f_rq^id^j$, because the three
conductors are coprime. The local factors at omitted primes are
therefore not silently retained in a primitive sum.

For the real-character assertion, a homomorphism from the units
modulo an odd prime power $p^b$ to $\{\pm1\}$ kills the subgroup
of order $p^{b-1}$ and factors modulo $p$. At $2$, every unit
congruent to $1$ modulo $8$ is a square modulo $2^b$. This follows
by induction: if an odd $v$ solves $v^2\equiv u\pmod{2^n}$,
$n\ge3$, replacing $v$ by $v+2^{n-1}$ toggles the next binary
digit of its square. Thus a real character at $2$ factors modulo
$8$. Apply the Chinese remainder theorem to $Q_hqd$. At $q,d$
the real local character is either trivial or the Legendre symbol;
the remaining primitive conductor divides $8M$.
\end{proof}

The characters requiring the square-root estimate in the construction
are among the first family with $(i,j)\ne(0,0)$: the untwisted powers
of $\psi$ are the retained characters. The real-character assertion
also verifies that the zero-free-box requirement involves a fixed
base family. Its primitive conductor is at most
$8M\exp(z+z')$ when $q\le e^z$, $d\le e^{z'}$.
Theorem~\ref{qs-selection} supplies these boxes at every height.

\subsection{Size bounds with the selection scale kept fixed}

Put $\alpha=2+8\epsilon$, where $0<\epsilon<1/52$.
Fix $J\ge1$ and write $U=X^{1/J}$ and $z=U^{\epsilon/5}$.
The prime selection is made at $U$, not at the later variable $x$.
For the selected primes one has
\begin{equation}\label{sa-prime-ranges}
 \frac z2<\log q\le z,\qquad
 \frac{z(\log U)^2}{2}<\log d\le z(\log U)^2.
\end{equation}
For all sufficiently large $X$, these imply
\begin{equation}\label{sa-size-comparison}
 \log(qd)\le X^{9\epsilon/(40J)},\qquad
 q,d>2X^{401}+1,
\end{equation}
and $q,d$ exceed any prescribed fixed constant.
Indeed $\log(qd)\le U^{\epsilon/5}(1+(\log U)^2)$, while
$9/40-1/5=1/40$. For the lower bounds,
$q,d>\exp(X^{\epsilon/(5J)}/2)$ once $\log U\ge1$;
this exceeds every fixed power of $X$.
One explicit verification uses
$\delta=\epsilon/(5J)$,
$n=\lceil402/\delta\rceil$ and $T=2^nn!$:
\[
 \exp(X^\delta/2)\ge X^{\delta n}/T\ge X^{402}/T.
\]
Taking $X>T(B+3)$ and $X\ge e^J$ gives all the bounds above
with an additional prescribed bound $B\ge1$.

If $U\le x\le X$ and $x^\alpha\le c2^h\le2x^\alpha$,
then $Q_h\le8Mx^\alpha$. Consequently, uniformly in such $x,h$,
\begin{equation}\label{sa-conductor}
 \log N_h\le2U^{\epsilon/5}(\log U)^2,
 \qquad
 \log(N_ht)\le2x^{\epsilon/4}\quad(1\le t\le x^{401})
\end{equation}
for sufficiently large $X$. For the first inequality, use
$\log x\le J\log U$ and absorb
$\log(8M)+\alpha J\log U+z$ into $z(\log U)^2$.
For the second use \eqref{sa-size-comparison}, $U\le x$, and
$9\epsilon/40<\epsilon/4$; the remaining logarithm is at most
$\log(8M)+(\alpha+401)\log x$.
Every threshold here depends only on the fixed data.

The congruences are also compatible without additional conditions:
$d\equiv7\pmod{16}$ and $d\equiv-1$ modulo $q$ and each odd
prime at most $z$ are reduced classes at pairwise coprime moduli.
For large $U$, $z\ge4M$ and $q>3$. Each odd prime $\ell\le z$
and the prime $q$ split in $\mathbb Q(\sqrt{-d})$, since
$-d\equiv1\pmod\ell$. The prime $2$ splits as well, since
$-d\equiv1\pmod8$. Requiring $d\equiv-1$ modulo higher prime
powers dividing $M$ is unnecessary for these splitting conditions.

\subsection{A simultaneous interval with the exact rescaling}

\begin{proposition}\label{sa-compatible-interval}
Assume \eqref{qs-qrh} and \eqref{ct-coarse-hypothesis}.
Fix $\psi$, a finite set $P$ of odd primes, and fixed abelian fields
$K_p\supset\mu_p$ for $p\in P$. Choose the generator bounds
$R_2,R_p\ge1$ of Proposition~\ref{ct-uniform-generators}, with
base field $K_\psi$, and put
\[
 \epsilon_0=\frac{9\epsilon}{10},\quad
 R_*=R_2+\sum_{p\in P}R_p,\quad
 \eta=\frac{\epsilon_0}{2C_K},\quad
 J_0=2^{R_*+1}(401/\eta)^{R_*}.
\]
Fix any $J\ge J_0$ and any integer $h_0\ge h_*$.
Set $\widetilde j_p=\lfloor\log_pR_p\rfloor+3$.
There are arbitrarily large $Z$ for which one can choose a single
$U>1$ and a single pair $q,d$ satisfying
\eqref{sa-prime-ranges}, the congruences of
Theorem~\ref{qs-selection}, and $(qd,2M)=1$, such that the following
hold simultaneously for every $x\in[Z,Z^{3/2}]$.

There is an integer $h=h(x)\ge h_0$ with
\begin{equation}\label{sa-rank-window}
 x^{1/J}\le U\le x,\qquad
 x^{2+8\epsilon}\le c2^h\le2x^{2+8\epsilon}.
\end{equation}
All twists specified in Lemma~\ref{sa-character-families} with a
nonzero quadratic factor obey \eqref{qs-bounds} at $y=x$.
The conductor bounds \eqref{sa-conductor} hold, and for every
integer $1\le t\le x^{400}$,
\begin{equation}\label{sa-partial-cohomology}
 \sum_{1\le e<t}\lambda_2(E_h,e)
 +\sum_{p\in P}\sum_{1\le e<t}
              \lambda_p((K_p)_{\widetilde j_p},e)
       \le \exp(x^{4\epsilon/5}).
\end{equation}
All lengths in this formula are finite. The primes $q,d$ exceed
$2x^{401}+1$ and any further fixed bound specified in advance.
\end{proposition}

\begin{proof}
All constants in the statement are fixed before $X$ is chosen.
Apply Theorem~\ref{qs-selection} at $U=X^{1/J}$ with its interval
exponent equal to $J$. The character bounds then hold throughout
$[U,U^J]=[U,X]$. By \eqref{sa-size-comparison},
\[
 \log(qd)\le X^{\epsilon_0/(4J)}
              \le X^{\epsilon_0/(4J_0)}.
\]
Proposition~\ref{ct-conditional-interval}, applied with
$\epsilon_0$, $C_t=401$, and $A=6$, gives $W$ with
$X^{1/J_0}\le W$, $W^2\le X$, and its cohomological estimates
on the whole interval $[W,W^2]$. In particular $W\ge U$.

Let $k=c^{1/\alpha}\ge1$, and for $s\ge1$ put
$e(s)=\lceil\alpha\log_2s\rceil$. Then
$s^\alpha\le2^{e(s)}\le2s^\alpha$.
Increase $X$ until $U\ge k^{400}$, $U\ge k^3$, and
$e(U)\ge h_0$. Set $Z=kW$. For $Z\le x\le Z^{3/2}$,
\[
 x\le(kW)^{3/2}\le W^2\le X,\qquad
 W\le x'=x/k\le W^2.
\]
Take $h=e(x')$. Then $h\ge h_0$ and
$x^\alpha=c(x')^\alpha\le c2^h\le2x^\alpha$.
Moreover $x^{1/J}\le X^{1/J}=U\le W\le x$.
Since $\alpha<3$ and $\log2>1/2$, one has
$h\le1+6\log x'\le6(1+\log x')$, as required by the
cohomological proposition. Finally
\[
                  x^{400}=k^{400}(x')^{400}\le(x')^{401}.
\]
Thus its estimate applies to every weight under consideration.

The layer used there is
$j_p=1+\lceil\log_pR_p\rceil$ (with $j_p=1$ for $R_p=1$).
One has $0\le\widetilde j_p-j_p\le2$.
Lemma~\ref{ct-growth} increases each odd-prime length by at
most $2R_p$ when passing to $\widetilde j_p$.
Divide the weighted estimate in
\eqref{ct-summed-bound} by $\log2$, the smallest of its prime
weights. Its left side dominates the corresponding unweighted
sum. After the layer changes the result is bounded by
\[
 \frac{\exp((x')^{18\epsilon/25})}{\log2}
       +2x^{400}\sum_{p\in P}R_p.
\]
Because $18/25<4/5$, this is at most
$\exp(x^{4\epsilon/5})$ for all sufficiently large $x$.
This proves \eqref{sa-partial-cohomology}, including finiteness.
The character and conductor conclusions follow from $[Z,Z^{3/2}]
\subset[U,X]$. Taking $X$ arbitrarily large makes $W$, hence $Z$,
arbitrarily large.
\end{proof}

The auxiliary exponent $\epsilon_0$ supplies a strict margin for the
fixed-layer change; it does not change the exponent in the prime
selection or in the rank window. The finite fields and generator bounds
are fixed before $J_0$. The potentially very large constant $J$ is
therefore an allowed fixed interval exponent in
Theorem~\ref{qs-selection}.

\subsection{The full cohomological bound and the remaining input}

For a CM field $E$ and an odd prime $p$, let $\iota$ be the action
of complex conjugation on $\mathcal S_p(E,e)$. Its two parts relevant
here are the eigenspaces with eigenvalues $(-1)^e$ and $(-1)^{e+1}$;
write their lengths as $\lambda_p^+(E,e)$ and
$\lambda_p^-(E,e)$. Since $2$ is a unit in $\mathbb Z_p$, the
projections $(1\pm(-1)^e\iota)/2$ give a direct sum and hence
\begin{equation}\label{sa-length-decomposition}
       \lambda_p(E,e)=\lambda_p^+(E,e)+\lambda_p^-(E,e).
\end{equation}
This identity also holds for infinite lengths. Finiteness must be
established before ordinary finite sums of lengths are used.

The precise additional estimates needed to complete the
cohomological input are, for every finite set $S$ of odd primes,
\begin{align}\label{sa-missing-parts}
 \sum_{p\in S}(\log p)\sum_{1\le e<t}\lambda_p^+(E_h,e)
       &\le (N_ht)^C+C\exp(x^{4\epsilon/5}),\notag\\
 \sum_{p\in S}(\log p)\sum_{1\le e<t}\lambda_p^-(E_h,e)
       &\le K\exp(8x^{\epsilon/4}),
\end{align}
with fixed $C,K\ge0$ and with each individual length finite.
The source obtains the first from its arithmetic counting assertions
using the fixed odd-prime bounds, and the second from its separate
estimate for the other eigenspaces. The first line is not a statement
about only the finite set $P$: it controls every odd prime.
Neither line is proved by Proposition~\ref{sa-compatible-interval}.

\begin{corollary}\label{sa-full-bound}
If \eqref{sa-missing-parts} holds for the fields and scales in
Proposition~\ref{sa-compatible-interval}, then, for every finite
set $S$ of primes and $1\le t\le x^{400}$,
\[
 \sum_{p\in S}(\log p)\sum_{1\le e<t}\lambda_p(E_h,e)
       \le \sqrt{\exp(x^\epsilon)}
\]
for all sufficiently large $x$ in those intervals. Every individual
length occurring here is finite.
\end{corollary}
\begin{proof}
Use \eqref{sa-length-decomposition},
\eqref{sa-partial-cohomology}, and \eqref{sa-conductor}.
The left side is at most
\[
 (\log2+C)\exp(x^{4\epsilon/5})
       +\exp(2C x^{\epsilon/4})+K\exp(8x^{\epsilon/4}).
\]
Put $b=x^{4\epsilon/5}$, $s=x^{\epsilon/5}$,
$D=9+2C$, and $A=2+C+K$.
For $x\ge1$ the last display is at most $A\exp(Db)$.
If $s\ge2(D+\log A+1)$, then
$\log A+Db\le bs/2=x^\epsilon/2$, since $b\ge1$.
For example $x\ge[2(D+\log A+1)]^{5/\epsilon}$ suffices.
The stated finiteness follows from the hypotheses on the two
parts and the two-adic estimate.
\end{proof}

There is a useful separate check of the original rescaling. Suppose
one already has the full cohomological interval statement, for a fixed
function $e(s)$ satisfying
$s^\alpha\le2^{e(s)}\le2s^\alpha$, with weights below $s^{401}$
and bound $\sqrt{\exp(s^\epsilon)}$. Its quantifiers are:
a fixed $J$ first; then every sufficiently large $X$; then every
permissible prime pair at $U=X^{1/J}$; finally an interval
$[W,W^2]$ with $U\le W$, $W^2\le X$.
Select the prime pair on $[U,X]$ before invoking this statement.
The substitution $k=c^{1/\alpha}$, $Z=kW$, $s=x/k$ used above
then proves the full cohomological bound for weights below $x^{400}$
throughout $[Z,Z^{3/2}]$: the weight inclusion is
$x^{400}\le s^{401}$, and
$\sqrt{\exp(s^\epsilon)}\le\sqrt{\exp(x^\epsilon)}$.
This is exactly the role of the exponents $401$, $400$, and $3/2$.
No separate intersection of prime-selection and cohomological
intervals is taken.

One tempting change of scale is false. If a pair satisfies
\eqref{sa-prime-ranges} at $U$ and $x=U^J$ with fixed $J>1$,
then, for sufficiently large $U$,
\[
          \log q\le U^{\epsilon/5}
                <\tfrac12 x^{\epsilon/5}.
\]
Thus the pair need not satisfy the required lower window at $x$.
The construction retains $U$ as its selection scale and uses only
$x^{1/J}\le U\le x$, exactly as in
\eqref{sa-rank-window}. Likewise, the character family is fixed
through $\psi$; replacing it by all characters of the growing
$2$-power modulus would be a stronger, unproved selection assertion.


% Conventional reconstruction of the actual coefficient-lattice duality.
% All labels in this fragment have the prefix coef:.
\section{Integral polynomial lattices and normalized Fock duality}
\label{coef:section}

We now construct the coefficient lattices used at the beginning of the
finite-group argument in Section~\ref{div:section}, and prove their
duality. The submitted construction uses polynomial dual-Weyl and Weyl
lattices, then transports them by local bases.\cite{alpo-uniform-source}
The proof below uses Gaussian elimination, divided row operations and
the elementary theory of lattices. In particular, the integral duality
will be proved here, rather than imported from a theorem on tableaux.

\subsection{The polynomial modules and their normalization}

Fix integers $0\le n\le r$ and a partition
$\nu=(\nu_1,\ldots,\nu_n)$, so
$\nu_1\ge\cdots\ge\nu_n\ge0$. Let $X=(x_{as})$ be an $r$ by $n$
matrix of independent variables. For a characteristic-zero domain $R$,
let $\nabla_\nu(R)$ consist of the polynomials in $R[X]$ which are
homogeneous of degree $\nu_s$ in the variables of each column $s$ and satisfy
\begin{equation}\label{coef:column-invariance}
 F\bigl(X(I+tE_{ij})\bigr)=F(X)\qquad(1\le i<j\le n)
\end{equation}
as polynomial identities in the additional indeterminate $t$.
Here $E_{ij}$ is the elementary matrix: the substitution adds $t$ times
column $i$ to column $j$. Identity \eqref{coef:column-invariance} is an
identity of polynomials, not merely a condition at the elements of a
residue field. Write $V_\nu=\nabla_\nu(\mathbb Q)$.

Let $\Delta_j(X)$ be the determinant of the first $j$ rows and first
$j$ columns, put $\nu_{n+1}=0$, and define
\[
 H_\nu(X)=\prod_{j=1}^n\Delta_j(X)^{\nu_j-\nu_{j+1}}.
\]
The polynomial $H_\nu$ belongs to $\nabla_\nu(\mathbb Z)$: every
initial minor is fixed by the indicated column additions, and the
degrees telescope to $\nu_s$. Its row degrees are
$(\nu_1,\ldots,\nu_n,0,\ldots,0)$.
Let $I_{r,n}$ denote the first $n$ columns of the identity matrix.
Then $H_\nu(I_{r,n})=1$.

For $1\le a,b\le r$ define the divided row operator
\[
 D_{ab}^{(k)}F=[t^k]F\bigl((I+tE_{ab})X\bigr),\qquad k\ge0.
\]
Thus the substitution changes row $a$ to row $a+t\,b$ when $a\ne b$.
For $a=b$ it multiplies row $a$ by $1+t$.
All these operators preserve integral coefficients. They preserve
$\nabla_\nu(R)$ because row substitutions commute with the right
column substitutions, and preserve column degrees. Let
\[
 \Delta_\nu(R)=\operatorname{span}_R
 \{D_1\cdots D_mH_\nu:m\ge0,\ D_i=D_{a_i b_i}^{(k_i)}\}.
\]
The empty product is allowed. In particular
$\Delta_\nu(R)\subset\nabla_\nu(R)$.
These are the polynomial Weyl and dual-Weyl lattices needed here.

For $F=\sum_\alpha F_\alpha X^\alpha$ and
$G=\sum_\alpha G_\alpha X^\alpha$ over $\mathbb Q$, define the
bilinear Fock pairing and its normalization by
\begin{equation}\label{coef:fock}
 [F,G]=\sum_\alpha\alpha!F_\alpha G_\alpha,
 \qquad C_\nu=[H_\nu,H_\nu],\qquad
 B_\nu(F,G)=\frac{[F,G]}{C_\nu},
 \qquad \alpha!=\prod_{a,s}\alpha_{as}!.
\end{equation}
The integer $C_\nu$ is positive: the Fock form is positive definite
on real polynomials, and $H_\nu(I_{r,n})=1$.
Thus the denominator in \eqref{coef:fock} is nonzero, including when
all parts of $\nu$ vanish. This explicit finite formula specifies the
normalizing scalar completely; a product formula for it is unnecessary.

\subsection{Integral coefficients from Gaussian elimination}

\begin{lemma}[Restriction to lower triangular matrices]\label{coef:restriction}
Let $F\in V_\nu$. Its restriction to matrices with $x_{as}=0$ for
$a<s$ determines $F$. If the restricted polynomial has integer
coefficients, then $F\in\mathbb Z[X]$.
\end{lemma}
\begin{proof}
The case $n=0$ concerns constant polynomials and is immediate. Assume
$n>0$ and put $\Delta=\prod_{j=1}^n\Delta_j$.
In the ring $A=\mathbb Z[X,\Delta^{-1}]$, Gaussian elimination gives
\[
                         X=LU,
\]
where $L$ is lower triangular in the rectangular sense
($L_{as}=0$ if $a<s$) and $U$ is upper triangular with diagonal entries
one. All entries belong to $A$.
For completeness, eliminate the entries above the diagonal in column
$j$, using the already treated columns $1,\ldots,j-1$ in this order.
The pivot in column $i$ is $\Delta_i/\Delta_{i-1}$, with
$\Delta_0=1$: these column operations leave each relevant initial
minor unchanged, and a lower triangular initial block has determinant
equal to the product of its pivots. Each pivot is a unit in $A$.
This proves the asserted factorization by induction, without dividing
by an integer.

An upper triangular unipotent matrix is a product of elementary
matrices $I+tE_{ij}$, $i<j$: successively clear its entries above the
diagonal by column additions. Therefore
\eqref{coef:column-invariance}, applied over the localized coefficient
ring, gives $F(X)=F(L)$. A zero restriction consequently implies $F=0$.
If the restriction has integer coefficients, it gives
$F\in A$.

We claim that $\mathbb Q[X]\cap A=\mathbb Z[X]$.
Write $F=G/m$ with $G\in\mathbb Z[X]$, $m>0$, and with no prime
dividing both $m$ and every coefficient of $G$. Membership in $A$
gives $\Delta^kG=mQ$ for some $k\ge0$ and $Q\in\mathbb Z[X]$.
For any prime $p\mid m$, the polynomial $\Delta$ is nonzero modulo
$p$: each initial determinant contains its diagonal monomial with
coefficient one, and a polynomial ring over $\mathbb F_p$ is a domain.
Reduction of the equality modulo $p$ therefore forces every coefficient
of $G$ to be divisible by $p$, a contradiction. Thus $m=1$.
\end{proof}

\begin{lemma}[The extremal row weight]\label{coef:extremal-weight}
The subspace of $V_\nu$ of row degree
$(\nu_1,\ldots,\nu_n,0,\ldots,0)$ is $\mathbb QH_\nu$.
Moreover, for every $F\in V_\nu$,
\begin{equation}\label{coef:evaluation-pairing}
                         B_\nu(H_\nu,F)=F(I_{r,n}).
\end{equation}
\end{lemma}
\begin{proof}
Projecting onto a fixed row degree preserves the column-invariance
identities, because every term of a column substitution preserves the
degree in each row. Let $F_0$ be the component of $F$ of the displayed
row degree.

A monomial with these row degrees and column degrees $\nu$, and
supported on entries with $a\ge s$, must be
\[
                         \prod_{s=1}^n x_{ss}^{\nu_s}.
\]
Indeed, row one can use only column one, so its prescribed degree
already exhausts that column; induction gives the assertion for every
row and column. The restriction of $H_\nu$ is the same diagonal
monomial. Lemma~\ref{coef:restriction} shows that
$F_0=cH_\nu$. Evaluation at $I_{r,n}$ leaves only the diagonal
monomial, so $c=F(I_{r,n})$.
Distinct row degrees are orthogonal for the Fock pairing, which is
diagonal on monomials. Since $H_\nu$ has the indicated row degree,
$[H_\nu,F]=[H_\nu,F_0]=cC_\nu$, proving
\eqref{coef:evaluation-pairing}.
\end{proof}

\begin{lemma}[Adjoint row operators]\label{coef:adjoint}
For a matrix $g\in\operatorname{Mat}_r(\mathbb Q)$, put
$T_gF(X)=F(gX)$. Then
\[
 [T_gF,G]=[F,T_{g^t}G],\qquad
 [D_{ab}^{(k)}F,G]=[F,D_{ba}^{(k)}G].
\]
The first identity holds over every characteristic-zero field after
extending the pairing, and applies independently to several disjoint
blocks of variables.
\end{lemma}
\begin{proof}
Constants satisfy the assertion directly. On polynomials of homogeneous
degree $m\ge1$ in any finite set of variables,
the multinomial theorem gives
\[
 [ (u\mathbin{\cdot}x)^m,(v\mathbin{\cdot}x)^m]
                         =m!(u\mathbin{\cdot}v)^m.
\]
The powers of linear forms span that space in characteristic zero.
One direct verification is the polarization identity
\[
 m!y_1\cdots y_m
  =\sum_{S\subset\{1,\ldots,m\}}(-1)^{m-|S|}
                               \left(\sum_{i\in S}y_i\right)^m:
\]
expanding the right side cancels every term missing one of the $m$
symbols and leaves the product term with coefficient $m!$.
Repetitions among the $y_i$ give every monomial. Substituting a matrix
in a linear form transposes its coefficient vector, so the first
identity follows from the displayed Fock formula and bilinearity.
Different total degrees are orthogonal. Apply this to the block matrix
which acts on the row variables of every column. Finally replace $g$
by $I+tE_{ab}$ and compare coefficients of $t^k$. This includes
$a=b$; in that case the divided operator multiplies a row-homogeneous
polynomial of row degree $d$ by $\binom dk$.
\end{proof}

For a word $D=D_1\cdots D_m$ let
$D^*=D_m^*\cdots D_1^*$, where
$(D_{ab}^{(k)})^*=D_{ba}^{(k)}$.
Lemmas~\ref{coef:extremal-weight} and \ref{coef:adjoint} give the
key identity
\begin{equation}\label{coef:word-evaluation}
                 B_\nu(F,DH_\nu)=(D^*F)(I_{r,n})
                       \qquad(F\in V_\nu).
\end{equation}

\subsection{Exact integral duality}

For a subgroup $L\subset V_\nu$ define
\[
 L^{\vee,B}=\{F\in V_\nu:B_\nu(F,L)\subset\mathbb Z\}.
\]

\begin{theorem}[Integral Weyl duality]\label{coef:integral-duality}
For every $0\le n\le r$ and every partition $\nu$ of length $n$,
both $\nabla_\nu(\mathbb Z)$ and $\Delta_\nu(\mathbb Z)$ are full
lattices in $V_\nu$, and
\begin{equation}\label{coef:duality}
 \Delta_\nu(\mathbb Z)^{\vee,B}=\nabla_\nu(\mathbb Z),
 \qquad
 \nabla_\nu(\mathbb Z)^{\vee,B}=\Delta_\nu(\mathbb Z).
\end{equation}
Thus $B_\nu$ gives a perfect integral pairing between these two
lattices. No restriction on the primes dividing $C_\nu$ is required.
\end{theorem}
\begin{proof}
Equation~\eqref{coef:word-evaluation} shows that the normalized pairing
of an integral polynomial in $\nabla_\nu$ with every generator of
$\Delta_\nu(\mathbb Z)$ is integral. This proves one inclusion in
the first equality.

Conversely, suppose $F\in V_\nu$ pairs integrally with
$\Delta_\nu(\mathbb Z)$. Then $(DF)(I_{r,n})\in\mathbb Z$ for
every word of divided row operators, since transposition and reversal
permute the set of all words. Consider the matrix $L(z)$ with diagonal
entries one, entries $z_{as}$ below the diagonal, and zero entries
above the diagonal. It is the first $n$ columns of a product of
elementary row matrices. More explicitly,
\[
 L(z)=\left(\prod_{s=1}^n\prod_{a>s}(I+z_{as}E_{as})\right)I_{r,n},
\]
where the factors are written from left to right in increasing $s$;
the order for fixed $s$ is immaterial. Products of two of the
off-diagonal summands vanish in this order: for a factor in column
$s$ followed by one in column $t\ge s$, a nonzero product would
require $s=b>t$. Consequently the displayed product has exactly the
specified entries. Expanding the substitutions one factor at a time
shows that every coefficient of $F(L(z))$ is $(DF)(I_{r,n})$ for an
appropriate word $D$, and hence is an integer.

Column homogeneity now recovers the restriction to an arbitrary lower
triangular matrix. Write its diagonal entries as $d_s$ and its lower
entries as $l_{as}$. In the Laurent polynomial ring in the $d_s$,
\[
 F((l_{as}))=\left(\prod_s d_s^{\nu_s}\right)
              F\bigl(L((l_{as}/d_s)_{a>s})\bigr).
\]
Every coefficient on the right is an integer. For each off-diagonal
monomial, the column degree determines its diagonal powers uniquely;
there is no collision of distinct terms. The left side is a polynomial,
so every surviving diagonal exponent is nonnegative. Its coefficients
are therefore integers. Lemma~\ref{coef:restriction} gives
$F\in\nabla_\nu(\mathbb Z)$, proving the first equality.

The space of polynomials of the fixed column degrees has a finite
monomial basis. Hence $\nabla_\nu(\mathbb Z)$ is a finitely generated
free abelian group. It spans $V_\nu$, since denominators in a rational
polynomial can be cleared and the defining identities have integral
coefficients. Its subgroup $\Delta_\nu(\mathbb Z)$ is also finitely
generated. It spans $V_\nu$ as well: otherwise nondegeneracy of $B_\nu$
would give a nonzero rational vector $F$ perpendicular to its rational
span. Then every rational multiple of $F$ would belong to
$\Delta_\nu(\mathbb Z)^{\vee,B}=\nabla_\nu(\mathbb Z)$, which is
impossible for a nonzero polynomial with integral coefficients.

Every full lattice in a finite-dimensional rational space with a
nondegenerate symmetric rational bilinear form equals its double dual.
Indeed, take a basis of the lattice; in these coordinates its dual is
the inverse Gram matrix applied to $\mathbb Z^q$, and dualizing once
more returns $\mathbb Z^q$ (the form here is symmetric).
Dualizing the first equality proves the second one. The same coordinate
calculation proves perfection. If $n=0$, both lattices are
$\mathbb Z$, $H_\nu=C_\nu=1$, and these assertions also hold directly.
\end{proof}

\begin{remark}[Saturation has two different roles]\label{coef:saturation}
The lattice $\nabla_\nu(\mathbb Z)$ is saturated in the group of
integral polynomials of column degree $\nu$: if a nonzero integer
multiple of an integral polynomial satisfies the defining identities,
the polynomial itself satisfies them. The smaller lattice
$\Delta_\nu(\mathbb Z)$ need not be saturated in
$\nabla_\nu(\mathbb Z)$.
For binary quadratics, $H=x^2$, $C=2$,
\[
 \nabla=\mathbb Zx^2\oplus\mathbb Zxy\oplus\mathbb Zy^2,
 \qquad
 \Delta=\mathbb Zx^2\oplus2\mathbb Zxy\oplus\mathbb Zy^2.
\]
Replacing $\Delta$ by its saturation would destroy the asserted
integral pairing, since $B(xy,xy)=1/2$.
Also, the class of $2xy$ in $\Delta/2\Delta$ is nonzero although its
image in $\nabla/2\nabla$ is zero. Reducing the inclusion of the
lattices modulo two is not the same as identifying their reductions.
The perfect pairing between the two abstract integral lattices remains
valid at this prime. This explains why a small residue characteristic
does not invalidate Theorem~\ref{coef:integral-duality}, and why the
torsion in $\nabla/\Delta$ must not be discarded.
\end{remark}

\subsection{Number fields, determinant characters and tensor products}

Let $K$ be a number field with ring of integers $\mathcal O$.
The defining conditions for $\nabla_\nu$ are finitely many homogeneous
linear equations with integer coefficients in the finite monomial
basis of column degree $\nu$. Passing from $\mathbb Z$ to $\mathcal O$
or $K$ commutes with their common kernel. One elementary proof puts
the integer matrix of these equations in Smith normal form: each
nonzero diagonal integer acts injectively on either characteristic-zero
domain, so the kernel is spanned by the same remaining basis columns.
Consequently
\[
 \nabla_\nu(\mathcal O)=\mathcal O\otimes_{\mathbb Z}
                 \nabla_\nu(\mathbb Z),\qquad
 \Delta_\nu(\mathcal O)=\mathcal O\otimes_{\mathbb Z}
                 \Delta_\nu(\mathbb Z).
\]
The second identity also follows immediately from the integral formulas
for the divided operators. Extending the two dual integral bases from
Theorem~\ref{coef:integral-duality} proves the same perfect pairing over
$\mathcal O$, with duality taken inside $\nabla_\nu(K)$.

The actual coefficient weights are dominant integral weights
$\lambda_1\ge\cdots\ge\lambda_n$; the $\lambda_i$ may be negative.
For $n>0$ put
\[
 a=\lambda_n,\qquad \nu_i=\lambda_i-a.
\]
Thus $\nu$ is a partition with last part zero. For $n=0$ put $a=0$.
On the square polynomial space define the representation
\begin{equation}\label{coef:representation}
 \rho_\lambda(g)F(X)=\det(g)^a F(g^tX),\qquad g\in\mathrm{GL}_n(K).
\end{equation}
It is a representation, since substitutions reverse matrix
multiplication and the transpose reverses it again. Negative powers
of $\det g$ are legitimate because $g$ is invertible.
Lemma~\ref{coef:adjoint} gives
\begin{equation}\label{coef:contravariant}
 B_\nu(\rho_\lambda(g)F,G)
       =B_\nu(F,\rho_\lambda(g^t)G),\qquad
 B_\nu(\rho_\lambda(g^{-t})F,\rho_\lambda(g)G)=B_\nu(F,G).
\end{equation}

Both standard lattices are preserved by
$\mathrm{GL}_n(\mathcal O)$ and by
$\mathrm{GL}_n(\mathcal O_{\mathfrak p})$ after localization.
For $\nabla$ this follows directly from integral substitution and the
fact that a positive or negative power of a unit determinant is a unit.
For $\Delta=\nabla^{\vee,B}$, it follows from
\eqref{coef:contravariant} and the stability of $\nabla$ under $g^t$;
apply the argument also to $g^{-1}$ to get equality.

For a finite index set $\mathcal T$, take a weight $\lambda_\tau$
for each $\tau$ and use disjoint variable blocks. Let
\[
 V=\bigotimes_{\tau\in\mathcal T}\nabla_{\nu_\tau}(K),\qquad
 N_0=\bigotimes_\tau\nabla_{\nu_\tau}(\mathcal O),\qquad
 D_0=\bigotimes_\tau\Delta_{\nu_\tau}(\mathcal O).
\]
The first tensor product is over $K$ and the latter two are over
$\mathcal O$. They are identified with products of polynomials in
the disjoint blocks. Monomial factorials factor over the blocks, so
the tensor-product pairing is exactly the Fock pairing divided by
\[
                             C=\prod_\tau C_{\nu_\tau}.
\]
Tensoring the dual bases proves
\begin{equation}\label{coef:standard-tensor-duality}
                  N_0^{\vee,B}=D_0,\qquad D_0^{\vee,B}=N_0.
\end{equation}
The representation $\rho(g)$ is the tensor product of
\eqref{coef:representation}. Formula~\eqref{coef:contravariant}
continues to hold componentwise. For an empty index set the space is
$K$, both lattices are $\mathcal O$, and $C=1$.

The integral bases constructed above give rational polynomial bases of
$V$. The matrix of $B$ in such a basis has rational entries. If all
entries of the matrices $g_\tau$ lie in a subfield $K_0\subset K$,
the matrix of $\rho(g)$ in the same basis also has entries in $K_0$:
the polynomial substitution and determinant powers have coefficients
in $K_0$, and expressing their outputs in a rational basis uses only
linear equations with rational coefficients. This gives the required
field of definition without making a special tableau-basis choice.

\subsection{Transport by the actual local bases}

Let $M_\tau\subset K^n$ be full $\mathcal O$-lattices. For every
nonzero prime $\mathfrak p$, choose matrices
$B_{\tau,\mathfrak p}\in\mathrm{GL}_n(K)$ whose columns form bases
of $(M_\tau)_{\mathfrak p}$ over $\mathcal O_{\mathfrak p}$.
Put $B_{\mathfrak p}=(B_{\tau,\mathfrak p})_\tau$ and define
\begin{equation}\label{coef:transport}
 N(M)=\bigcap_{\mathfrak p}\rho(B_{\mathfrak p})(N_0)_{\mathfrak p},
 \qquad
 D(M)=\bigcap_{\mathfrak p}\rho(B_{\mathfrak p})(D_0)_{\mathfrak p}
                         \quad\text{inside }V.
\end{equation}

\begin{lemma}[Well-defined transported lattices]\label{coef:transport-lattices}
The intersections in \eqref{coef:transport} are independent of the
chosen local bases. They are full lattices, their localizations are
the corresponding displayed local modules, and $D(M)\subset N(M)$.
For a tuple $g=(g_\tau)$ of invertible matrices,
\begin{equation}\label{coef:equivariance}
                  N(gM)=\rho(g)N(M),\qquad D(gM)=\rho(g)D(M).
\end{equation}
If $M_\tau^\vee$ denotes the ordinary bilinear dual lattice in $K^n$,
then
\begin{equation}\label{coef:transport-duality}
                     N(M)^{\vee,B}=D(M^\vee),
              \qquad D(M)^{\vee,B}=N(M^\vee).
\end{equation}
\end{lemma}
\begin{proof}
Two bases differ on the right by an element of
$\mathrm{GL}_n(\mathcal O_{\mathfrak p})$, which preserves the two
standard localized lattices. This proves independence.
Outside a finite set of primes the base lattices are
$\mathcal O_{\mathfrak p}^n$, so the identity matrices can be chosen.
The local modules in \eqref{coef:transport} therefore equal the
standard ones outside a finite set.

Here is the elementary patching argument. Given finitely many
exceptional local lattices and a reference lattice $L_0$, choose a
nonzero $b\in\mathcal O$ such that each lies between
$b(L_0)_{\mathfrak p}$ and $b^{-1}(L_0)_{\mathfrak p}$.
The finite module $b^{-1}L_0/bL_0$ is the direct sum of its primary
parts. In each part take the submodule specified by the corresponding
local lattice; at the additional primes dividing $b$ take the
submodule specified by $L_0$. Its inverse image in $b^{-1}L_0$ is a
full lattice with exactly the prescribed localizations. A lattice
equals the intersection of its localizations: in any finite generating
set, a failure of membership is detected in a primary component of
a finite quotient after clearing denominators. This proves both the
existence assertion and its equality with the intersections above.
The inclusion follows from $D_0\subset N_0$.
The matrices $g_\tau B_{\tau,\mathfrak p}$ are bases of
$(g_\tau M_\tau)_{\mathfrak p}$, proving
\eqref{coef:equivariance}.

Ordinary duality of a full lattice commutes with localization. For
example, an element pairing locally integrally with a finite generating
set can be multiplied by a single element outside $\mathfrak p$ so
that it pairs globally integrally with all generators. The bases of
$(M_\tau^\vee)_{\mathfrak p}$ are $B_{\tau,\mathfrak p}^{-t}$.
Equations~\eqref{coef:contravariant} and
\eqref{coef:standard-tensor-duality} therefore give
\[
 \bigl(\rho(B_{\mathfrak p})(N_0)_{\mathfrak p}\bigr)^{\vee,B}
     =\rho(B_{\mathfrak p}^{-t})(D_0)_{\mathfrak p}.
\]
Intersecting over all primes proves the first equality in
\eqref{coef:transport-duality}; exchange $N_0,D_0$ for the second.
\end{proof}

The definition in \eqref{coef:transport} is exactly the intersection
over all finite primes and all admissible local bases used in the
coefficient construction. The independence just proved is why those
extra quantifiers do not shrink the lattice further.

\subsection{The normalized Hermitian form and unimodular inputs}

Suppose $K$ has an involution $c\mapsto\bar c$.
It preserves $\mathcal O$ and acts coefficientwise on polynomials.
For each $\tau$ let $G_\tau\in\mathrm{GL}_n(K)$ be Hermitian,
$\bar G_\tau^{\,t}=G_\tau$. On $V$ define
\begin{equation}\label{coef:hermitian}
 h_G(F,F')=B(\bar F,\rho(G)F').
\end{equation}
Written in terms of polynomials, this is exactly
\[
 h_G(F,F')=
 \frac{\prod_\tau\det(G_\tau)^{a_\tau}}
      {\prod_\tau C_{\nu_\tau}}
 [\bar F,\ F'((G_\tau^tX_\tau)_\tau)].
\]
Thus both the factorial normalization and the determinant characters
are retained. Formula~\eqref{coef:contravariant} shows that $h_G$ is
Hermitian and nondegenerate. At a complex embedding where every
$G_\tau$ becomes positive definite, its numerical Gram matrix in a
rational polynomial basis is positive definite as well. Indeed the
ordinary complex Fock form is positive definite, the adjoint of
$\rho(g)$ is $\rho(g^*)$, and a factorization $G=SS^*$ gives
$\rho(G)=\rho(S)\rho(S)^*$.
The same adjoint identity gives, componentwise,
\[
 h_G(\rho(g)F,\rho(g)F')=h_{g^*Gg}(F,F').
\]
Thus a base isometry $g^*Gg=G$ acts by isometries on the coefficient
space. If it also preserves the input lattices, it preserves both
coefficient lattices by \eqref{coef:equivariance}. These are the
isometry and stability properties used for their finite stabilizers.

For a base lattice put
\[
 M_\tau^{\#,G_\tau}
   =\{y:\bar x^{\,t}G_\tau y\in\mathcal O\text{ for all }x\in M_\tau\}
   =G_\tau^{-1}(\bar M_\tau)^\vee.
\]
The symbol $\#$ on coefficient lattices below refers to $h_G$.

\begin{theorem}[The coefficient duality actually needed]\label{coef:hermitian-duality}
With all the preceding definitions,
\begin{equation}\label{coef:hermitian-duality-eq}
             N(M)^\#=D(M^{\#,G}),\qquad
             D(M)^\#=N(M^{\#,G}).
\end{equation}
In particular, if every input lattice is unimodular for its base
Hermitian form, then
\begin{equation}\label{coef:unimodular}
                         N(M)^\#=D(M)\subset N(M),\qquad
                         D(M)^\#=N(M).
\end{equation}
Consequently these coefficient lattices satisfy the initial dual
containments in the finite-group argument with $s=1$.
\end{theorem}
\begin{proof}
All standard lattice generators have integer coefficients. Conjugation
therefore gives $\overline{N(M)}=N(\bar M)$ and
$\overline{D(M)}=D(\bar M)$: conjugating a local basis at
$\overline{\mathfrak p}$ gives the local basis at $\mathfrak p$.
For $y\in V$, the defining condition for $N(M)^\#$ is
\[
 B(\overline{N(M)},\rho(G)y)\subset\mathcal O.
\]
By symmetry of $B$ and \eqref{coef:transport-duality}, this is
equivalent to $\rho(G)y\in D((\bar M)^\vee)$.
Equivariance \eqref{coef:equivariance} makes the resulting lattice
\[
 \rho(G^{-1})D((\bar M)^\vee)
                   =D(G^{-1}(\bar M)^\vee)=D(M^{\#,G}).
\]
The other equality follows by exchanging the two lattices.
The unimodular conclusion is immediate, including its inclusion,
which was proved in Lemma~\ref{coef:transport-lattices}.
\end{proof}

\begin{lemma}[Extension of a unimodular input]\label{coef:extension}
Let $E\hookrightarrow K$ be an embedding of number fields compatible
with their involutions. For a nondegenerate Hermitian $E$-space and a full
$\mathcal O_E$-lattice $L$, extension of scalars satisfies
\[
 (\mathcal O_K L)^\#=\mathcal O_K L^\#,
\]
where the form on the left is the scalar extension of the original
$E$-valued form. In particular a unimodular lattice remains unimodular,
even at ramified primes.
\end{lemma}
\begin{proof}
The module $L$ is finite projective over the Dedekind domain
$\mathcal O_E$. The torsion-free module $\mathcal O_K$ is flat over
$\mathcal O_E$. Duals of finite projective modules commute with scalar
extension: this is immediate for a finite free module, and follows
for a projective module by writing it as a direct summand of one.
Compatibility of the involutions also identifies the extension of
$\bar L$ with the conjugate of its extension. Use the formula
$L^\#=G^{-1}(\bar L)^\vee$ in any coordinates over $E$ and extend
scalars. The pairing is extended as a field-valued pairing; no trace
map is inserted, so there is no different-ideal factor.
\end{proof}

The first coefficient construction takes these scalar extensions of
the original unimodular base lattices. The second uses their ordinary
duals. The latter are unimodular for the inverse
transpose Gram form, because direct calculation gives
\[
 (M^\vee)^{\#,(G^{-1})^t}=(M^{\#,G})^\vee.
\]
Indeed, both sides equal $G^t\bar M$, using
$(\bar M^\vee)^\vee=\bar M$.
Thus Theorem~\ref{coef:hermitian-duality} supplies the exact duality
for both constructions once their stated base-lattice unimodularity is
available. The construction and uniform estimates for those underlying
arithmetic lattices remain separate requirements.

\begin{example}[A determinant character of negative weight]
For $n=1$ and weight $\lambda=-1$, the polynomial partition is zero
and $V=K$, but $\rho(g)=g^{-1}$.
Take $K=\mathbb Q(i)$, $M=\mathcal O$ and $G=2$.
Then $M^{\#,G}=\tfrac12\mathcal O$,
$N(M^{\#,G})=2\mathcal O$, and
$h_G(x,y)=\bar x y/2$ has dual lattice $2\mathcal O$, as required.
If the determinant factor in \eqref{coef:hermitian} were omitted,
the coefficient dual would instead be $\mathcal O$.
Negative weights therefore cannot be replaced by their nonnegative
polynomial parts without retaining the determinant character.
More generally, for any fractional ideal $I$ and integral rank-one
weight $a$, local transport gives $N(I)=D(I)=I^a$, including negative
$a$ and nonprincipal $I$. The Hermitian dual for Gram entry $g$ is
$g^{-a}\bar I^{-a}$, equal to $D(g^{-1}\bar I^{-1})$ as the theorem
asserts. This also checks the transport convention when no global
basis of the input lattice exists.
\end{example}

\subsection{Checks and the resulting scope of the reconstruction}

An accompanying exact computation constructs the two lattices
independently for twelve small cases. The dual-Weyl lattice is the
saturated integer kernel of the full column-invariance equations;
the Weyl lattice is the lattice generated from $H_\nu$ by closure
under all divided row operators. Smith and Hermite normal forms are
used over the integers. In every case the normalized matrix pairing
the two lattice bases is integral and has determinant $\pm1$.
For example, for square matrices the partitions $(2,1,0)$ and
$(3,1,0)$ have ranks $8$ and $15$, normalizing constants $3$ and $8$,
and indices $[\nabla:\Delta]=3$ and $4096$, respectively.
For four rows, two columns and partition $(2,1)$, the rank is $20$
and the index is $81$. These finite checks support the stated
conventions; the proof of the general result is the preceding argument.

The coefficient-lattice duality, its determinant twists, tensor
products, local transport and scalar extension are now proved in
ordinary mathematics with the actual definitions used by the divisor
construction. The proof uses characteristic zero, $n\le r$, a
nonincreasing nonnegative polynomial partition, invertible base Gram
matrices and full input lattices. The original integral weights may
be arbitrary dominant integers. No assumption that the residue
characteristic exceeds the degree is made. No counterexample to the
actual duality has been found; the explicit examples above instead
identify changes of normalization or saturation that would make it
false.



% Exact Fock normalizer and local equality of Weyl and dual-Weyl lattices.
% Labels in this fragment have the prefix fh-.
\section{The exact Fock normalizer and equality at large primes}
\label{fh-section}

We now compute the normalizing integer $C_\nu$ of
\eqref{coef:fock}. This gives an explicit integer which kills
$\nabla_\nu(\mathbb Z)/\Delta_\nu(\mathbb Z)$, and proves the
local equality of these lattices needed at the level primes.
The source obtains that equality from simplicity in the lowest
alcove.\cite{alpo-uniform-source} The calculation below instead
uses polynomial differential operators and the integral duality
already proved in Theorem~\ref{coef:integral-duality}.

\subsection{A differential determinant identity}

Let $X=(x_{ia})_{1\le i,a\le h}$ be a square matrix of variables,
write $\partial_{ia}=\partial/\partial x_{ia}$, and put
$E_{ij}=\sum_a x_{ia}\partial_{ja}$.
Composition of operators acts from right to left; in all products
below the factors are written from left to right in increasing $j$.

\begin{lemma}[The differential Capelli identity]\label{fh-capelli}
On $\mathbb Z[X]$ one has
\begin{equation}\label{fh-capelli-identity}
 \det(X)\det(\partial)=
 \sum_{\sigma\in\mathfrak S_h}\operatorname{sgn}(\sigma)
       \prod_{j=1}^h
           \bigl(E_{\sigma(j),j}+(h-j)\delta_{\sigma(j),j}\bigr).
\end{equation}
On the left, multiplication by $\det(X)$ follows differentiation.
\end{lemma}
\begin{proof}
Adjoin exterior generators $e_1,\ldots,e_h$, so
$e_ie_j=-e_je_i$ and $e_i^2=0$. Derivatives act only on the
polynomial variables. Set
\[
       \xi_a=\sum_i e_i x_{ia},\qquad
       A_j=\sum_a\xi_a\partial_{ja}=\sum_i e_iE_{ij}.
\]
We prove by descending induction on $j$ the operator identity
\begin{equation}\label{fh-exterior-ordering}
 \prod_{i=j}^h\bigl(A_i+(h-i)e_i\bigr)
   =\sum_{a_j,\ldots,a_h=1}^h
        \xi_{a_j}\cdots\xi_{a_h}
        \partial_{j,a_j}\cdots\partial_{h,a_h}.
\end{equation}
The empty tail is the identity. Suppose the assertion is known
for $j+1$, and denote its right side by $N_{j+1}$.
When $A_j$ acts on $N_{j+1}$, the terms in which its derivative
passes all $h-j$ exterior factors give exactly the right side
of \eqref{fh-exterior-ordering}. There is one additional term
for each factor hit by the derivative, since
\[
                         \partial_{ja}\xi_b=e_j\delta_{ab}.
\]
Each such term is $-e_jN_{j+1}$. To check the sign, if the
derivative hits the $t$th exterior factor of the tail, move the
resulting $e_j$ to the front, across $t$ factors, then move the
initial $\xi_a$ into its former position, across $t-1$ factors.
The combined sign is $(-1)^t(-1)^{t-1}=-1$.
Thus the total additional contribution is
$-(h-j)e_jN_{j+1}$, which the displayed shift cancels.
This proves the induction, using only integral operator identities.

Take the coefficient of $e_1\cdots e_h$ in
\eqref{fh-exterior-ordering} for $j=1$.
On the left it is the ordered sum on the right side of
\eqref{fh-capelli-identity}. On the right, only permutations of
$a_1,\ldots,a_h$ contribute: a repeated index makes the exterior
product zero. For $a_j=\tau(j)$ that coefficient is
$\operatorname{sgn}(\tau)\det(X)$.
The remaining sum of commuting partial derivatives is
$\det(\partial)$. This proves \eqref{fh-capelli-identity}.
\end{proof}

\subsection{The hook product}

Let $\nu$ be a partition. Its diagram consists of the pairs
$(i,j)$ with $1\le j\le\nu_i$. Write
$\nu'_j=\#\{i:\nu_i\ge j\}$ and define the hook length of
the pair $(i,j)$ by
\[
                  h_\nu(i,j)=\nu_i-j+\nu'_j-i+1.
\]
It is a positive integer: it counts that pair, the pairs to its
right in the same row, and those below it in the same column.
Zero rows and extra unused variables do not change the highest
polynomial $H_\nu$ or its Fock norm. Empty products are one.

\begin{theorem}[Exact Fock normalization]\label{fh-hook-product}
For the highest polynomial and the unnormalized Fock form of
Section~\ref{coef:section},
\begin{equation}\label{fh-hook-equality}
                   C_\nu=[H_\nu,H_\nu]
                          =\prod_{(i,j)\in\nu}h_\nu(i,j).
\end{equation}
If $\nu$ has $h>0$ positive parts, all factors on the right
are at most $\nu_1+h-1$.
\end{theorem}
\begin{proof}
For the zero partition, $H_\nu=1$ and $C_\nu=1$.
Otherwise remove the zero parts and work on the first $h$ rows
and columns. All $\nu_i$ are then positive. Define
$\nu^-=(\nu_1-1,\ldots,\nu_h-1)$, omitting any zero parts
when computing its highest polynomial. The defining products of
initial principal minors give
\[
                             H_\nu=\det(X)H_{\nu^-}.
\]
The operators $E_{ij}$ with $i<j$ kill $H_\nu$: they
differentiate addition of an earlier row to a later one, which
preserves every initial principal minor involving the later row;
the other factors do not involve that row. Also
$E_{ii}H_\nu=\nu_iH_\nu$ by row homogeneity.

In a term of the ordered determinant
\eqref{fh-capelli-identity} for a nonidentity permutation
$\sigma$, let $j$ be its greatest nonfixed index. Then
$\sigma(j)<j$. All factors to its right are diagonal and
act on $H_\nu$ by scalars, while the factor at $j$ is a
raising operator and kills it. Only the identity permutation
survives. Hence
\[
 \det(X)\det(\partial)H_\nu
                  =\left(\prod_{i=1}^h(\nu_i+h-i)\right)H_\nu.
\]
Cancellation of the nonzero polynomial $\det(X)$ yields
\begin{equation}\label{fh-lowering}
 \det(\partial)H_\nu
                  =\left(\prod_{i=1}^h(\nu_i+h-i)\right)H_{\nu^-}.
\end{equation}

For the Fock pairing, multiplication by $x_{ia}$ is adjoint to
$\partial_{ia}$; this follows immediately from the factorial
weight of a monomial. By composition and linearity,
multiplication by $\det(X)$ is adjoint to $\det(\partial)$.
Consequently \eqref{fh-lowering} gives the exact recurrence
\begin{equation}\label{fh-column-recurrence}
                 C_\nu=
                    \left(\prod_{i=1}^h(\nu_i+h-i)\right)C_{\nu^-}.
\end{equation}
The factors in this product are exactly the hook lengths in
the first column of $\nu$. Removing that column gives
$\nu^-$, and does not change the hook lengths of the remaining
pairs: their column indices decrease by one, their row lengths
decrease by one, and their column heights are the former heights
of the next columns. Induction on the number of pairs proves
\eqref{fh-hook-equality}. Finally
$h_\nu(i,j)\le\nu_1+h-1$, directly from its definition.
\end{proof}

\subsection{An explicit annihilator and the exact prime condition}

\begin{corollary}\label{fh-annihilator}
For every partition $\nu$,
\begin{equation}\label{fh-inclusions}
              C_\nu\nabla_\nu(\mathbb Z)
                    \subseteq\Delta_\nu(\mathbb Z)
                    \subseteq\nabla_\nu(\mathbb Z).
\end{equation}
In particular, if a prime $p$ divides no hook length of $\nu$,
then
\[
                \Delta_\nu(\mathbb Z_{(p)})
                          =\nabla_\nu(\mathbb Z_{(p)}).
\]
The equality remains true after extension to any
characteristic-zero coefficient domain in which $C_\nu$ is a unit.
\end{corollary}
\begin{proof}
For $F,G\in\nabla_\nu(\mathbb Z)$,
\[
                   B_\nu(C_\nu F,G)=[F,G]\in\mathbb Z.
\]
By Theorem~\ref{coef:integral-duality}, this says exactly that
$C_\nu F\in\Delta_\nu(\mathbb Z)$. The other inclusion was
part of the construction. Tensoring these inclusions with the
coefficient ring and using the perfect bases gives the assertion
whenever $C_\nu$ becomes a unit. Equation~\eqref{fh-hook-equality}
supplies the stated sufficient prime condition.
\end{proof}

Let $\lambda_1\ge\cdots\ge\lambda_N$ be a dominant integral
weight with $N\ge1$, and put $\nu_i=\lambda_i-\lambda_N$.
Then $\nu_N=0$. If $\nu$ is nonempty, its number $h$ of
positive parts is at most $N-1$, so every hook length is at most
\[
                          \lambda_1-\lambda_N+N-2.
\]
We therefore obtain the following precise conclusion.

\begin{corollary}\label{fh-large-prime-equality}
If $p$ is prime and
\begin{equation}\label{fh-prime-range}
                       \lambda_1-\lambda_N+N-1\le p,
\end{equation}
the standard Weyl and dual-Weyl lattices of weight $\lambda$
are equal after localization at $p$. The same is true at every
prime over $p$ in a number field, and after transport by a common
local basis. More explicitly, the single positive integer
$c=C_\nu$, independent of the polynomial, satisfies
\[
       p\nmid c,\qquad
       c\nabla_\nu(\mathbb Z)\subseteq\Delta_\nu(\mathbb Z).
\]
For a constant weight, including $N=1$, one can take $c=1$.
The module in rank zero also consists of constants and has $c=1$.
\end{corollary}
\begin{proof}
For a nonempty $\nu$, \eqref{fh-prime-range} makes every
positive hook length at most $p-1$. Apply
Corollary~\ref{fh-annihilator}. The determinant twist changes
the group action but not these standard polynomial lattices or
their normalized pairing. Localization, coefficient extension
and a common invertible transport preserve equality.
The constant cases follow from $H_\nu=1$.
\end{proof}

The actual source uses the stronger hypothesis
$\lambda_1-\lambda_N+N<p$: its quantity called the spread
includes $N$, not $N-1$. Thus Corollary~\ref{fh-large-prime-equality}
proves the precise local lattice assertion used there, with a
slightly larger allowed range and an explicit clearing integer.
It does not require a modular simplicity theorem or a passage
from equality of dimensions to equality of integral lattices.

The prime condition cannot be omitted in general. For
$\lambda=(2,0)$, the dual-Weyl lattice is
$\mathbb Zx^2\oplus\mathbb Zxy\oplus\mathbb Zy^2$, whereas
the Weyl lattice is
$\mathbb Zx^2\oplus\mathbb Z(2xy)\oplus\mathbb Zy^2$.
Here $C_{(2)}=2$, and the two lattices remain different at $p=2$.
This is the example of Section~\ref{coef:section}; it lies
outside \eqref{fh-prime-range}. No equality of the two arbitrary
integral lattices is being asserted.


% Integral initial-minor spanning, proved by an elementary leading-term argument.
% Labels in this fragment use the prefix sm-.
\section{Integral spanning by products of initial minors}
\label{sm-section}

The polynomial module in Section~\ref{coef:section} has an additional
integral description needed for branching: its elements are integral
linear combinations of products of initial-column minors. The submitted
proof obtains this description from a two-column argument and
semistandard tableaux.\cite{alpo-uniform-source}
We give a direct proof of the required spanning statement. The only
dimension calculation takes place over characteristic zero; the final
subtraction uses monic polynomials and therefore introduces no
denominators.

\subsection{A two-column leading-term lemma}

Let $d_1,\ldots,d_N$ be nonnegative integers. In the polynomial ring
$\mathbb Q[x_1,\ldots,x_N,y_1,\ldots,y_N]$, let $V$ be the
finite-dimensional space in which the total degree in $x_i,y_i$ is
$d_i$ for every $i$. Put
\[
 E=\sum_{i=1}^N x_i\frac{\partial}{\partial y_i},\qquad
 F=\sum_{i=1}^N y_i\frac{\partial}{\partial x_i},\qquad
 H=\sum_{i=1}^N
 \left(x_i\frac{\partial}{\partial x_i}
             -y_i\frac{\partial}{\partial y_i}\right).
\]
These operators preserve $V$ and satisfy $EF-FE=H$.
For an integer $h$, denote the $h$-eigenspace of $H$ by $V_h$.
Thus a monomial $x^\alpha y^\beta$ belongs to $V_h$ precisely when
\[
 \alpha_i+\beta_i=d_i\quad(1\le i\le N),
 \qquad \sum_i(\alpha_i-\beta_i)=h.
\]
Order monomials lexicographically with the variables in the order
\begin{equation}\label{sm-two-order}
                   x_1>\cdots>x_N>y_1>\cdots>y_N.
\end{equation}
The leading monomial is the greatest monomial with nonzero coefficient
in this order.

\begin{lemma}\label{sm-two-column}
Let $h\ge0$, and let $0\ne f\in V_h$ satisfy $Ef=0$.
If $x^\alpha y^\beta$ is its leading monomial, then
\begin{equation}\label{sm-ballot}
                \sum_{j=1}^i\beta_j
                    \le\sum_{j=1}^{i-1}\alpha_j
                    \qquad(1\le i\le N).
\end{equation}
The same assertion holds over any field of characteristic zero.
\end{lemma}
\begin{proof}
We first compute the dimension of the kernel. Complexify $V$ and make
its monomials orthogonal, with
\[
                \|x^\alpha y^\beta\|^2
                           =\prod_i\alpha_i!\,\beta_i!.
\]
Direct comparison of the coefficients of two monomials gives
$E^*=F$. For example, the contribution involving index $i$ uses
\[
 \beta_i(\alpha_i+1)!\,(\beta_i-1)!
                  =(\alpha_i+1)\alpha_i!\,\beta_i!.
\]
For $v\in V_{h+2}$, the commutator identity consequently gives
\[
            \|Fv\|^2=\|Ev\|^2+(h+2)\|v\|^2.
\]
Since $h+2>0$, the map $F:V_{h+2}\longrightarrow V_h$ is
injective. Its adjoint $E:V_h\longrightarrow V_{h+2}$ is therefore
surjective. The matrices have rational coefficients, so their ranks
over $\mathbb Q$ and $\mathbb C$ agree. It follows that
\begin{equation}\label{sm-kernel-dimension}
       \dim_{\mathbb Q}\ker(E|_{V_h})
                       =\dim_{\mathbb Q}V_h-\dim_{\mathbb Q}V_{h+2}.
\end{equation}

We count the difference on the right without a decomposition theorem
for representations. For $a\ge0$ define the Laurent polynomial
\[
                     \chi_a(T)=T^a+T^{a-2}+\cdots+T^{-a}.
\]
An elementary coefficient calculation gives
\begin{equation}\label{sm-character-product}
              \chi_a(T)\chi_b(T)
                 =\sum_{s=0}^{\min(a,b)}\chi_{a+b-2s}(T).
\end{equation}
Indeed, for $0\le t\le a+b$, the coefficient of $T^{a+b-2t}$
on either side is $1+\min(a,b,t,a+b-t)$; all other exponents have
zero coefficient. On the left this counts pairs $(j,k)$ with
$0\le j\le a$, $0\le k\le b$ and $j+k=t$. On the right it
counts $s$ with $0\le s\le\min(a,b)$ and $s\le t\le a+b-s$.

The dimension of $V_h$ is the coefficient of $T^h$ in
$\prod_i\chi_{d_i}(T)$. Iterating \eqref{sm-character-product},
starting from $h_0=0$, expresses this product as a sum of the
polynomials $\chi_{h_N}(T)$, one for each sequence
\begin{equation}\label{sm-path}
 0\le\beta_i\le\min(d_i,h_{i-1}),\qquad
 \alpha_i=d_i-\beta_i,\qquad
 h_i=h_{i-1}+d_i-2\beta_i
 \quad(1\le i\le N).
\end{equation}
Every $h_i$ is nonnegative. For $h\ge0$, the coefficient of $T^h$
minus the coefficient of $T^{h+2}$ in $\chi_k(T)$ is one if $k=h$
and zero otherwise. Therefore \eqref{sm-kernel-dimension} is the
number of sequences \eqref{sm-path} with $h_N=h$.
Since
\[
                    h_{i-1}=
                         \sum_{j<i}(\alpha_j-\beta_j),
\]
these sequences are exactly the nonnegative exponent arrays with the
prescribed row degrees and total weight $h$ which satisfy
\eqref{sm-ballot}.

For each such array we construct a monic polynomial in the kernel
with leading monomial $x^\alpha y^\beta$. Process the indices
$i=1,\ldots,N$ in order. Before processing $i$, retain a product of
brackets $x_jy_k-x_ky_j$ with $j<k<i$, together with
$h_{i-1}$ unpaired factors $x_j$, all with $j<i$.
Choose $\beta_i$ of the unpaired factors and replace each chosen
$x_j$ by
\[
                            x_jy_i-x_iy_j.
\]
Then adjoin $\alpha_i$ factors $x_i$. The choice is possible because
$\beta_i\le h_{i-1}$; the number of unpaired factors afterwards is
$h_{i-1}-\beta_i+\alpha_i=h_i$. This starts with the empty product
at $i=1$, when \eqref{sm-ballot} says $\beta_1=0$.

Every bracket and every unpaired $x_j$ is killed by $E$, so their
final product is killed by $E$. Its row degrees and weight are the
prescribed ones. Because $j<i$, the leading term of the displayed
bracket in \eqref{sm-two-order} is $x_jy_i$, with coefficient one.
Consequently the final product has leading monomial
$x^\alpha y^\beta$ and leading coefficient one.
Choose one such product for each exponent array satisfying
\eqref{sm-ballot}. Distinct choices of array give distinct leading
monomials, hence linearly independent polynomials. Their number is
the dimension in \eqref{sm-kernel-dimension}; they form a
$\mathbb Q$-basis of the kernel.

In a nonzero linear combination of these polynomials, the greatest
of their leading monomials with nonzero coefficient cannot cancel:
every monomial in the other participating polynomials is smaller.
This proves \eqref{sm-ballot} for the leading monomial of $f$.
Finally, the matrix of $E$ and all the constructed polynomials have
integer coefficients. The rank of that matrix is unchanged under
extension from $\mathbb Q$ to any characteristic-zero field.
The same products are therefore a basis of the corresponding kernel
over that field, and the leading-monomial conclusion is unchanged.
\end{proof}

\subsection{Several columns and integral subtraction}

Fix integers $N\ge r\ge1$ and a partition
$\mu_1\ge\cdots\ge\mu_r\ge0$. Let $X=(x_{ij})$ be an $N$ by $r$
matrix of indeterminates. For a characteristic-zero integral domain
$R$, define $\nabla_\mu(R)$ to be the set of polynomials in $R[X]$
which are homogeneous of degree $\mu_j$ in each column $j$ and
satisfy
\begin{equation}\label{sm-column-invariance}
          f\bigl(X(I+tE_{lj})\bigr)=f(X)
                    \qquad(1\le l<j\le r)
\end{equation}
as polynomial identities in the additional indeterminate $t$.
Here $E_{lj}$ is the elementary matrix, so this substitution adds
$t$ times column $l$ to column $j$.
Use the lexicographic order
\begin{equation}\label{sm-column-order}
 x_{11}>x_{21}>\cdots>x_{N1}>
 x_{12}>\cdots>x_{N2}>\cdots>x_{1r}>\cdots>x_{Nr}.
\end{equation}
For a strictly increasing row sequence
$I=(i_1<\cdots<i_s)$, $1\le s\le r$, put
\[
               \Delta_I(X)=
                 \det(x_{i_a,b})_{1\le a,b\le s}.
\]
This is an initial-column minor of height $s$.

\begin{lemma}\label{sm-leading-interlacing}
If $0\ne f\in\nabla_\mu(R)$ has leading monomial
$\prod_{i,j}x_{ij}^{a_{ij}}$, then for $j<r$ and $1\le i\le N$,
\begin{equation}\label{sm-interlacing}
                   \sum_{b\le i}a_{b,j+1}
                         \le\sum_{b<i}a_{b,j}.
\end{equation}
\end{lemma}
\begin{proof}
Pass to the fraction field of $R$. Differentiating
\eqref{sm-column-invariance} at $t=0$ shows that
\[
                       E_jf=0,\qquad
              E_j=\sum_i x_{ij}\frac{\partial}{\partial x_{i,j+1}}.
\]
Decompose $f$ according to the monomial in all columns other than
$j,j+1$ and according to the degrees
$d_i$ in $x_{ij},x_{i,j+1}$. The operator $E_j$ preserves each
of these finite components, and hence kills each one.
In particular it kills the nonzero component containing the leading
monomial of $f$. After removing its fixed monomial in the other
columns, this component has the two-column degrees $\mu_j,\mu_{j+1}$,
so its weight is $h=\mu_j-\mu_{j+1}\ge0$.
The restriction of \eqref{sm-column-order} to the two columns is
exactly \eqref{sm-two-order}. Its leading monomial is the part of
the leading monomial of $f$ in these two columns. Applying
Lemma~\ref{sm-two-column} gives \eqref{sm-interlacing}.
\end{proof}

\begin{lemma}\label{sm-matching-minors}
Suppose a nonnegative integer array $(a_{ij})$ has column sums
$\mu_j$ and satisfies \eqref{sm-interlacing}. There is a product of
initial-column minors with the same column degrees, leading monomial
$\prod_{i,j}x_{ij}^{a_{ij}}$, and leading coefficient one.
It has exactly $\mu_s-\mu_{s+1}$ factors of height $s$, where
$\mu_{r+1}=0$.
\end{lemma}
\begin{proof}
For each column $j$, form the nondecreasing list
\[
                       w_j(1),\ldots,w_j(\mu_j)
\]
in which the row index $i$ occurs $a_{ij}$ times. If
$1\le t\le\mu_{j+1}$ and $w_{j+1}(t)=i$, there are at least
$t$ entries of the list for column $j+1$ at most $i$.
By \eqref{sm-interlacing}, the list for column $j$ has at least
$t$ entries strictly less than $i$. Thus
\[
                         w_j(t)<w_{j+1}(t).
\]
For each $1\le t\le\mu_1$, put
\[
 s_t=\max\{j:t\le\mu_j\},\qquad
 I_t=(w_1(t),\ldots,w_{s_t}(t)).
\]
The preceding inequalities make $I_t$ a strictly increasing row
sequence. Consider $P=\prod_{t=1}^{\mu_1}\Delta_{I_t}(X)$.
In the determinant of any initial-column minor, its diagonal term
is the greatest term in \eqref{sm-column-order}: column one chooses
the smallest available row index, then column two chooses the
smallest remaining one, and so on. Its coefficient is one.
Leading monomials and coefficients multiply, so $P$ has the claimed
leading monomial and leading coefficient one. A factor has height
$s$ exactly when $\mu_{s+1}<t\le\mu_s$; this proves the last
assertion as well as the column degrees. When $\mu_1=0$, all lists
are empty and $P=1$ has the required properties.
\end{proof}

\begin{theorem}[Integral spanning by initial minors]
\label{sm-initial-minors}
For $N\ge r\ge1$ and a partition $\mu$ as above,
$\nabla_\mu(\mathbb Z)$ is spanned over $\mathbb Z$ by products
of initial-column minors having exactly $\mu_s-\mu_{s+1}$ factors
of height $s$, for $1\le s\le r$ and $\mu_{r+1}=0$.
The same equality of modules holds over any characteristic-zero
integral domain $R$, with the identical integral polynomials as
generators and coefficients in $R$.
\end{theorem}
\begin{proof}
An initial-column minor of height $s$ is unchanged under every
substitution in \eqref{sm-column-invariance}. If $j>s$, column $j$
does not occur. If $j\le s$, both columns $l<j$ occur, and the
substitution is a determinant-preserving column operation.
The indicated product of minors therefore belongs to
$\nabla_\mu(R)$, since its degree in column $j$ is
\[
                     \sum_{s\ge j}(\mu_s-\mu_{s+1})=\mu_j.
\]

Conversely, let $f\in\nabla_\mu(R)$ be nonzero, with leading
coefficient $c\in R$ and leading exponent array $(a_{ij})$.
Lemma~\ref{sm-leading-interlacing} and
Lemma~\ref{sm-matching-minors} supply a product $P$ of the required
minors with that leading monomial and leading coefficient one.
Then $f-cP$ still belongs to $\nabla_\mu(R)$ and, unless it is
zero, has a strictly smaller leading monomial.
There are only finitely many monomials with the fixed column degrees
$\mu_1,\ldots,\mu_r$. Repeating this subtraction must therefore
terminate at zero. It expresses $f$ as a linear combination over
$R$ of the specified integral products. Every subtraction used a
leading coefficient in $R$ and a monic polynomial; no division in
$R$ occurred. Taking $R=\mathbb Z$ proves the integral assertion.
\end{proof}

The case of no columns consists of constants and is represented by
the empty product. Zero parts of $\mu$ likewise cause no exception.
This argument proves exactly the integral spanning input and its
characteristic-zero extension. It does not assert that differentiating
the column identities suffices in positive characteristic: in
characteristic $p$, for example, $y^p$ is killed by
$x\partial_y$ but is not invariant under $y\mapsto y+tx$ as a
polynomial identity. The submitted source states a more general
result for domains of arbitrary characteristic; the present proof
does not replace that additional assertion. All coefficient lattices
to which we apply Theorem~\ref{sm-initial-minors} are in
characteristic zero.


\section{The compact branching law from polynomial representations}
\label{cb-section}

We prove the branching law needed both for the integral lowering map
and for the terminal radial operator.  The proof uses the polynomial
modules and their duality from Section~\ref{coef:section}.  In
particular, multiplicity one and the precise interlacing inequalities
will be proved, not assumed as additional representation-theoretic
inputs.

For a partition $\nu=(\nu_1,\ldots,\nu_r)$, use the polynomial
module $V_\nu=\nabla_\nu(\mathbb C)$ on an $r$ by $r$ matrix.
Zero last columns may be omitted: column degree zero means that
the polynomial does not involve that column.  We use the representation
\[
             (\rho(g)F)(X)=F(g^tX),\qquad g\in\operatorname{GL}_r(\mathbb C).
\]
The transpose makes $\rho(g)\rho(h)=\rho(gh)$.
Its infinitesimal operators are
\[
 E_{ab}=\sum_s x_{as}\frac{\partial}{\partial x_{bs}},
 \qquad
 [E_{ab},E_{cd}]=\delta_{bc}E_{ad}-\delta_{da}E_{cb}.
 \tag{CB1}
\]
For the Hermitian Fock form, linear in its first argument,
$E_{ab}^*=E_{ba}$.  This follows directly on monomials,
or by Lemma~\ref{coef:adjoint}.  The restriction of $\rho$ to
$U(r)$ is unitary.  The highest polynomial $H_\nu$ from
Section~\ref{coef:section} has row weight $\nu$ and satisfies
$E_{ab}H_\nu=0$ for $a<b$.

For a dominant integral tuple $\lambda_1\ge\cdots\ge\lambda_r$,
put $s=\lambda_r$ and define
\[
 V_\lambda=\det^{\,s}\otimes V_{\lambda-s(1,\ldots,1)}.
 \tag{CB2}
\]
Negative $s$ causes no difficulty on $\operatorname{GL}_r$ or $U(r)$.
The determinant character has absolute value one on the latter.
Its infinitesimal action adds $s$ to every diagonal weight.

\subsection{Highest weights and irreducibility}

\begin{lemma}\label{cb-polynomial-irreducible}
The representation $V_\lambda$ is irreducible, and its highest
weight is $\lambda$.
\end{lemma}
\begin{proof}
First suppose that $\lambda=\nu$ is a partition.  Theorem
\ref{coef:integral-duality} shows that the divided row operators
applied to $H_\nu$ span $\nabla_\nu(\mathbb Q)$.  The same is
true over $\mathbb C$.  Indeed these spaces have a fixed finite
monomial support, and their equations have rational coefficients:
column homogeneity and $E^{\rm col}_{ij}F=0$ for $i<j$ give their
defining equations.  In characteristic zero the latter differential
equation is equivalent to invariance under the corresponding column
addition, by the finite Taylor expansion.  Taking a kernel of a
finite rational matrix commutes with extension to $\mathbb C$.

Every divided off-diagonal row operator satisfies
$D_{ab}^{(k)}=E_{ba}^k/k!$;
the divided diagonal operator is $\binom{E_{aa}}k$ on each row
degree.  Thus $H_\nu$ is cyclic for the operators (CB1).
Suppose $W$ is a $U(r)$-invariant subspace of $V_\nu$ and let
$P$ be its orthogonal projection.  Its orthogonal complement is
invariant too, so $P$ commutes with $U(r)$ and hence with its
complexified infinitesimal operators (CB1).  It preserves row
weights.  Lemma~\ref{coef:extremal-weight} says that the row
weight $\nu$ space is exactly $\mathbb C H_\nu$.
Consequently $PH_\nu=cH_\nu$, and $P^2=P$ gives $c=0$ or $1$.
Cyclicity then gives $P=0$ or $1$.  This proves irreducibility.
Twisting by the one-dimensional determinant character preserves
irreducibility and shifts the highest weight as in (CB2).
\end{proof}

\begin{lemma}\label{cb-highest-classification}
Every irreducible finite-dimensional continuous unitary
representation of $U(r)$ is isomorphic to exactly one $V_\lambda$,
where $\lambda$ is a dominant integral tuple.  A highest vector
of weight $\lambda$ determines an intertwiner from $V_\lambda$
uniquely up to its prescribed nonzero scalar value on $H_\lambda$.
\end{lemma}
\begin{proof}
Smoothness follows from the convolution argument of
Lemma~\ref{rp-unitary-derivative}; that argument uses no
classification theorem.  The diagonal torus acts by commuting
unitary operators, which have a common orthonormal eigenbasis.
Each continuous character of a circle is $e^{it}\mapsto e^{imt}$
for an integer $m$: differentiating its smooth character gives
$e^{ict}$ for a real $c$, and period $2\pi$ gives $c\in\mathbb Z$.
Thus the torus weights are integral tuples, and only finitely
many occur.

Choose a nonzero vector $v$ of lexicographically maximal weight
$\lambda$.  For $a<b$, the operator $E_{ab}$ increases that
weight by $e_a-e_b$; the commutator with the diagonal operators
proves this assertion.  Hence $E_{ab}v=0$.  Moreover,
\[
 \|E_{ba}v\|^2
   =\langle E_{ab}E_{ba}v,v\rangle
   =(\lambda_a-\lambda_b)\|v\|^2\ge0.
\]
Here we used $E_{ab}^*=E_{ba}$ and
$[E_{ab},E_{ba}]=E_{aa}-E_{bb}$.
It follows that $\lambda$ is dominant.

We give the required uniqueness argument, including its
existence assertion.  In any unitary representation, the
Gram matrix of the vectors $Pv$, where $P$ ranges over words
in the $E_{ab}$, is determined by $\lambda$ and $\|v\|^2$.
To see this, write
\[
                    \langle Pv,Qv\rangle
                         =\langle Q^*Pv,v\rangle.
\]
Use (CB1) to order every word with lowering operators
$E_{ab}$, $a>b$, on the left, diagonal operators in the
middle, and raising operators, $a<b$, on the right.
This process terminates: swapping an adjacent pair removes
an inversion, and the commutator term has smaller word length.
Induct first on length, and then on the number of inversions.
A nonempty raising part annihilates $v$.  The diagonal part
acts by its polynomial in $\lambda$.  A nonempty lowering
part changes the weight by a nonzero sum of negative positive
roots and is therefore orthogonal to $v$.  Only the term
with neither raising nor lowering operators contributes.
This computes the Gram matrix using just $\lambda$.

Apply this calculation to $v$ and to the highest vector in
$V_\lambda$, normalizing both to have norm one.
The rule $PH_\lambda\mapsto Pv$ is well-defined: a linear
relation has squared norm zero on one side exactly when it
has squared norm zero on the other.  It is an isometry and
intertwines all $E_{ab}$.  Its image is the cyclic subspace
generated by $v$.  The group $U(r)$ is connected and is
generated by exponentials of its Lie algebra; explicitly,
unitary diagonalization writes any unitary matrix as an
exponential of a skew-Hermitian matrix.  The map is therefore
a group intertwiner.  In an irreducible representation the
cyclic subspace is the whole space.  Its uniqueness follows
from cyclicity of $H_\lambda$.

Finally the same normal ordering shows that every weight in
$V_\lambda$ is obtained from $\lambda$ by subtracting a sum
of positive roots: the module is cyclic and every word on
its highest vector reduces to lowering words.  A nonempty
such sum strictly decreases the lexicographic weight.
Thus $\lambda$ is its unique maximal weight.  An isomorphism
preserves weights, so different $\lambda$ give nonisomorphic
representations.
\end{proof}

\subsection{The highest vectors on restriction}

Let $N=n+1$ with $n\ge1$, and first let
$\lambda=(\lambda_1,\ldots,\lambda_n,0)$ be a partition.
Omit its unused last column and write the variable matrix as
$X=\left(\begin{smallmatrix}A\\z\end{smallmatrix}\right)$,
where $A$ is $n$ by $n$ and $z=(z_1,\ldots,z_n)$.
Let $\Delta_i$ be the first $i$ by $i$ minor of $A$, and let
$S_i$ be the minor with rows $1,\ldots,i-1,N$ and columns
$1,\ldots,i$.  Thus $S_1=z_1$.

\begin{lemma}\label{cb-joint-highest}
The vectors in $V_\lambda$ which have top-row weight
$\mu=(\mu_1,\ldots,\mu_n)$ and are annihilated by
$E_{ab}$ for $1\le a<b\le n$ form a space of dimension
one if
\[
                  \lambda_i\ge\mu_i\ge\lambda_{i+1}
                                      \qquad(1\le i\le n),
 \tag{CB3}
\]
and a zero space otherwise.  In the first case a generator is
\[
 \xi_{\lambda,\mu}(X)
       =\prod_{i=1}^n\Delta_i(X)^{\mu_i-\lambda_{i+1}}
                           S_i(X)^{\lambda_i-\mu_i}.
 \tag{CB4}
\]
Its last-row degree is $|\lambda|-|\mu|$.
\end{lemma}
\begin{proof}
The assertion that $E_{ab}F=0$ is equivalent, by finite
Taylor expansion, to invariance under adding a multiple
of top row $a$ to top row $b$.  These additions generate
the lower unipotent row subgroup.  By definition, $F$
is also invariant under the upper unipotent column subgroup.

On the nonempty open set $\prod_i\Delta_i\ne0$, Gaussian
elimination gives $A=LDU$, with $L$ lower unipotent,
$U$ upper unipotent, and $D=\operatorname{diag}(d_1,\ldots,d_n)$.
First multiply the columns by $U^{-1}$ and then the top
rows by $L^{-1}$.  Invariance gives
\[
                         F(A,z)=F(D,zU^{-1}).
 \tag{CB5}
\]
The row and column degrees force
\[
 F(D,w)=c\prod_i d_i^{\mu_i}w_i^{\lambda_i-\mu_i}.
 \tag{CB6}
\]
If one of these exponents is negative, the restriction is
zero, and (CB5) forces $F=0$.  In particular a nonzero
vector has $\lambda_i-\mu_i\ge0$.

On this diagonal matrix one has
$\Delta_i=d_1\cdots d_i$ and
$S_i=d_1\cdots d_{i-1}w_i$.  The Laurent product in
(CB4), allowing temporarily negative exponents
$\mu_i-\lambda_{i+1}$, restricts to the monomial in
(CB6).  It is itself invariant under both unipotent
groups.  Therefore (CB5) gives $F=c\xi_{\lambda,\mu}$
as a rational function.

The polynomials $\Delta_i,S_i$ are pairwise nonassociate
irreducibles.  Here is the elementary determinant fact
needed for this assertion.  A generic determinant is
irreducible by induction on its size: as a polynomial
in its last row it is linear, and its coefficients are
the maximal cofactors of the other rows.  Those cofactors
are irreducible by induction and no two are associates,
since they involve different column variables.  Their
greatest common divisor is one.  The polynomial is
therefore primitive and is irreducible over the fraction
field of the other-row variables, where it is a nonzero
linear polynomial.  Gauss's lemma proves irreducibility.
Size one is a variable.  Adjoining additional independent
variables preserves this conclusion.  Finally different
sizes have different degrees, and $S_i$ involves the last
row whereas $\Delta_i$ does not; hence none of the
displayed minors are associates.

Unique factorization now implies
$\mu_i-\lambda_{i+1}\ge0$ whenever $F\ne0$.
This proves the necessity of (CB3) and the upper bound
one on the dimension.  If (CB3) holds, (CB4) is a
nonzero polynomial.  Each factor is invariant under
the indicated column additions and top-row additions.
Its column degrees are $\lambda_i$, and its top-row
degrees telescope to $\mu_i$.
Each $S_i$ has last-row degree one, so the remaining
degree is $\sum_i(\lambda_i-\mu_i)$.
Thus (CB4) belongs to the required space and proves
the assertion.
\end{proof}

\begin{theorem}[Compact branching law]\label{ib-compact-branching}
For $N\ge2$ and a dominant integral highest weight
$\lambda=(\lambda_1,\ldots,\lambda_N)$, restriction
of $V_\lambda$ to $U(N-1)\times U(1)$ is the orthogonal
direct sum
\[
 V_\lambda\big|_{U(N-1)\times U(1)}
     \simeq\bigoplus_{\lambda_i\ge\mu_i\ge\lambda_{i+1}}
                           V_\mu\otimes t^{\,|\lambda|-|\mu|}.
 \tag{CB7}
\]
Each summand occurs once.  The character of the last
$U(1)$ is displayed, not suppressed.  For $N=1$ the
restriction to the trivial group has one one-dimensional
summand, of last-circle weight $\lambda_1$.
\end{theorem}
\begin{proof}
First subtract $\lambda_N$ from every weight, so that
the polynomial model of Lemma~\ref{cb-joint-highest}
applies.  A finite-dimensional unitary representation
is an orthogonal sum of irreducibles: choose a minimal
nonzero invariant subspace, take its invariant orthogonal
complement, and induct on dimension.
Lemma~\ref{cb-highest-classification} identifies every
summand of the restriction to $U(N-1)$ by a highest
vector.  Lemma~\ref{cb-joint-highest} allows exactly
the interlacing weights and at most one copy of each.
Conversely (CB4) provides a highest vector for every
interlacing weight.  The Gram-matrix argument in
Lemma~\ref{cb-highest-classification} identifies its
cyclic subspace with $V_\mu$, so every listed summand
does occur.  The last-row degree in that lemma proves
the $U(1)$ character.  Inequivalent unitary summands
are orthogonal: their cross-pairing defines an
intertwiner, which is zero by irreducibility and the
distinct highest weights.

Twisting back by $\det^{\lambda_N}$ adds $\lambda_N$
to all the top weights and to the last-circle weight.
It preserves (CB3) and changes the last exponent to
$|\lambda|-|\mu|$, as stated.  For $N=1$ (CB2) is
just the character $t\mapsto t^{\lambda_1}$.
\end{proof}

This proof also specifies the multiplicity-one normalization:
in the polynomial case there is a unique intertwiner
$J_{\lambda,\mu}:V_\mu\to V_\lambda$ with
$J_{\lambda,\mu}H_\mu=\xi_{\lambda,\mu}$.
The next section will address its integral denominator;
multiplicity one alone does not assert that this
normalization preserves the dual-Weyl lattice.


\section{The integral branching map and its exact normalization}
\label{ib-section}

Multiplicity one does not determine an integral normalization.
We prove integrality for the exact scalar used by the submitted
lowering map.  The ingredients are the preceding branching law,
the integral spanning theorem for initial minors
(Theorem~\ref{sm-initial-minors}), and the Capelli identities
of Section~\ref{rp-section}.  The central differential operators
will act on rows; this gives a direct projection argument.

First let $N=n+1$, $n\ge1$, let
$\lambda=(\lambda_1,\ldots,\lambda_n,0)$ be a partition,
and suppose $\lambda_i\ge\mu_i\ge\lambda_{i+1}$.
Use the polynomial modules and normalized Fock pairings
of Section~\ref{coef:section}.  Set
\[
 a_i=\mu_i-\lambda_{i+1},\quad b_i=\lambda_i-\mu_i,\quad
 d=\sum_i b_i,\quad D_0=\sum_i a_i,\quad
 D_1=\sum_i\min(\lambda_i-\lambda_{i+1},D_0).
 \tag{IB1}
\]
Define integer intervals and an integer $c$ by
\[
 \begin{split}
 R_i&=[\lambda_{i+1},\min(\lambda_i,\lambda_{i+1}+D_0)]\cap\mathbb Z,\\
 c&=\prod_{i=1}^n\ \prod_{v\in R_i\setminus\{\mu_i\}}
                  \prod_{j=1}^n\bigl((v-i)-(\mu_j-j)\bigr).
 \end{split}
 \tag{IB2}
\]
Empty products are one.  This is exactly $c_{\rm Br}$ in the
source: changing indices from zero to one cancels in each
difference.  Put $\xi=\xi_{\lambda,\mu}$ from (CB4), and let
$J:V_\mu\to V_\lambda$ be the intertwiner with $JH_\mu=\xi$.

\subsection{An integral lift of the required coefficient}

Write $\omega(\beta)=\sum_i i\beta_i$ for the weight of a
last-row monomial $z^\beta$, and write $[z^b]G$ for its
coefficient, a polynomial in the top block $A$.

\begin{lemma}\label{ib-integral-lift}
For every $f\in\nabla_\mu(\mathbb Z)$ there exists
$G\in\nabla_\lambda(\mathbb Z)$ of last-row degree $d$ with
\[
 G(A,z)=z^b f(A)+G_{<}(A,z),\qquad
 \omega(\beta)<\omega(b)
       \quad\text{for every last-row monomial }z^\beta\text{ of }G_{<}.
 \tag{IB3}
\]
In particular $[z^b]G=f$.
\end{lemma}
\begin{proof}
By Theorem~\ref{sm-initial-minors}, it suffices to lift a
product of initial-column minors of shape $\mu$; integer
linear combinations then give the assertion.
Induct on the number of nonconstant factors.
If there are none, $\mu=0$ and interlacing gives
$\lambda_2=\cdots=\lambda_N=0$.
The lift of $1$ is $z_1^{\lambda_1}$.

Choose a factor $M_I$ of height $s$, with rows
$I\subset\{1,\ldots,n\}$ and columns $1,\ldots,s$.
Thus $1\le s\le n$ and $\mu_s>\mu_{s+1}$, where
$\mu_{n+1}=0$.  Remove this factor to get shape
$\mu'=\mu-(1^s,0^{n-s})$.
If $\mu_s>\lambda_{s+1}$, subtract the same $s$
initial ones from $\lambda$ to obtain $\lambda'$.
Then $\lambda'$ interlaces $\mu'$: the only changed
lower boundary is $\lambda_{s+1}\le\mu_s-1$.
All subtracted entries are positive.  The other
inequalities follow by subtracting one from both
sides or leaving both sides unchanged.
Here $b'=\lambda'-\mu'=b$.
Multiply the inductively obtained lift by $M_I(A)$.

If $\mu_s=\lambda_{s+1}$, then $s<n$:
a height-$n$ factor would give $\mu_n\ge1>\lambda_N$.
Subtract $s+1$ initial ones from $\lambda$ to obtain
$\lambda'$.  Its upper inequality at position $s+1$ is
\[
              \lambda_{s+1}-1=\mu_s-1\ge\mu_{s+1}.
\]
At position $s$ the lower inequality is equality;
all other inequalities follow as before.
Thus $\lambda'$ interlaces $\mu'$ and
$b'=b-e_{s+1}\ge0$.
Append row $N$ to $I$ and take the initial-column
minor of height $s+1$.  Its last-row Laplace expansion is
\[
 \widetilde M_I(A,z)
   =z_{s+1}M_I(A)+
       \sum_{j=1}^{s}(-1)^{s+1+j}z_jM_{I,j}(A),
 \tag{IB4}
\]
where $M_{I,j}$ uses columns $1,\ldots,s+1$ with
column $j$ removed.  The leading last-row term has
weight $s+1$, and every other term has smaller weight.
Multiply this minor by the inductive lift.
Last-row weights add, proving (IB3).

Every factor used is an initial-column minor on $N$
rows, with integral coefficients and the required
column invariance.  Its shape adds to $\lambda$ and
its last-row degree adds to $d$.  This also verifies
membership and degree in both cases.
\end{proof}

\subsection{The exact integral projection}

\begin{lemma}\label{ib-capelli-projection}
The integer $c$ is nonzero.  On the last-row degree-$d$
subspace of $V_\lambda$, an operator with integer
differential coefficients preserves
$\nabla_\lambda(\mathbb Z)$ and equals $c$ times the
projection onto the $V_\mu$ summand in (CB7).
\end{lemma}
\begin{proof}
Use the top-row operators $E_{ab}$ of (CB1), with
$1\le a,b\le n$.  They preserve integral coefficients,
column degrees, column invariance, and last-row degree.
With indices starting at one write
\[
 C(u)=\sum_{\sigma\in\mathfrak S_n}\operatorname{sgn}(\sigma)
       \prod_{j=1}^n\bigl((u+j)\delta_{\sigma(j),j}
                                         -E_{\sigma(j),j}\bigr).
 \tag{IB5}
\]
The product is ordered by increasing $j$.
This is $B_n(u+1)$ from Corollary~\ref{rp-negative-convention}.
It is central and acts on $V_\eta$ by
$\chi_\eta(u)=\prod_j(u+j-\eta_j)$.
Put
\[
 Z=\prod_{i=1}^n\ \prod_{v\in R_i\setminus\{\mu_i\}} C(v-i).
 \tag{IB6}
\]
Its factors commute and preserve the integral polynomial module.

An interlacing summand of last-row degree $d$ has
$|\eta|=|\mu|$, so
$\sum_i(\eta_i-\lambda_{i+1})=D_0$.
Each summand is nonnegative, hence every $\eta_i\in R_i$.
If $\eta\ne\mu$, choose $i$ with $\eta_i\ne\mu_i$.
The included factor $C(\eta_i-i)$ has scalar zero on
$V_\eta$.  On $V_\mu$ the product has scalar exactly $c$.

Fix an included $v\in R_i\setminus\{\mu_i\}$.
The $j=i$ factor in (IB2) is nonzero.
If $j<i$, then
$\mu_j\ge\lambda_{j+1}\ge\lambda_i\ge v$, and
$v-i<\mu_j-j$.
If $j>i$, then
$v\ge\lambda_{i+1}\ge\lambda_j\ge\mu_j$, and
$v-i>\mu_j-j$.
Every factor is nonzero.
The branching law now proves the projection assertion.
\end{proof}

\begin{theorem}\label{ib-iota-integral}
The map $I=cJ$ is equivariant, sends $H_\mu$ to $c\xi$,
has last-row degree $d$, and satisfies
\[
                    I\nabla_\mu(\mathbb Z)\subset\nabla_\lambda(\mathbb Z).
 \tag{IB7}
\]
\end{theorem}
\begin{proof}
Take the integral lift $G$ of $f$ from
Lemma~\ref{ib-integral-lift}, and let $P_\mu$ be the
projection in Lemma~\ref{ib-capelli-projection}.
Coefficient extraction commutes with every top-row operator.
On $f\in V_\mu$, the operator (IB6) acts as $c$, so
\[
                    [z^b]P_\mu G=c^{-1}[z^b]ZG=f.
\]
The restriction of $[z^b]$ to $J(V_\mu)$ is the inverse
of $J$: it is equivariant and sends $\xi$ to $H_\mu$,
by (CB4) and its Laplace expansion.
The uniqueness in Lemma~\ref{cb-highest-classification}
therefore gives $[z^b]J=1$.
Consequently $P_\mu G=Jf$ and $If=ZG$, which is integral.
This also proves independence of the chosen lift.
The remaining properties follow from those of $J$.
\end{proof}

\subsection{Adjoints, the standard formula, and the norm}

Let $\Pi_\lambda,\Pi_\mu$ be the orthogonal Fock
projections onto the polynomial modules.  They have
rational matrices: choose rational bases and invert
their positive definite rational Gram matrices.
Set
\[
 C_\lambda=[H_\lambda,H_\lambda],\quad
 C_\mu=[H_\mu,H_\mu],\quad
 \kappa=\frac{[\xi,\xi]}{C_\lambda}>0,\qquad b!=\prod_i b_i!.
 \tag{IB8}
\]

\begin{theorem}\label{ib-standard-integral}
On the indicated polynomial modules define
\[
 \begin{split}
 \mathcal B F&=c\kappa\,\Pi_\mu([z^b]F),\\
 \mathcal I f&=\frac{c[\xi,\xi]}{b!C_\mu}\Pi_\lambda(z^b f).
 \end{split}
 \tag{IB9}
\]
Then $\mathcal I=I=cJ$ and $\mathcal B$ is its adjoint
for the normalized Fock pairings.  In particular
\[
 \mathcal I\nabla_\mu(\mathbb Z)\subset\nabla_\lambda(\mathbb Z),
 \qquad
 \mathcal B\Delta_\lambda(\mathbb Z)\subset\Delta_\mu(\mathbb Z).
 \tag{IB10}
\]
Both maps are nonzero and equivariant.  Their squared
operator norms for the normalized Hermitian Fock
metrics are exactly $c^2\kappa$.
These are the maps denoted $\operatorname{ioStd}$ and
$\operatorname{brStd}$ in the application.
\end{theorem}
\begin{proof}
Coefficient extraction commutes with $\operatorname{GL}_n$.
The projection $\Pi_\mu$ commutes with $U(n)$ by
unitarity, hence with its complexified Lie algebra
and with $\operatorname{GL}_n(\mathbb C)$.
Thus $\phi=\Pi_\mu[z^b]$ vanishes on every last-row
degree other than $d$; in degree $d$, multiplicity
one makes it vanish on every summand except
$J(V_\mu)$, on which $\phi J=1$.

The Gram-matrix calculation in
Lemma~\ref{cb-highest-classification} shows that the
pullback of the normalized Hermitian form under $J$
is a scalar multiple of the form on $V_\mu$.
Evaluation on $H_\mu$ gives the scalar $\kappa$.
Orthogonality of the summands therefore gives
$J^*=\kappa\phi$ and $I^*=\mathcal B$.
Their exact squared norms are $c^2\kappa>0$.

Comparing monomials in the top and bottom rows gives
\[
                         [F,z^b f]=b!\,[[z^b]F,f].
\]
After normalizing, the adjoint of $\phi$ is
$C_\lambda(b!C_\mu)^{-1}\Pi_\lambda(z^b\,\cdot)$.
Taking the adjoint of $c\kappa\phi$ proves the second
formula of (IB9), including every factorial.
All matrices are rational, so the same adjoint
identity holds for the rational bilinear forms.

The first containment in (IB10) is (IB7).
For $F\in\Delta_\lambda(\mathbb Z)$ and
$f\in\nabla_\mu(\mathbb Z)$ it follows that
\[
 B_\mu(\mathcal B F,f)=B_\lambda(F,\mathcal I f)\in\mathbb Z.
\]
Theorem~\ref{coef:integral-duality} gives
$\mathcal B F\in\Delta_\mu(\mathbb Z)$.
No prime dividing $b!$ or $C_\mu$ has been excluded.
\end{proof}

\begin{proposition}\label{ib-norm-bound}
Write $M=\lambda_1-\lambda_N$.  Then
\[
 |c|\le(M+n-1)^{nD_1},\qquad
 \kappa\le\dim V_\lambda,\qquad
 \|\mathcal B\|=\|\mathcal I\|
                  \le |c|\sqrt{\dim V_\lambda}.
 \tag{IB11}
\]
A zero base in the first power has exponent zero
and value $1$.  In particular these estimates imply
the source bound
\[
 \|\mathcal B\|,\|\mathcal I\|
       \le\bigl(2(M+N)\bigr)^{ND_1}\sqrt{\dim V_\lambda}.
 \tag{IB12}
\]
\end{proposition}
\begin{proof}
The interval $R_i$ contains $\mu_i$, and
$\#(R_i\setminus\{\mu_i\})
=\min(\lambda_i-\lambda_{i+1},D_0)$.
Both $v$ and $\mu_j$ in (IB2) lie between $\lambda_N$
and $\lambda_1$.  Every factor has absolute value
at most $M+n-1$, and there are $nD_1$ factors.
If $M+n-1=0$, then $n=1,M=0$ and $D_1=0$.

Normalize Haar measure on $U(N)$ to have mass one.
The highest vector $h=H_\lambda$ has normalized norm one.
Average the rank-one projection onto $\rho(g)h$.
The resulting self-adjoint operator commutes with
the group, so each of its eigenspaces is invariant.
Irreducibility makes it a scalar operator.
Its trace is one, hence its scalar is $(\dim V_\lambda)^{-1}$.
Consequently
\[
 \int_{U(N)}|\langle\xi,\rho(g)h\rangle_\lambda|^2\,dg
                         =\frac{\kappa}{\dim V_\lambda}.
\]
The evaluation identity of
Lemma~\ref{coef:extremal-weight} and unitarity identify
this matrix coefficient, up to conjugation and
inversion of $g$, with $\xi(g^t I_{N,n})$.
Each initial or replaced-row minor of a matrix with
orthonormal columns has absolute value at most one:
it is a coordinate of the unit exterior product of
those initial columns.  Formula (CB4) therefore
gives $|\xi(g^tI_{N,n})|\le1$.
The integral is at most one, proving the bound
on $\kappa$.  The exact norm formula proves the
last part of (IB11).  Finally $M+N\ge2$ and $n\le N$
give (IB12).
\end{proof}

\subsection{The local basis identity and a necessary scalar}

The statements extend to coefficient fields and tensor
products over embeddings.  The integral containments
also extend to characteristic-zero coefficient domains:
initial minors span the dual-Weyl module, the lift
and $Z$ have integer coefficients, and the perfect
Weyl duality extends with the coefficient ring.
In particular (IB10) imposes no exclusion of small primes.

For dominant integral $\lambda,\mu$ with the same
smallest entry $s$, subtract $s$ from both tuples.
Their canonical polynomial modules are those above;
the maps and integral pairings do not change.
The action acquires the common factor $\det^s$,
and the last-circle character becomes
$t^{d+s}=t^{|\lambda|-|\mu|}$.

\begin{corollary}\label{ib-diagonal-covariance}
Let $\lambda,\mu$ interlace and have the same smallest
entry.  Put $d_0=|\lambda|-|\mu|$ and
$\mathcal B_t=t^{-d_0}\mathcal B$ for $t\ne0$.
For $R_\lambda(g)=\rho_\lambda(g^{-t})$ and
$B=\operatorname{diag}(B',t^{-1})$ one has
\[
                R_\mu(B'^{-1})\mathcal B_t
                              =\mathcal B R_\lambda(B^{-1}).
 \tag{IB13}
\]
For tensor products the scalar is the product of
the corresponding scalars over the embeddings.
\end{corollary}
\begin{proof}
Equivariance and the last-circle character give
$\mathcal B\rho_\lambda(\operatorname{diag}(g,a))
=a^{d_0}\rho_\mu(g)\mathcal B$.
Now $R_\lambda(B^{-1})=\rho_\lambda(B^t)$ and
$B^t=\operatorname{diag}(B'^t,t^{-1})$.
Substitution proves the formula; tensor products
prove the final assertion.
In the arithmetic application $t=\overline{\pi_T}$,
as dictated by the Hermitian dual lattice.
\end{proof}

The scalar $c$ cannot simply be deleted.
Take $\lambda=(2,1,0)$ and $\mu=(2,0)$.
Write the two columns as $x=(x_1,x_2,x_3)^t$ and
$y=(y_1,y_2,y_3)^t$, and put $D_{ij}=x_i y_j-x_j y_i$.
Then $\xi=x_1D_{13}$ and
\[
 E_{21}\xi=x_1D_{23}+x_2D_{13}
                         =2x_2D_{13}-x_3D_{12}.
\]
Since $E_{21}(x_1^2)=2x_1x_2$, the unscaled map gives
\[
                  J(x_1x_2)=x_2D_{13}-\tfrac12x_3D_{12}.
\]
The coefficient of $x_1x_3y_2$ is $-1/2$.
Thus this is not an integral polynomial, although
$x_1x_2$ is in the integral dual-Weyl module.
Here $D_0=1$, $R_1=\{1,2\}$, $R_2=\{0,1\}$, and
(IB2) gives $c=(-2)(-2)=4$.
The actual map is $4J$, whose integrality is proved
in Theorem~\ref{ib-standard-integral}.
This refutes deletion of the normalization, not
the submitted integral branching assertion.

These results settle the polynomial branching input
and its local basis identity.  The actual twisted
automorphic classes, their coefficient lattices,
and their period sums are separate arithmetic statements.



% A mathematical section, to be included after the common preamble.
% Requires amsmath, amssymb, amsthm, booktabs and the theorem environments
% theorem, proposition, lemma, corollary, example, remark.
% All labels in this file have the prefix div:.
\section{Quotient metrics, local indices, and the divisor inequality}
\label{div:section}

We now prove the determinant and lattice assertions underlying the divisor
inequality. In particular, the inequality between a quotient covolume and
finite Gram determinants is derived, rather than taken as a hypothesis.
We also prove the operator estimates for a chain of rectangular matrices.
The source assertions being reconstructed occur in the quotient-metric
and local-to-global divisor constructions.\cite{alpo-uniform-source}
These results do not yet construct the particular automorphic lattices or
establish their uniform estimates. The remaining inputs are described at
the end of the section.

\subsection{Normalization and local indices}

Let $K$ be a totally imaginary number field, let $d=[K:\mathbb Q]$,
and put $r_2=d/2$. Write $\mathcal O=\mathcal O_K$ and let $v$ run over
the complex places, one from each pair of conjugate embeddings. For a
nonzero prime ideal $\mathfrak p$ put
\[
 w_{\mathfrak p}=\frac{\log N\mathfrak p}{d},\qquad
 \operatorname{ord}_{\mathfrak p}:K^\times\longrightarrow\mathbb Z.
\]
Valuations in this section are applied only to nonzero elements; a zero
determinant is never assigned a finite valuation. The product formula is
\begin{equation}\label{div:product-formula}
 \sum_{\mathfrak p}w_{\mathfrak p}\operatorname{ord}_{\mathfrak p}(a)
 =\frac2d\sum_{v\mid\infty}\log|v(a)|\qquad(a\in K^\times).
\end{equation}
Indeed, both sides multiplied by $d$ are
$\log|N_{K/\mathbb Q}(a)|$, computed from the factorization of the
principal fractional ideal and from the complex embeddings, respectively.

A full lattice in $K^q$ means a finitely generated $\mathcal O$-submodule
$L$ spanning $K^q$. Its determinant ideal $\mathfrak a(L)$ is the
fractional ideal generated by
$\det(x_1,\ldots,x_q)$ with $x_i\in L$, in the displayed coordinates.
Put $a_{\mathfrak p}(L)=\operatorname{ord}_{\mathfrak p}\mathfrak a(L)$.
For positive definite Hermitian matrices $G_v$ on $\mathbb C^q$, define
\begin{equation}\label{div:covolume}
 C_G(L)=N\mathfrak a(L)\prod_{v\mid\infty}\det G_v,
 \qquad c_G(L)=\frac2d\log C_G(L).
\end{equation}
Here the norm of a fractional ideal is positive and multiplicative.
The first expression is the real Minkowski covolume divided by
$2^{-r_2q}|\Delta_K|^{q/2}$, when the metric at a complex place is
the real part of its Hermitian form. Thus it removes exactly the
covolume of $\mathcal O^q$ with its standard metrics. The formula follows
first for a free lattice from the real determinant of a complex matrix,
which is the squared absolute value of its complex determinant, and then
for an arbitrary lattice by passing to a finite-index free sublattice.
Changing coordinates transforms the determinant ideal and the Gram
matrices by reciprocal factors; \eqref{div:product-formula} shows that
$C_G(L)$ is unchanged.

The usual normalized arithmetic degree and slope in this convention are
\[
 \widehat{\deg}(L)=-\tfrac12c_G(L),\qquad
 \mu(L)=-\frac{c_G(L)}{2q}\quad(q>0).
\]
This specifies the factor of two: a convention using logarithms of
squared norms has slope $2\mu(L)$. For the ordinary linear dual lattice
$L^\vee=\operatorname{Hom}_{\mathcal O}(L,\mathcal O)$ with dual metrics,
the determinant ideal and metric determinants are inverted. Consequently
\begin{equation}\label{div:ordinary-dual}
 c(L^\vee)=-c(L),\qquad 2\mu(L^\vee)=c(L)/q.
\end{equation}

For full lattices $A,B$ in the same $K$-space define the signed local index
\[
 i_{\mathfrak p}(A,B)=
 \operatorname{length}_{\mathcal O_{\mathfrak p}}
       \bigl(A_{\mathfrak p}/(A_{\mathfrak p}\cap B_{\mathfrak p})\bigr)
 -\operatorname{length}_{\mathcal O_{\mathfrak p}}
       \bigl(B_{\mathfrak p}/(A_{\mathfrak p}\cap B_{\mathfrak p})\bigr).
\]
The subscript denotes ordinary localization at the single prime.

\begin{lemma}[Index and covolume]\label{div:index}
For full lattices $A,B$ and fixed metrics $G$,
\begin{align}
 i_{\mathfrak p}(A,B)&=a_{\mathfrak p}(B)-a_{\mathfrak p}(A),
       \label{div:index-det}\\
 c_G(B)-c_G(A)&=2\sum_{\mathfrak p}w_{\mathfrak p}
                                  i_{\mathfrak p}(A,B).
       \label{div:index-covolume}
\end{align}
All sums are finite. The indices are additive:
$i_{\mathfrak p}(A,C)=i_{\mathfrak p}(A,B)+i_{\mathfrak p}(B,C)$.
\end{lemma}
\begin{proof}
Over the discrete valuation ring $\mathcal O_{\mathfrak p}$, choose a
basis of $A_{\mathfrak p}$. Elementary divisors give a basis in which
$B_{\mathfrak p}=\bigoplus_{j=1}^q\pi^{b_j}\mathcal O_{\mathfrak p}e_j$,
with $b_j\in\mathbb Z$. The two lengths in the definition are
$\sum_j\max(b_j,0)$ and $\sum_j\max(-b_j,0)$. Their difference is
$\sum_jb_j$, the valuation of the change-of-basis determinant. This
proves \eqref{div:index-det}. The determinant ideals have finite support,
and taking logarithms in \eqref{div:covolume} proves
\eqref{div:index-covolume}. Additivity follows from
\eqref{div:index-det}.
\end{proof}

\begin{lemma}[Determinantal ideals of rectangular matrices]\label{div:minors}
Let $A\in\operatorname{Mat}_{q\times n}(K)$ have rank $q$. Then
\[
 a_{\mathfrak p}(A\mathcal O^n)
 =\min_{\substack{I\subset\{1,\ldots,n\}\\|I|=q}}
           \operatorname{ord}_{\mathfrak p}\det A_I,
\]
where zero minors are omitted. More generally, for compatible matrices
$A,B$ and every $k$, the ideal generated by the $k$-minors of $AB$ is
contained in the product of the ideals generated by the $k$-minors of
$A$ and of $B$.
\end{lemma}
\begin{proof}
The exterior products of $q$ columns generate the top exterior power of
the image module, so its determinant ideal is precisely the ideal
generated by the maximal minors. An ideal over a discrete valuation
ring has valuation equal to the minimum valuation of any finite set of
nonzero generators. For the second assertion, Cauchy--Binet expresses
each $k$-minor of $AB$ as a sum of products of a $k$-minor of $A$ and a
$k$-minor of $B$. This proves the ideal inclusion, including cases in
which some minors vanish.
\end{proof}

\subsection{Hermitian duals and simultaneous localization}

Suppose that $K$ has an involution $a\mapsto\bar a$. This bar need not
act as complex conjugation at every embedding. Put
$\overline{\mathfrak p}=\{\bar a:a\in\mathfrak p\}$; then
$N\overline{\mathfrak p}=N\mathfrak p$ and
$\operatorname{ord}_{\mathfrak p}\bar a
 =\operatorname{ord}_{\overline{\mathfrak p}}a$.
Use $A^*=\bar A^{,t}$ for matrices over $K$ and ordinary conjugate
transpose for matrices over $\mathbb C$.

The ring relevant to a Hermitian local basis is
\[
 \mathcal O_{\{\mathfrak p,\overline{\mathfrak p}\}}
 =\mathcal O_{\mathfrak p}\cap
       \mathcal O_{\overline{\mathfrak p}}\subset K.
\]
It is obtained by inverting elements outside both primes and is stable
under the involution. A simultaneous local basis of $L$ at
$\mathfrak p,\overline{\mathfrak p}$ is a $K$-basis $e$ which is a basis
of both $L_{\mathfrak p}$ and $L_{\overline{\mathfrak p}}$.
Such a basis always exists. To see this, choose bases of the reductions
of the projective module $L$ modulo the two primes, choosing the same
basis if the primes coincide. The Chinese remainder theorem lifts these
to $q$ vectors of $L$. At each selected prime their determinant is a unit
in a local basis, so they form a basis there. This also proves their
$K$-linear independence. Equivalently, every finite projective module
of constant rank over this semilocal ring is free.

Let $H=H^*\in\operatorname{GL}_q(K)$, write
$h_H(x,y)=\bar x^{,t}Hy$, and define the Hermitian dual
\[
 L^\#=\{y\in K^q:h_H(x,y)\in\mathcal O\text{ for every }x\in L\},
 \qquad j_{\mathfrak p}(H,L)=i_{\mathfrak p}(L,L^\#).
\]
For a basis matrix $E$ with columns $e_1,\ldots,e_q$, its Gram
determinant is $g_H(e)=\det(E^*HE)\in K^\times$.

\begin{lemma}[The Hermitian index]\label{div:hermitian-index}
If $e$ is a simultaneous local basis of $L$ at
$\mathfrak p,\overline{\mathfrak p}$, then
\begin{align}
 \operatorname{ord}_{\mathfrak p}g_H(e)
 &=\operatorname{ord}_{\mathfrak p}\det H
       +a_{\mathfrak p}(L)+a_{\overline{\mathfrak p}}(L),
       \label{div:gram-valuation}\\
 j_{\mathfrak p}(H,L)&=-\operatorname{ord}_{\mathfrak p}g_H(e).
       \label{div:j-gram}
\end{align}
In particular $j_{\overline{\mathfrak p}}(H,L)=j_{\mathfrak p}(H,L)$.
\end{lemma}
\begin{proof}
The first identity is the determinant identity
$g_H(e)=\overline{\det E}\det H\det E$, using that $E$ represents $L$
at both primes. If $L^{\mathrm{lin},\vee}$ denotes the lattice dual
under the ordinary bilinear dot product, then
$L^\#=H^{-1}(\bar L)^{\mathrm{lin},\vee}$. It follows that
\[
 a_{\mathfrak p}(L^\#)
 =-\operatorname{ord}_{\mathfrak p}\det H
       -a_{\overline{\mathfrak p}}(L).
\]
Apply \eqref{div:index-det}. Finally $\det H$ is fixed by the
involution, so the right side of \eqref{div:gram-valuation} is invariant
under $\mathfrak p\leftrightarrow\overline{\mathfrak p}$.
\end{proof}

If $v(\bar H)=\overline{v(H)}$ for every embedding and every $v(H)$ is
positive definite, then the metrics $G_v=v(H)$ give the exact identity
\begin{equation}\label{div:rational-hermitian}
 c_H(L)=\sum_{\mathfrak p}w_{\mathfrak p}
                    \operatorname{ord}_{\mathfrak p}g_H(e_{\mathfrak p})
       =-\sum_{\mathfrak p}w_{\mathfrak p}j_{\mathfrak p}(H,L).
\end{equation}
Indeed, sum \eqref{div:gram-valuation}, use invariance of the weights
under conjugation, and apply \eqref{div:product-formula} to $\det H$.
The simultaneous bases may vary with $\mathfrak p$; the resulting
valuations do not.

\subsection{The dual exponent in the finite-group construction}

We can also derive the integer controlling the ambient dual lattice.
This uses the actual finite class sets, stabilizers and character
projectors occurring in the construction, rather than a bound for an
unspecified lattice.

\begin{lemma}[Dual of a saturated section]\label{div:cut-dual}
Let $(V,h)$ be a nondegenerate Hermitian $K$-space, let $N\subset V$
be a full lattice, and let $W\subset V$ be a nondegenerate subspace.
Let $\pi:V\to W$ be the orthogonal projection and put $N_W=N\cap W$.
For the restricted form on $W$ one has
\[
                       N_W^\#=\pi(N^\#).
\]
For the form $h/a$ on $W$, with $a\in\mathbb Q_{>0}$, the dual is
$a\pi(N^\#)$.
\end{lemma}
\begin{proof}
The inclusion from right to left follows from
$h(x,\pi z)=h(x,z)$ for $x\in W$.
For the reverse inclusion, $N_W$ is saturated in $N$: if a nonzero
$b\in\mathcal O$ satisfies $bx\in W$, then $x\in W$.
The quotient $N/N_W$ is consequently a finitely generated torsion-free
module over the Dedekind domain $\mathcal O$, hence projective. The
exact sequence splits. Given $y\in N_W^\#$, the conjugate-linear
functional $x\mapsto h(x,y)$ on $N_W$ extends, by zero on a chosen
complement, to an $\mathcal O$-valued conjugate-linear functional on
$N$. Extending scalars to $K$ and using nondegeneracy of $h$ gives a
vector $z\in N^\#$ which represents it. Thus
$h(x,\pi z-y)=0$ for all $x\in N_W$, and hence for all $x\in W$,
because $N_W$ spans $W$. Nondegeneracy on $W$ gives $\pi z=y$.
Finally $h(x,y)/a=h(x,y/a)$, proving the scaling assertion.
\end{proof}

For this construction fix an embedding $v_0:K\hookrightarrow\mathbb C$
which intertwines the involution with complex conjugation.
Let $\mathcal C$ be a finite set. For each $c\in\mathcal C$, let
$\Gamma_c$ be a finite group acting by isometries on a Hermitian
$K$-space $(V_c,h_c)$ which is positive definite at the distinguished
embedding $v_0$. Let $N_c,D_c$ be full lattices such that
\begin{equation}\label{div:class-lattices}
 D_c\subset N_c,\qquad \Gamma_cD_c=D_c,\qquad
                         sN_c^\#\subset D_c
\end{equation}
for one positive integer $s$. Form the space and its class-sum form
\[
 \mathcal V=\bigoplus_{c\in\mathcal C}V_c^{\Gamma_c},\qquad
 h(x,y)=\sum_{c\in\mathcal C}\frac{h_c(x_c,y_c)}{|\Gamma_c|},
\]
and the full lattices
\[
 N=\bigoplus_c(N_c\cap V_c^{\Gamma_c}),\qquad
 D'=\bigoplus_c(D_c\cap V_c^{\Gamma_c}).
\]
The restriction to each fixed subspace is nondegenerate, because it
is positive definite at $v_0$.

\begin{proposition}[Class-sum dual containment]\label{div:class-dual}
For the preceding form and lattices, $sN^\#\subset D'\subset N$.
\end{proposition}
\begin{proof}
The operator $\pi_c=|\Gamma_c|^{-1}\sum_{\gamma\in\Gamma_c}\gamma$
is the orthogonal projection onto $V_c^{\Gamma_c}$: it is idempotent,
its image is the fixed subspace, and its adjoint equals itself because
$\gamma^*=\gamma^{-1}$. Lemma~\ref{div:cut-dual} gives
\[
 (N_c\cap V_c^{\Gamma_c})^\#_{h_c/|\Gamma_c|}
          =|\Gamma_c|\pi_c(N_c^\#).
\]
Multiplying by $s$, the right side becomes
$\sum_{\gamma\in\Gamma_c}\gamma(sN_c^\#)$, which lies in $D_c$
by \eqref{div:class-lattices} and is fixed by $\Gamma_c$.
Taking the orthogonal direct sum proves the assertion. Notice that
no stabilizer-order factor survives in the result: the class-sum
weight cancels the denominator in the averaging projection.
\end{proof}

Suppose a finite group $G$ acts on $\mathcal V$ by isometries and
preserves $D'$. Let $\mathcal X$ be a finite set of distinct irreducible
complex characters of $G$ whose values belong to $K$, identified with
its image under $v_0$, let $P\subset G$ be a
subgroup, and let $\psi:P\to K^\times$ be a character. In $K[G]$ put
\begin{align*}
 e_{\mathcal X}
   &=\sum_{\chi\in\mathcal X}\frac{\chi(1)}{|G|}
                      \sum_{g\in G}\chi(g^{-1})g,\\
 e_{P,\psi}&=\frac1{|P|}\sum_{g\in P}\psi(g)^{-1}g,
       &e&=e_{\mathcal X}e_{P,\psi}.
\end{align*}
These formulas act on $\mathcal V$ through the specified representation.

\begin{lemma}[Integral denominators of the projectors]\label{div:projectors}
The operator $e$ is a self-adjoint idempotent. The coefficients of
$|G||P|e$ lie in $\mathcal O$. Its image is the intersection of the
sum of the $\chi$-isotypic subspaces, $\chi\in\mathcal X$, with the
subspace on which $P$ acts through $\psi$.
\end{lemma}
\begin{proof}
After extending scalars through $v_0$ to $\mathbb C$, for an irreducible
representation with character $\vartheta$,
the central element
$e_\chi=\chi(1)|G|^{-1}\sum_g\chi(g^{-1})g$ acts by a scalar,
by Schur's lemma. Taking its trace and using character orthogonality
shows that this scalar is $1$ for $\vartheta=\chi$ and $0$ otherwise.
One can obtain the needed orthogonality directly: averaging
$A\mapsto\rho_\vartheta(g)A\rho_\chi(g)^{-1}$ on the space of
homomorphisms between the two representations
is the projection onto the intertwining maps. Its trace is
$|G|^{-1}\sum_g\vartheta(g)\chi(g^{-1})$, and its rank is one or
zero by Schur's lemma. Complete reducibility follows by averaging a
projection onto any invariant subspace. This proves that the $e_\chi$
are mutually orthogonal central idempotents.

The finite sum for $e_{P,\psi}$ directly shows both idempotence and
that its image is the $\psi$-eigenspace. It commutes with the central
operator $e_{\mathcal X}$. At $v_0$, finite-group representations can
be made unitary by averaging a positive form. Hence
$\overline{\chi(g)}=\chi(g^{-1})$, as an identity in $K$, and
$\overline{\psi(g)}=\psi(g)^{-1}$. Since the given representation
is already unitary for $h$, the two idempotents are self-adjoint.
Their commuting product is the orthogonal projection onto the
intersection of their images.

Every $\chi(g)$ is an algebraic integer: the eigenvalues of a
finite-order matrix are roots of unity. The same is true of
$\psi(g)^{\pm1}$, and $\chi(1)$ is an integer.
Thus $|G|e_{\mathcal X}$ and $|P|e_{P,\psi}$ have integral
coefficients. Multiplication proves the asserted denominator bound.
\end{proof}

\begin{theorem}[The actual ambient dual exponent]\label{div:ambient-exponent}
Put $W=e\mathcal V$, $M=N\cap W$, and equip $W$ with the scaled form
$h_W=h|_W/|P|$. Then
\begin{equation}\label{div:ambient-containment}
                         |G|sM^\#\subset M.
\end{equation}
For any further algebraic quotient of $(W,h_W,M)$ as in
Proposition~\ref{div:local-bound}, its local Gram determinants therefore
satisfy
\[
 \operatorname{ord}_{\mathfrak p}g_\Gamma(e_{\mathfrak p})
             \le D\operatorname{ord}_{\mathfrak p}(|G|s).
\]
Consequently, for any finite set $S$ of primes,
\[
 \sum_{\mathfrak p\in S}w_{\mathfrak p}
             \operatorname{ord}_{\mathfrak p}g_\Gamma(e_{\mathfrak p})/D
                         \le\log(|G|s)\qquad(D>0).
\]
\end{theorem}
\begin{proof}
Lemma~\ref{div:cut-dual} gives $M^\#=|P|e(N^\#)$ for the scaled
form. Proposition~\ref{div:class-dual} and
Lemma~\ref{div:projectors} imply
\[
 |G|sM^\#=|G||P|e(sN^\#)\subset |G||P|e(D')\subset D'\subset N.
\]
The left side also lies in $W$, proving
\eqref{div:ambient-containment}. Apply
Proposition~\ref{div:local-bound}. Finally $|G|s$ is a positive
integer, all its finite valuations are nonnegative, and
\eqref{div:product-formula} gives
$\sum_{\mathfrak p}w_{\mathfrak p}
\operatorname{ord}_{\mathfrak p}(|G|s)=\log(|G|s)$.
This bounds every partial sum.
\end{proof}

These objects also admit the rational Gram matrices required above.
If the entries of the initial Gram matrices and representation matrices
lie in a subfield $K_0$ on which every embedding respects conjugation,
first enlarge $K_0$ by the finitely many character values and roots of
unity used in the projectors. These elements have the same property:
after any embedding
they are still finite-group character values and roots of unity, so
conjugation is inversion on their eigenvalues. The enlarged subfield
still has this property, since it is preserved by field operations.
The equations for the fixed subspaces and the projector matrices now
have entries in $K_0$. Gaussian elimination
therefore supplies bases over $K_0$, and restriction of the forms in
those bases has entries in $K_0$. Positivity at every embedding is
preserved under restriction to an independent set of columns and under
division by the positive integer $|P|$. This proves the requisite
properties of the resulting ambient matrices once they hold for the
initial coefficient forms and actions.

\begin{example}[A coefficient normalization cannot be omitted]
Let $K=\mathbb Q(i)$ and let
$N=\mathcal O x^2\oplus\mathcal O xy\oplus\mathcal O y^2$
in the space of binary quadratic forms. The Fock Hermitian form makes
these monomials orthogonal with squared norms $2,1,2$. Dividing the
form by the squared norm $2$ of the extremal monomial $x^2$ gives
Gram matrix $\operatorname{diag}(1,1/2,1)$, and direct calculation gives
\[
 N^\#=\mathcal O x^2\oplus2\mathcal O xy\oplus\mathcal O y^2
                  \subset N.
\]
For the unnormalized Fock form, however, the dual is
$\tfrac12\mathcal O x^2\oplus\mathcal O xy
 \oplus\tfrac12\mathcal O y^2$, which is not contained in $N$.
Indeed, its pairing with $x^2$ shows that $x^2/2$ lies in the dual,
whereas $x^2/2\notin N$.
Thus even this elementary coefficient space changes the admissible
dual exponent when the normalizing scalar is removed. This is a
counterexample to that simplification, not to the normalized
construction above.
\end{example}

\subsection{Quotient metrics and the Gram determinant bound}

Fix an embedding $v_0:K\hookrightarrow\mathbb C$ with
$v_0(\bar a)=\overline{v_0(a)}$ for every $a\in K$.
Let $H=H^*\in\operatorname{GL}_n(K)$ satisfy
\begin{equation}\label{div:H-condition}
 v(\bar H)=\overline{v(H)},\qquad v(H)>0
       \quad\text{for every complex embedding }v.
\end{equation}
The first condition holds, for example, when the entries lie in a
subfield on which all embeddings intertwine the involution with complex
conjugation. Let $P\in\operatorname{Mat}_{D\times n}(K)$ have rank $D$.
Define
\begin{equation}\label{div:quotient-matrices}
 G_v=\bigl(v(P)v(H)^{-1}v(P)^*\bigr)^{-1},\qquad
 \Gamma=(PH^{-1}P^*)^{-1}.
\end{equation}
The algebraic inverse defining $\Gamma$ exists: at $v_0$ its inverse
matrix is positive definite, hence has nonzero determinant. Also
$\Gamma=\Gamma^*$.

\begin{lemma}[The quotient metric]\label{div:quotient-metric}
For a positive definite Hermitian matrix $H_{\mathbb C}$ and a surjective
matrix $P_{\mathbb C}$, put
$G=(P_{\mathbb C} H_{\mathbb C}^{-1}P_{\mathbb C}^*)^{-1}$. Then
\[
 y^*Gy=\min_{P_{\mathbb C}x=y}x^*H_{\mathbb C}x.
\]
The unique minimizing vector is
$x_0=H_{\mathbb C}^{-1}P_{\mathbb C}^*Gy$ and is orthogonal to
$\ker P_{\mathbb C}$. A linear functional vanishing on this kernel
therefore has exactly the same operator norm before and after descent
to the quotient.
\end{lemma}
\begin{proof}
We have $P_{\mathbb C}x_0=y$, and for $z\in\ker P_{\mathbb C}$,
$z^*H_{\mathbb C}x_0=z^*P_{\mathbb C}^*Gy=0$.
Every other lift is $x_0+z$, whose squared norm is
$x_0^*H_{\mathbb C}x_0+z^*H_{\mathbb C}z\ge x_0^*H_{\mathbb C}x_0=y^*Gy$,
with equality only for $z=0$. The quotient map decreases norms, while
every vector has a lift of the same norm. These two facts give the
two inequalities between the operator norms of a functional and its
descent.
\end{proof}

\begin{proposition}[A local bound from the ambient dual]\label{div:local-bound}
Let $M\subset K^n$ be a full lattice, put $\Lambda=PM$, and let
$a\in\mathcal O\setminus\{0\}$ satisfy $aM^\#\subset M$, where the
dual is for $H$. Then, with the Hermitian dual of $\Lambda$ taken for
$\Gamma$,
\[
 a\Lambda^\#\subset\Lambda,
 \qquad
 \operatorname{ord}_{\mathfrak p}
       g_\Gamma(e_{\mathfrak p})
       \le D\operatorname{ord}_{\mathfrak p}(a)
       \quad\text{for every }\mathfrak p.
\]
In particular $a=\ell^g$, $\ell\ge1$, gives the bound
$gD\operatorname{ord}_{\mathfrak p}(\ell)$.
\end{proposition}
\begin{proof}
The algebraic map $J=H^{-1}P^*\Gamma$ satisfies $PJ=1$ and
\[
 h_H(x,Jy)=h_\Gamma(Px,y).
\]
Thus $y\in\Lambda^\#$ implies $Jy\in M^\#$. The hypothesis gives
$aJy\in M$, and applying $P$ gives $ay\in\Lambda$.
At $\mathfrak p$ the inclusion $a\Lambda^\#\subset\Lambda$ yields
\[
 D\operatorname{ord}_{\mathfrak p}(a)
       +a_{\mathfrak p}(\Lambda^\#)
       \ge a_{\mathfrak p}(\Lambda).
\]
Use Lemma~\ref{div:hermitian-index} to obtain the stated bound.
\end{proof}

\begin{lemma}[Determinant Cauchy--Schwarz]\label{div:det-cs}
For $D\times n$ complex matrices $U,V$,
\[
 |\det(UV^*)|^2\le\det(UU^*)\det(VV^*).
\]
\end{lemma}
\begin{proof}
The $D$th exterior power of $\mathbb C^n$ has the Hermitian inner
product determined by
$\langle u_1\wedge\cdots\wedge u_D,
v_1\wedge\cdots\wedge v_D\rangle
=\det(\langle u_i,v_j\rangle)$.
Apply ordinary Cauchy--Schwarz to the exterior products of the rows of
$U$ and $V$. Their squared norms are the two Gram determinants on the
right. This also proves the assertion when the row rank is deficient.
\end{proof}

\begin{theorem}[Covolume of an algebraic quotient]\label{div:quotient-covolume}
Under \eqref{div:H-condition}, for a full lattice $M\subset K^n$ and
$\Lambda=PM$, choose a simultaneous local basis $e_{\mathfrak p}$ of
$\Lambda$ at each conjugate pair of primes. Then
\begin{equation}\label{div:quotient-inequality}
 c_G(\Lambda)\le
 \sum_{\mathfrak p}w_{\mathfrak p}
       \operatorname{ord}_{\mathfrak p}g_\Gamma(e_{\mathfrak p}).
\end{equation}
The sum has finite support. The difference between the right and left
sides is independent of $M$ and is nonnegative. If the involution acts
as complex conjugation at every embedding of $K$, equality holds.
\end{theorem}
\begin{proof}
On the set of all complex embeddings define
$\tau(v)(a)=\overline{v(\bar a)}$. This is an involution commuting
with conjugation of embeddings. By \eqref{div:H-condition},
$\tau(v)(H)=v(H)$. Set
\[
 A_v=v(P)v(H)^{-1}v(P)^*,\qquad
 B_v=v(P)v(H)^{-1}\tau(v)(P)^*.
\]
Thus $G_v=A_v^{-1}$ and $v(\Gamma)=B_v^{-1}$.
The matrices $B_v$ are invertible, because the algebraic matrix inverted
in \eqref{div:quotient-matrices} is invertible. Apply
Lemma~\ref{div:det-cs} to
$U=v(P)v(H)^{-1/2}$ and $V=\tau(v)(P)v(H)^{-1/2}$:
\[
 |\det B_v|^2\le\det A_v\det A_{\tau(v)}.
\]
After inversion and logarithms this becomes
\[
 \log\det G_v+\log\det G_{\tau(v)}
       \le2\log|v(\det\Gamma)|.
\]
Sum over all embeddings. Since $\tau$ is a permutation, and conjugate
embeddings give equal real terms, we obtain
\begin{equation}\label{div:quotient-arch}
 \sum_{v\mid\infty}\log\det G_v
       \le\sum_{v\mid\infty}\log|v(\det\Gamma)|.
\end{equation}
By \eqref{div:gram-valuation}, the proposed upper bound equals
\[
 \sum_{\mathfrak p}w_{\mathfrak p}
       \operatorname{ord}_{\mathfrak p}\det\Gamma
       +2\sum_{\mathfrak p}w_{\mathfrak p}a_{\mathfrak p}(\Lambda).
\]
Both sums have finite support. Apply \eqref{div:product-formula} to
$\det\Gamma$ and then \eqref{div:quotient-arch}. This is exactly
\eqref{div:quotient-inequality}, by \eqref{div:covolume}.

More explicitly the difference is
\begin{equation}\label{div:quotient-defect}
 \delta(P,H)=\frac1{2d}\sum_{v:K\hookrightarrow\mathbb C}
 \log\frac{\det A_v\det A_{\tau(v)}}{|\det B_v|^2}\ge0.
\end{equation}
This expression contains no lattice. If $\tau$ fixes every embedding,
then $A_v=B_v$, so every summand is zero.
\end{proof}

\begin{example}[Strictness over a field which is not CM]\label{div:strict}
Let $\alpha=\sqrt[3]{2}$, let $\zeta$ be a primitive cube root of unity,
and take $K=\mathbb Q(\alpha,\zeta)$ with its usual complex conjugation
at the embedding where $\alpha>0$. Put $H=I_2$, $M=\mathcal O_K^2$,
and $P=(1,\alpha)$. Then $\Lambda=\mathcal O_K$ and
\[
 c_G(\Lambda)=-\log(1+2^{2/3}),\qquad
 \sum_{\mathfrak p}w_{\mathfrak p}
       \operatorname{ord}_{\mathfrak p}g_\Gamma(1)
       =-\tfrac13\log5.
\]
The inequality is strict.
\end{example}
\begin{proof}
The field has degree six, no real embeddings, and three complex places.
At every embedding $|v(\alpha)|=\alpha$, so
$G_v=(1+\alpha^2)^{-1}$. Algebraically
$\Gamma=(1+\alpha^2)^{-1}$, but its other embeddings are not all equal
to these positive metric matrices. Since
\[
 \prod_{k=0}^2(1+\alpha^2\zeta^{2k})=1+\alpha^6=5,
 \qquad |N_{K/\mathbb Q}(1+\alpha^2)|=25,
\]
the two formulas follow from \eqref{div:covolume} and
\eqref{div:product-formula}. Strictness follows from
$(1+\alpha^2)^3>1+\alpha^6=5$.
The fixed field of conjugation is $\mathbb Q(\alpha)$, which is not
totally real; hence $K$ is not a CM field. This example satisfies all
hypotheses of Theorem~\ref{div:quotient-covolume}. It disproves replacing
its inequality by an equality in the stated generality.
\end{proof}

\subsection{Patching the lattices and an exact divisor formula}

For a $K$-linear isomorphism $f$ between $q$-dimensional spaces with
metrics $G_v,G'_v$, write
\[
 d_v(f)=\frac{|\det v(f)|^2\det G'_v}{\det G_v},\qquad
 A(f)=\frac2d\sum_{v\mid\infty}\log d_v(f).
\]
These are coordinate independent. Directly from the product formula,
\begin{equation}\label{div:image-covolume}
 c_{G'}(fL)=c_G(L)+A(f).
\end{equation}

Here are the data for the formula. They are explicit finite-dimensional
objects, rather than an assumed divisor inequality.
There is a parameter set $\mathcal T$ with a distinguished element $0$.
For each $t$ there is a $D$-dimensional Hermitian $K$-space $Q_t$, a full
lattice $\Lambda_t\subset Q_t$, and a $K$-linear isomorphism
$\iota_t:Q_t\to E$ to one fixed Hermitian space $E$.
The space $Q_t$ and lattice $\Lambda_t$ arise from matrices and lattices
as in \eqref{div:quotient-matrices}: $\Lambda_t=P_tM_t$ with quotient
metrics $G_{t,v}$ and algebraic form $\Gamma_t$.
It is enough to require the global condition
\eqref{div:H-condition} for $t=0$; for other $t$ only the equality
$v_0(\Gamma_t)=G_{t,v_0}$, nondegeneracy and Hermitian symmetry are
needed below.

For each $j$ in a finite set $J$, there is a fixed Hermitian space $W_j$
of dimension $d_j$ and a full lattice $M_j\subset W_j$. For every $t$
there is a Hermitian space $T_{j,t}$ of dimension $d_j$, with full lattice
$L_{j,t}$, algebraic nondegenerate Hermitian form $H_{j,t}$, and
isomorphism $\theta_{j,t}:T_{j,t}\to W_j$.
At $v_0$ its metric is $v_0(H_{j,t})$. For $t=0$ the metric is
$v(H_{j,0})$ at every place and the entries of $H_{j,0}$ satisfy the
conjugation condition in \eqref{div:H-condition}.

Let $\kappa_v>0$, and let $\rho_t\in K^\times$ satisfy
$\rho_0=1$ and $v_0(\rho_t)>0$. The compatibility conditions on
determinants are
\begin{align}
 d_v(\iota_0)&=\kappa_v^D\prod_{j\in J}d_v(\theta_{j,0})
                       &&(v\mid\infty),\label{div:det-base}\\
 d_{v_0}(\iota_t)&=(\kappa_{v_0}v_0(\rho_t))^D
                         \prod_{j\in J}d_{v_0}(\theta_{j,t})
                       &&(t\in\mathcal T).\label{div:det-local}
\end{align}
Below we also prove the spectral calculation which produces such
determinant identities.

Choose parameters $t_{\mathfrak p}$ such that
\begin{equation}\label{div:patch-parameters}
 t_{\overline{\mathfrak p}}=t_{\mathfrak p},\qquad
 t_{\mathfrak p}=0\text{ except at finitely many primes},\qquad
 \operatorname{ord}_{\mathfrak p}(\rho_{t_{\mathfrak p}})=0.
\end{equation}
Put $R_t=\iota_t\Lambda_t$. There is a unique full lattice $R\subset E$
with $R_{\mathfrak p}=(R_{t_{\mathfrak p}})_{\mathfrak p}$ for every
$\mathfrak p$. For completeness, choose a nonzero integer $a$ such that
$aR_0\subset R_t\subset a^{-1}R_0$ for the finitely many lattices in
question. Enlarge the exceptional set to include the primes over $a$.
The finite module $a^{-1}R_0/aR_0$ is the direct sum of its primary
components. In each component select the submodule specified by the
desired local lattice, and take its inverse image in $a^{-1}R_0$.
This constructs $R$ and proves its asserted localizations. Uniqueness
follows because a finitely generated torsion module is zero if all its
localizations are zero.

Choose simultaneous local bases $e_{t,\mathfrak p}$ of $\Lambda_t$ and
$f_{j,t,\mathfrak p}$ of $L_{j,t}$. Set
\begin{align*}
 g_{t,\mathfrak p}&=g_{\Gamma_t}(e_{t,\mathfrak p}),\\
 I_{j,t,\mathfrak p}&=i_{\mathfrak p}(M_j,\theta_{j,t}L_{j,t}),
       &J_{j,t,\mathfrak p}&=j_{\mathfrak p}(H_{j,t},L_{j,t}),\\
 F_{t,\mathfrak p}&=\operatorname{ord}_{\mathfrak p}g_{t,\mathfrak p}
       +\sum_{j\in J}\bigl(I_{j,t,\mathfrak p}
              +I_{j,t,\overline{\mathfrak p}}+J_{j,t,\mathfrak p}\bigr),
       &\overline{\log\kappa}&=\frac1{r_2}\sum_{v\mid\infty}\log\kappa_v.
\end{align*}
The letter $F$ here denotes an explicitly defined local number, not an
additional hypothesis.

\begin{theorem}[Exact divisor formula]\label{div:patch-formula}
With the preceding data,
\begin{equation}\label{div:exact-divisor}
 c_E(R)=D\,\overline{\log\kappa}+\sum_{j\in J}c_{W_j}(M_j)
     +\sum_{\mathfrak p}w_{\mathfrak p}F_{t_{\mathfrak p},\mathfrak p}
     -\delta(P_0,H_0),
\end{equation}
where $\delta(P_0,H_0)\ge0$ is given explicitly by
\eqref{div:quotient-defect}. All sums over finite primes have finite
support. Dropping this nonnegative last term gives the divisor inequality.
\end{theorem}
\begin{proof}
Fix $\mathfrak p$ and put $t=t_{\mathfrak p}$. Let $\gamma$ be the
change-of-basis matrix from $\iota_0e_{0,\mathfrak p}$ to
$\iota_te_{t,\mathfrak p}$, and let $\delta_j$ be the corresponding
matrix from $\theta_{j,0}f_{j,0,\mathfrak p}$ to
$\theta_{j,t}f_{j,t,\mathfrak p}$.
Write $h_{j,t}=g_{H_{j,t}}(f_{j,t,\mathfrak p})$.
Changing Gram determinants at $v_0$ multiplies them by
$|v_0(\det\gamma)|^2$ and $|v_0(\det\delta_j)|^2$, respectively.
Apply \eqref{div:det-base} at $v_0$ and
\eqref{div:det-local}, substitute these change-of-basis identities,
and cancel the nonzero common factors. The resulting identity is
\[
 \overline{\det\gamma}\det\gamma\,g_{0,\mathfrak p}
                  \prod_jh_{j,t}
 =\rho_t^D g_{t,\mathfrak p}
                  \prod_j\bigl(\overline{\det\delta_j}
                                    \det\delta_j\,h_{j,0}\bigr).
\]
It is an identity in $K$: its image under the injective embedding $v_0$
has just been proved, and $v_0$ respects the involution.

Take $\operatorname{ord}_{\mathfrak p}$, use the last condition of
\eqref{div:patch-parameters}, and apply Lemmas~\ref{div:index} and
\ref{div:hermitian-index}. Because the bases are simultaneous at the
two primes, this gives
\begin{equation}\label{div:local-difference}
 F_{t,\mathfrak p}-F_{0,\mathfrak p}
 =i_{\mathfrak p}(R_0,R_t)
       +i_{\overline{\mathfrak p}}(R_0,R_t).
\end{equation}
By the first condition of \eqref{div:patch-parameters}, the two local
indices on the right are also those of $R_0,R$.
The terms $F_{0,\mathfrak p}$ have finite support, by
\eqref{div:gram-valuation} and the finite support of determinant ideals.
Equation \eqref{div:local-difference} and Lemma~\ref{div:index} therefore
give finite support for the selected terms as well. Summing and pairing
conjugate primes yields
\begin{equation}\label{div:patch-sum}
 \sum_{\mathfrak p}w_{\mathfrak p}F_{t_{\mathfrak p},\mathfrak p}
 =\sum_{\mathfrak p}w_{\mathfrak p}F_{0,\mathfrak p}
       +c_E(R)-c_E(R_0).
\end{equation}

At the base parameter, \eqref{div:rational-hermitian},
\eqref{div:index-covolume} and \eqref{div:image-covolume} give
\begin{align*}
 &\sum_{\mathfrak p}w_{\mathfrak p}
     (I_{j,0,\mathfrak p}+I_{j,0,\overline{\mathfrak p}}
                                   +J_{j,0,\mathfrak p})\\
 &\hspace{10mm}=c_{W_j}(\theta_{j,0}L_{j,0})-c_{W_j}(M_j)
                                     -c_{T_{j,0}}(L_{j,0})
  =A(\theta_{j,0})-c_{W_j}(M_j).
\end{align*}
Moreover \eqref{div:det-base} implies
$A(\iota_0)=D\overline{\log\kappa}+\sum_jA(\theta_{j,0})$.
Substitute these two identities and
$c_E(R_0)=c_{Q_0}(\Lambda_0)+A(\iota_0)$ into
\eqref{div:patch-sum}. The remaining difference
$\sum_{\mathfrak p}w_{\mathfrak p}
\operatorname{ord}_{\mathfrak p}g_{0,\mathfrak p}
-c_{Q_0}(\Lambda_0)$ is exactly $\delta(P_0,H_0)$ by
Theorem~\ref{div:quotient-covolume}. This proves the formula.
\end{proof}

\begin{corollary}[Insertion of local estimates]\label{div:scalar-divisor}
Assume $D>0$, and let $S$ contain the support of the selected local
terms. Abbreviate $I_{j,\mathfrak p}=I_{j,t_{\mathfrak p},\mathfrak p}$
and $J_{j,\mathfrak p}=J_{j,t_{\mathfrak p},\mathfrak p}$.
Suppose that real numbers satisfy
\begin{align*}
 \operatorname{ord}_{\mathfrak p}g_{t_{\mathfrak p},\mathfrak p}
       &\le Dg_{\mathfrak p},
 &\sum_{\mathfrak p\in S}w_{\mathfrak p}g_{\mathfrak p}&\le L_G,\\
 \sum_j(I_{j,\mathfrak p}+I_{j,\overline{\mathfrak p}})
       &\le D\left(\sum_j(2f_{j,\mathfrak p}+e_{j,\mathfrak p})
                                      +2/\ell_{\mathfrak p}^2\right),\\
 J_{j,\mathfrak p}&\le Dq_{j,\mathfrak p}+\eta_{j,\mathfrak p},
 &c_{W_j}(M_j)/D&\le L_j,\\
 \sum_{\mathfrak p\in S}w_{\mathfrak p}
                     (e_{j,\mathfrak p}+q_{j,\mathfrak p})&\le E_{1j},
 &\sum_{\mathfrak p\in S}w_{\mathfrak p}\eta_{j,\mathfrak p}/D&\le E_{2j},\\
 \sum_{\mathfrak p\in S}2w_{\mathfrak p}/\ell_{\mathfrak p}^2&\le E_3.
\end{align*}
Then
\begin{equation}\label{div:final-divisor}
 2\mu(R^\vee)\le\overline{\log\kappa}+\sum_jL_j
   +2\sum_j\sum_{\mathfrak p\in S}w_{\mathfrak p}f_{j,\mathfrak p}
   +L_G+\sum_j(E_{1j}+E_{2j})+E_3.
\end{equation}
\end{corollary}
\begin{proof}
Divide \eqref{div:exact-divisor} by $D$, discard $-\delta/D\le0$,
and insert the displayed estimates. Every finite-prime weight is
nonnegative, so the directions of all inequalities are preserved.
Interchange the two finite sums and use \eqref{div:ordinary-dual}.
\end{proof}

\subsection{A uniform estimate for the elementary prime error}

The term $E_3$ in the preceding corollary has a uniform numerical bound
which can be proved without any prime-distribution theorem.

\begin{proposition}\label{div:prime-error}
For a prime ideal $\mathfrak p$ above the rational prime $p$, let
$e_{\mathfrak p}=\operatorname{ord}_{\mathfrak p}(p)$ and
$\ell_{\mathfrak p}=p/\sqrt{e_{\mathfrak p}}$. Then
\[
 \sum_{\mathfrak p}\frac{2w_{\mathfrak p}}{\ell_{\mathfrak p}^2}
 =2\sum_{p\ {
m prime}}\frac{\log p}{p^2}
 <\frac{99841}{100000}<1.
\]
Every sum over a finite set of prime ideals obeys the same upper bound.
\end{proposition}
\begin{proof}
Write $N\mathfrak p=p^{f_{\mathfrak p}}$. The identity
$\sum_{\mathfrak p\mid p}e_{\mathfrak p}f_{\mathfrak p}=d$ gives the
first equality; all summands are positive.
For the numerical bound let $f(x)=\log x/x^2$, which is decreasing for
$x\ge2$. Every prime larger than $200$ is $1$ or $5$ modulo $6$.
For each of these two arithmetic progressions, comparison with its
first term and the subsequent integral gives
\[
 \sum_{p>200}f(p)
 \le2f(200)+\frac13\int_{200}^{\infty}f(x)\,dx
 =\frac{2\log200}{200^2}+\frac{\log200+1}{600}
 <\frac{2153}{200000},
\]
using $\log200<53/10$.
The finite sum is bounded by the following rational table. The last
column is an upper bound in units of $10^{-6}$.
\begin{center}\small
\begin{tabular}{lll}
\toprule
Interval&Primes&$10^6\sum\log p/p^2$\quad upper bound\\
\midrule
$[2,40]$&$2,3,5,7,11,13,17,19,23,29,31,37$&$468542$\\
$[41,80]$&$41,43,47,53,59,61,67,71,73,79$&$12964$\\
$[81,120]$&$83,89,97,101,103,107,109,113$&$3757$\\
$[121,160]$&$127,131,137,139,149,151,157$&$1753$\\
$[161,200]$&$163,167,173,179,181,191,193,197,199$&$1424$\\
\bottomrule
\end{tabular}
\end{center}
Here is a purely rational way to verify every entry, including the
bound for $\log200$. For $0\le z\le1/3$ put
\[
 U(z)=2\sum_{k=0}^{15}\frac{z^{2k+1}}{2k+1}
                   +\frac{2z^{33}}{33(1-z^2)}.
\]
The series obtained by integrating
$2/(1-z^2)=2\sum_{k\ge0}z^{2k}$ term by term shows that
$\log((1+z)/(1-z))\le U(z)$: bound every denominator in the remainder
by $33$ and sum the resulting geometric series.
For an integer $n\ge1$, choose $k$ with $2^k\le n<2^{k+1}$ and set
$z=(n-2^k)/(n+2^k)$. Then
\[
 \log n\le kU(1/3)+U(z).
\]
Substitution in each entry of the table proves the displayed rational
upper bound by integer arithmetic, and substitution of $n=200$ gives
$kU(1/3)+U(z)<53/10$. The accompanying short program uses precisely
these fractions, with no floating-point comparisons.
Adding the table gives
$\sum_{p\le200}\log p/p^2<488440/10^6=12211/25000$.
Together with the tail this is $99841/200000$; multiplication by two
proves the result.
\end{proof}

\subsection{Weighted operator matrices and rectangular chains}

We next give the norm bounds used by the chain construction, including
the dependence on the number of classes. This isolates the exact
pointwise estimate that must be proved for each actual kernel.

\begin{lemma}[A weighted operator matrix]\label{div:weighted-matrix}
Let $S,T$ be finite sets, let $U_s,V_t$ be complex normed spaces, and
let $\alpha_s,\beta_t>0$. For continuous linear maps
$A_{ts}:U_s\to V_t$ define
\[
 (\mathcal A x)_t=\sum_{s\in S}\alpha_s A_{ts}x_s,
 \qquad \|x\|_\alpha^2=\sum_s\alpha_s\|x_s\|^2,
 \qquad \|y\|_\beta^2=\sum_t\beta_t\|y_t\|^2.
\]
If $\|A_{ts}\|^2\le a_{ts}^2$, with $a_{ts}\ge0$, and
$M_{ts}=\sqrt{\beta_t\alpha_s}\,a_{ts}$, then
\begin{equation}\label{div:weighted-sharp}
 \|\mathcal A\|\le\|M\|_{\ell^2(S)\to\ell^2(T)}
       \le\left(\sum_{t,s}\beta_t\alpha_s a_{ts}^2\right)^{1/2}.
\end{equation}
In particular, if $\|A_{ts}\|^2\le A$, then
\begin{equation}\label{div:weighted-coarse}
 \|\mathcal A\|^2\le A\left(\sum_s\alpha_s\right)
                           \left(\sum_t\beta_t\right).
\end{equation}
If $S=\coprod_{a\in\mathcal A_0}S_a$ and
$T=\coprod_{b\in\mathcal A_1}T_b$, with
$\alpha_s,\beta_t\le1$, $|S_a|\le h$, $|T_b|\le h'$, this yields
$\|\mathcal A\|^2\le A|\mathcal A_0||\mathcal A_1|hh'$.
\end{lemma}
\begin{proof}
Put $u_s=\sqrt{\alpha_s}\|x_s\|$. The triangle inequality gives
\[
 \sqrt{\beta_t}\| (\mathcal A x)_t\|
       \le\sum_sM_{ts}u_s.
\]
Take the Euclidean norm in $t$ and use the definition of $\|M\|$.
The operator norm of a scalar matrix is at most its Hilbert--Schmidt
norm, proving \eqref{div:weighted-sharp}. The remaining claims follow
by summing the weights and the cardinality bounds. The argument includes
empty index sets: the corresponding map and sums are zero. Factors of
absolute value one in $A_{ts}$ do not affect any bound.
\end{proof}

The first bound in \eqref{div:weighted-sharp} is sometimes substantially
better than the class-count bound. Thus no enlargement of the number of
classes is intrinsically required in this part of the argument.

Suppose now that $V_0,\ldots,V_m$ are Hermitian $K$-spaces of possibly
different dimensions, and $K_j:V_{j-1}\to V_j$ are $K$-linear maps.
At each complex place suppose there are norm-preserving realizations
$\Phi_{j,v}$ into weighted spaces as above and exact identities
\[
 \Phi_{j,v}(K_{j,v}x)=\mathcal A_{j,v}\Phi_{j-1,v}(x).
\]
Linearity or surjectivity of $\Phi_{j,v}$ is not needed for the following
conclusion; preservation of norms and the displayed identity suffice.
Lemma~\ref{div:weighted-matrix} gives $\|K_{j,v}\|^2\le C_{j,v}$ with
\[
 C_{j,v}=A_{j,v}\left(\sum_s\alpha_{j-1,s}\right)
                          \left(\sum_t\alpha_{j,t}\right),
\]
or with the sharper scalar-matrix bound. Consequently
\begin{equation}\label{div:chain}
 \|K_{m,v}\cdots K_{1,v}\|^2\le\prod_{j=1}^m C_{j,v}.
\end{equation}
This follows by applying the bound successively to a vector. It requires
no square matrices, determinant nonvanishing, or equality of dimensions.
The empty chain is the identity.

\begin{lemma}[Spectral determinant and lower norm]\label{div:spectral}
Let $\iota:Q\to E$ be an isomorphism of finite-dimensional complex
Hermitian spaces, and let a finite nonempty family of subspaces $Q_p$
span $Q$. Suppose
$s_p>0$ and
\begin{equation}\label{div:spectral-cross}
 \langle\iota x,\iota y\rangle_E
       =s_p\langle x,y\rangle_Q
             \quad(x\in Q_p,\ y\in Q).
\end{equation}
Then
$\|\iota x\|_E^2\ge(\min_ps_p)\|x\|_Q^2$ for every $x\in Q$.
If the sum of the $Q_p$ is direct and $d_p=\dim Q_p$, then
\[
 d(\iota)=\prod_ps_p^{d_p}.
\]
In particular, if $s_p=\kappa\prod_j\lambda_{j,p}$ and
$\theta_j:T_j\to W_j$ has a direct-sum decomposition into subspaces
of dimensions $d_p$ on which its pullback metric has scalar
$\lambda_{j,p}$ in the sense of \eqref{div:spectral-cross}, then
\[
 d(\iota)=\kappa^{\dim Q}\prod_jd(\theta_j).
\]
\end{lemma}
\begin{proof}
Let $B=\iota^*\iota$, the positive self-adjoint operator on $Q$.
Equation \eqref{div:spectral-cross}, for all $y$, says exactly that
$Bx=s_px$ on $Q_p$. Since these spaces span $Q$, all eigenvalues of
$B$ occur among the $s_p$, proving the lower norm estimate by the
spectral theorem. If the sum is direct, concatenate bases of its
summands; the matrix of $B$ is diagonal by blocks, so
$\det B=\prod_ps_p^{d_p}$. This is the determinant ratio. The final
claim follows by substituting $s_p$ and applying the same calculation
to each $\theta_j$.
\end{proof}

\begin{proposition}[Height of a descended pairing]\label{div:pairing-height}
Let $P:V\to Q$ be surjective and give $Q$ its quotient metric. Suppose
$\iota:Q\to E$ satisfies Lemma~\ref{div:spectral} with
$s_{p,v}=\kappa_v\prod_j\lambda_{j,p,v}>0$ at each infinite place.
Let $K_m\cdots K_1:U\to U_m$ be a rectangular chain and let
$\mathcal P:U_m\times V\to K$ and $\mathcal B:U\times E\to K$
be bilinear pairings satisfying the exact identity
\begin{equation}\label{div:pairing-factorization}
 \mathcal B(x,\iota(Pz))
       =\mathcal P(K_m\cdots K_1x,z)
                         \quad(x\in U,z\in V).
\end{equation}
If $\|\mathcal P_v\|^2\le e^{h_{P,v}}$ and
$\|K_{j,v}\|^2\le e^{h_{j,v}}$, then
\begin{equation}\label{div:pairing-bound}
 \|\mathcal B_v\|^2
 \le\exp\left(h_{P,v}+\sum_jh_{j,v}
            -\log\kappa_v-\min_p\sum_j\log\lambda_{j,p,v}\right).
\end{equation}
For $\mathcal B\ne0$, define its squared-norm height by
$h_2(\mathcal B)=r_2^{-1}\sum_v\log\|\mathcal B_v\|^2$.
Then averaging \eqref{div:pairing-bound} gives
\[
 h_2(\mathcal B)\le\frac1{r_2}\sum_v
       \left(h_{P,v}+\sum_jh_{j,v}
                   -\min_p\sum_j\log\lambda_{j,p,v}\right)
       -\overline{\log\kappa}.
\]
\end{proposition}
\begin{proof}
Fix $v$, $x\in U_v$, $e\in E_v$, and put $q=\iota_v^{-1}e$.
By Lemma~\ref{div:quotient-metric} there is a lift $z$ of $q$ with
$\|z\|_V=\|q\|_Q$. Lemma~\ref{div:spectral} gives
\[
 \|q\|_Q\le\frac{\|e\|_E}
           {\sqrt{\kappa_v\min_p\prod_j\lambda_{j,p,v}}}.
\]
Extend \eqref{div:pairing-factorization} by complex bilinearity and
use the operator bounds and submultiplicativity:
\[
 |\mathcal B_v(x,e)|
 \le\frac{\|\mathcal P_v\|\prod_j\|K_{j,v}\|}
           {\sqrt{\kappa_v\min_p\prod_j\lambda_{j,p,v}}}
          \|x\|_U\|e\|_E.
\]
Square and use positivity to take logarithms of the denominator.
Since logarithm is increasing, the logarithm of the minimum of these
positive products is the minimum of their sums of logarithms.
If $\mathcal B\ne0$, a nonzero matrix entry remains nonzero under every
field embedding, so all its local operator norms are positive and the
height is defined. Averaging proves the last assertion. If
$\mathcal B=0$, the norm bound still holds, but no logarithm of zero is
needed.
\end{proof}

To use these bounds in the ordinary-slope comparison of
Section~\ref{cs-section}, the conventions must be converted together.
The bilinear pairing $\mathcal B$ and the induced linear map
$U\to E^\vee$ have the same operator norm. Since $n_v=2$ at a complex
place and $r_2=d/2$, the height used there satisfies
\[
 h_2(\mathcal B)=2h(\mathcal B),\qquad
 c(R)/D=2\mu(R^\vee).
\]
Accordingly, add the divisor inequality
\eqref{div:final-divisor} to the squared-norm height inequality above,
cancel the two occurrences of $\overline{\log\kappa}$, and divide
the whole inequality by two before applying the ordinary-slope
comparison. Converting only one of the two inequalities would give
an incorrect factor of two.

\subsection{Local divisibility gives an integral map}

\begin{proposition}\label{div:integral-pairing}
Let $N\subset U$ and $R\subset E$ be full lattices, with $\dim U>0$, and let
$\mathcal B:U\times E\to K$ be bilinear. For each prime let
$s_{\mathfrak p}\in\mathbb R$, with
$k_{\mathfrak p}=\lfloor s_{\mathfrak p}\rfloor$ of finite support.
Assume that for all $x\in N$, all $y\in R_{\mathfrak p}$, and every
prime $\mathfrak p$, either $\mathcal B(x,y)=0$ or
\[
 \operatorname{ord}_{\mathfrak p}\mathcal B(x,y)\ge s_{\mathfrak p}.
\]
Put $\mathfrak b=\prod_{\mathfrak p}\mathfrak p^{-k_{\mathfrak p}}$.
Then $\mathcal B$ induces an integral linear map
$\mathfrak bN\to R^\vee$, and
\[
 2\mu(\mathfrak bN)=2\mu(N)
             +2\sum_{\mathfrak p}w_{\mathfrak p}k_{\mathfrak p}.
\]
It is enough to establish the valuation hypothesis on any lattice with
the same localization as $R$ at the prime in question.
\end{proposition}
\begin{proof}
At $\mathfrak p$, denominators used to localize $N$ are units. Thus the
hypothesis extends to $x\in N_{\mathfrak p}$ by linearity and the
valuation inequality for sums. Multiplication by
$\mathfrak b_{\mathfrak p}=\pi^{-k_{\mathfrak p}}\mathcal O_{\mathfrak p}$
lowers a valuation by $k_{\mathfrak p}$, leaving it at least
$\lceil s_{\mathfrak p}\rceil-\lfloor s_{\mathfrak p}\rfloor\ge0$.
Hence $\mathcal B(\mathfrak bN,R)$ lies in every
$\mathcal O_{\mathfrak p}$, and therefore in $\mathcal O$.
This is precisely the asserted integral map.
If $a=\dim U>0$, then
$\mathfrak a(\mathfrak bN)=\mathfrak b^a\mathfrak a(N)$.
Taking logarithms in \eqref{div:covolume} and dividing by $a$ proves
the slope identity. In rank zero the map assertion remains true and
there is no slope assertion. Finally the valuation condition concerns
only the localization, so replacing a lattice by one with the same
localization does not change it.
\end{proof}

\subsection{Tests of the hypotheses and the remaining arithmetic work}

The following examples check three hypotheses which cannot be omitted.
They are counterexamples to weakened statements, not counterexamples
to the propositions proved above.

\begin{example}[One prime does not suffice for a Hermitian basis]
Take $K=\mathbb Q(i)$, $L=\mathbb Z[i]$, $H=(1)$,
$\mathfrak p=(2+i)$ and $\overline{\mathfrak p}=(2-i)$.
The vector $e=2-i$ is a basis of $L_{\mathfrak p}$, because it is a unit
there. Its Gram determinant is $5$, of valuation one at $\mathfrak p$.
But $L^\#=L$, so $j_{\mathfrak p}(H,L)=0$.
Thus \eqref{div:j-gram} would be false with an ordinary local basis at
only $\mathfrak p$. The simultaneous-basis hypothesis excludes this
example exactly because $e$ is not a unit at
$\overline{\mathfrak p}$.
\end{example}

\begin{example}[Bounds on spanning lines do not suffice]
Let $Q=E=\mathbb C^2$ with standard metrics,
$\iota=\operatorname{diag}(1,\eta)$, where $0<\eta<1$, and let
$Q_1=\mathbb Ce_1$, $Q_2=\mathbb C(e_1+\varepsilon e_2)$ with
$\varepsilon>0$. These lines span $Q$, and the squared norm ratios
on them are $1$ and $(1+\eta^2\varepsilon^2)/(1+\varepsilon^2)$.
Both are at least $(1+\varepsilon^2)^{-1}$, whereas the smallest
squared singular value of $\iota$ is $\eta^2$, which can be arbitrarily
small. The all-$y$ condition in \eqref{div:spectral-cross} is therefore
essential; knowing the norm identity only within each spanning
subspace would not prove Lemma~\ref{div:spectral}.
\end{example}

\begin{example}[Indices cannot be bounded above by multiplying minors]
Over a discrete valuation ring with uniformizer $\pi$, take
$A=(1,1)$ and $B=(1,-1+\pi^N)^t$. The first determinantal ideals of
$A$ and $B$ are both the unit ideal, while $AB=(\pi^N)$ has index
$N$. This holds for every $N\ge1$.
Thus Lemma~\ref{div:minors} gives a lower bound on the valuation of
the product ideal, and supplies no upper bound for it. The construction
must control the actual product or the actual local pairings when an
upper bound is needed.
\end{example}

The divisor inequality and the chain estimate have now been proved from
ordinary linear algebra, local algebra, the product formula and one
explicit finite rational calculation. The non-CM quotient issue,
the conjugate primes, all normalizing factors, and rectangular matrix
dimensions are accounted for. The construction required for the GRH
argument still has to supply the following facts for its particular
family, uniformly in the large parameters:
\begin{enumerate}
\item the ambient forms $H_t$, lattices $M_t$, quotient maps $P_t$,
      comparison maps $\iota_t,\theta_{j,t}$ and the simultaneous
      spectral decompositions proving
      \eqref{div:det-base}--\eqref{div:det-local};
\item the underlying unimodular arithmetic lattices that enter the
      coefficient construction, together with the remaining local
      index estimates and their error sums in
      Corollary~\ref{div:scalar-divisor};
\item the exact matrix factorization
      \eqref{div:pairing-factorization}, the isometric realizations of
      its individual steps, and bounds for the actual operator entries
      $A_{ts}$ and the final pairing $\mathcal P$;
\item the valuation inequality of
      Proposition~\ref{div:integral-pairing} for all vectors of the
      actual source and target lattices, together with the lower
      spectral bounds occurring in \eqref{div:pairing-bound}.
\end{enumerate}
The coefficient duality giving the initial containment
\eqref{div:class-lattices}, with $s=1$ for unimodular input lattices,
is proved in Section~\ref{coef:section}. The estimates concerning zeros in the supplied quasi-Riemann papers
do not by themselves imply any of these lattice containments or exact
factorizations. They may improve the auxiliary-prime choices used to
construct the family. They do not remove the need to construct the
operators and prove their uniform bounds. Establishing these remaining
facts is a mathematical reconstruction task; their absence from this
section is not evidence that the formal development assumes them
without proof.


% Concrete arithmetic coefficient maps, their divisibility and their spectra.
% All labels in this fragment have the prefix am-.
\section{Arithmetic coefficient maps and their singular values}
\label{am-section}

The coefficient lattices of Section~\ref{coef:section} give explicit
arithmetic maps whose image ideals and archimedean norms can be computed.
This section proves the estimates needed for these maps, including the
normalizing constant in their metrics. It also gives sharper bounds in
terms of elementary divisors and singular values. The comparison with the
actual class components is stated at the end. These are coefficient maps;
the different valuation inequalities for the iterated period pairing
require additional arguments.\cite{alpo-uniform-source}

\subsection{The actual map and its integral image}

Let $K$ be a number field, $\mathcal O$ its ring of integers, and $n\ge1$.
Fix a finite set $\mathcal T$. For each $\tau\in\mathcal T$ let
$\lambda_\tau=(\lambda_{\tau,1},\ldots,\lambda_{\tau,n})$ be a
nonincreasing tuple of integers. Put
\[
 a_\tau=\lambda_{\tau,n},\qquad
 \nu_{\tau,i}=\lambda_{\tau,i}-a_\tau,\qquad
 d_\tau=\sum_i\nu_{\tau,i}.
\]
Thus $\nu_\tau$ is a partition and $\nu_{\tau,n}=0$.
The arguments below also apply to a partition with a nonzero last part
and an independently specified integer $a_\tau$.

Use a separate $n$ by $n$ variable matrix $X_\tau$ for each $\tau$ and set
\[
 V=\bigotimes_{\tau\in\mathcal T}V_{\nu_\tau,K},\qquad
 \Delta_0=\bigotimes_\tau\Delta_{\nu_\tau}(\mathcal O),\qquad
 \nabla_0=\bigotimes_\tau\nabla_{\nu_\tau}(\mathcal O).
\]
The tensor products are realized as products of polynomials in the
disjoint variable blocks. Define
\[
 \rho_0(g)F((X_\tau)_\tau)=F((g_\tau^tX_\tau)_\tau),\qquad
 \rho_\lambda(g)=\left(\prod_\tau\det(g_\tau)^{a_\tau}\right)\rho_0(g).
 \tag{AM1}\label{am-action}
\]
Both are representations of $\prod_\tau\operatorname{GL}_n(K)$:
substitution reverses multiplication, and the transpose reverses it
again. The determinant powers in the second expression may be negative.

Let $\mathcal A$ be the finite set of arrays
$\alpha=(\alpha_{\tau,i,s})$ of nonnegative integers satisfying
\[
             \sum_{i=1}^n\alpha_{\tau,i,s}=\nu_{\tau,s}
                       \quad\text{for every $\tau,s$}.
\]
Every polynomial in $V$ is supported on $\mathcal A$.
For $g_\tau\in\operatorname{GL}_n(K)$ the coefficient map is
\begin{equation}\label{am-map}
 C_g:V\longrightarrow K^{\mathcal A},\qquad
 C_g(F)_\alpha=[X^\alpha]\rho_0(g)F.
\end{equation}
In particular the determinant character is \emph{not} inserted into
the definition of $C_g$. It will occur in its integral image ideal.
The map is injective: all the coefficients of $\rho_0(g)F$ occur in
\eqref{am-map}, and $\rho_0(g^{-1})$ is its inverse substitution.

The translates of $\prod_\tau H_{\nu_\tau}$ under $\rho_\lambda$ span
$V$. Here is a direct justification of the spanning statement. The
integral duality theorem proves that the divided-row-operation words
span $V$ over $K$. Each such word is a coefficient of a polynomial in
finitely many shear parameters. Polynomial interpolation recovers that
coefficient from finitely many values at nonnegative integers. Every
off-diagonal shear is invertible, and a diagonal substitution
$1+t$ is also invertible at these parameter values. Thus the relevant
values are invertible-matrix translates of the extremal polynomial.
Their determinant characters are nonzero scalars, which do not change
their $K$-span. Tensoring the argument proves the assertion for all blocks.

For full $\mathcal O$-lattices $M_\tau\subset K^n$, write
$\mathfrak d(M_\tau)\subset K$ for their determinant ideals, using
$e_1\wedge\cdots\wedge e_n$ to identify $\bigwedge^nK^n$ with $K$.
At a nonzero prime $\mathfrak p$, choose local bases
\[
       M_{\tau,\mathfrak p}=B_{\tau,\mathfrak p}\mathcal O_{\mathfrak p}^n.
\]
Define the transported coefficient lattice $\Delta_\lambda(M)$ by
\begin{equation}\label{am-local-lattice}
 \Delta_\lambda(M)_{\mathfrak p}
       =\rho_\lambda(B_{\mathfrak p})\Delta_{0,\mathfrak p}.
\end{equation}
The lattice $\nabla_\lambda(M)$ is defined in the same way with
$\nabla_0$. These definitions agree with the intersection over all
local bases used in Section~\ref{coef:section}.

To check independence of the bases, column invariance immediately
shows that $\rho_0(U)$ preserves $\nabla_0$ when
$U_\tau\in\operatorname{GL}_n(\mathcal O_{\mathfrak p})$.
It preserves $\Delta_0$ as well: the integral duality already proved in
Theorem~\ref{coef:integral-duality} and the transpose adjoint identity
identify $\Delta_0$ with the dual of $\nabla_0$ under the normalized
bilinear Fock pairing. Apply the adjoint identity also to $U^{-1}$ to
obtain equality. The determinant characters are units on this group,
even for negative $a_\tau$. Only finitely many local lattices in
\eqref{am-local-lattice} differ from the standard one, so their
intersection is a full lattice, by the elementary local patching result
of Section~\ref{div:section}.

\begin{theorem}[Integral coefficient map]\label{am-integral}
Suppose $g_\tau M_\tau\subset\mathcal O^n$ for every $\tau$, and put
\begin{equation}\label{am-ideal}
                  \mathfrak I=\prod_\tau\mathfrak d(M_\tau)^{a_\tau}.
\end{equation}
Then $\mathfrak I$ is a nonzero fractional ideal and
\[
 C_g(\Delta_\lambda(M))\subset\mathfrak I^{\mathcal A},
 \qquad
 C_g(\nabla_\lambda(M))\subset\mathfrak I^{\mathcal A}.
 \tag{AM2}\label{am-image}
\]
Equivalently, every coordinate of the image has valuation at least
$\sum_\tau a_\tau\operatorname{ord}_{\mathfrak p}\det B_{\tau,\mathfrak p}$
at each $\mathfrak p$. The assertion is valid at primes above $2$ and
does not require any $M_\tau$ to be free over $\mathcal O$.
\end{theorem}
\begin{proof}
Localize at $\mathfrak p$ and write
$F=\rho_\lambda(B)F_0$, where $F_0$ belongs to the appropriate standard
lattice. Its coefficients are integral over $\mathcal O_{\mathfrak p}$.
Since $g_\tau B_\tau$ has integral entries,
\[
 \rho_0(g)F
   =\left(\prod_\tau\det B_\tau^{a_\tau}\right)\rho_0(gB)F_0
\]
has every coefficient in
$\prod_\tau\det B_\tau^{a_\tau}\mathcal O_{\mathfrak p}$.
This is exactly $\mathfrak I_{\mathfrak p}$, because
$\mathfrak d(M_\tau)_{\mathfrak p}=\det B_\tau\mathcal O_{\mathfrak p}$.
Intersect these inclusions over $\mathfrak p$. The same description also
proves that the intersection defined by the conditions
\[
 \left(\prod_\tau\det B_\tau^{-a_\tau}\right)c
                              \in\mathcal O_{\mathfrak p}
\]
for all local bases is precisely the fractional ideal \eqref{am-ideal}.
\end{proof}

\subsection{Row weights and a sharper local divisibility statement}

For a polynomial in one variable block, its row degree is
$\beta=(\beta_1,\ldots,\beta_n)$ if every monomial has total degree
$\beta_i$ in row $i$. Projection to a row degree preserves the defining
column identities of $V_\nu$.

\begin{lemma}[Row weights]\label{am-weights}
If the row-degree-$\beta$ subspace of $V_\nu$ is nonzero, then
$\sum_i\beta_i=\sum_i\nu_i$, and for every subset
$S\subset\{1,\ldots,n\}$ with $|S|=k$,
\begin{equation}\label{am-majorization}
       \sum_{i\in S}\beta_i\le\sum_{i=1}^k\nu_i.
\end{equation}
In particular, for every $k$,
\[
 \sum_{i=n-k+1}^n\nu_i\le\sum_{i=1}^k\beta_i
                              \le\sum_{i=1}^k\nu_i.
\]
The weights $\nu$ and $(\nu_n,\ldots,\nu_1)$ both occur.
\end{lemma}
\begin{proof}
For $S=\{1,\ldots,k\}$, restrict a nonzero polynomial of row degree
$\beta$ to lower triangular matrices. The restriction is nonzero by
Lemma~\ref{coef:restriction}. In a monomial of this restriction, the
first $k$ rows use only the first $k$ columns. Their total degree is
therefore at most $\sum_{i\le k}\nu_i$. Row permutations preserve
column invariance and column degrees, which proves the assertion for
every $S$. The lower bound follows by applying the upper bound to the
complement of the first $k$ rows. The extremal polynomial $H_\nu$ has
row degree $\nu$; reverse its rows for the other endpoint.
\end{proof}

The lemma requires the column-invariance identities. Column homogeneity
alone would be insufficient. For example, in three rows the polynomial
$x_{11}x_{12}$ has column degrees $(1,1,0)$ and row degree $(2,0,0)$;
it violates \eqref{am-majorization} and is not column invariant.

\begin{theorem}[Exact local coefficient ideal]\label{am-smith}
Let $R$ be a discrete valuation ring of characteristic zero with fraction
field $L$, uniformizer $\pi$, and normalized valuation $v(\pi)=1$.
Let $M_\tau=B_\tau R^n$ and $g_\tau\in\operatorname{GL}_n(L)$, without
requiring $g_\tau M_\tau\subset R^n$. Let
$c_{\tau,1}\le\cdots\le c_{\tau,n}$ be the elementary-divisor
exponents of $g_\tau B_\tau$. The fractional ideal generated by all
coordinates $C_g(F)_\alpha$, as $F$ ranges over
$\Delta_\lambda(M)$, is exactly
\begin{equation}\label{am-sharp-local}
 \pi^eR,\qquad
 e=\sum_\tau\left(a_\tau v(\det B_\tau)
                         +\sum_i\nu_{\tau,i}c_{\tau,i}\right).
\end{equation}
The same statement holds with $\nabla_\lambda(M)$.
\end{theorem}
\begin{proof}
Smith reduction over $R$ gives
$g_\tau B_\tau=U_\tau D_\tau V_\tau$ with $U_\tau,V_\tau$ unimodular
and $D_\tau=\operatorname{diag}(\pi^{c_{\tau,1}},\ldots,
\pi^{c_{\tau,n}})$. This also holds for matrices with denominators:
first multiply the matrix by a power of $\pi$ and then subtract that
power from every elementary-divisor exponent.
For completeness, integral Smith reduction follows by moving an entry
of least valuation to the upper left. It divides every entry. Integral
row and column subtractions clear its row and column; induction treats
the remaining block, whose entries are still divisible by the chosen
pivot. Thus the successive pivot valuations are nondecreasing.

Integral unimodular substitutions preserve both standard lattices and
preserve the ideal generated by the coefficients of any polynomial:
one inclusion follows from integrality of the substitution, and the
reverse inclusion follows by applying its integral inverse.
It remains to treat the diagonal matrices $D_\tau$.
On a row-degree-$\beta$ term, $\rho_0(D_\tau)$ multiplies by
$\pi^{\sum_i c_{\tau,i}\beta_i}$. Summation by parts and
Lemma~\ref{am-weights} give
\[
 \sum_i c_{\tau,i}\beta_i
 =c_{\tau,n}d_\tau
   -\sum_{k<n}(c_{\tau,k+1}-c_{\tau,k})\sum_{i\le k}\beta_i
 \ge\sum_i c_{\tau,i}\nu_{\tau,i}.
\]
The determinant character of $B$ supplies the first term of $e$.
For equality take the product of the $H_{\nu_\tau}$ after undoing
the substitutions $V_\tau$. Each $H_{\nu_\tau}$ belongs to the Weyl
lattice, has row degree $\nu_\tau$, and has a coefficient equal to one:
its restriction to lower triangular matrices is the diagonal monomial
with coefficient one. Hence the lower bound is attained as the ideal
of coefficients, even after the substitutions $U_\tau$.
\end{proof}

In Theorem~\ref{am-integral} all $c_{\tau,i}$ are nonnegative, so
\eqref{am-sharp-local} implies the stated determinant-ideal bound.
The sharper formula also identifies exactly how much divisibility is
lost when the chosen linear functionals have denominators.

\subsection{The Fock metric and exact singular-value bounds}

Over $\mathbb C$ put
\[
 C_\tau=\sum_\alpha\alpha!\,|[X^\alpha]H_{\nu_\tau}|^2,
 \qquad C=\prod_\tau C_\tau>0,
 \qquad
 \|F\|_V^2=C^{-1}\sum_\alpha\alpha!\,|F_\alpha|^2.
\]
Give $\mathbb C^{\mathcal A}$ the diagonal metric
\begin{equation}\label{am-coordinate-metric}
 \|z\|_{\mathcal A}^2
       =\sum_{\alpha\in\mathcal A}w_\alpha|z_\alpha|^2,
       \qquad w_\alpha=\alpha!/C.
\end{equation}
Then $C_1$ is an isometric embedding. This explains the factorial weights
in the target metric; replacing them by the ordinary coordinate metric
would change the operator estimates.

\begin{theorem}[Largest and smallest singular values]\label{am-singular}
For each $\tau$ let
$s_{\tau,1}\ge\cdots\ge s_{\tau,n}>0$ be the singular values of
$g_\tau\in\operatorname{GL}_n(\mathbb C)$. Then
\begin{align}
 \|C_g\|_{\mathrm{op}}
       &=\prod_{\tau,i}s_{\tau,i}^{\nu_{\tau,i}},
       \label{am-op-exact}\\
 \inf_{F\ne0}\frac{\|C_gF\|_{\mathcal A}}{\|F\|_V}
       &=\prod_{\tau,i}s_{\tau,i}^{\nu_{\tau,n+1-i}}.
       \label{am-min-exact}
\end{align}
For $\rho_\lambda(g)$ multiply both expressions by
$\prod_\tau|\det g_\tau|^{a_\tau}$; equivalently replace $\nu$ by
$\lambda$ in the corresponding ordered products.
In particular
\begin{equation}\label{am-frobenius}
 \|C_gF\|_{\mathcal A}^2
 \le\left(\prod_\tau\|g_\tau\|_{\mathrm{op}}^{2d_\tau}\right)\|F\|_V^2
 \le\left(\prod_\tau\|g_\tau\|_{\mathrm{HS}}^{2d_\tau}\right)\|F\|_V^2.
\end{equation}
\end{theorem}
\begin{proof}
The Hermitian version of Lemma~\ref{coef:adjoint} says that the adjoint
of $\rho_0(g)$ for the Fock metric is $\rho_0(g^*)$.
Thus unitary substitutions act unitarily on $V$. A singular-value
decomposition reduces the problem to positive diagonal $g_\tau$.
Distinct row weights are orthogonal, and on the weight $\beta$ the
diagonal action has eigenvalue $\prod_i s_{\tau,i}^{\beta_i}$.
Writing $u_i=\log s_{\tau,i}$, summation by parts gives
\[
 \sum_i u_i\beta_i=u_nd_\tau
                   +\sum_{k<n}(u_k-u_{k+1})\sum_{i\le k}\beta_i.
\]
The two partial-sum bounds of Lemma~\ref{am-weights} prove the upper
and lower estimates. Its endpoint weights attain both estimates.
Their tensor products attain the endpoints simultaneously for all
$\tau$. Finally $C_1$ is isometric, so the singular values of $C_g$
are those of $\rho_0(g)$ on $V$. The scalar determinant twist gives
the last assertion. The two inequalities in \eqref{am-frobenius}
follow from $s_{\tau,i}\le s_{\tau,1}\le\|g_\tau\|_{\mathrm{HS}}$.
\end{proof}

No bound involving only $\det g$ could replace
\eqref{am-op-exact}: for $\nu=(1,0)$ and
$g=\operatorname{diag}(T,T^{-1})$, $T>1$, the determinant is one
and the operator norm is $T$. Also, the column-homogeneous polynomial
$x_{11}x_{12}$ above is multiplied by $T^2$ under
$\operatorname{diag}(T,1,1)$, whereas the claimed bound for partition
$(1,1,0)$ would be $T$. This confirms the necessity of column
invariance in the exact formula.

\begin{proposition}[Determinant of the substitution]\label{am-determinant}
Let $r=\dim_KV$. The numbers $\kappa_\tau=rd_\tau/n$ are integers, and
\[
 \det(\rho_0(g)|V)=\prod_\tau\det(g_\tau)^{\kappa_\tau},\qquad
 \det(\rho_\lambda(g)|V)
                 =\prod_\tau\det(g_\tau)^{\kappa_\tau+ra_\tau}.
 \tag{AM3}\label{am-det-formula}
\]
Over $\mathbb C$, in orthonormal source coordinates,
$\det(C_g^*C_g)=\prod_\tau|\det g_\tau|^{2\kappa_\tau}$.
\end{proposition}
\begin{proof}
Decompose $V$ into row-weight spaces in block $\tau$.
Row permutations identify the relevant weight spaces, so the sum of
the $i$th row degrees, counted with their dimensions, is independent
of $i$. Its sum over $i$ is $rd_\tau$; hence each sum is the integer
$\kappa_\tau$. This proves \eqref{am-det-formula} on diagonal matrices.
An elementary shear has unipotent action on the polynomial space of
degree $d_\tau$, since it changes one row degree in only one direction.
Its restriction to the invariant space $V$ has determinant one.
Gaussian elimination generates every invertible matrix by diagonal
matrices and elementary shears (a row interchange is a product of
three shears and a diagonal matrix). Thus the formula holds for every
$g$. The determinant twist acts by a scalar on an $r$-dimensional
space. The last formula follows because $C_1$ is isometric.
\end{proof}

\subsection{An elementary bound for the normalizing constant}

The preceding formulas retain the exact constant $C_\nu$.
Its required upper bound has a short proof by Gaussian integration;
no dimension formula for polynomial representations is needed.

\begin{theorem}[Norm of the extremal polynomial]\label{am-gaussian}
For every partition $\nu$ of length $n\ge1$,
\begin{equation}\label{am-normalization-bound}
 \prod_i\nu_i!\ \le\ C_\nu\ \le\
 \prod_{i=1}^n\frac{(\nu_i+n-i)!}{(n-i)!}
 \ \le\left(\prod_i\nu_i!\right)(\nu_1+n)^{n(n-1)/2}.
\end{equation}
\end{theorem}
\begin{proof}
The diagonal monomial in $H_\nu$ has coefficient one, so its contribution
to the Fock norm is $\prod_i\nu_i!$, proving the lower bound.

Let $Z$ be an $n$ by $n$ matrix of independent complex Gaussian entries,
each with density $\pi^{-1}e^{-|z|^2}$ on $\mathbb C$.
Polar integration gives
\[
 \mathbb E(z^a\overline{z}^{b})=
 \begin{cases}a!&a=b,\\0&a\ne b.\end{cases}
\]
Indeed the angular integral vanishes unless $a=b$; the remaining radial
integral is $\int_0^\infty t^ae^{-t}\,dt=a!$, by integration by parts.
Independence consequently proves
$C_\nu=\mathbb E|H_\nu(Z)|^2$.

Let $Z_j$ consist of the first $j$ columns of $Z$.
Cauchy--Binet gives
\[
           |\Delta_j(Z)|^2\le\det(Z_j^*Z_j).
\]
Apply Gram--Schmidt to these columns, writing $r_i>0$ for the length
of the component of column $i$ orthogonal to the preceding columns.
The columns are independent with probability one: the zero set of a
nonzero polynomial has Lebesgue measure zero, by induction on the
number of variables and Fubini's theorem. Thus
\[
 |H_\nu(Z)|^2
 \le\prod_j\det(Z_j^*Z_j)^{\nu_j-\nu_{j+1}}
 =\prod_i r_i^{2\nu_i},\qquad \nu_{n+1}=0.
\]
Conditionally on the preceding columns, the orthogonal component of
column $i$ is a standard complex Gaussian vector in dimension $n-i+1$.
This follows directly from unitary invariance of its density.
For such a vector in dimension $q$, the multinomial theorem and the
one-variable moments give
\[
 \mathbb E\left(\sum_{j=1}^q|z_j|^2\right)^m
 =\sum_{b_1+\cdots+b_q=m}\frac{m!}{b_1!\cdots b_q!}
                                      \prod_j b_j!
 =\frac{(m+q-1)!}{(q-1)!}.
\]
There are $\binom{m+q-1}{q-1}$ nonnegative $q$-tuples with sum $m$,
as follows by inserting $q-1$ dividers into a string of $m$ symbols.
Successive conditional integrations, starting with the last column,
therefore prove the middle inequality of \eqref{am-normalization-bound}.
Finally
\[
 \frac{(\nu_i+n-i)!}{(n-i)!\nu_i!}
          =\prod_{j=1}^{n-i}\frac{\nu_i+j}{j}
          \le(\nu_1+n)^{n-i}.
\]
Multiply over $i$.
\end{proof}

For $m\ge0$ write $m=nq+r$ with $0\le r<n$, and put
\[
                    b_n(m)=(q!)^{n-r}((q+1)!)^r.
\]
The least factorial weight in \eqref{am-coordinate-metric} is exactly
\begin{equation}\label{am-min-weight}
 w_{\min}=\min_{\alpha\in\mathcal A}w_\alpha
                    =\frac{\prod_{\tau,s}b_n(\nu_{\tau,s})}{C}.
\end{equation}
To prove this, minimize the product of the factorials separately in
each column. If two entries are $u\ge v+2$, replacing them by $u-1,v+1$
decreases the product by the factor $(v+1)/u<1$. Repetition leaves
entries differing by at most one, which are precisely the balanced
entries in the definition of $b_n$.

A less precise bound, but one expressed without factorial sums, is
\begin{equation}\label{am-weight-lower}
 w_{\min}\ge
 \left(n^{\sum_\tau d_\tau}
       \prod_\tau\prod_{i=1}^n
        \frac{(\nu_{\tau,i}+n-i)!}{(n-i)!\nu_{\tau,i}!}\right)^{-1}
 \ge
 \left(n^{\sum_\tau d_\tau}
       \prod_\tau(\nu_{\tau,1}+n)^{n(n-1)/2}\right)^{-1}.
\end{equation}
Indeed the multinomial theorem at $(1,\ldots,1)$ implies
$m!\le n^m\prod_i b_i!$ whenever $\sum_i b_i=m$.
Multiply this over the columns and apply Theorem~\ref{am-gaussian}.

\subsection{The resulting slope bound, with ordinary degree conventions}

Write $d_K=[K:\mathbb Q]$, and let $n_v=1$ or $2$ be the local degree
at a real or complex archimedean place. All metrics below are compatible
with complex conjugation. For a Hermitian lattice $L$ of rank $r>0$,
choose $0\ne\omega\in\bigwedge^r_K(KL)$ and write
$\bigwedge^r_{\mathcal O}L=\mathfrak a\omega$.
Our normalized degree and slope are
\begin{equation}\label{am-degree}
 \widehat{\deg}L=-\frac{\log N\mathfrak a}{d_K}
      -\frac1{d_K}\sum_{v\mid\infty}n_v\log\|\omega\|_v,
 \qquad \mu(L)=\widehat{\deg}L/r.
\end{equation}
Here $N\mathfrak a$ is the absolute norm of a fractional ideal, a
positive rational number. The product formula makes the definition
independent of $\omega$. Let $\mu_{\max}(L)$ be the supremum of the
slopes of its nonzero sublattices, with their induced metrics.

\begin{lemma}[A diagonal target]\label{am-diagonal-slope}
Suppose $\mathfrak I\subset K$ is a nonzero fractional ideal and
$K^A$ has diagonal metrics with constant positive diagonal weights
$w_a\ge w_0>0$ at every infinite place. Then every nonzero sublattice
$L\subset\mathfrak I^A$ satisfies
\[
                 \mu(L)\le-\frac{\log N\mathfrak I}{d_K}
                                      -\frac12\log w_0.
\]
\end{lemma}
\begin{proof}
Let $r$ be its rank, write $\bigwedge^rL=\mathfrak a\omega$, and
choose a nonzero coordinate $b$ of the exterior vector $\omega$.
Since all coordinates of $\bigwedge^rL$ lie in $\mathfrak I^r$,
$b\mathfrak a\subset\mathfrak I^r$ and hence
$\log N\mathfrak a+\log|N_{K/\mathbb Q}b|
\ge r\log N\mathfrak I$. In the exterior metric,
$\|\omega\|_v\ge w_0^{r/2}|v(b)|$, because its squared norm is the
sum of its squared coordinates times products of $r$ diagonal weights.
Insert these two inequalities into \eqref{am-degree} and use
$\sum_vn_v\log|v(b)|=\log|N_{K/\mathbb Q}b|$.
\end{proof}

For an injective $K$-linear map $f$, put
$h(f)=d_K^{-1}\sum_v n_v\log\|f\|_{v,\mathrm{op}}$.
For every sublattice $L$, the exterior operator estimate gives
$\widehat{\deg}L\le\widehat{\deg}f(L)+r h(f)$: use the same
fractional ideal in the determinant generators $\omega$ and $f\omega$,
and $\|f\omega\|_v\le\|f\|_{v,\mathrm{op}}^r\|\omega\|_v$.
This proves the following bound directly for the map constructed here.

\begin{corollary}[Slope of the coefficient lattice]\label{am-slope}
Give $\Delta_\lambda(M)$ the normalized Fock metric at every infinite
place and suppose $g_\tau M_\tau\subset\mathcal O^n$.
With $\mathfrak I$ as in \eqref{am-ideal} and the singular values of
$v(g_\tau)$ ordered decreasingly,
\begin{equation}\label{am-slope-bound}
 \mu_{\max}(\Delta_\lambda(M))
 \le-\frac{\log N\mathfrak I}{d_K}-\frac12\log w_{\min}
       +\frac1{d_K}\sum_v n_v\sum_{\tau,i}
                                  \nu_{\tau,i}\log s_{\tau,i}(v).
\end{equation}
The same bound holds for $\nabla_\lambda(M)$.
One may replace the last summand by
\[
 \frac1{2d_K}\sum_vn_v\sum_\tau d_\tau
                          \log\|v(g_\tau)\|_{\mathrm{HS}}^2
\]
and bound $w_{\min}$ by \eqref{am-weight-lower}.
\end{corollary}
\begin{proof}
Apply Theorem~\ref{am-integral}, Lemma~\ref{am-diagonal-slope}, and the
preceding exterior operator inequality to $C_g$. Its injectivity and
its exact operator norm were proved above. The weaker version is
\eqref{am-frobenius}.
\end{proof}

The use of ordinary slopes in \eqref{am-degree} matters. A convention
using logarithms of squared norms doubles both the slope and the height.
Under that convention the right side of \eqref{am-slope-bound} must
also be doubled; in particular the ideal term becomes
$-2\log N\mathfrak I/d_K$.

\subsection{The maps on actual class components}

We describe precisely the additional data needed to apply the result
to class components, and prove the resulting comparison.
In the application, class values are obtained by evaluating equivariant
functions on representatives of decorated-lattice isometry classes.
At the level in question the stabilizers are trivial. Evaluation then
identifies the space of such functions with the full product of their
coefficient spaces: assign arbitrary values at the representatives and
extend by the coefficient action. The extension is well defined because
two isometries with the same image differ by a stabilizer. Conversely,
equivariance determines the function from these values. The class-sum
metric has reciprocal stabilizer weights, which are one at this level.
For nontrivial stabilizers one must instead use their invariant
coefficient spaces; the full-product assertion would be false in general.

Let $\mathcal C$ be finite. For every $c\in\mathcal C$ and
$\tau\in\mathcal T$, let $M_{c,\tau}\subset K^n$ be a full lattice.
Consider the lattice in $\bigoplus_cV$ whose localization at
$\mathfrak p$ consists of the tuples $(F_c)_c$ satisfying
\begin{equation}\label{am-class-local}
 \rho_\lambda(B_{c,\mathfrak p}^{-1})F_c
                               \in\Delta_{0,\mathfrak p}\quad(c\in\mathcal C),
\end{equation}
where $B_{c,\tau,\mathfrak p}$ are local bases of $M_{c,\tau}$.
This lattice is exactly
$\mathcal L=\bigoplus_c\Delta_\lambda(M_c)$.
Consequently the projection onto any class is integral, and the
projection is injective on its class summand. The same assertions
remain true if the right side of \eqref{am-class-local} is replaced by
the integral span of the same divided-row words: this is the defining
Weyl lattice, as proved in Section~\ref{coef:section}.

At every infinite place choose invertible matrices $X_{c,\tau,v}$ and
give $K^n$ the metric
\[
       h_{c,\tau,v}(z,z')=(X_{c,\tau,v}^{-1}z)^*
                                      (X_{c,\tau,v}^{-1}z').
\]
Give $\mathcal L$ the direct-sum metric
\begin{equation}\label{am-class-metric}
       \|F\|_v^2=\sum_c
             \|\rho_\lambda(X_{c,v}^{-1})F_c\|_V^2.
\end{equation}
Fix a class $c$. Choose, for every $\tau$, a basis
$\ell_{\tau,1},\ldots,\ell_{\tau,n}$ of $(K^n)^*$ which is integral
on $M_{c,\tau}$, and let $g_\tau$ be its row matrix.
Define the actual map
\[
                     \Phi_c(F)=C_g(F_c).
\]
It is injective on the $c$th summand and takes that summand into
$\mathfrak I_c^{\mathcal A}$, where
$\mathfrak I_c=\prod_\tau\mathfrak d(M_{c,\tau})^{a_\tau}$.

At $v$, write $y_{\tau,j,v}$ for the Riesz vector of
$v(\ell_{\tau,j})$ for $h_{c,\tau,v}$, using the convention
$\ell(z)=h(y,z)$, conjugate-linear in the first argument.
Let $R_{\tau,v}$ be the matrix whose $j$th row is
$(X_{c,\tau,v}^{-1}y_{\tau,j,v})^*$. Then
\begin{equation}\label{am-riesz}
       v(g_\tau)X_{c,\tau,v}=R_{\tau,v},\qquad
       \|R_{\tau,v}\|_{\mathrm{HS}}^2
                         =\sum_j\|y_{\tau,j,v}\|_{h_{c,\tau,v}}^2.
\end{equation}
Both identities follow by writing
$h(y,z)=y^*(X^{-1})^*X^{-1}z$.
They retain the conjugates; replacing $y^*$ by $y^t$ would generally
give a different functional.

\begin{proposition}[Class coefficient bound]\label{am-class-bound}
Let $s_{\tau,1}(R_{\tau,v})\ge\cdots\ge s_{\tau,n}(R_{\tau,v})>0$
be the singular values in \eqref{am-riesz}, and put
$D_v=\prod_\tau|\det X_{c,\tau,v}|^{2a_\tau}$.
Then for all $F\in\bigoplus_cV$,
\begin{align}
 \|\Phi_c(F)\|_{\mathcal A}^2
 &\le D_v\left(\prod_{\tau,i}
                   s_{\tau,i}(R_{\tau,v})^{2\nu_{\tau,i}}\right)\|F\|_v^2
                          \label{am-class-sharp}\\
 &\le D_v\left(\prod_\tau
       \left(\sum_j\|y_{\tau,j,v}\|_{h_{c,\tau,v}}^2\right)^{d_\tau}\right)
                                                        \|F\|_v^2.
                          \label{am-class-hs}
\end{align}
The determinant contribution satisfies the exact identity
\begin{equation}\label{am-det-cancel}
 -\frac{\log N\mathfrak I_c}{d_K}
       +\frac1{2d_K}\sum_v n_v\log D_v
             =\sum_\tau a_\tau\widehat{\deg}(M_{c,\tau},h_{c,\tau}).
\end{equation}
\end{proposition}
\begin{proof}
Put $P_c=\rho_\lambda(X_{c,v}^{-1})F_c$ and
$d_v=\prod_\tau\det X_{c,\tau,v}^{a_\tau}$. The representation law gives
\[
             \rho_0(v(g))F_c=d_v\rho_0(R_v)P_c.
\]
Apply Theorem~\ref{am-singular} and
$\|P_c\|_V^2\le\sum_{c'}\|P_{c'}\|_V^2=\|F\|_v^2$.
Equation \eqref{am-riesz} proves the weaker bound.
Finally the exterior metric of the standard basis of $K^n$ is
$|\det X_{c,\tau,v}|^{-1}$. Definition \eqref{am-degree} therefore gives
\[
 \widehat{\deg}(M_{c,\tau},h_{c,\tau})
 =-\frac{\log N\mathfrak d(M_{c,\tau})}{d_K}
        +\frac1{d_K}\sum_v n_v\log|\det X_{c,\tau,v}|.
\]
Multiply by $a_\tau$ and sum. This also proves that the expression is
independent of all choices of coordinate bases and archimedean frames.
\end{proof}

\begin{corollary}[A bound from short integral functionals]\label{am-short}
Suppose
\[
 \widehat{\deg}(M_{c,\tau},h_{c,\tau})=-\tfrac12\log T
 \quad\text{for every $\tau$},\qquad T>0,
\]
and, for some $Q>0$,
\begin{equation}\label{am-short-assumption}
        \frac1{d_K}\sum_v n_v\|y_{\tau,j,v}\|_{h_{c,\tau,v}}^2
                                      \le Q\quad(\tau,j).
\end{equation}
Then the $c$th summand, with the metric \eqref{am-class-metric}, satisfies
\begin{equation}\label{am-final-slope}
 2\mu_{\max}(\Delta_\lambda(M_c))
 \le-\log w_{\min}+\sum_\tau d_\tau\log(nQ)
                                      -\sum_\tau a_\tau\log T.
\end{equation}
In particular, if $Q=n^2(2v_0)^2$ and
$c_{\nu}=C_\nu/\prod_i\nu_i!$, then
\begin{equation}\label{am-source-slope}
 2\mu_{\max}(\Delta_\lambda(M_c))
 \le\sum_\tau\left(
       d_\tau\log(4n^4v_0^2)+\log c_{\nu_\tau}-a_\tau\log T\right).
\end{equation}
One may further use
$c_\nu\le\prod_i\binom{\nu_i+n-i}{n-i}
\le(\nu_1+n)^{n(n-1)/2}$.
\end{corollary}
\begin{proof}
All $\|y_{\tau,j,v}\|^2$ are positive because the functionals form a
basis. The concavity bound $\log x\le x-1$, applied after division by
the mean, gives
\[
 \frac1{d_K}\sum_v n_v\log\left(\sum_j\|y_{\tau,j,v}\|^2\right)
 \le\log\left(\frac1{d_K}\sum_v n_v\sum_j\|y_{\tau,j,v}\|^2\right)
 \le\log(nQ).
\]
Use the integral image inclusion, Lemma~\ref{am-diagonal-slope},
\eqref{am-class-hs}, and the exact cancellation \eqref{am-det-cancel}.
This proves \eqref{am-final-slope}.
The multinomial bound preceding \eqref{am-weight-lower} gives
$w_{\min}^{-1}\le n^{\sum_\tau d_\tau}\prod_\tau c_{\nu_\tau}$.
For $Q=n^2(2v_0)^2$, combine the two logarithms of $n$ to obtain
\eqref{am-source-slope}. The bound on $c_\nu$ is
Theorem~\ref{am-gaussian}.
\end{proof}

\subsection{Base change, counterexamples, and the precise remaining input}

The data in the preceding corollary have their usual arithmetic meaning.
For a positive definite Hermitian space over a CM field $E$, integral
vectors in the Hermitian dual of a lattice give integral linear
functionals $z\mapsto h(y,z)$. A linearly independent $n$-tuple gives
an invertible matrix $g$, since nondegeneracy identifies a basis with
a basis of the dual. Upon an embedding $E\hookrightarrow K$, these
functionals extend to the $\mathcal O_K$-span of the lattice and remain
integral. Positivity and the Riesz identities persist at all embeddings.
No unimodularity of the selected $n$-tuple is required.

Normalized degree is unchanged by this base change. Indeed the
determinant ideal becomes $\mathfrak d(M)\mathcal O_K$, whose norm is
$N_E\mathfrak d(M)^{[K:E]}$; this follows prime by prime from
$\sum_{\mathfrak P\mid\mathfrak p}e_{\mathfrak P}f_{\mathfrak P}=[K:E]$.
Each archimedean term is repeated with total multiplicity $[K:E]$.
Both repetitions cancel the change from $[E:\mathbb Q]$ to
$[K:\mathbb Q]$ in \eqref{am-degree}. The same argument shows that
the mean of the squared norms in \eqref{am-short-assumption} is unchanged.
Thus a trace bound over the original CM field gives precisely the mean
bound used above, with no extra factor from the coefficient field.

Several possible simplifications are false, although none is used in
the actual coefficient-map statement.

\begin{itemize}
\item Negative determinant weights cannot be discarded. For $n=1$,
 $M=p\mathbb Z$, $a=-1$, and $\nu=0$, the transported lattice is
 $p^{-1}\mathbb Z$ and the coefficient map is the identity on constants.
 Its image is $p^{-1}\mathbb Z$, not $\mathbb Z$.
\item Integral functionals are essential for the determinant-ideal
 bound. For $M=\mathbb Z_p^2$, $a=0$, $\nu=(1,0)$, and
 $g=\operatorname{diag}(p^{-1},1)$, the polynomial $x_{11}$ acquires
 a coefficient $p^{-1}$. Formula \eqref{am-sharp-local}, which allows
 negative elementary divisors, still gives the correct answer.
\item The weighted coordinate metric is essential. For $\nu=(2,0)$
 in two rows, $C_\nu=2$ and the monomial $x_{11}x_{21}$ has normalized
 Fock norm $1/\sqrt2$. Its unweighted coefficient vector has norm one.
 Thus even $g=1$ would not give a contraction for the unweighted metric.
\item An integral independent $n$-tuple in a dual lattice need not be
 an integral basis. For $M=\mathbb Z_p^2$, the rows $(p,0),(0,1)$
 are independent and integral, but have index $p$ in the dual.
 The proof above uses only the inclusion $gM\subset\mathcal O^n$.
\item Stabilizers cannot be ignored when passing to class values. If
 a group of order two fixes the sole class representative and acts as
 $-1$ on a one-dimensional coefficient space in characteristic zero,
 every equivariant value is zero. It is not a full copy of the
 coefficient space. The actual full-product construction retains the
 trivial-stabilizer hypothesis.
\end{itemize}

The actual application supplies a decomposition into polynomial class
components, the local description \eqref{am-class-local}, and the
metric identity \eqref{am-class-metric}. Given these identifications,
this section proves the coefficient maps, their integral image,
their norm bounds, the determinant-frame cancellation, and the complete
deduction of \eqref{am-source-slope}. The existence of the decorated
base lattices, their specified determinant degree, and short independent
vectors in their Hermitian duals remain separate arithmetic inputs.
The estimates for the iterated period maps and their local vanishing
orders are also separate. In particular, the elementary-divisor formula
proved here must not be substituted for a valuation inequality for a
different map merely because both maps arise in the same construction.


\section{Decorated integral lattices and short vectors in their duals}
\label{al-section}

We prove three arithmetic inputs used by the coefficient maps of
Section~\ref{am-section}: the existence of the initial decorated
unimodular lattice, its kernel lattice's exact degree, and the existence
of sufficiently short independent vectors in the Hermitian dual.
We also prove the triviality of the actual stabilizers by their
unipotent reductions at two different rational primes.
These arguments use no zero-free region for an $L$-function.
The correspondence with the original construction is recorded
separately in the source audit~\cite{alpo-uniform-source}.

Throughout, $E$ is a CM field, $d=[E:\mathbb Q]=2f$, and
$\mathcal O=\mathcal O_E$ is its maximal order. Write
$\overline{\phantom{x}}$ for complex conjugation,
$\mathfrak D$ for the absolute different,
$\Delta=|\operatorname{disc}E|$, and $\rho=\Delta^{1/d}$.
Let $V$ have dimension $n\geq1$ over $E$ and let $h$ be a
Hermitian form, conjugate-linear in its first argument, positive
definite at every complex embedding. A full lattice means a
finitely generated $\mathcal O$-submodule spanning $V$ over $E$.
For a full lattice $L$, put
\[
 L^\#=\{y\in V:h(x,y)\in\mathcal O\text{ for every }x\in L\}.
\]
It is a full lattice. A lattice is integral if $L\subseteq L^\#$,
and unimodular if equality holds. In particular, integrality concerns
the $E$-valued form, not just the integrality of its rational trace.

\subsection{Residue decorations, indices, and determinant degree}

Let $S$ and $G$ be finite sets of nonzero prime ideals of $\mathcal O$.
The decorations used here consist of a unimodular lattice $N$, a
nonzero $\mathcal O$-linear map
\[
       \ell_P:N\longrightarrow\mathcal O/P\qquad(P\in S),
\]
and an isomorphism
\[
       b_Q:N/QN\longrightarrow(\mathcal O/Q)^n\qquad(Q\in G).
\]
Each map has a unique extension to the localization $N_P$ or $N_Q$.
For example, if $sx\in N$ and $s\notin P$, the extension of
$\ell_P$ is
\begin{equation}\label{al-extension}
                 \ell_P(x)=\bar s^{-1}\ell_P(sx).
\end{equation}
If also $tx\in N$ and $t\notin P$, linearity applied to
$tsx$ proves that this value is independent of $s$.
It proves additivity and linearity on the localization as well.
If a function on all of $V$ is wanted, extend it by zero outside
that localization. No assertion of linearity outside the localization
is made or needed.

\begin{proposition}[An explicit initial decoration]\label{al-base}
For the standard form $h(x,y)=\sum_{j=1}^n\bar x_jy_j$ on $E^n$,
take $N=\mathcal O^n$, let $\ell_P$ be the last coordinate modulo
$P$, and let $b_Q$ reduce all coordinates modulo $Q$.
These choices give decorations for every pair of finite sets $S,G$.
\end{proposition}
\begin{proof}
Testing $h(e_j,y)$ shows that $N^\#\subseteq\mathcal O^n$;
the reverse inclusion follows because conjugation preserves
$\mathcal O$. Thus $N$ is unimodular. The vector $e_n$ proves
that every $\ell_P$ is nonzero. Coordinate reduction is onto and
has kernel $QN$, which proves the assertion for $b_Q$.
Formula~\eqref{al-extension} gives the localized maps, including
their prescribed zero extensions. The construction works at
primes above two and does not require principal prime ideals.
\end{proof}

\begin{lemma}[Duals reverse finite indices]\label{al-dual-index}
If $L_1\subseteq L_2$ are full lattices, then
\[
 (L_1^\#)^\#=L_1,\qquad
 [L_1^\#:L_2^\#]=[L_2:L_1].
\]
Here brackets denote the cardinality of the indicated quotient.
\end{lemma}
\begin{proof}
A full lattice over the Dedekind domain $\mathcal O$ is projective.
The nondegenerate form identifies its Hermitian dual with its
conjugate module dual, so projective biduality proves the first
identity. Explicitly, after localization at a prime, choose a basis
of the free module. The dual basis for the nondegenerate form gives
the original module on taking the dual a second time; equality at
all primes proves the global equality.

For the index identity, localize $L_1\subseteq L_2$ at $P$ and
write it in elementary-divisor form
\[
 (L_2)_P=\bigoplus_j\mathcal O_Pe_j,\qquad
 (L_1)_P=\bigoplus_j\varpi_P^{a_j}\mathcal O_Pe_j,
                 \qquad a_j\geq0.
\]
The corresponding quotient of Hermitian duals at $\bar P$ has
the same length $\sum_ja_j$. This follows by taking the dual
basis: the exponents change from $a_j$ to $-a_j$, and scalar
conjugation changes $P$ to $\bar P$. Since
$N\bar P=NP$, multiplication of the local cardinalities proves
the second identity. The same argument also proves that every
Hermitian dual appearing here is a full lattice.
\end{proof}

\begin{proposition}[The common kernel]\label{al-kernel}
Let $N$ be unimodular and let $\ell_P:N\to\mathcal O/P$
be nonzero for each $P$ in a finite set $S$. Put
\[
 L=\{x\in N:\ell_P(x)=0\text{ for all }P\in S\},
       \qquad t=\prod_{P\in S}NP.
\]
Then $L$ is an integral full lattice and
\begin{equation}\label{al-indices}
 [N:L]=[L^\#:N]=t,\qquad [L^\#:L]=t^2.
\end{equation}
For the archimedean metrics induced by $h$, its normalized ordinary
determinant degree is
\begin{equation}\label{al-degree}
                         \widehat{\deg}L=-\frac{\log t}{d}.
\end{equation}
These conclusions do not require $S$ to be disjoint from $\bar S$.
\end{proposition}
\begin{proof}
A nonzero map to the one-dimensional residue field is onto.
Choose $u_P\in N$ with $\ell_P(u_P)=1$.
The Chinese remainder theorem gives $a_P\in\mathcal O$ such that
$a_P\equiv1\pmod P$ and $a_P\equiv0\pmod Q$ for $Q\ne P$ in $S$.
Thus $\ell_Q(a_Pu_P)$ is $1$ for $Q=P$ and $0$ otherwise.
The map $N\to\prod_{P\in S}\mathcal O/P$ is onto, with kernel
$L$, and hence $[N:L]=t$. Also
$(\prod_{P\in S}P)N\subseteq L\subseteq N$, so $L$ is full.
Since $N$ is integral, so is $L$. Duality gives
$L\subseteq N=N^\#\subseteq L^\#$, and
Lemma~\ref{al-dual-index} gives the remaining indices.

Use ordinary degree as defined in \eqref{am-degree}; the metrics
here are the metrics furnished by $h$ at each embedding.
Hermitian duality negates degree. Indeed it inverts the determinant
ideal up to conjugation and takes the reciprocal determinant norm
at every archimedean place, while conjugation does not change ideal
norm. Since $N=N^\#$, its degree is zero.
Passing to a sublattice of index $t$ changes the determinant ideal
norm by a factor $t$ and leaves the ambient metric unchanged.
This proves \eqref{al-degree}.
\end{proof}

In the actual application $|S|=f$ and $NP=T$ for every $P\in S$.
Consequently
\begin{equation}\label{al-actual-degree}
 [L^\#:L]=T^{2f},\qquad
 \widehat{\deg}L=-\tfrac12\log T,\qquad
 \mu(L)=-\frac{\log T}{2n}.
\end{equation}
The last expression becomes $-\log T/n$ in the doubled slope
convention of the original source. In particular,
\eqref{al-actual-degree} supplies the determinant-degree hypothesis
in Corollary~\ref{am-short} for every decorated representative.
It is not necessary to compute the full structure of $L^\#/L$
to obtain this degree.

\subsection{Unipotent reduction and trivial stabilizers}

\begin{lemma}\label{al-local-torsion}
Let $R$ be a discrete valuation ring of characteristic zero whose
residue field has characteristic $p>0$. If $A\in\mathrm{GL}_n(R)$
has finite order and its reduction is unipotent, then the order
of $A$ is a power of $p$. No condition on the ramification of
$R$ or on whether $p=2$ is needed.
\end{lemma}
\begin{proof}
First suppose $B\equiv I\pmod{\mathfrak m}$ and $B^r=I$ with
$p\nmid r$. Write $X=B-I$. If the entries of $X$ belong to
$\mathfrak m^k$ for $k\geq1$, the binomial identity gives
\[
  0=rX+\sum_{j=2}^r\binom rjX^j.
\]
The sum belongs entrywise to $\mathfrak m^{2k}$ and $r$ is a unit
of $R$, so $X$ belongs entrywise to $\mathfrak m^{2k}$.
Induction and $\bigcap_{k\geq1}\mathfrak m^k=0$ give $B=I$.

Now let the order of $A$ be $p^ar$ with $p\nmid r$ and put
$C=A^{p^a}$. Its order is $r$ and its reduction remains unipotent.
For any unipotent matrix $I+Z$ over a field of characteristic $p$,
$(I+Z)^{p^b}=I+Z^{p^b}=I$ when $p^b\geq n$.
Thus the order of $\bar C$ divides both $r$ and a power of $p$,
so $\bar C=I$. The first paragraph yields $C=I$ and hence $r=1$.
\end{proof}

\begin{lemma}\label{al-finite-isometries}
The isometries of a positive definite Hermitian space over $E$
which preserve a fixed full lattice form a finite group.
The same conclusion holds for the rational stabilizer of any
finite-adelic translate of the completed lattice.
\end{lemma}
\begin{proof}
The real form $B(x,y)=\operatorname{Tr}_{E/\mathbb Q}h(x,y)$
is positive definite on $V\otimes_{\mathbb Q}\mathbb R$.
Choose a $\mathbb Z$-basis of the lattice. Each basis vector can
be sent only to lattice points of its fixed $B$-norm. A discrete
lattice meets a bounded set in finitely many points. The images
of the basis determine the map, proving finiteness.

For the second assertion, let the adelic matrix be $x$ and put
$N_x=V\cap x\widehat N$. Choose a nonzero integer $a$ clearing
the finitely many denominators of both $x$ and $x^{-1}$.
Then $aN\subseteq N_x\subseteq a^{-1}N$, so $N_x$ is a full
lattice. A rational element stabilizing $x\widehat N$ stabilizes
$N_x$, and the first assertion applies. Only the containment
between these two global lattices is needed for finiteness;
no strong approximation statement is being used.
\end{proof}

\begin{proposition}[Two different residue characteristics]
\label{al-two-primes}
Let $\gamma$ be an isometry preserving a full lattice in $V$,
or a rational isometry stabilizing a finite-adelic translate.
Suppose that at places above two different rational primes $p,q$,
the matrices of $\gamma$ are conjugate to integral invertible
matrices with unipotent reductions. Then $\gamma=1$.
\end{proposition}
\begin{proof}
Lemma~\ref{al-finite-isometries} gives finite order.
The embeddings of $E$ in its completions are injective; thus
the order of the matrix is unchanged after these embeddings
and after conjugation. Lemma~\ref{al-local-torsion} makes this
order both a power of $p$ and a power of $q$. It is therefore one.
\end{proof}

This is the actual two-prime argument needed when the level
reduces to an upper triangular unipotent subgroup at the two
selected primes. Merely knowing that a matrix is unipotent after
reduction at one arbitrary prime would not prove the assertion.

There is also a useful single-prime alternative. If $p>dn+1$,
the hypotheses at $p$ alone imply $\gamma=1$.
Indeed a nontrivial finite element of $p$-power order has a power
$\delta$ of order $p$. On the rational vector space of dimension
$dn$, its minimal polynomial divides $X^p-1$, is squarefree, and
is not $X-1$. Thus it is divisible by
$\Phi_p(X)=1+X+\cdots+X^{p-1}$, whose irreducibility follows
from Eisenstein applied to $\Phi_p(X+1)$.
It follows that $p-1\leq dn$, a contradiction.
This recovers the separate large-prime decoration criterion.

If a decorated stabilizer consists of a rational isometry
$\gamma$ together with residue-frame matrices $a_Q$ and
$b_Q(\gamma y)=a_Qb_Q(y)$, then $\gamma=1$ also forces
$a_Q=1$ by surjectivity of $b_Q$. Thus these propositions prove
triviality of the entire decorated stabilizer, not just of its
rational projection.

\subsection{A Gaussian proof of the needed Euclidean transference}

We give the analytic step rather than importing a transference
constant. For a full lattice $\Lambda$ in Euclidean
$\mathbb R^D$, write
\[
 \Lambda^*=\{y:\langle y,x\rangle\in\mathbb Z\ (x\in\Lambda)\},
 \qquad
 \lambda_1(\Lambda^*)=\min_{0\ne y\in\Lambda^*}\|y\|.
\]

\begin{lemma}[Gaussian summation]\label{al-Gaussian-Poisson}
For $s>0$, $u\in\mathbb R^D$, and
$c=\operatorname{covol}(\Lambda)^{-1}$, one has
\begin{equation}\label{al-Gaussian-Poisson-eq}
 \sum_{x\in\Lambda}e^{-\pi s\|x-u\|^2}
 =c s^{-D/2}\sum_{y\in\Lambda^*}
       e^{-\pi\|y\|^2/s}\cos(2\pi\langle y,u\rangle).
\end{equation}
All sums converge absolutely, locally uniformly in $u$.
\end{lemma}
\begin{proof}
A lattice has $O((1+R)^D)$ points in a ball of radius $R$:
disjoint sufficiently small balls about its points lie in a fixed
enlargement of that ball. This proves the asserted convergence
and permits termwise integration on a fundamental parallelepiped.
The Fourier transform, with exponent $-2\pi i\langle y,z\rangle$,
of $e^{-\pi s\|z\|^2}$ is
$s^{-D/2}e^{-\pi\|y\|^2/s}$. In one dimension this follows by
differentiating the transform and integrating by parts, which
gives $G'(y)=-2\pi yG(y)/s$; the Gaussian integral gives
$G(0)=s^{-1/2}$. The latter integral follows by squaring it and
using polar coordinates. Multiplication proves the assertion in
dimension $D$.

Unfolding the periodization over a fundamental parallelepiped
now gives its Fourier coefficients as the terms on the right
of \eqref{al-Gaussian-Poisson-eq}, with the factor $c$.
The Fourier series is absolutely convergent. The difference
between it and the continuous periodization has all Fourier
coefficients zero and hence vanishes: convolution with products
of the one-dimensional Fej\'er kernels gives zero and converges
uniformly to the continuous difference. Pairing $y$ and $-y$
replaces the complex exponentials by the displayed cosines.
\end{proof}

\begin{theorem}[Covering and independent vectors]\label{al-transference}
If every nonzero point of $\Lambda^*$ has norm at least
$\lambda>0$, then every $u\in\mathbb R^D$ has some
$x\in\Lambda$ with
\begin{equation}\label{al-cover}
                         \|u-x\|\leq\frac D{2\lambda}.
\end{equation}
There are $D$ linearly independent vectors of $\Lambda$, each
of norm at most $D/\lambda$.
\end{theorem}
\begin{proof}
For $D=1$, write $\Lambda=a\mathbb Z$ in an orthonormal coordinate.
Then $\lambda\leq1/a$, and rounding to the nearest multiple of
$a$ proves both claims. Dimension zero is immediate.
Suppose $D\geq2$, put $m=D/2\geq1$, and suppose for a
contradiction that no lattice point is within $R=m/\lambda$
of some $u$. Define
\[
 A_s(u)=\sum_{x\in\Lambda}e^{-\pi s\|x-u\|^2},\quad
 Q_s=\sum_{y\in\Lambda^*}e^{-\pi\|y\|^2/s},\quad
 Q_s(u)=\sum_{y\in\Lambda^*}e^{-\pi\|y\|^2/s}
                              \cos(2\pi\langle y,u\rangle).
\]
For $0<\tau\leq1$, the distance assumption and Gaussian summation
give
\begin{align*}
 A_s(u)&\leq e^{-\pi s(1-\tau)R^2}A_{s\tau}(u),\\
 A_{s\tau}(u)&\leq c(s\tau)^{-m}Q_s.
\end{align*}
Indeed the cosines are at most one, and $Q_{s\tau}\leq Q_s$.
Consequently
\begin{equation}\label{al-Gauss-upper}
       Q_s(u)\leq e^{-\pi s(1-\tau)R^2}\tau^{-m}Q_s.
\end{equation}
For $0<\theta\leq1$, the lower bound on the nonzero dual vectors
gives
\[
 Q_s-Q_s(u)
 \leq2e^{-\pi(1-\theta)\lambda^2/s}Q_{s/\theta}.
\]
Also $A_{s/\theta}(0)\leq A_s(0)$ term by term, so Gaussian
summation at zero implies
$Q_{s/\theta}\leq\theta^{-m}Q_s$.
Since $Q_s>0$, combination with \eqref{al-Gauss-upper} yields
\begin{equation}\label{al-Gauss-contradiction}
 1\leq e^{-\pi s(1-\tau)R^2}\tau^{-m}
               +2e^{-\pi(1-\theta)\lambda^2/s}\theta^{-m}.
\end{equation}
Choose
\[
       s=\frac{18\lambda^2}{25m},\qquad
       \tau=\frac49,\qquad\theta=\frac14.
\]
The right side of \eqref{al-Gauss-contradiction} is
\[
 \left(\frac94e^{-2\pi/5}\right)^m
          +2\left(4e^{-25\pi/24}\right)^m<1.
\]
Here is an elementary numerical certificate. The perimeter of
an inscribed regular dodecagon gives
$\pi>3(\sqrt6-\sqrt2)>31/10$.
The first five terms of the exponential series at $31/25$
give $e^{31/25}>17/5$. Also
$e>163/60>19/7$ and
$e^{1/5}>1+1/5+1/50=61/50$, so
\[
 e^{25\pi/24}>e^{16/5}
   >(19/7)^3(61/50)>24.
\]
The two positive bases are therefore smaller than $45/68$ and
$1/6$. Since $m\geq1$, their indicated sum is smaller than
$45/68+1/3=203/204<1$. This contradiction proves
\eqref{al-cover}.

Put $R=D/(2\lambda)$. If the vectors of $\Lambda$ of norm at
most $2R$ failed to span $\mathbb R^D$, let $W$ be their span.
Discreteness gives $\varepsilon>0$ such that every lattice vector
of norm at most $2R+\varepsilon$ already has norm at most $2R$.
Choose $u\in W^\perp$ with $R<\|u\|<R+\varepsilon$.
By \eqref{al-cover} there is $x\in\Lambda$ with
$\|u-x\|\leq R$, and then $\|x\|<2R+\varepsilon$, so
$x\in W$. But orthogonality gives
$\|u-x\|\geq\|u\|>R$, a contradiction.
Select $D$ independent vectors from this spanning finite set.
\end{proof}

\subsection{The codifferent and short Hermitian dual vectors}

The codifferent is
\[
 \mathfrak D^{-1}
  =\{a\in E:\operatorname{Tr}_{E/\mathbb Q}(a\mathcal O)
                                                    \subseteq\mathbb Z\}.
\]
It is invariant under conjugation, because the trace is invariant
under every field automorphism. Its inverse has norm
$N\mathfrak D=\Delta$.

\begin{lemma}[A lower bound in the trace dual]\label{al-trace-dual}
Let $L$ be an integral full lattice and put $Y=L^\#$.
If $0\ne x\in V$ and
\[
       \operatorname{Tr}_{E/\mathbb Q}h(x,y)\in\mathbb Z
                                      \qquad(y\in Y),
\]
then
\begin{equation}\label{al-trace-lower}
             \operatorname{Tr}_{E/\mathbb Q}h(x,x)
                                      \geq\frac d{\rho^2}.
\end{equation}
\end{lemma}
\begin{proof}
Since $Y$ is an $\mathcal O$-module, the assumption implies
$h(x,y)\in\mathfrak D^{-1}$ for every $y\in Y$.
Since $L\subseteq Y$, for $\delta\in\mathfrak D$ and $l\in L$
we obtain
\[
 h(l,\delta x)=\delta\overline{h(x,l)}\in\mathcal O.
\]
Thus $\delta x\in Y$. Applying the previous inclusion with
$y=\delta x$ proves $\delta h(x,x)\in\mathfrak D^{-1}$.
Writing $\beta=h(x,x)$, this says
$\beta\mathfrak D^2\subseteq\mathcal O$.
It is a nonzero integral ideal, so
\[
       |N_{E/\mathbb Q}\beta|\,\Delta^2\geq1.
\]
All $d$ embeddings of $\beta$ are positive real numbers.
Their arithmetic-geometric mean inequality therefore gives
\[
 \operatorname{Tr}_{E/\mathbb Q}\beta
       \geq d\,|N_{E/\mathbb Q}\beta|^{1/d}
       \geq d\Delta^{-2/d}=d\rho^{-2}.
\]
This proves the claim without assuming that $L$ is unimodular
or prescribing the index $[L^\#:L]$.
\end{proof}

\begin{theorem}[Short independent dual vectors]\label{al-short-dual}
Every integral full lattice $L$ admits $n$ linearly independent
vectors $y_1,\ldots,y_n\in L^\#$ over $E$ such that
\begin{equation}\label{al-short-dual-bound}
       \frac1d\operatorname{Tr}_{E/\mathbb Q}h(y_j,y_j)
                                \leq n^2\rho^2\qquad(1\leq j\leq n).
\end{equation}
\end{theorem}
\begin{proof}
Equip $V\otimes_{\mathbb Q}\mathbb R$ with the positive definite
inner product $B(x,y)=\operatorname{Tr}_{E/\mathbb Q}h(x,y)$.
Its real dimension is $dn$, whereas the dimension of $V$ over
$E$ is $n$. The group $Y=L^\#$ is a full Euclidean lattice.
Every vector in its Euclidean dual belongs to $V$: a
$\mathbb Z$-basis of $Y$ has a rational nondegenerate Gram matrix,
whose inverse produces the dual basis inside $V$.
Lemma~\ref{al-trace-dual} therefore gives
\[
                     \lambda_1(Y^*)\geq\sqrt d/\rho.
\]
Theorem~\ref{al-transference} supplies $dn$ rationally independent
vectors of $Y$, each with
\[
 B(y,y)\leq\left(\frac{dn}{\sqrt d/\rho}\right)^2
                                                =dn^2\rho^2.
\]
They span $V$ over $\mathbb Q$, and hence also over $E$.
Select an $E$-basis from them. This proves
\eqref{al-short-dual-bound}. The selected tuple need not be an
$\mathcal O$-basis of $L^\#$; no such assertion is required.
\end{proof}

\begin{corollary}[The arithmetic input for the coefficient maps]
\label{al-coefficient-input}
For each decorated lattice of Proposition~\ref{al-kernel} with
$|S|=f$ and $NP=T$ for every $P\in S$, the degree and short-vector
inputs of Corollary~\ref{am-short} hold with
\[
                  Q=n^2\rho^2=n^2(2v_0)^2,
                         \qquad v_0=\rho/2.
\]
They remain true after any coefficient-field extension with
the induced metrics and normalized degree.
\end{corollary}
\begin{proof}
The degree is \eqref{al-actual-degree}. The mean over all
embeddings of $h(y_j,y_j)$ is its trace divided by $d$,
so Theorem~\ref{al-short-dual} is exactly the mean squared-norm
hypothesis in \eqref{am-short-assumption}. Hermitian pairing
with these dual vectors gives integral independent functionals.
The base-change assertions follow from the norm and archimedean
multiplicity calculation in Section~\ref{am-section}.
\end{proof}

\subsection{Finiteness of the actual decorated classes}

The short-vector theorem also gives a direct finiteness proof for
the class objects used by the coefficient maps. It does not require
a quantitative reduction theory or a class-number estimate.

\begin{theorem}[Finiteness of unimodular Hermitian lattice classes]
\label{al-finite-classes}
Fix the positive Hermitian space $(V,h)$ over $E$.
There are only finitely many isometry classes of full
$\mathcal O$-lattices $N\subset V$ satisfying $N=N^\#$.
\end{theorem}
\begin{proof}
The assertion is immediate when $n=0$, so assume $n\geq1$.
Apply Theorem~\ref{al-short-dual} to $L=N$. Since $N=N^\#$,
it supplies an $E$-basis $y_1,\ldots,y_n$ contained in $N$ with
\[
       \operatorname{Tr}_{E/\mathbb Q}h(y_i,y_i)
                      \leq C,\qquad C=dn^2\rho^2.
\]
Write $G=(h(y_i,y_j))$. Every entry belongs to $\mathcal O$.
At every embedding $\sigma:E\hookrightarrow\mathbb C$,
positivity and Cauchy--Schwarz give
\[
 |\sigma(G_{ij})|^2
       \leq \sigma(G_{ii})\sigma(G_{jj})\leq C^2,
\]
because each positive diagonal value is at most its full trace.
The Minkowski embedding of $\mathcal O$ is a lattice, hence
meets the compact set defined by $|\sigma(a)|\leq C$ in
only finitely many points. For completeness, if
$\omega_1,\ldots,\omega_d$ is an integral basis, the
Minkowski images of the $\omega_i$ form a real basis; its
inverse linear map sends a bounded set to a bounded set,
which contains only finitely many integer coordinate tuples.
Thus only finitely many Gram matrices $G$ can occur.

For each occurring $G$, choose one frame $z_1,\ldots,z_n$
with this Gram matrix, and put
$M=\sum_i\mathcal O z_i$. Any other frame with the same
Gram matrix is carried to this frame by the unique
$E$-linear map sending its $i$th vector to $z_i$; the Gram
identity proves that this map is an isometry. After applying
it, the original lattice therefore satisfies
\[
                         M\subseteq N\subseteq M^\#.
\]
The second inclusion follows from integrality of $N$.
The quotient $M^\#/M$ is finite: in the frame $z_i$, the
invertible matrix $G$ identifies $M^\#$ with a full
fractional lattice in $E^n$, and some nonzero integer
annihilates its quotient by $M$. There are only finitely
many subgroups of a finite abelian group, hence only
finitely many intermediate $\mathcal O$-lattices.
Taking the finite union over the possible $G$ proves the result.
\end{proof}

\begin{corollary}[Finiteness with the residue decorations]
\label{al-finite-decorated}
Fix finite sets $S,G$ of primes of $\mathcal O$. Decorate each
unimodular $N$ by a nonzero $\mathcal O$-linear map
$\ell_P:N\to\mathcal O/P$ for every $P\in S$, and by a
residue frame
\[
           N/QN\simeq(\mathcal O/Q)^n\qquad(Q\in G).
\]
Extend these maps to the appropriate localizations and then
by zero outside them, as above. The decorated objects have
finitely many orbits under the isometry group of $(V,h)$.
They consequently have finitely many orbits under its product
with any subgroup of the finite residue frame groups, acting
in the prescribed way. Their stabilizers in this product are
finite.
\end{corollary}
\begin{proof}
Choose representatives for the finitely many undecorated
classes from Theorem~\ref{al-finite-classes}. For a fixed
representative $N$, choose a finite set of $\mathcal O$-module
generators. A map $N\to\mathcal O/P$ is determined by its
values on these generators, so there are only finitely many
such maps. The module $N/QN$ is finite, so there are only
finitely many residue frames. Their localized extensions are
uniquely determined: if $s\notin P$ and $sx\in N$, the
extended value at $x$ is $\bar s^{-1}\ell_P(sx)$.
Thus adding the extensions introduces no further choices.
There are finitely many decorations of each chosen $N$.
Every decorated object is transported to one of these by an
isometry, which proves finiteness of the isometry orbits.
Passing to a larger group can only merge orbits.

An element stabilizing a decorated object has first component
in the finite group of isometries preserving $N$, by
Lemma~\ref{al-finite-isometries}. The second component belongs
to a finite product of finite matrix groups. The stabilizer
is therefore finite. Its triviality, when required, is the
separate two-prime conclusion of Proposition~\ref{al-two-primes}.
\end{proof}

\subsection{Counterexamples to omitted hypotheses and scope}

The following examples delimit the statements just proved.

\begin{itemize}
\item One residue characteristic is insufficient without a size
condition. For an odd prime $p$, multiplication by $\zeta_p$
on $E=\mathbb Q(\zeta_p)$ preserves $\mathcal O_E$ and the
standard positive Hermitian form, has order $p$, and reduces to
$1$ at the prime $(1-\zeta_p)$. At $p=2$, multiplication by $-1$
on $\mathbb Q(i)$ gives the same obstruction at $(1+i)$.
The two-prime statement retains different rational primes.
\item Positivity cannot be omitted from the finiteness step.
On $\mathbb Q(i)^2$ the indefinite Hermitian form
$h(x,y)=i(\bar x_1y_2-\bar x_2y_1)$ is preserved by the matrices
$\left(\begin{smallmatrix}1&m\\0&1\end{smallmatrix}\right)$
for $m\in\mathbb Z$. They preserve $\mathbb Z[i]^2$,
are unipotent at every prime, and need
not have finite order. The local torsion lemma always assumes
finite order, and the global argument obtains it from positivity.
\item Integrality is essential in Theorem~\ref{al-short-dual}.
Take $E=\mathbb Q(i)$, $h(x,y)=\bar xy$, and
$L=M^{-1}\mathbb Z[i]$ for an integer $M\geq3$.
Then $L^\#=M\mathbb Z[i]$, every nonzero dual vector has mean
squared norm at least $M^2$, whereas $n^2\rho^2=4$.
The omitted hypothesis $L\subseteq L^\#$ fails.
\item The index formula does not imply an arbitrary split
description of the discriminant module. For
$E=\mathbb Q(i)$, $N=\mathbb Z[i]$, and $S=\{P\}$ with
$P=(1+i)$, the kernel is $L=P$ and $L^\#=P^{-1}$.
Thus $L^\#/L\simeq\mathcal O/P^2$, which has cardinality
four but is not $(\mathcal O/P)^2$ as an $\mathcal O$-module:
the latter is killed by $P$ and the former is not.
The source's stronger split description retains the condition
that no selected prime is the conjugate of another.
\item The zero extension of a localized residue map is not
globally additive. In $\mathbb Q(i)$ at $P=(3)$,
both $1/3$ and $2/3$ lie outside $\mathcal O_P$ and have
zero extended value, while their sum $1$ has residue $1$.
Only additivity on the localization is used.
\end{itemize}

These are counterexamples to weakened hypotheses, not to the actual
statements. The arguments above close the standard decorated object's
existence, the kernel lattice's degree, the short-dual input, finiteness
of the decorated classes, and the two-prime stabilizer argument.
They do not identify all automorphic
class lattices with these decorated representatives, construct the
period maps, or prove the remaining global theta estimates.
Those identifications and analytic estimates require their separate
source arguments; no existence package is being substituted for them.


\section{Adelic classes and their integral polynomial coordinates}
\label{acl-section}

The coefficient estimates in Section~\ref{am-section} concern the
actual class lattices, rather than arbitrary lattices with the same
rank. We identify these objects explicitly. The argument below also
gives a short proof of the adelic existence statement needed at the
initial step of the construction. It uses the unramified quadratic
extension and the disjoint selected primes that occur in the
source~\cite{alpo-uniform-source}. Neither strong approximation for
the unitary group nor a class-number formula is required.

\subsection{Unit norms and integral orthonormal bases}

\begin{lemma}[Norms in an unramified quadratic extension]
\label{acl-local-norm}
Let $A$ be a complete discrete valuation ring with finite residue
field $\kappa$, and let $B$ be the ring of integers in an unramified
quadratic extension of its fraction field. Write $\bar x$ for the
nonidentity automorphism. Then the residue trace
$\kappa_B\to\kappa$ is surjective, and
\[
                         N_{B/A}:B^\times\longrightarrow A^\times
\]
is surjective. These statements include residue characteristic two.
\end{lemma}
\begin{proof}
Write $q=|\kappa|$. The trace is $x\mapsto x+x^q$ on
$\mathbb F_{q^2}$. In odd characteristic its value at $1$ is
nonzero. In characteristic two its value at any
$x\notin\mathbb F_q$ is nonzero, since $x^q\ne x$.
It is therefore a nonzero linear map from a two-dimensional
$\mathbb F_q$-space to a one-dimensional space, hence is onto.

The multiplicative group of a finite field is cyclic. One elementary
proof is to let $e$ be the exponent of this finite abelian group.
Combining elements of maximal prime-power orders gives an element
of order $e$. Every group element is a root of $X^e-1$, so the
group has at most $e$ elements; its exponent also divides its
order. The two numbers are equal. It follows that
$x\mapsto x^{q+1}$ maps $\mathbb F_{q^2}^\times$ onto
$\mathbb F_q^\times$.

Choose a uniformizer $\varpi$ of $A$, which is also a uniformizer
of $B$. Given $a\in A^\times$, first choose $u_1\in B^\times$
with $N(u_1)\equiv a\pmod\varpi$. Suppose inductively that
$N(u_j)\equiv a\pmod{\varpi^j}$. For $t\in B$,
\[
 N(1+\varpi^jt)
  =1+\varpi^j(t+\bar t)+\varpi^{2j}t\bar t
  \equiv1+\varpi^j\operatorname{Tr}(t)
                                     \pmod{\varpi^{j+1}}.
\]
Surjectivity of the residue trace permits a choice of $t_j$
such that $u_{j+1}=u_j(1+\varpi^jt_j)$ has
$N(u_{j+1})\equiv a\pmod{\varpi^{j+1}}$.
The $u_j$ converge in $B^\times$, because their successive
ratios belong to $1+\varpi^jB$. Their limit has norm $a$.
No division by two occurs.
\end{proof}

\begin{proposition}[Integral orthonormal bases]
\label{acl-local-basis}
Let $A$ be as above, and let $B$ be either its unramified quadratic
extension or $A\times A$, with the standard involution. Every
unimodular Hermitian form on a free $B$-module of rank $n$ has a
basis $e_1,\ldots,e_n$ with $h(e_i,e_j)=\delta_{ij}$.
\end{proposition}
\begin{proof}
First suppose $B$ is a discrete valuation ring. The Gram matrix
in an arbitrary basis belongs to $\operatorname{GL}_n(B)$.
If some diagonal entry is a unit, the corresponding vector has
unit norm. Otherwise some off-diagonal entry $a=h(e_i,e_j)$
is a unit. The reduction of
\[
 h(e_i+te_j,e_i+te_j)
   =h(e_i,e_i)+\operatorname{Tr}(ta)+t\bar t\,h(e_j,e_j)
\]
is nonzero for a suitable $t\in B$, by
Lemma~\ref{acl-local-norm}. Thus a vector of unit norm exists
also in this case. Multiplying it by a unit whose norm is the
inverse of its norm gives a vector $v$ with $h(v,v)=1$.

The integral projection $x\mapsto v h(v,x)$ splits off $Bv$.
Its orthogonal complement is a direct summand, hence free over
the discrete valuation ring. The Gram determinant shows that
the restricted form is again unimodular. Induction on $n$
proves the assertion, including the empty basis when $n=0$.

If $B=A\times A$, a free $B$-module is the product of two
free $A$-modules of rank $n$. A Hermitian form amounts to a
bilinear pairing between them, with the reverse pairing
determined by symmetry. Unimodularity says that this pairing
is perfect. Choose a basis on one factor and its dual basis
on the other. The corresponding $B$-basis has Gram matrix $1$.
\end{proof}

\subsection{Patching lattices and prescribing split-prime data}

Let $E/F$ be a CM quadratic extension, unramified at every finite
prime. Put $\mathcal O=\mathcal O_E$, and fix
\[
 V=E^n,\qquad h(x,y)=\sum_{i=1}^n\bar x_i y_i,
 \qquad N_0=\mathcal O^n,\qquad n\geq1.
\]
Let $U$ be its unitary group over $F$. At a finite place $v$ of
$F$, write $E_v=E\otimes_FF_v$ and $\mathcal O_v$ for its
maximal compact subring. The group $U(\mathbb A_{F,f})$ consists
of tuples in $U(F_v)$ preserving $(N_0)_v$ at all but finitely
many places.

\begin{lemma}[Patching full lattices]
\label{acl-patching}
Suppose a full lattice $M_P\subset E_P^n$ is prescribed at
each prime $P$ of $\mathcal O$, with $M_P=\mathcal O_P^n$
for all but finitely many $P$. Here $\mathcal O_P$ denotes the
completed discrete valuation ring. There is a unique full
$\mathcal O$-lattice $M\subset E^n$ with these completions:
\[
                         M=E^n\cap\prod_PM_P.
\]
In particular, for $g\in U(\mathbb A_{F,f})$,
\[
                         N_g=V\cap g\widehat N_0
\]
is unimodular and has completion $g_v(N_0)_v$ at every $v$.
\end{lemma}
\begin{proof}
Choose a nonzero rational integer $a$ such that, at every prime,
\[
                 a\mathcal O_P^n\subseteq M_P
                              \subseteq a^{-1}\mathcal O_P^n.
\]
Only finitely many local conditions need be imposed when choosing
$a$. The finite $\mathcal O$-module
$Q=a^{-1}\mathcal O^n/a\mathcal O^n$ is annihilated by $a^2$.
The Chinese remainder theorem decomposes it as the direct sum
of its localizations at the primes dividing $a$; these finite
modules are unchanged by completion. Each $M_P$ determines a
submodule of the corresponding local summand of $Q$. Take their
direct sum and its inverse image in $a^{-1}\mathcal O^n$.
The resulting module has exactly the required completions and
contains $a\mathcal O^n$, so it is full and finitely generated.

The same finite quotient shows that a lattice is the intersection
of its completions with $E^n$, proving uniqueness. Hermitian
duality commutes with completion over $F$: in local bases it is
given by the inverse transpose of the conjugate Gram matrix.
Each $g_v(N_0)_v$ is self-dual because $g_v$ is unitary.
Uniqueness therefore gives $N_g=N_g^\#$.
\end{proof}

Fix disjoint finite sets $S,G$ of primes of $\mathcal O$ such that
\begin{equation}\label{acl-primes}
               P\in S\cup G\quad\Longrightarrow\quad
                                      \bar P\notin S\cup G.
\end{equation}
Every selected prime lies over a split prime of $F$: in a quadratic
Galois extension an inert or ramified prime is fixed by conjugation.
Condition~\eqref{acl-primes} also says that no two selected primes
lie over the same prime of $F$.
At each $P\in S$, fix $c_P\in(\mathcal O/P)^\times$ and the
functional $\ell_{0,P}(x)=c_Px_n$ on $N_0$.
At each $Q\in G$, use the standard coordinate frame
$b_{0,Q}:N_0/QN_0\simeq(\mathcal O/Q)^n$.

A decorated lattice consists of a unimodular $N\subset V$, a
nonzero $\mathcal O$-linear map $\ell_P:N\to\mathcal O/P$ for each $P\in S$,
and an $\mathcal O/Q$-linear isomorphism
$b_Q:N/QN\to(\mathcal O/Q)^n$ for each
$Q\in G$. These are exactly the initial-step lattice and residue
data. The associated integral lattice is
\begin{equation}\label{acl-kernel}
                      L_d=\{x\in N:\ell_P(x)=0\text{ for all }P\in S\}.
\end{equation}
The maps on localizations, and their zero extensions when needed,
are uniquely determined by these restrictions, as in
Section~\ref{al-section}.

For $g\in U(\mathbb A_{F,f})$, transport the fixed functionals
and frames by $g^{-1}$ to obtain a decoration $d_g$ of $N_g$.
Local reduction agrees with $N_g/PN_g$ by the patching lemma.
In particular these functionals are nonzero and the frames are
isomorphisms.

\begin{theorem}[Every decorated lattice has an adelic representative]
\label{acl-points}
Every decorated lattice just defined equals $d_g$ for some
$g\in U(\mathbb A_{F,f})$.
\end{theorem}
\begin{proof}
Complete $N$ at every finite place $v$ of $F$. By
Proposition~\ref{acl-local-basis}, there is an isometry
$x_v:(N_0)_v\to N_v$. Outside finitely many places $N_v=(N_0)_v$,
and we take $x_v=1$ there. Hence $(x_v)$ is an adelic unitary
element.

At a selected split prime $v$, the local unitary group is
$\operatorname{GL}_n(F_v)$: after identifying
$E_v=F_v\times F_v$, its two components have the form
$A$ and $(A^{-1})^t$. The stabilizer of $(N_0)_v$ is
$\operatorname{GL}_n(\mathcal O_{F_v})$, whose reduction maps
onto $\operatorname{GL}_n(\kappa_v)$. To see the latter
surjectivity, lift the entries of any invertible residue matrix;
its determinant is a unit, so the lifted matrix is invertible.

If $P\in S$ lies over $v$, the pulled-back map $\ell_Px_v$
is a nonzero row vector over $\kappa_v$. The general linear
group acts transitively on nonzero row vectors: complete a
nonzero linear functional to a basis of the dual space.
Choose an integral local isometry $a_v$ such that
$\ell_Px_va_v=\ell_{0,P}$ modulo $P$.
If $Q\in G$ lies over $v$, choose instead $a_v$ with
$b_Qx_va_v=b_{0,Q}$ modulo $Q$, by lifting the inverse of
the residue matrix of $b_Qx_v$.
There is only one condition at each such $v$, by
\eqref{acl-primes}. Put $g_v=x_va_v$ at these places and
$g_v=x_v$ elsewhere. It has the same completed lattice as $x$,
and the prescribed functionals and frames. Thus $d_g=d$.
\end{proof}

Multiplying each $\ell_P$ by a fixed nonzero residue scalar
changes neither $L_d$ nor any stabilizer. This accounts for the
source's fixed diagonal rescaling of the level-prime decorations.
In particular, if $\pi\bar\pi=T$ and $(\pi,\bar\pi)=1$, then
$\bar\pi$ is a unit at each prime dividing $\pi$, so it is an
allowed choice of the scalar $c_P$. The proof does not assume
that $N$ has a global integral basis.

\subsection{Double cosets, stabilizers, and the inverse level character}

Let $K_0$ be the stabilizer of $\widehat N_0$, and let
$K_S\subseteq K_0$ preserve every $\ell_{0,P}$.
Reduction at the primes in $G$ gives a surjective homomorphism
\[
       r:K_S\longrightarrow\mathcal G,
       \qquad \mathcal G=\prod_{Q\in G}
                                \operatorname{GL}_n(\mathcal O/Q).
\]
Surjectivity follows by the same independent split-prime lifts
used in Theorem~\ref{acl-points}. Fix a subgroup $H\subseteq\mathcal G$
and put $K_H=r^{-1}(H)$ and $K_1=\ker r$.
These groups are compact and open: they impose finitely many
residue conditions inside the compact open group $K_0$.

The group $U(F)\times H$ acts on decorated lattices by
transporting the lattice and functionals with $\gamma\in U(F)$
and postcomposing the residue frames with $h\in H$.
Directly from the definition by inverse pullback,
\begin{equation}\label{acl-action}
 d_{\gamma g}=(\gamma,1)d_g,
 \qquad d_{gk}=(1,r(k)^{-1})d_g\quad(k\in K_S).
\end{equation}

\begin{proposition}\label{acl-double-cosets}
The maps above induce bijections
\[
 U(\mathbb A_{F,f})/K_1\simeq\{\text{decorated lattices}\},
 \qquad
 U(F)\backslash U(\mathbb A_{F,f})/K_H
       \simeq (U(F)\times H)\backslash\{\text{decorated lattices}\}.
\]
These double-coset sets are finite. For a representative $g$,
the stabilizer of $d_g$ in $U(F)\times H$ is isomorphic to
\begin{equation}\label{acl-rational-stabilizer}
                         \Gamma_g=U(F)\cap gK_Hg^{-1}
\end{equation}
by first projection. Its inverse sends $\gamma$ to
$(\gamma,r(g^{-1}\gamma g))$.
\end{proposition}
\begin{proof}
If $d_g=d_{g'}$, their completed lattices agree, so
$k=g^{-1}g'$ preserves $\widehat N_0$. Equality of the
functionals gives $k\in K_S$, and equality of the frames
gives $r(k)=1$. Conversely every such $k$ fixes $d_g$.
Theorem~\ref{acl-points} now proves the first bijection.
Equation~\eqref{acl-action}, together with the surjectivity
of $r$, proves the second.

Finiteness follows from Corollary~\ref{al-finite-decorated}.
For the stabilizer assertion, if
$\gamma\in\Gamma_g$ and $k=g^{-1}\gamma g$, then
\[
                  (\gamma,1)d_g=d_{\gamma g}=d_{gk}
                                      =(1,r(k)^{-1})d_g.
\]
Thus $(\gamma,r(k))$ fixes $d_g$. The converse follows by
the same computation and the first bijection. Once $\gamma$
is specified, the frame isomorphisms force the component in
$H$ uniquely, proving injectivity. These stabilizers are finite
by Lemma~\ref{al-finite-isometries}; at the two selected type
primes they are trivial under the unipotence hypotheses of
Proposition~\ref{al-two-primes}.
\end{proof}

Let $K$ be a number field containing the coefficient embeddings
of $E$ and the values of a character $\chi:H\to K^\times$.
Since $H$ is finite, these values are roots of unity. The
corresponding right level character is
\begin{equation}\label{acl-character}
                            \omega(k)=\chi(r(k))^{-1}.
\end{equation}
This inverse is necessary, as \eqref{acl-action} shows.

\subsection{Additional local bases and the row classes}

The row objects also occur in a description with actual local
bases at the auxiliary character primes. Forgetting a local basis
and remembering its residue frame does not change either the
class set or the stabilizer order. We give the argument explicitly,
since division by the group of basis changes would be incorrect:
that group is generally infinite.

\begin{lemma}\label{acl-free-fibers}
Let $p:\widetilde D\to D$ be a surjective map, equivariant for
a surjective homomorphism $\pi:\widetilde\Gamma\to\Gamma$.
Suppose $\ker\pi$ acts freely and transitively on each fiber
of $p$. Then $p$ induces a bijection of orbit sets, and for
every $x\in\widetilde D$ the homomorphism $\pi$ restricts
to an isomorphism
\[
           \operatorname{Stab}_{\widetilde\Gamma}(x)
                    \simeq\operatorname{Stab}_\Gamma(p(x)).
\]
\end{lemma}
\begin{proof}
If $p(y)=\gamma p(x)$, lift $\gamma$ to $g\in\widetilde\Gamma$.
Then $y$ and $gx$ lie in the same fiber, so $y=kgx$ for
some $k\in\ker\pi$. This proves injectivity on orbit sets;
surjectivity follows from that of $p$. For a stabilizer
element $\gamma$ downstairs, choose a lift $g$ and then $k$
with $kgx=x$. Thus the stabilizer map is onto. Its kernel
is trivial because $\ker\pi$ acts freely on the fiber.
\end{proof}

Let $Q$ be a finite set of additional primes disjoint from
$S\cup G$, and assume that $S\cup G\cup Q$ still contains
no conjugate pair. At $P\in Q$ put
$A_P=\mathcal O_{E,(P)}$, the localization inside $E$.
An enhanced decoration records a basis matrix
$B_P\in\operatorname{GL}_n(E)$ with
$B_PA_P^n=N\otimes_{\mathcal O}A_P$, rather than a residue
frame at $P$. Forget it by remembering
\[
                            b_P(x)=B_P^{-1}x\pmod P.
\]
Every residue frame has such a lift: first choose any basis
of the free $A_P$-module, then lift the required invertible
residue change of basis to $\operatorname{GL}_n(A_P)$.

For each $P\in Q$, choose a subgroup
$H_P\subseteq\operatorname{GL}_n(\mathcal O/P)$ and let
$\widetilde H_P$ be its inverse image in
$\operatorname{GL}_n(A_P)$. The acting groups are
\[
 \widetilde\Gamma=U(F)\times H\times\prod_{P\in Q}\widetilde H_P,
 \qquad
 \Gamma=U(F)\times H\times\prod_{P\in Q}H_P.
\]
The local action is $B_P\mapsto\gamma B_Pk_P^{-1}$.
Its residue frame is transformed by
$b_P\mapsto\bar k_Pb_P\gamma^{-1}$, so the forgetful map
is equivariant for the evident reduction homomorphism.
The kernel acts freely on each fiber: a basis matrix has no
nonidentity right stabilizer. It acts transitively as well,
since two bases with the same residue frame differ by a
unique matrix congruent to $1$ modulo $P$.
Lemma~\ref{acl-free-fibers} therefore applies.

\begin{corollary}[The row class set and its stabilizers]
\label{acl-row-classes}
Take $S=\varnothing$, so $L_d=N=N^\#$. The row descriptions
by enhanced decorations and by residue-decorated unimodular
lattices have canonically bijective class sets and isomorphic
stabilizers. The latter classes are the double cosets for
the compact open subgroup with the prescribed residue groups
at $G\cup Q$. All these sets are finite. The statement also
holds in rank zero: every group and every frame is then
trivial, and there is one zero lattice.
\end{corollary}
\begin{proof}
For positive rank, combine Lemma~\ref{acl-free-fibers} with
Proposition~\ref{acl-double-cosets}, using $G\cup Q$ in place
of $G$. The rank-zero assertion follows directly from the
definitions.
\end{proof}

Characters which factor through the residue groups pass through
this isomorphism with their exact values. If the row character
is denoted by $\chi_{\rm row}$ and the character on the local
basis description by $\chi_{\rm pt}$, the convention in the
source is
\begin{equation}\label{acl-row-character}
              \chi_{\rm pt}(g)
                        =\chi_{\rm row}(\pi g)^{-1}.
\end{equation}
Indeed, the former is the product of the chosen local
characters, while the latter is the product of their inverses.
This is compatible with the contragredient coefficient action
on one side of the period pairing and the covariant action
on the other. The equality of stabilizer orders above shows
that both descriptions have exactly the same weights in
their finite class sums. No index of the principal congruence
groups is introduced.

\subsection{The polynomial representation and the class lattice}

Fix embeddings $\tau:E\hookrightarrow K$ indexed by a finite
set $\Psi$, and dominant integral weights
$\lambda_\tau=(\lambda_{\tau,1},\ldots,\lambda_{\tau,n})$.
Write $a_\tau=\lambda_{\tau,n}$ and
$\nu_{\tau,i}=\lambda_{\tau,i}-a_\tau$. Let $W_\lambda(K)$
be the tensor product of the polynomial spaces already defined
in Section~\ref{coef:section}, with the determinant twists
recorded separately. In a block of variables $X_\tau$, the
covariant representation is
\[
 \rho_\lambda(A)F(X)
       =\left(\prod_\tau\det(A_\tau)^{a_\tau}\right)
                                         F((A_\tau^tX_\tau)_\tau).
\]
The representation used for evaluations on frames is its
contragredient realization
\begin{equation}\label{acl-contragredient}
 R_\lambda(\gamma)F(X)
    =\left(\prod_\tau\det(\tau\gamma)^{-a_\tau}\right)
                                  F(((\tau\gamma)^{-1}X_\tau)_\tau)
    =\rho_\lambda(((\tau\gamma)^{-1})^t)_\tau F.
\end{equation}
For an embedding $\sigma:K\hookrightarrow\mathbb C$, evaluate
$F$ on a tuple of invertible complex frames $Z=(Z_\tau)$ by
\begin{equation}\label{acl-evaluation}
               f_F(Z)=\prod_\tau\det(Z_\tau)^{a_\tau}
                                      \sigma(F)(Z).
\end{equation}
Negative $a_\tau$ are allowed, since the frames are invertible.
The defining identity is
\begin{equation}\label{acl-eval-action}
               f_{R_\lambda(\gamma)F}(Z)
                    =f_F(((\sigma\tau\gamma)^{-1}Z_\tau)_\tau).
\end{equation}

Evaluation is injective. Indeed, its determinant factor never
vanishes. If a polynomial vanishes on all invertible matrices,
fix any matrix $A$ and consider $A+tI$. Its determinant is a
nonzero polynomial in $t$, so the given polynomial vanishes
at $A+tI$ for all but finitely many $t$ and hence also at $t=0$.
Apply this argument to each block. A polynomial over an infinite
field which vanishes on every tuple has zero coefficients, by
induction on the number of variables.

For a full lattice $M\subset E^n$, denote by
$\Delta_{\lambda,\tau_\bullet}^{\mathrm{ctr}}(M)$ the
integral Weyl lattice transported by $R_\lambda$ from a local
basis of $M$ at every coefficient prime. Its independence of
these bases, finite generation, and compatibility with
localization were proved in Section~\ref{coef:section}.
Equivalently, at a prime $w$ of $K$, if $B_\tau$ is a basis
matrix of $M$ at $P_\tau=\tau^{-1}(w)$, this lattice is
\[
 \rho_\lambda\bigl(((\tau B_\tau)^{-1})^t\bigr)_\tau
                             \Delta_{\lambda}^{\mathrm{std}}(\mathcal O_{K,w}).
\]
This is the coefficient lattice in the actual class construction,
with $M=L_d^\#$.

\begin{lemma}[Conjugation in the coefficient lattice]
\label{acl-conjugate-lattice}
For the standard Hermitian form,
\begin{equation}\label{acl-conjugate-lattice-eq}
 \Delta_{\lambda,\tau_\bullet}^{\mathrm{ctr}}(L^\#)
                 =\Delta_{\lambda,\tau_\bullet\circ\bar{\ }}(L).
\end{equation}
The right-hand side is the covariant Weyl lattice evaluated
through the conjugate embeddings. The same identity holds
for the integral dual Weyl lattices.
\end{lemma}
\begin{proof}
Fix a coefficient prime $w$ and put $P_\tau=\tau^{-1}(w)$.
Choose a basis matrix $A_\tau$ of $L$ at the conjugate prime
$\bar P_\tau$. A basis of $L^\#$ at $P_\tau$ is
\[
                             B_\tau=(\bar A_\tau)^{-t}.
\]
This follows by writing the pairing in coordinates:
$h(Ax,y)=\bar x^t\bar A^ty$.
Substitution in the preceding local description gives
\[
        ((\tau B_\tau)^{-1})^t
                     =\tau(\bar A_\tau)
                     =(\tau\circ\bar{\ })(A_\tau).
\]
This is precisely the local description on the right of
\eqref{acl-conjugate-lattice-eq}, including every determinant
twist. Equality at all coefficient primes proves the global
identity. The argument applies to both integral coefficient
lattices because both have the same change-of-basis rule.
\end{proof}

Define $\mathcal F_K$ to be the space of functions $F$ from
decorated lattices to $W_\lambda(K)$ satisfying
\begin{equation}\label{acl-equivariant}
             F((\gamma,h)d)=\chi(h)R_\lambda(\gamma)F(d).
\end{equation}
Its integral lattice is
\[
 \mathcal F_{\mathcal O_K}
   =\{F\in\mathcal F_K:
             F(d)\in\Delta_{\lambda,\tau_\bullet}^{\mathrm{ctr}}(L_d^\#)
                                          \text{ for every }d\}.
\]
The condition is preserved by the action. Indeed,
$L_{(\gamma,h)d}=\gamma L_d$, the change-of-basis rule gives
the required coefficient-lattice transport, and roots of unity
are integral units.

\begin{proposition}[Evaluation at class representatives]
\label{acl-class-evaluation}
Choose representatives $d_c$ of the finite decorated class set,
and let $\Gamma_c$ be their stabilizers. Evaluation gives an
isomorphism
\begin{equation}\label{acl-evaluation-product}
 \mathcal F_K\simeq\bigoplus_c W_\lambda(K)^{\Gamma_c},
 \qquad
 \mathcal F_{\mathcal O_K}\simeq
   \bigoplus_c\left(
        \Delta_{\lambda,\tau_\bullet}^{\mathrm{ctr}}(L_{d_c}^\#)
                           \cap W_\lambda(K)^{\Gamma_c}\right),
\end{equation}
where $\Gamma_c$ acts by $\chi(h)R_\lambda(\gamma)$.
If the stabilizers are trivial, the invariant restrictions disappear.
\end{proposition}
\begin{proof}
An equivariant function takes a $\Gamma_c$-invariant value at
$d_c$, and all its other values are then determined. Conversely,
given invariant vectors $v_c$, define
\[
                     F((\gamma,h)d_c)
                               =\chi(h)R_\lambda(\gamma)v_c.
\]
Two choices differ by an element of $\Gamma_c$, so this is
well-defined. The representation law proves equivariance.
The lattice assertion follows from the transport property just
proved. In particular this argument neither averages over a
stabilizer nor introduces a denominator equal to its order.
\end{proof}

For clarity, localization introduces no such denominator either.
If $M_c$ is any finite collection of full coefficient lattices and
$W_0\subseteq\bigoplus_c W_\lambda(K)$ is a $K$-subspace, then
\begin{equation}\label{acl-localization}
 \left(W_0\cap\bigoplus_cM_c\right)_w
                     =W_0\cap\bigoplus_c(M_c)_w.
\end{equation}
For a vector on the right, each component has an integral
multiple with multiplier outside $w$. The product of the
finitely many multipliers works for every component and remains
outside $w$; it also preserves $W_0$. This proves the reverse
inclusion, while the forward inclusion is immediate. A full
lattice is the intersection of all its localizations: if
$x\notin M$, the proper annihilator of $x+M$ is contained in
some maximal ideal $w$, and $x\notin M_w$.
Consequently the displayed local coefficient conditions are
equivalent to the global lattice in
\eqref{acl-evaluation-product}.

\subsection{The adelic functions and the metric}

Over $\mathbb C$, associate to $F$ the function
\[
                              A_F(g,Z)=f_{F(d_g)}(Z).
\]
Equations~\eqref{acl-action}, \eqref{acl-eval-action}, and
\eqref{acl-equivariant} give exactly
\begin{equation}\label{acl-adelic-law}
 A_F(\gamma gk,((\sigma\tau(\gamma))Z_\tau)_\tau)
                           =\sigma(\omega(k))A_F(g,Z)
 \qquad(\gamma\in U(F),\ k\in K_H).
\end{equation}
Conversely, a function satisfying this law and having the prescribed
polynomial form at each $g$ determines $F$ uniquely.
Indeed, use Theorem~\ref{acl-points} to choose $g$ for each $d$;
another choice differs by $K_1$, whose character is one.
Injectivity of frame evaluation makes the polynomial independent
of the choice. To recover \eqref{acl-equivariant}, lift
$h^{-1}\in H$ to $k\in K_H$ and use
$d_{\gamma gk}=(\gamma,h)d_g$.
The rational subspace consists exactly of those functions whose
class polynomials have coefficients in $\sigma(K)$ and belong
to $\sigma(W_\lambda(K))$. The integral subspace is the one
specified in \eqref{acl-evaluation-product} and
\eqref{acl-localization}.

Give $W_\lambda(\mathbb C)$ its normalized Fock inner product,
as in Section~\ref{coef:section}. A unitary matrix acts
unitarily on the polynomial representation: the standard Gaussian
integral is invariant under its change of variables, and the
determinant character has absolute value one.
Every $\sigma\tau(\gamma)$ is unitary, since complex conjugation
on the CM field induces the given involution at every embedding.
Also $|\sigma\chi(h)|=1$. Hence the class norm is
\begin{equation}\label{acl-class-norm}
                  \|F\|_\sigma^2
                     =\sum_c\frac{\|\sigma(F(d_c))\|_{\mathrm{Fock}}^2}
                                             {|\Gamma_c|},
\end{equation}
independently of the representatives.

This is also the adelic quotient norm when the Haar measure is
normalized by $\operatorname{vol}(K_H)=1$. Indeed, the
double coset represented by $g_c$ is the quotient of $K_H$
by the finite group $g_c^{-1}\Gamma_{g_c}g_c$, and its quotient
mass is $1/|\Gamma_{g_c}|$. The coefficient norm is constant
on this double coset. The stabilizer identification in
Proposition~\ref{acl-double-cosets} gives exactly
\eqref{acl-class-norm}. One may equivalently take this finite
sum as the definition of the algebraic automorphic norm.
For an arbitrary fixed Haar measure $\mu$, the corresponding
integral is $\mu(K_H)$ times \eqref{acl-class-norm}.
Thus the algebraic class norm uses the measure divided by
$\mu(K_H)$; it does not assert that a separately chosen mass
or formal-degree normalization already has $\mu(K_H)=1$.
There is no additional factor $|H|$ in the finite class norm.
To justify the discrete quotient in this interpretation, a compact
open neighborhood meets $U(F)$ in a finite lattice-isometry group.
A smaller neighborhood excludes its nonidentity elements, proving
discreteness. The finite double-coset decomposition then reduces
the quotient integral to the displayed finite sum.

Under the actual two-prime unipotence hypotheses, every
$|\Gamma_c|$ equals one. Combining
\eqref{acl-conjugate-lattice-eq} with
\eqref{acl-evaluation-product} and \eqref{acl-class-norm}
therefore identifies the initial column lattice isometrically with
\begin{equation}\label{acl-actual-class-lattice}
                  \bigoplus_c
                         \Delta_{\lambda,\tau_\bullet\circ\bar{\ }}
                                      (L_{d_c}),
\end{equation}
with the standard coefficient metric on each summand.
The class lattices $L_{d_c}$ are integral kernels in unimodular
Hermitian lattices. Thus the degree and short-dual conclusions
of Section~\ref{al-section}, and the coefficient maps and slope
bound of Section~\ref{am-section}, apply to these actual objects.

\subsection{Necessary hypotheses and exact scope}

Several small examples guard the identifications above.

\begin{itemize}
\item The unramified hypothesis in Lemma~\ref{acl-local-norm}
cannot be dropped. In $\mathbb Q_2(i)/\mathbb Q_2$, a unit norm
is $a^2+b^2$ with exactly one of $a,b\in\mathbb Z_2$ odd,
so it is $1$ modulo $4$. In particular $-1$ is not a unit
norm. The unimodular rank-one forms with Gram entries $1$
and $-1$ are not integrally isometric. This does not affect
the unramified extension used above.
\item Arbitrary simultaneous decorations at conjugate split
primes cannot be prescribed. In rank one over
$\mathbb F_5\times\mathbb F_5$, a unitary matrix is
$(t,t^{-1})$. The two pulled-back functional coefficients
are $t^{-1}$ and $t$, with product one. The pair with
coefficients $1$ and $2$ cannot arise. Condition
\eqref{acl-primes} excludes precisely this obstruction.
\item Conjugation in \eqref{acl-conjugate-lattice-eq} matters.
In $E=\mathbb Q(i)$, let $P=(2+i)$ and take rank one and
weight $1$. Then $P^\#=\bar P^{-1}$, so its contragredient
coefficient lattice is $\bar P$, not $P$. The two ideals
are distinct.
\item Stabilizer invariants cannot be discarded without
triviality. A one-dimensional coefficient space with the
sign action of a stabilizer $\{\pm1\}$ has zero invariant
space, although the unrestricted coefficient space is
one-dimensional. The two-prime argument is used before
asserting the full direct sum in
\eqref{acl-actual-class-lattice}.
\end{itemize}

In the actual construction, the Hermitian spaces have standard
orthonormal frames and the standard integral lattice by their
definition. The class lattice is the inverse-coordinate image
of the selected decorated representative's kernel. Passing back
to coordinates gives exactly that kernel, and the frame
determinant is the ordinary matrix determinant. The selected
primes satisfy \eqref{acl-primes}: primes over $T$ divide
$\pi_T$ and do not divide $\bar\pi_T$, the type primes use
the prescribed conjugate half of the divisors of $g_1g_2$,
and the distinct rational primes keep these sets disjoint.
These are supplied hypotheses of the prime construction,
not new existence assumptions.

At the initial step the auxiliary coupling set is empty, the
index set has one element, and the reference character has no
auxiliary character primes. Thus no additional class label or
local-basis choice survives in \eqref{acl-actual-class-lattice}.
For the row step the coupling set is again empty; its extra
local bases are handled by Corollary~\ref{acl-row-classes}.

The present section proves the initial-step and row class
comparisons, the initial coefficient descriptions, and its
metric, together with the exact equality of stabilizer weights
needed for the row pairing. More elaborate intermediate steps also carry auxiliary
cosets and local character data. Their comparison with the
intermediate adelic level charts requires the corresponding
additional source arguments. The period-map valuation bound
and the global theta identities are not consequences of this
class comparison.


\section{Additional cosets at the penultimate step}
\label{acc-section}

The reference character used to construct the intermediate adelic
classes has no auxiliary character primes.  Consequently the only
additional datum at an intermediate step is a particular additive
coset, and this datum occurs only at the penultimate step.  We now
classify these cosets, prove their simultaneous existence on every
decorated lattice, and identify the resulting level subgroups.  This
proves the class comparison used at every step of the reference
construction, including its characters and stabilizer weights
\cite{alpo-uniform-source}.  The more general transfer between columns
with nontrivial auxiliary characters is a separate assertion.

\subsection{The selected rational primes are local uniformizers}

Write $\mathcal O=\mathcal O_E$, where $E$ is a CM field of degree
$2f$.  For a nonzero prime ideal $P$, put
$A_P=\mathcal O_{E,(P)}$ and $\kappa_P=\mathcal O/P$.
The valuation $v_P$ is normalized to have value one on a
uniformizer of the discrete valuation ring $A_P$.

\begin{lemma}\label{acc-uniformizer}
Let $p$ be a rational prime and let $\pi\in\mathcal O$ satisfy
\[
 \pi\overline\pi=p,\qquad (\pi)+(\overline\pi)=\mathcal O.
\]
Suppose that $f$ distinct prime ideals $P_1,\ldots,P_f$ divide
$(\overline\pi)$ and that $|\kappa_{P_i}|=p$ for every $i$.
Then
\[
 (\overline\pi)=\prod_{i=1}^f P_i,
 \qquad v_{P_i}(p)=1.
\]
In particular $p=|\kappa_{P_i}|$ is a uniformizer in $A_{P_i}$.
This assertion includes $p=2$.
\end{lemma}

\begin{proof}
Conjugation preserves the absolute norm of a principal ideal.
Taking ideal norms in $\pi\overline\pi=p$ therefore gives
\[
 N((\overline\pi))^2=p^{2f},
 \qquad N((\overline\pi))=p^f.
\]
The $P_i$ are distinct maximal ideals and hence pairwise coprime.
Thus $(\overline\pi)\subset\prod_iP_i$.
Both ideals have norm $p^f$, so the quotient of the larger ideal
by the smaller has one element.  They are equal.
Coprimality of $(\pi)$ and $(\overline\pi)$ gives
$v_{P_i}(\pi)=0$.  Taking valuations in the displayed product
identity now gives $v_{P_i}(p)=0+1=1$.
\end{proof}

The actual construction has two distinct rational primes $g_1,g_2$.
For each it supplies exactly the hypotheses of this lemma, with
$f$ selected prime ideals.  Let $\mathcal G$ be their disjoint union.
For $P\in\mathcal G$ write
\[
 t_P=|\kappa_P|\in\{g_1,g_2\}.
\]
We have proved $t_PA_P=PA_P$.  This equality is the precise
hypothesis needed below; membership $t_P\in P$ alone would not suffice.

\subsection{A quotient of intersections of local lattices}

The next statement applies to any number field, without a Hermitian
form.  Let $N\subset E^n$ be a full $\mathcal O$-lattice,
let $\mathcal G$ be a finite set of nonzero prime ideals, and
choose nonzero $t_P\in\mathcal O$ with $v_P(t_P)=1$.
Write $N_P=A_PN$, and fix a $\kappa_P$-linear isomorphism
\[
 b_P:N/PN\longrightarrow\kappa_P^n.
\]
It extends uniquely to $N_P/PN_P$ by inverting the reductions
of elements outside $P$.  We use $b_P$ also for this extended
reduction map on $N_P$.
Define $\mathcal O$-submodules of $E^n$ by
\[
 L(N)=\bigcap_{P\in\mathcal G}N_P,
 \qquad
 L^-(N)=\bigcap_{P\in\mathcal G}t_P^{-1}N_P.
\]
An intersection over an empty set is $E^n$.
These modules need not be finitely generated over $\mathcal O$.
Only their quotient, rather than either module separately, is finite.

\begin{theorem}\label{acc-residues}
The map
\begin{equation}\label{acc-reduction-map}
 L^-(N)\longrightarrow\mathcal A_n:=
       \prod_{P\in\mathcal G}\kappa_P^n,
 \qquad y\longmapsto \bigl(b_P(t_Py)\bigr)_P
\end{equation}
is surjective and has kernel $L(N)$.  It induces an isomorphism
of $\mathcal O$-modules
\[
                    L^-(N)/L(N)\simeq\mathcal A_n.
\]
In particular every element of $\mathcal A_n$ is realized on every
lattice $N$, and
\[
 |L^-(N)/L(N)|=\prod_{P\in\mathcal G}|\kappa_P|^n.
\]
\end{theorem}

\begin{proof}
Since $t_PA_P=PA_P$, the equality $b_P(t_Py)=0$ is equivalent to
$t_Py\in t_PN_P$, and hence to $y\in N_P$.
This proves the kernel assertion and also shows that
$L(N)\subset L^-(N)$.

For surjectivity, prescribe vectors $a_P\in\kappa_P^n$.
Choose $x_P\in N$ with $b_P(x_P)=a_P$.
For each fixed $P$, the Chinese remainder theorem gives
$e_P\in\mathcal O$ such that
\[
 e_P\equiv1\pmod P,\qquad
 e_P\in Q^{v_Q(t_P)}\quad
       (Q\in\mathcal G,\ Q\ne P).
\]
An exponent zero imposes no condition.  The moduli in these
conditions are pairwise coprime, so the theorem applies.
Put
\[
                         y=\sum_{P\in\mathcal G}
                                      \frac{e_P}{t_P}x_P.
\]
Fix $Q\in\mathcal G$.  The $P=Q$ term in $t_Qy$ is
$e_Qx_Q\in N_Q$ and has reduction $a_Q$.
For $P\ne Q$ the chosen divisibility gives
$e_P/t_P\in A_Q$, so
\[
             t_Q(e_P/t_P)x_P\in t_QN_Q=QN_Q.
\]
Thus $t_Qy\in N_Q$ and $b_Q(t_Qy)=a_Q$.
This proves surjectivity simultaneously at all the selected primes.
When $\mathcal G$ is empty, both the quotient and the product have
one element and $y=0$ suffices.  The same proof permits $n=0$.
The first isomorphism theorem and the cardinalities of the residue
vector spaces prove the remaining assertions.
\end{proof}

For the actual two rational primes the number in this theorem is
\begin{equation}\label{acc-number-cosets}
                         |\mathcal A_n|=(g_1g_2)^{fn}.
\end{equation}
A permitted additional coset is
\[
                          Y=y+L(N),\qquad y\in L^-(N).
\]
Its residue vector, denoted $a(Y)$, is the vector in
\eqref{acc-reduction-map}.  It does not depend on $y\in Y$.
The theorem says precisely that $Y\mapsto a(Y)$ is a bijection.

\subsection{Coordinates and change of residue frames}

Choose an $A_P$-basis matrix $F_P$ of $N_P$ that lifts the
inverse of $b_P$; explicitly,
\[
 F_P A_P^n=N_P,
 \qquad b_P(F_Px)=\overline x\quad(x\in A_P^n).
\]
Such a matrix exists: choose any local basis, then correct its
residue matrix by a lift in $\operatorname{GL}_n(A_P)$.
The reduction map on this general linear group is onto, since
any entrywise lift of an invertible residue matrix has unit determinant.
Define the coset in these coordinates by
\begin{equation}\label{acc-coordinate-coset}
                    C_P(Y)=F_P^{-1}y+A_P^n\subset E^n.
\end{equation}

\begin{proposition}\label{acc-coordinate-classification}
The coset \eqref{acc-coordinate-coset} is independent of the
representative $y\in Y$ and of the compatible basis matrix $F_P$.
For permitted cosets on two possibly different lattices with
residue frames,
\[
 \bigl(C_P(Y)\bigr)_{P\in\mathcal G}
   =\bigl(C_P(Y')\bigr)_{P\in\mathcal G}
                   \quad\Longleftrightarrow\quad a(Y)=a(Y').
\]
\end{proposition}

\begin{proof}
Changing $y$ by an element of $L(N)$ changes $F_P^{-1}y$ by
an element of $A_P^n$.  Moreover
$t_PF_P^{-1}y\in A_P^n$ and its reduction is $a_P(Y)$.
For $u,u'\in t_P^{-1}A_P^n$,
\[
 u+A_P^n=u'+A_P^n
   \ \Longleftrightarrow\ u-u'\in A_P^n
   \ \Longleftrightarrow\ \overline{t_Pu}=\overline{t_Pu'}.
\]
The last equivalence uses $t_PA_P=PA_P$.
This already proves the asserted comparison between two lattices.
For completeness, another compatible basis is $F'_P=F_Pk$,
where $k\in\operatorname{GL}_n(A_P)$ has reduction one.
Then $k^{-1}-1\in t_PM_n(A_P)$; hence
$(k^{-1}-1)F_P^{-1}y\in A_P^n$.
The coordinate coset is unchanged.
\end{proof}

Return to the positive Hermitian space and decorated lattices of
Section~\ref{acl-section}.  Here the selected residue-frame primes
are $\mathcal G$ and the functional primes are disjoint from them.
Let $\mathscr D$ be the set of decorated lattices and let
\[
                   \mathscr P=\{(d,Y):d\in\mathscr D,
                        \ Y\in L^-(N_d)/L(N_d)\}.
\]
The notation treats a quotient element as its additive coset.
The rational unitary group and the residue general linear groups act by
\[
 (\gamma,h)(d,Y)=((\gamma,h)d,\gamma Y),
 \qquad h=(h_P)_P\in\prod_P\operatorname{GL}_n(\kappa_P).
\]
Recall that the new residue frame is
$b'_P=h_Pb_P\gamma^{-1}$.
Since $t_P\in E$, it commutes with the $E$-linear map $\gamma$.
Consequently
\begin{equation}\label{acc-equivariance}
                a_P(\gamma Y)=h_Pa_P(Y).
\end{equation}
Theorems~\ref{acc-residues} and \ref{acc-coordinate-classification}
therefore give a bijection, compatible with both actions,
\begin{equation}\label{acc-product-description}
 \mathscr P\simeq\mathscr D\times\mathcal A_n,
 \qquad
 (\gamma,h)(d,a)=((\gamma,h)d,ha).
\end{equation}
In particular a rational unitary change alone preserves the coordinate
cosets, even when the chosen lifts of the residue frames are different.
A residue-frame change acts on their labels by the ordinary matrix
action.  No choice of global basis of $N_d$ is involved.

\subsection{The exact partition into adelic classes}

Let $B$ be a finite set of indices.  For $b\in B$ let
\[
 H_b=\prod_{P\in\mathcal G}H_{b,P}
      \subset\prod_{P\in\mathcal G}\operatorname{GL}_n(\kappa_P)
\]
be a product of subgroups, and let $\chi_b:H_b\to K^\times$
be a character in the coefficient field.  For $a\in\mathcal A_n$
put
\begin{equation}\label{acc-stabilizer-level}
 H_{b,a}=\{h\in H_b:ha=a\}
      =\prod_{P\in\mathcal G}
            \bigl(H_{b,P}\cap\operatorname{Stab}(a_P)\bigr).
\end{equation}
Here $\operatorname{Stab}(a_P)$ fixes the vector $a_P$, not merely
its one-dimensional span.  Choose one representative $a$ of each
$H_b$-orbit in $\mathcal A_n$.

\begin{theorem}\label{acc-class-partition}
There is a bijection
\begin{equation}\label{acc-class-bijection}
 \coprod_{b\in B}\ \coprod_{a\in H_b\backslash\mathcal A_n}
       (U(F)\times H_{b,a})\backslash\mathscr D
 \ \longrightarrow\
 \coprod_{b\in B}(U(F)\times H_b)\backslash\mathscr P,
 \qquad [d]\longmapsto[(d,a)].
\end{equation}
Every stabilizer on the left is isomorphic to the stabilizer of
its corresponding point on the right.  These isomorphisms preserve
the rational unitary component and the character, whose restriction
on the left is $\chi_b|_{H_{b,a}}$.

All class sets in \eqref{acc-class-bijection} are finite.
The number of displayed pieces is at most
$|B|\,|\mathcal A_n|$.
\end{theorem}

\begin{proof}
Use the product description \eqref{acc-product-description}.
Given $(d,a')$, choose $h\in H_b$ carrying $a'$ to the selected
representative $a$.  Applying $(1,h)$ puts its class in the indicated
piece.  Distinct $H_b$-orbits of vectors cannot meet.
If $(\gamma,h)(d,a)=(d',a)$, then $ha=a$, so
$h\in H_{b,a}$; this proves injectivity within the piece.
The same observation identifies the stabilizer groups themselves:
\[
 \operatorname{Stab}_{U(F)\times H_b}(d,a)
       =\operatorname{Stab}_{U(F)\times H_{b,a}}(d).
\]
It does not only identify their orders or their images in $U(F)$.
Both the rational component and the restricted character are unchanged.

Each group $H_{b,a}$ is a finite residue subgroup, so the class
finiteness theorem of Section~\ref{acl-section} applies.
There are finitely many $b$ and finitely many vector orbits.
There are at most $|\mathcal A_n|$ vector orbits for each $b$,
which proves the quantitative assertion.
\end{proof}

At every step other than the penultimate step the additional datum is
the empty set, rather than a nonempty additive coset.
For those steps replace $\mathcal A_n$ in the theorem by a one-element
set with trivial action.  Then $H_{b,a}=H_b$ and there is exactly
one piece for each $b$.  In the actual construction $m\geq2$,
so both the initial and the final step are of this kind.
Their index set $B$ has one element, giving a unique piece at each end.
At the penultimate step the number of pieces is bounded by
\[
                            |B|(g_1g_2)^{fn}.
\]

The source chooses representatives by first fixing the cosets in local
coordinates.  Proposition~\ref{acc-coordinate-classification} shows
that this is equivalent to fixing the vectors $a_P$ above.
The orbit of such a vector depends only on its $H_b$-orbit and not on
the underlying decorated lattice, by the simultaneous surjectivity
in Theorem~\ref{acc-residues}.  Thus the choice of preferred
representatives introduces no dependence on an unavailable lattice
or later parameter.  It selects exactly the pieces in
\eqref{acc-class-bijection}, up to renaming their indices.

\subsection{Level subgroups, coefficients, and all volume factors}

Use the notation of Section~\ref{acl-section}:
$K_S$ is the compact open subgroup preserving the base lattice and
its fixed residue functionals, and
\[
 r:K_S\longrightarrow
           \prod_{P\in\mathcal G}\operatorname{GL}_n(\kappa_P)
\]
is the surjective residue map.  Put
$K_{b}=r^{-1}(H_b)$ and $K_{b,a}=r^{-1}(H_{b,a})$.
Theorem~\ref{acl-double-cosets} identifies the $(b,a)$ piece with
\begin{equation}\label{acc-adelic-piece}
                         U(F)\backslash U(\mathbb A_{F,f})/K_{b,a}.
\end{equation}
The exact level character is
\begin{equation}\label{acc-character}
       \omega_{b,a}(k)=\chi_b(r(k))^{-1}
                              \qquad(k\in K_{b,a}).
\end{equation}
Indeed right multiplication by $k$ changes the residue frames by
$r(k)^{-1}$, so this follows directly from
\eqref{acc-equivariance} and the equivariance of coefficient functions.
No additional character comes from the additive coset.

More explicitly, suppose coefficient functions on $\mathscr P$ satisfy
\[
 F((\gamma,h)(d,a))=\chi_b(h)R(\gamma)F(d,a),
\]
where $R$ is the polynomial representation in
Section~\ref{acl-section}.  Restriction to a fixed representative $a$
gives exactly the corresponding equivariant coefficient function on
$\mathscr D$ for $U(F)\times H_{b,a}$.
Conversely, if $F_a$ is the function on the chosen piece, choose
$g=(\gamma,h)$ with $ha=a'$ and set
\[
 F(d,a')=\chi_b(h)R(\gamma)F_a(g^{-1}d).
\]
Two choices differ by an element of $U(F)\times H_{b,a}$,
so equivariance makes the result independent of the choice.
This proves the coefficient-space decomposition.

The integral lattice attached to a point depends on its functional
kernel lattice $L_d$ and on its fixed weight.  It does not depend on
$Y$.  Therefore the same restriction and transport identify the
integral coefficient lattices, as well as their rational and complex
spans.  In particular, the conjugate-prime coefficient identity of
Section~\ref{acl-section} remains valid on every piece.
If stabilizers are nontrivial their invariant subspaces are retained.
The weighted Hermitian norm is the orthogonal sum of the piece norms,
because corresponding stabilizers are the identical groups in
Theorem~\ref{acc-class-partition}.

\begin{proposition}\label{acc-index-volume}
For each vector $a$,
\[
 [K_b:K_{b,a}]=[H_b:H_{b,a}]=|H_ba|.
\]
For any one fixed Haar measure $\mu$ on the finite adelic group,
\[
                     \mu(K_{b,a})=\frac{\mu(K_b)}{|H_ba|}.
\]
The quotient integral for the piece
\eqref{acc-adelic-piece} is consequently
$\mu(K_{b,a})$ times its finite class sum with weights
$1/|\operatorname{Stab}|$.
\end{proposition}

\begin{proof}
Surjectivity of $r$ identifies the coset spaces
$K_b/K_{b,a}$ and $H_b/H_{b,a}$.
The map $hH_{b,a}\mapsto ha$ is a bijection from the latter
coset space to the orbit of $a$.
The finite disjoint union of $[K_b:K_{b,a}]$ translates of
$K_{b,a}$ has measure $\mu(K_b)$, giving the volume formula.
The quotient integral formula is the one proved in
Section~\ref{acl-section}, now at the subgroup $K_{b,a}$.
\end{proof}

This proof also applies to the source's local coset stabilizer:
let $A$ be the ring of integers of a nonarchimedean local field $k$,
with uniformizer $\varpi$.  If $J\subset\operatorname{GL}_n(A)$
preserves a lattice $\Lambda\subset k^n$, it acts on $k^n/\Lambda$; the index in $J$ of the
stabilizer of a quotient element is exactly its orbit size whenever
that orbit is finite.  One proves this by the same bijection from
stabilizer cosets to orbit elements.  For $\Lambda=\varpi^M A^n$, with $M\geq0$ an integer, and
vectors in $\varpi^{-M}A^n$, the quotient containing the orbit is finite.
Thus the relative-index factor in the additional local-level weight
is an orbit cardinality, with no omitted inverse.

There is a useful direct check on the stabilizer weights.
For a fixed decorated class represented by $d$, let
$\Gamma=\operatorname{Stab}_{U(F)\times H_b}(d)$.
Its inverse image in the enhanced class set consists of the
$\Gamma$-orbits in $\mathcal A_n$.
For any function $w$ constant on those orbits, orbit--stabilizer gives
\begin{equation}\label{acc-weighted-fiber}
 \sum_{[a]\in\Gamma\backslash\mathcal A_n}
           \frac{w(a)}{|\Gamma_a|}
      =\frac1{|\Gamma|}\sum_{a\in\mathcal A_n}w(a).
\end{equation}
To prove it, expand the sum on the right as a sum over orbits;
each orbit has $|\Gamma|/|\Gamma_a|$ elements.
In particular adjoining all the additional cosets multiplies the
weighted mass of a decorated class by $|\mathcal A_n|$.
Restricting to one piece instead uses its actual stabilizer and does
not insert this factor.  These two different operations must not be
confused.

\subsection{Checks on the necessary hypotheses}

\begin{example}\label{acc-norm-counterexample}
The induced map $L^-(N)/L(N)\to\mathcal A_n$ in
Theorem~\ref{acc-residues} need not be injective if
$v_P(t_P)=1$ is replaced by $t_P\in P$.
This can fail even when the number field is unramified at the prime.
Take $E=\mathbb Q(i)$, $P=(3)$, $A=A_P$ and $N_P=A$.
The polynomial $X^2+1$ is irreducible modulo $3$, so
$|\kappa_P|=9$ and $3$ is a uniformizer in $A$.
Set $t_P=|\kappa_P|=9$.
The cosets $A$ and $1/3+A$ are distinct, because $1/3\notin A$,
but both have residue $\overline{9y}=0$.
Thus the size of the residue field cannot be used as a uniformizer
at an arbitrary prime.  Lemma~\ref{acc-uniformizer} proves exactly
the stronger condition used at the actual selected primes.
\end{example}

\begin{example}\label{acc-ramified-counterexample}
A residue-degree-one hypothesis by itself also does not suffice.
In $E=\mathbb Q(\sqrt2)$ at $P=(\sqrt2)$, the residue field is
$\mathbb F_2$, but $2=(\sqrt2)^2$ has valuation two.
The distinct cosets $A_P$ and $1/\sqrt2+A_P$ both map to zero
under $y\mapsto\overline{2y}$.
The coprime factorization in Lemma~\ref{acc-uniformizer}
excludes this ramification at the selected primes.
\end{example}

\begin{example}\label{acc-line-counterexample}
Replacing the vector stabilizer by a line stabilizer changes the
level and its index.  On $\mathbb F_3$, let
$H=\mathbb F_3^\times$ act by multiplication and let $a=1$.
Its vector stabilizer is trivial and its orbit has two elements.
Its line stabilizer is all of $H$, since the vector space has
dimension one.  The resulting indices are two and one respectively.
The actual coordinate datum fixes an additive coset and hence a
vector, so it requires the former subgroup.
\end{example}

\begin{example}\label{acc-weight-counterexample}
In the same action of $\Gamma=\{\pm1\}$ on $\mathbb F_3$,
there are two orbits, $\{0\}$ and $\{1,-1\}$.
Their stabilizer orders are two and one.  Their weighted sum is
$1/2+1=3/2$, as \eqref{acc-weighted-fiber} predicts.
Using the original stabilizer order for both orbits would give one.
Thus an unchanged stabilizer denominator after adjoining a coset
would be incorrect; the actual chart comparison uses the smaller
stabilizer and has the stated exact equality.
\end{example}

These examples concern weakened statements, not flaws in the
actual source hypotheses.  The argument above proves the additional
coset classification, its simultaneous realization, the finite
partition into reference-step classes, and their exact integral
coefficients, characters, stabilizers and volume factors.
It uses the selected-prime hypotheses already verified in the
arithmetic construction.  It does not assert a comparison for
arbitrary nontrivial auxiliary character primes, nor does it prove
the remaining theta-period identities or the automorphic spectral
and mass formulas.


\section{Integral subspaces and the passage to period pairings}
\label{ip-section}

The local divisor argument applies a valuation bound to every integral
test vector after extending the coefficient field. This passage requires
an equality of integral submodules, not merely an equality of their
complex spans. We prove that equality and its valuation consequence.
The actual lower bounds for the theta periods remain separate inputs;
the following argument specifies exactly what their extension requires.

\subsection{Intersection with a rational subspace after extension}

Let $K_0\subseteq K$ be number fields with rings of integers $A$ and
$B$. Fix an embedding $\sigma:K\hookrightarrow\mathbb C$.
Let $V_0$ be a finite-dimensional $K_0$-vector space, let $M$ be a
full $A$-lattice in $V_0$, and identify
$V_0\otimes_{K_0,\sigma}\mathbb C$ with its specified complex
realization $V$. In particular a $K_0$-basis of $V_0$ is linearly
independent over $\mathbb C$ in this realization. Write $BM$ for
the $B$-span of $M$ in $V$.

\begin{lemma}\label{ip-intersection}
Let $e:V\to V$ be a complex-linear idempotent. Suppose there is
$0\ne a\in A$ such that $ae(M)\subseteq M$. Then
\begin{equation}\label{ip-cut}
                 BM\cap eV=B(M\cap eV).
\end{equation}
This equality holds at ramified primes as well as at unramified primes.
\end{lemma}
\begin{proof}
Since $M$ spans $V_0$, the denominator condition implies that $e$
preserves $V_0$. Set $q=a(1-e)$. It is an $A$-linear endomorphism
of $M$, and its complex-linear extension has kernel $eV$: the
conditions $e^2=e$ and $a\ne0$ show that $qv=0$ is equivalent to
$ev=v$. Thus $\ker(q|_M)=M\cap eV$.

We give the integral base-change argument explicitly. At a nonzero
prime $\mathfrak p$ of $A$, the module $M_{\mathfrak p}$ is free
of finite rank over the discrete valuation ring $A_{\mathfrak p}$.
Choose a basis. Smith normal form, with independent changes of bases
in the source and target, makes the matrix of $q$ diagonal, with
nonzero entries followed by zero entries. Its kernel is the span of
the basis vectors in those zero positions.
The ring $(A\setminus\mathfrak p)^{-1}B$ is a subring of $K$.
Multiplication by a nonzero element of $A_{\mathfrak p}$ remains
injective on it. Therefore the kernel of the same diagonal matrix
after tensoring is still exactly the span of those zero positions.
It follows that the natural map
\[
 B\otimes_A\ker(q|_M)\longrightarrow
          \ker\bigl(1\otimes q:B\otimes_A M\to B\otimes_A M\bigr)
\]
is an isomorphism after localization at every $\mathfrak p$.
It is an isomorphism globally: a nonzero element of its kernel or
cokernel would remain nonzero at a maximal ideal containing its
annihilator, contradicting the local calculation.

The map $B\otimes_A M\to V$ is injective as well. Locally over
$A_{\mathfrak p}$ its chosen basis of $M_{\mathfrak p}$ is a
$K_0$-basis and hence a complex-linearly independent tuple in $V$.
The localized evaluation map is consequently injective; the same
local criterion proves injectivity globally. Its image is $BM$.
Under this identification the last displayed kernel is $BM\cap eV$,
and the image of $B\otimes_A\ker(q|_M)$ is $B(M\cap eV)$.
This proves \eqref{ip-cut}. No assumption about an integral basis
of $B$ over $A$, or about prime splitting, has been made.
\end{proof}

The same proof applies to a finite matrix of maps with a common
denominator: after clearing that denominator, Smith normal form
computes the common kernel locally. This observation is useful if
the desired subspace is specified by linear equations instead of
by an idempotent.

\begin{example}[The subspace must be defined over the original field]
Take $K_0=\mathbb Q$, $K=\mathbb Q(\sqrt2)$,
$M=\mathbb Z^2\subset\mathbb C^2$, and
$e(x,y)=(x,\sqrt2x)$. Then $e^2=e$ and
$eV=\mathbb C(1,\sqrt2)$, but $M\cap eV=0$.
Indeed an integer pair on that line must have both coordinates zero.
On the other hand $(1,\sqrt2)\in BM\cap eV$, so
\eqref{ip-cut} fails. There is no nonzero integer $a$ for which
$ae(M)\subseteq M$. This is a counterexample to omitting the
rationality condition, not to Lemma~\ref{ip-intersection}.
\end{example}

The faithful complex realization is a separate condition. An arbitrary
torsion-free module inside a complex vector space need not satisfy it:
the map $\mathbb C\otimes_{\mathbb Z}(\mathbb Z+\mathbb Z\sqrt2)
\to\mathbb C$ has the nonzero kernel element
$1\otimes\sqrt2-\sqrt2\otimes1$. The source explicitly requires
injectivity of complexification in the hypothesis called an integral
form. Its subsequent intersection argument does not deduce this
injectivity from torsion-freeness alone.

\subsection{The integral denominator of the projectors}

Here is the application of the preceding lemma to the finite-group
subspaces. Let a finite group $G$ act on $V_0$, preserving $M$.
Assume that $K_0$ contains the values of the irreducible characters
in a chosen set $\mathcal X$ and of a one-dimensional character
$\psi:P\to K_0^\times$ of a subgroup $P\subseteq G$.
The two projectors are
\[
 e_{\mathcal X}
   =\frac1{|G|}\sum_{g\in G}\sum_{\chi\in\mathcal X}
                 \chi(1)\chi(g^{-1})\rho(g),
 \qquad
 e_\psi=\frac1{|P|}\sum_{p\in P}\psi(p)^{-1}\rho(p).
\]
The finite-group argument in Section~\ref{div:section} proves that
they are commuting idempotents: the first is central, and the second
projects to the $\psi$-subspace. Character values and roots of unity
are algebraic integers. Consequently, for $e=e_{\mathcal X}e_\psi$,
\begin{equation}\label{ip-denominator}
                     |G|\,|P|\,e(M)\subseteq M.
\end{equation}
Indeed, multiplying the two displayed sums by $|G||P|$ leaves
a finite sum of integral scalar multiples of the operators $\rho(gp)$,
each preserving $M$. Lemma~\ref{ip-intersection} therefore gives
the exact integral intersection with this simultaneous isotypic and
character subspace after extending the coefficient field.
The denominator in \eqref{ip-denominator} is used to make an
integral endomorphism. It is not an extra factor in every test
vector and is not an additional loss in the following valuation bound.

\subsection{Extension of a valuation bound}

For a nonzero prime $\mathfrak q$ of $B$, let
$v_{\mathfrak q}:K^\times\to\mathbb Z$ be the additive valuation
with a uniformizer of value one, and put $v_{\mathfrak q}(0)=+\infty$.
Thus $v_{\mathfrak q}(b)\ge0$ for $b\in B$ and
$v_{\mathfrak q}(x+y)\ge\min(v_{\mathfrak q}(x),v_{\mathfrak q}(y))$.

\begin{proposition}\label{ip-valuation}
Retain the notation of Lemma~\ref{ip-intersection}. Let $U$ be a
finite-dimensional $K$-vector space, let $\mathcal S\subseteq U$,
and let
\[
 \Phi:(U\otimes_{K,\sigma}\mathbb C)\times V\longrightarrow\mathbb C
\]
be complex-bilinear. Fix $k\in\mathbb Z$. Suppose that for every
$s\in\mathcal S$ and $m\in M\cap eV$ there is $a_{s,m}\in K$
with
\[
       \Phi(s,m)=\sigma(a_{s,m}),\qquad
                         v_{\mathfrak q}(a_{s,m})\ge k.
\]
Then for every $x\in B\mathcal S$ and $t\in BM\cap eV$ there is
a unique $a_{x,t}\in K$ with
\[
       \Phi(x,t)=\sigma(a_{x,t}),\qquad
                         v_{\mathfrak q}(a_{x,t})\ge k.
\]
The same assertion holds for the $B_{\mathfrak q}$-spans of both
sets, with the evident localization of the test lattice.
\end{proposition}
\begin{proof}
Write $x=\sum_i b_is_i$ with $b_i\in B$, and use
\eqref{ip-cut} to write $t=\sum_j c_jm_j$ with $c_j\in B$ and
$m_j\in M\cap eV$. Bilinearity gives
\[
       \Phi(x,t)=\sigma\left(\sum_{i,j}b_ic_ja_{s_i,m_j}\right).
\]
Each summand on the right has valuation at least $k$, since
both coefficients have nonnegative valuation. Their sum has
valuation at least $k$ as well. Injectivity of $\sigma$ proves
uniqueness, so the value does not depend on the chosen expansions.
For localization the same proof applies with coefficients in
$B_{\mathfrak q}$, whose valuations are still nonnegative.
\end{proof}

\begin{corollary}\label{ip-scaling}
Suppose a $K$-valued bilinear period pairing takes
$M_{\mathfrak q}\times T_{\mathfrak q}$ into
$\mathfrak q^{-b}B_{\mathfrak q}$, where $b\in\mathbb Z$.
If $x\in\mathfrak q^aM_{\mathfrak q}$ with $a\in\mathbb Z$,
then its pairing against every $t\in T_{\mathfrak q}$ has
valuation at least $a-b$.
\end{corollary}
\begin{proof}
Choose a uniformizer $\pi$ of $B_{\mathfrak q}$ and write
$x=\pi^a y$ with $y\in M_{\mathfrak q}$. The pairing equals
$\pi^a\Phi(y,t)$, so its valuation is at least $a-b$.
This argument uses only a local generator; it makes no assertion
that the global prime ideal is principal.
\end{proof}

These results are the ordinary algebra behind the source's transition
from integral period tests to all vectors in the final coefficient
lattice. In its notation, the projection is the level operator,
its integral denominator is the finite group order, and the cut
condition is \eqref{ip-cut}. The definitions and proofs are in
\path{Divisibility/TiedLemmas.lean} and
\path{GroupLattices/JointServices/ValIntNonneg.lean}, with paths
relative to \path{Formalism/LocalToGlobal}.
The exact declaration and read scopes are recorded in the audit.

What remains is mathematical content, not another scalar-extension
lemma: one must prove the initial valuation bound for the actual
theta periods and prove the asserted containment of their actual
source vectors in the required local lattices. The submitted
construction distinguishes the level prime, the two auxiliary
primes, and the other primes, with separate estimates in each case.
Proposition~\ref{ip-valuation} does not supply any of those estimates.


\section{Integral class pairings and the initial period bound}
\label{cp-section}

The initial period valuation has an elementary proof once the
actual polynomial lattices and their class coordinates have been
identified. The period is a finite class sum of the normalized
bilinear Fock pairing. Its only denominator at this stage is the
order of a finite stabilizer. We give the argument with arbitrary
finite stabilizers, so it also applies to the row classes with
additional level characters. The stronger divisibility of the
iterated columns is a separate assertion.

\subsection{Equivariant functions and integral pairings}

Let $K$ be a number field, $\mathcal O$ its ring of integers, and
let a group $G$ act on a set $\mathcal D$ with finitely many orbits.
Choose representatives $d_c$, indexed by a finite set $C$, and
suppose their stabilizers $\Gamma_c$ are finite. Put
$n_c=|\Gamma_c|\ge1$. Let $V^-$ and $V^+$ be finite-dimensional
$K$-spaces with representations $\rho^-$ and $\rho^+$ of $G$.
Suppose the nondegenerate bilinear pairing
\[
                    B:V^-\times V^+\longrightarrow K
\]
satisfies
\begin{equation}\label{cp-invariance}
 B(\rho^-(g)x,\rho^+(g)y)=B(x,y)\qquad(g\in G).
\end{equation}
For each $d\in\mathcal D$ let $M_d^\pm$ be full
$\mathcal O$-lattices, with
\[
 M_{gd}^\pm=\rho^\pm(g)M_d^\pm,
 \qquad B(M_d^-,M_d^+)\subseteq\mathcal O.
\]
Define the spaces and their integral lattices by
\[
 \begin{split}
 \mathcal A^\pm_K
   &=\{f:\mathcal D\to V^\pm:
                         f(gd)=\rho^\pm(g)f(d)\},\\
 \mathcal A^\pm_{\mathcal O}
   &=\{f\in\mathcal A^\pm_K:f(d)\in M_d^\pm
                                      \text{ for every }d\}.
 \end{split}
\]
Evaluation at the $d_c$ identifies these objects with
\[
 \bigoplus_c(V^\pm)^{\Gamma_c},\qquad
 \bigoplus_c\bigl(M_{d_c}^\pm\cap(V^\pm)^{\Gamma_c}\bigr),
\]
respectively. Indeed, an invariant value at a representative
extends uniquely by equivariance, independently of the element
used to transport that representative. This is the same
identification as in Proposition~\ref{acl-class-evaluation}.
Each displayed integral module is a full lattice in the
corresponding space: it is a submodule of a finite lattice, and
every vector in the invariant space has a nonzero integral
multiple in that lattice after clearing finitely many denominators.

Define
\begin{equation}\label{cp-period}
              P(f,h)=\sum_{c\in C}\frac{B(f(d_c),h(d_c))}{n_c}.
\end{equation}
Changing $d_c$ to $gd_c$ conjugates its stabilizer and preserves
its pairing by \eqref{cp-invariance}. Thus $P$ is independent
of all the representatives.

\begin{proposition}\label{cp-nondegenerate}
The pairing $P$ on $\mathcal A^-_K\times\mathcal A^+_K$
is nondegenerate. If $D$ is a positive integer divisible by
every $n_c$, then
\begin{equation}\label{cp-integrality}
       D\,P(\mathcal A^-_{\mathcal O},
                         \mathcal A^+_{\mathcal O})\subseteq\mathcal O.
\end{equation}
In fact one may take $D=\operatorname{lcm}_{c\in C}n_c$,
with the empty least common multiple defined to be $1$.
\end{proposition}
\begin{proof}
For a nonzero $x\in(V^-)^{\Gamma_c}$, choose $y\in V^+$
with $B(x,y)\ne0$. Its average
\[
                  y_0=\frac1{n_c}\sum_{\gamma\in\Gamma_c}
                                                       \rho^+(\gamma)y
\]
is invariant and satisfies $B(x,y_0)=B(x,y)$, by
\eqref{cp-invariance}. Thus the restriction to the invariant
spaces is nondegenerate on the left. The same argument on
the other side proves nondegeneracy on the right.
The direct sum of these pairings, multiplied by the nonzero
scalars $1/n_c$, remains nondegenerate. If $C$ is empty,
both spaces are zero and this statement has the same meaning.

For integral $f,h$, every $B(f(d_c),h(d_c))$ belongs to
$\mathcal O$, and $D/n_c$ is an integer. Multiplying
\eqref{cp-period} by $D$ proves \eqref{cp-integrality}.
The averaging used to prove nondegeneracy is not used to
prove this integral assertion.
\end{proof}

For a prime $w$ of $\mathcal O$, normalize $v_w$ by
$v_w(w)=1$ and put $v_w(0)=+\infty$. If $p$ is the
rational prime below $w$ and $e_w=v_w(p)$, then
\begin{equation}\label{cp-nat-valuation}
                        v_w(D)=e_w v_p(D).
\end{equation}
Indeed write $D=p^a b$ with $p\nmid b$. Bezout's identity
for $p,b$ shows that $b$ is a unit in $\mathcal O_w$.

\begin{corollary}\label{cp-local}
For integral $f,h$,
\[
                 v_w(P(f,h))\ge -e_wv_p(D).
\]
The same assertion holds if the first integral lattice is
replaced by its localization at $w$, and the second is
localized as well.
More generally, if $\mathfrak a$ is a nonzero fractional
ideal and $f\in\mathfrak a(\mathcal A^-_{\mathcal O})_w$,
$h\in(\mathcal A^+_{\mathcal O})_w$, then
\begin{equation}\label{cp-shift}
                  v_w(P(f,h))\ge v_w(\mathfrak a)-v_w(D).
\end{equation}
\end{corollary}
\begin{proof}
Localizing \eqref{cp-integrality} gives
$DP(f,h)\in\mathcal O_w$. This proves the first statements
and \eqref{cp-nat-valuation} gives their normalization.
The fractional ideal $\mathfrak a\mathcal O_w$ is principal;
write $f=a f_0$ with $v_w(a)=v_w(\mathfrak a)$ and
$f_0\in(\mathcal A^-_{\mathcal O})_w$. Bilinearity gives
\eqref{cp-shift}. Zero values cause no exception under the
chosen valuation convention.
\end{proof}

\subsection{The actual normalized polynomial pairing}

We now specify the preceding data for the row construction.
Use the dominant integral weights, determinant powers,
normalized bilinear Fock pairing $B$, and polynomial
representation $\rho_\lambda$ of Section~\ref{coef:section}.
For a tuple of full lattices $M=(M_\tau)_\tau$, put
\[
                M^+=N(M),\qquad M^-=D(M^\vee),
\]
in the notation of \eqref{coef:transport}. The integral
duality already proved in \eqref{coef:transport-duality} is
the exact equality
\begin{equation}\label{cp-polynomial-dual}
                     M^-=(M^+)^{\vee,B}.
\end{equation}
In particular $B(M^-,M^+)\subseteq\mathcal O$.
This uses the actual factorial normalization of the Fock
pairing, all determinant characters, and the unsaturated
Weyl lattice. It is valid at every coefficient prime,
including ramified primes and primes above $2$.

If the base change is $A=(A_\tau)_\tau$, the two coefficient
actions are $\rho_\lambda(A^{-t})$ and
$\rho_\lambda(A)$. Their pairing is invariant because
\[
 B(\rho_\lambda(A^{-t})x,\rho_\lambda(A)y)=B(x,y),
\]
by \eqref{coef:contravariant}. Multiplying the first action
by a finite-order level character and the second by its
inverse preserves this identity. Character values are
algebraic integral units, so they preserve the corresponding
integral lattices as well.

For the actual row objects the level-prime kernel condition
is absent: the underlying lattice is a unimodular Hermitian
lattice $N=N^\#$. The polynomial point lattice is the Weyl
lattice transported by the contragredient action on $N$;
the row lattice is the dual Weyl lattice transported
covariantly on $N$. Thus these are precisely
$D(M^\vee)$ and $N(M)$. The additional local bases in the
point description change neither classes nor stabilizers,
by Corollary~\ref{acl-row-classes}. The two level characters
are inverse, as in \eqref{acl-row-character}. Hence all
the hypotheses of Proposition~\ref{cp-nondegenerate} apply
to these actual row and point lattices.

For an embedding $\sigma:K\hookrightarrow\mathbb C$, let
$\operatorname{ev}(x_c)$ be the frame function associated
to a polynomial $x_c$ as in \eqref{acl-evaluation}.
Recovering its polynomial in the identity frame gives
exactly $\sigma(x_c)$: the determinant factor is removed
with the same exponent with which it was inserted.
The complex period in the construction is therefore
\begin{equation}\label{cp-actual-period}
 \operatorname{Per}(\operatorname{ev}(f),\operatorname{ev}(h))
      =\sum_c\frac{B_{\mathbb C}(\sigma(f(d_c)),
                                      \sigma(h(d_c)))}{n_c}
      =\sigma(P(f,h)).
\end{equation}
The last equality can be checked on monomials: the Fock
pairing is a finite sum of products of coefficients and
integer factorials, divided by the fixed nonzero rational
normalizing constant. An embedding commutes with all these
operations. No analytic comparison or unproved algebraicity
assertion is needed for this finite class sum.

Here the symbol $\operatorname{Per}$ denotes the final
polynomial period pairing. The iterated theta column is
substituted into its first argument later. Establishing
the arithmetic size of that column is a different step.

\subsection{Class restrictions and the stabilizer denominator}

For $c\in C$, let $e_c$ be restriction to the orbit of $d_c$,
followed by extension by zero. These operators preserve the
equivariant spaces, their integral lattices, and their
complexifications. They satisfy
\[
                 \sum_c e_c=1,\qquad e_ce_{c'}=0\quad(c\ne c').
\]
Give the complex point space its positive Hermitian form
\[
                   \langle x,y\rangle
                          =\sum_c n_c^{-1}\langle x_c,y_c\rangle_c,
\]
where the local forms are the normalized Fock forms. It is
conjugate-linear in its first variable. The $e_c$ are
orthogonal projections for this form, directly from the sum.

Fix a row vector $h$ and let $J(h)$ represent its period
functional, so
\[
       \langle J(h),x\rangle=\operatorname{Per}(x,\operatorname{ev}(h)).
\]
Such a vector exists uniquely because the Hermitian form is
positive definite on a finite-dimensional space. In coordinates
with Gram matrix $H$, solving this equation amounts to solving
the invertible linear system for $\overline{J(h)}^{\,t}H$.
In particular $J$ is conjugate-linear in $h$ under this
convention; no assertion of bilinearity of the Hermitian form
is being made.

\begin{proposition}\label{cp-class-contribution}
For integral $f,h$ the normalized class contribution
\[
                p_c(f,h)=n_c\langle e_cJ(h),e_c\operatorname{ev}(f)\rangle
\]
is $\sigma(B(f(d_c),h(d_c)))$, hence belongs to
$\sigma(\mathcal O)$. Moreover
\[
              \operatorname{Per}(\operatorname{ev}(f),\operatorname{ev}(h))
                           =\sum_c n_c^{-1}p_c(f,h).
\]
\end{proposition}
\begin{proof}
Orthogonality gives
$\langle e_cJ(h),e_cx\rangle=\langle J(h),e_cx\rangle$.
Apply \eqref{cp-actual-period} with $e_cf$ in place of $f$.
All terms other than $c$ vanish, leaving
$n_c^{-1}\sigma(B(f(d_c),h(d_c)))$. Multiplication by $n_c$
proves the first assertion and \eqref{cp-polynomial-dual}
proves its integrality. Summing the restrictions proves
the last identity.
\end{proof}

In the source, the row stabilizer embeds into the stabilizer
of its image in the last column class. This accounts for the
chosen common denominator. More generally, suppose a map
of objects is equivariant for an injective homomorphism of
acting groups. Its restriction to the stabilizer of an object
is injective into the stabilizer of the image, because
equivariance sends a fixing element to a fixing element.
Both groups here are finite, so Lagrange's theorem gives
divisibility of their orders. Changing the representative
of the image merely conjugates this injection.
Consequently the source's positive integer $D_\Gamma$,
divisible by the last-column stabilizer orders, is also
divisible by each $n_c$. Corollary~\ref{cp-local} gives
exactly
\begin{equation}\label{cp-source-bound}
 v_w\!\left(\sigma^{-1}\operatorname{Per}
                (\operatorname{ev}(f),\operatorname{ev}(h))\right)
                      \ge -e_wv_p(D_\Gamma).
\end{equation}
The inverse in this notation is applied only to a value
already proved to lie in $\sigma(K)$.

The same proof applies for every used finite-order twist.
If the row lattice is integral at all primes, it gives
\eqref{cp-source-bound} at all primes. If integrality has
only been supplied away from a prescribed exceptional set,
the conclusion holds away from that same set. No exceptional
prime is silently removed.

\subsection{Passage to columns and the remaining estimate}

Let $L$ be a lattice of column vectors and let
$T:L\otimes K\to\mathcal A^-_K$ be the rational map
which forms the iterated column and then its point value.
At a coefficient prime $w$, if one proves
\[
                  T(L_w)\subseteq\mathfrak a_w
                                      (\mathcal A^-_{\mathcal O})_w,
\]
then \eqref{cp-shift} immediately gives
\begin{equation}\label{cp-column-bound}
 v_w(P(Tx,h))\ge v_w(\mathfrak a_w)-v_w(D_\Gamma)
       \quad(x\in L_w,\ h\in\mathcal A^+_{\mathcal O}).
\end{equation}
This deduction retains the ramification factor in
\eqref{cp-nat-valuation}. The extension of the integral
cut and of the test-vector range is already proved in
Section~\ref{ip-section}. What remains here is the actual
containment for $T$, including the special estimates at
the level prime and the two type primes. The initial
period bound \eqref{cp-source-bound} alone does not give
that stronger containment.

For completeness, a finite set of residue-prime shifts has
an exact global interpretation. If $\mathfrak a=\prod_w w^{-b_w}$,
with integral $b_w$ of finite support, then a $K$-valued
bilinear pairing is integral on $\mathfrak a L\times M$
if and only if
\[
                     v_w(P(L,M))\ge b_w\quad\text{for every }w.
\]
Localize: $\mathfrak a_w$ has valuation $-b_w$, so the
condition is precisely that its product with every pairing
value have nonnegative valuation. An element of $K$ is
integral if and only if it is integral at every finite prime.
This proves both directions, including negative $b_w$.

Two examples show why the qualifications above are necessary.
Let a group of order two act trivially on a single object
and on $V^-=V^+=\mathbb Q$, with $B(x,y)=xy$ and both
lattices $\mathbb Z$. Then $P(1,1)=1/2$. Thus the
stabilizer denominator cannot be deleted in general, even
though the class contribution $p_c(1,1)$ is integral.
Also, over $K=\mathbb Q(\zeta_3)$, let an element act as
$\zeta_3$ on both one-dimensional coefficient spaces.
Then $B(\zeta_3x,\zeta_3y)=\zeta_3^2B(x,y)$, so a
class sum is not independent of representatives with those
two actions. Inverse characters restore invariance. Neither
example satisfies the corresponding qualified assertion
used in the submitted construction.


\section{The integral lowering map for every auxiliary character}
\label{atl-section}

We extend the class comparison to the nontrivial auxiliary
characters.  The relevant operation increases the rank by one,
then restricts a coefficient representation to the smaller rank.
Its integral normalization involves the conjugate of the selected
prime generator.  We give the local groups, the extension of
decorated lattices, the transport between class representatives,
and the descent of the resulting complex-valued functions.
The polynomial branching theorem is the separate result
Theorem~\ref{ib-standard-integral}; no automorphic decomposition
or period identity is needed for the arguments in this section.
The definitions compared here are those of the actual final
step for every auxiliary character in the source
\cite{alpo-uniform-source}.

\subsection{The selected primes and the character at the level}

Let $E/F$ be a CM quadratic extension, write
$\mathcal O=\mathcal O_E$, and write $x\mapsto\bar x$ for
its nontrivial automorphism.  A set $S$ of rational primes
contains the two type primes and all the primes $T$ used to
increase the rank.  For one such $T$ fix $\pi\in\mathcal O$
with
\begin{equation}\label{atl-prime-generator}
 \pi\bar\pi=T,\qquad
 (\pi)=\prod_{P\in\mathcal T}P,\qquad
 (\pi)+(\bar\pi)=\mathcal O.
\end{equation}
The primes in $\mathcal T$ are distinct.  The type primes
form another finite set $\mathcal G$ of primes of $E$, all
coprime to $T$.

The ideal equality in \eqref{atl-prime-generator} follows
from the actual prime data, rather than requiring a new
selection.  Namely, write $[E:\mathbb Q]=2f$.  The construction
has exactly $f$ distinct prime divisors of $\pi$, each of norm
$T$.  Conjugation preserves absolute ideal norms, so
$N((\pi))^2=N((T))=T^{2f}$ and $N((\pi))=T^f$.
The product of those $f$ prime divisors divides $(\pi)$
and has the same norm, hence equals $(\pi)$.  The coprimality
with its conjugate is also part of the selected prime data.

An auxiliary character determines a finite set $\mathcal Q$
of primes of $E$.  Each $Q\in\mathcal Q$ is the chosen prime
above a prime $v$ of $F$ such that
\begin{equation}\label{atl-good-primes}
 v\mathcal O=Q\bar Q,\quad Q\ne\bar Q,\quad
 N_{F/\mathbb Q}v=p\text{ is prime},\quad
 p\notin v^2,\quad p\notin S.
\end{equation}
The choice of one prime over each $v$ is fixed before the
auxiliary character is varied.  In particular no two selected
primes are conjugate.  The sets $\mathcal G$, $\mathcal T$,
and $\mathcal Q$ are disjoint, and $\mathcal Q$ is disjoint
from the conjugates of the other two sets.

\begin{lemma}\label{atl-prime-units}
For every $Q$ satisfying \eqref{atl-good-primes}, its residue
field has cardinality $p$, and $v_Q(p)=1$.  Both $\pi$ and
$\bar\pi$ are units in $\mathcal O_{(Q)}$.
These conclusions remain valid in residue characteristic two
whenever the displayed hypotheses hold.
\end{lemma}
\begin{proof}
The factorization of $v\mathcal O$ has ramification index one
at both its primes.  Since $[E:F]=2$, the equality
$\sum_{Q\mid v}e(Q/v)f(Q/v)=2$ gives residue degree one
at each prime.  Thus $NQ=Nv=p$.  Since $p$ belongs to $v$
and does not belong to $v^2$, its $v$-adic valuation is one;
extension to $Q$ leaves that valuation unchanged.

The rational prime $T$ differs from $p$, so $T$ is a unit
at $Q$.  Its two integral factors $\pi$ and $\bar\pi$
must then both have valuation zero.  This proves the last
assertion as well, without using the residue characteristic.
\end{proof}

Fix integers $m\geq2$, $r,t_0\geq0$, and put
$n=r+mt_0$.  At $Q\in\mathcal Q$, let
$\kappa_Q=\mathcal O/Q$ and let
$\xi_Q:\kappa_Q^\times\to K^\times$ be the residue
character, where the number field $K$ contains its values.
The lower-rank residue group and character are
\begin{equation}\label{atl-row-level}
 H_Q=\left\{
 \begin{pmatrix}A&B\\0&D\end{pmatrix}:
 A\in\operatorname{GL}_r(\kappa_Q),\quad
 D\in\operatorname{GL}_{mt_0}(\kappa_Q)\right\},
 \qquad
 \chi_Q\begin{pmatrix}A&B\\0&D\end{pmatrix}
                         =\xi_Q(\det D)^{-1}.
\end{equation}
An empty diagonal block has determinant one.  At a type prime
$P\in\mathcal G$ we keep its specified finite subgroup $H_P$
and character $\chi_P$.  The character of the action on
decorated objects is the product of these characters.  By
\eqref{acl-character}, the corresponding right character on
adelic functions is its inverse.  Thus the right character
at $Q$ is $\xi_Q(\det D)$, whereas the character in
\eqref{atl-row-level} is its reciprocal.

Here is the exact higher-rank local group.  Put
$A_Q=\mathcal O_{(Q)}$, let $\widetilde H_Q$ be the inverse
image of $H_Q$ in $\operatorname{GL}_n(A_Q)$, and set
\[
 K_{n+1}(C)=\ker\bigl(\operatorname{GL}_{n+1}(A_Q)
                 \longrightarrow\operatorname{GL}_{n+1}(A_Q/p^CA_Q)\bigr)
                    \qquad(C\geq1).
\]
Lemma~\ref{atl-prime-units} makes this the usual principal
congruence subgroup of depth $C$.

\begin{lemma}\label{atl-congruence-group}
The set
\begin{equation}\label{atl-top-group}
 \widetilde H_Q^{+}(C)
       =\operatorname{diag}(\widetilde H_Q,1)K_{n+1}(C)
\end{equation}
is a subgroup of finite index in $\operatorname{GL}_{n+1}(A_Q)$.
An element $k$ belongs to it if and only if there exists
$h\in\widetilde H_Q$ with
\[
                      k\equiv\operatorname{diag}(h,1)\pmod {p^C}.
\]
The formula
\begin{equation}\label{atl-top-character}
 \chi_Q^{+}(k)=\chi_Q(\bar h)
       \quad\text{if }k=\operatorname{diag}(h,1)u,
                         \quad u\in K_{n+1}(C),
\end{equation}
is a well-defined character.  It is trivial on $K_{n+1}(C)$,
and its restriction to $\operatorname{diag}(\widetilde H_Q,1)$
is exactly the lower-rank character.  Furthermore
\[
 \operatorname{diag}(h,1)\in\widetilde H_Q^{+}(C)
                       \quad\Longleftrightarrow\quad
 h\in\widetilde H_Q
 \qquad(h\in\operatorname{GL}_n(A_Q)).
\]
\end{lemma}
\begin{proof}
The principal congruence subgroup is normal, so its product
with a subgroup is a subgroup.  It has finite index because
the quotient ring $A_Q/p^CA_Q$ is finite.  Multiplication
by an integral invertible matrix preserves $p^CA_Q^{n+1}$,
so the membership criterion is equivalent to the displayed
factorization.

Reduction modulo $Q$ of every element of the product is
$\operatorname{diag}(\bar h,1)$ with $\bar h\in H_Q$.
It determines $\bar h$ uniquely.  Pulling $\chi_Q$ back
along this homomorphism proves well-definedness and all
the character assertions.  The same reduction proves the
forward implication in the last equivalence; the reverse
implication uses $u=1$.
\end{proof}

For the actual choice of depth, if $M\geq0$ is the local
constant used for this auxiliary character, then
\begin{equation}\label{atl-depth}
                  C=\max(4m+1,2m+2M).
\end{equation}
Indeed the preceding depth is $\max(4m,2m+2M-1)$, and
adding one gives \eqref{atl-depth}; $m\geq2$ ensures that
the subtraction in the second expression does not cause
an exceptional case.  The group in \eqref{atl-top-group}
is not replaced by its larger inverse image under reduction
modulo $Q$.  Its character is, however, already determined
by that reduction.  At the type primes the higher-rank
residue group is simply $\operatorname{diag}(H_P,1)$.

\subsection{Extension by the selected line}

Use the standard positive Hermitian form on $E^n$, with
the convention $h(v,w)=\bar v^{t}w$.  Let $N\subset E^n$
be a full self-dual $\mathcal O$-lattice.  A lower-rank
decorated object consists of $N$, its residue frames
\[
 b_P:N/PN\xrightarrow{\sim}\kappa_P^n
                       \quad(P\in\mathcal G),
\]
and at every $Q\in\mathcal Q$ a basis matrix
$F_Q\in\operatorname{GL}_n(E)$ with
$F_QA_Q^n=N_Q$.  Replacing all local bases by their
principal-congruence orbits gives the equivalent residue
description of Proposition~\ref{acl-row-classes}.

Extend this object to rank $n+1$ by
\begin{align}
 N^+&=N\oplus\mathcal O,\notag\\
 \ell_P(v,s)&=\bar\pi s\bmod P
                    &&(P\in\mathcal T),\label{atl-extended-object}\\
 b_P^+(v,s)&=(b_P(v),s/\pi\bmod P)
                    &&(P\in\mathcal G),\notag\\
 F_Q^+&=\operatorname{diag}(F_Q,\pi)
                    &&(Q\in\mathcal Q).\notag
\end{align}
Here $\ell_P$ is an $\mathcal O$-linear map to $\kappa_P$.
All divisions in the frames are by local units: at the
type primes this follows from $T$ being different from
the type primes, and at $Q$ it follows from
Lemma~\ref{atl-prime-units}.

\begin{proposition}\label{atl-extension}
The formulas \eqref{atl-extended-object} define a decorated
unimodular lattice.  Its lattice cut out by the functionals,
and its Hermitian dual, are respectively
\begin{equation}\label{atl-kernel-dual}
 L^+=N\oplus\pi\mathcal O,
               \qquad (L^+)^\#=N\oplus\bar\pi^{-1}\mathcal O.
\end{equation}
The construction is injective on decorated objects and
equivariant under
\[
 \gamma\longmapsto\operatorname{diag}(\gamma,1),\qquad
 h_P\longmapsto\operatorname{diag}(h_P,1),\qquad
 k_Q\longmapsto\operatorname{diag}(k_Q,1).
\]
These homomorphisms are injective, preserve the level
characters exactly, and send every lower-rank stabilizer
into the stabilizer of the extended object.  In the
parameter used for branching one has
\begin{equation}\label{atl-tz}
                            t=\bar\pi,\qquad \bar t t=T.
\end{equation}
\end{proposition}
\begin{proof}
The direct sum of two self-dual lattices is self-dual.
At $P\mid\pi$, the element $\bar\pi$ is a unit by
\eqref{atl-prime-generator}, so $\ell_P$ is nonzero.
Its kernel imposes precisely $s\in P$.  Intersecting
these conditions gives $s\in\prod_{P\in\mathcal T}P=(\pi)$.
This proves the first formula in \eqref{atl-kernel-dual}.
The pairing of $(0,\pi a)$ with $(v,s)$ is
$\bar\pi\bar a s$, so the last coordinate of the dual
is $\bar\pi^{-1}\mathcal O$; its first coordinates
form $N^\#=N$.

The displayed frames are invertible because their last
diagonal entries are local units.  The original lattice,
residue frames, and local bases are recovered by taking
the first $n$ coordinates, proving injectivity on objects.

For the residue-frame action
$b_P\mapsto h_Pb_P\gamma^{-1}$, its extension is
$b_P^+\mapsto\operatorname{diag}(h_P,1)b_P^+
\operatorname{diag}(\gamma,1)^{-1}$.
For a local basis the action is $F_Q\mapsto\gamma F_Qk_Q^{-1}$;
therefore
\[
 \operatorname{diag}(\gamma,1)\operatorname{diag}(F_Q,\pi)
       \operatorname{diag}(k_Q,1)^{-1}
                    =\operatorname{diag}(\gamma F_Qk_Q^{-1},\pi).
\]
The functionals at $\mathcal T$ are unchanged by the
unitary embedding.  These identities prove equivariance.
Injectivity of the group embeddings is immediate from
their first blocks.  Preservation of characters is
Lemma~\ref{atl-congruence-group} at $Q$, and the definition
of the diagonal residue character at the type primes.
Equivariance then proves the stabilizer assertion.
Finally \eqref{atl-tz} is exactly the last coordinate
of the dual lattice, together with $\pi\bar\pi=T$.
\end{proof}

This proposition does not assert that the induced map of
class sets is injective or surjective.  Both properties
are unnecessary.  At any class which meets the image one
may choose an extended object as representative; all
other classes keep arbitrary representatives.  There
is consequently no need to assume that every higher-rank
object comes from the smaller rank.

Both class sets at these levels are finite.  Indeed, the
finiteness theorem for positive definite unimodular lattices
gives finitely many underlying lattices up to rational
unitary isometry.  For each one, every residue functional
and residue frame has only finitely many possibilities.
The set of bases at $Q$ is a torsor under the integral
linear group.  Its quotient by the chosen level subgroup
is finite by Lemma~\ref{atl-congruence-group}, and likewise
for the lower-rank inverse image of $H_Q$.  Taking the
product over the finite set of primes leaves finitely
many possibilities.  This argument also applies to
higher-rank classes outside the image of the extension.

\subsection{Simultaneous bases and the scaled branching map}

Retain the coefficient field $K$, the finite family of
embeddings $\tau:E\hookrightarrow K$, and the representations
$R_\lambda$ and integral coefficient lattices of
\eqref{acl-contragredient}.  For a full lattice $M\subset E^d$
abbreviate its contragredient coefficient lattice by
$\Delta_\lambda(M)$.  It is the integral coefficient lattice used in the
class construction; all determinant twists are included.

\begin{lemma}\label{atl-semilocal-basis}
Let $\ell$ be a rational prime and $M\subset E^d$ a full
$\mathcal O$-lattice.  There exists
$B\in\operatorname{GL}_d(E)$ whose columns are a basis of
$M_P$ over $\mathcal O_{(P)}$ simultaneously at every
$P\mid\ell$.  If $w\mid\ell$ is a coefficient prime of $K$,
then
\begin{equation}\label{atl-local-coefficient}
 v\in\Delta_\lambda(M)_w
 \quad\Longleftrightarrow\quad
 R_\lambda(B^{-1})v\in
                 \Delta_\lambda^{\rm std}(\mathcal O_{K,w}).
\end{equation}
\end{lemma}
\begin{proof}
For every $P\mid\ell$, the module $M_P$ is free of rank
$d$ over the discrete valuation ring $\mathcal O_{(P)}$.
Choose a basis of $M/PM$.  The Chinese remainder theorem
for the pairwise comaximal submodules $PM$ lets us choose
$d$ elements of $M$ with these prescribed residue bases
at all the finitely many primes above $\ell$.
At each such prime their span has quotient $C$ with
$C/PC=0$.  The quotient is finitely generated; Nakayama's
lemma gives $C=0$.  Since $M_P$ has rank $d$, the selected
elements are a basis.  They are linearly independent over
$E$, and hence form the columns of an invertible matrix
$B$.  For $d=0$ take the unique empty matrix.

For each embedding $\tau$ the prime $\tau^{-1}(w)$ is
above $\ell$, so this same matrix supplies every required
local basis.  The defining change-of-basis rule from
Section~\ref{coef:section} gives
\[
 \Delta_\lambda(M)_w
     =R_\lambda(B)\Delta_\lambda^{\rm std}(\mathcal O_{K,w}).
\]
Applying $R_\lambda(B^{-1})$ proves the equivalence.
\end{proof}

Let $\lambda_\tau$ and $\sigma_\tau$ be dominant integer
weights in ranks $n+1$ and $n$, respectively.  We require
the actual interlacing and determinant-shift conditions
\begin{equation}\label{atl-weights}
 \lambda_{\tau,i}\geq\sigma_{\tau,i}
                       \geq\lambda_{\tau,i+1}\ (1\leq i\leq n),
               \qquad \sigma_{\tau,n}=\lambda_{\tau,n+1}.
\end{equation}
The ranks in the application are positive.  The standard
branching theorem gives a $K$-linear map
\[
 B_1:W_\lambda(K)\longrightarrow W_\sigma(K)
\]
with the specified polynomial normalization, such that
\begin{align}
 B_1R_\lambda(\operatorname{diag}(\gamma,1))
                &=R_\sigma(\gamma)B_1,
                    &&\gamma\in\operatorname{GL}_n(E),
                         \label{atl-standard-covariance}\\
 B_1\Delta_\lambda^{\rm std}(\mathcal O_{K,w})
                &\subseteq\Delta_\sigma^{\rm std}(\mathcal O_{K,w})
                    &&\text{for every }w.
                         \label{atl-standard-integral}
\end{align}
These are precisely the integral assertion of
Theorem~\ref{ib-standard-integral} and the diagonal
covariance of Corollary~\ref{ib-diagonal-covariance}, applied
to the contragredient realization.  Define
\begin{equation}\label{atl-scaled-map}
                 B_t=B_1R_\lambda(\operatorname{diag}(I_n,t))
                         \qquad(t\in E^\times).
\end{equation}
To check the normalization explicitly, set
\[
 b_{\tau,i}=\lambda_{\tau,i}-\sigma_{\tau,i}\geq0,
 \qquad a_\tau=\lambda_{\tau,n+1},\qquad
 N_\tau=\sum_{i=1}^{n+1}\lambda_{\tau,i}
                              -\sum_{i=1}^n\sigma_{\tau,i}.
\]
The standard polynomial formula first retains exactly the
monomials with exponent $b_{\tau,i}$ at the last-row variable
$X_{\tau,n+1,i}$, and exponent zero at every last-column
variable; it then applies a fixed projection and scalar.
Under $R_\lambda(\operatorname{diag}(I_n,t))$, each retained
monomial is multiplied by
$\prod_\tau\tau(t)^{-a_\tau-\sum_i b_{\tau,i}}$.
Since $a_\tau+\sum_i b_{\tau,i}=N_\tau$, this proves the
exact formula
\begin{equation}\label{atl-scaled-factor}
                      B_t=\left(\prod_\tau\tau(t)^{-N_\tau}\right)B_1.
\end{equation}
The exponents $N_\tau$ are integers and need not be
nonnegative.  This is the scaled map in the construction,
including the determinant contribution $a_\tau$.

\begin{proposition}\label{atl-scaled-integrality}
For $B'\in\operatorname{GL}_n(E)$ put
$B=\operatorname{diag}(B',t^{-1})$.  Then
\begin{equation}\label{atl-coordinate-identity}
           R_\sigma(B'^{-1})B_t=B_1R_\lambda(B^{-1}).
\end{equation}
The map $B_t$ intertwines the two representations under
$\gamma\mapsto\operatorname{diag}(\gamma,1)$.
For every full lattice $M\subset E^n$ it satisfies
\begin{equation}\label{atl-branch-integral}
 B_t\Delta_\lambda(M\oplus t^{-1}\mathcal O)
                           \subseteq\Delta_\sigma(M).
\end{equation}
All these statements hold at every finite coefficient
prime, including primes above $T$, primes above the
auxiliary set $\mathcal Q$, and primes above two.
\end{proposition}
\begin{proof}
By \eqref{atl-standard-covariance}, the left side of
\eqref{atl-coordinate-identity} is
\[
 B_1R_\lambda(\operatorname{diag}(B'^{-1},1))
       R_\lambda(\operatorname{diag}(I_n,t))
                 =B_1R_\lambda(\operatorname{diag}(B'^{-1},t)).
\]
This is the right side.  The same covariance, and the
commutation of the last-coordinate scaling with
$\operatorname{diag}(\gamma,1)$, prove the intertwining
assertion.

Fix a coefficient prime $w\mid\ell$.  Choose $B'$ as
in Lemma~\ref{atl-semilocal-basis} for $M$.  The matrix
$B=\operatorname{diag}(B',t^{-1})$ is a simultaneous
basis of $M\oplus t^{-1}\mathcal O$.  By
\eqref{atl-local-coefficient}, a vector in the localized
source lattice is carried by $R_\lambda(B^{-1})$ into
the standard integral lattice.  Apply
\eqref{atl-standard-integral} and
\eqref{atl-coordinate-identity}, and use the same local
criterion for the target.  This proves the localized
inclusion at $w$.  Intersecting over all $w$ proves
\eqref{atl-branch-integral}.  The argument never requires
that $t$ be a unit at $w$, or that the residue
characteristic exceed the degree of a polynomial.
\end{proof}

Adding an integer $a_\tau$ to every entry of a weight
does not change its normalized partition
$\lambda_{\tau,i}-\lambda_{\tau,n+1}$, and therefore
does not change its standard polynomial lattice as a
set of polynomials.  It does change the determinant
factor in $R_\lambda$.  In \eqref{atl-coordinate-identity}
the representations are always those of the actual
shifted weights; only the standard polynomial lattices
may be identified by equality of normalized partitions.

\subsection{Transport of class values and global descent}

Let $\mathcal D$ be the lower-rank decorated objects and
$\mathcal D^+$ the higher-rank objects, with their groups
$G$ and $G^+$ acting as above.  Write
$j:\mathcal D\to\mathcal D^+$ for the extension and
$\iota:G\to G^+$ for the group homomorphism.  Let
$\chi,\chi^+$ be the action characters.  We have proved
\begin{equation}\label{atl-action-compatibility}
 j(gd)=\iota(g)j(d),\qquad
 \chi^+(\iota(g))=\chi(g),\qquad
 \gamma(\iota(g))=\operatorname{diag}(\gamma(g),1).
\end{equation}
For a higher-rank equivariant coefficient function $F$,
define
\begin{equation}\label{atl-lowering}
                         (\mathcal L_tF)(d)=B_tF(j(d)).
\end{equation}

\begin{theorem}\label{atl-global-lowering}
The map $\mathcal L_t$ carries higher-rank equivariant
$K$-valued forms to lower-rank equivariant $K$-valued
forms.  With $t=\bar\pi$ and the actual coefficient
lattices, it carries the integral class lattice into
the integral lower-rank class lattice, and the same
holds after localization at every coefficient prime.
It agrees after complexification with evaluation on
the extended objects followed by the normalized
polynomial branching map.

More explicitly, choose representatives $d_c$ of the
lower-rank classes and $e_b$ of the higher-rank classes.
Let $b(c)$ be the class of $j(d_c)$, and choose
$a_c\in G^+$ with $a_ce_{b(c)}=j(d_c)$.  If $v_b=F(e_b)$,
then the value at $d_c$ is exactly
\begin{equation}\label{atl-representative-formula}
             (\mathcal L_tF)(d_c)
                 =B_t\bigl(\chi^+(a_c)
                          R_\lambda(\gamma(a_c))v_{b(c)}\bigr).
\end{equation}
It is independent of the choices $a_c$.
\end{theorem}
\begin{proof}
Equations \eqref{atl-action-compatibility} and the
intertwining assertion of
Proposition~\ref{atl-scaled-integrality} give
\begin{align*}
 (\mathcal L_tF)(gd)
  &=B_t\chi^+(\iota(g))
               R_\lambda(\operatorname{diag}(\gamma(g),1))F(j(d))\\
  &=\chi(g)R_\sigma(\gamma(g))(\mathcal L_tF)(d).
\end{align*}
Thus the function has the required equivariance.

At $j(d)$ the coefficient lattice is that of
$(L_{j(d)})^\#=N_d\oplus t^{-1}\mathcal O$,
whereas the lower-rank lattice is that of $N_d$.
Their integral compatibility is
\eqref{atl-branch-integral}.  Transport from $e_{b(c)}$
to $j(d_c)$ preserves these exact lattices: the
unitary action carries the underlying lattice to
the required one, its coefficient action is $R_\lambda$,
and $\chi^+(a_c)$ is a root of unity, hence an integral
unit at every coefficient prime.  This proves the
global and localized assertions.

If $a_c'$ is another choice, $a_c^{-1}a_c'$ stabilizes
$e_{b(c)}$.  Its coefficient action fixes $v_{b(c)}$
by Proposition~\ref{acl-class-evaluation}, so the
transported values are equal.  The same argument,
or the equivariance just proved, shows that the
right side of \eqref{atl-representative-formula}
is fixed by the lower-rank stabilizer.  No averaging
over a stabilizer occurs.  Finally all coefficient
operations and \eqref{acl-evaluation} commute with
the chosen embedding $K\hookrightarrow\mathbb C$.
This proves the complex comparison.
\end{proof}

For completeness, the integral coordinates used to
recognize the resulting complex form can also be
characterized without presupposing rationality.

\begin{proposition}\label{atl-complex-descent}
Fix $w\mid\ell$ in $K$ and an embedding
$\jmath:K\hookrightarrow\mathbb C$.  At every lower-rank
class representative $d_c$, choose a simultaneous
basis $B_c'$ of $N_{d_c}$ as in
Lemma~\ref{atl-semilocal-basis}.  Let $z$ be a complex
equivariant form in the coefficient representation.
Extract its unique polynomial coefficient $P_c$ from
$z(d_c)$ by \eqref{acl-evaluation}.  Then $z$ is the
complexification of an element of the localized integral
class lattice if and only if, for every $c$,
\begin{equation}\label{atl-complex-criterion}
 R_\sigma^{\mathbb C}(B_c'^{-1})P_c
      \in\jmath\bigl(\Delta_\sigma^{\rm std}
                                        (\mathcal O_{K,w})\bigr).
\end{equation}
The complex representation is obtained through the
embeddings $\jmath\tau$.
\end{proposition}
\begin{proof}
If the form is already the complexification of a
localized integral one, the criterion is
\eqref{atl-local-coefficient} at every class.
Conversely, the right side of
\eqref{atl-complex-criterion} supplies a unique
$K$-valued polynomial $q_c$ in the standard local
lattice.  The uniqueness follows from injectivity
of the field embedding on polynomial coefficients.
Undoing $R_\sigma(B_c'^{-1})$ gives
$v_c=R_\sigma(B_c')q_c\in\Delta_\sigma(N_{d_c})_w$
whose complex evaluation is $z(d_c)$.

For any element of the stabilizer of $d_c$, the
complex evaluations of $v_c$ and its coefficient
translate are equal, since $z$ is equivariant.
The injectivity proved after \eqref{acl-evaluation}
implies equality over $K$.  Hence $v_c$ is fixed
by the entire stabilizer.  These values determine
a $K$-valued equivariant form by
Proposition~\ref{acl-class-evaluation}.

There are finitely many classes.  For each $c$ choose
$s_c\in\mathcal O_K\setminus w$ with
$s_cv_c\in\Delta_\sigma(N_{d_c})$.  The product
$s=\prod_cs_c$ is still outside $w$ and makes every
class value integral.  Multiplication by $s$ preserves
all stabilizer equalities.  Thus the form belongs to
the localized integral class lattice, and its
complexification is $z$.  This also proves uniqueness.
\end{proof}

To see the exact coordinate formula for lowering,
put $B_c=\operatorname{diag}(B_c',t^{-1})$.  For the
transported value in \eqref{atl-representative-formula}
define
\begin{align*}
 \operatorname{co}_+(F,c)
  &=\jmath\left(R_\lambda(B_c^{-1})
       \chi^+(a_c)R_\lambda(\gamma(a_c))v_{b(c)}\right),\\
 \operatorname{co}_-(z,c)
  &=R_\sigma^{\mathbb C}(B_c'^{-1})P_c.
\end{align*}
Equation \eqref{atl-coordinate-identity} gives
\begin{equation}\label{atl-final-coordinate}
 \operatorname{co}_-(\mathcal L_tF,c)
                   =B_1^{\mathbb C}\operatorname{co}_+(F,c).
\end{equation}
The first coordinate map carries every localized
integral form into the standard integral polynomial
lattice.  The second recognizes precisely the
localized integral forms by
Proposition~\ref{atl-complex-descent}.  This is the
full coordinate comparison needed in the final
lowering step.

\subsection{Uniformity over the required characters and scope}

Fix the finite family of auxiliary characters used
in the construction.  For each member use its actual
finite support $\mathcal Q$ and its own constant $M$
in \eqref{atl-depth}.  The preceding proofs apply
to every member, and to every selected prime $T$,
with exactly the same weights for fixed $T$-index.
The support and depth change the finite class set and
its level subgroup; they do not change the formulas
\eqref{atl-kernel-dual}, \eqref{atl-tz}, or the
standard polynomial map $B_1$.  The simultaneous
local bases may depend on the class, the character,
and $w$.  Their existence is proved after these
parameters are fixed and requires no new prime
selection or conductor estimate.

In particular, no exclusion $w\notin\mathcal Q$ is
needed for the integral lowering map itself.  Any
such exclusion in a later period-valuation theorem
is a hypothesis of that theorem, not a consequence
of this coordinate comparison.  The exact integral
polynomial branching assertion remains the dependency
\eqref{atl-standard-integral}; the preceding proofs
do not infer it merely from irreducibility over $K$.
This section does not prove the transfer identities
for every intermediate nontrivially twisted theta
sum, nor the valuation improvement in the period
pairing.  It does prove the actual final class map,
its characters, its local integral coordinates, and
their global descent.

\subsection{Checks on essential hypotheses}

The following examples explain why several exact
features cannot be omitted.  They do not contradict
the hypotheses of the construction.

First, at a prime of residue characteristic $p$,
take $n=1$, $H_Q=\operatorname{GL}_1(\kappa_Q)$, and
$C=2$.  The matrix $I+pE_{21}$ reduces to the identity,
so it lies in the full inverse image of the diagonal
residue group.  It does not lie in
$\operatorname{diag}(\widetilde H_Q,1)K_2(2)$, because
its lower left entry is not divisible by $p^2$.
Thus replacing the actual group by that inverse
image changes the level, even when the character
is unchanged.

Second, in $E=\mathbb Q(i)$ take $\pi=2+i$.
For the standard Hermitian form on one dimension,
$(\pi\mathbb Z[i])^\#=\bar\pi^{-1}\mathbb Z[i]$.
The proposed replacement $\pi^{-1}\mathbb Z[i]$
is different: $\bar\pi/\pi=(3-4i)/5$ is not an
algebraic integer.  The conjugation in $t=\bar\pi$
therefore affects integral lattices.

Third, transport between class representatives
cannot be omitted from
\eqref{atl-representative-formula}.  Already for a
free transitive action of the group $\{1,s\}$ on
two points, the equivariance rule $F(sd)=-F(d)$
gives opposite values at the two representatives
of the same orbit.  A branching map equal to the
identity would preserve this discrepancy.

Finally the finiteness of the class set is used in
the descent proof.  In an infinite product over
distinct rational primes $p_j\ne\ell$, the vector
$(1/p_j)_j$ belongs componentwise to $\mathbb Z_{(\ell)}$.
There is no integer $s$ prime to $\ell$ which makes
all its components integers.  Hence it does not
belong to $(\prod_j\mathbb Z)\otimes\mathbb Z_{(\ell)}$,
although it belongs to $\prod_j\mathbb Z_{(\ell)}$.
The finite class theorem and the common denominator
in Proposition~\ref{atl-complex-descent} prevent
this failure in the actual construction.


\section{Integral formal series and bounded ordinary theta moments}
\label{ot-section}

This section proves the formal and functional-analytic parts of the
ordinary theta moment argument.  They are useful because bounded
ordinary derivatives alone do not define a bounded measure.  We give
the exact criterion, prove the descent from tensor powers, and prove
the integral functional equation used for the one-dimensional CM
theta function.  The geometric hypotheses needed to apply that
equation are stated explicitly at the end.

Fix a prime $p$.  Let $L$ be a complete nonarchimedean field containing
$\mathbb Q_p$, with its usual absolute value on $\mathbb Q_p$, and
write $\mathcal O=\{x\in L:|x|\le1\}$ and
$\mathfrak m=\{x\in L:|x|<1\}$.  A power series is called
\emph{integral} if all its coefficients belong to $\mathcal O$.
It is called \emph{bounded} if their absolute values have a finite
common bound.  No discreteness of the valuation is assumed.

\subsection{A complete measure criterion}

For $r\ge0$, put $X_r=\mathbb Z_p^r$.  For a multi-index
$n=(n_1,\ldots,n_r)\in\mathbb N^r$, set
\[
 B_n(x)=\prod_{j=1}^r\binom{x_j}{n_j},\qquad
 \binom Xk=\frac{X(X-1)\cdots(X-k+1)}{k!}.
 \tag{OT1}
\]
Empty products are one.  The polynomials $B_n$ take integral values
on $X_r$: this holds at nonnegative integers and then follows on
$\mathbb Z_p$ by continuity and density.
Let $S_jf(x)=f(x+e_j)$ and $\Delta_j=S_j-1$.

\begin{theorem}[Multivariable Mahler expansion]\label{ot-mahler}
Every continuous $f:X_r\to L$ has a unique uniformly convergent
expansion
\[
 f(x)=\sum_{n\in\mathbb N^r}a_n(f)B_n(x),\qquad
 a_n(f)=(\Delta_1^{n_1}\cdots\Delta_r^{n_r}f)(0).
 \tag{OT2}
\]
The coefficients tend to zero outside finite subsets of
$\mathbb N^r$, and
\[
                  \|f\|_\infty=\sup_n|a_n(f)|.
 \tag{OT3}
\]
When $r=0$, these assertions mean the single constant coefficient.
\end{theorem}
\begin{proof}
The ultrametric inequality gives $\|\Delta_j\|\le1$, hence
$|a_n(f)|\le\|f\|_\infty$.  We first prove their convergence
to zero.  A continuous function on the compact space $X_r$ can be
approximated uniformly, to any accuracy, by a function constant on
residue classes modulo $p^h$.  Indeed, uniform continuity gives
such an $h$, and one chooses one value on each of the finitely many
classes.

Let $g$ be such a function and put $N=p^h$.  On these functions
$S_j^N=1$.  In the ring $\mathbb Z[S]/(S^N-1)$ all coefficients
of $(S-1)^N$ are divisible by $p$: reduction modulo $p$ gives
$(S-1)^{p^h}=S^{p^h}-1=0$.  It follows that
\[
 \|\Delta_j^{kN}g\|_\infty\le |p|^k\|g\|_\infty.
 \tag{OT4}
\]
The same estimate holds with any additional differences in any
coordinates.  Thus $a_n(g)$ tends to zero whenever at least one
coordinate of $n$ tends to infinity, uniformly in the others.
Since $|a_n(f)-a_n(g)|\le\|f-g\|_\infty$, the asserted
cofinite convergence follows for $f$.

The series in (OT2) therefore converges uniformly: its tails have
norm at most the largest omitted coefficient, because
$|B_n(x)|\le1$.  Its sum is continuous.  For
$x\in\mathbb N^r$, the series is finite, and the ordinary Newton
identity gives its value as $f(x)$.  In one variable this identity
is the binomial inversion
\[
 f(m)=\sum_{n=0}^m\binom mn
          \sum_{j=0}^n(-1)^{n-j}\binom nj f(j).
\]
The coefficient of $f(j)$ is
$\binom mj\sum_{k=0}^{m-j}(-1)^k\binom{m-j}{k}$, which is
one for $j=m$ and zero otherwise.  Applying this in each coordinate
proves the multivariable identity.  Density of $\mathbb N^r$ in
$X_r$ proves (OT2) everywhere.  It also gives
$\|f\|_\infty\le\sup_n|a_n(f)|$, establishing (OT3).
Finally $\Delta_j\binom{x_j}{k}=\binom{x_j}{k-1}$ implies
$a_n(B_m)=1$ if $n=m$ and zero otherwise.  Applying the bounded
functional $a_n$ to any uniformly convergent expansion proves
uniqueness.
\end{proof}

By an $L$-valued bounded measure on $X_r$ we mean a continuous
$L$-linear functional on $C(X_r,L)$, with its supremum norm.  This
definition avoids any assumption of countable additivity in a
real-valued measure space.

\begin{corollary}[Exact moment criterion]\label{ot-measure}
A family $(c_n)_{n\in\mathbb N^r}$ is of the form
$c_n=\mu(B_n)$ for a bounded measure if and only if it is bounded.
The measure is unique, is given by
\[
             \mu(f)=\sum_n a_n(f)c_n,
             \qquad \|\mu\|=\sup_n|c_n|,
 \tag{OT5}
\]
and has norm at most one exactly when every $c_n$ is integral.
\end{corollary}
\begin{proof}
Boundedness of the $c_n$ and the cofinite convergence of $a_n(f)$
give convergence of (OT5) in the complete field $L$.  The
ultrametric inequality and (OT3) show that it is a linear
functional of norm at most $\sup|c_n|$.  Since
$\|B_n\|_\infty=1$ and $a_m(B_n)=\delta_{mn}$, the reverse
inequality holds.  Conversely a bounded measure has
$|\mu(B_n)|\le\|\mu\|$.  Applying it to (OT2) proves
uniqueness.
\end{proof}

Let $g\in L[[z_1,\ldots,z_r]]$ and write it in exponential
generating-function form
\[
 g(z)=\sum_{\alpha\in\mathbb N^r}
           M_\alpha\frac{z^\alpha}{\alpha!}.
 \tag{OT6}
\]
Here $\alpha!=\prod_j\alpha_j!$; it is nonzero in $L$.
Define the linear functional $\Lambda_g$ on $L[X_1,\ldots,X_r]$
by $\Lambda_g(X^\alpha)=M_\alpha$.

\begin{proposition}\label{ot-logarithm}
The coefficients of the formal series
\[
       g(\log(1+T_1),\ldots,\log(1+T_r))
       =\sum_n c_nT^n
 \tag{OT7}
\]
are exactly $c_n=\Lambda_g(B_n)$.  Consequently $\Lambda_g$
extends to a measure on $X_r$ of norm at most one if and only if
the series (OT7) is integral.
\end{proposition}
\begin{proof}
The formal identity
\[
 \exp(X\log(1+T))=\sum_{n\ge0}\binom Xn T^n
 \tag{OT8}
\]
holds in $L[X][[T]]$.  One proof differentiates in $T$: both sides
have constant term one and satisfy
$(1+T)\partial_T F=XF$, whose coefficients determine the series
uniquely.  Multiply this identity in the $r$ coordinates and apply
$\Lambda_g$ coefficient by coefficient.  Each coefficient uses
only finitely many monomials, so the operation is purely formal
and gives (OT7).  Corollary~\ref{ot-measure} proves the criterion.
Polynomial uniqueness follows also from the fact that the $B_n$
form a basis, since their leading monomials are $X^n/n!$.
\end{proof}

For a complex analytic function $G$ near zero in a finite-dimensional
complex vector space, choose vectors $v_1,\ldots,v_r$, an abstract
field embedding of its Taylor coefficients into $L$, and a period
$\omega\in L$.  Apply the proposition to the transported Taylor
series of $G(\omega\sum_j z_jv_j)$; equivalently set
$M_\alpha=\omega^{|\alpha|}\iota(D_{v_1}^{\alpha_1}
\cdots D_{v_r}^{\alpha_r}G(0))$.
Only algebraic operations on Taylor coefficients are involved; the
field embedding is not asserted to be continuous for the complex
topology.

\subsection{Descent from tensor powers}

\begin{theorem}\label{ot-tensor-root}
Let $G$ be as above, let $\Lambda$ be an additive subgroup of its
vector space, and let $h\ge1$.  Define the function in independent
blocks
\[
             G^{\otimes h}(z_1,\ldots,z_h)
                      =\prod_{j=1}^hG(z_j).
 \tag{OT9}
\]
If, for every finite family of directions in $\Lambda^h$, its
period-scaled Taylor functional is a measure of norm at most $C$,
then the corresponding functional of $G$ in every finite family
of directions in $\Lambda$ is a measure of norm at most $C^{1/h}$.
Here $C\ge0$.  In particular norm at most one descends from the
fourth tensor power, without any assumption on $G(0)$.
\end{theorem}
\begin{proof}
Fix $v_1,\ldots,v_r\in\Lambda$ and let $c_n$ be the coefficients
in (OT7) for these directions.  Use in the tensor product the
$hr$ directions obtained by placing $v_i$ in block $j$ and zero
in the other blocks.  Its logarithm-substituted Taylor series is
the product of $h$ copies of (OT7) in disjoint variable sets.
The coefficient indexed by $(n,\ldots,n)$ is therefore $c_n^h$.
The assumed norm bound and Corollary~\ref{ot-measure} give
$|c_n|^h\le C$.  Thus $|c_n|\le C^{1/h}$, and the same
corollary gives the claimed measure.  For $C=0$ every $c_n$ is
zero; this case is included.  For an empty direction family the
argument is simply the same inequality for the constant term.
\end{proof}

The independent blocks in this theorem matter.  A pointwise power
$G(z)^h$ combines coefficients by convolution and does not provide
the diagonal coefficient $c_n^h$.  The submitted fourth-power
argument uses the tensor product (OT9), so the theorem applies
with precisely its intended meaning.

\subsection{An integral functional equation}

\begin{theorem}\label{ot-functional-equation}
Let $\phi=(\phi_1,\ldots,\phi_r)\in
\mathcal O[[T_1,\ldots,T_r]]^r$ have zero constant term and
linear coefficient matrix in $\operatorname{GL}_r(\mathcal O)$.
Let $m\ge0$ be an integer with $|m|<1$, and let
$A,B\in\mathcal O[[T_1,\ldots,T_r]]$ have constant term one.
There exists a unique $f\in L[[T_1,\ldots,T_r]]$, with constant
term one, satisfying
\[
                     f(\phi(T))A(T)=f(T)^mB(T).
 \tag{OT10}
\]
This solution is integral.  If $m\ge2$ and
$0\ne\rho\in L[[T_1,\ldots,T_r]]$ satisfies
$\rho(\phi)=\rho^m$, then $\rho$ is a constant $c$ with
$c^{m-1}=1$.
\end{theorem}
\begin{proof}
Write the linear part of $\phi$ as $UT$.  Suppose the homogeneous
parts of $f$ below degree $n$ have been constructed.  The degree
$n$ part $f_n$ occurs on the left of (OT10) as $f_n(UT)$ and on
the right as $mf_n(T)$.  All other degree $n$ terms involve only
the already constructed parts.  Thus
\[
                       (S_n-m)f_n=R_n,
 \tag{OT11}
\]
where $S_n:P(T)\mapsto P(UT)$ is an automorphism of the finite
free $\mathcal O$-module of homogeneous polynomials of degree
$n$.  Its inverse substitutes $U^{-1}T$, so both $S_n$ and
$S_n^{-1}$ have integral entries.  The matrix
$1-mS_n^{-1}$ reduces to the identity modulo $\mathfrak m$;
its determinant is a unit of $\mathcal O$, and its inverse is
integral by the adjugate formula.  Hence $S_n-m$ is invertible
over $\mathcal O$.  If the lower parts of $f$ are integral, then
$R_n$ is integral, so (OT11) determines a unique integral $f_n$.
Starting with $f_0=1$ proves existence, integrality and uniqueness
degree by degree.  The same invertibility proves uniqueness among
all $L$-valued formal solutions, even if their coefficients were
not initially assumed bounded.  The case $r=0$ is the constant
identity $1=1$.

For the last assertion, the least degree of a nonzero homogeneous
part of $\rho$ is unchanged by substitution of $\phi$, since
$U$ is invertible.  Equation $\rho(\phi)=\rho^m$ forces this
degree to be zero.  Its constant $c$ is nonzero and satisfies
$c=c^m$.  Then $\rho/c$ satisfies (OT10) with $A=B=1$ and
constant term one.  Uniqueness makes it the constant one.
\end{proof}

In one variable the recurrence is especially transparent: if
$\phi(T)=aT+\cdots$, the coefficient of $f_n$ in (OT11) is
$a^n-m$, whose absolute value is one.  This proves exactly the
integrality step used in the ordinary CM theta construction and
also shows why the unit linear coefficient is essential.

\subsection{Bounded elliptic quotients}

The next elementary observation supplies the bounded-series step
in that construction.  We isolate the geometric inputs so that
integrality of a formal chart is not tacitly assumed.

\begin{lemma}\label{ot-bounded-unit}
Suppose $F,G\in L[[T]]$ are bounded, $FG=1$ and $F(0)=1$.
Then both $F$ and $G$ are integral.
\end{lemma}
\begin{proof}
For $0<r<1$ define $\|F\|_r=\max_n|F_n|r^n$; the maximum
exists because $F$ is bounded.  The same applies to $G$.
This norm is multiplicative.  Indeed, choose the smallest indices
$i,j$ where the two maxima occur.  In the coefficient of
$T^{i+j}$ in $FG$, every other summand has smaller weighted
absolute value: if its first index is less than $i$ it is already
smaller, and otherwise its second index is less than $j$.
The ultrametric inequality then gives equality for this
coefficient, while the reverse inequality for all coefficients
is immediate.  Thus $\|F\|_r\|G\|_r=1$.
Both norms are at least one, from their constant terms, hence
both are one.  Letting $r$ tend to one gives the assertion for
each coefficient.
\end{proof}

\begin{proposition}\label{ot-elliptic-bounded}
Let $E/L$ be a smooth Weierstrass elliptic curve, with origin $O$,
and suppose a formal parameter $T$ at $O$ gives expansions
\[
 x=T^{-2}X(T),\qquad y=T^{-3}Y(T),\qquad
 X,Y\in\mathcal O[[T]],\quad X(0)\in\mathcal O^\times.
 \tag{OT12}
\]
Let $R$ be a rational function on $E$, regular at $O$.
Assume every pole $P$ of $R$ is affine and satisfies $|x(P)|\le1$,
after passage to a complete algebraically closed extension if needed.
Then the expansion of $R$ in $T$ is bounded.
If $R(O)=1$ and the same pole condition holds for $R^{-1}$, then
the expansion of $R$ and its inverse are integral.
\end{proposition}
\begin{proof}
For each pole choose a sufficiently large multiplicity and form a
monic polynomial $D(X)$ whose roots are the finitely many $x(P)$.
Then $D(x)R$ has no affine pole, and $D$ has integral coefficients.
The affine coordinate ring of the smooth Weierstrass curve is
$L[x,y]$ modulo its Weierstrass equation.  A rational function
regular at every affine point belongs to this ring.  To recall
the local justification, smoothness makes each local ring a
discrete valuation ring, so absence of a pole means membership in
that local ring.  If an element of the fraction field belongs to
all those local rings but not to the affine ring $A$, its
denominator ideal $\{a\in A:af\in A\}$ is proper.  A maximal
ideal containing it contradicts membership in the corresponding
localization.  Thus
\[
                        D(x)R=A(x)+yB(x)
 \tag{OT13}
\]
for polynomials $A,B$, since the Weierstrass equation is monic of
degree two in $y$.

Put $d=\deg D$.  From (OT12),
$D(x)=T^{-2d}S(T)$, where $S$ is integral and
$S(0)=X(0)^d$ is a unit.  Its inverse is integral by the usual
coefficient recursion.  Multiplying the right side of (OT13)
by a sufficiently large power $T^k$ produces a bounded series:
it is a finite sum of scalar multiples of products of $X$ and
$Y$.  Therefore the Laurent expansion
$R=T^{2d-k}(T^k(A(x)+yB(x)))S^{-1}$ has bounded coefficients.
Regularity at $O$ removes its negative coefficients, yielding the
claimed power series.  Applying the same argument to $R^{-1}$
and Lemma~\ref{ot-bounded-unit} proves the last assertion.
All arguments may be made after the indicated extension; bounds
and integrality of the original coefficients then descend.
\end{proof}

\subsection{The precise one-dimensional CM consequence}

Here is a sufficient set of ordinary CM data, stated as conventional
mathematics.  Let $k$ be an imaginary quadratic field, let $p$ be
odd and split in $k$, and fix the embedding
$\mathcal O_k\hookrightarrow\mathbb Z_p$ associated with one
of its primes.  Suppose a complex elliptic curve with CM by
$\mathcal O_k$ has been transported to a good integral Weierstrass
model over $L$.  Assume its formal group has an integral
multiplicative coordinate $T$: its parameter at the origin is an
integral series in $T$ with unit linear coefficient, and the CM
endomorphism $\alpha$ acts as
\[
                         [\alpha](T)=(1+T)^{a}-1,
             \qquad a=\iota(\alpha)\in\mathbb Z_p.
 \tag{OT14}
\]
Assume the CM endomorphisms preserve the formal subgroup, and that
this subgroup consists precisely of $O$ and the affine points with
$|x|>1$.  Suppose a normalized entire theta function $\theta$
has $\theta(0)=1$, zeros at the translates of a nonzero point of
order two, and that
\[
              R_\alpha(z)=\frac{\theta(\alpha z)}
                                     {\theta(z)^{N(\alpha)}}
 \tag{OT15}
\]
is an elliptic rational function for each nonzero
$\alpha\in\mathcal O_k$ used below.  Finally suppose the
transported Taylor series in the formal coordinate is
\[
               f(T)=\widehat\theta(\omega\log(1+T))
 \tag{OT16}
\]
with the same period $\omega$ as in the differentiation
directions; its constant term is one.  These are hypotheses on
the actual model and normalization, not consequences of bounded
formal algebra.

\begin{corollary}\label{ot-CM-base}
Under these hypotheses, $f$ is integral.  Consequently the Taylor
functional of $\theta$ in the period-scaled one-dimensional
direction is a measure of norm at most one on $\mathbb Z_p$.
\end{corollary}
\begin{proof}
The Chinese remainder theorem gives
$\alpha\in\mathcal O_k$ which is one modulo the chosen prime
over $p$ and zero modulo its conjugate.  Thus $|a|=1$ and
$m=N_{k/\mathbb Q}(\alpha)$ is a positive integer divisible by
$p$.  In particular $m\ge2$ and $|m|<1$.

A good integral formal group has no nonzero prime-to-$p$ torsion.
In fact its multiplication-by-$n$ series is
$nU+U^2H(U)$ with $H$ integral; when $p\nmid n$ and
$0<|U|<1$, its absolute value is $|U|$, so it cannot vanish.
Since $p$ is odd, the theta zero of order two does not lie in the
formal subgroup.  No preimage of this zero under a CM
endomorphism lies there either, because that subgroup is
preserved.  The zeros and poles of $R_\alpha$ therefore avoid
the formal subgroup.  Also $R_\alpha(O)=1$.
Proposition~\ref{ot-elliptic-bounded} makes the expansion
$\epsilon(T)$ of $R_\alpha$ integral, with integral inverse.

The coefficients of $(1+T)^a-1$ are integral, by the
integer-valued binomial polynomials and density of $\mathbb N$
in $\mathbb Z_p$; its linear coefficient is the unit $a$.
Equations (OT14)--(OT16), as formal identities, give
\[
              f((1+T)^a-1)=f(T)^m\epsilon(T).
 \tag{OT17}
\]
Theorem~\ref{ot-functional-equation}, with $A=1$ and
$B=\epsilon$, proves that $f$ is integral.
Proposition~\ref{ot-logarithm} now gives the asserted measure.
\end{proof}

The extra hypotheses are necessary to identify this corollary with
the theta function used in a global construction.  In particular
one must prove existence and normalization of the good CM model
and multiplicative chart, that its period is a unit in the chosen
normalization, integrality for the required prime-to-$p$ torsion
translates and their finite linear combinations, and the passage
from the one-dimensional theta function to the higher-dimensional
lattice theta function.  The tensor descent in
Theorem~\ref{ot-tensor-root} closes the final extraction from the
fourth tensor power; it does not supply that tensor power's
integrality.  The original proof's treatment of $p=2$ requires
additional information on its specific point of order two, and
is not inferred from Corollary~\ref{ot-CM-base}.

\subsection{Checks against stronger false assertions}

\begin{enumerate}
\item Bounded power moments are insufficient.  The functional
$P\mapsto P'(0)$ has power moments zero except for the moment of
$X$, which is one, but takes $\binom X{p^a}$ to
$(-1)^{p^a-1}/p^a$ for every $a\ge1$.  These values have
unbounded absolute value on polynomials of supremum norm one,
so there is no bounded extension.  The binomial, rather than power, moments in
Corollary~\ref{ot-measure} are essential.
\item The invertible linear part in
Theorem~\ref{ot-functional-equation} cannot be discarded.  In
$\mathbb Q_p[[T]]$, put
$f(T)=\exp(T/p)$, $\phi(T)=pT$ and $m=p$.  Then
$f(\phi)=f^p$, with $A=B=1$, although $f$ is not integral.
\item The contraction condition $|m|<1$ is also essential.
For $\phi(T)=T$ and $m=1$, every series with constant term one
satisfies the equation with $A=B=1$, including $1+p^{-1}T$.
\item Boundedness of a series does not imply boundedness of its
reciprocal: $1+p^{-1}T$ is bounded, but its reciprocal has
coefficients $(-p^{-1})^n$.  Both bounds are used in
Lemma~\ref{ot-bounded-unit}.
\item At $p=2$ a formal group can have nonzero torsion of order
two: in the multiplicative formal group the parameter $T=-2$
has $|T|_2<1$ and $(1+T)^2-1=0$.  Thus the last proof's
prime-to-$p$ torsion argument genuinely requires odd $p$.
\end{enumerate}


\section{A classical ordinary CM input with the required normalization}
\label{ct-section}

The integrality theorem for reduced CM theta functions gives a shorter
route to part of the ordinary theta input.  We state precisely the
classical result used, match the even theta function and the period,
and prove the stronger norm assertion for finite combinations.  The
latter does not follow merely by bounding each theta function
separately.

Fix throughout an abstract field embedding $\iota:\mathbb C\hookrightarrow
\mathbb C_p$, compatible with the chosen embedding of the algebraic
coefficients.  The submitted construction allows an abstract field
isomorphism, which is more than is needed here.  Every transported
Taylor series, torsion character and complex period below uses this
same $\iota$.  No continuity for the complex topology, and no natural
embedding between the two fields, is asserted.

\subsection{The classical input and its scope}

We use the following classical facts in their elliptic-curve form.
For a good ordinary CM elliptic curve at a split prime $p\ge5$,
the reduced theta function with divisor $[O]$ and derivative one
has integral formal expansions after prime-to-$p$ Mumford translation;
the constant term is a unit unless the translated divisor contains
$O$.  The formal group has an integral isomorphism with
$\widehat{\mathbb G}_m$, whose period relative to a good invariant
differential is a unit.  The precise references are
Bannai--Kobayashi, Proposition~2.15, Corollary~2.16(i), and
\S3.1, in the version specified in the bibliography
\cite{bannai-kobayashi-2007}.

These statements use a good integral model, its invariant
differential, and reduced theta functions; arbitrary exponential
renormalizations are not included.  Prime-to-$p$ torsion includes
points of even order because $p\ge5$.  The full CM ring is the
maximal order.  The associated Hermitian form must be compatible
with complex multiplication; in one dimension it is a real
multiple of $z\bar w$, so this condition is automatic.
We import these classical theorems explicitly.  Their proofs are
not represented here as consequences of the finite moment argument.

Let $d$ be a positive squarefree integer with $d\equiv7\pmod8$,
put
\[
 k=\mathbb Q(\sqrt{-d}),\qquad
 \tau=\frac{1+i\sqrt d}{2},\qquad
 \Lambda=\mathcal O_k=\mathbb Z+\mathbb Z\tau,
 \qquad H(z,w)=\frac{2}{\sqrt d}z\bar w,
 \tag{CT1}
\]
and write $E(z,w)=\operatorname{Im}H(z,w)$.
For a function $F$ and $t\in\mathbb C$, define
\[
 (U_tF)(z)=
  e^{-\pi H(z,t)-\pi H(t,t)/2}F(z+t).
 \tag{CT2}
\]
Thus the Hermitian form is linear in its first variable.
Direct multiplication of the exponentials gives
\[
                   U_tU_s=e^{-\pi iE(t,s)}U_{t+s}.
 \tag{CT3}
\]
For rational lattice points $s,t\in\Lambda\otimes\mathbb Q$,
the factor in (CT3) is a root of unity, because
$E(\Lambda,\Lambda)\subset\mathbb Z$.

Fix a good integral CM model $\mathscr E$ over the valuation ring
$R$ of a finite extension of $\mathbb Q_p$, and a complex
uniformization compatible with a good invariant differential.
After a finite extension, take the full CM action to be defined
on this elliptic scheme.
If its period lattice is $\Gamma=c\Lambda$, its Hermitian form
satisfies
\[
                        H_\Gamma(cz,cw)=H(z,w).
 \tag{CT4}
\]
All the functions and torsion points used below can be defined
after a further finite extension.  We then embed the coefficients
in $\mathbb C_p$.  The chosen good model and its multiplicative
formal-group coordinate are part of the classical input.
The multiplicative coordinate may be defined over the completed
maximal unramified extension, or over $\mathcal O_{\mathbb C_p}$;
it need not be defined over the finite extension with valuation
ring $R$.  The finite ring $R$ will be used for the curve, its
line bundles and their finite free modules of global sections.
Integral changes of formal coordinate are then made after scalar
extension to $\mathbb C_p$.

\subsection{Identification of the even theta function}

Let $\vartheta_\Gamma$ be the reduced odd theta function for
$\Gamma$, with simple zeros exactly at $\Gamma$ and
$\vartheta'_\Gamma(0)=1$.  Put
$\vartheta(z)=\vartheta_\Gamma(cz)$ and $h=1/2$, and define
\[
                 \theta(z)=\frac{(U_h\vartheta)(z)}
                                      {(U_h\vartheta)(0)}.
 \tag{CT5}
\]
The denominator is nonzero because $h\notin\Lambda$.

\begin{lemma}\label{ct-even-identification}
The function (CT5) is entire and even, satisfies $\theta(0)=1$,
has simple zeros exactly at $h+\Lambda$, and satisfies
\[
 \theta(z+\ell)=(-1)^{N_{k/\mathbb Q}(\ell)}
            e^{\pi H(z,\ell)+\pi H(\ell,\ell)/2}\theta(z)
                   \qquad(\ell\in\Lambda).
 \tag{CT6}
\]
It is the unique entire function with this transformation law,
constant term one, and no zeros outside $h+\Lambda$.
\end{lemma}
\begin{proof}
The odd reduced theta function has semicharacter
$\alpha(u+v\tau)=(-1)^{u+v+uv}$.
Translation by $h$ multiplies it by
$e^{2\pi iE(h,u+v\tau)}=(-1)^v$, since
$E(h,u+v\tau)=-v/2$.  The translated semicharacter is therefore
$(-1)^{u+uv}$.  On the other hand
\[
 N(u+v\tau)=u^2+uv+\frac{d+1}{4}v^2\equiv u+uv\pmod2,
\]
because $(d+1)/4$ is even.  This proves (CT6).
Translation gives the stated zeros, and division by the value at
zero normalizes the constant term.

For completeness, every nonzero entire function with this
multiplier has exactly one zero, with multiplicity, in a
fundamental parallelogram.  Choose its boundary to avoid zeros.
The logarithmic derivative changes on translation by $\ell$ by
$\pi H(1,\ell)$.  Pairing the two horizontal and the two
vertical edges gives the boundary integral
\[
       \frac{\pi}{\operatorname{Im}\tau}(\tau-\bar\tau)
                              =2\pi i.
\]
The argument principle gives the claim.  Any other function in
the stated class consequently has one simple zero in the same
coset.  Its quotient by (CT5) is entire and $\Lambda$-periodic,
hence bounded and constant, and the value at zero fixes that
constant.  Finally $\theta(-z)$ has the same multiplier, the same
zero coset, and the same value at zero.  Uniqueness proves
evenness.
\end{proof}

Since $N(\ell)=|\ell|^2$ for $\ell\in k$, (CT6) is equivalent
to
\[
               e^{\pi i|\ell|^2}U_\ell\theta=\theta.
 \tag{CT7}
\]
Thus this is precisely the even theta normalization used in the
ordinary moment construction.  The condition
$d\equiv7\pmod{16}$ in that construction is stronger than the
congruence needed in this identification.

\subsection{One common unit period}

Let $\lambda(u)$ be the formal logarithm for the good invariant
differential on $\mathscr E$, normalized by $\lambda'(0)=1$ in
an integral parameter $u$ whose linear differential is that basis.
Choose the integral multiplicative isomorphism in the form
\[
                T=\exp(\lambda(u)/\Omega_p)-1,
                  \qquad |\Omega_p|_p=1.
 \tag{CT8}
\]
Its inverse is integral with unit linear coefficient.  In the
original variable $z$ for $\Lambda$, the same coordinate reads
\[
              z=\omega\log(1+T),\qquad
                        \omega=\Omega_p/\iota(c).
 \tag{CT9}
\]
This identity is an identity of transported formal Taylor series;
it is not a topological identification of complex and $p$-adic
neighborhoods.

Let $g_2,g_3$ be the Weierstrass invariants of $\Lambda$ and choose
$\delta\ne0$ with
\[
                       \delta^{12}=g_2^3-27g_3^2.
 \tag{CT10}
\]
The good model can be general Weierstrass.  With its invariants
$c_4,c_6,\Delta$, the change of differential gives
$c_4/12=\iota(g_2/c^4)$ and
$c_6/216=\iota(g_3/c^6)$.  Hence the identity
$c_4^3-c_6^2=1728\Delta$ gives
\[
               \iota(\delta)^{12}=\iota(c)^{12}\Delta.
 \tag{CT11}
\]
Since $|\Delta|_p=1$, one has
$|\iota(\delta)|_p=|\iota(c)|_p$.  Therefore
\[
             \Omega=\iota(\delta)\omega,
                    \qquad |\Omega|_p=1,
                    \qquad\omega=\Omega/\iota(\delta).
 \tag{CT12}
\]
The same $\Omega$ works for every torsion translate and every
finite family of directions: no period is chosen after selecting
the translate or a derivative order.

\begin{proposition}\label{ct-translates-integral}
For every $t\in\Lambda\otimes\mathbb Q$ of order prime to $p$
modulo $\Lambda$, the series
\[
                 F_t(T)=\widehat{U_t\theta}
                                  (\omega\log(1+T))
 \tag{CT13}
\]
is integral.  If $t+h\notin\Lambda$, it is a unit series.
\end{proposition}
\begin{proof}
For the good model, the classical input applies to the point $ch$
of order two, so the formal expansion of $U_h\vartheta$ has unit
constant term.  The division in (CT5) is therefore by a $p$-adic
unit.  By (CT3), $U_t\theta$ is this same unit inverse times a
root of unity times $U_{t+h}\vartheta$.
The point $c(t+h)$ still has order prime to $p$.
The classical theorem makes its reduced-theta formal expansion
integral, and a unit unless $t+h\in\Lambda$.  Substitution of
the integral inverse of (CT8) preserves integrality and unit
constant terms.  The change from $cz$ to $z$ is exactly (CT9),
which proves the displayed assertion.
\end{proof}

Proposition~\ref{ot-logarithm} consequently supplies norm-at-most-one
measures for all these one-dimensional period-scaled Taylor
functionals.  The exceptional translate has a zero constant term,
but its integrality and its measure remain valid.

\subsection{An integral expansion detects all global sections}

\begin{lemma}\label{ct-expansion-saturated}
Let $\mathscr X/R$ be a smooth proper curve with geometrically
integral special fiber, let $x\in\mathscr X(R)$, and let
$\mathscr L$ be a line bundle.  Choose an integral trivialization
of $\mathscr L$ in the formal neighborhood of $x$ and an integral
parameter there.  Then the expansion map
\[
 H^0(\mathscr X,\mathscr L)\longrightarrow R[[u]]
 \tag{CT14}
\]
is injective and is saturated: a global section over the fraction
field has integral expansion if and only if it is integral.
After scalar extension to $\mathbb C_p$, the coefficient
supremum norm of an expansion equals the maximum absolute value
of its coordinates in any $R$-basis of $H^0(\mathscr X,\mathscr L)$.
The same assertion holds if the expansion is multiplied by an
integral unit series, or its variable is changed by an integral
series with unit linear coefficient.
\end{lemma}
\begin{proof}
The module of global sections is finite over $R$ by properness
and torsion-free because the line bundle is flat.  It is therefore
free.  Choose an $R$-basis $s_1,\ldots,s_q$.
Reduction modulo a uniformizer injects this module modulo that
uniformizer into the space of sections on the special fiber:
this is the initial part of the cohomology sequence of
$0\to\mathscr L\xrightarrow{\pi}\mathscr L
\to\mathscr L|_{\mathscr X_k}\to0$.
On an integral smooth curve, a section which vanishes in the
completed local ring at a point vanishes globally.  Indeed its
germ vanishes in the local discrete valuation ring, hence on a
nonempty open set, and the generic value then vanishes.  The
local ring embeds in its power-series completion, since a
nonzero element has finite valuation.  Thus the reductions of
the expansions of the $s_i$ are linearly independent.

Some finite list of coefficient functionals already detects this
independence: their kernels form a decreasing sequence of
subspaces in a finite-dimensional vector space, whose full
intersection is zero.  Consequently the coefficient matrix has
a $q\times q$ minor which is a unit of $R$.
For any coefficients $a_i\in\mathbb C_p$, integrality of every
expansion coefficient of $\sum a_i s_i$ implies integrality of
the selected finite list and hence of all $a_i$, by the integral
inverse of that minor.  More generally the same inverse has
entries of absolute value at most one, so
\[
 \max_i|a_i|_p\le
   \sup_n|[u^n]\textstyle\sum_i a_i s_i|_p
                  \le\max_i|a_i|_p.
\]
This proves all the assertions about saturation and norms,
including injectivity.  Multiplication by an integral unit has
integral inverse, and so does a change of variable with unit
linear coefficient.  Applying each operation and its inverse
gives equality of norms.  The zero-dimensional space of sections
causes no exception.
\end{proof}

\begin{lemma}\label{ct-torsion-isometry}
In the situation of Lemma~\ref{ct-expansion-saturated}, let
$\mathscr X$ be an elliptic scheme.  Let $v$ be a torsion section
and suppose $\tau_v^*\mathscr L$ and $\mathscr L$ are isomorphic
on the generic fiber.  Suppose the chosen generic-fiber lift
$A_v$ of translation satisfies $nv=0$ and
$A_v^n=\zeta$ as a line-bundle lift, for some $n\ge1$ and a
root of unity $\zeta$.
Then $A_v$ preserves the integral module of global sections and
its inverse does also.
\end{lemma}
\begin{proof}
Enlarge the finite base field if necessary so that the point,
the lift and the root of unity are defined.  Properness extends
the point to an integral section.  The generic line-bundle
isomorphism is a rational section of the line bundle
$\operatorname{Hom}(\tau_v^*\mathscr L,\mathscr L)$.
Its divisor has no horizontal component.  Since the special
fiber is geometrically integral and reduced, its only possible
vertical divisor is an integral multiple of that fiber.
Multiplying by a power of a uniformizer removes that divisor
and gives an integral isomorphism $B_v$.
Thus $A_v=cB_v$ for a nonzero scalar $c$.
The relation $A_v^n=\zeta$ includes $nv=0$, and
$B_v^n$ is an automorphism of the line bundle over the identity.
It is multiplication by a unit of $R$: all global functions on
the proper connected elliptic scheme are scalars.
Taking valuations in $c^nB_v^n=\zeta$ gives $v(c)=0$.
Hence $A_v$ and its inverse are integral, as required.
\end{proof}

The hypotheses here concern actual line-bundle lifts.  An arbitrary
finite-order matrix need not preserve an arbitrarily chosen
lattice; the vertical-divisor argument supplies the necessary
integral comparison.

\subsection{Norm preservation for finite combinations}

For a bounded power series $F$, write
$\|F\|=\sup_n|[T^n]F|_p$.

\begin{theorem}\label{ct-combination-norm}
Let $a\ge0$, let $s_1,\ldots,s_m,t$ be rational lattice points
whose orders modulo $\Lambda$ are prime to $p$, and let
$b_1,\ldots,b_m\in\mathbb C_p$ be arbitrary.  Then
\[
 \left\|\sum_{j=1}^m b_j
       \widehat{U_t(U_{s_j}\theta)}
                         (\omega\log(1+T))^a\right\|
 =
 \left\|\sum_{j=1}^m b_j
       \widehat{U_{s_j}\theta}
                         (\omega\log(1+T))^a\right\|.
 \tag{CT15}
\]
In particular every common coefficient bound is preserved.
\end{theorem}
\begin{proof}
The cases $a=0$ or $m=0$ are immediate.  Choose an even positive
integer $N$, prime to $p$, such that $Ns_j,Nt\in\Lambda$ for
all $j$.  On the complex elliptic curve $\mathbb C/\Lambda$,
the functions
\[
                           (U_{s_j}\theta)(Nz)^a
\]
are reduced sections of the same symmetric line bundle
\[
              \mathscr L_N=[N]^*\mathcal O([h])^{\otimes a}.
 \tag{CT16}
\]
To check the multiplier, translation by $s_j$ changes the
semicharacter on $N\Lambda$ by
$e^{2\pi iE(s_j,N\ell)}=1$, since $Ns_j\in\Lambda$.
The Hermitian form is $aN^2H$ for every summand.

Use the natural integral model of the divisor
$a[N]^{-1}[h]$.  Its support consists of nonzero torsion points
of order prime to $p$.  They do not reduce to the origin, because
the formal group has no nonzero prime-to-$p$ torsion.  Thus its
canonical section gives an integral trivialization near the
origin.  The corresponding reduced function is
$\theta(Nz)^a$.  By Proposition~\ref{ct-translates-integral}
and the integral substitution $T\mapsto(1+T)^N-1$, its
formal series is an integral unit.  Therefore the actual
reduced Taylor expansion of any section of (CT16) is its
expansion in the integral local trivialization multiplied by
this unit series.  Lemma~\ref{ct-expansion-saturated} identifies
its coefficient norm with the integral global-section norm,
also for arbitrary $\mathbb C_p$ linear combinations.

Put $v=t/N$.  Since $N^2v=Nt\in\Lambda$, translation by $v$
preserves the line bundle of degree $aN^2$.  In reduced analytic
coordinates its lift is
\[
       (A_vF)(z)=e^{-\pi aN^2 H(z,v)-\pi aN^2H(v,v)/2}
                                                        F(z+v).
 \tag{CT17}
\]
The algebraic Mumford translation has exactly this reduced
formula, as in the cited construction.  The lift is algebraic
after a finite extension.  By (CT3), a positive power of $A_v$
is the action of a lattice translation, which is a root of
unity because the line bundle is symmetric.  Lemma~\ref{ct-torsion-isometry}
therefore proves that it is an integral automorphism of the
global-section module.

Using $Nv=t$ in (CT17) gives the exact identity
\[
 A_v\bigl((U_{s_j}\theta)(Nz)^a\bigr)
                         =(U_t(U_{s_j}\theta))(Nz)^a.
 \tag{CT18}
\]
Hence the two sides of (CT15), after substituting $Nz$ for $z$,
have equal norm.  On the formal multiplicative group this
substitution is $T\mapsto(1+T)^N-1$; its inverse is
$T\mapsto(1+T)^{1/N}-1$, and both are integral with unit
linear coefficient.  They preserve the coefficient norm by
Lemma~\ref{ct-expansion-saturated}, or directly by applying
composition and its inverse to bounded series.  Removing this
substitution proves (CT15).
\end{proof}

This proves the entire finite-combination norm assertion for the
one-dimensional theta function, with coefficients allowed to be
arbitrary $p$-adic numbers and with one period fixed beforehand.
It does not rely on estimating the coefficients $b_j$ separately,
so cancellations in their sum are retained.

\subsection{Matching the chosen torsion character}

The multiplicative coordinate in (CT8) is unique only up to
$T\mapsto(1+T)^u-1$ with $u\in\mathbb Z_p^\times$.
The next argument makes this freedom explicit and matches the
character used for the selected CM torsion points.

\begin{proposition}\label{ct-character-normalization}
Let $p\ge5$ split in $k$, and let $\iota$ induce one of its
two primes.  For $n\ge0$, let $C_n$ consist of the classes
\[
 v=\frac{a+b\tau}{p^n}\pmod\Lambda,
 \qquad a,b\in\mathbb Z,\qquad
              \iota(a+b\bar\tau)\in p^n\mathbb Z_p.
 \tag{CT19}
\]
These are precisely the $p^n$-torsion points in the formal
subgroup.  The compatible maps
\[
  v\longmapsto\iota\!\left(e^{2\pi i b/p^n}\right)
      =\iota\!\left(e^{-2\pi iE(1,v)}\right)
 \tag{CT20}
\]
are isomorphisms $C_n\simeq\mu_{p^n}$.  One unit change of
the multiplicative coordinate makes its value $1+T(v)$ equal
to (CT20) simultaneously for every $n$.  The normalized period
$\Omega$ in (CT12) remains a unit.
\end{proposition}
\begin{proof}
Because $p$ splits and is unramified, both
$\iota(\tau)$ and $\iota(\bar\tau)$ lie in $\mathbb Z_p$,
and their difference is a unit.  For each residue of $b$ modulo
$p^n$, (CT19) determines one residue of $a$; hence $|C_n|=p^n$.
The formula for $E(1,v)$ follows from
$\operatorname{Im}\tau=\sqrt d/2$.  It shows that (CT20)
is a well-defined injective group map, hence an isomorphism.
These maps commute with multiplication by $p$.

To identify the formal subgroup, choose
$\beta\in\mathcal O_k$ with
$\iota(\beta)\equiv0\pmod{p^n}$ and
$\iota(\bar\beta)\equiv1\pmod{p^n}$.
In the multiplicative coordinate its action is
$1+T\mapsto(1+T)^{\iota(\beta)}$.
This identity follows first on formal logarithms, where the CM
endomorphism scales the invariant differential by
$\iota(\beta)$, and then on formal series in characteristic
zero.  Thus $\beta$ kills every formal $p^n$-torsion point.
For a complex torsion class represented by $(a+b\tau)/p^n$,
this forces $\beta(a+b\tau)\in p^n\mathcal O_k$, and its
conjugate component gives (CT19).  Conversely the formal
multiplicative group has $p^n$ such points, so cardinalities
give equality with $C_n$.

The initial multiplicative coordinate and (CT20) are now two
compatible isomorphisms from the same groups $C_n$ to
$\mu_{p^n}$.  Their ratio is an automorphism of the inverse
system, given by compatible units modulo $p^n$, hence by one
$u\in\mathbb Z_p^\times$.  Replacing $1+T$ by $(1+T)^u$
matches them.  Both this series and its inverse are integral,
and the period is divided by the unit $u$.  Its absolute value
is unchanged.
\end{proof}

In particular the source's stronger condition
$\iota(\bar\nu)\in p^{2n}\mathbb Z_p$ for a representative
$y=s\nu/p^n$, with $p\nmid s$, implies (CT19).  It also implies
$|y|^2=s^2N(\nu)/p^{2n}\in\mathbb Z$, a fact used in its
additional Weyl-translation phase.  The normalized chart above
matches the root of unity required there.  We now prove the
translation identity, including that phase.

\subsection{The exact translation by selected $p$-power torsion}

For $\zeta\in\mathbb C_p$ with $|\zeta-1|_p<1$ and a bounded
series $F(T)=\sum_m F_mT^m$, define
\[
 (S_\zeta F)(T)=F(\zeta(1+T)-1),\qquad
 [T^n]S_\zeta F
       =\zeta^n\sum_{m\ge n}\binom mn F_m(\zeta-1)^{m-n}.
 \tag{CT21}
\]
The displayed sums converge.  They give a bounded series with
$\|S_\zeta F\|\le\|F\|$, and the inverse operation is
$S_{\zeta^{-1}}$.  Thus $S_\zeta$ preserves the norm.
It preserves sums and products.  These statements follow directly
by evaluating on $|T|_p\le r<1$: every substituted argument has
absolute value at most $\max(r,|\zeta-1|_p)<1$, so the defining
series and all products converge uniformly there.  Alternatively
one can use (CT21) and the ultrametric convergence criterion.

We will also use the following compatibility.  If $R$ is a
rational function on $\mathscr E$, without poles in the formal
subgroup, and $v$ is a point of that subgroup with multiplicative
coordinate $\zeta$, then
\[
                 S_\zeta\widehat R=\widehat{R(z+v)}.
 \tag{CT22}
\]
The hats here use the same formal coordinate on both sides.
The boundedness required on the left follows from
Proposition~\ref{ot-elliptic-bounded}.  To justify the identity,
the integral formal-group parameter converges on the open unit
disc and its addition law there is the elliptic-curve addition
law.  Substituting its convergent coordinate functions into the
rational expression for $R$ gives the same analytic function
after translation by $v$.  Its denominators do not vanish in
this disc.  Expanding this identity at zero, using (CT21), proves
(CT22).  This argument uses translation of points on the elliptic
curve; it does not incorrectly replace a torsion point by its
formal logarithm, which is zero.

For the original even theta function put
\[
           (W_yF)(z)=e^{\pi i|y|^2}(U_yF)(z),\qquad
              \kappa(y)=e^{-2\pi iE(1,y)}.
 \tag{CT23}
\]

\begin{theorem}\label{ct-power-translation}
Let $s\ge1$ be prime to $p$, let $st\in\Lambda$, and let
\[
 y=\frac{s\nu}{p^n},\qquad \nu\in\Lambda,\qquad
                     \iota(\bar\nu)\in p^{2n}\mathbb Z_p.
 \tag{CT24}
\]
For the coordinate normalized in
Proposition~\ref{ct-character-normalization}, set
$\zeta=\iota(\kappa(y))$.  Then $\zeta^{p^n}=1$ and
\[
 \widehat{W_y(U_t\theta)}(\omega\log(1+T))
       =S_\zeta\!\left(
                \widehat{U_t\theta}(\omega\log(1+T))\right).
 \tag{CT25}
\]
In particular all series on the right and left are bounded,
with the same coefficient norm.
\end{theorem}
\begin{proof}
We first prove the assertion for $t=0$ and any $y=\nu/p^n$
with the stronger conjugate divisibility in (CT24).
For $n=0$, $y\in\Lambda$, and (CT25) is just (CT7), with
$\zeta=1$.  Suppose $n\ge1$.  Write $\mathfrak p$ for the
prime of $k$ induced by $\iota$.  Choose by the Chinese
remainder theorem
\[
 \epsilon\in2\mathcal O_k,\qquad
 \epsilon\equiv1\pmod{\mathfrak p^n},\qquad
 \epsilon\equiv0\pmod{\bar{\mathfrak p}^{\,n}}.
 \tag{CT26}
\]
It has nonzero unit image at $\mathfrak p$.
The function
\[
                     R(z)=\frac{\theta(\epsilon z)}
                                      {\theta(\bar\epsilon z)}
 \tag{CT27}
\]
is elliptic: the two numerator and denominator multipliers
coincide because $\epsilon$ and $\bar\epsilon$ have the same
norm, both in (CT6) and in the Hermitian exponential.
Its zeros and poles avoid the formal subgroup.  Indeed each
arises from a preimage of the nonzero point $h$ of order two,
whereas CM endomorphisms preserve the formal subgroup and that
subgroup has no nonzero prime-to-$p$ torsion.
Thus $R$ and its inverse have bounded expansions.

Put $f(T)=\widehat\theta(\omega\log(1+T))$.
For a CM element $\alpha$ denote its integral formal
endomorphism by
$[\alpha](T)=(1+T)^{\iota(\alpha)}-1$.
The formal expansion of (CT27) is
\[
                   \widehat R(T)=
                   \frac{f([\epsilon](T))}{f([\bar\epsilon](T))}.
 \tag{CT28}
\]
Both denominator and numerator have constant term one and are
integral.  Since $y$ belongs to the formal $p^n$-torsion,
\[
 [\epsilon](T\mathbin\oplus y)
     =[\epsilon](T)\mathbin\oplus y,\qquad
 [\bar\epsilon](T\mathbin\oplus y)=[\bar\epsilon](T),
\]
where $\oplus y$ means $T\mapsto\zeta(1+T)-1$.

The points $\bar\epsilon y$ and $(\epsilon-1)y$ belong to
$\Lambda$.  This can be checked at $\mathfrak p$ and
$\bar{\mathfrak p}$ from (CT24)--(CT26); all other primes
see no denominator.  Moreover $\bar\epsilon y\in2\Lambda$,
since $\epsilon\in2\Lambda$ and $p$ is odd.
The multiplier in (CT6) at $\bar\epsilon y$ therefore has
semicharacter one.  Cancelling the equal Hermitian exponential
factors gives the exact complex identity
\[
 R(z+y)=\frac{(U_{\epsilon y}\theta)(\epsilon z)}
                                      {\theta(\bar\epsilon z)}.
 \tag{CT29}
\]
Apply (CT22) to (CT28) and (CT29), then cancel the unit series
$f([\bar\epsilon](T))$.  Since $[\epsilon]$ is an integral
formal automorphism, substituting its inverse gives
\[
                      S_\zeta f=\widehat{U_{\epsilon y}\theta}.
 \tag{CT30}
\]
The formal variable on the right is still $\omega\log(1+T)$.

It remains to check the phase.  Write $\epsilon=A+B\tau$,
with $A,B\in\mathbb Z$, and put $l=(\epsilon-1)y\in\Lambda$.
The strong conjugate divisibility gives
$|y|^2=N(\nu)/p^{2n}\in\mathbb Z$.
From (CT3) and (CT6),
\[
 U_{\epsilon y}\theta
   =e^{\pi iE(y,l)-\pi i|l|^2}U_y\theta.
 \tag{CT31}
\]
Now
\[
 E(y,l)=-B|y|^2,\qquad
 |l|^2=(N(\epsilon)-2A-B+1)|y|^2.
\]
Therefore the exponent coefficient in (CT31) differs from
$|y|^2$ by
$(2A-2-N(\epsilon))|y|^2$, an even integer because
$\epsilon\in2\mathcal O_k$ makes $N(\epsilon)$ divisible
by four.  This proves
$U_{\epsilon y}\theta=W_y\theta$, and establishes (CT25)
for the base theta function.

Return to general $t$ and (CT24).  Put $y'=\nu/p^n$ and
$\zeta'=\iota(\kappa(y'))$, so $y=sy'$ and
$\zeta=(\zeta')^s$.  Let
\[
              q(z)=\frac{(U_t\theta)(z)}{\theta(z)},
                       \qquad Q(z)=q(sz).
 \tag{CT32}
\]
The function $Q$ is elliptic: its multiplier on translating by
$l\in\Lambda$ is $e^{2\pi iE(t,sl)}=1$ because
$st\in\Lambda$.  Its poles are among the preimages of $h$
under multiplication by $s$, so none lies in the formal subgroup.
It is regular at zero.  Proposition~\ref{ot-elliptic-bounded}
shows that $\widehat Q$ is bounded.  Since
$\widehat Q=\widehat q\circ[s]$ and $p\nmid s$, the inverse
of $[s]$ is integral; hence $\widehat q$ is bounded as well.

Apply (CT22) to $Q$ and $y'$.  The identity
$[s](T\mathbin\oplus y')=[s](T)\mathbin\oplus y$ and
the complex identity $Q(z+y')=q(sz+y)$ give, after cancelling
the invertible substitution $[s]$,
\[
                         S_\zeta\widehat q
                                     =\widehat{q(z+y)}.
 \tag{CT33}
\]
The right side is regular at zero: $y$ has $p$-power order and
therefore is not the theta zero of order two.
Finally the definition (CT23) gives
\[
                    W_y(U_t\theta)(z)=q(z+y)W_y\theta(z).
\]
Multiply (CT33) by the already proved base-theta identity and
use multiplicativity of $S_\zeta$.  This is exactly (CT25).
The norm assertion follows from (CT21); integrality of the
right side follows from Proposition~\ref{ct-translates-integral}.
\end{proof}

\subsection{What this classical route supplies}

For any given good ordinary CM model at the actual large split
primes, this section proves the precise matching of the even
reduced theta function, all prime-to-$p$ translated integrality,
one common discriminant-normalized unit period, preservation of
the coefficient norm of finite theta combinations, and the
normalization of its torsion character together with the exact
selected-$p$-power translation identity.  These are all three
one-dimensional theta laws used in the moment construction.
The exact higher-dimensional
lattice-theta construction and its finite translation/support
identities remain to be connected.  The arguments do not assert
an integral chart at a supersingular prime, do not treat $p=2$
or $3$ through the cited elliptic integrality statement, and do
not replace a reduced theta function by an arbitrary quadratic
exponential multiple.


\section{Finite lattice translations and the normalization of theta measures}
\label{lm-section}

We now pass from the one-dimensional ordinary theta laws to the
finite operations which occur in the lattice construction.  These
operations introduce two kinds of constants: square roots of
lattice indices and an algebraic factor correcting a determinant.
Both must be retained.  An absolute value of one over the complex
numbers does not imply an absolute value of one at a finite place.
We prove the required arithmetic normalization and all the finite
measure arguments.  The precise analytic identity which identifies
the submitted higher-dimensional theta function is stated separately.

\subsection{Integral moments in a lattice of directions}

Fix $p$, a complete field $L\supset\mathbb Q_p$, an abstract
embedding $\iota$ of the complex coefficients into $L$, and a
period $\omega\in L$.  For an entire function $F$ on a
finite-dimensional complex space $V$ and an additive subgroup
$\Lambda\subset V$, say that its period-scaled moments are
integral on $\Lambda$ if, for every finite family
$v_1,\ldots,v_r\in\Lambda$, the formal series
\[
 \widehat F\left(\omega\sum_{i=1}^r v_i\log(1+T_i)\right)
                 \quad\hbox{belongs to }\mathcal O_L[[T_1,\ldots,T_r]].
 \tag{LM1}
\]
The notation means transport of Taylor coefficients and of the
coordinates of the $v_i$ by the fixed embedding.  It makes no
claim about convergence in the original complex topology.
By Proposition~\ref{ot-logarithm}, (LM1) is equivalent to
existence of the associated measures of norm at most one.

\begin{lemma}\label{lm-directions}
Suppose the series associated to $v_1,\ldots,v_r$ has coefficient
norm at most $C$.  Let
$w_i=\sum_j a_{ji}v_j$ for $1\le i\le s$, where
$\iota(a_{ji})\in\mathbb Z_p$.  Then the series associated to
$w_1,\ldots,w_s$ has norm at most $C$.
Consequently integral moments are preserved by finite products,
finite linear combinations with integral transported coefficients,
and pullback by a complex linear map $g:V'\to V$ satisfying
\[
 \forall v\in\Lambda'\quad
 \exists N\in\mathbb Z_{>0},\ p\nmid N,\quad Ngv\in\Lambda.
 \tag{LM2}
\]
\end{lemma}
\begin{proof}
In the original series substitute
\[
 T_j=\prod_{i=1}^s(1+S_i)^{\iota(a_{ji})}-1.
 \tag{LM3}
\]
These are integral series with zero constant term.  Every
coefficient of the substituted series is a finite sum of original
coefficients times integral scalars, so its absolute value is at
most $C$.  The formal identity
$\log\prod_i(1+S_i)^{a_i}=\sum_i a_i\log(1+S_i)$ follows by
differentiation and the constant term, and identifies the series
with the stated directions.  This proves the first assertion.
Products and finite integral linear combinations follow by
multiplying or adding the series for one common direction family.
For a finite family in $\Lambda'$, choose the finitely many
integers in (LM2).  Apply the first assertion to the vectors
$Ngv\in\Lambda$ with diagonal coefficients $1/N\in\mathbb Z_p$.
The chain rule identifies this result with the pullback.
\end{proof}

The same proof permits coefficients in the ring of integers of a
split imaginary quadratic field: their images lie in $\mathbb Z_p$.
In particular every finite family of directions in
$\mathcal O_k$ can be reduced to the direction $1$ in the
one-dimensional theorem of Section~\ref{ct-section}.  Products of
the one-dimensional theta functions therefore have integral
moments on the product lattice, with the same period.

\subsection{The finite translation operator}

Let $H$ be a positive definite Hermitian form on $V$, linear in
its first variable, and put $E=\operatorname{Im}H$.  For $x\in V$
define
\[
 (U_xF)(z)=e^{-\pi H(z,x)-\pi H(x,x)/2}F(z+x).
 \tag{LM4}
\]
Then $U_xU_y=e^{-\pi iE(x,y)}U_{x+y}$.
Let $\Lambda_0,\Lambda_1$ be commensurable full lattices on which
$E$ takes integer values.  Suppose a function $\alpha$ on their
rational span restricts to a semicharacter on each lattice:
\[
 \alpha(x+y)=\alpha(x)\alpha(y)e^{\pi iE(x,y)},\qquad
 |\alpha(x)|=1,\qquad x,y\in\Lambda_j.
 \tag{LM5}
\]
The values on the intersection are the same because a single
function $\alpha$ is used.  A theta function for
$(H,\Lambda_j,\alpha)$ is an entire function satisfying
$U_lF=\alpha(l)^{-1}F$ for $l\in\Lambda_j$.
Put $W_x=\alpha(x)U_x$ and
$d=[\Lambda_1:\Lambda_0\cap\Lambda_1]$.  With $R$ any set of
representatives of this finite quotient, define
\[
 I_{\Lambda_0,\Lambda_1}F
                 =\frac1{\sqrt d}\sum_{x\in R}W_xF,
 \tag{LM6}
\]
where $\sqrt d$ is the positive real square root before transport.

\begin{proposition}\label{lm-finite-intertwiner}
The expression (LM6) is independent of the representatives and
is a theta function for $(H,\Lambda_1,\alpha)$.
Suppose moreover $p\nmid d$, every $\alpha(x)$ on the rational
lattice is a root of unity, and the period-scaled moments of
$U_tF$ on $\Lambda_0$ are integral whenever
$Nt\in\Lambda_0$ for some $N>0$ prime to $p$.
Then $I_{\Lambda_0,\Lambda_1}F$ and all its prime-to-$p$
translates have integral moments on $\Lambda_1$.
\end{proposition}
\begin{proof}
For $x,y\in\Lambda_1$, (LM4)--(LM5) and the integrality of
$E(x,y)$ give $W_xW_y=W_{x+y}$.  If
$l\in\Lambda_0\cap\Lambda_1$, then $W_lF=F$.
Thus replacing $x$ by $x+l$ does not change $W_xF$.
Translation by any $y\in\Lambda_1$ permutes the quotient, so
$W_yI_{\Lambda_0,\Lambda_1}F=I_{\Lambda_0,\Lambda_1}F$.
This proves the first assertion, including entireness.

The finite group of order $d$ is killed by $d$, so
$d\Lambda_1\subset\Lambda_0$.  Each representative $x$ has
prime-to-$p$ order modulo $\Lambda_0$.  If
$Nt\in\Lambda_1$ with $p\nmid N$, then
$dN(t+x)\in\Lambda_0$ for every $x\in R$.  Also
\[
 U_t(W_xF)=\alpha(x)e^{-\pi iE(t,x)}U_{t+x}F.
 \tag{LM7}
\]
The pairing $E$ is rational on the common rational span, hence
the exponential in (LM7) is a root of unity.
The transported coefficient $1/\sqrt d$ has norm one since its
square is $1/d$; every root of unity also has norm one.
The assumed translated integrality, (LM2) with $g=1$ and
$N=d$, and Lemma~\ref{lm-directions} now prove the assertion.
\end{proof}

Two successive operators of this form are therefore finite
linear combinations of translates with integral coefficients
whenever both indices are prime to $p$.  Expanding the two sums
introduces only the cocycle in (LM4).  Multiplication by an
algebraic character value which is a $p$-adic unit preserves
the conclusion.  These observations apply to the submitted
construction's two successive lattice moves; their validity
does not depend on choosing representatives with small
archimedean size.

\subsection{An arithmetic factor correcting the determinant}

Let $k$ be an imaginary quadratic field, with class number $h$,
and let $(W,h_W)$ be a Hermitian $k$-space isometric to the
standard Hermitian space $k^m$.  Fix a full
$\mathcal O_k$-lattice $M\subset W$ and one rational isometry
$\gamma_0:W\to k^m$.  For each rational split prime $\ell$ in
a set $\mathcal S$, assume that there is an isometry
$\gamma_\ell$ carrying
$M\otimes\mathbb Z_\ell$ onto
$\mathcal O_k^m\otimes\mathbb Z_\ell$.
Choose one prime $\mathfrak l\mid\ell$, a generator
\[
       \mathfrak l^h=(\pi_\ell),\qquad
       r_\ell^{\,2h}=\frac{\pi_\ell}{\bar\pi_\ell},
 \tag{LM8}
\]
in an algebraic closure, and put
\[
 y_\ell=v_{\mathfrak l}
             \bigl(\det(\gamma_\ell\gamma_0^{-1})\bigr).
 \tag{LM9}
\]
The valuation is normalized to have value one on a uniformizer.
The existence of $\pi_\ell$ follows directly from the definition
of the class number, and the root $r_\ell$ is nonzero.

\begin{theorem}\label{lm-kappa}
The integer $y_\ell$ is independent of the locally integral
isometry $\gamma_\ell$.  Only finitely many of these integers
are nonzero, so
\[
                         \kappa=\prod_{\ell\in\mathcal S}
                                         r_\ell^{-y_\ell}
 \tag{LM10}
\]
is a well-defined nonzero algebraic number.
If $p\in\mathcal S$ and $\gamma$ is any isometry satisfying
$\gamma(M\otimes\mathbb Z_p)=\mathcal O_k^m\otimes\mathbb Z_p$
(so both maps between these specified lattices are integral), then,
for every embedding $\iota:\overline{\mathbb Q}\to\mathbb C_p$,
\[
 |\iota(\kappa)|_p^2
  \left|\iota\bigl(\det(\gamma\gamma_0^{-1})\bigr)\right|_p=1.
 \tag{LM11}
\]
Moreover every complex conjugate of $\kappa$ has absolute value one.
\end{theorem}
\begin{proof}
If $\gamma$ and $\gamma'$ both carry the completed lattice
onto the standard lattice, then $\gamma\gamma'^{-1}$ and its
inverse are integral matrices over
$\mathcal O_k\otimes\mathbb Z_\ell$.  Its determinant is a
unit at both primes above $\ell$.  Taking valuations proves
the independence in (LM9).  There is one nonzero integer $B$
such that both $\gamma_0M$ and $\mathcal O_k^m$ are contained
in $B^{-1}$ times the other.  For $\ell\nmid B$ their
completions are equal, and $\gamma_0$ itself is an allowable
local isometry.  Therefore $y_\ell=0$ away from the finite
set of prime divisors of $B$.

If $\ell\ne p$, the ideals of $\pi_\ell$ and $\bar\pi_\ell$
are supported above $\ell$.  Their images at every place
above $p$ are units.  Equation (LM8) gives
$|\iota(r_\ell)|_p=1$, so (LM10) yields
\[
                   |\iota(\kappa)|_p
                        =|\iota(r_p)|_p^{-y_p}.
 \tag{LM12}
\]
Set $a=\det(\gamma\gamma_0^{-1})$.
The isometry equation implies $a\bar a=1$.
At the chosen prime $\mathfrak p$ above $p$,
$v_{\mathfrak p}(a)=y_p$ by the independence already proved.
At its conjugate the valuation is $-y_p$.
Since $p$ splits, both ramification indices are one, while
$v_{\mathfrak p}(\pi_p)=h$ and
$v_{\bar{\mathfrak p}}(\pi_p)=0$.
If $\iota$ induces $\mathfrak p$, (LM8) consequently gives
\[
 |\iota(a)|_p=p^{-y_p}
                   =|\iota(r_p)|_p^{2y_p}.
 \tag{LM13}
\]
If it induces $\bar{\mathfrak p}$, both exponents change sign
and the same equality holds.  Multiplying by (LM12) squared
proves (LM11).
Finally every complex embedding of $k$ is either the given
one or its conjugate.  Thus $\pi_\ell/\bar\pi_\ell$ has
complex absolute value one in every embedding; every conjugate
of its $2h$-th root has absolute value one as well.  The finite
product (LM10) has the stated archimedean property.
\end{proof}

\begin{corollary}\label{lm-change-scalar}
Suppose $a$ is the determinant in (LM11), and $c$ is an
algebraic number satisfying $c^8=a^4$.  Then
$|\iota(\kappa c)|_p=1$.
\end{corollary}
\begin{proof}
Taking absolute values in the eighth-power identity gives
$|\iota(c)|_p^2=|\iota(a)|_p$.
The assertion is now (LM11).
\end{proof}

This is the exact exponent needed in the source's change of
theta normalization.  Merely knowing that $c$ or $\kappa$
has complex absolute value one would not give this conclusion.
The proof also explains why the choice of one prime over each
split rational prime and the sign in $r_\ell^{-y_\ell}$ matter.

\subsection{The precise fourth tensor criterion}

Let $\Lambda_0=\mathcal O_k^m\subset\mathbb C^m$, with its
product Hermitian form and product theta function
$\Theta_0(z)=\prod_{j=1}^m\theta(z_j)$ as in
Section~\ref{ct-section}.  Let $V$ have complex dimension
$m/4$, and let $\Lambda\subset V$ be a full lattice stable
under $\mathcal O_k$.  Suppose
$g:V^4\to\mathbb C^m$ is a complex linear isometry, and
$M_g=g(\Lambda^4)$ is commensurable with $\Lambda_0$.
Assume a common semicharacter on their rational span, with
root-of-unity values, and write
$j_g=[M_g:M_g\cap\Lambda_0]$.  Define the entire function
\[
 K_g(z)=j_g^{-1/2}
       \sum_{x\in M_g/(M_g\cap\Lambda_0)}
                       \alpha(x)(U_x\Theta_0)(gz).
 \tag{LM14}
\]
As above the expression uses any finite set of representatives.

\begin{theorem}\label{lm-fourth-criterion}
Suppose $p\ge5$ is split in $k$, the one-dimensional theta
function has the ordinary laws of Section~\ref{ct-section},
and $F$ is an entire function on $V$.  Suppose there are two
isometries $g_0,g_1$ as above and algebraic nonzero numbers
$\kappa,c$ such that
\[
 \begin{aligned}
 F^{\otimes4}&=\kappa K_{g_0},&
 K_{g_0}&=cK_{g_1},&
 c^8&=\det(g_1g_0^{-1})^4,\\
 p&\nmid j_{g_1},&
 |\iota(\kappa)|_p^2
      |\iota(\det(g_1g_0^{-1}))|_p&=1 .
 \end{aligned}
 \tag{LM15}
\]
Then every prime-to-$p$ translate $U_tF$ has integral
period-scaled moments on $\Lambda$.
\end{theorem}
\begin{proof}
Corollary~\ref{lm-change-scalar}, or its one-line norm
calculation, gives $|\iota(\kappa c)|_p=1$.
The product function $\Theta_0$ and all its prime-to-$p$
translates have integral moments on $\Lambda_0$, by the
one-dimensional theorem and Lemma~\ref{lm-directions}.
Proposition~\ref{lm-finite-intertwiner} proves the same assertion
for $I_{\Lambda_0,M_{g_1}}\Theta_0$ on $M_{g_1}$.

If $Nt\in\Lambda$, with $p\nmid N$, put
$t^{[4]}=(t,t,t,t)$.  Since $g_1$ is an isometry,
direct substitution in (LM4) gives
\[
 U_{t^{[4]}}(K_{g_1})
       =(U_{g_1t^{[4]}}
              I_{\Lambda_0,M_{g_1}}\Theta_0)\circ g_1.
 \tag{LM16}
\]
The vector $g_1t^{[4]}$ has prime-to-$p$ order modulo
$M_{g_1}$.  Pullback by $g_1$, which carries the direction
lattice $\Lambda^4$ exactly onto $M_{g_1}$, preserves integral
moments.  Applying (LM15) and (LM16) gives integral moments
for
\[
 U_{t^{[4]}}F^{\otimes4}=(U_tF)^{\otimes4}.
 \tag{LM17}
\]
Theorem~\ref{ot-tensor-root} now proves the assertion for $U_tF$.
\end{proof}

The arithmetic part of the last condition of (LM15) is precisely
Theorem~\ref{lm-kappa} when $g_0,g_1$ arise by complexifying the
rational isometries there.  The remaining analytic statement
$K_{g_0}=cK_{g_1}$ with its eighth-power determinant relation
is a substantive theta transformation theorem.  It is not proved
by the elementary moment criterion, by linearity of derivatives,
or by the complex norm of the two functions.

\subsection{Finite Fourier projection and support of measures}

Let $X=\mathbb Z_p^r$, let $\mu$ be a bounded $L$-valued
measure, and suppose $L$ contains all $p$-power roots of unity.
For a continuous function $w$, the measure $w\mu$ means
$(w\mu)(f)=\mu(wf)$.  For $n\ge1$ and $a\in X$, choose a
primitive $p^n$-th root $\zeta_n\in L$.

\begin{lemma}\label{lm-character-projection}
There is an exact equality of continuous functions
\[
 {\bf1}_{a+p^nX}(x)=p^{-nr}
     \sum_{y\in(\mathbb Z/p^n\mathbb Z)^r}
                       \zeta_n^{\,\langle y,x-a\rangle}.
 \tag{LM18}
\]
Consequently, restriction of $\mu$ to any residue class has
norm at most $\|\mu\|$, even though the displayed expression
contains a denominator divisible by $p$.
If
\[
       \chi_y\mu=\chi_y(a)\mu
       \quad\hbox{for every character }\chi_y
               \hbox{ of }X/p^nX,
 \tag{LM19}
\]
then $\mu$ is supported on $a+p^nX$.
\end{lemma}
\begin{proof}
The sum factors as a product of $r$ geometric sums.
A one-dimensional sum is $p^n$ if $x_i-a_i$ is zero
modulo $p^n$, and zero otherwise: in the latter case multiply
the sum by the nontrivial root $\zeta_n^{x_i-a_i}$ and subtract.
This proves (LM18).  Its left side has supremum norm one,
which proves the norm assertion.
Multiplying $\mu$ by (LM18) and applying (LM19) gives $\mu$
again.  For any continuous $f$ vanishing on the residue class,
$\mu(f)=\mu({\bf1}_{a+p^nX}f)=0$.  This is the stated support.
\end{proof}

Thus estimating the terms of the Fourier sum separately would
introduce a spurious loss of $p^{nr}$.  The projection itself
has integral values and has norm one.  The same proof applies
to the inverse image of a coset in any finite quotient of $X$,
using the characters of that quotient.

\begin{theorem}\label{lm-finite-approximation}
Let $Y\subset X$ be closed and let $D\subset Y$ be dense.
Suppose $\mu$ is supported on $Y$ and $\|\mu\|\le C$.
For every finite family $f_1,\ldots,f_s\in C(X,L)$ and
$\epsilon>0$, there are $d_1,\ldots,d_t\in D$ and
$a_1,\ldots,a_t\in L$ with $|a_i|\le C$ such that
\[
 \left|\mu(f_j)-\sum_{i=1}^t a_if_j(d_i)\right|<\epsilon
                             \qquad(1\le j\le s).
 \tag{LM20}
\]
If $C=1$, all coefficients are integral.
\end{theorem}
\begin{proof}
For $C=0$ the measure is zero and the empty sum suffices.
For $C>0$, uniform continuity gives $n$ for which every $f_j$
oscillates by less than $\epsilon/C$ on each residue class
modulo $p^n$.  There are finitely many such classes meeting
$Y$.  Each has a point $d_i\in D$, since its intersection
with $Y$ is relatively open and nonempty.  Let $B_i$ be the
class and set $a_i=\mu({\bf1}_{B_i})$.
Then $|a_i|\le\|\mu\|\le C$.
The function equal to $f_j(d_i)$ on each of these classes
differs from $f_j$ by less than $\epsilon/C$ on their union.
Outside that union, the measure vanishes by its support
condition.  Applying its norm bound proves (LM20).
The case $Y=\varnothing$ again gives the zero measure.
\end{proof}

This proves the finite simultaneous approximation required by a
polynomial-kernel argument as soon as the actual translation
identities establish the support and the rational lattice is
dense in that support.  It preserves the norm bound on each
coefficient and does not insert a factor equal to the number
of points in the finite sum.

\subsection{Limits of the reduction and counterexamples}

The preceding results prove every finite operation in the
passage from a locally normalized fourth tensor identity to
integral lattice moments, and the exact finite Fourier
projection used to obtain support.  To identify them with the
submitted global kernels one still has to supply the theta
transformation in (LM15), the locally standard isometries with
the required period lattices, and the actual higher-dimensional
translation equations which impose the stated kernel support.
Their names in the formal source do not discharge these
mathematical obligations.

Two false shortcuts are worth excluding.  First, over
$k=\mathbb Q(i)$ the number $a=(3+4i)/5$ has complex
absolute value one.  At the embedding into $\mathbb Q_5$
with $i\equiv2\pmod5$, its numerator is a unit, so
$|a|_5=5$.  The determinant correction in (LM11) cannot be
replaced by an archimedean norm calculation.
Second, multiplication of a theta function by $p^{-1}$
preserves its transformation law but destroys an integral
constant term of norm one.  The theta transformation law
alone therefore does not fix the arithmetic normalization.
These examples refute weakened claims, not the statements
with the normalizations retained above.


\section{The exact constant in a change of theta lattice}
\label{tc-section}

The remaining analytic normalization in Section~\ref{lm-section}
can be calculated explicitly.  We give the calculation, including
the finite quadratic Gauss sum on which its eighth power depends.
All square roots of positive real numbers in this section mean
the positive square root.  This convention is part of the
normalization.

Let $V$ be a complex vector space of dimension $r$, let $H$ be a
positive definite Hermitian form, linear in its first variable,
and put $E=\operatorname{Im}H$.  A full lattice $\Lambda$ is
self-dual if
\[
 \Lambda=\Lambda^*
   :=\{x\in V:E(x,\Lambda)\subset\mathbb Z\}.
 \tag{TC1}
\]
A sign semicharacter on $\Lambda$ is a function
$\alpha:\Lambda\to\{1,-1\}$ satisfying
\[
 \alpha(x+y)=\alpha(x)\alpha(y)e^{\pi iE(x,y)}.
 \tag{TC2}
\]
For several lattices we require their semicharacters to be
restrictions of one function, so they agree on intersections.
Write
\[
 U_xF(z)=e^{-\pi H(z,x)-\pi H(x,x)/2}F(z+x),
 \qquad W_x=\alpha(x)U_x.
 \tag{TC3}
\]
The theta space $\mathcal T(\Lambda,\alpha)$ consists of entire
functions $F$ satisfying $W_xF=F$ for $x\in\Lambda$.
The basic identities are
\[
 U_xU_y=e^{-\pi iE(x,y)}U_{x+y},\qquad
 U_xU_y=e^{-2\pi iE(x,y)}U_yU_x.
 \tag{TC4}
\]
They follow by expanding the three exponential factors and using
$H(y,x)=\overline{H(x,y)}$.  On any one of the lattices in
(TC2), the operators $W_x$ form a representation of its additive
group.  Also $\alpha(0)=1$ and $\alpha(-x)=\alpha(x)$.

\subsection{Self-dual lattices and their theta spaces}

Choose Lebesgue measure $dz$ by using any $H$-orthonormal complex
basis.  Thus
\[
                  \int_V e^{-\pi H(z,z)}\,dz=1.
 \tag{TC5}
\]
This measure does not depend on the chosen orthonormal basis.

\begin{lemma}\label{tc-covolume}
Every self-dual lattice has covolume one for this measure.
For every sign semicharacter on such a lattice,
$\dim_{\mathbb C}\mathcal T(\Lambda,\alpha)=1$.
\end{lemma}
\begin{proof}
In real orthonormal coordinates the matrix of $E$ has determinant
one.  If $B$ is a real basis matrix for $\Lambda$, its integral
alternating Gram matrix is $B^{\mathsf T}JB$.
Self-duality says that this integer matrix is unimodular.
Therefore $(\det B)^2=1$, which proves the covolume assertion.

For completeness, the dimension assertion has an elementary
Fourier proof.  A unimodular alternating integer form has a
symplectic basis $e_1,\ldots,e_r,f_1,\ldots,f_r$ with
$E(e_i,f_j)=-\delta_{ij}$.
Indeed, a primitive vector $e_1$ has a partner $f_1$ with
$E(e_1,f_1)=-1$ by unimodularity.  Subtracting integer multiples
of these two vectors splits off their orthogonal complement;
induction supplies the basis.
The $e_i$ are a complex basis of $V$: if
$\sum a_i e_i=i\sum b_i e_i$ with $a_i,b_i\in\mathbb R$,
pairing the two sides with $\sum b_i e_i$ under $E$ gives
$H(\sum b_i e_i,\sum b_i e_i)=0$.  Positivity then gives
$a_i=b_i=0$.

Use this complex basis, and write
$H(z,w)=z^{\mathsf T}G\overline w$ and
$(f_1,\ldots,f_r)=(e_1,\ldots,e_r)Z$.
The matrix $G$ is real symmetric and positive definite, since
$E(e_i,e_j)=0$.  The equations
$E(e_i,f_j)=-\delta_{ij}$ and $E(f_i,f_j)=0$ give
\[
             \operatorname{Im}Z=G^{-1}>0,\qquad Z=Z^{\mathsf T}.
 \tag{TC6}
\]
Choose $a_i,b_i\in\{0,\tfrac12\}$ with
$\alpha(e_i)=e^{2\pi ia_i}$ and
$\alpha(f_i)=e^{2\pi ib_i}$.
For $F\in\mathcal T(\Lambda,\alpha)$ put
$f(z)=e^{-\pi z^{\mathsf T}Gz/2}F(z)$.
The transformation equations become
\[
 \begin{split}
 f(z+u)&=e^{2\pi ia^{\mathsf T}u}f(z),\\
 f(z+Zv)&=e^{2\pi ib^{\mathsf T}v}
        e^{-2\pi iv^{\mathsf T}z-\pi iv^{\mathsf T}Zv}f(z)
                 \qquad(u,v\in\mathbb Z^r).
 \end{split}
 \tag{TC7}
\]
The first equation gives a Fourier expansion
$f(z)=\sum_{n\in\mathbb Z^r}c_n
 e^{2\pi i(n+a)^{\mathsf T}z}$.
Here the coefficients are obtained by integration on a real
unit cube after multiplication by $e^{-2\pi ia^{\mathsf T}z}$.
Cauchy's theorem on the coordinate rectangles makes the
appropriately weighted coefficient independent of the imaginary
part of $z$.  Shifting these rectangles in either direction
bounds the coefficients exponentially on every smaller closed
strip.  The resulting Fourier series therefore converges normally
on compact subsets and represents $f$.
Substitution in the second equation of (TC7) gives
\[
 c_{n+v}=c_n
  \exp\!\left(2\pi i(n+a)^{\mathsf T}Zv+
              \pi iv^{\mathsf T}Zv-2\pi ib^{\mathsf T}v\right).
 \tag{TC8}
\]
All coefficients are determined by $c_0$.
Conversely,
\[
 e^{\pi z^{\mathsf T}Gz/2}
 \sum_{n\in\mathbb Z^r}
  \exp\!\left(\pi i(n+a)^{\mathsf T}Z(n+a)
                 +2\pi i(n+a)^{\mathsf T}(z-b)\right)
 \tag{TC9}
\]
converges normally, because $\operatorname{Im}Z>0$, and
satisfies (TC7).  Its Fourier coefficients are nonzero, so it is
not the zero function.  This proves the dimension assertion.
The case $r=0$ consists simply of the constant functions.
\end{proof}

\subsection{Normalized finite averaging is unitary}

Suppose $\Lambda_i$ are commensurable self-dual lattices carrying
one common sign semicharacter.  Put
\[
 d_{ij}=[\Lambda_j:\Lambda_i\cap\Lambda_j]
        =[\Lambda_i:\Lambda_i\cap\Lambda_j],\qquad
 I_{ij}F=d_{ij}^{-1/2}
     \sum_{x\in\Lambda_j/(\Lambda_i\cap\Lambda_j)}W_xF.
 \tag{TC10}
\]
The equality of indices follows from Lemma~\ref{tc-covolume}.
On functions fixed by $\Lambda_i$ the summands are independent
of the choice of representatives.

\begin{lemma}\label{tc-unitary-intertwiner}
The operator $I_{ij}$ is a unitary isomorphism from the
$\Lambda_i$ theta space to the $\Lambda_j$ theta space, with
inverse $I_{ji}$.  The same inverse and norm assertions hold on
continuous functions with the corresponding transformation laws.
\end{lemma}
\begin{proof}
Let $\Gamma$ be the intersection of all lattices under
consideration.  For continuous functions satisfying the
$\Gamma$ transformation law define
\[
 \langle F,G\rangle_\Gamma
  =\frac1{\operatorname{vol}(V/\Gamma)}
    \int_{V/\Gamma}F(z)\overline{G(z)}
                        e^{-\pi H(z,z)}\,dz.
 \tag{TC11}
\]
The integrand is $\Gamma$-periodic.  Its norm is positive on
nonzero continuous functions.  For $x\in\Gamma^*$ the operator
$U_x$ preserves this space and is unitary, since
\[
 |U_xF(z)|^2e^{-\pi H(z,z)}
       =|F(z+x)|^2e^{-\pi H(z+x,z+x)}.
 \tag{TC12}
\]
All $\Lambda_i$ lie in $\Gamma^*$.
Consequently
\[
       P_i=\frac1{[\Lambda_i:\Gamma]}
                   \sum_{x\in\Lambda_i/\Gamma}W_x
 \tag{TC13}
\]
is an orthogonal projection onto the functions fixed by
$\Lambda_i$: it is the average of a finite unitary group
representation.

For $x\in\Gamma^*$, the commutator in (TC4) and finite
character orthogonality show that $P_iU_xP_i=0$ unless
$x\in\Lambda_i^*=\Lambda_i$.
Explicitly, the character being averaged is
$u\mapsto e^{-2\pi iE(u,x)}$ on $\Lambda_i/\Gamma$;
a nontrivial character has zero sum because translating its
sum by an element on which it is not one multiplies that sum
by a nontrivial scalar.
If $x\in\Lambda_i$, then $U_xP_i=\alpha(x)P_i$.
For $x$ also belonging to another specified lattice, the
common semicharacter therefore gives $W_xP_i=P_i$.
Thus
\[
 P_iP_jP_i=
 \frac{[(\Lambda_i\cap\Lambda_j):\Gamma]}
      {[\Lambda_j:\Gamma]}P_i=d_{ij}^{-1}P_i.
 \tag{TC14}
\]
On the image of $P_i$, regrouping the sum (TC13) gives
$I_{ij}=\sqrt{d_{ij}}P_j$.
Equation (TC14) gives $I_{ji}I_{ij}=1$ and
\[
 \|I_{ij}F\|^2=d_{ij}\langle F,P_jF\rangle
             =\|F\|^2.
\]
The reverse composition is the identity as well.
Finite translations preserve holomorphy, so the assertion
restricts to theta spaces.  For a function satisfying a
$\Lambda_i$ law, its normalized norm in (TC11) is independent
of the smaller common lattice $\Gamma$, by subdivision of
fundamental domains.
\end{proof}

\subsection{An eighth-power formula for finite Gauss sums}

A homogeneous quadratic function on a finite abelian group $M$
is a function $q:M\to\mathbb C^\times$ such that
$q(nx)=q(x)^{n^2}$ for integers $n$ and
\[
                b(x,y)=\frac{q(x+y)}{q(x)q(y)}
 \tag{TC15}
\]
is a bicharacter.  Its radical is
$R=\{x:b(x,M)=1\}$.
All values of $q$ are roots of unity, since
$q(x)^{|M|^2}=q(|M|x)=1$.

\begin{lemma}\label{tc-four-squares-congruence}
For every positive integer $N$ there are integers
$a,b,c,d$ with $a^2+b^2+c^2+d^2\equiv-1\pmod N$.
\end{lemma}
\begin{proof}
For an odd prime $\ell$, the sets of squares and of $-1$
minus squares in $\mathbb F_\ell$ both have size
$(\ell+1)/2$, so they intersect.  Thus $-1$ is a sum of
two squares modulo $\ell$.  At least one of the two
coordinates is nonzero.  Adjusting that coordinate at each
step lifts the solution to $\ell^s$, since its derivative
$2a$ or $2b$ is a unit modulo $\ell$.
For powers of two, start with
$1^2+1^2+1^2+2^2=7\equiv-1\pmod8$.
If the sum is correct modulo $2^s$, $s\ge3$, replacing
the first, odd coordinate $a$ by $a+2^{s-1}$ changes
the sum by $2^s$ modulo $2^{s+1}$.  It therefore corrects
the next bit when necessary.  The same initial tuple
also works modulo $2$ and $4$.
Apply the Chinese remainder theorem to each of the four
coordinates.  The assertion for $N=1$ is immediate.
\end{proof}

\begin{theorem}\label{tc-gauss-eight}
If $q$ is nondegenerate, then
\[
                    \left(\sum_{x\in M}q(x)\right)^8=|M|^4.
 \tag{TC16}
\]
More generally, if its sum $S$ is nonzero, then $q$ is trivial
on its radical $R$ and
\[
                         S^8=|R|^4|M|^4>0.
 \tag{TC17}
\]
\end{theorem}
\begin{proof}
First suppose $R=0$, and set $N=|M|^2$.
Choose $a,b,c,d$ as in Lemma~\ref{tc-four-squares-congruence}
and write $s=a^2+b^2+c^2+d^2$.
The integer quaternion matrix
\[
 T=\begin{pmatrix}
 a&-b&-c&-d\\
 b&a&-d&c\\
 c&d&a&-b\\
 d&-c&b&a
 \end{pmatrix}
 \quad\hbox{satisfies}\quad
                  T^{\mathsf T}T=TT^{\mathsf T}=sI_4.
 \tag{TC18}
\]
It permutes $M^4$, because $s\equiv-1\pmod N$ and
$N$ annihilates $M$; its inverse on this group is
$-T^{\mathsf T}$.
For $Q(x_1,\ldots,x_4)=\prod_iq(x_i)$ the quadratic
identity (TC15), applied to each row, gives
\[
                         Q(Tx)=Q(x)^s=Q(x)^{-1}.
 \tag{TC19}
\]
Indeed the exponents of $q(x_i)$ are the squared column
lengths $s$, while the exponents of each $b(x_i,x_j)$
are the column inner products, all zero for $i\ne j$.
The last equality uses $q(x_i)^N=1$.
Summing (TC19) gives $S^4=\overline S^{\,4}$.
On the other hand
\[
 |S|^2
  =\sum_{t\in M}q(t)\sum_{y\in M}b(y,t)=|M|.
 \tag{TC20}
\]
Here we substituted $x=y+t$ in $\sum_{x,y}q(x)/q(y)$;
nondegeneracy and character orthogonality leave only $t=0$.
Multiplying $S^4=\overline S^{\,4}$ by $S^4$ proves (TC16).

For a possibly degenerate $q$ with $S\ne0$, translation
of its sum by $r\in R$ gives $S=q(r)S$.
Thus $q(r)=1$.  The function descends to a nondegenerate
quadratic function $\bar q$ on $M/R$, and
$S=|R|\sum_{M/R}\bar q$.
Equation (TC16) gives
$S^8=|R|^8|M/R|^4=|R|^4|M|^4$, as claimed.
\end{proof}

\subsection{The phase in a composition of three lattice averages}

\begin{theorem}\label{tc-triple}
For three commensurable self-dual lattices with a common sign
semicharacter there is a scalar $\mu$ such that
\[
                      I_{23}I_{12}=\mu I_{13},
                         \qquad \mu^8=1
 \tag{TC21}
\]
on every function fixed by $\Lambda_1$.
\end{theorem}
\begin{proof}
Choose representatives $S_2$ for
$\Lambda_2/(\Lambda_1\cap\Lambda_2)$ and $S_3$ for
$\Lambda_3/(\Lambda_2\cap\Lambda_3)$.
For $z\in\Lambda_3$, $y\in\Lambda_2$, $z+y\in\Lambda_1$
put
\[
 \psi(z,y)=\alpha(z)\alpha(y)\alpha(z+y)e^{-\pi iE(z,y)},
 \quad
 B=\sum_{\substack{z\in S_3,\ y\in S_2\\z+y\in\Lambda_1}}
                         \psi(z,y).
 \tag{TC22}
\]
The projection argument of Lemma~\ref{tc-unitary-intertwiner}
gives, for $F$ fixed by $\Lambda_1$,
\[
 I_{31}I_{23}I_{12}F
       =\frac{\sqrt{d_{13}}}{\sqrt{d_{12}}\sqrt{d_{23}}}\,BF.
 \tag{TC23}
\]
In detail, expand the two rightmost averages and apply
$I_{31}=\sqrt{d_{13}}P_1$ on their output.
The projection kills $W_zW_yF$ unless $z+y\in\Lambda_1$,
by character orthogonality.  In the remaining cases
$W_zW_yF=\psi(z,y)F$ by (TC4).
Applying the inverse $I_{13}$ proves (TC21) with
$\mu=\sqrt{d_{13}}B/(\sqrt{d_{12}}\sqrt{d_{23}})$.
The unitarity of all three operators, applied to a nonzero
theta function from Lemma~\ref{tc-covolume}, gives $|\mu|=1$.

It remains to determine its phase.  Define subgroups
\[
 \begin{split}
 P&=\{(z,y)\in\Lambda_3\times\Lambda_2:z+y\in\Lambda_1\},\\
 T&=\{(t,u-t):t\in\Lambda_2\cap\Lambda_3,\
                         u\in\Lambda_1\cap\Lambda_2\}.
 \end{split}
 \tag{TC24}
\]
The pairs in the sum (TC22) form representatives for $P/T$.
To see existence, first reduce $z$ modulo
$\Lambda_2\cap\Lambda_3$ to $z_0\in S_3$, moving the
discarded element into $y$, and then reduce the new $y$
modulo $\Lambda_1\cap\Lambda_2$ to $y_0\in S_2$.
The sum $z_0+y_0$ stays in $\Lambda_1$.
For uniqueness, membership of the difference in $T$ first
forces equality of the $S_3$ representatives and then
equality of the $S_2$ representatives.  In particular $P/T$
is finite.

Direct use of (TC2) gives
\[
 \begin{split}
 \psi(z+z',y+y')&=\psi(z,y)\psi(z',y')
                                      e^{2\pi iE(y,z')},\\
 \psi(nz,ny)&=\psi(z,y)^{n^2}.
 \end{split}
 \tag{TC25}
\]
For the first equation, the excess exponent after expanding
the three semicharacters is
$2\pi i\{E(z,z')+E(y,y')+E(y,z')\}$.
The first two terms are integers and disappear.
For the second, $\alpha(nx)=\alpha(x)^n
=\alpha(x)^{n^2}$ on each lattice, since its values are signs.
The polar factor in (TC25) is a bicharacter.
For $(t,u-t)\in T$, formula (TC2) on $\Lambda_2$ gives
$\psi(t,u-t)=1$, and $E(y,t)\in\mathbb Z$ gives trivial
polar factor against every $(z,y)\in P$.
Thus $\psi$ descends to a homogeneous quadratic function
on $P/T$.  Its sum is $B$.
Since $|\mu|=1$, this sum is nonzero.  By
Theorem~\ref{tc-gauss-eight}, $B^8$ is positive real.
The multiplier relating $\mu$ to $B$ is positive real,
so $\mu^8$ is positive real and has absolute value one.
It is therefore one.
\end{proof}

\subsection{The Cayley phase and the Gaussian trace}

Let $\delta$ be a complex linear $H$-isometry with no nonzero
fixed vector.  Assume that $\Lambda$ and $\delta\Lambda$
are commensurable self-dual lattices, that the common function
$\alpha$ is a sign semicharacter on both, and that
\[
                         \alpha(\delta x)=\alpha(x)
                                      \quad(x\in\Lambda).
 \tag{TC26}
\]
This last equality is essential; agreement on an intersection
alone would not assert it.
Set
\[
 A=1-\delta^{-1},\quad C=A^{-1}-\tfrac12,\quad
 \Gamma=\Lambda\cap\delta\Lambda,\quad d=[\Lambda:\Gamma],
 \quad
 q(x)=\alpha(x)e^{-\pi H(Cx,x)}\quad(x\in\Lambda).
 \tag{TC27}
\]

\begin{lemma}\label{tc-cayley-quadratic}
The function $q$ descends to a nondegenerate homogeneous
quadratic function on $\Lambda/A\Gamma$.  Moreover
\[
                 N:=|\Lambda/A\Gamma|=d|\det_{\mathbb C}A|^2.
 \tag{TC28}
\]
\end{lemma}
\begin{proof}
Unitarity gives
$A^*=-\delta A$ and $A^{-1}+(A^*)^{-1}=1$,
so $C^*=-C$.
Thus $H(Cx,x)$ is purely imaginary.  For $x,y\in\Lambda$,
expanding $q(x+y)$ and using $E(x,y)\in\mathbb Z$ gives
\[
 \frac{q(x+y)}{q(x)q(y)}
                   =e^{-2\pi iE(A^{-1}x,y)}.
 \tag{TC29}
\]
For example the exponent before reducing modulo $2\pi i\mathbb Z$
is $\pi iE(x,y)-2\pi iE(Cx,y)$, equal to the displayed
exponent plus $2\pi iE(x,y)$.
This is a bicharacter; its symmetry follows also from its
definition as a polar factor.  Homogeneity follows from
$\alpha(nx)=\alpha(x)^{n^2}$.

If $t\in\Gamma$, then $t,\delta^{-1}t\in\Lambda$ and
$At=t-\delta^{-1}t\in\Lambda$.
Also
$H(CAt,At)=-iE(t,\delta^{-1}t)$.
Equations (TC2) and (TC26) give
\[
 q(At)=\alpha(t)\alpha(\delta^{-1}t)=1.
 \tag{TC30}
\]
The polar factor against $At$ is one, because $t\in\Lambda$.
Thus $q$ descends modulo $A\Gamma$.
The radical of (TC29) on $\Lambda$ consists precisely of
the $x$ with $A^{-1}x\in\Lambda^*=\Lambda$.
For such an $x$, write $t=A^{-1}x$.
The equation $\delta^{-1}t=t-x\in\Lambda$ implies
$t\in\Gamma$, hence $x\in A\Gamma$.
This proves nondegeneracy on the quotient.
Finally $A$ changes real volume by $|\det_{\mathbb C}A|^2$.
Lemma~\ref{tc-covolume} and $\operatorname{vol}(V/\Gamma)=d$
therefore give (TC28).
\end{proof}

\begin{theorem}\label{tc-unitary-constant}
For every $F\in\mathcal T(\Lambda,\alpha)$,
\[
 F(\delta^{-1}z)=
       c(\delta,\Lambda)\,I_{\Lambda,\delta\Lambda}F(z),
 \quad
 c(\delta,\Lambda)=
 \frac{|\det A|}{\det A}\,
         \frac1{\sqrt N}\sum_{x\in\Lambda/A\Gamma}q(x).
 \tag{TC31}
\]
In particular
\[
                         c(\delta,\Lambda)^8=(\det\delta)^4.
 \tag{TC32}
\]
\end{theorem}
\begin{proof}
Both sides of the first equation belong to the same
one-dimensional theta space; the left side has the required
transformation law by (TC26).  Lemma~\ref{tc-unitary-intertwiner}
makes its right-hand operator nonzero.
There is therefore a scalar $c$, independent of $F$.
We compute it, which is the step needed for its arithmetic
absolute value.

Put $K_w(z)=e^{\pi H(z,w)}$ and consider the periodized kernel
\[
                       P(z,w)=\sum_{x\in\Lambda}W_xK_w(z).
 \tag{TC33}
\]
For $z,w$ in compact sets, each summand and every derivative
is bounded by a polynomial in $\|x\|$ times
$e^{-\pi\|x\|^2/2+C\|x\|}$.
The lattice sum is consequently normally convergent.
Reindexing under $W_\lambda$ shows that
$P(\,\cdot\,,w)\in\mathcal T(\Lambda,\alpha)$.

We need a diagonal integral.  For $y\in\delta\Lambda$,
expand (TC33) and use (TC4).  After multiplication by
$e^{-\pi H(z,z)}$, its $x$-summand is
\[
 \alpha(y)\alpha(x)e^{-\pi iE(y,x)}
       e^{-\pi H(x+y,x+y)/2}e^{2\pi iE(x+y,z)}.
 \tag{TC34}
\]
The last factor is a character of the real torus $V/\Gamma$.
It has integral zero unless $x+y=0$; nondegeneracy of $E$
ensures that a nonzero $x+y$ gives a nontrivial character.
If $y\in\Lambda$ the unique surviving term is $x=-y$,
whose coefficient is one.  Otherwise no term survives.
Termwise integration is justified by the preceding Gaussian
bound on a compact fundamental parallelepiped.  Thus
\[
 \int_{V/\Gamma}(W_yP(\,\cdot\,,z))(z)
                        e^{-\pi H(z,z)}\,dz
     =\begin{cases}d,&y\in\Lambda,\\0,&y\notin\Lambda.\end{cases}
 \tag{TC35}
\]
The case $y=0$ shows that the kernel is not identically zero.
Averaging (TC35) over representatives of $\delta\Lambda/\Gamma$
gives
\[
 \int_{V/\Gamma}
   (I_{\Lambda,\delta\Lambda}P(\,\cdot\,,z))(z)
                          e^{-\pi H(z,z)}\,dz=\sqrt d.
 \tag{TC36}
\]

Set $g(z)=e^{-\pi H(Az,z)}$.
A direct expansion, using $C^*=-C$, gives the exact identity
\[
 \alpha(x)(U_xK_z)(\delta^{-1}z)e^{-\pi H(z,z)}
                  =q(x)\,g(z-A^{-1}x).
 \tag{TC37}
\]
For clarity, the linear terms agree because
$H(Az,A^{-1}x)=-H(\delta^{-1}z,x)$.
The constant terms agree because
$H(A^{-1}x,x)+H(x,A^{-1}x)=H(x,x)$.
Also
\[
                 |g(z)|=e^{-\pi H(Az,Az)/2}.
 \tag{TC38}
\]
This follows from $A+A^*=A^*A$ and is an integrable
positive Gaussian, since $A$ is invertible.
Decompose $\Lambda$ into the $N$ cosets of $A\Gamma$.
By Lemma~\ref{tc-cayley-quadratic}, the phase $q$ is constant
on each coset.  The integral of
$\sum_{\gamma\in\Gamma}g(z-A^{-1}x-\gamma)$ over
a fundamental domain is $\int_Vg$.
Absolute convergence follows from (TC38); hence both the
decomposition and this unfolding are justified.
We obtain
\[
 \int_{V/\Gamma}P(\delta^{-1}z,z)e^{-\pi H(z,z)}\,dz
       =\left(\sum_{\Lambda/A\Gamma}q(x)\right)\int_Vg(z)\,dz.
 \tag{TC39}
\]

Diagonalize the unitary map $\delta$ in an orthonormal basis.
The eigenvalues of $A$ are $\lambda_j=1-e^{-i\theta_j}$,
with $\operatorname{Re}\lambda_j>0$, since $\delta$ has no
fixed vectors.  Polar coordinates give, for such a $\lambda$,
\[
 \int_{\mathbb C}e^{-\pi\lambda|w|^2}\,dw
        =2\pi\int_0^\infty e^{-\pi\lambda t^2}t\,dt
        =\lambda^{-1}.
\]
Absolute convergence permits multiplication of these integrals,
so $\int_Vg=(\det A)^{-1}$.
Comparing (TC36) and (TC39) in the scalar identity
$P(\delta^{-1}z,w)=cI_{\Lambda,\delta\Lambda}P(\,\cdot\,,w)(z)$
now gives
\[
                c=\frac{\sum q(x)}{\sqrt d\,\det A}.
\]
Using $\sqrt N=\sqrt d\,|\det A|$ proves (TC31).

The normalized Gauss sum in (TC31) has eighth power one
by Theorem~\ref{tc-gauss-eight}.
Taking determinants in $A^*=-\delta A$ gives
\[
             \overline{\det A}=(-1)^r\det(\delta)\det A.
 \tag{TC40}
\]
Consequently $(|\det A|/\det A)^2=(-1)^r\det\delta$,
whose fourth power is $(\det\delta)^4$.
This proves (TC32).
\end{proof}

\subsection{Application to the actual fourth-power normalization}

Let $k=\mathbb Q(\sqrt{-D})$ with $D>0$, $D\equiv7\pmod8$,
and $\mathcal O_k=\mathbb Z[\tau]$,
$\tau=(1+i\sqrt D)/2$.
The same argument applies whenever the specified order is
$\mathbb Z[\tau]$; here $\mathcal O_k$ denotes the maximal
order used in the arithmetic construction.
Suppose $\Lambda\subset V$ is a full lattice stable under
$\mathcal O_k$ and self-dual for $E$.
Define on the whole complex space
\[
              \alpha_H(x)=
                    \exp\!\left(\pi i\frac{\sqrt D}{2}H(x,x)\right).
 \tag{TC41}
\]

\begin{lemma}\label{tc-norm-semicharacter}
The function (TC41) is a sign semicharacter on $\Lambda$.
Every isometry preserves this function, and it is multiplicative
under orthogonal direct sums.
On the standard lattice $\mathcal O_k^m$ with
$H_0(z,w)=(2/\sqrt D)\sum_jz_j\overline w_j$, it is
$(-1)^{\sum_jN_{k/\mathbb Q}(x_j)}$.
\end{lemma}
\begin{proof}
Put $Q(x)=(\sqrt D/2)H(x,x)$.
Since $\tau x\in\Lambda$,
$Q(x)=E(\tau x,x)\in\mathbb Z$.
For $x,y\in\Lambda$,
\[
 Q(x+y)-Q(x)-Q(y)+E(x,y)=2E(\tau x,y)\in2\mathbb Z.
 \tag{TC42}
\]
Subtracting $2E(x,y)\in2\mathbb Z$ gives (TC2).
The remaining assertions follow immediately from preservation
and additivity of the squared Hermitian norm.
\end{proof}

We also make explicit the nonprincipal elliptic base allowed
in the submitted construction.
Let
\[
 I=\mathbb Za+\mathbb Z\tau_b,\qquad
 \tau_b=\frac{b+i\sqrt D}{2},\qquad
 4ac=b^2+D,\qquad a>0\text{ odd},
 \tag{TC43}
\]
be an integral $\mathcal O_k$-ideal, and let $L=\lambda I$,
$\lambda\in\mathbb C^\times$.
Its Hermitian form is
$H_L(z,w)=2z\overline w/(a|\lambda|^2\sqrt D)$.

\begin{lemma}\label{tc-ideal-character}
The normalized translate by $\lambda a/2$ of the reduced odd
theta function of $L$ has semicharacter $\alpha_{H_L}$.
The lattice $L$ is self-dual for $\operatorname{Im}H_L$.
\end{lemma}
\begin{proof}
In the basis $\lambda a,\lambda\tau_b$, the alternating
pairing has entries $E(\lambda a,\lambda\tau_b)=-1$ and
$E(\lambda a,\lambda a)=E(\lambda\tau_b,\lambda\tau_b)=0$.
It is therefore unimodular.
For $x=\lambda(am+n\tau_b)$ the reduced odd theta
semicharacter is $(-1)^{m+n+mn}$.
Translation by $\lambda a/2$ multiplies it by
$e^{2\pi iE(x,\lambda a/2)}=(-1)^n$, yielding
$(-1)^{m+mn}$; the opposite sign in this exponential
gives the same value.
On the other hand
\[
 \frac{\sqrt D}{2}H_L(x,x)
       =\frac{N(am+n\tau_b)}a
       =am^2+bmn+cn^2.
 \tag{TC44}
\]
The relation in (TC43) makes $b$ odd, and then
$D\equiv7\pmod8$ and oddness of $a$ make $c$ even.
Thus (TC44) has parity $m+mn$, as required.
The translating half-period does not belong to $L$, so the
value used to normalize the theta function at zero is nonzero.
\end{proof}

Now put $\Lambda_0=L^{4r}$ with the product form $H_0$,
and let $\Theta_0$ be the product of these normalized
one-dimensional theta functions.
The case $I=\mathcal O_k$ and $\lambda=1$ is the base of
Section~\ref{ct-section}.  Lemma~\ref{tc-ideal-character}
handles the other ideal classes without changing the
analytic argument.
Let $g_0,g_1:V^4\to\mathbb C^{4r}$ be isometries from
$H^{\oplus4}$ to $H_0$, and put $M_i=g_i(\Lambda^4)$.
Assume that $M_0,M_1,\Lambda_0$ are commensurable and
$\det(g_1g_0^{-1})$ is algebraic.  These are satisfied
when the maps are obtained from isometries over $k$ in
the common scaled $k$-coordinates of the construction.
Use the same norm semicharacter $\alpha_{H_0}$ on all
three lattices.  They are self-dual by the isometries and
Lemma~\ref{tc-ideal-character}.
Define, as in Section~\ref{lm-section},
\[
          K_{g_i}(z)=(I_{\Lambda_0,M_i}\Theta_0)(g_i z).
 \tag{TC45}
\]

\begin{theorem}\label{tc-kernel-ratio}
Suppose $g_1-g_0$ is invertible.  Then there is a nonzero
algebraic number $c$ such that
\[
       K_{g_0}=cK_{g_1},\qquad
       c^8=\det(g_1g_0^{-1})^4.
 \tag{TC46}
\]
Thus the analytic identity required in the fourth tensor
criterion of Section~\ref{lm-section} holds with its exact
determinant exponent.
\end{theorem}
\begin{proof}
Put $\delta=g_1g_0^{-1}$ and $F=I_{\Lambda_0,M_0}\Theta_0$.
This is a nonzero theta function by
Lemma~\ref{tc-unitary-intertwiner}.
The hypothesis makes $\delta-1$ invertible.
The norm semicharacter is invariant under $\delta$ by
Lemma~\ref{tc-norm-semicharacter}, including the equality
(TC26).  Apply Theorem~\ref{tc-unitary-constant} on $M_0$
and then Theorem~\ref{tc-triple}:
\[
 F(\delta^{-1}w)
    =c_0 I_{M_0,M_1}I_{\Lambda_0,M_0}\Theta_0(w)
    =c_0\mu I_{\Lambda_0,M_1}\Theta_0(w),
 \quad c_0^8=(\det\delta)^4,\quad\mu^8=1.
\]
Evaluate at $w=g_1z$.  This gives (TC46) with $c=c_0\mu$.
Since $\det\delta$ is algebraic, that equation also proves
that $c$ is algebraic.  No algebraicity of the separate
complex Gaussian factors has been presumed.
\end{proof}

\begin{corollary}\label{tc-normalized-unit}
Use the reference isometry $\gamma_0=g_0$, and form
$\kappa$ as in (LM9)--(LM10), with the locally standard
maps carrying the given lattice onto the fixed target
$\Lambda_0$ in scaled $k$-coordinates.
If $p$ splits in $k$ and
$g_1(\Lambda^4\otimes\mathbb Z_p)
=\Lambda_0\otimes\mathbb Z_p$ in those coordinates, then,
for the same
embedding into $\mathbb C_p$ as in that theorem,
\[
                         |\iota(\kappa c)|_p=1.
 \tag{TC47}
\]
Moreover a nonzero theta function $F_\Lambda$ with
$F_\Lambda^{\otimes4}=\kappa K_{g_0}$ exists and is unique
up to a fourth root of unity.
\end{corollary}
\begin{proof}
The proof of Theorem~\ref{lm-kappa} applies to any fixed
target $\mathcal O_k$-lattice.  Indeed, two local maps
onto that target differ by an automorphism of its completed
lattice; its determinant has valuation zero at both primes.
The finite-support argument uses commensurability with the
fixed target in place of commensurability with
$\mathcal O_k^{4r}$, and every subsequent valuation
calculation is unchanged.
Equation (TC46) and Corollary~\ref{lm-change-scalar}
therefore give (TC47).  For existence, choose any nonzero element $\theta$
of $\mathcal T(\Lambda,\alpha_H)$ by
Lemma~\ref{tc-covolume}.
Both $\theta^{\otimes4}$ and $K_{g_0}$ are nonzero elements
of the one-dimensional theta space for $\Lambda^4$.
Thus $\kappa K_{g_0}=a\theta^{\otimes4}$ for $a\ne0$.
Choose a fourth root of $a$ and multiply $\theta$ by it.
Any two choices, and any two theta functions satisfying
the displayed normalization, differ by a fourth root of
unity, again by the one-dimensionality of the theta space.
\end{proof}

The determinant exponent in (TC46), the positive index
square roots in (TC10), and the invariant sign semicharacter
in (TC26) have all been retained.  The complex absolute
value of a theta normalization alone would not establish
(TC47).  The calculation above establishes the stronger
algebraic equality that supplies the finite-place bound.
It closes the analytic change-of-isometry input; identifying
the particular arithmetic column and its support is a
separate subsequent step.


\section{One ordinary period for the nonprincipal CM ideal bases}
\label{ci-section}

The source permits a nonprincipal elliptic lattice.  At every
prime where its moment theorem is applied it also assumes
$p\nmid N(I)$.  We verify the extension of
Section~\ref{ct-section} under this precise hypothesis, using a
prime-to-$p$ isogeny to keep the same period.  Choosing an unrelated
unit in $\mathbb C_p$ and making a binomial substitution with that
unit would not justify this step.

Keep $k=\mathbb Q(\sqrt{-D})$ and
$\mathcal O_k=\mathbb Z+\mathbb Z\tau$ as in
Section~\ref{tc-section}.  Throughout this section $p\geq5$
is a rational prime split in $k$; this is the scope of the
ordinary theta integrality theorem used here.
Let $I$ be the ideal (TC43),
so $[\mathcal O_k:I]=a$, $a$ is odd, and $p\nmid a$.
The source's curve ideals $J$ satisfy
$\gcd(N(J),2D)=1$, in particular $N(J)$ is odd.
Such an integral ideal $J$ can
first be written $J=gI$ with $g\in\mathbb Z_{>0}$ and
$I$ primitive of this form.  One obtains $g$ by taking the
content of a lattice basis; stability under $\tau$ shows
that the primitive Hermite normal form has second
$\tau$-coefficient one.
Then $N(J)=g^2a$, so $p\nmid N(J)$ implies $p\nmid a$
and $p\nmid g$.  The factor $g$ is included in the
scaling treated at the end of this section.
Put
\[
 H_I(z,w)=\frac{2}{a\sqrt D}z\overline w,\qquad
 E_I=\operatorname{Im}H_I,\qquad
 Q_I(z)=\frac{|z|^2}{a},\qquad h_I=\frac a2.
 \tag{CI1}
\]
Let $\theta_I$ be the reduced odd theta function translated
by $h_I$ and divided by its value at zero.  Lemma~\ref{tc-ideal-character}
gives its multiplier $(-1)^{Q_I(v)}$ for $v\in I$.
It is even and has its zero precisely at $h_I+I$, by the
same zero-count and uniqueness argument as in
Lemma~\ref{ct-even-identification}.  All translations below use
the Hermitian form $H_I$.

\subsection{Prime-to-$p$ isogenies and the common formal parameter}

Let $\mathscr E_0$ be the good ordinary model for
$\mathbb C/\mathcal O_k$ used in Section~\ref{ct-section}.
We choose its CM equation and differential over algebraic
numbers before completing at the selected place, as is
permitted in that classical CM input.  Further finite
extensions and the quotients below preserve this convention.
Write its period lattice relative to a good differential as
$c_0\mathcal O_k$, and normalize its multiplicative formal
parameter there.  If $\Omega_{p,0}$ is the associated unit
period, put
\[
                  \omega=\frac{\Omega_{p,0}}{\iota(c_0)}.
 \tag{CI2}
\]
The embedding $\iota$ is the fixed one from that section.

\begin{proposition}\label{ci-common-period}
There is a good ordinary model $\mathscr E_I$ for
$\mathbb C/I$ whose multiplicative formal parameter, pulled
back through the isogeny
\[
                    \psi:\mathbb C/I\longrightarrow
                         \mathbb C/\mathcal O_k,\qquad z\longmapsto z,
 \tag{CI3}
\]
has the same formal relation $z=\omega\log(1+T)$
in the unscaled complex variable.
If $\delta_I^{12}=g_2(I)^3-27g_3(I)^2$ and $\delta_0$
is the corresponding root for $\mathcal O_k$, then
\[
 \delta_I/\delta_0\in\overline{\mathbb Q}^{\,\times},
       \qquad |\iota(\delta_I/\delta_0)|_p=1.
 \tag{CI4}
\]
\end{proposition}
\begin{proof}
The inclusions $a\mathcal O_k\subset I\subset\mathcal O_k$
give (CI3) and its degree-$a$ dual, which in unscaled complex
coordinates sends $z$ to $az$.
Construct the latter as a quotient of $\mathscr E_0$ by its
subgroup $a^{-1}I/\mathcal O_k$.
After a finite extension this subgroup extends to a finite
\'etale subgroup of $\mathscr E_0[a]$: since $p\nmid a$,
multiplication by $a$ has invertible differential, and its
kernel on the smooth elliptic scheme is finite \'etale.
The translation action is free.  The finite locally free
quotient theorem \cite{stacks-07S7} makes the quotient a scheme
and its quotient map a torsor.  In this case it is finite
\'etale; smoothness and properness descend through that
cover.  Its geometrically connected genus-one fibers and
the image of the zero section make the quotient an elliptic
scheme.  Thus it is the required good model $\mathscr E_I$.
The dual isogeny (CI3) extends as well.  Both maps are
isomorphisms of formal groups: their composite is $[a]$,
whose formal inverse is integral, and each has a unit
linear term.  Ordinary reduction is preserved, for example
because they induce an isomorphism on the $p$-divisible
formal groups.

Write the period lattice of $\mathscr E_I$, relative to a
good differential, as $c_I I$.
The differential multiplier of $\psi$ in these coordinates
is $u=c_0/c_I$.  It is an algebraic number and a $p$-adic
unit, since $\psi$ is an algebraic isogeny between the
chosen algebraic models and its map on formal groups is an
integral isomorphism.
If $\lambda_I,\lambda_0$ are their normalized formal
logarithms, differentiation of the group homomorphism gives
\[
               \lambda_0(\psi(t))=\iota(u)\lambda_I(t).
 \tag{CI5}
\]
There is no extra constant because both sides vanish at zero.
Pull back the chosen multiplicative coordinate of
$\mathscr E_0$ along $\psi$.  Its period on $\mathscr E_I$
is
$\Omega_{p,I}=\Omega_{p,0}/\iota(u)
=\iota(c_I/c_0)\Omega_{p,0}$, again a unit.
In the variable $z$ for the lattice $I$, it therefore gives
\[
             \frac{\Omega_{p,I}}{\iota(c_I)}
                       =\frac{\Omega_{p,0}}{\iota(c_0)}=\omega.
 \tag{CI6}
\]
This proves the asserted equality of periods, rather than
merely equality of their absolute values.

Let $\Delta_0,\Delta_I$ be the discriminants of the two good
models.  The invariant-differential calculation (CT11) gives
$(\delta_I/c_I)^{12}=\Delta_I$ and
$(\delta_0/c_0)^{12}=\Delta_0$, after the fixed algebraic
identifications.  Their ratio shows that
$(\delta_I/\delta_0)^{12}
=(c_I/c_0)^{12}\Delta_I/\Delta_0$ is algebraic.
All its factors have $p$-adic absolute value one.
Taking a twelfth root proves both assertions in (CI4).
\end{proof}

\begin{theorem}\label{ci-integral-translates}
For every $t\in I\otimes\mathbb Q$ of order prime to $p$
modulo $I$,
\[
                     \widehat{U_t\theta_I}
                           \bigl(\omega\log(1+T)\bigr)
                         \in\mathcal O_{\mathbb C_p}[[T]].
 \tag{CI7}
\]
The series is a unit unless $t+h_I\in I$.
All finite-combination norm identities of
Theorem~\ref{ct-combination-norm} hold with $\theta_I$ and
this same $\omega$.
In particular every prime-to-$p$ translate has integral
period-scaled moments in every finite family of directions
from $I$.
\end{theorem}
\begin{proof}
Apply the reduced-theta integrality theorem cited in
Section~\ref{ct-section} to the good ordinary curve
$\mathscr E_I$ and to the nonzero two-torsion point $h_I$.
Its reduced-theta value at zero is a unit.  Translation
by $t$ composes with that translation up to the root of
unity $e^{-\pi iE_I(t,h_I)}$.
The sum $t+h_I$ still has prime-to-$p$ order.
The cited theorem therefore gives integrality, with unit
constant term except at the stated zero.
The inverse of the multiplicative coordinate constructed
in Proposition~\ref{ci-common-period} is integral with
unit linear coefficient, and (CI6) identifies its variable
with the one in (CI7).

The finite-combination assertion follows from the proof of
Theorem~\ref{ct-combination-norm} on this elliptic scheme:
choose an even common denominator $N$ prime to $p$, use the
line bundle $[N]^*\mathcal O([h_I])^{\otimes m}$, and use
the integral trivialization near zero given by its canonical
section.  Formula (CI7) makes the corresponding reduced
expansion an integral unit.  The saturated expansion lemma
identifies its coefficient norm with the integral
global-section norm.  Translation by $t/N$ has a power
equal to a root of unity and is therefore an integral
automorphism by the vertical-divisor proof of
Lemma~\ref{ct-torsion-isometry}.
The exact identity (CT18) uses only the Hermitian Weyl
formula, so it holds for $H_I$ as well.
Finally $T\mapsto(1+T)^N-1$ and its inverse are integral,
and removing that substitution gives the asserted norm
identity for arbitrary $\mathbb C_p$ coefficients.
Every element of $I$ belongs to $\mathcal O_k$, so its
image at this split prime lies in $\mathbb Z_p$.
Lemma~\ref{lm-directions} now gives the assertion for
arbitrary finite direction families.
\end{proof}

\subsection{The selected torsion character and its Weyl phase}

The chosen parameter is compatible with the required
character, not just with the invariant differential.

\begin{proposition}\label{ci-character}
The formal $p^n$-torsion of $\mathbb C/I$ consists of
\[
 y=\frac{am+b'\tau_b}{p^n}\pmod I,\qquad
 \iota(am+b'\overline{\tau_b})\in p^n\mathbb Z_p.
 \tag{CI8}
\]
In the parameter of Proposition~\ref{ci-common-period},
\[
                1+T(y)=\iota(e^{2\pi ib'/p^n})
                         =\iota(e^{-2\pi iE_I(a,y)}).
 \tag{CI9}
\]
If $st\in I$, $p\nmid s$, and
$y=s\nu/p^n$ with $\nu\in I$ and
$\iota(\bar\nu)\in p^{2n}\mathbb Z_p$, then
\[
 \widehat{e^{\pi iQ_I(y)}U_y(U_t\theta_I)}
                       \bigl(\omega\log(1+T)\bigr)
  =S_\zeta\!\left(\widehat{U_t\theta_I}
                       \bigl(\omega\log(1+T)\bigr)\right),
 \quad
 \zeta=\iota(e^{-2\pi iE_I(a,y)}).
 \tag{CI10}
\]
\end{proposition}
\begin{proof}
The isogeny (CI3) identifies the formal groups and their
$p$-power torsion.  Since
$\tau_b=\tau+(b-1)/2$, the image of the numerator in
(CI8), in the basis $1,\tau$, is
$am+b'(b-1)/2+b'\tau$.
Its conjugate-divisibility condition and its coefficient
of $\tau$ are exactly those in
Proposition~\ref{ct-character-normalization}.
This proves (CI8)--(CI9), including the simultaneous
normalization for every $n$.
Alternatively, $a$ is a unit modulo $p^n$, so (CI8)
has exactly one residue of $m$ for every $b'$; this
also proves directly that its character is an isomorphism.

We record the arithmetic changes in the proof of
Theorem~\ref{ct-power-translation}; they are the only
changes needed for (CI10).
Because $N(\nu)/a\in\mathbb Z$ and $p\nmid a$, the strong
conjugate divisibility gives
$Q_I(\nu/p^n)\in\mathbb Z$.
Choose $\epsilon\in2\mathcal O_k$ with
$\epsilon\equiv1\pmod{\mathfrak p^n}$ and
$\epsilon\equiv0\pmod{\bar{\mathfrak p}^{\,n}}$.
Then
\[
  (\epsilon-1)\nu/p^n\in I,\qquad
                  \bar\epsilon\nu/p^n\in2I.
 \tag{CI11}
\]
At the two primes over $p$ this follows from the two
congruences and the strong divisibility, using
$I\otimes\mathbb Z_p=\mathcal O_k\otimes\mathbb Z_p$.
At every other prime the denominator is a unit;
the factor two follows because $\epsilon$ is even and
$p$ is odd.
The rational elliptic function
$\theta_I(\epsilon z)/\theta_I(\bar\epsilon z)$ has no
zero or pole in the formal subgroup: all such points
map to the nonzero two-torsion point $h_I$.
Applying the formal translation identity (CT22) and
cancelling its unit denominator, exactly as in
(CT28)--(CT30), gives
$S_\zeta\widehat\theta_I
=\widehat{U_{\epsilon y}\theta_I}$ for $y=\nu/p^n$.

Write $\epsilon=A+B\tau$ and $l=(\epsilon-1)y$.
The two phase identities are now
\[
 E_I(y,l)=-B Q_I(y),\qquad
 Q_I(l)=(N(\epsilon)-2A-B+1)Q_I(y).
 \tag{CI12}
\]
Consequently
$e^{\pi iE_I(y,l)-\pi iQ_I(l)}
=e^{\pi iQ_I(y)}$, since their exponent difference
is $(2A-2-N(\epsilon))Q_I(y)$, an even integer.
This proves the base-theta assertion.
For $st\in I$, apply (CT22) to the elliptic rational
function
$z\mapsto(U_t\theta_I)(sz)/\theta_I(sz)$ and to
$\nu/p^n$.  Its poles avoid the formal subgroup.
The invertible substitution $[s]$ gives the ratio
identity (CT33), now for $I$.
Multiplying by the base-theta identity proves (CI10).
The operator $S_\zeta$ and its inverse preserve the
coefficient norm by (CT21).
\end{proof}

\subsection{Scaling the lattice and the actual restriction on primes}

For $L=\lambda I$ and its normalized theta function,
uniqueness gives $\theta_L(\lambda z)=\theta_I(z)$, and
the Hermitian forms give
\[
          (U_{\lambda t}^{L}\theta_L)(\lambda z)
                           =(U_t^{I}\theta_I)(z).
 \tag{CI13}
\]
Thus a direction $v_j=\lambda u_j\in L$ satisfies, identically
before transport,
\[
 (U_{\lambda t}^{L}\theta_L)
       \left(\omega\sum_j\lambda u_j\log(1+T_j)\right)
 =
 (U_t^{I}\theta_I)
       \left(\omega\sum_j u_j\log(1+T_j)\right).
 \tag{CI14}
\]
Here the displayed identity is understood first with an
indeterminate scalar $w$ in place of $\omega$, and then
after transporting its formal coefficients and setting
$w=\omega$.  This is why the same $\omega$ supplies the moments on
$L$: the complex scaling cancels in the function and
its directions.  No assertion that $|\iota(\lambda)|_p=1$
has been made or used.

Writing $\omega=\Omega/\iota(\delta_0)$, the number $\Omega$
is the single unit already fixed in Section~\ref{ct-section}.
Equation (CI4) implies that
$\Omega\,a\,\iota(\delta_I/\delta_0)$ is a unit as well.
This is the period adjustment that occurs when the source
changes from lattice directions to its arithmetic
trace-pairing coordinates.
The results here require $p\ge5$, $p$ split, and $p\nmid a$.
The inspected moment and local-normalization statements
explicitly impose the last condition; no claim for $p\mid a$
is needed for them.  Without it, the character in (CI9)
need not be a generator and the particular prime-to-$p$
isogeny argument above does not apply.


\section{Rational isometries with the required local theta normalization}
\label{ts-section}

We now produce the isometry required in
Theorem~\ref{tc-kernel-ratio}, with both its local lattice
condition and its absence of fixed vectors.  This removes
the corresponding existence hypothesis from the ordinary
moment criterion.  We give the four-square construction as
well, so that no existence of a rational standard form is
being presumed.

Let $k=\mathbb Q(\sqrt{-D})$, $\mathcal O_k$ and $I$ be
as in Sections~\ref{tc-section} and~\ref{ci-section}.
Write $\sigma=\sqrt{-D}$ for the element with positive
imaginary part in the fixed complex embedding.
Let $W$ be a $k$-vector space of dimension $r$, with
positive definite Hermitian form $h$, linear in its first
variable.  Its associated alternating rational form is
\[
       E_h(x,y)=
           \operatorname{Tr}_{k/\mathbb Q}\!\left(\frac{h(x,y)}{\sigma}\right)
          =\frac2{\sqrt D}\operatorname{Im}h(x,y).
 \tag{TS1}
\]
Let $M\subset W$ be a full $\mathcal O_k$-lattice,
self-dual for $E_h$.

\subsection{The four-square isometry}

\begin{lemma}\label{ts-rational-four-squares}
Every positive rational number is a sum of four rational
squares.
\end{lemma}
\begin{proof}
We include the elementary integer descent.
For an odd prime $\ell$, the pigeonhole argument in
Lemma~\ref{tc-four-squares-congruence} gives integers
$x,y$, chosen with $|x|,|y|\le(\ell-1)/2$, such that
$x^2+y^2+1\equiv0\pmod\ell$.
Thus some $m\ell$ is a sum of four squares with
$0<m<\ell$.
Choose the smallest positive $m$ for which
$m\ell=\sum_{i=1}^4a_i^2$.
If $m>1$, choose integers $b_i\equiv a_i\pmod m$
with $|b_i|\le m/2$.  Then
$\sum_i b_i^2=mm'$ with $0\le m'\le m$.
If $m'=0$, all $a_i$ are divisible by $m$, so
$m\mid\ell$, contradicting $1<m<\ell$.
If $m'=m$, equality in the square bound makes $m$
even and every $a_i=(m/2)(2c_i+1)$.
The sum of four odd squares is divisible by four,
so $m^2\mid\sum_i a_i^2=m\ell$, the same contradiction.
Hence $0<m'<m$.

Regard $a=(a_1,a_2,a_3,a_4)$ and $b=(b_1,b_2,b_3,b_4)$
as integer quaternions.
The four coordinates of $a\bar b$ are divisible by $m$:
modulo $m$, this product is $b\bar b=\sum_i b_i^2=0$.
The multiplicative norm identity, which is the matrix
identity (TC18), now writes
\[
                    m'\ell
                      =\left|\,a\bar b/m\,\right|^2
\]
as a sum of four integer squares, contradicting
minimality of $m$.  Therefore $m=1$.
The prime $2=1^2+1^2$ is immediate.
Multiplying quaternion norms proves the assertion for
every positive integer, by prime factorization.
Finally if $q=A/B>0$ with positive integers $A,B$,
write $AB$ as four integer squares and divide the
four coordinates by $B$.
\end{proof}

\begin{proposition}\label{ts-global-isometry}
There is a $k$-linear isometry
\[
        \gamma_0:(W^{\oplus4},h^{\oplus4})
                           \longrightarrow (k^{4r},h_{\rm st}),
 \qquad h_{\rm st}(x,y)=\sum_jx_j\overline y_j.
 \tag{TS2}
\]
\end{proposition}
\begin{proof}
Hermitian Gram--Schmidt over $k$ diagonalizes $h$ with
positive diagonal entries $q_j\in k\cap\mathbb R=\mathbb Q$.
At each step the orthogonal projection coefficients lie in
$k$, and positivity makes every pivot nonzero.
By Lemma~\ref{ts-rational-four-squares}, write
$q_j=a_j^2+b_j^2+c_j^2+d_j^2$ with rational coordinates.
The quaternion matrix (TC18) for these four coordinates
satisfies $T_j^*T_j=q_jI_4$, since its entries are rational.
After reordering the four copies of the diagonal basis,
the block diagonal matrix with blocks $T_j$ gives (TS2).
\end{proof}

\subsection{Local standard form at a split prime}

\begin{lemma}\label{ts-local-isometry}
If $p$ splits in $k$, there is a $k\otimes\mathbb Q_p$-linear
isometry $\gamma_p$ from $W\otimes\mathbb Q_p$ to the
standard Hermitian space carrying
\[
                     M\otimes\mathbb Z_p
                      \quad\hbox{onto}\quad
                     \mathcal O_k^r\otimes\mathbb Z_p.
 \tag{TS3}
\]
\end{lemma}
\begin{proof}
The split algebra is $\mathbb Q_p\times\mathbb Q_p$,
with conjugation interchanging the factors, and
$\mathcal O_k\otimes\mathbb Z_p=\mathbb Z_p\times\mathbb Z_p$.
Its two idempotents decompose $M\otimes\mathbb Z_p$
as $M_+\oplus M_-$, each free of rank $r$ over
$\mathbb Z_p$.
Hermitian sesquilinearity makes $h$ zero on each one
of the two summands paired with itself.
Write $\sigma=(s,-s)$; the number $s$ is a $p$-adic unit,
since a split prime is unramified.
For $x\in M_+$ and $y\in M_-$, if $h(x,y)=(u,0)$,
then $E_h(x,y)=u/s$.
The perfect alternating pairing on $M\otimes\mathbb Z_p$
therefore says that $h$ gives a perfect integral pairing
between $M_+$ and $M_-$.
Choose a basis of the former and its $h$-dual basis in
the latter.  Sending these to the corresponding standard
bases gives (TS3).  The conjugate cross-pairing is then
correct by the Hermitian identity.
\end{proof}

\begin{theorem}\label{ts-approximation}
Let $\mathcal S$ be a nonempty finite set of split primes
of $k$.  There is a $k$-linear isometry
$\gamma:W^{\oplus4}\to k^{4r}$ such that
\[
 \gamma(M^{\oplus4}\otimes\mathbb Z_p)
          =\mathcal O_k^{4r}\otimes\mathbb Z_p
       \quad(p\in\mathcal S),\qquad
                   \gamma-\gamma_0\ \hbox{is invertible}.
 \tag{TS4}
\]
\end{theorem}
\begin{proof}
Put $m=4r$ and use Lemma~\ref{ts-local-isometry} on
$M^{\oplus4}$ to choose each local $\gamma_p$.
Then $\delta_p=\gamma_p\gamma_0^{-1}$ is unitary on
the standard space.
At a split prime its first component can be any
$G_p\in\operatorname{GL}_m(\mathbb Q_p)$ and its other
component is $(G_p^{\mathsf T})^{-1}$.
The condition that $\delta_p\gamma_0$ carry the specified
lattice onto the standard one is open.
We may perturb $G_p$ within this open set so that
$\det(G_p-1)\det(G_p+1)\ne0$:
replace $G_p$ by $G_p+tI$ for arbitrarily small $t$,
avoiding the finitely many roots of the two determinant
polynomials.

Its Cayley inverse
\[
                     X_p=(\delta_p-1)(\delta_p+1)^{-1}
 \tag{TS5}
\]
is skew-Hermitian and invertible, and $1-X_p$ is
invertible.
The skew-Hermitian matrices over $k$ form a finite-dimensional
$\mathbb Q$-vector space, defined by the rational linear
equation $X^*=-X$.  A rational basis remains a basis after
extension to each $\mathbb Q_p$.
The Chinese remainder theorem, after clearing a finite
set of denominators, approximates the coefficients of
all the $X_p$ simultaneously by coefficients in $\mathbb Q$.
Choose a resulting matrix $X$ sufficiently close that
the determinants of $X$ and $1-X$ stay nonzero and that
\[
                       \delta=(1+X)(1-X)^{-1}
 \tag{TS6}
\]
has the required local lattice property.
To justify the approximation assertion directly, multiply
the finitely many prescribed rational $p$-adic coordinates
by an integer making them integral at the selected primes,
choose their required residues modulo sufficiently high
powers, solve the simultaneous integer congruences, and
divide by that integer.

Since $X^*=-X$, multiplication in (TS6) gives
$\delta^*\delta=1$.
Also $\delta-1=2X(1-X)^{-1}$ is invertible.
The map $\gamma=\delta\gamma_0$ consequently satisfies
both assertions in (TS4).
The zero-dimensional case is immediate with the
usual empty-determinant convention.
\end{proof}

\subsection{Ideal twists and the same determinant correction}

The trace lattices in the source are twisted by the
elliptic ideal.  This twist preserves the hypotheses
above and the local determinant comparison at its
allowed primes.

\begin{lemma}\label{ts-ideal-twist}
The lattice $MI$ is self-dual for $E_h/a$.
If $p\nmid a$, then
\[
                 MI\otimes\mathbb Z_p=M\otimes\mathbb Z_p,
             \qquad I\otimes\mathbb Z_p
                         =\mathcal O_k\otimes\mathbb Z_p.
 \tag{TS7}
\]
Every isometry in (TS2) or (TS4) remains an isometry
from $(h/a)^{\oplus4}$ to $h_{\rm st}/a$.
\end{lemma}
\begin{proof}
The inverse different of $\mathcal O_k$ is
$\sigma^{-1}\mathcal O_k$.
Thus an $\mathcal O_k$-stable lattice self-dual for
the trace form (TS1) is exactly a self-dual Hermitian
lattice: testing $E_h(x,\alpha M)$ for every
$\alpha\in\mathcal O_k$ is equivalent to
$h(x,M)\subset\mathcal O_k$.
This is also seen directly from the trace matrix in
the basis $1,\tau$, whose discriminant is $-D$.
For an invertible ideal,
$(MI)^\#=M^\#\bar I^{-1}$ with respect to $h$,
as follows locally by writing $I$ as a principal ideal.
In a quadratic field $I\bar I=(a)$.
The dual for $h/a$ is consequently
$aM^\#\bar I^{-1}=MI$.
The equivalence with alternating duality proves the
first assertion.
The finite quotient $\mathcal O_k/I$ has order $a$,
which proves (TS7) when $p\nmid a$.
Multiplying both Hermitian forms by $1/a$ does not
change their isometries.
\end{proof}

Let $\kappa$ be the determinant correction of (LM9)--(LM10)
formed from the untwisted lattice $M^{\oplus4}$, reference
isometry $\gamma_0$, and standard target $\mathcal O_k^{4r}$.
For the complex theta construction use the ideal-twisted
lattice $MI$, the Hermitian form $H=(2/(a\sqrt D))h$,
and the standard elliptic lattice $I$.
One may multiply both lattices by the same complex scalar
$\lambda$, replacing each Hermitian form by the transported
form $H_\lambda(\lambda x,\lambda y)=H(x,y)$;
the following conclusion is unchanged.

\begin{theorem}\label{ts-integral-theta}
Fix a split prime $p\ge5$ with $p\nmid a$.
Let $F_M$ be a theta function normalized by
\[
                         F_M^{\otimes4}=\kappa K_{\gamma_0}.
 \tag{TS8}
\]
Such a nonzero theta function exists and is unique up to a
fourth root of unity.
For the single ordinary period $\omega$ of
Proposition~\ref{ci-common-period}, every translate
$U_tF_M$ of order prime to $p$ modulo $MI$
has integral period-scaled moments on $MI$.
The statement also holds on the simultaneously scaled
lattices $\lambda MI$ and $\lambda I$ with the same
period in the lattice directions.
\end{theorem}
\begin{proof}
Apply Theorem~\ref{ts-approximation} with $\mathcal S=\{p\}$.
By (TS7), its isometry $\gamma$ is locally standard for
the ideal-twisted lattices too:
$\gamma((MI)^{\oplus4}\otimes\mathbb Z_p)
=I^{4r}\otimes\mathbb Z_p$.
Their intersection therefore has index prime to $p$.
For finite-index lattices this follows by tensoring the
finite quotient with $\mathbb Z_p$: it vanishes exactly
when the quotient's order is prime to $p$.
The determinant comparison in Theorem~\ref{lm-kappa}
uses the untwisted local equality in (TS4), and gives
\[
                 |\iota(\kappa)|_p^2
                 |\iota(\det(\gamma\gamma_0^{-1}))|_p=1.
 \tag{TS9}
\]
The analytic theorem applies to the ideal-twisted lattices
by Lemma~\ref{ts-ideal-twist}.  Their common norm
semicharacter is exactly the elliptic one by
Lemma~\ref{tc-ideal-character}.
The invertibility in (TS4) gives
$K_{\gamma_0}=cK_\gamma$ with
$c^8=\det(\gamma\gamma_0^{-1})^4$.
Equations (TS9) and this eighth-power equality give
$|\iota(\kappa c)|_p=1$.
Existence and uniqueness in (TS8) follow from the
theta-space argument in Corollary~\ref{tc-normalized-unit}.

By Theorem~\ref{ci-integral-translates}, the product
elliptic theta function and all its prime-to-$p$
translates have integral moments on $I^{4r}$.
The finite intertwiner of index prime to $p$ preserves
this property by Proposition~\ref{lm-finite-intertwiner}.
Pullback by $\gamma$ identifies the direction lattices
exactly.  Multiplying by the unit $\iota(\kappa c)$
therefore gives integral moments for every fourth
tensor power $(U_tF_M)^{\otimes4}$.
The tensor-root theorem, Theorem~\ref{ot-tensor-root},
gives the assertion for $U_tF_M$ itself.
For the scaled lattices, (CI13)--(CI14) identify the
elliptic expansions without any assumption on the
absolute value of $\iota(\lambda)$; the same finite
operator and pullback proof then applies.
\end{proof}

This theorem supplies the ordinary theta moments for the
normalized self-dual lattice vector itself, with the same
determinant correction as the untwisted trace lattice and
for every split prime $p\geq5$ prime to the elliptic
ideal norm.  This supplies all sufficiently large
selected split primes.  The source's generic moment
statement also permits $p=3$; that case is outside
the Bannai--Kobayashi input used in this reconstruction.
It does not yet identify the more general arithmetic
column, which includes further finite lattice moves,
local support conditions and incidence maps.


\section{Exact equivariance of the arithmetic theta family}
\label{te-section}

An eighth-power identity does not by itself give an identity of
theta functions with the prescribed scalar.  The arithmetic
hypotheses supply the additional information needed here.  We
prove the exact identities, including the determinant characters.
The argument uses unramifiedness of the CM extension at every
finite place and the even dimensions in the construction.

\subsection{The trace lattices and their semicharacter}

Let $k=\mathbb Q(\sqrt{-D})$, $\sigma=\sqrt{-D}=i\sqrt D$,
and $\tau=(1+\sigma)/2$, with
$\mathcal O_k=\mathbb Z[\tau]$ and $D\equiv7\pmod8$.
Let $E/F$ be a CM quadratic extension unramified at every
finite place, with $k\subset E$ and $F$ its maximal real
subfield.  Write $\mathcal O=\mathcal O_E$ and let $\Sigma$
be the embeddings of $E$ into $\mathbb C$ extending the
specified embedding of $k$; put $d=|\Sigma|$.
Choose a fractional $\mathcal O$-ideal $\mathfrak a$ with
\[
          \mathfrak a\overline{\mathfrak a}
                         =\mathfrak D_{E/k}^{-1}.
 \tag{TE1}
\]
This is the ideal denoted by the inverse of the chosen half
different in the construction.  Its existence is part of the
arithmetic datum; the assertion below retains the exact equality
(TE1).  Let $I$ be the elliptic ideal of norm $a$, and let
$\lambda\in\mathbb C^\times$ be its common complex scaling.

For unimodular Hermitian $\mathcal O$-lattices
$L\subset E^n$, $M\subset E^{n'}$, define
\[
 \begin{split}
 W_{n,n'}&=\prod_{\eta\in\Sigma}
                      \operatorname{Mat}_{n'\times n}(\mathbb C),\\
 \Lambda(L,M)&=\{(\lambda\eta(X))_{\eta\in\Sigma}:
             X\in I\mathfrak a\operatorname{Hom}_{\mathcal O}(L,M)\},\\
 H(X,Y)&=\frac{2}{a|\lambda|^2\sqrt D}
                  \sum_{\eta\in\Sigma}\operatorname{tr}(X_\eta Y_\eta^*),
       \qquad E_H=\operatorname{Im}H.
 \end{split}
 \tag{TE2}
\]
The symbol $E_H$ denotes the alternating form; the unadorned
$E$ continues to denote the number field.  The standard
Hermitian forms on the two vector spaces are used in (TE2).
Put $L_0=\mathcal O^n$ and $M_0=\mathcal O^{n'}$.

\begin{lemma}\label{te-trace-lattice}
Every lattice in (TE2) is self-dual for $E_H$.
All these lattices are commensurable and stable under
$\mathcal O$.  On each of them
\[
       Q(x)=\frac{\sqrt D}{2}H(x,x)\in\mathbb Z,
                    \qquad \alpha(x)=e^{\pi iQ(x)}\in\{1,-1\}
 \tag{TE3}
\]
is the same sign semicharacter.
\end{lemma}
\begin{proof}
The Hermitian trace dual of
$\operatorname{Hom}_{\mathcal O}(L,M)$ for the $E$-valued
form $\operatorname{tr}(XY^*)$ is itself.  To check this,
localize at a prime, use a basis of $L$ and a basis of $M$,
and use their Hermitian dual bases at the conjugate prime.
The matrix units and their duals give precisely the asserted
dual module.  This argument does not assume that the global
projective modules have bases.

For the $k$-valued trace form, taking a dual inserts
$\mathfrak D_{E/k}^{-1}$.  Thus the dual of
$I\mathfrak a\operatorname{Hom}(L,M)$ for that form divided
by $a$ is
\[
 a\overline I^{-1}\overline{\mathfrak a}^{-1}
       \mathfrak D_{E/k}^{-1}\operatorname{Hom}(L,M)
             =I\mathfrak a\operatorname{Hom}(L,M).
\]
Here $I\overline I=(a)$ and (TE1) were used.
The inverse different of $k/\mathbb Q$ is
$\sigma^{-1}\mathcal O_k$.  Since
\[
 \operatorname{Tr}_{k/\mathbb Q}(t/\sigma)
                              =2\operatorname{Im}(t)/\sqrt D,
\]
passing to the alternating form in (TE2) gives exactly
the same dual lattice.  The factors $\lambda$ cancel in
the numerator and denominator.  This proves self-duality.
Full lattices in the same rational vector space are
commensurable, and stability under $\mathcal O$ is immediate.
Finally $Q(x)=E_H(\tau x,x)$, so integrality and the
semicharacter law follow from (TC42).
\end{proof}

Use the normalized averaging operators $I_{\Lambda,\Lambda'}$
of (TC10), with positive square roots of their indices.
No change in that normalization will be made below.

\subsection{A common local hyperbolic decomposition}

\begin{lemma}\label{te-local-pair}
Let $A$ be a complete discrete valuation ring of
characteristic zero with finite residue field, and let $B$
be its unramified quadratic extension, with involution
$x\mapsto\bar x$.  Let $V$ be an even-dimensional
nondegenerate Hermitian space over $\operatorname{Frac}(B)$.
If $M,N$ are unimodular $B$-lattices in $V$, there is a
decomposition into hyperbolic planes with vectors $v_j,u_j$
and integers $a_j\geq0$ such that
\[
 \begin{split}
 M&=\bigoplus_j(Bv_j\oplus Bu_j),\\
 N&=\bigoplus_j(B\varpi^{-a_j}v_j\oplus B\varpi^{a_j}u_j),\\
 h(v_i,v_j)&=h(u_i,u_j)=0,\qquad h(v_i,u_j)=\delta_{ij}.
 \end{split}
 \tag{TE4}
\]
In particular, an idempotent $e$ preserves both lattices
and satisfies $e^*=1-e$.
For a split quadratic algebra $B=A\times A$, an idempotent
with the same properties exists in every dimension.
\end{lemma}
\begin{proof}
We use the convention that $h$ is conjugate-linear in its
first variable.  Trace and norm on units are surjective by
Lemma~\ref{acl-local-norm}, including in residue
characteristic two.

First a single unimodular even-dimensional lattice is a
sum of hyperbolic planes.  Proposition~\ref{acl-local-basis}
gives an orthonormal basis.  In any pair of its coordinates
choose $t\in B^\times$ with $t\bar t=-1$; the vector
$(1,t)$ is primitive and isotropic.  Unimodularity gives
a vector $z$ paired with it to one.  Replacing $z$ by
$z-c(1,t)$, where $c+\bar c=h(z,z)$, makes $z$ isotropic
without changing this pairing.  The resulting hyperbolic
plane splits integrally, and induction proves the assertion.

For two lattices, let $a\geq0$ be the least integer with
$\varpi^aN\subset M$.  Duality also gives
$\varpi^aM\subset N$.  If $a=0$, the lattices are equal
and the preceding argument applies.  If $a>0$, choose
$w\in\varpi^aN$ primitive in $M$; otherwise $a$ was not
minimal.  Its norm lies in $\varpi^{2a}A$.
Choose $z\in M$ with $h(w,z)=1$ and choose $c\in B$
whose trace is a unit.  The equation
\[
 h(w+\varpi^{2a}ctz,w+\varpi^{2a}ctz)=0
                    \qquad(t\in A)
 \tag{TE5}
\]
divided by $\varpi^{2a}$ has the form
$b+t\operatorname{Tr}(c)+\varpi^{2a}N(c)h(z,z)t^2=0$.
It has a root modulo $\varpi$, and its derivative is a
unit there.  Successively correcting the root modulo
$\varpi^r$ and taking its limit gives a root in $A$.
This also supplies the required elementary form of
Hensel's argument.

The resulting vector $v$ is primitive in $M$, is isotropic,
and still belongs to $\varpi^aN$, because
$\varpi^{2a}M\subset\varpi^aN$.
Choose $u\in M$ with $h(v,u)=1$ and make it isotropic
by subtracting $cv$ with $c+\bar c=h(u,u)$.
The vectors $v,u$ split a unimodular hyperbolic plane
off $M$.  The vectors $\varpi^{-a}v,\varpi^a u$
belong to $N$, are isotropic, and have pairing one;
they split the same plane off $N$.
The two orthogonal complements are again unimodular
lattices of even rank.  Induction gives (TE4).
Projection onto the span of the $v_j$ has the stated
properties.  In the split case use multiplication by
$(1,0)$; its adjoint is multiplication by $(0,1)$.
\end{proof}

\begin{proposition}\label{te-projector}
Let $M,N$ be unimodular $\mathcal O$-lattices of even
rank in one Hermitian $E$-space.  For every positive
rational integer $b$ there is an $E$-linear map $e$
preserving both lattices such that each of
\[
 e^2-e,\quad e+e^*-1,\quad e^*e,\quad ee^*,\quad
                 (1-e)^*(1-e),\quad(1-e)(1-e)^*
 \tag{TE6}
\]
maps $M+N$ into $b(M\cap N)$.
\end{proposition}
\begin{proof}
Put $\mathscr R=\operatorname{End}_{\mathcal O}(M)
\cap\operatorname{End}_{\mathcal O}(N)$.  It is a full
finitely generated $\mathbb Z$-lattice in
$\operatorname{End}_E(V)$ and a ring.  It is stable under
adjoint, since $M$ and $N$ are self-dual.
Choose a positive integer $c$ with
$c(M+N)\subset b(M\cap N)$.
At each finite place of $F$ above a divisor of $c$,
Lemma~\ref{te-local-pair} supplies an exact idempotent
preserving both completed lattices and satisfying
$e_v^*=1-e_v$.  All polynomials in (TE6) vanish at $e_v$.

Completion commutes with this intersection of two full
lattices: after clearing denominators it is the kernel
of a homomorphism to a finite module, and completion is
exact on finitely generated modules over $\mathbb Z$.
Thus these local maps belong to the corresponding
completions of $\mathscr R$.
The Chinese remainder theorem on the finite free
$\mathbb Z$-module $\mathscr R$ gives
$e\in\mathscr R$ congruent to all of them modulo
$c\mathscr R$ at the primes dividing $c$.
Every expression in (TE6) then lies in $c\mathscr R$;
membership is checked at these primes, and elsewhere
$c$ is a unit.  Its image of $M+N$ lies in
$c(M+N)\subset b(M\cap N)$, as claimed.
\end{proof}

\subsection{The finite quadratic sum has positive sign}

For three commensurable self-dual lattices
$\Lambda_1,\Lambda_2,\Lambda_3$, recall the finite group
and its quadratic function from the proof of
Theorem~\ref{tc-triple}:
\[
 \begin{split}
 P&=\{(z,y)\in\Lambda_3\times\Lambda_2:z+y\in\Lambda_1\},\\
 T&=\{(t,u-t):t\in\Lambda_2\cap\Lambda_3,
                              \ u\in\Lambda_1\cap\Lambda_2\},\\
 \psi(z,y)&=\alpha(z)\alpha(y)\alpha(z+y)e^{-\pi iE_H(z,y)}.
 \end{split}
 \tag{TE7}
\]
The function $\psi$ is a well-defined homogeneous quadratic
function on $P/T$.  Its sum differs from the scalar
in $I_{23}I_{12}=\mu I_{13}$ by a positive real factor,
and $|\mu|=1$.

\begin{lemma}\label{te-positive-gauss}
Suppose an endomorphism $e$ of the ambient rational space
preserves the three lattices, commutes with $k$, and
each of $e^2-e$, $e^*e$, and $(1-e)^*(1-e)$ carries
$\Lambda_1+\Lambda_2+\Lambda_3$ into
$\Lambda_1\cap\Lambda_2\cap\Lambda_3$.
Then $\mu=1$.
\end{lemma}
\begin{proof}
The map $(z,y)\mapsto(ez,ey)$ induces an idempotent on
$P/T$: its square minus itself has both coordinates in
the common intersection, and every such pair belongs
to $T$.  Consequently
$P/T=A\oplus B$, where $A$ and $B$ are the images of
$e$ and $1-e$.

Put
\[
 \beta(x,y)=\sqrt D\operatorname{Re}H(x,y)-E_H(x,y)
                         =2E_H((\tau-1)x,y).
 \tag{TE8}
\]
If $x,y,x+y$ lie in the indicated three lattices,
integrality of their $Q$-values gives
$\psi(x,y)=e^{\pi i\beta(x,y)}$.
For $(z,y)\in P$,
\[
       \beta(ez,ey)=\beta(z,e^*ey)\in2\mathbb Z,
\]
because $z,e^*ey\in\Lambda_3$ and $\Lambda_3$ is
$\mathcal O_k$-stable.  Hence $\psi=1$ on $A$.
The same argument with $1-e$ proves $\psi=1$ on $B$.
Writing $b_\psi$ for the polar bicharacter gives
\[
       \sum_{x\in P/T}\psi(x)
                    =\sum_{a\in A}\sum_{b\in B}b_\psi(a,b)\geq0.
 \tag{TE9}
\]
Indeed the inner sum of any character of a finite group
is either its order or zero, as follows by translating
the sum by an element where the character is nontrivial.
The scalar $\mu$ is therefore nonnegative real.  Its
absolute value one implies $\mu=1$.
\end{proof}

\begin{proposition}\label{te-exact-paths}
Let $n,n'$ be even.  Normalized averaging satisfies both
triangle identities with one of $L,M$ fixed, and the
square identity
\[
 \begin{split}
 I_{\Lambda(L_1,M),\Lambda(L_1,M_1)}
 I_{\Lambda(L,M),\Lambda(L_1,M)}
 \;=\;
 I_{\Lambda(L,M_1),\Lambda(L_1,M_1)}
 I_{\Lambda(L,M),\Lambda(L,M_1)}
 \end{split}
 \tag{TE10}
\]
on the theta line of $\Lambda(L,M)$.
\end{proposition}
\begin{proof}
For the square, choose $b$ with
$b(L+L_1)\subset L\cap L_1$, and apply
Proposition~\ref{te-projector} to $M,M_1$.
The induced map on matrices is $X\mapsto eX$.
It preserves each of the four lattices.
Every error in (TE6) sends each of them into their
common intersection.  In fact, if $A:L_i\to M_j$
and $v\in L_k$, then $bv\in L_i$ and
\[
       R A(v)=b^{-1}R A(bv)\in M\cap M_1
\]
for each error $R$, since $R(M+M_1)\subset b(M\cap M_1)$.
Multiplication by $I\mathfrak a$ preserves this inclusion.
Lemma~\ref{te-positive-gauss} applies to either of the
two triples from the lower left to the upper right
corner of the square.  Both composites are the same
direct averaging operator, proving (TE10).

For three varying domain lattices and a fixed target
$M$, choose $b$ carrying their sum into their
intersection, and apply the same argument to the pair
$M,M$.  This proves the left triangle identity.
For three varying target lattices and fixed $L$, use
Proposition~\ref{te-projector} for $L,L$, with $b$
carrying the sum of those target lattices into their
intersection.  Use the induced map $X\mapsto Xe^*$.
Its adjoint is $X\mapsto Xe$; the errors are precisely
the corresponding right multiplications by the
expressions in (TE6) or their adjoints.  They have the
required images by the same calculation.  This proves
the right triangle identity.  Zero-dimensional cases
have identity averaging and satisfy all the formulas.
\end{proof}

\subsection{The determinant in the unitary action}

We next remove the possible sign from the action on an
even number of copies of a theta line.

\begin{lemma}\label{te-negative-norm}
For every positive rational integer $m$ there is
$\ell\in\mathcal O$ with
\[
                         \ell\bar\ell\equiv-1\pmod{m\mathcal O}.
 \tag{TE11}
\]
\end{lemma}
\begin{proof}
At an inert place of $F$ dividing $m$, choose a unit
of norm $-1$ by Lemma~\ref{acl-local-norm}; at a split
place use $(1,-1)$.  Approximate all these units modulo
$m$ in the product of the completed integer rings.
The Chinese remainder theorem for $\mathcal O$ gives
an element $\ell$ with these residues.  Multiplication
and conjugation preserve the congruences, proving (TE11).
\end{proof}

\begin{lemma}\label{te-real-gauss}
For any three lattices of the form (TE2), or any three
commensurable $\mathcal O$-stable self-dual trace lattices
with the same form, their triple scalar $\mu$ is real,
and hence $\mu^2=1$.
For an $E$-linear unitary map $g$ with $1-g$ invertible,
the Cayley Gauss sum in (TC31) is real.
\end{lemma}
\begin{proof}
For the first assertion choose $r>0$ carrying the sum
of the three lattices into their intersection.
Choose $m$ divisible by $2r^2$ and by the exponent of
the finite group $P/T$, and choose $\ell$ from (TE11).
Multiplication by $\ell$ permutes $P/T$: multiplication
by $-\bar\ell$ is its inverse, because $m(P/T)=0$.
Write $\ell\bar\ell+1=2r^2s$ with
$s\in\mathcal O_F$.  The scalar action has adjoint
$\bar\ell$, so (TE8) gives
\[
 \beta(\ell z,\ell y)+\beta(z,y)
                       =2\beta(rsz,ry)\in2\mathbb Z.
\]
Here both arguments on the right lie in the common
intersection, and $\beta$ is even integral there.
It follows that
$\psi(\ell z,\ell y)=\overline{\psi(z,y)}$.
The sum of $\psi$ is real, and so is $\mu$.

For the second assertion use the notation $A,C,\Gamma,q$
of (TC27), on an $\mathcal O$-stable lattice $\Lambda$.
The maps $A,C$ are $E$-linear.
Choose $r>0$ with $rC\Lambda\subset\Lambda$, and choose
$m$ divisible by $2r$ and the exponent of
$\Lambda/A\Gamma$.
Multiplication by the corresponding $\ell$ again
permutes this quotient.  Write
$\ell\bar\ell+1=2rs$, $s\in\mathcal O_F$.
Since $C^*=-C$, the exponent in
\[
 q(\ell x)q(x)=
 \exp\!\left(\pi i E_H(\tau(\ell\bar\ell+1)x,x)
                   -\pi H(C(\ell\bar\ell+1)x,x)\right)
 \tag{TE12}
\]
is $2\pi i$ times an integer.
Indeed the first $E_H$-pairing is even integral, and
$C(\ell\bar\ell+1)x=2s(rCx)\in2\Lambda$.
The second Hermitian value is purely imaginary, so its
imaginary part is even integral as well.
Thus $q(\ell x)=\overline{q(x)}$, proving reality of
its sum.
\end{proof}

\begin{proposition}\label{te-unitary-square}
Let $\Lambda$ be an $\mathcal O$-stable self-dual trace
lattice of even complex dimension $r$.  Let $g$ be
an $E$-linear unitary map.  The nonzero scalar defined by
\[
          F(g^{-1}z)=c(g,\Lambda)I_{\Lambda,g\Lambda}F(z)
                        \qquad(F\in\mathcal T(\Lambda,\alpha))
 \tag{TE13}
\]
satisfies
\[
                         c(g,\Lambda)^2=\det_{\mathbb C}g.
 \tag{TE14}
\]
\end{proposition}
\begin{proof}
The scalar exists by the one-dimensionality of the theta
spaces and the invariance of the norm semicharacter.
If $1-g$ is invertible, its nondegenerate Cayley Gauss
sum has absolute value $\sqrt N$ by (TC16), and is real
by Lemma~\ref{te-real-gauss}.  Its normalized square is
therefore one.  Equations (TC31) and (TC40) give
$c^2=(-1)^r\det g=\det g$.

For a general $g$, choose a scalar $z\in k$ of norm one
such that $z\ne1$ and $1-zg$ is invertible.
There are infinitely many choices before these finitely
many exclusions: with $s=\sqrt{-D}$, the values
$(1+ts)/(1-ts)$ for $t\in\mathbb Q$ are distinct and
have norm one.
Write $g=g_1g_2$, where $g_1=zg$ and $g_2=z^{-1}$.
Both factors have no fixed vectors.
Pullback commutes with normalized averaging after
transforming both lattices; this follows by reindexing
the finite sum and using preservation of $H$ and $\alpha$.
Consequently the scalar in (TE13) is the product of the
two factor scalars and a triple scalar for
$\Lambda,g_1\Lambda,g\Lambda$.
The latter has square one by Lemma~\ref{te-real-gauss}.
Squaring and using the first case proves (TE14).
The zero-dimensional case has scalar and determinant one.
\end{proof}

\subsection{The exact family and its three transformation laws}

Let $\theta$ be the fixed nonzero normalized theta
function for the rank-one trace lattice
$\lambda I\mathfrak a$, as constructed in
Sections~\ref{tc-section}--\ref{ts-section}.
Let $F_0$ be its product in $nn'$ independent entries,
and define
\[
 \begin{split}
 T(L,M)&=I_{\Lambda(L,M_0),\Lambda(L,M)}
                 I_{\Lambda(L_0,M_0),\Lambda(L,M_0)}F_0,\\
 b_\chi(L,M)&=\chi(\det(L_0,L))^{n'/2}
                             \chi(\det(M_0,M))^{-n/2},\\
 e_\chi(L,M)&=b_\chi(L,M)T(L,M).
 \end{split}
 \tag{TE15}
\]
Here $\det(L,L_1)$ is the relative determinant ideal.
The character $\chi$ is multiplicative on fractional
ideals $\mathfrak c$ with
$\mathfrak c\overline{\mathfrak c}=\mathcal O$ and obeys
\[
             \chi((u))=\prod_{\eta\in\Sigma}\eta(u)^{-1}
                          \qquad(u\bar u=1).
 \tag{TE16}
\]
These are exactly the character requirements used in
this assertion.  In particular $\chi$ never vanishes
on such an ideal.

\begin{theorem}\label{te-equivariance}
For positive even $n,n'$, the functions in (TE15)
satisfy all three identities
\[
 \begin{split}
 e_\chi(L,M_1)&=\chi(\det(M,M_1))^{-n/2}
                         I_{\Lambda(L,M),\Lambda(L,M_1)}e_\chi(L,M),\\
 e_\chi(L_1,M)&=\chi(\det(L,L_1))^{n'/2}
                         I_{\Lambda(L,M),\Lambda(L_1,M)}e_\chi(L,M),\\
 e_\chi(L,M)((\gamma'^{-1}z_\eta\gamma)_\eta)
                         &=e_\chi(\gamma L,\gamma'M)(z).
 \end{split}
 \tag{TE17}
\]
In the last equation $\gamma,\gamma'$ are unitary
matrices over $E$, and each matrix in the argument is
understood through its embedding $\eta$.
These identities also hold in dimension zero with the
empty-product conventions.
\end{theorem}
\begin{proof}
Proposition~\ref{te-exact-paths} gives the first two
identities for $T$ with the character factors removed.
The relative determinant is multiplicative in its
three lattice arguments: at any prime it is the
quotient of the two determinant lattices.  Its
Hermitian norm is one by self-duality.
Thus multiplication of the factors in (TE15) gives
the first two identities of (TE17).

For the last identity put
$D(\gamma)=\prod_{\eta\in\Sigma}\eta(\det_E\gamma)$.
First let only $\gamma$ act, at the standard pair.
Group the entries of $F_0$ into its $n'$ rows.
On each row the forward map of period lattices is
$x\mapsto x\gamma^{-1}$, whose complex determinant
is $D(\gamma)^{-1}$ and whose dimension $dn$ is even.
The operator in (TE13) and the normalized averaging
both factor over this orthogonal product: the finite
quotient is a product, and its positive square root
is the product of the positive square roots.
Proposition~\ref{te-unitary-square} therefore gives
\[
 F_0((z_\eta\gamma)_\eta)
                   =D(\gamma)^{-n'/2}T(\gamma L_0,M_0)(z).
 \tag{TE18}
\]
Likewise, grouping by columns gives
\[
 F_0((\gamma'^{-1}z_\eta)_\eta)
                    =D(\gamma')^{n/2}T(L_0,\gamma'M_0)(z).
 \tag{TE19}
\]
Here the forward map is $x\mapsto\gamma'x$, giving
the opposite determinant exponent.
Pullback commutes with averaging, as in the proof of
Proposition~\ref{te-unitary-square}.
The square and triangle identities in (TE10) therefore
extend (TE18)--(TE19) to every pair:
\[
 T(L,M)((\gamma'^{-1}z_\eta\gamma)_\eta)
       =D(\gamma)^{-n'/2}D(\gamma')^{n/2}
                                      T(\gamma L,\gamma'M)(z).
 \tag{TE20}
\]
To see explicitly why no path ambiguity remains,
transport the two defining edges in (TE15) by
$(\gamma,\gamma')$, precede them by the two edges
from $(L_0,M_0)$ to $(\gamma L_0,\gamma'M_0)$, and
interchange adjacent horizontal and vertical edges
using (TE10).  Consecutive edges in one direction
collapse by the corresponding triangle identity.
The result is exactly the two edges defining the
right side of (TE20).

Finally
\[
 \det(L_0,\gamma L)=(\det_E\gamma)\det(L_0,L).
\]
The same formula holds for $M$, and (TE16) gives
\[
 b_\chi(\gamma L,\gamma'M)
            =D(\gamma)^{-n'/2}D(\gamma')^{n/2}b_\chi(L,M).
\]
Substitution in (TE20) proves the last identity of
(TE17), with its exact scalar.
\end{proof}

This proves the covariance and the two moves required
for the actual normalized family.  It uses no choice
of signs for individual lattice moves.  For comparison,
the homogeneous quadratic function
$q(x)=e^{2\pi ix^2/3}$ on $\mathbb Z/3\mathbb Z$ has
sum $i\sqrt3$, so a general normalized quadratic Gauss
sum need not be one.  The local hyperbolic decomposition
and the norm congruence (TE11) are the specific arithmetic
facts that remove this obstruction here.  The measure
support at the selected primes and its use in incidence
maps remain separate assertions.


\section{Exact ordinary moment laws for the full theta family}
\label{gm-section}

Integral Taylor moments alone do not identify translation by
$p$-power torsion with multiplication of a measure by a character.
We prove that identification for the actual theta functions.  A
finite expansion into elliptic theta functions supplies the
identity; the independent normalization already proved supplies
the bound on the measure.  No root of a translated identity is
taken.

\subsection{An orthogonal frame with prime-to-$p$ index}

Retain $k=\mathbb Q(\sqrt{-D})$, $D\equiv7\pmod8$, and the
fixed abstract isomorphism $\iota:\mathbb C\simeq\mathbb C_p$.
Assume $p\ge5$ splits in $k$.  Let $V_k$ be a positive definite
Hermitian $k$-space of dimension $g$, with form $h$ linear in its
first variable, and let $\Lambda$ be an integral unimodular
$\mathcal O_k$-lattice.  On its complexification use
\[
 H=\frac{2}{\sqrt D}h,\quad E=\operatorname{Im}H,\quad
 Q(x)=h(x,x),\quad \alpha(x)=e^{\pi iQ(x)},\quad
 W_x=\alpha(x)U_x.                                      \tag{GM1}
\]
The lattice is self-dual for $E$, and $\alpha$ restricts to
its norm semicharacter.  The trace lattices (TE2) have precisely
this description after choosing their rational $k$-structure;
their common complex scaling is included in that identification.

\begin{lemma}\label{gm-orthogonal-frame}
There are $x_1,\ldots,x_g\in\Lambda$ such that
\[
 h(x_i,x_j)=a_i\delta_{ij},\qquad
 a_i\in\mathbb Z_{>0},\quad p\nmid a_i.                 \tag{GM2}
\]
For $\Lambda_0=\bigoplus_i\mathcal O_kx_i$,
the finite index $[\Lambda:\Lambda_0]$ is prime to $p$,
and
\[
 \Lambda_0^\#=\bigoplus_i a_i^{-1}\mathcal O_kx_i,
                  \qquad \Lambda_0\subset\Lambda\subset\Lambda_0^\#.
                                                               \tag{GM3}
\]
Here the dual can be taken for $h$ or for $E$.
\end{lemma}
\begin{proof}
At $p$ the ring $\mathcal O_k\otimes\mathbb Z_p$ is
$\mathbb Z_p\times\mathbb Z_p$, with the factors interchanged
by conjugation.  A unimodular Hermitian pairing is a perfect
pairing of its two factors.  A basis of one factor and its dual
basis of the other give an orthonormal basis of
$\Lambda\otimes\mathbb Z_p$.

Approximate that basis sufficiently closely by vectors in
$\Lambda$, which is dense in its $p$-adic completion.  Apply
Hermitian Gram--Schmidt over $k$.  Every successive denominator
is a $p$-adic unit, because the approximating Gram matrix is
congruent to $1$ modulo $p$.  The resulting orthogonal vectors
belong to $\Lambda\otimes\mathbb Z_{(p)}$ and still form a
basis at $p$.  Multiply each by a positive integer prime to $p$
to put it in $\Lambda$.  Its norm is then in
$\mathcal O_k\cap\mathbb Q=\mathbb Z$, is positive by
definiteness, and is a $p$-adic unit.  This proves (GM2).
The local basis property implies that the finite quotient
$\Lambda/\Lambda_0$ has order prime to $p$.
The displayed formula for the dual follows from the diagonal
Gram matrix.  Integrality of $h$ on $\Lambda$ gives the second
inclusion in (GM3).  Hermitian and alternating duality agree
because the inverse different of $k/\mathbb Q$ is
$(\sqrt{-D})^{-1}\mathcal O_k$.
\end{proof}

Let $\theta$ be the normalized even elliptic theta function of
Section~\ref{ct-section}, for the lattice $\mathcal O_k$.
Use its single character-normalized ordinary period $\omega$.
Let $\mathcal T(\Lambda)$ denote the entire functions satisfying
$W_lF=F$ for every $l\in\Lambda$.

\begin{proposition}\label{gm-finite-expansion}
For every $F\in\mathcal T(\Lambda)$ there are finitely many
complex coefficients $c_j$ and $t_{ji}\in k$ such that
$a_it_{ji}\in\mathcal O_k$ and
\[
 F\left(\sum_i z_ix_i\right)
       =\sum_j c_j\prod_{i=1}^g(U_{t_{ji}}\theta)(z_i)^{a_i}.
                                                               \tag{GM4}
\]
The $U$ in each factor is the elliptic operator, with Hermitian
form $2z\bar w/\sqrt D$.  Thus all its translates have order
prime to $p$.  The coefficients need not be assumed integral.
\end{proposition}
\begin{proof}
Put $G(\sum_i z_ix_i)=\prod_i\theta(z_i)^{a_i}$.
It is nonzero, belongs to $\mathcal T(\Lambda_0)$, and
$G(0)=1$.  The commutator pairing
\[
 (\Lambda/\Lambda_0)\times(\Lambda_0^\#/\Lambda)
       \longrightarrow\mathbb C^\times,
             (l,w)\longmapsto e^{-2\pi iE(l,w)}              \tag{GM5}
\]
is perfect.  Indeed, vanishing of the pairing with all
$\Lambda_0^\#$ means membership in its dual $\Lambda_0$;
vanishing with all $\Lambda$ means membership in
$\Lambda^\#=\Lambda$.  The two finite groups have the
same order by their dual indices.

Choose representatives $l$ and $w$ of these groups, with
the zero coset represented by zero, and set
\[
                      G_w=\sum_{l\in\Lambda/\Lambda_0}W_lU_wG.
                                                               \tag{GM6}
\]
Since $w\in\Lambda_0^\#$, $U_wG$ has the same
$\Lambda_0$ multiplier as $G$.  The operators $W_l$ on
$\Lambda$ are additive, so (GM6) is well-defined and belongs
to $\mathcal T(\Lambda)$.  Commuting $U_{-w}$ past $W_l$
and using (GM5) gives
\[
 \sum_{w\in\Lambda_0^\#/\Lambda}U_{-w}G_w
    =\sum_l\left(\sum_w e^{-2\pi iE(l,w)}\right)W_lG
    =|\Lambda_0^\#/\Lambda|G\ne0.                            \tag{GM7}
\]
Hence some $G_w$ is nonzero.  By
Lemma~\ref{tc-covolume}, $\mathcal T(\Lambda)$ has dimension
one, so $F=cG_w$ for a scalar $c$ (including $c=0$).
For each summand in (GM6), the Heisenberg composition rule gives
\[
 W_lU_wG=\alpha(l)e^{-\pi iE(l,w)}U_{l+w}G.
\]
The coordinates of $l+w$ belong to $a_i^{-1}\mathcal O_k$
by (GM3).  Orthogonality of the frame identifies its translate
with the product in (GM4).  Absorb the displayed scalar into
$c_j$.  The phases are roots of unity, since $E(l,w)$ is
rational.  This proves the assertion.
For $g=0$ the empty product is $1$ and the assertion is immediate.
\end{proof}

\subsection{The measure and its torsion characters}

Write $\iota_-:k\to\mathbb Q_p$ for $b\mapsto\iota(\bar b)$,
and put
\[
                 X_\Lambda=\Lambda\otimes_{\mathcal O_k,\iota_-}
                                            \mathbb Z_p.        \tag{GM8}
\]
This is the conjugate ordinary component of the lattice, a free
$\mathbb Z_p$-module of rank $g$.  For $v\in\Lambda$ the
function
\[
              L_v(z)=\iota(h(v,z)),\qquad z\in\Lambda,
                                                               \tag{GM9}
\]
extends to a $\mathbb Z_p$-linear function on $X_\Lambda$:
in the frame (GM2), it is
$\sum_i a_i\iota(v_i)z_i$, where $z_i$ denotes the
$\iota_-$-coordinate.  This convention is essential; the
Hermitian adjoint in the point moments uses this conjugate
component.

For $y=S\kappa l/p^n$, with $l\in\Lambda$,
$S\in\mathbb Z_{>0}$ prime to $p$, and
\[
        \kappa\in\mathcal O_k,
                   \iota(\bar\kappa)\in p^{2n}\mathbb Z_p,
                                                               \tag{GM10}
\]
define on rational lattice representatives
\[
                  \chi_y(z)=\iota\bigl(e^{-2\pi iE(z,y)}\bigr).
                                                               \tag{GM11}
\]
It induces a character of $X_\Lambda/p^nX_\Lambda$.
For example, in an orthogonal frame in which all $y_i$ have
the strong divisibility (CT24),
\[
        \chi_y(z)=\prod_i
          \iota\bigl(e^{-2\pi iE_1(1,y_i)}\bigr)^{a_i z_i\bmod p^n},
             \qquad E_1=\operatorname{Im}(2z\bar w/\sqrt D).
                                                               \tag{GM12}
\]
To verify this formula, adjointness gives
$E_1(b,y_i)=E_1(1,\bar b y_i)$.
On the formal $p^n$-torsion, multiplication by $\bar b$ is
multiplication by $\iota(\bar b)$ modulo $p^n$; this is
Proposition~\ref{ct-character-normalization}.  It also applies
to $b\in\mathcal O_k\otimes\mathbb Z_{(p)}$, by inverting
an integer prime to $p$.  Equation (GM12) follows and proves
the assertion about factorization.

\begin{theorem}\label{gm-moment-law}
Let $F\in\mathcal T(\Lambda)$ and let $t\in V_k$ have
order prime to $p$ modulo $\Lambda$.
There is a unique bounded $\mathbb C_p$-valued measure
$\mu_{F,t}$ on $X_\Lambda$ such that, for all
$v_1,\ldots,v_r\in\Lambda$,
\[
 \widehat{U_tF}\left(\omega\sum_jv_j\log(1+T_j)\right)
     =\int_{X_\Lambda}\prod_j(1+T_j)^{L_{v_j}(z)}
                                                  \,d\mu_{F,t}(z).
                                                               \tag{GM13}
\]
There is an integer $S>0$ prime to $p$, independent of $n$,
$l$ and $\kappa$, such that for every $y$ as in (GM10),
\[
 \widehat{W_yU_tF}\left(\omega\sum_jv_j\log(1+T_j)\right)
     =\int_{X_\Lambda}\chi_y(z)
                \prod_j(1+T_j)^{L_{v_j}(z)}\,d\mu_{F,t}(z).
                                                               \tag{GM14}
\]
In particular the same measure supplies all derivatives and
all these torsion characters.  For each fixed $n$, the identity
also holds for the additive subgroup generated by these $y$.
If the period-scaled moments of
$U_tF$ are integral on $\Lambda$, then $\|\mu_{F,t}\|\le1$.
\end{theorem}
\begin{proof}
Use (GM4).  The coordinates $t_i$ of $t$ have prime-to-$p$
denominators.  In these coordinates $U_t$ replaces $t_{ji}$
by $t_i+t_{ji}$ and introduces only a root of unity into each
coefficient.  Every transported elliptic series is bounded
and integral by Proposition~\ref{ct-translates-integral}.
Their powers, products and the finite sum are bounded, by
$\max_j|\iota(c_j)|_p$.  By the multivariable moment theorem
of Section~\ref{ot-section}, this series is the Amice
transform of a bounded measure on $\mathbb Z_p^g$.
In the frame, $L_{x_i}(z)=a_i z_i$.  Since $a_i$ is a
$p$-adic unit, push this measure forward by
$(A_i)\mapsto(A_i/a_i)$ to obtain $\mu_{F,t}$.

For general directions $v_j$, the coefficients $\iota(v_{ji})$
lie in $\mathbb Z_p$.  Substitution (LM3) proves (GM13).
The frame directions already determine the measure uniquely,
because the binomial polynomials are dense in the continuous
functions on $\mathbb Z_p^g$, as proved in
Section~\ref{ot-section}.

Choose $s>0$ prime to $p$ so that every
$s(t_i+t_{ji})$ is in $\mathcal O_k$.  Choose $d>0$ prime
to $p$ so that $d\Lambda\subset\Lambda_0$, and put $S=sd$.
If $l=\sum l_ix_i$, then
\[
          y_i=\frac{s\nu_i}{p^n},\qquad
          \nu_i=d\kappa l_i\in\mathcal O_k,
                      \iota(\bar\nu_i)\in p^{2n}\mathbb Z_p.
                                                               \tag{GM15}
\]
Theorem~\ref{ct-power-translation} applies to every elliptic
factor with this same $s$.  Multiplying its identities,
taking the powers $a_i$, and adding the finitely many terms
proves the multivariable substitution identity
\[
      \widehat{W_yU_tF}(T_1,\ldots,T_g)
          =\widehat{U_tF}\bigl(\zeta_1(1+T_1)-1,
                         \ldots,\zeta_g(1+T_g)-1\bigr),
 \quad \zeta_i=\iota(e^{-2\pi iE_1(1,y_i)}).              \tag{GM16}
\]
In (GM16), the hats stand for the frame expansion
$\omega\sum_i x_i\log(1+T_i)$.
The Weyl phase is exactly the product of the elliptic phases
raised to $a_i$, since $Q(y)=\sum_i a_i|y_i|^2$.
Thus no root of unity was discarded in deriving (GM16).
The convergent substitutions (CT21) commute with finite
products and sums.  Under the Amice transform they multiply
the measure by $\prod_i\zeta_i^{A_i}$, which is precisely
(GM12) after $A_i=a_i z_i$.  This proves (GM14) for frame
directions, and (LM3) proves it for all directions.
For a finite sum or negative of the permitted lifts at one
fixed level $n$, each
coordinate in (GM15) still has an integral numerator with the
same conjugate divisibility.  Apply the elliptic law directly
to these summed numerators; the same proof gives (GM14) on
the indicated additive subgroup.
Finally, integrality on the frame directions says that the
Amice transform has coefficient norm at most one.  The norm
identity in the moment theorem and the invertible coordinate
change $A_i=a_i z_i$ give $\|\mu_{F,t}\|\le1$.
\end{proof}

The characters appearing here include all characters of
$X_\Lambda/p^nX_\Lambda$.  Indeed choose $\kappa$ with
$\iota(\kappa)\equiv1\pmod{p^{2n}}$ and
$\iota(\bar\kappa)\equiv0\pmod{p^{2n}}$ by the Chinese
remainder theorem.  The alternating pairing of a self-dual
lattice pairs its two ordinary components perfectly modulo
$p^n$.  In (GM2) this is the product of the elementary elliptic
pairings, multiplied by the units $a_i$.  Multiplication by
$S$ is invertible modulo $p^n$, and $l\in\Lambda$ surjects
onto either component modulo $p^n$.  Formula (GM11) therefore
runs through every character.

\subsection{The bound for the normalized arithmetic family}

\begin{theorem}\label{gm-arithmetic-unit-bound}
Retain the hypotheses and the even dimensions of
Section~\ref{te-section}.  Let $p\ge5$ split in $k$, and
assume $p$ is prime to the elliptic ideal norm $a$.
For all unimodular $\mathcal O_E$-lattices $L,M$, the
normalized arithmetic vector $e(L,M)$ of (TE15)
and every translate of order prime to $p$ have the exact
laws (GM13)--(GM14), with $\|\mu_{e(L,M),t}\|\le1$.
The period is the common period $\omega$ of
Proposition~\ref{ci-common-period}; it is not selected again
after choosing $L,M$, a derivative, or a torsion point.
\end{theorem}
\begin{proof}
First suppose $L$ and $M$ are standard at every prime over
$p$.  The base vector is the tensor product of the normalized
trace-lattice theta vector.  Theorem~\ref{ts-integral-theta}
and multiplication of Amice transforms give its integral
moments and those of every prime-to-$p$ translate.
Each of the two finite averaging operators defining $e(L,M)$
has index prime to $p$, because its lattices have equal
$p$-adic completions.  Proposition~\ref{lm-finite-intertwiner}
preserves the asserted integrality.

The determinant characters in its definition are $p$-adic
units.  Here is the exact argument for this character convention.
For a norm-one fractional ideal $J$, choose $b>0$ with
$J^b=(z)$, using finiteness of the ideal class group.  Then
$J^{2b}=(z/\bar z)$, with $z/\bar z$ of norm one.  The
principal-ideal convention (TE16) gives
\[
            \chi(J)^{2b}
                   =\prod_{\eta\in\Sigma}\eta(z/\bar z)^{-1}.
                                                               \tag{GM17}
\]
If $J$ is a unit above $p$, both $z$ and $\bar z$ are
units there.  Every factor on the right is a $p$-adic unit
at every transported place above $p$.  Taking valuations in
(GM17) proves the assertion.  The relative determinant ideals
of the standard local lattices have this property.

For arbitrary $L,M$, choose rational unitary transformations
$\gamma,\gamma'$ such that $\gamma L$ and $\gamma'M$
are standard at every prime over $p$.  This uses only weak
approximation: Proposition~\ref{acl-local-basis} gives a local
unitary transformation at each of those places.  Perturb it
within the open set preserving the specified lattice so that
$1+g$ is invertible.  Such a perturbation exists by multiplying
by local norm-one scalars near $1$ and avoiding the finitely
many roots of $\det(1+zg)$.  Its Cayley transform
$(1-g)(1+g)^{-1}$ is skew-Hermitian.  Weak approximation in
the finite-dimensional $F$-space of skew-Hermitian matrices,
followed by the inverse Cayley transform, gives the required
rational transformations simultaneously.  Nonvanishing of
the denominator and the prescribed lattice images are open
conditions, so they survive sufficiently close approximation.

The exact covariance proved in Section~\ref{te-section}
identifies $e(\gamma L,\gamma'M)$ with the pullback of
$e(L,M)$ by $z\mapsto\gamma'^*z\gamma$.  This unitary
map carries the corresponding period lattice onto the original
one.  Thus the integral moments just proved, including all
prime-to-$p$ translates, pull back to those on $\Lambda(L,M)$.
Apply Theorem~\ref{gm-moment-law} and its norm conclusion.
The equalities of periods under ideal isogenies and common
scaling were proved in Section~\ref{ci-section}; no assertion
about the $p$-adic absolute value of the complex scaling is
needed.
\end{proof}

\begin{corollary}\label{gm-algebraic-jets}
Let $\delta_0$ be the discriminant twelfth root for the
principal elliptic lattice.  For the family in
Theorem~\ref{gm-arithmetic-unit-bound}, every rational-lattice
translate $t\in\mathbb Q\Lambda(L,M)$ has algebraic jets
\[
       \delta_0^{-k}D_{v_1}\cdots D_{v_k}
                           (U_te(L,M))(0)\in\overline{\mathbb Q}
       \qquad(v_i\in\mathbb Q\Lambda(L,M)).                 \tag{GM18}
\]
The assertion also permits algebraic linear combinations of
these directions.
\end{corollary}
\begin{proof}
Fix an embedding of $\overline{\mathbb Q}$ into
$\mathbb C_p$, and let $\iota$ run through all its extensions
to an abstract isomorphism $\mathbb C\simeq\mathbb C_p$.
The previous proof applies to each one.  The ordinary period
can change with $\iota$, but always has the form
$\omega=\Omega/\iota(\delta_0)$ with $|\Omega|_p=1$.
Choose a sufficiently large split prime $p$ avoiding the
denominators of $t,v_1,\ldots,v_k$ and the elliptic ideal
norm.  Clearing those denominators, the moment formula and
the unit bound give an upper bound independent of $\iota$
for the absolute value of the transport of (GM18).  In fact
the bound is one if the directions belong to the lattice.

We recall why this forces algebraicity.  If a complex number
$c$ were transcendental, fix one such isomorphism $\iota_0$
and choose a transcendental $a\in\mathbb C_p$, over the
image of $\overline{\mathbb Q}$, with arbitrarily large
absolute value.  Such an $a$ exists because that algebraic
closure is countable whereas every nonempty open ball is
uncountable.  The isomorphism
$\overline{\mathbb Q}(c)\to
\overline{\mathbb Q}(\iota_0^{-1}a)$ sending $c$ to
$\iota_0^{-1}a$ extends to an automorphism of $\mathbb C$:
extend transcendence bases, then extend over their algebraic
closures.  Composing with $\iota_0$ contradicts the uniform
upper bound.  Thus the number in (GM18) is algebraic.
Multilinearity gives the final assertion.

Only existence of a split prime outside a finite list is
used here.  It also has an elementary proof in this quadratic
case: choose $m$ divisible by $2D$ and all primes in that
list.  Every prime divisor $q$ of $Dm^2+1>1$ is outside
the list, is prime to $2D$, and satisfies
$(Dm)^2\equiv-D\pmod q$.  Hence $q$ splits in $k$.
\end{proof}

This theorem supplies the character identities needed for finite
Fourier restriction of the actual measures.  It does not yet
assert that a particular local test function selects the two
matrix supports in (CV8); that identification also requires the
explicit finite translation operator defining the arithmetic
column.


\section{The two supports of the arithmetic theta column}
\label{csp-section}

We identify the finite operator defining the column with an
ordinary measure on each of the two required matrix supports.
The calculation keeps the normalized Fourier denominator,
the positive square root in the lattice move, and the
determinant character.  Their combination gives the precise
factor $p^{-fa}$ in the selected type.

\subsection{The selected prime and the actual moment coordinates}

Use the fields, ideal $I$ of norm $A$, inverse half different
$\mathfrak a$, and normalized family of
Section~\ref{te-section}; here the capital $A$ distinguishes
the elliptic ideal norm from the matrix rank $a$.
Let $a,b$ be positive even integers and let
$N\subset E^a$, $N'\subset E^b$ be unimodular Hermitian
$\mathcal O_E$-lattices.  Put
\[
 \mathcal M=I\mathfrak a\operatorname{Hom}_{\mathcal O_E}(N,N'),
 \qquad \Lambda=\delta_I\operatorname{emb}(\mathcal M),
                         \qquad F=e_\chi(N,N'),              \tag{CS1}
\]
where $\delta_I$ is the discriminant twelfth root used for
the elliptic ideal lattice.  All embeddings in $\operatorname{emb}$
belong to the CM type $\Sigma$, and $f=|\Sigma|$.

The selected prime data are
\[
 p\ge5,\quad \pi\in\mathcal O_E,\quad
 \pi\bar\pi=p,\quad (\pi,\bar\pi)=1,\quad
 \mathcal P=\{P:P\mid\pi\},\quad
                 |\mathcal P|=f,\quad N(P)=p.              \tag{CS2}
\]
Assume $p\nmid A$ and that $I\mathfrak a$ is a unit at
every prime over $p$.  These are the actual selected-prime
conditions.  In particular $p$ splits completely in $E$:
$N_{E/\mathbb Q}(\pi)=p^f$, so each of the $f$ primes in
$\mathcal P$ has multiplicity one; the conjugate primes
are disjoint and account for the other $f$ primes over $p$.
The inverse half different is a unit there because its inverse
is an integral ideal whose product with its conjugate is the
relative different, and this different is a unit at a completely
split prime.  Thus (CS1) is locally the ordinary Hom lattice
at these primes.

For each $P\in\mathcal P$ fix nonzero residue functionals
\[
 \ell_P:N_P\longrightarrow\mathbb F_p,\qquad
              \ell'_P:N'_P\longrightarrow\mathbb F_p.       \tag{CS3}
\]
They factor through reduction modulo $P$ and are surjective.
They define the finite local test function
\[
 \phi_P(X)=\mathbf1_{\{\ell'_PX=\ell_P\}},\qquad
             X\in\operatorname{Hom}_{\mathbb F_p}
                       (N_P/PN_P,N'_P/PN'_P).               \tag{CS4}
\]
Outside the integral Hom lattice its value is zero.
Let $\phi=\prod_{P\in\mathcal P}\phi_P$.

There are two types of transported place used below.
In the selected type, every map $\iota\eta$, $\eta\in\Sigma$,
induces a prime in $\mathcal P$.  The ordinary measure
variable in (GM8) then consists of the conjugate local
components at $\overline{\mathcal P}$.
In the conjugate type, every $\iota\eta$ induces a prime
in $\overline{\mathcal P}$ and the measure variable consists
of the components at $\mathcal P$.  These are precisely the
two place conditions used for the two column bounds.

We record the coordinate constants explicitly.  Let $B$ be a
number field containing all embeddings of $E$ in $\Sigma$,
and put
\[
 \mathcal L_B=\mathcal O_B\langle(\eta(x))_{\eta\in\Sigma}:
                                      x\in\mathfrak a\rangle,
 \qquad J_\eta=\{c\in B:c\,e_\eta\in\mathcal L_B\}.
                                                               \tag{CS5}
\]
The prescribed eigen-direction generator $c_\eta$ satisfies
$J_\eta=c_\eta\mathcal O_B$, after the finite extension in
the arithmetic datum.

\begin{lemma}\label{csp-coordinate-units}
At every transported place above $p$ one has
$|\iota(c_\eta)|_p=1$.
With the common period $\omega=\Omega_0/\iota(\delta_0)$,
put
\[
 \Omega_I=A\iota(\delta_I)\omega
       =A\Omega_0\iota(\delta_I/\delta_0),\qquad
            u_\eta=\frac{\iota(c_\eta)}{\Omega_I}.           \tag{CS6}
\]
All $u_\eta$ have absolute value one.  The moment of the
mixed derivative with direction $c_\eta$ in entry $(i,j)$
is exactly the corresponding product of the coordinates
\[
                         u_\eta\,\iota\eta((X^*)_{ji}).     \tag{CS7}
\]
There are no factorials in this mixed-derivative convention.
\end{lemma}
\begin{proof}
Every coordinate of $\mathcal L_B$ is integral at the chosen
place, since $\mathfrak a$ is a unit above $p$.  Hence
$|\iota(c_\eta)|_p\le1$.
For $\eta'\ne\eta$, complete splitting and the Chinese
remainder theorem give $x\in\mathcal O_E$ for which
$\iota\eta(x)-\iota\eta'(x)$ is a unit.  Start with the
vector $(1)_\eta\in\mathcal L_B$, which belongs to this
lattice because $1\in\mathfrak a$.  For each $\eta'\ne\eta$
apply the operator
\[
         (v_\rho)_\rho\longmapsto
                     ((\rho(x)-\eta'(x))v_\rho)_\rho.       \tag{CS8}
\]
It preserves $\mathcal L_B$, kills the $\eta'$ coordinate,
and multiplies the $\eta$ coordinate by a unit at the chosen
place.  After all the operators, the resulting vector is
$c e_\eta$ with $c\in J_\eta$ and $|\iota(c)|_p=1$.
Writing $c=c_\eta v$ with $v\in\mathcal O_B$ gives the
reverse inequality for $|\iota(c_\eta)|_p$.

The period and discriminant ratio in (CS6) are units by
Proposition~\ref{ci-common-period} and the ideal discriminant
comparison in that section; $p\nmid A$.  This proves the
unit assertions.
For a coordinate direction $v=c_\eta e_{\eta,i,j}$,
the Hermitian form in (TE2) gives, on rational lattice points,
\[
 h\bigl(v,\delta_I\operatorname{emb}(X)\bigr)
                 =\frac{c_\eta}{A\delta_I}\,
                                      \eta((X^*)_{ji}).      \tag{CS9}
\]
Multiply by $\omega^{-1}$ as in the derivative form of
(GM13).  This gives (CS7), with (CS6) exactly as stated.
Multilinearity proves the mixed-derivative formula.  These
coordinate directions lie in the algebraic linear span of
$\delta_0^{-1}\Lambda$: the embedding matrix is algebraic
and invertible, and $\delta_I/\delta_0$ is algebraic.
Corollary~\ref{gm-algebraic-jets} therefore also proves
algebraicity of these mixed derivatives, including all the
rational translates occurring below.
\end{proof}

\subsection{The finite operator, with its normalization}

Define the two commensurable self-dual lattices
$\mathcal A,\mathcal B\subset\mathcal M\otimes\mathbb Q$
by their localizations.  At a selected prime $P$, put
\[
 \mathcal A_P=p^{-1}\mathcal M_P,\quad
 \mathcal B_P=p\mathcal M_P,\qquad
 \mathcal A_{\bar P}=p\mathcal M_{\bar P},\quad
 \mathcal B_{\bar P}=p^{-1}\mathcal M_{\bar P};              \tag{CS10}
\]
at other primes both equal $\mathcal M$.
In formulas with $U,W,E$, identify rational matrices with
their scaled embedded vectors in (CS1).
The finite groups
\[
 T=(\mathcal M\cap\mathcal A)/(\mathcal A\cap\mathcal B),
 \qquad Y=\mathcal B/(\mathcal M\cap\mathcal B)             \tag{CS11}
\]
have order $p^{fab}$ and pair perfectly by
$\kappa(t,y)=e^{-2\pi iE(t,y)}$.
The first is the product of the matrix spaces in (CS4).
This follows directly from (CS10); the local trace duality
is perfect because $\mathcal M$ is self-dual.

For a function $\Phi$ on $T$, the finite theta operator is
\[
 \mathcal Q_\Phi F=
       \frac1{|T|}\sum_{y\in Y}\sum_{t\in T}
                  \Phi(t)\kappa(t,y)^{-1}W_yF.              \tag{CS12}
\]
Representatives in the indicated lattices are understood.
The expression is independent of them: on $\mathcal B$ the
operators $W_y$ are additive; representatives differing by
$\mathcal M\cap\mathcal B$ act identically on $F$; and
the characters factor through (CS11).

Here is the precise connection with the actual theta indicator.
At a general level $c$ and selected set of primes, use the
two half lattices defined as in (CS10), with $c$ in place
of $p$.  Put
$n=[\mathcal M\cap\mathcal A:\mathcal A\cap\mathcal B]$
and take representatives $R$ of
$\mathcal B/(\mathcal M\cap\mathcal B)$.
For a rational selected coset $x$, choose a lift $s_x$ in
some half lattice $\mathcal A(c')$, with $c'>0$, whose
selected components agree with $x$ modulo $c\mathcal M$.
The Chinese remainder theorem supplies such a lift.
The actual indicator is
\[
                  \Theta_x=\frac1n W_{-s_x}\sum_{y\in R}W_yF.
                                                               \tag{CS12a}
\]
Indeed its definition has an outside factor
$\mathrm{vol}^{-1/2}$.  Its Fourier coefficient at the
indicator of $\mathcal B$ is
$\mathrm{vol}^{-1/2}(\mathrm{vol}/n)$: only the zero
class of $\mathcal M\cap\mathcal A$ contributes.  The
two factors cancel exactly, giving $1/n$ in (CS12a).
Here $\mathrm{vol}>0$, since all lattice indices are
positive.  No assertion that this volume is a $p$-adic
unit is involved.

The lift in (CS12a) is independent of choice.  Two lifts
$s,t$ differ by an element of $\mathcal B$.  Moreover they
both belong to a common integral half lattice: if their
levels are $c',c''$, use $\mathcal A(\gcd(c',c''))$.
At a selected prime their difference is integral, so the
smaller of the two possible denominator bounds bounds
both; at the conjugate prime both have the corresponding
integral divisibility.  At other primes they belong to
$\mathcal M$.  Consequently $W_{-t}W_{t-s}=W_{-s}$
by the semicharacter law on that half lattice.  The sum
over $R$ is fixed by $W_{t-s}$, by permutation of its
classes.  This proves the claimed independence, including
its phase.

At $p$, the cosets with nonzero (CS4) are the elements
of $T$ in (CS11).  Choose their lifts in $\mathcal M$.
Commuting $W_{-s_x}$ through the sum in (CS12a), and using
$W_{s_x}F=F$, gives exactly (CS12).  Thus its denominator
is $|T|$, with no further $p$-primary factor.

\begin{lemma}\label{csp-fourier-fiber}
Choose $X_0\in\mathcal M$ satisfying
$\ell'_PX_0=\ell_P$ modulo $P$ for every selected $P$,
and put $K=\ker(\prod_P(\ell'_P\cdot))\subset T$.
Let $K^\perp\subset Y$ be its annihilator.  Then
\[
 |T/K|=|K^\perp|=p^{fa}=:h_0,
 \qquad
 \frac1{|T|}\sum_{t\in T}\phi(t)\kappa(t,y)^{-1}
       =\begin{cases}
          h_0^{-1}\kappa(X_0,y)^{-1},&y\in K^\perp,\\
          0,&y\notin K^\perp.
        \end{cases}                                       \tag{CS13}
\]
\end{lemma}
\begin{proof}
For each $P$, choose $v\in N'_P/PN'_P$ with
$\ell'_P(v)=1$.  The map $X=v\ell_P$ is a solution of
the affine equation.  Chinese remaindering lifts all these
solutions to one $X_0\in\mathcal M$.
The linear map $X\mapsto\ell'_PX$ onto the $a$-dimensional
dual of $N_P/PN_P$ is surjective, so its kernel has
codimension $a$.  Hence $|T/K|=p^{fa}$.
The support of $\phi$ is the coset $X_0+K$.
Summing a character on this coset gives zero unless it is
trivial on $K$, in which case it gives
$|K|\kappa(X_0,y)^{-1}$.  Division by $|T|$ proves (CS13).
The equality $|K^\perp|=|T/K|$ follows from perfectness
of the pairing in (CS11).
\end{proof}

\subsection{The conjugate type: exact restriction of the measure}

Suppose first that $\iota$ has the conjugate type.  Then
$Y$ is the formal $p$-torsion component, and the measure
variable is the selected component of $\mathcal M$.
Representatives of $Y$ can be chosen among the strong good
lifts in (GM10), with any prescribed additional
prime-to-$p$ multiple in $S$.  To see this explicitly,
choose $\kappa\in\mathcal O_k$ whose formal component is
$S^{-1}$ modulo $p^2$ and whose conjugate component is zero
modulo $p^2$.  The vectors $S\kappa l/p$, $l\in\Lambda$,
have selected component in $p\mathcal M_P$, lie in
$\mathcal B$, and realize every class of $Y$.
They differ from any other representatives by
$\mathcal M\cap\mathcal B$, so the operator is unchanged.
The exact law (GM14) therefore identifies (CS12) with the
measure
\[
 \left(\frac1{|T|}\sum_{y\in Y}\sum_{t\in T}
          \phi(t)\kappa(t,y)^{-1}\kappa(X,y)\right)d\mu_F(X)
                              =\phi(X)\,d\mu_F(X).          \tag{CS14}
\]
The equality is finite Fourier inversion.  In particular its
norm is at most one, with support
\[
 S_{\bar P}=\{X:\ XN_P\subset N'_P,\quad
                        \ell'_PX=\ell_P\text{ for all }P\in\mathcal P\}.
                                                               \tag{CS15}
\]
The Fourier denominator has disappeared by exact character
orthogonality.  Bounding its terms separately would give an
incorrect and unnecessarily large denominator.

\subsection{The selected type: the companion lattice}

Define a new unimodular lattice $N''$ by
\[
 N''_P=\ker\ell'_P\subset N'_P,\qquad
 N''_{\bar P}=(N''_P)^\#,\qquad
 N''_Q=N'_Q\quad(Q\nmid p).                               \tag{CS16}
\]
The dual at $\bar P$ is for the perfect Hermitian pairing
between the two conjugate completions.
Local patching, as in Lemma~\ref{acl-patching}, gives a
global projective lattice with these completions; their
duality proves global unimodularity.  Put
\[
           \mathcal M''=I\mathfrak a
                         \operatorname{Hom}_{\mathcal O_E}(N,N'').
                                                               \tag{CS17}
\]

\begin{lemma}\label{csp-companion}
The group $K^\perp$ in (CS13) is the image of
$\mathcal B\cap\mathcal M''$ in $Y$, and this image
identifies with
$\mathcal M''/(\mathcal M\cap\mathcal M'')$.
This quotient has order $h_0=p^{fa}$.  Moreover
\[
        \det(N',N'')=(\pi/\bar\pi),\qquad
        \pi N'\subset N'',\quad \bar\pi N''\subset N'.     \tag{CS18}
\]
\end{lemma}
\begin{proof}
At $P$, the kernel of $X\mapsto\ell'_PX$ consists of maps
whose image is in $\ker\ell'_P$.  Under the perfect matrix
trace pairing, its annihilator consists of adjoint matrices
that take $\ker\ell'_P$ integrally into $N_P$.
Equivalently, in bases with $\ell'_P$ the last coordinate,
the only possibly nonintegral row of the adjoint matrix is
its last row, with denominator $p$.  The other rows are
integral.  This is exactly the dual of
$\operatorname{Hom}(N_P,\ker\ell'_P)$ at $\bar P$,
which is $\mathcal M''_{\bar P}$.
It enlarges $\mathcal M_{\bar P}$ by an index $p^a$.

At $P$ the lattice $\mathcal M''_P$ is contained in
$\mathcal M_P$ and contains $p\mathcal M_P$.
Thus its selected component creates no quotient modulo
$\mathcal M$.  By the Chinese remainder theorem one may
choose representatives of $\mathcal M''/(\mathcal M\cap\mathcal M'')$
whose selected components are in $p\mathcal M_P$.
They then belong to $\mathcal B$ and realize the indicated
annihilator.  The quotient size is $\prod_Pp^a=h_0$.

The relative determinant has valuation $1$ at each $P$,
$-1$ at each $\bar P$, and zero elsewhere.  These are
the valuations of $(\pi/\bar\pi)$ by (CS2), proving the
first equality in (CS18).  The two inclusions follow
locally from (CS16): at $P$ the kernel contains $pN'_P$,
and at $\bar P$ its dual is contained in $p^{-1}N'_{\bar P}$.
At all other primes the lattices agree and $\pi,\bar\pi$
are units.
\end{proof}

The exact right-move identity (TE17) now gives
\[
 \sum_{y\in K^\perp}W_yF
      =\sqrt{h_0}\,\chi((\pi/\bar\pi))^{a/2}
                                      e_\chi(N,N'').        \tag{CS19}
\]
The positive real square root is the same one as in the
definition of the averaging operator.  Since $X_0\in\mathcal M$,
$W_{X_0}F=F$, and the Heisenberg commutator yields
\[
 \mathcal Q_\phi F
    =h_0^{-1}W_{-X_0}\sum_{y\in K^\perp}W_yF
    =h_0^{-1}\sqrt{h_0}\,
       \chi((\pi/\bar\pi))^{a/2}
                              W_{-X_0}e_\chi(N,N'').       \tag{CS20}
\]
This proves the needed identity of theta functions, before
taking a valuation.

Suppose $\iota$ has the selected type.  The character convention
(TE16) and (CS2) give
\[
 \left|\iota\chi((\pi/\bar\pi))\right|_p=p^f,\qquad
 \left|\iota\left(h_0^{-1}\sqrt{h_0}
                         \chi((\pi/\bar\pi))^{a/2}\right)\right|_p
                                                        =p^{fa}.
                                                               \tag{CS21}
\]
Indeed every $\iota\eta(\pi)$ has absolute value $p^{-1}$,
and every $\iota\eta(\bar\pi)$ has absolute value one.
The square root has absolute value $p^{-fa/2}$, cancelling
the absolute value $p^{fa/2}$ of the character factor.

The class of $X_0$ modulo $\mathcal M''$ is formal
$p$-torsion in this type: its selected component may have
order $p$, while its conjugate component already lies in
$\mathcal M''_{\bar P}$.  We spell out the harmless
choice of its representative.  If $r$ is any rational
formal $p^n$-torsion lift modulo a self-dual lattice, choose
a strong good lift $y$ congruent to $r$ modulo that lattice,
using the two ordinary idempotents and the Chinese remainder
theorem as above.  For $l=r-y$ in the lattice,
the Heisenberg formula and the lattice theta multiplier give
\[
                         W_rF'=\zeta W_yF',                 \tag{CS22}
\]
where $\zeta$ is a root of unity, since all quadratic and
alternating values at rational lattice points are rational.
The exact law for $y$ therefore gives a measure of norm
at most one for $W_rF'$, on the same ordinary lattice
component as for $F'$.  This observation needs only the
absolute value of $\zeta$, not an unspecified equality of
phases in (CS20).

Apply this to $r=-X_0$ and $F'=e_\chi(N,N'')$.
By (CS20)--(CS22), the selected-type measure for
$\mathcal Q_\phi F$ has norm at most $p^{fa}$ and is
supported on the conjugate ordinary component of
$\mathcal M''$.  In the original matrix language its
support is contained in
\[
 S_P=\{X:\ X^*(\ker\ell'_P)\subset N_P
                               \text{ for all }P\in\mathcal P\}.
                                                               \tag{CS23}
\]
To check the inclusion, let $z\in\ker\ell'_P=N''_P$.
For $v\in N_{\bar P}$, the image $Xv$ lies in $N''_{\bar P}$
on the conjugate component.  Perfect Hermitian duality
gives $h(v,X^*z)=h(Xv,z)\in\mathcal O_{E,P}$.
Since $N_P$ is its appropriate local dual, $X^*z\in N_P$.
This is exactly the condition in (CS23).

\subsection{The other primes and simultaneous moment approximation}

We verify that restoring the other local test functions does
not change either bound.  Let $S_0$ be the finite set of
auxiliary primes, disjoint from the primes over $p$, and
assume $S_0\cap\overline{S_0}=\varnothing$, as in the
selected local datum.  Let $c_A,g$ be positive integers
prime to $p$.
The full test function has the form
\[
                              \Phi=\phi\,\Phi_0,             \tag{CS24}
\]
where $\Phi_0$ is supported in the prescribed
$c_A^{-1}$ integral Hom lattice at $S_0$, is invariant under
$pg$ times that local Hom lattice, and all its transported
values have absolute value at most one.  Invariance means
that it is a function on the finite quotient at $S_0$;
it does not impose an additional condition at $p$.
These are the away-from-$p$ datum conditions in the column.

\begin{lemma}\label{csp-other-primes}
The full finite theta operator associated to (CS24), with
the same parameter $\mathrm{vol}=1$, is a finite sum
\[
                           \sum_j d_jW_{r_j}\mathcal Q_\phi F,
                                                               \tag{CS25}
\]
where $|\iota(d_j)|_p\le1$, each $r_j$ has order prime to
$p$ modulo both $\mathcal M$ and $\mathcal M''$, and the
representatives may be chosen arbitrarily divisible by $p$
at all primes over $p$.  In the conjugate type its measure
has norm at most one and support (CS15); in the selected
type it has norm at most $p^{fa}$ and support (CS23).
\end{lemma}
\begin{proof}
All representative quotients in the finite formula split
as the direct product of their $p$-primary part and their
part of order prime to $p$.  This is the elementary primary
decomposition of a finite abelian group, or equivalently
the Chinese remainder theorem on each localized lattice.
The alternating pairings between the factors are trivial:
a character pairing value has order dividing the orders
of both factors.  Each normalized Fourier denominator
therefore factors as a power of $p$ times an integer prime
to $p$.

Use the indicator formula (CS12a), and sum it with the
values of (CS24) over the finite selected coset quotient.
Its $p$-part is $T$; the $p$-part of the representative
group $R$ is $Y$; and the $p$-part of $n$ is $|T|$.
These assertions follow directly from (CS10), since
$p\nmid c_Ag$.  Choose $s_x=s_{x_p}+s_{x_0}$ according
to the primary decomposition, with $s_{x_p}\in\mathcal M$
and $s_{x_0}$ arbitrarily $p$-divisible.  Such representatives
exist by Chinese remaindering.  Expand the remaining
prime-to-$p$ finite sums in (CS12a).  Each coefficient is
a value of $\Phi_0$, times roots of unity and the inverse
of an integer prime to $p$, so has absolute value at most
one.  The remaining $p$-primary sum is exactly (CS12),
because (CS24) factors and its selected part is (CS4).

Here are explicit compatible choices of representatives.
If an away-from-$p$ class has order dividing $n$, choose
an integer $j\equiv1\pmod n$ with $j\equiv0\pmod{p^m}$,
and replace its representative $r$ by $jr$.
The class is unchanged, and its $p$-components are as
divisible as desired.  Choose the $p$-primary representatives
with the additional factor $S$ in (GM10) divisible by all
the finitely many away-from-$p$ denominators.  Their
pairings with the other representatives are then integral
at every prime.  Thus their $W$ operators commute exactly.
Changing a representative in its own integral lattice
quotient does not change the corresponding sum.  Combining
the remaining two operators introduces at most a root of
unity by the rational Heisenberg formula.  This proves
(CS25).  Taking $m\ge2$ ensures in addition that $r_j$
is integral in both lattices $\mathcal M,\mathcal M''$
at $p$, so its order in either quotient is prime to $p$.

For the conjugate type one may instead apply the $p$-primary
Fourier sum after each away-from-$p$ translate.  It then
restricts the norm-one measure of $W_{r_j}F$ by (CS14).
The chosen $r_j$ has zero reduction at the selected primes,
so the same affine condition (CS15) occurs.
For the selected type commute $W_{r_j}$ with $W_{-X_0}$
in (CS20).  The introduced factor is a root of unity.
Apply the ordinary-torsion law to the prime-to-$p$ translate
$W_{r_j}e_\chi(N,N'')$.  One can choose its strong lift
congruent to $-X_0$ modulo an additional prime-to-$p$
multiple of $\mathcal M''$, so the same argument (CS22)
applies.  Its measure still has norm at most one on the
ordinary component of $\mathcal M''$.
Equation (CS21) gives the stated bound and (CS23) the
support.  Finite sums with coefficients bounded by one
preserve these bounds by the ultrametric inequality.
\end{proof}

\begin{theorem}\label{csp-two-supports}
For the actual finite theta column with standard
representatives at the selected primes, its mixed-derivative
table has the simultaneous finite approximation required
in (CV9)--(CV10).  In the conjugate type the approximating
point matrices belong to $S_{\bar P}$ and their coefficients
have absolute value at most one.  In the selected type
they belong to $S_P$ and the table is approximated by
\[
       \sum_\nu b_\nu p^{-fa}
             \prod_{\eta,i,j}
                \bigl(u_\eta\iota\eta((X_\nu^*)_{ji})\bigr)^{m_{\eta,i,j}},
                        \qquad |b_\nu|_p\le1.               \tag{CS26}
\]
The $u_\eta$ are the fixed units (CS6), independent of the
derivative multi-index and the finite approximation.  Thus
the two ordinary support inputs in
Theorem~\ref{cv-two-kernels} are proved for this column.
\end{theorem}
\begin{proof}
Lemmas~\ref{csp-coordinate-units} and
\ref{csp-other-primes} give measures with the exact point
moments and the stated bounds.  Global rational lattice
points are dense in their ordinary completions, by the
Chinese remainder theorem.  In the conjugate type impose
the fixed residue equations (CS15); lifting those equations
and then arbitrarily many additional $p$-adic digits proves
density in the restricted support.  In the selected type
the global points of $\mathcal M''$ are dense in its
ordinary component and all satisfy (CS23).

Apply Theorem~\ref{lm-finite-approximation} to any finite
family of the point polynomials (CS7).  It approximates
all of them at once by a finite sum of point masses,
with coefficients bounded by the measure norm.
For the selected type multiply the measure by $p^{fa}$
first.  Its norm is then at most one, giving exactly
(CS26) after division by $p^{fa}$.

The arithmetic column also divides the theta jets by its
determinant character and inserts the auxiliary local
character values.  For standard representatives these
are $p$-adic units, by the principal-ideal argument
(GM17) and the prime-to-$p$ conductor condition.  The
normalization is therefore a scalar of absolute value
one and preserves both approximation assertions.
Where the coefficient table is descended to its specified
number field, apply the fixed embedding of that field;
algebraicity and the equality of the two descriptions of
the same jet identify the transported values.  This is
the precise descent condition in the actual table and
does not introduce an additional denominator.
\end{proof}

The two bounds above use the actual local affine equation,
the actual companion lattice, and the exact determinant
character.  A bounded measure with no character identities
would not give (CS14), and an eighth-power comparison of
theta vectors would not give (CS19).  Those dependencies
were separately proved in Sections~\ref{gm-section} and
\ref{te-section}.  The incidence identity between adjacent
columns is a further statement and is not asserted here.


\section{A finite coefficient field for the required theta jets}
\label{jcf-section}

We now pass from the algebraicity in Corollary~\ref{gm-algebraic-jets}
to one number field containing every coefficient used in a fixed
finite column construction. The distinction matters: a collection
of algebraic numbers need not lie in a single number field.
The proof uses the finiteness of positive definite lattice classes,
finite residue quotients, and the fixed bound on the orders of the
required derivatives. This is the finite-generation argument needed
for the submitted coefficient-field construction.\cite{alpo-uniform-source}

\subsection{The data fixed before adjoining coefficients}

Fix the CM field $E$, its CM type $\Sigma$, the finitely many
positive Hermitian spaces, their even ranks, and the normalized
theta family of Sections~\ref{te-section} and \ref{gm-section}.
Fix also the finite list of rational primes $T_u$ used in the
columns, the auxiliary primes $g_i$, and the finite family
$\Xi$ of finite-order character twists. These are arithmetic
data over $E$; they do not depend on an ambient coefficient
field. Let $\mathcal B$ be the finite set of primes of $E$
consisting of the primes above the $T_u,g_i$ and the primes
above the ramification loci of the characters in $\Xi$.
Taking all primes above these rational primes permits every
choice of one prime from a conjugate pair.

Let $m$ be the fixed length of the chain. For $R\subseteq
\mathcal B$ and an integer $b\in\{m,m+1\}$ put
\[
 q_R=\prod_{P\in R}NP,\qquad
                   M_{u,R,b}=T_uq_R^{2b}.                 \tag{JF1}
\]
The exponent $m$ is sufficient for the initial construction;
$m+1$ includes its prescribed increase of local level.
There are finitely many such positive integers. Fix a bound
$d$ on the total order of all derivatives in the finite
list of weights. One can instead take the exact order for
each step; taking their maximum only enlarges a finite set.
Fix the finitely many algebraic direction coefficients
$c_\eta$. The actual construction chooses its eigen-direction
lattices and their generators before the coefficient-field
enlargement; their values are included among the generators
below.

We shall allow all $S\subseteq\mathcal B$ for which the theta
indicator is defined, all positive divisors $c\mid M$ of
one of the integers in (JF1), and all integral self-dual
lattices in the fixed Hermitian spaces. Enlarging the finite
list in this way makes it independent of the later character
realizations and choices of prime representatives.

\subsection{Characters and finite residue quotients}

Let $\mathcal I^-$ be the group of fractional ideals $J$ of
$E$ satisfying $J\bar J=\mathcal O_E$, and let
$\mathcal P^-=\{(u):u\bar u=1\}$. The determinant character
used here satisfies
\[
                  \chi((u))=\prod_{\eta\in\Sigma}\eta(u)^{-1}
                         \quad(u\bar u=1).               \tag{JF2}
\]

\begin{lemma}\label{jcf-character-values}
The group $\mathcal I^-/\mathcal P^-$ is finite. Consequently
all values of $\chi$ belong to one number field.
\end{lemma}
\begin{proof}
The map to the ideal class group has finite image. Its kernel
consists of principal ideals $(z)$ with $z\bar z\in
\mathcal O_{E^+}^{\times,+}$. On that kernel the assignment
\[
 (z)\longmapsto z\bar z\pmod {N_{E/E^+}\mathcal O_E^\times}
\]
is well defined, because replacing $z$ by $z\varepsilon$ changes
its norm by $N\varepsilon$. Its kernel is precisely
$\mathcal P^-$: if $z\bar z=\varepsilon\bar\varepsilon$, then
$z/\varepsilon$ has norm one. The unit quotient on the right
is finite. Indeed the norm subgroup contains the squares of
$\mathcal O_{E^+}^\times$, and the unit theorem makes the
quotient by squares finite. This proves the first assertion.

Choose representatives $J_1,\ldots,J_r$ of this finite group
and let $h$ kill it. For each $i$ there is $u_i\in E^\times$
with $u_i\bar u_i=1$ and $J_i^h=(u_i)$. Formula (JF2) gives
\[
                  \chi(J_i)^h=
                         \prod_\eta\eta(u_i)^{-1}
                          \in\overline{\mathbb Q}^{\times}.
\]
Thus each $\chi(J_i)$ is algebraic. Every ideal $J$ has the
form $J_i(u)$ with $u\bar u=1$. Its value is
$\chi(J_i)\prod_\eta\eta(u)^{-1}$, which lies in the field
generated by the normal closure of $E$ and the finitely many
$\chi(J_i)$. This proves the assertion without adjoining
unboundedly many roots of arbitrary elements of $E$.
\end{proof}

A finite-order twist has finite image, contained in the roots
of unity of its order. Its ramification locus is finite.
Hence all its values, the roots of unity of orders
$(NP)^{2m+2}$ required at the finitely many local levels, and
the finitely many square roots of local norms lie in a number
field. Any prescribed finite list of algebraic frame
coefficients can be adjoined at the same time.
For algebraic Hecke characters with nonzero infinity type,
this statement concerns their normalized finite-idele
values; no assertion is made that their values on arbitrary
archimedean ideles are algebraic.

For two fixed self-dual lattices $L,L'$ write
\[
 H_S(L,L')=\{X\in\operatorname{Hom}_E(V,V'):
                         XL_P\subseteq L'_P\ (P\in S)\}.
\]

\begin{lemma}\label{jcf-finite-indicators}
For positive integers $M,c$ the quotient
\[
                  M^{-1}H_S(L,L')/cH_S(L,L')              \tag{JF3}
\]
is finite. Consequently the theta indicators
$\Theta_{L,L',S,c,x}$ with $x\in M^{-1}H_S(L,L')$ take
only finitely many distinct values as holomorphic functions.
\end{lemma}
\begin{proof}
Put $R_S=\bigcap_{P\in S}\mathcal O_{E,P}\subset E$.
If $S$ is empty, $R_S=E$ and the quotient is zero.
Otherwise $R_S$ is a semilocal Dedekind ring. An invertible
ideal of it is principal: choose, by the Chinese remainder
theorem, an element with its prescribed valuations at the
finitely many maximal ideals. Thus every finitely generated
projective $R_S$-module is free. In particular $H_S(L,L')$
is free of rank $ab$ over $R_S$, where $a,b$ are the ranks
of $L,L'$. Multiplication by $M$ identifies (JF3) with
$H_S/McH_S$. Finally
\[
 R_S/McR_S\simeq
     \prod_{P\in S}\mathcal O_E/P^{v_P(Mc)}
\]
with the factors of exponent zero omitted. The product is
finite. These facts also follow directly by choosing local
bases and applying the Chinese remainder theorem.

The finite theta indicator depends only on the residue class
of $x$ modulo $cH_S$. This is its coset definition: the two
lifts differ by the period of the characteristic function.
Equivalently, their Weyl operators differ by a stabilizer of
the averaged vector, with the same semicharacter phase. The
lift-independence calculation in Section~\ref{csp-section}
proves exactly this equality, including the phase. Thus the
number of distinct functions is bounded by the size of (JF3).
\end{proof}

\subsection{All isometric lattices use the same field}

Let $F_{L,L',S,c,x}$ denote the actual theta indicator divided
by its determinant-character factor. Its mixed derivatives
are the coefficients denoted by \emph{rational indices} in
the source. All of them are algebraic. Indeed the indicator
is a finite linear combination of rational translates of
the normalized theta vector; its coefficients are rational
numbers, square roots of positive integers, and roots of
unity. Corollary~\ref{gm-algebraic-jets}, with the coordinate
directions as in Lemma~\ref{csp-coordinate-units}, gives the
algebraicity of each derivative. Lemma~\ref{jcf-character-values}
permits division by its character factor.

\begin{proposition}\label{jcf-common-field}
There is a number field $K\subset\mathbb C$ containing all
mixed derivatives of order at most $d$ of
$F_{L,L',S,c,x}$ at zero, for every choice described above.
It also contains all the chosen direction, character, and
frame coefficients. It can be chosen Galois over $\mathbb Q$.
The original field $E$, primes, Hermitian lattices, weights,
and local moduli remain the fixed data in this assertion.
\end{proposition}
\begin{proof}
For each of the finitely many positive Hermitian spaces,
choose a finite set of representatives of integral self-dual
lattices, using Theorem~\ref{al-finite-classes}. For each
pair of representatives, each of the finitely many $S,M,c$,
and each class in (JF3), choose one theta indicator. There
are finitely many such functions. In a fixed finite number
of variables there are only finitely many monomials of total
degree at most $d$. Adjoin all the corresponding algebraic
mixed derivatives, all $c_\eta$, the normal closure of $E$,
and the finite character and frame data described above.
The resulting extension of $\mathbb Q$ is finite, since
adjoining a finite list of algebraic numbers gives a finite
extension by the tower formula.

Consider any other pair $\gamma L,\gamma'L'$ with
$\gamma,\gamma'$ rational unitary transformations. The
covariance in (TE17), applied term by term to the finite
indicator, carries its variable by
\[
                z_\eta\longmapsto
                             \eta(\gamma')^{-1}z_\eta\eta(\gamma).
                                                               \tag{JF4}
\]
It carries the lift of $x$ to the lift of $\gamma'x\gamma^{-1}$;
local quotients, their indices, and the Hermitian pairing
are preserved. Thus this is equality of the actual finite
averages, not merely equality up to a complex scalar.
After dividing by the determinant character, its remaining
scalar is a product of integer powers of
$\prod_\eta\eta(\det\gamma)$ and
$\prod_\eta\eta(\det\gamma')$. The principal-ideal rule
(JF2) gives this statement with no new roots. This scalar
and every entry of the linear substitution (JF4) belong to
the normal closure of $E$.

For clarity, if $G$ is holomorphic and $A$ is a linear
substitution whose entries lie in $K$, the ordinary chain
rule gives
\[
 D_{v_1}\cdots D_{v_j}(G\circ A)(0)
                    =D_{Av_1}\cdots D_{Av_j}G(0).        \tag{JF5}
\]
Expanding each $Av_i$ in the coordinate basis expresses the
right side as a finite sum of derivatives of order $j$ of
$G$, with coefficients in $K$. The factors $c_\eta$ in the
chosen directions and their inverses are already in $K$.
Therefore every required derivative for the new pair lies
in the same field. This covers all self-dual lattices, with
no new generators as the isometries vary.

Taking the normal closure of this finite field proves the
last field assertion. Complex conjugation then restricts
to an automorphism of $K$. This does not assert that $K$
is a CM field; no such assertion is needed here.
\end{proof}

\subsection{Order of choices and descent of the column identities}

The finite generator set just constructed depends only on
the data fixed before (JF1). In particular the primes in
$\mathcal B$ are primes of $E$, not primes of the enlarged
coefficient field. Field enlargement does not replace them
by a new list or change a lattice over $E$. The finite list
of sets $S\subseteq\mathcal B$ already includes every later
allowed selection, and (JF1) already includes every required
local modulus. The derivative bounds depend on the chosen
weights and ranks, which also precede the field enlargement.
An arbitrary locally constant $K$-valued test function on
one of the finite quotients is a $K$-linear combination of
its characteristic functions. Its values therefore do not
require further generators.

In particular there is no requirement that the selected
primes remain unramified in $K$. The ordinary moment and
valuation arguments use their specified splitting in $E$
and in the imaginary quadratic field. At any place of $K$
above such a prime, extending the embedding to $\mathbb C_p$
gives the same normalized absolute value and the same unit
bounds. Adjoining coefficients may ramify that place; it
does not alter those absolute values. Hence no further
prime selection is caused by Proposition~\ref{jcf-common-field}.

All equalities of theta coefficients and of finite class
sums which were proved over $\mathbb C$ now descend to $K$:
both sides have entries in $K$, and the embedding
$K\hookrightarrow\mathbb C$ is injective. In particular
this removes the coefficient-field proviso from the Hecke
commutation identity for these fixed columns. Every
required mixed derivative belongs to the source's set of
rational indices. The assertion is uniform over the fixed
finite construction; it does not assert that one number
field contains theta jets of unbounded order or at
unboundedly many independently chosen levels.


\section{The incidence identity and Hecke commutation}
\label{hec-section}

We prove the incidence identity for the actual affine matrix test
function and its passage to the theta columns. The scalar comes from
the change in the volume of the Hom lattice. The weights in the
global class sum are essential. All identities are first identities
of complex theta functions and their coefficient maps; the finite
coefficient-field construction then gives their arithmetic descent.

\subsection{Neighbors of a lattice with a residue functional}

Let $R$ be a discrete valuation ring with uniformizer $\varpi$,
fraction field $K_v$, and finite residue field $\mathbf k$ of order
$q$. A pair of rank $a$ consists of a free $R$-lattice $L\subset K_v^a$
and a nonzero $\mathbf k$-linear functional
$\ell:L/\varpi L\to\mathbf k$. We also write $\ell$ for its
composition with reduction on $L$. An upper neighbor of $(L,\ell)$ is
\[
 L^+=L+R\varpi^{-1}v,\qquad
 \ell^+(x)=\ell(\varpi x),\qquad v\in L,\quad\ell(v)=1.
 \tag{HC1}
\]
The displayed functional is well defined on $L^+/\varpi L^+$,
since $\varpi L^+\subset L$ and $\ell(\varpi y)=0$ for $y\in L$.
Its kernel in $L^+$ is exactly $L$.

\begin{lemma}\label{hec-neighbors}
There are $q^{a-1}$ distinct upper neighbors of $(L,\ell)$.
Every lower neighbor has underlying lattice $L^0=\ker\ell$.
Its functional $\ell^0:L^0/\varpi L^0\to\mathbf k$ is characterized
by
\[
             \ell^0(\varpi x)=\ell(x)\qquad(x\in L).
 \tag{HC2}
\]
There are $q^{a-1}$ such functionals, all nonzero.
\end{lemma}
\begin{proof}
Choose a basis $e_1,\ldots,e_a$ of $L$ with
$\ell(e_i)=0$ for $i<a$ and $\ell(e_a)=1$. The upper neighbors are
represented by $v=e_a+\sum_{i<a}u_i e_i$, $u_i\in\mathbf k$.
Indeed two such generators give the same lattice only if their
reductions span the same line in $L/\varpi L$; their last coordinates
are both one, so their reductions are equal. Formula (HC1) then fixes
the functional. The lattice $L^0$ has basis
$e_1,\ldots,e_{a-1},\varpi e_a$. Condition (HC2) prescribes the value
one on $\varpi e_a$ and leaves the other $a-1$ values arbitrary.
Conversely, for any such functional, the upper neighbor obtained
using $\varpi e_a$ is exactly $(L,\ell)$. The kernel identity in
(HC1) proves that every lower neighbor is of this form.
\end{proof}

For pairs $d=(L,\ell)$ and $d'=(M,\ell')$, of ranks $a,b$, and
$X\in\operatorname{Hom}_{K_v}(K_v^a,K_v^b)$, put
\[
 \phi_{d,d'}(X)=
  \mathbf1_{\{XL\subset M,\ \ell'X=\ell\text{ on }L\}}.
 \tag{HC3}
\]

\begin{proposition}\label{hec-local-count}
For every $X$,
\[
 \sum_{d_1:\ d'\to d_1}\phi_{d,d_1}(X)
       =\sum_{d_0:\ d_0\to d}\phi_{d_0,d'}(X).
 \tag{HC4}
\]
Both sums are either zero or one. Here arrows mean the upper-neighbor
relation (HC1).
\end{proposition}
\begin{proof}
Suppose $d'\to d_1=(M^+,\ell^+)$ and $\phi_{d,d_1}(X)=1$.
For $x\in L^0=\ker\ell$, the equality $\ell^+(Xx)=0$ gives
$Xx\in M$. Define $\ell^0(x)=\ell'(Xx)$ on $L^0$.
For $e\in L$ with $\ell(e)=1$,
\[
 \ell^0(\varpi e)=\ell'(\varpi Xe)=\ell^+(Xe)=1.
\]
More generally the same calculation gives (HC2), so $(L^0,\ell^0)$
is a lower neighbor and $\phi_{d_0,d'}(X)=1$.

Conversely suppose $d_0=(L^0,\ell^0)\to d$ and
$\phi_{d_0,d'}(X)=1$. Choose $e\in L$ with $\ell(e)=1$ and put
$v=X(\varpi e)\in M$. Then $\ell'(v)=1$. Define $M^+=M+R\varpi^{-1}v$
and its functional by (HC1). Since $L=L^0+Re$, the inclusion
$XL\subset M^+$ follows. For $x\in L^0$, $\ell^+(Xx)=0=\ell(x)$,
and for $x=e$ both values are one. Thus $\phi_{d,d_1}(X)=1$.

For uniqueness on the upper side, any solution satisfies
$Xe\in M^+$ and $\ell^+(Xe)=1$, hence $M^+=M+RXe$;
its functional is fixed by (HC1). On the lower side the underlying
lattice is $L^0$ and its functional must equal $\ell'X$.
The two constructions are inverse, which proves (HC4).
\end{proof}

For a self-dual Hermitian lattice over a CM field, changing the
component at a split prime $P$ forces the dual change at $\bar P$.
All other components and auxiliary frames remain fixed. Local patching
gives a unique global lattice, as in Lemma~\ref{acl-patching}.
Thus the same neighbor counts hold for the global pairs used here.
The test function is the product of (HC3) over the selected primes.
At a move at $P$, its other factors are identical on both sides of
(HC4); multiplying by them proves the global pointwise test identity.

\subsection{A finite theta calculation with a common reference lattice}

We explain the normalization in the lattice change, rather than
assuming that a finite identity of test functions automatically gives
an identity of differently normalized theta columns.
Fix positive even source and target ranks $a,b$.
Let $\mathcal S$ be a finite set of split primes, disjoint from its
conjugate. All period lattices in this paragraph agree outside
$\mathcal S\cup\overline{\mathcal S}$ and have dual components at
conjugate primes. Write $V_{\mathcal S}$ for the product of their
selected local vector spaces, with any one fixed additive Haar measure.
If $\mathcal L$ is one of these self-dual period lattices, write
$L=\prod_{P\in\mathcal S}\mathcal L_P$ for its selected component.
Write $v_{\mathcal L}$ for its normalized theta vector after removal
of the determinant character in (TE15). The exact lattice moves in
Theorem~\ref{te-equivariance} apply to these vectors.

For a compactly supported, locally constant function $\Phi$ on
$V_{\mathcal S}$, let $\mathcal T_{\mathcal L}(\Phi)$ be the finite
theta sum formed from the actual indicators (CS12a), with
$v_{\mathcal L}$ in place of $F$. Its support and its periods are
assumed to lie in the finite quotients used in that definition.
Increasing a support bound or refining a period does not change
this sum. The following proof includes that assertion.

\begin{lemma}\label{hec-common-theta}
For two such Hom period lattices $\mathcal L_1,\mathcal L_2$ and
one test function allowed for both,
\[
 \mathcal T_{\mathcal L_1}(\Phi)
   =\left(\frac{\operatorname{vol}L_2}
                  {\operatorname{vol}L_1}\right)^{1/2}
                                      \mathcal T_{\mathcal L_2}(\Phi).
 \tag{HC5}
\]
The square root is positive. This is an identity of entire functions.
\end{lemma}
\begin{proof}
Choose a sufficiently small selected Hom lattice $C$ contained in
$L_1\cap L_2$ and in the chosen periods of $\Phi$ for both lattices.
Choose a sufficiently
large selected Hom lattice $A$ containing $L_1,L_2$ and the support
of $\Phi$. These choices can be made within the Hom-lattice family:
enlarge the selected source and shrink the selected target to obtain
$C$, and do the reverse to obtain $A$. At conjugate primes impose
the dual changes and patch globally. Denote the resulting self-dual
period lattices by $\mathcal C,\mathcal A$.

Put $n_i=[L_i:C]$. The conjugate quotient dual to $L_i/C$ has order
$n_i$ and pairs perfectly with it by the Heisenberg commutator.
Consequently the projection of $v_{\mathcal L_i}$ onto the zero
selected coset modulo $C$ is
\[
 P_Cv_{\mathcal L_i}
       =\frac1{n_i}\sum_{y\in(L_i/C)^\vee}W_yv_{\mathcal L_i}
       =\frac1{\sqrt{n_i}}v_{\mathcal C}.
 \tag{HC6}
\]
The first equality is finite character orthogonality. The second is
the exact normalized lattice move, whose coefficient is
$n_i^{-1/2}$; moving the source and target consecutively gives the
same vector by the exact path independence proved in
Section~\ref{te-section}. This uses the positive normalization there.

Choose representatives of $A/C$ in $\mathcal A$ which are admissible
lifts for both original indicator systems. To do so, first choose
for each $i$ a positive denominator $c'_i$ clearing the selected
denominators of $A$ relative to $L_i$. At the conjugate selected
primes impose the additional divisibility by $c'_i$ required for
the half lattice in (CS12a), as well as membership in $\mathcal A$.
These are conditions at primes disjoint from the selected coset
conditions. The Chinese remainder theorem on $\mathcal A$ therefore
supplies all the required representatives simultaneously. For such a
representative $x$, the indicator of its coset applied to
$v_{\mathcal L_i}$ is
$n_i^{-1/2}W_{-x}v_{\mathcal C}$. These expressions are independent
of representatives: their difference lies in
$\mathcal A\cap\mathcal C$, the Weyl operators add on the integral
lattice $\mathcal A$, and this difference fixes $v_{\mathcal C}$.

For completeness, if $C\subset D\subset L_i$ is an intermediate
period (as for the actual scalar periods $cL_i$), the
identity of functions
$\mathbf1_{x+D}=\sum_{z\in D/C}\mathbf1_{x+z+C}$ gives the same
identity of these operators. Indeed the coefficient of a character
$y$ in the right side is multiplied by
\[
 \sum_{z\in D/C}e^{-2\pi iE(z,y)}
   =\begin{cases}|D/C|,&y\text{ is trivial on }D/C,\\
                 0,&\text{otherwise}.
     \end{cases}
\]
Together with the denominator $[L_i:C]$, this is exactly the
denominator and character sum for the coarser period $D$.
The lift independence in (CS12a) allows these common representatives
even when a support denominator exceeds the original period
denominator. Enlarging $A$ only inserts zero coefficients.
Thus the actual finite sum is
\[
 \mathcal T_{\mathcal L_i}(\Phi)
    =\frac1{\sqrt{n_i}}\sum_{x\in A/C}
                               \Phi(x)W_{-x}v_{\mathcal C}.
 \tag{HC7}
\]
Since $\operatorname{vol}L_i=n_i\operatorname{vol}C$, taking the
ratio proves (HC5). All sums are finite sums of translates of entire
theta functions, so the assertion holds as an identity of entire
functions, with no analytic convergence issue.
\end{proof}

Consider now ranks $a,b$ and lattices $N_1,N_2$ on the source and
$M_1,M_2$ on the target. They agree away from primes over a selected
rational prime $p$ which splits completely in $E$.
Thus $NP=p$ for each selected $P\mid p$, and $p$ is a local
uniformizer there. These are the actual selected-prime hypotheses.
For each such $P$, put
\[
 e_P=\operatorname{ord}_P\det N_2-\operatorname{ord}_P\det N_1,
 \qquad
 e'_P=\operatorname{ord}_P\det M_2-\operatorname{ord}_P\det M_1.
 \tag{HC8}
\]
The local period ideal is a unit. The determinant of
$\operatorname{Hom}(N,M)=M\otimes N^*$ therefore gives
\[
 \frac{\operatorname{vol}\operatorname{Hom}(N_2,M_2)_P}
      {\operatorname{vol}\operatorname{Hom}(N_1,M_1)_P}
                       =p^{b e_P-a e'_P}.
 \tag{HC9}
\]
This follows also from bases: an index $p^{e_P}$ change of $N$
changes the determinant of its dual $b$ times, and an index
$p^{e'_P}$ change of $M$ changes its determinant $a$ times.

The remaining character factor of the arithmetic column is
$D(N)/D(M)$. Away from the auxiliary level, its factor over a
prime of $F$ below $P$ is $w_P^{\operatorname{ord}_P\det N}$, where
$w_P$ is the nonzero unramified character value. If the selected
lift of that prime of $F$ is $\bar P$, both the determinant exponent
and the character are inverted, so the same factor results.
Here $w_{\bar P}=w_P^{-1}$ because the character is unramified at
both primes and trivial on the ground-field idele. Auxiliary frame
factors are unchanged under the present lattice move. Consequently
\[
 \frac{D(N_1)/D(M_1)}{D(N_2)/D(M_2)}
                         =\prod_{P\mid p}w_P^{e'_P-e_P}.
 \tag{HC10}
\]
Only the selected primes occur in this product.

\begin{proposition}\label{hec-kernel-move}
Let $\mathcal K_{N,M}(\Phi)$ be the actual theta coefficient map,
including the factors just described. Suppose the same test function
has the required support and periodicity for both pairs of lattices.
Then, over $\mathbb C$,
\[
 \mathcal K_{N_1,M_1}(\Phi)
   =\prod_{P\mid p}\left(w_P^{e'_P-e_P}
                    p^{(b e_P-a e'_P)/2}\right)
                                  \mathcal K_{N_2,M_2}(\Phi).
 \tag{HC11}
\]
Here $a,b$ are even, so the exponent is an integer.
\end{proposition}
\begin{proof}
Apply (HC5) and (HC9) to the determinant-normalized theta functions,
and then multiply by (HC10). The mixed derivatives defining the
column are linear operations on the same entire-function identity.
The polynomial coefficient extraction and the coefficient-module
maps are linear as well. Thus they preserve the scalar in (HC11).
No valuation estimate is substituted for this equality.
\end{proof}

\subsection{The pointwise column identity}

Fix pairs $d=(N,\ell)$ and $d'=(M,\ell')$, of even ranks $a,b$,
and a selected prime $P$. Write $N^0$ for the self-dual lattice
whose component at $P$ is $\ker\ell_P$, whose component at $\bar P$
is the dual one, and whose other components agree with $N$.
Let $\Phi_0$ be the prescribed auxiliary test function away from $p$.
The point kernel is
\[
          k(d,d')=\mathcal K_{N,M}(\phi_{d,d'}\Phi_0).
 \tag{HC12}
\]
All residue conditions at the other selected primes are included
in $\phi$. The auxiliary frames are fixed throughout each move.

\begin{theorem}\label{hec-point-identity}
Put $s=(b-a)/2$. Then
\[
 \sum_{d_1:\ d'\to d_1}k(d,d_1)
       =p^s\sum_{d_0:\ d_0\to d}k(d_0,d').
 \tag{HC13}
\]
\end{theorem}
\begin{proof}
If $d'\to d_1$ has underlying lattice $M_1$, the lower lattice
of $d_1$ is $M$. Change both lattices from $(N,M_1)$ to $(N^0,M)$.
At $P$ the two index exponents in (HC8) are $e_P=e'_P=1$,
and all other selected exponents vanish. Formula (HC11) therefore
has scalar exactly $p^{(b-a)/2}$; the character cancels.

The common test function is allowed for both pairs. If
$\phi_{d,d_1}(X)=1$, then $XN^0_P\subset M_P$ by the kernel
identity in (HC1). For periodicity, a period is divisible by $p$.
If $Z:N^0_P\to M_P$ is integral, then for $x\in N_P$ one has
$px\in N^0_P$ and
\[
                   pZx=Z(px)\in M_P=\ker\ell_{1,P}.
\]
Adding this period changes neither lattice inclusion nor the
residue equation. With the original pair, the same assertion
follows from $pM_{1,P}\subset\ker\ell_{1,P}$. Conditions at the
other selected primes are unchanged, and the lattices and frames
at the auxiliary primes are identical. Thus all support and
periodicity requirements in (HC11) hold. We obtain
\[
 k(d,d_1)=p^s\mathcal K_{N^0,M}(\phi_{d,d_1}\Phi_0).
\]
Every lower neighbor $d_0\to d$ has underlying lattice $N^0$,
so $k(d_0,d')=\mathcal K_{N^0,M}(\phi_{d_0,d'}\Phi_0)$.
Both sums are finite. Use linearity and (HC4) to conclude.
\end{proof}

\subsection{The weights in the global class sum}

Let a group $\Gamma$ act on a set $\mathscr X$ with finitely many
orbits and finite stabilizers. Let $\mathscr R\subset\mathscr X^2$
be an invariant relation with finite fibers in both directions.
Choose one representative $x_c$ of each orbit and put
$w_c=|\Gamma_{x_c}|^{-1}$.

\begin{lemma}\label{hec-incidence-mass}
For a simultaneous $\Gamma$-invariant vector-valued function $F$
on $\mathscr R$,
\[
 \sum_c w_c\sum_{y:\,(x_c,y)\in\mathscr R}F(x_c,y)
       =\sum_c w_c\sum_{x:\,(x,x_c)\in\mathscr R}F(x,x_c).
 \tag{HC14}
\]
\end{lemma}
\begin{proof}
There are finitely many $\Gamma$-orbits of related pairs, since
there are finitely many first-coordinate orbits and finite fibers.
For one such orbit, choose $(x,y)$ with first coordinate one of
the representatives. Its contribution to the left side is
\[
 \frac1{|\Gamma_x|}[\Gamma_x:\Gamma_x\cap\Gamma_y]F(x,y)
                 =\frac{F(x,y)}{|\Gamma_x\cap\Gamma_y|}.
\]
Computing the same orbit using the second coordinate gives exactly
the same value. Summing over these finitely many orbits proves (HC14).
\end{proof}

The coefficient spaces over $\mathscr X$ may carry representations
of $\Gamma$. Equivariant sections and equivariant kernels are allowed:
the expression $k(x,z)f(y)$ is simultaneously invariant in $(x,y)$,
because the source transformation of the kernel cancels that of $f$.
This is the only invariance required in (HC14).

\begin{theorem}\label{hec-global-identity}
On the actual finite spaces of equivariant forms, let $U_a,U_b$
sum over upper neighbors, and define the theta column by
\[
              (Kf)(d')=\sum_c w_c\,k(d_c,d')f(d_c).
 \tag{HC15}
\]
Then
\[
                  U_bK=p^{(b-a)/2}KU_a,
 \qquad U_b^rK=p^{r(b-a)/2}KU_a^r\quad(r\ge0).
 \tag{HC16}
\]
\end{theorem}
\begin{proof}
The class interpretation and its stabilizer weights were proved in
Sections~\ref{acl-section} and \ref{acc-section}. The local neighbors
there identify with (HC1); the residue functional records the
quotient map, so no further multiplicity occurs. Finiteness of the
fibers follows from Lemma~\ref{hec-neighbors}. Exact unitary
covariance in (TE17), together with the natural changes of auxiliary
frames, gives the equivariance of this kernel and of the forms.

Expand the left side of (HC16) using (HC15), interchange the two
finite sums, and use (HC13). The result, after division by $p^s$, is
\[
 \sum_c w_c\sum_{d_0:\,d_0\to d_c} k(d_0,d')f(d_c).
\]
Apply (HC14) with $F(x,y)=k(x,d')f(y)$. It equals
$\sum_c w_c k(d_c,d')\sum_{y:\,d_c\to y}f(y)=(KU_af)(d')$.
This proves the first equality, and induction proves the second,
including $r=0$.
\end{proof}

Proposition~\ref{jcf-common-field} constructs the number field $K$
containing the needed mixed derivatives, character values, and frame
coefficients for this fixed finite construction. All the coefficient
maps are therefore defined over $K$, and (HC16) descends to $K$ by
injectivity of $K\hookrightarrow\mathbb C$. This uses the finite
field construction, not merely the algebraicity of individual jets.
It discharges the commutation hypothesis in the finite valuation
iteration of Section~\ref{cv-section}. The exact source definitions
and the finite checks are identified in the accompanying trace
and manifest \cite{alpo-uniform-source}.

\begin{example}[Why the stabilizer weights are required]
Let $S_3$ act on one fixed point $*$ and on three permuted points.
Relate $*$ to each of the other three points and set $F=1$ on these
three pairs. There are two orbits of points. The two unweighted
class sums in (HC14) would be $3$ and $1$. The actual weights give
$3/6=1/2$ and $1/2$. Thus dropping the stabilizers would invalidate
the passage to class sums. They are present in the submitted proof.
\end{example}


\section{Local moments, integral kernels, and the column valuation bound}
\label{cv-section}

The integrality of the final lowering map does not by itself give the
stronger valuation estimate for a column of theta maps.  We prove the
finite algebraic passage from the required theta moments to that
estimate, with the actual determinant factors. The analytic input
is stated explicitly as a condition on the ordinary theta moments;
Section~\ref{tvc-section} assembles its proof for the actual kernel.
In particular, bounded ordinary power moments alone are insufficient.
The source of this part of the construction is the chain of local
kernel estimates and incidence operators in the submitted proof.
\cite{alpo-uniform-source}

\subsection{Polynomial substitution and its determinant factor}

Let $K$ be a number field, let $w$ be a finite prime of $K$, and let
$R=\mathcal O_{K,w}$.  We use the additive valuation $v_w$, with
$v_w(\varpi)=1$ for a uniformizer.  Let $I$ be a finite set.  Fix
$1\le a<b$ and dominant integral weights
$\lambda_\tau\in\mathbb Z^a$ and $\mu_\tau\in\mathbb Z^b$ for
$\tau\in I$.  Write $d_\tau=\lambda_{\tau,a}$ and
$d'_\tau=\mu_{\tau,b}$.  Suppose integers $s_\tau,s'_\tau$ and
nonnegative integers $n_\tau$ satisfy
\begin{align}
 \mu_{\tau,j}+s_\tau&=\lambda_{\tau,j}+s'_\tau
                  &&(1\le j\le a),\label{cv-weight-compatibility}\\
 \mu_{\tau,j}&=-s_\tau&&(a<j\le b),\notag\\
 n_\tau&=d_\tau+s'_\tau.\notag
\end{align}
Thus $d'_\tau=-s_\tau$.  The polynomial partitions underlying these
weights are $\nu_{\tau,j}=\lambda_{\tau,j}-d_\tau$ and
$\nu'_{\tau,j}=\mu_{\tau,j}-d'_\tau$; hence
\[
 \nu'_\tau=(\nu_{\tau,1}+n_\tau,\ldots,
              \nu_{\tau,a}+n_\tau,0,\ldots,0).
 \tag{CV1}
\]
Use the polynomial spaces and integral lattices of
Section~\ref{coef:section}, taking the tensor product over $\tau$.
We write $V_\lambda$ for the polynomial space with column degrees
$\nu_\tau$, and similarly for $V_\mu$.

If $M_\tau=B_\tau R^a$ is a full lattice, define its transported
dual-Weyl lattice $\mathcal N_\lambda(M)$ by the coefficient condition
\begin{equation}\label{cv-lattice-coordinates}
 F\in\mathcal N_\lambda(M)
 \quad\Longleftrightarrow\quad
 F_B(X):=\left(\prod_\tau\det(B_\tau)^{d_\tau}\right)
                       F((B_\tau X_\tau)_\tau)\in R[X].
\end{equation}
This is independent of the bases: a change by an element of
$\operatorname{GL}_a(R)$ and its inverse preserves integral
coefficients, and its determinant is a unit.  The corresponding Weyl
lattice is denoted $\mathcal D_\lambda(M)$; it is contained in
$\mathcal N_\lambda(M)$.  Formula
\eqref{cv-lattice-coordinates} includes the determinant twist even
when $d_\tau$ is negative.

For an $I$-tuple of $a$ by $b$ matrices $A=(A_\tau)$, define
\begin{equation}\label{cv-point-substitution}
 \begin{split}
 F^\sharp(X)&=F(X)\prod_\tau\det(X_\tau)^{n_\tau},\\
 Q_A F(Y)&=F^\sharp\bigl((A_\tau Y_{\tau,[1,a]})_\tau\bigr),
 \end{split}
\end{equation}
where $Y_{\tau,[1,a]}$ consists of the first $a$ columns of the $b$ by
$b$ matrix $Y_\tau$.  The matrices $A_\tau$ may have deficient rank.

\begin{lemma}\label{cv-substitution}
For every $A$, the polynomial $Q_AF$ belongs to $V_\mu$.
If $M'_{\tau}=B'_{\tau}R^b$, $A_\tau M'_\tau\subset M_\tau$,
and $F\in\mathcal N_\lambda(M)$, then
\begin{equation}\label{cv-point-integrality}
 Q_AF\in c(B,B')\mathcal N_\mu(M'),\qquad
 c(B,B')=\prod_\tau\det(B_\tau)^{s'_\tau}
                          \det(B'_\tau)^{-s_\tau}.
\end{equation}
The assertion remains valid after extension to any valued field,
with its valuation ring in place of $R$.
\end{lemma}
\begin{proof}
Multiplication by the determinant powers adds $n_\tau$ to every
column degree of $F$.  Substitution by $A_\tau$ commutes with right
column additions.  The unused columns have degree zero, so all
required column-invariance identities hold, including additions to an
unused column.  Equation (CV1) proves membership in $V_\mu$.

Put $A^0_\tau=B_\tau^{-1}A_\tau B'_\tau$.  The lattice inclusion is
equivalent to $A^0_\tau$ having entries in $R$.  Direct substitution,
using $n_\tau-d_\tau=s'_\tau$, gives
\begin{equation}\label{cv-exact-coordinate-factor}
 \left(\prod_\tau\det(B'_\tau)^{d'_\tau}\right)
                 (Q_AF)((B'_\tau Y_\tau)_\tau)
       =c(B,B')Q_{A^0}F_B(Y).
\end{equation}
The rightmost polynomial has integral coefficients: determinants,
multiplication, and substitution by integral matrices introduce no
denominators.  This proves \eqref{cv-point-integrality} and also proves
the assertion over an arbitrary valuation ring.
\end{proof}

\subsection{Finite moments and passage to a valuation ring}

Replace $A_{\tau,\alpha\beta}$ by independent variables
$Z_{\tau,\alpha\beta}$.  Expand the universal polynomial as
\begin{equation}\label{cv-universal-expansion}
                Q_ZF(Y)=\sum_m Z^m Q_m(F)(Y).
\end{equation}
Every occurring multi-index $m$ satisfies
\begin{equation}\label{cv-finite-degrees}
 \sum_{\alpha=1}^{a}\sum_{\beta=1}^{b}m_{\tau,\alpha\beta}
      =D_\tau:=\sum_{j=1}^{a}(\lambda_{\tau,j}+s'_\tau).
\end{equation}
In particular $D_\tau\ge0$, and the set $\mathcal M$ of such
multi-indices is finite.  For numbers $J_m\in K$, define the linear
operator
\[
                     \mathcal JF=\sum_{m\in\mathcal M}J_mQ_m(F).
 \tag{CV2}
\]
The column-invariance identities and degrees hold coefficient by
coefficient in the independent variables $Z$.  Thus
$\mathcal JF\in V_\mu$ for every choice of the $J_m$.

Fix an embedding $\iota:K\hookrightarrow L$ into a nonarchimedean
valued field such that $|\iota(x)|\le1$ if and only if $x\in R$.
Let $\mathcal O_L$ be its valuation ring.  For a family $\mathscr A$
of tuples of matrices over $L$, suppose each $A\in\mathscr A$
carries $M'_\tau\otimes_R\mathcal O_L$ into
$M_\tau\otimes_R\mathcal O_L$.  For $\kappa\in K$, impose the
following finite approximation condition:
\begin{equation}\label{cv-moment-approximation}
 \begin{gathered}
 \text{for every }\epsilon>0\text{ there are }A_1,\ldots,A_h\in
 \mathscr A\text{ and }c_1,\ldots,c_h\in\mathcal O_L\text{ such that}\\
 \left|\iota(J_m)-\iota(\kappa)\sum_{j=1}^h c_j A_j^m\right|
                  <\epsilon\quad(m\in\mathcal M).
 \end{gathered}
\end{equation}
The same finite sum must approximate every required moment.  Its
length and coefficients may depend on $\epsilon$.

\begin{theorem}\label{cv-moment-integrality}
Under \eqref{cv-moment-approximation},
\[
 \mathcal J\mathcal N_\lambda(M)
       \subset \kappa c(B,B')\mathcal N_\mu(M').
 \tag{CV3}
\]
If $\kappa=0$, then $\mathcal J=0$.
No completeness, compactness, or integral basis for $\mathcal O_L$
is needed in this assertion.
\end{theorem}
\begin{proof}
Fix $F\in\mathcal N_\lambda(M)$.  By
Lemma~\ref{cv-substitution}, each polynomial
\[
 \iota(\kappa)\sum_jc_jQ_{A_j}(\iota F)
\]
belongs to $\iota(\kappa c(B,B'))$ times the target integral
polynomial lattice over $\mathcal O_L$.
After applying the target coordinate transformation in
\eqref{cv-lattice-coordinates}, its coefficients converge to those
of the similarly transformed $\iota(\mathcal JF)$ as
$\epsilon\to0$.  Indeed, each coefficient is a fixed finite linear
combination of the errors in
\eqref{cv-moment-approximation}; the ultrametric inequality bounds
its absolute value by $C_F\epsilon$, where $C_F$ is the maximum of
the finitely many absolute values of its multipliers.

For a nonzero $q\in L$, the ball $q\mathcal O_L$ is closed in the
following elementary sense.  If $x$ can be approximated arbitrarily
closely by $qz$ with $|z|\le1$, choose an approximation with error at
most $|q|$; then $|x|\le|q|$.  For $q=0$, arbitrary approximation
means $x=0$.  Apply this observation coefficient by coefficient with
$q=\iota(\kappa c(B,B'))$.  When $\kappa\ne0$, divide the
coefficients by $q$ and use the defining property of $\iota$ to
descend integrality to $R$.  This proves (CV3).  When $\kappa=0$,
all coefficients vanish, proving the last assertion.
\end{proof}

\begin{lemma}[A bounded measure supplies the finite approximations]
\label{cv-measure-approximation}
Let $X$ be a compact space, let $L$ be complete, and let
$\ell:C(X,L)\to L$ be a continuous linear functional of norm at
most one.  Suppose $a:X\to\mathscr A$ is continuous, and
$J_m=\kappa\ell(a^m)$ for the finitely many relevant $m$.
Then \eqref{cv-moment-approximation} holds, provided $X$ has a
basis of clopen sets.
\end{lemma}
\begin{proof}
For any prescribed error, continuity and compactness give a finite clopen
partition $X=\coprod_{j=1}^h X_j$ such that all the functions $a^m$
oscillate by less than that error on each part.  Such a partition is
obtained by taking a finite clopen subcover and successively removing
earlier sets.  Omit empty parts, choose $x_j\in X_j$, and put
$c_j=\ell(1_{X_j})$.  Since $\|1_{X_j}\|_\infty=1$, one has
$|c_j|\le1$.  Applying $\ell$ to the step-function approximation
gives the required simultaneous estimates.  If $\kappa\ne0$,
divide the prescribed error by $|\kappa|$ first; if $\kappa=0$
the assertion is immediate.
\end{proof}

\subsection{The selected and conjugate primes}

Let $E/F$ be a CM extension with $[E:\mathbb Q]=2f$ and conjugation
$x\mapsto\bar x$.  Assume that a rational prime $p$ and an algebraic
integer $\pi$ satisfy
\[
 \pi\bar\pi=p,\qquad (\pi)+(\bar\pi)=\mathcal O_E,\qquad
 (\pi)=\prod_{P\in\mathcal P}P,\qquad
 |\mathcal P|=f,\quad\mathcal O_E/P=\mathbb F_p.
 \tag{CV4}
\]
Then $p$ is a uniformizer at every $P\in\mathcal P$ and at every
$\bar P$.  Let $N\subset E^a$ and $N'\subset E^b$ be self-dual
lattices for $h(x,y)=\bar x^{\,t}y$.  At every selected prime choose
nonzero linear functionals
\[
 \ell_P:N_P\to\mathbb F_p,\qquad
 \ell'_P:N'_P\to\mathbb F_p,
\]
linear over the local ring through its residue map.  Put
\[
 L_P=\ker\ell_P,\qquad L'_P=\ker\ell'_P,
\]
and let $L,L'$ agree with $N,N'$ at all other primes.  A superscript
$\#$ denotes the Hermitian dual lattice.

\begin{lemma}\label{cv-kernel-lattices}
Suppose $N_P=\mathcal O_{E,P}^a$ and
$N'_P=\mathcal O_{E,P}^b$ for all $P\in\mathcal P$.
Then
\[
 (L^\#)_P=N_P,\quad (L'^\#)_P=N'_P,
 \qquad v_P(\det L_P)=v_P(\det L'_P)=1,
 \tag{CV5}
\]
where determinant valuations are relative to the standard basis.
At conjugate primes,
\[
 v_{\bar P}(\det(L^\#)_{\bar P})
 =v_{\bar P}(\det(L'^\#)_{\bar P})=-1.
 \tag{CV6}
\]
In (CV5) a basis can be chosen with determinant exactly $p$.
\end{lemma}
\begin{proof}
Choose a coordinate $k$ on which $\ell_P$ is nonzero and divide its
residue row by that nonzero coefficient.  Lift the other coefficients
to $a_j\in\mathcal O_{E,P}$.  A basis of its kernel is
\[
                  p e_k,\qquad e_j-a_j e_k\quad(j\ne k).
\]
With the columns in their original order, its determinant is $p$.
The same proof applies to $L'_P$.
Hermitian duality exchanges $P$ and $\bar P$, and in bases sends
$B$ to $\bar B^{-t}$.  Since $L_{\bar P}=N_{\bar P}$ and $N$ is
self-dual, $(L^\#)_P=N_P$.  The inverse determinant gives (CV6).
These arguments use no division by two.
\end{proof}

Define two subsets of $\operatorname{Hom}_E(E^a,E^b)$:
\begin{align}
 \mathscr S_P&=\{x:x^*L'_P\subset N_P
                                \text{ for every }P\in\mathcal P\},
                         \label{cv-supports}\\
 \mathscr S_{\bar P}&=\{x:xN_P\subset N'_P,
                \ \ell'_P(xy)=\ell_P(y)
                \text{ for every }y\in N_P, P\in\mathcal P\}.
                         \notag
\end{align}
Here $x^*=\bar x^{\,t}$.

\begin{lemma}\label{cv-support-duality}
If $x\in\mathscr S_{\bar P}$, then
\[
               x^*(L'^\#)_{\bar P}\subset(L^\#)_{\bar P}
                          \quad(P\in\mathcal P).
 \tag{CV7}
\]
\end{lemma}
\begin{proof}
The functional identity implies $xL_P\subset L'_P$.  For
$z\in(L'^\#)_{\bar P}$ and $y\in L_P$, adjointness gives
$h(y,x^*z)=h(xy,z)\in\mathcal O_{E,\bar P}$.
The local description of the Hermitian dual therefore places
$x^*z$ in $(L^\#)_{\bar P}$.  This description follows directly
in a local basis from $L^\#_{\bar P}=\bar B^{-t}
\mathcal O_{E,\bar P}^a$.
\end{proof}

For the coefficient field $K$, take a set $I$ of $f$ embeddings
$\tau:E\hookrightarrow K$.  At the prime $w\mid p$ suppose either
all primes $\tau^{-1}w$ belong to $\mathcal P$, or all belong to
$\overline{\mathcal P}$.  These are respectively the two cases
below.  Write $e_w=v_w(p)$; for $z\in E$, at either corresponding
prime $Q$ one has
\[
                         v_w(\tau(z))=e_wv_Q(z),
 \tag{CV8}
\]
because $p$ is a uniformizer at $Q$.

\begin{quote}
\textbf{The required ordinary theta moments.}
\label{cv-theta-moments}
For the actual theta transition from rank $a$ to rank $b$, write its
polynomial kernel as $\eta\mathcal J$, where
$\eta\in K^\times$ is its algebraic character factor and
$v_w(\eta)=0$.  The indices used by $\mathcal J$ are exactly
\eqref{cv-finite-degrees}.  There are numbers
$u_\tau\in\mathbb C_p$ with $|u_\tau|=1$ such that its moments
$J_m$ satisfy \eqref{cv-moment-approximation} with matrices
\[
                         A_\tau(x)=u_\tau\,\iota(\tau(x^*))
 \tag{CV9}
\]
and with the following support and scalar:
\[
 \begin{array}{c|c|c}
 \tau^{-1}w\in\mathcal P&x\in\mathscr S_P&\kappa=p^{-fa}\\
 \tau^{-1}w\in\overline{\mathcal P}&x\in\mathscr S_{\bar P}
                                                  &\kappa=1.
 \end{array}
 \tag{CV10}
\]
The approximation is simultaneous for all needed indices, and uses
coefficients of absolute value at most one.
\end{quote}

This hypothesis is a finite consequence of the stronger measure
statement used in the source.  That statement constructs a measure
of norm at most one for the Taylor coefficients of a normalized CM
theta function at prime-to-$p$ torsion translates, after multiplying
the differentiation directions by one common $p$-adic unit period.
It then applies finite translations and the indicated support
conditions.  The measure statement has not been proved merely by
the finite-dimensional arguments above.  It must be supplied with
the actual normalization of the theta function and of its period.

\begin{theorem}[The two local kernel bounds]\label{cv-two-kernels}
Assume $N_P=\mathcal O_{E,P}^a$ and
$N'_P=\mathcal O_{E,P}^b$ for every $P\in\mathcal P$.
Assume the ordinary theta moment condition above and equality of the target
Weyl and dual-Weyl lattices at $w$.  Suppose
$s_\tau-s'_\tau=s_0\ge0$ is independent of $\tau$.
For $F\in\mathcal D_\lambda(L^\#)$, the actual kernel $K=\eta\mathcal J$
satisfies
\begin{align}
 p^{fa+\sum_\tau s_\tau}KF
       &\in\mathcal D_\mu(L')
       &&\text{if }\tau^{-1}w\in\mathcal P,\label{cv-P-bound}\\
 KF&\in p^{fs_0}\mathcal D_\mu(L'^\#)
       &&\text{if }\tau^{-1}w\in\overline{\mathcal P}.
                                                      \label{cv-bar-bound}
\end{align}
\end{theorem}
\begin{proof}
In the selected case the source lattice is standard at every
embedding.  Take its basis to be the identity.  By
Lemma~\ref{cv-kernel-lattices}, choose target bases of $L'$ with
determinant $p$.  Matrices (CV9) carry the target lattice into the
source lattice, by \eqref{cv-supports}; multiplication by $u_\tau$
does not change integrality.  The determinant factor of
Lemma~\ref{cv-substitution} is therefore
\[
                         c(B,B')=p^{-\sum_\tau s_\tau}.
\]
Theorem~\ref{cv-moment-integrality}, with $\kappa=p^{-fa}$,
proves \eqref{cv-P-bound} with target $L'$, first in the dual-Weyl
lattice.  The input Weyl lattice is contained in the source
dual-Weyl lattice, and the assumed target equality finishes the
argument.  The scalar $\eta$ is a unit.

In the conjugate case, Lemma~\ref{cv-support-duality} gives the
required matrix inclusions.  Both source and target dual lattices
have determinant of valuation $-e_w$ in each embedding.  Thus
\[
 v_w(c(B,B'))=e_w\sum_\tau(s_\tau-s'_\tau)=e_wfs_0.
\]
Here $\kappa=1$.  The same theorem gives
$KF\in c(B,B')\mathcal N_\mu(L'^\#)
=p^{fs_0}\mathcal N_\mu(L'^\#)$.  Equality of the two target
lattices gives \eqref{cv-bar-bound}.
\end{proof}

The target equality in this theorem is now supplied by
Corollary~\ref{fh-large-prime-equality}.  Its sufficient condition,
for a target weight $\mu$ of length $b$, is
\[
                       \mu_1-\mu_b+b-1\le p.
 \tag{CV11}
\]
The exact hook-product calculation in
Theorem~\ref{fh-hook-product} proves that
$[H_{\mu-\mu_b},H_{\mu-\mu_b}]$ is then a $p$-adic unit.  The
normalized pairing is integral on the dual-Weyl lattice itself,
which is therefore contained in its pairing dual, the Weyl lattice.
The reverse inclusion always holds.  The actual submitted proof
assumes the stronger inequality $\mu_1-\mu_b+b<p$, so its precise
input is covered.  No equality at arbitrary primes is asserted.

\subsection{Reduction to standard local lattices}

The standard-lattice assumption in Lemma~\ref{cv-kernel-lattices}
does not restrict the global class of a self-dual lattice.
We give the approximation argument needed for that assertion.

\begin{lemma}\label{cv-unitary-approximation}
Let $\mathcal P$ be a finite set of primes of a CM field $E$ such
that no member is the conjugate of another.  For a full lattice
$N\subset E^n$ there exists $g\in\operatorname{GL}_n(E)$ with
$g^*g=1$ and
\[
                     (gN)_P=\mathcal O_{E,P}^{n}
                          \quad(P\in\mathcal P).
 \tag{CV12}
\]
If $N$ is self-dual, the same equality holds at every $\bar P$.
\end{lemma}
\begin{proof}
At a split place of $F$ below $P$, a unitary matrix is determined by
an arbitrary matrix $G_P\in\operatorname{GL}_n(E_P)$; its conjugate
component is $(G_P^*)^{-1}$.  Choose $G_P$ carrying $N_P$ to the
standard lattice.  The property of carrying this lattice to the
standard lattice is open: after composing with $G_P^{-1}$, it is
membership in $\operatorname{GL}_n(\mathcal O_{E_P})$.

We prove simultaneous approximation by rational unitary matrices.
The Cayley transform
\[
              C(A)=(1-A)(1+A)^{-1},\qquad A^*=-A,
 \tag{CV13}
\]
is unitary wherever $1+A$ is invertible.  Its inverse is
$A=(1-G)(1+G)^{-1}$ on $\det(1+G)\ne0$.
At a split place the set of $G$ with $\det(1+G)\ne0$ is dense in
$\operatorname{GL}_n(E_P)$.  To see this, replace $G$ by $zG$ with
$z$ arbitrarily close to $1$; the polynomial $\det(1+zG)$ has only
finitely many roots and is not identically zero.  Thus we may choose
the local $G_P$ in this set while preserving its lattice property.

Approximate the entries of their Cayley inverses simultaneously by
entries of a matrix $B$ over $E$, with $B$ close to the prescribed
$A_P$ at $P$ and close to zero at every $\bar P$.  This is the usual
finite-prime approximation obtained by the Chinese remainder
theorem after clearing denominators.  Then $A=B-B^*$ is
skew-Hermitian and is close to $A_P$ at every selected prime.
The determinant $\det(1+A)$ is nonzero when the approximation is
sufficiently close.  The rational matrix $g=C(A)$ then has the
required local lattice property.  All continuity statements concern
finitely many entries and determinants.  This proof works also in
residue characteristic two, since its fields have characteristic zero.

For the last assertion, local Hermitian duality gives
$(gN)_{\bar P}=((gN)_P)^\#$, with conjugate coefficients.
The standard lattice is its own corresponding dual.
\end{proof}

\begin{corollary}[Transport to arbitrary representatives]
\label{cv-kernel-transport}
Suppose the ordinary theta moment condition holds for all representatives
which are standard at the selected primes.  Suppose also that the actual
kernel is equivariant under rational unitary changes of representatives:
if $g,g'$ act on its source and target lattices, and $T_g,T_{g'}$ are
the induced coefficient-space isomorphisms, then
\[
 T_{g'}K_{N,N'}=K_{gN,g'N'}T_g.
 \tag{CV13a}\label{cv-kernel-equivariance}
\]
Under the other hypotheses of Theorem~\ref{cv-two-kernels}, its two
conclusions hold for arbitrary self-dual $N,N'$.
\end{corollary}
\begin{proof}
Choose $g,g'$ by Lemma~\ref{cv-unitary-approximation}.
The changes of representatives preserve
\eqref{cv-supports}: replace $x$ by $g'xg^{-1}$, the two lattices
by $gN,g'N'$, and the residue functionals by composition with the
inverse.  Adjointness uses $g^{-1}=g^*$, and coefficient lattices
transform by the actions already proved in
Section~\ref{coef:section}.  In particular $T_g$ sends the source
Weyl lattice onto the transformed source Weyl lattice, and $T_{g'}$
sends each target Weyl lattice onto its transformed version.
Apply Theorem~\ref{cv-two-kernels} to $gN,g'N'$ and pull the two
inclusions back through $T_{g'}$, using \eqref{cv-kernel-equivariance}.
\end{proof}

\subsection{Algebraic character factors and stabilizer denominators}

\begin{lemma}\label{cv-character-unit}
Let $\chi$ be a Hecke character with conductor prime to
$p$, with values in a number field $K$, such that a positive power
of $\chi$ is algebraic.  An ideal or finite-idele
argument which is a unit at every prime over $p$ has character value
a unit at every $w\mid p$.  Finite-order character factors always
have this property.  The same assertion holds for a character on
norm-one fractional ideals $I\bar I=1$ whose value on a principal
norm-one ideal $(z)$ is a product of integral powers of the embeddings
of $z$.
\end{lemma}
\begin{proof}
First replace $\chi$ by its stated positive power; vanishing of
the valuation of a power implies vanishing of the original valuation.
For an ideal prime to the conductor and to $p$, raise its ray class
to a positive power which is principal, say generated by $z$ with
$z\equiv1$ at the conductor.  Since the ideal is prime to $p$,
$z$ is a unit at every prime over $p$.  The defining principal-ideal
formula for an algebraic Hecke character expresses this power of
its value as a product of integral powers of the conjugates of $z$.
All those factors are $w$-adic units.  Taking valuations shows that
the original value is also a unit.  Changing the ideal representative
or inserting local unit factors introduces only the same principal
formula and finite-order factors.  A root of unity has valuation
zero because a positive power is one.  This proves the assertion
for finite ideles as well.  The finiteness of the ray class group
used here follows from finiteness of the ordinary class group and
of the residue unit group at the conductor.

For a norm-one ideal $I$, choose $h>0$ with $I^h=(a)$.
Then $I^{2h}=(a/\bar a)$ and
$(a/\bar a)\overline{(a/\bar a)}=1$.  If $I$ is a unit above $p$,
then $a$ and $\bar a$ are units there.  Apply the stated
principal-ideal formula to $I^{2h}$ and take valuations.
\end{proof}

In the actual standard representatives the determinant ideals of
$N,N'$ are units above $p$.  The extra level primes are disjoint
from the selected primes, and the finite auxiliary twists are
unramified there.  The possibly half-integral unitary character types
become algebraic after squaring; equivalently their local values
satisfy the positive-power principal formulas just used.  The
character scalar $\eta$ is a ratio of the
corresponding determinant and local-frame character values.
Lemma~\ref{cv-character-unit} therefore gives $v_w(\eta)=0$.
This verification uses the prime-to-$p$ level condition; it cannot
be inferred solely from the absolute value of a complex character.

\begin{lemma}\label{cv-stabilizer-unit}
Let $\Gamma$ be a finite subgroup of $\operatorname{GL}_n(E)$, where
$[E:\mathbb Q]=d$.  Every prime divisor $p$ of $|\Gamma|$ satisfies
$p\le dn+1$.  Hence the stabilizer orders occurring at ranks at
most $N_m$ are units at $w\mid p$ whenever
\[
                              p>2fN_m+1.
 \tag{CV14}
\]
The same conclusion holds for stabilizers which inject into those
matrix groups.
\end{lemma}
\begin{proof}
Cauchy's theorem gives an element $g$ of order $p$.  Regard it as
a rational linear map on the $dn$-dimensional underlying vector
space.  Its minimal polynomial divides $X^p-1$ and is not $X-1$.
The polynomial $\Phi_p(X)=1+X+\cdots+X^{p-1}$ is irreducible over
$\mathbb Q$: after replacing $X$ by $X+1$, Eisenstein's criterion
at $p$ applies.  Thus the minimal polynomial has degree at least
$p-1$, and $p-1\le dn$.

The stabilizers of the decorated positive definite lattices are
finite by Section~\ref{al-section}.  Their projection to the rational
unitary matrices is injective: once the rational matrix and the
decorated lattice are fixed, every auxiliary level coordinate of a
stabilizing element is determined by its action on the prescribed
frame or residue marking.  This is the stabilizer identification of
Sections~\ref{acl-section} and~\ref{acc-section}.  Apply the first
assertion with $d=2f$ and $n\le N_m$.
\end{proof}

Thus a class-sum kernel with denominators $|\Gamma_c|$ preserves
any common local lattice containment satisfied by its individual
terms.  Multiplication by those denominators costs no valuation at
the selected prime.  This is stronger at that prime than the general
denominator bound in Section~\ref{cp-section}.

\subsection{Incidence sums and removal of one residue condition}

We give the local lattice argument that produces the useful gain.
For a local point $q=(N,\ell)$, let $L(q)=\ker\ell$.  An upper
neighbor $q^+=(N^+,\ell^+)$ has
\[
          N\subset N^+,\qquad L(q^+)=N,
 \tag{CV15}
\]
with the prescribed residue normalization.  The upper-neighbor
sum is an ordinary finite sum over these points.  The degree-one
residue field and the chosen normalization make the neighbor set
finite; the argument below needs no division by its cardinality.

For a subset $V\subset\mathcal P$, let $\mathcal L_i(V,q)$ be the
coefficient lattice obtained by using $L(q)_P$ at $P\in V$ and
$N(q)_P$ at $P\notin V$, with all other local conditions unchanged.
At conjugate primes use the dual lattices.  Fix a coefficient prime
$w$ of the selected type.  The map $\tau\mapsto P(\tau)=\tau^{-1}w$
is a bijection from $I$ to $\mathcal P$.  Indeed, (CV4) makes $p$
split completely in $E$, so the completion at every such prime is
$\mathbb Q_p$.  Two embeddings into $\mathbb C_p$ inducing the same
prime extend to the same embedding of $\mathbb Q_p$, and therefore
coincide.  The embeddings in $I$ are distinct and both sets have
cardinality $f$.

\begin{lemma}\label{cv-incidence-lattice}
At the fixed selected-type coefficient prime $w$, if $P\in V$,
replacing $q$ by a neighbor at $P$ gives
\[
             \mathcal L_i(V,q^+)=\mathcal L_i(V\setminus\{P\},q).
 \tag{CV16}
\]
Consequently the upper-neighbor sum at $P$ sends sections integral
in the first family of lattices to sections integral in the second.
\end{lemma}
\begin{proof}
At $P$ the equality is precisely $L(q^+)_P=N(q)_P$ from (CV15).
Every other selected-prime component is unchanged.  The conjugate
component of the global self-dual neighbor can also change, but none
of the embeddings defining the coefficient lattice at this $w$
sees a conjugate prime.  Transported polynomial lattices localized
at $w$ depend only on the indicated selected-prime components,
giving (CV16).
Each summand therefore lies in the same target lattice, and their
finite sum does also.  The equivariance factors used to express the
sum in class representatives preserve these transported lattices.
\end{proof}

\begin{lemma}\label{cv-erase-cost}
Let $M^0\subset M$ be a local sublattice of residue index $p$,
and let $\lambda$ be a dominant integral weight of length $n$.
Then
\[
                  p^{\lambda_1}\mathcal D_\lambda(M^0)
                              \subset\mathcal D_\lambda(M).
 \tag{CV17}
\]
For tensor factors at $w$, the exponent is the sum of
$\lambda_{\tau,1}$ over the embeddings which see that prime.
The statement permits negative $\lambda_1$.
\end{lemma}
\begin{proof}
Choose a basis in which $M=R^n$ and
$M^0=\operatorname{diag}(p,1,\ldots,1)R^n$; a unit determinant
change is harmless.  The round action on a row-weight component
of polynomial row degree $r_1$ multiplies it by
$p^{-d-r_1}$, where $d=\lambda_n$.
Every monomial of $V_{\lambda-d}$ has row degree
$r_1\le\lambda_1-d$, by the weight-majorization proof of
Section~\ref{am-section}.  Hence
$\lambda_1-d-r_1\ge0$.
The Weyl lattice is the sum of its integral row-weight components:
its generators are iterated divided row operators applied to the
highest vector, each of which has a definite row weight.  Multiplying
each component by the indicated nonnegative power of $p$ preserves
the Weyl lattice.  This proves (CV17).  Multiplying the powers in
the different tensor factors proves the last assertion.
\end{proof}

\subsection{The one-step bound on the actual class lattices}

Consider the transition from rank $N_i$ to $N_{i+1}$, with
$N_{i+1}-N_i=t+1=2s_0$.  Let $\mathscr S_i$ be its local lattice
of equivariant class sections.  The original coefficient lattice at
each point is obtained from its $L^\#$, as above.  Fix a set
$\mathcal T\subset I$ of distinguished embeddings and let
$V_w\subset\mathcal P$ be the corresponding selected primes.
Write $U_i$ for the product of their upper-neighbor sums at rank
$N_i$.  Operators at different primes commute, since they change
independent local components.

Suppose nonnegative integers $W_\tau$ satisfy, for every step,
\begin{equation}\label{cv-weight-bounds}
 0\le s_{i,\tau}\le W_\tau,\qquad
 (\mu_{i,\tau})_1\le W_\tau\quad(\tau\notin\mathcal T),
 \qquad N_i\le N_m.
\end{equation}
Assume the moment hypothesis and target Weyl/dual-Weyl equality
at each step, as well as \eqref{cv-kernel-equivariance} and (CV14).
Then the class-sum operator $K_i$
satisfies
\begin{equation}\label{cv-one-step}
 p^e U_{i+1}K_i\mathscr S_i\subset\mathscr S_{i+1},
 \qquad e=2\sum_\tau W_\tau+fN_m
\end{equation}
at selected-type primes.  At conjugate-type primes it satisfies
$K_i\mathscr S_i\subset\mathscr S_{i+1}$.

\begin{proof}
By \eqref{cv-P-bound}, multiplying the point kernel by
$p^{fN_i+\sum s_{i,\tau}}$ puts its value in the coefficient
lattice having every selected residue condition imposed.  The
class-sum denominator is a unit by
Lemma~\ref{cv-stabilizer-unit}; sums retain this containment.
Apply each upper-neighbor sum at a prime of $V_w$.
Lemma~\ref{cv-incidence-lattice} removes exactly that residue
condition.  For all remaining primes apply
Lemma~\ref{cv-erase-cost}.  The additional exponent is
$\sum_{\tau\notin\mathcal T}(\mu_{i,\tau})_1$; if some terms
are negative, the same inclusion holds and the upper bound below
only becomes easier.  The total exponent is at most
\[
 fN_i+\sum_\tau s_{i,\tau}
       +\sum_{\tau\notin\mathcal T}(\mu_{i,\tau})_1
                    \le fN_m+2\sum_\tau W_\tau=e.
\]
Any remaining nonnegative power of $p$ preserves the target
lattice.  This proves \eqref{cv-one-step}.
In the conjugate case, \eqref{cv-bar-bound} already gives a positive
power $p^{fs_0}$ times the original target lattice, and the same
stabilizer argument applies.
\end{proof}

\subsection{Accumulation along the column}

The preceding estimates must be combined with the correct direction
of the incidence commutation law.  Set $r=|V_w|$ and suppose
\begin{equation}\label{cv-commutation}
                         U_{i+1}K_i=p^{s_0r}K_iU_i.
\end{equation}
This is the finite incidence identity of the theta kernels; it is
an additional identity, not a consequence of their integral moments.
Let
$C_i=K_{i-1}\cdots K_0$, with $C_0=1$.

\begin{proposition}\label{cv-column}
Assume \eqref{cv-one-step} and \eqref{cv-commutation} for
$0\le i<m$.  Then
\begin{equation}\label{cv-column-bound}
 p^{b_i}C_iU_0^i\mathscr S_0\subset\mathscr S_i,\qquad
 b_i=ie+s_0r\frac{i(i+1)}2.
\end{equation}
In particular, if $A\ge0$ and
$\mathscr N=p^AU_0^m\mathscr S_0$, then for any $c\in K$ with
\[
                          v_w(c)\ge e_wb_m
 \tag{CV18}
\]
one has $cC_m\mathscr N\subset\mathscr S_m$.
\end{proposition}
\begin{proof}
Commuting one $U$ through the first $i$ transitions gives
\[
                         U_iC_i=p^{s_0ri}C_iU_0.
 \tag{CV19}
\]
This follows by induction directly from \eqref{cv-commutation}.
The assertion \eqref{cv-column-bound} is trivial for $i=0$.
Apply $p^eU_{i+1}K_i$ to its left side.  By
\eqref{cv-one-step}, the image is integral.  Equations
\eqref{cv-commutation} and (CV19) identify that image as
\[
 p^{b_i+e+s_0r(i+1)}C_{i+1}U_0^{i+1}\mathscr S_0.
\]
Since $b_{i+1}=b_i+e+s_0r(i+1)$, induction proves the formula.
Finally, (CV18) makes $cp^{A-b_m}$ an element of $R$.
Multiply the $i=m$ inclusion by this scalar.
\end{proof}

For the actual set of distinguished embeddings, $r=|\mathcal T|$.
Since $t+1=2s_0$, the final exponent can be written
\begin{equation}\label{cv-final-exponent}
 b_m=mE_p,\qquad
 E_p=2\sum_\tau W_\tau+fN_m
                    +|\mathcal T|\frac{(t+1)(m+1)}4.
\end{equation}
Although the last displayed summand can be a half-integer, $b_m$
is an integer because $m(m+1)$ is even.
The initial lattice in the submitted construction has exactly the
form $p^AU_0^m\mathscr S_0$, with
$A=m\sum_{P\in V_w}b_{\rm top}(\tau(P))\ge0$.
At conjugate-type primes $V_w$ is empty and the needed exponent is
zero.  Thus the finite argument gives the actual column valuation
bound, on its stated initial lattice, once the local moment and
incidence inputs have been supplied.  It does not assert the same
bound on every vector of the original unmodified class lattice.

The final all-twist lowering theorem of
Section~\ref{atl-section} preserves these local integral lattices.
Applying it introduces no further denominator.  Pairing with a row
then uses the finite-orbit integrality theorem of
Section~\ref{cp-section}; its stabilizer denominator is a unit at
$p$ by (CV14).  These two facts connect the local column bound to
the actual period matrix entry, without a stabilizer average in the
lowering map.

\subsection{Counterexamples to weakened assertions and remaining inputs}

The support and normalization conditions above have concrete force.
Here are elementary checks which prevent stronger, false replacements.

\begin{enumerate}
\item Take $a=1$, $b=2$, $M=\mathbb Z_p$, $M'=\mathbb Z_p^2$,
$\lambda=(0)$, $\mu=(1,0)$, $s=0$, and $s'=1$.
These weights satisfy the stated compatibility conditions, with
$n=1$ and the canonical source polynomial $F=1$.
For $A=(p^{-1},0)$, one has
$Q_AF=p^{-1}Y_{11}$, which is not integral.  Thus the matrix
support condition cannot be omitted from the finite moment theorem.
\item Bounding each ordinary power moment by one does not give a
measure of norm at most one on $\mathbb Z_p$.  For example, let
$J_1=1$ and $J_j=0$ for every $j\ne1$, including $j=0$.
The induced functional on polynomials is differentiation at zero.
It takes the integer-valued polynomial $\binom Xp$ to
$(-1)^{p-1}/p$.  Since $|\binom xp|_p\le1$ for every
$x\in\mathbb Z_p$, this functional cannot extend with norm at
most one.  Thus the ordinary theta measure input is materially
stronger than bounded individual derivatives.
\item The coefficient bound $|c_j|\le1$ in
\eqref{cv-moment-approximation} is essential: a single integral
evaluation multiplied by $p^{-1}$ has no such integrality property.
\item For $\lambda=(2,0)$, the polynomial Weyl lattice on one
column is generated by $x^2,2xy,y^2$, whereas the dual-Weyl lattice
contains $xy$.  They differ at $p=2$.  The representation-theoretic
prime bound therefore cannot be discarded.
\item A finite class sum with one class and stabilizer of order $p$
can equal $1/p$, even when its coefficient pairing is integral.
The unit-stabilizer assertion is indispensable for the loss-free
class-sum passage at the selected prime.
\item Reversing \eqref{cv-commutation} changes the column exponent.
For one-dimensional spaces over $\mathbb Q_p$, take $K_i=1$ and
$U_i=p^{s_0ri}$.  Then \eqref{cv-commutation} holds exactly, while
the reversed factor $p^{-s_0r}$ does not if $s_0r>0$.
\end{enumerate}

The finite approximation, substitution, lattice, determinant,
stabilizer, and accumulation arguments in this section are proved
in full. For the submitted theta construction, the normalized ordinary
theta moment theorem and its exact finite translation/support identities
are proved in Sections~\ref{gm-section} and \ref{csp-section}.
Section~\ref{hec-section} proves \eqref{cv-commutation};
Section~\ref{tvc-section} identifies the actual kernel with (CV2)
and assembles the selected-prime column bound.
Weyl/dual-Weyl equality at the target primes satisfying (CV11) has
been proved by the hook-norm argument of
Section~\ref{fh-section}. These inputs are proved independently
of the final branching map.


\section{The actual theta kernel and the selected-prime column estimate}
\label{tvc-section}

The moment and incidence hypotheses in Section~\ref{cv-section}
have now been proved for the normalized theta functions of the
construction. We make the remaining identification with the
polynomial kernel explicit. In particular, an arbitrary projection
onto the target representation has no effect, and ordinary mixed
derivatives introduce no factorial denominator. We then state exactly
which column estimate follows. This concerns the selected primes;
the estimates at the auxiliary level primes are separate.
\cite{alpo-uniform-source}

\subsection{From mixed derivatives to the polynomial map}

Retain the ranks $a<b$, embedding set $I$, dominant weights
$\lambda,\mu$, and nonnegative determinant exponents $n_\tau$
of \eqref{cv-weight-compatibility}. For the source variables
$X_{\tau,\alpha s}$, target variables $Y_{\tau,\beta j}$,
and new variables $Z_{\tau,\alpha\beta}$, put
\[
 \begin{split}
 F^\sharp(X)&=F(X)\prod_{\tau\in I}\det(X_\tau)^{n_\tau},\\
 \mathcal Q_F(Y,Z)
   &=F^\sharp\bigl((Z_\tau Y_{\tau,[1,a]})_{\tau\in I}\bigr)
     =\sum_m Z^m Q_m(F)(Y).
 \end{split}
                                                               \tag{TC1}
\]
Here $Y_{\tau,[1,a]}$ denotes the first $a$ columns of $Y_\tau$.
This is the universal polynomial of
\eqref{cv-universal-expansion}. The order of its two new matrix
indices is significant: $Z$ is an $a$ by $b$ matrix, whereas
the analytic theta variable is a $b$ by $a$ matrix. Thus
$Z_{\tau,\alpha\beta}$ corresponds to differentiation in
the analytic coordinate $(\tau,\beta,\alpha)$.

Let $G$ be holomorphic near zero in those analytic coordinates,
and let $c_\tau\ne0$ be the fixed direction constants. Define
\[
 J_m(G)=
  \left(\prod_{\tau,\alpha,\beta}
                  c_\tau^{m_{\tau,\alpha\beta}}\right)
  \left(\prod_{\tau,\alpha,\beta}
   \partial_{z_{\tau,\beta,\alpha}}^{m_{\tau,\alpha\beta}}
          G\right)(0).
                                                               \tag{TC2}
\]
Holomorphic mixed derivatives commute. Their multilinear definition
also shows that the result is independent of the order in which the
repeated directions are listed. In particular $J_m(G)$ is the
derivative, and is $m!$ times the corresponding Taylor coefficient
when all $c_\tau=1$.

\begin{lemma}\label{tvc-jet-map}
For any scalar family $(J_m)$, the operation on polynomials that
first multiplies by $\prod_\tau\det(X_\tau)^{n_\tau}$,
then substitutes
\[
              X_{\tau,\alpha s}
                   \longmapsto
                  \sum_\beta Y_{\tau,\beta s}Z_{\tau,\alpha\beta},
\]
and finally sends each monomial $Y^vZ^m$ to $J_mY^v$, is
precisely
\[
                       F\longmapsto\sum_m J_mQ_m(F).
                                                               \tag{TC3}
\]
Its image lies in $V_\mu$ for $F\in V_\lambda$. Consequently
composing it with any linear projection of the full target polynomial
ring onto $V_\mu$ which is the identity on $V_\mu$ leaves it unchanged.
\end{lemma}
\begin{proof}
After the first two operations the polynomial is exactly (TC1).
Applying the stated monomial rule gives (TC3) by linearity.
Lemma~\ref{cv-substitution} proves that $\mathcal Q_F(Y,Z)$,
as a polynomial in $Y$, satisfies the defining column degrees and
column-addition invariances of $V_\mu$. Comparing coefficients of
the independent variables $Z$ shows that each $Q_m(F)$ satisfies
these same identities. Hence every finite linear combination in
(TC3) belongs to $V_\mu$. A projection equal to the identity there
therefore fixes it. This argument makes no choice of a complement
and uses no integrality of a projection.
\end{proof}

For the actual lattices $N,N'$ and finite selected set $S$, put
\[
 H=\{x\in\operatorname{Hom}_E(E^a,E^b):xN_P\subseteq N'_P
                                                \ (P\in S)\}.
\]
For positive integers $c_A,c_B$, choose representatives $x$ of
$c_A^{-1}H/c_BH$. This quotient is finite by
Lemma~\ref{jcf-finite-indicators}. Let $G_x$ be its normalized
theta indicator, including division by the determinant-character
factor as in Section~\ref{jcf-section}. Let $\Xi(x)$ be the actual
finite test function, and let $\eta(N,N')$ be the ratio of the
remaining local frame and determinant characters. Set
\[
                 J_m=\sum_{x\in c_A^{-1}H/c_BH}
                                  \Xi(x)J_m(G_x).
                                                               \tag{TC4}
\]
All the quantities in this formula belong to the one field $K$
of Proposition~\ref{jcf-common-field}, after the fixed finite
choices have been made. Only the finitely many $m$ of
\eqref{cv-finite-degrees} are used.

\begin{proposition}\label{tvc-actual-kernel}
The actual polynomial theta kernel for this pair of decorated
lattices is
\[
                  K(N,N')F=\eta(N,N')\sum_m J_mQ_m(F).
                                                               \tag{TC5}
\]
Here $J_m$ is exactly (TC4), with the derivatives (TC2).
There is no additional factorial, lattice index, or projection
factor in (TC5).
\end{proposition}
\begin{proof}
For a fixed representative $x$, the polynomial theta operation is
the determinant multiplication and substitution of Lemma~\ref{tvc-jet-map},
followed by the derivative functional $Z^m\mapsto J_m(G_x)$.
That lemma identifies it with $\sum_mJ_m(G_x)Q_m$ and removes
any projection onto the target representation. Multiply by
$\Xi(x)$ and sum. Both sums are finite, so interchanging them
gives (TC4)--(TC5). The factor $\eta(N,N')$ is independent of $x$
and can be taken outside the sum.

The finite indicator $G_x$ already contains the Fourier average
and the positive square-root normalization proved in
Section~\ref{csp-section}. The present operation neither inserts
nor removes them. Likewise (TC2) is the ordinary repeated
derivative used there, not a divided Taylor coefficient.
This proves the normalization assertion. The admissibility of
the representatives is also exact: $c_Ax\in H$ implies
$c_Ac_Bx\in H$, so each term lies in the permitted level
$c_Ac_B$ and the indicator is evaluated at level $c_B$.
All equalities descend from the complex embedding to $K$ by
its injectivity.
\end{proof}

\subsection{The resulting local divisibility}

We retain all the geometric and coefficient data of
Section~\ref{cv-section}. Thus $[E:\mathbb Q]=2f$,
$\pi\bar\pi=p$, the $f$ primes dividing $\pi$ have norm $p$,
and the ranks are positive even integers
\[
                 N_{i+1}-N_i=t+1=2s_0\quad(0\le i<m).
\]
The CM elliptic ideal and the inverse half different are units at
$p$, and $p\ge5$. The source and target weights satisfy
\eqref{cv-weight-compatibility} at each transition. Choose
nonnegative integers $W_\tau$ with \eqref{cv-weight-bounds}, and suppose the
target hook bound \textup{(CV11)} holds at every transition and
\[
                              p>2fN_m+1.                       \tag{TC6}
\]
All other local levels in the test function are prime to $p$,
with integral values at the coefficient place in question, as in
the construction of Section~\ref{csp-section}. These conditions
are retained because they are used in the moment and character
unit calculations.

Let $\mathscr S_i$ be the local coefficient lattice of equivariant
class sections at rank $N_i$. Let $\mathcal T$ be the distinguished
set of embeddings, $r=|\mathcal T|$, and let $U_i$ be the product
of the upper-neighbor sums at its corresponding selected primes.
Let $K_i$ be the actual class-sum theta transition and put
$C_m=K_{m-1}\cdots K_0$.

\begin{theorem}\label{tvc-column}
At a coefficient place $w\mid p$ of selected type, put
\[
 e=2\sum_\tau W_\tau+fN_m,\qquad
 b_m=me+s_0r\frac{m(m+1)}2.
                                                               \tag{TC7}
\]
Then
\[
                  p^{b_m}C_mU_0^m\mathscr S_0
                                          \subseteq\mathscr S_m.
                                                               \tag{TC8}
\]
If the initial lattice is $\mathscr N=p^AU_0^m\mathscr S_0$
with $A\ge0$, and $v_w(c)\ge v_w(p)b_m$, then
\[
                              cC_m\mathscr N\subseteq\mathscr S_m.
                                                               \tag{TC9}
\]
At a coefficient place of conjugate type the column preserves
the original integral class lattice without a power of $p$.
Applying any of the integral lowering maps of
Section~\ref{atl-section}, and then pairing with an integral row,
introduces no further denominator at $w$.
\end{theorem}
\begin{proof}
Proposition~\ref{tvc-actual-kernel} identifies the actual operator
with the operator in \textup{(CV2)}. The moment laws of
Section~\ref{gm-section} and the two support calculations of
Section~\ref{csp-section} supply exactly \textup{(CV9)--(CV10)}:
the selected support has scalar $p^{-fN_i}$ and the conjugate
support has scalar one. Lemma~\ref{csp-coordinate-units} gives
the unit direction factors. The approximations use one integral
finite sum for all the required moments. Thus they are the
simultaneous approximations in
\eqref{cv-moment-approximation}, not merely separate bounds on
individual derivatives.

The character ratio is a unit by
Lemma~\ref{cv-character-unit}, and the hook bound identifies the
target Weyl and dual-Weyl lattices. Condition (TC6) makes every
class stabilizer order a unit by
Lemma~\ref{cv-stabilizer-unit}. Theorem~\ref{cv-two-kernels}
and the residue-removal argument therefore give
\[
                      p^eU_{i+1}K_i\mathscr S_i
                                                \subseteq\mathscr S_{i+1}.
\]
Theorem~\ref{hec-global-identity}, at each of the $r$ distinct
primes, gives
\[
                            U_{i+1}K_i=p^{s_0r}K_iU_i.
\]
The different local neighbor operators commute since they change
different local components. All these identities hold over the
same field $K$ by Proposition~\ref{jcf-common-field}.
Proposition~\ref{cv-column} now proves (TC8)--(TC9), with the
integer exponent (TC7). At conjugate-type places the second
bound of Theorem~\ref{cv-two-kernels} already preserves each
original integral lattice, so their composition does too.

The final lowering maps preserve the local lattices by
Section~\ref{atl-section}. The row pairing of
Section~\ref{cp-section} has only the stated class stabilizer
denominators, which are units by (TC6). This proves the last
assertion.
\end{proof}

The exponent (TC7) is the exponent $mE_p$ of
\eqref{cv-final-exponent}. The selected-prime part of the column
valuation argument is therefore complete under the displayed
construction hypotheses. This does not prove that
$U_0^m\mathscr S_0$ or its theta image is nonzero, does not supply
the estimates at the two auxiliary level primes, and does not
replace the archimedean norm and period comparison needed for
the final arithmetic inequality. Those are separate mathematical
assertions.


% Scalar rigidity for the actual Hermitian flags.
% Labels in this fragment use the prefix frg-.
\section{Scalar rigidity for Hermitian flags}
\label{frg-section}

We prove the three assertions which the source combines to
identify the scalar coefficients in its identity-cell calculation.
The conclusion concerns these coefficients and the negative norm;
it does not identify two ordered frames whose vectors can still
differ by lower flag terms. All conventions below agree with the
actual flag and infinitesimal-character definitions.
\cite{alpo-uniform-source}

\subsection{Flags and their infinitesimal characters}

Let $E$ be a number field with an involution $a\mapsto\bar a$,
and let $K\subseteq E$ be a subfield. Let $V$ be a
finite-dimensional $E$-space with a nondegenerate Hermitian form
$h$, conjugate linear in its first argument and linear in its
second. Thus
\[
 h(ax,by)=\bar a b\,h(x,y),\qquad
 h(y,x)=\overline{h(x,y)}.
\]
For $T\in E\setminus\{0\}$, a flag frame of length $r+1$
is a tuple
\[
 F=(e_0,\ldots,e_{r-1},z)
\]
with the $e_i$ linearly independent and
\[
 h(e_i,e_j)=0,\qquad h(e_i,z)=0,\qquad h(z,z)=-T.
\]
Write $f_i=e_i$ for $i<r$, $f_r=z$, and
\begin{gather*}
 X_a=\operatorname{span}_E(e_0,\ldots,e_{a-1})\quad(0\le a\le r),
 \\
 X=X_r,\qquad
 X^\perp=\{y:h(x,y)=0\text{ for every }x\in X\}.
\end{gather*}
For $y\in X$, let $[y]_i$ be its coefficient on $e_i$.
The set of $E$-linear endomorphisms used here is
\begin{equation}\label{frg-lie-definition}
 \mathfrak n_F=\left\{Y:
 \begin{array}{l}
 h(Yx,y)+h(x,Yy)=0\quad(x,y\in V),\\
 YX_a\subseteq X_{a-1}\quad(1\le a\le r),\\
 YX^\perp\subseteq X
 \end{array}\right\}.
\end{equation}
For $Y\in\mathfrak n_F$ define
\begin{equation}\label{frg-character}
 \ell_F(Y)=\sum_{i=0}^{r-1}[Yf_{i+1}]_i,
 \qquad d\psi_F(Y)=\operatorname{Tr}_{E/K}\ell_F(Y).
\end{equation}
Every vector whose coefficient occurs in this formula belongs
to $X$: for $i+1<r$ this follows from
\eqref{frg-lie-definition}, and for $i=r-1$ it follows from
$z\in X^\perp$. When $r=0$, the sum is zero.

Let $F'$ be another such frame, with negative norm parameter
$T'\ne0$. Given scalars $c_0,\ldots,c_r\in E$, the same-flag
condition is exactly
\begin{equation}\label{frg-same-flag}
                         f'_i-c_i f_i\in X_i
                                     \qquad(0\le i\le r).
\end{equation}
The relevance condition needed below is
\begin{equation}\label{frg-relevance}
 d\psi_F(Y)=d\psi_{F'}(Y)
                  \qquad(Y\in\mathfrak n_F\cap\mathfrak n_{F'}).
\end{equation}
It is equality of the actual linear functionals on the
intersection, not just equality on one chosen endomorphism.

The trace pairing has the nondegeneracy needed here without
an additional hypothesis. Indeed $E/K$ is a finite extension
of characteristic zero. If $b\ne0$, then
\[
 \operatorname{Tr}_{E/K}(b\,b^{-1})=[E:K]\ne0.
\]
Consequently
\begin{equation}\label{frg-trace-nondegenerate}
 \bigl(\operatorname{Tr}_{E/K}(ba)=0\text{ for all }a\in E\bigr)
                                      \quad\Longrightarrow\quad b=0.
\end{equation}
This proves the source's explicit trace-pairing assumption in
the present number-field setting.

\subsection{Relevance forces the flag scalars to agree}

\begin{lemma}\label{frg-triangular}
Under \eqref{frg-same-flag}, every $c_i$ is nonzero and
$X'_a=X_a$ for $0\le a\le r$.
\end{lemma}
\begin{proof}
The vectors $f_0,\ldots,f_r$ are independent. For if
$\sum_{i<r}a_i e_i+a_rz=0$, pairing with $z$ in the first
argument gives $-a_rT=0$. Then $a_r=0$, and independence
of the $e_i$ gives all other coefficients zero. The same
argument applies to $F'$.
Equation \eqref{frg-same-flag} writes the independent tuple
$F'$ as a triangular change of the independent tuple $F$,
with diagonal entries $c_i$. Its matrix is invertible,
so every $c_i\ne0$. Each initial triangular block is likewise
invertible, giving $X'_a=X_a$.
\end{proof}

\begin{proposition}\label{frg-relevant-scalars}
If \eqref{frg-same-flag} and \eqref{frg-relevance} hold, then
\begin{equation}\label{frg-all-scalars}
                              c_i=c_r\qquad(0\le i\le r).
\end{equation}
\end{proposition}
\begin{proof}
There is nothing to prove for $r=0$. Fix $0\le i<r$ and
$a\in E$. Nondegeneracy of $h$ and independence of the $f_j$
give a vector $v\in V$ such that
\begin{equation}\label{frg-dual-vector}
                     h(v,f_j)=a\,\delta_{j,i+1}
                                        \qquad(0\le j\le r).
\end{equation}
For completeness, extend the $f_j$ to a basis $b_1,\ldots,b_d$.
The Gram matrix $G=(h(b_j,b_k))$ is invertible, since a vector
orthogonal to the whole basis would contradict nondegeneracy.
Prescribe the desired values $h(v,b_k)$, using zero on the
additional basis vectors. If $v=\sum_j t_jb_j$, these equations
are $\bar t^{\,t}G=(h(v,b_k))_k$. Invert $G$ and apply the
involution to solve for every $t_j$. This proves
\eqref{frg-dual-vector}.

Define the $E$-linear endomorphism
\begin{equation}\label{frg-rank-two}
                       Y(y)=h(v,y)e_i-h(e_i,y)v.
\end{equation}
It is skew Hermitian. Indeed the two terms of $h(Yx,y)$ are
$h(x,v)h(e_i,y)-h(x,e_i)h(v,y)$, which cancel the two terms
of $h(x,Yy)$.
Also $h(v,X_{i+1})=0$ by \eqref{frg-dual-vector}.
For $x\in X_b$ we have $h(e_i,x)=0$. If $b\le i+1$, then
$Yx=0$; if $b\ge i+2$, then $Yx$ is a multiple of
$e_i\in X_{b-1}$. Finally, for $x\in X^\perp$ the second
term in \eqref{frg-rank-two} vanishes, and $Yx\in X$.
Thus $Y\in\mathfrak n_F$.

By Lemma~\ref{frg-triangular}, all the spaces $X'_b$ equal
$X_b$. The same argument, using $e_i\in X'_{i+1}$ and
$h(v,X'_{i+1})=0$, proves $Y\in\mathfrak n_{F'}$.
As $h(e_i,f_j)=0$ for every $j$, equation
\eqref{frg-dual-vector} gives immediately
\begin{equation}\label{frg-ell-unprimed}
                                  \ell_F(Y)=a.
\end{equation}

We compute the primed functional, retaining all triangular
terms. Since $f'_j\in X+Ez$ and $h(e_i,X+Ez)=0$,
\[
                            Yf'_{j+1}=h(v,f'_{j+1})e_i.
\]
If $j<i$, then $f'_{j+1}\in X_{i+1}$ and the scalar on
the right is zero. If $j>i$, then $e_i\in X'_{i+1}$
has zero coefficient on $e'_j$. For $j=i$, triangularity
gives
\[
 [e_i]'_i=c_i^{-1},\qquad
 h(v,f'_{i+1})=c_{i+1}a,
\]
because the error in $f'_{i+1}-c_{i+1}f_{i+1}$ lies in
$X_{i+1}$, which is orthogonal to $v$ in the required order.
Only this term survives. Hence
\begin{equation}\label{frg-ell-primed}
                         \ell_{F'}(Y)=\frac{c_{i+1}}{c_i}a.
\end{equation}
Relevance and \eqref{frg-ell-unprimed}--\eqref{frg-ell-primed}
give
\[
 \operatorname{Tr}_{E/K}\left(\left(\frac{c_{i+1}}{c_i}-1\right)a\right)=0
                                      \qquad(a\in E).
\]
Equation \eqref{frg-trace-nondegenerate} implies
$c_{i+1}=c_i$. This holds for each $i<r$, proving
\eqref{frg-all-scalars}.
\end{proof}

\subsection{An integral unit and the negative norm}

Assume now that the involution is complex conjugation at every
embedding:
\begin{equation}\label{frg-all-embeddings}
                 \sigma(\bar b)=\overline{\sigma(b)}
                    \quad(b\in E,\ \sigma:E\hookrightarrow\mathbb C).
\end{equation}
This includes the actual CM-field involution. Let $T,T'$ be
positive integers and suppose the final scalar is the image
of a unit $\zeta\in\mathcal O_E^\times$.

\begin{proposition}\label{frg-unit-norm-root}
Under \eqref{frg-same-flag} and \eqref{frg-all-embeddings},
if $c_r=\zeta\in\mathcal O_E^\times$, then $T=T'$ and
$\zeta$ is a root of unity. Relevance is not needed for this
assertion.
\end{proposition}
\begin{proof}
The final equation in \eqref{frg-same-flag} says
$z'=\zeta z+x$ with $x\in X$. The subspace $X$ is
totally isotropic and orthogonal to $z$, so
\[
                 -T'=h(z',z')=-\bar\zeta\zeta T,
 \qquad T'=\bar\zeta\zeta T.
\]
Let $d=[E:\mathbb Q]$. The norm of an algebraic unit is
$\pm1$: multiplication by the unit is an automorphism of
the free abelian group $\mathcal O_E$, so its determinant
is $\pm1$. The involution is a $\mathbb Q$-automorphism,
and conjugates the multiplication operators of $\zeta$
and $\bar\zeta$. Thus their norms agree. Taking norms in
the last display gives
\[
                 (T')^d=N_{E/\mathbb Q}(\zeta)^2T^d=T^d.
\]
Positivity gives $T'=T$, and cancellation gives
$\bar\zeta\zeta=1$. Equation \eqref{frg-all-embeddings}
now implies $|\sigma(\zeta)|=1$ for every embedding.

Here is the elementary finite-polynomial argument completing
the root-of-unity assertion. For every integer $m\ge1$,
the characteristic polynomial of multiplication by $\zeta^m$
on $\mathcal O_E$ is
\[
                P_m(t)=\prod_{\sigma:E\hookrightarrow\mathbb C}
                                  (t-\sigma(\zeta)^m)\in\mathbb Z[t].
\]
Its coefficient of $t^{d-j}$ has absolute value at most
$\binom dj$, since it is a sum of that many products of
complex numbers of absolute value one. Only finitely many
monic integer polynomials of degree $d$ have these coefficient
bounds. Their union of root sets is finite. In a fixed
embedding the powers $\sigma(\zeta)^m$ therefore take only
finitely many values. Two positive powers are equal;
injectivity and $\zeta\ne0$ give $\zeta^n=1$ for some
$n>0$.
\end{proof}

\subsection{Reduction at an odd unramified prime}

\begin{lemma}\label{frg-root-congruence}
Let $P$ be a nonzero prime ideal of $\mathcal O_E$, and
let $p\ge3$ be a rational prime such that
\begin{equation}\label{frg-prime-condition}
                         p\in P\quad\text{and}\quad p\notin P^2.
\end{equation}
If $\zeta\in\mathcal O_E$ is a root of unity and
$\zeta-1\in P$, then $\zeta=1$.
\end{lemma}
\begin{proof}
Suppose $\zeta\ne1$, and let its order be $m>1$.
Choose a prime $q\mid m$ and put $\eta=\zeta^{m/q}$.
Then $\eta$ has order exactly $q$, and $\eta-1\in P$.
The roots $1,\eta,\ldots,\eta^{q-1}$ are distinct, so
factoring $t^q-1$ over $E$, dividing by $t-1$, and
evaluating at $t=1$ gives
\begin{equation}\label{frg-prime-product}
                     q=\prod_{j=1}^{q-1}(1-\eta^j)\in P^{q-1}.
\end{equation}
Indeed each factor lies in $P$, because $\eta\equiv1\pmod P$.
In particular $q\in P$. If $p\ne q$, integers $a,b$
with $ap+bq=1$ would imply $1\in P$, a contradiction.
Thus $q=p$. Since $p-1\ge2$, equation
\eqref{frg-prime-product} implies $p\in P^2$, contrary
to \eqref{frg-prime-condition}.
\end{proof}

\begin{theorem}[The actual flag-scalar conclusion]
\label{frg-identity-cell}
Let $E,V,h,K,F,F',c$ satisfy the definitions above, with
$T,T'$ positive integers. Suppose
\eqref{frg-same-flag}, \eqref{frg-relevance} and
\eqref{frg-all-embeddings} hold. Suppose further that
$c_r=\zeta\in\mathcal O_E^\times$ and that, for a rational
prime $p\ge3$ and a prime ideal $P$,
\[
                         p\in P\setminus P^2,
                         \qquad\zeta\equiv1\pmod P.
\]
Then
\begin{equation}\label{frg-identity-conclusion}
                         T=T',\qquad c_0=\cdots=c_r=1.
\end{equation}
\end{theorem}
\begin{proof}
Proposition~\ref{frg-unit-norm-root} gives $T=T'$ and
makes $\zeta$ a root of unity. Lemma~\ref{frg-root-congruence}
gives $\zeta=1$. Finally Proposition~\ref{frg-relevant-scalars}
gives $c_i=c_r=\zeta$ for every $i$.
\end{proof}

The restrictions in the theorem cannot simply be discarded.
At the prime $2$, the unit $-1\in\mathbb Z$ satisfies
$-1\equiv1\pmod2$ and $2\notin(2)^2$, but is not $1$.
For ramified odd primes, take a primitive cube root $\omega$
in $E=\mathbb Q(\omega)$ and $P=(1-\omega)$.
Then $\omega\equiv1\pmod P$ and
$3=-\omega^{-1}(1-\omega)^2\in P^2$, so the claimed
root conclusion would again fail without the exclusion of
$P^2$. The unit condition in Proposition~\ref{frg-unit-norm-root}
is also necessary: for $V=\mathbb Q(i)$ with
$h(x,y)=-\bar x y$, take $r=0$, $z=1$, $z'=2$.
The same-flag relation has scalar $2$, but $T=1$ and $T'=4$.

Even the full theorem does not assert $F=F'$ as ordered
tuples. When $r\ge1$, replacing $z$ by $z+e_0$ preserves
its norm, all flag conditions and all scalars $c_i=1$.
It also preserves the infinitesimal character, since
$Ye_0=0$ for $Y\in\mathfrak n_F$, and the spaces
\eqref{frg-lie-definition} depend only on the $e_i$.
The source's precise conclusion is \eqref{frg-identity-conclusion},
which is what has been proved. No counterexample satisfying
those actual hypotheses was found.


% Ordinary mathematics for the vanishing of degenerate off-diagonal terms.
% Labels in this fragment have the prefix df-.
\section{Degenerate flag orbits and compact-group orthogonality}
\label{df-section}

The vanishing of a degenerate term between two different compact
types has a useful direct proof.  We give the flag calculation first,
including the case in which the two isotropic subspaces meet.  We then
identify a common compact unitary group and prove the required
orthogonality.  The argument uses no estimate for a lattice mass and
no asymptotic estimate for an orbital integral.  As usual, an orbital
integral below is considered only when it is absolutely convergent.
Its absolute convergence is a separate analytic assertion about the
functions used to construct the theta columns.

\subsection{Flags, characters, and the determinant}

Let $E/F$ be a CM extension, let a bar denote its nontrivial
$F$-automorphism, and let $V$ be an $E$-vector space of dimension
$k+2r+1$.  Suppose that $h$ is a nondegenerate Hermitian form on $V$,
conjugate linear in its first argument, of signature
$(k+r,r+1)$ at every real place of $F$.  The integers $k,r$ are
nonnegative.  An ordered flag with a negative vector consists of
linearly independent vectors $e_1,\ldots,e_r$ and a vector $z$ such
that
\[
 h(e_a,e_b)=0,\qquad h(e_a,z)=0,\qquad h(z,z)=-T,
 \qquad T\in F\text{ totally positive}.
\]
Write
\[
 X_a=\langle e_1,\ldots,e_a\rangle_E,\quad
 X=X_r,\quad X^+=X+Ez,\quad e_{r+1}=z,
 \qquad X_0=0.
\]
The Hermitian space
\begin{equation}\label{df-positive-quotient}
 Q=(X^+)^\perp/X
\end{equation}
is positive definite of dimension $k$ at every real place.  Indeed,
removing $r$ hyperbolic planes gives signature $(k,1)$ on
$X^\perp/X$, and removing the negative line generated by $z$ gives
\eqref{df-positive-quotient}.  In particular,
\begin{equation}\label{df-isotropic-in-perp}
 y\in(X^+)^\perp,\quad h(y,y)=0\quad\Longrightarrow\quad y\in X.
\end{equation}
All assertions about positivity may be checked at any one complex
embedding above a real place when only an algebraic equality is at issue.

Let $N$ be the subgroup of $U(V,h)$ defined by
\[
 (n-1)X_a\subseteq X_{a-1}\ (1\leq a\leq r),
 \qquad (n-1)X^\perp\subseteq X.
\]
It is a connected unipotent algebraic group over $F$.  In a Witt
basis extending the $e_a$, its matrices are triangular with all
diagonal entries equal to one.  Its Lie algebra consists of the
$E$-linear endomorphisms $Y$ satisfying
\begin{equation}\label{df-lie}
 h(Yv,w)+h(v,Yw)=0,\quad
 YX_a\subseteq X_{a-1},\quad YX^\perp\subseteq X.
\end{equation}
If $Y$ satisfies these conditions, let $c_a(Y)$ be the coefficient
of $e_a$ in $Ye_{a+1}$, and put
\begin{equation}\label{df-differential}
 \ell(Y)=\sum_{a=1}^r c_a(Y),\qquad
 d\psi(Y)=\operatorname{Tr}_{E/F}\ell(Y).
\end{equation}
These coefficients are defined because $Ye_{a+1}\in X_a$ for
$a<r$ and $Yz\in X$.  The function $d\psi$ vanishes on Lie
brackets.  For the coefficients with $a<r$, this follows from
$[Y_1,Y_2]X_{a+1}\subseteq X_{a-1}$.  For the last coefficient,
$[Y_1,Y_2]z\in X_{r-1}$.  Thus a fixed additive character of
$F\backslash\mathbb A_F$, composed with $d\psi\circ\log$, defines
a character of $N(\mathbb A_F)$.  Here exponential and logarithm
are finite polynomial expressions in a triangular matrix; the
commutator terms in their multiplication formula disappear under
$d\psi$.  When $r=0$, the group and the character are trivial.

Use primes for a second flag $(e'_1,\ldots,e'_r,z',T')$.  The flags
are called \emph{relevant} when
\begin{equation}\label{df-relevance}
 d\psi(Y)=d\psi'(Y)
 \qquad(Y\in\operatorname{Lie}N\cap\operatorname{Lie}N').
\end{equation}
This is equality of $F$-linear functionals.  It remains an equality
after extension to any completion of $F$.
Put
\[
 I=X\cap X',\qquad V_0=X^++X'^+,
 \quad
 D_a=\det\bigl(h(e_i,e'_j)\bigr)_{1\leq i,j\leq a},
 \qquad D_0=1,
\]
and define
\begin{equation}\label{df-G}
 G=D_{r+1}\overline{D_{r+1}}-TT'D_r\overline{D_r}.
\end{equation}
Conjugating every $D_a$ leaves $G$ unchanged.  Consequently the
opposite convention for which variable of a Hermitian form is linear
does not change the condition $G=0$.

\subsection{The consequence of relevance on a common isotropic vector}

For $0\ne x\in X$, its level is the unique integer $a$ such that
$x\in X_a\setminus X_{a-1}$; write
$x=\alpha e_a+x_0$ with $0\ne\alpha\in E$ and $x_0\in X_{a-1}$.

\begin{lemma}\label{df-successor}
Suppose that the flags are relevant.  If $0\ne x\in I$ has levels
$a,a'$ and leading coefficients $\alpha,\alpha'$ in the two flags,
then
\begin{equation}\label{df-successor-eq}
 \alpha e_{a+1}-\alpha'e'_{a'+1}\in X_a+X'_{a'}.
\end{equation}
\end{lemma}

\begin{proof}
For $v\in(X_a+X'_{a'})^\perp$, consider
\[
 A_{x,v}(y)=h(v,y)x-h(x,y)v.
\]
This endomorphism is skew Hermitian.  If $i\leq a$, it kills $e_i$.
If $i>a$, it sends $e_i$ into $Ex\subseteq X_a\subseteq X_{i-1}$.
It sends $X^\perp$ into $Ex$.  The same statements hold for the
primed flag.  Hence $A_{x,v}$ belongs to both Lie algebras in
\eqref{df-relevance}.  Directly from the coefficients in
\eqref{df-differential},
\[
 \ell(A_{x,v})=\alpha h(v,e_{a+1}),\qquad
 \ell'(A_{x,v})=\alpha'h(v,e'_{a'+1}).
\]
Set $d=\alpha h(v,e_{a+1})-\alpha'h(v,e'_{a'+1})$.
Replacing $v$ by $sv$ shows that
$\operatorname{Tr}_{E/F}(\bar s d)=0$ for every $s\in E$.
Taking $s=d$ gives $2\bar d d=0$, so $d=0$.
Thus $\alpha e_{a+1}-\alpha'e'_{a'+1}$ is orthogonal to
$(X_a+X'_{a'})^\perp$.  Nondegeneracy of $h$ identifies this double
orthogonal complement with $X_a+X'_{a'}$.
\end{proof}

\begin{lemma}\label{df-negative-not-sum}
Neither $z$ nor $z'$ belongs to $X+X'$.  There is consequently at
most one $c\in E$ with $z'-cz\in X+X'$.
\end{lemma}

\begin{proof}
If $z=x+x'$ with $x\in X$ and $x'\in X'$, then
$h(x,x')=h(x,z)=0$.  Both summands are isotropic, so $h(z,z)=0$,
contrary to $T\ne0$.  Interchanging the flags proves the assertion
for $z'$.  Subtracting two expressions $z'-cz\in X+X'$ proves
uniqueness of $c$.
\end{proof}

\begin{lemma}\label{df-common-levels}
If $I\ne0$ and the flags are relevant, there are a unique integer
$0\leq j<r$ and a unique $c\in E^\times$ such that
\begin{align}
 &\dim_E I=r-j,\qquad X=X_j\oplus I,
       \qquad X'=X'_j\oplus I,\label{df-common-splits}\\
 &z'-cz\in X+X',\label{df-z-relation}
\end{align}
and every nonzero vector of $I$ has the same level in both flags,
with leading coefficients related by $\alpha=c\alpha'$.
The levels that occur in $I$ are exactly $j+1,\ldots,r$.
\end{lemma}

\begin{proof}
Start with a common vector of levels $a,a'$.  As long as $a,a'<r$,
write \eqref{df-successor-eq} as
\[
 \alpha e_{a+1}-\alpha'e'_{a'+1}=s+s',
 \qquad s\in X_a,\quad s'\in X'_{a'}.
\]
Then $y=\alpha e_{a+1}-s=\alpha'e'_{a'+1}+s'$ is a common
vector of levels $a+1,a'+1$ with the same two leading coefficients.
If $a\ne a'$, this procedure reaches a stage where just one level
is $r$.  The next use of \eqref{df-successor-eq} puts either $z$ or
$z'$ in $X+X'$, contradicting Lemma~\ref{df-negative-not-sum}.
Thus $a=a'$.  Continuing to level $r$ gives
$z'-(\alpha/\alpha')z\in X+X'$.  Lemma~\ref{df-negative-not-sum}
makes this ratio independent of the original common vector.

Let $j+1$ be the least level occurring in $I$.  The construction
just given produces a vector of $I$ at every larger level.  No
nonzero vector of $I$ lies in $X_j$ or $X'_j$.  Conversely, choose
one vector at each level $j+1,\ldots,r$.  Successively canceling
the leading coefficient of a vector of $X$ shows that these vectors
together with $e_1,\ldots,e_j$ span $X$.  The same reasoning applies
to $X'$.  The direct sums in \eqref{df-common-splits} follow, and
their dimensions give $\dim_E I=r-j$.  This also proves the
asserted uniqueness of $j$.
\end{proof}

\subsection{The complete normal form needed for orthogonality}

\begin{theorem}\label{df-normal-form}
Suppose that the flags are relevant and $G=0$.  Then there are
$0\leq j\leq r$ and $c\in E^\times$ satisfying
\eqref{df-common-splits} and \eqref{df-z-relation}, with
\begin{equation}\label{df-normal-form-values}
 D_j\ne0,\qquad c\bar c T=T'.
\end{equation}
The radical of the restriction of $h$ to $V_0$ is $I$,
\[
 \dim_E V_0=r+j+1,
 \qquad
 \operatorname{sig}(V_0/I)=(j,j+1).
\]
Moreover,
\begin{equation}\label{df-common-perp}
 V_0^\perp\cap X=I=V_0^\perp\cap X',\quad
 V_0^\perp+X=(X^+)^\perp,\quad
 V_0^\perp+X'=(X'^+)^\perp.
\end{equation}
If $I\ne0$, the common-level and leading-coefficient assertions of
Lemma~\ref{df-common-levels} hold.  If $I=0$, one has $j=r$.
\end{theorem}

\begin{proof}
First suppose $I\ne0$, and take $j,c$ from
Lemma~\ref{df-common-levels}.  The pairing between $X_j$ and $X'_j$
is nondegenerate.  Indeed, a vector $x\in X_j$ orthogonal to $X'_j$
is orthogonal to all of $X'$, since $I\subseteq X$.  The relation
$z'-cz\in X+X'$ also makes it orthogonal to $z'$.  It is isotropic,
so \eqref{df-isotropic-in-perp}, applied to the primed flag, gives
$x\in X'$.  Thus $x\in X_j\cap I=0$.  This proves $D_j\ne0$,
including the convention $D_0=1$.

It remains to obtain the same conclusions when $I=0$.
The pairing between $X^+$ and $X'^+$ is then nondegenerate.
For if $x+az\in X^+$ is orthogonal to $X'^+$, its norm is
$-|a|^2T$ at every real place.  On $(X'^+)^\perp$ the form is
positive semidefinite with radical $X'$.  Therefore $a=0$ and
$x\in X'$, whence $x=0$.  Thus $D_{r+1}\ne0$.
Equation~\eqref{df-G} and $G=0$ imply $D_r\ne0$.

The space $X+X'$ is now nondegenerate of signature $(r,r)$.
There are unique $x\in X$, $x'\in X'$ such that
\[
 z_0=z-x\in(X+X')^\perp,\qquad
 z'_0=z'-x'\in(X+X')^\perp.
\]
They are obtained by solving the nondegenerate cross-pairing
systems.  Their norms are $-T$ and $-T'$.  Subtracting the
corresponding linear combinations of the first $r$ rows and columns
in the matrix defining $D_{r+1}$ gives
\[
 D_{r+1}=D_r h(z_0,z'_0).
\]
The space $(X+X')^\perp$ has signature $(k,1)$.  Write
$z'_0=cz_0+w$ with $h(z_0,w)=0$.  Its orthogonal complement to the
negative line $Ez_0$ is positive definite.  Consequently
\[
 |h(z_0,z'_0)|^2=TT'+T h(w,w)
\]
at every real place.  The equality $G=0$ forces $w=0$.
The scalar $c=h(z_0,z'_0)/(-T)$ belongs to $E$ and is nonzero.
This gives \eqref{df-z-relation}, with $j=r$.

We can finish both cases together.  Let
$H_0=X_j+X'_j$, which is a nondegenerate hyperbolic space of
signature $(j,j)$, and note that
$X+X'=H_0\oplus I$ with $I$ orthogonal to $H_0$.
Subtract from $z$ a vector of $X_j$ so that the resulting vector
$z_0$ is orthogonal to $H_0$.  This leaves its norm equal to $-T$;
it is also orthogonal to $I$.  Similarly subtract from $z'$ a
vector of $X'_j$ to obtain $z'_0\perp H_0$, still of norm $-T'$.
The relation \eqref{df-z-relation} implies
$z'_0-cz_0\in H_0+I$.  Since this difference is orthogonal to the
nondegenerate space $H_0$, it lies in $I$.
Comparing its norm with $-T'$ proves $c\bar c T=T'$.
We have the orthogonal decomposition, with the last summand radical,
\[
 V_0=H_0\mathbin{\perp}Ez_0\mathbin{\perp}I.
\]
It proves the assertions about $V_0$.  In particular
$V_0^\perp\cap X=I$ and similarly for $X'$.
The inclusion $V_0^\perp+X\subseteq(X^+)^\perp$ is immediate.
Both sides have dimension $k+r$, since
\[
 \dim(V_0^\perp+X)
 =(k+r-j)+r-(r-j)=k+r.
\]
This proves the other two equalities in \eqref{df-common-perp}.
The argument includes $r=0$, in which case it is just the equality
case of the reverse Cauchy--Schwarz inequality in signature $(k,1)$.
\end{proof}

\begin{corollary}\label{df-common-compact}
Under the hypotheses of Theorem~\ref{df-normal-form}, there is a
positive definite $E$-subspace $W\subseteq V_0^\perp$ of dimension
$k$ such that
\[
 V_0^\perp=W\oplus I.
\]
The maps induced by inclusion are isometries
\begin{equation}\label{df-quotient-isometries}
 W\simeq (X^+)^\perp/X,
 \qquad W\simeq (X'^+)^\perp/X'.
\end{equation}
The subgroup $U(W)$, acting as the identity on $W^\perp$, fixes
$V_0$ pointwise.  At infinity it is a compact group
$K_W\simeq\prod_{v\mid\infty}U(k)$, acting on the two quotients
in \eqref{df-quotient-isometries} by the same standard unitary
action, up to conjugation by a unitary change of basis.
\end{corollary}

\begin{proof}
The radical of $V_0^\perp$ is $I$.  In $I^\perp/I$ the form has
signature $(k+j,j+1)$, while the nondegenerate subspace $V_0/I$
has signature $(j,j+1)$.  Hence its orthogonal complement
$V_0^\perp/I$ is positive definite of dimension $k$.
Choose an $E$-linear complement $W$ of $I$ in $V_0^\perp$.
The form on $W$ is positive definite at every real place.
The injectivity and surjectivity of the maps in
\eqref{df-quotient-isometries} follow from
\eqref{df-common-perp}.  They preserve the Hermitian forms because
the subspaces being divided out are radicals.  Finally,
$V=W\perp W^\perp$, and $V_0\subseteq W^\perp$.
\end{proof}

\subsection{Averaging over the unipotent quotient}

We spell out a measure issue in applying the common compact group.
Its elements need not preserve the particular integral lattices
that define an arithmetic subgroup.  The appropriate operation is
first to integrate over the quotient of the \emph{real unipotent}
group by its arithmetic subgroup.

\begin{lemma}\label{df-unipotent-lattice}
Let $R$ be a unipotent algebraic group over a number field $F$ and
let $K_f\subseteq R(\mathbb A_{F,f})$ be a compact open subgroup.  Then
\[
 \Gamma=R(F)\cap K_f
\]
is a discrete cocompact subgroup of $R(F\otimes_{\mathbb Q}\mathbb R)$.
The group $R(F\otimes_{\mathbb Q}\mathbb R)$ is unimodular.
\end{lemma}

\begin{proof}
In characteristic zero, the exponential and logarithm of a
unipotent matrix are finite polynomials, inverse to one another.
The Lie algebra has a central series defined over $F$, with
additive $F$-vector spaces as successive quotients.  Choose an
$F$-basis compatible with a refinement of this series.  Ordered
exponential coordinates give polynomial coordinates on $R$; in
these coordinates multiplication is triangular, with the new
coordinate equal to the sum of the two coordinates plus a
polynomial in earlier coordinates.

Clearing the finitely many denominators in these polynomial
expressions, and shrinking the resulting fractional ideals if
necessary at the finitely many places where $K_f$ is prescribed,
gives coordinate lattices contained in $\Gamma$.  More explicitly,
each one-parameter subgroup $\exp(tY)$ contains
$\exp(\mathfrak a Y)\subseteq\Gamma$ for some nonzero fractional
ideal $\mathfrak a$; this follows by inserting $t\in\mathfrak a$
in the finite polynomial for $\exp(tY)$ and requiring membership
in $K_f$.  Successive multiplication by these elements reduces
each coordinate to a bounded fundamental parallelepiped for the
image of its ideal in $F\otimes\mathbb R$.  The triangular form
ensures that coordinates already reduced are unchanged.  The
resulting bounded closed coordinate box is compact and meets every
left $\Gamma$-coset.  Discreteness follows because $\log K_f$ is
contained in a product of fractional-ideal completions and the
diagonal image of each fractional ideal is discrete at infinity.
Finally, the adjoint action of a unipotent element on the Lie
algebra is unipotent, so its determinant is one.  This proves
unimodularity.
\end{proof}

Let $G_\infty=U(V)(F\otimes\mathbb R)$, and let
$R_\infty=N_\infty\cap N'_\infty$.  This is the real group of a
unipotent $F$-group.  Let $\Gamma\subset R_\infty$ be an arithmetic
subgroup obtained by imposing stabilizers of two integral lattices.
It has the form in Lemma~\ref{df-unipotent-lattice}, after taking
their compact open intersection.  Choose Haar measures $dg,dr$ and
the quotient measure $d\dot g$ characterized by
\[
 \int_{G_\infty}f(g)\,dg
 =\int_{R_\infty\backslash G_\infty}
       \int_{R_\infty}f(rg)\,dr\,d\dot g.
\]
Set $v_R=\operatorname{vol}(\Gamma\backslash R_\infty)$.

\begin{lemma}\label{df-average}
Suppose that a finite-dimensional vector-valued function $B$ on
$G_\infty$ satisfies $B(rg)=\chi(r)B(g)$, where $\chi$ is a
continuous unitary character of $R_\infty$, trivial on $\Gamma$,
and that $B$ is absolutely integrable on
$\Gamma\backslash G_\infty$.
If $\chi\ne1$, its integral is zero.  If $\chi=1$, then
\begin{equation}\label{df-disintegration}
 \int_{\Gamma\backslash G_\infty}B(g)\,dg
 =v_R\int_{R_\infty\backslash G_\infty}B(g)\,d\dot g.
\end{equation}
If a compact subgroup $K\subseteq G_\infty$ normalizes $R_\infty$,
then the quotient measure on $R_\infty\backslash G_\infty$ is
invariant under left multiplication by $K$.
\end{lemma}

\begin{proof}
Apply the displayed integration formula first to scalar absolute
values, and then to each coordinate of $B$, using a measurable
fundamental domain for $\Gamma$ in $R_\infty$.  This gives the
inner integral over $\Gamma\backslash R_\infty$.
If $\chi\ne1$, right translation by an element on which $\chi$
is nontrivial multiplies its character integral by a scalar other
than one.  The integral is therefore zero.  If $\chi=1$, the
inner integral is $v_R B(g)$, proving
\eqref{df-disintegration}.  Absolute convergence justifies these
uses of Fubini.

Conjugation by $k\in K$ multiplies $dr$ by a positive real scalar.
That scalar is a continuous homomorphism from the compact group
$K$ to $\mathbb R_{>0}$, hence is one.  Left multiplication by
$k$ preserves $dg$ as well.  Uniqueness of the quotient measure
in the integration formula proves the last assertion.
\end{proof}

\subsection{Vanishing of the actual compact-type orbital operator}

At each real place $v$ let $H_v,H'_v$ be finite-dimensional
Hilbert spaces carrying representations $\pi_v,\pi'_v$ of the
corresponding compact groups $U(k)$.  Put
\[
 H=\bigotimes_{v\mid\infty}H_v,\qquad
 H'=\bigotimes_{v\mid\infty}H'_v,
\]
with the external tensor-product actions
$\pi((k_v)_v)=\bigotimes_v\pi_v(k_v)$ and
$\pi'((k_v)_v)=\bigotimes_v\pi'_v(k_v)$.
Let
$\Phi:G_\infty\to H$, $\Phi':G_\infty\to H'$ be the functions
attached to the two flags.  A translate of the second function is
already included in the primed notation.  Suppose that
\[
 \Phi(ng)=\psi(n)\Phi(g),\qquad
 \Phi'(n'g)=\psi'(n')\Phi'(g).
\]
The representations of the compact stabilizers on $H,H'$ factor
through the positive quotients \eqref{df-positive-quotient}.
In particular, for the subgroup in Corollary~\ref{df-common-compact},
\[
 \Phi(kg)=\pi(k)\Phi(g),\qquad
 \Phi'(kg)=\pi'(k)\Phi'(g)\qquad(k\in K_W).
\]
For $v\in H$, $w\in H'$, write
$[v,w](\xi)=\langle w,\xi\rangle v$.  With inner products conjugate
linear in the first variable, define the orbital operator
\begin{equation}\label{df-orbital-operator}
 J=\int_{\Gamma\backslash G_\infty}
          [\Phi(g),\Phi'(g)]\,dg.
\end{equation}
The finite-place admissibility conditions on a pair of theta columns
make $\psi\overline{\psi'}$ trivial on $\Gamma$.  This is also the
precise condition needed for the integrand in
\eqref{df-orbital-operator} to descend to the quotient: rational
additive characters have product one, and their two finite-place
factors agree on the common integral stabilizer.  No normalization
of $v_R$ is needed for the following vanishing assertion.

\begin{theorem}\label{df-operator-vanishing}
Suppose that $G=0$ and that \eqref{df-orbital-operator} converges
absolutely.  At every real place let $\pi_v,\pi'_v$
be irreducible representations of $U(k)$, and suppose that they
are nonisomorphic at at least one place.  Then $J=0$.
\end{theorem}

\begin{proof}
If the flags are not relevant, the differential of
$\chi=\psi\overline{\psi'}$ is nonzero on the common real
unipotent group.  To see this globally, a nonzero $F$-valued
differential remains nonzero after applying
$\operatorname{Tr}_{F/\mathbb Q}(a\,\cdot)$ for some $a\in F$,
by nondegeneracy of the trace pairing.  Thus the character at
infinity is nontrivial.  Lemma~\ref{df-average} gives $J=0$.

Suppose the flags are relevant.  Equality of the differentials
implies equality of the characters on the common real unipotent
group, since exponential coordinates cover that connected group.
Use Theorem~\ref{df-normal-form} and
Corollary~\ref{df-common-compact} to choose $K_W$.
It fixes every vector $e_a,e'_a,z,z'$, so it normalizes both
unipotent groups and preserves their characters.  In particular,
it normalizes $R_\infty$.
By \eqref{df-disintegration} and the last part of
Lemma~\ref{df-average}, change of variables $g\mapsto kg$ gives
\[
 J=\pi(k)J\pi'(k)^{-1},\qquad k\in K_W.
\]
Here the conjugation on the second factor is essential:
$[Av,Bw]=A[v,w]B^{-1}$ for unitary $A,B$.
The isometries \eqref{df-quotient-isometries} identify the actions
with the two external tensor products of the specified $U(k)$
representations.  Changes of orthonormal basis only conjugate
one action by an inner automorphism of $U(k)$ and do not change
its isomorphism class.

For completeness, choose a place where the two irreducible
representations are nonisomorphic.  After fixing vectors and
covectors in all other tensor factors, the displayed intertwining
identity turns every matrix entry of $J$ in those factors into an
intertwiner between the two representations at that place.
Such an intertwiner is zero: its kernel and image are invariant,
so a nonzero intertwiner of irreducible representations would be
an isomorphism.  All matrix entries of $J$ vanish, proving $J=0$.
\end{proof}

\begin{corollary}\label{df-points}
Consider two distinct words $u,u'$ in the theta family, with their
specified dominant compact Harish--Chandra parameters
$\lambda_u(v),\lambda_{u'}(v)$.  Suppose these parameters differ
at one real place.  Every absolutely convergent term with $G=0$
in the operator from the $u'$-column space to the $u$-column space
is zero.  This includes the identity term whenever it occurs in
that off-diagonal expansion.
\end{corollary}

\begin{proof}
Number the diagonal coordinates from $0$ to $k-1$ and put
$\rho_a=(k-1)/2-a$.  The specified compact type has a nonzero
vector of diagonal-torus weight $\mu=\lambda-\rho$, and every
one of its weights $\nu$ satisfies
\[
 \mu-\nu=\sum_{a<b}c_{ab}(\varepsilon_a-\varepsilon_b),
 \qquad c_{ab}\in\mathbb Z_{\ge0},
\]
where $\varepsilon_a$ are the coordinate vectors.  This is the
highest-weight condition used to specify the compact types;
it suffices here without a representation-classification theorem.
Indeed, an isomorphism of the two representations would preserve
all diagonal-torus weights.  Hence both $\mu-\mu'$ and
$\mu'-\mu$ would be nonnegative combinations of the displayed
positive roots.  The linear functional
\[
             \ell(x)=\sum_{a=0}^{k-1}(k-a)x_a
\]
takes the strictly positive value $b-a$ on
$\varepsilon_a-\varepsilon_b$ for $a<b$.  Applying it to
these opposite combinations shows that every coefficient is
zero.  Thus $\mu=\mu'$, and consequently $\lambda=\lambda'$.
Different specified parameters therefore give nonisomorphic
representations.  Apply
Theorem~\ref{df-operator-vanishing}.  The corresponding term in
the theta operator is assembled from $J$ by finite-place scalar
factors and finite linear maps, all of which send the zero operator
to zero.
\end{proof}

In the source, the distinct-word hypothesis supplies different
compact parameters at every real place, so it is stronger than
the hypothesis of Corollary~\ref{df-points}.  The notation for the
source's determinants is the conjugate of one of the conventions
used in the intermediate normal-form calculation; as observed
after \eqref{df-G}, this has no effect on degeneracy or on
\eqref{df-normal-form-values}.  The calculation proves the
geometric and orthogonality content of the assertion called the
points lemma in the source~\cite{alpo-uniform-source}.
It does not prove the diagonal estimate for the sum of the
nonidentity degenerate terms.  That estimate additionally needs
the sizes of the orbital amplitudes, the count of their lattice
data, and the finite-level normalization.


\section{An exact integral for small Gram determinants}
\label{gv-section}

The geometric estimate for a degenerate orbital integral contains
the volume of tuples whose projected Gram determinant is small.
This volume admits a short proof by successive orthogonal projections.
We compute the relevant negative moments exactly, recover the submitted
estimate with a smaller constant, and also prove the optimal power
of the small determinant.  The passage from flag coordinates to this
Euclidean problem is a separate geometric step.
\cite{alpo-uniform-source}

Give $\mathbb C^n$ its usual Hermitian norm and real Lebesgue measure.
Let $P$ be an orthogonal projection of complex rank $d$, and let
$1\le j\le d\le n$.  For $v=(v_1,\ldots,v_j)$ put
\[
 \Gamma_P(v)=\Bigl[\det\bigl(\langle Pv_a,Pv_b\rangle\bigr)_{a,b=1}^j\Bigr]^{1/2}.
                                                               \tag{GV1}
\]
The determinant is real and nonnegative; its square root is the
nonnegative one.  Equivalently, $\Gamma_P$ is the norm of the
exterior product $Pv_1\wedge\cdots\wedge Pv_j$.  For $R>0$ let
\[
 \mathcal B_R=\{(v_1,\ldots,v_j):\|v_i\|\le R\text{ for every }i\}.
\]
We use the conventions $0!=1$ and that an empty product is one.

\subsection{One projected vector}

\begin{lemma}\label{gv-one-vector}
Let $Q$ be an orthogonal projection of rank $h$ on $\mathbb C^n$,
where $1\le h\le n$.  If $a$ is an integer with $0\le a<h$, then
\[
 \int_{\|v\|\le R}\|Qv\|^{-2a}\,dv
 =\frac{\pi^n R^{2(n-a)}}{(n-a)!}
       \frac{(h-a-1)!}{(h-1)!}.                               \tag{GV2}
\]
For $a\ge h$ the same integral is infinite.
\end{lemma}

\begin{proof}
The volume of a ball of radius $R$ in $\mathbb C^m$ is
$\pi^mR^{2m}/m!$.  One way to obtain the constant is to integrate
$e^{-\|x\|^2}$ first as a product of $m$ planar Gaussian integrals,
and then radially: the two values are $\pi^m$ and $m!$ times the
unit-ball volume.  Consequently radial integration in $\mathbb C^h$
uses the factor $2\pi^h r^{2h-1}/(h-1)!$.

After a unitary change of coordinates write $v=(u,w)$ with
$u\in\mathbb C^h$, $w\in\mathbb C^{n-h}$ and $Qv=(u,0)$.
Tonelli's theorem and the ball-volume formula give
\[
 \frac{2\pi^n}{(h-1)!(n-h)!}
   \int_0^R r^{2h-1-2a}(R^2-r^2)^{n-h}\,dr.
\]
This formula also covers $n=h$, since the zero-dimensional
ball has volume one.  Substituting $u=r^2/R^2$ changes the
last expression into
\[
 \frac{\pi^nR^{2(n-a)}}{(h-1)!(n-h)!}
   \int_0^1 u^{h-a-1}(1-u)^{n-h}\,du.
\]
For positive integers $b$ and nonnegative integers $c$,
successive integration by parts gives
\[
 \int_0^1 u^{b-1}(1-u)^c\,du
      =\frac{(b-1)!c!}{(b+c)!}.
\]
Use $b=h-a$, $c=n-h$ to obtain (GV2).
If $a\ge h$, restrict $w$ to a ball of radius $R/2$ and
$u$ to a sufficiently small ball.  The radial integral in
$u$ then contains a positive multiple of
$\int_0^{R/2}r^{2h-1-2a}\,dr$, which diverges.
\end{proof}

\subsection{Negative moments and a uniform sublevel estimate}

\begin{theorem}\label{gv-negative-moment}
If $a$ is an integer with $0\le a\le d-j$, then
\[
 \int_{\mathcal B_R}\Gamma_P(v)^{-2a}\,dv
 =\left(\frac{\pi^nR^{2(n-a)}}{(n-a)!}\right)^j
       \prod_{i=1}^j\frac{(d-i-a)!}{(d-i)!}.                  \tag{GV3}
\]
In particular the integral is finite.  The set $\Gamma_P=0$
has Lebesgue measure zero.  If $a\ge d-j+1$, the negative
moment in (GV3) is infinite.
\end{theorem}

\begin{proof}
Almost every projected tuple is linearly independent.  Indeed,
once $Pv_1,\ldots,Pv_{i-1}$ are independent, the condition
$Pv_i\in\operatorname{span}(Pv_1,\ldots,Pv_{i-1})$ describes a
proper complex linear subspace of $\mathbb C^n$, of codimension
$d-i+1\ge1$.  It has real Lebesgue measure zero.  Induction
and Fubini prove the assertion for the entire tuple.

For an independent partial tuple let $Q_i$ be the orthogonal
projection onto
\[
 \operatorname{im}P\cap
        \operatorname{span}(Pv_1,\ldots,Pv_{i-1})^\perp.
\]
Its rank is $d-i+1$.  Gram--Schmidt orthogonalization gives
\[
 \Gamma_P(v_1,\ldots,v_i)
   =\Gamma_P(v_1,\ldots,v_{i-1})\,\|Q_iv_i\|,
 \qquad \Gamma_P(\varnothing)=1.                              \tag{GV4}
\]
Integrate first in $v_j$.  Lemma~\ref{gv-one-vector} shows
that its contribution is
\[
 \frac{\pi^nR^{2(n-a)}}{(n-a)!}
       \frac{(d-j-a)!}{(d-j)!},
\]
independently of the preceding independent vectors.  Repeat
with $v_{j-1},\ldots,v_1$.  This proves (GV3).  All exchanges
are valid by Tonelli, even before finiteness is established;
the exceptional dependent tuples form a null set.  When $a=0$
the integrand is interpreted as one on that null set.

If $a\ge d-j+1$, then the conditional integral in $v_j$
is infinite for every independent preceding tuple, by the
last assertion of Lemma~\ref{gv-one-vector}.  Such preceding
tuples have full measure in a product of balls of positive
measure.  Tonelli therefore gives divergence.
\end{proof}

\begin{corollary}\label{gv-sublevel}
Let $1\le j<d\le n$, and write $s=d-j$.  Then for $t\ge0$,
\[
 \begin{split}
 &\operatorname{vol}\{v\in\mathcal B_R:\Gamma_P(v)\le t\}\\
 &\quad\le
  \frac{\pi^{nj}R^{2j(n-s)}t^{2s}}{((n-s)!)^j}
       \prod_{b=0}^{j-1}\frac{b!}{(s+b)!}.                    \end{split}
 \tag{GV5}
\]
The same upper bound holds if the rank of $P$ is at least $d$.
It also holds with every $v_i$ replaced by $Av_i$ both in the
norm conditions and in the Gram determinant, for any complex
linear automorphism $A$ with $|\det A|=1$.
\end{corollary}

\begin{proof}
For $t>0$, outside the null set $\Gamma_P=0$,
\[
        {\bf1}_{\{\Gamma_P\le t\}}
                    \le t^{2s}\Gamma_P^{-2s}.
\]
Integrating and applying (GV3) with $a=s$ proves (GV5).
For $t=0$ the null-set assertion proves the result directly.

If $\operatorname{rank}P\ge d$, choose a $d$-dimensional
subspace of its image and let $P_0$ be the orthogonal projection
onto it.  Then $\Gamma_{P_0}(v)\le\Gamma_P(v)$.  For example,
in coordinates adapted to the two images, Cauchy--Binet writes
the squared Gram determinant as a sum of squared absolute
values of minors; the smaller projection retains a subset
of these terms.  The sublevel set for $P$ is therefore
contained in the corresponding sublevel set for $P_0$.

Finally the real Jacobian of a complex linear map $A$ is
$|\det A|^2$.  The simultaneous substitution $w_i=Av_i$
preserves product Lebesgue measure when $|\det A|=1$.
\end{proof}

\subsection{The optimal power of the threshold}

The divergent moment at the next integer does not prevent
a stronger sublevel estimate.  To obtain it, integrate the
last perpendicular ball directly before applying a negative
moment to the preceding vectors.

\begin{theorem}\label{gv-optimal-power}
Let $1\le j\le d\le n$ and $q=d-j+1$.  Then for $t\ge0$,
\[
 \begin{split}
 &\operatorname{vol}\{v\in\mathcal B_R:\Gamma_P(v)\le t\}\\
 &\quad\le
  \frac{\pi^{nj}R^{2j(n-q)}t^{2q}}{q!\,((n-q)!)^j}
       \prod_{i=1}^{j-1}\frac{(j-i-1)!}{(d-i)!}.              \end{split}
 \tag{GV6}
\]
The extensions to rank at least $d$ and to $|\det A|=1$
in Corollary~\ref{gv-sublevel} remain valid.
For a projection of rank exactly $d$, no exponent larger
than $2q$ can replace $2q$ in an estimate with a finite
constant independent of $t\in(0,1)$.
\end{theorem}

\begin{proof}
First suppose $j\ge2$, and fix an independent preceding tuple
$v'=(v_1,\ldots,v_{j-1})$.  Write $\gamma=\Gamma_P(v')$.
The decomposition defined by $Q_j$ splits $v_j$ into a vector
in $\mathbb C^q$ and a vector in $\mathbb C^{n-q}$.
By (GV4) the first vector has norm at most $t/\gamma$.
The second has norm at most $R$.  Dropping their coupled
ball constraint gives the conditional volume bound
\[
 \frac{\pi^q}{q!}(t/\gamma)^{2q}
       \frac{\pi^{n-q}R^{2(n-q)}}{(n-q)!}.
\]
For the preceding $j-1$ vectors the exponent $a=q$ satisfies
$a=d-(j-1)$, so (GV3) applies.  Multiplication of its value
by the displayed factor proves (GV6).  For $j=1$, the same
two-ball estimate directly proves the formula, with the
product empty.  The case $t=0$ follows from the null-set
assertion.  The two extensions follow exactly as before.

To prove the optimality assertion, choose an orthonormal
basis of $\operatorname{im}P$.  When $j\ge2$, restrict
$v_1,\ldots,v_{j-1}$ to sufficiently small neighborhoods
of $(R/4)e_1,\ldots,(R/4)e_{j-1}$.  This gives a measurable
set of positive volume on which the projected vectors are
independent, all norms are less than $R/2$, and
\[
                  0<\gamma\le C=(R/2)^{j-1}.
\]
Decompose the last vector orthogonally into its components
in the preceding projected span, its $q$-dimensional
complement within $\operatorname{im}P$, and $\ker P$.
Restrict the first and third components to balls of radius
$R/4$, and the second to a ball of radius $t/C$.  For
$0<t\le CR/4$, all these last vectors have norm less
than $R$, and (GV4) gives $\Gamma_P\le t$.
The conditional volume is a fixed positive constant times
$t^{2q}$, independently of the chosen preceding tuple.
After integration over that tuple, the full sublevel
volume is at least $c\,t^{2q}$ for some $c>0$.
For $j=1$ use the same decomposition with the preceding
span absent and $C=1$.

If a bound $C' t^{2q+\delta}$ with $\delta>0$ held for
all sufficiently small $t$, division by $t^{2q}$ would
give $c\le C't^\delta$, a contradiction as $t\to0$.
\end{proof}

\subsection{The two volume bounds in the orbital calculation}

For direct comparison with the submitted estimates, put
$n=d+c$, $s=d-j\ge1$, and $R=M\ge1$.  Each factorial
ratio in (GV5) satisfies
\[
             \frac{b!}{(s+b)!}\le\frac1{s!}.
\]
Also $s!\ge(s/e)^s$: the term $s^s/s!$ of the exponential
series is at most $e^s$.  Hence (GV5) implies
\[
 \begin{split}
 &\operatorname{vol}\{v:\|Av_i\|\le M,\quad\Gamma_P(Av_1,\ldots,Av_j)\le t\}\\
 &\quad\le
   \frac{\pi^{j(j+c)}M^{2j(j+c)}}{((j+c)!)^j}
        \left(\frac{\pi e}{s}\right)^{js}t^{2s},
                                                               \end{split}
 \tag{GV7}
\]
provided $\operatorname{rank}P\ge d$ and $|\det A|=1$.

For the degenerate orbital integral, take
\[
 n=k+1,\qquad d=k-1,\qquad c=2,\qquad
 s=k-1-j,\qquad t=2y,
\]
where $1\le j$ and $j+2\le k$.  The source's middle
isometry $m$ satisfies $m^*H m=H$ for a nondegenerate
Hermitian matrix $H$.  Taking determinants gives
$|\det m|^2=1$, so $A=m^{-1}$ has the required Jacobian.
The sublevel volume is therefore at most
\[
 \frac{\pi^{j^2+2j}M^{2j^2+4j}}{((j+2)!)^j}
     \left(\frac{\pi e}{k-1-j}\right)^{j(k-1-j)}
     (2y)^{2(k-1-j)}.                                        \tag{GV8}
\]
The submitted bound has the same expression without the
factorial denominator and with an additional factor $2$.
Thus (GV8) improves its constant by at least
$2((j+2)!)^j$, with exactly the same hypotheses and powers
of $M,y$.  It also proves that the zero Gram locus is null.

The other orbital volume bound uses
$d=k-2$, $c=3$, $j=r$, and $s=k-2-r\ge1$.
The same argument gives
\[
 \frac{\pi^{r^2+3r}M^{2r^2+6r}}{((r+3)!)^r}
      \left(\frac{\pi e}{k-2-r}\right)^{r(k-2-r)}
      (2y)^{2(k-2-r)}.                                       \tag{GV9}
\]
This bounds the entire sublevel set and hence also its shell
$y<\Gamma_P\le2y$.  The submitted expression has a larger
factor $4M^{2r}$ and omits the factorial denominator.
Its assumed range $2r+6\le k$ is stronger than needed here.

The sharper exponent in (GV6) is also available for both
applications.  Using it in the complete orbital bound would
change the subsequent determinant powers and their comparison
with the lattice count, which must be tracked separately.
For the present reconstruction, (GV8)--(GV9) already supply
the exact volume estimates required by the existing argument.
They do not establish the flag-coordinate identity, the
coordinate-size bounds, or the normalization of the common
unipotent stabilizer; those remain separate inputs to the
orbital calculation.


\section{A sharper dyadic estimate for the degenerate orbital integral}
\label{ods-section}

The last summation in the degenerate amplitude estimate is a
measure-theoretic argument involving two size functions and two
Gram determinants. We prove it independently and improve both its
constant and its decay exponent. The geometric estimates supplying
its hypotheses remain separate; in particular the relation between
the determinant and the size functions cannot be omitted.
\cite{alpo-uniform-source}

Let $(X,\nu)$ be a measure space. All integrals of nonnegative
functions below are allowed initially to be infinite. Let
$A,B,S,S',Y,Y'$ be measurable nonnegative real functions, with
$S,S'\ge1$ and $Y,Y'>0$ almost everywhere. Fix integers
$r\ge1$, $j\ge0$, real numbers $m\ge1$, $e\ge0$, and
\[
 q\ge e+3r+1,\quad H,\Delta,\sigma\ge1,
 \qquad D,K,K'>0.
                                                               \tag{OD1}
\]
Suppose that almost everywhere
\begin{align}
 YY'D&\le2^j,& D&\le(SS')^r,\nonumber\\
 Y,Y'&\le\max(S,S')^r\Delta,&
 \sigma&\le\max\left(1,\frac{(SS')^r}{(2H)^{2r}}\right).
                                                               \tag{OD2}
\end{align}
For $i\ge0$ write
\[
 E_i(S)=\{S\le2^{i+1}H\}\cap
       \bigl(\{i=0\}\cup\{2^iH<S\}\bigr),
\]
where $\{i=0\}$ means all of $X$ if $i=0$ and the empty set
otherwise. Define $E_i(S')$ in the same way. Suppose that for all
$i,i'\ge0$ and $y>0$,
\begin{align}
 \int_{E_i(S)\cap\{S'\le2^{i'+1}H\}\cap\{y<Y\le2y\}}
                    A^2\,d\nu
 &\le K(2^{i+i'+1}H)^e y^{2m}2^{-2qi},\label{ods-first}\\
 \int_{E_{i'}(S')\cap\{S\le2^{i+1}H\}\cap\{y<Y'\le2y\}}
                    B^2\,d\nu
 &\le K'(2^{i+i'+1}H)^e y^{2m}2^{-2qi'}.
                                                               \label{ods-second}
\end{align}

\begin{theorem}\label{ods-dyadic}
Under these hypotheses, $AB$ is integrable and
\begin{equation}\label{ods-bound}
 \int_X AB\,d\nu
 \le 4\,2^{-j}\sqrt{KK'}(2H)^{e+4r}\Delta^2
        \left(\frac{2^j}{D}\right)^m
        \sigma^{-(q-e-3r-1)/r}.
\end{equation}
No local compactness or finiteness assumption on $\nu$ is required.
\end{theorem}

\begin{proof}
Discard a null set so that the pointwise hypotheses hold everywhere.
The sets $E_i(S)$ form a disjoint countable partition, and likewise
for $S'$. On their intersection put
\[
 u=2^{i+i'},\qquad R=2Hu,\qquad M=R^r\Delta,
 \qquad N=2^j/D.
\]
Then
\begin{equation}\label{ods-size-product}
 \max(S,S')\le R,\qquad SS'\le u(2H)^2,
 \qquad \sigma\le u^r.
\end{equation}
In particular $0<Y,Y'\le M$. Subdivide into the measurable sets
$P_{i,i',a,b}$, for $a,b\ge0$, on which
\[
       M2^{-a-1}<Y\le M2^{-a},\qquad
       M2^{-b-1}<Y'\le M2^{-b}.
\]
These form a disjoint countable partition of the intersection.
Set $y=M2^{-a-1}$, $y'=M2^{-b-1}$ and $w=yy'$.
On a nonempty piece the product condition in (OD2) gives
$w\le N$. Conditions \eqref{ods-first}--\eqref{ods-second},
followed by Cauchy--Schwarz on this piece, give
\begin{equation}\label{ods-piece-cs}
 \int_{P_{i,i',a,b}}AB\,d\nu
 \le \sqrt{KK'}R^e w^m u^{-q}.
\end{equation}
The square integrals in this use of Cauchy--Schwarz are finite by
the same two hypotheses. Empty pieces contribute zero.

We retain two factors that can be lost in a rougher summation.
First,
\[
 w=M^2\,2^{-a-b-2}
     =(2H)^{2r}u^{2r}\Delta^2\,2^{-a-b-2}.
\]
Second, on each nonempty piece the product bound in
\eqref{ods-size-product} gives
\[
                    N^{-1}=D2^{-j}
                          \le2^{-j}u^r(2H)^{2r}.
\]
As $m\ge1$ and $w\le N$, these imply
\begin{equation}\label{ods-power}
 \begin{split}
 w^m&\le N^{m-1}w\\
 &\le N^m\,2^{-j-2}(2H)^{4r}u^{3r}\Delta^2\,2^{-a-b}.
 \end{split}
\end{equation}
Let $n=q-e-3r-1\ge0$. Since $\sigma\le u^r$, one has
$u^{-n}\le\sigma^{-n/r}$. Substitution in
\eqref{ods-piece-cs} yields
\[
 \int_{P_{i,i',a,b}}AB\,d\nu
 \le C\,2^{-i-i'-a-b},
 \quad
 C=2^{-j-2}\sqrt{KK'}(2H)^{e+4r}\Delta^2N^m\sigma^{-n/r}.
\]
Finally, monotone convergence and the four geometric series give
\[
 \int_XAB\,d\nu
 \le C\left(\sum_{h=0}^{\infty}2^{-h}\right)^4=16C.
\]
This is \eqref{ods-bound}; its right side is finite. The proof
also covers $n=0$, so equality is allowed in the bound for $q$.
\end{proof}

\begin{corollary}\label{ods-source-bound}
Suppose $m,e,q$ are integers, $m\ge1$, and
$q\ge e+4r+2$, with all the other hypotheses above. Then
\[
 \int_X AB\,d\nu
 \le16\sqrt{KK'}(2H)^{e+4r}\Delta^2
          (2^j/D)^m\sigma^{-(q-e-4r-1)/r}.
                                                               \tag{OD3}
\]
Thus the exact dyadic summation used for the submitted degenerate
orbital estimate follows. In fact \eqref{ods-bound} improves its
constant by a factor $4\cdot2^j$ and supplies one additional power
of $\sigma^{-1}$ under these same hypotheses.
\end{corollary}
\begin{proof}
The hypothesis on $q$ implies (OD1). In the decay exponent,
$q-e-3r-1=(q-e-4r-1)+r$. Since $\sigma\ge1$ and $j\ge0$,
the right side of \eqref{ods-bound} is at most the right side of
(OD3), with precisely the stated factors.
\end{proof}

\subsection{Why the determinant--size condition is essential}

\begin{example}\label{ods-counterexample}
The assertion (OD3) is false if the condition $D\le(SS')^r$ is
deleted while its other hypotheses are retained. To see this,
let $L\ge512$ be an integer and give
$X=\{0,1,\ldots,2L\}$ counting measure. Put
\[
 \begin{gathered}
 r=m=H=\Delta=\sigma=K=K'=1,
 \quad j=e=0,
 \quad q=6,
 \quad N=2^{-2L},\quad D=N^{-1},\\
 S(a)=S'(a)=1,\quad
 Y(a)=2^{-a},\quad Y'(a)=2^{a-2L},
 \quad A(a)=Y(a)/2,\quad B(a)=Y'(a)/2.
 \end{gathered}
\]
Here $YY'D=1$, both $Y,Y'$ are at most one, and the condition
on $\sigma$ holds. Each interval $(y,2y]$ contains at most one
of the distinct powers of two $Y(a)$. At that point
$A(a)^2=Y(a)^2/4\le y^2$. Thus \eqref{ods-first} holds when
$i=0$, independently of $i'$, and its domain is empty when
$i\ge1$. The identical argument with $Y'$ proves
\eqref{ods-second}. All the retained assumptions are satisfied.
But
\[
                  \int_XAB\,d\nu=(2L+1)N/4>256N,
\]
whereas the right side of (OD3) is exactly $256N$.
The failed hypothesis is $D\le1$.
This example is not a counterexample to the submitted theorem,
which includes the determinant--size condition explicitly.
\end{example}

\subsection{Role in the orbital estimate}

In the degenerate orbital application, $m=k-1-j$, $e=16r^3$,
$D$ is the absolute value of the nonzero leading determinant,
and $Y,Y'$ are the corresponding normalized Gram determinants
on the two sides. The functions $S,S'$ are the maxima of the
flag lengths and their inverses. The two one-sided square-integral
estimates supply \eqref{ods-first}--\eqref{ods-second}; their
constants retain the unipotent Haar normalization and the stabilizer
measure. All of these are actual hypotheses, not consequences of
dyadic summation. Theorem~\ref{ods-dyadic} proves the final
measure-theoretic passage as soon as those geometric and
representation-theoretic estimates have been supplied. It does
not assert that a polynomial tail estimate by itself controls a
degenerate orbital integral.


\section{Local harmonic analysis in the theta estimates}
\label{sec:theta-local-audit}

This section reconstructs several analytic and local arithmetic steps that
enter the proposed estimates for the theta maps.  Every lemma proved here is
independent of the Riemann hypothesis.  The final estimate for the actual
global family is stated with its remaining arithmetic hypotheses displayed.
In particular, a local stationary phase formula is not a proof of the global
estimate for the sum of the orbital terms.

The source uses the normalized operators
\begin{equation}\label{th-normalized-K}
 K_{uv}=(c_uc_v)^{-1/2}P_u^*P_v,
 \qquad
 c_u=\frac{\kappa\,\operatorname{vol}(N_u)
                   \operatorname{vol}(H_u)}{\dim\pi_u}.
\end{equation}
Here $P_u:C_u\longrightarrow\mathcal H$ is a Poincar\'e-series map,
$C_u$ is a finite-dimensional Hilbert space, $\pi_u$ is a nonzero compact
stabilizer representation, and $\kappa$ is the volume of the finite-level
subgroup.  The source incorporates $\kappa$ into its symbol
$\operatorname{vol}(N_u)$.  All inner products below are conjugate linear
in the first variable.  Thus $K_{uv}:C_v\longrightarrow C_u$.

\subsection{The local rings and measures actually required}

Let $K$ be a finite product of finite extensions of $\mathbb Q_p$.  A
compact open additive subgroup of $K^n$ will be called a lattice.  For
$a\in K^\times$, write
\[
 |a|_K=\prod_{K_j\mid K}|a_j|_{K_j},
\]
where the absolute value in a field factor gives a uniformizer the value
$(\#k_j)^{-1}$.  For a matrix $B\in\operatorname{GL}_n(K)$,
linear change of variables multiplies additive Haar measure by
$|\det B|_K$.

Fix a continuous additive character $\psi:K\to\mathbb C^\times$ that is
nontrivial on every field factor.  Choose the self-dual Haar measure, and
use its product measure on $K^n$.  For a lattice $L\subset K^n$, put
\[
 L^\perp=\{y\in K^n:\psi(x^ty)=1\text{ for all }x\in L\}.
\]
Then $(L^\perp)^\perp=L$ and
\begin{equation}\label{th-dual-vol}
 \operatorname{vol}(L)\operatorname{vol}(L^\perp)=1,
 \qquad
 \int_L\psi(x^ty)\,dx=
 \operatorname{vol}(L)\,1_{L^\perp}(y).
\end{equation}
For completeness, the second identity follows by translating the integral
by an element $x_0\in L$: if the character is nontrivial, that translation
multiplies the integral by $\psi(x_0^ty)\ne1$.  The duality and volume
identities reduce by a linear change of basis to fractional ideals in a
local field.  If the annihilator of its ring of integers $O$ is $c^{-1}O$,
the measure for the character $x\mapsto\psi_0(cx)$ is
$|c|^{1/2}\,dx_0$, where $\psi_0$ has annihilator $O$ and
$\operatorname{vol}_{dx_0}(O)=1$.  The identities for fractional ideals
follow immediately.  A product of these calculations proves
\eqref{th-dual-vol}; a general additive lattice is a free
$\mathbb Z_p$-module, so the same calculation can alternatively be made
over $\mathbb Q_p$ using the trace pairing.

The two local orders in the source are as follows.

\begin{enumerate}
\item Away from the two distinguished split primes, $K$ is an unramified
quadratic algebra over a local field $F$, $R=O_K$, and the specified
character has annihilator $O_K$.  This includes the algebra $F\times F$.
The chosen measure has $\operatorname{vol}(R)=1$.
\item At a distinguished split prime, write $K=F\times F$, let $\varpi$
be a uniformizer of $F$, and let $\psi_F$ have annihilator $O_F$.  The
order and character are
\begin{equation}\label{th-split-order}
 R_\varpi=\{(x,y)\in O_F^2:x-y\in\varpi O_F\},
 \qquad
 \psi_R(x,y)=\psi_F((x-y)/\varpi).
\end{equation}
\end{enumerate}

\begin{lemma}\label{th-split-self-dual}
The order $R_\varpi$ is its own annihilator for $\psi_R$.
If $q=\#(O_F/\varpi O_F)$ and product Haar measure gives $O_F^2$
volume one, self-dual Haar measure for $\psi_R$ is $q$ times that
product measure.  In particular, it gives $R_\varpi$ volume one.
\end{lemma}

\begin{proof}
Write an element of $R_\varpi$ as $(y+\varpi z,y)$, with $y,z\in O_F$.
An element $(a,b)\in F^2$ annihilates this order if and only if
\[
 \psi_F\bigl(az+(a-b)y/\varpi\bigr)=1
 \quad(y,z\in O_F).
\]
This is equivalent to $a\in O_F$ and $(a-b)/\varpi\in O_F$, which
is precisely $(a,b)\in R_\varpi$.  The characters on the two factors
have scale factors $\varpi^{-1}$ and $-\varpi^{-1}$.  Their self-dual
measures therefore contribute $q^{1/2}$ each.  Finally
$[O_F^2:R_\varpi]=q$.
\end{proof}

Consequently it would be incorrect to replace all local orders by maximal
orders while retaining the same Haar measure.  The source does retain the
normalization in Lemma~\ref{th-split-self-dual}.  The proofs below allow
this nonmaximal order; they use the additive pairing, rather than a claim
that $R$ is a discrete valuation ring.

\subsection{A Gaussian integral with all normalizing factors}

Let $B\in\operatorname{GL}_n(K)$ be symmetric and put
$q_B(x)=\tfrac12x^tBx$.  Characteristic zero is essential here.
For a lattice $L\subset K^n$, impose the explicit condition
\begin{equation}\label{th-gauss-admissible}
 \psi\bigl(q_{B^{-1}}(\xi)\bigr)=1\qquad(\xi\in L^\perp).
\end{equation}
Such lattices exist: take $L$ sufficiently large, so that $L^\perp$
is contained in a neighborhood on which the indicated continuous
character-valued function is one.

\begin{lemma}\label{th-gaussian}
If \eqref{th-gauss-admissible} holds, then
\[
 \left|\int_L\psi(q_B(x))\,dx\right|=|\det B|_K^{-1/2}.
\]
The integral is independent of the lattice $L$ satisfying that condition.
Accordingly there is a well-defined complex number $\gamma_\psi(B)$,
of absolute value one, for which
\begin{equation}\label{th-gaussian-value}
 \int_L\psi(q_B(x))\,dx
   =\gamma_\psi(B)|\det B|_K^{-1/2}.
\end{equation}
It is unchanged under congruence:
$\gamma_\psi(C^tBC)=\gamma_\psi(B)$ for $C\in\operatorname{GL}_n(K)$.
\end{lemma}

\begin{proof}
Set $U=B^{-1}L^\perp$.  Condition \eqref{th-gauss-admissible} says
$\psi(q_B(u))=1$ for $u\in U$.  Polarization gives
$\psi(u^tBv)=1$ for $u,v\in U$.  Since $B$ is symmetric,
$U^\perp=BL$.  It follows that $BU=L^\perp\subset BL$, or $U\subset L$.
Writing $I_L=\int_L\psi(q_B(x))\,dx$ and making the change of variables
$x=y+u$ gives
\begin{align*}
 |I_L|^2
 &=\int_L\psi(q_B(u))
       \int_L\psi(y^tBu)\,dy\,du\\
 &=\operatorname{vol}(L)
       \int_{L\cap B^{-1}L^\perp}\psi(q_B(u))\,du\\
 &=\operatorname{vol}(L)\operatorname{vol}(B^{-1}L^\perp)
   =|\det B|_K^{-1}.
\end{align*}
All these integrals are over compact sets, so Fubini causes no
convergence issue.

If $a\notin L$, duality supplies $\xi\in L^\perp$ with
$\psi(\xi^ta)\ne1$.  Put $u=B^{-1}\xi\in U\subset L$.
Translation by $u$ multiplies the integral of $\psi(q_B)$ over $a+L$
by $\psi(\xi^ta)$: the other factors are one because
$\psi(q_B(u))=1$ and $\xi\in L^\perp$.  That coset integral is zero.
Thus enlarging $L$ to any lattice leaves $I_L$ unchanged.  For two
admissible lattices, enlarge both to their sum.  This proves independence.
Finally substitute $x=Cy$ in the integral.  The factor $|\det C|_K$
is canceled by the change in $|\det(C^tBC)|_K^{1/2}$.
\end{proof}

For $m\ge1$, let $J_m$ denote the matrix all of whose entries are one,
and put
\[
 H_m(a)=a^{-1}(I_m+J_m),\qquad
 \gamma_m(a)=\gamma_\psi(H_m(a))\quad(a\in K^\times).
\]
The determinant lemma, or the decomposition into the line spanned by
$(1,\ldots,1)$ and its sum-zero complement, gives
\begin{equation}\label{th-hessian-det}
 \det H_m(a)=(m+1)a^{-m},\qquad
 (I_m+J_m)^{-1}=I_m-\frac{J_m}{m+1}.
\end{equation}
The identities remain valid when the residue characteristic divides
$m+1$: division takes place in a characteristic-zero field.  The
congruence assertion in Lemma~\ref{th-gaussian} gives
$\gamma_m(ab^2)=\gamma_m(a)$.  Thus
\begin{equation}\label{th-gauss-H}
 \int_L\psi\bigl(\tfrac12x^tH_m(a)x\bigr)\,dx
 =\gamma_m(a)|m+1|_K^{-1/2}|a|_K^{m/2}
\end{equation}
whenever \eqref{th-gauss-admissible} is satisfied.

\begin{remark}\label{th-conductor}
There is no omitted conductor factor in \eqref{th-gauss-H}, because the
Haar measure is self-dual for the displayed character.  If a field factor
uses $\psi_c(x)=\psi_0(cx)$ and one insists on the measure for $\psi_0$,
the right side must instead be multiplied by $|c|^{-m/2}$.
Moreover the admissibility condition changes with $c$.  One may not
change the character and keep both the measure and the smallness
threshold unchanged.
\end{remark}

\subsection{An elementary analytic change of variables}

We record the analytic fact needed for stationary phase, including a
proof that does not assume that $2$ or $m+1$ is a unit of $O_F$.

\begin{lemma}\label{th-morse}
Let $F$ be a finite extension of $\mathbb Q_p$ and let $f$ be analytic
near $0\in F^m$.  Suppose $df(0)=0$ and the symmetric Hessian
$H=D^2f(0)$ is invertible over $F$.  On a sufficiently small ball
$B=(\varpi^aO_F)^m$ there is an analytic bijection
$\Phi:B\to B$ such that
\[
 f(y)=f(0)+\tfrac12\Phi(y)^tH\Phi(y).
\]
The map preserves Haar measure, and for every $b\ge a$ and $y,y'\in B$,
\[
 \Phi(y)-\Phi(y')\in(\varpi^bO_F)^m
 \quad\Longleftrightarrow\quad y-y'\in(\varpi^bO_F)^m.
\]
For an analytic family the ball can be chosen uniformly under the
following precise hypothesis: on a common polydisc of radius $r>0$,
the convergent power series $\sum_\alpha c_\alpha y^\alpha$ have
uniformly bounded Gauss norms
$\sup_\alpha |c_\alpha|r^{|\alpha|}$, and the inverse Hessians have
uniformly bounded matrix norms.  Shrinking the common radius is allowed.
\end{lemma}

\begin{proof}
Expanding the convergent power series and assigning each monomial of
degree at least two to two of its variables writes
\[
 f(y)-f(0)=\tfrac12y^t(H+C(y))y,
 \qquad C(0)=0,
\]
with $C(y)$ symmetric and analytic.  One explicit assignment is to put a
monomial into a diagonal entry if some variable occurs at least twice,
and otherwise into the two symmetric entries corresponding to its
first two variables.  Only the fixed factors $2$ and $1/2$ are
introduced; after shrinking the polydisc the resulting series converge.
Put $A(y)=H^{-1}C(y)$.  The matrix $A(y)$ is self-adjoint for the
bilinear form $H$, namely $A(y)^tH=HA(y)$.

The formal series
\[
 S(y)=\sum_{j\ge0}\binom{1/2}{j}A(y)^j
\]
converges on a sufficiently small polydisc.  Indeed the $p$-adic
absolute values of the coefficients are bounded by an exponential in
$j$, whereas the matrix norm of $A(y)$ tends to zero with $y$.
It satisfies $S(y)^2=I+A(y)$ and $S(y)^tH=HS(y)$.  Therefore
$S(y)^tHS(y)=H+C(y)$.  Set $\Phi(y)=S(y)y$.

Here are the required invertibility and measure details.  Shrink $B$
until $g(y)=\Phi(y)-y$ has Lipschitz constant strictly less than one,
and $g(0)=0$.  This is possible because its convergent power series
has no constant or linear term.  If $b\in B$, the map
$y\mapsto b-g(y)$ is a strict contraction of the complete ball $B$,
so it has a unique fixed point.  This proves bijectivity.
Furthermore
\[
 \|\Phi(y)-\Phi(y')\|=\|y-y'\|
\]
in the supremum norm, since the error has strictly smaller norm.
The map consequently permutes the residue classes of each smaller
ball.  These classes generate the Borel sets of $B$ and all have
the same volume; hence $\Phi$ preserves Haar measure.  The power-series
bounds just used can be made uniform when the coefficients and
$H^{-1}$ satisfy the stated uniform bounds.  Indeed on any strictly
smaller radius the uniform Gauss bound gives uniform convergence of
the tails and of their divided difference estimates.
\end{proof}

\subsection{The stationary phase formula used in the local induction}

For $D\in K^\times$ define, on $(K^\times)^m$,
\[
 f_D(x_1,\ldots,x_m)=x_1+\cdots+x_m+
                         \frac{D}{x_1\cdots x_m}.
\]
Fix a compact subset $Q\subset K^\times$, a compact open set
$C\subset(K^\times)^m$, and a lattice $L_1\subset K^m$.
The amplitudes considered below are all locally constant functions
$A:K^m\to\mathbb C$, supported in $C$ and invariant under translation
by $L_1$.  No bound on their values is required.  Choose $L_1$ smaller
if necessary; doing so enlarges the allowed class of amplitudes.

\begin{theorem}\label{th-stationary}
There is a product neighborhood $V$ of zero in $K$, depending on
$K,\psi,m,Q,C,L_1$, such that for $t\in V\cap K^\times$ and $D\in Q$,
\begin{align}\label{th-stationary-formula}
 &\int_{(K^\times)^m}A(x)\psi(f_D(x)/t)\,dx\notag\\
 &\quad=|t|_K^{m/2}|m+1|_K^{-1/2}
   \sum_{\substack{z\in K^\times\\z^{m+1}=D}}
    \gamma_m(zt)|z|_K^{m/2}
    \psi((m+1)z/t)A(z,\ldots,z).
\end{align}
The sum is finite.  If $K$ is a field it has at most $m+1$ terms;
if $K$ has $d$ field factors it has at most $(m+1)^d$ terms.
\end{theorem}

\begin{proof}
We first take $K=F$ to be a field.  Direct differentiation gives
\[
 \frac{\partial f_D}{\partial x_i}
   =1-\frac{D}{x_i\prod_jx_j}.
\]
All derivatives vanish if and only if $x_1=\cdots=x_m=z$ and
$z^{m+1}=D$.  At that point the diagonal Hessian entries are $2/z$
and the off-diagonal entries are $1/z$.  Thus the Hessian is $H_m(z)$,
which is invertible by \eqref{th-hessian-det}, and
$f_D(z,\ldots,z)=(m+1)z$.

The set of possible roots $z$, for $D\in Q$, is compact in $F^\times$:
their absolute values lie in finitely many valuation classes, bounded
above and away from zero.  The coefficients of the Taylor expansion
at these roots, and the inverse Hessians, are therefore bounded
uniformly.  More precisely, choose a common expansion radius strictly
smaller than the least absolute value of these roots.  On that radius
the reciprocal factors in $f_D$ have geometrically convergent series
with uniformly bounded Gauss norms.  This is the uniform analytic
hypothesis of Lemma~\ref{th-morse}, which supplies a common sufficiently small
ball $B=(\varpi^aO_F)^m\subset L_1$.  We may also arrange that distinct
critical points have disjoint translates of $B$.  To see uniform
separation, otherwise there would be distinct roots $z_j,z'_j$ with
$z_j-z'_j\to0$; after taking a convergent subsequence, their common
limit $z$ is nonzero.  The function $w\mapsto w^{m+1}$ has derivative
$(m+1)z^m\ne0$, so is injective on a neighborhood of $z$, a
contradiction.  Its local injectivity follows by exactly the contraction
argument in Lemma~\ref{th-morse}, applied after its nonzero linear term
has been divided out.

On a critical ball, $A$ is constant.  Applying the measure-preserving
change of variables and then \eqref{th-gauss-H} gives
\[
 \int_{(z,\ldots,z)+B}A(x)\psi(f_D(x)/t)\,dx
 =A(z,\ldots,z)\psi((m+1)z/t)
   \gamma_m(zt)|m+1|^{-1/2}|zt|^{m/2}.
\]
The use of the Gaussian formula is justified uniformly for small $t$.
In fact
\[
 (H_m(z)/t)^{-1}
       =tz\left(I_m-\frac{J_m}{m+1}\right).
\]
On the fixed compact lattice $B^\perp$, the associated quadratic
form tends uniformly to zero as $t\to0$, with $z$ in its compact
set.  It is consequently in the kernel of $\psi$ for all sufficiently
small $t$.  This proves the admissibility condition, including at
primes dividing $2(m+1)$.

We give the uniform vanishing argument on the complement.  Near any
pair $(D_0,x_0)$ with $\nabla f_{D_0}(x_0)\ne0$, choose $j$ for which
$|a|=|\partial_jf_{D_0}(x_0)|$ is the largest absolute value among
the partial derivatives.  In particular $a\ne0$.  On sufficiently small product balls
in $D$ and $x$, the map
\[
 x\longmapsto(f_D(x),x_1,\ldots,\widehat{x_j},\ldots,x_m)
\]
is a bijection onto a product of balls.  Its first ball is a translate
of $a\varpi^bO_F$; the other balls have radius $|\varpi|^b$.
Its Jacobian has constant absolute value $|a|$.  These facts follow
by dividing the first coordinate by $a$ and applying the same strict
contraction argument: the nonlinear remainder and the variation in
the derivative can both be made smaller than the linear term.
For fixed other coordinates, this is a one-variable bijection of
balls, preserving their normalized Haar measures.  Shrink the
balls further so that $A$ is constant there.  Integration in the
first coordinate now vanishes as soon as
$\psi$ is nontrivial on $t^{-1}a\varpi^bO_F$.

Only finitely many such bounds are needed.  Here is the compactness
point needed for that assertion.  A root $z_0^{m+1}=D_0$ continues
uniquely, in a fixed small root ball, for all $D$ sufficiently near
$D_0$, by the preceding local inverse argument.  Roots outside those
balls cannot suddenly enter the fixed compact set under consideration:
a convergent subsequence would give an additional root at $D_0$.
Thus after shrinking a neighborhood of each $D_0$, the critical
balls have fixed disjoint locations and their centers vary within
those balls.  The remaining part of $C$ is a compact open set without
critical points, uniformly for these $D$.  It is a finite disjoint
union of balls of the kind just described.  A finite cover of $Q$
provides finitely many choices of $a,b$.  Choose $V$ small enough to
make all corresponding character integrals vanish.  Summing the
critical-ball integrals proves \eqref{th-stationary-formula}.

For a product $K=\prod F_j$, the additive character and self-dual
measure factor over $j$, as do the phase and the critical equations.
A compactly supported function invariant under a fixed product
lattice is a finite linear combination of products of characteristic
functions of balls in the factors.  Apply the field result to each
such product and use linearity and Fubini.  Enclose the projections
of $Q$ in compact subsets of $F_j^\times$ and take the product of the
resulting neighborhoods.  This proves the full assertion, with
uniformity for all the stated amplitudes.
\end{proof}

\begin{remark}\label{th-uniformity-limit}
The theorem gives a common neighborhood for the displayed fixed local
data.  It does not give a neighborhood independent of $m$, $p$, the
field, or the character conductor.  In particular the inverse Hessian
contains $(m+1)^{-1}$, and the analytic square root used in
Lemma~\ref{th-morse} depends on the valuation of $2$.  This dependence
cannot be discarded when passing to a varying global family.
The source's local theorem likewise quantifies its neighborhood after
the local field, character, dimension, amplitude scale, and compact
parameter set have been fixed.  Its later passage from these local
identities to a uniform positive majorant requires a separate argument.
\end{remark}

\subsection{Finite cokernels and positivity of the local kernel}

Let $R\subset K$ be either of the self-dual compact open orders above.
For $Z\in M_n(R)\cap\operatorname{GL}_n(K)$, define
\[
 C_Z=R^n/ZR^n,\qquad e_Z=[e_n],\qquad
 \lambda_Z([v])=e_n^tZ^{-1}v\pmod R\quad\text{in }K/R.
\]
For $n=0$ take $C_Z=0$, $e_Z=0$, and $\lambda_Z=0$.
The functional $\lambda_Z$ is well-defined because replacing $v$ by
$v+Zw$ adds $e_n^tw\in R$.  Call $(C_Z,e_Z,\lambda_Z)$ the marked
finite module associated to $Z$.  An isomorphism of two such modules
is an $R$-linear isomorphism carrying the marked element and the
functional to their counterparts.

\begin{lemma}\label{th-coker}
The module $C_Z$ is finite, and
\[
 \#C_Z=|\det Z|_K^{-1}.
\]
These assertions do not require $R$ to be a principal ideal domain.
\end{lemma}

\begin{proof}
The submodule $ZR^n$ is open and compact, since $Z$ is an invertible
continuous linear map.  It is contained in $R^n$.  The quotient is
therefore compact and discrete, hence finite.  All its cosets have
volume $\operatorname{vol}(ZR^n)=|\det Z|_K$, while
$\operatorname{vol}(R^n)=1$.
\end{proof}

Here is the precise counting mechanism for positivity.  Let $\mathcal W$
be a finite set.  For each $w\in\mathcal W$, choose a finite family
of marked finite $R$-modules $C_{w,q}$ and weights $a_{w,q}\in\mathbb C$.
Put
\begin{equation}\label{th-coker-kernel}
 \mathcal K(w,w')=
  \sum_{q,q'}a_{w,q}\overline{a_{w',q'}}
       \#\operatorname{Iso}(C_{w,q},C_{w',q'}).
\end{equation}
\begin{lemma}\label{th-coker-gram}
The matrix $\mathcal K$ is positive semidefinite.  More precisely, if
$C_\alpha$ runs through the isomorphism classes occurring, and
\[
 B_w(\alpha)=\sum_q a_{w,q}
                  \#\operatorname{Iso}(C_{w,q},C_\alpha),
\]
then
\begin{equation}\label{th-coker-gram-formula}
 \mathcal K(w,w')=
 \sum_\alpha
 \frac{B_w(\alpha)\overline{B_{w'}(\alpha)}}
      {\#\operatorname{Aut}(C_\alpha)}.
\end{equation}
\end{lemma}

\begin{proof}
If two modules are not isomorphic, both sides of the corresponding
identity for their isomorphism count are zero.  If they are isomorphic
to $C_\alpha$, each set of isomorphisms to $C_\alpha$ is a torsor
under the finite nonempty group $\operatorname{Aut}(C_\alpha)$.
Consequently
\[
 \#\operatorname{Iso}(C,D)=
 \sum_\alpha
 \frac{\#\operatorname{Iso}(C,C_\alpha)
       \#\operatorname{Iso}(D,C_\alpha)}
      {\#\operatorname{Aut}(C_\alpha)}.
\]
Substitution gives \eqref{th-coker-gram-formula}.  Its quadratic form
at coefficients $b_w$ is
$\sum_\alpha|\sum_w \overline{b_w} B_w(\alpha)|^2/
\#\operatorname{Aut}(C_\alpha)$, which is nonnegative.
\end{proof}

In the source, $w$ consists of positive block sizes $s_1,\ldots,s_l$
and scalars $t_1,\ldots,t_l\in K^\times$.  The matrix $Z(w)$ is block
diagonal with blocks $t_j$ times the reversing permutation matrix of
size $s_j$.  The index $q$ ranges over the integral matrices in its
specified lower- and upper-unipotent orbit, modulo integral
unipotent changes.  The weight is
$\psi(\sum x_{i+1,i}+\sum y_{i,i+1})$.
For every finite collection of $w$, the source proves that only
finitely many such cosets occur.  Once that finiteness and the
invariance under the integral changes have been established, its
local kernel is exactly \eqref{th-coker-kernel}.  Thus the positivity
step is Lemma~\ref{th-coker-gram}, with no appeal to a global
automorphic identity.

There is also an immediate support restriction.  All matrices in
the orbit of $Z(w)$ have determinant $\det Z(w)$, because the
unipotent factors have determinant one.  By Lemma~\ref{th-coker},
\begin{equation}\label{th-det-support}
 \mathcal K(w,w')\ne0\quad\Longrightarrow\quad
 |\det Z(w)|_K=|\det Z(w')|_K.
\end{equation}
The factor multiplying the kernel after adjoining a block of size
$s$ is
$|tt'|_K^{s(s-1)/2}$.  It is a product of a positive factor depending
on $t$ and one depending on $t'$, so preserves positive
semidefiniteness.  This establishes the positivity and determinant
support portions of that construction.  The stronger support
condition $t^{-1}-t'^{-1}\in R$ uses the additional averaging over
successive coordinate functionals; it does not follow merely from
\eqref{th-det-support}.

\subsection{Finite Fourier averaging and the local limit}

Two further steps of the local construction admit short proofs that
do not require an estimate for an oscillatory integral.  For Hermitian
matrices $M,N$, the notation $M\le N$ means that $N-M$ is positive
semidefinite.

\begin{lemma}\label{th-cut-contraction}
Let $G$ and $G_c$ be Hermitian matrices on a finite index set, with
$0\le G_c\le G$.  There are vectors $u_i$ in a finite-dimensional
Hilbert space and a linear contraction $T$ such that
\[
 G_{ij}=\langle u_i,u_j\rangle,\qquad
 (G_c)_{ij}=\langle Tu_i,Tu_j\rangle.
\]
\end{lemma}

\begin{proof}
Equip the vector space of coefficients with the positive semidefinite
form $a^*Gb$ and quotient by its null space.  The form $a^*G_cb$
descends to that quotient: a vector of zero $G$-norm has zero
$G_c$-norm, and the Cauchy--Schwarz inequality for $G_c$ makes its
pairing with every vector zero.  On the quotient the second form is
$\langle a,Bb\rangle$ for a self-adjoint operator $0\le B\le I$.
Take $T=B^{1/2}$ and let $u_i$ be the images of the coordinate vectors.
The spectral theorem gives $\|T\|\le1$ and both stated identities.
This also covers the zero-dimensional quotient.
\end{proof}

\begin{proposition}\label{th-fourier-average}
Let $A$ be a finite nonempty set, let $B$ be a finite abelian group,
and put $a=\#A$, $b=\#B$.  Let $I$ be a finite index set and
$u_{\rho,t,i}$ vectors in a Hilbert space $\mathcal E$, indexed by
$\rho\in A$, $t\in B$, and $i\in I$.  Suppose that
\[
 \langle u_{\rho,t,i},u_{\rho',t',j}\rangle=0
 \quad\text{if }\rho\ne\rho'.
\]
For each character $\chi\in\widehat B$ choose numbers
$\omega_\rho(\chi)$ of modulus one, and put
\[
 v_{\rho,\chi,i}=\frac1b\sum_{t\in B}
             \overline{\chi(t)}u_{\rho,t,i},\qquad
 v_{\chi,i}=\frac1a\sum_{\rho\in A}
             \omega_\rho(\chi)v_{\rho,\chi,i},
\]
\[
 \Theta_\chi(i,j)=\frac1a\sum_{\rho\in A}
          \langle v_{\rho,\chi,i},v_{\rho,\chi,j}\rangle.
\]
If $T:\mathcal E\to\mathcal E$ is a linear contraction and
$N_\chi(i,j)=a\langle Tv_{\chi,i},Tv_{\chi,j}\rangle$, then
\begin{equation}\label{th-fourier-domination}
 0\le N_\chi\le\Theta_\chi,
 \qquad
 \sum_{\chi\in\widehat B}\Theta_\chi(i,j)
 =\frac1{ab}\sum_{\rho\in A}\sum_{t\in B}
       \langle u_{\rho,t,i},u_{\rho,t,j}\rangle.
\end{equation}
Consequently, if the matrix on the right is at most a fixed matrix
$M$, then $\sum_\chi\Theta_\chi\le M$.
\end{proposition}

\begin{proof}
For coefficients $c_i$, write
$v_{\rho,\chi,c}=\sum_i c_i v_{\rho,\chi,i}$ and similarly
for $v_{\chi,c}$.  Orthogonality between distinct $\rho$ gives
\[
 \sum_{i,j}\overline{c_i}c_j\Theta_\chi(i,j)
 =\frac1a\sum_\rho\|v_{\rho,\chi,c}\|^2
 =a\|v_{\chi,c}\|^2
 \ge a\|Tv_{\chi,c}\|^2.
\]
This proves the first assertion, including positivity.  Character
orthogonality gives
\[
 \sum_{\chi\in\widehat B}\chi(t)\overline{\chi(t')}
 =b\,1_{t=t'}.
\]
This follows either by decomposing $B$ into cyclic
groups and summing geometric series, or by translating this
character sum.  Expanding the inner products and applying that
identity proves the second assertion.  The last assertion follows
directly by comparison of quadratic forms.
\end{proof}

Here is the exact role of this proposition in the local induction.
One chooses additive lattices
$L_1\subset R\subset L_2$ and representatives $A$ of $L_2/R$
and $B$ of $R/L_1$.  The latter quotient is the group in
Proposition~\ref{th-fourier-average}; representatives of $A$ are
only a finite indexing set.  Each shift $\sigma=\rho+t$ changes
the small parameter $\tau$ to $\tau_\sigma$ by
$\tau_\sigma^{-1}=\tau^{-1}+\sigma$.
The source's unrestricted orbital sum supplies a Gram matrix on the
indices $(\rho,t,i)$.  Restricting the orbital sum to a subset
supplies a smaller positive semidefinite matrix, so
Lemma~\ref{th-cut-contraction} supplies $T$.
The support condition in the preceding paragraph gives the required
orthogonality when $\rho\ne\rho'$.
Self-duality identifies $\widehat{R/L_1}$ with $L_1^\perp/R$.
The phases $\omega_\rho$ account for the representatives of
$L_2/R$.  The crucial remaining identity says that the matrix
$a\langle Tv_{\chi,i},Tv_{\chi,j}\rangle$ equals the desired
kernel after adjoining the larger block, for parameters in the
chosen compact annulus.  Once that identity and the bound for the
averaged unrestricted diagonal have been established,
\eqref{th-fourier-domination} proves both the positive majorant and
its total-mass bound.  Thus the averaging inequality itself adds
neither an unspecified constant nor a loss depending on $a$ or $b$.

The following compactness argument removes the annulus.  A function
on $K$ is called $R$-periodic if it is constant on every additive
coset of $R$.

\begin{proposition}\label{th-local-limit}
Let $L_c\subset K$ be a compact open additive subgroup containing
$R$, and let $I$ be a finite index set.  Let $M$ be a fixed positive
semidefinite matrix on $I$.  Suppose that $N(t)$ is a Hermitian
matrix for every $t\in K^\times$, with $N(t)=0$ outside $L_c$.
Let $\Omega_j$ be increasing subsets of $K^\times$ whose union is
$K^\times$.  Suppose that for every $j$ there is an $R$-periodic
matrix-valued function $\Theta_j$ such that
\[
 \Theta_j(x)\ge0,\qquad
 \Theta_j(x)=0\ (x\notin L_c),\qquad
 \sum_{x\in K/R}\Theta_j(x)\le M,
\]
and $\Theta_j(t)\ge N(t)$ for $t\in\Omega_j$.
Then there is an $R$-periodic function $\Theta$ with the same
positivity, support, and total-mass properties and
$\Theta(t)\ge N(t)$ for every $t\in K^\times$.
Any fixed collection of zero matrix entries shared by all
$\Theta_j(x)$ is shared by $\Theta(x)$.  If $L_c=R$, then
\begin{equation}\label{th-unit-head-bound}
 0\le\Theta(x)\le1_R(x)M\qquad(x\in K).
\end{equation}
\end{proposition}

\begin{proof}
The quotient $L_c/R$ is compact and discrete, hence finite.  Thus
the matrices $\Theta_j$ are described by a fixed finite number of
complex entries.  For every coset $x$ one has
$0\le\Theta_j(x)\le M$, because all other summands are positive
semidefinite.  The $2$ by $2$ principal minors therefore give
\[
 |\Theta_j(x)_{il}|^2
 \le\Theta_j(x)_{ii}\Theta_j(x)_{ll}
 \le M_{ii}M_{ll}.
\]
A subsequence converges in this finite-dimensional space.  Define
$\Theta$ by its limits on these cosets and set it to zero outside
$L_c$.  The cone of positive semidefinite matrices is closed:
for every vector $c$, the inequality $c^*Ac\ge0$ passes to an
entrywise limit.  This proves positivity and the total-mass bound,
since the latter is a finite sum.  For fixed $t\in K^\times$,
one has $t\in\Omega_j$ for all sufficiently large $j$, so the
same closedness argument gives $\Theta(t)\ge N(t)$.
Zero entries also pass to the limit.  Finally, if $L_c=R$, the
sum over $K/R$ has just one possibly nonzero term, proving
\eqref{th-unit-head-bound}.
\end{proof}

For the present product of local fields one may take $\Omega_j$ to
be the product of the annuli on which each component valuation lies
between $-j$ and $j$.  These sets are compact, open, exhaust the
units of $K$, and are stable under multiplication by any root of
unity.  If the source's additional support assertion
\[
 N(t)_{hh'}\ne0\quad\Longrightarrow\quad
       (\det Z(h))t\in R
\]
has been proved, a common permissible subgroup is the compact open
$R$-module
$L_c=R+\sum_h(\det Z(h))^{-1}R$, where $h$ ranges over the fixed
finite family of previous block configurations and $Z(h)$ is their
block reversal matrix defined above.  If all $Z(h)$ and their inverses have
entries in $R$, then each determinant is a unit of $R$, so
$L_c=R$.  In this case the upper bound in
\eqref{th-unit-head-bound} is precisely the final extra condition
in the local majorant statement, after substituting its matrix
$M$.  No rate of convergence, uniform stationary-phase radius,
or infinite-dimensional compactness theorem is needed for this
last passage.  The subsequent local reconstruction proves that
the truncated orbital Gram matrix equals the integral of the
larger-block kernel.  It establishes the actual orbital germ
formula and the cancellation between distinct stationary points.

\subsection{Reproducing constants and convolution on a quotient}

The possible scalar in the local reproducing formula can be determined
without evaluating a quotient integral point by point.  We state the
needed functional-analytic argument explicitly.

Let $G$ be a second countable unimodular locally compact group with a
fixed Haar measure.  Let $\pi$ be a unitary representation on a Hilbert
space $V$, and suppose $d_\pi>0$ satisfies the formal-degree identity
\begin{equation}\label{th-schur}
 \int_G\langle a,\pi(g)b\rangle
       \overline{\langle a',\pi(g)b'\rangle}\,dg
   =d_\pi^{-1}\langle a,a'\rangle\langle b',b\rangle
\end{equation}
for all vectors.  The existence of this identity for an irreducible
discrete series is a representation-theoretic input; it is not being
inferred from unitarity alone.  Choose $w\in V$ of norm one and put
\[
 p(g)=d_\pi\overline{\langle w,\pi(g)w\rangle}.
\]
In addition to \eqref{th-schur}, assume $p\in L^1(G)$.

\begin{lemma}\label{th-projection}
For convolution $(p*q)(x)=\int_Gp(g)q(g^{-1}x)\,dg$ one has
$p*p=p$ and $p^*(g)=\overline{p(g^{-1})}=p(g)$.
In any strongly continuous unitary representation $R$ of $G$, the
bounded operator
\[
 R(p)v=\int_Gp(g)R(g)v\,dg
\]
is an orthogonal projection.  In particular, if $v\ne0$ and
$R(p)v=cv$ with $c\ne0$, then $c=1$.
\end{lemma}

\begin{proof}
The coefficient belongs to $L^2$ by \eqref{th-schur}, so the
convolution integrand is absolutely integrable by Cauchy--Schwarz.
Unitarity and \eqref{th-schur} give
\begin{align*}
 (p*p)(x)
 &=d_\pi^2\int_G
    \overline{\langle w,\pi(g)w\rangle}
    \langle\pi(x)w,\pi(g)w\rangle\,dg\\
 &=d_\pi\langle\pi(x)w,w\rangle=p(x).
\end{align*}
The identity $p^*=p$ follows directly from unitarity.
The integral defining $R(p)v$ exists in norm and
$\|R(p)\|\le\|p\|_1$.  Fubini is justified by the majorant
$|p(g)q(h)|\|v\|$, and gives $R(p)R(q)=R(p*q)$.
Invariance under inversion gives $R(p)^*=R(p^*)$.
Thus $R(p)$ is self-adjoint and idempotent, hence an orthogonal
projection.  Applying idempotence to $v$ gives $c^2v=cv$;
$c\ne0$ and $v\ne0$ imply $c=1$.
\end{proof}

For application to the source, let $S=N\rtimes H$ be the local
stabilizer, with $N$ unipotent and $H$ compact.  Its action on $N$
preserves Haar measure: its modular factor is a continuous
homomorphism from the compact group $H$ to $\mathbb R_{>0}$ and
therefore is one.  Thus $S$ is unimodular.  Let $\rho$ be its strongly
continuous unitary representation, formed from the specified character of $N$ and
the compact representation of $H$.  The induced Hilbert space is
\[
 \mathcal I=\left\{F:G\to V_\rho:
 F(sg)=\rho(s)F(g),\quad
 \int_{S\backslash G}\|F(g)\|^2\,d\bar g<\infty\right\},
\]
where the functions are measurable and functions equal almost
everywhere are identified.  Right
translation $R(g)F(x)=F(xg)$ is unitary, by the quotient integration
formula.  Lemma~\ref{th-projection} applies to this representation.
It proves that the nonzero scalar in a reproducing formula for a
nonzero $F\in\mathcal I$ must be one.

The nonvanishing, membership in $\mathcal I$, and the $L^1$ condition
are substantive hypotheses.  The last does not follow from square
integrability of discrete-series coefficients.  In the source it is
obtained from a weighted coefficient estimate under a large
Harish--Chandra parameter gap.  The argument above replaces the
subsequent scalar-normalization argument, not that weighted estimate.

\subsection{Unfolding the Gram matrix}

Let $\Gamma\subset G$ be a discrete subgroup, and let
$\Gamma_u\subset\Gamma$ be subgroups.  Let $C_u$ be finite-dimensional
Hilbert spaces and let $\phi_{u,v}(g)$ depend linearly on $v\in C_u$,
with $\phi_{u,v}$ left $\Gamma_u$-invariant.  Define
\[
 P_uv(g)=\sum_{\delta\in\Gamma_u\backslash\Gamma}
                         \phi_{u,v}(\delta g).
\]
Assume the series converge absolutely almost everywhere, give
$L^2(\Gamma\backslash G)$ functions, and that for every pair of
vectors from chosen bases,
\begin{equation}\label{th-unfold-integrable}
 \int_{\Gamma_u\backslash G}|\phi_{u,v}(g)|
     \sum_{\delta\in\Gamma_{u'}\backslash\Gamma}
                  |\phi_{u',w}(\delta g)|\,d\bar g<\infty.
\end{equation}

\begin{lemma}\label{th-unfold}
Choose representatives $\eta$ for
$\Gamma_{u'}\backslash\Gamma/\Gamma_u$, and put
$J_\eta=\Gamma_u\cap\eta^{-1}\Gamma_{u'}\eta$.
With compatible quotient measures,
\begin{equation}\label{th-unfold-formula}
 \langle P_uv,P_{u'}w\rangle
 =\sum_\eta\int_{J_\eta\backslash G}
       \overline{\phi_{u,v}(g)}\phi_{u',w}(\eta g)\,d\bar g.
\end{equation}
The right side is absolutely convergent.
\end{lemma}

\begin{proof}
Expand the first Poincar\'e series and integrate on a measurable
fundamental domain for $\Gamma$.  Translating that domain by the
representatives of $\Gamma_u\backslash\Gamma$ tiles a fundamental
domain for $\Gamma_u$.  Since $P_{u'}w$ is left $\Gamma$-invariant,
the result is
$\int_{\Gamma_u\backslash G}
\overline{\phi_{u,v}(g)}P_{u'}w(g)\,d\bar g$.
For each double coset $\Gamma_{u'}\eta\Gamma_u$, its left cosets
are $\Gamma_{u'}\eta\alpha$, with
$\alpha\in J_\eta\backslash\Gamma_u$.
Expanding this part of the second series and unfolding once more
gives the integral over $J_\eta\backslash G$ in
\eqref{th-unfold-formula}.  Tonelli applied to the absolute values,
using \eqref{th-unfold-integrable}, justifies both rearrangements
and proves the asserted convergence.  Finite-dimensionality extends
the basis assumption to all vectors.
\end{proof}

By \eqref{th-normalized-K}, division of each term by $\sqrt{c_uc_{u'}}$
gives the corresponding expansion of $K_{uu'}$.  This explains which
kernel is being bounded: it is the Gram operator of the actual
Poincar\'e maps, rather than an unrelated matrix with a similar name.
Further decomposition of the double cosets into flags, lattices,
and local orbits, and the estimates for their absolute sums, remain
arithmetic work.  The hypotheses in Lemma~\ref{th-unfold} must be
proved for the proposed seeds; they are not consequences of the
formal definition of a Poincar\'e series.

\subsection{From the orbital estimates to the pairwise bound}

The last combination of the orbital estimates admits a short proof.
We give it because it specifies exactly what remains to be estimated.
Let $a_t$ be an absolutely summable family in a Hilbert space
$\mathcal A$, and let $v_t$ be an absolutely summable family in a
Hilbert space $\mathcal V$.  Let $r(t)$ be a row label, and let
$i(r)$ belong to a finite set $I$.  Suppose that for each of two
disjoint sets of terms $S=\mathcal G,\mathcal N$,
\begin{equation}\label{th-row-norm}
 \left\|\sum_{t\in S}a_t\right\|^2
  =\operatorname{Re}\sum_{i,j\in I}\Gamma_{ij}
       \left\langle\sum_{t\in S:i(r(t))=i}v_t,
                       \sum_{t\in S:i(r(t))=j}v_t\right\rangle.
\end{equation}
Assume $P\gamma I-\Gamma$ is positive semidefinite, with
$P\ge0$, $\gamma>0$, and assume
\begin{equation}\label{th-row-orth}
 \langle v_t,v_{t'}\rangle=0
 \quad\text{if }i(r(t))=i(r(t'))\text{ and }r(t)\ne r(t').
\end{equation}
In the application, $\mathcal A$ is a space of Hilbert--Schmidt
operators.  The sets $\mathcal G$ and $\mathcal N$ consist of the
generic and nongeneric terms with nonzero Gram determinant.

For equal-row pairs of nonzero-determinant terms, suppose an integer
$s(t,t')\ge1$ and a scalar/non-scalar distinction at $s=1$ are given.
For a set $\mathcal Q$ of ordered pairs define
\[
 E(\mathcal Q)=\gamma\sum_{(t,t')\in\mathcal Q}
                    \operatorname{Re}\langle v_t,v_{t'}\rangle,
 \qquad
 E_{\mathrm{abs}}(\mathcal Q)=\gamma\sum_{(t,t')\in\mathcal Q}
                    |\langle v_t,v_{t'}\rangle|.
\]
All these sums converge absolutely.  Put $\mathcal B$ for the
equal-row pairs with $s\ge2$, $\mathcal C$ for the generic
equal-row pairs with $s=1$ that are not scalar, $\mathcal M$ for
the scalar generic pairs with $s=1$, and $\mathcal U$ for the
nongeneric pairs with $s=1$.

\begin{proposition}\label{th-pair-bound}
Let $d\ge0$ and $A,\epsilon,\epsilon'\ge0$.  Suppose
\begin{equation}\label{th-four-estimates}
 E_{\mathrm{abs}}(\mathcal B)\le\epsilon d,\quad
 E_{\mathrm{abs}}(\mathcal C)\le\epsilon d,\quad
 E(\mathcal M)\le A d,\quad E(\mathcal U)\le\epsilon'd.
\end{equation}
Then
\begin{equation}\label{th-pair-output}
 \left\|\sum_{t\in\mathcal G}a_t+
             \sum_{t\in\mathcal N}a_t\right\|^2
 \le P\left(\sqrt{A+2\epsilon}
                  +\sqrt{\epsilon+\epsilon'}\right)^2d.
\end{equation}
\end{proposition}

\begin{proof}
The positive semidefinite inequality and
\eqref{th-row-norm} bound the squared norm for a set $S$ by
\[
 P\gamma\sum_{i\in I}
       \left\|\sum_{t\in S:i(r(t))=i}v_t\right\|^2.
\]
Expand the squares.  Equation \eqref{th-row-orth} removes pairs
with different row labels.  For $S=\mathcal G$, the remaining
pairs partition into $s\ge2$, scalar $s=1$, and nonscalar $s=1$;
for $S=\mathcal N$, they partition into $s\ge2$ and $s=1$.
Using \eqref{th-four-estimates} gives bounds
$P(A+2\epsilon)d$ and $P(\epsilon+\epsilon')d$ respectively.
The triangle inequality gives \eqref{th-pair-output}.  Absolute
summability justifies every expansion and restriction of the pair sums.
\end{proof}

\begin{lemma}\label{th-small-errors}
If $b\ge16$ and
\[
 0\le P\le e^{6b^{-7}},\quad
 0\le A\le(1+b^{-7})e^{b^{-14}},\quad
 0\le\epsilon,\epsilon'\le b^{-14},
\]
then the multiplier in \eqref{th-pair-output} is at most
\begin{equation}\label{th-sharper-error}
 1+14b^{-7}\le1+b^{-6}.
\end{equation}
\end{lemma}

\begin{proof}
Write $u=b^{-7}\le1/1000$.  The elementary bound
$e^x\le1+x+2x^2$ for $0\le x\le1/2$ follows from
$e^x(1-x)\le1$ and $1/(1-x)\le1+x+2x^2$.
It gives $P\le1+7u$ and $e^{u^2}\le1+2u^2$.
Consequently
\[
 A+2\epsilon\le1+u+5u^2\le1+2u,
 \qquad P(A+2\epsilon)\le1+10u.
\]
Also $P(\epsilon+\epsilon')\le\tfrac94u^2$.
The sum of the two square roots is therefore at most
$1+5u+\tfrac32u=1+\tfrac{13}{2}u$.  Squaring and using
$u\le1/1000$ bounds it by $1+14u$.
Finally $14u\le bu=b^{-6}$.
\end{proof}

The slightly stronger bound $14b^{-7}$ is available without changing
the orbital estimates.  It improves the last numerical combination;
it does not eliminate any of the four estimates in
\eqref{th-four-estimates}.

Here is a second useful simplification of the numerical part.  Suppose
local quantities $q_{\mathfrak p}$ satisfy
$q_{\mathfrak p}\ge p$ at a place above the rational prime $p$, and
there are at most $f$ such places above each $p$.  For $a>1$,
\begin{equation}\label{th-density-bound}
 \prod_{\mathfrak p}\exp(4q_{\mathfrak p}^{-a})
 \le\exp\left(4f\,2^{-a}\left(1+\frac2{a-1}\right)\right).
\end{equation}
Indeed
$\sum_{\mathfrak p}q_{\mathfrak p}^{-a}
\le f\sum_{n\ge2}n^{-a}
\le f(2^{-a}+\int_2^\infty t^{-a}\,dt)$.
This also proves convergence, and applies to every finite subproduct.
The source uses $a=k-2r-2\ge2$ and the weaker upper bound
$\exp(12f2^{-a})$.  Thus its local-density accumulation does not
conceal a constant depending on the number of places.

Likewise, for $q\ge3$,
\[
 \prod_{j=1}^k(1-q^{-j})^{-1}
 \le\exp\left(\frac{q}{(q-1)^2}\right)
 \le\exp\left(\frac9{4q}\right).
\]
Use $-\log(1-z)\le z/(1-z)$ and sum the geometric series.
If two split primes have sizes $g_1,g_2\ge(kb)^8$ and each has $f$
places, the product of these bounds is at most
$\exp(9f/(2(kb)^8))$.  This is smaller than the coarser
$\exp(6b^{-7})$ used in Lemma~\ref{th-small-errors} when
$f\le k$, $k\ge1$, and $b\ge1$.
For clarity, the actual prefactor $P$ in the source's row estimate
is the different product
\[
 P=\prod_{v\mid g_1g_2}(1+3/q_v)
   =\bigl((1+3/g_1)(1+3/g_2)\bigr)^f.
\]
The equality uses complete splitting at these two rational primes.
It satisfies the required bound directly:
\[
 P\le\exp(3f/g_1+3f/g_2)
   \le\exp\bigl(6f/(kb)^8\bigr)
   \le\exp(6b^{-7}).
\]
Thus the separate Euler-product estimate above is not needed to
identify this prefactor.

\subsection{The remaining global statement, with its scale visible}

We first specify all parameters in the following conditional numerical
combination.  The integers $k,f\ge1$ and $r\ge0$ have the following
meaning in the source.  The CM extension is $E/F$, with $F$ totally
real of degree $f$ over $\mathbb Q$.  The relevant Hermitian space has
dimension $k+2r+1$ over $E$ and signature $(k+r,r+1)$ at each real
place of $F$.  The integer $r$ is the length of the isotropic flag,
and the compact stabilizer has dimension $k$ in its defining
representation.  The number of successive lifts is $2r+2$.  The
positive real number $b$ is the height scale for the parameters of
the representations; in the growing family it is $\exp(x^\varepsilon)$,
where $x>1$ tends to infinity and $\varepsilon>0$ is fixed.

The real numbers $c_5>0$ and $C_n\ge0$ are respectively the decay
constant in the large-corank and nonscalar bounds and the growth
constant in the nongeneric bound.  The concrete source takes
$c_5=1/3$ and chooses $C_n$ after the fixed auxiliary constants,
independently of the subsequent growing parameters.  The positive
real number $C_\Gamma$ is chosen after the constants in the diagonal
estimates, and
\[
 \Gamma=\exp\bigl(C_\Gamma(2r+2)^5f\log b\bigr).
\]
The real number $C_0\ge1$ measures the assumed comparison between
the dimensions $d_u$ and $D_*/\Gamma^2$: the source requires
$C_0^{-1}\le d_u\Gamma^2/D_*\le C_0$ for every $u$.
Only its upper bound enters the diagonal estimate below.

Here $D_*>0$ is the proposed main term for the dimension of the
ambient automorphic space.  More precisely, write $G$ for the unitary
group just specified, $d_\infty>0$ for the product of the
archimedean formal degrees with the chosen Haar measures, and
$e_{\mathrm{fin}}$ for the finite-adelic convolution idempotent
specifying the finite-level condition.  Its normalization is
\[
 D_*=d_\infty\,
 \operatorname{vol}\bigl(G(F)\backslash G(\mathbb A_F)\bigr)
 \operatorname{Re} e_{\mathrm{fin}}(1).
\]
Both the volume and the idempotent use the measures fixed in the
automorphic construction.  This definition does not itself assert
an asymptotic formula for the actual dimension.  In the numerical
argument that follows, $k,f,r,b,c_5,C_n,C_\Gamma,C_0,D_*$ are explicit
input numbers with the domains just stated and the displayed
inequalities; their realization by the particular global family
remains a separate assertion.

For distinct indices $u,u'$, put
$d=d_ud_{u'}/D_*$, where $d_u=\dim C_u$ and $D_*>0$.
The actual four estimates in the source have the following scales:
\begin{align}\label{th-global-scales}
 \epsilon&=\exp(-c_5kf\log k),\notag\\
 \epsilon'&=2^{-(k-2r-5)}
       \exp\bigl(C_n(2r+2)^5f\log b\bigr),\notag\\
 A&=(1+\delta)e_1,\qquad
 0\le\delta\le b^{-7},\qquad
 1\le e_1\le\exp(12f2^{-(k-2r-2)}).
\end{align}
The dimensional restrictions include
\begin{equation}\label{th-global-parameters}
 \begin{gathered}
 b\ge16,\quad k\ge32,\quad f\le k,\quad 4r+12\le k,
 \quad70\log b\le k,\\
 5c_5f\log k\ge1,\qquad
 8C_n(2r+2)^5f\log b\le k.
 \end{gathered}
\end{equation}
They imply $\epsilon,\epsilon'\le e^{-k/5}\le b^{-14}$ and
$12f2^{-(k-2r-2)}\le e^{-k/5}$.  To verify the latter bound, use
$k-2r-2\ge k/2+4$, so that the first expression is at most
$f2^{-k/2}\le k2^{-k/2}\le e^{-k/5}$ for $k\ge32$.
The last elementary inequality follows because
$\log k\le k((\log2)/2-1/5)$ at $k=32$ and the right side has
larger derivative thereafter.  For $\epsilon'$, use
$k-2r-5\ge k/2$ and
$C_n(2r+2)^5f\log b\le k/8$; then
$-(k/2)\log2+k/8\le-k/5$.

Assuming the actual orbital decomposition satisfies
\eqref{th-row-norm}--\eqref{th-four-estimates}, and assuming the
degenerate and identity terms vanish for $u\ne u'$, the already
proved proposition and lemma give
\begin{equation}\label{th-off-final}
 \|K_{uu'}\|_{\mathrm{HS}}^2
 \le(1+14b^{-7})\frac{d_ud_{u'}}{D_*}.
\end{equation}
The source states the weaker $(1+b^{-6})$ version.

On the diagonal write $K_{uu}=I+Z_u+R_u$, where $Z_u$ is the sum
of the nonidentity terms with zero Gram determinant.  The separate
estimates required there are
\begin{equation}\label{th-diag-input}
 \|Z_u\|_{\mathrm{HS}}^2\le d_u e^{-kf\log k},\qquad
 \|R_u\|_{\mathrm{HS}}^2\le\Gamma\frac{d_u^2}{D_*},
 \qquad \frac{d_u\Gamma^2}{D_*}\le C_0.
\end{equation}
They imply the slightly sharper diagonal defect
\begin{equation}\label{th-diag-final}
 \|K_{uu}-I\|_{\mathrm{HS}}^2
 \le d_u\left(e^{-kf\log k/2}+\sqrt{C_0/\Gamma}\right)^2.
\end{equation}
This is just the triangle inequality together with the last
inequality in \eqref{th-diag-input}.  All constants and their roles
are explicit at this stage.

The Hilbert--Schmidt norm of a matrix representing a linear map in
orthonormal bases is the operator's Hilbert--Schmidt norm.
Consequently, if $M_u$ is the orthonormal matrix of $P_u$ and
$A_u=c_u^{-1/2}M_u$, then
\[
 A_u^*A_{u'}=K_{uu'}.
\]
Equations \eqref{th-off-final} and \eqref{th-diag-final} therefore
give exactly the pairwise and diagonal estimates needed by the
previous slope comparison.  Compression by an isometric inclusion
of the ramified subspace is handled by the compression lemma proved
there.  This last identification requires no additional estimate.

\paragraph{Local statements established in the following sections.}
The local reconstruction proves the
finiteness and representative independence of the actual local
sums, their local constancy uniformly against all comparison
configurations, and the strong inverse-parameter support.  It also
proves the root-fiber change of variables and the cancellation of
terms from distinct stationary points.  The fixed-initial-block
construction and the Hessenberg change of variables establish the
actual kernel germ in Theorem~\ref{ga-actual-germ}, hence the exact
truncated orbital identity.  The finite character identity and its
averaged diagonal bound then complete the positive-majorant
induction in Theorem~\ref{tp-majorant}.
Propositions~\ref{th-fourier-average} and \ref{th-local-limit}
supply its averaging and limiting steps with the asserted constants.

\paragraph{What has not yet been proved in this reconstruction.}
The following are concrete remaining assertions about the proposed
automorphic family, rather than consequences of the local lemmas.
First, the actual seeds must have the weighted integrability and
nonvanishing properties used above; the source obtains these from
discrete-series and Bessel-model theorems with large parameter gaps.
Second, the global recursion after Poisson summation must
bound the scalar generic pair sum by $(1+\delta)e_1d$ with
$\delta\le b^{-7}$.  Its arithmetic inputs include the lattice
indices and the sums over the dual lattices, with constants uniform
in the stated growing parameters.  Third, the larger-corank,
nonscalar, and nongeneric sums in \eqref{th-global-scales}, the
degenerate diagonal bound in \eqref{th-diag-input}, and the required
ambient-dimension estimate must be established for these particular
maps.  Positivity and stationary phase alone imply none of those
global estimates.

No counterexample satisfying all the actual hypotheses of the local
results proved here has been found.  This is not an assertion of an
exhaustive audit of the theta energy recursion.  The source trace
accompanying this section distinguishes the statements and proofs
read in full from the still unresolved arithmetic steps.

\subsection{Exact small controls and excluded shortcuts}

The accompanying program uses only integer arithmetic and reduction
modulo cyclotomic polynomials.  For $a=p^e$, $z=1$, and
$L=\mathbb Z_p^m$, it checks the finite sums
\[
 G_{p,m,e}=\sum_{x\in(\mathbb Z/p^e\mathbb Z)^m}
 \exp\left(\frac{2\pi i}{p^e}
       \left(\sum_i x_i^2+\sum_{i<j}x_ix_j\right)\right).
\]
The Gaussian lemma predicts
$|G_{p,m,e}|^2=p^{em+v_p(m+1)}$ when the admissibility condition
holds.  The checked cases are
$(p,m,e)=(2,1,2),(2,3,3),(3,2,2),(3,3,1),(5,2,1),(5,4,2)$.
The checks include residue characteristics dividing both $2$ and
$m+1$.  It also verifies the full stationary phase identity for
$m=1$, $D=1$, and amplitude $1_{\mathbb Z_p^\times}$ at
$(p,e)=(3,2),(3,3),(5,2),(2,8)$, using critical balls of radius
$p^{-1}$ for odd $p$ and $2^{-3}$ for $p=2$.

Two exact failures explain hypotheses that must be retained.
For $(p,m,e)=(2,3,1)$ one has $G_{2,3,1}=-4$, so the normalized
squared integral is $1/4$, whereas the determinant expression
without admissibility would be $1/2$.  The dual quadratic form is
not trivial on $L^\perp$ in this example.  This is a counterexample
to dropping \eqref{th-gauss-admissible}, not to
Lemma~\ref{th-gaussian} or to the source's corresponding statement.
For the split order \eqref{th-split-order}, with $F=\mathbb Q_p$
and odd $p$, the integral of $\psi_R(x^2,y^2)$ over
$\mathbb Z_p^2$ is $1/p$ for product Haar measure and $1$ for the
correct self-dual measure.  Thus the order's factor $p$ cannot be
omitted.  Both values are checked exactly for $p=3,5$.

Finally, a degenerate Hessian really is excluded by characteristic
zero, not by a condition on residue characteristic.  Over a field
of characteristic $p$, the same phase at $m=p-1$ has Hessian
determinant zero at its critical point.  The Gaussian argument then
does not apply.  The actual source assumes that every positive
integer is invertible in the local algebra, so this example does
not satisfy its hypotheses.



\section{A local change of variables and cancellation of mixed critical terms}
\label{sec:theta-orbital-identity}

We continue with the same local data as in the preceding section:
$K$ is a finite product of finite extensions of $\mathbb Q_p$,
$R\subset K$ is a compact open subring, and $\psi$ is a continuous
additive character whose annihilator of $R$ is $R$ itself.  Haar
measure is normalized by $\operatorname{vol}(R)=1$.  This includes
the nonmaximal split order already described.  Absolute values and
small neighborhoods in $K$ are understood componentwise.  In
particular, a parameter tends to zero in every field factor.

Fix an integer $m\ge1$, put $s=m+1$, and let $\Omega\subset K^\times$
be compact and open.  Suppose that $\Omega$ is invariant under
multiplication by every $s$th root of unity in $K$.
For $t\in K^\times$, define
\begin{equation}\label{orb-cut}
 \mathcal C_t=\{(-t)^{-m}z^s:z\in\Omega\}.
\end{equation}
Fix an additive lattice $\Lambda$ such that
$\Omega+\Lambda=\Omega$.  A function on $\Omega$ is called
$\Lambda$-invariant when $f(z+\lambda)=f(z)$ for
$z\in\Omega$ and $\lambda\in\Lambda$.
We may always replace $\Lambda$ by a smaller lattice.  None of the
neighborhoods below depends on the values of $f$ or $g$.

\subsection{Uniform choices near one}

We record the elementary local facts used to keep the dependence on
residue characteristic and conductor explicit.

\begin{lemma}\label{orb-near-one}
For each positive integer $s$, the power series
$\beta_s(X)=(1+X)^{1/s}$ converges on a sufficiently small
neighborhood of zero in $K$.  On a smaller neighborhood it satisfies
\[
 \beta_s(X)^s=1+X,\qquad
 \beta_s(X)=1+\frac{X}{s}b_s(X),\qquad b_s(X)\in R,
 \quad b_s(0)=1.
\]
It is the unique $s$th root sufficiently close to one.
Given $\Omega$ and $\Lambda$ as above, there is a neighborhood
$U$ of one in $K^\times$, invariant under inversion, such that
\[
 u\Omega=\Omega,\qquad (u-1)\Omega\subset\Lambda,
 \qquad |u_j|=1
\]
for every $u\in U$ and every field component $j$.
\end{lemma}

\begin{proof}
In a field factor the coefficients of the binomial series have at
most exponential growth.  Indeed, for a fixed rational number $a$,
\[
 \left|\binom a n\right|
 \le\max(1,|a|)^n|n!|^{-1},
\]
and the valuation formula for $n!$ gives an exponential upper bound
for its inverse absolute value.  The series therefore converges on
a sufficiently small ball.  Its formal identity
$\beta_s(X)^s=1+X$ is an identity of convergent series there.
The function $b_s(X)=s(\beta_s(X)-1)/X$, extended by $b_s(0)=1$,
is continuous.  Since $R$ is a neighborhood of one, it takes
values in $R$ after another shrink.  Two nearby roots differ by an
$s$th root of unity.  There are finitely many such roots, and one
can choose the neighborhood to exclude all but one.  Applying these
facts in the finitely many field factors proves the assertion for
$K$.

For the second assertion, compactness of $\Omega$ and continuity
of multiplication give $(u-1)\Omega\subset\Lambda$ for every
$u$ in a sufficiently small neighborhood of one.  Then
$u\Omega\subset\Omega$.  Intersect this neighborhood with its
inverse to obtain equality.  The component absolute values of a
unit near one are all one.
\end{proof}

The neighborhood in this lemma is allowed to depend on $s$, on the
fields, and on $R$.  In particular, it is not asserted to be
independent of the valuations of $s$ or of $2$.  The lemma remains
valid when the residue characteristic divides either integer.

Let $\Sigma\subset K$ be finite.  For $t$ sufficiently small and
$\sigma\in\Sigma$, set
\begin{equation}\label{orb-shift}
 t_\sigma=(t^{-1}+\sigma)^{-1}=\frac{t}{1+\sigma t}.
\end{equation}
The denominator is then a unit, $t_\sigma/t$ tends uniformly to one,
and $|t_\sigma|=|t|$.  Taking the near-one $s$th root of
$(t_\sigma/t)^m$ and using Lemma~\ref{orb-near-one} gives
\begin{equation}\label{orb-cut-shift}
 \mathcal C_{t_\sigma}=\mathcal C_t
 \qquad(\sigma\in\Sigma).
\end{equation}
The same argument shows that $\mathcal C_t+R=\mathcal C_t$
when $t$ is small: if $c=(-t)^{-m}z^s$ and $r\in R$, then
$r/c$ tends uniformly to zero, and
$z\beta_s(r/c)\in\Omega$ is an $s$th root of $(-t)^m(c+r)$.
Compactness of $R$ and of the set of inverses of elements of
$\Omega$ makes this uniform in $c$ and $r$.

\subsection{An exact formula for the sum over root fibers}

Choose any representatives $c_x\in K^\times$ of the additive
cosets $x\in K/R$.  Such representatives exist: the compact open
ring contains a product of sufficiently small open balls, so each
coordinate of a representative can be perturbed away from zero.
The following identity is independent of these choices.

\begin{theorem}\label{orb-change-variable}
For all sufficiently small $t\in K^\times$ and every
$\Lambda$-invariant function $f:\Omega\to\mathbb C$,
\begin{equation}\label{orb-change-formula}
 \sum_{\substack{x\in K/R\cr c_x\in\mathcal C_t}}
       \ \sum_{z^s=(-t)^m c_x} f(z)
 =\frac{|s|}{|t|^m}\int_\Omega |z|^m f(z)\,dz.
\end{equation}
All roots in the inner sum belong to $\Omega$, and both sums are
finite.  The required small neighborhood depends only on
$K,R,m,\Omega,\Lambda$.
\end{theorem}

\begin{proof}
If one root of $z^s=(-t)^mc$ lies in $\Omega$, every root does,
because their ratios are $s$th roots of unity.  Put
\[
 H(c)=\sum_{\substack{z\in\Omega\cr z^s=(-t)^mc}}f(z).
\]
This function has compact support $\mathcal C_t$.  For
$c\in\mathcal C_t$ and $r\in R$, multiplication by
$\omega=\beta_s(r/c)$ bijects the roots above $c$ and $c+r$.
For small $t$, uniformly over these choices,
$z(\omega-1)\in\Lambda$.  Thus it preserves $f$, and
$H(c+r)=H(c)$.  The invariance of $\mathcal C_t$ under $R$
proves the same equality outside $\mathcal C_t$.  Consequently
\[
 \sum_{x\in K/R}H(c_x)=\int_K H(c)\,dc,
\]
since each additive coset has volume one and only finitely many
cosets meet the compact support.

The map $z\mapsto c=z^s/(-t)^m$ has nonzero derivative
$s z^m/(-t)^m$ at every point of $\Omega$.  On sufficiently small
product balls it is an analytic bijection onto a product of balls,
and multiplies Haar measure by $|s||z|^m/|t|^m$.  To justify this
without a smoothness assumption on $f$, divide the local map by
its derivative: on a smaller ball it is a linear map followed by
the identity plus a strict contraction, exactly as in the proof
of the local Morse lemma.  It therefore maps each sufficiently
small coset partition bijectively to the corresponding partition.
The Jacobian factor is constant on that ball.  A finite disjoint
partition of $\Omega$ by such balls proves the change-of-variables
formula, with every root counted once.  This gives
\eqref{orb-change-formula}.
\end{proof}

The source writes the integral on the right as a finite lattice
sum and additionally assumes that $|z|$ is constant on the
$\Lambda$-cosets in $\Omega$.  Under that assumption its sum is
exactly the Haar integral in \eqref{orb-change-formula}.  The
Haar-integral version itself needs no such extra assumption.

\subsection{Vanishing of the mixed critical terms}

\begin{theorem}\label{orb-cross-vanish}
Fix $\zeta\in K^\times$ with $\zeta^s=1$ and $\zeta\ne1$,
and fix a finite set $\Sigma\subset K$.  There are neighborhoods
$E$ of one and $V$ of zero such that the following holds.
Let $t\in V\cap K^\times$, let
$t_1=t_\sigma$, $t_2=t_{\sigma'}$ with
$\sigma,\sigma'\in\Sigma$, and let $\kappa\in E$.
Suppose that $g:\Omega\to\mathbb C$ is $\Lambda$-invariant
and satisfies
\begin{equation}\label{orb-cross-support}
 g(z)\ne0\quad\Longrightarrow\quad
 z^{-1}-(\zeta z)^{-1}\in R.
\end{equation}
Then, for every choice of the representatives $c_x$,
\begin{equation}\label{orb-cross-sum}
 \sum_{\substack{x\in K/R\cr c_x\in\mathcal C_t}}
 \ \sum_{z^s=(-t_1)^m c_x}
 g(z)\psi\left(sz/t_1-s\zeta\kappa z/t_2\right)=0.
\end{equation}
In particular this holds with the source's further condition
$\kappa^s=(t_2/t_1)^m$; that condition is unnecessary for this
vanishing assertion alone.
\end{theorem}

\begin{proof}
Write
\[
 a=\frac{s}{t_1}-\frac{s\zeta\kappa}{t_2},\qquad
 \rho=\kappa\frac{t_1}{t_2},\qquad
 T=(-t_1)^m.
\]
We may make $\rho$ as close to one as needed by shrinking $E$
and $V$.  The truncation sets for $t_1$ and $t$ agree by
\eqref{orb-cut-shift}.  Define
\[
 H(c)=\sum_{\substack{z\in\Omega\cr z^s=Tc}}
             g(z)\psi(az).
\]
The important point is that $H$ is still $R$-periodic, despite
the oscillating factor.  For $r\in R$, multiplication by
$\omega=\beta_s(r/c)$ matches the roots above $c$ and $c+r$,
preserves $\Omega$, and preserves $g$.  Moreover
\[
 az(\omega-1)=\frac{aT}{s z^m}\,r\,b_s(r/c),\qquad
 \frac{aT}{s z^m}
    =(-1)^m\frac{t_1^{m-1}}{z^m}(1-\zeta\rho).
\]
The second expression belongs to $R$ on the support of $g$,
uniformly for sufficiently small parameters.  If $m\ge2$, it
tends uniformly to zero, since $\Omega$ is compact in $K^\times$.
If $m=1$, then $s=2$ and $\zeta^{-1}=\zeta$, so
\[
 -\frac{1-\zeta\rho}{z}
 =-\frac{1-\zeta}{z}+\frac{\zeta(\rho-1)}{z}\in R.
\]
The first term belongs to $R$ by \eqref{orb-cross-support};
the second does so by the uniform smallness of $\rho-1$.
Since $r,b_s(r/c)\in R$ and $\psi$ is trivial on $R$, the
phase is unchanged.  Terms where $g(z)=0$ cause no difficulty,
because the matched value of $g$ is also zero.  Therefore
$H(c+r)=H(c)$, including outside its $R$-stable support.

As in the preceding theorem, the sum in \eqref{orb-cross-sum}
is consequently
\begin{equation}\label{orb-cross-integral}
 \frac{|s|}{|t_1|^m}
       \int_\Omega |z|^m g(z)\psi(az)\,dz.
\end{equation}
Choose a fixed smaller lattice $\Lambda_0\subset\Lambda$
such that $\Omega$ is a finite union of its cosets and $|z|$
is constant on each such coset.  This is possible because the
absolute value is locally constant on $K^\times$ and $\Omega$
is compact.  The amplitude $|z|^m g(z)$ is constant on each coset,
independently of the values of $g$.

There is at least one field component $j$ with $\zeta_j\ne1$.
After shrinking $E$ and $V$, the number $1-\zeta_j\rho_j$
has the same nonzero absolute value as $1-\zeta_j$.
Thus $|a_j|=|s(1-\zeta_j\rho_j)/t_{1,j}|$ tends to infinity.
The annihilator $\Lambda_0^\perp$ is compact, so
$a\notin\Lambda_0^\perp$ once $t$ is sufficiently small.
Character orthogonality on each coset of $\Lambda_0$ now makes
the integral in \eqref{orb-cross-integral} zero.  This proves
\eqref{orb-cross-sum}.
\end{proof}

The proof isolates the role of the support condition: it is needed
for invariance under a change of representative when $m=1$.
For $m\ge2$ the same conclusion holds without that condition.
All choices remain local choices made after $K,R,\psi,m,\Omega,
\Lambda,\Sigma,\zeta$ are fixed.  No constant uniform over
all residue characteristics or over all block sizes has been
asserted.

\begin{remark}[The support condition cannot simply be omitted when $m=1$]
Take $K=\mathbb Q_3$, $R=\mathbb Z_3$, and a character with
annihilator $R$.  Let $m=1$, $s=2$, $\zeta=-1$,
$\Sigma=\{0\}$, $\kappa=1$, and set
\[
 \Omega=(3+9\mathbb Z_3)\cup(-3+9\mathbb Z_3),\qquad
 \Lambda=9\mathbb Z_3,\qquad g=1_{3+9\mathbb Z_3}.
\]
All the assumptions of Theorem~\ref{orb-cross-vanish}, other than
\eqref{orb-cross-support}, hold.  Put $t=3^N$ with $N\ge4$.
One may choose $c=-9/t$ as the representative of its coset in
$K/R$.  The roots of $z^2=-tc$ are $3$ and $-3$, and the
contribution of this coset is $\psi(12/t)$.
Now replace only this representative by $c+1$.  Its root in
$3+9\mathbb Z_3$ is
\[
 z'=3\sqrt{1-3^{N-2}}.
\]
The near-one square root exists, and its binomial series gives
\[
 \frac{4(z'-3)}{t}
       =-\frac23+O(3^{N-3})
       \equiv-\frac23\pmod{\mathbb Z_3}.
\]
Indeed the first-order term is $-2/3$ and, for $j\ge2$, the
$j$th term has valuation at least $N-3$; the binomial
coefficients of $(1+X)^{1/2}$ lie in $\mathbb Z_3$.
Thus this coset's contribution is multiplied by
$\psi(-2/3)\ne1$, while all other contributions are unchanged.
The two total sums differ, so they cannot both vanish.
Such a failure occurs for arbitrarily small $t$.
The omitted hypothesis fails at $z=3$, since
$z^{-1}-(-z)^{-1}=2/3\notin R$.
This is a counterexample to removing a hypothesis, not to the
statement in the source.
\end{remark}

\subsection{The exact truncated orbital identity and its remaining algebraic inputs}

A block configuration $h$ consists of a finite list of block sizes
and one scalar in $K^\times$ for each block.  Write $h[s,z]$ for
the configuration obtained by adjoining the block $zJ_s$, where
$J_s$ is the reversing permutation matrix of size $s$.  Zero-sized
blocks contribute empty matrices and can be omitted.  The matrix
$Z(h)$ is the block diagonal matrix of the configuration.

Let $\mathcal K(h,h')$ denote the local marked-cokernel kernel of
the preceding section.  For convenience define
\begin{align}
 N_s(h,h';z,z')
   &=|zz'|^{s(s-1)/2}
                 \mathcal K(h[s,z],h'[s,z']),\\
 W_h(t,c)&=h[m,t][1,c].
\end{align}
The truncated Gram sum is
\begin{equation}\label{orb-cut-gram}
 U_t(h,h';t_1,t_2)=|t_1t_2|^{m(m+1)/2}
       \sum_{\substack{x\in K/R\cr c_x\in\mathcal C_t}}
          \mathcal K(W_h(t_1,c_x),W_{h'}(t_2,c_x)).
\end{equation}

The next theorem isolates the exact algebraic inputs still needed
for this particular kernel.  The hypotheses are stated in
ordinary mathematical terms rather than concealed in a general
existence assertion.

\begin{theorem}\label{orb-cut-identity}
Fix a finite collection $\mathcal F$ of block configurations of
the same total size, a compact open $\Omega$ as above, and a
finite $\Sigma\subset K$.  Suppose that the following properties
hold for the kernel $\mathcal K$.
\begin{enumerate}
\item The kernel is Hermitian.
\item For each $h\in\mathcal F$ there is an additive lattice
$L_h$ such that, whenever $z,z'\in K^\times$, $z\in\Omega$,
and $z'-z\in L_h$, one has
\[
 \mathcal K(h[s,z'],w)=\mathcal K(h[s,z],w)
 \quad\text{for every comparison configuration }w.
\]
\item For $h,h'\in\mathcal F$ and $z,z'\in K^\times$,
\begin{equation}\label{orb-inverse-support}
 \mathcal K(h[s,z],h'[s,z'])\ne0
 \quad\Longrightarrow\quad z^{-1}-z'^{-1}\in R.
\end{equation}
\item With $\gamma_m$ the Gaussian phase from the preceding
section, the following germ formula holds.  For every compact
$Q\subset K^\times$ there is a sufficiently small neighborhood
of zero such that, for $h\in\mathcal F$, $t$ in that neighborhood,
$ct^m\in Q$, and every comparison configuration $w$,
\begin{align}\label{orb-germ-input}
 \mathcal K(W_h(t,c),w)
  ={}&|t|^{-m^2/2}|s|^{-1/2}
       \sum_{z^s=(-t)^m c}\gamma_m(zt)|z|^{m^2/2}
             \psi(sz/t)\mathcal K(h[s,z],w).
\end{align}
The neighborhood is uniform in $c$ and $w$ as stated.
\end{enumerate}
Then for all sufficiently small $t\in K^\times$, all
$\sigma,\sigma'\in\Sigma$, and all $h,h'\in\mathcal F$,
with $t_1=t_\sigma$ and $t_2=t_{\sigma'}$, one has
\begin{equation}\label{orb-exact-cut}
 U_t(h,h';t_1,t_2)
   =\int_\Omega N_s(h,h';z,z)
                         \psi(z(\sigma-\sigma'))\,dz.
\end{equation}
The small neighborhood is independent of the representatives $c_x$.
\end{theorem}

\begin{proof}
All choices below involve finitely many configurations, shifts,
and $s$th roots of unity, so the required neighborhoods and lattices
can be intersected.  Choose the unique root near one satisfying
\[
 \kappa^s=(t_2/t_1)^m.
\]
For each root $z^s=(-t_1)^mc$, the roots for the second germ
formula are exactly $\zeta\kappa z$, with $\zeta^s=1$.
The truncation sets for $t,t_1,t_2$ agree.  In particular all these roots
are in $\Omega$.

Take
$Q=\{(-1)^m z^s:z\in\Omega\}$ in the germ hypothesis.
Apply \eqref{orb-germ-input} first to the first argument of the
kernel in \eqref{orb-cut-gram}, and then to its second argument
by Hermitian symmetry.  These are finite sums.  Since
$|t_1|=|t_2|=|t|$ and $|\zeta|=|\kappa|=1$, the result is
\begin{align}\label{orb-expanded-cut}
 U_t(h,h';t_1,t_2)
 ={}&\frac{|t|^m}{|s|}
 \sum_{\zeta^s=1}
 \sum_{\substack{x\in K/R\cr c_x\in\mathcal C_t}}
 \ \sum_{z^s=(-t_1)^m c_x}
 |z|^{-m}\notag\\
 &\quad\cdot\gamma_m(zt_1)
       \overline{\gamma_m(\zeta\kappa zt_2)}
       N_s(h,h';z,\zeta\kappa z)\notag\\
 &\quad\cdot\psi(sz/t_1-s\zeta\kappa z/t_2).
\end{align}
For completeness, the power of $|t|$ comes from
$m(m+1)-m^2=m$.  The power of $|z|$ comes from
$m^2-s(s-1)=-m$.  These equalities account for the whole
normalization in this calculation.

The factor $\kappa t_2/t_1$ is a square when $t$ is sufficiently
small.  Square invariance of the Gaussian phase gives
\[
 \gamma_m(\zeta\kappa zt_2)=\gamma_m(\zeta zt_1).
\]
The local constancy hypothesis, applied in the second argument
using Hermitian symmetry, gives
\[
 N_s(h,h';z,\zeta\kappa z)=N_s(h,h';z,\zeta z).
\]
Here $\zeta z\in\Omega$, and
$(\kappa-1)\zeta z$ is uniformly small.  The positive factor in
the definition of $N_s$ is unchanged because $|\kappa|=1$.

Choose a fixed lattice $\Lambda_0$ small enough that $\Omega$
is a union of its cosets, that $|z|$ is constant on those cosets,
and that $(z+\lambda)/z$ is a square for
$z\in\Omega$, $\lambda\in\Lambda_0$.  Also require
$\Lambda_0\subset L_h$ and $\zeta\Lambda_0\subset L_h$
for every $h\in\mathcal F$ and $\zeta^s=1$.
These choices are possible by compactness and
Lemma~\ref{orb-near-one}.  The amplitude
\[
 g_{\zeta,t_1}(z)=|z|^{-m}\gamma_m(zt_1)
           \overline{\gamma_m(\zeta zt_1)}
                       N_s(h,h';z,\zeta z)
\]
is then $\Lambda_0$-invariant, with that lattice independent of
$t_1$.  Square invariance handles the two Gaussian factors, and
the local constancy hypothesis handles both kernel arguments.
If this amplitude is nonzero, \eqref{orb-inverse-support}
gives exactly \eqref{orb-cross-support}.
Theorem~\ref{orb-cross-vanish} therefore removes every term in
\eqref{orb-expanded-cut} with $\zeta\ne1$.

For $\zeta=1$ the two Gaussian factors multiply to one.
The remaining phase satisfies
\begin{equation}\label{orb-diagonal-phase}
 \psi(sz/t_1-s\kappa z/t_2)
               =\psi(z(\sigma-\sigma'))\qquad(z\in\Omega).
\end{equation}
Here is a direct verification, including the necessary uniformity.
As a convergent power series in the small parameter $t$,
\[
 \kappa=1+\frac{m}{s}(\sigma-\sigma')t+t^2q(t),
\]
where $q$ is analytic near zero.  This follows by substituting
$t_2/t_1=(1+\sigma t)/(1+\sigma't)$ in its defining
near-one root.  Consequently
\[
 s/t_1-s\kappa/t_2
     =\sigma-\sigma'+tq_1(t)
\]
for another analytic function $q_1$.  Its values are bounded on
a smaller neighborhood.  Multiplication by $z\in\Omega$ makes
the remainder lie in $R$ for all sufficiently small $t$.
There are finitely many pairs of shifts, so the choice is uniform.
Triviality of $\psi$ on $R$ proves \eqref{orb-diagonal-phase}.

After shrinking $\Lambda_0$ again to annihilate all differences
$\sigma-\sigma'$, the remaining amplitude
\[
 |z|^{-m}N_s(h,h';z,z)\psi(z(\sigma-\sigma'))
\]
is $\Lambda_0$-invariant.  Theorem~\ref{orb-change-variable}
turns its root sum into
\[
 \frac{|s|}{|t_1|^m}
       \int_\Omega N_s(h,h';z,z)
                      \psi(z(\sigma-\sigma'))\,dz.
\]
This cancels the factor $|t|^m/|s|$ in
\eqref{orb-expanded-cut}, proving \eqref{orb-exact-cut}.
\end{proof}

The theorem proves the full analytic passage from the orbital germ
formula to the truncated orbital identity, including the mixed critical terms.
It does not establish the germ formula or the inverse-parameter
support for an arbitrary Hermitian kernel.  For the actual marked
cokernel kernel those properties are additional algebraic facts;
their exact content is examined next.

\subsection{Finiteness of the actual sums}

We make the kernel definition precise enough for the algebraic
support argument.  For a configuration $w$ of total size $n$, let
$L_w(K)$ be the block lower unipotent group, with identity matrices
on its diagonal blocks and arbitrary blocks strictly below them.
Let $U_n(K)$ be the upper unipotent group.  For either group,
placing $R$ in place of $K$ means that all its entries lie in $R$.
The source's lower group is exactly $L_w(K)$: conjugating a lower
triangular diagonal block by its reversing permutation matrix
makes it upper triangular, so a diagonal block that remains lower
unipotent must be the identity.

For
\[
 q=(x,y)\in L_w(R)\backslash L_w(K)
                    \times U_n(K)/U_n(R),
\]
choose representatives and put
\[
 Z_q=xZ(w)y,\qquad
 a_q=\psi\left(\sum_{i=0}^{n-2}x_{i+1,i}
                       +\sum_{i=0}^{n-2}y_{i,i+1}\right).
\]
Let $\mathcal I_w$ be the set of those cosets for which
$Z_q\in M_n(R)$.  The marked module attached to $Z_q$ is
$C_{Z_q}=R^n/Z_qR^n$, with marked vector $[e_{n-1}]$ and
functional $[v]\mapsto (Z_q^{-1}v)_{n-1}\pmod R$.
For $n=0$, take the zero module with zero markings.  Then the actual kernel is
\begin{equation}\label{orb-actual-kernel}
 \mathcal K(w,w')=
 \sum_{q\in\mathcal I_w}\sum_{q'\in\mathcal I_{w'}}
 a_q\overline{a_{q'}}\,
       \#\operatorname{Iso}(C_{Z_q},C_{Z_{q'}}).
\end{equation}
Isomorphisms here preserve both markings.

\begin{lemma}\label{orb-finite-orbits}
The set $\mathcal I_w$ is finite, and
\eqref{orb-actual-kernel} is independent of every representative
choice.  In particular it is the finite Gram kernel described
earlier.  For $w=h[s,z]$, the union of the integral orbit cosets
is still finite when $z$ ranges over a fixed compact subset of
$K^\times$.
\end{lemma}

\begin{proof}
First consider all pairs $(x,y)$ for which $xZ(w)y$ has entries
in $R$.  Write $xZ(w)y=LU$ in block form, with $L=x$ and
$U=Z(w)y$.  The determinant of the $i$th diagonal block of $U$
is the fixed nonzero number $\det(t_iJ_{s_i})$.
Block Gaussian elimination recovers $L$ and $U$ from the product
$Z$ by
\[
 U_{ij}=Z_{ij}-\sum_{a<i}L_{ia}U_{aj}\quad(j\ge i),\qquad
 L_{ji}=\left(Z_{ji}-\sum_{a<i}L_{ja}U_{ai}\right)U_{ii}^{-1}
 \quad(j>i).
\]
At each step $U_{ii}^{-1}$ is its adjugate divided by its fixed
nonzero determinant.  Thus the entries of $L,U$, and
$y=Z(w)^{-1}U$ are bounded when the entries of $Z$ lie in the
compact set $R$.  All conditions on $x,y$ are closed.  The set
of such pairs is therefore a closed bounded subset of a
finite-dimensional vector space over the local fields, hence
compact.  Its image in the displayed coset space is finite because
the integral unipotent subgroups are open and that coset space
is discrete.  The case $n=0$ has just one pair and the same
conclusion.
The same proof is uniform for $z$ in a compact subset of
$K^\times$: the diagonal-block determinants and their inverses
then range over compact sets.  Apply the recovery formulas on
the compact set of pairs $(Z,z)$ instead.  The group $L_w$
depends only on the block sizes, so the quotient space is fixed.

Changing representatives replaces $Z_q$ by $aZ_qb$, with
$a\in L_w(R)$ and $b\in U_n(R)$.  Both matrices are invertible
over $R$.  The map $[v]\mapsto[av]$ identifies the cokernels.
It preserves the marked vector because $ae_{n-1}=e_{n-1}$,
and preserves the functional because
$e_{n-1}^tb^{-1}=e_{n-1}^t$.
The sums of the adjacent entries are additive on unipotent
groups.  Their changes lie in $R$, so the phase $a_q$ is
unchanged.  This proves representative independence and the
claimed identification.
\end{proof}

\begin{proposition}\label{orb-local-constancy-proved}
Fix $h$, an integer $s\ge1$, and a compact set
$C\subset K^\times$.  There is an additive lattice $L$ such
that for $z\in C$ and $z'\in K^\times$ with $z'-z\in L$,
\[
 \mathcal K(h[s,z'],w)=\mathcal K(h[s,z],w)
 \quad\text{for every configuration }w.
\]
The lattice is independent of the comparison configuration.
\end{proposition}

\begin{proof}
Enlarge $C$ to a compact open subset of $K^\times$ and choose
a preliminary lattice so that its small translates remain
inside this larger set.  By Lemma~\ref{orb-finite-orbits}, a
single finite list of orbit cosets contains all the possible
integral terms.  Choose fixed representatives $(x_q,y_q)$
for them.  Write
$Z_q(z)=x_qZ(h[s,z])y_q$.
These matrices are invertible, and their inverses range over
a compact set as $q,z$ vary.  Choose an integer $b\ge0$
such that every entry of every $Z_q(z)^{-1}$ lies in
$p^{-b}R$.

Shrink the lattice $L$ so that, for every listed coset,
\[
 \Delta_q=Z_q(z')-Z_q(z)\in p^{2b+1}M_n(R)
 \quad\text{whenever }z'-z\in L.
\]
This is possible because the difference is a fixed linear
function of $z'-z$ for each of finitely many $q$.
It implies that $Z_q(z')$ is integral exactly when $Z_q(z)$
is integral.  Put $M=Z_q(z)^{-1}\Delta_q$.
Then $M\in p^{b+1}M_n(R)$, and the geometric series for
$(I+M)^{-1}$ converges inside $M_n(R)$, since $R$ is compact
and hence complete.  Thus
\[
 Z_q(z')R^n=Z_q(z)R^n,
 \qquad Z_q(z')^{-1}-Z_q(z)^{-1}\in pM_n(R).
\]
The first identity identifies the cokernels by the identity
map of $R^n$; the second preserves every inverse-coordinate
functional modulo $R$.  The marked vector is unchanged,
and the phase depends only on $x_q,y_q$, not on $z$.
Consequently every marked module and weight in the finite
sum is unchanged.  Counting isomorphisms to any comparison
marked module gives the asserted equality, uniformly in $w$.
\end{proof}

\subsection{The finite-module character argument}

We use the following elementary duality fact.  Its proof also
explains why a nonmaximal order causes no difficulty.

\begin{lemma}\label{orb-dual-span}
Let $C$ be a finite $R$-module, and let
$\lambda_1,\ldots,\lambda_l,\nu:C\to K/R$ be $R$-linear.
Then $\nu$ is an $R$-linear combination of the $\lambda_i$
if and only if it vanishes on their common kernel.
\end{lemma}

\begin{proof}
Only the reverse implication needs proof.  The character
$c\mapsto\psi(\nu(c))$ then factors through the image of
$\lambda=(\lambda_1,\ldots,\lambda_l)$.
That image is a finite additive subgroup of $(K/R)^l$, contained
in $(p^{-N}R/R)^l$ for a sufficiently large integer $N$.
Every character of a subgroup of a finite abelian group extends
to the group: adjoin one element at a time, choosing a root of
the prescribed value of its least positive multiple lying in
the subgroup.  Self-duality of $R$ gives a perfect pairing
\[
 R/p^NR\ \times\ p^{-N}R/R\longrightarrow\mathbb C^\times,
 \qquad(a,z)\longmapsto\psi(az).
\]
Its kernels are exactly $p^NR$ and $R$; the two finite groups
have the same order by Haar-volume scaling.  Thus the extended
character has the form
$\psi(\sum_i a_i\lambda_i(c))$ with $a_i\in R$.
Apply this equality to $rc$ for every $r\in R$.  Self-duality
then implies
$\nu(c)-\sum_i a_i\lambda_i(c)=0$ in $K/R$ for every $c$.
\end{proof}

For an integral invertible $n$ by $n$ matrix $Z$, put
\[
 \lambda_j^Z([v])=(Z^{-1}v)_j\pmod R
 \qquad(0\le j<n).
\]
These functionals are well-defined on $C_Z$.  If
$\phi:C_Z\to C_{Z'}$ is an isomorphism, let $T_j(\phi)$
mean that for every $i\ge j$ there are $a_{il}\in R$ such that
\begin{equation}\label{orb-triangular}
 \lambda_i^{Z'}\phi=\lambda_i^Z+
                                  \sum_{l>i}a_{il}\lambda_l^Z.
\end{equation}
The marked-module isomorphisms already satisfy $T_{n-1}$.
By Lemma~\ref{orb-dual-span}, assuming $T_{k+1}$, the additional
condition $T_k$ is equivalent to vanishing of
$\lambda_k^{Z'}\phi-\lambda_k^Z$ on the common kernel of the
functionals $\lambda_i^Z$ with $i>k$.

\begin{proposition}\label{orb-triangular-average}
Let $w=h[s,t]$ and $w'=h'[s,t']$, where the two previous
configurations have the same total size $r$, and let
$n=r+s$.  If $r\le k\le n-2$, then the sum in
\eqref{orb-actual-kernel} is unchanged when, among the
isomorphisms counted in each summand, the condition $T_{k+1}$
is strengthened to $T_k$.
\end{proposition}

\begin{proof}
It suffices to show that the weighted contribution of the
isomorphisms satisfying $T_{k+1}$ but not $T_k$ is zero.
We partition the upper-unipotent variables into smaller cosets.
Put $a=k+1$ and $l=n-a$.  Let $P_k\subset U_n(K)$ consist of
matrices whose two diagonal blocks of sizes $a,l$ have entries
in $R$; the upper-right rectangle is unrestricted.  This is a
subgroup containing $U_n(R)$.  Every coset in $P_k/U_n(R)$ has
a representative
\[
 p_B=\begin{pmatrix}I_a&B\\0&I_l\end{pmatrix},
 \qquad B\in M_{a,l}(K)/M_{a,l}(R).
\]
Indeed one can remove both integral diagonal blocks by right
multiplication.  Two displayed representatives define the same
coset exactly when their difference $B-B'$ is integral.

Fix the lower-unipotent cosets and representatives $y_0,y'_0$
of the upper cosets modulo $P_k$.  Write
$Z_0=xZ(w)y_0$ and $Z'_0=x'Z(w')y'_0$.
Only finitely many $B,B'$ make $Z_0p_B,Z'_0p_{B'}$ integral,
by Lemma~\ref{orb-finite-orbits}.  Empty collections contribute
zero and may be discarded.  The first $a$ columns of $Z_0$
and $Z'_0$ are then integral, since multiplication by $p_B$
does not change them.  Define
\[
 S_0=Z_0(R^a\oplus0)\subset R^n,\quad
 P=R^n/S_0,\quad
 \Lambda:P\to K^l,\quad
 \Lambda([v])=\operatorname{pr}_{>k}(Z_0^{-1}v),
\]
and define $S'_0,P',\Lambda'$ in the same way.
The maps are well-defined.  For
$L_B=Z_0p_BR^n\subset R^n$, the restriction
\begin{equation}\label{orb-tail-section}
 \Lambda:L_B/S_0\longrightarrow R^l
\end{equation}
is an isomorphism.  It is surjective by the last $l$ columns
of $Z_0p_B$; its kernel is zero because a vector with last
$l$ inverse coordinates zero belongs to the first $a$ columns.
The analogous statement holds for $L'_{B'}$.

The marked vector $e=e_{n-1}$ satisfies
$\Lambda([e])=\Lambda'([e])=0$.
To check this, use
$Z_0^{-1}e=y_0^{-1}Z(w)^{-1}x^{-1}e
=t^{-1}y_0^{-1}e_r$.
An upper triangular matrix times $e_r$ has zero coordinates
greater than $r$, hence greater than $k$.  The same calculation
applies to $Z'_0$.

We now describe the isomorphisms in this finite sum by data
independent of $B,B'$.  The higher inverse-coordinate maps of
$Z_0p_B$ are just $\Lambda$ modulo $R^l$.
An isomorphism
$\phi:R^n/L_B\to R^n/L'_{B'}$ satisfying $T_{k+1}$ therefore
satisfies
\[
 \Lambda'(v')\equiv D\Lambda(v)\pmod{R^l}
 \quad\text{whenever }\phi([v])=[v'],
\]
for some upper unipotent $D\in U_l(R)$.
Choose once and for all representatives for the equivalence
relation
\[
 D_1\sim D_2\quad\Longleftrightarrow\quad
 (D_1-D_2)\Lambda(P)\subset R^l.
\]
The isomorphism $\phi$ determines exactly one such class.
Replacing $D$ by its chosen representative preserves the
congruence above.

There is a unique lift
\begin{equation}\label{orb-partial-lift}
 \Psi:P\longrightarrow P',\qquad
 \Lambda'\Psi=D\Lambda,
\end{equation}
that induces $\phi$.  Indeed, start with any representative
$v'$ of $\phi([v])$.  Its error
$\Lambda'(v')-D\Lambda(v)$ lies in $R^l$ and can be removed
by subtracting the unique element of $L'_{B'}/S'_0$ with
that image under \eqref{orb-tail-section}.  This defines the
lift independently of all choices and proves linearity.
It is injective because an element in its kernel lies in
$L_B/S_0$ and has zero $\Lambda$-image.  It is surjective by
the same adjustment in the source, using $D^{-1}R^l=R^l$.
The lift carries the marked vector to its counterpart, since
both have zero higher coordinates.

Conversely, a marked isomorphism $\Psi:P\to P'$ satisfying
\eqref{orb-partial-lift}, with $D$ a chosen representative,
descends to exactly one isomorphism $\phi$ whenever
\begin{equation}\label{orb-related-lattices}
 \Psi(L_B/S_0)=L'_{B'}/S'_0.
\end{equation}
Its last inverse-coordinate functional is preserved because
the last row of $D$ is the last coordinate vector; hence it
is an isomorphism of the marked modules.  It satisfies
$T_{k+1}$ by construction.

The failure of $T_k$ is also intrinsic to $\Psi$.  It is
equivalent to the existence of $v,v'\in R^n$ such that
\begin{equation}\label{orb-bad-pair}
 \Psi([v])=[v'],\qquad \Lambda(v)=\Lambda'(v')=0,
 \qquad (Z_0'^{-1}v')_k-(Z_0^{-1}v)_k\notin R.
\end{equation}
To see this, apply Lemma~\ref{orb-dual-span} to the common
kernel of the higher coordinate functionals.  A representative
whose higher coordinates lie in $R^l$ can be changed, using
\eqref{orb-tail-section}, to one whose higher coordinates
are exactly zero.  The corresponding change in the target
is made using $\Psi$.  These changes do not alter any
inverse-coordinate difference modulo $R$.  On exact zero
higher coordinates, multiplication by $p_B^{-1}$ does not
alter the $k$th coordinate, proving \eqref{orb-bad-pair}.
The converse follows by applying the same criterion to that
pair of representatives.

Group all bad isomorphisms by their lift $\Psi$.  Each resulting
set of pairs $(B,B')$ is finite.  Fix such a $\Psi$, choose
$v,v'$ as in \eqref{orb-bad-pair}, and write the last
nonintegral difference there as $d$.  By self-duality there
is $r_0\in R$ with $\psi(r_0d)\ne1$.
Put $v_0=r_0v$, $v'_0=r_0v'$, and
\[
 \xi=Z_0^{-1}v_0,\quad \xi'=Z_0'^{-1}v'_0,
 \quad\beta=(1,0,\ldots,0)\in R^{1\times l},
 \quad\beta'=\beta D^{-1}.
\]
Extend $\beta,\beta'$ by zeros in the first $a$ coordinates.
Both $\xi$ and $\xi'$ vanish in the last $l$ coordinates.
The matrices
\[
 \sigma=I+\xi\beta,\qquad
 \sigma'=I+\xi'\beta'
\]
are therefore upper unipotent shears supported in the
upper-right rectangle.  Their inverses have the corresponding
minus signs.

Replacing $p_B,p_{B'}$ by $\sigma p_B,\sigma'p_{B'}$
preserves integrality: their products with $Z_0,Z'_0$ change
by $v_0\beta,v'_0\beta'$, whose entries lie in $R$.
It also preserves \eqref{orb-related-lattices}.
Here are the details of that assertion.  For $x\in R^n$,
put $b(x)=\beta\Lambda(x)$.  Then
\[
 x\in Z_0\sigma p_BR^n
 \quad\Longleftrightarrow\quad
 b(x)\in R\text{ and }x-b(x)v_0\in L_B.
\]
This follows directly by applying the inverse shear; on
$L_B$ the higher inverse coordinates are integral, and
$\beta\xi=0$.
For representatives $x,y$ with $\Psi([x])=[y]$, the analogous
scalar in the target is
$\beta'\Lambda'(y)=\beta'D\Lambda(x)=b(x)$.
Moreover $\Psi([v_0])=[v'_0]$.  The equivalence of the two
membership conditions now follows from
\eqref{orb-related-lattices}.  Applying the inverse shears
shows that the resulting transformation is a bijection of
the finite set of pairs grouped under $\Psi$.

Only one adjacent upper entry occurs in the shears, the
$(k,k+1)$ entry.  Since $\beta_{k+1}=\beta'_{k+1}=1$, the
change in the difference of the upper phases is
\[
 \xi_k-\xi'_k=-r_0d.
\]
Thus this permutation multiplies every weight in the finite
sum for $\Psi$ by the same number $\psi(-r_0d)\ne1$.
That sum must be zero.  Summing over the finitely many lifts
and over the outer cosets proves the proposition.

Finally, using canonical representatives $p_B$ has not
changed the original counts.  A right integral upper
unipotent change replaces the inverse-coordinate functionals
by an integral upper unipotent combination.  It preserves
the last functional and every condition $T_j$, in both
source and target.  Its phase is trivial.  Hence all the
partitions and changes of representatives used above are
legitimate in \eqref{orb-actual-kernel}.
\end{proof}

\begin{corollary}\label{orb-support-proved}
For the actual kernel \eqref{orb-actual-kernel}, if $s\ge1$,
the previous block configurations $h,h'$ have equal total
size, and
$\mathcal K(h[s,t],h'[s,t'])\ne0$, then
\[
 t^{-1}-t'^{-1}\in R.
\]
\end{corollary}

\begin{proof}
Put $r=\dim Z(h)=\dim Z(h')$ and $n=r+s$.
Marked isomorphisms already satisfy $T_{n-1}$.
Repeated use of Proposition~\ref{orb-triangular-average}
shows that the same kernel is obtained by counting only
isomorphisms satisfying $T_r$.
If its value is nonzero, at least one such isomorphism occurs.
For $Z=xZ(h[s,t])y$,
\[
 Z^{-1}e_{n-1}=t^{-1}y^{-1}e_r.
\]
Its $r$th coordinate is $t^{-1}$ and all coordinates greater
than $r$ are zero.  The corresponding coordinates for $Z'$
are $t'^{-1}$ and zeros.  Evaluate the relation
\eqref{orb-triangular} with $i=r$ on the marked vector,
which the isomorphism preserves.  It gives
$t'^{-1}\equiv t^{-1}\pmod R$, as asserted.
When $s=1$ no averaging step is needed: $r=n-1$ already.
\end{proof}

\paragraph{Exact finite controls.}
An accompanying integer-arithmetic program checks the averaging
identity for a single block of size two over $\mathbb Q_p$,
with $p=3,5$, $t=p$, and $t'=vp$ for $v=1,2,p+1$.
For a block with scalar $pv$, writing the upper unipotent
parameter as $b/p$ with $0\le b<p$ gives the integral orbit matrix
\[
 Z_b=\begin{pmatrix}0&pv\\pv&vb\end{pmatrix}.
\]
All marked-module isomorphisms are enumerated directly in
$(\mathbb Z/p^2\mathbb Z)^2/Z_b(\mathbb Z/p^2\mathbb Z)^2$,
using the possible images of its two generators.  The program
independently tests the additional condition $T_0$ and reduces
the phase sums modulo $1+X+\cdots+X^{p-1}$.
For $v=2$, the unrestricted phase polynomial is
$p(1+X+\cdots+X^{p-1})$, whereas the restricted polynomial is
zero: cancellation is essential.  For $v=1,p+1$, both values
are $p^2$.  These six exact controls agree with the support
corollary.  Four additional exact congruence checks, for
$N=4,5,6,7$, verify the representative change in the preceding
$\mathbb Q_3$ counterexample.  The general assertions rest on
the displayed proofs, not on these finite computations.

\paragraph{Scope of the reconstructed truncated orbital identity.}
For the actual finite marked-cokernel kernel, its Gram formula
proves Hermitian symmetry, Proposition~\ref{orb-local-constancy-proved}
proves the required local constancy uniformly against all comparison
configurations, and Corollary~\ref{orb-support-proved} proves the
inverse-parameter support.  The following finite-type argument
reduces comparison-configuration uniformity to finitely many
fixed marked modules.  Section~\ref{hr-section} constructs the
actual matrix test functions, and
Section~\ref{sec:theta-scalar-germ} proves their algebraic
reduction to the oscillatory integral.  Together with the previously
proved stationary-phase formula, these give the actual kernel germ
in Theorem~\ref{ga-actual-germ}.  Thus all hypotheses of
Theorem~\ref{orb-cut-identity} are established for this kernel.
The averaged diagonal estimate needed for the positive-majorant
induction is a separate finite character calculation, proved in
Section~\ref{tp-section}.  The quantitative global theta energy
estimates remain separate statements; they do not follow merely
from these local identities.


\subsection{Reduction of the germ formula to finitely many fixed marked modules}
\label{gu-section}

The comparison configuration in the remaining orbital germ formula is
arbitrary. Its uniformity can nevertheless be reduced to finitely many
fixed targets. The reason is that the determinant stays in a compact
subset of $K^\times$. This observation does not prove the scalar germ
formula, but removes a separate uniformity requirement from that task.
We retain the local field product $K$, compact open subring $R$, and
marked cokernels used in Section~\ref{sec:theta-orbital-identity}.

\begin{lemma}\label{gu-bounded-types}
Let $n\ge0$ and let $D\subset K^\times$ be compact. There is an
integer $N\ge0$, depending only on $R$ and $D$, such that every
$Z\in M_n(R)\cap\mathrm{GL}_n(K)$ with $\det Z\in D$ satisfies
\[
             p^N R^n\subset ZR^n\subset R^n.
\]
Consequently the finite modules $C_Z=R^n/ZR^n$, with any specified
vector $e\in C_Z$ and any $R$-linear functional
$\lambda:C_Z\to K/R$, belong to finitely many isomorphism classes.
Writing $a=\#(R/p^NR)$ and $B=a^n$, their number is at most
\[
                         2^B B a^B.
\]
These assertions hold for nonmaximal compact open orders as well.
\end{lemma}

\begin{proof}
The inverse set $D^{-1}$ is compact. Since $R$ is a neighborhood
of zero, multiplication by a sufficiently large power of $p$ sends
all of $D^{-1}$ into $R$; choose such a power $p^N$.
The adjugate of an integral matrix has entries in $R$. Therefore
\[
 p^N Z^{-1}=\frac{p^N}{\det Z}\operatorname{adj}(Z)
                                      \in M_n(R),
\]
which proves the first inclusion. The second is integrality of $Z$.
The case $n=0$ has the zero module and also satisfies the conclusion.

The group $R/p^NR$ is finite: it is compact and discrete, because
$p^NR$ is an open subgroup of $R$. Put $A=R/p^NR$.
The module $C_Z$ is a quotient of the finite module $A^n$ and has
cardinality at most $B$. There are at most $2^B$ choices for its
defining submodule, at most $B$ choices for its marked vector,
and at most $a^B$ choices for its functional. For the last claim,
$p^NC_Z=0$ implies that the functional takes values in
\[
          (K/R)[p^N]=p^{-N}R/R,
\]
which has cardinality $a$ by multiplication by $p^N$.
Counting all functions into this finite set only increases the
number of $R$-linear ones. These finite presentations cover all
the marked modules under consideration, so also bound their
number of isomorphism classes. No principal-ideal property of
$R$ is used.
\end{proof}

For a finite $R$-module with specified vector and functional,
write $T=(C,e,\lambda)$ and define its orbital coefficient by
\begin{equation}\label{gu-orbital-coefficient}
 \mathcal B(w;T)=
       \sum_{q\in\mathcal I_w}a_q\,
                     \#\operatorname{Iso}(C_{Z_q},T).
\end{equation}
Isomorphisms preserve both specifications. This is a finite sum by
Lemma~\ref{orb-finite-orbits}. It depends only on the isomorphism
class of $T$. Directly from the definition of the kernel,
\begin{equation}\label{gu-kernel-coefficients}
 \mathcal K(w,w')=
       \sum_{q'\in\mathcal I_{w'}}\overline{a_{q'}}\,
                        \mathcal B(w;C_{Z_{q'}}).
\end{equation}
The notation on the right includes the specified vector and
functional of each cokernel.

\begin{proposition}\label{gu-uniform-germ}
Fix $m\ge1$, put $s=m+1$, and fix a finite collection
$\mathcal F$ of configurations whose total size is $r$.
Put $n=r+s$ and fix a compact set $Q\subset K^\times$.
For a critical root $z^s=(-t)^mc$, write
\[
 b(t,c;z)=|t|^{-m^2/2}|s|^{-1/2}
              \gamma_m(zt)|z|^{m^2/2}\psi(sz/t).
\]
Suppose that for each $h\in\mathcal F$ and each fixed
$Z'\in M_n(R)\cap\mathrm{GL}_n(K)$ there is a neighborhood
$V_{h,Z'}$ of zero such that, whenever $t,c\in K^\times$,
$t\in V_{h,Z'}$ and $ct^m\in Q$, one has
\begin{equation}\label{gu-pointwise-germ}
 \mathcal B(W_h(t,c);C_{Z'})
       =\sum_{z^s=(-t)^mc}
                 b(t,c;z)\mathcal B(h[s,z];C_{Z'}).
\end{equation}
Then there is one neighborhood $V$ of zero such that, for every
$h\in\mathcal F$, $t,c\in K^\times$ with $t\in V$ and
$ct^m\in Q$, and every comparison configuration $w'$, one has
\begin{equation}\label{gu-uniform-germ-conclusion}
 \mathcal K(W_h(t,c),w')
       =\sum_{z^s=(-t)^mc}
                    b(t,c;z)\mathcal K(h[s,z],w').
\end{equation}
The total size and parameters of $w'$ are unrestricted. In fact,
only finitely many fixed matrices $Z'$ are needed in the
hypothesis.
\end{proposition}

\begin{proof}
Let $H=Z(h)$. The determinant of a reversing matrix of size $j$
is $(-1)^{j(j-1)/2}$. Hence
\[
 \det Z(W_h(t,c))
      =\det H\,(-1)^{m(m-1)/2}t^mc.
\]
For every critical root,
\begin{align*}
 \det Z(h[s,z])
  &=\det H\,(-1)^{m(m+1)/2}z^{m+1}\\
  &=\det H\,(-1)^{m(m+1)/2+m}t^mc
    =\det Z(W_h(t,c)).
\end{align*}
The last two exponents differ by $2m$. Multiplication by the
lower and upper unipotent matrices in an orbital term does not
change the determinant.

Thus the determinants of all integral matrices occurring on
either side, for every allowed parameter and every
$h\in\mathcal F$, lie in one compact subset $D\subset K^\times$.
Lemma~\ref{gu-bounded-types} gives a finite list of possible marked
cokernel types. Keep only those which actually occur on one of
the two sides for some allowed parameter, and choose a realizing
integral matrix $Z_T\in M_n(R)\cap\mathrm{GL}_n(K)$ for each type
$T$. This list and its representatives depend on the already
fixed $\mathcal F,m,Q$, not on any later $t,c$ or $w'$.
If the list is empty, all relevant orbital coefficients vanish
and the result is immediate.

Intersect the finitely many neighborhoods $V_{h,Z_T}$.
For $t$ in their intersection, \eqref{gu-pointwise-germ} holds
for every possible target type $T$. It also holds with any target
of a different type: both sides are zero, because an isomorphism
can exist only between equal marked types.

Apply these identities to the finitely many marked modules
$C_{Z_{q'}}$ occurring in an arbitrary comparison configuration
$w'$. Multiply by $\overline{a_{q'}}$ and sum. Formula
\eqref{gu-kernel-coefficients} and interchange of the finite sums
give \eqref{gu-uniform-germ-conclusion}. The chosen neighborhood
is independent of $w'$, as required.
\end{proof}

The proposition supplies precisely the comparison-configuration
uniformity in \eqref{orb-germ-input}. The scalar identity
\eqref{gu-pointwise-germ} is supplied by the initial-block
reduction in Section~\ref{hr-section} and the scalar germ in
Section~\ref{sec:theta-scalar-germ}. The present proposition
shows that, for fixed $m$, $\mathcal F$ and $Q$,
one need not establish bounds simultaneously for infinitely many
unrelated target modules. The finite list is determined before
the small parameter is chosen, so no exchange of an infinite
intersection with a neighborhood occurs.


\section{The scalar orbital germ from a Hessenberg determinant}
\label{sec:theta-scalar-germ}

We retain $K,R,\psi$ and the Haar normalization of the preceding
sections.  In particular $K$ is a finite product of finite extensions
of $\mathbb Q_p$, $R$ is a compact open subring, and the annihilator
of $R$ is $R$.  An absolute value $|a|$ denotes the modulus of
multiplication by $a$ on the additive group of $K$.  It is the
product of the usual local-field moduli.  All product spaces carry
product Haar measure.  The integer $p$ tends to zero in every
factor, and the lattices $p^aR$, $a\in\mathbb Z$, are cofinal among
neighborhoods of zero.

The purpose of this section is to prove the scalar two-block
orbital germ.  The essential algebraic step is a change from
unipotent coordinates to the first row of an upper Hessenberg
matrix.  Its determinant determines the one remaining diagonal
entry.  Only its expansion through first order is needed; local
constancy makes the resulting formula exact.

\subsection{Test functions and orbital sums}

Fix $m\ge1$, and write $s=m+1$ and $d=m(m+1)/2$.
Let $L,\Lambda\subset K$ be compact open additive subgroups,
and let $Q\subset K^\times$ be compact.  A permitted test function
$\Phi:M_s(K)\to\mathbb C$ satisfies the following conditions:
\begin{enumerate}
\item $\Phi(aZb)=\Phi(Z)$ for integral lower unipotent $a$ and
integral upper unipotent $b$ of size $s$.
\item If $\Phi(Z)\ne0$, then every entry of $Z$ belongs to $L$.
\item If $\det Z\in Q\cup(-Q)$ and $Y\in M_s(\Lambda)$, then
$\Phi(Z+Y)=\Phi(Z)$.
\end{enumerate}
The neighborhoods constructed below depend on $K,R,\psi,m,L,
\Lambda,Q$, and not on the values of $\Phi$.

Let $J_j$ denote the reversing permutation matrix of size $j$.
For a block configuration $w$ of size $s$, define
\[
 \mathcal O_\Phi(w)=
 \sum_{q\in L_w(R)\backslash L_w(K)\times U_s(K)/U_s(R)}
       \Phi(x_qZ(w)y_q)
       \psi\left(\sum_i(x_q)_{i+1,i}
                    +\sum_i(y_q)_{i,i+1}\right).
\]
Here the lower group and representatives have the definitions
used in \eqref{orb-actual-kernel}.  Every sum in this section is
finite.  Indeed the block Gaussian elimination proof of
Lemma~\ref{orb-finite-orbits} applies with the compact set $L$
in place of $R$: bounded product entries and fixed nonzero
diagonal-block determinants bound both unipotent factors.
Invariance of $\Phi$ and triviality of $\psi$ on $R$ prove
independence of representatives.

We abbreviate
\[
 \mathcal O_\Phi(t,c)=\mathcal O_\Phi(\operatorname{diag}(tJ_m,c)),
 \qquad
 \mathcal O_\Phi(zJ_s)=\mathcal O_\Phi(\text{one block }zJ_s).
\]
The integrands below are locally constant on their relevant
determinant sets.  Condition (3), continuity of the matrix maps,
and the compact support just proved show that all Haar integrals
are absolutely convergent.  No global regularity of $\Phi$
outside these sets is needed.

\subsection{The first-row determinant calculation}

Use indices $0,1,\ldots,m$.  For $\rho\in K$, $x=(x_1,\ldots,x_m)$,
and $y=(y_{ij})_{0\le i<j\le m}$, let
$H_t(\rho,x,y)$ be the matrix with diagonal
$(\rho,x_1,\ldots,x_m)$, entries $y_{ij}$ above the diagonal,
entries $t$ immediately below it, and all other entries zero.
Thus $H_t$ is upper Hessenberg.  Let $C_t(x,y)$ be its trailing
$m$ by $m$ principal submatrix; it does not depend on $\rho$ or
on the first-row variables $y_{0j}$.

For $x\in(K^\times)^m$ and $D\in K^\times$, put
\begin{align}
 P(x)&=\prod_{i=1}^m x_i,&
 x_0(D,x)&=D/P(x),\\
 g_D(x,y)&=\frac{y_{01}}{x_1}
     +x_0(D,x)\sum_{i=1}^{m-1}
                       \frac{y_{i,i+1}}{x_ix_{i+1}},&
 f_D(x)&=\sum_{i=1}^m x_i+x_0(D,x).
\end{align}
The empty sum is zero when $m=1$.

\begin{lemma}\label{sg-determinant-expansion}
Fix compact sets of allowed $D\in K^\times$,
$x\in(K^\times)^m$, and $y\in K^d$.  For all sufficiently small
$t$, uniformly over these sets,
\begin{equation}\label{sg-tail-modulus}
 C_t(x,y)\text{ is invertible},\qquad
                  |\det C_t(x,y)|=|P(x)|.
\end{equation}
There is then a unique $\rho=\rho_t(D,x,y)$ with
$\det H_t(\rho,x,y)=D$, and
\begin{equation}\label{sg-rho-expansion}
 \rho_t(D,x,y)=x_0(D,x)+t g_D(x,y)+t^2 E_t(D,x,y),
\end{equation}
where the values of $E_t$ are contained in one compact subset
of $K$.  The neighborhood and this bound are independent of
the individual parameters in the specified compact sets.
\end{lemma}

\begin{proof}
The determinant is linear in $\rho$, with coefficient
$\det C_t$.  At $t=0$ that coefficient is $P(x)$, which is
bounded away from zero in every field factor.  The ratio
$\det C_t/P(x)$ is therefore uniformly in $1+pR$ for small $t$.
Such a ratio is a unit of $R$, with inverse given by a convergent
geometric series, and hence has modulus one.  This proves
\eqref{sg-tail-modulus} and uniqueness of $\rho$.

The coefficient of $t$ in the Hessenberg determinant is
\begin{equation}\label{sg-det-linear}
 -\sum_{i=0}^{m-1}y_{i,i+1}
                 \prod_{j\ne i,i+1}d_j,
 \qquad (d_0,d_1,\ldots,d_m)=(\rho,x_1,\ldots,x_m).
\end{equation}
For completeness, in the permutation expansion a term using
exactly one entry below the diagonal must use an adjacent
transposition.  All other positions must be diagonal: the
sum of column minus row indices is zero, so the sole negative
displacement $-1$ permits exactly one positive displacement
$+1$.  The sign of that transposition is negative, proving
\eqref{sg-det-linear}.  Terms using two or more subdiagonal
entries are divisible by $t^2$.

Put $S=\sum_{i=1}^{m-1}y_{i,i+1}/(x_ix_{i+1})$.
There are polynomial remainders, uniformly bounded on the
compact parameter sets for small $t$, such that
\[
 \det C_t=P(1-tS)+t^2 A_t,
 \qquad
 \det H_t(0,x,y)=-t y_{01}P/x_1+t^2 B_t.
\]
Solving the linear equation for $\rho$ gives
\[
 \rho_t=
 \frac{D+t y_{01}P/x_1-t^2B_t}{P(1-tS)+t^2A_t}.
\]
The denominator divided by $P$ is uniformly near one.
Expanding its inverse by the geometric series, through the
linear term, gives exactly \eqref{sg-rho-expansion}.
All remainders are bounded uniformly because $P^{-1}$, the
polynomial coefficients and the inverse of that near-one
denominator are bounded on the stated compact sets.  This also
proves the assertion when $m=1$, where the higher remainders
in the determinant vanish.
\end{proof}

The next observation ensures that this calculation covers every
contributing orbit, including ones for which the diagonal entries
were not initially assumed invertible.

\begin{lemma}\label{sg-nonzero-diagonals}
Fix a compact $Q_0\subset K^\times$ and a compact open additive
subgroup $L$.  There is a compact open set $E\subset K^\times$
and a neighborhood of zero such that the following holds.
If $t$ belongs to that neighborhood, all entries of
$\rho,x,y$ belong to $L$, and
\[
                    \det H_t(\rho,x,y)\in Q_0,
\]
then every $x_i$ belongs to $E$.
The set $E$ can also be chosen to contain every $x_i$ for which
$x_0(D,x),x_1,\ldots,x_m\in L$ and $D\in Q_0$.
\end{lemma}

\begin{proof}
Choose $a\ge0$ such that $L$, $Q_0$ and $Q_0^{-1}$ are
contained in $p^{-a}R$.  Write the determinant as
\[
             D=\rho\prod_i x_i+t B.
\]
For $t\in R$ all values of the polynomial $B$ lie in one
fixed compact subset of $K$.  Consequently, for small $t$,
$tB/D\in pR$, uniformly.  It follows that
$\rho\prod_i x_i=D(1-tB/D)$ is invertible, and hence so is
every $x_i$.  Moreover
\[
 x_i^{-1}=\frac{\rho\prod_{j\ne i}x_j}{D}
                           (1-tB/D)^{-1}
             \in p^{-(m+1)a}R.
\]
The same bound holds when $\rho=x_0(D,x)$ and $t=0$.
Thus one may take
\[
 E=\{u\in L\cap K^\times:u^{-1}\in p^{-(m+1)a}R\}.
\]
This set is open and compact in $K^\times$: both $u$ and its
inverse are bounded, so no sequence in it can approach a
nonunit in any field component.  It is a finite union of
cosets of some sufficiently small additive lattice.
\end{proof}

\subsection{Construction of the amplitude}

Let $Q_0=Q\cup(-Q)$.  For $D\in Q_0$, define
\begin{equation}\label{sg-amplitude}
 A_D(x)=|P(x)|^{-1}\int_{K^d}
   \Phi\bigl(J_sH_0(x_0(D,x),x,y)\bigr)
                          \psi(g_D(x,y))\,dy
 \quad(x\in(K^\times)^m),
\end{equation}
and set $A_D(x)=0$ otherwise.

\begin{lemma}\label{sg-amplitude-uniform}
There are a compact open $E\subset K^\times$ and an additive
lattice $\Lambda_1\subset K$, depending only on the fixed
data, such that every amplitude \eqref{sg-amplitude} is
supported in $E^m$ and invariant under translations by
$\Lambda_1^m$.  The inner integrands are supported in $L^d$.
All these assertions are uniform in $D\in Q_0$ and in every
permitted $\Phi$.
\end{lemma}

\begin{proof}
If the matrix value of $\Phi$ is nonzero, all its diagonal
and upper entries before reversal belong to $L$.  Their
diagonal product is $D$.  Lemma~\ref{sg-nonzero-diagonals},
at $t=0$, bounds $x$ in $E^m$, and condition (2) bounds $y$
in $L^d$.

Choose an additive lattice small enough to preserve $E$ and
$L$ under translation.  On the compact set
$Q_0\times E^m\times L^d$, the rational functions $P^{-1}$,
$x_0$ and $g_D$ have bounded denominators and are uniformly
continuous.  Shrink the lattice so that translating any of
the $x$ or $y$ coordinates by it has the following effects:
$|P|$ is unchanged; every entry of the change in
$J_sH_0(x_0,x,y)$ lies in $\Lambda$; and the change in
$g_D$ lies in $R$.  These choices are uniform because the
parameter set is compact, and $D$ is held fixed during each
translation.  More explicitly, the differences in reciprocal
coordinates satisfy
\[
 (x_i+u)^{-1}-x_i^{-1}=-u\bigl(x_i(x_i+u)\bigr)^{-1};
\]
their denominators remain in fixed bounded sets.  Products
and sums of these identities give each asserted uniform
continuity bound.

The determinant of the unperturbed matrix is
$(-1)^{s(s-1)/2}D$, which belongs to $Q\cup(-Q)$.
Thus condition (3) makes the matrix factor unchanged.
The character factor is unchanged because $\psi$ is trivial
on $R$.  Outside $E^m\times L^d$, both values vanish because
that set is stable under the selected translations.  This
proves a common translation invariance for the full inner
integrand.  Integrating in $y$ proves the claim for $A_D$.
\end{proof}

\subsection{The exact oscillatory integral}

\begin{theorem}\label{sg-oscillatory-reduction}
For all sufficiently small $t\in K^\times$, uniformly in
$c\in K^\times$ with $ct^m\in Q$ and in permitted $\Phi$,
put $D=(-t)^mc$.  Then
\begin{equation}\label{sg-oscillatory-identity}
 \mathcal O_\Phi(t,c)
   =|t|^{-d}\int_{(K^\times)^m}
                         A_D(x)\psi(f_D(x)/t)\,dx.
\end{equation}
For every $z\in K^\times$ with $z^s=D$,
\begin{equation}\label{sg-critical-amplitude}
 A_D(z,\ldots,z)
       =|z|^{m(m-1)/2}\mathcal O_\Phi(zJ_s).
\end{equation}
\end{theorem}

\begin{proof}
For the two-block matrix $W=\operatorname{diag}(tJ_m,c)$,
the lower group consists of matrices with identity diagonal
blocks and a single arbitrary last row below the diagonal.
Write this row as $(\eta_{m-1},\ldots,\eta_0)$.
Its simple-root phase is $\psi(\eta_0)$.
Write an upper unipotent matrix as $U=I+(u_{ij})_{i<j}$.
Product additive Haar measure in the $u_{ij}$ is Haar measure
on this group, normalized by $\operatorname{vol}(U_s(R))=1$.
Indeed left and right translations are triangular affine
changes in these coordinates with diagonal entries one;
successive one-variable translations preserve product Haar
measure.  The last-row coordinates likewise have measure
normalized by $\operatorname{vol}(R^m)=1$.

The summand is invariant under the corresponding integral
subgroups.  Its support lies in finitely many cosets and each
has volume one.  Hence the orbital sum equals the double
Haar integral of that summand.  Reverse the rows of $LWU$.
The resulting matrix has first row
\[
 t\sum_{i=0}^{m-1}\eta_i U_{i,*}+c e_m^{\,t},
\]
and row $i+1$ equal to $tU_{i,*}$ for $0\le i<m$.
Its determinant is $(-t)^mc=D$: subtract the indicated
linear combination of the last $m$ rows from the first row
and expand along the remaining entry $c$.

Now put $n_{ij}=tu_{ij}$.  This contributes the Jacobian
$|t|^{-d}$.  Set
\[
 x_i=n_{i-1,i}\quad(1\le i\le m),\qquad
 y_{ij}=n_{i-1,j}\quad(1\le i<j\le m).
\]
These are all the scaled upper entries, merely reordered.
The remaining first-row entries are denoted
$\tau_j=y_{0j}$, $1\le j\le m$, and its first diagonal
entry is $\rho=t\eta_0$.  Thus the reversed matrix is
$H_t(\rho,x,y)$ and its phase is
\begin{equation}\label{sg-coordinate-phase}
 \psi\left(\eta_0+\sum_{i=0}^{m-1}u_{i,i+1}\right)
              =\psi\left((\rho+\sum_{i=1}^m x_i)/t\right).
\end{equation}

If the matrix value of $\Phi$ is nonzero, all entries
$\rho,x,y$ belong to $L$.  Its determinant is in $Q_0$.
Lemma~\ref{sg-nonzero-diagonals} therefore restricts $x$ to
a fixed compact open $E^m$, and the upper entries below
the first row lie in $L$.  The same restrictions contain
the support of the amplitude in \eqref{sg-amplitude}.
Outside these restrictions both relevant integrands vanish.

Fix such $x$ and the upper entries below the first row.
The affine map $\eta\mapsto\tau$ has coefficient matrix
\[
             \bigl(tU_{i,j}\bigr)_{1\le j\le m,\ 0\le i<m}
                     =C_t(x,y)^{\,t}.
\]
For small $t$ it is invertible by
Lemma~\ref{sg-determinant-expansion}, and its determinant
has modulus $|P(x)|$.  It is therefore a bijection of all
$K^m$, with $d\eta=|P(x)|^{-1}d\tau$.
The entry $\rho$ after this change is the unique
$\rho_t(D,x,y)$ in that lemma, since the determinant remains
$D$.  In particular, this is a change on the whole first-row
space, not a selection of a branch or a restriction to an
unspecified neighborhood.

When one of the first-row coordinates lies outside $L$,
both matrix test functions vanish.  For all the remaining
$y\in L^d$, the expansion \eqref{sg-rho-expansion} is uniform.
Make $t$ small enough that
\[
 t\in\Lambda,\qquad
 \rho_t-x_0\in\Lambda,\qquad tE_t\in R.
\]
The first two conditions and test-function condition (3),
applied at $J_sH_0(x_0,x,y)$, give
\[
 \Phi(J_sH_t(\rho_t,x,y))
                  =\Phi(J_sH_0(x_0,x,y)).
\]
The third and \eqref{sg-coordinate-phase} give
\[
 \psi\left((\rho_t+\sum_i x_i)/t\right)
                 =\psi(f_D(x)/t)\psi(g_D(x,y)).
\]
Substitute these equalities into the changed integral.
Its two Jacobian factors are $|t|^{-d}$ and $|P(x)|^{-1}$.
Fubini's theorem, applicable to these compactly supported
locally constant functions, gives exactly
\eqref{sg-oscillatory-identity}.

Finally let $x_1=\cdots=x_m=z$ with $z^s=D$.  Then
$x_0=z$ and $g_D(x,y)=z^{-1}\sum_i y_{i,i+1}$.
Substitute $y_{ij}=zv_{ij}$ in \eqref{sg-amplitude}.
The factor $|P|^{-1}$ is $|z|^{-m}$ and the Jacobian is
$|z|^d$.  The matrix becomes $zJ_sU$, where $U$ is upper
unipotent with upper entries $v_{ij}$, and its phase is the
usual $\psi(\sum_i v_{i,i+1})$.  For a single block the lower
group is trivial.  The remaining Haar integral is therefore
$\mathcal O_\Phi(zJ_s)$, by the same volume-one coset argument.
As $d-m=m(m-1)/2$, this proves
\eqref{sg-critical-amplitude}.
\end{proof}

\subsection{The scalar germ}

\begin{corollary}\label{sg-scalar-germ}
For all sufficiently small $t\in K^\times$, uniformly in
$ct^m\in Q$ and in all permitted test functions,
\begin{align}\label{sg-germ-formula}
 \mathcal O_\Phi(t,c)
 ={}&|t|^{-m^2/2}|s|^{-1/2}
   \sum_{z^s=(-t)^mc}\gamma_m(zt)|z|^{m^2/2}
                     \psi(sz/t)\mathcal O_\Phi(zJ_s).
\end{align}
There is no exclusion of residue characteristics dividing
$2s$; all neighborhoods may depend on $m$, the local field
product, the order, the character and the stated compact sets.
\end{corollary}

\begin{proof}
Lemma~\ref{sg-amplitude-uniform} supplies one support and one
translation lattice for every $A_D$.  Apply
Theorem~\ref{th-stationary} with these data and $D\in Q_0$.
Then apply \eqref{sg-oscillatory-identity} and
\eqref{sg-critical-amplitude}.  The exponents combine as
\[
 -d+\frac m2=-\frac{m^2}{2},\qquad
 \frac m2+\frac{m(m-1)}2=\frac{m^2}{2}.
\]
This gives precisely \eqref{sg-germ-formula}, with every
normalization unchanged.  Uniformity follows because both
the oscillatory reduction and stationary phase used
neighborhoods chosen from the fixed data before $\Phi,t,c$.
\end{proof}

This proves the scalar two-block orbital germ and its full
compact-parameter uniformity.  Section~\ref{hr-section} supplies
the test function for the marked-cokernel coefficient with fixed
initial blocks.  The following assembly applies the scalar
formula to that function and then to the actual kernel.


\subsection{Removing the fixed initial blocks}
\label{hr-section}

We now prove the reduction of a scalar orbital coefficient to a
function of the last blocks. This is the reduction required before
the stationary-phase calculation. All sums below are finite on their
effective support. The proof retains the specified vector and
inverse-coordinate functional of each cokernel.

Fix a configuration $h$ of total size $a$, put $H=Z(h)$ and
$d=\det H\in K^\times$, and fix a finite marked module
$T=(C,e,\lambda)$. For an integral matrix $B$ of size $b$, let
$T_B$ denote $R^b/BR^b$ with its usual last-vector and
last-inverse-coordinate specifications when $B$ is invertible.
Define
\[
 F_T(B)=
 \begin{cases}
 \#\operatorname{Iso}(T_B,T),&B\in M_b(R)\cap\mathrm{GL}_b(K),\\
 0,&\text{otherwise}.
 \end{cases}
\]
In size zero this uses the zero module and zero specifications.
For an integral singular matrix its cokernel cannot be a finite
module: extension from $R$ to $K$ leaves a nonzero cokernel in at
least one field factor, whereas extension of a finite module is zero:
its cardinality annihilates it and is a nonzero integer invertible in $K$.
Here $K=R[1/p]$, since $R$ is open. Thus assigning zero in the
singular case agrees with counting isomorphisms to the finite target.

Fix $s\ge1$. For $A\in\mathrm{GL}_a(K)$ and rectangular matrices
$X\in M_{s,a}(K)$, $Y\in M_{a,s}(K)$, put
\begin{equation}\label{hr-block-matrix}
 B(A,X,Y,G)=
 \begin{pmatrix}A&AY\\ XA&XAY+G\end{pmatrix}
 =\begin{pmatrix}I&0\\X&I\end{pmatrix}
   \begin{pmatrix}A&0\\0&G\end{pmatrix}
   \begin{pmatrix}I&Y\\0&I\end{pmatrix}.
\end{equation}
Write $\ell(X)=X_{1,a}$ and $r(Y)=Y_{a,1}$ when $a>0$,
and put both expressions equal to zero when $a=0$.

For each integral orbital coset $q\in\mathcal I_h$, choose its
representative matrix $A_q=x_qHy_q$ and phase $a_q$. Define
\[
 \mathcal X_q=
  \{X\bmod M_{s,a}(R):XA_q\in M_{s,a}(R)\},\qquad
 \mathcal Y_q=
  \{Y\bmod M_{a,s}(R):A_qY\in M_{a,s}(R)\}.
\]
These are finite sets. For example, the first is the quotient
$M_{s,a}(R)A_q^{-1}/M_{s,a}(R)$ of two compact open additive
lattices; the denominator is contained in the numerator because
$A_q$ is integral. The second is treated in the same way. Put
\begin{equation}\label{hr-Phi}
 \Phi_{h,T,s}(G)=
 \sum_{q\in\mathcal I_h}a_q
    \sum_{X\in\mathcal X_q}\sum_{Y\in\mathcal Y_q}
        \psi(\ell(X)+r(Y))F_T(B(A_q,X,Y,G)).
\end{equation}
The sets indexing this sum do not depend on $G$ or on $T$.

\begin{lemma}\label{hr-well-defined}
Formula \eqref{hr-Phi} is independent of all the indicated
representatives. If $v\in U_s^-(R)$ and $u\in U_s^+(R)$ are
lower and upper unipotent matrices, respectively, then
\begin{equation}\label{hr-invariance}
                  \Phi_{h,T,s}(vGu)=\Phi_{h,T,s}(G).
\end{equation}
Moreover,
\begin{equation}\label{hr-support}
 \Phi_{h,T,s}(G)\ne0
       \quad\Longrightarrow\quad G\in d^{-1}M_s(R).
\end{equation}
\end{lemma}

\begin{proof}
For any integral lower unipotent matrix $v_0$ and upper
unipotent matrix $u_0$, the replacement $B\mapsto v_0Bu_0$
preserves integrality and invertibility. It preserves the marked
module: the induced cokernel map sends $[x]$ to $[v_0x]$,
the last vector is fixed by $v_0$, and the last row of $u_0^{-1}$
is the last coordinate row. Hence $F_T(v_0Bu_0)=F_T(B)$.

Replacing $X,Y$ by $X+X_0,Y+Y_0$ with integral $X_0,Y_0$
multiplies \eqref{hr-block-matrix} on the left and right by
integral unipotent matrices. The corner phase changes by an
element of $R$, on which $\psi$ is trivial. This proves
independence of the rectangular representatives.
If the representative of $q$ changes $A_q$ to $v_0A_qu_0$,
the bijection $X\mapsto Xv_0^{-1}$, $Y\mapsto u_0^{-1}Y$
identifies the summands. The block matrix is then multiplied
by $\operatorname{diag}(v_0,I_s)$ and
$\operatorname{diag}(u_0,I_s)$. The corner entries and the
phase $a_q$ are unchanged modulo $R$.

For \eqref{hr-invariance}, use the bijections
$X\mapsto vX$ and $Y\mapsto Yu$. Then
\[
 B(A,vX,Yu,vGu)=\operatorname{diag}(I_a,v)
                  B(A,X,Y,G)\operatorname{diag}(I_a,u).
\]
The first row of a lower unipotent matrix and the first column
of an upper unipotent matrix are coordinate vectors. Therefore
$\ell(vX)=\ell(X)$ and $r(Yu)=r(Y)$, which proves the identity.

If the finite sum is nonzero, some block matrix
$B=B(A_q,X,Y,G)$ is integral. Its upper-left block $A=A_q$,
both $XA$ and $AY$, and its lower-right block $XAY+G$ are
integral. Also $\det A=d\in R$ and $\operatorname{adj}(A)$
is integral. The identity
\[
 dG=d(XAY+G)-(XA)\operatorname{adj}(A)(AY)
\]
shows that $dG$ is integral. For $a=0$, $d=1$ and the same
assertion follows directly from $B=G$.
\end{proof}

We next give the perturbation estimate that preserves the
functional as well as the underlying module.

\begin{lemma}\label{hr-perturbation}
Let $B\in\mathrm{GL}_b(K)$, write $\delta=\det B$, and
suppose that $P\in M_b(R)$ and
$P\in\delta^2pM_b(R)$. Then
\[
                         F_T(B+P)=F_T(B)
\]
for every fixed finite marked module $T$.
If $B$ is integral, the identification is the identity on the
ambient module $R^b$: the two image lattices are equal and
$(B+P)^{-1}-B^{-1}$ is integral.
\end{lemma}

\begin{proof}
If $B$ is not integral, adding the integral matrix $P$ cannot
make it integral, so both sides are zero. Assume $B$ integral
and write $P=\delta^2pP_0$ with $P_0$ integral. Put
\[
 A=I+B^{-1}P=I+p\delta\operatorname{adj}(B)P_0.
\]
Since $\delta\in R$, this is a matrix in $I+pM_b(R)$.
Every such matrix is invertible over $R$: its geometric-series
inverse converges, because the powers of a matrix in $pM_b(R)$
have entries in $p^jR$ and these sets tend to zero. Compactness
of $R$ makes it complete. This argument also applies at $p=2$
and to a nonmaximal order.

Now $B+P=BA$, so $(B+P)R^b=BR^b$. Moreover
\[
 (B+P)^{-1}-B^{-1}
   =-pA^{-1}\operatorname{adj}(B)P_0\operatorname{adj}(B)
                                                     \in M_b(R).
\]
The identity on the common cokernel preserves the specified
vector and every inverse-coordinate functional modulo $R$.
It consequently identifies the two sets of marked
isomorphisms to $T$.
\end{proof}

\begin{corollary}\label{hr-local-constancy}
Let $Q\subset K^\times$ be compact. There is an additive
lattice $L\subset R$, depending only on $R,d,Q$, such that
\[
 \Phi_{h,T,s}(G+E)=\Phi_{h,T,s}(G)
 \quad\text{if }\det G\in Q\cup(-Q),\quad E\in M_s(L).
\]
The choice is independent of $T$ and $s$.
\end{corollary}

\begin{proof}
Choose $k\ge0$ such that $p^kd^{-2}q^{-2}\in R$ for every
$q\in Q$, using compactness of the inverse set; take
$L=p^{k+1}R$. In every term of \eqref{hr-Phi}, the full block
matrix has determinant $d\det G$. Replacing $G$ by $G+E$
adds only the lower-right block $E$, an integral perturbation
whose entries divided by $(d\det G)^2$ belong to $pR$.
Lemma~\ref{hr-perturbation} therefore applies term by term.
The indexing sets and phases do not change.
\end{proof}

\begin{proposition}\label{hr-sum-identity}
Let $g$ be any configuration of total size $s$, and let $h\,g$
denote concatenation of their lists of blocks. Let
\[
 \mathcal Q_g=L_g(R)\backslash L_g(K)\times U_s(K)/U_s(R)
\]
be its full orbital coset space. For $q\in\mathcal Q_g$, write
$G_q=x_qZ(g)y_q$ and let $a_q$ be its phase. Then
\begin{equation}\label{hr-orbital-sum}
 \mathcal B(h\,g;T)
           =\sum_{q\in\mathcal Q_g}a_q\Phi_{h,T,s}(G_q).
\end{equation}
The support of the sum on the right is finite. The matrices
$G_q$ in it are not required to be integral.
\end{proposition}

\begin{proof}
Split the coordinates into the first $a$ and last $s$.
Every lower matrix for the concatenated configuration has
a unique expression
\[
 \begin{pmatrix}v_h&0\\Xv_h&v_g\end{pmatrix},
 \qquad v_h\in L_h(K),\quad v_g\in L_g(K),
\]
and every upper matrix has a unique expression
\[
 \begin{pmatrix}u_h&u_hY\\0&u_g\end{pmatrix},
 \qquad u_h\in U_a(K),\quad u_g\in U_s(K).
\]
Choose representatives for the four diagonal cosets. Left
integral multiplication on the lower side then leaves exactly
$X\bmod M_{s,a}(R)$; right integral multiplication on the
upper side leaves exactly $Y\bmod M_{a,s}(R)$.
For clarity, once the diagonal representatives agree, the
quotient of two lower matrices is
$\left(\begin{smallmatrix}I&0\\X'-X&I\end{smallmatrix}\right)$;
the corresponding quotient on the upper side is
$\left(\begin{smallmatrix}I&Y'-Y\\0&I\end{smallmatrix}\right)$.
This proves both existence and uniqueness of the resulting
coset parametrization.

Their orbital matrix is exactly
$B(v_hHu_h,X,Y,v_gZ(g)u_g)$.
The adjacent-entry phase factors into the phases of the first
and last configurations, multiplied by
$\psi(\ell(X)+r(Y))$. Indeed the only adjacent entry crossing
the block boundary is the corner entry, and
$\ell(Xv_h)=\ell(X)$, $r(u_hY)=r(Y)$ by triangularity.
No determinant or measure factor enters this bijection of
discrete cosets.

If the full orbital matrix is integral, its first block is
integral, and $XA$ and $AY$ are integral. Thus the nonzero
terms are exactly among those defining \eqref{hr-Phi}.
There are only finitely many full integral orbital cosets by
Lemma~\ref{orb-finite-orbits}. The bijection just proved
therefore has finite effective support, so it can be regrouped
without an infinite-sum interchange. Grouping first by the
coset of the last configuration gives
\eqref{hr-orbital-sum}.
\end{proof}

Applying Proposition~\ref{hr-sum-identity} with
$g=[m,t][1,c]$ and with $g=[m+1,z]$ gives the two scalar
identities required to remove the fixed initial blocks. The
same function $\Phi_{h,T,m+1}$ occurs in both. It is invariant
under integral lower and upper unipotent multiplication, is
supported in $d^{-1}M_{m+1}(R)$, and is locally constant on
the determinant range $Q\cup(-Q)$ with a lattice independent
of the target module. These are precisely the required
hypotheses on the matrix function in the subsequent
oscillatory reduction.

\begin{remark}[Nonintegral matrices in the last blocks]
Their inclusion in \eqref{hr-orbital-sum} is essential.
Take $K=\mathbb Q_3$, $R=\mathbb Z_3$, the standard character
with $\psi(1/3)=\zeta_3$, one initial block $H=(3)$,
one last block $G=(-1/3)$, and the zero marked target module.
The full orbital matrices are
\[
 \begin{pmatrix}3&3y\\3x&3xy-1/3\end{pmatrix}.
\]
Integrality forces $x=u/3$, $y=v/3$ modulo $R$, with
$uv\equiv1\pmod3$. Their determinant is $-1$, so every
integral matrix has zero cokernel and contributes its phase.
The only pairs are $(u,v)=(1,1),(2,2)$; their sum is
\[
             \zeta_3^2+\zeta_3^4=-1.
\]
The last scalar $-1/3$ is not integral. Restricting the right
side of \eqref{hr-orbital-sum} to integral last-block matrices
would incorrectly give zero. The actual source retains the
full coset sum.
\end{remark}

\begin{remark}[Preserving the module does not preserve the functional]
Take $R=\mathbb Z_p$, $B=(p^2)$ and $P=(p^3)$.
The image ideals of $B$ and $B+P=p^2(1+p)$ are equal, but
their inverse entries differ by
$-1/(p(1+p))\notin R$. The two cyclic cokernels with marked
vector $[1]$ are the same underlying module; any isomorphism
preserving this generator is the identity. Their specified
functionals differ on that generator, so they are not
isomorphic as marked modules. This explains why a perturbation
of size merely $p\det(B)R$ is insufficient. The squared
determinant condition in Lemma~\ref{hr-perturbation} retains
the information that is actually used in the germ formula.
\end{remark}


\subsection{The germ formula for the actual marked-cokernel kernel}
\label{ga-section}

We now combine the initial-block reduction and the scalar
oscillatory calculation.  All objects in the statement are
the actual finite orbital sums defined above.

\begin{theorem}\label{ga-actual-germ}
Fix a block configuration $h$, an integer $m\ge1$, and a
compact set $Q\subset K^\times$.  Put $s=m+1$.
There is a neighborhood $V$ of zero such that for all
$t,c\in K^\times$ with $t\in V$ and $ct^m\in Q$, and every
comparison configuration $w'$, one has
\begin{align}\label{ga-actual-kernel-formula}
 \mathcal K(W_h(t,c),w')
 ={}&|t|^{-m^2/2}|s|^{-1/2}
       \sum_{z^s=(-t)^mc}\gamma_m(zt)|z|^{m^2/2}
          \psi(sz/t)\mathcal K(h[s,z],w').
\end{align}
The size, block decomposition and parameters of $w'$ are
unrestricted.  The neighborhood is chosen before $t,c,w'$.
For a fixed finite collection of possible $h$, one neighborhood
works for every member of that collection.
\end{theorem}

\begin{proof}
Write $d_h=\det Z(h)$.  For each fixed finite marked target
module $T$, Section~\ref{hr-section} constructs the test function
$\Phi_{h,T,s}$.  Lemma~\ref{hr-well-defined} makes it invariant
under integral lower and upper unipotent multiplication and
supports it on $d_h^{-1}M_s(R)$.  Corollary~\ref{hr-local-constancy}
provides one additive lattice on which it is locally constant
when the determinant belongs to $Q\cup(-Q)$.
These bounds are independent of the target module.
They are precisely the three test-function assumptions of
Section~\ref{sec:theta-scalar-germ}.

Apply Corollary~\ref{sg-scalar-germ}.  Proposition~\ref{hr-sum-identity}
identifies its two orbital coefficients, giving
\begin{align}\label{ga-marked-formula}
 \mathcal B(W_h(t,c);T)
 ={}&|t|^{-m^2/2}|s|^{-1/2}
       \sum_{z^s=(-t)^mc}\gamma_m(zt)|z|^{m^2/2}
               \psi(sz/t)\mathcal B(h[s,z];T).
\end{align}
All sums here include every orbital coset on their finite
effective supports, including nonintegral last-block matrices
in the initial-block reduction.

Proposition~\ref{gu-uniform-germ} already shows that these
identities against fixed marked targets imply the germ
uniformly against all comparison configurations.  In the
present proof one can also see the uniformity directly:
the test-function bounds and the scalar-germ neighborhood
were independent of $T$.  For any $w'$, apply
\eqref{ga-marked-formula} to each marked module $C_{Z_{q'}}$
in its finite integral orbital expansion, multiply by
$\overline{a_{q'}}$, and sum.  Formula
\eqref{gu-kernel-coefficients} gives
\eqref{ga-actual-kernel-formula}.  The root sum is finite,
so the interchange introduces no convergence requirement.
This argument does not require the size of $w'$ to equal
the size of $W_h(t,c)$.

Finally, intersect the neighborhoods for finitely many $h$.
This retains a neighborhood of zero and proves the last assertion.
\end{proof}

\begin{corollary}\label{ga-actual-truncation}
For the actual marked-cokernel kernel, the truncated orbital
identity \eqref{orb-exact-cut} holds under the local hypotheses
stated in Theorem~\ref{orb-cut-identity}, without an additional
germ assumption.
\end{corollary}

\begin{proof}
The Gram formula gives Hermitian symmetry.
Proposition~\ref{orb-local-constancy-proved} supplies the
required local constancy, Corollary~\ref{orb-support-proved}
supplies the inverse-parameter support, and
Theorem~\ref{ga-actual-germ} supplies the last hypothesis.
Apply Theorem~\ref{orb-cut-identity}.
\end{proof}

This closes the local orbital germ and truncated orbital identity.
It does not assert a global theta energy bound: the global
arithmetic and automorphic estimates remain distinct statements.


\section{The single-block identity and the local majorant induction}
\label{tp-section}

We retain the product of local fields $K$, compact open self-dual
subring $R$, character $\psi$, and orbital kernel $\mathcal K$ of
Section~\ref{sec:theta-orbital-identity}. Haar measure gives $R$
volume one. Write $[a]$ for the class of $a$ in $K/R$, and also
write $\psi$ for the induced character of that quotient. For an
integral invertible matrix $Z$ of size $n$, put
\[
 C_Z=R^n/ZR^n,\qquad e_Z=[e_n],\qquad
 \lambda_Z([v])=[(Z^{-1}v)_n].
\]
In size zero the module, marked vector and functional are zero.
Isomorphisms of marked modules preserve both the vector and the
functional; an unqualified isomorphism below preserves only the
$R$-module structure. All sums over finite modules are ordinary
finite sums. For a block configuration $h$, define
\[
 c(h)=|\det Z(h)|_K^{-1},\qquad N_0(h,h')=\mathcal K(h,h'),
\]
\[
 N_s(h,h';t,t')=|tt'|_K^{s(s-1)/2}
           \mathcal K(h[s,t],h'[s,t'])\quad(s\ge1).
\]
Here $h[s,t]$ means that the block $tJ_s$ is appended to $h$,
where $J_s$ reverses the coordinates. These are the source's
normalizations\cite{alpo-uniform-source}; in particular $N_1$ has no additional modulus
factor. We prove the exact single-block identity and then deduce
the local majorant theorem from the orbital germ proved in
Theorem~\ref{ga-actual-germ}.

\subsection{Duality and a bordered cokernel}

\begin{lemma}\label{tp-dual}
For $C=R^n/ZR^n$ as above, put $C^*=\operatorname{Hom}_R(C,K/R)$.
Then $\#C^*=\#C=|\det Z|_K^{-1}$, and
\[
 \sum_{\alpha\in C^*}\psi(\alpha(x))
       =\#C\,1_{x=0},\qquad
 \sum_{x\in C}\psi(\alpha(x))=\#C\,1_{\alpha=0}.
\]
\end{lemma}

\begin{proof}
An $R$-linear functional on $R^n$ with values in $K/R$ is
represented by a row $\xi\in K^n$ modulo $R^n$. It vanishes on
$ZR^n$ exactly when $\xi Z\in R^n$. Consequently
$C^*=R^nZ^{-1}/R^n$ as an additive group. Multiplication by $Z$
identifies this with the quotient of $R^n$ by $R^nZ$, whose
order is $|\det Z|_K^{-1}$ by the Haar-volume calculation in
Lemma~\ref{th-coker}. The same calculation gives the order of $C$.

If $x=[v]\ne0$, then $Z^{-1}v$ has a coordinate outside $R$.
Self-duality of $R$ gives an $r\in R$ on which its pairing by
$\psi$ is nontrivial. The corresponding row $r e_i^tZ^{-1}$
represents an $\alpha\in C^*$ with $\psi(\alpha(x))\ne1$.
Translation by that $\alpha$ makes the first sum zero. If
$\alpha\ne0$, choose $x$ with $\alpha(x)\ne0$ in $K/R$.
Self-duality gives $r\in R$ with $\psi(\alpha(rx))\ne1$;
translation by $rx$ makes the second sum zero. In the zero cases
every summand is one. This proof uses no principal-ideal property.
\end{proof}

For rows $\xi$, columns $\eta$, and $c\in K^\times$, set
\begin{equation}\label{tp-border}
 B(Z,\xi,\eta,c)=
 \begin{pmatrix}Z&Z\eta\\ \xi Z&\xi Z\eta+c\end{pmatrix}.
\end{equation}
Suppose $Z$ is integral and invertible. The conditions
$\xi Z\in R^n$ and $Z\eta\in R^n$ identify $\xi\bmod R^n$
with $\alpha\in C_Z^*$ and $\eta\bmod R^n$ with
$\beta=[Z\eta]\in C_Z$. The remaining condition for
\eqref{tp-border} to be integral is
\begin{equation}\label{tp-integral-condition}
                         \alpha(\beta)=-[c].
\end{equation}
Every pair $(\alpha,\beta)$ has such representatives. These
assertions follow either by applying $Z^{-1}$, or by evaluating
the functional on the standard basis of $R^n$.

\begin{lemma}\label{tp-border-isomorphisms}
Let $Z,Z'$ be integral invertible matrices, possibly of different
sizes. Fix $c\in K^\times$ and pairs $(\alpha,\beta)$ and
$(\alpha',\beta')$ satisfying \eqref{tp-integral-condition}.
There is a bijection between marked isomorphisms of the two
bordered cokernels and the unqualified isomorphisms
$\theta:C_Z\to C_{Z'}$ satisfying
\[
                 \alpha'\theta=\alpha,\qquad
                 \theta\beta=\beta'.
\]
\end{lemma}

\begin{proof}
For $C=C_Z$, form the $R$-module
\[
 E_\alpha=\{(x,k)\in C\times K:[k]=-\alpha(x)\}.
\]
The map $(v,r)\mapsto([v],r-\xi v)$ from $R^n\oplus R$
onto $E_\alpha$ has kernel consisting exactly of
$(Zu,\xi Zu)$ with $u\in R^n$. Surjectivity follows by lifting
$x$ to $v$ and using $[k]=-\alpha(x)$ to see that $k+\xi v\in R$.
Thus it identifies the quotient by the first $n$ columns of
\eqref{tp-border} with $E_\alpha$. The last column maps to
$g=(\beta,c)$, and the last standard vector maps to $e=(0,1)$.
Block inversion gives the last inverse coordinate
$c^{-1}(r-\xi v)$. Hence the bordered marked module is
\[
 E_\alpha/Rg,\qquad [e],\qquad [(x,k)]\longmapsto[k/c].
\]
Here $g\in E_\alpha$ by \eqref{tp-integral-condition}, and
$Rg$ is a free copy of $R$ because $c\in K^\times$.

An isomorphism of these marked quotients lifts uniquely to an
isomorphism $E_\alpha\to E_{\alpha'}$ preserving $k$. Indeed,
given a representative of the image class, preservation of the
functional says that its second coordinate differs from $k$ by
an element of $cR$. Subtracting the corresponding multiple of
$g'=(\beta',c)$ gives a unique lift with second coordinate $k$.
Uniqueness proves additivity and $R$-linearity, and the same
construction for the inverse proves bijectivity. The lift sends
$g$ to $g'$ and $e$ to $e'$, by their quotient classes and their
respective second coordinates $c$ and $1$.

Quotienting $E_\alpha$ by $Re$ recovers $C$, with functional
$-\alpha$ determined by $k\bmod R$. The lifted isomorphism
therefore induces $\theta$ with $\alpha'\theta=\alpha$ and
$\theta\beta=\beta'$. Conversely such a $\theta$ gives
$(x,k)\mapsto(\theta x,k)$ on the two $E$'s. It preserves
$e,g$ and descends to a marked isomorphism of the bordered
quotients. These constructions are inverse, proving the claim.
\end{proof}

\subsection{The exact single-block sum}

For marked modules $(C,e,\lambda)$ and $(C',e',\lambda')$,
define
\begin{equation}\label{tp-finite-inner}
 I_c(C,C')=
 \sum_{\theta:C\simeq C'}\ 
 \sum_{\substack{\alpha\in C^*,\ \beta\in C\\
                         \alpha(\beta)=-[c]}}
 \psi\bigl(\alpha(e-\theta^{-1}e')
                    +\lambda(\beta)-\lambda'(\theta\beta)\bigr).
\end{equation}
The outer isomorphisms are unqualified. The finite expression
depends only on $c\bmod R$ and is zero for all but finitely many
such classes.

\begin{proposition}\label{tp-explicit-single}
For configurations $h,h'$ and $c\in K^\times$,
\begin{equation}\label{tp-single-formula}
 N_1(h,h';c,c)=
 \sum_{q\in\mathcal I_h}\sum_{q'\in\mathcal I_{h'}}
       a_q\overline{a_{q'}}\,I_c(C_{Z_q},C_{Z_{q'}}),
\end{equation}
where each cokernel on the right retains its original marked
vector and last inverse-coordinate functional.
\end{proposition}

\begin{proof}
Fix lower and upper representatives $(x,y)$ of an orbital class
for $h$, so that $Z=xZ(h)y$. Representatives over this class
after appending the size-one block are
\[
 x_+=\begin{pmatrix}x&0\\ \xi x&1\end{pmatrix},\qquad
 y_+=\begin{pmatrix}y&y\eta\\0&1\end{pmatrix}.
\]
The rows $\xi$ and columns $\eta$ range independently over
$K^n/R^n$. To verify this parametrization, project a lower or
upper representative to its upper left block. Use an integral
change to put that block equal to the chosen representative.
The remaining integral changes add an arbitrary integral row
to $\xi$, or an arbitrary integral column to $\eta$; conversely
only such changes keep the upper left representative fixed.
This proves both existence and uniqueness of the indicated
coset coordinates.

Matrix multiplication gives
$x_+\operatorname{diag}(Z(h),c)y_+=B(Z,\xi,\eta,c)$.
Since the last column of $x$ and last row of $y$ are the last
standard vector and covector, the additional phase is
\[
 \psi(\xi_n+\eta_n)
       =\psi\bigl(\alpha(e_Z)+\lambda_Z(\beta)\bigr).
\]
For size zero both extra terms are zero. Integrality is precisely
\eqref{tp-integral-condition}; if the bordered matrix is integral,
its upper left matrix $Z$ is integral as well. Lemma~\ref{tp-border-isomorphisms}
now replaces the marked-isomorphism count by the number of
$\theta$ satisfying $\alpha'=\alpha\theta^{-1}$ and
$\beta'=\theta\beta$. The difference of the two phases is
exactly the argument of $\psi$ in \eqref{tp-finite-inner}.
Summing over the remaining $\alpha,\beta$ proves
\eqref{tp-single-formula}. The base orbital sets and all modules
are finite, so every rearrangement is finite.
\end{proof}

\begin{theorem}\label{tp-parseval}
The function $c\mapsto N_1(h,h';c,c)$ is $R$-periodic on
$K^\times$ and extends to a finitely supported function on
$K/R$. For any choice of unit representative $c_x$ of each
additive class $x\in K/R$,
\begin{equation}\label{tp-parseval-eq}
 \sum_{x\in K/R}N_1(h,h';c_x,c_x)
                       =c(h)c(h')N_0(h,h').
\end{equation}
\end{theorem}

\begin{proof}
Periodicity and finite support follow from
\eqref{tp-finite-inner} and \eqref{tp-single-formula}.
Each additive class has a representative in $K^\times$:
adding a suitable element of the open subgroup $R$ avoids
the finitely many coordinate hyperplanes of the product $K$.
For a fixed $\theta$ in \eqref{tp-finite-inner}, summing over
$c\bmod R$ removes the condition $\alpha(\beta)=-[c]$, since
each pair $(\alpha,\beta)$ determines exactly one such class.
The remaining sum factors as
\[
 \left(\sum_{\alpha\in C^*}
       \psi(\alpha(e-\theta^{-1}e'))\right)
 \left(\sum_{\beta\in C}
       \psi\bigl((\lambda-\lambda'\theta)(\beta)\bigr)\right).
\]
Lemma~\ref{tp-dual} makes this $(\#C)^2$ if $\theta$ preserves
both marks, and zero otherwise. An isomorphism $C\simeq C'$
implies $\#C=\#C'$. Therefore
\[
 \sum_{c\bmod R}I_c(C,C')
       =\#C\,\#C'\,\#\operatorname{Iso}((C,e,\lambda),(C',e',\lambda')).
\]
Every integral orbital matrix for $h$ has cokernel order $c(h)$,
and similarly for $h'$. Substitution in \eqref{tp-single-formula}
gives \eqref{tp-parseval-eq}.
\end{proof}

\begin{example}\label{tp-factor-example}
The factors $c(h)c(h')$ cannot be deleted. Take
$K=\mathbb Q_2$, $R=\mathbb Z_2$, and $h=h'=[1,2]$.
There is just one base orbital matrix, namely $(2)$; its
marked cokernel is $C=\mathbb Z/2\mathbb Z$, with $e=1$
and $\lambda(x)=[x/2]$. It has only the identity module
automorphism. Write $\alpha_r(x)=[rx/2]$, $r=0,1$.
The phase in \eqref{tp-finite-inner} is always one.
Three pairs $(r,\beta)$ have $r\beta=0$ and one has
$r\beta=1$. Thus $N_1$ takes values $3$ on the zero coset,
$1$ on $1/2+\mathbb Z_2$, and zero elsewhere. Its sum is
$4=c(h)^2N_0(h,h)$, whereas $N_0(h,h)=1$.
This is a counterexample to an incorrectly normalized identity,
not to \eqref{tp-parseval-eq} or to the submitted theorem.
\end{example}

\subsection{Support and the unrestricted diagonal}

\begin{lemma}\label{tp-minor-support}
For $s\ge1$, a nonzero $N_s(h,h';t,t')$ implies
\[
                 (\det Z(h))t\in R,
                 \qquad(\det Z(h'))t'\in R.
\]
If both $Z(h)$ and its inverse are integral, and likewise for
$h'$, then $N_0(h,h')=1$.
\end{lemma}

\begin{proof}
Let $r$ be the total size of $h$ and $n=r+s$. In an integral
orbital matrix $x\operatorname{diag}(Z(h),tJ_s)y$, take the
minor in the first $r+1$ rows and columns $1,\ldots,r,n$.
Multiplication on the left restricts to a lower unipotent
matrix on these rows and has determinant one. Before that
multiplication, the last selected row is zero in the first $r$
columns and has last entry $t$; the upper left $r$ by $r$
block is $Z(h)$ times an upper unipotent matrix. The minor is
therefore exactly $(\det Z(h))t$. Integrality puts it in $R$.
If the kernel is nonzero, at least one integral orbital matrix
occurs on each side, proving both assertions of support.

For the last assertion every reversing block of $Z(h)$ is
invertible over $R$. An integral matrix $xZ(h)y$ recovers
integral $x,y$ by successive block elimination. Indeed, the
first diagonal block is $t_1J_{s_1}$ times an upper unipotent
matrix, so the latter is integral and invertible over $R$.
The other blocks in its row and column recover the first
upper and lower off-diagonal blocks integrally. Their Schur
complement is integral and repeats the argument on the
remaining blocks. Thus only the integral orbital coset
occurs. Its phase is one and its cokernel is zero. The same
holds for $h'$, and the zero marked modules have one
isomorphism. This also covers empty configurations.
\end{proof}

For equal-size configurations $h,h'$ put
\[
 U_m(h,h';t,t')=|tt'|_K^{m(m+1)/2}
   \sum_{c\bmod R}\mathcal K(h[m,t][1,c],h'[m,t'][1,c]),
\]
using unit representatives. Theorem~\ref{tp-parseval} applies
to the configurations $h[m,t]$ and $h'[m,t']$. Since
$c(h[m,t])=c(h)|t|_K^{-m}$, it gives the exact identity
\begin{equation}\label{tp-unrestricted}
             U_m(h,h';t,t')=c(h)c(h')N_m(h,h';t,t').
\end{equation}

Fix a finite collection $\mathcal F$ of equal-size configurations.
Write $D_s=\operatorname{diag}_{h\in\mathcal F}(c(h)^s)$,
and regard kernels as matrices on $\mathcal F$.
Suppose a positive $R$-periodic majorant $\Theta_m$ satisfies
\[
 N_m(t,t)\le\Theta_m(t),\qquad
 \sum_{x\in K/R}\Theta_m(x)\le D_mN_0D_m,
\]
with finite support on $K/R$. Each positive summand is at most
the whole sum. Conjugating by $D_1$ and using
\eqref{tp-unrestricted} proves
\begin{equation}\label{tp-diagonal}
                     U_m(t,t)\le D_{m+1}N_0D_{m+1}
                         \quad(t\in K^\times).
\end{equation}
In particular the same bound holds for every convex average
of these unrestricted diagonals. No new analytic estimate is
required for this step.

\subsection{Completion of the induction from the germ}

\begin{theorem}\label{tp-majorant}
For every $s\ge1$ and finite family $\mathcal F$ of equal-size
configurations there is an $R$-periodic matrix-valued function
$\Theta_s$ with
\begin{equation}\label{tp-majorant-properties}
 0\le N_s(t,t)\le\Theta_s(t)\quad(t\in K^\times),\qquad
 \sum_{x\in K/R}\Theta_s(x)\le D_sN_0D_s.
\end{equation}
It is zero outside the fixed compact open $R$-module
\[
                 L_c=R+\sum_{h\in\mathcal F}(\det Z(h))^{-1}R.
\]
Its $(h,h')$ entry vanishes when
$|\det Z(h)|_K\ne|\det Z(h')|_K$.
If all $Z(h)$ are invertible over $R$, then
\[
                      0\le\Theta_s(x)\le1_R(x)N_0.
\]
The construction uses the germ uniformly on each fixed compact
parameter set. Theorem~\ref{ga-actual-germ} supplies exactly
that uniformity, independently of the comparison configuration.
\end{theorem}

\begin{proof}
For $s=1$, take $\Theta_1$ to be the periodic extension of
$N_1(t,t)$. Positivity is the Gram identity, its exact mass is
Theorem~\ref{tp-parseval}, and its support is
Lemma~\ref{tp-minor-support}. Determinant support gives the
claimed zero entries. This proves the base case.

Suppose the result is known for $s=m\ge1$. Fix a compact open
set $\Omega\subset K^\times$ invariant under $(m+1)$st roots
of unity. Choose additive lattices $L_1\subset R\subset L_2$
so that $\Omega,L_c\subset L_1^\perp$, and $\Omega$ and every
entry of $N_{m+1}(z,z)$ on it are invariant under translation
by $L_2^\perp$. Such lattices exist by compactness and the
uniform local constancy already proved in
Section~\ref{sec:theta-orbital-identity}; shrinking a lattice
to a power of $p$ times $R$ preserves these conditions.
Choose representatives $A$ of $L_2/R$, $B$ of $R/L_1$, and
put $a=\#A$, $b=\#B$.

For all shifts $\sigma=\rho+v$, with $\rho\in A$ and $v\in B$,
choose one sufficiently small $t\in K^\times$ and set
$t_\sigma^{-1}=t^{-1}+\sigma$. Their absolute values equal
$|t|_K$. The germ consequence in
Section~\ref{sec:theta-orbital-identity} identifies the truncated
Gram matrix by
\begin{equation}\label{tp-truncated-average}
 G_c((\sigma,h),(\sigma',h'))
       =\int_\Omega N_{m+1}(h,h';z,z)
                             \psi(z(\sigma-\sigma'))\,dz.
\end{equation}
The unrestricted Gram matrix is
$G((\sigma,h),(\sigma',h'))=U_m(h,h';t_\sigma,t_{\sigma'})$.
Both are finite sums of positive Gram matrices, with
$0\le G_c\le G$. Finiteness follows also directly from
Theorem~\ref{tp-parseval}. Thus Lemma~\ref{th-cut-contraction}
realizes them as $\langle u_{\sigma,h},u_{\sigma',h'}\rangle$
and $\langle Tu_{\sigma,h},Tu_{\sigma',h'}\rangle$ for a
contraction $T$.

By \eqref{tp-unrestricted} and the inverse-parameter support
proved earlier, these vectors are orthogonal when their
$\rho$ indices differ: a nonzero entry would imply
$\sigma-\sigma'\in R$, contrary to the choice of $A$.
Use the convention that the inner product is antilinear in
its first variable, and define
\[
 v_{\rho,z,h}=\frac1b\sum_{v\in B}\psi(zv)u_{\rho+v,h},\qquad
 v_{z,h}=\frac1a\sum_{\rho\in A}\psi(z\rho)v_{\rho,z,h},
\]
\[
 \widetilde\Theta(z)_{hh'}=
       \frac1a\sum_{\rho\in A}
             \langle v_{\rho,z,h},v_{\rho,z,h'}\rangle.
\]
These expressions are considered for $z\in L_1^\perp$, and
$\widetilde\Theta$ is extended by zero to the other classes
of $K/R$. The values of $\widetilde\Theta$ are $R$-periodic,
since $v\in R$ and $\psi$ is trivial on $R$.

For $z_1\in\Omega$, substitute \eqref{tp-truncated-average}
in $a\langle Tv_{z_1,h},Tv_{z_1,h'}\rangle$. The resulting
factor in the integrand is
\[
 \frac{a}{(ab)^2}
  \left|\sum_{\sigma\in A+B}\psi((z-z_1)\sigma)\right|^2.
\]
The set $A+B$ is a system of representatives of $L_2/L_1$.
Because $z,z_1\in L_1^\perp$, character orthogonality makes
this factor $a$ on $z-z_1\in L_2^\perp$ and zero otherwise.
The whole coset $z_1+L_2^\perp$ lies in $\Omega$ and the
kernel is constant there. Its volume is $1/a$, since
$\operatorname{vol}(L_2)=a$ and the measure is self-dual.
Consequently
\[
 a\langle Tv_{z_1,h},Tv_{z_1,h'}\rangle
                         =N_{m+1}(h,h';z_1,z_1).
\]
Proposition~\ref{th-fourier-average}, with characters
$v\mapsto\psi(-zv)$, now gives
\[
 \widetilde\Theta(z_1)\ge N_{m+1}(z_1,z_1),\qquad
 \sum_{z\in L_1^\perp/R}\widetilde\Theta(z)
   =\frac1{ab}\sum_{\rho\in A}\sum_{v\in B}
                        U_m(t_{\rho+v},t_{\rho+v})
   \le D_{m+1}N_0D_{m+1},
\]
where the final inequality is \eqref{tp-diagonal}.
Replace $\widetilde\Theta$ by $1_{L_c}\widetilde\Theta$.
This keeps positivity and only decreases its total mass.
It still dominates the kernel on $\Omega$, because that
kernel is zero outside $L_c$ by Lemma~\ref{tp-minor-support}.
Every Gram entry between heads with different determinant
absolute values vanishes, since all $|t_\sigma|_K$ are equal.
Thus the required zero entries persist through the averaging.

Finally exhaust $K^\times$ by compact open products of annuli,
stable under all roots of unity. The fixed module $L_c$ works
for every annulus. Proposition~\ref{th-local-limit} gives
one majorant on all units with the same mass and support,
and with the same zero entries. This completes the induction.
If all heads are invertible over $R$, then $c(h)=1$ and
$L_c=R$; the one possible support coset gives the last
inequality directly from the total-mass bound. The finite
target reduction is Proposition~\ref{gu-uniform-germ}.
\end{proof}

The scalar germ, averaged diagonal estimate and all finite
averaging steps of the local induction are therefore proved for
the actual kernel. A local majorant alone does not prove the
four global orbital estimates \eqref{th-four-estimates}; their
uniform arithmetic and automorphic estimates remain distinct
tasks in the complete argument.


\section{Global ideal sums in the theta estimates}
\label{gt-section}

The local results in the preceding sections leave several global
estimates to be proved.  We begin with a complete reconstruction of the
sum over the dual lattices in the scalar orbital calculation.  This
part admits a shorter proof and a stronger estimate than the one
stated in the source~\cite{alpo-uniform-source}.  All the constants in this section are explicit
and independent of the number field and its discriminant.

\subsection{Ideal norms and separation of algebraic integers}

Let $E$ be a totally complex number field of degree $2f$, where $f\geq1$.
Write $\mathcal O_E$ for its ring of integers and choose embeddings
$\sigma_1,\ldots,\sigma_f:E\longrightarrow\mathbb C$, one from each
complex conjugate pair.  A nonzero fractional ideal $\mathfrak a$ is a
nonzero finitely generated $\mathcal O_E$-submodule of $E$ for which
$m\mathfrak a\subseteq\mathcal O_E$ for some positive integer $m$.
Its absolute norm is
\[
 N\mathfrak a=\frac{[\mathcal O_E:m\mathfrak a]}{m^{2f}}.
\]
This does not depend on $m$.  Equivalently it is the ratio of the
covolumes of $\mathfrak a$ and $\mathcal O_E$ under the embedding
$\alpha\mapsto(\sigma_v\alpha)_{v=1}^f$ into $\mathbb C^f$.
The choice of Lebesgue measure cancels in this ratio.

We use the usual multiplicativity of ideal norms.  To recall its
content, every nonzero fractional ideal has a unique factorization
$\prod_{\mathfrak p}\mathfrak p^{n_{\mathfrak p}}$, with finitely many
nonzero integers $n_{\mathfrak p}$, and its norm is
$\prod_{\mathfrak p}|\mathcal O_E/\mathfrak p|^{n_{\mathfrak p}}$.
For integral ideals this follows by the Chinese remainder theorem
and the filtration by successive powers of each prime ideal; each
successive quotient has the cardinality of its residue field.
Clearing denominators gives the fractional-ideal assertion.

\begin{lemma}\label{gt-product-minimum}
If $0\ne\alpha\in\mathfrak a$, then
\begin{equation}\label{gt-product-minimum-eq}
 \prod_{v=1}^f|\sigma_v\alpha|\geq\sqrt{N\mathfrak a}.
\end{equation}
The same inequality applies to $\alpha-\beta$ for distinct
$\alpha,\beta\in\mathfrak a$.
\end{lemma}

\begin{proof}
The containment $\alpha\mathcal O_E\subseteq\mathfrak a$ has a
positive integral index.  Multiplication by $\alpha$ has real
determinant
$|N_{E/\mathbb Q}\alpha|=\prod_v|\sigma_v\alpha|^2$ on
$\mathbb C^f$.  Comparing covolumes gives
\[
 |N_{E/\mathbb Q}\alpha|
   =N\mathfrak a\,[\mathfrak a:\alpha\mathcal O_E]
   \geq N\mathfrak a.
\]
Taking square roots proves the first assertion.  The second follows
because $\mathfrak a$ is an additive subgroup.
\end{proof}

\begin{lemma}[A box count]\label{gt-box-count}
For positive real numbers $R_1,\ldots,R_f$, let
\[
 B_{\mathfrak a}(R)=
 \#\{0\ne\alpha\in\mathfrak a:
                  |\sigma_v\alpha|\leq R_v\text{ for all }v\}.
\]
Then
\begin{equation}\label{gt-box-count-eq}
 B_{\mathfrak a}(R)
       \leq9^f\frac{\prod_vR_v^2}{N\mathfrak a}.
\end{equation}
Moreover $B_{\mathfrak a}(R)=0$ if
$\prod_vR_v<\sqrt{N\mathfrak a}$.
\end{lemma}

\begin{proof}
The last assertion is Lemma~\ref{gt-product-minimum}.
If the set is empty, the estimate is immediate.  Otherwise put
\[
 a=\left(\frac{\sqrt{N\mathfrak a}}{\prod_vR_v}\right)^{1/f},
 \qquad 0<a\leq1.
\]
After dividing the $v$th coordinate by $R_v$, the points lie in the
closed unit polydisc in $\mathbb C^f$.  Distinct resulting points
$z,z'$ satisfy
\[
 \prod_v|z_v-z'_v|\geq a^f,
 \qquad \max_v|z_v-z'_v|\geq a.
\]
The open polydiscs of radius $a/2$ about these points are therefore
pairwise disjoint.  They are all contained in the polydisc of radius
$1+a/2$.  Comparing their $2f$-dimensional volumes gives
\[
 B_{\mathfrak a}(R)
 \leq\left(1+\frac2a\right)^{2f}
 \leq\left(\frac3a\right)^{2f}
 =9^f\frac{\prod_vR_v^2}{N\mathfrak a}.
\]
One may apply this argument first to every finite subset.  The
uniform finite bound then also proves finiteness of the full set,
so no separate lattice-point asymptotic is being used.
\end{proof}

\subsection{A stronger smoothed sum with its exact logarithmic factor}

For $x=(x_v)_{v=1}^f$ with $x_v>0$, and $0\ne\alpha\in E$, set
\[
 q_x(\alpha)=\prod_{v=1}^f
       \min\left(1,\frac{x_v}{|\sigma_v\alpha|}\right),
 \qquad
 W_{K,x}(\alpha)=q_x(\alpha)^K
 =\prod_{v=1}^f
       \min\left(1,\left(\frac{x_v}{|\sigma_v\alpha|}\right)^K\right).
\]
Here $K$ is a positive real number.  Define $\Theta>0$ by
\begin{equation}\label{gt-theta}
 \Theta^f=\frac{\sqrt{N\mathfrak a}}{\prod_vx_v}.
\end{equation}

\begin{theorem}\label{gt-smoothed}
If $K>2$ and $\Theta\geq1$, then the sum over all nonzero elements
of $\mathfrak a$ converges and satisfies
\begin{equation}\label{gt-smoothed-exact}
 \sum_{0\ne\alpha\in\mathfrak a}W_{K,x}(\alpha)
 \leq
 \left(\frac{9K}{K-2}\right)^f\Theta^{-Kf}
 \sum_{j=0}^{f-1}\frac{((K-2)f\log\Theta)^j}{j!}.
\end{equation}
For $K\geq3$ and $\Theta\geq1$, the following simpler estimates
also hold:
\begin{align}
 \sum_{0\ne\alpha\in\mathfrak a}W_{K,x}(\alpha)
    &\leq\bigl(27\Theta^{-(K-1)}\bigr)^f,
       \label{gt-smoothed-simple}\\
 \sum_{0\ne\alpha\in\mathfrak a}W_{K,x}(\alpha)
    &\leq\bigl(81(1+\log\Theta)\Theta^{-K}\bigr)^f.
       \label{gt-smoothed-log}
\end{align}
\end{theorem}

\begin{proof}
For $r>0$ and $x>0$ the elementary identity
\[
 \min(1,x/r)^K
       =K\int_1^\infty 1_{\{r\leq xt\}}t^{-K-1}\,dt
\]
follows by integrating from $\max(1,r/x)$ to infinity.
Multiply this identity over the embeddings and sum over the ideal.
All functions are nonnegative, so the interchange of sums and
integrals is justified by monotone convergence.  The box count gives
\begin{align*}
 \sum_{0\ne\alpha\in\mathfrak a}W_{K,x}(\alpha)
 &=K^f\int_{[1,\infty)^f}
        B_{\mathfrak a}((x_vt_v)_v)
                 \prod_vt_v^{-K-1}\,dt_1\cdots dt_f\\
 &\leq9^fK^f\Theta^{-2f}
       \int_{\substack{t_v\geq1\\\prod_vt_v\geq\Theta^f}}
                         \prod_vt_v^{-(K-1)}\,dt_1\cdots dt_f.
\end{align*}
The restriction on the product is justified by the vanishing
assertion of Lemma~\ref{gt-box-count}.  Substitute $t_v=e^{u_v}$.
Writing $a=K-2>0$ and $L=f\log\Theta\geq0$, the last integral is
\[
 \int_{\substack{u_v\geq0\\\sum_vu_v\geq L}}
                      e^{-a\sum_vu_v}\,du_1\cdots du_f
   =\frac1{(f-1)!}\int_L^\infty e^{-au}u^{f-1}\,du.
\]
The equality follows by first computing the volume
$u^f/f!$ of $\{u_v\geq0:\sum_vu_v\leq u\}$ and then integrating
with respect to $u$.  Repeated integration by parts gives
\[
 \frac1{(f-1)!}\int_L^\infty e^{-au}u^{f-1}\,du
       =a^{-f}e^{-aL}\sum_{j=0}^{f-1}\frac{(aL)^j}{j!}.
\]
Substituting $e^{-aL}=\Theta^{-(K-2)f}$ proves
\eqref{gt-smoothed-exact}, including convergence.  These formulas
also hold for $f=1$, when the polynomial is $1$.

For the two simpler estimates, the same integral calculation without
the product restriction shows, for every $s>2$,
\begin{equation}\label{gt-moment-bound}
 \sum_{0\ne\alpha\in\mathfrak a}q_x(\alpha)^s
       \leq\left(\frac{9s}{s-2}\right)^f\Theta^{-2f}.
\end{equation}
Lemma~\ref{gt-product-minimum} also gives
\begin{equation}\label{gt-q-bound}
 q_x(\alpha)\leq
       \frac{\prod_vx_v}{\prod_v|\sigma_v\alpha|}
       \leq\Theta^{-f}.
\end{equation}
Consequently, if $s\leq K$, multiplication of
\eqref{gt-moment-bound} by $\Theta^{-f(K-s)}$ yields
\begin{equation}\label{gt-interpolation}
 \sum_{0\ne\alpha\in\mathfrak a}q_x(\alpha)^K
       \leq\left(\frac{9s}{s-2}
                         \Theta^{-(K-s+2)}\right)^f.
\end{equation}
Taking $s=3$ proves \eqref{gt-smoothed-simple}.
For \eqref{gt-smoothed-log}, take
$s=2+\eta$, where $\eta=(1+\log\Theta)^{-1}$.
Then $0<\eta\leq1$, so $s\leq3\leq K$, and
\[
 \frac{9s}{s-2}=9(3+2\log\Theta)
       \leq27(1+\log\Theta),\qquad
 \Theta^\eta\leq e<3.
\]
Equation~\eqref{gt-interpolation} now gives the stated bound.
\end{proof}

\begin{corollary}\label{gt-subgroup-sum}
Let $L\subseteq\mathfrak a$ be an additive subgroup, let $A>0$,
and suppose that $x_v\geq0$ and
\begin{equation}\label{gt-product-size}
 \prod_vx_v\leq(A/\Theta)^f\sqrt{N\mathfrak a},
 \qquad \Theta\geq A.
\end{equation}
For $K>2$, the right side of \eqref{gt-smoothed-exact}, with
$\Theta$ replaced by $\Theta/A$, bounds the sum of $W_{K,x}$ over
$L\setminus\{0\}$.  For $K\geq3$ one may also use
\begin{equation}\label{gt-subgroup-simple}
 \bigl(27(A/\Theta)^{K-1}\bigr)^f
 \quad\hbox{or}\quad
 \bigl(81(1+\log(\Theta/A))(A/\Theta)^K\bigr)^f.
\end{equation}
Every assertion holds, in particular, for every finite partial sum.
\end{corollary}

\begin{proof}
If some $x_v=0$, every summand is zero: a nonzero field element has
nonzero image under every embedding.  Suppose therefore that all
$x_v>0$.  Define $T$ by
$T^f=\sqrt{N\mathfrak a}/\prod_vx_v$.  Then
$T\geq\Theta/A\geq1$, and Theorem~\ref{gt-smoothed} applies.
The exact expression in \eqref{gt-smoothed-exact} is decreasing in
$T$: before its evaluation it is the product of $T^{-2f}$ and a
positive integral over
$\{u_v\geq0:\sum_vu_v\geq f\log T\}$.
The other two expressions are also decreasing for $T\geq1$ and
$K\geq3$.  For the logarithmic expression this follows directly from
\[
 \frac{d}{dT}\bigl((1+\log T)T^{-K}\bigr)
       =T^{-K-1}\bigl(1-K(1+\log T)\bigr)<0.
\]
The sum over $L$ is at most the sum over $\mathfrak a$.
\end{proof}

\subsection{The actual period and dual lattices}

We now assume that $E$ is a CM field and write
$\alpha\mapsto\bar\alpha$ for its complex conjugation.  Put
\[
 \Delta_E=\operatorname{disc}(E),\qquad
 \rho_E=|\Delta_E|^{1/(2f)},\qquad
 \mathfrak D_E^{-1}
   =\{x\in E:\operatorname{Tr}_{E/\mathbb Q}
                          (x\mathcal O_E)\subseteq\mathbb Z\}.
\]
Thus $\mathfrak D_E$ is the different, and $\rho_E$ is the root
discriminant.  The trace form is nondegenerate in characteristic
zero, so the displayed set is a nonzero fractional ideal.

\begin{lemma}\label{gt-different}
For every nonzero fractional ideal $\mathfrak a$,
\begin{align}
 \{x\in E:\operatorname{Tr}_{E/\mathbb Q}(x\mathfrak a)
                    \subseteq\mathbb Z\}
       &=\mathfrak D_E^{-1}\mathfrak a^{-1},
          \label{gt-trace-dual}\\
 N\mathfrak D_E&=|\Delta_E|.\label{gt-different-norm}
\end{align}
The ideals $\mathfrak D_E^\nu$ are preserved by complex
conjugation for every integer $\nu$.
\end{lemma}

\begin{proof}
If $x$ belongs to the set on the left of \eqref{gt-trace-dual},
then for every $y\in\mathfrak a$ and $b\in\mathcal O_E$ one has
$by\in\mathfrak a$ and $\operatorname{Tr}(xyb)\in\mathbb Z$.
Thus $x\mathfrak a\subseteq\mathfrak D_E^{-1}$.  Conversely
this containment implies the trace condition by taking $b=1$.
Multiplication by the inverse ideal gives
\eqref{gt-trace-dual}.

Let $e_1,\ldots,e_{2f}$ be an integral basis of $\mathcal O_E$
and let $T=(\operatorname{Tr}(e_ie_j))_{i,j}$.  In these coordinates
the codifferent is $T^{-1}\mathbb Z^{2f}$.  Its covolume relative
to $\mathcal O_E$ is $|\det T|^{-1}=|\Delta_E|^{-1}$.
Multiplicativity of norms gives \eqref{gt-different-norm}.
Finally the trace is unchanged by complex conjugation.  Hence
conjugation preserves the codifferent, its inverse, and all their
integer powers.
\end{proof}

Fix a positive integer $g$, an integer $\nu$, and a nonnegative
integer $s$.  The period lattice and the lift lattice in the global
orbital calculation are, respectively,
\begin{align}
 P_{\nu,s}
   &=\{y\in g\mathfrak D_E^\nu:
                    y-(-1)^s\bar y\in g^2\mathfrak D_E^\nu\},
           \label{gt-period-lattice}\\
 L_{\nu,s}
   &=\{x\in g^{-1}\mathfrak D_E^{-\nu}:
                    x-(-1)^s\bar x\in\mathfrak D_E^{-\nu}\}.
           \label{gt-lift-lattice}
\end{align}
The dual of the period lattice is taken for the pairing
$\operatorname{Tr}_{E/\mathbb Q}(xy/g)$:
\begin{equation}\label{gt-period-dual-definition}
 P_{\nu,s}^{\vee,g}
   =\{x\in E:\operatorname{Tr}_{E/\mathbb Q}(xy/g)
                    \in\mathbb Z\text{ for every }y\in P_{\nu,s}\}.
\end{equation}
The superscripts specify this normalization; using the unscaled
trace pairing would give a different power of $g$.

\begin{proposition}\label{gt-lattice-containments}
For the lattices just defined,
\begin{align}
 g^2\mathfrak D_E^\nu
       &\subseteq P_{\nu,s}\subseteq g\mathfrak D_E^\nu,
             \label{gt-period-between}\\
 \mathfrak D_E^{-\nu-1}
       &\subseteq P_{\nu,s}^{\vee,g}
                  \subseteq g^{-1}\mathfrak D_E^{-\nu-1},
             \label{gt-dual-between}\\
 \mathfrak D_E^{-\nu}
       &\subseteq L_{\nu,s}
                  \subseteq g^{-1}\mathfrak D_E^{-\nu}.
             \label{gt-lift-between}
\end{align}
In particular these groups are full lattices.  The containing ideals
in \eqref{gt-dual-between} and \eqref{gt-lift-between} have norms
\begin{align}
 \sqrt{N(g^{-1}\mathfrak D_E^{-\nu-1})}
       &=(g\rho_E^{\nu+1})^{-f},
             \label{gt-dual-norm}\\
 \sqrt{N(g^{-1}\mathfrak D_E^{-\nu})}
       &=(g\rho_E^\nu)^{-f}.
             \label{gt-lift-norm}
\end{align}
These assertions need no unramifiedness assumption on $E$ and no
restriction on the primes dividing $2g$.
\end{proposition}

\begin{proof}
The different and all its powers are conjugation invariant.  Thus
$y\in g^2\mathfrak D_E^\nu$ implies both
$y\in g\mathfrak D_E^\nu$ and
$y-(-1)^s\bar y\in g^2\mathfrak D_E^\nu$.
This proves \eqref{gt-period-between}.
If $x\in P_{\nu,s}^{\vee,g}$, its trace pairing with
$g^2\mathfrak D_E^\nu$ is integral, so
\[
 \operatorname{Tr}_{E/\mathbb Q}(gx\mathfrak D_E^\nu)
        \subseteq\mathbb Z.
\]
Lemma~\ref{gt-different} gives
$gx\in\mathfrak D_E^{-\nu-1}$, proving the upper containment in
\eqref{gt-dual-between}.  For its lower containment, pair
$x\in\mathfrak D_E^{-\nu-1}$ with
$y\in P_{\nu,s}\subseteq g\mathfrak D_E^\nu$.
Then $xy/g\in\mathfrak D_E^{-1}$, whose trace is integral.
The upper containment in \eqref{gt-lift-between} is part of the
definition, and the lower one follows from conjugation invariance
in the same way as \eqref{gt-period-between}.

The lower and upper groups in these containments are nonzero
fractional ideals, so each intervening additive group is a full
lattice.  Finally $N(g\mathcal O_E)=g^{2f}$ and
$N\mathfrak D_E=\rho_E^{2f}$.  Multiplicativity proves
\eqref{gt-dual-norm} and \eqref{gt-lift-norm} for every integer
$\nu$, including negative $\nu$.
\end{proof}

For the blocks used in the theta calculation, the exponent is
\begin{equation}\label{gt-block-exponent}
 \nu(r,b,s)=2r+3-2b-s.
\end{equation}
Here $r$ is a nonnegative integer, $b$ is the first block index, and
$s\geq1$ is its length.  Define
\[
 \mathfrak a_{b,s}=g^{-1}\mathfrak D_E^{-(\nu(r,b,s)+1)},
 \qquad
 \mathfrak l_{b,s}=g^{-1}\mathfrak D_E^{-\nu(r,b,s)}.
\]
Then
\begin{equation}\label{gt-adjacent-ideals}
 \mathfrak l_{b,s}=\mathfrak a_{b,s+1},\qquad
 \sqrt{N\mathfrak a_{b,s}}
    =(g\rho_E^{2r+4-2b-s})^{-f},\qquad
 \sqrt{N\mathfrak l_{b,s}}
    =(g\rho_E^{2r+3-2b-s})^{-f}.
\end{equation}
This is an identity of ideals, not an estimate; it accounts for the
one-step change in the discriminant exponent in the recursion.

\begin{corollary}[The actual dual-lattice error]\label{gt-actual-error}
Let $\Theta>4$.  Suppose that the nonnegative widths
$x^{D}_{b,v}$ and $x^{L}_{b,s,v}$ satisfy
\begin{align}
 \prod_vx^{D}_{b,v}
       &\leq\Theta^{-f}(g\rho_E^{2r+3-2b})^{-f},
                 \label{gt-dual-width}\\
 \prod_vx^{L}_{b,s,v}
       &\leq(2/\Theta)^f
                       (g\rho_E^{2r+3-2b-s})^{-f}.
                 \label{gt-lift-width}
\end{align}
For each indicated block, the sums
\begin{equation}\label{gt-actual-sums}
 \sum_{0\ne\alpha\in P_{\nu(r,b,1),1}^{\vee,g}}
                  W_{5,x^D_b}(\alpha),
 \qquad
 \sum_{0\ne\alpha\in L_{\nu(r,b,s),s}}
                  W_{5,x^L_{b,s}}(\alpha)
\end{equation}
are bounded by each of the following quantities:
\begin{align}
 &15^f(2/\Theta)^{5f}
       \sum_{j=0}^{f-1}
                  \frac{(3f\log(\Theta/2))^j}{j!},
                  \label{gt-error-exact}\\
 &(432\Theta^{-4})^f,
                  \label{gt-error-simple}\\
 &\bigl(2592(1+\log(\Theta/2))\Theta^{-5}\bigr)^f.
                  \label{gt-error-log}
\end{align}
The bounds are uniform over the block indices and over every finite
family of width vectors satisfying
\eqref{gt-dual-width}--\eqref{gt-lift-width}.
\end{corollary}

\begin{proof}
Proposition~\ref{gt-lattice-containments} places the two summation
groups in $\mathfrak a_{b,1}$ and $\mathfrak l_{b,s}$,
respectively.  Their norm formulas are
\eqref{gt-adjacent-ideals}.  Thus the assumptions are precisely
\eqref{gt-product-size}, with $A=1$ for the first group and $A=2$
for the second.  Corollary~\ref{gt-subgroup-sum}, with $K=5$, gives
the asserted common bounds by using $A=2$ for both.  The constants
are $9K/(K-2)=15$, $27\cdot2^4=432$, and
$81\cdot2^5=2592$.  No union bound over the blocks is taken:
the same bound holds separately for every required sum.
\end{proof}

\subsection{Exact covolumes, including primes above two}

Write $F$ for the maximal totally real subfield of $E$, so
$[F:\mathbb Q]=f$.  In this subsection suppose that $E/F$ is
unramified at every finite prime of $F$ dividing $g$.
This is weaker than assuming that $E/F$ is unramified at every
finite prime.  It also holds when every rational prime dividing
$g$ splits completely in $E$, as at the selected levels in the
application.  We give the elementary local calculation, since
division by $2$ would be invalid at some of these primes.

\begin{lemma}\label{gt-quadratic-count}
Let $A$ be the ring of integers of a nonarchimedean local field,
let $B$ be a quadratic \emph{\'etale} $A$-algebra which is
unramified, and let $\tau$ be its nontrivial involution.
Thus $B=A\times A$, with the factors exchanged, or $B$ is the
ring of integers of the unramified quadratic field extension.
For a nonzero nonunit $a\in A$ and $\epsilon\in\{1,-1\}$,
\begin{equation}\label{gt-local-kernel-card}
 \#\ker\bigl(1-\epsilon\tau:B/aB\longrightarrow B/aB\bigr)
       =\#(A/aA).
\end{equation}
\end{lemma}

\begin{proof}
In the split case the kernel consists of the pairs
$(x,\epsilon x)$, so the assertion is immediate.
In the field case choose an element of the quadratic residue
field which generates it over the residue field of $A$, and lift
it to $\omega\in B$.  Nakayama's lemma shows that
$1,\omega$ is an $A$-basis of $B$.  Its conjugate has a different
residue, since finite fields are perfect, so
$\delta=\omega-\tau\omega$ is a unit of $B$.
Put $t=\omega+\tau\omega\in A$.

For $\epsilon=1$, write an element as $x+y\omega$.
Its difference with its conjugate is $y\delta$, which vanishes
modulo $aB$ if and only if $y\in aA$.  The kernel is therefore
the copy of $A/aA$ represented by $x$.
For $\epsilon=-1$ the map is the trace
$x+y\omega\mapsto2x+ty$ into $A/aA$.
This map is surjective: if the residue characteristic is odd,
$2$ is a unit; if it is $2$, the equality
$\delta=2\omega-t$ shows that $t$ is a unit.
Its domain has cardinality $\#(A/aA)^2$, so its kernel has
cardinality $\#(A/aA)$.  This proves both signs without
inverting $2$.
\end{proof}

\begin{proposition}\label{gt-exact-covolumes}
Under the unramifiedness assumption of this subsection,
\begin{align}
 [g\mathfrak D_E^\nu:P_{\nu,s}]
    &=[P_{\nu,s}:g^2\mathfrak D_E^\nu]=g^f,
                 \label{gt-period-indices}\\
 [g^{-1}\mathfrak D_E^{-\nu}:L_{\nu,s}]&=g^f.
                 \label{gt-lift-index}
\end{align}
For standard Lebesgue measure on $\mathbb C^f$, the covolumes
under $\sigma=(\sigma_v)_v$ are
\begin{align}
 \operatorname{covol}\sigma(P_{\nu,s})
    &=2^{-f}g^{3f}\rho_E^{(2\nu+1)f},
                 \label{gt-period-covol}\\
 \operatorname{covol}\sigma(L_{\nu,s})
    &=2^{-f}g^{-f}\rho_E^{(1-2\nu)f}.
                 \label{gt-lift-covol}
\end{align}
In particular, if $\Lambda_{\nu,s}=\rho_E^{-\nu}
\sigma(P_{\nu,s})$, then
\begin{equation}\label{gt-normalized-covol}
 \operatorname{covol}\Lambda_{\nu,s}
      =(\rho_E/2)^f g^{3f},
\end{equation}
independently of $\nu$ and $s$.
\end{proposition}

\begin{proof}
At a finite prime $\mathfrak p$ of $F$ dividing $g$, put
$A=\mathcal O_{F,\mathfrak p}^{\wedge}$ and
$B=\mathcal O_E\otimes_{\mathcal O_F}A$.
The trace pairing of $B$ over $A$ is perfect because $B/A$ is
unramified.  This can be seen directly from the basis in
Lemma~\ref{gt-quadratic-count}: its trace determinant is
$(\omega-\tau\omega)^2$, a unit; in the split case use the
two standard idempotents.  The absolute trace pairing is the
composition of this relative pairing with the trace of $F$.
Consequently the localization of $\mathfrak D_E$ is
$\mathfrak D_F B$.  More explicitly, choose a basis of $B$
over $A$ and its relative trace-dual basis, both integral.
An element is dual to $B$ for the absolute trace exactly when
all its coefficients in that dual basis belong to the absolute
codifferent of $A$.  This proves the assertion about the
different and its inverse, and hence about every integer power.
Thus locally $\mathfrak D_E^\nu=dB$ for a nonzero $d\in F_{\mathfrak p}$
fixed by $\tau$.

For $\epsilon=(-1)^s$, reduction modulo $g\mathfrak D_E^\nu$
therefore identifies
\[
 \{z\in\mathfrak D_E^\nu:z-\epsilon\bar z
                            \in g\mathfrak D_E^\nu\}
       /g\mathfrak D_E^\nu
\]
with a product, over $\mathfrak p\mid g$, of the kernels in
Lemma~\ref{gt-quadratic-count}.  Its cardinality is
$\#(\mathcal O_F/g\mathcal O_F)=g^f$.
The whole quotient
$\mathfrak D_E^\nu/g\mathfrak D_E^\nu$ has cardinality
$g^{2f}$.  Multiplication by $g$ gives both indices in
\eqref{gt-period-indices}.  Applying the same argument with
$\mathfrak D_E^{-\nu}$ and then multiplying by $g^{-1}$ proves
\eqref{gt-lift-index}.  When $g=1$ the quotients are trivial
and the same formulas hold.

For completeness,
\begin{equation}\label{gt-ideal-covol}
 \operatorname{covol}\sigma(\mathfrak a)
       =2^{-f}|\Delta_E|^{1/2}N\mathfrak a
       =2^{-f}\rho_E^fN\mathfrak a.
\end{equation}
Indeed the matrix of all $2f$ complex embeddings on an integral
basis has determinant of absolute value $|\Delta_E|^{1/2}$.
Replacing each conjugate pair of rows by its real and imaginary
parts divides the determinant's absolute value by $2$.
This gives the formula for $\mathcal O_E$, and the definition of
$N\mathfrak a$ gives it for every fractional ideal.
The two index formulas and
$N(g^a\mathfrak D_E^b)=g^{2af}\rho_E^{2bf}$ now give
\eqref{gt-period-covol} and \eqref{gt-lift-covol}.
Multiplication by the positive real number $\rho_E^{-\nu}$
on $\mathbb C^f$ scales volume by $\rho_E^{-2\nu f}$,
which proves \eqref{gt-normalized-covol}.
\end{proof}

Use the real inner product
$\langle z,w\rangle_{\mathbb R}
   =\operatorname{Re}\sum_v\bar z_vw_v$ on $\mathbb C^f$.
For a full lattice $\Lambda$, write
\[
 \Lambda^*=\{\xi\in\mathbb C^f:
       \langle\xi,\lambda\rangle_{\mathbb R}\in\mathbb Z
                       \text{ for every }\lambda\in\Lambda\}.
\]

\begin{proposition}\label{gt-Fourier-dual}
For every nonzero real number $c$, without any unramifiedness
assumption, one has
\begin{equation}\label{gt-Fourier-dual-eq}
 (c\sigma(P_{\nu,s}))^*
       =\left\{\frac2{gc}\overline{\sigma(x)}:
                             x\in P_{\nu,s}^{\vee,g}\right\}.
\end{equation}
Under the unramifiedness assumption used in
Proposition~\ref{gt-exact-covolumes}, one further has
\begin{align}
 P_{\nu,s}^{\vee,g}&=L_{\nu+1,s+1},
                 \label{gt-exact-trace-dual}\\
 \operatorname{covol}\sigma(P_{\nu,s}^{\vee,g})
       &=2^{-f}g^{-f}\rho_E^{-(2\nu+1)f},
                 \label{gt-dual-covol}\\
 \operatorname{covol}\bigl(\rho_E^{\nu+1}
                       \sigma(P_{\nu,s}^{\vee,g})\bigr)
       &=(\rho_E/(2g))^f.
                 \label{gt-normalized-dual-covol}
\end{align}
\end{proposition}

\begin{proof}
For $x,y\in E$ the normalizations satisfy exactly
\begin{equation}\label{gt-pairing-identity}
 \left\langle\frac2{gc}\overline{\sigma(x)},
                         c\sigma(y)\right\rangle_{\mathbb R}
       =\frac2g\operatorname{Re}\sum_v\sigma_v(xy)
       =\operatorname{Tr}_{E/\mathbb Q}(xy/g).
\end{equation}
This gives one inclusion in \eqref{gt-Fourier-dual-eq}.
For the other, choose a $\mathbb Z$-basis $y_1,\ldots,y_{2f}$
of $P_{\nu,s}$.  It is a $\mathbb Q$-basis of $E$.
If $\xi$ belongs to the Fourier dual, the prescribed values
$\langle\xi,c\sigma(y_i)\rangle_{\mathbb R}$ are integers.
Nondegeneracy of the rational trace form gives a unique $x\in E$
with those values equal to $\operatorname{Tr}(xy_i/g)$.
The $2f$ vectors $c\sigma(y_i)$ form a real basis, so
\eqref{gt-pairing-identity} implies
$\xi=2\overline{\sigma(x)}/(gc)$.
The integral values show that $x\in P_{\nu,s}^{\vee,g}$.

For any real full lattice, its Euclidean dual has reciprocal
covolume.  In coordinates, if the basis matrix is $M$, the dual
basis matrix is $(M^{-1})^{\mathsf T}$, which proves this assertion.
Taking $c=1$ in \eqref{gt-Fourier-dual-eq} and using
\eqref{gt-period-covol} therefore gives
\[
 (2/g)^{2f}\operatorname{covol}\sigma(P_{\nu,s}^{\vee,g})
       =\bigl(2^{-f}g^{3f}\rho_E^{(2\nu+1)f}\bigr)^{-1}.
\]
This is \eqref{gt-dual-covol}.

It remains to identify the lattice, rather than just its volume.
Put $\epsilon=(-1)^s$.  If $x\in P_{\nu,s}^{\vee,g}$, then
$x\in g^{-1}\mathfrak D_E^{-\nu-1}$ by
Proposition~\ref{gt-lattice-containments}.
For every $z\in\mathfrak D_E^\nu$ the element
$g(z+\epsilon\bar z)$ belongs to $P_{\nu,s}$.
Pairing with it and using invariance of trace under conjugation
gives
\[
 \operatorname{Tr}\bigl((x+\epsilon\bar x)z\bigr)\in\mathbb Z.
\]
Hence $x+\epsilon\bar x\in\mathfrak D_E^{-\nu-1}$ by
Lemma~\ref{gt-different}, and so
$P_{\nu,s}^{\vee,g}\subseteq L_{\nu+1,s+1}$.
Equation~\eqref{gt-lift-covol} shows that the latter lattice has
the covolume in \eqref{gt-dual-covol}.  The inclusion has index
one, proving \eqref{gt-exact-trace-dual}.  Finally scaling by
$\rho_E^{\nu+1}$ gives \eqref{gt-normalized-dual-covol}.
\end{proof}

\subsection{The global Poisson inequality with its actual regularity}

The summation step in the source uses continuous nonnegative
functions, not merely Schwartz functions.  We prove the needed
version.  This also fixes the Fourier and volume normalizations.
For an integrable function $\phi:\mathbb C^f\to\mathbb C$, let
\[
 \widehat\phi(\xi)=\int_{\mathbb C^f}\phi(z)
          e^{-2\pi i\langle\xi,z\rangle_{\mathbb R}}\,dz,
\]
where $dz$ is standard Lebesgue measure.

\begin{theorem}\label{gt-Poisson}
Let $\Lambda$ be a full lattice in $\mathbb C^f$ and
$V=\operatorname{covol}\Lambda$.
Suppose that $\phi:\mathbb C^f\to\mathbb C$ is continuous and
that, for some finite $C\geq0$,
\begin{equation}\label{gt-Poisson-decay}
 |\phi(z)|\leq C\prod_{v=1}^f(1+|z_v|)^{-3}.
\end{equation}
If $\sum_{\xi\in\Lambda^*}|\widehat\phi(\xi)|<\infty$, then
for every $z_0\in\mathbb C^f$,
\begin{equation}\label{gt-Poisson-eq}
 \sum_{\lambda\in\Lambda}\phi(z_0+\lambda)
       =V^{-1}\sum_{\xi\in\Lambda^*}
          \widehat\phi(\xi)
                 e^{2\pi i\langle\xi,z_0\rangle_{\mathbb R}}.
\end{equation}
Both series converge absolutely, and the first converges locally
uniformly in $z_0$.
If $\phi$ is real and nonnegative, then every finite subset
$S\subseteq z_0+\Lambda$ satisfies
\begin{equation}\label{gt-Poisson-inequality}
 \sum_{z\in S}\phi(z)
       \leq V^{-1}\int_{\mathbb C^f}\phi(z)\,dz
          +V^{-1}\sum_{0\ne\xi\in\Lambda^*}
                                  |\widehat\phi(\xi)|.
\end{equation}
\end{theorem}

\begin{proof}
Write $h(z)=\prod_v(1+|z_v|)^{-3}$.
This function is integrable because each complex coordinate has
real dimension two and
$2\pi\int_0^\infty r(1+r)^{-3}\,dr<\infty$.
Choose $\delta>0$ so that the open polydiscs
$\lambda+\{z:\max_v|z_v|<\delta\}$, $\lambda\in\Lambda$,
are pairwise disjoint.  Such a choice exists because a full
lattice is discrete.  On each of these polydiscs,
\[
 h(\lambda)\leq(1+\delta)^{3f}h(z).
\]
Integration over the disjoint polydiscs proves
$\sum_{\lambda\in\Lambda}h(\lambda)<\infty$.
For $z_0$ in any fixed compact set, a further fixed factor bounds
$h(z_0+\lambda)$ by $h(\lambda)$, by the triangle inequality.
Thus \eqref{gt-Poisson-decay} and the Weierstrass test imply that
\[
 F(z_0)=\sum_{\lambda\in\Lambda}\phi(z_0+\lambda)
\]
is a continuous $\Lambda$-periodic function.

Integrate over a half-open fundamental parallelepiped $Q$ of
$\Lambda$.  Absolute integrability permits unfolding, and for
$\xi\in\Lambda^*$ the corresponding Fourier coefficient is
\[
 V^{-1}\int_Q F(z)e^{-2\pi i\langle\xi,z\rangle_{\mathbb R}}\,dz
       =V^{-1}\widehat\phi(\xi).
\]
The assumed absolute convergence implies that the right side of
\eqref{gt-Poisson-eq}, as a Fourier series in $z_0$, converges
uniformly to a continuous function with exactly those
coefficients.  The difference from $F$ has all Fourier
coefficients zero and therefore vanishes.  One elementary
justification of this last assertion is to identify
$\mathbb C^f/\Lambda$ with
$(\mathbb R/\mathbb Z)^{2f}$ using a lattice basis and convolve
with products of the one-dimensional Fej\'er kernels
\[
 F_N(t)=\frac1N\left|\sum_{j=0}^{N-1}e^{2\pi ijt}\right|^2.
\]
These kernels are nonnegative, have integral one, and their
mass outside any fixed neighborhood of $0$ tends to zero:
away from the integers the geometric-sum formula bounds
$F_N(t)$ by a constant divided by $N$.
Consequently their convolution with a continuous function on
the torus converges uniformly to that function.
All those convolutions are zero when all its Fourier
coefficients vanish.  This proves \eqref{gt-Poisson-eq}.
For real nonnegative $\phi$, the finite sum is at most the full
sum, and the triangle inequality bounds the nonzero Fourier
terms.  The zero coefficient is $\int\phi$, giving
\eqref{gt-Poisson-inequality}.
\end{proof}

\begin{corollary}\label{gt-actual-Poisson-error}
Let $P=P_{\nu,s}$, $c\ne0$, $\Lambda=c\sigma(P)$, and
$V=\operatorname{covol}\Lambda$.
Let $\phi$ be continuous, real, nonnegative, and satisfy
\eqref{gt-Poisson-decay}.
Suppose $B\geq0$ and nonnegative widths $x_v$ satisfy
\[
 V^{-1}\left|\widehat\phi
              \left(\frac2{gc}\overline{\sigma(\alpha)}\right)\right|
       \leq B W_{5,x}(\alpha)
       \quad(0\ne\alpha\in P^{\vee,g}),
 \qquad
 \prod_vx_v\leq(A/\Theta)^f(g\rho_E^{\nu+1})^{-f},
\]
where $A>0$ and $\Theta\geq A$.
For every $y_0\in E$ and finite subset $S\subseteq y_0+P$,
\begin{equation}\label{gt-actual-Poisson-error-eq}
 \sum_{y\in S}\phi(c\sigma(y))
       \leq V^{-1}\int\phi
                   +B\bigl(27(A/\Theta)^4\bigr)^f.
\end{equation}
The exact logarithmic-polynomial bound of
Corollary~\ref{gt-subgroup-sum} can replace the last factor.
If $E/F$ is unramified above $g$, then
\[
 V=|c|^{2f}2^{-f}g^{3f}\rho_E^{(2\nu+1)f}.
\]
In particular $c=\rho_E^{-\nu}$ gives
$V=(\rho_E/2)^fg^{3f}$.
\end{corollary}

\begin{proof}
The containing ideal for $P^{\vee,g}$ and its norm are exactly
\eqref{gt-dual-between} and \eqref{gt-dual-norm}.
Corollary~\ref{gt-subgroup-sum} therefore bounds the sum of
the nonzero Fourier coefficients divided by $V$.
Their sum is finite; the zero coefficient is finite because
$\phi$ is integrable.  Proposition~\ref{gt-Fourier-dual}
identifies all Fourier frequencies, with no omitted factor
of $2$, $g$, or $c$.  Theorem~\ref{gt-Poisson} now proves
\eqref{gt-actual-Poisson-error-eq}, and
Proposition~\ref{gt-exact-covolumes} gives the displayed volume.
\end{proof}

\subsection{Necessary hypotheses and the limits of this improvement}

\begin{example}\label{gt-critical-exponent}
The condition $K>2$ cannot be dropped from a uniform assertion of
convergence.  Take $E=\mathbb Q(i)$,
$\mathfrak a=\mathbb Z[i]$, and $x=1/\Theta<1$.  For $K=2$,
\[
 \sum_{0\ne\alpha\in\mathbb Z[i]}W_{2,x}(\alpha)
       =x^2\sum_{0\ne\alpha\in\mathbb Z[i]}|\alpha|^{-2}
       =\infty.
\]
Indeed, the $a$ points $a+bi$, $0\leq b<a$, have squared modulus
at most $2a^2$, so their contribution to the latter sum is at least
$1/(2a)$.  Summing over $a\geq1$ proves divergence.  The same
conclusion holds for $0<K\leq2$, since $0<q_x(\alpha)\leq1$.
The exponent in the actual theta application is $K=5$.
\end{example}

\begin{example}\label{gt-log-necessary}
The logarithmic factor at the full decay power cannot in general be
replaced by a constant independent of $\Theta$.  Take
$E=\mathbb Q(\sqrt2,i)$, so $f=2$, and
$\mathfrak a=\mathcal O_E$.  The algebraic integer
$u=1+\sqrt2$ is a unit, with inverse $\sqrt2-1$.
At the two complex places its absolute values are $u$ and $u^{-1}$.
Set $x_1=x_2=\Theta^{-1}$, where $\Theta>1$.
For every integer $n$ satisfying
$|n|\leq\log\Theta/\log u$, both absolute values of $u^n$
lie between $\Theta^{-1}$ and $\Theta$.  Hence
\[
 q_x(u^n)=\Theta^{-2},\qquad
 W_{K,x}(u^n)=\Theta^{-2K}.
\]
The powers are distinct, and therefore
\[
 \sum_{0\ne\alpha\in\mathcal O_E}W_{K,x}(\alpha)
 \geq
 \left(2\left\lfloor\frac{\log\Theta}{\log(1+\sqrt2)}
                         \right\rfloor+1\right)\Theta^{-2K}.
\]
This rules out a bound $C^2\Theta^{-2K}$ with a fixed $C$ for this
one field.  In degree four the degree-one logarithmic polynomial
in \eqref{gt-smoothed-exact} has the correct order of growth.
\end{example}

Two further restrictions are essential in the elementary count.
The zero element is excluded: for $\mathfrak a=\mathbb Z[i]$
and $0<R<1/3$, including zero would give at least one point,
whereas the proposed right side $9R^2$ is smaller than one.
Also the product lower bound must apply to differences of points.
A bound only on each individual point would allow arbitrarily many
distinct points on the complex unit circle and would not imply a
packing estimate.  The actual summation sets lie in fractional
ideals, so both restrictions are satisfied.

The unramifiedness assumption for the exact indices also has
content, although it was unnecessary for the containing ideals.
For $E=\mathbb Q(i)$, $g=2$, $\nu=0$, and $s=0$, one has
\[
 P_{0,0}=\{2(a+bi):4bi\in4\mathbb Z[i]\}=2\mathbb Z[i].
\]
Thus $[2\mathbb Z[i]:P_{0,0}]=1$, whereas the formula
$g^f$ would give $2$.  This prime is ramified in
$\mathbb Q(i)/\mathbb Q$, so it is excluded by
Proposition~\ref{gt-exact-covolumes} and by the actual source
hypotheses.  The example also explains why a proof that
diagonalizes the involution by dividing by $2$ would not
justify the asserted range of that proposition.

The source's stated smoothed-sum estimate has the weaker form
$(C\Theta^{-(K-5/2)})^f$; its proof already obtains an intermediate
bound with exponent $K-1$ before weakening it.  Thus retaining that
intermediate exponent is an improvement of the stated result, not
a correction of a false theorem.  The direct integration above
also gives the explicit logarithmic polynomial at exponent $K$,
smaller constants, and the exact convergence range $K>2$.
The actual application has $K=5$ and the width factor $2$ recorded
in Corollary~\ref{gt-actual-error}.

This proves the ideal-summation component and the global Poisson
step of the scalar orbital estimate, including its actual period
lattices, covolumes, trace and Fourier normalizations, and
discriminant exponents.  It does not by itself prove the width
inequalities \eqref{gt-dual-width}--\eqref{gt-lift-width}: those
come from the chosen archimedean test functions and their parameter
bounds.  The Fourier-coefficient bound for those particular
test functions must also be established.  Nor does the section
supply the finite-place mass identities or the remaining
larger-corank, nonscalar, and nongeneric orbital estimates.  The
degenerate diagonal estimate and the comparison with the ambient
automorphic dimension remain separate obligations.  No
counterexample to the actual ideal-sum or lattice-containment
statements has been found.


\section{Local unitary densities and the exact mass ratio}
\label{md-section}

We verify the local densities and the mass normalization used in
Section~\ref{as-section}. The finite-field calculations and all
normalization changes are proved below, including residue
characteristic $2$. The global mass formula is used as an explicitly
stated classical theorem. This is sufficient to check the source's
mass-ratio assertion; it is not a new proof of the general mass formula.

\subsection{Counting primitive frames over a finite field}

Let $k=\mathbb F_q$, where $q$ is a prime power, and let
$k'=\mathbb F_{q^2}$ with involution $a\mapsto\overline a=a^q$.
On $(k')^n$ put $h(x,y)=\overline x^{\,t}y$.
For a matrix $G=\overline G^{\,t}\in M_m(k')$ of rank $m-c$, let
$\mathcal E_-(n,G)$ be the set of matrices $X\in M_{n,m}(k')$
of column rank $m$ satisfying $\overline X^{\,t}X=G$.
In the split case, let $G\in M_m(k)$ have rank $m-c$, and let
$\mathcal E_+(n,G)$ consist of pairs $(X,Y)$ of matrices in
$M_{n,m}(k)$, both of column rank $m$, satisfying $X^tY=G$.
In both cases $0\leq c\leq m$.
For $\eta\in\{1,-1\}$ define
\begin{equation}\label{md-density}
 d_\eta(n,m,c;q)=
 \begin{cases}
 \displaystyle\prod_{j=n-m-c+1}^{n}(1-\eta^j q^{-j}),&m+c\leq n,\\
 0,&m+c>n.
 \end{cases}
\end{equation}
An empty product is $1$. Equivalently, one may start the product at
$\max(0,n-m-c+1)$: the extra cases then contain the zero factor
$1-\eta^0q^0$.

\begin{lemma}\label{md-unitary-order}
For $n\geq0$,
\[
 \#U_n(k)=q^{n(n-1)/2}\prod_{j=1}^n(q^j-(-1)^j).
\]
For $n\geq1$, the number of nonzero isotropic vectors in $(k')^n$ is
\[
                  (q^n-(-1)^n)(q^{n-1}-(-1)^{n-1}).
\]
Every nondegenerate Hermitian form on $(k')^n$ is isometric to $h$.
These assertions hold also when $q$ is even.
\end{lemma}

\begin{proof}
The norm $N:k'^\times\to k^\times$, $N(a)=a^{q+1}$, is
surjective and every fiber has $q+1$ elements: the multiplicative
group of a finite field is cyclic. The trace $a+\overline a$ is
a nonzero $k$-linear map $k'\to k$, hence is also surjective.
It is nonzero even in characteristic $2$, since the fixed field
of the involution is exactly $k$.

A nonzero nondegenerate Hermitian form has a vector of nonzero
length. Otherwise, expanding $h(ax+y,ax+y)=0$ for all $a\in k'$
would give $\operatorname{Tr}(\overline a h(x,y))=0$ for all $a$.
The trace pairing is nondegenerate, so the form would be zero.
Normalize a vector of nonzero length using the norm map, take its
nondegenerate orthogonal complement, and induct. This proves the
last assertion and extends every orthonormal frame to a basis.

Let $A_n$ count vectors of length $0$ and let $B_n$ count vectors
of length a prescribed nonzero element of $k$. The latter count
is independent of that element by norm surjectivity. We have
$A_0=1$, $B_0=0$, and, by separating the last coordinate,
\[
 A_n=A_{n-1}+(q^2-1)B_{n-1},\qquad
 A_n+(q-1)B_n=q^{2n}.
\]
Equivalently $A_n-B_n=-q(A_{n-1}-B_{n-1})=(-q)^n$.
Solving these two equations gives
\[
 B_n=q^{n-1}(q^n-(-1)^n),\qquad
 A_n=q^{2n-1}+(q-1)(-1)^nq^{n-1}.
\]
Thus $A_n-1$ is the second expression in the lemma.
Choosing the first vector of an orthonormal basis in $B_n$ ways
and then a basis of its orthogonal complement gives
$\#U_n(k)=B_n\#U_{n-1}(k)$, proving the product formula.
\end{proof}

\begin{theorem}\label{md-frame-count}
For the two sets of primitive frames defined above,
\[
 \#\mathcal E_-(n,G)=q^{2nm-m^2}d_{-1}(n,m,c;q),\qquad
 \#\mathcal E_+(n,G)=q^{2nm-m^2}d_1(n,m,c;q).
\]
The same formulas hold for any nondegenerate target Hermitian
form in the first case and any perfect target bilinear pairing
in the second case.
\end{theorem}

\begin{proof}
Changes of target bases remove the target form or pairing.
If $m>n$ both sets are empty, and the displayed products are zero.
The assertion is immediate if $m=0$. We may therefore assume
$1\leq m\leq n$.

First consider the Hermitian case. A basis of the radical and
Lemma~\ref{md-unitary-order} put $G$ into the form
$\operatorname{diag}(I_r,0_c)$, where $r=m-c$.
Choose the first $r$ columns as an orthonormal frame. There are
$\#U_n(k)/\#U_{n-r}(k)$ choices. The last $c$ columns are a
linearly independent ordered isotropic frame in the orthogonal
complement of dimension $N=n-r$; all their pairings must vanish.

Let $F(N,c)$ count such ordered isotropic frames. Their span is
contained in its own orthogonal complement, so $2c\leq N$ is
necessary. If $c\geq1$, choose the first nonzero isotropic vector
$v$. The quotient $v^\perp/k'v$ is nondegenerate of dimension
$N-2$. To see this directly, choose $w$ with $h(v,w)=1$ and
replace $w$ by $w+av$ so that it is isotropic; the required
choice is possible because the trace is surjective. The span
of $v,w$ is then a nondegenerate plane.
Every isotropic $(c-1)$-frame in the quotient has
$q^{2(c-1)}$ lifts to $v^\perp$, all of which form with $v$
an independent isotropic frame. Therefore
\[
 F(N,c)=(q^N-(-1)^N)(q^{N-1}-(-1)^{N-1})
                         q^{2(c-1)}F(N-2,c-1).
\]
Together with $F(N,0)=1$, this proves
\[
 F(N,c)=q^{c(c-1)}
           \prod_{j=N-2c+1}^{N}(q^j-(-1)^j)
                  \quad(2c\leq N).
\]
The necessary condition $2c\leq n-r$ is precisely $m+c\leq n$.
Multiplying by the orthonormal-frame count makes the two products
adjacent, with indices $n-m-c+1,\ldots,n$. The total exponent of
$q$ after writing each factor as $q^j(1-(-1)^jq^{-j})$ is
\[
 \frac{n(n-1)-(n-r)(n-r-1)}2+c(c-1)
             +\sum_{j=n-m-c+1}^{n}j
                       =2nm-m^2.
\]
This proves the Hermitian formula, including the impossible cases.

For the split case, the number of possible matrices $X$ is
$\prod_{i=0}^{m-1}(q^n-q^i)$. For a fixed $X$, a change of
dual target bases reduces the equation to saying that the first
$m$ rows of $Y$ are $G$. Write the remaining $n-m$ rows as $Z$.
The stacked matrix $Y$ is injective precisely when $Z$ is
injective on $\ker G$, which has dimension $c$. After choosing
a complement of this kernel, the number of choices is
\[
           q^{(n-m)(m-c)}\prod_{i=0}^{c-1}(q^{n-m}-q^i).
\]
It is zero if $c>n-m$. Otherwise, multiply by the count of $X$
and extract powers of $q$. The remaining product has exactly
the indices $n-m-c+1,\ldots,n$, and the extracted exponent is
$2nm-m^2$. This proves the split formula.
\end{proof}

\subsection{Lifting primitive frames, including the dyadic case}

Let $A$ be a complete discrete valuation ring with uniformizer
$\varpi$ and residue field $k$ of size $q$. Let $B$ be either
the integers in the unramified quadratic extension of its fraction
field or $A\times A$ with its interchange involution. A Hermitian
matrix over $B$ is a matrix $G=\overline G^{\,t}$. In the split
case it is equivalent to one arbitrary matrix over $A$.
Let $\Phi\in M_n(B)$ be an invertible Hermitian matrix, and let
$G\in M_m(B)$ be Hermitian with reduction of corank $c$.
A primitive solution modulo $\varpi^a$ to
\[
                         \overline X^{\,t}\Phi X=G
\]
means one whose reduction has full column rank at both residue
factors, if there are two. Denote their number by $N_a$.

\begin{proposition}\label{md-lifting}
If $N_1\ne0$, then for every $a\geq1$,
\[
 N_a=q^{(a-1)(2nm-m^2)}N_1.
\]
The normalized primitive density is consequently
\[
 \lim_{a\to\infty}q^{-a(2nm-m^2)}N_a
                        =d_\eta(n,m,c;q),
\]
with $\eta=-1$ in the inert case and $\eta=1$ in the split case.
If $N_1=0$, all the counts and the density are zero.
\end{proposition}

\begin{proof}
Suppose first that the residue algebra is $k'$. A tentative lift
of a solution at level $a$ may be corrected by $\varpi^aH$.
Modulo $\varpi^{a+1}$ its Gram matrix changes by
\[
        \varpi^a(\overline X^{\,t}\Phi H+
                                      \overline H^{\,t}\Phi X).
\]
All matrices on the right can be reduced modulo $\varpi$.
Since $X$ is injective and $\Phi$ is invertible,
$\overline X^{\,t}\Phi$ has rank $m$. Thus
$T=\overline X^{\,t}\Phi H$ ranges over all of $M_m(k')$.
The map $T\mapsto T+\overline T^{\,t}$ is onto the
$k$-space of Hermitian matrices: choose upper off-diagonal
entries freely and choose each diagonal entry using the
surjective trace $k'\to k$. This argument uses no division by $2$.
The source space of corrections has $k$-dimension $2nm$ and the
target has dimension $m^2$, so every allowed Gram correction has
$q^{2nm-m^2}$ preimages.

In the split case write the two components as $X,Y$.
The linearized Gram map is
$\delta X^t\Phi Y+X^t\Phi\delta Y$. Already its second
term is onto $M_m(k)$, since $X$ has rank $m$.
The same dimension count gives the same number of lifts.
Iterating proves the first formula; division by the indicated
power of $q$ and Theorem~\ref{md-frame-count} prove the second.
Primitivity is preserved under lifting.
\end{proof}

In particular the unitary stabilizer density of a unimodular
rank-$n$ lattice is
\begin{equation}\label{md-stabilizer-density}
                 \beta_{n,v}=
                       \prod_{j=1}^n(1-\eta_v^j q_v^{-j}).
\end{equation}
This follows by taking $m=n$, $c=0$: a matrix with the prescribed
invertible Gram matrix is automatically invertible. The argument
does not require the local base field to be unramified over
$\mathbb Q_2$. It only requires that its quadratic algebra $B$
be unramified or split.

\begin{example}\label{md-degenerate-test}
For $q=2$, $n=2$, $m=1$, and zero Gram matrix, there are $9$
primitive Hermitian vectors and $3$ split pairs. Their normalized
densities are $9/8$ and $3/8$, respectively. For nonzero Gram
matrix the two counts are both $6$, and the two densities are
$3/4$. These values follow directly from the formulas above and
show why the corank $c$ must occur in the lower limit of the product.
Also, in rank $n=m=1$, a zero Gram matrix has no primitive
Hermitian solution. If primitivity is deleted the zero vector is
a solution. The source's density retains both corank and
primitivity, so neither example contradicts it.
\end{example}

\subsection{The precise global input}

Let $E/F$ be a CM extension: $F$ is totally real of degree
$f=[F:\mathbb Q]$ and $E$ is a totally imaginary quadratic
extension. Write $\mathcal O_F,\mathcal O_E$ for their rings of
integers, $d_F>0$ for the absolute discriminant of $F$, and
$\mathfrak d_{E/F}$ for the relative discriminant. Assume that
$E/F$ is unramified at every finite prime. For a finite prime $v$
of $F$ put $q_v=\#(\mathcal O_F/v)$ and
$\eta_v=1$ if $v$ splits in $E$, $\eta_v=-1$ if it is inert.
Let $\chi$ be the quadratic Hecke character of $E/F$.
For $j\geq2$ put
\[
                L_j=L_F(j,\chi^j)
                   =\prod_v(1-\eta_v^j q_v^{-j})^{-1}.
\]
Here $\chi^j$ means the trivial character for even $j$ and
$\chi$ for odd $j$. Write $L_1=L_F(1,\chi)>0$ for the value
defined by analytic continuation, or equivalently by the ratio
of the residues of $\zeta_E$ and $\zeta_F$ at $1$.

For a positive definite Hermitian space $V$ over $E$ and a full
$\mathcal O_E$-lattice $L$, its genus is the set of lattices
isometric to $L$ at every finite completion, modulo global
isometry. Its mass is
\[
                \operatorname{Mass}(L)
                   =\sum_{[L']\in\operatorname{gen}(L)}
                                  \frac1{\#\operatorname{Aut}(L')}.
\]
The class set and stabilizers are finite. The class interpretation
and finiteness in the setting used here were proved in
Sections~\ref{al-section} and \ref{acl-section}.
For $n\geq1$, let $M_n$ be the mass of the standard lattice
$\mathcal O_E^n$ with Hermitian form $\overline x^{\,t}y$.

We use the classical Siegel mass formula in its unitary form.
With unimodular local lattices, it states
\begin{equation}\label{md-classical-mass}
 \operatorname{Mass}(L)=
  2d_F^{n^2/2}\,N_{F/\mathbb Q}(\mathfrak d_{E/F})^{n(n+1)/4}
   \left(\prod_{j=1}^n\frac{(j-1)!}{(2\pi)^j}\right)^f
                            \prod_v\beta_{n,v}^{-1}.
\end{equation}
The factor $2$ is the unitary Tamagawa number; the compact
archimedean group has Haar probability measure when interpreting
the class mass. The product has the usual $L(1,\chi)$ convergence
normalization. The formula in this generality is recorded in
Kirschmer, Theorem 4.2.3, from Gan--Yu, \S10.
Terakado--Yu, equation (2.6), gives the same constants.
\cite{kirschmer-herm,terakado-yu-mass}
The latter article imposes an additional dyadic assumption for
its ramified local factors. We do not use those factors: the
unramified local formula needed here has been proved above.
Thus no unramifiedness condition on $F/\mathbb Q$ is being added
to the source hypotheses.

\begin{theorem}\label{md-mass-ratio}
Under the preceding hypotheses, $M_n>0$ and
\[
 M_n=2d_F^{n^2/2}
             \prod_{j=1}^n
                  \left(\frac{(j-1)!}{(2\pi)^j}\right)^f L_j.
\]
In particular, for $n\geq2$,
\begin{equation}\label{md-ratio}
 \frac{M_n}{M_{n-1}}
       =\left((n-1)!\left(\frac{\rho}{2\pi}\right)^n\right)^f
                   d_F^{-1/2}L_n,
             \qquad \rho=d_F^{1/f}.
\end{equation}
Equivalently, with $v=\rho/2$,
\[
 \frac{M_n}{M_{n-1}}
     =\frac{\bigl((n-1)!(v/\pi)^n\bigr)^f}
                  {\sqrt{(2v)^f}}\,L_n.
\]
This is the source mass ratio, including its archimedean
normalization and discriminant power.
\end{theorem}

\begin{proof}
The relative discriminant is the unit ideal.
Substitute \eqref{md-stabilizer-density} in the classical
formula \eqref{md-classical-mass}. Reorganizing the finitely
many rank factors gives the displayed expression for $M_n$.
Equivalently, one first uses the usual $L(1,\chi)$
convergence factor and then the absolutely convergent products
with $j\geq2$; no rearrangement of an unregularized product at
$1$ is needed. Positivity also follows directly because the
genus contains the standard lattice and all summands are positive.

On division by the expression in rank $n-1$, the factor $2$
and all factors with $j<n$ cancel. The remaining discriminant
power is $d_F^{(2n-1)/2}$. Since $\rho^f=d_F$, this is
$\rho^{nf}d_F^{-1/2}$, which proves \eqref{md-ratio}.
Replacing $\rho$ by $2v$ gives its last form.
\end{proof}

\subsection{The source field and Euler-product conventions}

The source initially specifies a totally imaginary number field $E$
with an involution whose image under every complex embedding is complex
conjugation. Its fixed field is indeed a totally real field $F$ with
$[E:F]=2$. The involution cannot be trivial, since that would make
every embedding real. Choose $\alpha$ not fixed by it. For $x\in E$,
the element
\[
 b=\frac{x-\overline x}{\alpha-\overline\alpha}
\]
is fixed, and $x-b\alpha$ is fixed as well. Thus $E=F+F\alpha$,
and the sum has dimension two. Every embedding of $F$ extends to
an embedding of $E$; its values on $F$ are real by the specified
conjugation identity. This proves the asserted CM interpretation.

\begin{proposition}\label{md-source-conventions}
Let complex conjugation on $E$ be denoted by a bar.
The following condition is equivalent to finite unramifiedness
of $E/F$:
for every maximal ideal $P\subset\mathcal O_E$ there exists
$\theta\in\mathcal O_E$ with $\theta-\overline\theta\notin P$.
Under this condition,
\[
 [E:\mathbb Q]=2f,\qquad |d_E|=d_F^2,
 \qquad |d_E|^{1/(2f)}=\rho.
\]
The product indexed by primes $P$ of $E$ with local factor
\[
 \begin{cases}
 (1-(-1)^n(NP)^{-n/2})^{-1},&P=\overline P,\\
 (1-(NP)^{-n})^{-1/2},&P\ne\overline P
 \end{cases}
\]
equals $L_n$ for $n\geq2$, where the positive square root is used.
\end{proposition}

\begin{proof}
At a ramified prime of a quadratic extension, conjugation acts
trivially on the residue field, so every integral
$\theta-\overline\theta$ belongs to $P$.
At an inert unramified prime, conjugation is the nontrivial
automorphism of its quadratic residue-field extension. Choose
a residue outside the base field and lift it to $\mathcal O_E$.
At a split prime, $P$ and $\overline P$ are distinct, and the
Chinese remainder theorem gives $\theta\equiv0\pmod P$ and
$\theta\equiv1\pmod{\overline P}$. In particular
$\overline\theta\equiv1\pmod P$. This proves both implications.

The discriminant formula in a tower gives
$|d_E|=d_F^2N(\mathfrak d_{E/F})=d_F^2$, and the degree
formula gives $[E:\mathbb Q]=2f$.
At an inert prime $v$, there is one $P$ and $NP=q_v^2$.
Its displayed factor is $(1-(-1)^nq_v^{-n})^{-1}$.
At a split prime there are two primes of norm $q_v$, and
their positive square-root factors multiply to
$(1-q_v^{-n})^{-1}$. These are exactly the $v$-factors of $L_n$.
Absolute convergence for $n\geq2$ justifies grouping them.
\end{proof}

\begin{corollary}\label{md-euler-bound}
For every integer $n\geq2$,
\[
                         |\log L_n|\leq f\log\zeta(n).
\]
Thus the recurrence and Euler-factor bound assumed in
Section~\ref{as-section} follow from the stated classical
mass formula with the actual source hypotheses.
\end{corollary}

\begin{proof}
For $0<x<1$ and $\epsilon\in\{1,-1\}$,
$|\log(1-\epsilon x)|\leq-\log(1-x)$; this follows from
$1+x\leq(1-x)^{-1}$ and monotonicity of the logarithm.
Apply this to every local factor to obtain
$|\log L_n|\leq\log\zeta_F(n)$.
Above a rational prime $p$, there are at most $f$ prime ideals
of $F$, each of norm at least $p$. Therefore
\[
 \prod_{v\mid p}(1-q_v^{-n})^{-1}
                             \leq(1-p^{-n})^{-f}.
\]
Multiplying the absolutely convergent products proves
$\zeta_F(n)\leq\zeta(n)^f$ and the assertion.
\end{proof}

The finite-field formula proves the source primitive-frame density
including all Gram coranks. The mass ratio removes the mass input
from the archimedean size calculation, subject only to the classical
global mass theorem explicitly identified above. The source-space
dimension and formal-degree inputs of that calculation are distinct
representation-theoretic assertions and are not inferred from the
mass ratio.


% Exact column dimensions and finite-group fixed spaces.
% Labels in this fragment use the prefix cd-.
\section{Column dimensions and finite-group fixed spaces}
\label{cd-section}

We prove the dimension formula for the actual decorated source
space, including its finite-level factors and its coefficient
dimension. We also compute the fixed spaces and simultaneous
kernels of the norm operators at the auxiliary primes.
These are finite-dimensional algebraic statements. The mass which
appears below is defined as a weighted sum of lattice classes;
its separate analytic product formula is not used in this section.
The definitions and prime conditions are those of the submitted
column construction.\cite{alpo-uniform-source}

\subsection{Evaluation of columns on double cosets}

Let $\Gamma,G_\infty,G_f$ be groups, with homomorphisms
$\iota_\infty:\Gamma\to G_\infty$ and
$\iota_f:\Gamma\to G_f$. Let $K\subseteq G_f$ be a subgroup,
$k$ a field, $\psi:K\to k^\times$ a character, and
$\pi:G_\infty\to\operatorname{GL}(V)$ a representation on a
finite-dimensional $k$-space. A column is a linear map
$T:V\to\{\text{functions }G_\infty\times G_f\to k\}$ such that
\begin{align}
 (Tv)(\iota_\infty(\gamma)g,\iota_f(\gamma)x)
       &=(Tv)(g,x),\notag\\
 (Tv)(g,x\kappa)&=\psi(\kappa)(Tv)(g,x),\notag\\
 (T\pi(h)v)(g,x)&=(Tv)(gh,x)
 \label{cd-column-equivariance}
\end{align}
for all indicated group elements. Denote this vector space by
$\mathcal C(\pi,K,\psi)$.

\begin{proposition}\label{cd-column-decomposition}
Suppose the double-coset set
$\mathcal X=\iota_f(\Gamma)\backslash G_f/K$ is finite and
\begin{equation}\label{cd-trivial-stabilizers}
     x^{-1}\iota_f(\gamma)x\in K\quad\Longrightarrow\quad\gamma=1
                            \qquad(x\in G_f,\ \gamma\in\Gamma).
\end{equation}
Choose representatives $s_c$, $c\in\mathcal X$. Evaluation gives
a linear isomorphism
\[
 \mathcal C(\pi,K,\psi)\longrightarrow\bigoplus_{c\in\mathcal X}V^*,
 \qquad T\longmapsto\bigl(v\mapsto(Tv)(1,s_c)\bigr)_c.
\]
In particular
\begin{equation}\label{cd-column-rank}
                  \dim_k\mathcal C(\pi,K,\psi)
                               =|\mathcal X|\dim_k V.
\end{equation}
The same conclusion holds if $K$ is contained in a larger subgroup
for which \eqref{cd-trivial-stabilizers} is known.
\end{proposition}
\begin{proof}
Every $x\in G_f$ can be written
$x=\iota_f(\gamma)s_c\kappa$. This expression is unique once
$c$ and its representative are fixed. Indeed, two such expressions
give
\[
 s_c^{-1}\iota_f(\gamma'^{-1}\gamma)s_c=\kappa'\kappa^{-1}\in K,
\]
so \eqref{cd-trivial-stabilizers} gives $\gamma=\gamma'$ and
then $\kappa=\kappa'$.
For any collection $L_c\in V^*$ define
\begin{equation}\label{cd-column-extension}
       (Tv)(g,\iota_f(\gamma)s_c\kappa)
          =\psi(\kappa)L_c\bigl(\pi(\iota_\infty(\gamma)^{-1}g)v\bigr).
\end{equation}
Uniqueness makes this well-defined. The homomorphism laws show
directly that it satisfies all three identities
\eqref{cd-column-equivariance}. Those identities also force
\eqref{cd-column-extension} for any given column. Thus extension
and evaluation are inverse. Containment in a larger subgroup
preserves \eqref{cd-trivial-stabilizers}.
\end{proof}

There is a useful convention check over $\mathbb C$. Suppose $V$
has a Hermitian form linear in its first argument, $\pi$ is
unitary, and $|\psi|=1$. The unique vectors $F(x)$ satisfying
$(Tv)(1,x)=\langle v,F(x)\rangle$ obey
\begin{equation}\label{cd-riesz-character}
 F(\iota_f(\gamma)x\kappa)
            =\psi(\kappa)^{-1}\pi(\iota_\infty(\gamma))F(x).
\end{equation}
Indeed \eqref{cd-column-equivariance} gives
\[
 (Tv)(1,\iota_f(\gamma)x\kappa)
     =\psi(\kappa)\langle\pi(\iota_\infty(\gamma)^{-1})v,F(x)\rangle,
\]
and unitarity and conjugate linearity in the second argument give
\eqref{cd-riesz-character}. The passage $T\mapsto F$ is
conjugate linear, not complex linear. With the opposite Hermitian
convention, the same identity is written
$(Tv)(1,x)=\langle F(x),v\rangle$. This is the source's convention.

Without \eqref{cd-trivial-stabilizers}, the evaluated covectors
must satisfy the actual stabilizer conditions
\[
 L_c\bigl(\pi(\iota_\infty(\gamma))v\bigr)
   =\psi(s_c^{-1}\iota_f(\gamma)s_c)^{-1}L_c(v)
 \quad\text{if }s_c^{-1}\iota_f(\gamma)s_c\in K.
\]
They cannot in general be replaced by all of $V^*$.

\subsection{The fixed-space count with its character hypotheses}

Let now $\Gamma,K,K_0,D$ be subgroups of a group $G$, with
$K,D\subseteq K_0$. Assume
\begin{equation}\label{cd-neat}
                         \Gamma\cap xK_0x^{-1}=\{1\}
                                             \qquad(x\in G).
\end{equation}
Let $W$ be a finite-dimensional $k$-representation $\rho$ of
$\Gamma$, and let $\chi:K\to k^\times$ be a character.
Write $\mathcal A_K$ for the functions $f:G\to W$ satisfying
\[
                         f(\gamma x\kappa)
                                =\chi(\kappa)\rho(\gamma)f(x).
\]
Suppose $D$ normalizes $K$, $\chi(\delta\kappa\delta^{-1})
=\chi(\kappa)$ for $\delta\in D$, and $\chi$ is trivial on
$D\cap K$. Denote by $\mathcal A_K^D$ those functions which
also satisfy $f(x\delta)=f(x)$ for all $\delta\in D$.

\begin{proposition}\label{cd-fixed-dimension}
If $\Gamma\backslash G/K$ is finite, then $[D:D\cap K]$ is
finite and
\begin{equation}\label{cd-fixed-dimension-eq}
                   \dim_k\mathcal A_K^D\,[D:D\cap K]
                                     =\dim_k\mathcal A_K.
\end{equation}
This holds in every characteristic, without dividing by the index.
\end{proposition}
\begin{proof}
The product $L=DK=KD$ is a subgroup of $K_0$. Define
$\widetilde\chi(\delta\kappa)=\chi(\kappa)$.
If $\delta\kappa=\delta'\kappa'$, then
$\kappa'\kappa^{-1}=\delta'^{-1}\delta\in D\cap K$,
so this definition is independent of the decomposition.
For multiplication, use
\[
 (\delta\kappa)(\delta'\kappa')
         =\delta\delta'(\delta'^{-1}\kappa\delta')\kappa'.
\]
The conjugation hypothesis makes $\widetilde\chi$ a character.
It extends $\chi$ and is trivial on $D$. The transformation
rules therefore identify $\mathcal A_K^D$ with the corresponding
space $\mathcal A_L$ of level $L$ and character
$\widetilde\chi$.

The projection
$\Gamma\backslash G/K\to\Gamma\backslash G/L$ is onto.
For a representative $x$ of a target class, its fiber consists of
the classes of $x\delta$, $\delta\in D$. Two such classes
coincide exactly when $\delta_1^{-1}\delta_2\in D\cap K$.
In fact, an equality gives
$\gamma x\delta_1\kappa=x\delta_2$ and hence
$x^{-1}\gamma x\in L\subseteq K_0$; equation
\eqref{cd-neat} forces $\gamma=1$.
The converse is immediate. Every fiber thus has cardinality
$[D:D\cap K]$, which is finite because the source is finite and
the target is nonempty. Evaluation on classes, exactly as in
\eqref{cd-column-extension} but with values in $W$, gives
\[
 \dim_k\mathcal A_K=\dim_k W\,|\Gamma\backslash G/K|,
 \qquad
 \dim_k\mathcal A_L=\dim_k W\,|\Gamma\backslash G/L|.
\]
The fiber count proves \eqref{cd-fixed-dimension-eq}.
\end{proof}

\subsection{The coefficient dimension from compact branching}

For a dominant integral tuple $\lambda$ of length $k$, let
$V_\lambda$ be the representation of Section~\ref{cb-section}.
Put $s_k=\prod_{j=0}^{k-1}j!$, with $s_0=1$.

\begin{proposition}[Weyl dimension formula]\label{cd-weyl-dimension}
For $k\ge1$,
\begin{equation}\label{cd-weyl-formula}
        \dim V_\lambda
            =\frac{\prod_{i<j}(\lambda_i-\lambda_j+j-i)}{s_k}.
\end{equation}
Equivalently, for $u_i=\lambda_i+(k+1)/2-i$,
$\dim V_\lambda=\prod_{i<j}(u_i-u_j)/s_k$.
The dimension is the same over any characteristic-zero coefficient
field in which the polynomial representation is defined.
\end{proposition}
\begin{proof}
A common determinant twist changes neither side. Subtract the
smallest entry of $\lambda$ so that all entries are nonnegative.
For $k=1$ the module is one-dimensional and the formula holds.
By the already proved compact branching law,
\[
 \dim V_\lambda
    =\sum_{\lambda_i\ge\mu_i\ge\lambda_{i+1}}\dim V_\mu.
\]
Assume the formula for length $k-1$. Put
$a_i=\lambda_i+k-i$ and $b_i=\mu_i+k-1-i$.
The summation ranges are independent integer intervals
\[
                          a_{i+1}\le b_i<a_i
                                        \qquad(1\le i<k).
\]
The $a_i$ are strictly decreasing nonnegative integers.
For $r=k-1$, the induction hypothesis equals
\[
       \dim V_\mu
            =\det\left[\binom{b_i}{r-j}\right]_{1\le i,j\le r}.
\]
To verify this determinant formula, the polynomial
$\binom{t}{m}$ has degree $m$ and leading coefficient $1/m!$.
Changing the columns from these polynomials to monomials is
triangular. The Vandermonde determinant then gives
$\prod_{i<j}(b_i-b_j)/\prod_{j=0}^{r-1}j!$, with the displayed
sign because the degrees run in decreasing order.

Multilinearity in the rows and the telescoping identity
\[
 \sum_{b=A}^{B-1}\binom bq
                         =\binom B{q+1}-\binom A{q+1}
\]
show that the sum over $\mu$ is
\[
 \det\left[
      \binom{a_i}{r-j+1}-\binom{a_{i+1}}{r-j+1}
              \right]_{1\le i,j\le r}.
\]
This is exactly
$\det[\binom{a_i}{k-j}]_{1\le i,j\le k}$: subtract row
$i+1$ from row $i$ for $i<k$, leaving the last row unchanged.
The row transformation has determinant one; the last column
then has zero entries except for its final entry one.
Applying the same Vandermonde calculation proves
\eqref{cd-weyl-formula}. The assertion about $u$ follows by
subtracting two of its entries. Finally the polynomial modules
are kernels of finite rational linear systems, so extension of a
characteristic-zero coefficient field preserves their dimensions.
\end{proof}

\subsection{Counting the actual decorated lattice classes}

Let $E$ be the CM field of degree $2f$ used in the construction,
and fix a positive definite Hermitian space of dimension $k\ge1$
over $E$. Consider its unimodular $\mathcal O_E$-lattices.
Choose one lattice $N_c$ in each isometry class and put
\begin{equation}\label{cd-mass-definition}
                   M_k=\sum_c\frac1{|\operatorname{Aut}(N_c)|}.
\end{equation}
The class set and the automorphism groups are finite by the
results of Section~\ref{al-section}.

In the actual application, $E/F$ is unramified at every finite
prime and the fixed space is $E^k$ with
$h(x,y)=\sum_i\bar x_i y_i$. Lemma~\ref{acl-local-norm} and
Proposition~\ref{acl-local-basis} show that every unimodular
lattice has an orthonormal basis at every finite completion,
including the split places. All these lattices therefore belong
to the genus of the standard lattice $\mathcal O_E^k$.
Consequently $M_k$ in \eqref{cd-mass-definition} is exactly
the standard genus mass used in Section~\ref{md-section}.

At each of $f$ selected primes of residue cardinality $g_1$ and
each of $f$ selected primes of residue cardinality $g_2$, choose
an isomorphism $N/PN\simeq\mathbb F_{g_i}^{k}$, modulo left
composition by the upper triangular unipotent group
$U_k(\mathbb F_{g_i})$. At each of $f$ further selected primes
of residue cardinality $T$, choose a nonzero linear functional
\[
                         N/PN\longrightarrow\mathbb F_T.
\]
All these primes are distinct. The actual parameters have
$g_1,g_2,T$ prime, with $g_1\ne g_2$ and $T\ne g_i$.
Each selected collection has $f$ elements, and contains only
the chosen primes from the conjugate pairs. A local basis at an
auxiliary level prime is taken modulo the indicated group, not
modulo a Borel subgroup or the scalar matrices.

Let $\mathcal D_T$ be the resulting isometry classes of decorated
lattices, and put
\[
 G_k(q)=\prod_{j=0}^{k-1}(q^k-q^j)
       =q^{k^2}\prod_{j=1}^k(1-q^{-j}),\qquad
 c_k(q)=\frac{G_k(q)}{q^{k(k-1)/2}}.
\]

\begin{lemma}\label{cd-decoration-count}
All decorated stabilizers are trivial, and
\begin{equation}\label{cd-classes-mass}
             |\mathcal D_T|
                 =M_k\bigl(c_k(g_1)c_k(g_2)\bigr)^f(T^k-1)^f.
\end{equation}
\end{lemma}
\begin{proof}
For a full projective rank-$k$ lattice, localization at $P$
is free of rank $k$. Reduction therefore gives
$N/PN\simeq\mathbb F_q^k$, regardless of whether $N$ is
globally free. Its ordered bases form a torsor for
$\operatorname{GL}_k(\mathbb F_q)$. Choosing independent
columns successively gives $G_k(q)$ bases. The group $U_k$
has $q^{k(k-1)/2}$ elements, because its entries strictly above
the diagonal are arbitrary and all other entries are fixed.
Its action on the bases is free, giving exactly $c_k(q)$
choices modulo $U_k$. The dual vector space has $T^k-1$
nonzero vectors, giving that number of nonzero functionals.
Choices at the distinct primes are independent. Thus every
$N_c$ has
\[
                 A=\bigl(c_k(g_1)c_k(g_2)\bigr)^f(T^k-1)^f
\]
possible decorations.

An isometry preserving one decoration has, in the corresponding
local bases, unipotent reductions at both residue characteristics
$g_1,g_2$. It has finite order because it preserves a positive
definite lattice. Proposition~\ref{al-two-primes} makes it the
identity. The same conclusion follows from just the first prime
under the source's additional inequality $2fk<g_1-1$, by the
single-prime argument following that proposition.
The action of $\operatorname{Aut}(N_c)$ on the decorations is
therefore free. There are $A/|\operatorname{Aut}(N_c)|$ orbits
over $N_c$. Summing over $c$ proves \eqref{cd-classes-mass}.
\end{proof}

To incorporate a character of the upper unipotent groups, one
can instead retain the full local bases and let the product of
these groups act along with the rational isometries, with that
character on the coefficient space. The action on the full
decorated objects is free: its rational part is the identity
by the preceding proof, and then the local matrices are the
identity because the bases are isomorphisms. Evaluation therefore
still gives one copy of the coefficient space per class. Thus
the finite-level character does not change the dimension.
The character orientation remains the inverse level convention
proved in Section~\ref{acl-section}; no invariant restriction has
been discarded before proving stabilizer triviality.

\begin{theorem}[The exact source-space dimension]\label{cd-source-dimension}
At each of the $f$ coefficient embeddings let
$\lambda^{(v)}$ be the dominant integral highest weight and
$u_{v,i}=\lambda^{(v)}_i+(k+1)/2-i$. For the actual column
space $\mathcal M_{\epsilon,T}$ with these coefficient modules
and the decorations just specified,
\begin{equation}\label{cd-source-formula}
 \dim\mathcal M_{\epsilon,T}
   =M_k\bigl(c_k(g_1)c_k(g_2)\bigr)^f
           \prod_v\frac{\prod_{i<j}(u_{v,i}-u_{v,j})}{s_k}
                           (T^k-1)^f.
\end{equation}
The same formula holds over the number field of coefficients.
\end{theorem}
\begin{proof}
The class evaluation of Proposition~\ref{acl-class-evaluation},
or equivalently Proposition~\ref{cd-column-decomposition},
gives the number of decorated classes times the dimension of
the coefficient space. That space is the tensor product over
the $f$ embeddings; its dimension is the product of the individual
dimensions. Apply Proposition~\ref{cd-weyl-dimension} and
Lemma~\ref{cd-decoration-count}. Field extension preserves the
dimension, as in Proposition~\ref{cd-weyl-dimension}.
\end{proof}

For the actual weights, the source identifies $u_{v,i}$ with its
Harish--Chandra parameters. The source's zero-based sequence
$\lambda_i-i$ differs from $u_{v,i}$ by one common constant,
which cancels from all differences. Equation
\eqref{cd-source-formula} is therefore precisely
\eqref{as-column-assumption}, with the mass defined by
\eqref{cd-mass-definition}. The integrality of the highest weights
is part of the actual construction; an arbitrary real tuple in
a scalar size estimate need not define a representation.

\subsection{The auxiliary residue-field actions}

Let
$D=\prod_{P\in S_T}(\mathcal O_E/P)^\times$
act by scaling the chosen nonzero functionals, leaving the lattice
and every other decoration unchanged. This action commutes with
isometries and with the changes of local bases.

\begin{lemma}\label{cd-diamond-free}
The action of $D$ on $\mathcal D_T$ is free. After choosing one
class in each $D$-orbit, the column space is, as a $D$-module,
the space of functions $D\to W_0$ with translation action,
for a finite-dimensional coefficient space $W_0$.
\end{lemma}
\begin{proof}
If scaling by $a\in D$ preserves a class, some rational
isometry relates the original decoration to the scaled one.
It preserves the lattice and the local-basis cosets at $g_1,g_2$,
so the preceding stabilizer argument makes it the identity.
Each remaining equality is $a_P\ell_P=\ell_P$ for a nonzero
functional $\ell_P$, forcing $a_P=1$. Thus $a=1$.

Choose representatives of the $D$-orbits and, for every $a\in D$,
use the actual object obtained by scaling such a representative
by $a$. Freeness makes these distinct isometry classes. Evaluation
identifies a column with an arbitrary coefficient vector at each
of them. Scaling permutes these values by translation in $D$;
the lattice and the coefficient vector spaces themselves do not
change. Taking $W_0$ to be the direct sum of the coefficient
spaces at the chosen orbit representatives proves the assertion.
\end{proof}

For a factor $D_P$, let $R_a$ be translation by $a$ on this
function space and put $N_P=\sum_{a\in D_P}R_a$.

\begin{proposition}\label{cd-norm-kernel-dimensions}
The fixed space for one $D_P$ and the simultaneous kernel satisfy
\begin{align}
       \dim\mathcal M_{\epsilon,T}^{D_P}\,(T-1)
                         &=\dim\mathcal M_{\epsilon,T},
                                      \label{cd-one-fixed}\\
       \dim\bigcap_{P\in S_T}\ker N_P\,(T-1)^f
                         &=\dim\mathcal M_{\epsilon,T}\,(T-2)^f.
                                      \label{cd-all-kernels}
\end{align}
The counting argument is valid over every coefficient field.
\end{proposition}
\begin{proof}
On functions on a finite group $D_P$, translation-invariant
functions are exactly the constants. The norm operator sends
a function to the constant function whose value is the sum
of all its values. The sum functional is onto in every
characteristic: a function supported at one point can have
arbitrary sum. Its kernel therefore has dimension $|D_P|-1$.

For several factors, identify the space of functions with
$W_0\otimes_k\bigotimes_P k^{D_P}$. Write
\[
 k^{D_P}=\ker\left(\sum:k^{D_P}\to k\right)\oplus k\delta_1,
\]
where $\delta_1$ is the indicator of the identity element and
has sum one. Each $N_P$ has the same kernel as this sum map
on its tensor factor. The simultaneous kernel is consequently
$W_0$ tensored with the product of all the displayed kernels.
Its dimension is $\dim W_0\prod_P(|D_P|-1)$, whereas the whole
space has dimension $\dim W_0\prod_P|D_P|$.
Finally $|D_P|=T-1$ and $|S_T|=f$. This proves both formulas.
For $T=2$, each factor is trivial and the norm kernel is zero,
exactly as the second formula states.
\end{proof}

The trivial-stabilizer and character conditions in these results
are substantive. For example, take $G=\Gamma=K_0=D$ cyclic of
order two, $K=\{1\}$, and the trivial one-dimensional coefficient
representation. Then $\mathcal A_K$ and its $D$-fixed space both
have dimension one, but $[D:D\cap K]=2$; formula
\eqref{cd-fixed-dimension-eq} would be false without
\eqref{cd-neat}. Separately, take $G=K=K_0=D$ cyclic of order
two, $\Gamma=\{1\}$, and the nontrivial character over
$\mathbb Q$. Now stabilizers are trivial, but $\mathcal A_K$
has dimension one and its $D$-fixed space is zero, so triviality
of $\chi$ on $D\cap K$ cannot be omitted. Both examples are
excluded by the actual hypotheses.

These arguments prove the source-space dimension and the
finite-group fixed-space counts. They do not evaluate the mass
\eqref{cd-mass-definition}, its recurrence in $k$, an
archimedean formal degree, or an ambient automorphic dimension.


\section{The formal degree with the specified Haar measure}
\label{fd-section}

We determine the formal degree in the normalization used by the mass
formula. This also checks the meaning of the representation parameter:
the source specifies a smallest compact type, rather than writing an
infinitesimal character into its definition of a discrete series.
The calculation below uses classical discrete-series theorems, stated
explicitly. The measure comparison, scalar integral, and the application
to the source's precise definition are proved here.

For an integer $a\ge0$, write
\[
                     s_a=\prod_{j=0}^{a-1}j!,\qquad s_0=1.
\]
Let $p,q\ge1$, put $n=p+q$, and set
\[
 G=U(p,q)=\{g\in\operatorname{GL}_n(\mathbb C):g^*Jg=J\},
 \quad J=\operatorname{diag}(I_p,-I_q),\quad
 K=U(p)\times U(q).
\]
All Lebesgue measures on complex matrices below are the products
of $d(\Re z)\,d(\Im z)$ over their entries. Haar measure on a compact
group is a probability measure unless stated otherwise.

\subsection{Two elementary matrix integrals}

Put
\[
 \mathcal B_{p,q}=\{Z\in M_{p,q}(\mathbb C):I_q-Z^*Z>0\},
\]
where $>0$ means positive definite.

\begin{lemma}\label{fd-matrix-integrals}
For $\alpha>-1$ and $u>p+q-1$, respectively,
\begin{align}
 \int_{\mathcal B_{p,q}}\det(I_q-Z^*Z)^\alpha\,dZ
   &=\pi^{pq}\prod_{j=0}^{q-1}
                  \frac{\Gamma(\alpha+j+1)}
                       {\Gamma(\alpha+p+j+1)},\label{fd-ball-integral}\\
 \int_{M_{p,q}(\mathbb C)}\det(I_q+Z^*Z)^{-u}\,dZ
   &=\pi^{pq}\prod_{j=0}^{q-1}
                  \frac{\Gamma(u-p-j)}{\Gamma(u-j)}.
                                                   \label{fd-affine-integral}
\end{align}
Here $\Gamma(v)=\int_0^\infty t^{v-1}e^{-t}\,dt$ for $v>0$.
\end{lemma}
\begin{proof}
First, for a vector $z\in\mathbb C^p$, radial integration gives
\begin{align*}
 \int_{|z|<1}(1-|z|^2)^\alpha\,dz
       &=\pi^p\frac{\Gamma(\alpha+1)}
                         {\Gamma(\alpha+p+1)},\\
 \int_{\mathbb C^p}(1+|z|^2)^{-u}\,dz
       &=\pi^p\frac{\Gamma(u-p)}{\Gamma(u)}\quad(u>p).
\end{align*}
Indeed the radial volume element is
$2\pi^p t^{2p-1}dt/\Gamma(p)$. Its constant follows by
integrating $e^{-|z|^2}$ in Cartesian and polar coordinates.
The substitution $x=t^2$ reduces the two integrals to the
beta integral. The identity
$B(a,b)=\Gamma(a)\Gamma(b)/\Gamma(a+b)$ follows by setting
$r=x+y$, $v=x/(x+y)$ in the product of the two defining
gamma integrals. All integrands are nonnegative, so the changes
of variables and repeated integration follow from Tonelli's theorem.

Write a matrix as $Z=[W,z]$, with $q-1$ columns in $W$.
The Schur complement gives
\[
 \det(I_q\pm Z^*Z)
  =\det(I_{q-1}\pm W^*W)
       \bigl(1\pm z^*(I_p\pm WW^*)^{-1}z\bigr).
 \tag{FD1}
\]
For the minus sign, membership in $\mathcal B_{p,q}$ is
equivalent to $W\in\mathcal B_{p,q-1}$ and
$z^*(I_p-WW^*)^{-1}z<1$. The equivalence of
$I_{q-1}-W^*W>0$ and $I_p-WW^*>0$ follows by comparing
their nonzero singular values.
In either case substitute $z=(I_p\pm WW^*)^{1/2}v$.
The real Jacobian is $\det(I_p\pm WW^*)$, which equals
$\det(I_{q-1}\pm W^*W)$. The equality follows, for example,
by expanding $\det\bigl(\begin{smallmatrix}I&W\\
\mp W^*&I\end{smallmatrix}\bigr)$ in its two Schur complements.

If the two integrals are denoted by $I_{p,q}(\alpha)$ and
$J_{p,q}(u)$, this gives
\begin{align*}
 I_{p,q}(\alpha)
   &=\pi^p\frac{\Gamma(\alpha+1)}{\Gamma(\alpha+p+1)}
                                      I_{p,q-1}(\alpha+1),\\
 J_{p,q}(u)
   &=\pi^p\frac{\Gamma(u-p)}{\Gamma(u)}
                                      J_{p,q-1}(u-1).
\end{align*}
At $q=0$ both integrals are one. Iteration proves the formulas,
including convergence under exactly the indicated sufficient
conditions.
\end{proof}

\subsection{The compact-dual normalization}

For a block matrix
$g=\bigl(\begin{smallmatrix}A&B\\C&D\end{smallmatrix}\bigr)\in G$,
the action on $\mathcal B_{p,q}$ is
\[
                 g\cdot Z=(AZ+B)(CZ+D)^{-1}.
\]
The denominator is invertible: otherwise a nonzero vector
$(Zv,v)$ of strictly negative norm would be carried to a
vector with second component zero, whose norm is nonnegative.
The stabilizer of $0$ is $K$, and the action is transitive.
For instance, the positive Hermitian block matrix
\[
 g_Z=
 \begin{pmatrix}
 (I_p-ZZ^*)^{-1/2}&Z(I_q-Z^*Z)^{-1/2}\\
 Z^*(I_p-ZZ^*)^{-1/2}&(I_q-Z^*Z)^{-1/2}
 \end{pmatrix}
 \tag{FD2}
\]
belongs to $G$ and carries $0$ to $Z$. The needed identities
between its off-diagonal blocks follow by a singular-value
decomposition of $Z$.

\begin{lemma}\label{fd-quotient-measure}
The measure
\[
                   d\nu_0(Z)=\det(I_q-Z^*Z)^{-n}\,dZ
 \tag{FD3}
\]
is $G$-invariant. It defines a Haar measure $\mu_0$ on $G$ by
\[
       \int_G F(g)\,d\mu_0(g)
          =\int_{\mathcal B_{p,q}}\int_K F(g_Zk)\,dk\,d\nu_0(Z).
 \tag{FD4}
\]
The pushforward of Haar probability measure on $U(n)$ to
$U(n)/(U(p)\times U(q))$ has, in its graph chart
$M_{p,q}(\mathbb C)$, density
\[
             c_{p,q}\det(I_q+Z^*Z)^{-n}\,dZ,\qquad
             c_{p,q}=\frac{s_n}{\pi^{pq}s_ps_q}.
 \tag{FD5}\label{fd-cpq}
\]
\end{lemma}
\begin{proof}
Put $Z'=g\cdot Z$ and $R=CZ+D$. Directly from $g^*Jg=J$,
\[
                  I_q-Z'^*Z'=R^{-*}(I_q-Z^*Z)R^{-1}.
 \tag{FD6}
\]
The complex derivative sends a variation $H$ to
$(A-Z'C)HR^{-1}$. The block determinant identity
\[
                         \det(A-Z'C)\det R=\det g
\]
follows by eliminating the upper-right block after replacing
the second block column of $g$ by its product with
$\bigl(\begin{smallmatrix}Z\\I_q\end{smallmatrix}\bigr)$.
The complex determinant of the derivative is
$(\det g)^q\det R^{-n}$; its real Jacobian is
$|\det R|^{-2n}$, since $|\det g|=1$.
Together with (FD6) this proves invariance of (FD3).
Changing the section in (FD4) multiplies it on the right by an
element of $K$, which does not change the inner integral.
Left translation by $G$ changes the section in this way and
preserves $\nu_0$, proving the Haar assertion.

The compact quotient is the Grassmannian of $q$-planes
in $\mathbb C^{p+q}$. Its graph chart sends $Z$ to the
column space of $\bigl(\begin{smallmatrix}Z\\I_q\end{smallmatrix}\bigr)$.
The same calculation for $g\in U(n)$ has plus signs:
\[
                   I_q+Z'^*Z'=R^{-*}(I_q+Z^*Z)R^{-1}.
 \tag{FD7}
\]
It therefore determines the invariant density in the chart
up to a positive constant. Haar measure gives a smooth
positive density in local coordinates. The complement of
the graph chart is the zero set of a nonzero Pl\"ucker
coordinate, and has measure zero. To justify this last
assertion without a measure convention, cover the
Grassmannian by its finitely many graph charts; a nonzero
complex polynomial in each chart vanishes on a set of
Lebesgue measure zero, by induction on the number of
variables and Fubini's theorem.
Formula \eqref{fd-affine-integral} at $u=n$ gives total mass
\[
 \pi^{pq}\prod_{j=0}^{q-1}\frac{(q-j-1)!}{(n-j-1)!}
                        =\pi^{pq}\frac{s_ps_q}{s_n}.
\]
Normalizing the total mass to one proves (FD5).
\end{proof}

We now specify the Lie-algebra density used in the source.
For either real form $\mathfrak u(p,q)$ or $\mathfrak u(n)$
inside $M_n(\mathbb C)$, let $\omega_{\mathfrak g}$ be the
positive Lebesgue measure characterized by
\[
 (X,Y)\longmapsto X+iY
       \quad\hbox{carrying}\quad
 \omega_{\mathfrak g}\otimes\omega_{\mathfrak g}
       \quad\hbox{to ordinary Lebesgue measure on }M_n(\mathbb C).
 \tag{FD8}
\]
This map is a real-linear isomorphism. On the diagonal write
$X_{ii}=ix_i$ with $x_i\in\mathbb R$; off the diagonal
use one complex coordinate $z_{ij}$ for each $i<j$.
The other entry is $-\varepsilon_i\varepsilon_j\bar z_{ij}$,
where $\varepsilon_i$ is the sign of the Hermitian form.
For each pair $i<j$, the real determinant in (FD8) has
absolute value $4$; on each diagonal coordinate it has
absolute value $1$. Consequently
\[
 \omega_{\mathfrak g}
    =2^{n(n-1)/2}\prod_i dx_i
                  \prod_{i<j}d(\Re z_{ij})\,d(\Im z_{ij}).
 \tag{FD9}
\]
This also proves uniqueness. In particular the density is
the same in these coordinates for the two real forms.

Call a Haar measure on $G$ compact-dual normalized if its
density at the identity relative to $\omega_{\mathfrak u(p,q)}$
equals that of Haar probability measure on $U(n)$ relative
to $\omega_{\mathfrak u(n)}$. Density may equivalently be
defined by the limit of the ratio of the measures of
exponential images of shrinking coordinate balls: the
Jacobian is continuous and the derivative of the exponential
map at zero is the identity.

\begin{proposition}\label{fd-measure-match}
There is exactly one compact-dual normalized Haar measure, and it is
\[
                             \mu_{\rm cd}=c_{p,q}\mu_0.
 \tag{FD10}\label{fd-cd-measure}
\]
This is the source's measure specified by its compact
probability and equal exponential-density conditions.
\end{proposition}
\begin{proof}
For both groups the decomposition at the identity into
the block-diagonal compact algebra and its complementary
off-diagonal coordinates has
\[
          \omega_{\mathfrak g}
                =2^{pq}\omega_{\mathfrak k}\,dZ.
 \tag{FD11}
\]
Indeed (FD9) contributes one additional factor $2$ for
each of the $pq$ cross-block complex entries.
Let $b_K$ be the density of Haar probability measure on
$K$ relative to $\omega_{\mathfrak k}$.
For $G$, the local product map $(Z,k)\mapsto g_Zk$ has
derivative the direct-sum identification at $(0,1)$.
By (FD3)--(FD4), the identity density of $\mu_0$ is
$b_K/2^{pq}$.
For the compact group use the analogous section obtained
by orthonormalizing the graph of $Z$; its derivative has
off-diagonal blocks $Z,-Z^*$ and the same coordinate
Jacobian. Formula (FD5) gives identity density
$c_{p,q}b_K/2^{pq}$.
Their ratio proves (FD10). All Haar measures on a fixed
Lie group are positive scalar multiples of one another,
so the density also proves uniqueness.
\end{proof}

\subsection{Classical representation-theoretic inputs}

For a strongly continuous irreducible unitary discrete-series
representation $\pi$ of $G$, its formal degree with respect to
a Haar measure $\mu$ is the positive number $d_\mu(\pi)$ in
Schur orthogonality; in particular
\[
  \int_G|\langle\pi(g)v,w\rangle|^2\,d\mu(g)
      =d_\mu(\pi)^{-1}\|v\|^2\|w\|^2.
 \tag{FD12}
\]
It follows directly that
$d_{a\mu}(\pi)=a^{-1}d_\mu(\pi)$ for $a>0$.
A discrete series is an irreducible unitary representation
occurring as a closed subrepresentation of $L^2(G,\mu)$.

We use the following classical facts. They are background
theorems, not new deductions or unproved assertions from
the submitted archive.

\begin{enumerate}
\item Harish--Chandra's classification and formal-degree theorem.
Discrete series have regular integral Harish--Chandra parameters,
considered modulo the compact Weyl group.
For fixed $G$ and $\mu$ there is a positive constant $b_\mu$,
independent of the parameter and its chamber, such that
\[
                    d_\mu(\pi_\lambda)
                       =b_\mu\prod_{i<j}|\lambda_i-\lambda_j|.
 \tag{FD13}\label{fd-classical-hc}
\]
Here the integrality convention for $U(p,q)$ is
$\lambda_i-(n-1)/2\in\mathbb Z$.
The constant is left undetermined at this stage.
\cite[Theorem 16, \S41]{harish-chandra-ds-ii}
\item Blattner's formula: choose the positive roots making
$\lambda$ dominant, let $\rho_{\rm c}$ and $\rho_{\rm n}$
be the half sums of its compact and noncompact positive
roots, and let $P_{\rm n}(\beta)$ count the expressions of
$\beta$ as a sum of noncompact positive roots. For a
compact-dominant integral highest weight $\mu$, its
multiplicity is
\[
 m_\mu(\pi_\lambda)=
  \sum_{w\in W_K}\operatorname{sgn}(w)
       P_{\rm n}\bigl(w(\mu+\rho_{\rm c})
                                -\lambda-\rho_{\rm n}\bigr).
 \tag{FD14}\label{fd-blattner}
\]
The partition function is finite because these roots lie
in an open half-space.
\cite[III.1, theorem (2)]{duflo-heckman-vergne}
The stated highest-weight convention is also checked
against the explicit $U(p,q)$ convention in
\cite[\S2.2]{baldoni-vergne}.
\item The characterization by a compact type
\cite[Theorem 1.1]{parthasarathy-minimal}:
put $\nu=\lambda-\rho$, with $\rho=\rho_{\rm c}+\rho_{\rm n}$,
and suppose $\nu$ is dominant integral. An irreducible
unitary representation containing the compact type of
highest weight $\nu+2\rho_{\rm n}$, and containing no
compact-dominant type $\nu+2\rho_{\rm n}-\alpha$ for a
noncompact positive root $\alpha$, is $\pi_\lambda$.
\end{enumerate}

The semisimple versions apply to $U(p,q)$ through the
finite central covering $SU(p,q)\times U(1)\to U(p,q)$.
An irreducible unitary representation has a scalar
central character. Its restriction to $SU(p,q)$ is
irreducible, since a subspace invariant under that group
is then invariant under the whole group.
The character on the compact center has absolute value one;
probability measure on that center changes no
matrix-coefficient norm. The possible fixed measure factor
in the covering is absorbed in $b_\mu$ in (FD13).
The compact-type and root statements restrict to the
trace-zero hyperplane; all root differences have sum zero,
so adding the fixed central weight leaves every comparison
unchanged. The complexification of $SU(p,q)$ is $SL_n(\mathbb C)$,
which is simply connected, as required in the third input.

\subsection{An explicit scalar discrete series fixes the constant}

Let $s$ be an integer with $s>n-1$, and put
$\alpha=s-n>-1$. Consider the Hilbert space
\[
 \mathcal H_s=\left\{f\text{ holomorphic on }\mathcal B_{p,q}:
        \|f\|_s^2=
         \int_{\mathcal B_{p,q}}|f(Z)|^2
                         \det(I_q-Z^*Z)^{s-n}\,dZ<\infty\right\}.
 \tag{FD15}
\]
It is a closed subspace of the indicated $L^2$ space.
To see this, the mean-value estimate on a small ball about
any interior point, where the weight is bounded below,
controls point evaluation uniformly on compact subsets.
An $L^2$-Cauchy sequence of holomorphic functions therefore
converges locally uniformly to a holomorphic function,
which represents its $L^2$ limit. In particular every
point evaluation is continuous.

For $g\in G$ define $j(g,Z)=\det(CZ+D)$, and set
\[
       (\pi_s(g)f)(Z)=j(g^{-1},Z)^{-s}f(g^{-1}\cdot Z).
 \tag{FD16}
\]
The exponent is integral, so no choice of branch is involved.

\begin{lemma}\label{fd-scalar-model}
The representation $\pi_s$ is irreducible, unitary and
belongs to the discrete series. If
$I_s=I_{p,q}(s-n)$ is the integral in
\eqref{fd-ball-integral}, then
\[
                    d_{\mu_0}(\pi_s)=I_s^{-1}
                      =\pi^{-pq}
                        \prod_{i=1}^p\prod_{j=1}^q(s-i-j+1).
 \tag{FD17}\label{fd-scalar-degree}
\]
\end{lemma}
\begin{proof}
The identity
$j(gh,Z)=j(g,h\cdot Z)j(h,Z)$ proves the representation
law. Equations (FD6) and the Jacobian calculation prove
unitarity after a change of variables in (FD15).
The representation is strongly continuous: on a compact
subset of the ball the transformed holomorphic function
converges uniformly as $g\to1$. In the inner product with
a fixed $f$, the integral over the complementary tail is
uniformly bounded by its $L^2$ norm times $\|f\|_s$.
Letting that tail tend to zero proves
$\langle\pi_s(g)f,f\rangle\to\|f\|_s^2$, and unitarity
then gives norm convergence.

The constant function $1$ has norm squared $I_s$.
Averaging over the scalar rotations $Z\mapsto e^{it}Z$
and using the one-variable Cauchy formula gives
\[
                         \langle f,1\rangle_s=I_s f(0),
 \tag{FD18}
\]
with the inner product linear in the first argument.
Fubini is legitimate because $I_s<\infty$ and
$\int|f|\,\det(I-Z^*Z)^{s-n}\le\sqrt{I_s}\|f\|_s$.
Put $v=I_s^{-1/2}1$. For $Z=g\cdot0$, (FD16)--(FD18) give
\[
 |\langle f,\pi_s(g)v\rangle_s|^2
       =I_s|j(g,0)|^{-2s}|f(Z)|^2
       =I_s\det(I_q-Z^*Z)^s|f(Z)|^2.
 \tag{FD19}
\]
Thus the translates of $v$ span a dense subspace: a function
orthogonal to all of them vanishes at every point.

For $k=\operatorname{diag}(A,D)\in K$, the constant
has character $\det(D)^s$. The subspace with this
character is one-dimensional. Indeed its functions
are invariant under $Z\mapsto e^{it}Z$, and their
holomorphic power series have no terms of positive degree.
An orthogonal projection onto a closed invariant
subspace commutes with $K$, hence sends $v$ to either
$0$ or $v$. It then has the same alternative on every
translate of $v$, proving irreducibility by density.

Integrating (FD19) with (FD3)--(FD4) gives
\[
               \int_G|\langle f,\pi_s(g)v\rangle_s|^2\,d\mu_0(g)
                                  =I_s\|f\|_s^2.
 \tag{FD20}
\]
Consequently $f\mapsto I_s^{-1/2}
\langle f,\pi_s(\,\cdot\,)v\rangle_s$ is an isometric
intertwining map into the left regular representation,
with closed image. This proves the discrete-series claim.
Taking $f=v$ in (FD20) and comparing with (FD12) gives
$d_{\mu_0}(\pi_s)=I_s^{-1}$.
Finally, reversal of the factors in
\eqref{fd-ball-integral} gives
\[
 I_s=\pi^{pq}\prod_{j=0}^{q-1}
                       \frac{\Gamma(s-p-j)}{\Gamma(s-j)}.
\]
The gamma recurrence, obtained by integration by parts,
now proves (FD17).
\end{proof}

Its Harish--Chandra parameter is also explicit. In the
complexified Lie algebra with upper triangular positive
root spaces, differentiating (FD16) shows that the
constant vector is killed by every positive root vector
and has highest weight
\[
                 \Lambda=(\underbrace{0,\ldots,0}_{p},
                            \underbrace{s,\ldots,s}_{q}).
 \tag{FD21}
\]
For a cross-block upper root vector the derivative is a
negative coordinate derivative, so it kills a constant;
compact upper root vectors kill it as well.
The diagonal calculation gives $0$ in the first block
and $s$ in the second. The usual highest-weight
infinitesimal character is $\Lambda+\rho$, where
$\rho_i=(n+1)/2-i$. This follows from the
Harish--Chandra isomorphism, or from the ordered Capelli
eigenvalue calculation on a highest vector.
Only the unordered tuple is needed here.
Its within-block Vandermonde factors are $s_p,s_q$, and
its cross-block factors have absolute values
$s-p-j+i$ for $1\le i\le p$, $1\le j\le q$.
All are positive because $s>n-1$. Hence
\[
 \prod_{i<j}|(\Lambda+\rho)_i-(\Lambda+\rho)_j|
       =s_ps_q\prod_{i=1}^p\prod_{j=1}^q(s-i-j+1).
 \tag{FD22}
\]

\begin{theorem}\label{fd-formula}
For every discrete series of $U(p,q)$ with
Harish--Chandra parameter $\lambda$,
\begin{equation}\label{fd-exact}
            d_{\mu_{\rm cd}}(\pi_\lambda)
                 =\frac{\prod_{i<j}|\lambda_i-\lambda_j|}{s_n}.
\end{equation}
No extra lower bound on the noncompact root gaps is needed
in this classical argument.
\end{theorem}
\begin{proof}
Comparing (FD17), (FD22) and the parameter-independent
constant in (FD13) gives
$b_{\mu_0}=\pi^{-pq}/(s_ps_q)$.
Divide by $c_{p,q}$ using (FD10) and the scaling
rule following (FD12). The result is $b_{\mu_{\rm cd}}=1/s_n$.
\end{proof}

\paragraph{Compact signatures.}
The same formula holds when $pq=0$. In this case the group
is $U(n)$, and its compact-dual normalized Haar measure is
its probability measure. An irreducible unitary representation
of a compact group is finite dimensional. Schur orthogonality
for probability measure makes its formal degree its dimension.
If $\mu$ is its dominant integral highest weight and
$\lambda=\mu+\rho$, Weyl's dimension formula gives
\[
 \dim V_\mu
   =\prod_{i<j}\frac{\mu_i-\mu_j+j-i}{j-i}
   =\frac{\prod_{i<j}(\lambda_i-\lambda_j)}{s_n}.
\]
The equality $\prod_{i<j}(j-i)=s_n$ follows by grouping
the factors with fixed $j$. This is also the dimension
formula proved in Section~\ref{cd-section}.
For $n=0$ the group is trivial, its irreducible representation
is one dimensional, and both the empty Vandermonde and
$s_0$ are one. Thus its formal degree is one.

\subsection{Checking the source's compact-type definition}

Now allow the signs $\varepsilon_i\in\{1,-1\}$ to be in
any order, and suppose
\[
             \lambda_1>\cdots>\lambda_n,\qquad
             \lambda_i-(n-1)/2\in\mathbb Z.
 \tag{FD23}
\]
Positive roots are $e_i-e_j$ for $i<j$. A positive root
is compact if $\varepsilon_i=\varepsilon_j$ and
noncompact otherwise. Use their half sums
$\rho_{\rm c},\rho_{\rm n}$ and put
\[
 \rho=\rho_{\rm c}+\rho_{\rm n},\quad
 \rho_i=(n+1)/2-i,\quad
 \mu_0=\lambda+\rho_{\rm n}-\rho_{\rm c}.
 \tag{FD24}
\]
The compact Weyl group permutes the indices of each sign.
A compact highest weight $\mu$ is dominant when
$\mu_i\ge\mu_j$ for $i<j$ of the same sign.
The squared norm used in the source is
$\|\mu+2\rho_{\rm c}\|^2=\sum_i(\mu_i+2\rho_{{\rm c},i})^2$.

\begin{proposition}\label{fd-lowest-type}
The classical representation $\pi_\lambda$ contains the
compact type $\mu_0$ with multiplicity one. Every other
compact highest weight $\mu$ occurring in it satisfies
\[
             \|\mu+2\rho_{\rm c}\|^2
                        >\|\mu_0+2\rho_{\rm c}\|^2.
 \tag{FD25}
\]
Conversely, an irreducible unitary representation containing
the type $\mu_0$ and satisfying the corresponding weak
minimum inequality for all its compact types is
$\pi_\lambda$.
\end{proposition}
\begin{proof}
If $pq=0$, then $G=K$ and $\mu_0=\lambda-\rho$.
The corresponding compact representation has just this
one compact type, and any irreducible representation
containing it is that representation. This includes the
trivial group when $n=0$. We may therefore assume $pq>0$.
The strict integral gaps imply $\lambda_i-\lambda_j\ge j-i$.
For indices of the same sign let $b_{ij}$ count the
opposite-sign indices strictly between them. Directly
from the half sums,
\[
 \rho_{{\rm c},i}-\rho_{{\rm c},j}=j-i-b_{ij},\qquad
 \rho_{{\rm n},i}-\rho_{{\rm n},j}=b_{ij}.
\]
Thus
$\mu_{0,i}-\mu_{0,j}=\lambda_i-\lambda_j-(j-i)+2b_{ij}\ge0$.
Also $\mu_0=\lambda+\rho-2\rho_{\rm c}$ is integral:
$\lambda_i+\rho_i\in\mathbb Z$ and
$2\rho_{{\rm c},i}\in\mathbb Z$.
The tuple $\xi=\lambda+\rho_{\rm n}=\mu_0+\rho_{\rm c}$
is strictly decreasing within each sign.

For two tuples strictly decreasing within each sign, their
scalar product strictly decreases when one tuple is
permuted by a nonidentity element of $W_K$.
This is the rearrangement inequality: swapping an inverted
pair changes the scalar product by
$(a_i-a_j)(b_i-b_j)>0$, and repeated swaps sort each block.
Consequently, for $w\ne1$,
\[
                       \langle\lambda,w\xi-\xi\rangle<0.
\]
Every nonzero sum of positive noncompact roots instead
has positive scalar product with $\lambda$.
In (FD14) at $\mu=\mu_0$, only the identity permutation
can therefore contribute, and it contributes
$P_{\rm n}(0)=1$. This proves occurrence and multiplicity.

If $m_\mu(\pi_\lambda)>0$, some term of (FD14) is nonzero.
For its permutation $w$, write
\[
          w(\mu+\rho_{\rm c})=\lambda+\rho_{\rm n}+\beta,
 \qquad \beta\in\sum_{\alpha>0,\ {\rm noncompact}}\mathbb N\alpha.
\]
Put $x=\mu+\rho_{\rm c}$. The weak rearrangement
inequality applied within each sign gives
\begin{align*}
 \|\mu+2\rho_{\rm c}\|^2
   &=\|x\|^2+2\langle x,\rho_{\rm c}\rangle+\|\rho_{\rm c}\|^2\\
   &\ge\|wx+\rho_{\rm c}\|^2
     =\|\lambda+\rho+\beta\|^2
     \ge\|\lambda+\rho\|^2.
\end{align*}
The final inequality is strict if $\beta\ne0$, since
$\lambda+\rho$ has positive scalar product with every
positive root. If $\beta=0$, the two compact-dominant
tuples $x,\xi$ in the same $W_K$-orbit are equal, so
$\mu=\mu_0$. This proves (FD25).

For the converse, put $\nu=\lambda-\rho$. It is dominant
integral by (FD23). For a positive noncompact root
$\alpha=e_i-e_j$,
\[
 \|\mu_0-\alpha+2\rho_{\rm c}\|^2
       -\|\mu_0+2\rho_{\rm c}\|^2
       =-2\langle\lambda+\rho,\alpha\rangle+2<0.
 \tag{FD26}
\]
Indeed $\lambda_i-\lambda_j\ge j-i\ge1$, and
$\rho_i-\rho_j=j-i\ge1$.
Whenever $\mu_0-\alpha$ is compact dominant, its occurrence
would contradict the assumed minimum. All the hypotheses
of the stated compact-type characterization apply, with
$\nu+2\rho_{\rm n}=\mu_0$. It follows that the representation
is $\pi_\lambda$.
\end{proof}

\begin{corollary}\label{fd-source}
Suppose the source's Haar and representation data satisfy
their stated identifications: its measure is compact-dual
normalized; its representation is nonzero, irreducible,
unitary, strongly continuous and square integrable; and
it contains the compact type $\mu_0$ of (FD24), with
minimum squared norm as above. Then its formal degree
$d_\lambda$ satisfies
\[
                  s_n d_\lambda
                      =\prod_{i<j}(\lambda_i-\lambda_j).
 \tag{FD27}
\]
In particular the formal-degree assumption
\eqref{as-formal-degree-assumption} holds with equality
before its factor $1/2$ is inserted.
\end{corollary}
\begin{proof}
The source's equal-density measure conditions are exactly
(FD8)--(FD10), so Proposition~\ref{fd-measure-match}
identifies its Haar measure. Its compact-type conditions
are exactly those in Proposition~\ref{fd-lowest-type}.
Apply Theorem~\ref{fd-formula}; strict decrease removes
the absolute values.

There is also a direct way to apply the source's numerical
identification of its degree. That identification quantifies
over every representation with the specified properties
and every positive formal degree satisfying Schur
orthogonality. The classical $\pi_\lambda$ exists and
has those properties by (FD14) and
Proposition~\ref{fd-lowest-type}. Substituting this one
representation gives the same value for the stipulated
degree. Thus no inference from an empty class of
representations is used.
\end{proof}

\subsection{Normalization checks and scope}

For $p=q=1$, one has $c_{1,1}=1/\pi$, and the scalar
parameter difference is $s-1$. Equations (FD17) and
\eqref{fd-exact} give
\[
               d_{\mu_0}=(s-1)/\pi,\qquad
               d_{\mu_{\rm cd}}=s-1.
\]
More generally the constant is $s_n/(\pi^{pq}s_ps_q)$,
not just $\pi^{-pq}$. A claim that (FD27) holds for an
arbitrary Haar measure is false: replacing the measure
by twice itself divides its formal degree by two.
The source explicitly fixes the measure and is not
subject to this objection.

The compact-type condition also has real content.
For $U(1,1)$ the two discrete series in a regular
infinitesimal-character packet have different smallest
torus weights. Replacing the stipulated type by an
arbitrary occurring type would not justify the
parameter identification above. The proof retains
both its occurrence and its minimum property.

This section supplies the precise formal-degree value
and its normalization from the three named classical
theorems. It does not reprove those deep theorems or
claim a line-by-line reconstruction of the source's
independent discrete-series construction. It closes
the formal-degree input to the elementary size estimates.
The global multiplicity bound, the integrated Jacquet
identities and the arithmetic theta moments remain
separate questions.


% Rational orbital estimates and the complex multiplicity comparison.
% Labels in this fragment have the prefix amb-.
\section{Rational orbital estimates and the ambient multiplicity}
\label{amb-section}

The dimension of the decorated, definite source space in
Section~\ref{cd-section} is different from the multiplicity in the
indefinite automorphic space. Here we isolate that second count,
prove its rational-element and distant-orbit estimates, and give the
exact analytic reduction needed for the final comparison. The
coefficients throughout this section are complex. In particular, the
large auxiliary prime below imposes a condition on the finite level;
it is not a coefficient characteristic. These are the conventions of
the submitted multiplicity theorem.\cite{alpo-uniform-source}

\subsection{The multiplicity and its proposed main term}

Let $E/F$ be a CM extension, with $F$ totally real of degree $f$,
$[E:F]=2$, and $E/F$ unramified at all finite places. Let $V$ be an
$n$-dimensional Hermitian space over $E$, with a self-dual
$\mathcal O_E$-lattice $\Lambda_0$. Write $G=\operatorname U(V)$,
$\Gamma=G(F)$ and
\[
 G_\infty=\prod_{v\mid\infty}G(F_v),\qquad
 G_f=G(\mathbb A_{F,f}),\qquad
 \mathcal X=\Gamma\backslash(G_\infty\times G_f).
\]
At every real place the signature is $(p,q)$, where $p+q=n$.
Choose coordinates in which the Hermitian matrix is a diagonal
matrix $J_v$ with entries $1$ and $-1$. No ordering of these signs
is assumed. The associated positive norm is the usual Euclidean
norm in these coordinates. Put
\[
                       r_E=|D_E|^{1/(2f)}.
\]
Let $K_f$ stabilize $\widehat\Lambda_0$, and let $K_f(g)$ be its
principal congruence subgroup of positive integer level $g$.

The finite test function $e:G_f\to\mathbb C$ is locally constant,
compactly supported in $K_f$, bi-invariant under $K_f(g)$, and
satisfies
\begin{equation}\label{amb-idempotent}
 e*e=e,\qquad e(y^{-1})=\overline{e(y)},\qquad
 \sum_{i,j}\overline{c_i}c_j e(y_i^{-1}y_j)\ge0.
\end{equation}
The last inequality holds for every finite collection of $y_i,c_i$.
There is a rational prime $P>4n$ which splits completely in $E$,
such that
\begin{equation}\label{amb-unipotent-level}
 e(y)\ne0\quad\Longrightarrow\quad
             (y-1)^n\widehat\Lambda_0\subseteq
                                      P\widehat\Lambda_0.
\end{equation}
Fix Haar measures $dg_v$ on $G(F_v)$ and $dg_f$ on $G_f$.

For each real place let $\sigma_v$ be the discrete-series
representation with the specified Harish--Chandra parameter
$\lambda_v$, and choose a unit vector $w_v$ in its lowest
$K_v$-type. Let $d_v>0$ be its formal degree for $dg_v$ and
$\sigma=\widehat\bigotimes_v\sigma_v$. With $R$ the right regular
representation on $L^2(\mathcal X)$, define
\begin{equation}\label{amb-D-definition}
 D=\dim_{\mathbb C}
   \bigl\{R(e)T:T\in
       \operatorname{Hom}^{\mathrm{cont}}_{G_\infty}
                         (\sigma,L^2(\mathcal X))\bigr\}.
\end{equation}
The dimension in \eqref{amb-D-definition} is used as a finite number
only after finite-dimensionality has been established. Write
\begin{equation}\label{amb-main-term}
 m_0=\Bigl(\prod_v d_v\Bigr)e(1),\qquad
 D_*=m_0\operatorname{vol}(\mathcal X).
\end{equation}
Lemma~\ref{amb-positive-type} below shows that $e(1)$ is real and
nonnegative. The source's identification of its row multiplicity
with \eqref{amb-D-definition} is the identification used to write
$D=D_{\rm all}$; it does not change $D_*$.

The claimed comparison is
\begin{equation}\label{amb-target}
                 |D-D_*|\le
                  \exp\bigl(-nf\log n/16\bigr)D_*.
\end{equation}
The source proves this for sufficiently large $n$, uniformly under
the following quantitative assumptions. Here $C_g\ge0$, $b\ge2$,
and $h$ is a lower bound for every difference
$|\lambda_{v,i}-\lambda_{v,j}|$ whose two indices have opposite
signs in $J_v$:
\begin{align}
 q^2&\le n,& \log f&\le\sqrt n,& \log r_E&\le e^n,
                                                        \notag\\
 |\lambda_{v,i}|&\le b^2,&1\le g&\le b^{C_g},&
 \log b&\le\exp(nf\log n/16),                         \label{amb-parameters}\\
 h&\ge8\cdot10^{14}(C_g+4)n^{19}f^5\log(br_E).
                                                        \notag
\end{align}
The size of the threshold for $n$ may depend on $C_g$.
The statements proved below specify exactly which additional
analytic inputs are still needed for \eqref{amb-target}.

\subsection{Unipotence of nearby rational isometries}

For an invertible complex matrix $A$ with singular values
$s_1,\ldots,s_n$, set
\[
             t(A)=\Bigl(\sum_i(\log s_i)^2\Bigr)^{1/2}.
\]
Its polar decomposition is $A=k\exp X$, where $k$ is unitary,
$X$ is Hermitian, and $\|X\|_{\rm HS}=t(A)$. In particular,
$\|A\|_{\rm op}\le e^{t(A)}$.

\begin{lemma}\label{amb-archimedean-traces}
Suppose $n\ge1$, $\epsilon\ge0$, $n\epsilon\le1/4$ and
$t(A)\le\epsilon$. Then, for $0\le j\le n$,
\begin{equation}\label{amb-power-traces}
                  |\operatorname{tr}(A^j)-n|\le3n.
\end{equation}
If $\operatorname{tr}A=n$, then
\begin{equation}\label{amb-close-to-one}
                  \|A-1\|_{\rm HS}\le3\sqrt{n\epsilon}.
\end{equation}
\end{lemma}
\begin{proof}
The operator-norm bound gives
$|\operatorname{tr}(A^j)|\le n\|A^j\|_{\rm op}
\le ne^{j\epsilon}\le ne^{1/4}<3n/2$.
The last strict bound follows, for example, by comparing the
exponential series with $\sum_{k\ge0}(1/4)^k=4/3$.
The triangle inequality proves \eqref{amb-power-traces}.
If the trace is $n$, then
\[
 \|A-1\|_{\rm HS}^2
   =\operatorname{tr}(A^*A)+n-2\Re\operatorname{tr}A
   =\operatorname{tr}(e^{2X})-n
   \le n(e^{2\epsilon}-1).
\]
Here $\epsilon\le1/4$. The exponential series gives
$e^{2\epsilon}-1\le2\epsilon/(1-2\epsilon)
\le4\epsilon$. Taking square roots proves the stronger bound
$2\sqrt{n\epsilon}$, hence \eqref{amb-close-to-one}.
\end{proof}

\begin{lemma}\label{amb-integral-traces}
Let $E$ be a number field, $\Lambda$ a full
$\mathcal O_E$-lattice in $E^n$, and $A\Lambda\subseteq\Lambda$.
Suppose that $P\mathcal O_E$ is radical and
$(A-1)^n\Lambda\subseteq P\Lambda$. Then
\[
               \operatorname{tr}(A^j)-n\in P\mathcal O_E
                                      \qquad(j\ge0).
\]
No freeness of $\Lambda$ over $\mathcal O_E$ is required.
\end{lemma}
\begin{proof}
At every nonzero prime $\mathfrak p$, the localization
$\Lambda_{\mathfrak p}$ is a free module of rank $n$ over the
discrete valuation ring $\mathcal O_{E,\mathfrak p}$.
Thus the characteristic polynomial of $A$ has coefficients in
every $\mathcal O_{E,\mathfrak p}$, hence in their intersection
$\mathcal O_E$. In particular all the traces in question are
integral. At a prime $\mathfrak p\mid P$, the reduction of $A$
on $\Lambda_{\mathfrak p}/\mathfrak p\Lambda_{\mathfrak p}$ is
unipotent. Its $j$th power has all eigenvalues equal to one over
an algebraic closure of the residue field, so its trace is $n$.
Consequently $\operatorname{tr}(A^j)-n$ belongs to every such
$\mathfrak p$. Radicality gives
$\bigcap_{\mathfrak p\mid P}\mathfrak p=P\mathcal O_E$.
\end{proof}

\begin{proposition}\label{amb-near-unipotent}
Let $\Lambda,A,P$ be as in Lemma~\ref{amb-integral-traces}, with
$P>3n$. For every complex embedding $\iota:E\hookrightarrow
\mathbb C$, suppose that $\iota(A)$ is similar to a matrix
$A_\iota$ satisfying $t(A_\iota)\le\epsilon$, where
$n\epsilon\le1/4$. Then
\begin{equation}\label{amb-unipotence}
                   \det(T-A)=(T-1)^n,
                   \qquad(A-1)^n=0.
\end{equation}
Moreover \eqref{amb-close-to-one} holds for each $A_\iota$.
\end{proposition}
\begin{proof}
For $1\le j\le n$ put $a_j=\operatorname{tr}(A^j)-n$.
Lemma~\ref{amb-integral-traces} writes $a_j=P b_j$ with
$b_j\in\mathcal O_E$. If $a_j\ne0$, then
\[
 |N_{E/\mathbb Q}(a_j)|
       =P^{[E:\mathbb Q]}|N_{E/\mathbb Q}(b_j)|
       \ge P^{[E:\mathbb Q]}.
\]
But similarity invariance of trace and
\eqref{amb-power-traces} bound the same norm by
$(3n)^{[E:\mathbb Q]}$, a contradiction. Thus
$\operatorname{tr}(A^j)=n$ for $1\le j\le n$.

For completeness these power traces determine the characteristic
polynomial in characteristic zero. If $\alpha_i$ are its roots
with multiplicity, then the polynomial
$Q(z)=\prod_i(1-\alpha_i z)$ satisfies the formal identity
\[
             -zQ'(z)=Q(z)\sum_{j\ge1}
                              \operatorname{tr}(A^j)z^j.
\]
Writing $Q(z)=\sum_{j=0}^n c_jz^j$ with $c_0=1$, comparison
of degree $j$ determines $j c_j$ from $c_0,\ldots,c_{j-1}$ and
the first $j$ traces. Those traces agree with the identity matrix,
so $Q(z)=(1-z)^n$. Cayley--Hamilton now gives
$(A-1)^n=0$. Finally the first trace is $n$, and
Lemma~\ref{amb-archimedean-traces} gives the asserted bound.
\end{proof}

For an adelic point $x=(x_\infty,x_f)$, define its rational
lattice and positive metrics by
\begin{equation}\label{amb-point-lattice}
 \Lambda_x=V\cap x_f\widehat\Lambda_0,
 \qquad\|u\|_{x,v}=\|x_v^{-1}u_v\|_2.
\end{equation}
Here $u_v$ is expressed in the chosen $J_v$-coordinates.
The lattice is full and self-dual. Indeed a common integer clears
the denominators of $x_f$ and $x_f^{-1}$, which bounds
$\Lambda_x$ between two homothetic copies of $\Lambda_0$.
The Chinese remainder theorem identifies its completions with
$x_f\widehat\Lambda_0$; duality can then be checked at each
finite prime. This also proves
\[
 (x^{-1}\gamma x)_f\in K_f\Longrightarrow
                        \gamma\Lambda_x=\Lambda_x,
\]
and \eqref{amb-unipotent-level} gives
$(\gamma-1)^n\Lambda_x\subseteq P\Lambda_x$ whenever the
finite test function is nonzero.

Complete splitting of $P$ makes $P\mathcal O_E$ radical. The
$f$ selected complex embeddings and their conjugates are all the
embeddings of $E$. Therefore Proposition~\ref{amb-near-unipotent}
applies to every rational summand with
$e((x^{-1}\gamma x)_f)\ne0$ and
$t((x^{-1}\gamma x)_v)\le\epsilon$ for all $v$.
In particular it applies at the exact radius
\begin{equation}\label{amb-small-radius}
                  \epsilon=10^{-6}n^{-8}f^{-2}.
\end{equation}
The source's further radius bounds follow directly: if
$\delta=3\sqrt{n\epsilon}$ and $q^2\le n$, then
\[
             \delta<\tfrac14,
             \qquad\delta qnf
               =3\cdot10^{-3}\frac q{n^{5/2}}
               \le3\cdot10^{-3}<\tfrac1{100}.
\]

\subsection{A separation bound from normalized covolumes}

We prove the separation estimate used in the orbit count without
requiring the construction of a Harder--Narasimhan filtration.
Put $c_E=2^{-f}|D_E|^{1/2}$. For a nonzero
$\mathcal O_E$-lattice $M$ of rank $r$, with its positive metrics
at the $f$ complex places, define
\begin{equation}\label{amb-size}
 H(M)=c_E^{-r}\operatorname{covol}_{\mathbb R}(M),\qquad
 s(M)=H(M)^{1/(2fr)},\qquad
 \lambda(\Lambda)=\inf_{0\ne M\subseteq\Lambda}s(M).
\end{equation}
Covolumes are measured in the real span of the lattice, with the
Euclidean metric which is the sum of the selected complex metrics.
For a torsion-free quotient, use the orthogonal quotient metrics.
These are exactly the normalized covolumes and exponents in the
submitted least-size function.

\begin{lemma}\label{amb-size-duality}
For every full $\mathcal O_E$-lattice $\Lambda$ with positive
metrics, $\lambda(\Lambda)>0$. Let $\Lambda^*$ be its
$\mathcal O_E$-linear dual, with the dual metrics. Then
\begin{equation}\label{amb-height-dual}
 H(\Lambda^*)=H(\Lambda)^{-1},\qquad
 \sup_{\substack{S\subsetneq\Lambda\\\Lambda/S\ \text{torsion-free}}}
              s(\Lambda/S)=\lambda(\Lambda^*)^{-1}.
\end{equation}
If $\Lambda$ is self-dual for a Hermitian form whose matrix is
$J_v$ in an orthonormal basis of each positive metric, then
$\lambda(\Lambda^*)=\lambda(\Lambda)$.
\end{lemma}
\begin{proof}
For $1\le r\le n$, a $\mathbb Z$-basis of any rank-$r$
$\mathcal O_E$-sublattice gives a nonzero vector in the fixed
lattice $\bigwedge^{2fr}_{\mathbb Z}\Lambda$. Its Euclidean
length is the covolume of the sublattice. A lattice has a positive
minimum length for nonzero vectors: an invertible real linear
map from $\mathbb Z^N$ bounds every such length below by its
smallest singular value. Taking the minimum over the finitely
many ranks proves positivity in \eqref{amb-size}.

If $M_0=\bigoplus_{i=1}^r\mathcal O_E b_i$ is free and
$G_v=(\langle b_i,b_j\rangle_v)$, the discriminant determinant
for the Minkowski embedding gives
\begin{equation}\label{amb-free-height}
                       H(M_0)=\prod_v\det G_v.
\end{equation}
Explicitly the covolume of $\mathcal O_E$ in its $f$ complex
coordinates is $c_E$; the real determinant of the complex change
of basis at $v$ is $|\det B_v|^2=\det G_v$. Multiplying these
determinants proves \eqref{amb-free-height}. The dual Gram
matrices are $G_v^{-1}$, so the dual-height assertion follows
for $M_0$.

For arbitrary $M$, choose such a free sublattice $M_0\subset M$
of finite index $a$. Then $H(M)=H(M_0)/a$ and
$[M_0^*:M^*]=a$. The latter index equality can be checked at
each prime: diagonalizing the inclusion over the discrete
valuation ring changes its exponents to their negatives in the
dual inclusion, with the same lengths of the two finite
quotients. Their global orders are the products of the residue
norms to these lengths. Thus
$H(M^*)=aH(M_0^*)=H(M)^{-1}$, without assuming $M$ is free.

If $S$ is saturated in $\Lambda$, then $Q=\Lambda/S$ is
projective over the Dedekind domain $\mathcal O_E$.
Dualizing the split module sequence identifies $Q^*$ with the
saturated submodule of $\Lambda^*$ annihilating $S$, with its
induced metric. Conversely every saturated submodule of
$\Lambda^*$ arises in this way, by dualizing its projective
quotient. Saturating a nonzero sublattice decreases its covolume
and preserves its rank; consequently the infimum in
\eqref{amb-size} is unchanged if it is restricted to saturated
sublattices. Since $s(Q^*)=s(Q)^{-1}$, this proves the supremum
identity.

Finally the Hermitian form identifies $\Lambda$ conjugate
linearly with $\Lambda^*$. At each complex place the norm of the
functional $u\mapsto\langle y,J_vu\rangle$ is $\|y\|$, because
$J_v$ is unitary. This identification is an isometry, takes
$\mathcal O_E$-sublattices to such sublattices and preserves their
real covolumes. Hence it preserves the infimum defining
$\lambda$.
\end{proof}

\begin{lemma}\label{amb-map-separation}
Under the self-duality hypothesis of Lemma~\ref{amb-size-duality},
every nonzero $E$-linear map $T$ with
$T\Lambda\subseteq\Lambda$ satisfies
\begin{equation}\label{amb-geometric-separation}
            \Bigl(\prod_v\|T_v\|_{\rm op}\Bigr)^{1/f}
                                     \ge\lambda(\Lambda)^2.
\end{equation}
\end{lemma}
\begin{proof}
Let $S=\Lambda\cap\ker T$, $Q=\Lambda/S$, and $M=T\Lambda$,
of common positive rank $r$. The first is saturated, and the
induced map $Q\to M$ is an isomorphism of
$\mathcal O_E$-modules. At each complex place its singular
values are at most $\|T_v\|_{\rm op}$: identify the quotient
with the orthogonal complement of $\ker T_v$ and restrict $T_v$.
The change-of-volume formula, applied first to a finite-index
free sublattice and then cancelling that common index, gives
\[
                 H(M)\le H(Q)\prod_v\|T_v\|_{\rm op}^{2r}.
\]
Take $2fr$th roots. By Lemma~\ref{amb-size-duality},
$s(M)\ge\lambda(\Lambda)$ and
$s(Q)\le\lambda(\Lambda)^{-1}$. These inequalities give
\eqref{amb-geometric-separation}.
\end{proof}

For the point lattice \eqref{amb-point-lattice}, put
\begin{equation}\label{amb-tau}
          \lambda_x=\lambda(\Lambda_x),\qquad
          \tau_x=\min(1/2,\lambda_x)>0.
\end{equation}
The positive metrics in \eqref{amb-point-lattice} are compatible
with the original Hermitian form, since $x_v$ is an isometry of
$J_v$. Thus Lemma~\ref{amb-map-separation} applies. If distinct
rational elements $\gamma,\eta$ both have
$(x^{-1}\gamma x)_f,(x^{-1}\eta x)_f\in K_f$, their difference
preserves $\Lambda_x$, and
\begin{equation}\label{amb-two-matrices}
 \Bigl(\prod_v\|(x^{-1}\gamma x)_v
                         -(x^{-1}\eta x)_v\|_{\rm op}\Bigr)^{1/f}
                                                    \ge\tau_x^2.
\end{equation}

\subsection{Counting rational matrices and summing the tail}

\begin{lemma}\label{amb-packing}
Let $W_1,\ldots,W_f$ be real normed spaces of dimension $N$,
where $f,N\ge1$. Suppose a set $S\subset\prod_vW_v$ satisfies
\[
       \Bigl(\prod_v\|s_v-t_v\|\Bigr)^{1/f}\ge a>0
                       \quad(s,t\in S,\ s\ne t).
\]
For positive $R_v$, the subset with $\|s_v\|\le R_v$ for all
$v$ is finite, of cardinality at most
\begin{equation}\label{amb-packing-bound}
              \left(1+\frac{2(\prod_vR_v)^{1/f}}a\right)^{Nf}.
\end{equation}
\end{lemma}
\begin{proof}
Rescale the $v$th coordinate by $R/R_v$, where
$R=(\prod_vR_v)^{1/f}$. The product of the scale factors is one,
so the separation hypothesis is unchanged. In the maximum norm
on $\prod_vW_v$, the new set is $a$-separated and lies in the
closed ball of radius $R$. The open balls of radius $a/2$ about
any finite selection of its points are disjoint and lie in the
ball of radius $R+a/2$. Lebesgue volume in dimension $Nf$ gives
\eqref{amb-packing-bound} for every finite selection. An infinite
set has finite subsets of arbitrarily large cardinality, so the
same bound proves finiteness as well.
\end{proof}

\begin{corollary}\label{amb-rational-ball}
For $r_v\ge0$, the rational elements with finite component in
$xK_fx^{-1}$ and $t((x^{-1}\gamma x)_v)\le r_v$ form a finite
set of cardinality at most
\begin{equation}\label{amb-ball-bound}
 A_x\exp\left(2n^2\sum_vr_v\right),\qquad
                   A_x=\left(\frac{3\sqrt n}{\tau_x^2}\right)^{2n^2f}.
\end{equation}
\end{corollary}
\begin{proof}
Use \eqref{amb-two-matrices} and Lemma~\ref{amb-packing} in the
spaces of complex $n\times n$ matrices with operator norm,
each of real dimension $2n^2$, with $a=\tau_x^2$ and
$R_v=e^{r_v}$. The map from rational elements to their matrices
is injective, since an embedding of $E$ is injective.
Writing $R=e^{\sum_vr_v/f}\ge1$, we have
$1+2R/\tau_x^2\le3R/\tau_x^2\le3\sqrt nR/\tau_x^2$.
Raise this inequality to $2n^2f$. The proof actually permits
omitting $\sqrt n$; we keep it to match the source bound.
\end{proof}

\begin{lemma}[A dyadic sum]\label{amb-shells}
Let $S$ be a set, $d_v:S\to[0,\infty)$, and
$w:S\to[0,\infty)$, with $f\ge1$ indices $v$. Suppose
\begin{align*}
 \#\{s:d_v(s)\le r_v\text{ for all }v\}
                 &\le A\exp(\kappa\sum_vr_v),\qquad r_v\ge0,\\
 w(s)&\le\prod_v\exp\left[-\frac h8
                              \min(d_v(s)^2,d_v(s)/c)\right],
\end{align*}
where the counted sets are finite and $A,\kappa\ge0$.
Assume
\begin{equation}\label{amb-shell-hypotheses}
 0<\epsilon\le1/c,\quad c>0,\quad
 \kappa\epsilon\le\log2,\quad
 32\kappa\le h\epsilon,\quad
                     T:=h\epsilon^2/16\ge1.
\end{equation}
Then the sum over $s$ for which some $d_v(s)>\epsilon$ is
absolutely convergent and at most
\begin{equation}\label{amb-shell-bound}
                          2Af3^f e^{-T}.
\end{equation}
\end{lemma}
\begin{proof}
Partition each coordinate into $[0,\epsilon]$ and
$(2^{j-1}\epsilon,2^j\epsilon]$, $j\ge1$.
For $d\ge\epsilon$ the inequality $\epsilon\le1/c$ gives
$\min(d^2,d/c)\ge\epsilon d$. Thus the weight in the $j$th
interval, $j\ge1$, is at most $e^{-T2^j}$; in interval zero it
is at most one. The number of points in a product of intervals
is bounded using the encompassing box. Set
\[
 W_0=e^{\kappa\epsilon},\qquad
 W_j=\exp\bigl((\kappa\epsilon-T)2^j\bigr)\quad(j\ge1).
\]
Their products, multiplied by $A$, bound the contributions of
the corresponding boxes. We have $W_0\le2$ and
$\kappa\epsilon\le T/2$. Hence
\[
       \sum_{j\ge1}W_j
           \le\sum_{i\ge0}e^{-T2^i}\le2e^{-T}\le1.
\]
For the middle inequality use $2^i\ge1+i$ and $T\ge1$, so
$e^{-T2^i}\le e^{-T}2^{-i}$. Write the last sum as $\beta$
and $W_0$ as $\alpha$. Summing over all multi-indices except
zero gives at most
\[
 A\bigl((\alpha+\beta)^f-\alpha^f\bigr)
       \le Af\beta(\alpha+\beta)^{f-1}
       \le2Af3^f e^{-T}.
\]
Every box is finite, so the partition also proves that the
index set is countable. The finite bounds and the convergent
nonnegative majorant justify the sum and its convergence.
\end{proof}

\begin{lemma}\label{amb-positive-type}
For the test function in \eqref{amb-idempotent},
\begin{equation}\label{amb-test-bound}
                         e(1)\ge0,\qquad |e(y)|\le e(1).
\end{equation}
If $e(1)=0$, then $e$ is identically zero.
\end{lemma}
\begin{proof}
The one-point positive-type condition makes $e(1)$ real and
nonnegative. The two-point condition says that
$\begin{pmatrix}e(1)&e(y)\\\overline{e(y)}&e(1)\end{pmatrix}$
is positive semidefinite. Its determinant is nonnegative, which
is \eqref{amb-test-bound}. This also proves the last assertion.
\end{proof}

With inner products conjugate linear in the first argument, set
\begin{equation}\label{amb-kernel-function}
 \Phi(g)=\Bigl(\prod_vd_v\Bigr)
          \overline{\prod_v\langle w_v,\sigma_v(g_v)w_v\rangle}
                 e(g_f),\qquad
 K(x,y)=\sum_{\gamma\in\Gamma}\Phi(x^{-1}\gamma y).
\end{equation}
In particular $\Phi(1)=m_0$. The representation-theoretic
estimate used for the next theorem is the precise assertion
\begin{equation}\label{amb-decay-input}
 |\langle w_v,\sigma_v(a)w_v\rangle|
       \le\exp\left[-\frac h8
                 \min\bigl(t(a)^2,t(a)/\sqrt{2q}\bigr)\right].
\end{equation}
It is needed for all $a\in G(F_v)$ and the chosen lowest-type
vectors. We now derive it, with a stronger constant, from a
classical integral formula for the lowest compact type.

\subsection{Deriving decay from the lowest compact-type formula}

We first give the elementary estimate for the scalar integrand.
We then identify the exact integral formula in a classical
representation-theoretic theorem, including its character,
measure, and projection conventions.

For a finite real random variable $Y$ supported in $[-M,M]$,
write
\[
 \varphi(u)=\log\mathbb E(e^{uY}),\qquad
 \operatorname{Var}(Y)=\mathbb E(Y^2)-\mathbb E(Y)^2.
\]
The expectation is a finite sum with nonnegative weights of sum one.

\begin{lemma}\label{amb-variance}
Suppose $\mathbb E(Y)=0$ and $0<\theta\le1$ satisfies
$2M\theta\le1$. Then
\[
                     \log\mathbb E(e^Y)
                  \ge\frac{\theta}{\mathrm e}\operatorname{Var}(Y),
\]
where $\mathrm e=\exp(1)$.
\end{lemma}
\begin{proof}
Differentiation of a finite sum shows that $\varphi''(u)$ is
the variance for the probability measure whose weights are
multiplied by $e^{uY}/\mathbb E(e^{uY})$.
For $u\ge0$ each multiplier is at least $e^{-2Mu}$.
Variance is $\min_{c\in\mathbb R}\mathbb E((Y-c)^2)$.
Taking this minimum after the pointwise comparison gives
\[
                \varphi''(u)\ge e^{-2Mu}\operatorname{Var}(Y).
\]
In particular $\varphi$ is convex. Since
$\varphi(0)=\varphi'(0)=0$ and $2M\le1/\theta$,
\[
 \varphi(\theta)
   =\int_0^\theta(\theta-u)\varphi''(u)\,du
   \ge\operatorname{Var}(Y)
             \int_0^\theta(\theta-u)e^{-u/\theta}\,du
   =\frac{\theta^2}{\mathrm e}\operatorname{Var}(Y).
\]
Convexity and $\varphi(0)=0$ give
$\varphi(1)\ge\varphi(\theta)/\theta$, proving the claim.
\end{proof}

Let $J$ be a diagonal signature matrix with $q\ge1$ negative
entries, and let
\[
             X=X^*,\qquad JX+XJ=0,\qquad t=\|X\|_{\rm HS}.
\]
The signs may occur in any order. For $1\le j<n$ define
\[
 \Delta_j(X)=\det\bigl((e^X)_{a b}\bigr)_{1\le a,b\le j}>0.
\]
For a strictly decreasing real tuple $\lambda$, put
\begin{equation}\label{amb-phase}
 \mathcal P_\lambda(X)=
       \sum_{j=1}^{n-1}(\lambda_j-\lambda_{j+1}+1)
                                      \log\Delta_j(X).
\end{equation}
Suppose that $\lambda_a-\lambda_b\ge h\ge0$ whenever
$a<b$ and the corresponding signs differ.

\begin{proposition}\label{amb-principal-minor-bound}
With these conventions,
\begin{equation}\label{amb-phase-bound}
 \mathcal P_\lambda(X)\ge
       \frac{h}{2\mathrm e}
                \min\left(t^2,\frac{t}{\sqrt{2q}}\right).
\end{equation}
\end{proposition}
\begin{proof}
The claim is immediate if $X=0$, so suppose $t>0$.
The relation $JXJ=-X$ makes the spectrum symmetric about zero.
After grouping the signs, $X$ has blocks
$\left(\begin{smallmatrix}0&B\\B^*&0\end{smallmatrix}\right)$.
It has at most $q$ positive eigenvalues, and the sum of their
squares is $t^2/2$. If $M$ is their sum, Cauchy--Schwarz gives
\[
                          M\le\sqrt{q/2}\,t.
\]

On $\bigwedge^j\mathbb C^n$, let $D_j(X)$ act by applying
$X$ to each factor and adding. Its eigenvalues are sums of
$j$ eigenvalues of $X$, hence lie in $[-M,M]$.
For the unit vector $v_j=e_1\wedge\cdots\wedge e_j$,
the identity $\bigwedge^j(e^X)=e^{D_j(X)}$ gives
\[
 \Delta_j(X)=\langle v_j,e^{D_j(X)}v_j\rangle.
\]
Expand $v_j$ in an orthonormal eigenbasis of $D_j(X)$.
The squared absolute values of its coefficients define a
finite probability measure, and its exponential moment is
the last expression. Its mean is
$\sum_{a\le j}X_{aa}=0$. Its variance is
\[
 Q_j(X)=\|D_j(X)v_j\|^2
                  =\sum_{a\le j<b}|X_{ba}|^2.
\]
Indeed the nonzero replacements of an occupied basis vector
by an unoccupied one are distinct orthogonal exterior basis
vectors; the diagonal contribution is zero.

Choose
$\theta=\min(1,1/(\sqrt{2q}\,t))$.
Then $2M\theta\le1$, so Lemma~\ref{amb-variance} gives
$\log\Delta_j(X)\ge\theta Q_j(X)/\mathrm e$.
The coefficients in \eqref{amb-phase} are positive. Summing
and interchanging two finite sums,
\[
 \begin{split}
 \sum_{j=1}^{n-1}(\lambda_j-\lambda_{j+1}+1)Q_j(X)
  &=\sum_{a<b}(\lambda_a-\lambda_b+b-a)|X_{ba}|^2\\
  &\ge h\sum_{a<b}|X_{ba}|^2
   =\frac h2\,t^2.
 \end{split}
\]
Here entries with equal signs vanish. Multiplication by
$\theta/\mathrm e$ proves \eqref{amb-phase-bound}.
\end{proof}

To state the representation-theoretic identity,
let $W$ be the lowest $K$-type in $\sigma_\lambda$,
$j:W\hookrightarrow H$ its isometric inclusion, $\tau$ its
compact representation, and $w_0$ a unit highest-weight vector.
Let $P_0$ be orthogonal projection onto $\mathbb Cw_0$ and
$d_W=\dim W$.

\begin{theorem}[The lowest compact-type integral]\label{amb-compression}
Suppose
$\lambda_1>\cdots>\lambda_n$ and
$\lambda_i-(n-1)/2\in\mathbb Z$.
Then
\begin{equation}\label{amb-compression-input}
 j^*\sigma_\lambda(e^X)j
   =d_W\int_K e^{-\mathcal P_\lambda(k^{-1}Xk)}
                 \tau(k)P_0\tau(k)^{-1}\,dk ,
 \qquad X=X^*,\ JX+XJ=0,
\end{equation}
with probability Haar measure on $K$.
\end{theorem}
\begin{proof}
We specify the two classical representation-theoretic inputs
and then prove the change of coordinates in the formula.
Let $\mathfrak b$ be the upper triangular Borel subalgebra of
$\mathfrak{gl}_n(\mathbb C)$, with strictly upper triangular
radical $\mathfrak n$, and let $T$ be the diagonal torus
of $G$. The Cartan involution is conjugation by $J$, so
$\mathfrak b$ is stable under it and its real Levi subgroup
is the compact torus $T$. Put
\[
 \rho_i=(n+1)/2-i,\qquad \nu=\lambda-\rho,\qquad
 \mu_0=\nu+2\rho_{\rm n}
                  =\lambda+\rho_{\rm n}-\rho_{\rm c}.
\]
The differences of consecutive coordinates of $\lambda$
are positive integers. Thus $\nu$ is an integral dominant
weight and defines a unitary character of $T$.
The classical realization of discrete series by
cohomological induction identifies the representation
induced from this character, denoted $A_{\mathfrak b}(\nu)$,
with $\sigma_\lambda$. Its lowest compact type has highest
weight $\mu_0$ and multiplicity one.
For the hypotheses in this identification, the admissibility
conditions of \cite[(5.1), Theorem 5.3]{vogan-zuckerman}
are exactly that $\nu$ be a unitary torus character and
$\langle \nu,e_a-e_b\rangle\ge0$ for $a<b$.
The identifying criterion
\cite[Proposition 6.1]{vogan-zuckerman}
applies to $\sigma_\lambda$: its lowest type is $\mu_0$,
and its Casimir eigenvalue is
$\|\lambda\|^2-\|\rho\|^2$. The additional condition on a
noncompact part of the Cartan subalgebra is empty because
$T$ is compact. These semisimple statements apply through
$SU(p,q)\times U(1)\longrightarrow U(p,q)$; the extra
central factor acts by the prescribed character and all
root shifts have coordinate sum zero.

We use Sun's integral formula for the lowest compact type
of $A_{\mathfrak b}(\nu)$
\cite[Theorem 6.4 and \S6.4, p.~43]{sun-eisenstein}.
In its precise specialization here, for $x\in\exp(\mathfrak p_0)$,
\begin{equation}\label{amb-sun}
 j^*\sigma_\lambda(x^2)j
 =d_W\int_K(\delta_{\mathfrak n}\nu)(\widetilde H(lx))
                       \tau(l^{-1})P_0\tau(l)\,dl.
\end{equation}
Here $\mathfrak p_0=\{X=X^*:JXJ=-X\}$ and
$\delta_{\mathfrak n}$ is the determinant of the diagonal
adjoint action on $\mathfrak n$. It has weight $2\rho$.
To define $\widetilde H$, factor $g=n a u$ in the complex
group, where $n$ is upper triangular unipotent, $a$ is
positive real diagonal, and $u\in U(n)$. Then
$\widetilde H(g)=a^2$.
Indeed Sun's definition is $y\overline{\theta(y^{-1})}$,
where the bar is conjugation with respect to the real
group $U(p,q)$; its composition with the Cartan involution
is the compact conjugation $g\mapsto(g^*)^{-1}$.
It therefore sends $a^{-1}$ to $a$.
The Haar measure in \eqref{amb-sun} is a probability
measure by the conventions of \cite[\S1.2]{sun-eisenstein}.
Sun's matrix coefficient is the inclusion followed by
the representation and projection onto the compact
isotypic component. Multiplicity one and unitarity make
it exactly $j^*\sigma_\lambda(g)j$, with value $1_W$ at
the identity. Thus no scalar factor has been suppressed.

Here is an elementary verification of the matrix conversion.
Set $x=e^{X/2}$ and $B=l e^X l^{-1}$, a positive definite
Hermitian matrix of determinant one. There is a unique
factorization
\[
                 B=n a^2 n^*,\qquad
       n\text{ upper triangular unipotent},\quad
       a=\operatorname{diag}(a_1,\ldots,a_n)>0.
\]
It follows inductively from the Schur complement of the
last diagonal entry: if
$B=\left(\begin{smallmatrix}B_0&b\\b^*&c\end{smallmatrix}\right)$,
then $c>0$ and $B_0-bb^*/c>0$; factor the latter matrix,
take $a_n^2=c$, and take the last column of $n$ above
the diagonal to be $b/c$. This also proves uniqueness.
The matrix $a^{-1}n^{-1}lx$ is unitary, so this is the
factorization used to define $\widetilde H(lx)=a^2$.

Write $z_i=a_i^2$ and
$T_i=\det B_{\{i+1,\ldots,n\},\{i+1,\ldots,n\}}$.
Triangularity gives
$T_i=\prod_{j>i}z_j$, and $\prod_j z_j=\det B=1$.
The character $\delta_{\mathfrak n}\nu$ has weight
$\kappa=\nu+2\rho=\lambda+\rho$. Summation by parts gives
\[
 \log(\delta_{\mathfrak n}\nu)(a^2)
    =\sum_j\kappa_j\log z_j
    =-\sum_{i=1}^{n-1}
             (\lambda_i-\lambda_{i+1}+1)\log T_i.
\]
The complementary-principal-minor identity, with $\det B=1$,
is
\[
       T_i=\det (B^{-1})_{\{1,\ldots,i\},\{1,\ldots,i\}}.
\]
For clarity, partition $B$ after its first $i$ rows
and columns. If its lower right block is $D$, its
Schur complement is $S$, and the upper left block of
$B^{-1}$ is $S^{-1}$. Consequently
$\det(B^{-1})_{\le i,\le i}=(\det S)^{-1}
 =\det D/\det B$, which proves the identity.
Since $B^{-1}=l e^{-X}l^{-1}$, we have established
\[
  (\delta_{\mathfrak n}\nu)(\widetilde H(lx))
                  =e^{-\mathcal P_\lambda(-lXl^{-1})}.
\]
Finally $J$ lies in the center of $K$ and $JXJ=-X$.
Replacing $l$ by $lJ$ in \eqref{amb-sun} removes the minus
sign in the phase. The projection factor is unchanged,
because $\tau(J)$ is scalar by irreducibility of $\tau$.
Replace $l$ by $k^{-1}$ to obtain
\eqref{amb-compression-input}.
\end{proof}

Thus the character shifts and dimension factor in the
source's integral formula agree with a classical theorem.
The proof uses the cited results on cohomological induction;
it does not reprove those representation-theoretic theorems.

\begin{corollary}\label{amb-decay-from-compression}
Under the hypotheses of Theorem~\ref{amb-compression},
suppose $h\ge0$ is a lower bound for all positive
opposite-sign root differences. Then for any unit
vectors $u,v$ in the lowest compact type and every $a\in G$,
\[
 |\langle u,\sigma_\lambda(a)v\rangle|
 \le\exp\left[-\frac{h}{2\mathrm e}
           \min\left(t(a)^2,\frac{t(a)}{\sqrt{2q}}\right)\right].
\]
Consequently \eqref{amb-decay-input} holds.
\end{corollary}
\begin{proof}
Compact Schur orthogonality gives
$d_W\int_K\tau(k)P_0\tau(k)^{-1}\,dk=1_W$:
the integral commutes with $\tau(K)$, and its trace is one.
The norm and the off-diagonal signature conditions on $X$
are unchanged by conjugation by $K$. Proposition
\ref{amb-principal-minor-bound} bounds every positive scalar
in \eqref{amb-compression-input} by the asserted exponential.
All the conjugates of $P_0$ are positive semidefinite.
The integral therefore lies between zero and that scalar
times $1_W$ in the order on self-adjoint operators; its norm
has the same bound.

For completeness, write $a=k e^X$ using polar decomposition.
The unitary identity gives
$J(a^*a)J=(a^*a)^{-1}$, so
$X=\frac12\log(a^*a)$ satisfies $JXJ=-X$.
The unitary factor $k$ lies in $K$, and
$\|X\|_{\rm HS}=t(a)$ by the definition using singular values.
Since $j(W)$ is $K$-stable, the compression at $a$ is
$\tau(k)$ times the compression at $e^X$ and has the same
norm. This proves the coefficient estimate.
Finally $\mathrm e<4$, for example by bounding the exponential
series by $1+1+\sum_{m\ge2}2^{-(m-1)}=3$, with strict
inequality. Thus $1/(2\mathrm e)>1/8$.
\end{proof}

The restriction to the lowest compact type is essential.
For unrestricted unit vectors in a unitary representation,
fix $a$ and a unit vector $v$, and take $u=\sigma_\lambda(a)v$.
Then $|\langle u,\sigma_\lambda(a)v\rangle|=1$.
If $h>0$ and $t(a)>0$, the displayed exponential is less
than one. Thus the corresponding assertion for all
unit vectors would be false; neither the source nor the
corollary makes that assertion.
Interchanging $J$ and $-J$ changes neither the group,
the compact subgroup, nor the phase. Hence $q$ can be
replaced by $\min(p,q)$ in the bound. In the range
$n\ge3$, $q^2\le n$ used below, one already has $q\le p$.

\subsection{The distant rational terms}

\begin{theorem}[The distant rational terms]\label{amb-far}
Assume $n\ge3$, $1\le q$, $q^2\le n$, $f\ge1$,
$r_E\ge1$, $b\ge2$, $C_g\ge0$, $1\le g\le b^{C_g}$,
$\log f\le\sqrt n$, and the lower bound for $h$ in
\eqref{amb-parameters}. Assume \eqref{amb-decay-input} and put
\begin{equation}\label{amb-truncation}
                X=nfr_Egb,\qquad \tau_0=X^{-11q}.
\end{equation}
For every $x$ with $\lambda_x\ge\tau_0$, the rational sum
\begin{equation}\label{amb-far-sum}
 \sum_{\substack{\gamma\in\Gamma\\
       t((x^{-1}\gamma x)_v)>\epsilon\text{ for some }v}}
                     |\Phi(x^{-1}\gamma x)|
                              \le m_0e^{-n^2f}
\end{equation}
is absolutely convergent, where $\epsilon$ is
\eqref{amb-small-radius}. Consequently its integral over any
measurable subset of $\{x:\lambda_x\ge\tau_0\}$ in a
finite-volume fundamental domain is at most $e^{-n^2f}D_*$.
Terms with $\Phi(x^{-1}\gamma x)=0$ may freely be included.
\end{theorem}
\begin{proof}
If $m_0=0$, Lemma~\ref{amb-positive-type} gives $e=0$, and all
assertions hold. Otherwise divide the summands by $m_0$.
Their nonzero finite components lie in $K_f$.
Lemma~\ref{amb-positive-type} and \eqref{amb-decay-input} give
the weight hypothesis of Lemma~\ref{amb-shells}, while
Corollary~\ref{amb-rational-ball} gives its counting hypothesis
with $A=A_x$, $\kappa=2n^2$ and $c=\sqrt{2q}$.

Here are all the numerical estimates needed to apply that lemma
and obtain its asserted strength. Put $L=\log(br_E)$ and
$s=\sqrt n$. We use $L\ge\log2\ge2/3$; the last inequality
follows by integrating $2/(1-u^2)\ge2$ from $0$ to $1/3$.
The lower bound for $h$ gives exactly
\begin{equation}\label{amb-T-lower}
 T=\frac{h\epsilon^2}{16}
                      \ge50(C_g+4)n^3fL.
\end{equation}
It also gives $h\epsilon\ge64n^2$ and $T\ge1$.
For example
$h\epsilon\ge8\cdot10^8(C_g+4)n^{11}f^3L$, which is
larger than $64n^2$ under the stated assumptions.
Moreover $\epsilon\le1/(2q)\le1/\sqrt{2q}$ and
$2n^2\epsilon\le2\cdot10^{-6}\le\log2$.
These prove \eqref{amb-shell-hypotheses}.

Since $X\ge2n\ge6$, we have $\tau_0<1/2$; therefore
$\lambda_x\ge\tau_0$ implies $\tau_x\ge\tau_0$.
The elementary bound $\log n\le\sqrt n$ follows from
$\log(s/2)\le s/2-1$ and $\log2\le1$. Thus
\[
 \log X\le2s+(C_g+1)L,\qquad q\le s,
 \qquad \log(2f3^f)\le4f.
\]
Using $\log3\le2$, $s\le n$ and $n\ge3$ now gives
\begin{align*}
 \log A_x+\log(2f3^f)+n^2f
 &\le2n^2f\bigl(\log3+\tfrac12\log n+22q\log X\bigr)
                                                        +4f+n^2f\\
 &\le n^2f\bigl(5+s+88n+44s(C_g+1)L+4/n^2\bigr)\\
 &\le n^3f\bigl(92+44(C_g+1)L\bigr)\\
 &\le n^3f(44C_g+182)L\\
 &\le50(C_g+4)n^3fL\le T.
\end{align*}
In the third line use $5+4/n^2\le3n$; in the fourth use
$92\le138L$. Lemma~\ref{amb-shells} therefore bounds the
normalized sum by $A_x2f3^fe^{-T}\le e^{-n^2f}$.
Multiplication by $m_0$ proves \eqref{amb-far-sum}.
The sum is measurable, since it is a countable sum of measurable
nonnegative terms. Integrating the pointwise bound and using
the volume of the full domain proves the final assertion.
\end{proof}

\subsection{The trace identity and the remaining comparison}

The following statement records the exact analytic conclusion
needed from the automorphic-kernel construction. It is not a
consequence of the lattice packing argument: the convolution
operator $R(\Phi)$ must be a finite-rank orthogonal projection
of rank \eqref{amb-D-definition}, and its continuous kernel must
have diagonal
\begin{equation}\label{amb-projection-diagonal}
                  K(x,x)=\sum_{j=1}^D|F_j(x)|^2
\end{equation}
for an orthonormal basis $F_1,\ldots,F_D$ of its image, with
the series in \eqref{amb-kernel-function} representing this
kernel. With these assertions supplied, finite summation and
orthonormality immediately give
\begin{equation}\label{amb-trace}
                 K(x,x)\ge0,\qquad
                 \int_{\mathcal X}K(x,x)\,dx=D.
\end{equation}
In particular, the diagonal is integrable. A finite-rank kernel
on $L^2$ must be given by such representatives, or an equivalent
proved diagonal identity; equality of measurable kernels only
almost everywhere on $\mathcal X\times\mathcal X$ would not
by itself determine the diagonal.

\begin{proposition}\label{amb-final-reduction}
Assume $n\ge e^{15}$ and $f\ge1$. Suppose the actual kernel
satisfies \eqref{amb-projection-diagonal} and its diagonal
series is absolutely convergent. Let
\[
 \mathcal M=\{x:\lambda_x\ge\tau_0\},\qquad
 \mathcal B=\mathcal X\setminus\mathcal M,
 \qquad a=nf\log n,\qquad u=n^2f.
\]
Assume the three further estimates
\begin{align}
 \int_{\mathcal M}
   \sum_{\substack{\gamma\ne1\\
          t((x^{-1}\gamma x)_v)\le\epsilon\ (\forall v)}}
                      |\Phi(x^{-1}\gamma x)|\,dx
                  &\le e^{-a/4}D_*,                    \label{amb-near-input}\\
 \operatorname{vol}(\mathcal B)
                  &\le e^{-a/4}\operatorname{vol}(\mathcal X),
                                                        \label{amb-rare-input}\\
 \int_{\mathcal B}K(x,x)\,dx&\le e^{-u}D.             \label{amb-deep-input}
\end{align}
Under the hypotheses of Theorem~\ref{amb-far}, these assertions
imply \eqref{amb-target}.
\end{proposition}
\begin{proof}
On $\mathcal M$, separate the identity term $m_0$, the nonidentity
terms with all distances at most $\epsilon$, and the remaining
terms. Absolute convergence, \eqref{amb-near-input} and
Theorem~\ref{amb-far} give
\[
 \left|\int_{\mathcal M}K(x,x)\,dx
                      -m_0\operatorname{vol}(\mathcal M)\right|
                  \le(e^{-a/4}+e^{-u})D_*.
\]
On $\mathcal B$, positivity in \eqref{amb-trace} gives
\[
 \left|\int_{\mathcal B}K(x,x)\,dx
                      -m_0\operatorname{vol}(\mathcal B)\right|
 \le e^{-u}D+m_0\operatorname{vol}(\mathcal B).
\]
Adding and using \eqref{amb-rare-input} proves the exact inequality
\begin{equation}\label{amb-error-before-absorption}
             |D-D_*|\le(2e^{-a/4}+e^{-u})D_*+e^{-u}D.
\end{equation}
Set $z=e^{-a/16}$ and $Y=|D-D_*|$. Since $\log n\le n$,
$u\ge a$, and $n\ge e^{15}$ gives $a\ge240$.
Consequently
\[
        e^{-a/4}\le z/8,\qquad e^{-u}\le z/8,\qquad z\le1;
\]
for the first two bounds it is enough that
$3a/16\ge\log8$ and $15a/16\ge\log8$.
Use $D\le D_*+Y$ in \eqref{amb-error-before-absorption}.
It follows that
\[
                    Y\le(z/2)D_*+Y/2,
\]
whence $Y\le zD_*$. This is \eqref{amb-target}, including the
case $D_*=0$.
\end{proof}

\subsection{Hypothesis checks and precise scope}

The small-distance condition does not force a rational element
to be the identity in the indefinite group. For example, in a
hyperbolic plane with Hermitian matrix
$\begin{pmatrix}0&1\\1&0\end{pmatrix}$, choose a nonzero
$a\in\mathcal O_E$ with $\bar a=-a$. Then
\[
 \gamma=\begin{pmatrix}1&Pa\\0&1\end{pmatrix},\qquad
 x_R=\begin{pmatrix}R&0\\0&R^{-1}\end{pmatrix}
\]
are isometries, $\gamma$ is a nonidentity integral unipotent
congruent to one modulo $P$, and
$x_R^{-1}\gamma x_R=\begin{pmatrix}1&Pa/R^2\\0&1\end{pmatrix}$
tends to one at every infinite place as $R\to\infty$.
One can take the same $R$ at every place and adjoin a positive
orthogonal summand of any dimension. This satisfies the actual
unimodularity and congruence conditions. It explains why
\eqref{amb-near-input} is a separate counting problem, even
after Proposition~\ref{amb-near-unipotent} is proved.

If the congruence condition is omitted, $\gamma=-1$ has distance
zero at every place and is not unipotent. If the small-distance
condition is omitted, the product
\[
 \begin{pmatrix}1&Pa\\0&1\end{pmatrix}
 \begin{pmatrix}1&0\\Pa&1\end{pmatrix}
 =\begin{pmatrix}1+P^2a^2&Pa\\Pa&1\end{pmatrix}
\]
is an integral isometry congruent to one modulo $P$, but has
trace $2+P^2a^2\ne2$, so is not unipotent. These are
counterexamples to weakened statements, not to the actual
propositions above. In a definite group a unitary unipotent is
the identity, by the spectral theorem. Thus, once the actual
kernel and trace identity are supplied, the definite case gives
the exact equality $D=D_*$.

The present section proves the rational unipotence and separation
arguments, the orbit packing, the full distant-term estimate
using the classical integral formula and its elementary
decay estimate, and the exact deduction
of \eqref{amb-target} from the displayed analytic inputs.
It retains the explicitly cited classical representation
theorems. It does not supply the automorphic
finite-rank kernel construction, the integrated nonidentity
unipotent estimate \eqref{amb-near-input}, the small-volume bound
\eqref{amb-rare-input}, or the cusp estimate
\eqref{amb-deep-input}. These are separate substantive inputs
in the actual source chain. No counterexample satisfying the
actual hypotheses was found in the components proved here.
The Galois-cohomology bounds at other coefficient primes are a
different dependency and are not consequences of this complex
multiplicity argument.


\section{The automorphic projection and its pointwise trace}
\label{ak-section}

We supply the kernel construction required in
Section~\ref{amb-section}.  The diagonal identity is proved
pointwise.  Equality of measurable kernels almost everywhere
on a product would not suffice for this purpose.
The only additional arithmetic theorem used here is the
classical finiteness theorem for automorphic forms of fixed
level, compact type and infinitesimal character; its precise
use is stated below.

\subsection{A matrix coefficient gives the required projection}

All Hilbert space inner products in this section are linear
in the second variable.  Let \(G\) be a unimodular locally
compact group, let \(\sigma\) be an irreducible unitary
representation on a Hilbert space \(H\), and let \(w\in H\)
have norm one.  Suppose its formal degree is \(d>0\).
Polarization of the formal-degree identity gives
\[
 d\int_G\overline{\langle a,\sigma(g)b\rangle}
                 \langle c,\sigma(g)v\rangle\,dg
                       =\langle c,a\rangle\langle b,v\rangle .
 \tag{AK1}
\]
Every integral in (AK1) is absolutely convergent by
Cauchy--Schwarz.  Put
\[
 a_w(g)=\langle w,\sigma(g)w\rangle,\qquad
                   \phi(g)=d\,\overline{a_w(g)}.
 \tag{AK2}
\]
Assume \(\phi\in L^1(G)\).

\begin{proposition}\label{ak-projection}
For every strongly continuous unitary representation \(\pi\)
on a Hilbert space \(\mathcal H\), the bounded operator
\[
                  P=\int_G\phi(g)\pi(g)\,dg
\]
is an orthogonal projection.  Moreover
\[
                  P\pi(x)P=a_w(x)P \qquad(x\in G).
 \tag{AK3}
\]
Evaluation at \(w\) is a linear bijection
\[
 \operatorname{Hom}^{\rm cont}_G(H,\mathcal H)
                       \longrightarrow\operatorname{ran}P,
                  \qquad T\longmapsto Tw .
 \tag{AK4}
\]
The intertwiner corresponding to a vector \(v\) has operator
norm \(\|v\|\).
\end{proposition}
\begin{proof}
The operator integral exists on each vector, by strong
continuity and the bound \(\|\pi(g)v\|=\|v\|\).
Its norm is at most \(\|\phi\|_1\).
Unitarity gives
\(\phi(g^{-1})=\overline{\phi(g)}\), so \(P=P^*\).
Apply (AK1) with \(a=b=w\), \(c=\sigma(y)w\), and
\(v=\sigma(x)w\).  The result is
\[
 \int_G\phi(g)\phi(x^{-1}g^{-1}y)\,dg
                                      =a_w(x)\phi(y).
 \tag{AK5}
\]
Indeed the second factor in the integrand is
\(d\langle\sigma(y)w,\sigma(g)\sigma(x)w\rangle\).
For \(x=1\), this proves \(\phi*\phi=\phi\).
Integrating (AK5) in \(\pi\), using Fubini for the
\(L^1\) functions, proves both \(P^2=P\) and (AK3).

If \(v\in\operatorname{ran}P\) has norm one, (AK3) gives
\(\langle v,\pi(g)v\rangle=a_w(g)\).  Therefore the map
\[
       \sum_i c_i\sigma(g_i)w\longmapsto
                          \sum_i c_i\pi(g_i)v
\]
preserves the squared norm: expand it as a finite sum
of the coefficients at \(g_i^{-1}g_j\).
It is well defined and extends to an isometry \(T_v\)
on \(H\), since the orbit of the nonzero vector \(w\)
is dense by irreducibility.  The same formula proves
equivariance.  Scaling treats every \(v\), including zero.
Conversely, any continuous intertwiner satisfies
\(P(Tw)=T\sigma(\phi)w=Tw\), by (AK1).
Density of the orbit of \(w\) proves uniqueness.
The isometry construction gives the asserted operator norm.
\end{proof}

We apply this to the product of the archimedean
representations in Section~\ref{amb-section}.
Write \(d_\infty=\prod_vd_v\) and
\(w=\bigotimes_v w_v\).  The finite tensor product has
formal degree \(d_\infty\): multiply (AK1) for elementary
tensors and use their density.  Let \(E_f=R(e)\) on
\(L^2(\mathcal X)\).  The finite convolution identities
in \eqref{amb-idempotent} show that \(E_f\) is an
orthogonal projection.  It commutes with every
archimedean right translation.  Hence
\[
 Q=E_fP=R(\Phi),\qquad
 \Phi(g_\infty,g_f)=d_\infty
       \overline{\langle w,\sigma(g_\infty)w\rangle}\,e(g_f)
 \tag{AK6}
\]
is also an orthogonal projection.  Formula (AK4) restricts
to a linear bijection between its image and the intertwiner
space in \eqref{amb-D-definition}: the condition on an
intertwiner is \(E_fT=T\).  Thus the rank of \(Q\), once
proved finite, is exactly the required \(D\).

\subsection{Two-point counting and convergence of the series}

Retain \(n,f,q,h\) and \(\tau_x\) from
Section~\ref{amb-section}, with \(n\ge3\) and \(q\le n/2\),
as in its large-\(n\) application.  We first record two extensions
of its lattice argument.

\begin{lemma}\label{ak-count}
If \(T:\Lambda_y\to\Lambda_x\) is a nonzero
\(\mathcal O_E\)-linear map, then
\[
              \left(\prod_v\|T_v\|_{\rm op}\right)^{1/f}
                                      \ge\lambda_x\lambda_y.
 \tag{AK7}
\]
Consequently the number of rational \(\gamma\) for which
\((x^{-1}\gamma y)_f\in K_f\) and
\(t((x^{-1}\gamma y)_v)\le r_v\), with all \(r_v\ge0\),
is at most
\[
 \left(\frac{3\sqrt n}{\tau_x\tau_y}\right)^{2n^2f}
                               \exp\left(2n^2\sum_v r_v\right).
 \tag{AK8}
\]
The function \(x\mapsto\tau_x\) is positive and continuous.
\end{lemma}
\begin{proof}
In the proof of Lemma~\ref{amb-map-separation}, use
the quotient of \(\Lambda_y\) by the kernel, and the
image in \(\Lambda_x\).  The quotient size is at most
\(\lambda_y^{-1}\), and the image size is at least
\(\lambda_x\).  The same covolume inequality proves (AK7).
The difference of two rational matrices counted in (AK8)
is such a map.  Lemma~\ref{amb-packing}, in the
\(2n^2\)-dimensional real matrix spaces, proves (AK8)
exactly as it proved \eqref{amb-ball-bound}.

The finite point lattice is locally constant modulo \(K_f\).
At infinity, a change of positive metric by maps whose
norms and inverse norms are at most \(e^{r_v}\)
changes every sublattice size by a factor between
\(e^{-\sum_vr_v/f}\) and \(e^{\sum_vr_v/f}\).
Taking infima preserves these bounds.  Taking the minimum
with \(1/2\) preserves them as well.  Since the maps tend
to the identity when the point tends to a fixed point,
these inequalities prove continuity.  Positivity was proved
in Lemma~\ref{amb-size-duality}.
\end{proof}

For \(q>0\), assume
\[
                         H=\frac{h}{8\sqrt{2q}}>2n^2.
 \tag{AK9}
\]
This is much weaker than \eqref{amb-parameters}.
The proved compact-type bound implies, at each place,
\[
 \bigl|\langle w_v,\sigma_v(g)w_v\rangle\bigr|
 \le \exp\!\left(-\frac h8
                \min\{t(g)^2,t(g)/\sqrt{2q}\}\right)
                         \le e^H e^{-Ht(g)}.
 \tag{AK10}
\]
For \(t\le1\), the last right side is at least one;
for \(t\ge1\), the minimum in the exponent is
\(t/\sqrt{2q}\).  This proves the second inequality.

The coefficient is integrable.  Here is the growth
calculation needed to justify that assertion.  In the
Cartan coordinates \(a_1\ge\cdots\ge a_q\ge0\) on
\(U(n-q,q)\), one has \(t^2=2\sum_i a_i^2\).
The positive restricted roots have multiplicities
\[
 e_i\pm e_j:2,\qquad e_i:2(n-2q),\qquad 2e_i:1.
\]
The Jacobian of the polar map, up to a fixed positive
constant depending on the Haar normalization, is
the product of \(\sinh(\alpha(a))^{m_\alpha}\).
This follows by differentiating the polar map on each
restricted-root subspace: its two opposite directions
give the factor \(\sinh(\alpha(a))\).
Bounding each such factor by its exponential gives
\[
 J(a)\le C\exp\left(2\sum_i(n-2i+1)a_i\right)
                                \le C e^{\sqrt2 n\sqrt q\,t}.
\]
Since \(q\le n\), condition (AK9) implies
\(H>\sqrt2 n\sqrt q\).  Integrating (AK10) over the
positive chamber, and then over the two compact factors,
is finite; the remaining radial volume is at most a
constant times \(t^{q-1}\,dt\).  When \(q=0\) the
archimedean group is compact, so integrability is immediate.
The finite function \(e\) is compactly supported.
Thus \(\Phi\in L^1(G_\infty\times G_f)\), as required
for (AK6).

\begin{theorem}\label{ak-kernel}
Under (AK9), or in the compact case \(q=0\), the series
\[
                    K(x,y)=\sum_{\gamma\in\Gamma}
                                        \Phi(x^{-1}\gamma y)
 \tag{AK11}
\]
converges absolutely and locally uniformly.
It defines a continuous kernel on
\(\mathcal X\times\mathcal X\), and for some finite
constant \(C_\Phi\),
\[
 \sum_\gamma|\Phi(x^{-1}\gamma y)|
                  \le C_\Phi(\tau_x\tau_y)^{-2n^2f}.
 \tag{AK12}
\]
It represents \(Q\) on continuous compactly supported
test functions.
\end{theorem}
\begin{proof}
Positivity of \(e\) gives \(|e(g_f)|\le e(1)\), by its
two-by-two positive matrix, as in
Lemma~\ref{amb-positive-type}.
For \(q>0\), separate the series into the shells
\(j_v\le t((x^{-1}\gamma y)_v)<j_v+1\),
where all \(j_v\) are nonnegative integers.
Equations (AK8)--(AK10) bound its absolute value by
\[
 m_0\left(\frac{3\sqrt n}{\tau_x\tau_y}\right)^{2n^2f}
 e^{f(H+2n^2)}
 \sum_{(j_v)\in\mathbb Z_{\ge0}^{\,f}}
                             e^{-(H-2n^2)\sum_vj_v}.
\]
The sum is \((1-e^{-(H-2n^2)})^{-f}\).
This proves (AK12), including \(e=0\).
For \(q=0\), all distances are zero and (AK8)
directly bounds the finite sum.

On a compact set of pairs, \(\tau_x\tau_y\) has a positive
lower bound.  The same geometric sum bounds the tails
uniformly.  The union of the rational elements with all
distances bounded on such a compact set is finite:
their finite components lie in one fixed compact adelic
set and their archimedean matrices are bounded.
After a fixed denominator is cleared, their entries
belong to a fixed lattice in the Minkowski space of \(E\),
whose bounded subsets are finite.  Thus these are tails
of the original series, not only tails of a varying
partition.  Local uniform convergence and continuity
of each term prove continuity of \(K\).
Reindexing \(\Gamma\) proves its invariance in both
quotient variables, and inversion proves
\(K(y,x)=\overline{K(x,y)}\).

For \(F\in C_c(\mathcal X)\), unfolding gives
\[
 \int_{\mathcal X}K(x,y)F(y)\,dy
       =\int_{G_\infty\times G_f}\Phi(g)F(xg)\,dg .
 \tag{AK13}
\]
Absolute convergence follows from
\(\|F\|_\infty\|\Phi\|_1<\infty\).
The latter integral represents the bounded \(L^2\)
operator \(R(\Phi)\), by its defining operator integral
and Fubini against a compactly supported test function.
The left side is continuous locally in \(x\), by
continuity of \(K\) on compact sets.  Thus (AK13)
gives a continuous representative of \(QF\).
\end{proof}

\subsection{A continuous kernel determines its diagonal}

\begin{lemma}\label{ak-diagonal}
Let \(X\) be a locally compact second countable space
with a positive Radon measure of full support.
Let \(Q\) be an orthogonal projection on \(L^2(X)\),
represented on \(C_c(X)\) by a continuous Hermitian
kernel \(K\).  Then every vector of \(\operatorname{ran}Q\)
has a unique continuous representative.  For an orthonormal
basis \((F_i)\) of that range, with these representatives,
\[
 K(x,y)=\sum_iF_i(x)\overline{F_i(y)},\qquad
 K(x,x)=\sum_i|F_i(x)|^2\ge0
 \tag{AK14}
\]
at every point.  In particular
\[
                \int_X K(x,x)\,dx
                    =\dim\operatorname{ran}Q
 \tag{AK15}
\]
as an equality of nonnegative extended real numbers.
\end{lemma}
\begin{proof}
Choose nonnegative \(u_{x,j}\in C_c(X)\), of integral
one, supported in successively smaller neighborhoods
of \(x\).  Such functions exist by local compactness,
full support and regularity of the measure.  The kernel
formula and continuity give
\[
 \langle Qu_{x,j},Qu_{y,l}\rangle
   =\langle u_{x,j},Qu_{y,l}\rangle
                   \longrightarrow K(x,y).
\]
Taking \(y=x\) shows that \(Qu_{x,j}\) is a Cauchy
sequence in \(L^2\).  Let its limit be \(k_x\).
It lies in the closed range of \(Q\), and
\[
                  \langle k_x,k_y\rangle=K(x,y).
 \tag{AK16}
\]
In particular \(x\mapsto k_x\) is continuous in Hilbert
norm, by expanding \(\|k_x-k_y\|^2\).

For \(F\in\operatorname{ran}Q\), define
\(\widetilde F(x)=\langle k_x,F\rangle\).
This is continuous.  It represents \(F\) in \(L^2\).
To verify this last assertion without assigning point
values to an \(L^2\) class, first let \(u\in C_c(X)\).
Pairing with \(v\in C_c(X)\) in the defining limit gives
\(\langle v,k_x\rangle=\overline{(Qv)(x)}\).
Integrating this identity and using \(Q=Q^*\) gives
\[
                Qu=\int_X u(x)k_x\,dx
\]
as a Hilbert space integral; the integrand is continuous
and bounded on its compact support.
Pairing with \(F\) gives
\(\int\overline{u(x)}\widetilde F(x)\,dx
=\langle Qu,F\rangle=\langle u,F\rangle\).
Thus \(\widetilde F=F\) almost everywhere.
Two continuous representatives agreeing almost everywhere
are equal, because a nonempty open set has positive measure.

Expand \(k_x\) and \(k_y\) in an orthonormal basis of
\(\operatorname{ran}Q\).  Parseval and Cauchy--Schwarz
give both formulas in (AK14), with an absolutely convergent
series at each pair of points.  Separability gives a
countable basis.  Tonelli and \(\|F_i\|_2=1\) prove
(AK15), including an infinite-dimensional range.
\end{proof}

\subsection{Finiteness and the exact trace}

We use the following classical arithmetic input.
For a reductive group over a number field, the space
of smooth automorphic forms of fixed compact open
level, fixed finite collection of archimedean compact
types, and annihilated by a fixed finite-codimensional
ideal of the center of the universal enveloping algebra,
is finite dimensional.  Moderate growth is required.
This is Harish--Chandra's finiteness theorem; the adelic
version follows from Borel's finite class decomposition.
The precise statements and the passage from the
classical to the adelic quotient are given in
\cite[Theorems 6.2.3 and 6.3.2]{getz-hahn-2022}.
These classical theorems are imported here.

\begin{theorem}\label{ak-finite-trace}
In the setting of Section~\ref{amb-section}, under (AK9)
or \(q=0\), the operator \(Q=R(\Phi)\) has finite rank
\(D\) as defined in \eqref{amb-D-definition}.
Its continuous kernel satisfies
\[
       K(x,x)=\sum_{i=1}^D|F_i(x)|^2,\qquad
                      \int_{\mathcal X}K(x,x)\,dx=D
 \tag{AK17}
\]
for a continuous orthonormal basis of its image.
\end{theorem}
\begin{proof}
It remains to check the hypotheses of the stated
finiteness theorem; neither finite rank nor an
integrable diagonal is assumed in this check.
By Proposition~\ref{ak-projection}, every vector
of the range is \(Tw\) for a continuous intertwiner
\(T\) satisfying \(E_fT=T\).
It is fixed by \(K_f(g)\), because \(e\) is
bi-invariant under that subgroup.
Its archimedean compact translates lie in the
image under \(T\) of the fixed lowest compact type.
The center of the universal enveloping algebra
acts on that image by the infinitesimal character
of \(\sigma\).

We justify smoothness directly.  The lowest compact
type occurs with multiplicity one in \(\sigma\).
Compact Schur orthogonality supplies a matrix
coefficient measure \(p_w\) on \(K_\infty\)
whose action on \(\sigma\) is projection onto
\(\mathbb Cw\).  Choose \(b\in C_c^\infty(G_\infty)\)
with \(c=\langle w,\sigma(b)w\rangle\ne0\);
an approximate identity supplies such \(b\).
Then
\[
                  a=c^{-1}p_w*b*p_w\in C_c^\infty(G_\infty),
                  \qquad \sigma(a)w=w.
\]
Consequently \(R(a)Tw=Tw\).
Convolution by \(a\), together with normalized Haar
measure on \(K_f(g)\), gives smooth representatives
of these \(L^2\) vectors.  To see this in local
coordinates, its kernel is a locally finite sum of
smooth functions: on compact sets all possible
rational matrices lie in a fixed compact adelic set.
For \(x\) in a fixed compact neighborhood the
other variable is restricted to the image of
that neighborhood times the fixed compact support.
Every derivative of the kernel is bounded there,
so Cauchy--Schwarz justifies differentiation under
the integral.  These representatives agree with
the continuous representatives in
Lemma~\ref{ak-diagonal}.
The compact-type and infinitesimal-character
identities therefore hold as identities of smooth
functions, not only of Hilbert space vectors.

Finally (AK12) and (AK14) give
\[
                 |(Tw)(x)|\le
                  C_\Phi^{1/2}\tau_x^{-2n^2f}\|Tw\|_2 .
 \tag{AK18}
\]
On a fixed finite adelic component, let \(\Lambda_*\)
be its point lattice at archimedean identity.
Changing its metric by \(g_\infty^{-1}\) gives
\[
       \tau_{(g_\infty,x_f)}
             \ge\tau_{(1,x_f)}
                   \left(\prod_v\|g_v\|_{\rm op}\right)^{-1/f}.
\]
The operator norms are at least one, since the
matrices are unitary for an involutory Hermitian form.
Thus (AK18) is polynomial growth in matrix coordinates
on every classical component.  There are finitely many
such components at the fixed level by Borel's class
decomposition.  Apply the classical finiteness theorem
on each component, with the same compact type and
the kernel of the fixed infinitesimal character.
Their finite direct sum contains the range of \(Q\).
This proves finite rank.
The bijection (AK4) identifies the rank with \(D\).
Now (AK17) follows from Lemma~\ref{ak-diagonal}.
\end{proof}

This proves the kernel and trace input in
Proposition~\ref{amb-final-reduction}.  Its nearby
unipotent, small cusp-volume and deep cusp-mass bounds
are separate quantitative statements; the qualitative
finiteness theorem gives none of those bounds.
The construction also explains the role of continuity:
on a nonatomic space an arbitrary measurable kernel
may be changed on the diagonal without changing its
integral operator.  No such change preserves the
continuous kernel constructed above.


% A direct proof of the volume estimate for small least size.
% All labels in this fragment have prefix rv-.
\section{The volume of the region with small least size}
\label{rv-section}

We prove the volume estimate needed in
\eqref{amb-rare-input}.  A sublattice of least normalized size
replaces the complete Harder--Narasimhan flag.  If that size is
less than one, the sublattice is totally isotropic.  Unfolding
the number of primitive isotropic sublattices of a fixed rank
then gives the required estimate.

The global classical inputs are the unitary Tamagawa-number
theorem, the Tamagawa number one of the special linear group,
the analytic class-number formula, the local-global principle
for Hermitian spaces, and Louboutin's residue bound.  We specify
their use and normalize all measures below.  We do not import
the submitted flag-volume estimate.

\subsection{Statement and the sublattice of least size}

Use the notation of Section~\ref{amb-section}.  Thus $E/F$ is
CM and unramified at every finite place,
$[F:\mathbb Q]=f$, $r_E=|D_E|^{1/(2f)}$, and $V$ has dimension
$n=p+q$ and signature $(p,q)$ at every real place.
The lattice $\Lambda_0$ is self-dual.  For
$x\in\mathcal X=G(F)\backslash G(\mathbb A_F)$ use the point
lattice and metrics in \eqref{amb-point-lattice}, and write
$\lambda_x$ for its least size \eqref{amb-size}.
Let $\mathbb P$ be Haar measure on $\mathcal X$, divided by
its total volume.

\begin{theorem}\label{rv-main}
Suppose
\[
       \log n\ge15,\qquad q^2\le n,\qquad
       \log r_E\le e^n .
\]
Then
\begin{equation}\label{rv-probability}
 \mathbb P\{x:\lambda_x<1/2\}
       \le \exp(-nf\log n/3).
\end{equation}
Consequently, for $b\ge2$, $g\ge1$, and
$\tau_0=(nfr_Egb)^{-11q}$, the complement of
$\{x:\lambda_x\ge\tau_0\}$ has probability at most
$\exp(-nf\log n/3)$.  This proves the weaker estimate
\eqref{amb-rare-input}.  No bound on $\log f$ is needed here.
\end{theorem}

All positive metrics below are compatible with the Hermitian
form: its matrix in an orthonormal basis has entries $1,-1$
on the diagonal.  A sublattice is called primitive if its
quotient in the ambient lattice is torsion-free.

\begin{lemma}\label{rv-minimum}
For every full $\mathcal O_E$-lattice $\Lambda$ the infimum
$\lambda(\Lambda)$ is attained by a nonzero primitive sublattice
$M$.  Put $s=s(M)$.  Every nonzero sublattice of $M$ has size
at least $s$, and every nonzero torsion-free quotient of $M$
has size at most $s$.
If $\Lambda$ is self-dual and $s<1$, then $M$ is totally
isotropic.  Its rank is therefore between one and $q$.
\end{lemma}
\begin{proof}
Saturation preserves rank and decreases covolume.  At a fixed
rank $j$, a primitive sublattice determines, up to sign, a
primitive decomposable vector in
$\bigwedge_{\mathbb Z}^{2fj}\Lambda$; its norm is the real
covolume.  A bounded ball contains only finitely many vectors
of this fixed exterior lattice.  Thus there are only finitely
many primitive sublattices of bounded covolume at each rank.
There are only $n$ possible ranks.  Taking, for example,
$s(\Lambda)$ as an initial upper bound proves attainment.

The assertion for sublattices follows from minimality.
For a primitive submodule $N\subset M$, of rank $a<j$,
orthogonal quotient metrics give
\[
                  H(M)=H(N)H(M/N).
\]
This identity follows by taking a basis of the underlying
$\mathbb Z$-lattice, extending a basis of $N$, and projecting
the remaining vectors orthogonally; the determinant becomes
block triangular.  The normalizing powers of $c_E$ multiply
because ranks add.  Since
$H(N)\ge s^{2fa}$ and $H(M)=s^{2fj}$, it gives
$s(M/N)\le s$.

Let $\overline{M}^{\,*}$ denote the conjugate
$\mathcal O_E$-linear dual, with its dual metrics.  The
Hermitian form restricted to $M$ defines a homomorphism
\[
                         T:M\longrightarrow\overline{M}^{\,*}.
\]
Its operator norm at each complex place is at most one:
it is the restriction of the unitary matrix $J_v$, followed
by orthogonal projection onto the dual of the subspace.
Suppose that $T\ne0$.  Let $Q=M/\ker T$ and $I=T(M)$, of
common positive rank $a$.  By the duality of sublattices and
quotients in Lemma~\ref{amb-size-duality}, the quotient
assertion just proved gives
\[
                 s(Q)\le s,\qquad s(I)\ge s^{-1}.
\]
The change-of-volume argument of
Lemma~\ref{amb-map-separation}, now with these domain and
target lattices, gives
\[
 s(I)\le s(Q)\Bigl(\prod_{v\mid\infty}
                                  \|T_v\|_{\rm op}\Bigr)^{1/f}.
\]
It would follow that $1\ge s^{-2}>1$, a contradiction.
Hence the restricted pairing is zero.  An isotropic subspace
of a Hermitian space of signature $(p,q)$ has dimension at
most $\min(p,q)$: projection to either definite summand is
injective on that subspace.  In particular its dimension
is at most $q$.
\end{proof}

For self-dual $\Lambda$, duality also gives $H(\Lambda)=1$.
If $q=0$, the lemma excludes $\lambda(\Lambda)<1$, while
$s(\Lambda)=1$ gives the reverse inequality.  Thus
$\lambda_x=1$ everywhere in that case.  In the rest of the
proof assume $q\ge1$.

\subsection{The quotient volumes and their Haar measures}

At a real place put on $\operatorname U(a,b)$ the Haar
measure whose density at the identity is Euclidean volume
for
\[
       \langle X,Y\rangle=\Re\operatorname{tr}(X^*Y)
\]
on its real Lie algebra.  We call this the
Hilbert--Schmidt normalization.  At finite places normalize
the stabilizer of a self-dual lattice to have volume one.
Let $\mathcal V_{a,b}$ be the resulting adelic quotient
volume, with signature $(a,b)$ at all real places, and set
$r=a+b>0$.
Write $d_F=|D_F|$ and
$L_\ell=L_F(\ell,\chi^\ell)$ for the factors in
Section~\ref{md-section}, where $\chi$ is the quadratic
character of $E/F$.

\begin{lemma}\label{rv-unitary-volume}
For $r>0$,
\begin{equation}\label{rv-unitary-volume-formula}
        \mathcal V_{a,b}
          =2d_F^{r^2/2}\prod_{\ell=1}^{r}L_\ell .
\end{equation}
Consequently, if $m=n-2j>0$, the ratio for signatures
$(p-j,q-j)$ and $(p,q)$ is
\begin{equation}\label{rv-volume-ratio}
 \frac{\mathcal V_{p-j,q-j}}{\mathcal V_{p,q}}
       =r_E^{-2fj(n-j)}
                      \prod_{\ell=m+1}^{n}L_\ell^{-1}.
\end{equation}
\end{lemma}
\begin{proof}
We use the classical unitary Tamagawa-number theorem in the
following form: over a fixed quadratic extension, the
Tamagawa quotient volume of a nondegenerate Hermitian
space is independent of its form and of its positive
dimension, and equals two in the commutative unitary case.
Weil proves the independence in \cite[\S54]{weil-siegel};
his unitary case has parameter $e=1/2$, as specified in
\cite[\S27]{weil-siegel}.  The threshold $r>4e-2$ is
therefore precisely $r>0$.  The absolute normalization
agrees with the unitary mass formula already stated in
\eqref{md-classical-mass}.

Here are the measure comparisons needed to apply this
theorem to our normalization.  At finite places the density
of the self-dual stabilizer is
$\prod_{\ell=1}^{r}(1-\eta_v^\ell q_v^{-\ell})$ by
\eqref{md-stabilizer-density}, independent of the Hermitian
form, including at dyadic places.  At infinity the conversion
from an algebraic invariant differential to
Hilbert--Schmidt volume is independent of the signature.
One can check the latter directly on the Lie algebra.
For a Hermitian matrix $H$, the map $X\mapsto HX$
identifies $X^*H+HX=0$ with the skew-Hermitian matrices.
After writing $H=AJA$ with $A$ positive definite and $J$
diagonal with entries $1,-1$, put $Y=AXA^{-1}$.
Then $HX=AJYA$.  Multiplication by $J$ is an isometry.
On the real space of skew-Hermitian matrices the map
$Z\mapsto AZA$ has determinant $|\det H|^r$:
diagonalize $A$, count its real diagonal coordinates and
its complex off-diagonal coordinates.  The corresponding
finite-place determinant is a unit in an integral
self-dual basis.  In arbitrary rational bases the
remaining determinant factors cancel by the product
formula.  Thus changing the signature introduces no
new constant in the normalized quotient volume.

It suffices to compute that volume for signature $(r,0)$.
In Hilbert--Schmidt normalization the compact group has
volume
\begin{equation}\label{rv-compact-volume}
 U_r:=\operatorname{vol}\operatorname U(r)
      =2^{r(r-1)/2}\prod_{\ell=1}^{r}
                        \frac{2\pi^\ell}{(\ell-1)!}.
\end{equation}
For completeness, in the fibration
$\operatorname U(r-1)\subset\operatorname U(r)\to S^{2r-1}$
the off-diagonal tangent directions have squared norm
twice their sphere norm; the remaining tangent direction
has unchanged norm.  The quotient metric volume is
$2^{r-1}$ times the standard sphere volume
$2\pi^r/(r-1)!$.  Induction, starting with the circle
of length $2\pi$, proves \eqref{rv-compact-volume}.
Multiplying the mass in Theorem~\ref{md-mass-ratio} by
$U_r^f$ cancels all factorials and powers of $2\pi$ and
gives \eqref{rv-unitary-volume-formula}.
Finally $d_F=r_E^f$ and
$n^2-(n-2j)^2=4j(n-j)$ prove \eqref{rv-volume-ratio}.
\end{proof}

We also need a general linear quotient volume.  At the
$f$ complex places use
\[
 da=\prod_{v\mid\infty}
              \frac{\prod_{\alpha,\beta}
                     d\Re A_{v,\alpha\beta}\,
                     d\Im A_{v,\alpha\beta}}
                   {|\det A_v|^{2j}},
 \qquad
 u(A)=\prod_{v\mid\infty}|\det A_v|^2.
\]
For a projective rank-$j$ $\mathcal O_E$-module $P$ embedded
as a lattice in $E^j$, let $\operatorname{Aut}(P)$ act on
$\operatorname{GL}_j(\mathbb C)^f$ by left multiplication,
and let $\mathcal G(P)$ be the volume of the part of that
quotient with $1\le u(A)<e$.
The number $u(A)$ is invariant under $\operatorname{Aut}(P)$,
since the absolute norm of the determinant of an integral
automorphism is one.

\begin{lemma}\label{rv-gl-volume}
With one representative of each isomorphism class of
projective rank-$j$ $\mathcal O_E$-modules,
\begin{equation}\label{rv-gl-formula}
 \mathcal G_j:=\sum_{[P]}\mathcal G(P)
     =c_E^{j^2}\operatorname{Res}_{s=1}\zeta_E(s)
                         \prod_{\ell=2}^{j}\zeta_E(\ell).
\end{equation}
For any nonnegative measurable function $h$,
\begin{equation}\label{rv-gl-radial}
 \int_{\operatorname{Aut}(P)\backslash
                         \operatorname{GL}_j(\mathbb C)^f}
                  h(u(A))\,da
      =\mathcal G(P)\int_0^\infty h(u)\,\frac{du}{u}.
\end{equation}
\end{lemma}
\begin{proof}
Here the classical input is the Tamagawa number one of
$\operatorname{SL}_j$ over a number field, recorded, for
example, in \cite[\S57]{weil-siegel}; for $j=1$ the special
linear group is trivial.  Together with the analytic
class-number formula for $E$, applied to the determinant
torus, it says that the norm-one adelic quotient of
$\operatorname{GL}_j$ has volume one with the following
normalization:
\begin{equation}\label{rv-tamagawa-gl}
 |D_E|^{-j^2/2}
       \bigl(\operatorname{Res}_{s=1}\zeta_E(s)\bigr)^{-1}
       \prod_{v\mid\infty}2^{j^2}\,da_v
       \prod_{w\nmid\infty}(1-q_w^{-1})^{-1}\,da_w.
\end{equation}
At a finite place, $da_w=|\det A|_w^{-j}dA$, with
$M_j(\mathcal O_{E,w})$ of additive volume one.
The measure on the norm-one group is obtained by
disintegration against $d\log|\det|_{\mathbb A_E}$.
To explain the compatibility, the invariant differential
on $\operatorname{GL}_j$ decomposes into the invariant
differential on its determinant-one fiber and $dt/t$.
The determinant sequence is split; applying Fubini to it
multiplies the special-linear Tamagawa number by the
idele-class norm-one volume.  The analytic class-number
formula gives the latter as one after the residue factor
in \eqref{rv-tamagawa-gl} is included.  The factor two in
each complex additive coordinate and the discriminant
factor are the usual additive adelic normalization.
This also checks \eqref{rv-tamagawa-gl} directly when $j=1$.

Counting ordered independent columns over the residue
field gives
\[
 \operatorname{vol}_{da_w}\operatorname{GL}_j(\mathcal O_{E,w})
                 =\prod_{\ell=1}^{j}(1-q_w^{-\ell}).
\]
The convergence factor in \eqref{rv-tamagawa-gl} cancels
the term with $\ell=1$.  Dividing out the finite compact
group and sorting the remaining components by the
associated projective module therefore gives
\[
 1=c_E^{-j^2}
       \bigl(\operatorname{Res}_{s=1}\zeta_E(s)\bigr)^{-1}
       \left(\prod_{\ell=2}^{j}\zeta_E(\ell)^{-1}\right)
       \sum_{[P]}\mathcal G(P).
\]
This proves \eqref{rv-gl-formula}; no assumption that
$P$ is free has been imposed.

For \eqref{rv-gl-radial}, multiplication at every complex
place by $u^{1/(2fj)}I_j$ gives a central one-parameter
subgroup with determinant norm $u$.  Haar disintegration
along this subgroup has radial measure $du/u$ and a
constant fiber volume.  The slab $1\le u<e$ has radial
volume one, so its volume $\mathcal G(P)$ is exactly that
constant.  This proves the formula, first for intervals
and then for all nonnegative measurable functions.
\end{proof}

\subsection{Primitive isotropic sublattices and the unipotent radical}

Fix a rank $j$ with $1\le j\le\min(p,q)$ and put
$m=n-2j>0$ and $d=2mj+j^2$.
Let $P$ be a primitive totally isotropic sublattice of a
self-dual Hermitian lattice $\Lambda$.

\begin{lemma}\label{rv-integral-splitting}
There is an integral orthogonal decomposition
\[
                     \Lambda=\mathbb H(P)\perp\Lambda',
\]
where $\mathbb H(P)$ is the hyperbolic module
$P\oplus\overline P^{\,*}$ and $\Lambda'$ is self-dual
of rank $m$ and signature $(p-j,q-j)$ at each real place.
The stabilizer of $P$ in $\operatorname{Aut}(\Lambda)$
is a semidirect product of its integral unipotent radical
with
$\operatorname{Aut}(P)\times\operatorname{Aut}(\Lambda')$.
\end{lemma}
\begin{proof}
Finite unramifiedness implies
$\operatorname{Tr}_{E/F}(\mathcal O_E)=\mathcal O_F$.
Indeed the local trace is surjective in an unramified
quadratic extension, and it is the sum map in the split
case.  The image of the global trace is an ideal of
$\mathcal O_F$ whose localization at every finite prime
is the unit ideal.  Choose $\theta\in\mathcal O_E$ with
$\theta+\bar\theta=1$.

Primitivity implies that $\Lambda/P$ is projective over
the Dedekind domain $\mathcal O_E$.  Restriction of the
perfect pairing on $\Lambda$ thus gives a surjection
$\Lambda\to\overline P^{\,*}$, and this surjection has
an $\mathcal O_E$-linear section.  The pairing of its
image with $P$ is perfect.  If $B$ is the Hermitian
Gram map of the section, subtract from the section
the map into $P$ corresponding to $\theta B$.
The new Gram map is
$B-\theta B-\bar\theta B=0$.
This calculation may be made after localization in
bases; it is intrinsic and hence holds for projective
modules.  We have obtained an integral isotropic
section, whose sum with $P$ is $\mathbb H(P)$.
Its pairing is perfect, so it is an orthogonal direct
summand.  Its orthogonal complement is the required
self-dual $\Lambda'$.

Every automorphism of $P$ extends to $\mathbb H(P)$
by its inverse conjugate dual, and every automorphism
of $\Lambda'$ extends by the identity.  Conversely,
an isometry preserving $P$ acts on $P$ and on
$P^\perp/P=\Lambda'$ by such automorphisms.
After removing these actions it lies in the unipotent
radical.  This proves the stated semidirect product.
The argument used a trace-one element, rather than
division by two, so it remains integral at dyadic places.
\end{proof}

At infinity choose hyperbolic coordinates
$E^j\oplus E^m\oplus E^j$ in which the Hermitian matrix is
\[
 \begin{pmatrix}0&0&I_j\\0&J'&0\\I_j&0&0\end{pmatrix},
 \qquad J'=\operatorname{diag}(1,\ldots,1,-1,\ldots,-1),
\]
and the compatible positive metric is the ordinary
Euclidean metric.  The unipotent radical consists of
\[
 n(Z,B)=
 \begin{pmatrix}
 I_j&-Z^*J'&B\\
 0&I_m&Z\\
 0&0&I_j
 \end{pmatrix},
                  \qquad B+B^*=-Z^*J'Z.
\]
Its tangent vectors are parametrized by $Z\in M_{m,j}(\mathbb C)$
and skew-Hermitian $C\in M_j(\mathbb C)$, and have
squared Hilbert--Schmidt norm
$2\|Z\|_{\rm HS}^2+\|C\|_{\rm HS}^2$.
Give the radical the Haar measure with this tangent density.

\begin{lemma}\label{rv-radical-covolume}
In the product of the radicals at the $f$ complex places,
the integral subgroup has covolume
\begin{equation}\label{rv-radical-formula}
                  r_E^{fd/2}H(P)^{n-j}.
\end{equation}
Here $H(P)$ is taken in the positive metrics just chosen.
\end{lemma}
\begin{proof}
Projection to $Z$ identifies the quotient by the center
with $\operatorname{Hom}_{\mathcal O_E}
(\overline P^{\,*},\Lambda')$.  It is surjective on
integral points: for any integral $Z$ one may take
$B=-\theta Z^*J'Z$.  The fibers are translates of the
integral skew-Hermitian central lattice.  Multiplication
in these coordinates is triangular with diagonal
Jacobian one, so its Haar measure is the product of
the two additive measures.  A fundamental domain may
be chosen by reducing $Z$ first and then the central
coordinate.  Its volume is the product of the two
additive covolumes, even though the group need not
be abelian.

For the $Z$ coordinate, the elementary determinant
identity for tensor products gives normalized height
\[
 H\bigl(\operatorname{Hom}(\overline P^{\,*},\Lambda')\bigr)
                =H(P)^m H(\Lambda')^j=H(P)^m.
\]
This identity follows in free bases from the determinant
of the Gram matrices and, for projective modules, by
localizing their inclusions in free lattices and
multiplying the resulting indices.
Its ordinary additive covolume is
$c_E^{mj}H(P)^m$.
The factor two in the squared $Z$-norm multiplies this
by $2^{fmj}$, giving $r_E^{fmj}H(P)^m$.

For clarity we also calculate the center, including its
diagonal coordinates.  Put
$\mathcal O_E^-=\{z\in\mathcal O_E:\bar z=-z\}$.
In the selected complex embeddings,
the real projection of $\mathcal O_E$ is
$\tfrac12\mathcal O_F$, by trace surjectivity.
Since $\operatorname{covol}(\mathcal O_E)
=2^{-f}d_F$, orthogonal projection gives
\[
              \operatorname{covol}(\mathcal O_E^-)
                         =\sqrt{d_F}=r_E^{f/2}.
\]
An off-diagonal skew-Hermitian entry has a second,
conjugate entry and hence has squared norm twice its
ordinary complex squared norm.  Its integral
covolume is $2^fc_E=r_E^f$.  For $P=\mathcal O_E^j$
the total central covolume is therefore
$r_E^{f j^2/2}$.

For general $P$, write it, in an $E$-basis, as
$\bigoplus_i\mathfrak a_i$.  An off-diagonal coordinate
belongs to $\mathfrak a_i\bar{\mathfrak a}_k$ and
contributes the additional factor
$N\mathfrak a_i\,N\mathfrak a_k$.
The invariant ideal $\mathfrak a_i\bar{\mathfrak a}_i$
is $\mathfrak b_i\mathcal O_E$ with
$N_{F/\mathbb Q}\mathfrak b_i=N_{E/\mathbb Q}\mathfrak a_i$.
This follows prime by prime, both in the split and in
the unramified inert case.  Its diagonal skew part is
$\mathfrak b_i\mathcal O_E^-$ and contributes
$N\mathfrak a_i$.  Thus the total ideal factor is
$(\prod_iN\mathfrak a_i)^j$.
For an arbitrary positive metric on $P$, a complex
change of basis $A$ acts on the center by $C\mapsto ACA^*$,
whose real determinant is $|\det A|^{2j}$.
Multiplication over the complex places supplies exactly
the remaining metric factor in $H(P)^j$.
The central covolume is
$r_E^{fj^2/2}H(P)^j$.
Multiplying it by the $Z$ covolume proves
\eqref{rv-radical-formula}, since $m+j=n-j$.
\end{proof}

\subsection{The real Iwasawa factor}

Put
\[
 S_a=\prod_{\ell=1}^{a}\frac{2\pi^\ell}{(\ell-1)!},
                   \qquad S_0=1.
\]
The stabilizer of the isotropic $j$-plane has Levi
subgroup
\[
       \operatorname{GL}_j(\mathbb C)\times
                          \operatorname U(p-j,q-j).
\]
Its embedding is
$(A,b)\mapsto\operatorname{diag}(A,b,A^{-*})$.
Let $K=\operatorname U(p)\times\operatorname U(q)$
carry probability Haar measure.

\begin{lemma}\label{rv-iwasawa}
For the Hilbert--Schmidt measures on the unitary group,
its radical, and the smaller unitary group, and for the
ordinary complex matrix Haar measure $da$ used above,
\begin{equation}\label{rv-iwasawa-formula}
 \begin{split}
 \int_{\operatorname U(p,q)}F(g)\,dg
   =c_j\,2^{j^2}
       \int_N\int_{\operatorname{GL}_j(\mathbb C)}
       \int_{\operatorname U(p-j,q-j)}\int_K
       &F\bigl(n\operatorname{diag}(A,b,A^{-*})k\bigr)\\[-1mm]
       &\hspace{-18mm}\cdot|\det A|^{-2(n-j)}
                         \,dk\,db\,da\,dn ,
 \end{split}
\end{equation}
for every nonnegative measurable $F$, where
\begin{equation}\label{rv-iwasawa-constant}
 c_j=
 2^{-d/2}\frac{U_pU_q}
           {2^{j^2/2}U_jU_{p-j}U_{q-j}}
 =2^{-j(j+1)/2}
                \frac{S_pS_q}{S_jS_{p-j}S_{q-j}} .
\end{equation}
\end{lemma}
\begin{proof}
The compact group $K$ is transitive on isotropic
$j$-planes: the projections of an orthonormal isotropic
frame to its positive and negative summands are scaled
orthonormal frames with the same scale; extend each
to a unitary basis.  This gives $G=NMK$.
The redundancy is the compact subgroup
$K_M\simeq\operatorname U(j)\times
\operatorname U(p-j)\times\operatorname U(q-j)$.
On its first factor the embedding duplicates a tangent
matrix, so its metric is twice the ordinary
Hilbert--Schmidt metric.  Its volume is
$2^{j^2/2}U_jU_{p-j}U_{q-j}$.
Likewise the full embedded general linear factor has
twice its usual tangent metric, and real dimension
$2j^2$, accounting for $2^{j^2}$ in
\eqref{rv-iwasawa-formula}.

It remains to compute the transverse Jacobian.
The radical Lie algebra is orthogonal to its Cartan
involutive image.  Its projection onto the orthogonal
complement of the compact Lie algebra is
$(X-\vartheta X)/2$, of norm $\|X\|/\sqrt2$.
It is orthogonal to the noncompact Levi directions.
There are $d$ radical directions, so this projection
contributes $2^{-d/2}$.
The quotient of the compact volumes is $U_pU_q/\operatorname{vol}K_M$,
giving the first expression for $c_j$.

Finally conjugation by the Levi acts on the complex
$Z$ coordinate with real determinant
$|\det A|^{2m}$ and on the skew-Hermitian central
coordinate with determinant $|\det A|^{2j}$.
Its modular character is consequently
$|\det A|^{2(n-j)}$.
Left invariance therefore requires its reciprocal in
the $NMK$ integral.  This determines the measure
globally from the just computed density at the identity.
Inserting \eqref{rv-compact-volume} into the first
expression in \eqref{rv-iwasawa-constant} gives the
second, by cancellation of the powers of two.
\end{proof}

\subsection{Unfolding the number of isotropic sublattices}

Let $N_j(x,T)$ count the primitive totally isotropic
rank-$j$ sublattices of $\Lambda_x$ whose size in its
point metrics is less than $T$.  This number is finite:
the exterior-lattice argument in Lemma~\ref{rv-minimum}
gives finiteness for any bound on the covolume.
It is measurable, being a countable sum of indicators
of continuous height inequalities on the arithmetic
components of the adelic quotient.

\begin{proposition}\label{rv-mean-value}
For $1\le j\le\min(p,q)$ and $m=n-2j>0$,
\begin{equation}\label{rv-mean-value-formula}
 \int_{\mathcal X}N_j(x,T)\,d\mathbb P(x)
       =\frac{\kappa_j}{n-j}\,
                        T^{2fj(n-j)}\qquad(T>0),
\end{equation}
where
\begin{equation}\label{rv-kappa}
 \kappa_j=
 r_E^{-fd/2}\,2^{-fj(j+1)/2}
 \left(\frac{S_pS_q}{S_{p-j}S_{q-j}S_j}\right)^f
 \operatorname{Res}_{s=1}\zeta_E(s)
 \frac{\prod_{\ell=2}^{j}\zeta_E(\ell)}
      {\prod_{\ell=m+1}^{n}L_\ell},
 \qquad d=2mj+j^2 .
\end{equation}
\end{proposition}
\begin{proof}
We give the orbit and measure details of the unfolding.
After division by $K_f$, the adelic quotient is the
disjoint union of
$\operatorname{Aut}(\Lambda)\backslash G_\infty$,
over the isometry classes in the genus of $\Lambda_0$.
Finite compact volume is one, so the measures on these
components are precisely the archimedean quotient
measures.  For each component, split the sum defining
$N_j$ into orbits under $\operatorname{Aut}(\Lambda)$.
The usual change of variables in a fundamental domain
gives
\[
 \int_{\operatorname{Aut}(\Lambda)\backslash G_\infty}
                    \sum_P {\bf1}_{s_x(P)<T}\,dx
  =\sum_{[P]}
       \int_{\operatorname{Aut}(\Lambda,P)\backslash G_\infty}
                                  {\bf1}_{s_x(P)<T}\,dx .
\]
Here the sum on the right is over primitive isotropic
sublattices modulo the indicated automorphism group.
All summands are nonnegative, so Tonelli justifies
this equality before integrability is known.

The isometry classes of marked pairs $(\Lambda,P)$
are exactly the pairs consisting of an isomorphism
class of projective rank-$j$ modules $P$ and an
isometry class of self-dual lattices $\Lambda'$ of
signature $(p-j,q-j)$.
Indeed Lemma~\ref{rv-integral-splitting} gives this
description and its converse by
$\mathbb H(P)\perp\Lambda'$.
For the converse to lie in the prescribed rational
space, we use the classical local-global principle
for Hermitian spaces
\cite[Theorem 2.4.1]{kirschmer-herm}.
At every finite place all unimodular Hermitian
lattices over the unramified quadratic algebra are
isometric: first use the finite-field classification
in Lemma~\ref{md-unitary-order}, and then the successive
Gram-matrix lifting in Proposition~\ref{md-lifting}.
The split case is the corresponding perfect-pairing
calculation.  At infinity the signatures agree.
Thus the rational spaces are isometric and the
resulting lattice lies in the same genus.
An isometry of marked pairs induces an isomorphism
of $P$ and an isometry on $P^\perp/P$; conversely
these extend to the hyperbolic orthogonal sums.
No class is lost by requiring a free basis for $P$,
because no such requirement has been made.

For one marked pair, its stabilizer is the semidirect
product in Lemma~\ref{rv-integral-splitting}.  Apply
\eqref{rv-iwasawa-formula} at the $f$ complex places.
Its radical quotient contributes
$r_E^{fd/2}H(P)^{n-j}$ by
Lemma~\ref{rv-radical-covolume}.  Denote by
$V(\Lambda')$ the volume of
$\operatorname{Aut}(\Lambda')\backslash
\operatorname U(p-j,q-j)^f$.
If $u=\prod_v|\det A_v|^2$, the size of the fixed
$P$ in the point metric is
\[
                         s_x(P)=(H(P)/u)^{1/(2fj)}.
\]
Indeed $n$ acts trivially on $P$, the Levi acts there
as $A$, and the right compact factor preserves
the positive norm.  Thus the integral of this
marked pair equals
\[
 \begin{split}
 c_j^f2^{fj^2}\,r_E^{fd/2}H(P)^{n-j}V(\Lambda')
       \int_{\operatorname{Aut}(P)\backslash
                           \operatorname{GL}_j(\mathbb C)^f}
       u^{-(n-j)}
       {\bf1}_{u>H(P)T^{-2fj}}\,da .
 \end{split}
\]
This also describes a fundamental domain in the
semidirect quotient: first reduce the Levi
coordinate and then the radical coordinate.
Integral Levi automorphisms preserve the radical
lattice and have global modular character one.
Consequently the radical covolume just used has no
extra index or stabilizer factor.

By \eqref{rv-gl-radial}, the last integral is
\[
 \mathcal G(P)\int_{H(P)T^{-2fj}}^\infty
                                u^{-(n-j)}\,\frac{du}{u}
  =\frac{\mathcal G(P)}{n-j}
                    H(P)^{-(n-j)}T^{2fj(n-j)}.
\]
The power of $H(P)$ cancels exactly.  Summing over
the two sets of lattice classes and dividing by
$\mathcal V_{p,q}$ proves
\[
 \int N_j(x,T)\,d\mathbb P(x)
 =c_j^f2^{fj^2}r_E^{fd/2}\mathcal G_j
      \frac{\mathcal V_{p-j,q-j}}{\mathcal V_{p,q}}
                  \frac{T^{2fj(n-j)}}{n-j}.
\]
Inserting \eqref{rv-gl-formula},
\eqref{rv-volume-ratio}, and
\eqref{rv-iwasawa-constant}, and using
$c_E=(r_E/2)^f$, yields \eqref{rv-kappa}.
In particular the computed integral is finite,
justifying all ordinary integral assertions as well.
\end{proof}

\subsection{An elementary bound for the constant}

\begin{lemma}\label{rv-kappa-bound}
Under the hypotheses of Theorem~\ref{rv-main},
for every $1\le j\le q$,
\begin{equation}\label{rv-kappa-small}
                 \kappa_j\le
                    \exp(-jnf\log n/2).
\end{equation}
\end{lemma}
\begin{proof}
The inequalities imply $p=n-q\ge q$ and
$n-2j\ge n-2\sqrt n\ge2$, so all preceding formulas
apply.  The first two powers in \eqref{rv-kappa}
are at most one, since $r_E\ge1$.

The factorials give a particularly large saving.
Direct cancellation shows
\[
 \frac{S_q}{S_{q-j}S_j}
   =\pi^{j(q-j)}
          \prod_{a=1}^{j}\frac{(a-1)!}{(q-j+a-1)!}
   \le\pi^{j(q-j)}.
\]
Moreover
$S_p/S_{p-j}\le[2\pi^p/(p-j)!]^j$.
Put $h=p-j\ge n-2\sqrt n$.  Hence
\[
 \log\frac{S_pS_q}{S_{p-j}S_{q-j}S_j}
       \le j\bigl(\log2+(n-j)\log\pi-\log(h!)\bigr).
\]
The integral of $\log x$ from one to $h$ gives
$\log(h!)\ge h\log h-h+1$.
Since $\log2<1$, $\log\pi<2$, and $h\le n$,
the last displayed right side is at most
$j(3n-h\log h)$.
For $\sqrt n\ge64$,
\[
 \log h\ge\log n+\log(1-2/\sqrt n)
                          \ge\log n-4/\sqrt n .
\]
Thus
\[
 3n-h\log h
 \le -n\log n+2\sqrt n\log n+4\sqrt n+3n .
\]
The ratio of the last three terms to $n\log n$
is at most
$1/32+1/240+1/5<1/4$, using $\log n\ge15$.
Consequently
\begin{equation}\label{rv-sphere-saving}
 \log\frac{S_pS_q}{S_{p-j}S_{q-j}S_j}
                         \le-\frac34jn\log n.
\end{equation}

For the residue we use the following unconditional
classical bound, valid for every number field $K$
of degree $d>1$:
\[
 \operatorname{Res}_{s=1}\zeta_K(s)
       \le\left(\frac{e\log|D_K|}{2(d-1)}\right)^{d-1}.
 \qquad\text{\cite[Theorem 1(1)]{louboutin-residue}}
\]
Here $d=2f$ and $\log|D_E|=2f\log r_E$, so it gives
\begin{equation}\label{rv-residue-bound}
       \log\operatorname{Res}_{s=1}\zeta_E(s)
                           \le 2f(n+1).
\end{equation}
Indeed $f/(2f-1)\le1$ and
$\log\log r_E\le n$.  The discriminant is greater
than one for a CM field, so these logarithms are
defined; if the logarithm of the residue bound is
negative, the same upper estimate still holds.

For completeness, the remaining Euler products
require only
\[
 \sum_{\ell=2}^{\infty}\log\zeta(\ell)
 \le\sum_{\ell=2}^{\infty}(\zeta(\ell)-1)
 =\sum_{a=2}^{\infty}\frac1{a(a-1)}=1.
\]
Since $\zeta_E(\ell)\le\zeta(\ell)^{2f}$, the numerator
in \eqref{rv-kappa} contributes at most $2f$ to its
logarithm.  Each inverse factor $L_\ell^{-1}$ is at
most $\zeta_F(\ell)\le\zeta(\ell)^f$: compare
$1-\eta_v^\ell q_v^{-\ell}$ with
$(1-q_v^{-\ell})^{-1}$ at every finite prime.
The denominator therefore contributes at most $f$.
Combining these estimates with
\eqref{rv-sphere-saving} and \eqref{rv-residue-bound},
\[
        \log\kappa_j
          \le-\frac34jnf\log n+2fn+5f
          \le-\frac12jnf\log n .
\]
The last inequality follows from
$\tfrac14jn\log n\ge\tfrac{15}{4}n\ge2n+5$.
\end{proof}

\begin{proof}[Proof of Theorem~\ref{rv-main}]
The case $q=0$ was already proved.
For $q\ge1$, if $\lambda_x<1/2$,
Lemma~\ref{rv-minimum} gives a primitive isotropic
sublattice of some rank $1\le j\le q$ and size
less than $1/2$.  Hence, pointwise,
\[
              {\bf1}_{\lambda_x<1/2}
                     \le\sum_{j=1}^{q}N_j(x,1/2).
\]
Put $A=nf\log n$.  The mean-value formula and
\eqref{rv-kappa-small} give
\[
 \mathbb P\{\lambda_x<1/2\}
 \le\sum_{j=1}^{q}
        \frac{2^{-2fj(n-j)}}{n-j}\kappa_j
 \le\sum_{j\ge1}e^{-jA/2}
 \le2e^{-A/2}\le e^{-A/3}.
\]
Here $A\ge15$, so both final inequalities hold.
This proves \eqref{rv-probability}.
If $q\ge1$, then $nfr_Egb\ge2n>2$ and
$\tau_0<1/2$; its exceptional set is smaller.
If $q=0$, then $\tau_0=1$ and $\lambda_x=1$,
so the exceptional set is empty.
\end{proof}

\paragraph{Relation to the source and remaining estimates.}
The constant \eqref{rv-kappa} is exactly the source's
\texttt{kapBd} for a flag with one nonzero rank increment:
rewriting $r_E/2$ as its \texttt{vInv} gives its two
displayed powers of two.  The source's set
$\operatorname{stepLt}(1/2)\ne0$ is the set
$\lambda_x<1/2$.  In one direction, a lattice of largest
quotient size less than $1/2$ itself has size less than
$1/2$.  In the other, the minimizing lattice in
Lemma~\ref{rv-minimum} has all quotient sizes at most
its size, and therefore occurs in that sum of submodules.
Thus \eqref{rv-probability} also proves the stated
\texttt{NonzeroFlagBound}, by a different argument.
This section settles the rare-volume input only.
The integrated count of nearby nonidentity unipotents,
the deep-cusp mass bound, and the finite-rank kernel
identity in the ambient comparison remain separate.


% A direct proof of the integrated count of nearby rational isometries.
% Labels in this fragment have the prefix nu-.
\section{The integrated contribution of nearby rational isometries}
\label{nu-section}

We prove the estimate \eqref{amb-near-input}.  The argument
has three parts.  A canonical flag of small sublattices
puts every sufficiently small integral isometry in one
unipotent radical.  Successive matrix blocks give a
packing estimate in that radical.  Iteration of the
one-sublattice unfolding in Section~\ref{rv-section}
then integrates the estimate, with all constants.
The only classical global inputs are those already
specified in that section.

\subsection{Statement and elementary inequalities for lattices}

Retain the notation of Sections~\ref{amb-section}
and~\ref{rv-section}.  In particular $E/F$ is a
CM extension unramified at every finite place,
$[F:\mathbb Q]=f$, $n=p+q$, and the Hermitian lattices
are self-dual.  The positive metrics are compatible
with their Hermitian forms.  Haar probability on the
adelic quotient is denoted by $\mathbb P$.
Set
\begin{equation}\label{nu-parameters}
 \begin{gathered}
 A=nf\log n,\qquad X=nfr_Egb,\qquad
 \tau_0=X^{-11q},\\
 \epsilon=10^{-6}n^{-8}f^{-2},\qquad
 \delta=3\sqrt{n\epsilon}
                =3\cdot10^{-3}n^{-7/2}f^{-1}.
 \end{gathered}
\end{equation}
For a point $x$, let $N_\delta(x)$ be the number of
nonidentity integral isometries of its point lattice
whose matrices, in its orthonormal frames, satisfy
$\|\gamma_v-I\|_{\rm HS}\le\delta$ at every complex place.

\begin{theorem}\label{nu-main}
Suppose $C_g\ge0$, $b\ge2$, $1\le g\le b^{C_g}$,
$1\le q$, $q^2\le n$, and
\begin{equation}\label{nu-hypotheses}
 \begin{gathered}
 \log n\ge20,\qquad n\ge4\log(C_g+3),\\
 \log r_E\le e^n,\qquad \log f\le\sqrt n,\qquad
                       \log b\le e^{A/16}.
 \end{gathered}
\end{equation}
Then
\begin{equation}\label{nu-count-integral}
   \int_{\lambda_x\ge\tau_0}N_\delta(x)\,d\mathbb P(x)
                                      \le e^{-A/4}.
\end{equation}
Consequently, under the finite-level congruence
hypotheses of Proposition~\ref{amb-near-unipotent},
\begin{equation}\label{nu-coefficient-integral}
 \int_{\lambda_x\ge\tau_0}
  \sum_{\substack{\gamma\ne1\\
       t((x^{-1}\gamma x)_v)\le\epsilon\ (\forall v)}}
      |\Phi(x^{-1}\gamma x)|\,dx
                                      \le e^{-A/4}D_*.
\end{equation}
Terms where $\Phi$ vanishes may be included in the sum.
\end{theorem}

Here integral means preservation of the point lattice
at all finite places; it does not require a free
$\mathcal O_E$-basis.  All the counts are finite:
distinct integral matrices are separated by
\eqref{amb-two-matrices}, and their archimedean norms
are bounded.  They are measurable, being countable
sums of indicators of measurable conditions.

For a nonzero metrized projective lattice $L$, write
$\lambda(L)$ for its least sublattice size, and write
$\mu(L)$ for the supremum of the sizes of its nonzero
projective quotients, with quotient metrics.
The size $s$ and height $H$ are those of
\eqref{amb-size}.  Thus
$H(L)=s(L)^{2f\operatorname{rank}L}$.
A lattice is semistable if every nonzero sublattice
has size at least $s(L)$.

\begin{lemma}\label{nu-map-size}
If $T:L\to M$ is a nonzero integral map, then
\begin{equation}\label{nu-map-size-inequality}
 \left(\prod_v\|T_v\|_{\rm op}\right)^{1/f}
                              \ge\frac{\lambda(M)}{\mu(L)}.
\end{equation}
For a semistable lattice of size $s$,
$\lambda(L)=\mu(L)=s$.
\end{lemma}
\begin{proof}
Let $Q=L/(L\cap\ker T)$ and $I=TL$.  They have the
same positive rank $r$, and the quotient-to-image
map is an isomorphism of projective modules.
Its singular values at $v$ are at most
$\|T_v\|_{\rm op}$.  As in
Lemma~\ref{amb-map-separation}, change of covolume gives
\[
 H(I)\le H(Q)\prod_v\|T_v\|_{\rm op}^{2r}.
\]
Take $2fr$th roots and use
$s(I)\ge\lambda(M)$ and $s(Q)\le\mu(L)$.
If $L$ is semistable and $Q=L/S$ is a nonzero
projective quotient, the exact covolume identity
$H(L)=H(S)H(Q)$ gives $s(Q)\le s(L)$.
The identity quotient and the sublattice $L$ itself
give the reverse inequalities.
\end{proof}

\subsection{A canonical flag and its small automorphisms}

\begin{lemma}\label{nu-largest-minimum}
Every nonzero positive metrized lattice $L$ has a
unique primitive sublattice $M$ of maximum rank
among those with size $\lambda(L)$.  It is
semistable.  If $M\ne L$, then
\begin{equation}\label{nu-quotient-strict}
                        \lambda(L/M)>\lambda(L).
\end{equation}
\end{lemma}
\begin{proof}
The minimum is attained by the exterior-lattice
argument of Lemma~\ref{rv-minimum}, which applies to
any positive metric, and saturation preserves a
minimizer.  Choose one of maximum rank, and put
$s=\lambda(L)$.  Every sublattice of $M$ has size
at least $s$, so $M$ is semistable.

For a nonzero primitive sublattice $N$ of $L/M$,
let $\widetilde N$ be its inverse image in $L$.
The quotient metrics give the exact identity
\[
 H(\widetilde N)=H(M)H(N).
\]
If $s(N)<s$, the size of $\widetilde N$ is less than
$s$.  If $s(N)=s$, it is another minimizer of
larger rank.  Both are impossible.  The least
size of the nonzero lattice $L/M$ is attained,
which proves the strict inequality.

If $M=L$, uniqueness of the maximum-rank
primitive minimizer is immediate.  Otherwise,
if $B$ is another minimum-size sublattice, it too
is semistable of size $s$.  The quotient map
$B\to L/M$ has operator norm at most one.
If nonzero, Lemma~\ref{nu-map-size} and
\eqref{nu-quotient-strict} would give a geometric
mean operator norm greater than one.
Thus $B\subset M$.  This proves uniqueness of
the maximum-rank minimizer.
\end{proof}

\begin{proposition}\label{nu-canonical-flag}
For every self-dual point lattice $\Lambda$ there
is a uniquely specified flag of primitive
isotropic sublattices
\[
          0=P_0\subsetneq P_1\subsetneq\cdots
                                      \subsetneq P_a
\]
with the following properties.  Put
$Q_i=P_i/P_{i-1}$, $j_i=\operatorname{rank}Q_i$,
$J=\sum_i j_i$, and $M=P_a^\perp/P_a$.
Then
\begin{equation}\label{nu-flag-properties}
 \begin{gathered}
 j_i\ge1,\quad J\le q,\quad
 Q_i\text{ is semistable of size }\tau_i,\quad
             0<\tau_1<\cdots<\tau_a<1/2,\\
 \lambda(M)\ge1/2,\qquad \mu(M)\le2.
 \end{gathered}
\end{equation}
The flag is empty exactly when $\lambda(\Lambda)\ge1/2$.
If nonempty, $\tau_1=\lambda(\Lambda)$.
Every integral isometry $\gamma$ with
$\|\gamma_v-I\|_{\rm op}<1/4$ at every place acts
as the identity on all $Q_i$ and on $M$.
It therefore lies in the unipotent radical of the
parabolic stabilizing this flag.
\end{proposition}
\begin{proof}
If $\lambda(\Lambda)<1/2$, take the sublattice in
Lemma~\ref{nu-largest-minimum}.  It is isotropic
by Lemma~\ref{rv-minimum}.  The lattice
$P_1^\perp/P_1$ is self-dual with its compatible
quotient metric, by the integral hyperbolic
splitting in Lemma~\ref{rv-integral-splitting}.
It embeds isometrically as a sublattice of
$\Lambda/P_1$.  Its least size is consequently
strictly greater than $\tau_1$, by
\eqref{nu-quotient-strict}.  Repeat this construction
in the quotient for as long as the least size is
less than $1/2$.
Each step reduces both signature entries by a
positive rank, so there are at most $q$ steps.
Their inverse images give the stated flag.
The construction is unique at each step.
Self-duality of the last quotient and
Lemma~\ref{amb-size-duality} give $\mu(M)\le2$.

We check the automorphism assertion at the first
step; the same argument then applies to every
quotient.  The map
\[
        P_1\xrightarrow{\gamma-I}\Lambda
                                  \longrightarrow\Lambda/P_1
\]
is integral and has operator norm less than $1/4$.
If nonzero, Lemma~\ref{nu-map-size} gives a lower
bound greater than one, since
$\mu(P_1)=\tau_1<\lambda(\Lambda/P_1)$.
Thus $\gamma$ preserves the $E$-span of $P_1$.
It preserves $P_1$ itself, with equality: an
invertible linear map preserving this span
preserves it also under its inverse, and
$P_1=\Lambda\cap EP_1$.
The endomorphism $(\gamma-I)|_{P_1}$ must vanish,
since a nonzero endomorphism of the semistable
lattice $P_1$ has geometric mean norm at least one.

Unitarity now gives the identity action on
$\Lambda/P_1^\perp$, the conjugate dual of $P_1$,
and gives an isometry on $P_1^\perp/P_1$.
Its difference from the identity has operator
norm less than $1/4$, because restriction and
orthogonal quotient do not increase the norm.
Induction applies.  On the final self-dual
quotient a nonzero integral endomorphism has
geometric mean norm at least
$\lambda(M)^2\ge1/4$, by
Lemma~\ref{amb-map-separation}; hence the induced
isometry is the identity there too.

Relative to the flag and its perpendicular flag,
the resulting matrix is upper block triangular
with identity diagonal blocks.  These are exactly
the elements of the unipotent radical.  For an
empty flag the same last argument says $\gamma=I$.
\end{proof}

\subsection{Packing successive matrix blocks}

For a sequence $\mathbf j=(j_1,\ldots,j_a)$ of
positive integers, let $J=\sum_i j_i$, $m=n-2J$,
and put
\begin{equation}\label{nu-dimension}
 d(\mathbf j)=4\sum_{i<k}j_ij_k+2Jm+\sum_i j_i^2
             =2nJ-2J^2-\sum_i j_i^2\le2nJ.
\end{equation}
For positive $\boldsymbol\tau$ and a function $F$
on the positive real numbers, define
\begin{equation}\label{nu-layer-product}
 \begin{split}
 L_{\mathbf j}(\boldsymbol\tau;F)
 ={}&\prod_{i<k}
       F(\tau_i/\tau_k)^{2fj_ij_k}
       F(\tau_i\tau_k)^{2fj_ij_k}\\
    &{}\times\prod_i
       F(\tau_i)^{2fj_im}F(\tau_i^2)^{fj_i^2}.
 \end{split}
\end{equation}
This is the source's function \texttt{layerProd},
with the rank and measure conventions made explicit.

\begin{lemma}\label{nu-packing}
At a point with the nonempty canonical flag
of Proposition~\ref{nu-canonical-flag}, and for
$0<\delta<1/4$,
\begin{equation}\label{nu-packing-bound}
 1+N_\delta(x)\le
 L_{\mathbf j}\!\left(\boldsymbol\tau;
                         r\longmapsto1+\frac{4\delta}{r}\right).
\end{equation}
\end{lemma}
\begin{proof}
At each place choose orthonormal coordinates
adapted to the flag and its perpendicular flag.
The ordered diagonal blocks have sizes
\[
                j_1,\ldots,j_a,m,j_a,\ldots,j_1.
\]
The Hermitian form pairs block $i$ with block
$i^*=2a+2-i$ by the identity, and has a
signature matrix on the middle block.
Such coordinates follow by pairing an orthonormal
basis of the isotropic space with its image under
the unitary involution representing the form.
The remaining orthogonal complement has the
induced compatible metric.
The corresponding graded integral lattices are
\[
              Q_1,\ldots,Q_a,M,
                   \overline{Q_a}^{\,*},\ldots,
                   \overline{Q_1}^{\,*},
\]
with their quotient metrics.  These are integral
identifications: dualize the split projective
sequences using self-duality of $\Lambda$.

Write an element of the radical as $I+Z$, with
$Z$ strictly upper block triangular.  Order the
upper blocks by their height, namely the difference
between column and row indices.  The unitary
equation is
\[
              Z^\dagger+Z+Z^\dagger Z=0,
\]
where $\dagger$ is adjunction for the displayed
Hermitian form.  Adjunction preserves block
height.  A block of height $h$ in
$Z^\dagger Z$ depends only on blocks of smaller
positive height.  Thus, after those earlier
blocks are fixed, the blocks of height $h$ are
paired by a linear adjoint relation.  The
independent blocks, with their real dimensions
and separation scales, are as follows:
\[
\begin{array}{c|c|c}
 \text{map}&\text{real dimension}&\text{scale}\\ \hline
 Q_k\longrightarrow Q_i,\quad i<k
       &2j_ij_k&\tau_i/\tau_k\\
 \overline{Q_k}^{\,*}\longrightarrow Q_i,\quad i<k
       &2j_ij_k&\tau_i\tau_k\\
 M\longrightarrow Q_i
       &2j_im&\tau_i/2\\
 \overline{Q_i}^{\,*}\longrightarrow Q_i
       \text{ with skew-Hermitian difference}
       &j_i^2&\tau_i^2
\end{array}
\]
In the last row the Hermitian part of the block
is determined by earlier heights; its free part
is skew-Hermitian.  In every other row the whole
complex block is free.  The list selects one
from each reflected pair
$(r,s)\leftrightarrow(s^*,r^*)$; the fixed pairs
are precisely the last row.

Here is why the separation statements are
integral statements, although the orthogonal
splittings need not be defined over $E$.
Extend the isotropic flag by its perpendicular
flag to a filtration $F_\bullet\Lambda$ with the
displayed graded pieces.  For two integral
matrices agreeing at all smaller heights,
their difference $T$ satisfies
$T(F_s)\subset F_{s-h}$ at the first remaining
height $h$.  It therefore induces integral maps
$F_s/F_{s-1}\to F_{s-h}/F_{s-h-1}$.
Their matrices are exactly the differences of
the corresponding orthogonal blocks.
After fixing one reflected pair at height $h$,
its other member is fixed too.  One may hence
process the independent blocks at that height
in any order.  Lemma~\ref{nu-map-size} gives
the tabulated scales, since the $Q_i$ and their
duals are semistable of sizes $\tau_i$ and
$\tau_i^{-1}$, and $\mu(M)\le2$.

For each fixed set of earlier coordinates,
the next coordinate has norm at most $\delta$
at every place.  This is immediate for a block
of $Z$, since $\|Z\|_{\rm op}\le\delta$.
For a skew-Hermitian coordinate apply the
contractive projection $B\mapsto(B-B^*)/2$;
its differences are the original block
differences, because the Hermitian parts
have already been fixed.
The packing lemma \ref{amb-packing}, in the
product of these real normed spaces, bounds
the number of possibilities by
$(1+2\delta/r)^{fD}$ for a row of dimension
$D$ and scale $r$.
For the middle-block row this is
$(1+4\delta/\tau_i)^{2fj_im}$.
The other rows are bounded by replacing
$2\delta$ with $4\delta$.
Multiplication of the successive fibre
bounds proves \eqref{nu-packing-bound}.
All the small isometries are in this radical
by Proposition~\ref{nu-canonical-flag}, and
the identity accounts for the additional one.
\end{proof}

\subsection{The measure of flags by repeated unfolding}

Let $\mathcal F_{\mathbf j}(x)$ be the set of all
primitive isotropic flags with rank increments
$\mathbf j$.  No semistability condition is
imposed on these flags.  For such a flag let
$\tau_i=s_x(P_i/P_{i-1})$, with the induced
quotient metric.
Define
\begin{equation}\label{nu-kappa}
 \begin{split}
 \kappa_{\mathbf j}
  ={}&r_E^{-f d(\mathbf j)/2}
       2^{-\frac f2\sum_i j_i(j_i+1)}
       \left(\frac{S_pS_q}
        {S_{p-J}S_{q-J}\prod_i S_{j_i}}\right)^f\\
     &{}\times
       \bigl(\operatorname{Res}_{s=1}\zeta_E(s)\bigr)^a
       \frac{\displaystyle\prod_i\prod_{\ell=2}^{j_i}
                         \zeta_E(\ell)}
            {\displaystyle\prod_{\ell=n-2J+1}^{n}L_\ell}.
 \end{split}
\end{equation}
The factors $S_r$ and $L_\ell$ have the definitions
in Section~\ref{rv-section}.  Rewriting $r_E/2$
as the source's \texttt{vInv} shows that
\eqref{nu-kappa} is exactly its \texttt{kapBd}:
\[
 \frac{d(\mathbf j)}2-nJ+J^2-\frac J2
                      =-\frac12\sum_i j_i(j_i+1).
\]

\begin{proposition}\label{nu-flag-measure}
If $J\le q$ and $n-2J>0$, then every nonnegative
Borel function $B$ on $(0,\infty)^a$ satisfies
\begin{equation}\label{nu-unfolded-flags}
 \begin{split}
 &\int_{\mathcal X}
       \sum_{\mathcal P\in\mathcal F_{\mathbf j}(x)}
           B(\boldsymbol\tau(\mathcal P,x))\,d\mathbb P(x)\\
 &\quad=\kappa_{\mathbf j}
      \int_{(0,\infty)^a}
       B(\boldsymbol\tau)
       L_{\mathbf j}(\boldsymbol\tau;\operatorname{id})
                     \prod_i2fj_i\,\frac{d\tau_i}{\tau_i}.
 \end{split}
\end{equation}
The equality is an equality of nonnegative
integrals and does not assume their finiteness.
\end{proposition}
\begin{proof}
The one-sublattice unfolding
of Proposition~\ref{rv-mean-value} gives more
than its displayed cumulative formula.
It also disintegrates the measure over the
self-dual quotient $P^\perp/P$.
To see this directly, insert in that proof a
nonnegative function of this quotient and its
induced metric.  In the Levi coordinates
$n\,\operatorname{diag}(A,b,A^{-*})\,k$,
the size is $(H(P)/u)^{1/(2fj)}$, whereas the
quotient metric depends only on $b$.
The radical acts trivially on the quotient
and the compact factor is an isometry.
The radical covolume, the determinant
disintegration on $\operatorname{GL}_j$, and
the cancellation of $H(P)^{n-j}$ therefore
remain exactly the same.
The last unitary integral is simply the
integral of the inserted function.

After division by the original unitary volume,
the quotient variable consequently has its
Haar probability measure, and the size
variable has density
\[
       \kappa_j(E;p,q)\,
               2fj\,\tau^{\,2fj(n-j)}\,\frac{d\tau}{\tau}.
\]
The coefficient agrees with
\eqref{rv-mean-value-formula}, since integrating
this density over $(0,T)$ gives
$\kappa_jT^{2fj(n-j)}/(n-j)$.
The projective-module and quotient-lattice
classes are summed independently, exactly as
in the proof there.  This proves the
disintegration for nonnegative simple
functions, and monotone convergence proves
it for nonnegative measurable functions.

A flag above its first member is precisely
a flag in $P_1^\perp/P_1$.  Repeatedly apply
this disintegration with
\[
 t_{i-1}=j_1+\cdots+j_{i-1},\quad
 p_i=p-t_{i-1},\quad q_i=q-t_{i-1},\quad
 n_i=n-2t_{i-1}.
\]
Its coefficient is
$\prod_i\kappa_{j_i}(E;p_i,q_i)$.
The compact volumes and the $L$-factors
telescope to \eqref{nu-kappa}.  The
discriminant exponents telescope because
\[
 \sum_i\bigl(2j_in_i-3j_i^2\bigr)
                         =d(\mathbf j).
\]
Finally the radial powers satisfy the exact
identity
\[
 L_{\mathbf j}(\boldsymbol\tau;\operatorname{id})
          =\prod_i\tau_i^{\,2fj_i(n_i-j_i)}.
\]
Indeed, for each $i<k$ the ratio and product
factors together contribute
$\tau_i^{4fj_ij_k}$ and no power of $\tau_k$.
This proves \eqref{nu-unfolded-flags}.
All changes of order are justified by Tonelli.
\end{proof}

\subsection{Uniform constants and the truncated radial integral}

\begin{lemma}\label{nu-kappa-bound}
If $\log n\ge15$, $q^2\le n$,
$\log r_E\le e^n$, and $1\le J\le q$, then
\begin{equation}\label{nu-kappa-small}
                         \kappa_{\mathbf j}\le e^{-JA/2}.
\end{equation}
\end{lemma}
\begin{proof}
In the product of one-step compact factors
from the preceding proof, each integer
$p_i-j_i$ is at least $n-q-J\ge n-2\sqrt n$,
and $p_i+q_i\le n$.
The factorial proof of \eqref{rv-sphere-saving}
therefore applies with the original $n$
at every step.  Explicitly,
\[
 \log\frac{S_{p_i}S_{q_i}}
          {S_{p_i-j_i}S_{q_i-j_i}S_{j_i}}
 \le j_i\bigl(3n-h_i\log h_i\bigr)
 \le-\tfrac34j_in\log n,\qquad h_i=p_i-j_i.
\]
Sum over $i$.  The first two factors in
\eqref{nu-kappa} are at most one.
Equation~\eqref{rv-residue-bound} bounds the
$a$ residues by a logarithm of at most
$2af(n+1)$.
The elementary Euler-product estimates in
that same proof bound the numerator zeta
products by a logarithm of at most $2af$,
and all inverse $L$-factors together by $f$.
Here $n-2J\ge2$, so their indices are at least
two.  Since $a\le J$, we obtain
\[
 \log\kappa_{\mathbf j}
 \le-\tfrac34JA+2Jfn+4Jf+f
 \le-\tfrac34JA+Jf(2n+5)
 \le-\tfrac12JA.
\]
The last inequality uses
$n\log n/4\ge 15n/4\ge2n+5$.
\end{proof}

\begin{lemma}\label{nu-radial-bound}
Under the hypotheses of Theorem~\ref{nu-main},
let
\[
 \mathcal S_a=\{\boldsymbol\tau:
                 \tau_0\le\tau_1<\cdots<\tau_a<1/2\}.
\]
Then
\begin{equation}\label{nu-one-type}
 \begin{split}
 \kappa_{\mathbf j}\int_{\mathcal S_a}
   L_{\mathbf j}\!\left(\boldsymbol\tau;
                 r\longmapsto1+\frac{4\delta}{r}\right)
   L_{\mathbf j}(\boldsymbol\tau;\operatorname{id})
             \prod_i2fj_i\,\frac{d\tau_i}{\tau_i}
                  \le e\,e^{-3JA/8}.
 \end{split}
\end{equation}
\end{lemma}
\begin{proof}
On the ordered simplex, every argument
$\tau_i/\tau_k$ with $i<k$, $\tau_i\tau_k$,
$\tau_i$, or $\tau_i^2$ lies in $(0,1)$.
Multiplying the two layer products replaces
each such argument $r$ by $r+4\delta$.
Their product is therefore at most
$(1+4\delta)^{fd(\mathbf j)}$.
Equations~\eqref{nu-parameters} and
\eqref{nu-dimension} give
\[
 \log(1+4\delta)^{fd(\mathbf j)}
 \le4\delta f(2nJ)
 \le 24\cdot10^{-3}n^{-2}<1,
\]
since $J\le\sqrt n$.

Enlarging the ordered simplex to
$[\tau_0,1/2]^a$ bounds the remaining radial
integral by
\[
 \prod_i2fj_i\log\frac1{2\tau_0}
        \le\bigl(22f q^2\log X\bigr)^a.
\]
The logarithm is positive because $X\ge2n$
and $q\ge1$.  We verify explicitly that
\begin{equation}\label{nu-log-absorption}
                   22f q^2\log X\le e^{A/8}.
\end{equation}
The hypotheses give
\[
 \log X\le\log n+\sqrt n+e^n+(C_g+1)e^{A/16}
                         \le(C_g+3)e^{A/16}.
\]
Indeed $\log n+\sqrt n\le e^n$ and
$A/16\ge n$, since $\log n\ge20$.
Using $q^2\le n$,
\[
 \log(22f q^2\log X)
       \le A/16+\log22+\sqrt n+\log n+\log(C_g+3).
\]
For $n\ge e^{20}$ we have
$\log22\le\sqrt n$, $\log n\le\sqrt n$,
and $3\sqrt n\le n/4$.
The last term is at most $n/4$ by
\eqref{nu-hypotheses}.  Thus the right side
is at most $A/16+n/2\le A/8$.
This proves \eqref{nu-log-absorption}.
Combine it with $a\le J$ and
\eqref{nu-kappa-small} to obtain
\eqref{nu-one-type}.
\end{proof}

\begin{proof}[Proof of Theorem~\ref{nu-main}]
The parameter $\delta$ in \eqref{nu-parameters}
is less than $1/4$.  If the canonical flag of
$x$ is empty, $N_\delta(x)=0$.
If $\lambda_x\ge\tau_0$ and the flag is
nonempty, its size tuple belongs to
$\mathcal S_a$ by
Proposition~\ref{nu-canonical-flag}.
The packing bound therefore gives, pointwise,
\[
 \begin{split}
 {\bf1}_{\lambda_x\ge\tau_0}N_\delta(x)
 \le\sum_{J=1}^{q}\ 
       \sum_{\substack{\mathbf j\text{ a composition}\\
                       \text{of }J}}\ 
       \sum_{\mathcal P\in\mathcal F_{\mathbf j}(x)}
       {\bf1}_{\boldsymbol\tau(\mathcal P,x)\in\mathcal S_a}
       L_{\mathbf j}\!\left(\boldsymbol\tau(\mathcal P,x);
                      r\longmapsto1+\frac{4\delta}{r}\right).
 \end{split}
\]
Only the one canonical flag is needed to
dominate the left side; all other flags add
nonnegative terms.  Their possible lack of
semistability is immaterial, since the
packing estimate was applied only to the
canonical flag.
Apply Proposition~\ref{nu-flag-measure}
and Lemma~\ref{nu-radial-bound}.
There are $2^{J-1}$ compositions of $J$:
choose which of the $J-1$ spaces between
$J$ ordered units are separators.
Consequently
\[
 \int_{\lambda_x\ge\tau_0}N_\delta(x)\,d\mathbb P(x)
 \le e\sum_{J=1}^{q}2^{J-1}e^{-3JA/8}
 \le\frac{e\,e^{-3A/8}}{1-2e^{-3A/8}}
 \le2e\,e^{-3A/8}\le e^{-A/4}.
\]
Both final inequalities hold for $A\ge16$,
as is certainly the case here.

Finally, Proposition~\ref{amb-near-unipotent}
puts every nonzero summand in
\eqref{nu-coefficient-integral} among the
integral isometries counted by $N_\delta(x)$.
The unit-vector and positive-type bounds
give $|\Phi|\le m_0$.  Multiply
\eqref{nu-count-integral} by
$m_0\operatorname{vol}(\mathcal X)=D_*$.
This proves \eqref{nu-coefficient-integral},
including the case $m_0=0$.
\end{proof}

\paragraph{Scope.}
For fixed $C_g$, the two lower bounds on $n$
in \eqref{nu-hypotheses} give an explicit
sufficient threshold.  Thus the estimate has
the quantifiers of the source's
\texttt{NearModerate}; the stronger block
packing constant used here is not needed
for that conclusion.  In the compact case
$q=0$, a unipotent unitary matrix is the
identity, so the corresponding near sum is
zero.  This section does not assert the
deep-cusp mass estimate
\eqref{amb-deep-input}.  That estimate
requires control of automorphic functions,
in addition to the lattice counts proved here.


% A quantitative proof of the mass estimate in the deepest part of the cusp.
% Labels in this fragment have the prefix dc-.
\section{The mass in the deepest part of the cusp}
\label{dc-section}

We now prove the remaining estimate \eqref{amb-deep-input}.
The argument uses the zero constant term of a discrete-series
automorphic form, a small fundamental domain for one unipotent
radical, and an explicit local Sobolev inequality.  We import
the classical Euclidean transference theorem of Banaszczyk,
the cuspidality theorem for discrete-series matrix coefficients,
and the usual analytic continuation and functional equation
of a finite-order Hecke character.  These imports are stated
precisely where they are used.  The quantitative estimates
and their application to the actual kernel are proved below.

\subsection{Statement and an absolute quotient-volume bound}

Retain the notation of Sections~\ref{amb-section},
\ref{ak-section}, \ref{rv-section}, and \ref{nu-section}.
In particular, \(E/F\) is CM and unramified at every finite
place, \([F:\mathbb Q]=f\), and the Hermitian lattice is
self-dual.  Write
\[
 N=n^2f,\qquad X=nfr_Egb,\qquad \tau_0=X^{-11q}.
\]
Let \(K(x,y)\) be the orthogonal-projection kernel constructed
in Section~\ref{ak-section}, of finite rank \(D\).

\begin{theorem}\label{dc-main}
Suppose \(\log n\ge20\), \(1\le q\), \(q^2\le n\),
\(b\ge2\), \(g\ge1\), and
\(\lvert\lambda_{v,i}\rvert\le b^2\).
Suppose also that the hypotheses of
Section~\ref{ak-section} hold, so its absolutely locally
uniformly convergent projection kernel and finite-rank
identity are available.  Then
\begin{equation}\label{dc-main-bound}
        \int_{\lambda_x<\tau_0}K(x,x)\,dx
                                    \le e^{-N}D.
\end{equation}
No upper bound on \(f,r_E,g\), or \(\log b\) is needed
in the estimates of this section.
\end{theorem}

We may normalize finite Haar measure by
\(\operatorname{vol}(K_f)=1\), and archimedean Haar measure
by the Hilbert--Schmidt density on each real Lie algebra.
Changing Haar measure multiplies an integral kernel and
its diagonal by the reciprocal factor and leaves its
operator, rank, and integrated diagonal unchanged.
Thus this normalization does not change
\eqref{dc-main-bound}.

Let \(\chi\) be the quadratic Hecke character of \(E/F\),
and put \(d_F=|D_F|=r_E^f\).  It has finite conductor one
and is odd at every real place.

\begin{lemma}\label{dc-volume}
In the preceding normalization,
\begin{equation}\label{dc-volume-bound}
 L_F(1,\chi)\le3^f d_F^{1/4},\qquad
 \operatorname{vol}(\mathcal X)\le X^N,\qquad
 \operatorname{vol}(K_f(g))^{-1}\le g^{2N}.
\end{equation}
\end{lemma}
\begin{proof}
We use the classical Hecke theorem in the following form:
for a nontrivial primitive finite-order character, the
completed \(L\)-function is entire, satisfies its functional
equation with root number of absolute value one, and the
uncompleted function has polynomial growth on each fixed
vertical strip.  The completion and functional equation
are recorded in \cite[\S2.1, (2.1)--(2.3)]{zaman-hecke-estimates};
the qualitative growth needed here also follows from
the convexity estimate there, Lemma~2.1, together with
the Euler product and the functional equation.
For our totally odd character the completion is
\[
 \Lambda(s)=d_F^{s/2}
       \left\{\pi^{-(s+1)/2}
                    \Gamma\left(\frac{s+1}{2}\right)\right\}^{f}
                    L_F(s,\chi).
\]
Since \(\chi\) is real, its functional equation is
\(\Lambda(s)=W\Lambda(1-s)\), where \(|W|=1\).
The Euler product and the integral test give
\[
 |L_F(3/2+it,\chi)|\le\zeta(3/2)^f\le3^f.
\]
On the other boundary of the strip, the functional equation
and \(\Gamma(z+1)=z\Gamma(z)\) therefore give
\[
 |L_F(-1/2+it,\chi)|
       \le d_F\left(\frac{3}{\pi}
                            |1/4-it/2|\right)^f.
\]
Indeed the two gamma arguments are \(5/4-it/2\) and
\(1/4+it/2\), whose quotient in absolute value is
\(|1/4-it/2|\).

Put \(F(s)=L_F(s,\chi)/(s+2)^f\).  On
\(\Re s=3/2\), \(|F(s)|\le1\); on \(\Re s=-1/2\),
\(|F(s)|\le d_F\), because
\[
 \frac{|1/4-it/2|}{|3/2+it|}\le\frac12,\qquad
                         \frac{3}{2\pi}<1.
\]
The maximum principle on this strip consequently yields
\[
 |F(\sigma+it)|\le d_F^{(3/2-\sigma)/2}.
\]
For clarity, the unbounded-strip step follows by applying
the ordinary maximum principle on rectangles to
\[
 F(s)d_F^{(s-3/2)/2}
                  \exp\{\varepsilon(s^2-9/4)\}.
\]
Its vertical boundary values have absolute value at most
one.  Polynomial vertical growth makes its horizontal
boundary values tend to zero as the rectangle grows.
Then let \(\varepsilon\) decrease to zero.  Evaluation at
\(s=1\) proves the first assertion.

By \eqref{rv-unitary-volume-formula},
\[
 \operatorname{vol}(\mathcal X)
     =2d_F^{n^2/2}\prod_{\ell=1}^n L_F(\ell,\chi^\ell).
\]
For the remaining factors the Euler products give
\[
 \prod_{\ell=2}^n |L_F(\ell,\chi^\ell)|
 \le\left(\prod_{\ell=2}^{\infty}\zeta(\ell)\right)^f
 \le e^f.
\]
The last inequality follows from
\[
 \sum_{\ell=2}^{\infty}\log\zeta(\ell)
 \le\sum_{\ell=2}^{\infty}\sum_{a=2}^{\infty}a^{-\ell}
 =\sum_{a=2}^{\infty}\frac1{a(a-1)}=1.
\]
Thus the volume is at most
\(18^f r_E^{N/2+f/4}\).
Since \(X\ge18\), \(r_E\le X\), and \(n\ge2\), this is
at most \(X^{N/2+5f/4}\le X^N\).

Finally, reduction modulo \(g\) embeds \(K_f/K_f(g)\)
in the endomorphisms of
\(\widehat\Lambda_0/g\widehat\Lambda_0\).
The latter additive group has order \(g^{2fn^2}\).
This follows locally from freeness over a discrete
valuation ring and hence holds for a projective lattice.
The index, and therefore the reciprocal level volume,
is at most \(g^{2N}\).
\end{proof}

\subsection{A small fundamental domain for a radical}

We use the least sublattice size \(\lambda(L)\) and
greatest quotient size \(\mu(L)\) of
Section~\ref{nu-section}, always with induced and quotient
positive metrics.  Conjugate duality gives
\[
            \lambda(\overline L^{\,*})=\mu(L)^{-1}.
\]
If \(\mathfrak a\) is a fractional ideal of \(E\), then
\begin{equation}\label{dc-ideal-size}
 \lambda(\mathfrak a L)
     =|N_{E/\mathbb Q}\mathfrak a|^{1/(2f)}\lambda(L),
 \qquad
 \mu(\mathfrak a L)
     =|N_{E/\mathbb Q}\mathfrak a|^{1/(2f)}\mu(L).
\end{equation}
Indeed \(M\mapsto\mathfrak aM\) is a bijection between
sublattices, and between projective quotients, in the
two modules.  On rank \(r\) covolumes it multiplies
the normalized height by
\(|N\mathfrak a|^r\), as follows by localization.
Taking \(2fr\)-th roots proves the identities.

For a real Euclidean lattice \(L\), denote its covering
radius by \(\operatorname{covrad}(L)\), its first vector
minimum by \(\lambda_1(L)\), and its Euclidean integral
dual by \(L^\vee\).  We import Banaszczyk's theorem
\begin{equation}\label{dc-transference}
       \operatorname{covrad}(L)\lambda_1(L^\vee)
                         \le\frac{\dim_{\mathbb R}L_{\mathbb R}}2.
\end{equation}
This is \cite[Theorem~2.2]{banaszczyk-transference},
also stated with these Euclidean conventions in
\cite[Theorem~3.7.2, (3.54)]{bost-kunnemann-extensions}.
It is used for arbitrary lattices, without semistability.

\begin{lemma}\label{dc-hom-cover}
Let \(A,B\) be positive metrized projective
\(\mathcal O_E\)-lattices of ranks \(a,b\ge1\).
Give \(\operatorname{Hom}_{\mathcal O_E}(A,B)\) the
Euclidean norm
\((\sum_v\|T_v\|_{\rm HS}^2)^{1/2}\).  Then
\begin{equation}\label{dc-hom-cover-bound}
 \operatorname{covrad}\operatorname{Hom}(A,B)
          \le\frac{ab\sqrt f\,r_E}{2}
                                      \frac{\mu(B)}{\lambda(A)}.
\end{equation}
\end{lemma}
\begin{proof}
Write \(\mathfrak d_E^{-1}\) for the absolute
codifferent.  The trace pairing identifies the integral
dual of \(\operatorname{Hom}(A,B)\) with
\(\mathfrak d_E^{-1}\operatorname{Hom}(B,A)\).
Under the Euclidean Riesz identification, a map \(U\)
in this module corresponds to \(2U^*\), because
\[
 \operatorname{Tr}_{E/\mathbb Q}\operatorname{tr}(UT)
                  =2\sum_v\Re\operatorname{tr}(U_vT_v).
\]
The trace pairing is perfect for projective modules:
one checks this in free local bases at every finite
prime.  Formula \eqref{dc-ideal-size} and
\(|N\mathfrak d_E^{-1}|=|D_E|^{-1}=r_E^{-2f}\)
give
\(\lambda(\mathfrak d_E^{-1}A)=r_E^{-1}\lambda(A)\).
For \(U\ne0\), Lemma~\ref{nu-map-size} and the
arithmetic--geometric mean inequality give
\[
 \left(\sum_v\|2U_v^*\|_{\rm HS}^2\right)^{1/2}
       \ge2\sqrt f\,r_E^{-1}\frac{\lambda(A)}{\mu(B)}.
\]
Apply \eqref{dc-transference} in real dimension \(2fab\).
\end{proof}

Fix a point \(x\) with \(0<\tau=\lambda_x<1/2\).
Use the canonical isotropic flag of
Proposition~\ref{nu-canonical-flag}, with successive
semistable quotient sizes
\[
             0<\tau_1<\cdots<\tau_a<1/2,\qquad
             \tau_1=\tau,\qquad a\le q.
\]
Set \(\tau_{a+1}=1/2\).

\begin{lemma}\label{dc-radical-domain}
There is a nonzero primitive isotropic sublattice \(P\)
of rank \(j\le q\) such that the following holds.
Let \(U\) be the unipotent radical of its stabilizer,
and put
\[
 \Gamma_U(g)=
 \{\gamma\in U(F):\gamma\Lambda_x=\Lambda_x,\
                         (\gamma-1)\Lambda_x\subseteq g\Lambda_x\}.
\]
In orthonormal Witt coordinates at infinity, every
class in \(\Gamma_U(g)\backslash U_\infty\) has a
representative \(\exp A\exp C\), where \(C\) is central,
\([A,C]=0\), and
\begin{equation}\label{dc-radius}
 \|A\|_{\rm HS}+\|C\|_{\rm HS}\le\rho,\qquad
 \rho=\frac{n^2\sqrt f\,g r_E}{4}(2\tau)^{1/q}.
\end{equation}
Here and below a norm on a product of real places is
the square root of the sum of the squared norms.
\end{lemma}
\begin{proof}
An extension of lattices whose greatest quotient sizes
are at most \(s\) again has greatest quotient size at
most \(s\).  To see this, take any projective quotient,
saturate the image of the sublattice, and use the exact
covolume identity between this saturated image and
the remaining quotient.  Saturation decreases
covolume; both induced maps are contractions for
quotient metrics.  Each of the two factors therefore
has size at most \(s\).  Apply this successively to the
semistable grades of \(P_i\) to obtain
\(\mu(P_i)\le\tau_i\).
The self-dual core \(P_i^\perp/P_i\) has least size
\(\tau_{i+1}\) if \(i<a\), and at least \(1/2\)
if \(i=a\), by construction of the flag.

The product of the \(a\) ratios
\(\tau_i/\tau_{i+1}\) is \(2\tau<1\).
For some \(i\),
\[
 \frac{s}{t}:=\frac{\tau_i}{\tau_{i+1}}
       \le(2\tau)^{1/a}\le(2\tau)^{1/q}=:R.
\]
Take \(P=P_i\), \(j=\operatorname{rank}P\), and
\(M=P^\perp/P\), of rank \(m=n-2j\).
Then
\[
 \mu(P)\le s,\quad
 \lambda(\overline P^{\,*})\ge s^{-1},\quad
 \lambda(M)\ge t,\quad \mu(M)\le t^{-1},\quad
                         s^2\le R/4.
\]
The last assertion uses \(s/t\le R\) and \(t\le1/2\).

Use the integral hyperbolic splitting proved in
Section~\ref{rv-section}.  At infinity use orthogonal
Witt coordinates \(P\oplus M\oplus\overline P^{\,*}\),
so the Hermitian matrix is
\[
       \begin{pmatrix}0&0&I\\0&J_0&0\\I&0&0\end{pmatrix}.
\]
The first and central Lie coordinates of the radical are
\[
 A_Z=\begin{pmatrix}0&-Z^*J_0&0\\0&0&Z\\0&0&0\end{pmatrix},
 \qquad
 C_0=\begin{pmatrix}0&0&C\\0&0&0\\0&0&0\end{pmatrix},
 \qquad C^*=-C.
\]
Thus \(\|A_Z\|_{\rm HS}=\sqrt2\|Z\|_{\rm HS}\).
Projection of \(\Gamma_U(g)\) onto the \(Z\)-coordinate
is exactly
\[
               g\operatorname{Hom}(\overline P^{\,*},M).
\]
For surjectivity, choose \(\theta\in\mathcal O_E\)
with \(\theta+\bar\theta=1\), available because
\(E/F\) is unramified at finite places.  For an integral
\(Z\) divisible by \(g\), the group matrix with top
right block \(B=-\theta Z^*J_0Z\) is integral and
congruent to the identity modulo \(g\).
This statement is intrinsic to the integral splitting;
its expression in orthonormal real coordinates need not
be integral.  It nevertheless identifies the projected
lattice and its induced metric exactly.
Lemma~\ref{dc-hom-cover} now gives first-coordinate
covering radius, in the radical norm, at most
\begin{equation}\label{dc-first-radius}
                  \frac{jm\sqrt f\,g r_E}{\sqrt2}\frac{s}{t}.
\end{equation}

We give the corresponding duality for the center,
including the prime two.  Let \(\mathcal C(P)\) be the
skew-Hermitian maps \(\overline P^{\,*}\to P\).
The \(\mathcal O_F\)-pairing
\[
 \mathcal C(P)\times\mathcal C(\overline P^{\,*})
       \longrightarrow\mathcal O_F,\qquad
                         (C,U)\longmapsto-\operatorname{tr}(CU)
\]
is perfect.  In local free bases, an off-diagonal pair
is the trace pairing for the unramified quadratic
algebra.  A diagonal coordinate is its trace-zero
rank-one module.  At an inert prime choose a generator
\(\alpha\) of the quadratic algebra with unit
discriminant; then
\(2\alpha-\operatorname{Tr}(\alpha)\) generates its
trace-zero module and its square is a unit.
This holds also at a dyadic inert prime: the trace of
such an \(\alpha\) is a unit, and the equation
\(2a+\operatorname{Tr}(\alpha)b=0\) forces
\(b\in2\mathcal O_F\).
At a split prime the generator is \((1,-1)\).
These checks prove perfectness at every finite prime,
and hence globally.

The Euclidean integral dual of \(\mathcal C(P)\) is
therefore represented by
\(\mathfrak d_F^{-1}\mathcal C(\overline P^{\,*})\).
There is no factor two here: the matrix trace already
lies in \(F\), and the real inner product is
\(-\sum_v\operatorname{tr}(C_vU_v)\).
Because
\[
 |N_{E/\mathbb Q}(\mathfrak d_F^{-1}\mathcal O_E)|
                        =d_F^{-2}=|D_E|^{-1},
\]
every nonzero dual vector, regarded as a map
\(P\to\mathfrak d_F^{-1}\overline P^{\,*}\),
has Euclidean norm at least
\(\sqrt f/(r_Es^2)\), by Lemma~\ref{nu-map-size}.
Its real dimension is \(fj^2\); transference gives
central covering radius at most
\begin{equation}\label{dc-central-radius}
                        \frac{j^2\sqrt f\,g r_E}{2}s^2.
\end{equation}

Reduce the first coordinate by an element of
\(\Gamma_U(g)\), using \eqref{dc-first-radius}, and
then reduce the remaining central coordinate by
\(g\mathcal C(P)\), using \eqref{dc-central-radius}.
Group multiplication is triangular, so the second
operation leaves the first coordinate unchanged.
An affine translation of the central fiber has the
same covering radius.  The resulting representative
is \(\exp A_Z\exp C_0\); the two matrices commute.
Its sum of norms is at most
\[
 \sqrt f\,g r_ER\left(\frac{jm}{\sqrt2}+\frac{j^2}{2}\right)
 \le\sqrt f\,g r_ER\,j(n-j)
 \le\frac{n^2\sqrt f\,g r_E}{4}R.
\]
This proves the assertion.
\end{proof}

\subsection{Zero constant terms for the actual projection}

We use the following classical Harish--Chandra theorem.
A matrix coefficient between compact-finite vectors
of a discrete-series representation of a connected
semisimple real Lie group with finite center belongs
to its Harish--Chandra cusp Schwartz space.  In particular,
for every proper parabolic radical \(U\) its integrals
\[
                   \int_{U(\mathbb R)}c(aub)\,du
\]
converge absolutely and vanish, for all fixed \(a,b\).
The Schwartz assertion is in
\cite[\S33, Lemma~65, Corollary~1]{harish-chandra-ds-ii};
the cusp assertion and its constant-term application
are explicitly recalled in
\cite[\S4.1]{khare-larsen-savin}.
We import this local theorem, rather than a statement
about the global cuspidal spectrum.

It applies to \(U(p,q)\) as follows.
The map \(U(1)\times SU(p,q)\to U(p,q)\),
\((z,g)\mapsto zg\), is onto with finite kernel.
The pullback of an irreducible representation is a
unitary central character times an irreducible
representation of \(SU(p,q)\).  Square integrability
and compact finiteness pass across this cover because
the added center is compact.  A parabolic radical
is contained in \(SU(p,q)\); its constant-term
integral is thus covered by the stated theorem.
The same reasoning applies to the contragredient
coefficient and to a finite product of real places.

\begin{lemma}\label{dc-cuspidal}
Write the projection identity of Section~\ref{ak-section}
as
\[
 K(x,y)=\sum_{i=1}^D F_i(x)\overline{F_i(y)},\qquad
 F_i=T_iw,
\]
where each \(T_i:\sigma\to L^2(\mathcal X)\) is an
isometric intertwiner and \(w\) is the fixed unit
lowest-compact-type vector.  Every \(F_i\), every
archimedean right translate of it, and every one of
its right Lie derivatives is cuspidal.
\end{lemma}
\begin{proof}
For fixed \(y\), integrate the kernel column along
\(U(F)\backslash U(\mathbb A_F)\).
This quotient is compact.  Absolute locally uniform
convergence of the kernel series allows termwise
integration, including absolute values.  Decomposing
\(G(F)\) into left \(U(F)\)-cosets unfolds this integral
to
\[
 \sum_{\gamma\in U(F)\backslash G(F)}
       \int_{U(\mathbb A_F)}
                   \Phi(x^{-1}u^{-1}\gamma y)\,du .
\]
The total integral of the absolute values is finite
by the preceding compact-quotient calculation.
Fubini is therefore applicable to the finite and
infinite factors.  Each archimedean radical integral
vanishes by the local theorem, so every kernel
column has zero constant term.

These columns span \(\operatorname{span}\{F_1,\ldots,F_D\}\).
Indeed, if their coefficient vectors did not span
\(\mathbb C^D\), a nonzero linear combination of the
\(F_i\) would vanish at every point, contradicting
their orthonormality in \(L^2\).
Consequently all \(F_i\) are cuspidal.
Right translations preserve constant terms.
For Lie derivatives, differentiation under the
constant-term integral is valid because its domain
is compact and the functions are smooth.
\end{proof}

For the particular radical of Lemma~\ref{dc-radical-domain},
cuspidality also means zero average on
\(\Gamma_U(g)\backslash U_\infty\), with \(x_f\) fixed.
Here is the finite-place justification.
If
\[
 K_{U,f}(g)=U(\mathbb A_{F,f})\cap x_fK_f(g)x_f^{-1},
\]
then
\[
                 U(\mathbb A_{F,f})=U(F)K_{U,f}(g).
\]
This follows directly from the two-step coordinates:
approximate the first additive coordinate by a rational
one, lift it by the trace-one formula above, and then
approximate the central coordinate.  Density of a
finite-dimensional \(F\)-vector space in its finite
adeles follows by clearing denominators and the
Chinese remainder theorem.  Thus the adelic
constant-term quotient, after imposing right level
\(K_{U,f}(g)\), is precisely
\(\Gamma_U(g)\backslash U_\infty\).
The first-coordinate and central lattices also show
directly that this real quotient is compact.
Normalize its Haar measure to have total mass one.

\subsection{Derivatives on a bounded range of compact types}

At one real place use the Hilbert--Schmidt inner
product on the real Lie algebra \(\mathfrak g\) of
\(U(p,q)\), and let \(\mathfrak g=\mathfrak k\oplus
\mathfrak p\) be its orthogonal Cartan decomposition.
Put
\[
                \beta(P)=2\sqrt n\,b^2+4n^{3/2}+4P.
\]
In the product representation, let \(V_P\) be the
span of all products of at most \(P\) complexified
Lie algebra operators applied to the full lowest
compact type \(W=\bigotimes_v W_v\).
It is a finite-dimensional compact-invariant space
of smooth vectors.

\begin{lemma}\label{dc-derivative}
For a Hilbert--Schmidt unit vector \(Y\) in the Lie
algebra of one real place and \(u\in V_P\),
\[
                  \|d\sigma(Y)u\|\le\beta(P)\|u\|.
\]
For a general product Lie algebra vector \(Y\),
\begin{equation}\label{dc-derivative-bound}
       \|d\sigma(Y)u\|
                   \le\sqrt f\,\beta(P)\|Y\|_{\rm HS}\|u\|.
\end{equation}
\end{lemma}
\begin{proof}
We include the Casimir computation to fix its sign.
For orthonormal bases of \(\mathfrak k,\mathfrak p\)
the invariant form is \(B(A,B)=\Re\operatorname{tr}(AB)\),
negative on \(\mathfrak k\) and positive on
\(\mathfrak p\).  Thus
\[
 \Omega_G=\sum_{\mathfrak p}Y^2-\sum_{\mathfrak k}Y^2,
 \qquad \Omega_K=-\sum_{\mathfrak k}Y^2,
 \qquad
 \sum_{\mathfrak g}\|d\sigma(Y)u\|^2
                    =\langle u,(2\Omega_K-\Omega_G)u\rangle.
\]
The Casimir eigenvalues in this normalization are
\[
 c_G=\|\lambda_v\|^2-\|\rho\|^2,\qquad
 c_K(\mu)=\|\mu+\rho_K\|^2-\|\rho_K\|^2.
\]
These formulas follow by writing the complexified
Casimir in matrix units \(E_{ij}\): move each positive
root operator to the right of its negative root
operator and apply it to a highest vector.  The
commutators contribute \(2\rho\), and the diagonal
terms contribute the squared weight.  The same
calculation for the two compact blocks gives
\(\rho_K\).  On the noncompact representation the
Harish--Chandra parameter is, by definition, the
shifted parameter for this central character.

The lowest compact highest weight is
\(\kappa=\lambda_v+\rho-2\rho_K\), as checked in
Section~\ref{fd-section}.  We use the elementary bounds
\[
 \|\rho\|,\|\rho_K\|\le\frac12 n^{3/2},\qquad
              \|\kappa\|\le\sqrt n\,b^2+\frac32n^{3/2}.
\]
All weights of the compact module of highest weight
\(\kappa\) have norm at most \(\|\kappa\|\).
For completeness, the highest-weight construction
places each weight below \(\kappa\) by a sum of
positive roots.  The set of weights is Weyl invariant.
Testing an arbitrary real linear functional, and
moving it into the dominant chamber, therefore
places every weight in the convex hull of the Weyl
orbit of \(\kappa\).  That convex hull lies in the
closed ball of radius \(\|\kappa\|\).

A highest weight \(\mu\) occurring in \(V_P\) occurs
in \(W_v\otimes\mathfrak g_{\mathbb C}^{\otimes r}\)
for some \(r\le P\).  The weights of the adjoint
module are zero or \(e_i-e_j\), of norm \(\sqrt2\).
Consequently
\[
 \|\mu\|\le\|\kappa\|+\sqrt2P,\qquad
 c_K(\mu)\le
             \bigl(\sqrt n\,b^2+2n^{3/2}+2P\bigr)^2.
\]
If the expression inside parentheses is \(a_P\),
the energy above is at most
\[
 \left(2a_P^2+\frac14 n^3\right)\|u\|^2
                                      \le4a_P^2\|u\|^2.
\]
This is \(\beta(P)^2\|u\|^2\).
The calculation applies to the compact-isotypic
decomposition of every \(u\in V_P\), including in
the product representation.  Cauchy--Schwarz for a
linear combination of the orthonormal Lie basis
proves the single-place bound.  Sum over the \(f\)
places and apply Cauchy--Schwarz once more to obtain
\eqref{dc-derivative-bound}.
\end{proof}

\subsection{An explicit local Sobolev estimate}

Write
\[
 {\bf F}(u,y)=\bigl((T_1u)(y),\ldots,(T_Du)(y)\bigr)
                         \in\mathbb C^D,\qquad
                  \tau_y=\min\{\lambda_y,1/2\}.
\]
For \(u\in V_P\) these are smooth representatives
obtained by right differentiation and compact
translation of the \(F_i\).  They satisfy pointwise
equivariance and are right \(K_f(g)\)-invariant.
Each \(T_i\) is isometric, so
\begin{equation}\label{dc-L2}
       \int_{\mathcal X}\|{\bf F}(u,y)\|^2\,dy=D\|u\|^2.
\end{equation}

\begin{lemma}\label{dc-sobolev}
For \(u\in V_P\),
\begin{equation}\label{dc-sobolev-bound}
 \|{\bf F}(u,y)\|^2
 \le
  \bigl(200N^3\beta(P+N)^2\bigr)^N
       \left(\frac7{\tau_y^2}\right)^{2N}
          \frac{D}{\operatorname{vol}(K_f(g))}\|u\|^2.
\end{equation}
\end{lemma}
\begin{proof}
Choose a Hilbert--Schmidt orthonormal basis
\(Y_1,\ldots,Y_N\) of the product real Lie algebra,
each basis vector supported at one place.  Set
\[
 r=\frac1{20N^2},\qquad
 \Psi(t)=\exp(t_1Y_1)\cdots\exp(t_NY_N),
                         \quad t\in[-r,r]^N.
\]
We record explicit chart bounds.  Put \(S=\sum_j|t_j|\),
so \(S\le1/(20N)\).  In the ambient product of
matrix spaces, the \(j\)-th derivative column differs
from \(Y_j\) by norm at most \(e^S-1\le2S\).
The operator norm of the difference between the full
derivative and the linear isometry
\(t\mapsto\sum t_jY_j\) is at most
\(2S\sqrt N\le1/10\).
Integrating along a straight segment in the cube
shows that \(\Psi\) is injective.
Its left Maurer--Cartan columns are
\[
 A_j(t)=\operatorname{Ad}
       \bigl(\exp(-t_NY_N)\cdots\exp(-t_{j+1}Y_{j+1})\bigr)Y_j.
\]
They differ from \(Y_j\) by norm at most
\(e^{2S}-1\le3S\).  Their full derivative matrix
has smallest singular value at least
\(1-3S\sqrt N>1/2\).
Thus Haar measure in the chart has density at least
\(2^{-N}\), and \(\|A_j(t)\|_{\rm HS}<2\).

The map of the chart times \(K_f(g)\) into
\(\mathcal X\), translated by \(y\), has multiplicity
at most
\[
                    M_y=\left(\frac7{\tau_y^2}\right)^{2N}.
\]
Indeed any two lifts differ by a rational isometry
preserving \(\Lambda_y\), of operator norm less than
\(e^{2Nr}<2\) at every real place.
The difference of two such isometries is a nonzero
integral endomorphism.  Lemma~\ref{amb-map-separation}
gives geometric mean operator norm at least
\(\lambda_y^2\).  The packing argument of
Lemma~\ref{amb-packing}, in real matrix dimension
\(2n^2\) at each place, bounds their number by
\[
                  (1+4/\tau_y^2)^{2N}\le M_y.
\]
Consequently, for any fixed smooth vector \(v\)
in the span under consideration,
\begin{equation}\label{dc-chart-integral}
 \int_{[-r,r]^N}\|{\bf F}(v,y\Psi(t))\|^2\,dt
       \le\frac{2^N M_yD}{\operatorname{vol}(K_f(g))}\|v\|^2.
\end{equation}
This follows from the chart Jacobian, the multiplicity
bound, right level invariance, and \eqref{dc-L2}.

Here is an elementary Sobolev inequality sufficient
for this chart.  The one-dimensional fundamental
theorem of calculus, averaged over \([-r,r]\),
expresses \(h(0)\) as its average plus an integral
of \(h'\) with kernel of absolute value at most
\(1/2\).  Iterating in the \(N\) coordinates gives
\[
 h(0)=\sum_{I\subseteq\{1,\ldots,N\}}
          \int_{[-r,r]^N}k_I(t)\partial_Ih(t)\,dt,
 \qquad |k_I(t)|\le(2r)^{-N}r^{|I|}.
\]
The same identity holds for a finite-dimensional
Hilbert-space-valued function.  Cauchy--Schwarz gives
\[
 \|h(0)\|^2\le(2r)^{-N}(1+r^2)^N
                     \sum_I\int\|\partial_Ih(t)\|^2\,dt
              \le r^{-N}\sum_I\int\|\partial_Ih(t)\|^2\,dt.
\]
Use \(h(t)={\bf F}(u,y\Psi(t))\).
For \(I=\{i_1<\cdots<i_k\}\), direct differentiation
of the ordered product gives
\[
 \partial_Ih(t)=
 {\bf F}\bigl(d\sigma(A_{i_1}(t))\cdots
                       d\sigma(A_{i_k}(t))u,\ y\Psi(t)\bigr).
\]
Each coordinate is differentiated only once, so this
identity has no additional derivative terms.
Expand the \(A_j(t)\) in the fixed orthonormal basis.
The squared sum of coefficients of their tensor
product is at most \(4^k\); pointwise Cauchy--Schwarz
therefore bounds the squared norm by \(4^k\) times
the sum for all \(N^k\) fixed words of length \(k\).
Lemma~\ref{dc-derivative} bounds each such word by
\(\beta(P+N)^k\|u\|\).
Applying \eqref{dc-chart-integral} word by word yields
\[
 \|{\bf F}(u,y)\|^2
 \le r^{-N}\frac{2^NM_yD}{\operatorname{vol}(K_f(g))}
                 (1+4N\beta(P+N)^2)^N\|u\|^2.
\]
Since \(N\ge1\), \(\beta(P+N)\ge1\), and
\(r^{-1}=20N^2\), this proves the stated constant.
\end{proof}

\subsection{Iteration of the zero-average identity}

Fix \(x\) with \(\tau=\lambda_x<\tau_0\), and use the
radical and \(\rho\) of Lemma~\ref{dc-radical-domain}.
Let \(\Omega_x\) consist of the points
\(x\exp A\exp C\) with first and central radical
coordinates satisfying
\(\|A\|_{\rm HS}+\|C\|_{\rm HS}\le\rho\);
the finite component remains \(x_f\).
It contains \(x\), the paths
\(x\exp(t(A+C))\), \(0\le t\le1\), and representatives,
up to left rational and right \(K_f(g)\) action,
of all points on the radical quotient through \(x\).

If \(\rho\le\log2\), then throughout this set
\begin{equation}\label{dc-size-nearby}
                   \tau_y\ge\tau/2.
\end{equation}
Indeed at place \(v\) the positive metrics change
by a map and inverse of norm at most
\(\exp(\|A_v\|_{\rm HS}+\|C_v\|_{\rm HS})\).
Taking the normalized product over the \(f\) places
changes every sublattice size by at most
\(\exp(\rho/\sqrt f)\).  Since \(\tau<1/2\),
the same lower bound survives truncation at \(1/2\).

Set
\[
 P_0=8qN,\qquad \beta_*=\beta(P_0+N).
\]
Our parameters give the following simple bounds:
\begin{equation}\label{dc-numerical}
 \beta_*\le X^3,\qquad N\le X^2,\qquad
 156800\le X,\qquad
 \rho\le\tfrac12X^{-9}<\log2.
\end{equation}
For the first, \(q\le\sqrt n\) implies
\[
 \beta_*\le2\sqrt n\,b^2+4n^{3/2}+36n^{5/2}f.
\]
Divide by \(n^3f^3b^3\le X^3\); the quotient is at most
\[
        n^{-5/2}+\tfrac12n^{-3/2}+\tfrac92n^{-1/2}<1
                         \qquad(n\ge100).
\]
The other elementary bounds follow from \(X\ge2n\),
\(\log n\ge20\), and
\[
 \rho\le\frac{n^2\sqrt f\,gr_E}{2}\tau^{1/q}
                          \le\tfrac12X^2X^{-11}.
\]
Equations \eqref{dc-sobolev-bound},
\eqref{dc-size-nearby}, and \eqref{dc-numerical}
give, uniformly for \(y\in\Omega_x\) and \(u\in V_{P_0}\),
\begin{equation}\label{dc-initial-evaluation}
 \|{\bf F}(u,y)\|^2
   \le X^{13N}\tau^{-4N}
       \frac{D}{\operatorname{vol}(K_f(g))}\|u\|^2.
\end{equation}
Here \(200\cdot28^2=156800\):
the Sobolev constant contributes \(N^3\le X^6\),
\(\beta_*^2\le X^6\), and the remaining constant
contributes at most \(X\).

\begin{lemma}\label{dc-iteration}
Suppose \(M_{P+1}\) bounds
\(\|{\bf F}(u,y)\|/\|u\|\) for all nonzero
\(u\in V_{P+1}\), \(y\in\Omega_x\).
For \(0\le P<P_0\), the same quantity on \(V_P\)
is bounded by
\[
                   2\rho\sqrt f\,\beta_* M_{P+1}.
\]
\end{lemma}
\begin{proof}
For \(y=x\exp A\exp C\in\Omega_x\), differentiate
along \(x\exp(t(A+C))\).  The two coordinates commute,
and \(d\sigma(A+C)u\in V_{P+1}\).
Lemma~\ref{dc-derivative} gives
\[
 \|{\bf F}(u,y)-{\bf F}(u,x)\|
                  \le\rho\sqrt f\,\beta_*M_{P+1}\|u\|.
\]
The same inequality holds for every point of the
radical quotient, by its representatives in
\(\Omega_x\) and the rational and level invariance.
Its average is zero by Lemma~\ref{dc-cuspidal} and
the finite-place argument following it.
Averaging the difference therefore bounds
\(\|{\bf F}(u,x)\|\) by the same right side.
The triangle inequality proves the assertion at
all \(y\in\Omega_x\).
\end{proof}

\begin{proof}[Proof of Theorem~\ref{dc-main}]
If \(D=0\), the finite-rank identity makes the kernel
zero, and there is nothing to prove.
Otherwise iterate Lemma~\ref{dc-iteration} \(P_0\)
times, starting with \eqref{dc-initial-evaluation}.
Since \(w\in V_0\) has norm one, the diagonal identity
gives
\[
 K(x,x)\le(2\rho\sqrt f\,\beta_*)^{16qN}
       X^{13N}\tau^{-4N}
                   \frac{D}{\operatorname{vol}(K_f(g))}.
\]
Moreover
\[
 2\rho\sqrt f\,\beta_*
   =\frac{n^2fgr_E}{2}\,2^{1/q}\beta_*\tau^{1/q}
                         \le X^5\tau^{1/q},
\]
because \(q\ge1\), \(n^2fgr_E\le X^2\),
and \(\beta_*\le X^3\).
It follows that
\[
 K(x,x)\le
 X^{(80q+13)N}\tau^{12N}
                   \frac{D}{\operatorname{vol}(K_f(g))}
 \le X^{(-52q+13)N}
                   \frac{D}{\operatorname{vol}(K_f(g))}.
\]
The diagonal is nonnegative.  Integrating this
uniform bound over the deep subset, bounding its
volume by the whole quotient volume, and using
Lemma~\ref{dc-volume}, gives
\[
 \int_{\lambda_x<\tau_0}K(x,x)\,dx
 \le X^{(-52q+16)N}D
 \le X^{-36N}D
 \le e^{-N}D.
\]
This is \eqref{dc-main-bound}.
\end{proof}

\begin{corollary}\label{dc-ambient-comparison}
Assume the field, lattice, finite test-function,
congruence, and discrete-series hypotheses in
Section~\ref{amb-section}.  Suppose \(C_g\ge0\),
\(b\ge2\), \(1\le q\), all the quantitative conditions
\eqref{amb-parameters} hold, and
\[
              \log n\ge20,\qquad n\ge4\log(C_g+3).
\]
Then the actual ambient multiplicity is finite and
\[
             |D-D_*|\le
                    \exp(-nf\log n/16)\,D_*.
\]
\end{corollary}
\begin{proof}
The compact-type estimate of Section~\ref{amb-section}
supplies the local decay required in
Section~\ref{ak-section}.  The lower bound for \(h\)
in \eqref{amb-parameters} in particular implies its
condition \(h/(8\sqrt{2q})>2n^2\).
That section therefore constructs the finite-rank
orthogonal projection, its absolutely convergent
continuous kernel, and its pointwise diagonal.
Theorem~\ref{nu-main} proves \eqref{amb-near-input}.
The rare-cusp estimate \eqref{rv-probability} implies
\eqref{amb-rare-input}, since
\(\tau_0<1/2\) and \(e^{-nf\log n/3}\le e^{-nf\log n/4}\).
Theorem~\ref{dc-main} proves \eqref{amb-deep-input},
and Theorem~\ref{amb-far} supplies the distant terms.
All hypotheses of Proposition~\ref{amb-final-reduction}
are now proved, so it yields the asserted comparison.
\end{proof}

This corollary concerns the complex ambient multiplicity.
It retains the explicitly stated classical imports of
the preceding sections.  It does not assert that
the remaining arithmetic and functorial steps of
the submitted GRH argument have been verified.


\section{Archimedean polynomial bounds and the global scale}
\label{as-section}

We now reconstruct the explicit polynomial and size estimates used
to choose the widths in Section~\ref{gt-section}.  The parameters
and polynomials below are the ones in the source, with ordinary
mathematical notation~\cite{alpo-uniform-source}.  The resulting boundary estimate improves
its numerical constant $160$ to $80$.  The mass recurrence, column
dimension formula, and formal-degree lower bound are stated
explicitly when they are used. Section~\ref{md-section} proves the
mass recurrence from the stated classical mass formula. The dimension
and formal-degree assertions are separate inputs. The coefficient
identities underlying the Fourier and off-diagonal estimates are
proved in Sections~\ref{ca-section} and \ref{cg-section}; their
application to the generated Jacquet functions has additional
analytic requirements.

\subsection{Vandermonde determinants under binary perturbations}

For a strictly decreasing real tuple $z=(z_1,\ldots,z_J)$, put
\[
 \mathcal V(z)=\prod_{1\leq i<j\leq J}(z_i-z_j),
 \qquad H_n=\sum_{j=1}^n\frac1j,\qquad H_0=0.
\]
Empty products have value one.

\begin{lemma}\label{as-vandermonde}
Let $J\geq1$ and $\epsilon_1,\ldots,\epsilon_J\in\{0,1\}$.
Then
\[
 R_J(\epsilon)=\prod_{i<j}
       \left(1+\frac{\epsilon_j-\epsilon_i}{4(j-i)}\right)>0.
\]
If
\[
 L_J(\epsilon)=\frac14\sum_{i=1}^J
               \epsilon_i(H_{i-1}-H_{J-i}),
\]
then
\begin{equation}\label{as-log-vand}
 L_J(\epsilon)-\frac{\pi^2J}{144}
       \leq\log R_J(\epsilon)\leq L_J(\epsilon),
 \qquad |L_J(\epsilon)|\leq\frac{J\log2}{4}.
\end{equation}
In particular $R_J(\epsilon)\leq2^{J/4}\leq e^{J/4}$.
\end{lemma}

\begin{proof}
Each fraction in the product has absolute value at most $1/4$.
For $|x|\leq1/4$, integration of
$x-\log(1+x)=\int_0^x t/(1+t)\,dt$ gives
\[
 0\leq x-\log(1+x)\leq\frac{x^2}{2(1-|x|)}
                                      \leq\frac23x^2.
\]
The sum of the fractions is $L_J(\epsilon)$: the coefficient
of $\epsilon_i$ is $(H_{i-1}-H_{J-i})/4$.
Also
\[
 \frac23\sum_{i<j}
       \left(\frac{\epsilon_j-\epsilon_i}{4(j-i)}\right)^2
 \leq\frac1{24}\sum_{d=1}^{J-1}\frac{J-d}{d^2}
 \leq\frac{\pi^2J}{144},
\]
using $\sum_{d\geq1}d^{-2}=\pi^2/6$.
This proves the first two inequalities in \eqref{as-log-vand}.

The coefficients $H_{i-1}-H_{J-i}$ are increasing and
antisymmetric about the middle index.  Thus the largest possible
absolute sum against binary $\epsilon_i$ is the sum of their
positive terms.  If $J=2p$, this sum is
$2p(H_{2p}-H_p)$; if $J=2p+1\geq3$, it is
$(2p+1)(H_{2p}-H_p)$.  These identities follow by summing
$\sum_{j=0}^nH_j=(n+1)H_n-n$.
Finally
$H_{2p}-H_p\leq\int_p^{2p}dt/t=\log2$.
For $J=1$ the sum is zero.  This proves the last inequality.
\end{proof}

\subsection{The spectral parameters and the critical scale}

Fix integers $k\geq4$ and $m\geq2$, a real $\Delta\geq1$,
and $v_*\geq0$.  Set
\[
 v_i=v_*+4\Delta(k-1-i)\quad(1\leq i\leq k-1).
\]
Suppose that real numbers $b_1,\ldots,b_m,a$ satisfy
\begin{equation}\label{as-parameters}
 v_1\leq4k\Delta,\quad b_1\geq1,\quad
 b_{j+1}-b_j\geq1,\quad a-v_1\geq4k\Delta.
\end{equation}
There are two cases.  In the first put $u_k=v_*-1/2$.
In the second choose $b_*$ with
$b_1/2\leq b_*\leq b_1-3$ and put $u_k=-b_*$.
For a binary sequence $\epsilon=(\epsilon_1,\ldots,\epsilon_{k-1})$
with $\epsilon_{k-1}=0$, complete the tuple by
\[
 u_1=a,\qquad
 u_{i+1}=v_i-\tfrac12-\epsilon_i\Delta
                      \quad(1\leq i\leq k-2).
\]
Write $u(\epsilon)$ for this tuple, and put
\[
 \lambda=(v_1,\ldots,v_{k-1},-b_1,\ldots,-b_m),
 \qquad n=k+m-1,\qquad B_0=\prod_{j=1}^m b_j.
\]
Define
\begin{equation}\label{as-effective-a}
 a_0=
 \begin{cases}a-v_1&\text{in the first case},\\
 a b_* /(k\Delta)&\text{in the second case},\end{cases}
 \qquad
 \eta_*=
 \begin{cases}0&\text{in the first case},\\
 4k\Delta/b_*&\text{in the second case}.
 \end{cases}
\end{equation}
All these tuples are strictly decreasing.  For the interior entries,
successive differences are $\Delta(4+\epsilon_{i+1}-\epsilon_i)$,
and the two end inequalities follow from \eqref{as-parameters}.
The last nonnegative entry of $\lambda$ is $v_*\geq0$, whereas
its remaining entries are negative.  These observations also show
that every Vandermonde determinant below is positive.

\begin{lemma}\label{as-vand-comparison}
With this notation, $a_0\geq1$ and
\begin{align}
 \mathcal V(\lambda)&\geq
       \mathcal V(v_1,\ldots,v_{k-1})B_0^{k-1},
                    \label{as-vand-lambda}\\
 \mathcal V(u(\epsilon))&\leq
       e^{k/4}\bigl(2a_0(1+\eta_*)\bigr)^{k-1}
                    \mathcal V(v_1,\ldots,v_{k-1}).
                    \label{as-vand-u}
\end{align}
\end{lemma}

\begin{proof}
Every mixed factor $v_i+b_j$ is at least $b_j$, and every
factor $b_j-b_i$ for $i<j$ is at least one.
Factoring the Vandermonde determinant proves
\eqref{as-vand-lambda}.
In the first case the last $k-1$ entries of $u$ are the
binary perturbations of the arithmetic progression $v_i$.
Their Vandermonde determinant is exactly
$R_{k-1}(\epsilon)\mathcal V(v_1,\ldots,v_{k-1})$.
Furthermore each $a-u_i$, $2\leq i\leq k$, is at most
\[
 a-v_*+\tfrac12
   =(a-v_1)+4\Delta(k-2)+\tfrac12
   \leq2(a-v_1).
\]
Lemma~\ref{as-vandermonde} proves \eqref{as-vand-u}.

In the second case put $q=k-2\geq2$.
The interior $q$ entries have Vandermonde determinant at most
$e^{q/4}\mathcal V(v_1,\ldots,v_q)$.
Each such entry is positive and at most $v_1$, whence
\[
 a-u_i\leq a,\qquad u_i+b_*\leq b_*(1+\eta_*).
\]
Also
\[
 \frac{a+b_*}{a_0}=\frac{k\Delta}{a}+\frac{k\Delta}{b_*}
                  \leq\frac{1+\eta_*}{4}.
\]
The missing factors in the progression are exactly
\[
 \mathcal V(v_1,\ldots,v_{k-1})
       =\mathcal V(v_1,\ldots,v_q)(4\Delta)^q q!.
\]
Combining these facts gives the stronger estimate
\[
 \frac{\mathcal V(u(\epsilon))}
      {\mathcal V(v_1,\ldots,v_{k-1})}
 \leq\frac{e^{q/4}}4
       \bigl(a_0(1+\eta_*)\bigr)^{k-1}
                  \frac{k^q}{4^q q!}.
\]
Here $4^q q!\geq(q+2)^q=k^q$ for $q\geq2$.
For $2\leq q\leq5$ this is checked by multiplication.
For $q\geq6$, the integral bound
$\log(q!)\geq\int_1^q\log t\,dt\geq q\log q-q$
and $e<3$ give $4^q q!\geq(4q/3)^q\geq(q+2)^q$.
This proves \eqref{as-vand-u} also in the second case.
Finally $a_0\geq4k\Delta$ in the first case, and
$a_0\geq4b_*\geq2$ in the second.
\end{proof}

Let $\rho\geq1$, $g\geq2\pi$, and $\kappa=2\pi/\rho$.
Let $V_\lambda>0$ satisfy the explicit formal-degree inequality
\begin{equation}\label{as-formal-degree-assumption}
 V_\lambda\geq\tfrac12\mathcal V(\lambda).
\end{equation}
In the application $V_\lambda$ is the formal degree of the
specified archimedean representation, multiplied by
$\prod_{j=0}^{n-1}j!$, for the Haar measure in the mass formula.
Corollary~\ref{fd-source} proves the stronger equality
$V_\lambda=\mathcal V(\lambda)$ with precisely that normalization,
using the classical representation theorems stated there.
For each binary sequence define its positive critical scale
\begin{equation}\label{as-critical-scale}
 T_\epsilon^k=
 \kappa^{-((m-1)k+m(m-1)/2)}\rho^{-(m-1)/2}
 g^{mk+m(m-1)/2}
       \frac{V_\lambda}{\mathcal V(u(\epsilon))}.
\end{equation}

Suppose in addition that a real number $B\geq1$ satisfies
\begin{equation}\label{as-small-parameters}
 \eta=\frac{4mk\Delta}{b_1}\leq\frac1{10},\qquad
 b_j\leq B^2\ (1\leq j\leq m),\qquad
 \mu=\frac{2m\log B}{k}\leq\frac1{10}.
\end{equation}
In the second case $b_*\geq b_1/2$ and $m\geq2$ imply
$0\leq\eta_*\leq\eta\leq1/10$.

\begin{proposition}\label{as-critical-lower}
For every allowed sequence,
\begin{equation}\label{as-critical-lower-eq}
 T_\epsilon\geq
 c_0 g^m\kappa^{-(m-1)}\frac{B_0}{a_0},
 \qquad
 c_0=\frac{19}{50}\frac{e^{-\mu}}{1+\eta_*}.
\end{equation}
\end{proposition}

\begin{proof}
The factor
$((g/\kappa)^{m(m-1)/2}\rho^{-(m-1)/2})^{1/k}$
is at least one, because $g/\kappa\geq\rho\geq1$.
Equations~\eqref{as-formal-degree-assumption} and
\eqref{as-vand-u} therefore give
\[
 T_\epsilon\geq g^m\kappa^{-(m-1)}
       \frac{e^{-1/4}}2
       \left(\frac{B_0}{a_0(1+\eta_*)}\right)^{1-1/k}.
\]
We have $B_0\leq B^{2m}$ and $a_0(1+\eta_*)\geq1$.
Thus the last power is at least
$e^{-\mu}B_0/(a_0(1+\eta_*))$.
Finally $e^{-1/4}/2>19/50$.
For example the exponential series gives
$e^{1/4}\leq1+1/4+(1/4)^2/(2(1-1/4))=31/24$,
and $12/31>19/50$.
\end{proof}

\subsection{The explicit boundary polynomial estimate}

For two allowed sequences $\epsilon,\epsilon'$ and positive
real numbers $T,T'$, put
\[
 P(Y)=\prod_{j=1}^n(Y+\lambda_j),\qquad
 Q(Y)=-TT'\prod_{i=1}^k
          (Y+\tfrac12+u_i(\epsilon))
          (Y+\tfrac12+u_i(\epsilon')),
\]
and, whenever the denominator is nonzero, define
\[
 \beta(Y)=g^m\kappa^{-(m-1)}
             \sqrt{\left|\frac{P(Y)P(Y+1)}{Q(Y)}\right|},
 \qquad t=\sqrt{\frac{T_\epsilon T_{\epsilon'}}{TT'}}.
\]

\begin{theorem}\label{as-boundary}
Suppose that $1\leq l\leq k-2$,
$\epsilon_l\ne\epsilon'_l$, and
$\tau=v_l+Y\in[5/2,\Delta/2]$.
Then $P(Y)$, $P(Y+1)$, and $Q(Y)$ are all nonzero, and
\begin{equation}\label{as-boundary-eq}
 0<\beta(Y)\leq80\,t\sqrt{k}\sqrt{\tau/\Delta}.
\end{equation}
\end{theorem}

\begin{proof}
Put $y=-Y=v_l-\tau$.  Then $0<y\leq v_1\leq4k\Delta$;
more precisely $y\geq v_*+7\Delta/2$.
Grouping corresponding factors of $P$ and $Q$ gives
\begin{equation}\label{as-beta-factorization}
 \beta(Y)=\frac{g^m\kappa^{-(m-1)}}{\sqrt{TT'}}
       N(y)U(y)\prod_{i=1}^{k-1}R_i(y),
\end{equation}
where
\[
 N(y)=\prod_{j=1}^m\sqrt{(b_j+y)(b_j+y-1)},\qquad
 U(y)=\frac1{a+1/2-y},
\]
and, except for the last index in the second case,
\[
 R_i(y)=
 \frac{\sqrt{|(v_i-y)(v_i+1-y)|}}
      {\sqrt{|(v_i-\epsilon_i\Delta-y)
                    (v_i-\epsilon'_i\Delta-y)|}}.
\]
At the last index in the second case replace this by
\[
 R_{k-1}(y)=
       \frac{\sqrt{|(v_*-y)(v_*+1-y)|}}{b_*+y-1/2}.
\]
All the displayed denominators are positive.  At $i=l$ the
two factors in the denominator have absolute values $\tau$
and $\Delta-\tau$.  At $i<l$ they are positive, and at
$i>l$ they are negative before absolute values are taken.
This proves the nonvanishing assertions and justifies the
factorization.

Since $1+x\leq e^x$,
\[
 N(y)\leq\prod_j(b_j+y)
       \leq B_0\exp\left(y\sum_jb_j^{-1}\right)
       \leq B_0e^\eta.
\]
For $i<l$, put $d=l-i$ and $X=4d\Delta+\tau$.
Then
\[
 R_i(y)\leq\frac{X+1}{X-\Delta}
       \leq1+\frac2{4d-1}.
\]
For $i>l$ outside the exceptional last index,
$X=v_i-y\leq-7\Delta/2$, so the numerator is at most
$-X$ and each denominator factor has absolute value at
least $-X$.  Consequently $R_i(y)\leq1$.
At the differing index the exact value and bound are
\[
 R_l(y)=\sqrt{\frac{\tau+1}{\Delta-\tau}}
              \leq\sqrt{\frac{14\tau}{5\Delta}}.
\]

The finite product before this index has the useful sharper bound
\begin{equation}\label{as-product-sqrt}
 \prod_{d=1}^p\left(1+\frac2{4d-1}\right)
                 \leq\frac53\sqrt p\quad(p\geq1).
\end{equation}
It is equality at $p=1$.  For the induction step, the inequality
\[
 \left(\frac{4p+5}{4p+3}\right)^2
                    \leq\frac{p+1}{p}
\]
reduces to $8p+9\geq0$.
Thus the product over $i<l$ is at most $(5/3)\sqrt k$;
this also holds for the empty product when $l=1$.

In the first case $U(y)\leq1/a_0$ and the last factor is
at most one.  In the second case
\[
 R_{k-1}(y)\leq\frac{4k\Delta}{b_*},\qquad
 U(y)\leq\frac1{a-v_1}\leq\frac2a,
\]
so their product is at most $8/a_0$.
Combining all these estimates in
\eqref{as-beta-factorization}, and applying
Proposition~\ref{as-critical-lower} to both sequences, gives
\[
 \beta(Y)\leq
 \frac{8(5/3)\sqrt{14/5}\,e^\eta}{c_0}
               t\sqrt k\sqrt{\tau/\Delta}.
\]
Finally $e^{-\mu}\geq1-\mu\geq9/10$,
$1+\eta_*\leq11/10$, and
$e^\eta\leq(1-\eta)^{-1}\leq10/9$.
Since $\sqrt{14/5}<8367/5000$, the remaining constant is at most
\[
 8\frac53\frac{8367}{5000}\frac{10}{9}
       \frac{11/10}{(19/50)(9/10)}
       =\frac{122716}{1539}<80.
\]
\end{proof}

\subsection{The exact mass assumptions and their algebraic consequence}

We next specify the global input that controls the product of the
numbers $t$ in Theorem~\ref{as-boundary}.  Fix an integer
$r\geq0$, put $m=2r+2$ and $n=k+m-1$, and let $f\geq1$
be the number of complex places.
At each place use spectral data as above, with the same
$k,m,\Delta,g,\rho$, but possibly different $a,b_j,v_*$ and
choice of the last entry.  The tuple $\lambda$ is fixed when
the binary sequence varies.  Write $V_{\lambda,v}$ and
$\mathcal V(u_v(\epsilon))$ for the corresponding quantities.
Put $\kappa=2\pi/\rho$ and
$s_d=\prod_{j=0}^{d-1}j!$.

Let $g=g_1g_2$, with integers $g_1,g_2\geq2$.
For $q>1$ define
\[
 G_d(q)=q^{d^2}\prod_{j=1}^d(1-q^{-j}),\qquad
 c_k(q)=\frac{G_k(q)}{q^{k(k-1)/2}},\qquad
 r_{k,r}(q)=\frac{G_n(q)}
                     {q^{r(2k+2r+1)+k(k-1)/2}}.
\]
At a prime power $q$, $G_d(q)$ is the order of
$\operatorname{GL}_d(\mathbb F_q)$, obtained by choosing
successive linearly independent columns.  The other two
expressions specify the exact finite-level factors used here.

Assume that positive masses $M_j$ and positive factors $L_j$
are given in the indicated ranges, with
\begin{align}
 M_j&=\bigl((j-1)!\kappa^{-j}\bigr)^f
              \rho^{-f/2}L_jM_{j-1}\quad(k<j\leq n),
                  \label{as-mass-assumption}\\
 |\log L_j|&\leq f\log\zeta(j),
                  \label{as-L-assumption}\\
 D_*&=M_n\bigl(r_{k,r}(g_1)r_{k,r}(g_2)\bigr)^f
                        \prod_v\frac{V_{\lambda,v}}{s_n},
                  \label{as-Dstar-assumption}\\
 d_{\epsilon,T}&=M_k\bigl(c_k(g_1)c_k(g_2)\bigr)^f
            \prod_v\frac{\mathcal V(u_v(\epsilon))}{s_k}
                              (T^k-1)^f,\qquad T\geq2.
                  \label{as-column-assumption}
\end{align}
Here $\zeta(j)=\sum_{a\geq1}a^{-j}$ in its convergent range.
The $M_j$ stand for the masses with the specified Haar
normalization, and $d_{\epsilon,T}$ for the chosen source-space
dimension. The mass recurrence and the Euler-factor bound are proved
in Section~\ref{md-section}, the source-space dimension in
Theorem~\ref{cd-source-dimension}, and the formal-degree
identification in Corollary~\ref{fd-source}. The expression $D_*$
is the prescribed leading term. Its quantitative comparison with
the actual ambient multiplicity is a separate estimate; the
identities here do not assert that comparison.

\begin{proposition}\label{as-mass-ratio}
Under \eqref{as-mass-assumption}--\eqref{as-column-assumption},
\begin{equation}\label{as-mass-ratio-eq}
 \frac{D_*}{d_{\epsilon,T}}
   =E_T\prod_v\left(\frac{T_{\epsilon,v}}T\right)^k,
\end{equation}
where
\[
 E_T=\left(\prod_{j=k+1}^nL_j\right)
       \left(\prod_{j=k+1}^n
                    (1-g_1^{-j})(1-g_2^{-j})\right)^f
                         (1-T^{-k})^{-f}.
\]
Moreover $E_T>0$ and
\begin{equation}\label{as-E-bound}
 |\log E_T|\leq8f2^{-k}.
\end{equation}
\end{proposition}

\begin{proof}
Directly from the finite products,
\[
 \frac{r_{k,r}(q)}{c_k(q)}
   =q^{km+m(m-1)/2}\prod_{j=k+1}^n(1-q^{-j}).
\]
The exponent identity uses $m=2r+2$ and $n=k+2r+1$.
Multiplying the mass recurrence gives
\[
 \frac{M_n}{M_k}
   =\left(\kappa^{-((m-1)k+m(m-1)/2)}
             \rho^{-(m-1)/2}\prod_{j=k}^{n-1}j!\right)^f
                   \prod_{j=k+1}^nL_j.
\]
Since $s_n/s_k=\prod_{j=k}^{n-1}j!$, the factorials cancel
in $D_*/d_{\epsilon,T}$.  The remaining factors are exactly
\eqref{as-critical-scale} at each place and $E_T$.
This proves \eqref{as-mass-ratio-eq}.

For $j\geq3$, monotonicity of $x^{-j}$ gives
\[
 \zeta(j)-1\leq2^{-j}+\int_2^\infty x^{-j}\,dx
       =2^{-j}\left(1+\frac2{j-1}\right)\leq2^{1-j}.
\]
Hence $\log\zeta(j)\leq2^{1-j}$.
For $0\leq x\leq1/2$,
$0\leq-\log(1-x)\leq2x$, by integrating $(1-t)^{-1}$.
The logarithm of the first product in $E_T$ is therefore
bounded in absolute value by $2f2^{-k}$.
Each of the two finite-level products contributes at most
$2f2^{-k}$, and the last factor contributes at most
$2f2^{-k}$.  This proves \eqref{as-E-bound}.
\end{proof}

\begin{corollary}\label{as-product-critical}
Suppose $8f2^{-k}\leq\log2$, $C_0\geq1$, $\Gamma\geq1$, and
\begin{equation}\label{as-choice-dimension}
 \frac{D_*}{C_0\Gamma^2}\leq d_{\epsilon,T}
                      \leq\frac{C_0D_*}{\Gamma^2}.
\end{equation}
Then
\[
 (2C_0)^{-1/k}\Gamma^{2/k}
  \leq\prod_v\frac{T_{\epsilon,v}}T
  \leq(2C_0)^{1/k}\Gamma^{2/k}.
\]
In particular, if $\log\Gamma\leq k/8$, the upper bound
is at most $C_t=2C_0e^{1/4}$.
If the same assumptions hold for $(\epsilon',T')$, then
\begin{equation}\label{as-product-t}
 \prod_v\sqrt{\frac{T_{\epsilon,v}T_{\epsilon',v}}{TT'}}
                              \leq C_t.
\end{equation}
\end{corollary}

\begin{proof}
Equation~\eqref{as-E-bound} gives $1/2\leq E_T\leq2$.
Substitute \eqref{as-choice-dimension} in
\eqref{as-mass-ratio-eq} and take positive $k$th roots.
The bound on $\Gamma$ gives $\Gamma^{2/k}\leq e^{1/4}$,
and $(2C_0)^{1/k}\leq2C_0$.  Multiplying the two upper
bounds and taking a square root proves
\eqref{as-product-t}.
\end{proof}

\subsection{Actual widths and a stronger lower bound for their scale}

Assume the hypotheses of Theorem~\ref{as-boundary} at every
place, with one fixed differing index $l$, and suppose
\eqref{as-product-t} holds.  Let
\[
 N=10m,\qquad I=\Delta/2-5/2,
 \qquad N^3\leq\Delta/2,
 \qquad C\geq1,\qquad
 H=\frac{C\pi N}{I}\leq1.
\]
These assumptions imply $I>0$.
Let $B_v$ be the supremum of $\beta_v(Y)$ on
$5/2\leq v_l+Y\leq\Delta/2$, and let $B_{1,v}$ be its
supremum on $5/2\leq v_l+Y\leq N^3$.
The polynomials have no zeros in these denominators, so these
are finite positive suprema of continuous functions on nonempty
compact intervals.  Theorem~\ref{as-boundary} gives, with
$A_0=80$ and $t_v=\sqrt{T_{\epsilon,v}T_{\epsilon',v}/(TT')}$,
\begin{equation}\label{as-B-bounds}
 B_v\leq A_0t_v\sqrt k,
 \qquad B_{1,v}\leq A_0t_v\sqrt k\sqrt{N^3/\Delta}.
\end{equation}

Put $h=C\kappa N/(gI)$.  Since $\rho\geq2$ in the
application, $\kappa\leq\pi$; assume $\rho\geq2$ in the
rest of this subsection.  Then $0<h\leq H\leq1$.
For $b,s\geq1$ with $b+s\leq r+2$ set
\[
 A_{b,s,v}=C B_{b,s,v}h^{2(b-1)+s-1},
 \qquad
 B_{b,s,v}=
 \begin{cases}B_{1,v}&(b,s)=(1,1),\\B_v&\text{otherwise},\end{cases}
\]
and for $1\leq b\leq r+1$ put
\[
 A'_{b,v}=\sqrt{\frac{C B_v\kappa}{\sqrt{TT'}}}\,h^{b-1}.
\]
Define the positive number $\Theta$ by the finite maximum
\begin{equation}\label{as-scale-definition}
 \Theta^{-f}=
 \max\left(\left\{\prod_vA_{b,s,v}:b,s\geq1,
                                           b+s\leq r+2\right\}
       \cup\left\{\prod_vA'_{b,v}:1\leq b\leq r+1\right\}\right).
\end{equation}
This is precisely the scale used for the source's block widths;
there is no choice of a favorable subsequence of blocks.

\begin{proposition}\label{as-widths}
Write $U=C_t^{1/f}$.  Then
\begin{equation}\label{as-theta-lower}
 \Theta\geq\min\left\{
 \frac{\sqrt{\Delta/(kN^3)}}{UCA_0},\quad
 \frac{I}{\pi UC^2A_0\sqrt k\,N},\quad
 \sqrt{\frac{\sqrt{TT'}}{\pi UCA_0\sqrt k}}
                            \right\}.
\end{equation}
For $\nu(r,b,s)=2r+3-2b-s$, define the unnormalized widths
\[
 x_{b,s,v}=\frac{A_{b,s,v}}{g\rho^{\nu(r,b,s)+1}},
 \qquad
 x'_{b,v}=\frac{A'_{b,v}}{g\rho^{\nu(r,b,r+2-b)}}.
\]
Then the actual dual widths $x^D_{b,v}=x_{b,1,v}$ and lift widths
\begin{equation}\label{as-actual-lift-widths}
 x^L_{b,s,v}=
 \begin{cases}
 x_{b,s+1,v}&b+s\leq r+1,\\
 2x'_{b,v}&b+s=r+2
 \end{cases}
\end{equation}
satisfy \eqref{gt-dual-width}--\eqref{gt-lift-width}, with
$\rho=\rho_E$.  Also
\[
 \frac{(\prod_vx_{b,s,v})^2}{N\mathfrak a_{b,s}}
                                          \leq\Theta^{-2f}.
\]
\end{proposition}

\begin{proof}
Multiplying \eqref{as-B-bounds} over the places shows that
the $f$th roots of the products in
\eqref{as-scale-definition} are bounded by, respectively,
\[
 UCA_0\sqrt{kN^3/\Delta},\qquad
 UCA_0\sqrt k\,H^2,\qquad
 UCA_0\sqrt k\,H,\qquad
 \sqrt{\frac{\pi UCA_0\sqrt k}{\sqrt{TT'}}}.
\]
The four cases are $(b,s)=(1,1)$, $s=1,b\geq2$,
$s\geq2$, and the terminal widths $A'_b$.
Since $H\leq1$, the second bound is at most the third.
Taking the reciprocal of the maximum proves
\eqref{as-theta-lower}.

The remaining claims follow directly from
\eqref{as-scale-definition} and
\eqref{gt-adjacent-ideals}.  For the interior lift widths use
$\nu(r,b,s+1)+1=\nu(r,b,s)$.
For the terminal lift widths the factor $2$ contributes $2^f$
to the product, exactly as in \eqref{gt-lift-width}.
Finally the squared ratio in the statement is
$(\prod_vA_{b,s,v})^2$ by the ideal norm formula.
\end{proof}

To verify that these definitions are the original widths,
let $d=2(b-1)+s-1$, $\kappa_g=2\pi/g$, and put
$\mathcal A_v=B_{b,s,v}/(g^m\kappa^{-(m-1)})$.
Their scalar polynomial formula is
\[
 C\kappa_g^{-(m-1)}\mathcal A_v
               (C\kappa_gN/I)^d.
\]
Multiplication by $g\rho^{\nu(r,b,s)+1}$ gives
$CB_{b,s,v}h^d=A_{b,s,v}$, since
$\nu(r,b,s)+1=m-1-d$.
For the terminal width use instead
$\mathcal A_v=B_v/(g^m\kappa^{-(m-1)})$; its squared formula is
\[
 C\kappa_g^{-(m-1)}\mathcal A_v
       \kappa_g^{2b-1}(TT')^{-1/2}(CN/I)^{2(b-1)}.
\]
Multiplication by $(g\rho^{r+1-b})^2$ gives
$(A'_{b,v})^2$.  All quantities are nonnegative, so taking
square roots preserves the equality.  This verifies both
normalizations, including the terminal square root.

\begin{corollary}\label{as-power-scale}
Fix positive constants $C,C_t,C_\Delta,c_T,C_b$, with
$C\geq1$, $C_t\geq1$, $C_\Delta>4$, and $c_TC_b>1$.
Suppose $m=2r+2\leq f\leq k$ and
\begin{equation}\label{as-growing-parameters}
 B\geq k^{C_b/2},\qquad
 \Delta\geq(kf\log B)^{C_\Delta},\qquad
 T,T'\geq B^{c_Tm}.
\end{equation}
Under the preceding spectral, mass, dimension, and
formal-degree hypotheses, every fixed
\begin{equation}\label{as-improved-exponent}
 0<c<\min\left\{\frac{C_\Delta-1}{2},
                         \frac{2c_TC_b-1}{4}\right\}
\end{equation}
satisfies $\Theta\geq k^c$ for all sufficiently large $k$.
The threshold is uniform in $m,f,B,\Delta,T,T'$ and in the
individual spectral parameters subject to the displayed
assumptions.
\end{corollary}

\begin{proof}
The constants are fixed, while $2\leq m\leq f\leq k$.
Thus $N=10m\leq10f$ and
$\log B\geq(C_b/2)\log k$.
For sufficiently large $k$, uniformly in the other parameters,
the lower bound for $\Delta$ gives
$N^3\leq\Delta/2$, $I=\Delta/2-5/2\geq\Delta/3$,
and $C\pi N/I\leq1$.
Since $1\leq U=C_t^{1/f}\leq C_t$, the first term of
\eqref{as-theta-lower} is at least a positive fixed constant times
\[
 k^{(C_\Delta-1)/2}
       f^{(C_\Delta-3)/2}(\log B)^{C_\Delta/2},
\]
and the second is at least a positive fixed constant times
\[
 k^{C_\Delta-1/2}
       f^{C_\Delta-1}(\log B)^{C_\Delta}.
\]
The positive powers of $f$ can be bounded below by constants.
The third term is at least a positive fixed constant times
\[
 (TT')^{1/4}k^{-1/4}
    \geq B^{c_Tm/2}k^{-1/4}
    \geq k^{(2c_TC_b-1)/4}.
\]
Each exponent is strictly larger than $c$, so its fixed
multiplicative constant is absorbed by increasing the threshold
for $k$.  This proves the uniform claim.
\end{proof}

\begin{corollary}\label{as-final-smallness}
Let $C_1\geq1$, $c>0$, and $C_f>0$ be fixed, with $cC_f>7$.
If $\Theta\geq k^c$, $2r+2\leq f$, and
$C_f\log B\leq f\log k$, then, for sufficiently large $k$,
\[
 r(C_1/\Theta)^f\leq1,\qquad
 C_1(r+1)(C_1/\Theta)^f\leq B^{-7}.
\]
\end{corollary}

\begin{proof}
Since $r+1\leq f\leq2^f$, the prefactor
$C_1(r+1)C_1^f$ is at most $(2C_1^2)^f$.
Also $B^7\leq k^{7f/C_f}$.
Therefore the second asserted product, after multiplication
by $B^7$, is at most
\[
 \bigl(2C_1^2k^{-(c-7/C_f)}\bigr)^f,
\]
which is at most one for large $k$.
The first assertion follows in the same way from
$r\leq2^f$ and $k^c\geq2C_1$.
\end{proof}

The boundary constant $80$, the factor $C_t^{1/f}$ in
\eqref{as-theta-lower}, and the exponent range
\eqref{as-improved-exponent} strengthen the corresponding
size estimates stated in the source.  The displayed hypotheses
include the source-space dimension formula and formal-degree lower
bound; the mass recurrence is now supplied by Section~\ref{md-section}.
The polynomial coefficient identities are proved in
Sections~\ref{ca-section} and \ref{cg-section}. Their application to
the particular Jacquet functions, the remaining orbital energy
estimates, and the ambient-dimension assertion are not consequences
of the elementary size estimates above.


\section{Polynomial identities behind the archimedean coefficients}
\label{ca-section}

The next step concerns the differential operators that produce
the archimedean Fourier estimates.  We prove their leading
polynomial identity directly, including the two different orders
of the matrix indices, and then prove the coefficient extraction
and its positive matrix estimate.  We also bound the explicit
rational correction contributed by the preceding operators.
These are actual algebraic steps in the source's coefficient
argument~\cite{alpo-uniform-source}.  The Fourier integral
identity for the radial Hermitian form, and positivity of that
form, are identified separately at the end.

\subsection{The differential operators and their leading coefficients}

Let $k,r$ be nonnegative integers and $n=2r+k+1$.
Use the ordered basis
\[
 e_0,\ldots,e_{r-1},w_0,\ldots,w_k,f_{r-1},\ldots,f_0
\]
of $\mathbb C^n$.  Let $H$ be the Hermitian matrix specified by
\[
 He_a=f_a,\quad Hf_a=e_a,\quad Hw_b=w_b\ (b<k),
 \quad Hw_k=-w_k.
\]
Thus $H^*=H$ and $H^2=1$.  Put
\[
 G=\{g\in\operatorname{GL}_n(\mathbb C):g^*Hg=H\}.
\]
Its real Lie algebra consists of the matrices $X$ satisfying
$X^*H+HX=0$.

For a matrix $Z$ set
\[
 Z_1=\tfrac12(Z-HZ^*H),\qquad
 Z_2=\tfrac1{2i}(Z+HZ^*H).
\]
Both matrices belong to this real Lie algebra, and
$Z=Z_1+iZ_2$.  If $F$ is a smooth function on $G$ with
values in a complex vector space of finite dimension, define
\[
 (D_XF)(g)=\left.\frac{d}{dt}F(g\exp(tX))\right|_{t=0},
 \qquad E_{ab}=D_{(E^{ab})_1}+iD_{(E^{ab})_2},
\]
where $E^{ab}$ is the matrix unit in row $a$, column $b$.
The same definitions apply to smooth functions on an open subset
of $G$, by locality.  These are complex-linear differential
operators with smooth coefficients.  In particular, their
finite compositions preserve smoothness and do not enlarge
support.

For ordered injective lists $a=(a_0,\ldots,a_{m-1})$ and
$b=(b_0,\ldots,b_{m-1})$ of basis indices, define the column
determinant
\begin{equation}\label{ca-column}
 Q_{a,b}(u)=
 \sum_{\sigma\in S_m}\operatorname{sgn}(\sigma)
 \prod_{j=0}^{m-1}
 \bigl((u-j)\delta_{a_{\sigma(j)},b_j}\,\mathrm{Id}
                                  +E_{a_{\sigma(j)},b_j}\bigr).
\end{equation}
The product is in the displayed order; when it acts on a
function, the rightmost operator acts first.  No commutation
of differential operators is understood.
For $m=0$ the value is $\mathrm{Id}$.

\begin{lemma}\label{ca-column-leading}
The expression \eqref{ca-column} is a polynomial in $u$ of
degree at most $m$, with leading coefficient
\[
 [u^m]Q_{a,b}(u)=
 \det\bigl(\delta_{a_i,b_j}\bigr)_{0\leq i,j<m}\,\mathrm{Id}.
\]
This identity remains valid for arbitrary endomorphisms
$E_{ab}$ of any complex vector space; no Lie algebra
relations are needed.
\end{lemma}

\begin{proof}
The scalar $u$ commutes with every endomorphism.
On expanding each ordered product, a term of degree $m$
can only arise by choosing the scalar $u$ from every factor.
Its coefficient is
$\prod_j\delta_{a_{\sigma(j)},b_j}\,\mathrm{Id}$.
Summing over $\sigma$ gives the displayed determinant.
All other terms have degree at most $m-1$.
\end{proof}

For $0\leq m\leq r$ set
\[
 X_{a,c}(g)=(Hg)_{e_a,c},\qquad
 Y_{a,c}(g)=(g^{-1})_{c,e_a}\quad(0\leq a<m).
\]
There are two orders on the basis indices.  The first is the
displayed order used to define $H$.  The second begins with
$f_0,f_1,\ldots,f_{r-1}$ and is followed by
$e_0,\ldots,e_{r-1},w_0,\ldots,w_k$.
Write $c$ increasing with respect to the indicated order
in the following two sums:
\begin{align}
 (A_m(u)F)(g)
   &=\sum_{c\text{ increasing in the first order}}
          \det(X_{a,c_b}(g))_{a,b<m}\,
                 (Q_{(e_0,\ldots,e_{m-1}),c}(u)F)(g),
                    \label{ca-Aplus}\\
 (\widetilde A_m(u)F)(g)
   &=\sum_{c\text{ increasing in the second order}}
          \det(Y_{a,c_b}(g))_{a,b<m}\,
        (Q_{c,(f_0,\ldots,f_{m-1})}(u+m-1)F)(g).
                    \label{ca-Aminus}
\end{align}
Each determinant multiplies the result of the differential
operator; its entries are not differentiated in these formulas.
Put
\[
 d_m(g)=\det\bigl((Hg)_{e_a,e_b}\bigr)_{a,b<m},
 \qquad d_0(g)=1.
\]
These are the actual principal minors in the source at levels
$m\leq r$.  The additional vectors used for the final principal
minor do not occur at these levels.

\begin{proposition}\label{ca-A-leading}
For every smooth $F$ and every $g\in G$,
\[
 [u^m](A_m(u)F)(g)=d_m(g)F(g),\qquad
 [u^m](\widetilde A_m(u)F)(g)=\overline{d_m(g)}F(g).
\]
Both expressions have degree at most $m$.
\end{proposition}

\begin{proof}
In \eqref{ca-Aplus}, Lemma~\ref{ca-column-leading} gives a
nonzero leading coefficient only if the set of indices in
$c$ is $\{e_0,\ldots,e_{m-1}\}$.  The order condition then
forces $c=(e_0,\ldots,e_{m-1})$, with determinant of the
Kronecker matrix equal to one.  This proves the first formula.
The same argument in the second order forces
$c=(f_0,\ldots,f_{m-1})$.  The shift $u+m-1$ has no effect
on its leading coefficient.

Since $g^{-1}=Hg^*H$, one has
\[
 (g^{-1})_{f_b,e_a}
     =\overline{g_{f_a,e_b}}
     =\overline{(Hg)_{e_a,e_b}}.
\]
Taking determinants proves the second formula.  The empty
determinant case $m=0$ also gives both identities.
\end{proof}

Fix $0\leq j\leq r$ and work on the open set
$G_j=\{g\in G:d_j(g)\ne0\}$.
Let $c=(2r+k)/2$.  The two components of the operators
at a real parameter $y$ are
\begin{align}
 L^+_{j,e}(y)F&=d_j^{-1}A_{j-e}(y+c)F,
                    \label{ca-Lplus}\\
 L^-_{j,e}(y)F&=\overline{d_j}^{-1}
                       \widetilde A_{j-e}(y-c)F,
                    \label{ca-Lminus}
\end{align}
for $e=0,1$ and $e\leq j$.  Define
$L^\pm_{0,1}(y)F=0$.
These formulas include the special level-zero convention;
they never invoke an operator with a negative index.

\begin{corollary}\label{ca-normalized-leading}
As operators on smooth functions on $G_j$,
\[
 L^\pm_{j,0}(y)=y^j\mathrm{Id}
                       +\sum_{\ell<j}y^\ell B^\pm_\ell,\qquad
 L^\pm_{j,1}(y)=\sum_{\ell<j}y^\ell C^\pm_\ell,
\]
for differential operators with smooth coefficients independent
of $y$.  A translation $y\mapsto y+a$ by a fixed real number
does not change these leading coefficients.
At $j=0$, these formulas mean
$L^\pm_{0,0}=\mathrm{Id}$ and $L^\pm_{0,1}=0$.
\end{corollary}

\begin{proof}
Apply Proposition~\ref{ca-A-leading} with $m=j$ or $j-1$.
For $e=0$, multiplication by $d_j^{-1}$ or
$\overline{d_j}^{-1}$ cancels the leading coefficient
exactly.  For $e=1$ the degree is at most $j-1$.
The denominators are smooth and nonzero on $G_j$.
Translations preserve the leading coefficient of a polynomial.
\end{proof}

This proof also gives the needed regularity of every coefficient.
If a test function is smooth with compact support in $G_j$,
then each coefficient has support in the same compact set.
More generally, applying any fixed complex-linear restriction
or pullback map preserves the polynomial identity.  If a
Hermitian form is only defined on a subspace of functions,
all coefficient functions must belong to that subspace.
It suffices to check membership for $j+1$ distinct real
values of $y$: inversion of the scalar Vandermonde matrix
then expresses every coefficient as a finite linear combination
of those values.  This observation prevents any implicit
exchange of differentiation with an unproved integral.

\subsection{The leading coefficient in the mixed pairing}

In this subsection a sesquilinear form is linear in its first
argument and conjugate-linear in its second, which is the
convention of the source's coefficient functions.
Let $\mathcal B:V\times V'\to\mathbb C$ be such a form.
Suppose four vector polynomials in a real variable $y$ satisfy
\[
 X_0(y)=y^j x+\text{terms of degree less than }j,\quad
 Y_0(y)=y^j x'+\text{terms of degree less than }j,
\]
\[
 \deg X_1<j,\qquad \deg Y_1<j.
\]
For $j=0$ the last two polynomials are zero.
Let $U:V\to V$ be complex-linear and let $s\in\mathbb C$.
Then
\begin{equation}\label{ca-mixed-polynomial}
 F(y)=
 \mathcal B\bigl(U X_1(y)+sU X_0(y),Y_0(y)\bigr)
       -\mathcal B\bigl(U X_0(y),Y_1(y)\bigr)
\end{equation}
is a polynomial in the real variable $y$ of degree at most $2j$,
and
\begin{equation}\label{ca-mixed-leading}
 [y^{2j}]F(y)=s\,\mathcal B(Ux,x').
\end{equation}
Indeed, in the first term only
$\mathcal B(sy^j Ux,y^j x')$ contributes to that degree,
and the second term has degree at most $2j-1$.
Reality of $y$ is used when conjugating its powers.

For completeness the extra summand $sU X_0$ comes from an
explicit operation on the two-component indices, not from
an estimate.  Given a finite sequence of signs
$\varepsilon_\ell\in\{+,-\}$, regard its vectors as indexed by
$\{0,1\}^a$.  On the $\ell$th factor apply the matrix
\[
 \begin{pmatrix}1&0\\ s&1\end{pmatrix}
 \quad\text{if }\varepsilon_\ell=+,\qquad
 \begin{pmatrix}1&0\\ \overline s&1\end{pmatrix}
 \quad\text{if }\varepsilon_\ell=-.
\]
The matrices on distinct factors commute.  One may additionally
multiply every component by the same function
$\exp(sz-\overline s\,\overline z)$.
The latter has absolute value one, but its preservation of
a particular differential Hermitian form is an additional
identity, not a consequence of its pointwise absolute value.

On adding a new positive component, the transformed components
are precisely $UX_0$ and $UX_1+sUX_0$.
Contracting the new positive and negative components by
\[
 \mathcal B_1((x_0,x_1),(y_0,y_1))
           =\mathcal B(x_1,y_0)-\mathcal B(x_0,y_1)
\]
therefore gives \eqref{ca-mixed-polynomial}.
Repeated contraction gives a sesquilinear form on any number
of new components.
Applying Corollary~\ref{ca-normalized-leading} to a new pair
$L^+_{j,e}(y)$ and $L^-_{j,e}(y+1)$ proves
\eqref{ca-mixed-leading} for the actual normalized differential
operators.  The shift by one on the negative side changes no
leading coefficient.  No positivity is required for this
algebraic assertion.

\subsection{Exact extraction by interpolation}

For distinct real nodes $y_0,\ldots,y_d$, put
\[
 c_i=\prod_{\substack{0\leq h\leq d\\h\ne i}}
                    (y_i-y_h)^{-1},
 \qquad C=\sum_{i=0}^d|c_i|.
\]
Each $c_i$ is nonzero and $C>0$.
Lagrange interpolation of a vector polynomial $P$ of degree
at most $d$ gives
\begin{equation}\label{ca-Lagrange-leading}
 [y^d]P(y)=\sum_{i=0}^d c_i P(y_i).
\end{equation}
To prove this for any complex vector space, subtract the
usual interpolation expression from $P$.  Its coefficients
lie in a finite-dimensional subspace, so the scalar
root-counting argument applied to coordinates proves that
the difference vanishes.  Comparing coefficients of $y^d$
then gives \eqref{ca-Lagrange-leading}.
For $d=0$, it simply says that a constant polynomial equals
its value at its single node.

\begin{proposition}\label{ca-iterated-extraction}
Let $K,d$ be nonnegative integers.  Suppose that quantities
$W_a(y_1,\ldots,y_a)$, $0\leq a\leq K$, take values in a
complex vector space, with $W_0$ independent of the variables.
Assume that, for each fixed $y_1,\ldots,y_{a-1}$,
$W_a$ is a polynomial in its last real variable of degree
at most $d$, with leading coefficient $s_aW_{a-1}$.
For each $a$ choose $d+1$ distinct real nodes $y_{a,i}$ and
let $c_{a,i}$ be their coefficients above.  Then
\begin{equation}\label{ca-extraction}
 \left(\prod_{a=1}^K s_a\right)W_0
  =\sum_{i_1,\ldots,i_K}
       \left(\prod_{a=1}^K c_{a,i_a}\right)
           W_K(y_{1,i_1},\ldots,y_{K,i_K}).
\end{equation}
No assumption about simultaneous polynomial degree in all
variables is needed.
\end{proposition}

\begin{proof}
Extract the coefficient in the last variable using
\eqref{ca-Lagrange-leading}, then extract successively in
the preceding variables.  At each step the extracted
coefficient is the asserted scalar multiple of $W_{a-1}$,
independent of the variable just removed.
The finite sums commute with scalar multiplication.
The case $K=0$ has one empty tuple and is the identity.
\end{proof}

In the mixed pairing above one takes $d=2j$ and every
$s_a=s$.  Newly inserted operators may act on the previous
ones; they are not commuted past them.
Their leading term is the identity, so the same one-variable
argument applies to each successive insertion.
Thus \eqref{ca-extraction} proves the source's algebraic
degree and extraction identity on the common domain of its
forms and restrictions.

The natural probability weights are
\[
 \theta_{i_1,\ldots,i_K}
   =\frac{\prod_a|c_{a,i_a}|}{\prod_a C_a},
 \qquad C_a=\sum_i|c_{a,i}|.
\]
Their sum is one, because a product of finite sums expands
as a sum over all tuples.  Consequently, for vectors in a
normed space,
\begin{align}
 \left(\prod_a|s_a|\right)\|W_0\|
   &\leq\left(\prod_aC_a\right)
         \sum_{\boldsymbol i}\theta_{\boldsymbol i}
                    \|W_K(\boldsymbol y_{\boldsymbol i})\|,
                    \label{ca-extraction-norm}\\
 \left(\prod_a|s_a|\right)^2\|W_0\|^2
   &\leq\left(\prod_aC_a\right)^2
         \sum_{\boldsymbol i}\theta_{\boldsymbol i}
                    \|W_K(\boldsymbol y_{\boldsymbol i})\|^2.
                    \label{ca-extraction-square}
\end{align}
The first is the triangle inequality in
\eqref{ca-extraction}; the second follows from
Cauchy--Schwarz for the probability weights.

\begin{lemma}\label{ca-Chebyshev}
If $d\geq1$ and $\alpha<\beta$, the nodes
\[
 y_i=\frac{\alpha+\beta}{2}
       +\frac{\beta-\alpha}{2}\cos\frac{i\pi}{d},
       \qquad 0\leq i\leq d,
\]
satisfy the exact identity
\[
 \sum_i|c_i|=\frac12\left(\frac4{\beta-\alpha}\right)^d.
\]
This is the smallest possible coefficient sum among all
choices of $d+1$ distinct nodes in $[\alpha,\beta]$.
\end{lemma}

\begin{proof}
The polynomial $T_d$ determined by
$T_d(\cos\theta)=\cos(d\theta)$ has degree $d$, leading
coefficient $2^{d-1}$, and absolute value at most one on
$[-1,1]$.  These facts follow from
$T_0=1$, $T_1=x$, and $T_{d+1}=2xT_d-T_{d-1}$.
The scaled polynomial
\[
 R(y)=T_d\left(\frac{2y-\alpha-\beta}{\beta-\alpha}\right)
\]
has leading coefficient
$2^{2d-1}(\beta-\alpha)^{-d}$ and value $(-1)^i$ at $y_i$.
The nodes decrease with $i$, so $c_i$ also has sign $(-1)^i$.
Equation~\eqref{ca-Lagrange-leading} gives the equality.
For any other nodes in the interval, applying the same
identity to $R$ and using $|R|\leq1$ gives the asserted
lower bound.
\end{proof}

The actual construction needs its nodes separated from the
parameters already used.  The exact Chebyshev identity alone
does not supply those simultaneous separation conditions.
The source uses grids whose coefficient sums are bounded
by $(28N/I)^{2j}$.  Section~\ref{cg-section} supplies the
complete construction and improves $28$ to $24$.
For those actual grids,
\eqref{ca-extraction-norm}--\eqref{ca-extraction-square}
apply with that bound, with no loss from the number of grid
tuples.

\subsection{The explicit rational correction}

Let $q_0$ be a real-valued function, and let
$w=((\varepsilon_1,a_1),\ldots,(\varepsilon_L,a_L))$
be a finite sequence, with $\varepsilon_\ell\in\{+,-\}$
and $a_\ell\in\mathbb R$.  The modified function occurring
in the source is exactly
\begin{equation}\label{ca-qword}
 q_w(y)=q_0(y)
   \prod_{\varepsilon_\ell=+}
           \frac{y-a_\ell+1}{y-a_\ell}
   \prod_{\varepsilon_\ell=-}
           \frac{a_\ell-y}{a_\ell-y-1}.
\end{equation}
Only the displayed evaluation at $y$ is needed below.

\begin{lemma}\label{ca-q-distortion}
Suppose $D>1$, $q_0(y)\ne0$, and
\[
 |y-a_\ell|\geq D\quad(\varepsilon_\ell=+),\qquad
 |a_\ell-y-1|\geq D\quad(\varepsilon_\ell=-).
\]
Then $q_w(y)/q_0(y)>0$ and
\begin{equation}\label{ca-q-distortion-eq}
 \left|\log\frac{q_w(y)}{q_0(y)}\right|
                      \leq\frac{L}{D-1}.
\end{equation}
In particular, if $N\geq20$, $L\leq N$ and $D\geq N^2/8$,
then
\[
 e^{-20/49}\leq\frac{q_w(y)}{q_0(y)}\leq e^{20/49}.
\]
If also $q_0(y)>0$ and
$\sqrt{|P(y)P(y+1)|/q_0(y)}\leq A$, then
\begin{equation}\label{ca-node-majorant}
 q_w(y)>0,\qquad
 \sqrt{\frac{|P(y)P(y+1)|}{q_w(y)}}
                                      \leq e^{10/49}A.
\end{equation}
\end{lemma}

\begin{proof}
Every factor in \eqref{ca-qword} is $1+u$ with
$|u|\leq1/D<1$.  It is positive, and
\[
 |\log(1+u)|
  =\left|\int_0^u\frac{dt}{1+t}\right|
  \leq\frac{|u|}{1-|u|}\leq\frac1{D-1}.
\]
Adding these inequalities proves
\eqref{ca-q-distortion-eq}.
Moreover
\[
 \frac{L}{D-1}\leq\frac{8N}{N^2-8}
                          \leq\frac{20}{49}\quad(N\geq20).
\]
The last function decreases for $N>\sqrt8$; at $N=20$
its value is $20/49$.
Taking exponentials gives the first uniform assertion,
and taking the reciprocal square root gives
\eqref{ca-node-majorant}.
\end{proof}

In the actual construction $N=20(r+1)$.
The separation conditions proved for its new and previous
parameters give precisely the denominator bounds in this
lemma, with $D=N^2/8$ in the exceptional short interval
and a larger value in the other intervals.
The polynomial $q_0$ is $Q$ from Section~\ref{as-section}.
On an interval belonging to a differing binary position it
is positive: exactly one of the paired factors in $Q$ has
opposite signs, and its defining initial minus sign reverses
that sign.  Thus \eqref{ca-node-majorant} applies there.
The source uses the weaker factor $eA$.
This improvement concerns the rational modification alone;
it does not assert that the separated grids or positive
Hermitian forms exist for arbitrary parameters.

\subsection{The exact positive matrix estimate}

We next give the finite-dimensional step that turns the
interpolation identity into the normalized local coefficient
estimate.  All its hypotheses are stated for arbitrary
linear combinations, as required by positivity of a block
matrix.

Fix $K,d\geq0$ and the nodes of
Proposition~\ref{ca-iterated-extraction}.
For a tuple $\boldsymbol i=(i_1,\ldots,i_K)$ put
\[
 c_{\boldsymbol i}=\prod_a c_{a,i_a},\quad
 C_*=\prod_a\sum_i|c_{a,i}|,\quad
 \theta_{\boldsymbol i}=|c_{\boldsymbol i}|/C_*.
\]
Let $I_+$ and $I_-$ be finite sets.  For each tuple suppose
$R^+_{\boldsymbol i}$ and $R^-_{\boldsymbol i}$ are
positive semidefinite Hermitian forms on
$\mathbb C^{I_+}$ and $\mathbb C^{I_-}$, and that
$B_{\boldsymbol i}$ is a sesquilinear form between these
two spaces.  Let $P$ and $q$ be real-valued at the nodes,
with $P(y_{a,i})P(y_{a,i}+1)\ne0$ and $q(y_{a,i})>0$.
Assume
\begin{equation}\label{ca-pairing-assumption}
 |B_{\boldsymbol i}(x,y)|^2
       \prod_{a=1}^Kq(y_{a,i_a})
       \leq R^+_{\boldsymbol i}(x,x)
                 R^-_{\boldsymbol i}(y,y)
 \quad\text{for every }x,y.
\end{equation}
Fix positive $\Gamma,S$ and nonnegative $T,A$ such that
\[
 \frac{|P(y_{a,i})P(y_{a,i}+1)|}{q(y_{a,i})}
                                                  \leq A^2.
\]
Define normalized forms
\[
 G^+_{\boldsymbol i}
       =\frac{R^+_{\boldsymbol i}}
                 {\Gamma S\prod_a|P(y_{a,i_a})|},\qquad
 G^-_{\boldsymbol i}
       =\frac{R^-_{\boldsymbol i}}
                 {\Gamma S\prod_a|P(y_{a,i_a}+1)|},
\]
and the mixed form
\[
 F=\frac{T^K}{\Gamma S}
            \sum_{\boldsymbol i}c_{\boldsymbol i}B_{\boldsymbol i}.
\]
Write $\overline G^\pm=\sum_{\boldsymbol i}
\theta_{\boldsymbol i}G^\pm_{\boldsymbol i}$.

\begin{proposition}\label{ca-positive-block}
For all $x\in\mathbb C^{I_+}$ and $y\in\mathbb C^{I_-}$,
\begin{equation}\label{ca-positive-bound}
 |F(x,y)|\leq C_*T^KA^K
       \sqrt{\overline G^+(x,x)\,\overline G^-(y,y)}.
\end{equation}
Consequently, if $X\geq0$ and
$C_*T^KA^K\leq X^K$, the block Hermitian form with
diagonal forms $X^K\overline G^+$, $X^K\overline G^-$
and mixed form $F$ is positive semidefinite.
\end{proposition}

\begin{proof}
Divide \eqref{ca-pairing-assumption} by $(\Gamma S)^2$.
The normalizations give exactly
\[
 \left|\frac{B_{\boldsymbol i}(x,y)}{\Gamma S}\right|
 \leq
 \left(\prod_a
       \frac{|P(y_{a,i_a})P(y_{a,i_a}+1)|}{q(y_{a,i_a})}
       \right)^{1/2}
       \sqrt{G^+_{\boldsymbol i}(x,x)
                         G^-_{\boldsymbol i}(y,y)}.
\]
In particular the factor on the right is at most $A^K$.
The triangle inequality and
$|c_{\boldsymbol i}|=C_*\theta_{\boldsymbol i}$ yield
\[
 |F(x,y)|\leq C_*T^KA^K
     \sum_{\boldsymbol i}\theta_{\boldsymbol i}
          \sqrt{G^+_{\boldsymbol i}(x,x)
                            G^-_{\boldsymbol i}(y,y)}.
\]
Cauchy--Schwarz for the probability weights proves
\eqref{ca-positive-bound}.
If $a=\overline G^+(x,x)$ and
$b=\overline G^-(y,y)$, the quadratic value of the block
form is at least
$X^K(a+b)-2|F(x,y)|\geq X^K(\sqrt a-\sqrt b)^2$.
This is nonnegative, including when one of the diagonal
values is zero.  For $K=0$ all empty products are one and
the same proof applies.
\end{proof}

In the source the original size $S$ is the product
$\prod_{\ell\in w}|P(a_\ell)|$.
Adding positive parameters $y_1,\ldots,y_K$ multiplies it by
$\prod_a|P(y_a)|$; adding the corresponding negative
parameters $y_a+1$ multiplies it by
$\prod_a|P(y_a+1)|$.  Thus the displayed normalization is
exactly the one in its Gram sums.  The factor $\Gamma$
cancels in the estimate.  It has not been estimated or
replaced by one.

For the Fourier coefficient at level $j$, the interpolation
degree is $d=2j$.  If each actual grid has
$C_a\leq(28N/I)^{2j}$, and $A_{\rm pol}$ is the unscaled
boundary majorant
$\sup\sqrt{|P(y)P(y+1)|/q_0(y)}$ on the relevant interval,
one may take
\begin{equation}\label{ca-X-exact}
 X=\frac{g\lambda}{2\pi}
         \left(\frac{28N}{I}\right)^{2j}
                         e^{10/49}A_{\rm pol}.
\end{equation}
Here $g>0$ is the level parameter and $\lambda>0$ is the
actual scalar change of the complex integration variable.
This conclusion uses \eqref{ca-pairing-assumption} for the
particular Hermitian forms.  Its proof has not supplied
that representation-theoretic hypothesis.

\subsection{Fourier normalization and the remaining analytic input}

The source's additive character at a complex place is
\[
 \psi_g(l,c)=\exp\left(\frac{4\pi i}{g}\Re(lc)\right).
\]
If the radial variable is $z=\lambda c$, put
\[
 s=\frac{2\pi i\,l}{g\lambda},\qquad
 T=\frac{g\lambda}{2\pi}.
\]
Then the exact phase and scalar identities are
\begin{equation}\label{ca-phase}
 e^{sz-\overline s\,\overline z}=\psi_g(l,c),
 \qquad T^Ks^K=i^Kl^K.
\end{equation}
These follow by substituting $z=\lambda c$ and using that
$g,\lambda$ are real.  In the usual Fourier transform
$\widehat f(\xi)=\int_{\mathbb C}f(c)
e^{-2\pi i\Re(\overline\xi c)}\,dc$, the same character
therefore has frequency $\xi=-2\overline l/g$.

If the actual radial Fourier identity represents its
normalized coefficient as
\[
 \widehat K(l)
     =(\Gamma S)^{-1}\,
               \mathcal B(U_sx,x'),
\]
then the proved extraction identity identifies the mixed
form in Proposition~\ref{ca-positive-block} with
$i^Kl^K\widehat K(l)$.
Here $U_s$ consists of the explicit phase and the component
matrices described above.  There is a second necessary
identity: $U_s$ must preserve the original nonnegative
Hermitian form.  If it does, Cauchy--Schwarz supplies the
coefficient bound with no power of $l$ as well.
The diagonal forms in this latter bound need not be the
averaged diagonal forms in Proposition~\ref{ca-positive-block}.
If both diagonal products are bounded by a common
nonnegative majorant $M^2$, taking the stronger estimate
gives
\[
 |\widehat K(l)|\leq
       M\min\bigl(1,(X/|l|)^K\bigr)\qquad(l\ne0).
\]
At $l=0$ one uses the bound $M$ from the original form.
Obtaining this common majorant is a separate requirement
when assembling the global estimate.

Thus the differential-operator leading terms, real polynomial
extraction, rational correction, finite averaging, and exact
normalizations have now been proved.  The remaining actual
inputs are the radial Fourier identity, nonnegativity and
the mixed inequality \eqref{ca-pairing-assumption} for the
specified differential Hermitian forms, preservation of those
forms by $U_s$, and the analytic function-domain conditions.
The finite grid and separation conditions are discharged in
Section~\ref{cg-section}.  The mass,
source-dimension, and formal-degree hypotheses in
Section~\ref{as-section} also remain separate.

\begin{remark}\label{ca-counterexamples}
Two possible shortcuts are false.  First, entrywise bounds
do not imply the block positivity used here.  Let the two
diagonal matrices be $I_2$ and the mixed matrix have every
entry equal to one.  Each mixed entry has absolute value
at most the geometric mean of its two diagonal entries.
For $x=(1,1)$ the block quadratic form at $(x,-x)$ is
$4-8=-4$.  This is why \eqref{ca-pairing-assumption} is
required for all linear combinations.

Second, the real parameter in
\eqref{ca-mixed-polynomial} cannot be replaced by an
arbitrary complex parameter while retaining the same
polynomial identity.  With the ordinary form on
$\mathbb C$, $X_0(y)=Y_0(y)=y$ gives
$\mathcal B(X_0(y),Y_0(y))=|y|^2$, not $y^2$.
At $y=i$ these values are $1$ and $-1$.
The actual source uses real interpolation nodes and
tests the mixed form on all linear combinations, so
neither example contradicts its stated hypotheses.
\end{remark}


\section{Separated interpolation sets in the actual spectral interval}
\label{cg-section}

We now supply the separated interpolation sets left open in
Section~\ref{ca-section}.  The construction below is the one
used in the source: Chebyshev nodes in disjoint intervals,
together with a separate set in a shorter interval near the
left endpoint~\cite{alpo-uniform-source}.  Keeping the exact
Chebyshev coefficient sum improves the constant $28$ to $24$
at every degree needed in the argument.

\subsection{The spectral interval}

Use the parameters and polynomials $P,Q$ of
Section~\ref{as-section}.  Fix
$1\leq l\leq k-2$ with $\epsilon_l\ne\epsilon'_l$ and put
\[
 \alpha=\frac52-v_l,\qquad
 \beta=\frac{\Delta}{2}-v_l,\qquad
 I=\beta-\alpha=\frac{\Delta}{2}-\frac52.
\]
Let $r\geq0$ be an integer, $N=20(r+1)$ and $R=2r$.
For the construction assume
\begin{equation}\label{cg-length}
                        I\geq6N^3+6N.
\end{equation}
In particular $\Delta>5$.

\begin{lemma}\label{cg-spectral}
For every $Y\in[\alpha,\beta]$, $P(Y)\ne0$ and $Q(Y)>0$.
Also $-Y$ lies between the two differing spectral entries
with a margin of two:
\begin{equation}\label{cg-OnCut}
 \min(u_{l+1}(\epsilon),u_{l+1}(\epsilon'))+2
       \leq -Y
       \leq\max(u_{l+1}(\epsilon),u_{l+1}(\epsilon'))-2.
\end{equation}
\end{lemma}

\begin{proof}
Write $\tau=v_l+Y$, so $5/2\leq\tau\leq\Delta/2$.
The factors $Y+v_i=\tau+4\Delta(l-i)$ are positive
for $i\leq l$ and negative for $i>l$.
Furthermore
\[
 Y=\tau-v_l\leq\Delta/2-v_*-4\Delta(k-1-l)<0,
\]
so every $Y-b_j$ is negative.  This proves $P(Y)\ne0$.
The factor $Y+1/2+a$ is positive by
$a-v_1\geq4k\Delta$.
For $1\leq i\leq k-2$, the two factors of $Q$
associated with $u_{i+1}$ are
\[
              \tau+4\Delta(l-i)-\epsilon_i\Delta,
       \qquad \tau+4\Delta(l-i)-\epsilon'_i\Delta.
\]
They are both positive if $i<l$, and both negative if
$i>l$.  At $i=l$ they are $\tau$ and $\tau-\Delta$,
in some order, so their product is negative.
The two factors associated with $u_k$ are equal and
nonzero.  In the first case their common value is
$Y+v_*<0$; in the second it is $Y+1/2-b_*<0$.
Thus the product defining $Q$ has exactly one negative
paired factor.  Its initial factor $-TT'$ is negative,
so $Q(Y)>0$.

The two entries in \eqref{cg-OnCut} are
$v_l-1/2-\Delta$ and $v_l-1/2$.
The upper inequality is $\tau\geq5/2$.
The lower is $\tau\leq\Delta-3/2$, which follows from
$\tau\leq\Delta/2$ and $\Delta\geq3$.
\end{proof}

In particular $P(y)$ and $P(y+1)$ are nonzero whenever
$y,y+1\in[\alpha,\beta]$.  The same margin condition
holds at both parameters.

\subsection{The finite sets}

Put
\[
 q=\frac{I}{6N},\qquad \ell=q-1,\qquad
 A_a=\alpha+(2N+2a-2)q\quad(1\leq a\leq N).
\]
Condition~\eqref{cg-length} gives $q\geq N^2+1$ and
$\ell\geq N^2>0$.  Define
\[
                    J_a=[A_a,A_a+\ell].
\]
For $d\geq1$, let
\[
 c_{a,d,i}
   =A_a+\frac{\ell}{2}
                \left(1+\cos\frac{i\pi}{d}\right)
                       \quad(0\leq i\leq d),
\]
and set
\[
 \mathcal G_a=
       \bigcup_{1\leq d\leq\max(1,R)}
                    \{c_{a,d,0},\ldots,c_{a,d,d}\}.
\]
For positive $a$ outside $\{1,\ldots,N\}$ let
$\mathcal G_a=\varnothing$.
Independently choose arbitrary
$\delta_0,\ldots,\delta_{N-1}\in[0,1/2]$ and put
\[
 z_j=\alpha+jN^2+\delta_j\quad(0\leq j<N),\qquad
 \mathcal G_0=\{z_0,\ldots,z_{N-1}\}.
\]
All choices of the $\delta_j$ have the properties below.
The source's short-interval nodes, which additionally avoid
the finitely many zeros of $P(Y)P(Y+1)$ by construction,
are among these choices.  Here that avoidance also follows
from Lemma~\ref{cg-spectral}.

\begin{proposition}\label{cg-sets}
The set $\mathcal G_0$ has exactly $N$ elements.
Every $y$ in any of the displayed sets satisfies
$y,y+1\in[\alpha,\beta]$, and every $y\in\mathcal G_0$
also satisfies
\[
                 \frac52\leq v_l+y,\quad
                              v_l+y+1\leq N^3.
\]
For $e,e'\in\{0,1\}$, the following separations hold:
\begin{align}
 |(y+e)-(y'+e')|&\geq N^2+1
       &&(y\in\mathcal G_a,\ y'\in\mathcal G_{a'},\
                   1\leq a,a'\leq N,\ a\ne a'),\label{cg-long-sep}\\
 |(y+e)-(y'+e')|&\geq N^2+1
       &&(y\in\mathcal G_0,\ y'\in\mathcal G_a,\
                   1\leq a\leq N),\label{cg-cross-sep}\\
 |(y+e)-(y'+e')|&\geq N^2-\frac32\geq\frac{N^2}{8}
       &&(y,y'\in\mathcal G_0,\ y\ne y').\label{cg-short-sep}
\end{align}
\end{proposition}

\begin{proof}
Every long-interval point satisfies
\[
 \alpha+2Nq\leq y,\qquad
             y+1\leq\alpha+(4N-1)q<\alpha+6Nq=\beta.
\]
For a short-interval point,
\[
 \alpha\leq y,\qquad
 y+1\leq\alpha+(N-1)N^2+\frac32
                      \leq\alpha+N^3-\frac52.
\]
The last inequality uses $N^2\geq4$, and the last upper
bound is less than $\beta$ by \eqref{cg-length}.
It also proves the stated shorter-interval bound.

If $a<a'$, the least possible difference
$(y'+e')-(y+e)$ is at least
\[
             A_{a'}-(A_a+\ell+1)
                    =\bigl(2(a'-a)-1\bigr)q\geq q.
\]
This proves \eqref{cg-long-sep}.
The least difference between a long point with its shift
and a short point with its shift is at least
\[
          2Nq-N^3+\frac52
                 \geq N^3+2N+\frac52>N^2+1,
\]
proving \eqref{cg-cross-sep}.
If $j<j'$, then
$z_{j'}-z_j\geq N^2-1/2>0$.
Thus the $z_j$ are distinct, and allowing the two shifts
decreases their distance by at most one.
Finally $N^2-3/2\geq N^2/8$ for $N\geq20$.
\end{proof}

\begin{proposition}[Interpolation with the actual separation]
\label{cg-interpolation}
For every $1\leq a\leq N$ and every integer $0\leq d\leq R$,
there are $d+1$ distinct points $y_0,\ldots,y_d$ in
$\mathcal G_a$ such that their Lagrange leading-coefficient
weights satisfy
\begin{equation}\label{cg-weight}
 \sum_{i=0}^d\left|
       \prod_{h\ne i}(y_i-y_h)^{-1}\right|
                       \leq\left(\frac{24N}{I}\right)^d.
\end{equation}
\end{proposition}

\begin{proof}
For $d=0$, choose the endpoint $A_a+\ell$, which occurs
in the degree-one set.  The empty product is one, so
\eqref{cg-weight} holds.
For $d\geq1$ use $y_i=c_{a,d,i}$.
Lemma~\ref{ca-Chebyshev} gives the exact coefficient sum
$\frac12(4/\ell)^d$.  Since $q=I/(6N)$ and
$q\geq N^2+1$,
\[
 \frac4\ell=\frac{24N}{I}\frac{q}{q-1}
              \leq\frac{24N}{I}\left(1+\frac1{N^2}\right).
\]
For an integer $d\geq1$ and $x\geq0$ with $dx<1$,
\[
 (1+x)^d
     =\sum_{h=0}^d\binom dh x^h
     \leq\sum_{h=0}^{\infty}(dx)^h
     =\frac1{1-dx}.
\]
Here $d\leq R=2r<N/10$, so in particular
$d/N^2\leq1/2$ and the last quantity is at most two.
The factor $1/2$ in the exact coefficient sum absorbs
this loss.  This proves \eqref{cg-weight}.
The same proof works for any $d\leq N^2/2$, if the set
is constructed with nodes of that degree.
\end{proof}

\subsection{Extension of signed parameter sequences}

A signed parameter is a pair $(\varepsilon,t)$ with
$\varepsilon\in\{+,-\}$ and $t\in\mathbb R$.
Its underlying node is $t$ for sign $+$ and $t-1$
for sign $-$.
Call a finite sequence
$w=((\varepsilon_1,t_1),\ldots,(\varepsilon_h,t_h))$
admissible if $h\leq N$, its $i$th underlying node belongs to
$\mathcal G_{h+1-i}\cup\mathcal G_0$, and
\begin{equation}\label{cg-admissible}
       |t_i-t_j+\delta|>2
       \quad(i\ne j,\ \delta\in\mathbb Z,\ |\delta|\leq2).
\end{equation}
These are the actual per-place conditions used in the
source.  Every parameter belongs to $[\alpha,\beta]$
and hence satisfies \eqref{cg-OnCut}.

For a list $\boldsymbol y=(y_1,\ldots,y_K)$, write
\[
 w_+(\boldsymbol y)
       =((+,y_1),\ldots,(+,y_K),w),\qquad
 w_-(\boldsymbol y)
       =((-,y_1+1),\ldots,(-,y_K+1),w).
\]
The notation means concatenation with the original sequence.

\begin{theorem}\label{cg-extension}
Let $w$ be admissible of length $h$ and let $h+K\leq N$.
For every $0\leq d\leq R$ there are $K$ grids
$y_{a,0},\ldots,y_{a,d}$, $1\leq a\leq K$, satisfying
\eqref{cg-weight}, such that every selection
$\boldsymbol y=(y_{1,i_1},\ldots,y_{K,i_K})$
has the following properties:

\begin{enumerate}
\item Both $w_+(\boldsymbol y)$ and $w_-(\boldsymbol y)$
are admissible, have length $h+K$, and satisfy the
spectral margin \eqref{cg-OnCut} at every parameter.
\item The new nodes are pairwise at distance greater
than four.  For every old parameter $(\varepsilon,t)$,
every new node $y$ and every $e\in\{0,1\}$,
\[
                         |y+e-t|\geq N^2+1.
\]
\item The Lagrange product weights, normalized by their
sum of absolute values as in
Section~\ref{ca-section}, define a probability measure
on the resulting finite family of extended sequences.
\end{enumerate}
At degree $d=0$ one may instead choose $K$ distinct
nodes from $\mathcal G_0$ not already underlying old
parameters.  All conclusions still hold, with the
distance lower bound $N^2/8$ in place of $N^2+1$,
and every new node lies in the shorter interval above.
\end{theorem}

\begin{proof}
For $a=1,\ldots,K$ take the degree-$d$ grid from
$\mathcal G_{h+K+1-a}$ in
Proposition~\ref{cg-interpolation}.
These block indices are distinct and are all greater
than $h$.  Each old underlying node either belongs to
$\mathcal G_0$ or to a block with index at most $h$.
The separations \eqref{cg-long-sep} and
\eqref{cg-cross-sep} therefore apply to every new pair
and every new--old pair, for all shifts by zero or one.
In particular all those parameter distances exceed four.
Adding an integer of absolute value at most two leaves
a distance greater than two, proving
\eqref{cg-admissible} for the extended sequences.
The previously existing pairs already satisfy it.

For the block membership condition, the $a$th new
underlying node lies in
$\mathcal G_{h+K+1-a}$ by construction.
An old entry moves from position $i$ to $K+i$, and
\[
                  (h+K)+1-(K+i)=h+1-i.
\]
Thus its required block is unchanged.  Both signs
of extension use the same underlying nodes.  The
interval and spectral assertions follow from
Proposition~\ref{cg-sets} and Lemma~\ref{cg-spectral}.

Every Lagrange weight is nonzero, because the nodes
within its grid are distinct.  The product of the
normalized absolute weights sums to one by expansion
of a product of finite sums.  Assign the sum of these
weights to each sequence if several choices represent
the same sequence.  This gives the stated probability
measure without any counting factor.

For degree zero, at most $h$ elements of $\mathcal G_0$
are underlying nodes of old parameters.  Since
$|\mathcal G_0|=N\geq h+K$, at least $K$ unused
nodes remain.  Choose any $K$ of them.
Equation~\eqref{cg-short-sep} supplies all new--new
and relevant new--old separations within the short
set, and \eqref{cg-cross-sep} handles old long nodes.
Here $N^2/8\geq50>4$, so the same admissibility
argument applies.  Every grid has one node and weight
one.  For $K=0$ both constructions give the original
sequence and the probability measure of mass one on it.
\end{proof}

\begin{corollary}\label{cg-actual-constant}
The actual Fourier interpolation degree $d=2j$,
$0\leq j\leq r$, admits the construction above.
The denominators in \eqref{ca-q-distortion-eq} are
bounded below by $N^2/8$ for the short-interval
construction and by $N^2+1$ for the long one.
Consequently the algebraic barrier in
\eqref{ca-X-exact} may be replaced by
\[
 \frac{g\lambda}{2\pi}
        \left(\frac{24N}{I}\right)^{2j}
                             e^{10/49}A_{\rm pol}.
\]
This improvement retains the particular Hermitian-form
inequality and Fourier identity assumed in
Section~\ref{ca-section}.
\end{corollary}

\begin{proof}
For an old positive parameter, the denominator is
$y-t$; for an old negative parameter it is
$t-y-1$.  The latter is the difference between its
underlying node and $y$.  Both are covered by the
proved separation inequalities, using the appropriate
zero or one shift.  Apply
Lemma~\ref{ca-q-distortion} and
Proposition~\ref{ca-positive-block}.
\end{proof}

The interval condition \eqref{cg-length} also follows
eventually from the numerical range in
Section~\ref{as-section}.  Indeed, there $m=2r+2\leq k$,
so $N=10m\leq10k$.  If
$\Delta\geq(kf\log B)^{C_\Delta}$ with $C_\Delta>4$,
$f\geq2$ and $\log B\geq1$, then
$\Delta\geq k^{C_\Delta}$.
For all sufficiently large $k$, depending only on
$C_\Delta$, this gives
$\Delta/2-5/2\geq6000k^3+60k\geq6N^3+6N$.
Thus this construction introduces no new asymptotic
hypothesis in that range.

This closes the finite-set construction, all simultaneous
separation conditions, and the interpolation probabilities.
For a finite family of complex places, perform the same
construction at each place.  Its constants depend only on
$N,I$ and the degree, not on the translated interval
endpoint $\alpha$, the place, or the choices $\delta_j$.
The source bounds the lengths of the parameter sequences
by their maximum over the places.  Replacing the sequence
at one place by any extension of length at most $N$
therefore preserves the global length restriction.
It does not prove the radial Fourier identity, the
positivity and mixed inequality for the specified
differential forms, their wave isometry, or the analytic
domain and limiting assertions.  The separate mass,
source-dimension and formal-degree hypotheses remain
as stated earlier.


\section{The terminal radial operator and its positive density}
\label{rp-section}

We prove the finite-dimensional part of the radial positivity
argument, including the identity for the actual terminal
differential operator.  All ordered products below have the
same order and shifts as the operators in the source
\cite{alpo-uniform-source}.  The proof uses exterior powers
to establish the required identities for Capelli minors.
It then replaces a further interpolation over the spectrum
of a matrix by an orthogonal decomposition using real
interpolation polynomials.  The passage from this pointwise
density to the integrated radial forms still requires the
analytic assertions specified at the end of the section.

\subsection{Ordered minors and exterior powers}

Let $\mathcal A$ be a unital complex algebra and let
$E_{ab}\in\mathcal A$, $0\leq a,b<N$, satisfy
\begin{equation}\label{rp-gl}
 [E_{ab},E_{cd}]=\delta_{bc}E_{ad}-\delta_{da}E_{cb}.
\end{equation}
For increasing lists $I=(i_0,\ldots,i_{n-1})$ and
$J=(j_0,\ldots,j_{n-1})$, with $0\leq n\leq N$, put
\begin{equation}\label{rp-minor}
 Q_{I,J}(u)=
 \sum_{\sigma\in\mathfrak S_n}\operatorname{sgn}(\sigma)
 \prod_{a=0}^{n-1}
 \bigl((u-a)\delta_{i_{\sigma(a)},j_a}
                  +E_{i_{\sigma(a)},j_a}\bigr).
\end{equation}
The factors are multiplied from left to right in increasing
$a$.  An empty product is $1$.  Write $Q_N(u)$ for the
minor with both lists equal to $(0,\ldots,N-1)$.
The identities in this subsection are algebraic, so they
also apply to differential operators acting on a common
space of smooth functions.

\begin{lemma}\label{rp-capelli}
For $P\in\operatorname{GL}_N(\mathbb C)$ and $R=P^{-1}$,
define
\[
 E'_{ab}=\sum_{c,d}P_{ca}R_{bd}E_{cd}.
\]
Then
\begin{equation}\label{rp-minor-covariance}
 Q_{I,J}(E';u)=
 \sum_{\substack{|A|=n\\|B|=n}}
       \det P_{A,I}\,\det R_{J,B}\,Q_{A,B}(E;u).
\end{equation}
Here subsets are always written in increasing order.
If $\mathcal A=\operatorname{End}_{\mathbb C}(W)$ for a
finite-dimensional Hilbert space $W$ and
$E_{ab}^*=E_{ba}$, then
\begin{equation}\label{rp-minor-adjoint}
                      Q_{I,J}(u)^*=Q_{J,I}(\overline u).
\end{equation}
In particular $Q_N(u)$ is self-adjoint for real $u$.
\end{lemma}

\begin{proof}
Let $V=\mathbb C^N$ and let $L_a(v)$ act on the $a$th
factor of $V^{\otimes n}$, with matrix entries
$v\delta_{ij}+E_{ij}$.  Operators in $\mathcal A$ act in
their displayed order.  Let $P_{ab}$ interchange auxiliary
factors $a,b$, with indices now $0,\ldots,n-1$.
Expanding one matrix entry using \eqref{rp-gl} gives
\begin{equation}\label{rp-rtt}
 ((u-v)1-P_{ab})L_a(u)L_b(v)
 =L_b(v)L_a(u)((u-v)1-P_{ab}).
\end{equation}
Indeed the coefficient identity, after subtracting the
commuting scalar terms, is
\[
 (u-v)[E_{ij},E_{kl}]
 =(u-v)(\delta_{kj}E_{il}-\delta_{il}E_{kj}).
\]
Thus this formula requires precisely \eqref{rp-gl}, with
the indicated sign.

Let $A_n$ be the antisymmetrizer and set
\[
 A_n=\frac1{n!}\sum_{\sigma}\operatorname{sgn}(\sigma)P_\sigma,
 \qquad L=L_0(u)L_1(u-1)\cdots L_{n-1}(u-n+1).
\]
We claim that
\begin{equation}\label{rp-fusion}
                           A_nL=A_nLA_n.
\end{equation}
For adjacent positions $a,a+1$ their spectral parameters
differ by one.  Formula \eqref{rp-rtt} therefore moves
$(1-P_{a,a+1})/2$ through that pair, reversing their
order.  This projector commutes with every other factor
$L_b$, and $A_n(1-P_{a,a+1})/2=A_n$.  It follows that
$A_nL(1+P_{a,a+1})=0$.  Hence $A_nLP_\sigma
=\operatorname{sgn}(\sigma)A_nL$ for every permutation,
which proves \eqref{rp-fusion}.

Use the orthonormal exterior basis
\[
 \omega_I=\frac1{\sqrt{n!}}
      \sum_{\sigma}\operatorname{sgn}(\sigma)
      e_{i_{\sigma(0)}}\otimes\cdots\otimes e_{i_{\sigma(n-1)}}.
\]
The coefficient of $\omega_I$ in $A_nL(e_{j_0}\otimes
\cdots\otimes e_{j_{n-1}})$ is $Q_{I,J}(u)/\sqrt{n!}$.
On the other hand the antisymmetrization of the input is
$\omega_J/\sqrt{n!}$.  Equation \eqref{rp-fusion} proves
that the matrix entries of $A_nLA_n$ on $\bigwedge^nV$
are exactly $Q_{I,J}(u)$, with no additional factorial.

The auxiliary matrix $E'$ is $P^tE(P^t)^{-1}$.
Conjugate every auxiliary factor by $P^t$.
The antisymmetrizer commutes with this simultaneous change
of basis.  The entries of $\bigwedge^nP^t$ and
$\bigwedge^nR^t$ are respectively $\det P_{A,I}$ and
$\det R_{J,B}$.  This proves
\eqref{rp-minor-covariance}.

For the adjoint assertion first take $u$ real.  Each
$L_a(u-a)$ is self-adjoint because $E_{ij}^*=E_{ji}$.
Inside an expression $A_n(\cdots)A_n$ one may interchange
two neighboring factors whose labels are $a<b$.
To see this, use \eqref{rp-rtt} with parameter difference
$b-a$.  Its factor $(b-a)1-P_{ab}$ commutes with all the
other $L_c$ and acts on either exterior space as the
nonzero scalar $b-a+1$.  Move this factor across the
expression and cancel that scalar.  Reversing the ordered
list $0,\ldots,n-1$ by adjacent interchanges only exchanges
pairs with $a<b$ in this direction.  Consequently
\[
 A_nL_0(u)\cdots L_{n-1}(u-n+1)A_n
 =
 A_nL_{n-1}(u-n+1)\cdots L_0(u)A_n.
\]
The right side is the adjoint of the left side.
Comparison of exterior matrix entries proves
\eqref{rp-minor-adjoint} for real $u$.  Both sides are
polynomials in $\overline u$, so equality for all real
$u$ proves the general assertion.  For $n=0$ all these
identities reduce to $1=1$.
\end{proof}

\begin{corollary}\label{rp-central-character}
The full determinant $Q_N(u)$ commutes with every $E_{ab}$.
Suppose a module for these operators contains a vector $v$
such that
\[
 E_{ab}v=0\quad(a<b),\qquad E_{aa}v=\lambda_a v.
\]
Then
\begin{equation}\label{rp-highest-eigenvalue}
                 Q_N(u)v=\prod_{a=0}^{N-1}(u-a+\lambda_a)v.
\end{equation}
The same scalar identity holds on the submodule generated
by $v$.
\end{corollary}

\begin{proof}
For $n=N$, \eqref{rp-minor-covariance} says $Q_N(E';u)
=Q_N(E;u)$, because $\det P\det P^{-1}=1$.
Take $P(t)=\exp(tX)$ and differentiate at $0$.
The derivative of $E'_{ab}$ is
$[\sum_{c,d}X_{cd}E_{cd},E_{ab}]$ by \eqref{rp-gl}.
The product rule therefore gives
$[\sum X_{cd}E_{cd},Q_N(u)]=0$, proving centrality.
This argument may equally be read as an identity of
formal power series, so no topology on $\mathcal A$ is
needed.

In a summand of \eqref{rp-minor} belonging to a nonidentity
permutation, let $a$ be its largest nonfixed column.
Then $\sigma(a)<a$.  All factors to its right are diagonal
and act on $v$ by scalars; its factor is $E_{\sigma(a),a}$
and annihilates $v$.  Only the identity permutation
remains, yielding \eqref{rp-highest-eigenvalue}.
Centrality extends the identity from $v$ to every vector
obtained by applying the $E_{ab}$, and thence to their
linear span.
\end{proof}

\begin{corollary}\label{rp-negative-convention}
For the alternative convention
\[
 B_N(u)=\sum_{\sigma}\operatorname{sgn}(\sigma)
           \prod_{a=0}^{N-1}
             \bigl((u+a)\delta_{\sigma(a),a}-E_{\sigma(a),a}\bigr)
\]
one has $B_N(u)=(-1)^NQ_N(-u)$.  It is central,
self-adjoint for real $u$ under the preceding adjoint
hypothesis, and acts on the highest-weight submodule
above by $\prod_a(u+a-\lambda_a)$.
\end{corollary}

\begin{proof}
Factor $-1$ out of every factor in the ordered product.
The other conclusions follow from the two preceding
results.
\end{proof}

\subsection{The complementary contraction}

For $N\geq2$, vectors $a,a'\in\mathbb C^N$ and covectors
$\beta,\beta'\in(\mathbb C^N)^*$, define
\begin{align}
 \mathcal C_E(u;a,a',\beta,\beta')
 ={}&\sum_{\substack{i<i'\\j<j'}}
       (-1)^{i+i'+j+j'}
       (\beta_i\beta'_{i'}-\beta_{i'}\beta'_i)
       (a_ja'_{j'}-a_{j'}a'_j)\notag\\
 &\hspace{32mm}\cdot Q_{\{i,i'\}^{c},\{j,j'\}^{c}}(u).
 \label{rp-contraction}
\end{align}
All coordinate factors in this formula are scalars placed
to the left of the differential operators, when these
are the chosen algebra $\mathcal A$.

\begin{lemma}\label{rp-adapted-contraction}
Let $P=(a,a',b_1,\ldots,b_{N-2})$ be an invertible matrix.
Suppose $\beta(b_i)=\beta'(b_i)=0$ for all $i$, and let
$E'$ be its change of generators from Lemma~\ref{rp-capelli}.
Then
\begin{equation}\label{rp-contraction-adapted}
 \mathcal C_E(u;a,a',\beta,\beta')
 =\bigl(\beta(a)\beta'(a')-\beta(a')\beta'(a)\bigr)
     Q_{\{2,\ldots,N-1\},\{2,\ldots,N-1\}}(E';u).
\end{equation}
\end{lemma}

\begin{proof}
We first verify covariance of \eqref{rp-contraction}.
For two-element subsets $I,J$ put
\[
       C_{I,J}(E)=(-1)^{\sum I+\sum J}Q_{I^c,J^c}(E;u).
\]
Complementary minors of an invertible scalar matrix obey
\[
 \det P_{A^c,I^c}
   =(-1)^{\sum A+\sum I}\det P\,\det(P^{-1})_{I,A}.
\]
For completeness this is the coefficient identity obtained
by applying $\bigwedge^NP$ to
$e_A\wedge e_{A^c}$ and comparing the wedge pairing
\[
 \bigwedge\nolimits^2\mathbb C^N\times
 \bigwedge\nolimits^{N-2}\mathbb C^N
                   \longrightarrow\bigwedge\nolimits^N\mathbb C^N.
\]
The sign of $e_A\wedge e_{A^c}$ is
$(-1)^{\sum A-1}$; the two constant signs cancel.
Equivalently, expanding that wedge pairing is the usual
Laplace expansion, which proves the displayed identity
without any assumption on the complementary minor.

Apply \eqref{rp-minor-covariance} and this identity twice.
Since $\det P\det P^{-1}=1$, one obtains
\[
 C(E')=(\bigwedge\nolimits^2P^{-1})\,C(E)\,
                                     (\bigwedge\nolimits^2P).
\]
The vector $a\wedge a'$ changes by $\bigwedge^2P^{-1}$,
and the covector $\beta\wedge\beta'$ by $\bigwedge^2P$.
Their contraction with $C$ is consequently unchanged.
In the basis formed by the columns of $P$ the two vectors
are $e_0,e_1$, and both covectors vanish in positions
$2,\ldots,N-1$.  Only $I=J=\{0,1\}$ survives in
\eqref{rp-contraction}; its sign is $+1$.  This gives
\eqref{rp-contraction-adapted}, including the case $N=2$,
whose remaining determinant is the empty determinant $1$.
\end{proof}

\subsection{The actual terminal differential operator}

Let $V$ be a complex Hermitian space of signature $(k,1)$,
$k\geq1$, with its Hermitian form $h$ linear in the first
argument.  Let $z,z'\in V$ satisfy
$h(z,z)=-T$, $h(z',z')=-T'$ with $T,T'>0$.
For $g\in U(V,h)$ put
\[
 y=g^{-1}z,\qquad x=z',\qquad
 \zeta=h(gz',z)=h(x,y),\qquad
 G(g)=|\zeta|^2-TT'.
\]
We work at points where $G(g)>0$.  On the terminal group
of the source, the previously fixed isotropic vectors
have been removed; the remaining space is exactly this
$(k+1)$-dimensional space $V$.

For $A\in\operatorname{End}_{\mathbb C}(V)$ write
$A^\dagger$ for its $h$-adjoint and define
\[
 A_1=\tfrac12(A-A^\dagger),\qquad
 A_2=\tfrac1{2i}(A+A^\dagger).
\]
Both belong to the real Lie algebra $\mathfrak u(V,h)$,
and $A=A_1+iA_2$.  If $F$ is a smooth function on an
open part of $U(V,h)$, define
\[
 R_XF(g)=\left.\frac{d}{dt}\right|_{t=0}F(g\exp(tX)),
 \qquad E_AF=R_{A_1}F+iR_{A_2}F .
\]
In a fixed basis of $V$ set $E_{ab}=E_{e_a e_b^*}$,
where $e_b^*$ is the coordinate covector, not the
Hermitian dual.  The commutator formula for
left-invariant vector fields gives \eqref{rp-gl}.
It also follows directly by differentiating the group
commutator
$\exp(tX)\exp(sY)\exp(-tX)\exp(-sY)$ at $(0,0)$
and extending complex linearly.

At the fixed point $g$, put
$\beta(v)=h(v,y)$ and $\beta'(v)=h(v,x)$.
Define the terminal $(2,2)$ entry by
\begin{equation}\label{rp-actual-terminal}
 (\mathcal T_{22}(Y)F)(g)
 =-\bigl(\mathcal C_E(Y+(k-2)/2;y,x,\beta,\beta')F\bigr)(g).
\end{equation}
The point-dependent scalar coordinates of $y,\beta$
multiply the operators on the left; they are not
differentiated in this definition.
This is the terminal entry of the original two-by-two
operator: its full spectral shift is $Y+k/2$, and the
complementary contraction uses that shift minus one.

\begin{proposition}\label{rp-terminal-identity}
Choose an $h$-orthonormal basis $b_1,\ldots,b_{k-1}$
of $\langle y,x\rangle^\perp$, and put
$U_{ab}(v)=h(v,b_b)b_a$.  Let
$Q_{k-1}(E_U;u)$ be the ordered full determinant of
the operators $E_{U_{ab}}$.  Then
\begin{equation}\label{rp-terminal-action}
 (\mathcal T_{22}(Y)F)(g)
       =G(g)\bigl(Q_{k-1}(E_U;Y+(k-2)/2)F\bigr)(g).
\end{equation}
The basis in this formula is held fixed while the
differential operators are applied.
\end{proposition}

\begin{proof}
The Gram matrix of $y,x$ is
$\left(\begin{smallmatrix}-T&\overline\zeta\\
\zeta&-T'\end{smallmatrix}\right)$, with determinant
$TT'-|\zeta|^2=-G(g)<0$.  It is nondegenerate of
signature $(1,1)$; its orthogonal complement is positive
definite of dimension $k-1$.  Gram--Schmidt therefore
gives the stated basis, including the empty basis when
$k=1$.

Use the basis $P=(y,x,b_1,\ldots,b_{k-1})$ of $V$.
The covectors $\beta,\beta'$ vanish on its last $k-1$
columns, and
\[
 \beta(y)\beta'(x)-\beta(x)\beta'(y)
                    =TT'-|\zeta|^2=-G(g).
\]
For a last-column basis index $b$, the corresponding
row of $P^{-1}$ is $v\mapsto h(v,b_b)$.
Consequently the transformed generator $E'_{ab}$ in
the last $(k-1)\times(k-1)$ block is $E_{U_{ab}}$:
the rank-one matrix underlying it is
$P_{\cdot,a}(P^{-1})_{b,\cdot}=b_a h(\,\cdot\,,b_b)$.
Lemma~\ref{rp-adapted-contraction}, followed by the
minus sign in \eqref{rp-actual-terminal}, proves
\eqref{rp-terminal-action}.
\end{proof}

Let $K_{z'}$ be the stabilizer of $z'$ in $U(V,h)$.
It is the compact unitary group of the positive space
$(z')^\perp$.  Let $W,W'$ be finite-dimensional complex
Hilbert spaces and let $\pi':K_{z'}\to U(W)$ be a
continuous representation.  Suppose that
$F:U(V,h)\supset D\to\operatorname{Hom}_{\mathbb C}(W,W')$
is smooth and satisfies
\begin{equation}\label{rp-equivariance}
                     F(gh)=F(g)\pi'(h)
\end{equation}
whenever $h\in K_{z'}$ and the relevant points lie in $D$.
It suffices that this holds in a neighborhood of each
point where a derivative is taken.

\begin{lemma}\label{rp-unitary-derivative}
A continuous finite-dimensional unitary representation
of a compact matrix Lie group is smooth.  Its
complexified derivative is a Lie algebra representation.
For the matrices $U_{ab}$ above, its operators
$D_{ab}=d\pi'_{\mathbb C}(U_{ab})$ satisfy
\[
 [D_{ab},D_{cd}]=\delta_{bc}D_{ad}-\delta_{da}D_{cb},
                      \qquad D_{ab}^*=D_{ba}.
\]
\end{lemma}

\begin{proof}
Here is a direct smoothness argument.  Normalize Haar
measure on the compact group to have mass one.  For a
smooth nonnegative function $f$ of integral one supported
near the identity, the operator
$\pi'(f)=\int f(h)\pi'(h)\,dh$ is arbitrarily close to
the identity, by continuity, and hence invertible.
For every $v$ the function
\[
 \pi'(g)\pi'(f)v=\int f(g^{-1}h)\pi'(h)v\,dh
\]
is smooth in $g$, by differentiation under the compact
integral.  Since $\pi'(f)$ is onto, every orbit map is
smooth.  Such functions $f$ exist by a smooth coordinate
chart and a nonnegative smooth bump function.  Haar
measure here may be obtained from the invariant
Riemannian volume on the compact unitary group.

Differentiating the representation identity and its
group commutator gives a real Lie algebra homomorphism;
its complex-linear extension has the same property.
Differentiating preservation of the inner product gives
$d\pi'(X)^*=-d\pi'(X)$ for a skew-Hermitian $X$.
Since $U_{ab}^\dagger=U_{ba}$, the decomposition into
the two skew-Hermitian parts above gives
$D_{ab}^*=D_{ba}$.  Finally
$U_{ab}U_{cd}=\delta_{bc}U_{ad}$ by orthonormality,
which gives the displayed commutator.
\end{proof}

\begin{corollary}\label{rp-finite-terminal}
For a fixed $g$ and its fixed basis $(b_a)$ put
\[
                 C_g(Y)=Q_{k-1}(D;Y+(k-2)/2).
\]
For every smooth equivariant $F$ as above,
\begin{equation}\label{rp-evaluation-transfer}
             G(g)^{-1}(\mathcal T_{22}(Y)F)(g)=F(g)C_g(Y).
\end{equation}
The operator $C_g(Y)$ on $W$ is self-adjoint for real
$Y$.  No choice of a smooth family of bases in $g$ is
needed for either assertion.
\end{corollary}

\begin{proof}
Every $U_{ab}$ annihilates $z'$; its two real Lie algebra
parts lie in the Lie algebra of $K_{z'}$.
Differentiating \eqref{rp-equivariance} gives
$E_{U_{ab}}F(g)=F(g)D_{ab}$.  Repeated differentiation
preserves the order: for example
$E_{U_{ab}}E_{U_{cd}}F(g)=F(g)D_{ab}D_{cd}$.
Thus every ordered minor transfers to the corresponding
minor in $D$.  Proposition~\ref{rp-terminal-identity}
proves \eqref{rp-evaluation-transfer}, and
Lemma~\ref{rp-capelli} proves self-adjointness.
\end{proof}

\subsection{A pointwise orthogonal decomposition}

We state explicitly the algebraic part of the normal
form that the positivity argument uses.  This does not
assume that the algebraic involution on differential
operators is already an adjoint for an integrated
radial form.

\begin{proposition}\label{rp-idempotent-positive}
Let $\mathcal S$ be a nonempty finite set and let
$p_\rho(Y)\in\mathbb R[Y]$, $\rho\in\mathcal S$, be
pairwise distinct.  Let $\mathcal V$ be a complex vector
space, let
$\kappa:\mathbb C^{\mathcal S}\to
\operatorname{End}_{\mathbb C}(\mathcal V)$ be a unital
algebra homomorphism, and let
$\operatorname{ev}:\mathcal V\to\operatorname{Hom}(W,W')$
be linear.  Suppose that for each real $Y$ there is a
self-adjoint $C(Y)\in\operatorname{End}(W)$ such that
\begin{equation}\label{rp-normal-evaluation}
 \operatorname{ev}\bigl(\kappa((p_\rho(Y))_\rho)\Phi\bigr)
                         =(\operatorname{ev}\Phi)C(Y)
                  \qquad(\Phi\in\mathcal V).
\end{equation}
Write $P_\rho=\kappa(1_{\{\rho\}})$ and use the
Hilbert--Schmidt form linear in the first argument,
$\langle A,B\rangle_{\rm HS}=\operatorname{tr}(B^*A)$.
Then the maps $\operatorname{ev}(P_\rho\Phi)$ are
pairwise orthogonal and
\begin{equation}\label{rp-positive-decomposition}
 \left\langle\operatorname{ev}(\kappa(f)\Phi),
                    \operatorname{ev}\Phi\right\rangle_{\rm HS}
       =\sum_{\rho\in\mathcal S}
              f(\rho)\|\operatorname{ev}(P_\rho\Phi)\|_{\rm HS}^2
                  \qquad(f\in\mathbb C^{\mathcal S}).
\end{equation}
In particular it is nonnegative for real $f\geq0$,
and is bounded between
$\min f\,\|\operatorname{ev}\Phi\|_{\rm HS}^2$ and
$\max f\,\|\operatorname{ev}\Phi\|_{\rm HS}^2$.
\end{proposition}

\begin{proof}
Choose a real $Y_0$ at which all $p_\rho(Y_0)$ are
distinct.  Such a point exists because each nonzero
difference polynomial has finitely many roots.
Put $\xi_\rho=p_\rho(Y_0)$ and form the real
interpolation polynomials
\[
 \ell_\rho(X)=
       \prod_{\sigma\ne\rho}\frac{X-\xi_\sigma}{\xi_\rho-\xi_\sigma}.
\]
Since $\kappa$ is an algebra homomorphism,
$P_\rho=\ell_\rho(\kappa((\xi_\sigma)_\sigma))$.
Induction on the degree of a polynomial in
\eqref{rp-normal-evaluation} therefore gives
\[
       \operatorname{ev}(P_\rho\Phi)
          =(\operatorname{ev}\Phi)\ell_\rho(C(Y_0)).
\]
The matrix $\ell_\rho(C(Y_0))$ is self-adjoint.
For any $\Phi,\Psi\in\mathcal V$, cyclicity of the
finite-dimensional trace gives
\[
 \langle\operatorname{ev}(P_\rho\Phi),\operatorname{ev}\Psi\rangle_{\rm HS}
 =\langle\operatorname{ev}\Phi,\operatorname{ev}(P_\rho\Psi)\rangle_{\rm HS}.
\]
On $\mathcal V$ the $P_\rho$ are mutually annihilating
idempotents whose sum is the identity.  Apply this last
identity to $P_\rho\Phi$ and $P_\sigma\Phi$ to see that
their evaluations are orthogonal for $\rho\ne\sigma$.
Their evaluations sum to $\operatorname{ev}\Phi$.
Finally $\kappa(f)=\sum_\rho f(\rho)P_\rho$, which gives
\eqref{rp-positive-decomposition} and its bounds.
No assertion about the spectrum of $C(Y_0)$ on vectors
outside these evaluated functions is needed.
\end{proof}

\subsection{The actual interlacing parameters and signs}

Here are the precise scalar parameters used in this
application.  Let $k\geq1$, $T,T'>0$, and let
$u_1,\ldots,u_k,u'_1,\ldots,u'_k$ be real numbers with
\begin{equation}\label{rp-interlacing}
 \max(u_{i+1},u'_{i+1})+1
              \leq\min(u_i,u'_i)\qquad(1\leq i<k).
\end{equation}
The source also requires these numbers to lie in a
common coset of $\mathbb Z$; the following arguments
do not require that extra condition.
Set
\[
 a_i=\max(u_{i+1},u'_{i+1})+\tfrac12,\qquad
 b_i=\min(u_i,u'_i)-\tfrac12,
\]
and let $\mathcal S$ be the product of the finite sets
$\{a_i,a_i+1,\ldots,a_i+\lfloor b_i-a_i\rfloor\}$,
$1\leq i<k$.  Thus a member $\rho$ is a tuple
$(\rho_1,\ldots,\rho_{k-1})$ with $\rho_i\in[a_i,b_i]$.
The set is nonempty; for $k=1$ it consists of the empty
tuple.  Put
\begin{equation}\label{rp-spectral-polynomials}
 p_\rho(Y)=\prod_{i=1}^{k-1}(Y+\rho_i),\qquad
 q(Y)=-TT'\prod_{i=1}^{k}(Y+\tfrac12+u_i)
                                  (Y+\tfrac12+u'_i).
\end{equation}

\begin{lemma}\label{rp-parameter-signs}
The polynomials $p_\rho$ are pairwise distinct.
Fix $1\leq l\leq k$ and a real $Y$ with
\begin{equation}\label{rp-cut-margin}
          \min(u_l,u'_l)+2\leq-Y\leq\max(u_l,u'_l)-2.
\end{equation}
Then, for every $\rho\in\mathcal S$,
\[
 (-1)^{k-l}p_\rho(Y+\delta)>0\quad(-1\leq\delta\leq1),
               \qquad q(Y)>0,\quad q(Y-1)>0.
\]
\end{lemma}

\begin{proof}
For $i<k-1$,
$\rho_i\geq\max(u_{i+1},u'_{i+1})+1/2$ and
$\rho_{i+1}\leq\min(u_{i+1},u'_{i+1})-1/2$.
Hence $\rho_i-\rho_{i+1}\geq1$.
Equality of the monic polynomials $p_\rho,p_\sigma$
gives equality of their multisets of real roots, and
the decreasing order gives $\rho=\sigma$.

For $i<l$, \eqref{rp-interlacing} implies
$\rho_i\geq\max(u_l,u'_l)+1/2$, so
$Y+\rho_i\geq5/2$.
For $i\geq l$ it gives
$\rho_i\leq\min(u_l,u'_l)-1/2$, so
$Y+\rho_i\leq-5/2$.
These signs persist after a shift of absolute value at
most one.  There are exactly $k-l$ negative factors,
proving the first assertion, also for the empty tuple.

For $i<l$, both $u_i,u'_i$ are at least
$\max(u_l,u'_l)+1$; for $i>l$ both are at most
$\min(u_l,u'_l)-1$.  Thus the two factors belonging to
each $i\ne l$ in $q(Y)$ and $q(Y-1)$ have the same,
nonzero sign.  At $i=l$, one entry is at most $-Y-2$
and the other is at least $-Y+2$ by
\eqref{rp-cut-margin}.  The two factors have opposite
signs both after adding $1/2$ and after subtracting
$1/2$.  The product in \eqref{rp-spectral-polynomials}
is negative, and its factor $-TT'$ makes it positive.
\end{proof}

For clarity we give the exact recursion for the scalar
diagonal coefficients in the source.  Fix a finite set
of positions $S$ and a list $w$ of triples
$\ell=(\varepsilon,a,y)$, where
$\varepsilon\in\{+,-\}$, $a\in S$, and $y\in\mathbb R$.
Positions are distinct and any two parameters have
distance greater than two.  Fix
$e:S\to\{0,1\}$.  For a parameter $Y$ set
\[
 \phi_\ell(Y)=
 \begin{cases}(Y-y)/(Y-y-1),&\varepsilon=+,\\
               (y-Y)/(y-Y-1),&\varepsilon=-,
 \end{cases}
 \qquad
 \eta_\ell(Y)=
 \begin{cases}(Y-y+1)/(Y-y),&\varepsilon=+,\\
               (y-Y)/(y-Y-1),&\varepsilon=-.
 \end{cases}
\]
Define $q_\varnothing(Y)=q(Y)$,
$x_{\rho,\varnothing,e}(Y)=(-1)^{k-l}p_\rho(Y)$
and $c_{\rho,\varnothing,e}=1$.  On prepending a triple
$\ell=(\varepsilon,a,y)$, set
\begin{align*}
 q_{\ell w}(Y)&=q_w(Y)\eta_\ell(Y),\\
 x_{\rho,\ell w,e}(Y)&=
       x_{\rho,w,e}(Y)
           \begin{cases}\phi_\ell(Y),&e(a)=0,\\1,&e(a)=1,\end{cases}\\
 c_{\rho,\ell w,e}&=c_{\rho,w,e}
   \begin{cases}
      q_w(y)/x_{\rho,w,e}(y+1),&\varepsilon=+,\,e(a)=0,\\
      q_w(y-1)/x_{\rho,w,e}(y-1),&\varepsilon=-,\,e(a)=0,\\
      x_{\rho,w,e}(y),&e(a)=1 .
   \end{cases}
\end{align*}

\begin{lemma}\label{rp-recursive-positive}
If every parameter $y$ in the list satisfies
\eqref{rp-cut-margin}, then all the coefficients
$c_{\rho,w,e}$ are strictly positive, and every
denominator appearing in their definition is nonzero.
\end{lemma}

\begin{proof}
If $|s|>1$, all of $s/(s-1)$ and $(s+1)/s$ are
positive.  The same statement with $s$ replaced by
$-s$ proves that $\phi_\ell(Y)$ and $\eta_\ell(Y)$
are positive whenever $|Y-y|>1$.
When a new parameter $y$ is inserted, each previously
inserted parameter $y'$ satisfies $|y-y'|>2$.
Thus $|y+\delta-y'|>1$ for $\delta=-1,0,1$.
Lemma~\ref{rp-parameter-signs} and the product formulas
give
$x_{\rho,w,e}(y),x_{\rho,w,e}(y-1),
x_{\rho,w,e}(y+1),q_w(y),q_w(y-1)>0$.
Every case of the recursion for $c$ is consequently a
positive product or quotient.  Induction on the length
proves both assertions.
\end{proof}

\begin{corollary}\label{rp-terminal-density-positive}
Take the actual terminal evaluation
$\operatorname{ev}_{c,g}\Phi=\Phi(c,g)$, with
$G(g)>0$, on the space of smooth radial functions.
Suppose that its algebraic normal form supplies a
unital algebra homomorphism $\kappa$ with
\begin{equation}\label{rp-required-normal-form}
 \kappa((p_\rho(Y))_\rho)=G^{-1}\mathcal T_{22}(Y)
                         \quad\text{on that function space}.
\end{equation}
Then \eqref{rp-positive-decomposition} holds at every
such point, with the actual parameters
\eqref{rp-spectral-polynomials}.

In particular, let $\Psi_e$ be any family of functions
in that space, indexed by $e:S\to\{0,1\}$, and let
$h$ be the length of a list satisfying the hypotheses
of Lemma~\ref{rp-recursive-positive}.  The terminal
density
\begin{equation}\label{rp-terminal-density}
 \mathcal D(c,g)=
  \left(\frac{|\zeta(g)|^2}{G(g)}\right)^h
  \sum_e\left\langle
    \operatorname{ev}_{c,g}
       \bigl(\kappa((c_{\rho,w,e})_\rho)\Psi_e\bigr),
            \operatorname{ev}_{c,g}\Psi_e\right\rangle_{\rm HS}
\end{equation}
is nonnegative.  The choice $\Psi_e=(F_w\Phi)_e$
is exactly the terminal density in the source.
If the radial functions are extended by zero off their
prescribed domain, the same conclusion holds there.
\end{corollary}

\begin{proof}
Corollary~\ref{rp-finite-terminal} gives
\eqref{rp-normal-evaluation} from
\eqref{rp-required-normal-form}.
Lemma~\ref{rp-parameter-signs} supplies distinct
polynomials, so Proposition~\ref{rp-idempotent-positive}
applies.  Each coefficient $c_{\rho,w,e}$ is positive
by Lemma~\ref{rp-recursive-positive}.
Also $|\zeta|^2=G+TT'>0$.  Every summand in
\eqref{rp-terminal-density} is therefore nonnegative.
Off the radial domain both evaluated functions are
zero, by the specified extension.
\end{proof}

\subsection{Scope and checks on the hypotheses}

The terminal identity
\eqref{rp-terminal-action} and its compact
representation interpretation
\eqref{rp-evaluation-transfer} are proved here for
the actual operators, with their exact shifts and signs.
For the subsequent positivity step the only normal-form
property used is \eqref{rp-required-normal-form};
the relations involving the additional raising and
lowering operators are unnecessary for this step.
The source derives this property as the $(2,2)$ entry
of a larger normal form.  Section~\ref{rd-section}
constructs exactly the required property directly
from compact branching.  The orthogonal decomposition
also strengthens the source's positivity conclusion to
the exact sum of squares \eqref{rp-positive-decomposition}
and its minimum and maximum bounds.

This pointwise result does not yet identify the
integrated radial form with the integral of
\eqref{rp-terminal-density}.  That identification needs
the adjoint relations on the actual class of functions,
integrability after the successive differential
operations, and the successive one-variable integral
identities with their boundary terms.  The wave
identities and the common positive majorants used in
Section~\ref{ca-section} consequently remain subject to
those analytic steps.

Two simple checks explain essential hypotheses without
giving counterexamples to the source.  A scalar operator
is not positive merely because it is self-adjoint:
for $N=1$, $E_{00}=0$, the determinant is $Q_1(u)=u$,
which is negative for real $u<0$.
The positivity above uses the positive interlacing
weights, not positivity of an arbitrary Capelli
polynomial.  Likewise replacing the strict separation
$|y-y'|>2$ by merely distinct parameters can make the
recursion undefined: with two plus parameters $y'=0$
and $y=1$, the factor $\phi_{\ell'}(y)=y/(y-1)$
has a pole.  The interpolation sets in
Section~\ref{cg-section} satisfy the stronger separation,
and Lemma~\ref{rp-parameter-signs} supplies all the
actual sign conditions.


\section{The finite radial decomposition without the larger normal form}
\label{rd-section}

We now construct the algebra homomorphism required in
\eqref{rp-required-normal-form}.  Only the terminal
operator proved in Section~\ref{rp-section} and the
compact branching law of Theorem~\ref{ib-compact-branching}
are needed.  In particular, the construction does not
require the additional raising and lowering operators
in the larger source normal form
\cite{alpo-uniform-source}.

\subsection{The function space and the actual matrix field}

Retain the Hermitian space $V$ of signature $(k,1)$,
the vectors $z,z'$ and the positive constants $T,T'$
from Section~\ref{rp-section}.  Let
\[
 K_z=\{a\in U(V,h):az=z\},\qquad
 K_{z'}=\{b\in U(V,h):bz'=z'\}.
\]
Let $\pi:K_z\to U(W')$ and
$\pi':K_{z'}\to U(W)$ be nonzero irreducible,
finite-dimensional continuous unitary representations.
Let $r\geq0$ and set
\[
 \Omega=(\mathbb C^\times)^r,\qquad
 D=\{g\in U(V,h):G(g)>0\}.
\]
For $r=0$, $\Omega$ is a one-point space.
Define $\mathscr R$ to be the space of smooth functions
\[
 \Phi:\Omega\times D\longrightarrow\operatorname{Hom}(W,W')
\]
with
\begin{equation}\label{rd-biequivariance}
 \Phi(c,agb)=\pi(a)\Phi(c,g)\pi'(b)
              \quad(a\in K_z,\ b\in K_{z'}).
\end{equation}
Extend these functions by zero outside $\Omega\times D$.
Smoothness is required on this open submanifold only,
not across its boundary.  This is the source's radial
function space after identifying the terminal group
with $U(V,h)$.  Its formulation using local smooth
extensions to matrix space is equivalent: local
coordinates for an embedded Lie group extend a smooth
function by making it constant in the normal
coordinates.

For $g\in D$, write $x=z'$, $y=g^{-1}z$, and let
\[
                       E_g=\langle y,x\rangle^\perp.
\]
This is a positive space of dimension $k-1$.
Using any orthonormal basis of $E_g$, define the
matrix $C_g(Y)$ on $W$ by
Corollary~\ref{rp-finite-terminal}.

\begin{lemma}\label{rd-smooth-matrix}
The matrix $C_g(Y)$ is independent of the chosen
orthonormal basis of $E_g$.  It is smooth in $g$,
polynomial in $Y$, and obeys
\begin{equation}\label{rd-matrix-equivariance}
 C_{agb}(Y)=\pi'(b)^{-1}C_g(Y)\pi'(b).
\end{equation}
Consequently
\begin{equation}\label{rd-actual-X}
             (X(Y)\Phi)(c,g)=\Phi(c,g)C_g(Y)
\end{equation}
defines an endomorphism of $\mathscr R$.  This is the
actual differential operator
$G^{-1}\mathcal T_{22}(Y)$, and it does not enlarge
the support of a function.
\end{lemma}

\begin{proof}
A change of orthonormal basis in $E_g$ conjugates its
general-linear generators exactly as in
Lemma~\ref{rp-capelli}.  Full determinant covariance
makes $C_g(Y)$ unchanged.

To check smoothness without choosing a global frame,
put
\[
 v_g=y+\frac{h(y,x)}{T'}x.
\]
Then $v_g\in x^\perp$ and direct expansion gives
$h(v_g,v_g)=G(g)/T'>0$.  The space $E_g$ is the
orthogonal complement of $v_g$ in the fixed positive
space $x^\perp$.  Its orthogonal projection is
\[
       v\longmapsto v-\frac{h(v,v_g)}{h(v_g,v_g)}v_g
                                      \quad(v\in x^\perp),
\]
which varies smoothly in $g$.  Near a fixed point,
project a fixed basis of $E_g$ by these projections
and apply positive-definite Gram--Schmidt.  The
resulting local orthonormal frame is smooth, since
all denominators and the required squared norms
are positive near that point.  The definition of
$C_g(Y)$ is polynomial in the associated rank-one
matrices and their images under $d\pi'_{\mathbb C}$,
so it is smooth in that neighborhood.  Frame
independence makes these descriptions agree.

For $g'=agb$ one has $y'=b^{-1}y$, while
$b^{-1}x=x$.  Thus $E_{g'}=b^{-1}E_g$.
Conjugating the rank-one matrices by $b^{-1}$
conjugates their representation operators by
$\pi'(b)^{-1}$.  Ordered products give
\eqref{rd-matrix-equivariance}.
Combining it with \eqref{rd-biequivariance} proves
that \eqref{rd-actual-X} again satisfies
\eqref{rd-biequivariance}.  Joint smoothness and
the zero extension are preserved.
Corollary~\ref{rp-finite-terminal}, applied to
$g\mapsto\Phi(c,g)$ for fixed $c$, identifies it
with the actual differential operator.  The
displayed multiplication formula proves the
support assertion.
\end{proof}

\subsection{Matching the source parameters with highest weights}

Let the spectral parameters $u,u'$ satisfy
\eqref{rp-interlacing}, and suppose the compact types
are specified by the source's Capelli equations:
on the representation with parameters $u$ the
full $k$-dimensional determinant is
\begin{equation}\label{rd-source-Capelli}
       Q_k(w)=\prod_{i=1}^k
           \left(w-\frac{k-1}{2}+u_i\right)\operatorname{Id},
\end{equation}
and the analogous equation holds for $u'$.
The common integral coset condition is as in
Section~\ref{rp-section}.

\begin{lemma}\label{rd-highest-parameters}
With highest weights indexed from $1$ to $k$, the
representations $\pi,\pi'$ have respectively
\begin{equation}\label{rd-lambda-u}
 \lambda_i=u_i+i-1-\frac{k-1}{2},
 \qquad
 \lambda'_i=u'_i+i-1-\frac{k-1}{2}.
\end{equation}
In particular these numbers are integers.
The common irreducible constituents on restriction
to a common positive $(k-1)$-dimensional subspace
are indexed exactly by the tuples $\rho\in\mathcal S$
from Section~\ref{rp-section}, under
\begin{equation}\label{rd-mu-rho}
                   \rho_i=\mu_i-i+1+\frac{k-2}{2}.
\end{equation}
On such a constituent the operator $C_g(Y)$ acts by
$p_\rho(Y)$.
\end{lemma}

\begin{proof}
The highest-weight classification in
Lemma~\ref{cb-highest-classification} gives integral
dominant weights $\lambda,\lambda'$.
Corollary~\ref{rp-central-character} gives the
full Capelli eigenvalue
$\prod_{i=1}^k(w-i+1+\lambda_i)$.
Its constants $\lambda_i-i+1$ are strictly
decreasing.  The $u_i$ are also strictly decreasing
by \eqref{rp-interlacing}.  Comparison with
\eqref{rd-source-Capelli} and equality of ordered
real roots prove \eqref{rd-lambda-u}.

By Theorem~\ref{ib-compact-branching}, a highest
weight $\mu$ occurs on restriction from the left
type precisely when
$\lambda_i\geq\mu_i\geq\lambda_{i+1}$, and similarly
for the right type.  Applying \eqref{rd-lambda-u}
and \eqref{rd-mu-rho}, the simultaneous conditions
become
\[
 \max(u_{i+1},u'_{i+1})+\tfrac12
          \leq \rho_i
          \leq \min(u_i,u'_i)-\tfrac12 .
\]
Moreover
$\rho_i-\max(u_{i+1},u'_{i+1})-1/2
=\mu_i-\max(\lambda_{i+1},\lambda'_{i+1})$
is an integer.  These are exactly the finite sets
defining $\mathcal S$.
The inequalities automatically make $\mu$ dominant.
Finally the $(k-1)$-dimensional Capelli eigenvalue
at $w=Y+(k-2)/2$ is
\[
 \prod_{i=1}^{k-1}(Y+(k-2)/2-i+1+\mu_i)
                                      =p_\rho(Y).
\]
For $k=1$ the common subgroup is trivial, each
original representation is one-dimensional,
$\mathcal S$ has the empty tuple, and this last
determinant is $1$.
\end{proof}

\begin{lemma}\label{rd-common-types}
For every $\Phi\in\mathscr R$ and $g\in D$,
the map $\Phi(c,g)$ vanishes on all right
constituents whose type does not also occur on
the left.  Thus its nonzero components have
precisely the possible Capelli polynomials
$p_\rho(Y)$, $\rho\in\mathcal S$.
\end{lemma}

\begin{proof}
Let $H_g$ be the subgroup that acts arbitrarily
on $E_g$ and as the identity on $\langle y,x\rangle$.
It is isomorphic to $U(k-1)$.
For $h\in H_g$, one has $h\in K_{z'}$ and
$ghg^{-1}\in K_z$.  Evaluating
\eqref{rd-biequivariance} in these two ways gives
\[
             \Phi(c,g)\pi'(h)
                 =\pi(ghg^{-1})\Phi(c,g).
\]
Both restrictions are orthogonal sums of the
irreducibles in the compact branching theorem.
An intertwiner between two nonisomorphic
irreducibles is zero: its kernel and image
are invariant, so a nonzero such map would
be an isomorphism.  Projection onto the
orthogonal left summands therefore shows that
every noncommon right summand is annihilated.
Lemma~\ref{rd-highest-parameters} identifies the
remaining labels and eigenvalues.
\end{proof}

\subsection{Construction of the projections}

\begin{theorem}\label{rd-constructed-kappa}
There is a unique unital algebra homomorphism
\[
       \kappa:\mathbb C^{\mathcal S}
                    \longrightarrow
                         \operatorname{End}_{\mathbb C}(\mathscr R)
\]
such that
\begin{equation}\label{rd-kappa-identity}
       \kappa((p_\rho(Y))_\rho)
                =X(Y)=G^{-1}\mathcal T_{22}(Y)
                            \qquad(Y\in\mathbb C).
\end{equation}
It can be chosen as a polynomial in one actual
terminal operator.  For
$P_\rho=\kappa(1_{\{\rho\}})$, every radial
function has the decomposition
\begin{equation}\label{rd-decomposition}
 \Phi=\sum_{\rho\in\mathcal S}P_\rho\Phi,\qquad
 \mathcal T_{22}(Y)P_\rho\Phi
                        =G\,p_\rho(Y)P_\rho\Phi.
\end{equation}
At every point the summands are orthogonal in
the Hilbert--Schmidt form.  In particular
\begin{equation}\label{rd-pointwise-contraction}
 \sum_\rho\|(P_\rho\Phi)(c,g)\|_{\rm HS}^2
                      =\|\Phi(c,g)\|_{\rm HS}^2 .
\end{equation}
All projections preserve the stated smooth
equivariance and do not enlarge supports.
\end{theorem}

\begin{proof}
Choose a real $Y_0$ for which
$\xi_\rho=p_\rho(Y_0)$ are distinct, as in
Proposition~\ref{rp-idempotent-positive}.
Let
\[
 q_0(Z)=\prod_{\rho\in\mathcal S}(Z-\xi_\rho),
 \qquad
 \ell_\rho(Z)=
       \prod_{\sigma\ne\rho}
           \frac{Z-\xi_\sigma}{\xi_\rho-\xi_\sigma}.
\]
We first prove $q_0(X(Y_0))=0$ on $\mathscr R$.
At a fixed $(c,g)$, induction on the degree shows
that a polynomial $a$ obeys
\[
 (a(X(Y_0))\Phi)(c,g)
                         =\Phi(c,g)a(C_g(Y_0)).
\]
This uses \eqref{rd-actual-X} for every function
to which the next operator is applied.
Lemma~\ref{rd-common-types} shows that
$\Phi(c,g)q_0(C_g(Y_0))=0$: on common types
the factor is zero, while on other types
$\Phi(c,g)$ itself is zero.
Outside the prescribed domain all functions
are zero.  Pointwise equality proves
$q_0(X(Y_0))=0$.

In the quotient algebra $\mathbb C[Z]/(q_0)$,
the $\ell_\rho$ are mutually annihilating
idempotents and their sum is $1$.
Indeed each claimed polynomial identity
vanishes at every distinct $\xi_\rho$ and
is therefore divisible by $q_0$.
Set
\[
     P_\rho=\ell_\rho(X(Y_0)),\qquad
     \kappa(f)=\sum_\rho f(\rho)P_\rho.
\]
The same identities hold as operators, proving
that $\kappa$ is a unital algebra homomorphism.
For a right constituent of common type $\sigma$,
its evaluation satisfies
\[
            (P_\rho\Phi)(c,g)
              =\Phi(c,g)\ell_\rho(C_g(Y_0)),
 \qquad
 \ell_\rho(C_g(Y_0))|_{W_\sigma}
                          =\delta_{\rho,\sigma}\operatorname{Id}.
\]
Noncommon constituents are killed by $\Phi(c,g)$.
Consequently both sides of
\eqref{rd-kappa-identity}, evaluated at $(c,g)$,
act as $p_\sigma(Y)$ on every surviving
constituent, and agree.  This proves the identity
for all complex $Y$, without an analytic
continuation argument.

The idempotent identities now imply
\eqref{rd-decomposition}.  Since $C_g(Y_0)$ is
self-adjoint, the proof of
Proposition~\ref{rp-idempotent-positive} gives
pointwise orthogonality and
\eqref{rd-pointwise-contraction}.
Lemma~\ref{rd-smooth-matrix} shows that every
polynomial in $X(Y_0)$ acts within $\mathscr R$.
Its pointwise matrix formula also preserves
the support of the original function.

Finally any unital algebra homomorphism satisfying
\eqref{rd-kappa-identity} must send $1_{\{\rho\}}$
to $\ell_\rho(X(Y_0))$, since
$1_{\{\rho\}}=\ell_\rho((\xi_\sigma)_\sigma)$
in the finite function algebra.  This proves
uniqueness.  In particular its action agrees
with the action of the source's larger normal
form whenever that form is supplied.
\end{proof}

No estimate on the separation between the
$\xi_\rho$ is needed.  They are used to construct
the projections, but the resulting estimate
\eqref{rd-pointwise-contraction} has constant one.
In particular there is no hidden dependence on
the size of the type set or on the interpolation
denominators in this estimate.

\begin{corollary}\label{rd-integrated-projections}
Let $\nu$ be any positive measure on
$\Omega\times D$.  Whenever the functions in
question are square integrable, the projections
$P_\rho$ are mutually orthogonal contractions
for the Hilbert--Schmidt $L^2(\nu)$ norm, and
\[
 \|\Phi\|_{L^2(\nu)}^2
                  =\sum_\rho\|P_\rho\Phi\|_{L^2(\nu)}^2.
\]
For $f\in\mathbb C^{\mathcal S}$,
\[
 \|\kappa(f)\Phi\|_{L^2(\nu)}
              \leq \max_\rho|f(\rho)|\,\|\Phi\|_{L^2(\nu)}.
\]
The same assertions hold after inserting any
nonnegative measurable scalar weight into the
measure.  Each $P_\rho$ commutes with scalar
multiplication and with differentiation in
the variables $c\in\Omega$.
\end{corollary}

\begin{proof}
Integrate \eqref{rd-pointwise-contraction};
the sum is finite.  More generally, its
orthogonal components give
$\|\kappa(f)\Phi(c,g)\|_{\rm HS}^2
=\sum_\rho|f(\rho)|^2
\|(P_\rho\Phi)(c,g)\|_{\rm HS}^2$.
Integration proves the bound.
The projections are pointwise right
multiplications by matrices depending only
on $g$, which proves the last assertions
whenever both sides are defined on the
smooth radial space.
\end{proof}

\begin{corollary}\label{rd-actual-positive-density}
Under the actual compact-type and spectral-margin
hypotheses, the terminal density
\eqref{rp-terminal-density} is nonnegative
without an additional normal-form premise.
For a fixed $c$, it is integrable against Haar
measure on $U(V,h)$ if all its finitely many
functions $\Psi_e(c,\cdot)$ have compact
support in $D$.
\end{corollary}

\begin{proof}
Theorem~\ref{rd-constructed-kappa} supplies
exactly \eqref{rp-required-normal-form}, so
Corollary~\ref{rp-terminal-density-positive}
applies.  On the finite union of the compact
supports, $G$ has a positive minimum and
$|\zeta|^2/G$ is bounded.  The functions,
their projections, and all displayed
coefficients are continuous there.
The density is therefore bounded and
compactly supported.  Haar measure is
finite on compact sets, proving
integrability.
\end{proof}

The construction proves the entire finite
spectral decomposition needed for terminal
positivity, not just a conditional implication
from the larger normal form.  It also proves
an exact contraction estimate for arbitrary
positive measures.  The remaining analytic
task is to identify the original integrated
differential forms with these densities on
the source's generated function classes,
including their noncompact members and
successive radial integrations.  Compact
support at one fixed value of $c$ is not
asserted for that whole class.


\section{Noncompact integration of the radial derivative terms}
\label{rf-section}

The remaining radial integration takes place on the whole complex
plane.  Both the point $t=0$ and infinity must therefore be treated.
We prove the Fourier regularization and the actual elementary
differential identities used in this passage
\cite{alpo-uniform-source}.  Two simplifications are useful.
Only the Fourier transform of the weighted density is needed.
Furthermore, for an admissible word the differential expansion
can be written explicitly as a sum over subsets.
The estimates for the actual Jacquet integrals needed to apply
these facts are identified separately at the end.

\subsection{Fourier normalization and extension across zero}

We identify $\mathbb C$ with $\mathbb R^2$ and use the measure
$dt=dx\,dy$ for $t=x+iy$.  The Fourier transform in this section is
\begin{equation}\label{rf-transform}
 \widehat E(s)=\int_{\mathbb C}
        \exp(st-\overline s\,\overline t)\,E(t)\,dt.
\end{equation}
For $s=a+ib$, the exponential in \eqref{rf-transform} is
$\exp(2i(bx+ay))$.  In particular its absolute value is one.
All derivatives below are real derivatives.  The norm
$\|D^jF(t)\|$ is the operator norm of the $j$-linear derivative
on the Euclidean space $\mathbb R^2$; for $j=0$ it is $|F(t)|$.

\begin{lemma}\label{rf-inversion}
Let $E\in L^1(\mathbb C)$ and suppose
$\widehat E\in L^1(\mathbb C)$.  Then
\begin{equation}\label{rf-inverse}
 E_0(t)=\frac1{\pi^2}\int_{\mathbb C}
       \exp(-st+\overline s\,\overline t)\widehat E(s)\,ds
\end{equation}
is continuous and equals $E$ almost everywhere.
It equals $E$ at every point at which $E$ is continuous.
If $|s|^j\widehat E(s)$ is integrable for $0\leq j\leq m$,
then $E_0$ is $C^m$, and differentiation under the integral is valid.
\end{lemma}

\begin{proof}
The absolute convergence and continuity of \eqref{rf-inverse}
follow from dominated convergence.  To fix the normalization
and prove the assertion about $E$, put
\[
 \eta_a(t)=\frac1{\pi a}e^{-|t|^2/a}\qquad(a>0).
\]
The one-dimensional Gaussian integral, applied separately to
the two real coordinates of $s$, gives
\[
 \widehat{\eta_a}(s)=e^{-a|s|^2},\qquad
 \frac1{\pi^2}\int_{\mathbb C}e^{-a|s|^2}
        e^{st-\overline s\,\overline t}\,ds=\eta_a(t).
\]
Fubini's theorem is applicable because $E$ is integrable and
the Gaussian in $s$ is integrable.  It therefore gives
\[
 (E*\eta_a)(t)=\frac1{\pi^2}\int_{\mathbb C}
      e^{-st+\overline s\,\overline t}
      e^{-a|s|^2}\widehat E(s)\,ds.
\]
The right side tends to $E_0(t)$ uniformly in $t$ as $a\to0$.

On the other hand, $E*\eta_a\to E$ in $L^1$.
Indeed,
\[
 \|E*\eta_a-E\|_1
 \leq\int_{\mathbb C}\eta_a(v)\,
                    \|E(\,\cdot-v)-E\|_1\,dv.
\]
Translations are continuous in $L^1$: this follows first for
continuous functions of compact support by uniform continuity,
and then for an arbitrary integrable function by approximation
in $L^1$.  Split the last integral into a fixed small disk
and its complement, use this translation continuity on the
disk and the bound $2\|E\|_1$ on the complement.
The Gaussian mass of the complement tends to zero.
An $L^1$-convergent subsequence converges almost everywhere,
so the uniform limit $E_0$ equals $E$ almost everywhere.
At a point of continuity of $E$, their difference is continuous.
If its value were nonzero, its absolute value would be bounded
away from zero throughout a small disk, contradicting equality
almost everywhere.

For a unit real direction $v$, differentiating the exponential
in \eqref{rf-inverse} gives the factor
$-sv+\overline s\,\overline v$, whose absolute value is at most
$2|s|$.  The stated moment assumptions permit repeated
differentiation by dominated convergence and prove the last
assertion.
\end{proof}

\begin{theorem}\label{rf-weighted-extension}
Let $m\geq0$ be an integer and let $D:\mathbb C\to\mathbb C$
be measurable and continuous on $\mathbb C^\times$.
Put
\[
                   E(t)=(1+|t|^2)^2D(t).
\]
Suppose that $E\in L^1(\mathbb C)$ and, for some $C\geq0$,
\begin{equation}\label{rf-weighted-decay}
             |\widehat E(s)|\leq C(1+|s|)^{-m-3}
                      \qquad(s\in\mathbb C).
\end{equation}
Then there is a unique $C^m$ function $D_0$ on $\mathbb C$
equal to $D$ on $\mathbb C^\times$.  It satisfies
\begin{align}
 \int_{\mathbb C}D_0(t)\,dt
      &=\int_{\mathbb C}D(t)\,dt, \label{rf-same-integral}\\
 \|D^jD_0(t)\|&\leq A_m C(1+|t|)^{-4}
       \qquad(0\leq j\leq m), \label{rf-all-derivatives}\\
 \int_{\mathbb C}D^jD_0(t)\,dt&=0
       \qquad(1\leq j\leq m). \label{rf-zero-derivatives}
\end{align}
In the last line the integral is the componentwise integral
of the multilinear derivative.  The constant $A_m$ depends
only on $m$, and the following completely explicit choice is valid:
\begin{align*}
 B_0&=1,\qquad B_j=\prod_{\nu=0}^{j-1}(5\nu+4),\\
 M_{m,j}&=\frac{2^{j+1}}{\pi}
                   \frac{(j+1)!(m-j)!}{(m+2)!}
                       \qquad(0\leq j\leq m),\\
 A_m&=\max_{0\leq j\leq m}
       \sum_{\nu=0}^j\binom j\nu
                    16\,2^{\nu/2}B_\nu M_{m,j-\nu}.
\end{align*}
No separate estimate for $\widehat D$ is required.
\end{theorem}

\begin{proof}
The hypothesis $E\in L^1$ implies $D\in L^1$ because
$(1+|t|^2)^{-2}\leq1$.  For $0\leq j\leq m$, polar coordinates
and the substitution $r=u/(1-u)$ give
\begin{align*}
 \int_{\mathbb C}|s|^j(1+|s|)^{-m-3}\,ds
 &=2\pi\int_0^\infty r^{j+1}(1+r)^{-m-3}\,dr\\
 &=2\pi\frac{(j+1)!(m-j)!}{(m+2)!}.
\end{align*}
The last integral can also be evaluated by repeated integration
by parts; its endpoint terms vanish since $j\leq m$.
Lemma~\ref{rf-inversion} supplies a $C^m$ extension $E_0$ and gives
\begin{equation}\label{rf-E-derivative}
                       \|D^jE_0(t)\|\leq M_{m,j}C.
\end{equation}
Since $E$ is continuous away from zero, $E_0=E$ there.
Define $D_0(t)=(1+|t|^2)^{-2}E_0(t)$.  This proves the asserted
extension and uniqueness, because $\mathbb C^\times$ is dense.

For clarity, a uniform derivative bound for the weight can
be obtained without an unspecified differentiation constant.
For every coordinate multi-index $\alpha$ with $|\alpha|=j$,
\[
 \partial^\alpha(1+x^2+y^2)^{-2}
     =\frac{P_\alpha(x,y)}{(1+x^2+y^2)^{2+j}},
 \qquad \deg P_\alpha\leq j.
\]
The sum of the absolute values of the coefficients of
$P_\alpha$ is at most $B_j$.  To see this, differentiating once
replaces $P$ by
\[
  (1+x^2+y^2)\partial_iP-(4+2j)x_iP.
\]
The coefficient sum increases by a factor at most
$3j+4+2j=5j+4$.  This proves the bound by induction.
On $|t|\leq1$ and $|t|\geq1$, respectively, the displayed
rational expression then gives
\[
 |\partial^\alpha(1+|t|^2)^{-2}|
                      \leq16B_j(1+|t|)^{-4}.
\]
Passing from coordinate derivatives to the operator norm
costs at most $2^{j/2}$, since
$|v_x|+|v_y|\leq\sqrt2\,|v|$.
The multilinear product rule, together with
\eqref{rf-E-derivative}, proves \eqref{rf-all-derivatives}
with the stated $A_m$.

These bounds imply integrability of $D_0$ and all of its
derivatives through order $m$.  Equality almost everywhere
with $D$ proves \eqref{rf-same-integral}.
To prove \eqref{rf-zero-derivatives}, it is enough to consider
a coordinate derivative $\partial_i H$, where
$H=\partial^\alpha D_0$ has order at most $m-1$.
Choose a smooth compactly supported function $\chi$ equal to
one on the unit disk and put $\chi_R(t)=\chi(t/R)$.
Integration by parts on its compact support gives
\[
 \int_{\mathbb C}\chi_R\,\partial_i H
        =-\int_{\mathbb C}H\,\partial_i\chi_R.
\]
The right side has absolute value at most
$R^{-1}\|\partial_i\chi\|_\infty\|H\|_1$, hence tends to zero.
Dominated convergence on the left gives
$\int\partial_iH=0$.  The argument applies to every component
of every positive-order derivative.  It treats zero and
infinity simultaneously, without discarding an annular
boundary term.
\end{proof}

\begin{corollary}\label{rf-finite-combination}
Let $D_1,\ldots,D_N$ satisfy
Theorem~\ref{rf-weighted-extension} with a common order $m$,
and constants $C_i$ that may depend on $i$.  Let $0\leq j_i\leq m$,
$a_i\in\mathbb C$ and fixed real directions
$v_{i,1},\ldots,v_{i,j_i}\in\mathbb C$.
If, for $t\ne0$,
\[
 f(t)=\sum_{i=1}^N a_i
      D^{j_i}D_i(t)[v_{i,1},\ldots,v_{i,j_i}],
\]
then $f$ extends continuously to $\mathbb C$, its extension
is $O((1+|t|)^{-4})$, and
\[
        \int_{\mathbb C}f(t)\,dt
           =\sum_{\{i:j_i=0\}}a_i\int_{\mathbb C}D_i(t)\,dt.
\]
The implicit constant can be taken to be
$A_m\sum_i |a_i|C_i\prod_{\nu=1}^{j_i}|v_{i,\nu}|$.
The assertion includes $N=0$.
\end{corollary}

\begin{proof}
Replace $D_i$ by its extension from the theorem.
They agree on an open neighborhood of every nonzero point,
so their derivatives agree there.  The finite sum of the
extended derivatives is continuous and has the stated
bound.  Integrate term by term and use
\eqref{rf-same-integral}--\eqref{rf-zero-derivatives}.
\end{proof}

\begin{example}\label{rf-critical-power}
The extra integrable Fourier moment cannot simply be omitted.
For $t\ne0$ set
\[
 E(t)=\int_0^1\frac1{\pi a}e^{-|t|^2/a}\,da,\qquad
 D(t)=\frac{E(t)}{(1+|t|^2)^2},
\]
and give both functions any finite values at zero.
Tonelli's theorem gives $\|E\|_1=1$, and differentiation
under the integral shows that $E$ is smooth away from zero.
Fubini's theorem and the Gaussian calculation above give
\[
 \widehat E(s)=
 \begin{cases}
 (1-e^{-|s|^2})/|s|^2,&s\ne0,\\
 1,&s=0.
 \end{cases}
\]
Thus $|\widehat E(s)|\leq4(1+|s|)^{-2}$.
For $0<|t|<1$, restricting the integral to
$|t|^2\leq a\leq1$ gives
$E(t)\geq(\pi e)^{-1}\log(|t|^{-2})$.
Consequently $D$ has no continuous extension at zero.
This is a counterexample to replacing the power $-3$
by $-2$ in the case $m=0$; the source uses the sufficient
power $-m-3$.
\end{example}

\subsection{The actual two-component differential identity}

Let $\mathcal H$ be a complex Hilbert space, with inner product
linear in the first argument.  Write
\[
 \partial_t=\tfrac12(\partial_x-i\partial_y),\qquad
 \partial_{\overline t}=\tfrac12(\partial_x+i\partial_y),
 \qquad p=t\partial_t,\quad
 \overline p=\overline t\,\partial_{\overline t}.
\]
Here $\overline p$ denotes this displayed differential
operator, not an instruction to conjugate its argument.
For $u,v:\mathbb C^\times\to\mathcal H^2$ put
\[
 D(u,v)(t)=|t|^{-2}
              \bigl(\langle u_0(t),v_0(t)\rangle+
                    \langle u_1(t),v_1(t)\rangle\bigr).
\]
For a real number $Y$, define
\begin{align}
 V_+(Y)v&=(Yv_0+pv_0-tv_1,\ t^{-1}v_0),\nonumber\\
 I_+(Y)u&=(-\overline t^{-1}u_1,\
             \overline t\,u_0+(Y+1)u_1-\overline p\,u_1),
                    \label{rf-plus-operators}\\
 V_-(Y)v&=(Yv_0-\overline p\,v_0-\overline t\,v_1,\
                                      \overline t^{-1}v_0),\nonumber\\
 I_-(Y)u&=(-t^{-1}u_1,\
                    t u_0+(Y-1)u_1+p u_1).
                    \label{rf-minus-operators}
\end{align}
These are the source's actual one-coordinate operators, including
their shifts by $1$.  Let $Ev=(0,v_0)$.

\begin{lemma}\label{rf-one-letter}
For $C^1$ Hilbert-space-valued functions $u,v$ on
$\mathbb C^\times$ and real $Y$,
\begin{align}
 D(I_+(Y)u,V_+(Y)v)
      &=D(u,v)-\partial_{\overline t}D(u,Ev),\label{rf-plus-density}\\
 D(I_-(Y)u,V_-(Y)v)
      &=D(u,v)+\partial_tD(u,Ev).\label{rf-minus-density}
\end{align}
\end{lemma}

\begin{proof}
Put $A(t)=\langle u_1(t),v_0(t)\rangle$.
Direct expansion of \eqref{rf-plus-operators}, using
linearity in the first argument and $Y=\overline Y$, gives
\[
 \langle I_+u,V_+v\rangle
 =\langle u,v\rangle+\overline t^{-1}A
       -\langle\overline t^{-1}\overline p\,u_1,v_0\rangle
       -\langle u_1,t^{-1}p v_0\rangle.
\]
The last two terms sum to $-\partial_{\overline t}A$:
\[
 \partial_{\overline t}A
    =\langle\partial_{\overline t}u_1,v_0\rangle
       +\langle u_1,\partial_t v_0\rangle.
\]
Since
$\partial_{\overline t}|t|^{-2}
       =-\overline t^{-1}|t|^{-2}$,
multiplication by $|t|^{-2}$ proves \eqref{rf-plus-density}.

The analogous expansion of \eqref{rf-minus-operators} is
\[
 \langle I_-u,V_-v\rangle
 =\langle u,v\rangle-t^{-1}A
       +\langle t^{-1}p u_1,v_0\rangle
       +\langle u_1,\overline t^{-1}\overline p\,v_0\rangle.
\]
The last two terms sum to $\partial_tA$.
Using $\partial_t|t|^{-2}=-t^{-1}|t|^{-2}$ proves
\eqref{rf-minus-density}.  All calculations are valid in a
Hilbert space by the differentiability of its inner product.
\end{proof}

\subsection{An explicit expansion for an admissible word}

Let $S$ be a finite set and let
$\mathcal B=\{0,1\}^S$.  Replace $\mathcal H^2$ by
$\mathcal H^{\mathcal B}$ and sum its inner product over
$\mathcal B$.  For $a\in S$, each of the operators
\eqref{rf-plus-operators}--\eqref{rf-minus-operators} acts
on the $a$-th two-component factor and as the identity
on the other factors.  Denote by $E_a$ the corresponding
elementary matrix $E_{10}$:
\[
 (E_a v)_b=
 \begin{cases}
 v_{b[a\leftarrow0]},& b(a)=1,\\
 0,&b(a)=0.
 \end{cases}
\]

A word of length $h$ is a sequence
$\ell_j=(a_j,\sigma_j,Y_j)$, where
$a_j\in S$, $\sigma_j\in\{+,-\}$ and $Y_j\in\mathbb R$.
Assume that the $a_j$ are pairwise distinct.
This is the part of the source's admissibility condition
used in the present identity.  Define the ordered products
\[
 I_w=I_{\sigma_1,a_1}(Y_1)\cdots I_{\sigma_h,a_h}(Y_h),
 \qquad
 V_w=V_{\sigma_1,a_1}(Y_1)\cdots V_{\sigma_h,a_h}(Y_h).
\]
The rightmost factor acts first.  Put
\[
 (\epsilon_j,\delta_j)=
 \begin{cases}
 (-1,\partial_{\overline t}),&\sigma_j=+,\\
 (1,\partial_t),&\sigma_j=-.
 \end{cases}
\]
For $J\subseteq\{1,\ldots,h\}$, let
$E_J=\prod_{j\in J}E_{a_j}$,
$\epsilon_J=\prod_{j\in J}\epsilon_j$, and
$\delta_J=\prod_{j\in J}\delta_j$.
The factors in these three products commute.
For the empty set they are the identity, $1$, and the
identity differential operator, respectively.

\begin{proposition}\label{rf-subset-expansion}
For smooth $u,v:\mathbb C^\times\to\mathcal H^{\mathcal B}$,
\begin{equation}\label{rf-explicit-expansion}
             D(I_wu,V_wv)
       =\sum_{J\subseteq\{1,\ldots,h\}}
                         \epsilon_J\delta_JD(u,E_Jv).
\end{equation}
There are at most $2^h$ terms, their coefficients have
absolute value one, and the derivative order of the
$J$-term is $|J|$.
\end{proposition}

\begin{proof}
The case $h=0$ is immediate.  Write $w=\ell_1w'$.
Sum Lemma~\ref{rf-one-letter} over the other component
indices to obtain
\[
 D(I_wu,V_wv)=D(I_{w'}u,V_{w'}v)
     +\epsilon_1\delta_1
          D(I_{w'}u,E_{a_1}V_{w'}v).
\]
Since $a_1$ does not occur in $w'$, $E_{a_1}$ commutes
with every constant component matrix in $V_{w'}$.
It also commutes with multiplication by $t,\overline t$
and with their derivatives.  Hence
$E_{a_1}V_{w'}=V_{w'}E_{a_1}$.
Apply the induction hypothesis first to $(u,v)$ and
then to $(u,E_{a_1}v)$.  The resulting terms correspond,
respectively, to subsets not containing $1$ and to
subsets containing $1$.  Smoothness permits the
indicated differentiations of these identities.
\end{proof}

\begin{theorem}\label{rf-noncompact-step}
In Proposition~\ref{rf-subset-expansion}, suppose for every
$J\subseteq\{1,\ldots,h\}$ that
\[
 E^{(J)}(t)=(1+|t|^2)^2D(u,E_Jv)(t)
\]
is integrable and
\[
             |\widehat {E^{(J)}}(s)|
                    \leq C_J(1+|s|)^{-h-3}
                    \qquad(s\in\mathbb C)
\]
for some $C_J\geq0$.
Then $D(I_wu,V_wv)$ has a continuous integrable extension
to the whole plane, bounded in absolute value by
$A_h(\sum_J C_J)(1+|t|)^{-4}$, and
\begin{equation}\label{rf-full-plane-identity}
       \int_{\mathbb C}D(I_wu,V_wv)(t)\,dt
                       =\int_{\mathbb C}D(u,v)(t)\,dt.
\end{equation}
No support condition in $t$ is imposed.
\end{theorem}

\begin{proof}
Apply Theorem~\ref{rf-weighted-extension} to the finitely
many functions $D(u,E_Jv)$, with $m=h$.
An iterated Wirtinger derivative of order $j$ is a sum
of $2^j$ coordinate derivatives, each with coefficient
of absolute value $2^{-j}$.  Thus \eqref{rf-all-derivatives}
bounds it by the same constant, and its integral is
zero for $j\geq1$ by \eqref{rf-zero-derivatives}.
Insert these extended functions in
\eqref{rf-explicit-expansion}.  The empty subset contributes
$D(u,v)$; every other subset integrates to zero.
\end{proof}

\begin{example}\label{rf-repeated-index}
The distinct-index assumption in
Proposition~\ref{rf-subset-expansion} is essential.
Take $\mathcal H=\mathbb C$, one two-component factor,
two identical plus letters with $Y=0$, and
$u=(0,1)$, $v=(1,0)$.
Then $I_+^2u=(-\overline t^{-1},0)$ and
$V_+^2v=(-1,0)$, so
\[
                D(I_+^2u,V_+^2v)=\frac1{t\overline t^2}.
\]
The proposed subset expression would instead be
$-2\partial_{\overline t}|t|^{-2}
       =2/(t\overline t^2)$, since $E^2=0$.
The source's admissible words have distinct indices,
so this does not contradict its statement.
\end{example}

\subsection{Relation to the radial form and remaining analysis}

The measure $|t|^{-2}dt$ explains the shifts in the
operators.  On compactly supported smooth functions
away from zero, integration by parts gives
\[
       p^*=-\overline p,\qquad
       \overline p^{\,*}=-p,\qquad t^*=\overline t.
\]
These formulas define the formal adjoint of every
matrix differential operator considered here.

\begin{lemma}\label{rf-inverse-adjoint}
As differential operators on $\mathbb C^\times$,
\[
             I_\sigma(Y)V_\sigma(Y)^*
             =V_\sigma(Y)^*I_\sigma(Y)=1
                      \qquad(Y\in\mathbb R).
\]
Consequently $I_wV_w^*=1$ for an ordered word.
If matrix differential operators $R,R'$ satisfy
$R=V_w^*R'V_w$, then
\begin{equation}\label{rf-mirror-identity}
                         R'V_w=I_wR.
\end{equation}
\end{lemma}

\begin{proof}
In the plus case,
\[
 V_+(Y)^*=
  \begin{pmatrix}Y-\overline p&\overline t^{-1}\\
                          -\overline t&0\end{pmatrix},
 \qquad
 I_+(Y)=
  \begin{pmatrix}0&-\overline t^{-1}\\
                  \overline t&Y+1-\overline p\end{pmatrix}.
\]
Their products are the identity because
$[\overline p,\overline t]=\overline t$ and
$[\overline p,\overline t^{-1}]=-\overline t^{-1}$.
In the minus case the corresponding matrices are
\[
 V_-(Y)^*=
  \begin{pmatrix}Y+p&t^{-1}\\-t&0\end{pmatrix},
 \qquad
 I_-(Y)=
  \begin{pmatrix}0&-t^{-1}\\t&Y-1+p\end{pmatrix},
\]
and the calculation uses $[p,t]=t$ and
$[p,t^{-1}]=-t^{-1}$.
In the word product $I_wV_w^*$ these identities cancel
adjacent factors, starting with the last letter.
Multiplying $R=V_w^*R'V_w$ on the left by $I_w$
proves \eqref{rf-mirror-identity}.
These are algebraic identities; they assert no
unproved adjoint-domain statement on a noncompact space.
\end{proof}

\begin{lemma}\label{rf-L2-smoothness}
Let $(Z,\nu)$ be a $\sigma$-finite measure space and let
$F:\mathbb C^\times\times Z\to\mathbb C^d$ be measurable.
Suppose that, outside one null subset of $Z$, it is smooth
in $t$ and all its coordinate derivatives are jointly
measurable.  Suppose also that, for every compact
$K\subset\mathbb C^\times$ and every multi-index $\alpha$,
\begin{equation}\label{rf-local-L2-bounds}
       \sup_{t\in K}
          \|\partial^\alpha_tF(t,\cdot)\|_{L^2(Z,\nu)}
                                                     <\infty.
\end{equation}
Then $t\mapsto F(t,\cdot)$ is a smooth map into
$L^2(Z,\nu;\mathbb C^d)$, with the expected derivatives.
A pointwise dominating $L^2$ function is not required.
\end{lemma}

\begin{proof}
Fix a closed disk $K$ inside $\mathbb C^\times$ and a
smaller concentric disk.  Every sufficiently short
segment from the smaller disk lies in $K$.
The fundamental theorem of calculus, applied outside
the common null set, gives
\[
 F(t+h,z)-F(t,z)
       =\int_0^1\bigl(h_x\partial_xF+h_y\partial_yF\bigr)
                                                    (t+uh,z)\,du.
\]
Cauchy--Schwarz in $u$, followed by Tonelli's theorem,
or equivalently the integral triangle inequality in
$L^2$, bounds the $L^2$ norm of the difference by
$\sqrt2|h|$ times a common bound for the two first
derivatives on $K$.  This proves norm continuity.
Applying the same argument to every
$\partial^\alpha F$ proves their norm continuity.

The second-order Taylor formula along that segment
gives
\[
 F(t+h,\cdot)-F(t,\cdot)
       -h_x\partial_xF(t,\cdot)-h_y\partial_yF(t,\cdot)
 =\int_0^1(1-u)
       \sum_{i,j\in\{x,y\}}h_i h_j\partial_{ij}F(t+uh,\cdot)\,du.
\]
If $B$ bounds the $L^2$ norms of the second derivatives
on $K$, the norm of the right side is at most
$B(|h_x|+|h_y|)^2/2\leq B|h|^2$.
Thus the derivative exists in $L^2$ and has the asserted
value.  Repeating this proof for each
$\partial^\alpha F$ proves smoothness of all orders.
The same estimates justify every integral displayed
above; no assumption of finite measure on $Z$ is used.
\end{proof}

For the source radial form, $\mathcal H$ is the $L^2$ space
of the remaining radial coordinates and the terminal
unitary group, with its specified positive measure
and Hilbert--Schmidt coefficient space.  Its measure
is independent of the coordinate $t$ being integrated.
When the radial functions and their derivatives are
smooth as maps into this Hilbert space, all the preceding
identities apply directly.  Lemma~\ref{rf-L2-smoothness}
shows that it suffices to bound their $L^2$ derivative
norms locally uniformly.  The source's coordinate
operators include multiplication by $t^{\pm1}$ and
$\overline t^{\pm1}$ together with $p,\overline p$.
On a compact subset of $\mathbb C^\times$, every
coordinate derivative of order $j$ is a finite sum
of words in these operators, of length at most $2j$.
For example $\partial_t=t^{-1}p$ and
$\partial_{\overline t}=\overline t^{-1}\overline p$,
and one repeats these identities by composition.
Consequently locally uniform $L^2$ bounds for all
coordinate-operator words imply
\eqref{rf-local-L2-bounds}.  This also justifies
differentiating all the integrated pairings in
Lemma~\ref{rf-one-letter}, without constructing a
separate pointwise dominating function.

With $u=Rx$ and $v=y$, identity
\eqref{rf-mirror-identity} and Theorem~\ref{rf-noncompact-step}
give
\[
 \int_{\mathbb C}D(R'V_wx,V_wy)(t)\,dt
                  =\int_{\mathbb C}D(Rx,y)(t)\,dt,
\]
provided the displayed weighted Fourier bounds hold for
the densities $D(Rx,E_Jy)$.
This is precisely the one-coordinate integration step
after the source's algebraic relation
$R=V_w^*R'V_w$ has been established.
The argument uses the actual shifts and the full
noncompact coordinate domain.

The present section proves the elementary differential
expansion, the inverse formal-adjoint identity, and the
analytic removal of every derivative term.  It does
not yet prove that all actual Jacquet-generated functions
satisfy the required $L^2$ differentiability and weighted
Fourier estimates, nor does it identify the preceding
algebraic relation with all of the source's recursively
defined frame operators.  The source obtains the
Fourier estimates from four specified differential
operators preserving the class of coefficients and
from an identity that transfers multiplication by the
Fourier parameter to those operators.  Their construction
and the precise Jacquet pairing and measure identities
remain substantive dependencies.  In particular,
pointwise terminal positivity alone is not being used
as an integrated noncompact theorem.


\section{A regular coefficient for the inverse terminal operator}
\label{rl-section}

We next reconstruct an actual full-group coefficient used in the
four-operator Fourier argument.  Radial coordinates involve ratios
of consecutive minors, which may be undefined inside the relevant
open subset of the full group.  The construction below has only the
positive Gram polynomial in its denominator.  It gives the
zero-order spectral factor used in the source
\cite{alpo-uniform-source}, including covariance, growth of every
derivative, and preservation of the coefficient class.

\subsection{The Hermitian space, minors and cofactor vector}

Let $k\geq1$ and $r\geq0$.  Fix a Witt decomposition
\[
 V=U\oplus W\oplus U^\vee,\qquad
 U=\langle e_1,\ldots,e_r\rangle,\quad
 U^\vee=\langle e_1^\vee,\ldots,e_r^\vee\rangle,
\]
for a Hermitian form $h$, linear in its first argument.
The subspaces $U,U^\vee$ are isotropic,
$h(e_i,e_j^\vee)=\delta_{ij}$, and $W$ is perpendicular
to their sum and has signature $(k,1)$.
Choose a basis of $W$ in which the form is
$\operatorname{diag}(1,\ldots,1,-1)$.
In the resulting basis of $V$ its matrix $H$ satisfies
$H^*=H$ and $H^2=1$.
Write $\mathbf G=U(V,h)$.

Fix $z,z'\in W$ with
\[
                 h(z,z)=-T,\qquad h(z',z')=-T',
                         \qquad T,T'>0.
\]
Put $f_i=e_i$, $f'_i=e_i$ for $1\leq i\leq r$, and
$f_{r+1}=z$, $f'_{r+1}=z'$.
For $g\in\mathbf G$ define
\begin{align}
 Z(g)_{ab}&=h(gf'_b,f_a),\qquad 1\leq a,b\leq r+1,
                                      \label{rl-Z}\\
 D(g)&=\det Z(g),\qquad
 d(g)=\det Z(g)_{\{1,\ldots,r\},\{1,\ldots,r\}},\nonumber\\
 \mathcal G(g)&=|D(g)|^2-TT'|d(g)|^2. \label{rl-Gram}
\end{align}
When $r=0$, $d(g)=1$.
These are the last two minors and the actual Gram polynomial
in the source, before passage to radial coordinates.
Define
\begin{equation}\label{rl-cofactor-vector}
 L(g)=\sum_{a=1}^{r+1}
        \overline{\operatorname{adj}(Z(g))_{r+1,a}}\,f_a,
 \qquad y(g)=g^{-1}L(g).
\end{equation}
Let $\omega(g)$ be the projection of the $W$-component
of $y(g)$ onto the positive space
$W'_+=W\cap(z')^\perp$.
Fix an orthonormal basis $B_1,\ldots,B_k$ of $W'_+$ and put
\[
                 \Omega_j(g)=h(y(g),B_j).
\]
Thus $\omega(g)=\sum_j\Omega_j(g)B_j$.

\begin{lemma}\label{rl-cofactor-norm}
For every $g\in\mathbf G$,
\begin{align}
 h(y(g),f'_b)&=\overline{D(g)}\,\delta_{b,r+1},
                                                  \label{rl-dual-pairing}\\
 h(y(g),y(g))&=-T|d(g)|^2,\label{rl-y-norm}\\
 \sum_{j=1}^k|\Omega_j(g)|^2
             &=\frac{\mathcal G(g)}{T'}.\label{rl-omega-norm}
\end{align}
In particular $\mathcal G(g)\geq0$.
Each $\Omega_j$ is the restriction to $\mathbf G$ of a
polynomial of degree $r+1$ in the conjugates of the
matrix entries of $g$.
No minor is required to be nonzero in these assertions.
\end{lemma}

\begin{proof}
Since $g$ is unitary,
\[
 h(y,f'_b)=h(L,gf'_b)
   =\sum_a
     \overline{\operatorname{adj}(Z)_{r+1,a}}\,
     \overline{Z_{ab}}
   =\overline{D}\,\delta_{b,r+1}.
\]
This is the adjugate identity, valid also when $Z$ is singular.
All $f_a$ with $a\leq r$ are mutually perpendicular
and perpendicular to $z$, so
\[
 h(y,y)=h(L,L)
       =-T|\operatorname{adj}(Z)_{r+1,r+1}|^2
       =-T|d|^2.
\]
The first $r$ equations in \eqref{rl-dual-pairing}
say that $y\in U^\perp=U\oplus W$.
Its component along $z'$ is
$-\overline D\,z'/T'$.  Consequently
\[
           y=u-\frac{\overline D}{T'}z'+\omega
                          \qquad(u\in U,\ \omega\in W'_+).
\]
Taking its norm gives
$-T|d|^2=-|D|^2/T'+h(\omega,\omega)$,
which proves \eqref{rl-omega-norm}.

For the polynomial assertion, $Z$ is linear in the
entries of $g$, so every cofactor has degree $r$.
On $\mathbf G$ one has $g^{-1}=Hg^*H$.
Replacing $g^{-1}$ by this formula in
\eqref{rl-cofactor-vector} gives a polynomial of
degree $r+1$ in the conjugate entries.
Pairing with the fixed $B_j$ is linear and does not
change that degree.
\end{proof}

\subsection{Covariance, including the unipotent subgroups}

Let $N\subset\mathbf G$ consist of the elements satisfying
\[
 n e_i-e_i\in\langle e_1,\ldots,e_{i-1}\rangle,
       \qquad n w-w\in U\quad(w\in W).
\]
This is the source's upper unipotent subgroup in the
chosen Witt basis.  It preserves $U^\perp$ and acts
as the identity on $U^\perp/U$.  The same assertions
hold for its inverse, by triangular inversion.
Let $K_z$ and $K_{z'}$ fix every $e_i,e_i^\vee$
and fix $z$ and $z'$, respectively.  They act on
$W$ and are the compact groups isomorphic to $U(k)$.

\begin{lemma}\label{rl-cofactor-covariance}
Let $h_L\in K_z$, $h_R\in K_{z'}$ and
$n_L,n_R\in N$, and set
\[
                  g'=h_L^{-1}n_L^{-1}g n_Rh_R.
\]
Every leading principal minor of $Z(g')$ equals the
corresponding minor of $Z(g)$.  Moreover
\begin{equation}\label{rl-omega-covariance}
                \omega(g')=h_R^{-1}\omega(g),
       \qquad \mathcal G(g')=\mathcal G(g).
\end{equation}
These statements hold also when intermediate minors vanish.
\end{lemma}

\begin{proof}
Write $f$ for the matrix with columns $f_a$, and
similarly $f'$ for $f'_b$.  There are upper triangular
unipotent matrices $A,B$ of size $r+1$ with
$n_Lf=fA$ and $n_Rf'=f'B$.
Since the compact factors fix these columns,
unitarity gives
\[
                         Z(g')=A^*Z(g)B.
\]
For a leading block, upper triangularity of $A,B$
means that multiplication uses only the corresponding
leading blocks.  Their determinants are one.
This proves the minor identities, and hence the
identity for $\mathcal G$.

For completeness, the cofactor transformation can be
obtained without assuming that $Z$ is invertible.
The identity
\[
 \operatorname{adj}(A^* Z B)
       =B^{-1}\operatorname{adj}(Z)(A^*)^{-1}
\]
holds for invertible $Z$ by the inverse formula;
both sides are polynomial in the entries of $Z$,
so it holds for every $Z$ by continuity and the
density of invertible matrices.  The last row of
$B^{-1}$ is the last coordinate row.  Transposing
and conjugating that row identity in
\eqref{rl-cofactor-vector} yields
\[
                 L(g')=n_L^{-1}L(g)
                       =h_L^{-1}n_L^{-1}L(g).
\]
The second equality holds because
$n_L^{-1}L(g)$ lies in $\langle f_1,\ldots,f_{r+1}\rangle$,
which $h_L$ fixes pointwise.  Thus
$y(g')=h_R^{-1}n_R^{-1}y(g)$.
The vector $y(g)$ belongs to $U^\perp$ by
Lemma~\ref{rl-cofactor-norm}, and $n_R^{-1}$
does not change its $W$-component.
Finally $h_R$ preserves $W$, fixes $z'$, and
acts unitarily on $W'_+$.  Projecting onto $W'_+$
proves \eqref{rl-omega-covariance}.
\end{proof}

\subsection{The radial cell and the compact Capelli contraction}

For $c=(c_1,\ldots,c_r)\in(\mathbb C^\times)^r$,
let $b(c)\in\mathbf G$ fix $W$ and satisfy
\[
 b(c)e_i=c_i e_i^\vee,\qquad
 b(c)e_i^\vee=\overline c_i^{-1}e_i.
\]
Let $m\in U(W,h|_W)$, extended as the identity on
$U\oplus U^\vee$.  Put
\[
             a(c)=\prod_{i=1}^r c_i,\qquad
             \zeta(m)=h(mz',z),\qquad
             G(m)=|\zeta(m)|^2-TT'.
\]

\begin{lemma}\label{rl-cell}
For $g=b(c)m$ one has
\begin{align}
 Z(g)&=\operatorname{diag}(c_1,\ldots,c_r,\zeta(m)),\nonumber\\
 \mathcal G(g)&=|a(c)|^2G(m),\label{rl-cell-Gram}\\
 \Omega_j(g)&=\overline{a(c)}\,h(m^{-1}z,B_j).
                                                 \label{rl-cell-omega}
\end{align}
\end{lemma}

\begin{proof}
Pair the images of the indicated basis vectors.
The decomposition into $W$ and its perpendicular
hyperbolic planes gives the diagonal matrix.
Its last cofactor row has only the last entry
$a(c)$, so $L(g)=\overline{a(c)}z$.
Since $b(c)$ fixes $z$, it follows that
$y(g)=\overline{a(c)}m^{-1}z$.
The remaining statements follow from the definitions.
\end{proof}

Let $\pi':K_{z'}\to U(W_{\pi'})$ be a finite-dimensional
continuous unitary representation.  Its complexified
derivative in the fixed basis $B_1,\ldots,B_k$
gives the operators $E_{ij}$ of
Lemma~\ref{rp-capelli}.  For a unit vector
$\nu=(\nu_1,\ldots,\nu_k)\in\mathbb C^k$ define
\begin{equation}\label{rl-cap-vector}
 C_\nu(Y)=\sum_{i,j=1}^k(-1)^{i+j}
       \overline{\nu_i}\nu_j\,
       Q_{\{1,\ldots,k\}\setminus\{i\},
         \{1,\ldots,k\}\setminus\{j\}}
          \left(Y+\frac{k-2}{2}\right).
\end{equation}
The indexing convention in the minor is that of
Lemma~\ref{rp-capelli}, with positions starting at zero
inside each ordered tuple.  The sign is unchanged
when the ambient indices are shifted from zero to one.

\begin{lemma}\label{rl-compressed-capelli}
The operator $C_\nu(Y)$ is the compact Capelli
determinant on the orthogonal complement of
$\sum_i\nu_iB_i$, with shift $Y+(k-2)/2$.
In particular it is independent of an orthonormal
basis of that complement, and
\[
 C_\nu(Y)^*=C_\nu(\overline Y).
\]
If $h\in K_{z'}$ and $\nu'$ is the coordinate vector
of $h^{-1}\sum_i\nu_iB_i$, then
\begin{equation}\label{rl-cap-covariance}
          C_{\nu'}(Y)=\pi'(h)^{-1}C_\nu(Y)\pi'(h).
\end{equation}
Also $C_{\lambda\nu}(Y)=C_\nu(Y)$ when $|\lambda|=1$.
\end{lemma}

\begin{proof}
Complete $\nu$ to a unitary matrix
$P=(v_1,\ldots,v_{k-1},\nu)$.
Apply all-minor covariance from Lemma~\ref{rp-capelli}
to the determinant on its first $k-1$ columns.
The complementary-minor formula gives
\[
 \det P_{\widehat i,\{1,\ldots,k-1\}}
                =(-1)^{i+k}\det P\,\overline{\nu_i}.
\]
The conjugate factor from the inverse matrix is
$(-1)^{j+k}\overline{\det P}\,\nu_j$.
Their product is
$(-1)^{i+j}\overline{\nu_i}\nu_j$, which proves
\eqref{rl-cap-vector} and its interpretation.
Frame independence and the adjoint formula follow
from the corresponding Capelli statements.
Conjugation by $h^{-1}$ carries the complement and
all its rank-one Lie algebra generators to the
new complement, giving \eqref{rl-cap-covariance}.
The scalar phase cancels directly in
\eqref{rl-cap-vector}.  When $k=1$ this is the
empty determinant $1$, so that case is included.
\end{proof}

On the open set $\mathbf G^+=\{g:\mathcal G(g)>0\}$,
define
\[
     \nu(g)=\sqrt{T'/\mathcal G(g)}\,\Omega(g),
                   \qquad C_g(Y)=C_{\nu(g)}(Y).
\]
Lemma~\ref{rl-cofactor-norm} makes $\nu(g)$ a unit vector.
In particular the square root is the positive real one.
Substitution into \eqref{rl-cap-vector} gives the
formula without square roots
\begin{equation}\label{rl-rational-matrix}
 C_g(Y)=\frac{T'}{\mathcal G(g)}
        \sum_{i,j=1}^k(-1)^{i+j}
            \overline{\Omega_i(g)}\Omega_j(g)
            Q_{\widehat i,\widehat j}
                   \left(Y+\frac{k-2}{2}\right).
\end{equation}
It is smooth on $\mathbf G^+$, including at the
zeros of intermediate principal minors.

\begin{proposition}\label{rl-actual-terminal-lift}
For a radial function $\Phi$ of Section~\ref{rd-section},
at $g=b(c)m$ with $G(m)>0$,
\begin{equation}\label{rl-terminal-match}
   G(m)^{-1}(\mathcal T_{22}(Y)\Phi)(c,m)
                              =\Phi(c,m)C_g(Y).
\end{equation}
Moreover, with $g'$ as in Lemma~\ref{rl-cofactor-covariance},
\[
                C_{g'}(Y)=\pi'(h_R)^{-1}C_g(Y)\pi'(h_R).
\]
\end{proposition}

\begin{proof}
By \eqref{rl-cell-omega}, $\nu(g)$ differs by the
unit scalar $\overline{a(c)}/|a(c)|$ from the
normalized projection of $m^{-1}z$ onto $W'_+$.
Its orthogonal complement in $W'_+$ is therefore
$\langle m^{-1}z,z'\rangle^\perp$.
Lemma~\ref{rl-compressed-capelli} and the actual terminal
identity of Corollary~\ref{rp-finite-terminal} prove
\eqref{rl-terminal-match}.
The covariance assertion follows from
\eqref{rl-omega-covariance} and
\eqref{rl-cap-covariance}.
\end{proof}

\subsection{The inverse polynomial and its regular full-group lift}

Retain the common finite type set $\mathcal S$ and
$p_\rho(Y)=\prod_{i=1}^{k-1}(Y+\rho_i)$ from
Section~\ref{rd-section}.  It is nonempty under the
interlacing hypotheses there.
Let $\sigma\in\{1,-1\}$; for the actual signed
source operator it is $\sigma=(-1)^{k-l_0}$.
Choose $Y\in\mathbb C$ such that $p_\rho(Y)\ne0$
for every $\rho\in\mathcal S$.  Define
\begin{align}
 P_Y(X)&=\prod_{\rho\in\mathcal S}
                         (X-\sigma p_\rho(Y)),\nonumber\\
 q_Y(X)&=\frac{1-P_Y(X)/P_Y(0)}{X}. \label{rl-inverse-polynomial}
\end{align}
The numerator has zero constant term, so $q_Y$ is
a polynomial, of degree at most $|\mathcal S|-1$.
Repeated values of $p_\rho(Y)$ cause no difficulty.
For every $\rho$,
$q_Y(\sigma p_\rho(Y))=(\sigma p_\rho(Y))^{-1}$.
For real $Y$ this polynomial has real coefficients.
An explicit simultaneous choice avoiding the zeros
at $Y-1,Y,Y+1$ is
\[
 Y=2+\max\bigl(\{0\}\cup
          \{|\rho_i|:\rho\in\mathcal S,\ 1\leq i<k\}\bigr).
\]

Let $W_\pi$ be the left coefficient space and set
$\mathcal H=\operatorname{Hom}(W_{\pi'},W_\pi)$ with
the Hilbert--Schmidt inner product.
For $g\in\mathbf G^+$ define the endomorphism
\begin{equation}\label{rl-zero-coefficient}
       M_Y(g):A\longmapsto
          \mathcal G(g)^{-1}A\,q_Y(\sigma C_g(Y)).
\end{equation}

\begin{theorem}\label{rl-regular-inverse}
The coefficient \eqref{rl-zero-coefficient} is smooth
on $\mathbf G^+$ and has the following properties.

\begin{enumerate}
\item On a radial cell it realizes exactly the signed
terminal inverse multiplied by the reciprocal full
Gram polynomial:
\[
 M_Y(b(c)m)\Phi(c,m)
    =\frac1{|a(c)|^2G(m)}
       \left(\kappa\bigl(
            (\sigma p_\rho(Y))^{-1}\bigr)\Phi\right)(c,m).
\]
Thus it is the actual coefficient used for the product
of this inverse and the reciprocal Gram factor.
\item It commutes with the full compact and unipotent
action on coefficient functions.  More explicitly,
for the elements in Lemma~\ref{rl-cofactor-covariance},
\[
 M_Y(g)\bigl(\pi(h_L)A\pi'(h_R)^{-1}\bigr)
    =\pi(h_L)\bigl(M_Y(g')A\bigr)\pi'(h_R)^{-1}.
\]
\item Its pointwise Hilbert--Schmidt adjoint is smooth
and has the same derivative growth bounds as $M_Y$.
For real $Y$ it equals $M_Y$.
\item For every $\varepsilon>0$ and integer $m\geq0$,
there is a finite constant $C_{\varepsilon,m}$ such that
every word $D_1\cdots D_j$ of at most $m$ left or right
Lie derivatives in matrices whose entries have absolute
value at most one satisfies
\begin{equation}\label{rl-growth}
 \|D_1\cdots D_j M_Y(g)\|
       \leq C_{\varepsilon,m}
                      (1+\|g\|_{\rm F})^{(r+1)m}
         \quad(\mathcal G(g)\geq\varepsilon).
\end{equation}
The same bound holds for its adjoint.
The constant depends on the fixed Hermitian data,
the finite representations, $Y,\sigma,\varepsilon,m$,
and not on $g$ or on the selected unit Lie directions.
\end{enumerate}
\end{theorem}

\begin{proof}
The finite decomposition of Theorem~\ref{rd-constructed-kappa}
and \eqref{rl-terminal-match} identify
$q_Y(\sigma C_g(Y))$ on the evaluated radial image
with the specified inverse.
No assertion that it is the inverse on unused
right constituents is needed.  This proves the first
part, using \eqref{rl-cell-Gram}.
Smoothness follows from \eqref{rl-rational-matrix}.
The covariance of $C_g$ and invariance of $\mathcal G$
give the second part by moving $\pi'(h_R)^{-1}$
through the polynomial.
The adjoint of right multiplication by a matrix $B$
on the Hilbert--Schmidt space is right multiplication
by $B^*$.  The adjoint formula in
Lemma~\ref{rl-compressed-capelli} therefore proves
the third part; the final assertion for real $Y$
also uses that $q_Y$ has real coefficients.

Here is a direct derivative estimate, which avoids
differentiating any orthonormal frame.  Write
$d_0=r+1$, $n=2r+k+1$, and use the fixed polynomial
extension of $\Omega$ from Lemma~\ref{rl-cofactor-norm}.
There is a constant $C_\Omega$ such that any word
$L$ of $a$ left or right unit Lie derivatives obeys
\begin{equation}\label{rl-polynomial-derivative}
               \|L\Omega(g)\|
       \leq C_\Omega(nd_0)^a(1+\|g\|_{\rm F})^{d_0}.
\end{equation}
Indeed each such vector field replaces an entry
of $g$ by a sum of at most $n$ entries with coefficients
of absolute value at most one, and likewise for a
conjugate entry.  On a monomial of degree $d_0$
there are $d_0$ choices of a differentiated factor.
Thus each derivative multiplies the coefficient
sum by at most $nd_0$, without increasing the
degree.  The Euclidean sum of the component coefficient
norms gives a possible $C_\Omega$.

For $v\in\mathbb C^k\setminus\{0\}$ set
\[
 \mathcal F_Y(v):A\longmapsto
       \frac{1}{T'\|v\|^2}
          A\,q_Y\bigl(\sigma C_{v/\|v\|}(Y)\bigr).
\]
This is a smooth function with values in
$\operatorname{End}(\mathcal H)$, and
$\mathcal F_Y(tv)=t^{-2}\mathcal F_Y(v)$ for real $t>0$.
Consequently
\[
 \|D^\ell\mathcal F_Y(v)\|
       \leq M_\ell\|v\|^{-2-\ell},\qquad
 M_\ell=\max_{\|u\|=1}\|D^\ell\mathcal F_Y(u)\|<\infty.
\]
The finiteness follows from smoothness and compactness
of the unit sphere.  These formulas define the
constants entirely in the fixed finite coefficient space.
By \eqref{rl-omega-norm},
$M_Y(g)=\mathcal F_Y(\Omega(g))$, and
$\|\Omega(g)\|\geq\delta=\sqrt{\varepsilon/T'}$
on the set under consideration.

The iterated chain rule for a word of $j$ vector fields
expresses its derivative of $\mathcal F_Y(\Omega)$
as a sum over partitions of those $j$ positions.
A partition with $\ell$ blocks contributes
$D^\ell\mathcal F_Y$ applied to $\ell$ derivatives
of $\Omega$, using the original order within each
block.  This statement follows by induction:
the next derivative either differentiates
$\mathcal F_Y$ or one of its arguments.
It does not require the vector fields to commute.
Using \eqref{rl-polynomial-derivative}, its norm is
at most
\[
 (nd_0)^j
 \sum_{\ell=0}^j
    \left\{\begin{matrix}j\\\ell\end{matrix}\right\}
       M_\ell\delta^{-2-\ell}C_\Omega^\ell
                    (1+\|g\|_{\rm F})^{d_0\ell},
\]
where the braces count partitions into $\ell$ nonempty
blocks, with the usual value $1$ at $(j,\ell)=(0,0)$.
Take the maximum of the corresponding fixed constants
over $0\leq j\leq m$, and use $\ell\leq j\leq m$.
This proves \eqref{rl-growth}.
Taking an adjoint is a fixed real-linear isometry
on the finite coefficient space, so the same argument
and bound apply to the adjoint coefficient.
\end{proof}

\subsection{Preservation of the actual rapidly decreasing class}

For $\varepsilon>0$, let $\mathscr S_\varepsilon$
consist of the smooth $\mathcal H$-valued functions
on $\mathbf G$ which vanish on $\mathcal G<\varepsilon$
and for which every mixed left and right Lie derivative
decays faster than every power of $1+\|g\|_{\rm F}$,
uniformly over unit Lie directions.
This is the coefficient class in the source's
smoothness, decay and support hypotheses.
Its formulation with the norm $\|g\|_{\rm F}$ alone
is equivalent: $|\det g|=1$, so Hadamard's determinant
inequality implies $\|g\|_{\rm F}\geq1$.
Hence $1+\|g\|_{\rm F}\leq2\|g\|_{\rm F}$.

\begin{corollary}\label{rl-class-preserved}
Extend $M_YF$ by zero on $\mathcal G\leq0$.
Both $F\mapsto M_YF$ and its pointwise adjoint
preserve $\mathscr S_\varepsilon$.
They are zero-order differential operators with
smooth coefficients on $\mathbf G^+$, they commute
with the source compact and unipotent actions, and
all their coefficient derivatives have polynomial
growth on $\mathcal G\geq\varepsilon$.
\end{corollary}

\begin{proof}
The functions in question vanish on the neighborhood
$\mathcal G<\varepsilon$ of $\mathcal G=0$,
so their zero extensions are smooth there.
On $\mathbf G^+$ use the product rule.
For a derivative of order at most $m$, every
coefficient factor is bounded by
\eqref{rl-growth}.  The differentiated $F$ can be
bounded by an arbitrarily large negative power;
choosing that power larger by $(r+1)m$ gives the
required decay.  The support is not enlarged.
Covariance and the adjoint properties were proved
in Theorem~\ref{rl-regular-inverse}.
\end{proof}

The exponent $(r+1)m$ in \eqref{rl-growth} is sufficient
for the actual qualitative coefficient condition.
It also improves the exponent $3(r+1)m$ used in
the corresponding source provider for this zero-order
factor.  The constants are not claimed to be uniform
as $Y$ approaches a forbidden spectral value.
They can be chosen uniformly on a fixed compact set
of $Y$ avoiding those finitely many values, since
the coefficients of $q_Y$ and all sphere derivatives
above then vary continuously on compact sets.

\begin{example}\label{rl-vanishing-middle-minor}
The domain $\mathbf G^+$ really contains points
where a radial coordinate is unavailable.
Take $k=r=1$, a basis $(e,a,z,e^\vee)$ with
$h(e,e^\vee)=1$, $h(a,a)=1$, $h(z,z)=-1$,
and all other required pairings zero.
Set $z'=z$ and define
\[
 ge=a+z,\quad
 ga=\tfrac12e+e^\vee,\quad
 gz=\tfrac12e-e^\vee,\quad
 ge^\vee=\tfrac12(a-z).
\]
The displayed images have exactly the same pairings
as the basis, so $g\in\mathbf G$.
Here
\[
             Z(g)=\begin{pmatrix}0&-1\\-1&0\end{pmatrix},
       \qquad d(g)=0,\quad D(g)=-1,\quad\mathcal G(g)=1.
\]
Thus a proposed full-group lift involving $d(g)^{-1}$
would be singular at an allowed point.
The cofactor coefficient above is regular there:
$L(g)=e$, $y(g)=a+z$, and $\Omega(g)=1$.
This is not a counterexample to the source, whose
construction uses the regular Gram denominator.
\end{example}

\subsection{What this completes}

The preceding construction proves the actual
zero-order spectral factor, rather than assuming
a smooth matrix field with suitable bounds.
The polynomial cofactor identity proves its full-group
definition; the radial-cell calculation proves that
it is the correct operator; and covariance and
\eqref{rl-growth} prove the analytic coefficient
conditions needed to compose it with other
differential operators.

The remaining full-group operators in the Fourier
identity include two weighted radial transfer
operators and terminal off-diagonal contractions.
Their restriction identities and formal-adjoint
passage through the Jacquet transform are not
proved merely by the zero-order construction here.
Likewise the integrated Jacquet pairing and the
actual local uniform $L^2$ bounds remain separate.
No full weighted Fourier decay theorem is asserted
until those remaining identities are established.


\section{Polynomial differential operators and the Toda radial part}
\label{rt-section}

We now construct two of the actual differential operators used
in the Fourier argument.  Their coefficients are polynomial on
the full unitary group.  We also prove their precise restrictions
to the radial cell, including the scalar conjugation and the
conjugate-coordinate operator.  The proof of the radial identity
uses a last-column expansion of a Capelli determinant and the
two unipotent equivariances; no general theorem about radial
parts is required.  These are the operators denoted by
the weighted isotropic Capelli minors in the source
\cite{alpo-uniform-source}.

\subsection{Right derivatives and the two polynomial operators}

Keep the Hermitian space and notation of Section~\ref{rl-section}.
In particular $V=U\oplus W\oplus U^\vee$,
$\dim U=r$, and $\dim W=k+1$.
Fix any ordering of its chosen Witt basis.
For $X$ in the real Lie algebra of $\mathbf G=U(V,h)$ put
\[
 (R_XF)(g)=\left.\frac{d}{dt}\right|_{t=0}F(g\exp(tX)).
\]
For an arbitrary complex matrix $A$ define the complexified
right derivative by
\begin{equation}\label{rt-complex-derivative}
 R_AF=R_{(A-HA^*H)/2}F
            +iR_{(A+HA^*H)/(2i)}F.
\end{equation}
Both matrices on the right belong to the real Lie algebra.
This is complex linear in $A$, and
$[R_A,R_B]=R_{[A,B]}$.  Indeed the real right derivatives
are the left-invariant vector fields with tangent vectors
$gA$, and differentiating these vector fields gives this
bracket; complexification preserves it.

Write $E_{ab}=R_{e_ae_b^*}$, where $e_b^*$ here means
the coordinate covector, not the Hermitian dual.
For ordered lists $I=(i_1,\ldots,i_m)$ and
$J=(j_1,\ldots,j_m)$ of distinct basis indices define
\begin{equation}\label{rt-minor}
 Q_{I,J}(u)=
 \sum_{\sigma\in S_m}\operatorname{sgn}(\sigma)
 \prod_{b=1}^{m}
 \bigl((u-b+1)\delta_{i_{\sigma(b)},j_b}
                         +E_{i_{\sigma(b)},j_b}\bigr).
\end{equation}
Products act on a function with the rightmost factor first.
The empty minor is the identity.  These are the minors
of Theorem~\ref{rp-capelli}.

Set $I=(e_1,\ldots,e_r)$ and
$I^\vee=(e_1^\vee,\ldots,e_r^\vee)$.  If
$J=(j_1,\ldots,j_r)$ runs over the increasing $r$-element
subsets of the chosen basis, put
\begin{align}
 \beta_J(g)&=\det\bigl(h(ge_{j_b},e_a)\bigr)_{a,b=1}^{r},
                                                        \label{rt-beta}\\
 \alpha_J(g)&=\det\bigl((g^{-1}e_a)_{j_b}\bigr)_{a,b=1}^{r}.
                                                        \label{rt-alpha}
\end{align}
The transpose in the second displayed determinant does
not change its value.  For $Y\in\mathbb C$ set
\begin{align}
 \mathcal A_YF(g)
   &=\sum_J\beta_J(g)\,
       Q_{I,J}\left(Y+\frac{2r+k}{2}\right)F(g),
                                                   \label{rt-A}\\
 \overline{\mathcal A}_YF(g)
   &=\sum_J\alpha_J(g)\,
       Q_{J,I^\vee}\left(Y-\frac{2r+k}{2}+r-1\right)F(g).
                                                   \label{rt-Av}
\end{align}
The bar in $\overline{\mathcal A}$ is a label for the
conjugate-coordinate operator; it does not conjugate the
argument $F$ or the parameter $Y$.
Changing the order of one summation list changes both
its scalar determinant and its quantum minor by the
same sign.  Hence any fixed rule for ordering the subsets
gives the same operators.  This explains the different
ordering of the dual indices used in the source.
For $r=0$ both operators are the identity.

\begin{proposition}\label{rt-polynomial}
The operators \eqref{rt-A} and \eqref{rt-Av} have
differential order at most $r$ and polynomial coefficients
of degree $r$ in $g$ and $\overline g$, respectively.
They are smooth on the whole group, including the zeros
of every intermediate minor.

Every derivative of a scalar coefficient, of any fixed
length $m$ in left or right Lie derivatives with directions
of norm at most one, is bounded by
$C_m(1+\|g\|)^r$.  The constant is independent of $g$
and of those directions.  Their formal adjoints for
Haar measure also have polynomial coefficients of degree
at most $r$ and order at most $r$.

Both the operators and their formal adjoints preserve the
class $\mathscr S_\varepsilon$ of
Corollary~\ref{rl-class-preserved}.
\end{proposition}

\begin{proof}
The first coefficients in \eqref{rt-beta} are polynomial
of degree $r$.  On $\mathbf G$ one has $g^{-1}=Hg^*H$,
so the second coefficients are polynomial of degree $r$
in $\overline g$.  Each term in \eqref{rt-minor} is a
constant linear combination of words of length at most
$r$ in right derivatives.  Thus no division occurs in
either operator.

A left or right Lie derivative sends a matrix entry of
$g$ to a linear combination of its entries, and likewise
for conjugate entries.  On the finite-dimensional space
of polynomials of degree at most $r$, this derivative
therefore acts by a linear map whose coefficients are
bounded when the direction has norm at most one.
Iterating $m$ times bounds the sum of the absolute
coefficients by a fixed constant $C_m$.
Evaluating the resulting polynomials proves the
stated bound.

For clarity, Haar measure on this group is unimodular.
The real bilinear form
$B(X,Z)=\operatorname{Re}\operatorname{tr}(XZ)$ on its
Lie algebra is nondegenerate and invariant under the
adjoint action of every group element.  Nondegeneracy
follows by complexifying: the complexified Lie algebra
is the full matrix algebra, whose trace pairing is
nondegenerate.  An adjoint transformation preserves
$B$, so its real determinant has absolute value one.
The absolute volume form defined at the identity
therefore extends to a left and right invariant
measure.  Integration by parts for compactly supported
functions gives $R_X^*=-R_X$ for real $X$.
Taking adjoints of \eqref{rt-A} and \eqref{rt-Av}
reverses each word and conjugates its coefficients.
Leibniz's rule differentiates those polynomials but
does not increase their degree.

Finally every derivative of an output is a finite sum
of a polynomially bounded coefficient times a derivative
of the input.  Arbitrarily fast polynomial decay of
all input derivatives implies the same for the output.
A differential operator does not enlarge support.
This proves preservation of $\mathscr S_\varepsilon$
for the operators and their adjoints.
\end{proof}

\subsection{Equivariance on the full group}

Let $N,K_z,K_{z'}$ be as in
Lemma~\ref{rl-cofactor-covariance}.
The following identities hold for arbitrary smooth
functions, before imposing a character on $N$.

\begin{proposition}\label{rt-full-covariance}
If $s_L\in K_zN$ and $s_R\in K_{z'}N$, then
\begin{align*}
 \mathcal A_Y\bigl(F(s_L\,\cdot\,s_R)\bigr)(g)
       &=(\mathcal A_YF)(s_Lgs_R),\\
 \overline{\mathcal A}_Y\bigl(F(s_L\,\cdot\,s_R)\bigr)(g)
       &=(\overline{\mathcal A}_YF)(s_Lgs_R).
\end{align*}
The same assertions hold for vector-valued functions
component by component.
\end{proposition}

\begin{proof}
Both $s_L$ and $s_R$ preserve $U$ and act on it with
determinant one.  They also preserve $U^\perp$, and
the induced action on $V/U^\perp$, in the basis
given by $I^\vee$, has determinant one.

Left translation commutes with right derivatives.
For the first operator,
\[
 h(s_Lge_j,e_a)=h(ge_j,s_L^{-1}e_a).
\]
The basis change on $U$ has determinant one, so every
$r$-row determinant in \eqref{rt-beta} is unchanged.
For the second operator,
$\bigwedge^r(s_L^{-1})$ fixes
$e_1\wedge\cdots\wedge e_r$, so
$\alpha_J(s_Lg)=\alpha_J(g)$.
This proves left covariance.

For right translation by $s$, complexified derivatives
transform by
\[
 R_A\bigl(F(\,\cdot\,s)\bigr)(g)
          =(R_{s^{-1}As}F)(gs).
\]
The exterior covariance of the quantum minors gives
\begin{equation}\label{rt-right-minor}
 Q'_{I,J}(u)
    =\sum_{K,L}\det(s^{-1})_{K,I}\,
                   \det(s)_{J,L}\,Q_{K,L}(u),
\end{equation}
where $Q'$ is formed using $A\mapsto R_{s^{-1}As}$.
This is precisely the all-minor covariance in
Theorem~\ref{rp-capelli}; no coefficient is
differentiated in this constant change of basis.

Since $\bigwedge^r(s^{-1})$ fixes $\bigwedge^rU$,
the first determinant on the right of
\eqref{rt-right-minor}, when $I$ is the fixed list
of $U$-indices, equals one for $K=I$ and zero otherwise.
Cauchy--Binet then gives
\[
 \sum_J \beta_J(g)\det(s)_{J,L}=\beta_L(gs).
\]
This proves right covariance of $\mathcal A_Y$.

For the second operator, Cauchy--Binet gives instead
\[
 \sum_J\det(s^{-1})_{K,J}\alpha_J(g)=\alpha_K(gs).
\]
The remaining coefficients in
\eqref{rt-right-minor} are $\det(s)_{I^\vee,L}$.
The last $r$ coordinate rows of $s$ vanish on
$U^\perp$ and have determinant one on the quotient.
Consequently these coefficients are one for
$L=I^\vee$ and zero otherwise.  This proves right
covariance of $\overline{\mathcal A}_Y$.
The argument is unaffected by any fixed reordering
of the dual basis.
\end{proof}

\subsection{The two unipotent characters and elementary derivative rules}

For $n\in N$, define
\[
 \ell_z(n)=
 \sum_{i=1}^{r-1}(ne_{i+1})_{e_i}+(nz)_{e_r},
 \qquad
 \psi_z(n)=\exp\bigl(\ell_z(n)-\overline{\ell_z(n)}\bigr).
\]
For $r=0$ set $\ell_z=0$.
The same definition with $z'$ gives $\psi_{z'}$.
These are the actual normalizations of the source's
characters.  Only their values on the roots inside
$U$ and $U^\vee$ enter the proof below.

Let $D$ be an open subset of $\mathbf G$ invariant
under multiplication by $N$ on both sides.
Let $F:D\to E$ be smooth, where $E$ is a
finite-dimensional complex vector space, and suppose
\begin{equation}\label{rt-whittaker}
 F(n_Lgn_R)=\psi_z(n_L)\psi_{z'}(n_R)F(g).
\end{equation}
Fix $m\in U(W)$ and consider the points $b(c)m\in D$,
with all $c_i\ne0$, as in Lemma~\ref{rl-cell}.
Write
\[
 F_m(c)=F(b(c)m),\qquad
 p_i=c_i\partial_{c_i},\qquad
 \bar p_i=\overline c_i\partial_{\overline c_i}.
\]
Wirtinger derivatives are meant in these formulas.

\begin{lemma}\label{rt-root-rules}
For the ordered indices $e_1,\ldots,e_r$,
right differentiation on a function satisfying
\eqref{rt-whittaker} gives
\begin{align}
 E_{e_i,e_j}F&=\delta_{j,i+1}F
                                      &&(i<j),\label{rt-right-rule}\\
 (E_{e_i,e_j}F)(b(c)m)
       &=\delta_{i,j+1}\frac{c_i}{c_j}F_m(c)
                                      &&(i>j),\label{rt-left-rule}\\
 (E_{e_i,e_i}F)(b(c)m)&=p_iF_m(c).
                                                   \label{rt-torus-rule}
\end{align}
Equation \eqref{rt-left-rule} holds also with $F$
replaced by any word of right derivatives applied
to $F$.

For the reversed dual list
$v_i=e_{r+1-i}^\vee$, the corresponding rules are
\begin{align}
 E_{v_i,v_j}F&=\delta_{j,i+1}F&&(i<j),\label{rt-dual-right}\\
 (E_{v_i,v_j}F)(b(c)m)
   &=\delta_{i,j+1}
       \frac{\overline c_{r+1-j}}{\overline c_{r+1-i}}
         F_m(c)&&(i>j),\label{rt-dual-left}\\
 (E_{v_i,v_i}F)(b(c)m)
   &=-\bar p_{r+1-i}F_m(c).\label{rt-dual-torus}
\end{align}
The dual lower-root rule again holds after any
word of right derivatives.
\end{lemma}

\begin{proof}
For $i<j$, the real root subgroup inside $N$ is
\[
 n_{ij}(w)=
 \exp\bigl(w e_ie_j^*-\overline w\,
                         e_j^\vee(e_i^\vee)^*\bigr).
\]
Its character is $\exp(w-\overline w)$ when
$j=i+1$ and is one otherwise.
Complexifying its two real tangent directions as
in \eqref{rt-complex-derivative} amounts to taking
$\partial_w$ at zero.  The right character law
therefore gives \eqref{rt-right-rule}.

At $g=b(c)m$ one has, for $i>j$,
\[
 \operatorname{Ad}(g)(e_ie_j^*)
          =\frac{c_i}{c_j}
                      e_i^\vee(e_j^\vee)^*.
\]
The matrix on the right is an upper root in the
dual part of $N$.  Its character derivative is
one when $i=j+1$ and zero otherwise.  For example,
if its root parameter is $w$, the corresponding
upper entry in the $U$-block is $-\overline w$,
whose contribution to
$\ell-\overline\ell$ is again $w-\overline w$.
Transferring the right derivative at $g$ to this
left derivative proves \eqref{rt-left-rule}.
Every right derivative commutes with left
translations.  Thus every right-differentiated
function retains the left character law, which
proves the asserted strengthened version.

The two real directions in
\eqref{rt-complex-derivative} for $e_ie_i^*$
scale $c_i$ respectively by $e^{t/2}$ and
$e^{-it/2}$.  Their complex combination is
$c_i\partial_{c_i}$, proving \eqref{rt-torus-rule}.

In the reversed dual ordering, the upper root
$e_a^\vee(e_b^\vee)^*$ has $a>b$ and has
the same character derivative, equal to one
precisely when $a=b+1$.  For a lower root in
that ordering, conjugation by $b(c)m$ gives
\[
 \operatorname{Ad}(b(c)m)
       \bigl(e_a^\vee(e_b^\vee)^*\bigr)
           =\frac{\overline c_b}{\overline c_a}
                          e_ae_b^* .
\]
This proves \eqref{rt-dual-right} and
\eqref{rt-dual-left}.  The real torus directions
for $e_i^\vee(e_i^\vee)^*$ scale $c_i$ by
$e^{-t/2}$ and $e^{-it/2}$; their complex
combination is $-\bar c_i\partial_{\bar c_i}$.
This proves the last identity.
\end{proof}

\subsection{An elementary determinant recurrence}

The following argument is the essential radial calculation.
It explains why only adjacent roots survive.

\begin{lemma}\label{rt-cofactor-reduction}
Let $E_{ab}$ satisfy the $\mathfrak{gl}_s$ commutator
relations, and let $s\geq2$.
In its universal enveloping algebra, put
\[
 C_s(u)=Q_{(1,\ldots,s-2,s),(1,\ldots,s-1)}(u).
\]
Then
\begin{equation}\label{rt-cofactor-ideal}
 C_s(u)=E_{s,s-1}Q_{(1,\ldots,s-2),(1,\ldots,s-2)}(u)
                  +\sum_{j=1}^{s-2}E_{sj}V_j
\end{equation}
for some elements $V_j$ of the enveloping algebra.
No assertion about the values of the $V_j$ is needed.
\end{lemma}

\begin{proof}
In each term of the determinant, there is exactly
one factor from row $s$, necessarily an off-diagonal
factor $E_{sj}$.  If $j<s-1$, move this factor to
the far left.  For a factor from an earlier column
$b<j$ one has
\[
 \bigl((u-b+1)\delta_{ab}+E_{ab}\bigr)E_{sj}
  =E_{sj}\bigl((u-b+1)\delta_{ab}+E_{ab}\bigr)
                                  -\delta_{aj}E_{sb}.
\]
Every new last-row generator has an even smaller
column index.  Repeating this identity expresses
the term as a sum of elements with leftmost factor
$E_{s\ell}$, where $\ell\leq j<s-1$.

If $j=s-1$, every preceding row index is at most
$s-2$.  The same commutator is then zero.
Moving $E_{s,s-1}$ to the far left leaves exactly
the determinant on the first $s-2$ rows and
columns, with parameter $u$ and positive sign.
Summing all terms proves \eqref{rt-cofactor-ideal}.
\end{proof}

\begin{proposition}\label{rt-raw-toda}
For $0\leq s\leq r$, let $K_s(u)$ be the
coordinate differential operator defined by
\[
 K_0(u)=1,\qquad K_1(u)=u+p_1,
\]
and, for $s\geq2$,
\begin{equation}\label{rt-K-rec}
 K_s(u)=(u-s+1+p_s)K_{s-1}(u)
                      -\frac{c_s}{c_{s-1}}K_{s-2}(u).
\end{equation}
For every $F$ satisfying \eqref{rt-whittaker},
\begin{equation}\label{rt-raw-identity}
 \bigl(Q_{(e_1,\ldots,e_s),(e_1,\ldots,e_s)}(u)F\bigr)
                (b(c)m)=K_s(u)F_m(c).
\end{equation}
\end{proposition}

\begin{proof}
We prove the slightly stronger inductive assertion in which
the full left character law is retained, but on the right
only the upper-root character laws within the first $s$
isotropic indices are required.  The stated theorem is
its special case.
The assertions for $s=0,1$ are immediate from the
empty determinant and \eqref{rt-torus-rule}.
Expand the size $s$ determinant in its last column.
The last factor in each product acts first.
By \eqref{rt-right-rule}, all terms vanish except
the last-row diagonal term and the term with row
$s-1$.  Thus its value equals that of
\[
 Q_{s-1}(u)(u-s+1+E_{ss})F-C_s(u)F.
\]
Here the abbreviated indices belong to the first
isotropic list.

The function $E_{ss}F$ retains left equivariance,
and it retains right equivariance for all upper
roots within the first $s-1$ indices, since those
generators commute with $E_{ss}$.
The induction hypothesis therefore applies to
this function.  The resulting diagonal term is
$K_{s-1}(u)(u-s+1+p_s)F_m$.
The operator $K_{s-1}(u)$ uses only
$c_1,\ldots,c_{s-1}$, so it commutes with $p_s$.

Apply Lemma~\ref{rt-cofactor-reduction} to the
other term.  The strengthened lower-root rule
in Lemma~\ref{rt-root-rules} kills every summand
$E_{sj}V_jF$ with $j<s-1$.  The remaining term
equals
\[
 \frac{c_s}{c_{s-1}}\,
       \bigl(Q_{s-2}(u)F\bigr)(b(c)m).
\]
The induction hypothesis gives exactly
\eqref{rt-K-rec}.  This proof uses only the root
laws actually needed at each stage; it does not
assume that arbitrary right derivatives preserve
all right character laws.
\end{proof}

\subsection{The exact scalar conjugation and the conjugate operator}

Define
\begin{equation}\label{rt-gauge}
 \gamma(c)=\prod_{i=1}^{r}
       |c_i|^{\,2(r-i+1)+k},
 \qquad
 \Phi(c,m)=\gamma(c)F(b(c)m).
\end{equation}
This is exactly the source's scalar radial
normalization.  Let
\[
 T_i(Y)=
 \begin{pmatrix}Y+p_i&-c_i\\ c_i^{-1}&0\end{pmatrix},
 \qquad
 \overline T_i(Y)=
 \begin{pmatrix}Y-\bar p_i&-\bar c_i\\
                         \bar c_i^{-1}&0\end{pmatrix},
\]
where entries act on functions of $c$.
Write their ordered products as
\[
 \begin{pmatrix}A_r(Y)\\ B_r(Y)\end{pmatrix}
       =T_r(Y)\cdots T_1(Y)\binom10,\qquad
 \begin{pmatrix}\overline A_r(Y)\\ \overline B_r(Y)\end{pmatrix}
       =\overline T_r(Y)\cdots\overline T_1(Y)\binom10.
\]
For $r=0$ both columns are $(1,0)^{\mathsf t}$.
Set $a(c)=\prod_i c_i$.

\begin{theorem}\label{rt-actual-lifts}
For every smooth $F$ satisfying
\eqref{rt-whittaker}, and every allowed cell point,
\begin{align}
 \gamma(c)(\mathcal A_YF)(b(c)m)
       &=a(c) A_r(Y)\Phi(c,m),\label{rt-lift-first}\\
 \gamma(c)(\overline{\mathcal A}_YF)(b(c)m)
       &=\overline{a(c)}\,\overline A_r(Y)\Phi(c,m).
                                             \label{rt-lift-second}
\end{align}
Both identities hold for every complex $Y$.
The full-group operators on the left are the
polynomial operators already constructed in
\eqref{rt-A} and \eqref{rt-Av}.
\end{theorem}

\begin{proof}
At a cell point, $h(b(c)me_j,e_a)$ is zero unless
$j=e_a$, and its nonzero value is $c_a$.
Only one scalar minor in \eqref{rt-A} survives.
Consequently
\[
 (\mathcal A_YF)(b(c)m)
      =a(c)\bigl(Q_{I,I}(u)F\bigr)(b(c)m),
             \qquad u=Y+\frac{2r+k}{2}.
\]
For $n_i=2(r-i+1)+k$,
\[
 \gamma p_i\gamma^{-1}=p_i-\frac{n_i}{2},
 \qquad u-i+1-\frac{n_i}{2}=Y.
\]
Multiplication by $c_s/c_{s-1}$ commutes with
$\gamma$.  Conjugating \eqref{rt-K-rec} therefore
gives the recurrence
\[
 A_s=(Y+p_s)A_{s-1}
                -\frac{c_s}{c_{s-1}}A_{s-2},
\]
with $A_0=1$, $A_1=Y+p_1$.
The first column of $T_s\cdots T_1$ obeys the
same recurrence, since $B_{s-1}=c_{s-1}^{-1}A_{s-2}$.
This proves \eqref{rt-lift-first}.

For the second operator,
$(b(c)m)^{-1}e_a=\bar c_a e_a^\vee$.
Thus \eqref{rt-Av} reduces to
\[
 \overline{a(c)}\,
 Q_{I^\vee,I^\vee}(w)F(b(c)m),
        \qquad w=Y-\frac{2r+k}{2}+r-1 .
\]
Simultaneously reversing both index lists of a
full Capelli determinant does not change it,
by its covariance under a permutation matrix.
Apply the proof of Proposition~\ref{rt-raw-toda}
to the reversed dual list in
Lemma~\ref{rt-root-rules}.
At the vertex corresponding to the original
index $i$, the diagonal factor is
\[
                w-(r-i)-\bar p_i,
\]
and the edge between $i$ and $i+1$ has factor
$-\bar c_{i+1}/\bar c_i$.
Conjugation by $\gamma$ changes that diagonal
factor to $Y-\bar p_i$, because
\[
 w-(r-i)+\frac{n_i}{2}=Y.
\]

For completeness, the resulting recurrence can
be read in either direction along the ordered
set of vertices $1,\ldots,r$.
Its expansion is a sum over sets of pairwise
disjoint adjacent edges.  An edge $(i,i+1)$
contributes $-\bar c_{i+1}/\bar c_i$; every
uncovered vertex $j$ contributes $Y-\bar p_j$.
All factors in a fixed term commute: an
uncovered derivative involves none of the
variables in the edge factors.  Splitting these
sets according to whether the last vertex is
uncovered or lies in its last edge gives the
recurrence for $\overline T_r\cdots\overline T_1$.
Splitting at the first vertex gives the reversed
dual recurrence.  They are therefore equal.
This proves \eqref{rt-lift-second}, including
$r=0$ and $r=1$.
\end{proof}

\paragraph{Scope.}
The two polynomial full-group operators, their
translation covariance, their action on the
coefficient class, and their exact radial
identities are now proved.  Together with
Section~\ref{rl-section}, this supplies three
concrete ingredients of the four-operator
Fourier construction.  The remaining coupled
operators involving the off-diagonal terminal
entries require their own identities.  The
passage from formal adjoints to the actual
Jacquet pairing, the uniform square-integrability
bounds for the generated class, and the final
weighted Fourier estimate are not consequences
of the two polynomial radial identities alone.


\section{The two differential operators involving the terminal matrix}
\label{rc-section}

We next construct the remaining two polynomial differential
operators in the four-operator argument.  The calculation
extends Section~\ref{rt-section} by retaining a block of
terminal indices.  A general bordered-minor identity gives
both the ordinary recurrence and its boundary term.
All signs and normalizations below refer to the actual
terminal operators of the source \cite{alpo-uniform-source}.

\subsection{The terminal entries and a bordered-minor identity}

Use the notation of Section~\ref{rt-section}, with the
ordered basis
\[
 e_1,\ldots,e_r,\ w_0,\ldots,w_k,\
                         e_r^\vee,\ldots,e_1^\vee .
\]
Here $h(w_i,w_j)=\epsilon_i\delta_{ij}$, where
$\epsilon_0=\cdots=\epsilon_{k-1}=1$ and
$\epsilon_k=-1$.  Put $d=k+1$.
For $m\in U(W)$ set
\[
 a_j=(m^{-1}z)_j,\qquad b_i=h(w_i,m^{-1}z),
 \qquad a'_j=z'_j,\qquad b'_i=h(w_i,z').
\]
These four vectors have exactly the coefficient conventions
in the terminal matrix.  Let $Q^W$ denote the quantum
minors made from right derivatives on $U(W)$.
Write $\widehat i$ for the increasing list
$\{0,\ldots,k\}\setminus\{i\}$, and similarly
$\widehat{i,i'}$ for two omissions.  Set $v=Y+k/2$.
The three entries needed here are
\begin{align}
 t_{21}(Y)&=
 -\sum_{i,j}(-1)^{i+j}b'_i a_j
                    Q^W_{\widehat i,\widehat j}(v),
                                                     \label{rc-t21}\\
 t_{12}(Y)&=
 -\sum_{i,j}(-1)^{i+j}b_i a'_j
                    Q^W_{\widehat i,\widehat j}(v-1),
                                                     \label{rc-t12}\\
 t_{22}(Y)&=
 -\sum_{\substack{i<i'\\j<j'}}(-1)^{i+i'+j+j'}
     (b_i b'_{i'}-b_{i'}b'_i)
     (a_j a'_{j'}-a_{j'}a'_j)
                 Q^W_{\widehat{i,i'},\widehat{j,j'}}(v-1).
                                                     \label{rc-t22}
\end{align}
Every scalar coefficient in these formulas multiplies
the result of the indicated differential operator.
It is not differentiated by that operator.
The last entry is the one already identified with a
positive compact Capelli contraction in
Proposition~\ref{rp-finite-terminal}.

We record two algebraic forms of the same determinant
calculation.  Lists $I=(i_1,\ldots,i_q)$ and
$J=(j_1,\ldots,j_q)$ in the following lemma have distinct
entries, none equal to a distinguished index $e$.
They need not be equal to one another.
All indices in signs inside this lemma start at one.

\begin{lemma}\label{rc-border}
Let $E_{ab}$ satisfy the general linear Lie algebra
relations, and let $\Lambda$ be a complex linear map
from their enveloping algebra to a complex vector space.

Suppose that for every $X$ and all $i\in I,j\in J$,
\[
 \Lambda(E_{ie}X)=\ell_i\Lambda(X),\qquad
 \Lambda(XE_{ej})=r_j\Lambda(X).
\]
Then
\begin{align}
 \Lambda\bigl(Q_{(e,I),(e,J)}(u)\bigr)
 &=\Lambda\bigl((u+E_{ee})Q_{I,J}(u-1)\bigr)\nonumber\\
 &\quad-\sum_{\alpha,\beta=1}^{q}
     (-1)^{\alpha+\beta}\ell_{i_\alpha}r_{j_\beta}
       \Lambda\bigl(Q_{I\setminus i_\alpha,J\setminus j_\beta}
                                                   (u-2)\bigr).
                                                     \label{rc-border-first}
\end{align}
If instead
\[
 \Lambda(E_{ej}X)=\ell_j\Lambda(X),\qquad
 \Lambda(XE_{ie})=r_i\Lambda(X),
\]
then
\begin{align}
 \Lambda\bigl(Q_{(I,e),(J,e)}(u)\bigr)
 &=\Lambda\bigl(Q_{I,J}(u)(u-q+E_{ee})\bigr)\nonumber\\
 &\quad-\sum_{\alpha,\beta=1}^{q}
     (-1)^{\alpha+\beta}r_{i_\alpha}\ell_{j_\beta}
       \Lambda\bigl(Q_{I\setminus i_\alpha,J\setminus j_\beta}
                                                     (u)\bigr).
                                                     \label{rc-border-last}
\end{align}
The formulas include $q=0$, with empty sums.
\end{lemma}

\begin{proof}
Besides its column expression, a quantum minor has
the reverse row expression
\begin{equation}\label{rc-row-form}
 Q_{I,J}(u)=
 \sum_{\tau\in S_q}\operatorname{sgn}(\tau)
 \prod_{\alpha=q}^{1}
 \bigl((u-\alpha+1)\delta_{i_\alpha,j_{\tau(\alpha)}}
                                  +E_{i_\alpha,j_{\tau(\alpha)}}\bigr).
\end{equation}
This follows from the anti-involution
$E_{ab}^*=E_{ba}$ and the adjoint identity in
Theorem~\ref{rp-capelli}: apply that identity to
$Q_{J,I}(\overline u)$ and reverse each product.
It is an identity in the enveloping algebra and
does not require a Hilbert-space realization.

Expand the determinant on the left of
\eqref{rc-border-first} in its first column.
The first term is
$(u+E_{ee})Q_{I,J}(u-1)$.
For row $i_\alpha$ the sign is $(-1)^\alpha$
and the first factor is $E_{i_\alpha,e}$.
Remove it using the left rule for $\Lambda$.
In the remaining minor the first row is $e$.
Expanding \eqref{rc-row-form} in that first row
places $E_{e,j_\beta}$ at the far right and
leaves the minor with parameter $u-2$.
Its sign is $(-1)^{\beta-1}$.
Remove the rightmost factor by the right rule.
The combined sign is
$(-1)^{\alpha+\beta-1}=-(-1)^{\alpha+\beta}$,
which proves \eqref{rc-border-first}.

For \eqref{rc-border-last}, expand in the last
column.  The diagonal term is
$Q_{I,J}(u)(u-q+E_{ee})$.
For row $i_\alpha$ the remaining minor has last
row $e$ and is followed by $E_{i_\alpha,e}$.
The column sign is $(-1)^{\alpha+q+1}$.
Remove the final factor by the right rule.
Expansion of \eqref{rc-row-form} in the last
row puts $E_{e,j_\beta}$ at the far left,
with sign $(-1)^{q+\beta}$; the minor left
over has parameter $u$.
The product of the two signs is again
$-(-1)^{\alpha+\beta}$.
This proves the second identity.
\end{proof}

\subsection{Mixed roots and the boundary contraction}

Let $F$ satisfy the smoothness and two character
laws \eqref{rt-whittaker}.  At a fixed cell point
$g=b(c)m$, evaluation of a right differential word
applied to $F$ gives a map $\Lambda$ as in
Lemma~\ref{rc-border}.
In addition to Lemma~\ref{rt-root-rules}, the
character laws give, for $1\leq l\leq r$,
\begin{align}
 E_{e_l,w_j}F&=\delta_{l,r}a'_jF,&
 E_{w_i,e_l^\vee}F&=\delta_{l,r}b'_iF,
                                                     \label{rc-mixed-right}\\
 (E_{w_i,e_l}XF)(b(c)m)
     &=\delta_{l,r}c_l^{-1}b_i(XF)(b(c)m),&
 (E_{e_l^\vee,w_j}XF)(b(c)m)
     &=\delta_{l,r}\bar c_l^{-1}a_j(XF)(b(c)m)
                                                     \label{rc-mixed-left}
\end{align}
for every right differential word $X$.

Here is a direct verification of the normalizations.
For the root $e_r w_j^*$ the real unitary root
matrix is
$t e_r w_j^*-\bar t\,\epsilon_jw_j(e_r^\vee)^*$.
Its contribution to $\ell_{z'}$ is $t z'_j$,
so its complex character derivative is $a'_j$.
For $w_i(e_r^\vee)^*$ the contribution is
$-\bar t\,\epsilon_i z'_i$, whose character
derivative is $\epsilon_i\overline{z'_i}=b'_i$.
Earlier isotropic indices give zero.
At a cell point conjugation gives
\[
 \operatorname{Ad}(b(c)m)(w_i e_l^*)
      =c_l^{-1}(mw_i)(e_l^\vee)^*,\qquad
 \operatorname{Ad}(b(c)m)(e_l^\vee w_j^*)
      =\bar c_l^{-1}e_l(w_j^*m^{-1}).
\]
The left character derivatives are respectively
$c_l^{-1}h(mw_i,z)=c_l^{-1}b_i$
and $\bar c_l^{-1}(m^{-1}z)_j$ when $l=r$.
Right differentiation preserves left equivariance,
which proves \eqref{rc-mixed-left} with the
arbitrary word $X$.

There is one sign in the boundary term which is
useful to make explicit.  For a fixed omitted
index $s$, write
$p_s(i)=i-\mathbf1_{i>s}$ for the position of
$i$ in the increasing list $\widehat s$,
starting at zero.  Grouping a double sum by
unordered row and column pairs gives
\begin{align}
 &\sum_{s,t}(-1)^{s+t}b'_s a_t
       \sum_{\substack{i\ne s\\j\ne t}}
        (-1)^{p_s(i)+p_t(j)}b_i a'_j
            Q^W_{\widehat{s,i},\widehat{t,j}}(v-1)
       =t_{22}(Y),                                  \label{rc-boundary-sign}\\
 &\sum_{s,t}(-1)^{s+t}b_s a'_t
       \sum_{\substack{i\ne s\\j\ne t}}
        (-1)^{p_s(i)+p_t(j)}b'_i a_j
            Q^W_{\widehat{s,i},\widehat{t,j}}(v-1)
       =t_{22}(Y).                                  \label{rc-boundary-sign-dual}
\end{align}
The lists with two omissions are always put in
increasing order.
For \eqref{rc-boundary-sign}, the two row orders
give $b'_s b_i-b'_i b_s$, which is the negative
of the row factor in \eqref{rc-t22}.
The two column orders give its column factor.
For \eqref{rc-boundary-sign-dual} the row
factor has the correct order and the column
factor changes sign.  In each case the minus
sign is exactly the one in \eqref{rc-t22}.
For $s<i,t<j$, the position corrections
subtract two from $s+i+t+j$ and hence do not
change its parity.  These observations verify
all four terms for every pair and prove
both identities.

\subsection{The ordinary and conjugate recurrences}

For $0\leq l\leq r$, let
\[
 I_{l,s}=(e_{l+1},\ldots,e_r,
                             w_0,\ldots,\widehat w_s,\ldots,w_k),
\]
and
\[
 J_{l,s}=(w_0,\ldots,\widehat w_s,\ldots,w_k,
                             e_r^\vee,\ldots,e_{l+1}^\vee).
\]
Let $\gamma$ and $\Phi$ be as in \eqref{rt-gauge}.
Define
\begin{align}
 F_l(c,m)
  &=\gamma(c)\sum_{s,t}(-1)^{s+t}b'_s a_t\,
       \bigl(Q_{I_{l,s},I_{l,t}}(u_l)F\bigr)(b(c)m),
       &u_l&=Y+r-l+\frac k2,                       \label{rc-F}\\
 F_l^-(c,m)
  &=\gamma(c)\sum_{s,t}(-1)^{s+t}b_s a'_t\,
       \bigl(Q_{J_{l,s},J_{l,t}}(v-1)F\bigr)(b(c)m).
                                                     \label{rc-Fminus}
\end{align}
The coefficients $a,b$ here are evaluated at $m$.
They are independent of $c$.

\begin{proposition}\label{rc-recurrences}
The terminal values are
\[
                F_r=-t_{21}(Y)\Phi,\qquad
                F_r^-=-t_{12}(Y)\Phi.
\]
For $r\geq1$ the boundary formulas are
\begin{align}
 F_{r-1}&=(Y+p_r)F_r-c_r^{-1}t_{22}(Y)\Phi,
                                                     \label{rc-first-boundary}\\
 F_{r-1}^-&=(Y-\bar p_r)F_r^-
                              -\bar c_r^{-1}t_{22}(Y)\Phi .
                                                     \label{rc-last-boundary}
\end{align}
For $0\leq l\leq r-2$,
\begin{align}
 F_l&=(Y+p_{l+1})F_{l+1}
                       -\frac{c_{l+2}}{c_{l+1}}F_{l+2},
                                                     \label{rc-first-bulk}\\
 F_l^-&=(Y-\bar p_{l+1})F_{l+1}^-
                       -\frac{\bar c_{l+2}}{\bar c_{l+1}}F_{l+2}^- .
                                                     \label{rc-last-bulk}
\end{align}
\end{proposition}

\begin{proof}
Right derivatives with indices in $W$ act on
$F(b(c)m)$ as the corresponding right derivatives
in $m$.  The scalar $\gamma$ is independent of
$m$.  The definitions \eqref{rc-t21} and
\eqref{rc-t12} therefore give the terminal values.

For \eqref{rc-first-bulk}, apply
\eqref{rc-border-first} with distinguished
index $e_{l+1}$, before inserting the outside
coefficients in \eqref{rc-F}.
All left and right root constants vanish except
for the next index $e_{l+2}$.  The product of
its two constants is $c_{l+2}/c_{l+1}$.
The remaining parameters are $u_l-1=u_{l+1}$
and $u_l-2=u_{l+2}$.
The diagonal right derivative becomes $p_{l+1}$
on the cell.  Conjugating it by $\gamma$ changes
$u_l+p_{l+1}$ into $Y+p_{l+1}$, since
$p_{l+1}\log\gamma=r-l+k/2$.
The outside coefficients are independent of this
coordinate, so the derivative does not create
an additional term.

For $l=r-1$, the remaining indices are the
$W$-indices with one omission.
Equations \eqref{rc-mixed-right} and
\eqref{rc-mixed-left} give constants
$c_r^{-1}b_i$ and $a'_j$.
The contracted double sum is exactly
$c_r^{-1}t_{22}(Y)$ by
\eqref{rc-boundary-sign}.  The same scalar
conjugation proves \eqref{rc-first-boundary}.

For the conjugate recurrence apply
\eqref{rc-border-last}, with distinguished
index $e_{l+1}^\vee$ appended to the two lists.
Its diagonal generator commutes with every generator
in the remaining minor, since that minor has no row
or column with the distinguished index.  We may therefore
move the diagonal factor to the left before restricting
to the cell.
The right derivative in its diagonal is
$-\bar p_{l+1}$.
If the original size of the minor is
$r-l+k$, its diagonal scalar is
\[
       (v-1)-(r-l+k-1)=Y-r+l-\frac k2 .
\]
Conjugating $-\bar p_{l+1}$ by $\gamma$
adds $r-l+k/2$, leaving $Y-\bar p_{l+1}$.
All parameters of the smaller minors remain
$v-1$, exactly as in \eqref{rc-Fminus}.

Away from the boundary the only surviving
pair of root constants is the adjacent dual
pair, with product
$\bar c_{l+2}/\bar c_{l+1}$.
At the boundary their values are
$b'_i$ and $\bar c_r^{-1}a_j$ by
\eqref{rc-mixed-right} and
\eqref{rc-mixed-left}.
Identity \eqref{rc-boundary-sign-dual}
identifies the double contraction with
$\bar c_r^{-1}t_{22}(Y)$.
This proves the conjugate formulas.
\end{proof}

\begin{corollary}\label{rc-chain}
With the transfer columns of
Theorem~\ref{rt-actual-lifts},
\begin{align}
 F_0&=-t_{21}(Y)A_r(Y)\Phi
                              -t_{22}(Y)B_r(Y)\Phi,\label{rc-chain-plus}\\
 F_0^-&=-t_{12}(Y)\overline A_r(Y)\Phi
                              -t_{22}(Y)\overline B_r(Y)\Phi .
                                                     \label{rc-chain-minus}
\end{align}
\end{corollary}

\begin{proof}
All terminal operators act only on $m$ and commute
with coordinate operators in $c$.
Expand the recurrence \eqref{rc-first-bulk}
over the path of vertices $1,\ldots,r$.
The coefficient of $F_r$ is the sum over
disjoint adjacent edges, with diagonal factors
$Y+p_i$ at uncovered vertices and edge
factors $-c_{i+1}/c_i$.
As in the proof of Theorem~\ref{rt-actual-lifts},
this is $A_r(Y)$.
The coefficient of the other boundary term is
the same sum on $1,\ldots,r-1$, multiplied by
$c_r^{-1}$; it equals
$c_r^{-1}A_{r-1}(Y)=B_r(Y)$.
Substituting the terminal values proves
\eqref{rc-chain-plus}.
The same expansion, with $p_i$ replaced by
$-\bar p_i$ and $c_i$ by $\bar c_i$, gives
\eqref{rc-chain-minus}.
For $r=0$ the assertions reduce to the terminal
identities and $B_0=\overline B_0=0$.
\end{proof}

\subsection{Polynomial coefficients on the full group}

Let $\operatorname{vol}$ be the alternating
form taking value one on the displayed ordered
basis of $V$.  Let $M=r+k$.
For an increasing $M$-element list $J$, define
\begin{align}
 B_J(g)&=\operatorname{vol}
       (g^{-1}e_r,\ldots,g^{-1}e_1,g^{-1}z,e_J),
                                                     \label{rc-global-B}\\
 A_J(g)&=\det
       \begin{pmatrix}
        h(g\,\cdot,e_1)\\[-2pt]
        \vdots\\
        h(g\,\cdot,e_r)\\
        h(g\,\cdot,z)\\
        e_J^*
       \end{pmatrix}.
                                                     \label{rc-global-A}
\end{align}
The last $M$ rows in \eqref{rc-global-A}
are the coordinate covectors in the list $J$.
All matrix sizes are $2r+k+1$.
Set $\sigma=(-1)^{r(k+1)}$ and define
\begin{align}
 \mathcal R_Y F(g)
  &=-\sigma\sum_{i,J}(-1)^i b'_i B_J(g)\,
                 Q_{I_{0,i},J}(u_0)F(g),
                                                     \label{rc-global-R}\\
 \mathcal R_Y^- F(g)
  &=-\sum_{t,J}(-1)^t a'_t A_J(g)\,
                 Q_{J,J_{0,t}}(v-1)F(g).
                                                     \label{rc-global-Rminus}
\end{align}
These are the actual full-group coefficient
determinants in the source, with their constant
signs chosen explicitly.

\begin{lemma}\label{rc-cell-coefficients}
At $g=b(c)m$, all terms in
\eqref{rc-global-B} vanish except
$J=I_{0,t}$, and
\[
                 B_{I_{0,t}}(b(c)m)
                       =\sigma(-1)^t\overline{a(c)}\,a_t .
                                                     \tag{RC1}
\]
All terms in \eqref{rc-global-A} vanish except
$J=J_{0,s}$, and
\[
                 A_{J_{0,s}}(b(c)m)
                       =(-1)^s a(c)b_s .             \tag{RC2}
\]
\end{lemma}

\begin{proof}
In \eqref{rc-global-B} the first $r$ columns
are $\bar c_r e_r^\vee,\ldots,\bar c_1e_1^\vee$,
and the next column is $m^{-1}z$.
Any nonzero complementary coordinate determinant
must therefore contain all $U$-indices and all
but one of the $W$-indices.
Move the first $r$ columns past the other
$r+k+1$ columns, and then move the $m^{-1}z$
column past the $r$ coordinate columns of $U$.
These moves give the sign
$(-1)^{r(r+k+1)+r}=(-1)^{r(k+1)}$.
Expansion in its omitted $W$-coordinate gives
the extra sign $(-1)^t$.
This proves (RC1).

In \eqref{rc-global-A} the first $r$ rows are
$c_i e_i^*$, while the next row vanishes on
$U\oplus U^\vee$ and has $W$-coordinates $b_s$.
The complementary rows must therefore be the
dual indices and all but one $W$-index.
Expansion of that next row in column $r+s$
has sign $(-1)^{r+(r+s)}=(-1)^s$.
The remaining coordinate determinant is one.
This proves (RC2).
\end{proof}

\begin{theorem}\label{rc-actual-lifts}
The operators \eqref{rc-global-R} and
\eqref{rc-global-Rminus} have order at most
$r+k$ and polynomial coefficients of degree
$r+1$.  They are defined on the entire group.
They commute with left translation by $K_zN$
and right translation by $K_{z'}N$.
They and their Haar formal adjoints preserve
$\mathscr S_\varepsilon$.

For every $F$ satisfying \eqref{rt-whittaker},
their exact radial identities are
\begin{align}
 \gamma(c)(\mathcal R_YF)(b(c)m)
   &=\overline{a(c)}
       \bigl(t_{21}(Y)A_r(Y)+t_{22}(Y)B_r(Y)\bigr)\Phi,
                                                     \label{rc-radial-R}\\
 \gamma(c)(\mathcal R_Y^-F)(b(c)m)
   &=a(c)
       \bigl(t_{12}(Y)\overline A_r(Y)
                    +t_{22}(Y)\overline B_r(Y)\bigr)\Phi .
                                                     \label{rc-radial-Rminus}
\end{align}
The statements hold for every complex $Y$.
\end{theorem}

\begin{proof}
The coefficients $B_J$ have degree $r+1$
in $\overline g$, since $g^{-1}=Hg^*H$.
The coefficients $A_J$ have degree $r+1$
in $g$.  The order and polynomial assertions
follow from the sizes of the quantum minors.
The proof of Proposition~\ref{rt-polynomial}
also gives derivative bounds
$C_m(1+\|g\|)^{r+1}$ for these coefficients
and proves preservation of
$\mathscr S_\varepsilon$ by both operators
and their formal adjoints.

We give the full covariance argument to account
for a determinant factor which cannot simply
be omitted.
The vector
$\eta=e_r\wedge\cdots\wedge e_1\wedge z$
is fixed by $\bigwedge^{r+1}s_L$ for
$s_L\in K_zN$: the $U$-block is unipotent
and $s_Lz-z\in U$.
Hence $B_J(s_Lg)=B_J(g)$.
The corresponding wedge of covectors in
\eqref{rc-global-A} is likewise unchanged,
because unitarity replaces each fixed vector
by its image under $s_L^{-1}$.
Left translations commute with right
derivatives, proving left covariance.

Now let $s\in K_{z'}N$.
Put $P=U\oplus W$ and
$\ell(x)=h(x,z')$ for $x\in P$.
The action of $s$ preserves $P$ and satisfies
$\ell(sx)=\ell(x)$.
Its determinant on $P$ equals its determinant
on $V$: the actions on $U$ and on
$V/P$ have determinant one.
Thus the fixed row vector in the exterior
power,
\[
       \rho=\sum_i(-1)^i b'_i e_{I_{0,i}}
          =(-1)^r\iota_\ell
                  (e_1\wedge\cdots\wedge e_r
                         \wedge w_0\wedge\cdots\wedge w_k),
\]
satisfies
$\bigwedge^M(s^{-1})\rho=\det(s)^{-1}\rho$.
On the other side, the volume definition gives
\[
 \sum_J B_J(g)\det(s)_{J,L}
                  =\det(s) B_L(gs).
                                                     \tag{RC3}
\]
Indeed the left side is
$\operatorname{vol}(g^{-1}\eta\wedge s e_L)$;
apply $s^{-1}$ to every vector in this
volume.  Formula \eqref{rt-right-minor}
and (RC3) cancel the two determinant
factors and prove right covariance of
$\mathcal R_Y$.

For the second operator, the elementary
dual volume identity is
\[
 \sum_J\det(s^{-1})_{K,J}A_J(g)
                      =\det(s)^{-1}A_K(gs).
                                                     \tag{RC4}
\]
It follows either by applying $s^{-1}$ to
the coordinate rows in the full determinant,
or by the alternating Cauchy--Binet formula.
The fixed column covector is
\[
                  \kappa=\sum_t(-1)^t a'_t e_{J_{0,t}}^*.
 \]
On $V/U$, it is contraction of its standard
top-degree form by the vector $z'$.
The action of $s$ fixes this vector modulo
$U$ and has determinant $\det(s)$ on the
quotient.  Therefore
$\kappa\circ\bigwedge^M s=\det(s)\kappa$.
This factor cancels the one in (RC4) when
\eqref{rt-right-minor} is applied.
It proves right covariance of
$\mathcal R_Y^-$.

Finally (RC1), (RC2), and $\sigma^2=1$ give
\[
 \gamma(\mathcal R_YF)(b(c)m)
       =-\overline{a(c)}F_0,\qquad
 \gamma(\mathcal R_Y^-F)(b(c)m)
       =-a(c)F_0^- .
 \]
Corollary~\ref{rc-chain} proves
\eqref{rc-radial-R} and
\eqref{rc-radial-Rminus}.
The source sometimes writes the conjugate
transfer recurrence with its new coordinate
factor on the right.  Earlier-coordinate
operators commute with that factor, so it
is exactly the column used here.
\end{proof}

\paragraph{Scope.}
The four required polynomial differential
lifts and the regular spectral inverse have
now been constructed and matched with their
actual radial actions.  This is an algebraic
and differential statement.  The terminal
quantum relation used to cancel the two
mixed terms in the Fourier identity, and
the analytic passage to the Jacquet pairing
and its uniformly bounded generated class,
remain distinct obligations.  No integration
by parts for that pairing is inferred from
the ordinary Haar formal-adjoint calculation.


\section{The determinant identity for the terminal differential operators}
\label{rq-section}

The terminal matrix is not a matrix of commuting operators.
Its determinant identity has a spectral shift, and its
coefficients are functions which the derivatives can differentiate.
We prove the identity with both features retained.
The proof uses quantum minors, elementary exterior algebra,
and a polynomial cancellation argument.  It applies directly
to smooth functions and does not require an approximation by
regular functions.  The resulting identity is the source's
terminal determinant formula \cite{alpo-uniform-source}.

\subsection{Signed complementary minors}

Let $n\geq2$, let $\mathcal A$ be a unital complex algebra, and let
$E_{ab}\in\mathcal A$, $1\leq a,b\leq n$, satisfy
\[
 [E_{ab},E_{cd}]=\delta_{bc}E_{ad}-\delta_{da}E_{cb}.
\]
Use the ordered minors $Q_{I,J}(u)$ from
\eqref{rt-minor}.  Put
\[
 T_{ab}(u)=u\delta_{ab}+E_{ab},\qquad
 D(u)=Q_{\{1,\ldots,n\},\{1,\ldots,n\}}(u),
 \qquad C_{ij}(u)=(-1)^{i+j}Q_{\widehat i,\widehat j}(u).
\]
For distinct pairs define
\[
 N_{iqjp}(u)=
 \epsilon(i,q)\epsilon(j,p)(-1)^{i+q+j+p}
 Q_{\widehat{i,q},\widehat{j,p}}(u),
 \qquad
 \epsilon(i,q)=
 \begin{cases}1&i<q,\\-1&i>q.\end{cases}
\]
Set $N_{iqjp}=0$ when $i=q$ or $j=p$.
All omission lists are increasing.  Empty minors are one.
Here $C$ is indexed by the omitted row and omitted column;
it has not been transposed into an adjugate matrix.

We use the alternating and change-of-basis identities
for quantum minors proved in Theorem~\ref{rp-capelli}.
In particular $D(u)$ commutes with every $E_{ab}$:
the full determinant is unchanged by conjugation of the
generators, and differentiating this identity at the
identity matrix gives the asserted commutators.
It therefore commutes with every quantum minor.

\begin{lemma}\label{rq-cofactor-expansions}
For all indices and all $u\in\mathbb C$,
\begin{align}
 \sum_p C_{qp}(u)T_{cp}(u+1)
       &=\delta_{cq}D(u+1),                 \label{rq-cofactor-one}\\
 \sum_p N_{iqjp}(u)T_{cp}(u+1)
       &=\delta_{cq}C_{ij}(u+1)
                   -\delta_{ci}C_{qj}(u+1). \label{rq-cofactor-two}
\end{align}
\end{lemma}

\begin{proof}
The column expression for $Q$ and its reverse row
expression \eqref{rc-row-form} give the usual four
end-row and end-column expansions, with their displayed
shifts.  For example, expansion in the last row, written
with the last row factor on the right, gives
\[
 Q_{I,J}(v)=
 \sum_{\beta=1}^{s}(-1)^{s+\beta}
 Q_{I\setminus i_s,J\setminus j_\beta}(v-1)
 T_{i_sj_\beta}(v).
\]
An arbitrary row can be moved to the last position;
alternation contributes its permutation sign.
Equivalently, one may put an arbitrary row first in the
increasing-product row expression
\[
 Q_{I,J}(v)=
 \sum_{\tau\in S_s}\operatorname{sgn}(\tau)
 \prod_{\alpha=1}^{s}
 T_{i_\alpha,j_{\tau(\alpha)}}(v-s+\alpha).
\]
This expression follows by reversing the row list in
\eqref{rc-row-form} and using alternation.
These expansions are identities in $\mathcal A$;
they do not commute adjacent factors.

For \eqref{rq-cofactor-one}, append row $c$ to the list
with row $q$ removed, and expand along the appended row.
If $c\ne q$, a row is repeated and the minor is zero.
If $c=q$, moving that row back to its original position
gives $D(u+1)$ with the sign already included in
$C_{qp}$.

For \eqref{rq-cofactor-two}, start instead with the row
list omitting $i,q$ and the column list omitting $j$,
append row $c$, and expand in that row.
The result is zero unless $c=i$ or $c=q$.
In those two cases moving the appended row into its
increasing position gives respectively
$-C_{qj}(u+1)$ and $C_{ij}(u+1)$.
For precision, removing $q$ changes the position of
an index $i$ to $i-\mathbf1_{i>q}$, and removing $j$
changes the position of $p$ to $p-\mathbf1_{p>j}$.
The two resulting signs are exactly
$\epsilon(i,q)\epsilon(j,p)(-1)^{i+q+j+p}$.
This verifies every sign in
\eqref{rq-cofactor-two}, including the zero terms.
\end{proof}

\begin{theorem}\label{rq-jacobi}
The signed quantum minors satisfy
\begin{equation}\label{rq-jacobi-identity}
 C_{ij}(u+1)C_{qp}(u)-C_{qj}(u+1)C_{ip}(u)
                  =D(u+1)N_{iqjp}(u).
\end{equation}
No assumption that $\mathcal A$ is a domain is needed.
\end{theorem}

\begin{proof}
Fix $i,q,j$ and let $Z_p$ be the difference of the
two sides of \eqref{rq-jacobi-identity}.  Multiply
$Z_p$ on the right by $T_{cp}(u+1)$ and sum over $p$.
Lemma~\ref{rq-cofactor-expansions} gives
\begin{align*}
 \sum_p Z_pT_{cp}(u+1)
 &=\delta_{cq}C_{ij}(u+1)D(u+1)
       -\delta_{ci}C_{qj}(u+1)D(u+1)\\
 &\quad-D(u+1)
       \bigl(\delta_{cq}C_{ij}(u+1)
                       -\delta_{ci}C_{qj}(u+1)\bigr)=0.
\end{align*}
The last equality uses the centrality of $D$.
It remains to justify cancellation without assuming
invertibility of the matrix $T$.

Adjoin a central indeterminate $t$ and replace
$E_{ab}$ by $E_{ab}+t\delta_{ab}$.
These new generators satisfy the same Lie relations.
The preceding calculation gives polynomials
$Z_p(t)\in\mathcal A[t]$ satisfying
\[
 \sum_p Z_p(t)\bigl(E_{cp}+(u+1+t)\delta_{cp}\bigr)=0
                       \qquad\text{for every }c.
\]
If some $Z_p(t)$ is nonzero, let $d$ be the largest
degree of all of them.  The coefficient of $t^{d+1}$
in the equation with index $c$ is exactly the
degree-$d$ coefficient of $Z_c(t)$.  It vanishes
for every $c$, a contradiction.  Thus all $Z_p(t)$
are zero.  Substituting $t=0$ proves the theorem.
\end{proof}

\subsection{Moving the variable coefficients through a minor}

Suppose $\mathcal A$ also contains elements
$y_1,\ldots,y_n,\eta_1,\ldots,\eta_n$ such that
\begin{equation}\label{rq-vector-laws}
 [E_{ab},y_c]=-\delta_{ac}y_b,\qquad
 [E_{ab},\eta_c]=\delta_{bc}\eta_a.
\end{equation}
Assume that all the $y$'s and $\eta$'s commute with
one another.  They need not commute with the $E$'s.

\begin{lemma}\label{rq-moving}
Under these assumptions,
\begin{align}
 D(w)y_j
       &=y_jD(w)-\sum_l y_lC_{jl}(w),          \label{rq-det-vector}\\
 [C_{qp}(u),\eta_i]
       &=\sum_a N_{aqip}(u)\eta_a,             \label{rq-minor-covector}\\
 \sum_i\eta_iC_{ij}(u-1)
       &=\sum_iC_{ij}(u)\eta_i,               \label{rq-shift-one}\\
 \sum_a\eta_aN_{aqcp}(u-1)
       &=\sum_aN_{aqcp}(u)\eta_a.             \label{rq-shift-two}
\end{align}
\end{lemma}

\begin{proof}
For \eqref{rq-det-vector}, use the row expansion of
$D(w)$ which puts row $j$ in the far-left factor:
\[
 D(w)=\sum_l(-1)^{j+l}
                 T_{jl}(w-n+1)Q_{\widehat j,\widehat l}(w).
\]
Every generator in the smaller minor commutes with
$y_j$, because its row is different from $j$.
Commuting $y_j$ through the remaining factor gives
$[T_{jl}(w-n+1),y_j]=-y_l$.
This proves \eqref{rq-det-vector}.

For \eqref{rq-minor-covector}, the assertion is zero
if $i=p$, since column $p$ is absent.
Otherwise expand the minor defining $C_{qp}(u)$
in column $i$, putting that column factor on the right.
All factors in the smaller minors commute with
$\eta_i$ because their columns differ from $i$.
The remaining commutator is
$[T_{ai}(u-n+2),\eta_i]=\eta_a$.
The position of $a$ after removing $q$ is
$a-\mathbf1_{a>q}$ and that of $i$ after removing
$p$ is $i-\mathbf1_{i>p}$.
The signs from these positions, together with
$(-1)^{q+p}$, are precisely those in
$N_{aqip}(u)$.  This proves the second identity.

Here is an exterior-algebra proof of both shift
identities, including their parameter change.
Adjoin symbols $\theta_1,\ldots,\theta_n$ commuting
with $\mathcal A$ and satisfying
$\theta_a\theta_b=-\theta_b\theta_a$ and $\theta_a^2=0$.
Set
\[
 \beta_j(v)=\sum_a\theta_aT_{aj}(v),\qquad
                    \eta=\sum_a\theta_a\eta_a.
\]
Relation \eqref{rq-vector-laws} gives, term by term,
\[
                   \beta_j(v)\eta=-\eta\beta_j(v-1).
\]
Indeed the commutator contribution is
$\sum_a\theta_a\theta_j\eta_a=-\theta_j\eta$.
Repeatedly using this equality yields
\[
 \beta_{j_1}(u)\beta_{j_2}(u-1)\cdots
       \beta_{j_s}(u-s+1)\eta
   =(-1)^s\eta\,\beta_{j_1}(u-1)\cdots\beta_{j_s}(u-s).
\]
The coefficient of
$\theta_{i_1}\cdots\theta_{i_{s+1}}$ is
\begin{equation}\label{rq-general-shift}
 \sum_{\alpha=1}^{s+1}(-1)^{\alpha-1}
       \eta_{i_\alpha}Q_{I\setminus i_\alpha,J}(u-1)
  =
 \sum_{\alpha=1}^{s+1}(-1)^{\alpha-1}
       Q_{I\setminus i_\alpha,J}(u)\eta_{i_\alpha}.
\end{equation}
This follows directly from the column definition
of the minors.  Taking $I$ to be the full index
list and $J$ to omit $j$ proves
\eqref{rq-shift-one}.  Taking $I$ to omit $q$
and $J$ to omit $c,p$, then multiplying by the
fixed omission signs, proves \eqref{rq-shift-two}.
\end{proof}

\subsection{The contracted determinant identity}

For fixed complex vectors $x=(x_j)$ and
$\xi=(\xi_i)$ define
\begin{align*}
 S(u;x,\xi)&=\sum_{i,j}\xi_i x_j C_{ij}(u),\\
 H_y(u;\xi)&=\sum_{i,j}\xi_i y_j C_{ij}(u),\\
 H_\eta(u;x)&=\sum_{i,j}x_j\eta_i C_{ij}(u),\\
 H_{y\eta}(u)&=\sum_{i,j}\eta_i y_j C_{ij}(u),\\
 H_2(u;x,\xi)&=
 \sum_{\substack{i<i'\\j<j'}}(-1)^{i+i'+j+j'}
  (\eta_i\xi_{i'}-\eta_{i'}\xi_i)
  (y_jx_{j'}-y_{j'}x_j)
            Q_{\widehat{i,i'},\widehat{j,j'}}(u).
\end{align*}
In every expression the coefficient precedes the minor.

\begin{theorem}\label{rq-contracted}
One has the identity in $\mathcal A$
\begin{equation}\label{rq-contracted-identity}
 D(u+1)H_2(u-1;x,\xi)
     +H_y(u+1;\xi)H_\eta(u-1;x)
                    =H_{y\eta}(u)S(u;x,\xi).
\end{equation}
\end{theorem}

\begin{proof}
Both sides are bilinear in $x,\xi$, so fix $p,q$
and take $x_j=\delta_{jp}$, $\xi_i=\delta_{iq}$.
For brevity write
\[
 C^+_{ij}=C_{ij}(u+1),\quad
 C^0_{ij}=C_{ij}(u),\quad C^-_{ij}=C_{ij}(u-1),
 \quad N^0_{aqcp}=N_{aqcp}(u).
\]
Expanding the two alternating factors in $H_2$
and grouping their four terms gives
\[
 H_2(u-1;x,\xi)=\sum_{a,c}\eta_a y_cN_{aqcp}(u-1).
\]
By \eqref{rq-shift-two} and the commutativity of
$\eta_a$ with $y_c$, this is
$\sum_c y_c\sum_aN^0_{aqcp}\eta_a$.
Likewise \eqref{rq-shift-one} gives
\[
 H_\eta(u-1;x)=\sum_iC^0_{ip}\eta_i,\qquad
 H_{y\eta}(u)=\sum_l y_l\sum_i C^+_{il}\eta_i.
\]
Let
\[
                  \mathcal E=
       \sum_{l,c,a}y_l C^+_{cl}N^0_{aqcp}\eta_a .
\]
Using \eqref{rq-det-vector} in the first summand
of the proposed identity gives
\[
 D(u+1)H_2(u-1;x,\xi)
       =\sum_{l,a}y_lD(u+1)N^0_{aqlp}\eta_a-\mathcal E.
\]
The second summand is
$\sum_{l,a}y_l C^+_{ql}C^0_{ap}\eta_a$.
Their sum is therefore
\[
 \sum_{l,a}y_l
   \bigl(D(u+1)N^0_{aqlp}+C^+_{ql}C^0_{ap}\bigr)\eta_a
                       -\mathcal E
   =\sum_{l,a}y_l C^+_{al}C^0_{qp}\eta_a-\mathcal E
\]
by Theorem~\ref{rq-jacobi}.
On the other hand,
\eqref{rq-minor-covector} gives
\[
 H_{y\eta}(u)C^0_{qp}
       =\sum_{l,a}y_l C^+_{al}C^0_{qp}\eta_a-\mathcal E.
\]
This is the same expression.  Bilinearity now proves
\eqref{rq-contracted-identity}.
\end{proof}

\subsection{Application to the actual terminal group}

Return to $W$ of signature $(k,1)$ and its group
$U(W,h)$, with $h$ linear in the first argument.
For a smooth function on this group let $E_{ab}$
be its complexified right derivatives.
Multiplication by
\[
 y_j(m)=(m^{-1}z)_j,\qquad
       \eta_i(m)=h(mw_i,z)=h(w_i,m^{-1}z)
\]
defines commuting operators.
Differentiation of $m^{-1}z$ and of the row
$z^*Hm$ gives exactly
\eqref{rq-vector-laws}.
Take $x_j=z'_j$, $\xi_i=h(w_i,z')$ and
$u=Y+k/2$.  Equations
\eqref{rc-t21}--\eqref{rc-t22} become
\[
 t_{11}(Y)=D(u),\quad
 t_{21}(Y)=-H_y(u;\xi),\quad
 t_{12}(Y)=-H_\eta(u-1;x),\quad
 t_{22}(Y)=-H_2(u-1;x,\xi).
\]
Thus Theorem~\ref{rq-contracted} proves
\begin{equation}\label{rq-terminal-factorization}
 t_{11}(Y+1)t_{22}(Y)-t_{21}(Y+1)t_{12}(Y)
                         =-H_{y\eta}(u)S(u;x,\xi).
\end{equation}
All compositions here are compositions of genuine
differential operators with variable coefficients.

Choose orthonormal bases $B_1,\ldots,B_k$ of $z^\perp$
and $B'_1,\ldots,B'_k$ of $(z')^\perp$ in $W$.
Write $C^R_{B'}(u)$ for the degree-$k$ Capelli
determinant of the right derivatives along
$B'_a h(\,\cdot\,,B'_b)$.
Write $C^L_B(u)$ for the same polynomial in left
derivatives, with every product of derivatives reversed.
This reversal is part of the definition.

\begin{theorem}\label{rq-terminal-determinant}
For every smooth function $F$ on the terminal group,
\begin{equation}\label{rq-exact-terminal}
 \bigl(t_{11}(Y+1)t_{22}(Y)-t_{21}(Y+1)t_{12}(Y)\bigr)F
          =-TT'\,C^L_B(Y+k/2)C^R_{B'}(Y+k/2)F.
\end{equation}
The formula holds for every complex $Y$ and needs
neither nonvanishing intermediate minors nor a
positivity restriction on the Gram polynomial.
\end{theorem}

\begin{proof}
In the ordered basis
$B'_1,\ldots,B'_k,z'/\sqrt{T'}$,
the vector $x$ has only its last entry nonzero,
equal to $\sqrt{T'}$, and the covector $\xi$
has last entry $-\sqrt{T'}$.
The change-of-basis rule for minors therefore gives
\[
                       S(u;x,\xi)=-T'C^R_{B'}(u).
\]
At a fixed $m$, apply the same calculation to
$m^{-1}z$, whose Hermitian norm is $-T$.
An orthonormal frame of its perpendicular space is
$m^{-1}B_1,\ldots,m^{-1}B_k$.  Since the
coefficients in $H_{y\eta}$ precede the derivatives,
they are frozen at this evaluation point.  Consequently
\[
 (H_{y\eta}(u)G)(m)
                    =-T(C^R_{m^{-1}B}(u)G)(m).
\]
There is no derivative of the frozen frame in this
identity.

For completeness, if $X_a=\operatorname{Ad}(m^{-1})Y_a$
are fixed tangent matrices at this point, then
\[
 (R_{X_1}\cdots R_{X_s}G)(m)
                       =(L_{Y_s}\cdots L_{Y_1}G)(m).
\]
Indeed both sides are the mixed derivative at zero
of
$G(\exp(t_1Y_1)\cdots\exp(t_sY_s)m)$.
Complexification gives the same identity for the
matrix units defining the Capelli determinant.
This proves
$(C^R_{m^{-1}B}(u)G)(m)=(C^L_B(u)G)(m)$,
including the required reversal of each product.
Substitution into
\eqref{rq-terminal-factorization} proves
\eqref{rq-exact-terminal}.
\end{proof}

\begin{corollary}\label{rq-type-scalar}
Suppose
$F(hmh')=\pi(h)F(m)\pi'(h')$ for
$h\in U(z^\perp)$ and $h'\in U((z')^\perp)$,
where $\pi,\pi'$ are irreducible finite-dimensional
unitary representations with highest weights
$\lambda=(\lambda_1\geq\cdots\geq\lambda_k)$ and
$\lambda'=(\lambda'_1\geq\cdots\geq\lambda'_k)$.
Then the left side of \eqref{rq-exact-terminal} equals
\[
 -TT'\prod_{j=1}^k
   (Y+k/2-j+1+\lambda_j)
   (Y+k/2-j+1+\lambda'_j)\,F.
\]
In the source's parameters
$u_j=\lambda_j+(k+1)/2-j$ and
$u'_j=\lambda'_j+(k+1)/2-j$, this scalar is
\[
                 -TT'\prod_{j=1}^k
                    (Y+\tfrac12+u_j)(Y+\tfrac12+u'_j).
\]
\end{corollary}

\begin{proof}
The right derivatives act through the derived
representation $\pi'$ on the right of $F$.
Left derivatives act in the reverse order through
$\pi$ on its left; the reversal in $C^L_B$
therefore restores the ordinary Capelli polynomial.
Its value in highest weight $\lambda$ is
$\prod_j(u-j+1+\lambda_j)$ by
Theorem~\ref{rp-capelli}.
Apply this on both sides of
\eqref{rq-exact-terminal} and substitute $u=Y+k/2$.
\end{proof}

\paragraph{Scope.}
This proves the terminal determinant formula and its
actual finite-type scalar, including all variable
coefficient corrections.  It does not assert the
remaining commutation relations of the complete
terminal matrix, the recursive frame identity, or
integration by parts for the Jacquet pairing.
Those assertions have their own hypotheses and are
not consequences of this determinant identity alone.

The spectral shift cannot be suppressed.  In dimension
two, take the trivial representation $E_{ab}=0$ and
$u=1$.  Then $C_{11}(1)C_{22}(1)=1$, while
$D(1)N_{1212}(1)=0$.  The unshifted version of
\eqref{rq-jacobi-identity} is therefore false even in
this elementary representation.  The submitted
identity has the shift proved above, so this is a
counterexample only to that stronger, altered statement.


\section{Commutation relations of the terminal matrix}
\label{rr-section}

We prove the full commutation relation of the actual terminal
matrix by adding one isotropic pair.  A full Capelli determinant
on the larger unitary group restricts to a product of three
$2\times2$ matrices.  The full determinants commute.
An elementary independence argument then recovers all sixteen
commutation relations of the middle matrix.
The smooth extension needed for this argument is given explicitly.

\subsection{The open cell with one isotropic pair}

Put $d=k+1$, let $J=\operatorname{diag}(1,\ldots,1,-1)$
on $W=\mathbb C^d$, and let
\[
 V=\mathbb C e\oplus W\oplus\mathbb C f,\qquad
 H=\begin{pmatrix}0&0&1\\0&J&0\\1&0&0\end{pmatrix}.
\]
For $v\in W$, $t\in\mathbb R$, and $c\in\mathbb C^\times$,
write
\[
 n(v,t)=
 \begin{pmatrix}
 1&-v^*J&it-\tfrac12v^*Jv\\
 0&I&v\\0&0&1
 \end{pmatrix},
 \qquad
 w(c)=
 \begin{pmatrix}0&0&\bar c^{-1}\\0&I&0\\c&0&0\end{pmatrix}.
\]
Both matrices preserve $H$.  The matrices $n(v,t)$ form
the unipotent group $N$.
For $z\in W$ set
\[
 \ell_z(n(v,t))=-v^*Jz,\qquad
              \chi_z(n)=\exp(\ell_z(n)-\overline{\ell_z(n)}).
\]
The coordinate $v$ is additive under multiplication in
$N$, so $\chi_z$ is a character.  This is the character
normalization used in the preceding radial calculations.

\begin{lemma}\label{rr-cell}
The map
\[
 N\times\mathbb C^\times\times U(W,J)\times N
   \longrightarrow \{g\in U(V,H):g_{fe}\ne0\},\qquad
 (n_L,c,m,n_R)\longmapsto n_Lw(c)m n_R
\]
is a smooth bijection with smooth inverse.
Here $m$ acts as the identity on $\mathbb Ce\oplus\mathbb Cf$.
\end{lemma}

\begin{proof}
Let $g$ be in the displayed open set and put $c=g_{fe}$.
Multiplication of the four displayed block matrices shows
that the only possible parameters satisfy
\[
 v_L=c^{-1}g_{We},\qquad
 v_R^*J=-c^{-1}g_{fW},\qquad
 a_L=c^{-1}g_{ee},\qquad a_R=c^{-1}g_{ff},
\]
where $a_\nu=it_\nu-\tfrac12v_\nu^*Jv_\nu$.
The zero Hermitian norm of column $e$ of $g$ gives
$2\operatorname{Re}a_L=-v_L^*Jv_L$.
The analogous row identity, obtained from $gHg^*=H$,
gives $2\operatorname{Re}a_R=-v_R^*Jv_R$.
Thus $t_L,t_R$ are uniquely the imaginary parts of
$a_L,a_R$.

Remove the two resulting unipotent factors.
The matrix $n_L^{-1}gn_R^{-1}$ has first column $cf$
and last row $ce^*$.
Its other columns are orthogonal to that first column,
so their $e$-coordinates vanish except in the last
column.  The analogous row orthogonality makes the
remaining off-diagonal blocks zero.  The pairing of
the first and last columns then fixes the $(e,f)$
entry to be $\bar c^{-1}$.
The remaining $W$-block $m$ satisfies $m^*Jm=J$.
It equals
$g_{WW}+v_Lc\,v_R^*J$ by multiplication of the factors.
This proves both existence and uniqueness.
All formulas use smooth operations and division only
by the nonzero number $c$, which proves smoothness.
\end{proof}

\begin{corollary}\label{rr-extension}
Every smooth function $\phi(c,m)$ on
$\mathbb C^\times\times U(W,J)$ extends uniquely to a
smooth function on this open cell with
\[
 F(n_Lw(c)m n_R)
                  =\chi_z(n_L)\chi_{z'}(n_R)\phi(c,m).
\]
It has the two prescribed unipotent character laws.
\end{corollary}

\begin{proof}
Use the coordinates of Lemma~\ref{rr-cell}.
Multiplication on either side by $N$ stays in the
cell, and its factor is multiplied by the corresponding
unipotent element.  The character law gives the assertion.
\end{proof}

\subsection{The full Capelli determinant in these coordinates}

Let $p=c\partial_c$, $\bar p=\bar c\partial_{\bar c}$
and $\gamma(c)=|c|^{d+1}$.  Thus
$p\gamma=\bar p\gamma=(d+1)\gamma/2$.
Define
\[
 T(Y)=\begin{pmatrix}Y+p&-c\\c^{-1}&0\end{pmatrix},
 \qquad
 L(Y)=\begin{pmatrix}Y-\bar p&\bar c^{-1}\\-\bar c&0\end{pmatrix}.
\]
Let $t(Y)$ be the actual terminal matrix from
Section~\ref{rq-section}.  Its entries act only on $m$
and commute with operators in $c$.

\begin{lemma}\label{rr-radial-capelli}
Let $\mathcal D(Y)$ be the full Capelli determinant on
$U(V,H)$ with parameter $Y+(d+1)/2$.
For every smooth function $F$ with the two unipotent
character laws, put
$\Phi(c,m)=\gamma(c)F(w(c)m)$.
Then
\begin{equation}\label{rr-sandwich}
 \gamma(c)(\mathcal D(Y)F)(w(c)m)
                         =(L(Y)t(Y)T(Y))_{11}\Phi(c,m).
\end{equation}
Consequently the operators
$\mathcal C(Y)=(L(Y)t(Y)T(Y))_{11}$ commute
for all complex parameters.
\end{lemma}

\begin{proof}
We give the bordered determinant calculation explicitly.
Order the full basis as $(e,1,\ldots,d,f)$ and put
$U=Y+(d+1)/2$.  For an omission in the middle block,
write $Q_{\widehat i,\widehat j}$ for its quantum minor.
Expansion first in column $e$, then in column $f$,
and then in the two exposed end rows gives the following
four terms:
\begin{align*}
 &T_{ee}(U)D(U-1)T_{ff}(U-d-1),\\
 &-\sum_{i,j}(-1)^{i+j}
      E_{fj}T_{ee}(U)Q_{\widehat i,\widehat j}(U-1)E_{if},\\
 &-\sum_{i,j}(-1)^{i+j}
      E_{ie}Q_{\widehat i,\widehat j}(U-2)
                        T_{ff}(U-d-1)E_{ej},\\
 &-\sum_{\substack{i\ne i'\\j\ne j'}}
   \epsilon(i,i')\epsilon(j,j')(-1)^{i+i'+j+j'}
   E_{ie}E_{fj}Q_{\widehat{i,i'},\widehat{j,j'}}(U-2)
                                      E_{ej'}E_{i'f}.
\end{align*}
The remaining terms are
\begin{align*}
 &(-1)^{d+1}E_{fe}Q_{(e,1,\ldots,d),(1,\ldots,d,f)}(U-1),\\
 &\sum_i(-1)^{i+d}E_{ie}
       Q_{(1,\ldots,\widehat i,\ldots,d,f),(1,\ldots,d)}(U-1)E_{ef}.
\end{align*}
These formulas are obtained using the end-row and
end-column expansions of Lemma~\ref{rq-cofactor-expansions}.
For the fourth line, the position of $i'$ after removing
$i$ is $i'-\mathbf1_{i'>i}$, and similarly for $j'$.
These position changes give the two factors $\epsilon$.
When an end-row factor is moved past an end-column
factor, their possible extra commutator is a multiple
of $E_{ef}$ or $E_{fe}$ and is exactly one of the last
two displayed terms.  Thus this expansion retains the
noncommutative correction terms.

At the point $w(c)m$, the left and right character
rules of \eqref{rc-mixed-right}--\eqref{rc-mixed-left}
are
\[
 \begin{gathered}
 E_{ej}F=x_jF,\qquad E_{if}F=\xi_iF,\\
 (E_{ie}XF)(w(c)m)=c^{-1}\eta_i(XF)(w(c)m),\\
 (E_{fj}XF)(w(c)m)=\bar c^{-1}y_j(XF)(w(c)m).
 \end{gathered}
\]
They hold for every right differential word $X$ in
the left rules.  The right central root $E_{ef}$
annihilates $F$, and the left central root $E_{fe}$
annihilates the evaluation of every $XF$.
The last two terms therefore vanish.

The diagonal derivatives restrict to $p$ and $-\bar p$.
Conjugation by $\gamma$ changes the two diagonal
factors to $Y+p$ and $Y-\bar p$, respectively.
The middle parameter $U-1$ is $Y+k/2$.
Grouping the four sums by their omitted pairs,
as in \eqref{rc-boundary-sign}, leaves precisely
\[
 (Y-\bar p)t_{11}(Y)(Y+p)
 +(Y-\bar p)t_{12}(Y)c^{-1}
 +\bar c^{-1}t_{21}(Y)(Y+p)
 +\bar c^{-1}t_{22}(Y)c^{-1}.
\]
This is the $(1,1)$ entry of $L(Y)t(Y)T(Y)$.

The full Capelli determinants commute, by their
centrality.  They preserve both unipotent character
laws: left translations commute with right
derivatives, and invariance under conjugation gives
commutation with right translations.  Corollary
\ref{rr-extension} allows an arbitrary smooth radial
function as their restriction.  Applying
\eqref{rr-sandwich} twice therefore proves
$[\mathcal C(U),\mathcal C(V)]=0$ on every such function.
\end{proof}

\subsection{An elementary matrix calculation}

For a $2\times2$ matrix $A$ over an algebra, let
$A_1=A\otimes I$ and $A_2=I\otimes A$, with their
entries ordered as written.  Let $P$ interchange the
two factors of $\mathbb C^2\otimes\mathbb C^2$ and put
$R(s)=sI-P$.
The identity
\[
 R(u-v)A_1(u)A_2(v)=A_2(v)A_1(u)R(u-v)
\]
is equivalent, entry by entry, to
\begin{equation}\label{rr-RTT}
 (u-v)[A_{ab}(u),A_{cd}(v)]
       =A_{cb}(u)A_{ad}(v)-A_{cb}(v)A_{ad}(u).
\end{equation}

\begin{lemma}\label{rr-Weyl}
If $cc^{-1}=c^{-1}c=1$ and $[p,c]=c$, then
\[
                  \begin{pmatrix}u+p&-c\\c^{-1}&0\end{pmatrix}
\]
satisfies \eqref{rr-RTT}.
Both $T$ and $L$ above therefore satisfy this identity.
\end{lemma}

\begin{proof}
One also has $[p,c^{-1}]=-c^{-1}$.
Commutators of two diagonal entries, or of $c$ with
$c^{-1}$, vanish.  The only nonzero pairs, up to
interchanging them, are
$[u+p,-c]=-c$ and $[u+p,c^{-1}]=-c^{-1}$.
Substitution into \eqref{rr-RTT} gives respectively
$-(u-v)c$ and $-(u-v)c^{-1}$ on both sides.
Pairs involving the zero entry give zero on both
sides.  This exhausts the four entries and proves
the identity.
For $L$, take the Weyl pair
$p'=-\bar p$, $c'=-\bar c^{-1}$,
$(c')^{-1}=-\bar c$.
\end{proof}

\begin{lemma}\label{rr-defect}
Let $Q(u)$ be any polynomial $2\times2$ operator
matrix acting only on $m$, and put $A(u)=L(u)Q(u)T(u)$.
Set
\[
 X(u,v)=R(u-v)Q_1(u)Q_2(v)
                         -Q_2(v)Q_1(u)R(u-v).
\]
Then
\begin{equation}\label{rr-projected-defect}
 (u-v-1)[A_{11}(u),A_{11}(v)]
  =\bigl(L_2(v)L_1(u)X(u,v)T_1(u)T_2(v)\bigr)_{11,11}.
\end{equation}
\end{lemma}

\begin{proof}
The two Weyl pairs commute with one another and with
all entries of $Q$.  Move $R$ through the two $L$
factors by Lemma~\ref{rr-Weyl}, and then through the
two $T$ factors by the same identity.  Ordinary
associative matrix multiplication gives
\[
 R A_1(u)A_2(v)-A_2(v)A_1(u)R
                       =L_2(v)L_1(u)X(u,v)T_1(u)T_2(v).
\]
The row and column indexed by $(1,1)$ are eigenvectors
of $P$ with eigenvalue one.  Taking that entry gives
\eqref{rr-projected-defect}.
\end{proof}

\subsection{Recovering every coefficient}

\begin{lemma}\label{rr-independence}
Suppose the left side of
\eqref{rr-projected-defect} vanishes on all smooth
functions of $(c,m)$.  If $u-v\ne1$, then
$X(u,v)$ is the zero operator matrix on smooth
functions of $m$.
\end{lemma}

\begin{proof}
Use indices $0,1$ in this calculation only.
Fix a smooth function $f(m)$, a point $m$, a complex
number $\beta$ and an integer $n$, and apply
\eqref{rr-projected-defect} to
\[
                     |c|^{2\beta}c^n f(m).
\]
Here $|c|^{2\beta}=\exp(2\beta\log|c|)$ is a smooth
single-valued function on $\mathbb C^\times$.
It is an eigenfunction of $p$ with eigenvalue
$\beta+n$ and of $\bar p$ with eigenvalue $\beta$.
Put $z=\beta+n$.

For $A=(a_1,a_2)$ and $B=(b_1,b_2)$ the first-row
and first-column coefficients in the right side of
\eqref{rr-projected-defect}, after removing the
common nonzero factor $|c|^{2\beta}c^n$, are
\[
      L_A(\beta)R_B(z)c^{-(b_1+b_2)}\bar c^{-(a_1+a_2)},
\]
where
\[
\begin{array}{c|cccc}
A&00&01&10&11\\ \hline
L_A(\beta)&(u-\beta)(v-\beta)&u-\beta&v-\beta+1&1
\end{array}
\]
and
\[
\begin{array}{c|cccc}
B&00&01&10&11\\ \hline
R_B(z)&(u+z)(v+z)&u+z-1&v+z&1.
\end{array}
\]
The two shifts by one follow from
$p(c^{-1}h)=c^{-1}(p-1)h$ and
$\bar p(\bar c^{-1}h)=\bar c^{-1}(\bar p-1)h$.
No entry of $X$ differentiates these scalar factors.

The nine monomials
$c^{-b}\bar c^{-a}$, $0\leq a,b\leq2$,
are linearly independent on every nonempty open subset
of $\mathbb C^\times$.  Indeed multiplying by
$c^2\bar c^2$ gives a polynomial in $c,\bar c$;
the invertible change $c=x+iy$, $\bar c=x-iy$
reduces the assertion to polynomial uniqueness on
an open subset of $\mathbb R^2$.
Thus for each pair $(a,b)$,
\[
 \sum_{\substack{a_1+a_2=a\\b_1+b_2=b}}
 L_A(\beta)R_B(\beta+n)(X_{A,B}(u,v)f)(m)=0.
\]
This holds for every complex $\beta$ and every
integer $n$.  Polynomial uniqueness first in
$\beta$ and then in $n$ makes it a polynomial
identity in two independent variables $\beta,z$.
The argument remains valid for vector-valued
coefficients by finite Vandermonde inversion.

When $a=0$ or $2$, there is just one nonzero row
polynomial.  When $a=1$, the two row polynomials
are $u-\beta$ and $v-\beta+1$.
They are independent because their difference is
the nonzero constant $u-v-1$.
The same assertion holds for the two column
polynomials $u+z-1$ and $v+z$ when $b=1$.
Products of these independent polynomials in the
two separate variables are independent.
Every coefficient $(X_{A,B}(u,v)f)(m)$ therefore
vanishes.  Since $f,m$ were arbitrary, $X(u,v)=0$.
\end{proof}

\begin{theorem}\label{rr-terminal-RTT}
The actual terminal differential matrix satisfies,
on every smooth function on $U(W,J)$,
\[
 (u-v)[t_{ab}(u),t_{cd}(v)]
      =t_{cb}(u)t_{ad}(v)-t_{cb}(v)t_{ad}(u)
                \qquad(a,b,c,d\in\{1,2\}).
\]
The statement holds for every pair of complex
parameters, including $u-v=1$.
\end{theorem}

\begin{proof}
Lemma~\ref{rr-radial-capelli} gives the commuting
$(1,1)$ entries required by Lemma~\ref{rr-independence}.
Hence the defect matrix vanishes for $u-v\ne1$.
All its entries are polynomials in $u,v$ with
differential-operator coefficients, since the
terminal entries are finite sums of quantum minors.
Applying an entry to any fixed smooth function and
evaluating at a point gives a polynomial which
vanishes for $u-v\ne1$; it therefore vanishes
identically.  This includes the exceptional parameters
and proves the theorem.
\end{proof}

\begin{corollary}\label{rr-terminal-adjugate}
Put $J_2=\begin{pmatrix}0&1\\-1&0\end{pmatrix}$ and
\[
 q(Y)=t_{11}(Y+1)t_{22}(Y)-t_{21}(Y+1)t_{12}(Y).
\]
As operator matrices on smooth terminal functions,
\[
 t(Y+1)^{\mathsf T}J_2t(Y)=q(Y)J_2,\qquad
 t(Y)J_2t(Y+1)^{\mathsf T}=q(Y)J_2.
\]
On the compact types of Corollary~\ref{rq-type-scalar},
$q(Y)$ is the scalar polynomial computed there.
At a value where that scalar is nonzero,
\[
             t(Y)^{-1}=-q(Y)^{-1}
                              J_2t(Y+1)^{\mathsf T}J_2
\]
is both a left and a right inverse.
\end{corollary}

\begin{proof}
Write $S=t(Y+1)$ and $T=t(Y)$.
Use Theorem~\ref{rr-terminal-RTT} with parameter
difference one.  It gives
\[
 S_{11}T_{21}=S_{21}T_{11},\qquad
 S_{12}T_{22}=S_{22}T_{12},
\]
and, on adding its $(12,21)$ and $(11,22)$ entries,
\[
 S_{12}T_{21}-S_{22}T_{11}
                       =-S_{11}T_{22}+S_{21}T_{12}.
\]
These are the two zero diagonal entries and the
negative lower off-diagonal entry of the first
matrix identity; its upper off-diagonal entry
is the definition of $q(Y)$.

The same commutation relation, in its $(11,12)$,
$(21,12)$, $(11,22)$, and $(21,22)$ entries, gives
\begin{align*}
 T_{11}S_{12}-T_{12}S_{11}&=0,\\
 T_{11}S_{22}-T_{12}S_{21}&=q(Y),\\
 T_{21}S_{12}-T_{22}S_{11}&=-q(Y),\\
 T_{21}S_{22}-T_{22}S_{21}&=0.
\end{align*}
This proves the second matrix identity.
Apply Corollary~\ref{rq-type-scalar} for the scalar
value.  Finally $J_2^2=-I$, so the two identities
give the displayed two-sided inverse whenever
that scalar is nonzero.
\end{proof}

\paragraph{Scope.}
The theorem proves the full commutation relation of
the terminal matrix itself, without assuming a
finite compact type or a special Jacquet origin
for the smooth function.  The recursive frame
identities and the estimates justifying integration
with the Jacquet pairing remain separate.


\section{Construction of the recursive radial frame}
\label{re-section}

The terminal commutation and determinant identities suffice to
construct the change of variables in the positivity argument.
We first construct its initial connection directly from a
polynomial coefficient.  We then prove the two recursive changes
of variables by \(2\times2\) matrix multiplication.  In particular,
the larger normal form involving individual raising operators is
not a premise of this construction.

\subsection{The initial connection}

Let \(k\geq1\), let \(W\) have signature \((k,1)\), and retain
the negative vectors \(z,z'\), with
\[
 h(z,z)=-T,\qquad h(z',z')=-T',\qquad
 \zeta(m)=h(mz',z),\qquad G=\zeta\bar\zeta-TT'.
\]
We work on the open set \(G>0\).  In particular, \(G,\zeta,\bar\zeta\)
are invertible smooth multiplication operators.  The space of
smooth functions has the fixed left and right compact equivariance
of Section~\ref{rd-section}; it can also depend on independent
variables \(c_1,\ldots,c_r\in\mathbb C^\times\).
The formal adjoint is taken for terminal Haar measure, the
Hilbert--Schmidt inner product, and
\(\prod_i |c_i|^{-2}\,dc_i\).

Write
\[
 t(Y)=
 \begin{pmatrix} a(Y)&b(Y)\\c(Y)&\tau(Y)\end{pmatrix}
\]
for the actual terminal matrix of
\eqref{rc-t21}--\eqref{rc-t22}, including its full Capelli
entry \(a(Y)\).  The letter \(c(Y)\) in this display denotes
an operator, and is unrelated to the independent coordinates
\(c_i\).  Sections~\ref{rq-section} and \ref{rr-section} give
\begin{align}
 (u-v)[t_{ij}(u),t_{kl}(v)]
 &=t_{kj}(u)t_{il}(v)-t_{kj}(v)t_{il}(u),       \label{re-rtt}\\
 a(Y+1)\tau(Y)-c(Y+1)b(Y)&=q(Y),                \label{re-one-det}\\
 q(Y)&=-TT'\prod_{i=1}^k
       (Y+\tfrac12+u_i)(Y+\tfrac12+u'_i).       \label{re-q}
\end{align}
The real numbers \(u_i,u'_i\) are the shifted compact highest
weights used in the terminal determinant calculation.
All these equalities are equalities of differential operators,
not merely of their principal symbols.

\begin{lemma}\label{re-base-leading}
Put \(h=k-1\).  There are operators \(B,C\) such that
\begin{align}
 \tau(Y)&=Gx_0(Y),&
 x_0(Y)&=Y^h+B Y^{h-1}+O(Y^{h-2}),             \label{re-x-leading}\\
 c(Y)&=-\bar\zeta Y^{h+1}+C Y^h+O(Y^{h-1}).    \label{re-c-leading}
\end{align}
Here \(x_0(Y)=\sum_\rho p_\rho(Y)P_\rho\), with
\(p_\rho(Y)=\prod_{i=1}^{k-1}(Y+\rho_i)\), is the actual
finite spectral decomposition of
Theorem~\ref{rd-constructed-kappa}.
When \(h=0\), the first formula means \(x_0=1\), we set
\(B=0\), and the second formula has no remainder.
The operator \(C\) has differential order at most one.
For real \(Y\), the formal adjoint of \(t_{ij}(Y)\) is
\(t_{ji}(Y)\).
\end{lemma}

\begin{proof}
The spectral formula, including the fact that its coefficients
have order zero, was proved in
Theorem~\ref{rd-constructed-kappa}.
There is also a direct check on its leading normalization.
Put \(y=m^{-1}z\), \(\eta_i=h(mw_i,z)\),
\(x=z'\), and \(\xi_i=h(w_i,z')\).
The leading coefficient of the double cofactor contraction is
\[
 -\sum_{i<j}(\eta_i\xi_j-\eta_j\xi_i)(y_i x_j-y_jx_i)
 =-\{(\eta y)(\xi x)-(\eta x)(\xi y)\}=G.
\]
The leading coefficient of the single cofactor contraction is
\(-\xi y=-\bar\zeta\).  A cofactor of size \(k\) has degree
\(k\) in its parameter; its coefficient of degree \(k-1\)
contains at most one Lie algebra generator.  The shift
\(Y\mapsto Y+k/2\) changes this coefficient only by an order
zero term.  This proves \eqref{re-c-leading} and the order
assertion.

For completeness, the adjoint assertion follows directly from
the same minors.  In an orthonormal basis put
\(\epsilon_i=h(w_i,w_i)\in\{1,-1\}\).
Complexified right differentiation satisfies
\(E_{ij}^*=\epsilon_i\epsilon_jE_{ji}\).
The row and column forms of a quantum minor give
\[
 Q_{I,J}(u)^*=\epsilon_I\epsilon_J Q_{J,I}(\bar u).
\]
Indeed, taking adjoints reverses the order of all factors;
use the row form with both index lists reversed.  Reversing
both lists contributes the square of the same permutation sign,
which is one.  For the signed cofactors this reads
\(C_{ij}(u)^*=\epsilon_i\epsilon_jC_{ji}(\bar u)\).
Consequently, for real \(u\),
\[
 \left(-\sum_{i,j}\xi_i y_j C_{ij}(u)\right)^*
 =-\sum_{i,j}C_{ij}(u)\eta_i x_j
 =-\sum_{i,j}\eta_i x_j C_{ij}(u-1).
\]
The last equality is precisely the covector shift
\eqref{rq-shift-one}.  Thus \(c(Y)^*=b(Y)\).
The full determinant is self-adjoint by the same minor rule.
Finally \(\tau(Y)=Gx_0(Y)\) is self-adjoint on the radial
space: \(G\) is real and each \(P_\rho\) is a pointwise
orthogonal projection.  This proves every matrix entry of the
assertion.
\end{proof}

\begin{lemma}\label{re-two-dets}
The relations \eqref{re-rtt} and \eqref{re-one-det} imply
\begin{equation}
 t(Y+1)^t Jt(Y)=q(Y)J,\qquad
 t(Y)Jt(Y+1)^t=q(Y)J,\qquad
 J=\begin{pmatrix}0&1\\-1&0\end{pmatrix}.       \label{re-dets}
\end{equation}
They also imply that the members of each fixed entry family
\(t_{ij}(Y)\) commute with one another.
\end{lemma}

\begin{proof}
For a fixed entry, \eqref{re-rtt} gives
\((u-v-1)[t_{ij}(u),t_{ij}(v)]=0\).
Interchanging \(u,v\) gives the same assertion with
\(u-v+1\); subtraction proves commutation, without a restriction
on \(u-v\).
At \(u=Y+1,v=Y\), the remaining entries of
\eqref{re-rtt} give
\[
 a_+c=c_+a,\quad b_+\tau=\tau_+b,\quad
 ab_+=ba_+,\quad c\tau_+=\tau c_+,
\]
where a subscript \(+\) means evaluation at \(Y+1\).
They further give
\[
 a_+\tau-c_+b=a\tau_+-bc_+
                         =\tau a_+-cb_+
                         =\tau_+a-b_+c.
\]
For example, the first equality is the relation for \(c_+,b\);
the second follows from that for \(a_+,\tau\).
The fourth expression follows in the same way from that for
\(b_+,c\).  Multiplying the two displayed \(2\times2\)
matrices now gives \eqref{re-dets} entry by entry.
\end{proof}

\begin{theorem}\label{re-initial-connection}
Define
\begin{equation}
 \nabla=-G^{-1}C-(B+\tfrac12)\bar\zeta G^{-1},
 \qquad
 \nabla^*=\text{its formal adjoint}.
                                                               \label{re-connection}
\end{equation}
Then \(\nabla\) is a first-order operator and
\begin{gather}
 [\nabla,\zeta]=1,\quad[\nabla,\bar\zeta]=0,\quad
 [\nabla,G]=\bar\zeta,\qquad
 [x_0(Y),\nabla]=0,                               \label{re-connection-laws}\\
 [\nabla^*,\bar\zeta]=-1,\quad[\nabla^*,\zeta]=0,\quad
 [\nabla^*,G]=-\zeta.                            \label{re-star-laws}
\end{gather}
The connection and its adjoint commute with the independent
coordinates \(c_i,\bar c_i\), their inverses, and their Euler
derivatives.

Let \(\sigma\in\{1,-1\}\), replace \(t\) by \(\sigma t\),
and put \(x(Y)=\sigma x_0(Y)\).  Define
\begin{equation}
 I=\{Y\in\mathbb C:p_\rho(Y)\ne0\text{ for every }\rho\},\quad
 x(Y)^{-1}=\sigma\sum_\rho p_\rho(Y)^{-1}P_\rho\quad(Y\in I).
                                                               \label{re-initial-domain}
\end{equation}
On \(I\) put
\begin{equation}
 \ell(Y)=(Gx(Y))^{-1}(\sigma c(Y)),\qquad
 \ell_0(Y)=-\nabla-(Y+\tfrac12)\bar\zeta G^{-1},\qquad
 e(Y)=\zeta G\bar\zeta^{-1}\{\ell(Y)-\ell_0(Y)\}.
                                                               \label{re-initial-e}
\end{equation}
Each \(e(Y)\) commutes with \(\zeta,\bar\zeta,G\) and their
inverses.  These operators supply the actual initial data of
the recursive construction, with
\begin{equation}
 \begin{aligned}
 &R=F=1,\quad N=0,\quad d=c_{\rm wt}=1,\quad t_{22}(Y)=Gx(Y),\\
 &t_{21}(Y)=Gx(Y)\{\ell_0(Y)+\bar\zeta\zeta^{-1}G^{-1}e(Y)\}.
 \end{aligned}
                                                               \label{re-initial-frame}
\end{equation}
The determinant \(q\) is unchanged by the sign \(\sigma\).
\end{theorem}

\begin{proof}
Compare coefficients of the polynomial parameter \(v\) in
\eqref{re-rtt}.  Commutation of the family \(c(v)\) gives
\([c(u),\bar\zeta]=0\).
The coefficient of \(v^{h+2}\) in the relation for
\(\tau(u),c(v)\) gives \([\tau(u),\bar\zeta]=0\).
Commutation of the family \(\tau(v)\) gives
\([\tau(u),G]=0\).  Since \(G=\zeta\bar\zeta-TT'\) and
\(\bar\zeta\) is invertible, also \([\tau(u),\zeta]=0\).
Finally the coefficient of \(v^{h+1}\) in
\[
 (u-v)[c(u),\tau(v)]
                    =c(u)\tau(v)-c(v)\tau(u)
\]
is
\([c(u),G]=-\bar\zeta\tau(u)\).
It follows that
\begin{equation}
 [c(u),\zeta]=-\tau(u),\qquad [c(u),\bar\zeta]=0.
                                                               \label{re-zeta-comm}
\end{equation}
Comparing the coefficient of \(u^h\) proves
\([C,\zeta]=-G\), \([C,\bar\zeta]=0\).  The coefficient \(B\)
commutes with scalar multiplication by the spectral
decomposition.  Thus \eqref{re-connection} gives the first
three identities in \eqref{re-connection-laws}.

It remains to check the commutation with \(x_0(u)\).
Work in the algebra of formal Laurent series in \(v^{-1}\)
with differential-operator coefficients.  The leading term
of \(\tau(v)\) is \(Gv^h\), with invertible coefficient, so
its formal inverse exists.  Equations
\eqref{re-x-leading}--\eqref{re-connection} give
\begin{equation}
 \ell(v)=\tau(v)^{-1}c(v)
      =-v\bar\zeta/G-\nabla-\tfrac12\bar\zeta/G+O(v^{-1}).
                                                               \label{re-laurent}
\end{equation}
Dividing the RTT relation for \(\tau(u),c(v)\) by the
commuting \(\tau(v)\) gives
\begin{equation}
 [\tau(u),\ell(v)]
       =-\tau(u)\frac{\ell(u)-\ell(v)}{u-v}.       \label{re-tau-ell}
\end{equation}
This identity may initially be read with \(\tau(u)\) also
formally invertible; multiplication by it cancels its inverse.
The coefficient of \(v^0\) on its right is
\(\bar\zeta x_0(u)\).  On its left the coefficient is
\[
 [Gx_0(u),-\nabla]
        =-G[x_0(u),\nabla]+\bar\zeta x_0(u).
 \]
Hence \([x_0(u),\nabla]=0\).  This is a polynomial
identity in \(u\); it therefore holds for every complex \(u\).
The interpolation formula for \(P_\rho\) then also gives
\([P_\rho,\nabla]=0\) for every \(\rho\).
Taking formal adjoints gives \eqref{re-star-laws}.
All coefficients used here act only in the terminal variable,
which proves the assertions about the independent \(c_i\).

The finite spectral formula proves the stated two-sided
inverse on \(I\), commutation of the \(x(Y)\), their
commutation with all scalar functions, and their
self-adjointness for real \(Y\).
Equation \eqref{re-zeta-comm} shows
\([\ell(Y),\zeta]=-1\) and
\([\ell(Y),\bar\zeta]=0\).
The same is true of \(\ell_0(Y)\); their difference therefore
commutes with \(\zeta,\bar\zeta,G\), proving the assertion about
\(e\).  The last displayed identities are now definitions.
The sign disappears from \(\ell\), and its square in the
determinant is one.  For \(k=1\), \eqref{re-laurent} is an
exact equality, so \(e=0\).  No limiting argument at a spectral
root has been used.
\end{proof}

The construction stays within the actual radial function space.
The cofactor covariance in Section~\ref{rc-section} shows that
the terminal entries preserve the two compact equivariances.
Each of their polynomial coefficients is a finite linear
combination of evaluations, so \(C\) preserves them as well.
The spectral projections preserve them by
Theorem~\ref{rd-constructed-kappa}; multiplication by
\(\zeta,\bar\zeta,G\) and their inverses does so because these
functions are invariant.  All these coefficients are smooth on
\(G>0\).  Formal adjoints preserve the equivariance because
the compact actions are unitary.  Consequently every operator
just constructed acts on the stated smooth radial space and
preserves supports.

\subsection{Two elementary Gauss identities}

The following calculations take place in an arbitrary
associative complex algebra.  Every inverse mentioned is a
two-sided inverse.  Suppose \(t(u)\) satisfies
\eqref{re-rtt} and \eqref{re-dets}, and that \(\tau(u)\) is
invertible for \(u\) in a set \(I\).  Write
\[
 \ell_u=\tau_u^{-1}c_u,\qquad
 m_u=\tau_u^{-1}b_u,\qquad
 H_u=q(u)\tau_u^{-1}\tau_{u+1}^{-1}.
\]
The last expression is used only when \(u,u+1\in I\).

\begin{lemma}\label{re-gauss}
Whenever the indicated adjacent points belong to \(I\),
\begin{align}
 m_{u+1}&=b_u\tau_u^{-1},&
 \ell_u&=c_{u+1}\tau_{u+1}^{-1},               \label{re-shifts}\\
 a_u-b_u\tau_u^{-1}c_u&=q(u)\tau_{u+1}^{-1},&
 a_u-c_u\tau_u^{-1}b_u&=q(u-1)\tau_{u-1}^{-1},\label{re-schur}\\
 t(u)&=
 \begin{pmatrix}1&m_{u+1}\\0&1\end{pmatrix}
 \begin{pmatrix}q(u)\tau_{u+1}^{-1}&0\\0&\tau_u\end{pmatrix}
 \begin{pmatrix}1&0\\\ell_u&1\end{pmatrix},     \label{re-gauss-plus}\\
 t(u)^t&=
 \begin{pmatrix}1&\ell_{u-1}\\0&1\end{pmatrix}
 \begin{pmatrix}q(u-1)\tau_{u-1}^{-1}&0\\0&\tau_u\end{pmatrix}
 \begin{pmatrix}1&0\\m_u&1\end{pmatrix}.        \label{re-gauss-minus}
\end{align}
If \(u\ne v\), put
\(\eta_{uv}=(\ell_u-\ell_v)/(u-v)\) and
\(\mu_{uv}=(m_u-m_v)/(u-v)\).  Then
\begin{align}
 [\tau_u,\ell_v]&=-\tau_u\eta_{uv},&
 [\tau_u,m_v]&=\tau_u\mu_{uv},                 \label{re-gauss-comm}\\
 [\ell_u,\ell_v]&=\eta_{uv}(\ell_u-\ell_v),    \label{re-ell-comm}\\
 (u-v-1)[m_u,\ell_v]&=H_v-H_{u-1}.            \label{re-mixed}
\end{align}
For \eqref{re-mixed} assume \(u-1,u,v,v+1\in I\).
If \(d=d^*\) is invertible and \(dt(u)\) is Hermitian for real
\(u\), then
\begin{equation}
             dm_{u+1}=\ell_u^*d                 \label{re-adj-gauss}
\end{equation}
for real \(u,u+1\in I\).
\end{lemma}

\begin{proof}
The zero diagonal entries of the two identities
\eqref{re-dets} give
\(b_{u+1}\tau_u=\tau_{u+1}b_u\) and
\(c_u\tau_{u+1}=\tau_uc_{u+1}\).
These are \eqref{re-shifts}.  The off-diagonal entries,
with these substitutions, give \eqref{re-schur}.
Multiplying the three factors proves the two factorizations.

The RTT relation for \(\tau_u,c_v\), after multiplication
by \(\tau_v^{-1}\), gives the first identity in
\eqref{re-gauss-comm}.  For clarity, the second can be
derived without dividing by \(u-v+1\).
Put \(D=u-v\) and \(K=[\tau_u,m_v]\).
The relations for \(\tau_u,b_v\), together with the same
relation with \(u,v\) interchanged, give
\[
 DK=\tau_v^{-1}b_u\tau_v-m_v\tau_u,\qquad
 \tau_v^{-1}b_u\tau_v=b_u-K.
\]
Thus \((D+1)K=b_u-m_v\tau_u
=b_u-\tau_um_v+K\), proving \(DK=\tau_u(m_u-m_v)\).
Next expand \(c_uc_v=c_vc_u\), use
\(c_w=\tau_w\ell_w\), and move \(\tau_u,\tau_v\)
to the left using \eqref{re-gauss-comm}.
Cancellation of their invertible product gives
\eqref{re-ell-comm}.

The relation for \(b_u,c_v\) gives, by the same movement,
\[
 D[m_u,\ell_v]-(m_u-m_v)\ell_v-(\ell_u-\ell_v)m_u
       =H_v-H_{u-1}+m_v\ell_v-\ell_um_u.
\]
Here use the first Schur expression at \(v\) and the second
at \(u\), so that no inverse at \(u+1\) is required.
The left side is
\((D-1)[m_u,\ell_v]+m_v\ell_v-\ell_um_u\).
Cancellation proves \eqref{re-mixed}.
Finally weighted Hermiticity says \(c_u^*d=db_u\) and
\(\tau_u^*d=d\tau_u\).  Hence
\[
 \ell_u^*d=c_u^*(\tau_u^*)^{-1}d
                    =db_u\tau_u^{-1}=dm_{u+1}.
\]
\end{proof}

\subsection{The recursive construction and its domains}

We now allow an auxiliary tensor product of finitely many copies
of \(\mathbb C^2\).  All matrices below have differential-operator
coefficients on the radial space and act on this tensor product.
A new factor always means a factor on which every existing
coefficient acts as the identity.  This condition is essential.
Let
\[
 w=\zeta\bar\zeta/G,\qquad
 \kappa=\bar\zeta/(\zeta G),\qquad t_z=TT'/\bar\zeta.
\]
A frame consists of sets \(I\subset\Omega\), a matrix \(t(Y)\)
defined on \(\Omega\), an invertible \(F\), commuting operators
\(x(Y)\), an invertible self-adjoint \(c_{\rm wt}\), and
operators \(e(Y)\) for \(Y\in I\).  Its scalar determinant
\(q(Y)\) is defined on the adjacent-point domain
\(\Omega^{(2)}=\{Y:Y,Y+1\in\Omega\}\).  Put
\begin{equation}
 d=w^N c_{\rm wt},\quad
 \ell_Y=-\nabla-(Y+\tfrac12)\bar\zeta/G+\kappa e(Y),\quad
 m_Y=d^{-1}\ell_{Y-1}^*d
       \quad(Y\in\mathbb R,\ Y-1\in I).
                                                               \label{re-frame-data}
\end{equation}
Here \(\ell_Y\) is defined for \(Y\in I\).  For real \(Y\)
with \(Y-1,Y\in I\), the displayed \(m_Y\) agrees with
the upper Gauss coordinate by the adjoint identity
\eqref{re-adj-gauss}.
The frame identities are
\begin{align}
 R&=F^*dF,&
 Ft_{22}(Y)F^{-1}&=Gx(Y)\quad(Y\in\Omega),       \label{re-frame-main}\\
 Ft_{21}(Y)F^{-1}&=Gx(Y)\ell_Y
                    \quad(Y,Y+1\in I).        \label{re-frame-lower}
\end{align}
We require \(x(Y)^{-1}\) on \(I\), that the \(x(Y)\)
commute with one another, \(c_{\rm wt}\), \(\nabla\), and
\(\zeta,\bar\zeta,G\), that \(x(Y)^*=x(Y)\) for real \(Y\),
and that \(c_{\rm wt}\) and \(e(Y)\) commute with the latter
three multiplication operators.  In addition, \(R=R^*\),
\(Rt(Y)\) is Hermitian at real \(Y\in\Omega\), and
\eqref{re-rtt} holds on \(\Omega\), with
\eqref{re-dets} whenever \(Y,Y+1\in\Omega\).
These conditions hold for the initial frame constructed above.

For a new factor let
\begin{equation}
 P_{ij}=E_{ji},\qquad Q_{ij}=E_{ij},\qquad
 R_P(D)=\frac{DI-P}{D-1},\qquad
 R_Q(d_0)=\frac{d_0I-Q}{d_0-1}.
                                                               \label{re-fresh-matrices}
\end{equation}
Here \(i,j\in\{1,2\}\) are the matrix indices of \(t\),
and \(E_{ij}\) acts on the new factor.  The operation
\(\operatorname{ins}t(a)\) means that the two indices of
\(t(a)\) act on that new factor.

\begin{theorem}\label{re-recursive-frame}
Suppose the preceding frame conditions hold.

For a real number \(a\), assume
\begin{equation}
                a,a+1\in I,\qquad q(a)\ne0.     \label{re-plus-hyp}
\end{equation}
There is a new frame with
\begin{equation}
 \begin{aligned}
 \Omega_+&=\Omega\setminus\{a+1\},&
 I_+&=I\setminus\{a,a+1\},\\
 t_+(Y)&=t(Y)R_P(Y-a),&
 q_+(Y)&=q(Y)\frac{Y-a+1}{Y-a},\\
 R_+&=R\,\operatorname{ins}t(a),&
 N_+&=N+1.
 \end{aligned}                                                   \label{re-plus-module}
\end{equation}
Its other entries are
\begin{equation}
 F_+=C_+L_aF,\quad
 C_+=\begin{pmatrix}\zeta^{-1}&0\\0&G/\bar\zeta\end{pmatrix},\quad
 L_a=\begin{pmatrix}1&0\\\ell_a&1\end{pmatrix}.                 \label{re-plus-F}
\end{equation}
\begin{align}
 x_+(Y)&=x(Y)
       \begin{pmatrix}(Y-a)/(Y-a-1)&0\\0&1\end{pmatrix},\nonumber\\
 c_{{\rm wt},+}&=c_{\rm wt}
       \begin{pmatrix}q(a)x(a+1)^{-1}&0\\0&x(a)\end{pmatrix},
                                                               \label{re-plus-xc}
\end{align}
With \(\Delta_{Ya}e=(e(Y)-e(a))/(Y-a)\), put
\begin{equation}
 e_+(Y)=
 \begin{pmatrix}
 e(Y)-\Delta_{Ya}e+t_z&-(Y-a)^{-1}\\
 0&x(a)^{-1}e(Y)x(a)
 \end{pmatrix}.                                \label{re-plus-e}
\end{equation}

For a real number \(b\), assume
\begin{equation}
             b-1,b,b+1\in I,\qquad q(b-1)\ne0.  \label{re-minus-hyp}
\end{equation}
There is a new frame with
\begin{equation}
 \begin{aligned}
 \Omega_-&=\Omega\setminus\{b-1\},&
 I_-&=I\setminus\{b-1,b\},\\
 t_-(Y)&=R_Q(b-Y)t(Y),&
 q_-(Y)&=q(Y)\frac{b-Y}{b-Y-1},\\
 R_-&=R\,\operatorname{ins}t(b)^t,&
 N_-&=N+1.
 \end{aligned}                                                   \label{re-minus-module}
\end{equation}
Its other entries are
\begin{equation}
 F_-=C_-M_bF,\quad
 C_-=\begin{pmatrix}\bar\zeta^{-1}&0\\0&G/\zeta\end{pmatrix},\quad
 M_b=\begin{pmatrix}1&0\\m_b&1\end{pmatrix}.                  \label{re-minus-F}
\end{equation}
\begin{align}
 x_-(Y)&=x(Y)
       \begin{pmatrix}(b-Y)/(b-Y-1)&0\\0&1\end{pmatrix},\nonumber\\
 c_{{\rm wt},-}&=c_{\rm wt}
       \begin{pmatrix}q(b-1)x(b-1)^{-1}&0\\0&x(b)\end{pmatrix},
                                                               \label{re-minus-xc}
\end{align}
\begin{equation}
 e_-(Y)=
 \begin{pmatrix}
 e(Y)&0\\
 -\dfrac{q(b-1)x(b)^{-1}x(b-1)^{-1}}{b-Y-1}
       &x(b)^{-1}e(Y)x(b)-\Delta_{Yb}e+t_z
 \end{pmatrix}.                                  \label{re-minus-e}
\end{equation}
All displayed rational formulas are asserted on their indicated
domains.  In particular, \(q_\pm\) is defined on
\(\Omega_\pm^{(2)}=\{Y:Y,Y+1\in\Omega_\pm\}\).
For \(q_+\), this condition excludes \(Y=a\), so its displayed
formula never evaluates its apparent pole.
\end{theorem}

\begin{proof}
Conjugate the old matrix by \(F\), and abbreviate it by \(t\)
during the calculation.  Its lower-right entry is
\(\tau_Y=Gx(Y)\); the old density is \(d\).
Similarity preserves \eqref{re-rtt} and \eqref{re-dets}.
The Gauss coordinates coincide with \(\ell_Y\) when
\(Y,Y+1\in I\), and \eqref{re-adj-gauss} identifies the upper
coordinate with \(m_Y\) at the real parameters used below.

First we check the conjugation of the last row.  For the plus
step put \(D=Y-a\) and
\(\eta=(\ell_Y-\ell_a)/D\).  Before conjugating by \(L_a\),
the new entries are
\begin{equation}
 (t_+)_{22}=\frac{\tau_Y}{D-1}
      \begin{pmatrix}D&0\\-\ell_Y&D-1\end{pmatrix},\qquad
 (t_+)_{21}=\frac{\tau_Y}{D-1}
      \begin{pmatrix}(D-1)\ell_Y&-1\\0&D\ell_Y\end{pmatrix}.
                                                               \label{re-plus-blocks}
\end{equation}
Set the two matrices on the right equal to \(U,V\).
The first identity of \eqref{re-gauss-comm} gives
\begin{equation}
 L_a\tau_Y U
       =\tau_Y\operatorname{diag}(D,D-1)L_a.
                                                               \label{re-plus-tau-calc}
\end{equation}
Moreover
\begin{equation}
 V L_a^{-1}
 =U L_a^{-1}
       \begin{pmatrix}\ell_Y-\eta&-D^{-1}\\0&\ell_Y+\eta\end{pmatrix}.
                                                               \label{re-plus-low-calc}
\end{equation}
Three entries follow by \(D\eta=\ell_Y-\ell_a\).
For the remaining lower-left entry the difference is a scalar
multiple of
\((\ell_Y-\ell_a)^2-D[\ell_Y,\ell_a]\), which is zero
by \eqref{re-ell-comm}.
Thus conjugation by \(L_a\) gives the asserted new
\(\tau\) and the Gauss coordinate in
\eqref{re-plus-low-calc}.

The assertion about \(\tau\) holds on all of \(\Omega_+\),
including points where \(x(Y)\) is not invertible.
To see this without division by \(D\), use the RTT identity
in the form
\begin{equation}
 D\ell_a\tau_Y-c_Y=(D-1)\tau_Y\ell_a .
                                                               \label{re-plus-all-reg}
\end{equation}
It requires only the inverse of \(\tau_a\), and proves
\eqref{re-plus-tau-calc} directly with \(c_Y\) in place of
\(\tau_Y\ell_Y\).

For the minus step put \(d_0=b-Y\).
Before conjugating by \(M_b\), the two new entries, multiplied
by \(d_0-1\), are
\[
 A=\begin{pmatrix}d_0\tau_Y&0\\-b_Y&(d_0-1)\tau_Y\end{pmatrix},
 \qquad
 B_0=\begin{pmatrix}d_0c_Y&0\\-a_Y&(d_0-1)c_Y\end{pmatrix}.
\]
The first Schur formula gives
\begin{equation}
 B_0=A
       \begin{pmatrix}\ell_Y&0\\-H_Y/(d_0-1)&\ell_Y\end{pmatrix}.
                                                               \label{re-minus-block}
\end{equation}
The second identity of \eqref{re-gauss-comm} gives
\[
 M_b A=\operatorname{diag}(d_0\tau_Y,(d_0-1)\tau_Y)M_b.
\]
The mixed identity \eqref{re-mixed}, with \(u=b,v=Y\), then gives
\begin{equation}
 M_b
 \begin{pmatrix}\ell_Y&0\\-H_Y/(d_0-1)&\ell_Y\end{pmatrix}
 =
 \begin{pmatrix}\ell_Y&0\\-H_{b-1}/(d_0-1)&\ell_Y\end{pmatrix}M_b.
                                                               \label{re-minus-low-calc}
\end{equation}
For the lower-right entry on all of \(\Omega_-\), use instead
the undivided RTT equality
\begin{equation}
 d_0m_b\tau_Y-b_Y=(d_0-1)\tau_Ym_b.             \label{re-minus-all-reg}
\end{equation}
It only uses \(\tau_b^{-1}\).

It remains to conjugate by \(C_+\) or \(C_-\).
The level relations give
\begin{equation}
 \eta=-\bar\zeta/G+\kappa\Delta_{Ya}e,\qquad
 \bar\zeta/G-\zeta^{-1}=\kappa t_z .
                                                               \label{re-coordinate-simplify}
\end{equation}
The first relation in \eqref{re-gauss-comm}, applied to
\(\tau_a=Gx(a)\), also gives
\begin{equation}
               x(a)^{-1}e(Y)x(a)=e(Y)+\Delta_{Ya}e.
                                                               \label{re-e-conj}
\end{equation}
Here \(Y,Y+1,a,a+1\in I\) and \(Y\ne a\), exactly as in the
new lower-entry domain.
For the plus step, the two diagonal entries after conjugation
are
\[
 \ell_Y-\eta-\zeta^{-1},\qquad
 \ell_Y+\eta+\bar\zeta/G,
\]
and the upper entry is \(-\kappa/D\).
Equations \eqref{re-coordinate-simplify} and
\eqref{re-e-conj} give exactly \eqref{re-plus-e}.
For the minus step the diagonal entries are
\(\ell_Y\) and \(\ell_Y+\bar\zeta/G-\zeta^{-1}\).
Its lower entry is
\[
 -\frac{G}{\zeta}\frac{H_{b-1}}{d_0-1}\bar\zeta
 =-\kappa\,\frac{q(b-1)x(b)^{-1}x(b-1)^{-1}}{d_0-1}.
\]
Using \eqref{re-e-conj} with \(a=b\) gives
\eqref{re-minus-e}.  The requirement \(b+1\in I\) is used
in this last application.

We next prove the congruence for the form.  The weighted
Gauss factorization gives
\begin{align*}
 d\,t(a)
 &=L_a^*\operatorname{diag}
       \bigl(dq(a)\tau_{a+1}^{-1},d\tau_a\bigr)L_a,\\
 d\,t(b)^t
 &=M_b^*\operatorname{diag}
       \bigl(dq(b-1)\tau_{b-1}^{-1},d\tau_b\bigr)M_b.
\end{align*}
All \(x\)'s commute with \(d\) and with the coordinate
multiplications.  Since
\(w/(\zeta\bar\zeta)=G^{-1}\), direct diagonal multiplication
gives
\begin{align*}
 C_+^*w^{N+1}c_{{\rm wt},+}C_+
 &=\operatorname{diag}(dq(a)\tau_{a+1}^{-1},d\tau_a),\\
 C_-^*w^{N+1}c_{{\rm wt},-}C_-
 &=\operatorname{diag}(dq(b-1)\tau_{b-1}^{-1},d\tau_b).
\end{align*}
Undoing the old conjugation therefore proves
\(R_\pm=F_\pm^*w^{N+1}c_{{\rm wt},\pm}F_\pm\).
The inverse of \(F_+\), for example, is
\(F^{-1}L_a^{-1}C_+^{-1}\); the minus inverse is identical
with \(M_b,C_-\) substituted.

We verify the remaining matrix relations explicitly.
The matrices in \eqref{re-fresh-matrices} satisfy
\eqref{re-rtt}: their numerators have, respectively, the
generator matrices \(-E_{ji}\) and \(E_{ij}\), and
\([E_{ij},E_{kl}]=\delta_{jk}E_{il}-\delta_{il}E_{kj}\)
gives each of the sixteen identities.
Multiplication gives their two shifted determinants as
\((D+1)/D\) and \(d_0/(d_0-1)\), respectively.
Before division the numerators are
\((D-1)(D+1)\) and \(d_0(d_0-2)\); the successive
denominators are \((D-1)D\) and
\((d_0-1)(d_0-2)\).
Products of matrices with mutually commuting coefficient
algebras preserve RTT: in the two auxiliary tensor copies,
move each old coefficient past each fresh coefficient and
apply the two RTT equations in succession.
They also preserve the two determinant equations, since
\((AB)^t=B^tA^t\) for such coefficient algebras.
This proves the claimed \(q_\pm\).

The new weighted Hermiticity follows from the old RTT
relation at \(a,Y\), or \(b,Y\).  More explicitly, after
using \(dt(u)=(dt(u))^*\), the plus assertion is the exchange
identity
\[
 \operatorname{ins}t(a)\,t(Y)R_P(Y-a)
       =R_P(Y-a)t(Y)\operatorname{ins}t(a);
\]
the minus assertion is
\[
 \operatorname{ins}t(b)^t R_Q(b-Y)t(Y)
       =t(Y)R_Q(b-Y)\operatorname{ins}t(b)^t.
\]
The first is RTT on the old auxiliary and new factors.
For the second, expand \(Q_{ij}=E_{ij}\): its
\((i,j;k,l)\) entry is the same RTT identity with the two
indices belonging to the new factor interchanged.
This is an entry calculation; it does not use a rule
reversing arbitrary products under partial transpose.
The self-adjointness of \(R_\pm\) also follows from the
congruences just proved.  The reality of \(q_\pm\) follows
from the reality of their parameters.
Commutation within the new lower-left family follows from
RTT as in Lemma~\ref{re-two-dets}.

Finally the inverses in the new \(x,c_{\rm wt}\) exist on
the specified sets by \eqref{re-plus-hyp} or
\eqref{re-minus-hyp}.  Their commutations and adjoints
follow entry by entry from the old ones: every new
diagonal factor is a real scalar rational function times
one of the commuting \(x\)'s or their inverses.
The formulas for \(e_\pm\) show that they commute with
the three coordinate multiplications.
All frame conditions have now been checked.
\end{proof}

\subsection{Iteration, scalar densities and the empty word}

Start with the actual frame of
Theorem~\ref{re-initial-connection}, with
\(\sigma=(-1)^{k-l}\), \(1\le l\le k\), as in the application.
Take an ordered sequence of plus and minus steps on distinct
auxiliary factors.  At each step require exactly
\eqref{re-plus-hyp} or \eqref{re-minus-hyp} for the current
domain and current determinant.
Theorem~\ref{re-recursive-frame} then applies by induction.
This formulation allows the sharper recursively admissible
domains of the source.

A convenient sufficient condition, used for its original
admissible words, is the following.  All parameters \(a_j\)
are real, \(a_j-1,a_j,a_j+1\) belong to the initial set \(I\),
\(q(a_j),q(a_j-1)\ne0\), and
\begin{equation}
             a_i-a_j\notin\{-2,-1,0,1,2\}\quad(i\ne j).
                                                               \label{re-word-domain}
\end{equation}
Indeed an earlier step removes only its parameter and a
neighbor from \(I\).  Its determinant multiplier has its
zero and pole at such neighbors.  Condition
\eqref{re-word-domain} prevents any later required evaluation
from meeting either one.  Distinct auxiliary factors supply
the required commutation with each new factor.
This proves sufficiency for every order of the steps.

\begin{lemma}\label{re-shifted-domain}
For each step define \(y_j=a_j\) in the plus case and
\(y_j=b_j-1\) in the minus case.
It is enough to require that the auxiliary factors are distinct,
that all \(y_j\) are real, that
\[
 y_j,y_j+1\in I,\qquad q(y_j)\ne0,\qquad
 y_j+2\in I\ \hbox{for a minus step},
\]
in the initial frame, and that
\[
 y_i-y_j\notin\{-2,-1,0,1,2\}\quad(i\ne j).
\]
The resulting invertibility domain after any subset of steps is
the initial \(I\) with the pairs \(\{y_j,y_j+1\}\) removed.
If, in addition,
\begin{equation}
 |y_i-y_j|>2,\qquad
 \sigma p_\rho(y_j)>0,\quad
 \sigma p_\rho(y_j+1)>0,\quad q(y_j)>0
                                                               \label{re-shifted-signs}
\end{equation}
for every \(\rho,j\), all the scalar weights in
\eqref{re-plus-scalar}--\eqref{re-minus-scalar} are positive.
\end{lemma}

\begin{proof}
Both signs remove exactly \(\{y_j,y_j+1\}\) from the
current \(I\).  The required points for a later plus step are
\(y_i,y_i+1\); for a later minus step they are
\(y_i,y_i+1,y_i+2\).  An intersection with a removed pair
would give one of the excluded integer differences.
The determinant multipliers from earlier steps, at a general
point \(Y\), are respectively
\[
 \frac{Y-y_j+1}{Y-y_j},\qquad
 \frac{Y-y_j-1}{Y-y_j}.
\]
At \(Y=y_i\) neither has a zero or a pole.  This proves
recursive admissibility.

For positivity, the two possible nontrivial multipliers of
\(x(Y)\) are
\[
 \frac{Y-y_j}{Y-y_j-1},\qquad
 \frac{Y-y_j-1}{Y-y_j}.
\]
When \(Y=y_i\) or \(y_i+1\), both are positive: the difference
\(Y-y_j\) is either greater than one or less than zero.
The two determinant multipliers at \(Y=y_i\) are positive
as well.  In both signs of the recursion, the required values
of \(x\) are exactly \(x(y_i)\) and \(x(y_i+1)\), and the
required value of \(q\) is \(q(y_i)\).
Thus every new scalar weight is a product or quotient of
positive numbers, proving the assertion inductively.
\end{proof}

This also gives the version with shifted spectral cuts used
in the source.  Namely, it is enough that for each \(j\)
either \(y_j\) or \(y_j+1\) satisfy
\eqref{rp-cut-margin}, and that the \(y_j\) are separated
by more than two.  Lemma~\ref{rp-parameter-signs} gives
\eqref{re-shifted-signs} in both cases, using \(q(Y-1)>0\)
when the cut is imposed at \(Y=y_j+1\).
For the additional domain point \(y_j+2\), its proof gives
slightly more than was stated: every factor \(Y+\rho_i\)
has absolute value at least \(5/2\) with a fixed sign,
so that sign persists for every shift of absolute value at
most two.  Consequently \(p_\rho(y_j+2)\ne0\).
No stronger separation between the unshifted parameters
\(a_j,b_j\) is needed for this version.

\begin{corollary}\label{re-density}
Let \(n\) be the number of steps.  If \(e\) is a choice
of \(0\) or \(1\) on each auxiliary factor, then
\begin{equation}
 x_n(Y)=\sum_{\rho,e}x_{\rho,e}(Y)P_\rho\otimes E_{ee},
 \qquad
 c_{{\rm wt},n}=\sum_{\rho,e}c_{\rho,e}P_\rho\otimes E_{ee},
 \qquad
 R_n=F_n^*w^n c_{{\rm wt},n}F_n .              \label{re-final-congruence}
\end{equation}
Initially \(x_{\rho,e}(Y)=\sigma p_\rho(Y)\) and
\(c_{\rho,e}=1\).
At a plus step of parameter \(a\), multiply
\(x_{\rho,e}(Y)\) by
\((Y-a)/(Y-a-1)\) when its new bit is \(0\), and by \(1\)
when it is \(1\).  Multiply \(c_{\rho,e}\) respectively by
\begin{equation}
                  q(a)/x_{\rho,e}(a+1),\qquad x_{\rho,e}(a).
                                                               \label{re-plus-scalar}
\end{equation}
At a minus step of parameter \(b\), the corresponding
multipliers are \((b-Y)/(b-Y-1)\), \(1\), and
\begin{equation}
                  q(b-1)/x_{\rho,e}(b-1),\qquad x_{\rho,e}(b).
                                                               \label{re-minus-scalar}
\end{equation}
Every quantity on the right is evaluated before the step.
Under the spectral-margin and separation hypotheses of
Corollary~\ref{rp-terminal-density-positive}, or under
the shifted hypotheses of Lemma~\ref{re-shifted-domain}, all
\(c_{\rho,e}>0\).
Consequently, for every smooth compactly supported radial
vector \(\Phi\),
\begin{equation}
 \int\langle R_n\Phi,\Phi\rangle
       =\int w^n\sum_{\rho,e}c_{\rho,e}
                       \|P_\rho(F_n\Phi)_e\|^2\geq0.             \label{re-positive-form}
\end{equation}
The integration measure is the radial measure specified at
the start of the section.
\end{corollary}

\begin{proof}
All initial projections commute, and every new formula for
\(x,c_{\rm wt}\) is diagonal on the new factor.
Induction therefore proves the spectral formulas and the
scalar recurrences.  They are exactly the scalar expressions
whose signs were proved in
Corollary~\ref{rp-terminal-density-positive}.
On a compact subset of \(G>0\), the inverses of \(G,\zeta\)
and \(\bar\zeta\), and all coefficients of the finitely many
operators, are bounded with all needed derivatives.
Differential operators do not enlarge supports.
Formal integration by parts is therefore legitimate in
\eqref{re-final-congruence}, giving
\eqref{re-positive-form}.  The pointwise orthogonality of
the \(P_\rho\) justifies its last expansion.
\end{proof}

For the empty sequence, \(n=0\), \(R_0=F_0=1\), and
\eqref{re-positive-form} is simply the ordinary squared norm.
When \(r=0\), there are no independent \(c_i\), and their
product measure and all assertions about their Euler
derivatives disappear.  When \(k=1\), the finite type has
one element, \(p_\rho=1\), and the initial correction
\(e(Y)\) is zero.  These cases require no additional
invertibility assumption.

This proves the actual algebraic frame and its compactly
supported positive form.  Passing to the specific noncompact
vectors used later still requires their regularity, decay,
support and operator-commutation hypotheses to be checked
against the analytic theorems of
Sections~\ref{jg-section} and \ref{jp-section}.
The algebraic congruence alone does not assert those hypotheses.
The correspondence of the formulas with the submitted source
is recorded separately; see \cite{alpo-uniform-source}.


\section{The coordinate integrals in the Jacquet estimates}
\label{jk-section}

This section proves the coordinate estimates used in the
construction of the integrable majorant for the Jacquet pairing.
All the integrals are Lebesgue integrals of nonnegative functions.
They are therefore well defined, possibly with value \(+\infty\),
before convergence is established.  Every interchange below follows
from Tonelli's theorem.  We give an exact value and a necessary and
sufficient convergence criterion for the basic integral; the
weighted estimate then follows from elementary averaging.
The notation and the precise correspondence with the source are
recorded at the end of the section.

\subsection{The basic integral and its exact convergence region}

Let \(p,m\) be nonnegative integers.  On
\[
 X_{p,m}=\mathbb C^p\times\mathbb C^p\times
                  \mathbb C^m\times\mathbb R
\]
write \(x=(a,b,z,t)\), with indices \(0\leq j<p\).
Lebesgue measure on \(\mathbb C\) means
\(d(\operatorname{Re}z)\,d(\operatorname{Im}z)\);
in dimension zero it has total mass one.
For \(0\leq l\leq p\), put
\begin{align}
 L_l(x)&=1+\sum_{j<l}|a_j|^2,\nonumber\\
 B_l(x)&=-\tfrac12\|z\|^2+it-\sum_{j\geq l}b_j\overline{a_j},
       \label{jk-coordinates}\\
 R_l(x)&=L_p(x)+\|z\|^2+\sum_{j\geq l}|b_j|^2,\nonumber\\
 d_l(x)&=\bigl(L_l(x)R_l(x)+|B_l(x)|^2\bigr)^{1/2}.
       \nonumber
\end{align}
In particular \(d_l\geq1\).  Define
\[
 \sigma_l=\begin{cases}2&l<p,\\m+1&l=p,\end{cases}
 \qquad
 I_{p,m}(\nu)=
       \int_{X_{p,m}}\prod_{l=0}^{p}d_l(x)^{-\sigma_l-\nu_l}\,dx,
       \qquad \nu\in\mathbb R^{p+1}.
\]
Write \(S_l=\nu_0+\cdots+\nu_l\).

\begin{theorem}\label{jk-exact}
The integral \(I_{p,m}(\nu)\) is finite if and only if
\begin{equation}\label{jk-cone}
 S_l>0\quad(0\leq l\leq p),\qquad
 S_l+2\sum_{j=l+1}^{p}\nu_j>0\quad(0\leq l<p).
\end{equation}
On this region it equals
\begin{equation}\label{jk-product}
 I_{p,m}(\nu)=
 \mathcal B_m(S_p)
 \prod_{l=0}^{p-1}
   \frac{(2\pi)^2}{S_l\bigl(S_l+2\sum_{j>l}\nu_j\bigr)},
\end{equation}
where, for \(v>0\),
\begin{equation}\label{jk-base-value}
 \mathcal B_m(v)=
  2^m\pi^{m+1/2}
  \frac{\Gamma((m+v)/2)\Gamma(v)}
       {\Gamma((m+1+v)/2)\Gamma(m+v)} .
\end{equation}
An empty product has value one.
\end{theorem}

\begin{proof}
First take \(p=0\).  Directly from \eqref{jk-coordinates},
\[
 d_0^2=(1+\|z\|^2/2)^2+t^2.
\]
For \(q>1/2\) and \(A>0\), integration of
\[
 u^{-q}=\frac{1}{\Gamma(q)}
                   \int_0^\infty s^{q-1}e^{-su}\,ds
\]
against \(u=A^2+t^2\) gives
\[
 \int_{\mathbb R}(A^2+t^2)^{-q}\,dt
  =\sqrt\pi\,\frac{\Gamma(q-\tfrac12)}{\Gamma(q)}
                      A^{1-2q}.
\]
Here we used \(\int_{\mathbb R}e^{-st^2}\,dt=\sqrt{\pi/s}\).
The same calculation in \(m\) complex coordinates gives, for \(h>m\),
\[
 \int_{\mathbb C^m}(1+\|z\|^2/2)^{-h}\,dz
                    =(2\pi)^m\frac{\Gamma(h-m)}{\Gamma(h)}.
\]
These identities follow directly from the displayed Laplace
integral and Tonelli; the Gaussian identity follows by squaring the
one dimensional integral and using polar coordinates in
\(\mathbb R^2\).
Apply them with \(q=(m+1+v)/2\) and \(h=m+v\).
This proves \eqref{jk-base-value} when \(v>0\).
If \(m+v\leq0\), the integral in \(t\) already diverges.
If \(m+v>0\) but \(v\leq0\), integration in \(t\) leaves
\((1+\|z\|^2/2)^{-m-v}\), whose integral diverges:
for \(m\geq1\), polar coordinates give a tail comparable to
\(\int_1^\infty r^{-1-2v}\,dr\).
For \(m=0\) the first alternative applies.
Thus the base integral is finite exactly when \(v>0\).

Now suppose \(p\geq1\), and set \(L=1+|a_0|^2\).
Completing the square gives the identity
\begin{equation}\label{jk-square}
 d_0^2=
 L\left|b_0-\frac{a_0B_1}{L}\right|^2+\frac{d_1^2}{L}.
\end{equation}
Indeed \(B_0=B_1-b_0\overline{a_0}\),
\(R_0=R_1+|b_0|^2\), and \(L_0=1\);
expansion of both sides gives the same expression.
The remaining \(d_l\), \(l\geq1\), do not depend on \(b_0\).
For fixed remaining coordinates, integration in \(b_0\) is finite
exactly when \(\nu_0>0\).  In that case polar coordinates give
\begin{equation}\label{jk-first-integration}
 \int_{\mathbb C}d_0^{-2-\nu_0}\,db_0
       =\frac{2\pi}{\nu_0}L^{\nu_0/2-1}d_1^{-\nu_0}.
\end{equation}
For \(\nu_0\leq0\), the two dimensional tail diverges, independently
of all the fixed coordinates.  The full integral is then infinite.

For fixed \(a_0\), make the substitution
\[
 a_j=\sqrt L\,\widetilde a_{j-1},\quad
 b_j=\sqrt L\,\widetilde b_{j-1}\quad(j\geq1),\qquad
 z=\sqrt L\,\widetilde z,\quad t=L\widetilde t.
\]
Its Jacobian is \(L^{2(p-1)+m+1}\).  The definitions show that
\[
 L_{l+1}=L\widetilde L_l,\quad
 B_{l+1}=L\widetilde B_l,\quad
 R_{l+1}=L\widetilde R_l,\quad
 d_{l+1}=L\widetilde d_l.
\]
The sum of the remaining critical exponents is
\(2(p-1)+m+1\), exactly the Jacobian exponent.
Consequently the remaining factor in \(a_0\) is
\[
 (1+|a_0|^2)^{-1-\nu_0/2-\sum_{j\geq1}\nu_j}.
\]
Its integral is finite exactly when
\(\nu_0+2\sum_{j\geq1}\nu_j>0\), and then equals
\(2\pi/(\nu_0+2\sum_{j\geq1}\nu_j)\).
Thus, whenever these two inequalities hold,
\begin{equation}\label{jk-recursion}
 I_{p,m}(\nu)=
 \frac{(2\pi)^2}
      {\nu_0(\nu_0+2\sum_{j\geq1}\nu_j)}
 I_{p-1,m}(\nu_0+\nu_1,\nu_2,\ldots,\nu_p).
\end{equation}
This is also an identity of extended nonnegative integrals.
The lower dimensional integral is strictly positive, so divergence
of either factor cannot be concealed by a zero factor.
Induction now proves necessity and sufficiency of all the
inequalities in \eqref{jk-cone}, and multiplication of the factors
in \eqref{jk-recursion} proves \eqref{jk-product}.
\end{proof}

\begin{corollary}\label{jk-uniform}
On every compact subset of the open region \eqref{jk-cone},
the integrands defining \(I_{p,m}\) have a common integrable
majorant.  In particular the integral depends continuously on
\(\nu\) there.  The same assertion holds after any fixed number of
derivatives with respect to the exponents.
\end{corollary}

\begin{proof}
At any point \(\nu^*\) of the region, choose \(\varepsilon>0\)
such that \(\nu^*-2\varepsilon(1,\ldots,1)\) remains in the region.
For \(|\nu_l-\nu_l^*|<\varepsilon\), the inequality \(d_l\geq1\)
gives an integrable majorant with exponents
\(\nu_l^*-\varepsilon\).
For derivatives, each additional factor is a product of powers
of \(\log d_l\).  The elementary bound
\((\log u)^N\leq C_{N,\varepsilon}u^\varepsilon\), \(u\geq1\),
gives the majorant with exponents
\(\nu_l^*-2\varepsilon\), with a smaller \(\varepsilon\) if needed.
A finite covering of the compact set supplies a finite sum of
such majorants.  Dominated convergence proves the assertions.
\end{proof}

\subsection{An integral with one coordinate removed}

For \(a\in\mathbb C^{p+1}\), \(b\in\mathbb C^p\),
\(z\in\mathbb C^m\), \(t\in\mathbb R\), regard
\[
 x=(a,(0,b_0,\ldots,b_{p-1}),z,t)\in X_{p+1,m}.
\]
Use the quantities \(L_l,d_l,\sigma_l\) for this larger space.
For \(\alpha\in\mathbb R\) and \(\nu\in\mathbb R^{p+1}\), define
\[
 T_{p,m}(\alpha,\nu)=
 \int L_1^{(\alpha-2)/2}d_1^{-\alpha}
                \prod_{i=0}^{p}d_{i+1}^{-\sigma_{i+1}-\nu_i}
                                      \,da\,db\,dz\,dt .
\]

\begin{corollary}\label{jk-tail}
This integral is finite exactly when
\[
 \alpha+2\sum_{i=0}^{p}\nu_i>0
 \quad\hbox{and}\quad
 (\nu_0+\alpha,\nu_1,\ldots,\nu_p)
                 \text{ satisfies \eqref{jk-cone}}.
\]
Under these conditions,
\begin{equation}\label{jk-tail-value}
 T_{p,m}(\alpha,\nu)=
       \frac{2\pi}{\alpha+2\sum_i\nu_i}
          I_{p,m}(\nu_0+\alpha,\nu_1,\ldots,\nu_p).
\end{equation}
No sign condition on \(\alpha\) is required.
\end{corollary}

\begin{proof}
Set \(L=1+|a_0|^2\) and use the same scaling as in the
proof of Theorem~\ref{jk-exact}.
The Jacobian is now \(L^{2p+m+1}\) and cancels the sum
of the critical exponents of \(d_1,\ldots,d_{p+1}\).
The residual power is
\(L^{-1-\alpha/2-\sum_i\nu_i}\).
Its integral over \(a_0\in\mathbb C\) gives the displayed factor
and its exact convergence condition.  The other factor is the
strictly positive integral on the right of \eqref{jk-tail-value}.
Tonelli proves both directions, including divergence cases.
\end{proof}

\subsection{Three elementary estimates}

\begin{lemma}\label{jk-sublevel}
Let \(P(u)=c_0+c_1u+c_2u^2\), with complex coefficients, and let
\(M=\max_j|c_j|\).  If \(M>0\) and \(\varepsilon>0\), then
\begin{equation}\label{jk-sublevel-measure}
 \bigl|\{u\in[1,2]:|P(u)|\leq\varepsilon M\}\bigr|
                                  \leq \sqrt{432\varepsilon}.
\end{equation}
Consequently, for \(0<\delta\leq1/2\),
\begin{equation}\label{jk-sublevel-integral}
 \int_1^2(1+|P(u)|)^{-\delta}\,du
                         \leq30(1+M)^{-\delta/2}.
\end{equation}
\end{lemma}

\begin{proof}
Let \(E\) be the set in \eqref{jk-sublevel-measure} and
\(\mu=|E|>0\).  Select three points \(u_1,u_2,u_3\) of \(E\)
at pairwise distance at least \(\mu/6\): after selecting each
of the first two points, delete its interval of radius
\(\mu/6\); the deleted measure is at most \(2\mu/3\).
Lagrange interpolation gives
\[
 P(u)=\sum_{i=1}^{3}P(u_i)
              \prod_{j\ne i}\frac{u-u_j}{u_i-u_j}.
\]
Each coefficient of the numerator \(\prod_{j\ne i}(u-u_j)\)
has absolute value at most \(4\), since \(1\leq u_j\leq2\).
Each denominator has absolute value at least \(\mu^2/36\).
Comparing coefficients yields
\(M\leq432\varepsilon M/\mu^2\), proving
\eqref{jk-sublevel-measure}.

For \(M\geq1\), split at \(|P|=\sqrt M\).
The contribution from the first set is at most
\(\sqrt{432}M^{-1/4}\); on its complement the integrand
is at most \(M^{-\delta/2}\).
Since \(\delta\leq1/2\), their sum is at most
\((\sqrt{432}+1)M^{-\delta/2}
 \leq30(1+M)^{-\delta/2}\).
For \(M\leq1\), including \(M=0\), the integral is at most one,
which satisfies the same bound.
\end{proof}

\begin{lemma}\label{jk-riesz}
If \(0<e<\gamma<2\), there is a finite constant
\(C_{e,\gamma}\), independent of \(c\in\mathbb C\) and \(D>0\),
such that
\begin{equation}\label{jk-riesz-bound}
 \int_{\mathbb C}
 (|w-c|^2+D^2)^{-1+e/2}|w|^{-\gamma}\,dw
       \leq C_{e,\gamma}(|c|^2+D^2)^{-(\gamma-e)/2}.
\end{equation}
\end{lemma}

\begin{proof}
Scale by \(r=(|c|^2+D^2)^{1/2}\).  The resulting factor is
\(r^{e-\gamma}\), and the new parameters satisfy
\(|c|^2+D^2=1\).
For \(|w|\geq3\), the integrand is at most a constant times
\(|w|^{-2+e-\gamma}\); its integral is uniformly finite because
\(\gamma>e\).
It remains to consider \(|w|\leq3\).
If \(|c|\leq1/2\), then \(D\geq\sqrt3/2\), and the only
possible singularity is bounded by a constant times
\(|w|^{-\gamma}\), integrable because \(\gamma<2\).
If \(|c|>1/2\), the discs \(|w|<1/4\) and
\(|w-c|<1/4\) are disjoint.
On the first, the other factor is bounded uniformly.
On the second, \(|w|^{-\gamma}\) is bounded uniformly, while
\[
 (|w-c|^2+D^2)^{-1+e/2}\leq |w-c|^{-2+e},
 \qquad w\ne c,
\]
which is integrable because \(e>0\).
On the remainder both distances are at least \(1/4\), so
the integrand is bounded uniformly.
This proves the normalized estimate and hence the result.
\end{proof}

\begin{lemma}\label{jk-slow}
Let \(0<\kappa<1\), \(R>0\), and
\[
 w(\beta)=\prod_{i=1}^{n}
                (A_i^2+(\beta-b_i)^2)^{-t_i/2},
 \qquad
 A_i\geq R,\quad t_i\geq0,\quad \sum_i t_i\leq T .
\]
For every \(c\in\mathbb R\),
\begin{equation}\label{jk-slow-bound}
 \int_{\mathbb R}w(\beta)|\beta-c|^{-\kappa}\,d\beta
 \leq
 \left(1+\frac{2\,19^{\max(T,0)/2}}{1-\kappa}\right)
 R^{-\kappa}\int_{\mathbb R}w(\beta)\,d\beta .
\end{equation}
The assertion is valid for extended nonnegative integrals.
\end{lemma}

\begin{proof}
Translate to \(c=0\).  The case \(\int w=+\infty\) is
immediate, so suppose that integral is finite.
For \(|\beta|\leq R\) and \(R\leq s\leq2R\),
\[
 A_i^2+(s-b_i)^2
 \leq A_i^2+2(\beta-b_i)^2+18R^2
 \leq19\bigl(A_i^2+(\beta-b_i)^2\bigr).
\]
Therefore \(w(\beta)\leq19^{\max(T,0)/2}w(s)\).
Averaging over \(s\in[R,2R]\) gives
\[
 w(\beta)\leq19^{\max(T,0)/2}R^{-1}
                       \int_{\mathbb R}w(s)\,ds
             \qquad (|\beta|\leq R).
\]
Integration of \(|\beta|^{-\kappa}\) on \([-R,R]\) gives
the second term in the constant.
On \(|\beta|\geq R\), use
\(|\beta|^{-\kappa}\leq R^{-\kappa}\).
\end{proof}

\subsection{A finite integral below the critical exponents}

\begin{proposition}\label{jk-weighted}
Let \(P\geq1\), \(m\geq0\), \(0<\gamma\leq1\), and
\(e_0,\ldots,e_P>0\).  Suppose that
\begin{equation}\label{jk-loss-condition}
 e_0+2\sum_{l=1}^{P}e_l<\gamma/2.
\end{equation}
Use \eqref{jk-coordinates} on \(X_{P,m}\), and put
\(s_l=\sigma_l-e_l\).  Then
\begin{equation}\label{jk-weighted-integral}
 \int_{X_{P,m}}\prod_{l=0}^{P}d_l(x)^{-s_l}
                                  |b_0|^{-\gamma}\,dx<\infty .
\end{equation}
\end{proposition}

\begin{proof}
The hypotheses give \(0<e_0<\gamma<2\) and \(0<e_l<1\)
for all \(l\); hence \(s_l\geq0\).
Fix all variables except \(b_0\), and let \(L=1+|a_0|^2\).
Write \(D_0=d_0|_{b_0=0}\) and \(D_l=d_l\) for \(l\geq1\);
the latter quantities are independent of \(b_0\).
With \(c=a_0B_1/L\), identity \eqref{jk-square} gives
\[
 d_0^2=L|b_0-c|^2+D_1^2/L,\qquad
 D_0^2=L|c|^2+D_1^2/L .
\]
Applying Lemma~\ref{jk-riesz}, with its positive parameter
\(D=D_1/L\), yields
\begin{equation}\label{jk-b-integration}
 \int_{\mathbb C}d_0^{-2+e_0}|b_0|^{-\gamma}\,db_0
              \leq C L^{-1+\gamma/2}D_0^{-(\gamma-e_0)}.
\end{equation}
The powers of \(L\) here follow by factoring \(L\) from
\(d_0^2\) and using
\(|c|^2+(D_1/L)^2=D_0^2/L\).

Put \(\alpha=\gamma/2-e_0>0\) and \(\kappa=\gamma/2\).
For the point with \(b_0=0\), let
\[
 \theta_l=\operatorname{Im}\sum_{j\geq l}b_j\overline{a_j}.
\]
Then
\[
 D_l^2=A_l^2+(t-\theta_l)^2,\qquad
 A_l^2=L_lR_l+
       \left(-\|z\|^2/2-
                   \operatorname{Re}\sum_{j\geq l}
                                     b_j\overline{a_j}\right)^2.
\]
For \(l\geq1\), both \(L_l\) and \(R_l\) are at least \(L\);
thus \(A_l\geq L\).
Moreover \(B_0=B_1\) and \(R_0=R_1\) at \(b_0=0\).
It follows that
\[
 D_1\leq\sqrt L\,D_0,\qquad |t-\theta_1|\leq D_0.
\]
Apart from the single value \(t=\theta_1\), this implies
\[
 D_0^{-(\gamma-e_0)}
   \leq L^{\alpha/2}D_1^{-\alpha}|t-\theta_1|^{-\kappa}.
\]
Apply Lemma~\ref{jk-slow} to
\[
 w(t)=D_1^{-\alpha}\prod_{l=1}^{P}D_l^{-s_l}.
\]
Its exponents are nonnegative and have sum at most
\(\alpha+\sum_{l=1}^{P}\sigma_l\); all its parameters
\(A_l\) are at least \(L\).
The constant is therefore independent of the fixed coordinates.
After integrating in \(t\), the product of the powers of \(L\) is
\[
 L^{-1+\gamma/2}L^{\alpha/2}L^{-\kappa}
                                      =L^{(\alpha-2)/2}.
\]
By \eqref{jk-b-integration}, the integral in
\eqref{jk-weighted-integral} is consequently at most a constant
times
\[
 T_{P-1,m}\bigl(\alpha,(-e_1,\ldots,-e_P)\bigr).
\]
Condition \eqref{jk-loss-condition} says
\(\alpha>2\sum_{l=1}^{P}e_l\).
The first condition of Corollary~\ref{jk-tail} follows.
Every partial sum for the shifted vector
\((\alpha-e_1,-e_2,\ldots,-e_P)\) is positive.
Every other required expression is of the form
\[
 \alpha-\sum_{j\leq l}e_j-2\sum_{j>l}e_j
                       \geq\alpha-2\sum_{j=1}^{P}e_j>0,
\]
with the appropriate initial segment starting at \(j=1\).
Thus all conditions of that corollary hold, proving finiteness.
\end{proof}

\subsection{The estimate with a linear character}

For \(x=(a,b,z,t)\in X_{P,m}\), define the complex linear form
\[
 \ell(x)=\sum_{j=0}^{P-1}(a_j\eta_j+b_j\eta'_j)
                                      +\sum_{i=1}^{m}z_i\xi_i,
\]
where \(\eta,\eta'\in\mathbb C^P\), \(\xi\in\mathbb C^m\).

\begin{theorem}\label{jk-oscillatory}
Let \(P\geq1\), \(m\geq0\),
\[
 0<\delta\leq\tfrac12,\qquad
 0<\gamma\leq\delta/8,\qquad e_l>0,
 \qquad e_0+2\sum_{l=1}^{P}e_l<\gamma/2 .
\]
Set \(s_l=\sigma_l-e_l\).
There is a finite constant depending only on
\(P,m,\delta,\gamma,e\) such that, for every
\(\eta,\eta',\xi,c_0\) with \(\eta'_0\ne0\) and \(c_0\in\mathbb R\),
\begin{equation}\label{jk-os}
 \int_{X_{P,m}}\prod_{l=0}^{P}d_l(x)^{-s_l}
          \bigl(1+|2\operatorname{Im}\ell(x)+c_0|\bigr)^{-\delta}
             \,dx
 \leq C|\eta'_0|^{-\gamma}(1+|c_0|)^{-\delta/4}.
\end{equation}
In particular the bound is uniform in all the other coefficients
of the linear form.
\end{theorem}

\begin{proof}
Write \(F(x)=\prod_l d_l(x)^{-s_l}\).
We use two changes of variables preserving Lebesgue measure.
All the averaging is over \([1,2]\), an interval of length one.

First put
\[
 \rho(u)=\frac{(1+iu)^2}{1+u^2},\qquad
 \mathcal R_u(a,b,z,t)=(\rho(u)a,\rho(u)b,\rho(u)z,t).
\]
Since \(|\rho(u)|=1\), this map preserves measure.
All \(L_l,R_l,B_l\), and hence \(F\), are unchanged, whereas
\(\ell(\mathcal R_u x)=\rho(u)\ell(x)\).
For \(X\in\mathbb C\), direct multiplication gives
\[
 2\operatorname{Im}(\rho(u)X)+c_0
     =\frac{(2\operatorname{Im}X+c_0)
                 +4\operatorname{Re}X\,u
                 +(c_0-2\operatorname{Im}X)u^2}{1+u^2}.
\]
Let \(M\) be the maximum of the absolute values of the three
coefficients of the numerator.  Then
\[
 |c_0|\leq M,\qquad
 |\operatorname{Re}X|\leq M/4,\qquad
 |\operatorname{Im}X|\leq M/2,\qquad |X|\leq M.
\]
As \(1+u^2\leq5\), Lemma~\ref{jk-sublevel} shows that
\begin{equation}\label{jk-rotation-bound}
 \int_1^2
 (1+|2\operatorname{Im}(\rho(u)X)+c_0|)^{-\delta}\,du
 \leq30\,5^\delta
              (1+|c_0|)^{-\delta/4}(1+|X|)^{-\delta/4}.
\end{equation}
Indeed \((1+M)^2\geq(1+|c_0|)(1+|X|)\).
Changing variables by \(\mathcal R_u\), and then averaging,
bounds the left side of \eqref{jk-os} by
\begin{equation}\label{jk-after-rotation}
 30\,5^\delta(1+|c_0|)^{-\delta/4}
                 \int F(x)(1+|\ell(x)|)^{-\delta/4}\,dx .
\end{equation}

The second change of variables is
\[
 \mathcal S_u(a,b,z,t)=
 ((ua_0,a_1,\ldots),(b_0/u,b_1,\ldots),z,t).
\]
The real Jacobian is \(u^2u^{-2}=1\).
For \(1\leq u\leq2\), each \(L_l\) can only increase,
each \(R_l\) is at least one quarter of its original value,
and each \(B_l\) is unchanged.  Thus
\[
 d_l(\mathcal S_u x)\geq d_l(x)/2,\qquad
 F(\mathcal S_u x)\leq2^{\sum_l s_l}F(x).
\]
Here \(s_l\geq0\), as already observed in
Proposition~\ref{jk-weighted}.
Consequently change of variables and averaging give
\begin{align}
 \int F(x)(1+|\ell(x)|)^{-\delta/4}\,dx
 &\leq2^{\sum_l s_l}
   \int F(x)\int_1^2
       (1+|\ell(\mathcal S_u x)|)^{-\delta/4}\,du\,dx.
 \label{jk-after-scaling}
\end{align}
Write
\[
 A=a_0\eta_0,\qquad C=b_0\eta'_0,\qquad
 B=\sum_{j\geq1}(a_j\eta_j+b_j\eta'_j)+\sum_i z_i\xi_i.
\]
Then \(u\ell(\mathcal S_u x)=C+Bu+Au^2\).
Since \(u\leq2\), another application of
Lemma~\ref{jk-sublevel}, now with exponent \(\delta/4\), gives
\[
 \int_1^2(1+|\ell(\mathcal S_u x)|)^{-\delta/4}\,du
 \leq30\,2^{\delta/4}(1+|C|)^{-\delta/8}
 \leq30\,2^{\delta/4}|b_0\eta'_0|^{-\gamma}
\]
whenever \(b_0\eta'_0\ne0\).
The last inequality uses \(0<\gamma\leq\delta/8\).
As \(\eta'_0\ne0\), the excluded set \(b_0=0\) has measure
zero.  Alternatively the bound holds everywhere with value
\(+\infty\) assigned to the singular factor.
Insert this inequality into \eqref{jk-after-scaling}.
Proposition~\ref{jk-weighted} bounds the resulting integral of
\(F(x)|b_0|^{-\gamma}\).
Together with \eqref{jk-after-rotation}, this proves
\eqref{jk-os}.
\end{proof}

\subsection{Sharp convergence restrictions and scope}

The inequalities in \eqref{jk-cone} are necessary, not merely
convenient sufficient assumptions.  The proof of
Theorem~\ref{jk-exact} gives a divergent complex coordinate integral
or the divergent final real or radial integral whenever any of
them fails.  For example, \(p=0,\nu_0=0\) gives logarithmic
divergence; for \(p=1\), the vectors
\((0,1)\) and \((2,-1)\) respectively violate the first
complex integration condition and the second one at equality,
although the remaining inequalities are positive.

The extra weight in Theorem~\ref{jk-oscillatory} is also essential.
For \(P\geq1\) and \(e_0>0\), the integral of
\(\prod_l d_l^{-\sigma_l+e_l}\) without that weight is infinite,
because its first exponent parameter is \(\nu_0=-e_0<0\).
If all the coefficients \(\eta,\eta',\xi\) vanish, the weight
is a positive constant, so the same divergence persists.
This supplies a counterexample to omitting the condition on the
linear form; it is not a counterexample under the stated
condition \(\eta'_0\ne0\).
In Lemma~\ref{jk-riesz}, \(\gamma\geq2\) makes the origin
nonintegrable, and \(\gamma\leq e\) makes infinity nonintegrable.
The restriction \(e>0\) is also required for a bound uniform
as \(D\downarrow0\) with \(c\ne0\): a disc around \(c\)
has a divergent limiting integral when \(e\leq0\).

The definitions \eqref{jk-coordinates} and the critical exponents
are exactly those in the source's coordinate integral
\(\mathrm{CHInt}\).  Theorem~\ref{jk-exact} proves its assertion
\(\mathrm{CH}\) and strengthens it to an exact evaluation and a
necessary convergence criterion.  Corollary~\ref{jk-tail}
proves its assertion \(\mathrm{TailFin}\), identity
\eqref{jk-square} is its \(\mathrm{C1}\), and the three elementary
lemmas prove \(\mathrm{SUB}\), \(\mathrm{KER}\), and
\(\mathrm{SLOW}\).  Theorem~\ref{jk-oscillatory}, with \(P=p+1\),
is exactly its assertion \(\mathrm{OS}\).
These identifications refer to the definitions in
\texttt{TailIntegrals/TailInt.lean} in the source archive
\cite{alpo-uniform-source}; the accompanying source trace records
the other files and the parts read.

This section supplies actual estimates, with proofs, rather than
assuming the existence of a Jacquet majorant.
To obtain that majorant one still has to prove the geometric
identities expressing the unipotent integral in these coordinates,
the corresponding induction on the rank, and the group integral
used for the matrix coefficient bound.
The convergence and differentiability of the final Jacquet
pairing also require the subsequent unfolding and domination
arguments.  None of those further assertions is inferred solely
from the coordinate estimates proved here.


\section{Weighted integrability on the unipotent group}
\label{ju-section}

We now pass from the coordinate estimates of
Section~\ref{jk-section} to the unipotent integral that occurs in
the Jacquet majorant.  We prove the required coordinate identities,
including their Jacobians, and the complete induction.
The negative norm of the vector defining the character is used at
a specific point: it prevents the coefficient needed by
Theorem~\ref{jk-oscillatory} from vanishing in the last step.

\subsection{The Hermitian spaces, flags, and unipotent groups}

Fix \(k\geq1\).  Let \(V_0=\mathbb C^{k+1}\), with its usual
Euclidean norm and Hermitian form
\[
 h_0(v,w)=\sum_{j=1}^{k}v_j\overline{w_j}
                         -v_-\overline{w_-}.
\]
All Hermitian forms here are linear in the first variable.
Inductively set
\[
 V_r=\mathbb Ce_0\oplus V_{r-1}\oplus\mathbb Cf_0,\qquad
 h_r(\alpha e_0+v+\beta f_0,\alpha'e_0+w+\beta'f_0)
       =\alpha\overline{\beta'}+\beta\overline{\alpha'}
                                      +h_{r-1}(v,w).
\]
The three summands are Euclidean orthogonal, and \(e_0,f_0\)
have Euclidean norm one.  Relabel the vectors of the smaller
space as \(e_1,\ldots,e_{r-1}\), \(f_1,\ldots,f_{r-1}\).
Thus \(h_r(e_i,f_j)=\delta_{ij}\).
Write \(h_r(v,w)=w^*H_rv\); then \(H_r^*=H_r\) and \(H_r^2=1\).
Let
\[
 G_r=U(V_r,h_r),\qquad K_r=G_r\cap U(V_r)
\]
where the second unitary group uses the Euclidean norm.
Every element of \(K_r\) commutes with \(H_r\).
The subgroup \(N_r\) consists of the elements \(n\in G_r\) with
\[
 ne_j-e_j\in\operatorname{span}(e_0,\ldots,e_{j-1})
                 \quad(0\leq j<r),\qquad
 nw-w\in\operatorname{span}(e_0,\ldots,e_{r-1})
                 \quad(w\in V_0).
\]
In particular \(N_0=\{1\}\).

Let \(u_1\) denote the first positive unit vector in \(V_0\),
and put
\[
 \varepsilon=(u_1+u_-)/\sqrt2,\qquad
 (\phi_0,\ldots,\phi_r)=(f_0,\ldots,f_{r-1},\varepsilon),
 \qquad \phi'_j=H_r\phi_j .
\]
The vectors \(\phi_j,\phi'_j\), together with
\(u_2,\ldots,u_k\), form a Euclidean orthonormal basis.
Their Hermitian products are
\[
 h_r(\phi_i,\phi_j)=h_r(\phi'_i,\phi'_j)=0,\qquad
 h_r(\phi_i,\phi'_j)=\delta_{ij}.
\]
For \(g\in G_r\), define
\[
 \Delta_l(g)=
       \|g\phi_0\wedge\cdots\wedge g\phi_l\|>0
                         \quad(0\leq l\leq r),
 \qquad
 \Xi_s(g)=\prod_{l=0}^{r}\Delta_l(g)^{-s_l}.
\]
The exterior norm is the Euclidean one: its square is the
determinant of the Gram matrix.
The critical exponents are
\[
 s_l^0=\begin{cases}2&l<r,\\k&l=r.\end{cases}
\]

\begin{lemma}\label{ju-heisenberg}
For \(r\geq1\), let \(W=V_{r-1}\) and \(J=H_{r-1}\).
Set
\[
 n(v,t)=
 \begin{pmatrix}
 1&-v^*J&it-\tfrac12 h_{r-1}(v,v)\\
 0&I&v\\0&0&1
 \end{pmatrix},\qquad
 \iota(m)=\operatorname{diag}(1,m,1).
\]
Every element of \(N_r\) has a unique expression
\(n(v,t)\iota(m)\), with \(v\in W\), \(t\in\mathbb R\),
and \(m\in N_{r-1}\).  These are global smooth coordinates.
There is a Haar measure on \(N_r\) which in these coordinates is
\(dv\,dt\,dm\), where \(dm\) is the recursively chosen Haar
measure on \(N_{r-1}\), starting with mass one on \(N_0\).
\end{lemma}

\begin{proof}
The matrix \(n(v,t)\) preserves the Hermitian form by direct
multiplication.  It satisfies the defining flag conditions.
An element \(g\in N_r\) fixes \(e_0\).  Since it preserves the
form, its last row is \((0,0,1)\); consequently
\[
 g=\begin{pmatrix}1&\alpha&\beta\\0&m&v\\0&0&1\end{pmatrix}.
\]
Pairing its middle columns with each other gives \(m^*Jm=J\).
Pairing them with its last column gives
\(\alpha=-v^*Jm\), and the norm of the last column gives
\(2\operatorname{Re}\beta=-h_{r-1}(v,v)\).
The flag conditions on the middle block are exactly those of
\(N_{r-1}\).  Thus \(t=\operatorname{Im}\beta\) gives the
claimed factorization, its uniqueness, and its smooth inverse.

Multiplication of the displayed matrices gives
\[
 n(v,t)n(v',t')=
 n(v+v',\,t+t'-\operatorname{Im}h_{r-1}(v',v)),\qquad
 \iota(m)n(v,t)\iota(m)^{-1}=n(mv,t).
\]
Left and right translations in the first group have real
Jacobian one.  Also \(|\det_{\mathbb C}m|=1\), by taking
determinants in \(m^*Jm=J\), so the real Jacobian of
\((v,t)\mapsto(mv,t)\) is one.
It follows directly from these two formulas and induction that
\(dv\,dt\,dm\) is invariant under both left and right translations.
It is positive on nonempty open sets and finite on compact sets,
so it is a Haar measure.
\end{proof}

\subsection{The exact exterior norm formula}

\begin{lemma}\label{ju-coordinate-formula}
On \(W=V_{r-1}\), write
\[
 v=\sum_{j=0}^{r-1}(a_j\phi'_j+b_j\phi_j)
                        +\sum_{i=2}^{k}z_i u_i,\qquad
 \widetilde t=t+\operatorname{Im}\sum_jb_j\overline{a_j}.
\]
The map \((v,t)\mapsto(a,b,z,\widetilde t)\) preserves
Lebesgue measure.  With \(p=r\), \(m=k-1\), and the
quantities of \eqref{jk-coordinates}, one has
\begin{equation}\label{ju-formula}
 \Delta_l(n(v,t))=d_l(a,b,z,\widetilde t)
                                      \quad(0\leq l\leq r).
\end{equation}
\end{lemma}

\begin{proof}
The change from \(v\) to \((a,b,z)\) is unitary for the
Euclidean norm.  The subsequent change in \(t\) is a translation
depending on those coordinates, and hence has Jacobian one.
As
\[
 h_{r-1}(v,v)=
         2\operatorname{Re}\sum_jb_j\overline{a_j}+\|z\|^2,
\]
the upper right entry of \(n(v,t)\) is precisely \(B_0\).
Its relevant column vectors are
\[
 n(v,t)f_0=B_0e_0+v+f_0,\qquad
 n(v,t)\phi_j=\phi_j-\overline{a_j}e_0.
\]
Here the \(\phi_j\) on the right belong to \(W\), and correspond
to indices \(j+1\) in \(V_r\).
In the wedge with \(l\) such middle columns, subtract
\(\sum_{j<l}b_jn(v,t)\phi_j\) from the first column.
The first column becomes
\[
 y=B_le_0+f_0+\sum_j a_j\phi'_j+
                   \sum_{j\geq l}b_j\phi_j+\sum_i z_i u_i.
\]
The remaining columns are \(u_j=\phi_j-\overline{a_j}e_0\),
\(j<l\).
Let \(a_{<l}=(a_0,\ldots,a_{l-1})^T\).
Their Gram matrix, with \(y\) first, is
\[
 \begin{pmatrix}
 R_l+|B_l|^2&-\overline{B_l}\,a_{<l}^*\\
 -B_l a_{<l}&I+a_{<l}a_{<l}^*
 \end{pmatrix}.
\]
The lower block has determinant \(L_l\) and inverse
\(I-a_{<l}a_{<l}^*/L_l\).
The inverse formula follows by multiplication; the determinant
follows by using a unitary basis whose first vector is proportional
to \(a_{<l}\), or immediately when this vector is zero.
Taking the Schur complement now gives
\[
 L_l\left(R_l+|B_l|^2
       -|B_l|^2\frac{\|a_{<l}\|^2}{L_l}\right)
                                  =L_lR_l+|B_l|^2.
\]
For \(l=0\) the same identity uses the empty lower block.
This proves \eqref{ju-formula}.
\end{proof}

\subsection{Orthogonalization of an isotropic flag}

\begin{lemma}\label{ju-isotropic-qr}
Let \(Y_0,\ldots,Y_r\in V_r\) be linearly independent and
span a totally isotropic subspace.  There exist
\(\kappa\in K_r\) and an upper triangular matrix
\(\tau=(\tau_{ij})_{0\leq i,j\leq r}\), with positive real
diagonal entries, such that
\[
 Y_j=\sum_{i\leq j}\tau_{ij}\kappa\phi_i,\qquad
 \|Y_0\wedge\cdots\wedge Y_l\|
                            =\prod_{j=0}^{l}\tau_{jj}.
\]
\end{lemma}

\begin{proof}
Apply Euclidean Gram--Schmidt to the \(Y_j\).
It gives Euclidean orthonormal vectors \(q_j\) in their span
and the indicated upper triangular coefficients with positive
diagonal.  Their span is still totally isotropic.
Let \(V_r=V_+\oplus V_-\) be the orthogonal decomposition
into the eigenspaces of \(H_r\), and let \(P_\pm\) be the
corresponding projections.  The two identities
\[
 \langle q_i,q_j\rangle=\delta_{ij},\qquad h_r(q_j,q_i)=0
\]
imply
\[
 \langle P_+q_i,P_+q_j\rangle
   =\langle P_-q_i,P_-q_j\rangle=\delta_{ij}/2.
\]
The same is true for the \(\phi_j\).
There are \(r+1=\dim V_-\) negative projected vectors,
and \(\dim V_+=k+r\geq r+1\).
Extend the positive projected frames to orthonormal bases;
the negative projected frames are already bases.
The two resulting unitary changes of basis give
\(\kappa=\kappa_+\oplus\kappa_-\in K_r\) with
\(\kappa\phi_j=q_j\).
Taking the wedge of the first \(l+1\) triangular expansions
proves the norm formula.
Since \(H_r\in K_r\) interchanges \(\phi_j\) and \(\phi'_j\),
replacing \(\kappa\) by \(\kappa H_r\) gives the identical
assertion with the primed reference frame, exactly as in the
source's assertion \(\mathrm{IW}\).
\end{proof}

\begin{lemma}\label{ju-levi-factor}
For \(m\in G_{r-1}\), one can choose
\(\kappa\in K_{r-1}\) such that
\[
 m\phi_0=\Delta_0(m)\kappa\phi_0 .
\]
For that choice, putting \(v'=\kappa^{-1}v\), one has
\begin{align}
 \Delta_0(n(v,t)\iota(m))&=\Delta_0(n(v',t)),\nonumber\\
 \Delta_{j+1}(n(v,t)\iota(m))
     &=\Delta_j(m)\Delta_{j+1}(n(v',t))
                                    \quad(0\leq j<r).
 \label{ju-levi-norms}
\end{align}
\end{lemma}

\begin{proof}
Apply Lemma~\ref{ju-isotropic-qr} on \(W\) to \(m\phi_j\).
Conjugation gives
\(n(v,t)\iota(\kappa)=\iota(\kappa)n(\kappa^{-1}v,t)\).
The first column \(f_0\) is fixed by \(\iota(m)\).
In the first \(j+2\) columns of the wedge, the triangular
expansion of the other \(j+1\) columns therefore contributes
the factor \(\prod_{i\leq j}\tau_{ii}=\Delta_j(m)\).
The exterior action of \(\iota(\kappa)\) is Euclidean unitary,
so it does not change the norm.
The assertion for the first column follows by the same
unitary invariance.  Finally
\(\tau_{00}=\|m\phi_0\|=\Delta_0(m)\).
\end{proof}

\subsection{The character and the nonvanishing coefficient}

Fix a vector \(z\in V_0\) with \(h_0(z,z)<0\).
For \(r\geq1\), put
\[
 \ell_z(g)=\sum_{a=0}^{r-2}(ge_{a+1})_{e_a}
                                  +(gz)_{e_{r-1}},
 \qquad \lambda_z(g)=2\operatorname{Im}\ell_z(g).
\]
For \(r=0\), put \(\ell_z=\lambda_z=0\).
This is the character normalization in the source.
The flag conditions show that each consecutive entry, and the
last coefficient involving \(z\), is additive under multiplication
in \(N_r\): in multiplying consecutive entries, all possible
intermediate terms lie strictly earlier in the flag.
For the last coefficient use
\(n'z-z\in\operatorname{span}(e_0,\ldots,e_{r-1})\)
and the fact that the coefficient of \(e_{r-1}\) in \(ne_j\)
is \(\mathbf1_{\{j=r-1\}}\).
Thus \(\lambda_z\) is additive on \(N_r\), and
\(\exp(i\lambda_z)\) is a unitary character.
For \(W=V_{r-1}\), let \(f_z=e_0\) if \(r-1\geq1\),
and \(f_z=z\) if \(r-1=0\).
Inspection of the first row and the middle block gives
\begin{equation}\label{ju-lambda}
 \lambda_z(n(v,t)\iota(m))
       =2\operatorname{Im}h_{r-1}(v,mf_z)+\lambda_z(m).
\end{equation}
Indeed the first summand of \(\ell_z\) is
\(-h_{r-1}(mf_z,v)\), and its remaining summands are
exactly \(\ell_z(m)\).

The fixed number
\[
 c_z=h_{r-1}(\phi_0,f_z)
\]
is nonzero.  If \(r-1\geq1\), it equals
\(h_{r-1}(f_0,e_0)=1\).
If \(r-1=0\), it equals \(h_0(\varepsilon,z)\).
Were this zero, the first positive coordinate of \(z\)
would equal its negative coordinate; consequently
\(h_0(z,z)=\sum_{i=2}^{k}|z_i|^2\geq0\), a contradiction.

\begin{proposition}\label{ju-inner}
Let \(0<\delta\leq1/2\), \(0<\gamma\leq\delta/8\),
and \(s_l=s_l^0-e_l\), where \(e_l>0\) and
\(e_0+2\sum_{l=1}^{r}e_l<\gamma/2\).
Define
\[
 s'_j=s_{j+1}-\gamma\,\mathbf1_{\{j=0\}}
                                      \quad(0\leq j<r).
\]
There is a constant independent of \(m\in G_{r-1}\) such that
\begin{align}
 &\int_{W\times\mathbb R}
  \Xi_s(n(v,t)\iota(m))
  (1+|\lambda_z(n(v,t)\iota(m))|)^{-\delta}\,dv\,dt
  \nonumber\\
 &\hspace{15mm}\leq
 C|c_z|^{-\gamma}\Xi_{s'}(m)
                            (1+|\lambda_z(m)|)^{-\delta/4}.
 \label{ju-inner-bound}
\end{align}
\end{proposition}

\begin{proof}
Choose \(\kappa\) as in Lemma~\ref{ju-levi-factor}, set
\(v'=\kappa^{-1}v\), and put \(w=\kappa^{-1}mf_z\).
This change of variables preserves Lebesgue measure.
By \eqref{ju-levi-norms}, the factor independent of \(v',t\)
is \(\prod_{j=0}^{r-1}\Delta_j(m)^{-s_{j+1}}\).
Lemma~\ref{ju-coordinate-formula} gives the remaining
coordinate kernel of Theorem~\ref{jk-oscillatory}, with
\(P=r\), \(m\) there equal to \(k-1\).
In those coordinates,
\[
 h_{r-1}(v',w)=
 \sum_j a_j h_{r-1}(\phi'_j,w)
 +\sum_j b_j h_{r-1}(\phi_j,w)
 +\sum_i z_i h_{r-1}(u_i,w).
\]
Thus the coefficient denoted by \(\eta'_0\) in that theorem
is \(h_{r-1}(\phi_0,w)\).
Since \(m\phi_0=\Delta_0(m)\kappa\phi_0\), invariance of the
Hermitian form gives
\[
 \Delta_0(m)h_{r-1}(\phi_0,w)
            =h_{r-1}(m\phi_0,mf_z)=c_z .
\]
This coefficient is nonzero, and its negative \(\gamma\)-power
is \(|c_z|^{-\gamma}\Delta_0(m)^\gamma\).
Apply Theorem~\ref{jk-oscillatory} with \(c_0=\lambda_z(m)\)
and use \eqref{ju-lambda}.  Combining the powers of
\(\Delta_j(m)\) gives exactly \(\Xi_{s'}(m)\).
No measurable choice of \(\kappa\) as a function of \(m\)
is needed: the argument gives a bound for each fixed \(m\)
by the displayed continuous function of \(m\).
\end{proof}

\begin{theorem}\label{ju-integrability}
For every \(r\geq0\) and \(\delta>0\), there is
\(\epsilon_r(\delta)>0\) such that
\begin{equation}\label{ju-final}
 \int_{N_r}\Xi_s(n)(1+|\lambda_z(n)|)^{-\delta}\,dn<\infty
\end{equation}
whenever \(0<s_l^0-s_l\leq\epsilon_r(\delta)\) for all \(l\).
The same exponent window works for every negative norm vector
\(z\in V_0\).  The integral is finite for every Haar measure on
\(N_r\).
\end{theorem}

\begin{proof}
It suffices to prove the assertion for \(0<\delta\leq1/2\),
because increasing \(\delta\) decreases the integrand.
Take \(\epsilon_0(\delta)=1/10\); the base group is a singleton.
Suppose the result is known at rank \(r-1\) for exponent
\(\delta/4\), and put
\[
 E=\epsilon_{r-1}(\delta/4),\qquad
 \gamma=\min(\delta/8,E/2),\qquad
 \epsilon_r(\delta)=\min\bigl(E/2,\gamma/(4r+4)\bigr).
\]
If \(0<e_l\leq\epsilon_r(\delta)\), then
\[
 e_0+2\sum_{l=1}^{r}e_l
 \leq(2r+1)\frac{\gamma}{4r+4}<\gamma/2.
\]
The lower exponents in Proposition~\ref{ju-inner} have losses
\[
 s_j^0-s'_j=e_{j+1}+\gamma\mathbf1_{\{j=0\}},
\]
where \(s_j^0\) now refers to rank \(r-1\).
These losses are positive and at most \(E\).
Use the Haar coordinates of Lemma~\ref{ju-heisenberg},
Tonelli, and \eqref{ju-inner-bound}.  The induction hypothesis
makes the resulting integral over \(N_{r-1}\) finite.
The constants involving \(z\) are finite since \(c_z\ne0\);
the choice of \(\epsilon_r(\delta)\) does not depend on \(z\).
Finally every Haar measure is a positive scalar multiple of
the one constructed in Lemma~\ref{ju-heisenberg}, by uniqueness
of Haar measure, so finiteness is independent of normalization.
\end{proof}

\subsection{The exact scope of the conclusion}

Theorem~\ref{ju-integrability} proves the source's assertion
\(\mathrm{BAt}(k,r,\delta)\); its proof also establishes
\(\mathrm{FormulaF}\) and the compact orthogonalization assertion
\(\mathrm{IW}\).

These are the actual assertions in the source archive
\cite{alpo-uniform-source}; the source trace identifies their
definitions precisely.
The proof uses the actual flag, character, Hermitian form,
and critical exponents.
The Gram determinant calculation provides a short direct
replacement for the coordinate computations in
\texttt{DelFormula.lean}.
The induction retains precisely the exponent change
\(\delta\mapsto\delta/4\) and the subtraction of \(\gamma\)
from the first lower exponent.

The group integral defining \(\mathrm{AAt}\), the separated
matrix coefficient estimate, and the final dominated
unfolding argument for the Jacquet pairing are additional
steps.  They are not asserted by \eqref{ju-final} alone.


\section{The full group integral and the Jacquet majorant}
\label{jm-section}

We retain the Hermitian spaces and flags of
Section~\ref{ju-section}.  We prove the group integral required
for the matrix coefficient estimate, and then construct the
majorant with its two required properties.
The proof includes the rank zero case, the Haar density on the
open cell, and the integration in each Heisenberg variable.

\subsection{Singular minors and compact changes of flag}

Write \(M(g)=\|g\|_{\mathrm F}\) for the Frobenius norm and set
\[
 D_l(g)=\det\bigl(h_r(g\phi'_b,\phi'_a)\bigr)_{0\leq a,b\leq l},
 \qquad
 \mathcal F_s(g)=\prod_{l=0}^{r}|D_l(g)|^{-s_l}.
\]
We use positive exponents \(s_l\), assigning value \(+\infty\)
when a minor vanishes.  Define also
\[
 \widehat\Delta_l(g)=
       \|g\phi'_0\wedge\cdots\wedge g\phi'_l\|,
 \qquad
 \widehat\Xi_s(g)=\prod_l\widehat\Delta_l(g)^{-s_l}.
\]
These last quantities are finite and positive on \(G_r\).
For a fixed Haar measure \(dg\), let
\[
 J_r(A,s)=\int_{G_r}M(g)^{-A}\mathcal F_s(g)\,dg,\qquad
 B(k,r)=2(k+r)(r+1)+1.
\]
All assertions about these integrals are independent of the
normalization of Haar measure.

The groups \(G_r\) are unimodular.  One way to see this directly
is to use the invariant bilinear form
\(\operatorname{Re}\operatorname{tr}(XY)\) on their real Lie
algebras.  It is nondegenerate: the complexification of the
Lie algebra is the full matrix algebra, on which the trace
form is nondegenerate, and \(\operatorname{tr}(XY)\) is real
for the two elements of this real form.
Preservation of this form gives
\(|\det_{\mathbb R}\operatorname{Ad}g|=1\).
A left invariant volume form is therefore right invariant
as well.

\begin{lemma}\label{jm-compact-factor}
For every nonnegative measurable function \(\Phi\) invariant
under left and right multiplication by \(K_r\), and every
\(x,y\in G_r\),
\begin{equation}\label{jm-both}
 \int_{G_r}\Phi(g)\mathcal F_s(x^{-1}gy)\,dg
  =\widehat\Xi_s(x)\widehat\Xi_s(y)
                        \int_{G_r}\Phi(g)\mathcal F_s(g)\,dg .
\end{equation}
If only left \(K_r\)-invariance is assumed, the corresponding
identity with \(y=1\) still holds.
\end{lemma}

\begin{proof}
Apply Lemma~\ref{ju-isotropic-qr} to the primed frame
\(x\phi'_j\).  Thus
\(x\phi'_j=\sum_{i\leq j}\tau_{ij}\kappa\phi'_i\),
with \(\kappa\in K_r\).
Invariance of \(h_r\) gives
\[
 h_r(x^{-1}g\phi'_b,\phi'_a)
                         =h_r(g\phi'_b,x\phi'_a).
\]
Taking the first \(l+1\) rows of the triangular change
therefore multiplies the absolute value of the determinant
by \(\prod_{j\leq l}\tau_{jj}=\widehat\Delta_l(x)\).
Consequently
\[
 \mathcal F_s(x^{-1}g)
          =\widehat\Xi_s(x)\mathcal F_s(\kappa^{-1}g).
\]
Change variables by left multiplication by \(\kappa\).
This proves the assertion for \(y=1\), including when the
integral is infinite.
The same triangular calculation on the columns of \(y\phi'_j\)
gives \(\mathcal F_s(gy)=\widehat\Xi_s(y)\mathcal F_s(g\kappa')\)
for some \(\kappa'\in K_r\).
Right invariance of Haar measure and of \(\Phi\) proves
\eqref{jm-both}.
\end{proof}

\subsection{The base case}

\begin{lemma}\label{jm-ball}
On \(G_0=U(k,1)\),
\[
                       \int_{G_0}M(g)^{-A}\,dg<\infty
                              \qquad\text{if }A>2k.
\]
Every \(g\in G_0\) is a product of two elements of \(K_0\)
and a matrix \(a_t\), \(t\geq0\), which is the identity on
\(\operatorname{span}(u_2,\ldots,u_k)\) and has matrix
\[
 \begin{pmatrix}\cosh t&\sinh t\\
                 \sinh t&\cosh t\end{pmatrix}
\]
on \(\operatorname{span}(u_1,u_-)\).
\end{lemma}

\begin{proof}
The quotient \(G_0/K_0\) is the unit ball
\(\mathbb B^k\subset\mathbb C^k\).
Indeed a block matrix
\(g=\left(\begin{smallmatrix}A&b\\c^*&d\end{smallmatrix}\right)\)
acts by \(w\mapsto(Aw+b)/(c^*w+d)\).
It preserves the ball, and the stabilizer of zero is exactly
\(K_0=U(k)\times U(1)\).
For \(w=\rho u\), \(\|u\|=1\), \(0\leq\rho<1\), the matrix
\[
 b_w=
 \begin{pmatrix}
 I+(a-1)uu^*&a\rho u\\
 a\rho u^*&a
 \end{pmatrix},
 \qquad a=(1-\rho^2)^{-1/2},
\]
preserves the form and sends zero to \(w\).
At zero use \(b_0=I\).
Rotating \(u\) by \(U(k)\) shows that \(b_w\in K_0a_tK_0\)
with \(\tanh t=\rho\), proving the second assertion.

For completeness, the invariant measure on the ball is,
up to a positive constant,
\[
                       (1-\|w\|^2)^{-k-1}\,dw.
\]
The form identity gives
\(1-\|gw\|^2=(1-\|w\|^2)/|c^*w+d|^2\).
Differentiating the fractional linear map and taking a
block determinant gives its complex Jacobian
\(\det g/(c^*w+d)^{k+1}\).
Since \(|\det g|=1\), its real Jacobian is the squared
absolute value, which proves invariance of the displayed
density.
Disintegration of Haar measure over the compact stabilizer
therefore gives this quotient measure, up to normalization.
Finally
\[
 M(b_w)^2=k-1+\frac{2(1+\rho^2)}{1-\rho^2}
                         \geq\frac2{1-\rho^2}.
\]
The required integral is bounded by a constant times
\[
 \int_0^1 \rho^{2k-1}
               (1-\rho^2)^{A/2-k-1}\,d\rho,
\]
which is finite for \(A>2k\).
\end{proof}

\begin{lemma}\label{jm-sphere}
Let \(d\omega\) be probability measure on the unit sphere in
\(\mathbb C^k\).  For \(0<s<k\),
\begin{equation}\label{jm-angular}
 \int |1-\zeta_1|^{-s}\,d\omega(\zeta)
       =\frac{\Gamma(k)\Gamma(k-s)}
                    {\Gamma(k-s/2)^2}<\infty.
\end{equation}
For \(0\leq\rho<1\),
\begin{equation}\label{jm-poisson}
 \int |1-\rho\zeta_1|^{-2k}\,d\omega(\zeta)
                                      =(1-\rho^2)^{-k}.
\end{equation}
\end{lemma}

\begin{proof}
Rotational invariance and a complex Gaussian integral give
\[
 \int\zeta_1^a\overline{\zeta_1}^{\,b}\,d\omega=0
       \quad(a\ne b),\qquad
 \int|\zeta_1|^{2j}\,d\omega
                         =\frac{j!}{(k)_j}.
\]
Here \((v)_j=v(v+1)\cdots(v+j-1)\).
To check the second formula, integrate
\(|z_1|^{2j}e^{-\|z\|^2}\) over \(\mathbb C^k\) first by
coordinates and then in polar coordinates.
The coordinate value is \(\pi^k j!\), while comparison with
the case \(j=0\) gives the stated quotient of gamma factors.
This also covers \(k=1\).

The binomial series, absolutely convergent for \(\rho<1\),
now gives, for \(a>0\),
\[
 \int|1-\rho\zeta_1|^{-2a}\,d\omega
                  =\sum_{j=0}^\infty
                    \frac{(a)_j^2}{j!(k)_j}\rho^{2j}.
\]
Taking \(a=k\) proves \eqref{jm-poisson}.
For \(a=s/2\), the sum at \(\rho=1\) is evaluated by
inserting the beta integral
\[
 \frac{(a)_j}{(k)_j}
 =\frac{\Gamma(k)}{\Gamma(a)\Gamma(k-a)}
               \int_0^1 u^{a+j-1}(1-u)^{k-a-1}\,du.
\]
Summing the remaining binomial series inside this nonnegative
integral gives the beta integral with parameters \(a,k-2a\),
and hence the right side of \eqref{jm-angular}.
The beta integral itself follows from
\(\Gamma(\alpha)\Gamma(\beta)\) by the change
\((x,y)=(vu,v(1-u))\).
Its parameters here are positive because \(0<s<k\).
Fatou's lemma, as \(\rho\uparrow1\), first proves finiteness
of the integral in \eqref{jm-angular}.
For \(1/2\leq\rho<1\) and \(|z|\leq1\),
\[
 |1-\rho z|^2-\rho|1-z|^2
                         =(1-\rho)(1-\rho|z|^2)\geq0.
\]
Thus the integrands are dominated by
\(2^{s/2}|1-\zeta_1|^{-s}\).
Dominated convergence proves the equality.
\end{proof}

\begin{proposition}\label{jm-base}
For \(0<s<k\) and \(A>2k\), the integral \(J_0(A,s)\) is finite.
\end{proposition}

\begin{proof}
On this group the single primed vector is
\(\phi'_0=(u_1-u_-)/\sqrt2\).
Its orbit under probability Haar measure on \(K_0\) has
the distribution \((\zeta,-v)/\sqrt2\), where \(\zeta\)
is uniform on the unit sphere and \(v\) is uniform on the
unit circle, independently.
It follows from Lemma~\ref{jm-sphere} that
\[
 C_s:=\int_{K_0}|h_0(\kappa\phi'_0,\phi'_0)|^{-s}\,d\kappa
             =2^s\frac{\Gamma(k)\Gamma(k-s)}
                          {\Gamma(k-s/2)^2}<\infty .
\]
The compact group is transitive on the Euclidean unit
isotropic vectors, by Lemma~\ref{ju-isotropic-qr}.
For every nonzero isotropic \(y\), therefore,
\[
 \int_{K_0}|h_0(\kappa y,\phi'_0)|^{-s}\,d\kappa
                                      =C_s\|y\|^{-s}.
\]
For \(a_t\) of Lemma~\ref{jm-ball},
\(\|a_t\kappa\phi'_0\|\) has the distribution
\(|\cosh t-\sinh t\,\zeta_1|\).
Equation \eqref{jm-poisson}, with \(\rho=\tanh t\), gives
\[
             \int_{K_0}\|a_t\kappa\phi'_0\|^{-2k}\,d\kappa=1.
\]
Jensen's inequality for the concave power \(s/(2k)\in(0,1)\)
implies
\(\int_{K_0}\|a_t\kappa\phi'_0\|^{-s}\,d\kappa\leq1\).
The decomposition \(G_0=K_0a_tK_0\) then shows that
\[
 \int_{K_0\times K_0}\mathcal F_s(\kappa g\kappa')\,
                          d\kappa\,d\kappa'\leq C_s
                                     \qquad(g\in G_0).
\]
Average the definition of \(J_0\) on both sides by \(K_0\).
Tonelli and the invariance of Haar measure and \(M\) give
\[
 J_0(A,s)\leq C_s\int_{G_0}M(g)^{-A}\,dg<\infty
\]
by Lemma~\ref{jm-ball}.
\end{proof}

\subsection{Haar measure on the open cell}

For \(r\geq1\), set \(W=V_{r-1}\), \(d=\dim_{\mathbb C}W
=2r+k-1\), and write
\[
 w(c)=
 \begin{pmatrix}0&0&\bar c^{-1}\\0&I&0\\c&0&0\end{pmatrix},
                              \qquad c\in\mathbb C^\times.
\]
The matrices \(n(v,t)\) and \(\iota(m)\) are those of
Lemma~\ref{ju-heisenberg}.

\begin{lemma}\label{jm-cell-haar}
The map
\[
 (n_L,c,m,n_R)\longmapsto n_Lw(c)\iota(m)n_R
\]
is a smooth bijection, with smooth inverse, onto
\(\{g\in G_r:g_{f_0e_0}\ne0\}\).
Its complement has Haar measure zero.
For a positive constant depending only on the Haar
normalizations, the measure in these coordinates is
\begin{equation}\label{jm-haar}
                  C\,|c|^{2d}\,dn_L\,dc\,dm\,dn_R .
\end{equation}
Here \(dn_L,dn_R\) are the Lebesgue Haar measures on
\(W\times\mathbb R\), and \(dc\) is planar Lebesgue measure.
\end{lemma}

\begin{proof}
The coordinates follow by solving the block product.
If \(c=g_{f_0e_0}\ne0\), the left vector is
\(v_L=g_{We_0}/c\), the right vector is determined by
\(v_R^*J=-g_{f_0W}/c\), and the two upper right parameters
are \(g_{e_0e_0}/c\) and \(g_{f_0f_0}/c\).
Their imaginary parts determine the two real parameters;
the real parts follow from isotropy.
Removing these factors leaves \(w(c)\iota(m)\), with
\(m=g_{WW}+v_Lc\,v_R^*J\).
These formulas prove the assertion about coordinates.

We give the density calculation.  Translation invariance
reduces it to \(n_L=n_R=1\), \(m=1\).
Conjugation by a fixed \(m\in G_{r-1}\) preserves the
Lebesgue measure on \(W\times\mathbb R\), as in
Lemma~\ref{ju-heisenberg}, so this reduction introduces no
extra factor.
For
\[
 X(v,t)=\begin{pmatrix}0&-v^*J&it\\0&0&v\\0&0&0\end{pmatrix},
\]
the contribution of the left variable to \(g^{-1}dg\) is
\[
 w(c)^{-1}X(v,t)w(c)=
 \begin{pmatrix}
 0&0&0\\cv&0&0\\i|c|^2t&-\bar c v^*J&0
 \end{pmatrix}.
\]
Thus its real Jacobian is \(|c|^{2d+2}\).
The \(c\)-variable contributes the diagonal matrix with entries
\(dc/c\), \(0\), and \(-d\bar c/\bar c\), whose real Jacobian
relative to planar \(dc\) is \(|c|^{-2}\).
The middle Lie algebra and the right Heisenberg Lie algebra
complete these to a direct sum of the real Lie algebra of
\(G_r\), and their Jacobians are one.
The product is \(|c|^{2d}\), proving \eqref{jm-haar}.

Finally \(G_r\) is connected.
For its polar decomposition \(g=\kappa P\), the identity
\(H_r(g^*g)H_r=(g^*g)^{-1}\) implies
\(H_rPH_r=P^{-1}\).  Thus \(P^t\in G_r\) for \(0\leq t\leq1\)
and \(\kappa\in K_r\); the latter group is
\(U(k+r)\times U(r+1)\), which is connected.
The real analytic function \(|g_{f_0e_0}|^2\) is not
identically zero, since it is nonzero at \(w(1)\).
Its zero set has measure zero in analytic local coordinates.
To recall the elementary justification, at each zero some
finite directional derivative is nonzero, unless all
derivatives vanish and analytic continuation makes the
function identically zero.  In a sufficiently small box this
derivative stays nonzero; on each line in that direction
Rolle's theorem bounds the number of zeros by its order.
Fubini and a countable covering give measure zero.
Haar measure has a smooth positive density in these coordinates,
which proves the assertion.
\end{proof}

\subsection{Integration of one Heisenberg factor}

\begin{lemma}\label{jm-e-int}
Put \(D=2d+1\).  If \(A>D\), then
\begin{equation}\label{jm-e-bound}
 \int_{W\times\mathbb R}
        M(w(c)\iota(m)n(v,t))^{-A}\,dv\,dt
 \leq C_{A,d}\min(|c|^{-D},|c|^{-1})M(m)^{D-A}
\end{equation}
for all \(c\ne0\) and \(m\in G_{r-1}\).
\end{lemma}

\begin{proof}
Writing \(b=it-h_{r-1}(v,v)/2\), direct block multiplication gives
\[
 M(w(c)\iota(m)n(v,t))^2
 =|c|^{-2}+M(m)^2+\|mv\|^2
                       +|c|^2(1+\|v\|^2+|b|^2).
\]
In particular this is at least
\(M(m)^2+|c|^2(\|v\|^2+t^2)\).
Scale \(v,t\) by \(|c|\), then by \(M(m)\).
The resulting bound is
\[
 |c|^{-D}M(m)^{D-A}
                    \int_{\mathbb R^D}(1+\|u\|^2)^{-A/2}\,du .
\]
The same norm is at least
\(M(m)^2+\|mv\|^2+|c|^2t^2\).
Now substitute \(v'=mv\), whose real Jacobian is
\(|\det m|^2=1\), and \(t'=|c|t\).
This gives the identical bound with \(|c|^{-1}\) in place
of \(|c|^{-D}\).
The remaining Euclidean integral is finite because \(A>D\);
polar coordinates, or the Gaussian calculation in
Theorem~\ref{jk-exact}, prove this directly.
Taking the smaller of the two bounds gives the result.
\end{proof}

\begin{lemma}\label{jm-left-factor}
Fix \(n_L=n(v,t)\), and set
\[
 Y_0=n_Lf_0,\qquad Y_{j+1}=n_L\phi'_j\quad(0\leq j<r),
 \qquad \delta_l=\|Y_0\wedge\cdots\wedge Y_l\|,
 \qquad Q=\delta_0>0.
\]
For every \(c\ne0\), there exist
\(\kappa\in K_r\), \(u\in G_{r-1}\), and a Heisenberg
element \(n_0\) such that
\begin{equation}\label{jm-left-decomposition}
 n_Lw(c)=\kappa w(cQ)\iota(u)n_0,\qquad
               \widehat\Delta_j(u)=\delta_{j+1}/Q.
\end{equation}
\end{lemma}

\begin{proof}
The vector \(Y_0\) is nonzero and isotropic.
Compact transitivity supplies \(\kappa f_0=Y_0/Q\).
Consequently
\((\kappa w(cQ))^{-1}n_Lw(c)\) fixes \(e_0\).
The block argument in Lemma~\ref{ju-heisenberg}, with
\(m\in G_{r-1}\) instead of \(N_{r-1}\), identifies its
stabilizer as the products \(\iota(u)n_0\).
This proves the factorization.
Apply it to the primed flag, whose first vector is \(e_0\).
The rightmost \(n_0\) changes its columns by an upper triangular
matrix with diagonal one.
The norm of the first \(j+2\) transformed columns on the left
is \(|c|\delta_{j+1}\).
On the right it is
\(|c|Q\,\widehat\Delta_j(u)\), since \(f_0\) is Euclidean
orthogonal to \(W\) and \(\kappa\) is Euclidean unitary.
Cancellation gives the claimed norm identities.
\end{proof}

\begin{theorem}\label{jm-group-integral}
For every \(k\geq1\), \(r\geq0\), and
\[
                0<s_l^0-s_l\leq1/10\quad(0\leq l\leq r),
\]
one has
\begin{equation}\label{jm-AAt}
                       J_r(A,s)<\infty
                                   \qquad(A\geq B(k,r)).
\end{equation}
\end{theorem}

\begin{proof}
The case \(r=0\) follows from Proposition~\ref{jm-base},
since \(0<s_0<k\) and \(B(k,0)=2k+1\).
Suppose \(r\geq1\), and use induction at rank \(r-1\).
Write
\[
 e_l=s_l^0-s_l,\quad
 \sigma=\sum_{l=0}^{r}s_l,\quad
 a=2d-\sigma,\quad D=2d+1,\quad s^-=(s_1,\ldots,s_r).
\]
Since
\(B(k,r)-B(k,r-1)=2k+4r=D+1\), one has
\(A>D\) and \(A-D\geq B(k,r-1)\).
The induction hypothesis gives
\(J_{r-1}(A-D,s^-)<\infty\).

Both left and right multiplication by a Heisenberg element
preserve all the minors \(D_l\).
Indeed its action on the primed flag is upper triangular
with diagonal one, and the left action gives the same row
operations by Hermitian invariance.
The block matrix \(w(c)\iota(m)\) then gives
\[
 D_0=c,\qquad D_l=c\,D_{l-1}(m)\quad(l\geq1).
\]
Thus in the coordinates of Lemma~\ref{jm-cell-haar},
\[
 \mathcal F_s(n_Lw(c)\iota(m)n_R)
                  =|c|^{-\sigma}\mathcal F_{s^-}(m).
\]
For fixed \(n_L,c\), use \eqref{jm-left-decomposition}.
The change
\(n_0\iota(m)n_R=\iota(m)(\iota(m)^{-1}n_0\iota(m))n_R\)
is a translation in the right Heisenberg integral, and
introduces no factor.  Lemma~\ref{jm-e-int} therefore gives
\[
 \int_{n_R}M(n_Lw(c)\iota(m)n_R)^{-A}\,dn_R
 \leq C_{A,d}\min(|cQ|^{-D},|cQ|^{-1})M(um)^{D-A}.
\]
Integration over \(m\), followed by \(g=um\), is evaluated
using the left version of Lemma~\ref{jm-compact-factor}:
\begin{align*}
 \int_{G_{r-1}}\mathcal F_{s^-}(m)M(um)^{-(A-D)}\,dm
       &=\widehat\Xi_{s^-}(u)J_{r-1}(A-D,s^-),\\
 \widehat\Xi_{s^-}(u)
       &=Q^{\sum_{l=1}^{r}s_l}
                             \prod_{l=1}^{r}\delta_l^{-s_l}.
\end{align*}
The second expression depends on \(n_L\), and not on the
particular factorization or on \(c\).

The integral in \(c\) has an elementary exact value.
For \(1<\sigma<2d+1\), polar coordinates give
\[
 \int_{\mathbb C}|c|^{2d-\sigma}
             \min(|c|^{-2d-1},|c|^{-1})\,dc
          =2\pi\left(\frac{1}{2d-\sigma+1}
                                  +\frac{1}{\sigma-1}\right).
\]
The hypotheses ensure these inequalities:
\(\sigma\geq2r+k-(r+1)/10>1\), and
\(\sigma<2r+k<4r+2k-1=2d+1\).
The substitution \(c'=Qc\) consequently extracts the
factor \(Q^{-a-2}\).

It remains to integrate
\[
 \delta_0^{-a-2+\sum_{l=1}^{r}s_l}
                                  \prod_{l=1}^{r}\delta_l^{-s_l}
\]
over \(n_L\).
Lemma~\ref{ju-coordinate-formula} applies also to the
reference frame \(\phi'_j\) on \(W\), with its paired frame
\(\phi_j\): its proof uses only the displayed orthogonality
and pairing identities.  Hence this is the coordinate integral
\(I_{r,k-1}(\nu)\), where
\[
           \nu_0=e_0+2\sum_{j=1}^{r}e_j,\qquad
           \nu_l=-e_l\quad(1\leq l\leq r).
\]
For its partial sums,
\[
 S_l=e_0+\sum_{j=1}^{l}e_j+2\sum_{j=l+1}^{r}e_j>0,\qquad
 S_l+2\sum_{j>l}\nu_j=e_0+\sum_{j=1}^{l}e_j>0.
\]
Every inequality in Theorem~\ref{jk-exact} holds.
The remaining integral is therefore finite, with the exact
gamma-product value given there.
All substitutions and integrations involved nonnegative
functions; Tonelli justifies them before finiteness is known.
This completes the induction.
\end{proof}

\subsection{The separated bound for matrix coefficients}

\begin{proposition}\label{jm-coefficient}
Under the hypotheses of Theorem~\ref{jm-group-integral},
there is a finite constant \(C\) such that, for any functions
\(F,F'\) into a normed space that are strongly measurable
almost everywhere for Haar measure and satisfy
\[
 \|F(g)\|\leq cM(g)^{-A},\qquad
 \|F'(g)\|\leq c'M(g)^{-A}\qquad(g\in G_r),
\]
where \(c,c'\geq0\), one has
\begin{equation}\label{jm-coefficient-bound}
 \int_{G_r}\|F(x^{-1}gy)\|\,\|F'(g)\|\,dg
                         \leq Ccc'\Xi_s(x)\Xi_s(y)
                                  \qquad(x,y\in G_r).
\end{equation}
One may take \(C=J_r(A,s)\) with the present norm conventions.
\end{proposition}

\begin{proof}
Left and right translations preserve Haar null sets, so the
integrand is measurable for the completed Haar measure.
The two reference exterior vectors defining \(D_l\) have
Euclidean norm one.  Cauchy--Schwarz on the exterior power
therefore gives \(|D_l(g)|\leq M(g)^{l+1}\).
Also \(M(g)\geq1\), since \(|\det g|=1\).
As
\[
 \sum_{l=0}^{r}(l+1)s_l
   <\sum_{l=0}^{r}(l+1)s_l^0
       =(r+1)(k+r)<B(k,r)\leq A,
\]
it follows that
\begin{equation}\label{jm-determinant-majorization}
                         M(g)^{-A}\leq\mathcal F_s(g).
\end{equation}
Let \(H=H_r\).  Conjugation by \(H\) preserves both Haar
measure and \(M\).  Moreover
\[
                  \widehat\Xi_s(HxH)=\Xi_s(x),
\]
because \(H\phi'_j=\phi_j\), and left multiplication by
\(H\) is Euclidean unitary.
Apply \eqref{jm-determinant-majorization} to
\(H(x^{-1}gy)H\), multiply by \(cc'M(g)^{-A}\), and integrate.
After the substitution \(g'=HgH\), formula \eqref{jm-both}
with \(x'=HxH\), \(y'=HyH\), and \(\Phi(g')=M(g')^{-A}\)
gives exactly \eqref{jm-coefficient-bound}.
Theorem~\ref{jm-group-integral} proves that its constant is finite.
\end{proof}

\subsection{The compact stabilizers and the required majorant}

For the negative vector \(z\in V_0\), let \(K_z\) be the
subgroup of \(G_r\) fixing \(z\) and each \(e_j,f_j\).
It is compact: on \(V_0\) it acts trivially on the negative
line \(\mathbb Cz\) and unitarily on the positive definite
space \(z^\perp\), so is isomorphic to \(U(k)\).
It normalizes \(N_r\), preserves its Haar measure under
conjugation, and satisfies
\(\lambda_z(h^{-1}nh)=\lambda_z(n)\).
The flag definition proves normalization.
The formula for \(\ell_z\), using the fixed vectors \(e_j,z\),
proves character invariance.
For Haar measure, use the recursion of
Lemma~\ref{ju-heisenberg}: conjugation acts on each
Heisenberg parameter by \(v\mapsto hv\), \(t\mapsto t\),
with real determinant \(|\det h|^2=1\).
Induction gives measure preservation.

\begin{lemma}\label{jm-compact-stabilizer}
Suppose \(s_l\geq0\) and
\(\int_{N_r}\Xi_s(n)(1+|\lambda_z(n)|)^{-1/2}\,dn<\infty\).
Then, for Haar measure \(dh\) on \(K_z\),
\begin{equation}\label{jm-stabilizer-integral}
 \int_{K_z\times N_r}
       \Xi_s(nh)(1+|\lambda_z(n)|)^{-2}\,dh\,dn<\infty.
\end{equation}
\end{lemma}

\begin{proof}
Compactness gives a uniform bound \(L\geq1\) on
\(\|h^{-1}\|_{\mathrm{op}}\), \(h\in K_z\).
The exterior norm inequality
\(\|\bigwedge^{l+1}h\,v\|\geq L^{-l-1}\|v\|\)
implies
\[
          \Xi_s(hg)\leq L^{\sum_l(l+1)s_l}\Xi_s(g).
\]
For fixed \(h\), write \(nh=h(h^{-1}nh)\), and change
variables \(n'=h^{-1}nh\).  The measure and character are
unchanged.  Since
\((1+|u|)^{-2}\leq(1+|u|)^{-1/2}\), the integral over \(n\)
is at most the assumed finite integral times the displayed
constant, independently of \(h\).
Haar measure on the compact group \(K_z\) is finite.
Tonelli proves \eqref{jm-stabilizer-integral}.
\end{proof}

\begin{theorem}\label{jm-majorant}
Let \(k\geq1\), \(r\geq0\), \(A\geq B(k,r)\), and let
\(z,z'\in V_0\) have negative Hermitian norm.
There exists a positive continuous function \(\Xi\) on \(G_r\)
with the following two properties:
\begin{enumerate}
 \item For all \(F,F'\) satisfying the decay bounds in
 Proposition~\ref{jm-coefficient}, the integral in
 \eqref{jm-coefficient-bound} is at most
 \(Ccc'\Xi(x)\Xi(y)\).
 \item Both integrals
 \[
 \int_{K_z\times N_r}\Xi(nh)(1+|\lambda_z(n)|)^{-2}\,dh\,dn,
 \qquad
 \int_{K_{z'}\times N_r}
             \Xi(nh)(1+|\lambda_{z'}(n)|)^{-2}\,dh\,dn
 \]
 are finite.
\end{enumerate}
One may take
\[
 \Xi=\Xi_s,\qquad
 s_l=s_l^0-\epsilon,\qquad
 0<\epsilon\leq\min\bigl(1/10,\epsilon_r(1/2)\bigr).
\]
\end{theorem}

\begin{proof}
The exterior Gram determinants are continuous and positive
on the group, so \(\Xi_s\) is positive and continuous.
Its extension by zero outside the closed matrix group \(G_r\)
is Borel measurable on the ambient matrix space; all the Haar
measures in these assertions are supported on \(G_r\).
Theorem~\ref{ju-integrability} gives the two unipotent
integrals with exponent \(1/2\), using the same \(\epsilon\)
for \(z\) and \(z'\).
Theorem~\ref{jm-group-integral} and
Proposition~\ref{jm-coefficient} give the first property.
Lemma~\ref{jm-compact-stabilizer} gives the second.
\end{proof}

Theorem~\ref{jm-group-integral} is the actual assertion
\(\mathrm{AAt}\) in the source, and
Theorem~\ref{jm-majorant} proves the actual existence
assertion \(\mathrm{XiExists}\), including both
\(\mathrm{CoefBound}\) and \(\mathrm{XiInt}\).
The norm is exactly the Frobenius norm called \(gN\), and
the exponent threshold \(B(k,r)\) is exactly \(AB\).
The source uses \(n_L^{-1}\) in its open cell; the change
\((v,t)\mapsto(-v,-t)\) preserves the Lebesgue measure and
identifies that convention with the one used here.
These statements and conventions are in
\texttt{StatementB.lean} and the
\texttt{TailIntegrals} files recorded in the source trace
\cite{alpo-uniform-source}.

The construction of the majorant is now complete.
Its application to differentiation, regularized pairings,
and the final Jacquet adjointness identity still requires
the subsequent dominated integration and unfolding arguments.


\section{The integration formula and the dilation measure}
\label{jg-section}

The normalization of Haar measure in the Jacquet pairing matters.
We now derive the actual integration formula, including the terminal
compact quotient, and prove the integrability of the weight used in
the differentiability estimates.  No formula for integration over
double cosets is assumed.

\subsection{The compact quotient at rank zero}

On \(V_0=\mathbb C^k\oplus\mathbb C e_-\), use the form
\[
 h_0((u,w),(u',w'))=\langle u,u'\rangle-w\bar w',
 \qquad k\geq1,
\]
linear in the first variable.  Write \(G_0=U(h_0)\), and let
\(K_-= \{\operatorname{diag}(u,1):u\in U(k)\}\).
The orbit of \(e_-\) is the hyperboloid
\[
 X=\{(u,w):|w|^2-\|u\|^2=1\}.
\]
Transitivity follows by completing a negative unit vector to an
orthonormal basis for the Hermitian form.  Its stabilizer is \(K_-\).

\begin{lemma}\label{jg-hyperboloid}
Up to a positive constant, the invariant measure on \(X\), in the
coordinates
\[
 (u,\theta)\longmapsto
       (u,\sqrt{1+\|u\|^2}\,e^{i\theta}),
 \qquad u\in\mathbb C^k,\quad \theta\in\mathbb R/(2\pi\mathbb Z),
\]
is \(du\,d\theta\).  In particular, if \(dk\) is probability Haar
measure on \(K_-\), then for every nonnegative measurable \(f\) on
\(G_0\),
\begin{equation}\label{jg-base-integration}
 \int_{G_0}f(g)\,dg
 =C\int_{|w|>1}(|w|^2-1)^{k-1}
     \int_{K_-\times K_-}
       f(k_1^{-1}b(w)k_2)\,dk_1\,dk_2\,dw ,
\end{equation}
where \(C>0\) and \(b(w)\) acts as the identity on
\(\operatorname{span}(e_1,e_-)^\perp\) and on
\(\operatorname{span}(e_1,e_-)\) has matrix
\[
 b(w)=
 \begin{pmatrix}\bar w&\sqrt{|w|^2-1}\\
                 \sqrt{|w|^2-1}&w\end{pmatrix}.
\]
\end{lemma}

\begin{proof}
Put \(Q(u,w)=|w|^2-\|u\|^2\).
Every element of \(G_0\) preserves \(Q\) and real Lebesgue measure:
its complex determinant has absolute value one.
Consequently the measure on \(Q=1\) obtained by restricting
Lebesgue measure to \(1-\eta<Q<1+\eta\), dividing by \(2\eta\),
and letting \(\eta\downarrow0\), is invariant.
For compactly supported continuous test functions this limit exists
by the ordinary change of variables
\[
       (u,q,\theta)\longmapsto
               (u,\sqrt{q+\|u\|^2}\,e^{i\theta}).
\]
Planar polar measure in the last coordinate is
\(\rho\,d\rho\,d\theta=\tfrac12\,dq\,d\theta\).
Thus the resulting measure on \(X\) is
\(\tfrac12\,du\,d\theta\), proving invariance and the first assertion.

For completeness, this quotient measure determines Haar measure on
\(G_0\) without any choice of density.  Choose local smooth sections
of \(G_0\to X\), combine them into a measurable section \(\sigma\),
and set
\[
 \int f\,d\widetilde\mu
     =\int_X\int_{K_-}f(\sigma(x)k)\,dk\,d\nu(x).
\]
Changing \(\sigma(x)\) multiplies it on the right by an element of
\(K_-\), so this definition is independent of the section.
For \(a\in G_0\), the two representatives
\(a\sigma(x)\) and \(\sigma(ax)\) differ by an element of \(K_-\).
Invariance of \(\nu\) and of \(dk\) proves left invariance of
\(\widetilde\mu\).  It is nonzero and finite on compact sets, by
local trivializations and compactness of \(K_-\).  It is therefore
a positive multiple of the fixed Haar measure.
Local sections can also be obtained directly by smooth Hermitian
Gram--Schmidt near any negative unit vector; a countable cover and
disjoint measurable refinements give the stated measurable section.

Now write \(u=r\omega\), where \(\omega\) is a unit vector in
\(\mathbb C^k\).  If \(\rho=|w|=\sqrt{1+r^2}\), then
\[
 r^{2k-1}\,dr=(\rho^2-1)^{k-1}\rho\,d\rho .
\]
The pushforward of probability Haar measure on \(U(k)\) by
\(k\mapsto k e_1\) is normalized spherical measure: it is invariant,
and transitivity and compact averaging give uniqueness.
The matrix \(b(w)\) belongs to \(G_0\) and sends \(e_-\) to
\((\sqrt{|w|^2-1}\,e_1,w)\).
Substitution in the quotient formula gives
\eqref{jg-base-integration}.  The omitted set \(u=0\) has zero
\(du\,d\theta\)-measure, also when \(k=1\).
\end{proof}

Let \(z,z'\in V_0\) be negative, and write
\[
           T=-h_0(z,z)>0,\qquad T'=-h_0(z',z')>0,
           \qquad K_z=\operatorname{Stab}_{G_0}(z).
\]
Choose \(a,a'\in G_0\) with
\(a e_-=z/\sqrt T\) and \(a'e_-=z'/\sqrt{T'}\).
Then
\[
 \mu(\zeta)=a\,b\!\left(-\frac{\zeta}{\sqrt{TT'}}\right)(a')^{-1}
       \qquad (|\zeta|^2>TT')
\]
is a smooth section satisfying \(h_0(\mu(\zeta)z',z)=\zeta\).

\begin{corollary}\label{jg-terminal-measure}
For every such smooth section and every choice of Haar measures,
there is a constant \(C>0\) such that
\begin{equation}\label{jg-terminal}
 \int_{G_0}f(g)\,dg
 =C\int_{|\zeta|^2>TT'}(|\zeta|^2-TT')^{k-1}
      \int_{K_z\times K_{z'}}
        f(k_1^{-1}\mu(\zeta)k_2)\,dk_1\,dk_2\,d\zeta .
\end{equation}
\end{corollary}

\begin{proof}
Apply \eqref{jg-base-integration} to \(a^{-1}g a'\).
Unimodularity gives invariance of Haar measure under these
translations.  The substitution
\(w=-\zeta/\sqrt{TT'}\) contributes the positive factor
\((TT')^{-k}\), which is absorbed in \(C\).
The conjugate compact subgroups are \(K_z\) and \(K_{z'}\).

The formula does not depend on the particular section.
For a fixed \(\zeta\) in the exterior region, \(K_z\) acts
transitively on the vectors \(v\) satisfying
\(h_0(v,v)=-T'\) and \(h_0(v,z)=\zeta\):
after the change by \(a^{-1}\), the last coordinate is fixed
and the first \(k\) coordinates run over a sphere of positive
radius.  Once the image of \(z'\) has been matched, two group
elements differ on the right by an element of \(K_{z'}\).
Thus every fiber is one \(K_z\times K_{z'}\) orbit.
Bi-invariance of the two compact Haar measures proves that its
average is independent of the representative.
\end{proof}

\subsection{Iteration of the open cell}

Embed \(V_0\) as the middle summand of
\[
 V_r=\bigoplus_{a=0}^{r-1}\mathbb C e_a
          \oplus V_0\oplus
             \bigoplus_{a=0}^{r-1}\mathbb C f_a,\qquad
 h_r(e_a,f_b)=\delta_{ab}.
\]
The \(f_a\) are ordered in reverse order in the matrix coordinates.
The vector spaces and Euclidean norms are those in
Section~\ref{ju-section}.  Put
\[
       M_a=2(r-a)+k+1,\qquad
       b(c,m)e_a=c_a f_a,\quad
       b(c,m)f_a=\bar c_a^{-1}e_a,\quad
       b(c,m)|_{V_0}=m .
\]
Here \(c\in(\mathbb C^\times)^r\) and \(m\in G_0\).
Extend \(K_z,K_{z'}\) by the identity on all the outer pairs.
These groups normalize \(N_r\) and preserve its Haar measure,
as proved in Section~\ref{jm-section}.

\begin{proposition}\label{jg-full-cell}
For any choices of Haar measure there is a positive constant
\(C_r\) such that
\begin{equation}\label{jg-iterated}
 \int_{G_r}f(g)\,dg
 =C_r\int_{\mathbb C^r}
      \prod_{a=0}^{r-1}|c_a|^{2(M_a-2)}
      \int_{N_r\times N_r\times G_0}
          f(n_L^{-1}b(c,m)n_R)\,dn_L\,dn_R\,dm\,dc .
\end{equation}
This holds for every nonnegative measurable \(f\).
Values on any \(c_a=0\) can be assigned arbitrarily.
\end{proposition}

\begin{proof}
At the first step Lemma~\ref{jm-cell-haar} has middle dimension
\(d=2r+k-1=M_0-2\) and density \(|c_0|^{2d}\).
Apply the same formula to the middle \(G_{r-1}\) variable.
The outer Weyl matrix is the identity on that middle space,
so it commutes with every embedded middle matrix.
The product of the two Weyl matrices is the corresponding
restriction of \(b(c,m)\).

The left Heisenberg factor followed by the middle unipotent
factor gives exactly the coordinates
\(n(v,t)\iota(n')\) of Lemma~\ref{ju-heisenberg}.
On the right, reorder
\(\iota(n')n(v,t)\) as
\(n(n'v,t)\iota(n')\).
The real Jacobian in \(v\) is
\(|\det n'|^2=1\); the central coordinate is unchanged.
Thus the two product measures are Haar measures on \(N_r\),
up to constant factors.  Inversion also preserves Haar measure
on \(N_r\), because this nilpotent group is unimodular.
This accounts for the inverse in \eqref{jg-iterated}.

Repeat until reaching \(G_0\).
At stage \(a\) the exponent is \(2(M_a-2)\).
Each discarded complement has Haar measure zero by
Lemma~\ref{jm-cell-haar}; Fubini in the preceding coordinates
therefore discards only a null set at every stage.
Tonelli permits all the substitutions for nonnegative \(f\).
The case \(r=0\) has empty products and is immediate.
\end{proof}

Let
\[
 Q=(\mathbb C^\times)^r\times
                    \{\zeta\in\mathbb C:|\zeta|^2>TT'\},
 \qquad g_q=b(c,\mu(\zeta)),
\]
and define
\begin{equation}\label{jg-radial-measure}
 d\nu(q)=
    \prod_{a=0}^{r-1}|c_a|^{2(M_a-2)}
          (|\zeta|^2-TT')^{k-1}\,dc\,d\zeta .
\end{equation}
For \(p=((k_1,n_1),(k_2,n_2))\), put
\[
 dp=dk_1\,dn_1\,dk_2\,dn_2,\qquad
                s_p(g)=k_1^{-1}n_1^{-1}g n_2k_2 .
\]

\begin{theorem}\label{jg-intformula}
There is a unique normalization of Haar measure \(dg\) on
\(G_r\), for the fixed measures \(dn,dk_1,dk_2\), for which
\begin{equation}\label{jg-integral}
       \int_{G_r}f(g)\,dg
        =\int_Q\int f(s_p(g_q))\,dp\,d\nu(q)
\end{equation}
for every nonnegative measurable \(f\).
\end{theorem}

\begin{proof}
Insert \eqref{jg-terminal} in \eqref{jg-iterated}.
The compact middle matrices commute with \(b(c,1)\).
Move each compact factor to the outside, conjugating the
adjacent unipotent variable.  These conjugations preserve
Haar measure, so the result is a positive constant times
the right side of \eqref{jg-integral}.
Divide \(dg\) by that constant.
For uniqueness, apply the formula to a nonnegative compactly
supported continuous function of positive integral.

All measures in this argument are locally finite and
\(\sigma\)-finite.  The density in \eqref{jg-radial-measure}
is continuous on \(Q\), and the section is smooth there.
The excluded complex coordinate hyperplanes and terminal
circle have planar product measure zero.  Consequently the
same formula holds with the archive's nonnegative density on
all of \(\mathbb C^r\times\mathbb C\), including its convention
at the boundary when \(k=1\).
\end{proof}

\subsection{Integrability at the exact decay threshold}

\begin{theorem}\label{jg-lone}
For \(k\geq1\), \(r\geq0\), and
\(A>2(k+r)(r+1)\), one has
\begin{equation}\label{jg-weight-integrable}
                     \int_{G_r}M(g)^{-A}\,dg<\infty .
\end{equation}
In particular every almost everywhere strongly measurable \(F\) satisfying
\(\|F(g)\|\leq C M(g)^{-A}\) is integrable when
\(A\geq B(k,r)=2(k+r)(r+1)+1\).
\end{theorem}

\begin{proof}
The case \(r=0\) is the ball integral proved in
Section~\ref{jm-section}: its radial upper bound is a constant
times
\(\int_0^1 \rho^{2k-1}(1-\rho^2)^{A/2-k-1}\,d\rho\),
which is finite for \(A>2k\).

Suppose \(r\geq1\), and set \(d=2r+k-1\), \(D=2d+1\).
Retain the two constant terms in the exact Frobenius formula
in the proof of Lemma~\ref{jm-e-int}.  The same two changes
of variables now give
\begin{equation}\label{jg-sharp-e}
 \int_{\mathbb C^d\times\mathbb R}
    M(w(c)\iota(m)n(v,t))^{-A}\,dv\,dt
 \leq C_{A,d}\min(|c|^{-D},|c|^{-1})
       \bigl(M(m)^2+|c|^2+|c|^{-2}\bigr)^{(D-A)/2}.
\end{equation}
Indeed the squared norm is at least either
\[
 B_c^2+|c|^2(\|v\|^2+t^2)
 \quad\hbox{or}\quad B_c^2+\|mv\|^2+|c|^2t^2,
 \qquad B_c^2=M(m)^2+|c|^2+|c|^{-2}.
\]
The real Jacobian of \(v\mapsto mv\) is one.
The Euclidean integral after scaling is finite since \(A>D\).

For the left Heisenberg variable write
\[
 Q_L=\|n(v_L,t_L)f_0\|
  =\bigl(1+\|v_L\|^2+t_L^2+
                       h_{r-1}(v_L,v_L)^2/4\bigr)^{1/2}.
\]
Lemma~\ref{jm-left-factor} gives, for each fixed \(n_L,c\),
\[
             n_Lw(c)=\kappa w(cQ_L)\iota(u)n_0.
\]
After multiplication by \(\iota(m)n_R\), move \(n_0\)
past \(\iota(m)\), change the right Heisenberg variable
by translation, and change \(m\) to \(um\).
These preserve the relevant Haar measures; left multiplication
by \(\kappa\) preserves \(M\).
Thus the integral over \(m,n_R\) is exactly its value with
\(n_L=1\) and \(c\) replaced by \(cQ_L\).
This equality holds pointwise in \(n_L,c\), so no measurable
choice of \(\kappa,u,n_0\) is needed.

In the Haar density \(|c|^{2d}\,dc\), the substitution
\(c'=cQ_L\) contributes \(Q_L^{-2d-2}\).
This factor is integrable in \(n_L\).  In fact
\[
 \begin{split}
 \int Q_L^{-2d-2}\,dv_L\,dt_L
 &\leq \int_{\mathbb C^d\times\mathbb R}
                   (1+\|v_L\|^2+t_L^2)^{-d-1}\,dv_L\,dt_L\\
 &<\infty ,
 \end{split}
\]
since \(2d+2>2d+1\).
It remains to integrate \eqref{jg-sharp-e} against
\(|c|^{2d}\,dc\,dm\).

Put \(b=A-D>1\) and \(R=M(m)\geq1\).
For \(|c|=\rho\geq1\), the radial factor
\(\rho^{2d}\rho\,\rho^{-D}\) is one, and
\[
 \int_1^\infty(R^2+\rho^2+\rho^{-2})^{-b/2}\,d\rho
 \leq R^{1-b}\int_0^\infty(1+s^2)^{-b/2}\,ds.
\]
For \(0<\rho\leq1\), the radial factor is \(\rho^{2d}\), and
the integral is at most \(R^{-b}/(2d+1)\leq R^{1-b}/(2d+1)\).
The remaining middle integral is bounded by
\[
              C\int_{G_{r-1}}M(m)^{-(A-D-1)}\,dm.
\]
Finally
\[
 2(k+r)(r+1)-2(k+r-1)r=D+1.
\]
Our assumption therefore gives precisely the induction
hypothesis at rank \(r-1\).  Tonelli justifies the argument
before finiteness has been established.
The final assertion is the dominated integrability criterion
for an almost everywhere strongly measurable function,
applied to the integrable scalar majorant \(C M^{-A}\).
\end{proof}

\subsection{Dilation and its two Jacobians}

For \(t>0\), define
\[
 a_t e_a=t^{r-a}e_a,\qquad
 a_t f_a=t^{-(r-a)}f_a,\qquad a_t|_{V_0}=1,
 \qquad
 \kappa_r=\sum_{a=0}^{r-1}2(M_a-1)(r-a).
\]
These elements belong to \(G_r\) and commute with \(K_z,K_{z'}\).
The character \(\lambda_z=2\operatorname{Im}\ell_z\) is the one
defined in Section~\ref{ju-section}.

\begin{proposition}\label{jg-dilation}
Conjugation \(\delta_t(n)=a_tna_t^{-1}\) preserves \(N_r\) and
satisfies
\begin{equation}\label{jg-diln}
 \lambda_z(\delta_t n)=t\lambda_z(n),\qquad
                 (\delta_t)_*(dn)=t^{-\kappa_r}\,dn .
\end{equation}
For \(s>0\), put
\(S_s(c,\zeta)=((s^{r-a}c_a)_a,\zeta)\).
Then
\begin{equation}\label{jg-dilq}
 (S_s)_*\nu=s^{-\kappa_r}\nu,\qquad
       a_{t_1}^{-1}g_q a_{t_2}=g_{S_{t_1t_2}q}.
\end{equation}
\end{proposition}

\begin{proof}
The defining triangular conditions for \(N_r\) are preserved
by diagonal conjugation.
Every entry in \(\ell_z\) has dilation weight one: the
successive \(e_a,e_{a+1}\) weights differ by one, and the
last \(e_{r-1}\) has weight one while \(V_0\) has weight zero.
This proves the character identity, including \(r=0\).

In the first Heisenberg factor,
\[
       (v,u)\longmapsto(t^r a_t^{\mathrm{mid}}v,t^{2r}u).
\]
The middle complex determinant has absolute value one.
The real Jacobian is consequently
\(t^{2r(d+1)}\), with \(d=2r+k-1\).
The remaining \(N_{r-1}\) has its own dilation Jacobian.
The recurrence
\(\kappa_r=2r(2r+k)+\kappa_{r-1}\), starting with zero,
gives the stated sum.  A map whose coordinate Jacobian is
\(t^{\kappa_r}\) pushes the measure forward to
\(t^{-\kappa_r}\) times the measure.  This fixes the sign
of the exponent in \eqref{jg-diln}.

In the radial measure, \(c_a\mapsto s^{r-a}c_a\) multiplies
its density by \(s^{2(M_a-2)(r-a)}\) and its planar
volume by \(s^{2(r-a)}\).
The total exponent is \(\kappa_r\); the terminal factor is
unchanged.  This proves the first identity in \eqref{jg-dilq}.
The second follows by applying both sides to \(e_a,f_a,V_0\).
\end{proof}

Theorems~\ref{jg-intformula} and~\ref{jg-lone}, together with
Proposition~\ref{jg-dilation} and the compact conjugation
invariance already proved, establish the geometric and
\(L^1\) hypotheses used by the source's Jacquet pairing
argument.  The exponents here are the actual exponents in
that argument; in particular the normalization constant is
absorbed in Haar measure, and the two dilation pushforwards
have negative exponents.  The corresponding source
definitions and the precise read scopes are recorded in
the accompanying manifest \cite{alpo-uniform-source}.


\section{Regularization, unfolding, and Jacquet adjointness}
\label{jp-section}

We prove the analytic passage from the group operators to the
Jacquet pairing.  The geometric integration formula and the
integrable majorant are the theorems of
Sections~\ref{jm-section} and~\ref{jg-section}.
In particular, none of the estimates in this section assumes
absolute convergence of an unweighted integral over the
unipotent group.

\subsection{The transform and the differentiability class}

Fix finite dimensional complex Hilbert spaces \(W',W\), continuous
unitary representations \(\pi,\pi'\) of \(K_z,K_{z'}\), and
\(E=\operatorname{Hom}_{\mathbb C}(W',W)\).
Use the Hilbert--Schmidt inner product on \(E\), linear in the
first variable.  Replacing its norm by the operator norm changes
the estimates below only by constants depending on the dimensions.

Let \(S_z=N_rK_z\) and \(S_{z'}=N_rK_{z'}\).
Each element has a unique expression \(nk\); the two factors have
trivial intersection because \(N_r\) is unipotent and \(K_z\)
fixes all the outer vectors.  Both groups are unimodular, and
\(dn\,dk\) is Haar measure: compact conjugation preserves \(dn\).
We write \(P=S_z\times S_{z'}\), with product Haar measure \(dp\).
The coordinates \(p=((k_1,n_1),(k_2,n_2))\) are those of
Theorem~\ref{jg-intformula}; their order is retained in formulas.
Set
\[
 \begin{split}
 s_p(g)&=k_1^{-1}n_1^{-1}g n_2k_2,\\
 \rho(p)F(g)&=\pi(k_1)F(s_p(g))\pi'(k_2^{-1}),\\
 \chi_\tau(p)&=\exp\!\bigl(i(\tau_1\lambda_z(n_1)
                              -\tau_2\lambda_{z'}(n_2))\bigr),
                 \qquad \tau\in\mathcal Q=(0,\infty)^2 .
 \end{split}
\]
The characters \(\lambda_z,\lambda_{z'}\) are additive on \(N_r\)
and invariant under their compact stabilizers.  Thus \(\chi_\tau\)
is a unitary character of \(P\).
For the product inherited from \(S_z\times S_{z'}\),
\[
                       \rho(p)\rho(p')=\rho(pp').
\]
This follows directly by multiplying \(n_1k_1\) and \(n'_1k'_1\),
and similarly on the right.  The transformations of \(E\) supplied
by the compact representations preserve its Hilbert--Schmidt norm.
Each \(\rho(p)\) is therefore unitary on \(L^2(G_r,E)\).

For \(X\) in the real Lie algebra of \(G_r\), define
\[
 L_XF(g)=\left.\frac d{dt}F(e^{-tX}g)\right|_{t=0},
 \qquad
 R_XF(g)=\left.\frac d{dt}F(ge^{tX})\right|_{t=0}.
\]
For words \(u=(X_1,\ldots,X_i)\), \(v=(Y_1,\ldots,Y_j)\),
let \(L_u=L_{X_1}\cdots L_{X_i}\) and
\(R_v=R_{Y_1}\cdots R_{Y_j}\), with the empty word equal to
the identity.  Left and right derivatives commute.
Choose the compact unit ball in the Lie algebra consisting of
matrices whose entries have absolute value at most one, and set
\[
 \|F\|_{A,m}
 =\max_{i+j\leq m}\ 
      \sup_{\substack{g\in G_r\\X_a,Y_b\text{ in that ball}}}
        M(g)^A\|L_uR_vF(g)\|.
\]
Write \(\mathcal C_A^m\) for the \(C^m\) functions for which this
quantity is finite.  This is the differentiability and decay class
in the source: a \(C^m\) function on the embedded matrix Lie group
has local \(C^m\) ambient extensions, by extending in local
submanifold coordinates, and conversely such extensions give the
intrinsic derivatives used here.
Throughout the section,
\[
                         A\geq B(k,r).
\]
Theorem~\ref{jg-lone} implies
\(\mathcal C_A^0\subset L^1(G_r,E)\).

Define, wherever the integral is absolutely convergent,
\[
 J_\tau F(g)=\int_P\chi_\tau(p)\rho(p)F(g)\,dp,\qquad
 \Phi_{F,F'}(p)=\int_{G_r}
                  \langle\rho(p)F(g),F'(g)\rangle\,dg.
\]
Let \(X(p)=\Xi(n_1k_1)\Xi(n_2k_2)\), for the function
\(\Xi\) constructed in Section~\ref{jm-section}.
Its two proved properties give
\begin{align}
 \int_{G_r}\|\rho(p)F(g)\|\|F'(g)\|\,dg
       &\leq C\|F\|_{A,0}\|F'\|_{A,0}X(p),
                                      \label{jp-coef}\\
 \int_P\frac{X(p)}
       {(1+|\lambda_z(n_1)|)^2(1+|\lambda_{z'}(n_2)|)^2}\,dp
       &<\infty .                     \label{jp-xi}
\end{align}
In particular \(\Phi\) is an ordinary absolutely convergent
group integral for every \(p\).

\begin{lemma}\label{jp-orbital}
For \(F\in\mathcal C_A^0\), the integral \(J_\tau F(g_q)\)
is absolutely convergent for almost every \(q\) with respect to
\(\nu\), on one common set valid for every \(\tau\in\mathcal Q\).
On that set it is continuous in \(\tau\).
It has a jointly measurable version in \((\tau,q)\).
\end{lemma}

\begin{proof}
The integration formula and unitarity on \(E\) give
\[
 \int_Q\int_P\|\rho(p)F(g_q)\|\,dp\,d\nu(q)
                         =\int_{G_r}\|F(g)\|\,dg<\infty .
\]
Tonelli supplies the common full measure set.
The integrand's norm is independent of \(\tau\), so dominated
convergence gives continuity there.
The map \((\tau,q,p)\mapsto\chi_\tau(p)\rho(p)F(g_q)\) is
measurable.  Integration of a measurable finite dimensional
vector valued function, with value zero assigned where its
norm integral is infinite, is measurable; this follows
componentwise from truncation and the scalar integration
theorem.  This proves the final assertion.
\end{proof}

We use this version to define
\[
 N_\tau(F)=\int_Q\|J_\tau F(g_q)\|^2\,d\nu(q)\in[0,\infty],
 \qquad
 \beta_\tau(F,F')=
       \int_Q\langle J_\tau F(g_q),J_\tau F'(g_q)\rangle\,d\nu(q)
\]
when both squared norms are finite.
The eventual pairing in the source is \(\beta=\beta_{(1,1)}\).

\subsection{Dilation estimates and the regularized form}

For \(t=(t_1,t_2)\in\mathcal Q\), put
\[
 D_tF(g)=F(a_{t_1}^{-1}g a_{t_2}),\qquad
 U_tF(g)=F(a_{t_1}g a_{t_2}^{-1}).
\]
Here \(U_t=D_{(t_1^{-1},t_2^{-1})}\).

\begin{lemma}\label{jp-dilation-seminorm}
On every compact rectangle in \(\mathcal Q\),
\(\|D_tF\|_{A,m}\leq C\|F\|_{A,m}\).
On \([1/2,2]^2\) one also has
\begin{equation}\label{jp-dilation-cont}
 \|D_tF-F\|_{A,4}
       \leq C(|t_1-1|+|t_2-1|)\|F\|_{A,5}.
\end{equation}
For \(|\alpha|\leq m\), the ordinary parameter derivatives
\(\partial_t^\alpha U_tF\) have decay norm of order zero at
most \(C_\alpha\|F\|_{A,m}\), uniformly on a compact rectangle.
\end{lemma}

\begin{proof}
The Frobenius norm is submultiplicative, and the matrices
\(a_t,a_t^{-1}\) are uniformly bounded on such a rectangle.
Consequently \(M(a_{t_1}^{-1}g a_{t_2})\) and \(M(g)\) are
comparable there, uniformly in \(g\).
The chain rule gives
\[
 L_XD_tF=D_t L_{\operatorname{Ad}(a_{t_1}^{-1})X}F,\qquad
 R_YD_tF=D_t R_{\operatorname{Ad}(a_{t_2}^{-1})Y}F.
\]
Iterating these identities, and using multilinearity in the
Lie algebra elements, proves the first bound: the indicated
adjoint maps have uniformly bounded coefficients.

Let \(H\) be the diagonal matrix with entries
\(r,r-1,\ldots,1,0,\ldots,0,-1,\ldots,-r\).
Then \(a_t=e^{(\log t)H}\), \(H\) belongs to the real Lie
algebra, and
\[
       \partial_{t_1}D_tF=t_1^{-1}D_tL_HF,\qquad
       \partial_{t_2}D_tF=t_2^{-1}D_tR_HF.
\]
The corresponding formulas for \(U_t\) have the opposite
signs and the same bounded scalar factors.
Higher parameter derivatives are finite sums of translated
Lie derivatives, with smooth bounded coefficients.
This proves their asserted decay estimates.

Apply one parameter derivative also to the displayed formulas
for \(L_uR_vD_tF\), with \(|u|+|v|\leq4\).
It differentiates either an adjoint coefficient or the
translated function.  The latter term uses at most five
Lie derivatives of \(F\).  All terms are bounded, after
multiplication by \(M(g)^A\), by \(C\|F\|_{A,5}\).
Integrate first along \(t_1\) and then along \(t_2\) from
\((1,1)\) to \(t\).  Taking the suprema defining the
seminorm gives \eqref{jp-dilation-cont}.
\end{proof}

Choose a real smooth function \(\chi\), supported in
\((1/2,2)\), with \(\int_{\mathbb R}\chi(t)\,dt=1\).
It need not be nonnegative.  Write
\[
 \delta_t p=((k_1,a_{t_1}n_1a_{t_1}^{-1}),
             (k_2,a_{t_2}n_2a_{t_2}^{-1})),\qquad
 \lambda_1=\lambda_z(n_1),\quad\lambda_2=\lambda_{z'}(n_2).
\]
Define the compactly supported \(C^4\) function
\[
 a_{F,F',p}(t)=
       \chi(t_1)\chi(t_2)(t_1t_2)^{\kappa_r}
                                  \Phi_{F,F'}(\delta_t p),
\]
extended by zero outside \(\mathcal Q\).
For \(F,F'\in\mathcal C_A^4\), define
\begin{equation}\label{jp-reg-definition}
 B_\tau(F,F')
 =\int_P\frac{1}{(1+\tau_1^2\lambda_1^2)
                    (1+\tau_2^2\lambda_2^2)}
    \int_{\mathbb R^2}
      (1-\partial_{t_1}^2)(1-\partial_{t_2}^2)
          a_{F,F',p}(t)\,
      e^{i(\tau_1t_1\lambda_1-\tau_2t_2\lambda_2)}\,dt\,dp .
\end{equation}
This is exactly the source's regularized pairing.

\begin{proposition}\label{jp-reg-bound}
The integral \eqref{jp-reg-definition} is absolutely convergent.
It is sesquilinear in \(F,F'\), continuous in \(\tau\in\mathcal Q\),
and satisfies, locally uniformly in \(\tau\),
\begin{equation}\label{jp-Bbound}
       |B_\tau(F,F')|\leq C_\tau\|F\|_{A,4}\|F'\|_{A,4}.
\end{equation}
\end{proposition}

\begin{proof}
The substitution \(g=a_{t_1}h a_{t_2}^{-1}\) in the group
integral gives
\begin{equation}\label{jp-phase-translate}
 \Phi_{F,F'}(\delta_t p)
                  =\Phi_{U_tF,U_tF'}(p).
\end{equation}
It uses both unimodularity of \(G_r\) and commutation of
\(a_t\) with the compact stabilizers.
For fixed \(p\), the derivatives in \(t\) on the right may
be taken under the group integral.  Indeed every derivative
of total order at most four is bounded, on the fixed compact
rectangle, by a constant times
\[
                 M(s_p(g))^{-A}M(g)^{-A}.
\]
This is integrable by \eqref{jp-coef}.
The same bound applies to the continuous highest derivatives,
so repeated dominated convergence proves the asserted \(C^4\)
regularity, including the mixed derivative of order \((2,2)\).

After differentiation, apply \eqref{jp-coef} to each of the
translated derivatives, and then the preceding lemma.
For \(0\leq i,j\leq2\) this yields
\[
 |\partial_{t_1}^i\partial_{t_2}^j a_{F,F',p}(t)|
 \leq C\|F\|_{A,4}\|F'\|_{A,4}X(p)
                      {\bf1}_{[1/2,2]^2}(t).
\]
On a compact subset of \(\mathcal Q\),
\[
 (1+\tau_j^2\lambda_j^2)^{-1}
                         \leq C(1+|\lambda_j|)^{-2}.
\]
Equation~\eqref{jp-xi} proves absolute convergence and the
bound.  It also supplies a common dominating function for
continuity in \(\tau\).  Sesquilinearity follows from that
of the integrands, with all the integrals now justified.
\end{proof}

\subsection{Averaging in the two characters}

For a real \(C^2\) function \(w\) of compact support in
\((0,\infty)\), write
\[
                \widehat w(x)=\int_{\mathbb R}w(\tau)e^{i\tau x}\,d\tau.
\]
Twice integrating by parts gives
\[
          |\widehat w(x)|\leq C_w(1+|x|)^{-2}.
\]
For example combine the bounds \(\|w\|_1\) and
\(|x|^{-2}\|w''\|_1\); the boundary terms vanish by compact
support.
For two such functions set
\(\Omega(p)=\widehat w_1(\lambda_1)\widehat w_2(-\lambda_2)\).

\begin{proposition}\label{jp-reg-average}
For \(F,F'\in\mathcal C_A^4\),
\begin{equation}\label{jp-average}
 \int_{\mathcal Q}w_1(\tau_1)w_2(\tau_2)B_\tau(F,F')\,d\tau
                           =\int_P\Omega(p)\Phi_{F,F'}(p)\,dp .
\end{equation}
The integral on the right is absolutely convergent.
\end{proposition}

\begin{proof}
Its absolute convergence follows from
\eqref{jp-coef}, \eqref{jp-xi}, and Fourier decay.
For each fixed \(p\), integration by parts in \(t_1,t_2\)
gives
\[
 \begin{split}
 &\frac{1}{(1+\tau_1^2\lambda_1^2)(1+\tau_2^2\lambda_2^2)}
    \int (1-\partial_{t_1}^2)(1-\partial_{t_2}^2)a_p(t)
                         e^{i(\tau_1t_1\lambda_1-\tau_2t_2\lambda_2)}\,dt\\
 &\hspace{25mm}
       =\int a_p(t)e^{i(\tau_1t_1\lambda_1-\tau_2t_2\lambda_2)}\,dt .
 \end{split}
\]
The right side is asserted here for fixed \(p\); no unweighted
interchange of \(p\) and \(t\) is being made.
The absolute bound used in Proposition~\ref{jp-reg-bound}
allows averaging the left side in \(\tau\) under the \(p\)
integral.  For fixed \(p\), all \(t,\tau\) integrals are
over compact sets, so their interchange gives
\[
 \int_P\int_{\mathbb R^2}
       \chi(t_1)\chi(t_2)(t_1t_2)^{\kappa_r}
       \Phi_{F,F'}(\delta_t p)
       \widehat w_1(t_1\lambda_1)
       \widehat w_2(-t_2\lambda_2)\,dt\,dp .
\]
This new double integral is absolutely convergent.
To see this, for each \(t\) in the support change variables
\(p'=\delta_t p\).
Proposition~\ref{jg-dilation} gives
\(dp=(t_1t_2)^{-\kappa_r}dp'\), and its character identity
replaces \(t_j\lambda_j\) by the corresponding \(\lambda_j'\).
The absolute integral is therefore bounded by
\[
 \|\chi\|_1^2\int_P|\Omega(p')\Phi_{F,F'}(p')|\,dp'<\infty .
\]
Fubini now applies.  The same substitution without absolute
values and \((\int\chi)^2=1\) prove \eqref{jp-average}.
\end{proof}

\begin{proposition}\label{jp-unfold}
For \(F\in\mathcal C_A^4\) and nonnegative \(w_1,w_2\) as above,
\begin{equation}\label{jp-unfold-identity}
 \int_{\mathcal Q}w_1(\tau_1)w_2(\tau_2)N_\tau(F)\,d\tau
                  =\int_P\Omega(p)\Phi_{F,F}(p)\,dp<\infty .
\end{equation}
The equality regards the nonnegative real number on the left
as a complex number on the right.
\end{proposition}

\begin{proof}
Fix \(q\) in the common full measure set of
Lemma~\ref{jp-orbital} and put \(f_q(p)=\rho(p)F(g_q)\).
This is an \(E\)-valued \(L^1(P)\) function.
Expanding the squared norm of its character integral gives
an absolutely convergent double integral.
In that integral set the second variable equal to \(p\)
and the first to \(pp'\).
Left invariance of Haar measure, the character identity,
and unitarity on \(E\) give
\begin{equation}\label{jp-fixed-q}
 \|J_\tau F(g_q)\|^2
 =\int_P\int_P \chi_\tau(p')
     \langle\rho(p')F(s_p(g_q)),F(s_p(g_q))\rangle\,dp'\,dp .
\end{equation}
The compact factors attached to \(p\) act unitarily.
They transform the two entries of the inner product on
the right into \(\rho(pp')F(g_q)\) and \(\rho(p)F(g_q)\).
The absolute double integral before inserting the character
is at most \(\|f_q\|_1^2\).
Thus averaging \eqref{jp-fixed-q} in \(\tau\) gives the same
formula with \(\chi_\tau(p')\) replaced by \(\Omega(p')\).

It remains to justify integration in \(q\); the preceding
bound need not be integrable in \(q\).
Instead use the Fourier decay already obtained.
The integration formula for \(G_r\), followed by Tonelli, gives
\[
 \begin{split}
 &\int_Q\int_P\int_P|\Omega(p')|
       \|\rho(p')F(s_p(g_q))\|\,\|F(s_p(g_q))\|
                        \,dp'\,dp\,d\nu(q)\\
 &\quad=\int_P|\Omega(p')|
               \int_{G_r}\|\rho(p')F(g)\|\,\|F(g)\|\,dg\,dp'
 <\infty ,
 \end{split}
\]
by \eqref{jp-coef} and \eqref{jp-xi}.
This proves that the averaged right side of
\eqref{jp-fixed-q} is integrable in \(q\).
On the left, \(w_1w_2\geq0\), so Tonelli applies even before
finiteness is known.  Fubini on the absolutely integrable
right side now proves \eqref{jp-unfold-identity}.
\end{proof}

\subsection{Equality at every positive character}

\begin{theorem}\label{jp-pairing-equality}
For every \(F,F'\in\mathcal C_A^5\) and every
\(\tau\in\mathcal Q\), the two squared norms
\(N_\tau(F),N_\tau(F')\) are finite and
\begin{equation}\label{jp-beta-equals}
                         \beta_\tau(F,F')=B_\tau(F,F').
\end{equation}
In particular
\[
       |\beta(F,F')|\leq C\|F\|_{A,4}\|F'\|_{A,4}.
\]
\end{theorem}

\begin{proof}
First take only \(F\in\mathcal C_A^4\), and put
\(b(\tau)=B_\tau(F,F)\).
The two averaging propositions say
\[
       \int w_1w_2N_\tau(F)\,d\tau=\int w_1w_2b(\tau)\,d\tau
\]
for every pair of nonnegative test functions.
The left side is finite.
On any compact rectangle in \(\mathcal Q\), choose test
functions positive on that rectangle; this shows that
\(\tau\mapsto N_\tau(F)\) is locally integrable and finite
almost everywhere.
At its Lebesgue points, use products of nonnegative smooth
approximate identities in the displayed equality.
Their support can be kept inside \(\mathcal Q\).
The Lebesgue differentiation theorem on the left and continuity
of \(b\) on the right give
\[
                           N_\tau(F)=b(\tau)
                           \quad\hbox{for almost every }\tau .
\]
In particular \(b\) is real and nonnegative almost everywhere,
and hence everywhere by continuity.

For any fixed \(\tau_0\), choose a sequence of such equality
points tending to \(\tau_0\).
The pointwise continuity in Lemma~\ref{jp-orbital} and
Fatou's lemma give
\begin{equation}\label{jp-diagonal-domination}
                    N_{\tau_0}(F)\leq b(\tau_0).
\end{equation}
Thus the norm is finite at every positive character even
for \(F\in\mathcal C_A^4\).  Moreover, locally uniformly in
\(\tau\),
\[
              \|J_\tau F(g_\cdot)\|_{L^2(Q,\nu)}
                                      \leq C\|F\|_{A,4}.
\]
Linearity holds because the defining orbital integrals are
absolutely convergent on common full measure sets.

We next use the fifth derivative.
Changes of variables \(p=\delta_t p'\), and then
\(q'=S_{t_1t_2}q\), give
\begin{align}
 J_\tau(D_tF)(g)
     &=(t_1t_2)^{\kappa_r}
           J_{(\tau_1t_1,\tau_2t_2)}F(a_{t_1}^{-1}g a_{t_2}),
                                                   \label{jp-J-dilation}\\
 N_\tau(D_tF)
     &=(t_1t_2)^{\kappa_r}
                           N_{(\tau_1t_1,\tau_2t_2)}(F).
                                                   \label{jp-N-dilation}
\end{align}
These equalities use precisely the two pushforward exponents
in Proposition~\ref{jg-dilation}.
They are valid on common full measure sets in \(q\);
the radial scaling preserves null sets.

For \(F\in\mathcal C_A^5\), the estimate just proved and
\eqref{jp-dilation-cont} show
\[
 \|J_\tau(D_tF)(g_\cdot)-J_\tau F(g_\cdot)\|_{L^2}
 \leq C\|D_tF-F\|_{A,4}
                 \longrightarrow0\qquad(t\longrightarrow(1,1)).
\]
Consequently \(N_\tau(D_tF)\to N_\tau(F)\).
Equation~\eqref{jp-N-dilation} now proves continuity of
\(\tau'\mapsto N_{\tau'}(F)\) at \(\tau\).
Since \(\tau\) was arbitrary, it is continuous on
\(\mathcal Q\), and its almost everywhere equality with
\(b\) is therefore equality everywhere.

Finally both sides of \eqref{jp-beta-equals} are
sesquilinear on \(\mathcal C_A^5\).
The integral defining the left side is absolutely convergent
by Cauchy--Schwarz and the finite diagonal norms.
Equality of the diagonal values, followed by complex
polarization, proves \eqref{jp-beta-equals}.
The last bound is \eqref{jp-Bbound} at \((1,1)\).
\end{proof}

\subsection{Adjoints without differentiation of an orbital integral}

Let \(U\subset G_r\) be open and stable under left \(S_z\)
and right \(S_{z'}\).
Consider the finite order differential operator
\[
              \Theta F=\sum_T A_T(g)R_{w_T}F,
 \qquad
 \Theta^\dagger F'
       =\sum_T(-1)^{|w_T|}
                     R_{\operatorname{rev}(w_T)}(A_T(g)^*F').
\]
The coefficients and their pointwise adjoints are smooth
on \(U\), and \(m=\max_T|w_T|\).
Assume that \(\Theta\rho(p)F=\rho(p)\Theta F\) on \(U\)
for \(C^m\) functions \(F\).
The support of a function is its closed support in \(G_r\).

\begin{lemma}\label{jp-ibp}
If a \(C^1\) scalar function \(h\) and \(R_Xh\) are in
\(L^1(G_r)\), then \(\int R_Xh\,dg=0\).
Consequently integration by parts against a real Lie
derivative has no boundary term whenever the three products
\(\langle f,f'\rangle\),
\(\langle R_Xf,f'\rangle\), and
\(\langle f,R_Xf'\rangle\) are integrable.
\end{lemma}

\begin{proof}
The fundamental theorem of calculus gives, for \(t>0\),
\[
       h(ge^{tX})-h(g)=\int_0^t(R_Xh)(ge^{sX})\,ds.
\]
The absolute double integral on the right is
\(t\|R_Xh\|_1\), by right invariance of Haar measure.
Fubini therefore applies.  Integration of the left side
gives zero, and the right side gives \(t\int R_Xh\,dg\).
For the second assertion use the product rule for the
inner product, which is differentiated in a real parameter.
\end{proof}

\begin{theorem}\label{jp-adjoint}
Suppose \(F,F'\) are \(C^{m+5}\), supported in \(U\), and
\[
       F,F',\Theta F,\Theta^\dagger F'\in\mathcal C_A^5 .
\]
For every \(p\in P\), every term \(T\), and
\(0\leq i\leq j\leq|w_T|\), assume the actual integrability
condition
\begin{equation}\label{jp-intclause}
 \int_{G_r}
   \bigl\|R_{\operatorname{rev}(w_T[0:i])}(A_T^*F')(g)\bigr\|
   \bigl\|R_{w_T[j:]}(\rho(p)F)(g)\bigr\|\,dg<\infty .
\end{equation}
Then, for every \(\tau\in\mathcal Q\),
\begin{equation}\label{jp-adjoint-identity}
            \beta_\tau(\Theta F,F')
                         =\beta_\tau(F,\Theta^\dagger F').
\end{equation}
\end{theorem}

\begin{proof}
The product \(A_T^*F'\), extended by zero outside \(U\),
is \(C^{m+5}\) on all of \(G_r\).
At every point outside \(U\), the closed support hypothesis
gives a neighborhood on which \(F'\) is zero, so no
regularity of \(A_T\) outside \(U\) is needed.

Fix \(p,T\), and transfer the letters of \(w_T\) one at a
time from \(\rho(p)F\) to \(A_T^*F'\).
At the \(j\)-th transfer, the three product norms needed in
Lemma~\ref{jp-ibp} are exactly the cases
\((i,j)=(j,j+1),(j,j),(j+1,j+1)\) of
\eqref{jp-intclause}.  Thus each transfer is legitimate,
and
\[
 \int\langle R_{w_T}\rho(p)F,A_T^*F'\rangle\,dg
   =(-1)^{|w_T|}
      \int\langle\rho(p)F,
                 R_{\operatorname{rev}(w_T)}(A_T^*F')\rangle\,dg .
\]
The reversal of the word and the sign are both necessary.
Summing and using commutation on \(U\), while \(F'\) is
zero outside \(U\), proves
\[
                  \Phi_{\Theta F,F'}(p)
                               =\Phi_{F,\Theta^\dagger F'}(p).
\]
The same identity holds at \(\delta_t p\).
It therefore holds for their regularized brackets and their
parameter derivatives, and gives equality of the two
\(B_\tau\) forms.  Theorem~\ref{jp-pairing-equality} now
gives \eqref{jp-adjoint-identity}.
No derivative of the unweighted orbital integral was used.
\end{proof}

\subsection{The polynomial coefficient criterion}

To check \eqref{jp-intclause} in the setting used by the
source, put \(v_a=e_a\) for \(a<r\), \(v_r=z\), and use
\(v'_r=z'\) on the right.
Let
\[
 D_l(g)=\det\bigl(h_r(gv'_b,v_a)\bigr)_{0\leq a,b<l},
 \quad D_0=1,\qquad
 \mathcal G(g)=|D_{r+1}(g)|^2-TT'|D_r(g)|^2 .
\]
Here this notation for \(D_l\) refers to these Bessel
minors, not the flag minors of Section~\ref{jm-section}.
The function \(\mathcal G\) is a real polynomial in the
real and imaginary matrix entries.
Every \(S_z,S_{z'}\) change of the indicated frames is
triangular with diagonal one.  Hence every \(D_l\),
and therefore \(\mathcal G\), is invariant under those
left and right actions.
It follows that \(U=\{\mathcal G>0\}\) is stable and open.

\begin{corollary}\label{jp-poly-adjoint}
Let \(d\geq0\), \(\epsilon>0\), and suppose that for every
coefficient \(A_T,A_T^*\), every left and right word of
total length at most \(m+5\), and every
\(\mathcal G(g)\geq\epsilon\), its derivative has norm at
most \(C M(g)^d\), uniformly for Lie algebra letters in
the fixed unit ball.
Assume that \(\Theta\) is smooth and commutes with \(\rho\)
on \(\{\mathcal G>0\}\).
If \(F,F'\in\mathcal C_{A+d}^{m+5}\) vanish whenever
\(\mathcal G<\epsilon\), then
\[
                         \beta(\Theta F,F')
                                      =\beta(F,\Theta^\dagger F').
\]
\end{corollary}

\begin{proof}
Continuity of \(\mathcal G\) places both closed supports
in \(\{\mathcal G\geq\epsilon\}\subset U\).
The Leibniz rule gives
\[
                  F,F',\Theta F,\Theta^\dagger F'
                                           \in\mathcal C_A^5 .
\]
Indeed each differentiated term is bounded by
\(C M^d M^{-(A+d)}\), and uses at most \(m+5\)
derivatives of either factor.  Fixed letters occurring in
\(\Theta\) are handled by scaling them into the unit ball.
All derivatives vanish locally where the function is
outside its closed support.

The first factor in \eqref{jp-intclause} is consequently
bounded by \(C M(g)^{-A}\).
For fixed \(p\), derivatives of the second factor are
finite linear combinations of derivatives of \(F\) at
\(s_p(g)\), with fixed conjugated Lie algebra letters.
Submultiplicativity of \(M\) gives
\(M(s_p(g))^{-(A+d)}\leq C_pM(g)^{-(A+d)}\).
Their product is therefore at most
\(C_p M(g)^{-2A-d}\leq C_pM(g)^{-A}\).
This is integrable by Theorem~\ref{jg-lone}, proving
every case of \eqref{jp-intclause}.
The preceding theorem applies.
\end{proof}

In particular, polynomial matrix coefficients satisfy this
criterion: invariant differentiation acts by linear vector
fields on matrix entries and does not increase polynomial
degree.  The criterion also covers coefficients
\(P(g)\mathcal G(g)^{-q}\) on \(U\), for a polynomial \(P\)
and integer \(q\geq0\).  After \(j\) derivatives, the
Leibniz and chain rules give denominators
\(\mathcal G^{q+i}\), \(i\leq j\), and numerator degree at
most \(\deg P+j\deg\mathcal G\).
On \(\mathcal G\geq\epsilon\) the denominators are bounded
away from zero, giving the required polynomial growth for
all the finitely many derivatives.

This proves the source's pairing bound, equality with the
regularized pairing, adjointness with its full integrability
clause, and its polynomial coefficient version.
The fifth derivative has a specific role: it gives continuity
of dilations in the norm using four derivatives and upgrades
the almost everywhere identity in the characters to an
identity at \((1,1)\).
Absolute convergence of \(J_\tau F(g)\) at every group
element has neither been assumed nor concluded.
All source correspondences are recorded in the accompanying
trace and manifest \cite{alpo-uniform-source}.


\section{Fourier decay for the actual Jacquet transforms}
\label{jf-section}

We now apply the group operators constructed above to actual
Jacquet integrals.  In particular, the differential operators,
their adjoints, and their decay estimates in this section are
not additional hypotheses.  A separate final paragraph specifies
the further operator identities needed to apply the conclusion
to every recursively generated vector in the source
\cite{alpo-uniform-source}.

\subsection{Smooth Jacquet integrals and their radial adjoints}

Retain the Hermitian spaces, measures and compact representations
of Sections~\ref{rl-section}, \ref{jg-section} and
\ref{jp-section}.  Let \(k\geq1\), \(r\geq0\),
\(G_r=U(k+r,r+1)\), and
\(E=\operatorname{Hom}_{\mathbb C}(W',W)\), with its
Hilbert--Schmidt inner product linear in the first variable.
Let \(J_rF=J_{(1,1)}F\) be the Jacquet integral of
Section~\ref{jp-section}.  Put
\[
 \mathscr S_\varepsilon(G_r,E)=
 \{F\in C^\infty(G_r,E):
       \|F\|_{A,m}<\infty\text{ for all }A,m,
       \ F=0\text{ on }\mathcal G<\varepsilon\},
 \qquad \varepsilon>0.
\]
Here \(\mathcal G\) is the full Gram polynomial
\eqref{rl-Gram}, and \(\|\cdot\|_{A,m}\) is the mixed
left and right derivative seminorm defined in
Section~\ref{jp-section}.  These are the actual coefficient
functions used in the Fourier argument.  Write
\(\mathscr S_+=\bigcup_{\varepsilon>0}\mathscr S_\varepsilon\).

Number the outer coordinates from \(0\) to \(r-1\), and set
\[
 d(c)=\prod_{j=0}^{r-1}c_j,\qquad
 \gamma_r(c)=\prod_{j=0}^{r-1}|c_j|^{2(r-j)+k},\qquad
 d\eta_r(c,m)=\prod_{j=0}^{r-1}|c_j|^{-2}\,dc_j\,dm.
\]
Empty products are one.  The terminal variable \(m\) belongs
to \(G_0=U(k,1)\), restricted to
\(G(m)=|h(mz',z)|^2-TT'>0\).  On this open radial domain put
\begin{equation}\label{jf-radial-transform}
 \Phi_F(c,m)=\gamma_r(c)J_rF(b(c)m).
\end{equation}
The integration formula and Theorem~\ref{jp-pairing-equality}
show that \(\Phi_F\in L^2(d\eta_r,E)\) for every
\(F\in\mathscr S_+\).  More precisely, for fixed compatible
Haar measures there is a positive constant \(C_r\) such that
\begin{equation}\label{jf-radial-pairing}
 \langle\Phi_F,\Phi_H\rangle_{L^2(d\eta_r)}
                          =C_r\beta(F,H).
\end{equation}
Indeed, \(\gamma_r^2\prod|c_j|^{-2}\) is exactly the outer
density in the open-cell Haar formula.  The compact-terminal
integration formula of Section~\ref{jg-section} then gives
\eqref{jf-radial-pairing}; compact equivariance preserves the
Hilbert--Schmidt inner product.  The constant is independent
of \(F,H\), and can be made one by the allowed normalization
of terminal Haar measure.

\begin{lemma}\label{jf-pointwise-smooth}
For \(F\in\mathscr S_+\), the Jacquet integral is absolutely
convergent and smooth at every point of the open cell with
\(\mathcal G>0\).  It may be differentiated there under the
integral sign.  If \(\Theta\) is one of the group operators
of Theorems~\ref{rt-actual-lifts}, \ref{rc-actual-lifts},
or a composition of these with the coefficient of
Theorem~\ref{rl-regular-inverse}, then
\begin{equation}\label{jf-commute}
             \Theta J_rF=J_r(\Theta F),\qquad
             \Theta^\dagger J_rF=J_r(\Theta^\dagger F).
\end{equation}
Both \(\Theta F\) and \(\Theta^\dagger F\) belong to the
same \(\mathscr S_\varepsilon\) as \(F\).
\end{lemma}

\begin{proof}
We give the uniform majorant needed for the pointwise statement.
For one outer pair, the inverse coordinates in
Lemma~\ref{jm-cell-haar} are rational functions of the entries
of the group matrix and of its nonzero lower-left entry.
More explicitly, for
\(g=n_Lw(t)\iota(g')n_R\), let \(J\) be the matrix of
the form on the smaller space.  Then
\[
 t=g_{fe},\quad v_L=g_{We}/t,\quad
 v_R^*J=-g_{fW}/t,\quad
 g'=g_{WW}+v_Lt v_R^*J.
\]
The two central parameters are the imaginary parts of
\(g_{ee}/t\) and \(g_{ff}/t\); their real parts are
determined by isotropy.  Iteration proves that the inverse
coordinates on the full open cell are rational, with
denominators products of the nonzero leading minors and
their conjugates.

Fix a compact subset of the radial domain and allow both
compact factors of the Jacquet integral to vary.  The leading
minors are invariant under the integrating groups and are
bounded away from zero on this compact set.  The preceding
inverse formulas imply
\[
  1+\|n_L\|_{\rm coord}+\|n_R\|_{\rm coord}
              \leq C M(k_L^{-1}n_L^{-1}b(c)m n_Rk_R)^a
\]
for some \(C,a>0\).  Coordinates on the nilpotent groups
have Lebesgue Haar measure.  Every fixed derivative of the
integrand is a finite sum of left or right derivatives of
\(F\), multiplied by polynomials in these nilpotent
coordinates with coefficients bounded on the chosen compact
set.  Choosing the decay exponent in \(\|F\|_{A,m}\)
larger than this polynomial degree and the dimension of the
integration space gives an integrable majorant.  Dominated
convergence proves absolute convergence and every derivative
asserted in the lemma.  It also proves smooth compact
equivariance, either directly or from the smoothness of
finite dimensional continuous unitary representations.

The covariance proved in Sections~\ref{rt-section} and
\ref{rc-section}, and for the spectral coefficient in
Section~\ref{rl-section}, says precisely that these operators
commute with the action \(\rho(p)\).  Moving the operator
inside the now justified integral proves the first identity
in \eqref{jf-commute}.  The action \(\rho(p)\) is unitary
for Haar measure, so its commutation with \(\Theta\)
also implies commutation with \(\Theta^\dagger\).
The polynomial bounds for all derivatives of the coefficients
and their adjoints, on \(\mathcal G\geq\varepsilon\),
prove preservation of every seminorm after increasing its
decay exponent and derivative order.  Differential operators
do not enlarge support.  The resulting functions extend
smoothly by zero across the omitted region, since they vanish
on the open set \(\mathcal G<\varepsilon\).
The same dominated argument proves the second identity.
\end{proof}

\begin{lemma}\label{jf-radial-adjoint}
Suppose a differential operator \(\Theta\) commutes with
the left and right unipotent translations and has radial
operator \(\theta\) after multiplication by \(\gamma_r\).
Then its Haar formal adjoint has radial operator
\(\theta^*\), the formal adjoint for \(d\eta_r\).
This applies to all four group operators used below.
\end{lemma}

\begin{proof}
Use the coordinates \((n_L,c,m,n_R)\) before dividing
by the compact groups.  Haar measure is
\[
 dn_L\,dn_R\,
     \gamma_r(c)^2\prod_j|c_j|^{-2}\,dc_j\,dm.
\]
Take a coordinate chart in \((c,m)\).  In this chart any
finite order differential operator has a unique expansion
in ordered products of a fixed invariant frame on each
nilpotent factor and ordinary coordinate derivatives in
\((c,m)\), with smooth coefficient functions.  This follows
by replacing the ordinary derivatives in the nilpotent
coordinates by the invertible invariant frame and proceeding
from the highest order to the lowest order.  Translation
invariance forces the coefficients in this expansion to be
independent of \(n_L,n_R\).

The nilpotent groups are unimodular.  Each real invariant
vertical vector field consequently has formal adjoint equal
to its negative.  On an equivariant function it acts by a
purely imaginary scalar, because the character is unitary;
evaluation on that character therefore preserves the adjoint
operation.  Horizontal integration by parts uses exactly
the remaining measure
\(\gamma_r^2\prod|c_j|^{-2}dc_j\,dm\).
Thus reduction to the character commutes with formal
adjoints for this measure.  Finally multiplication by the
positive function \(\gamma_r\) is an isometry from that
weighted space to \(L^2(d\eta_r)\), and conjugates both
the operator and its adjoint.  This proves the lemma locally,
hence throughout the open cell.

For the actual operators, both reduced adjoints preserve
the compact equivariance conditions: the group operator
commutes with the unitary compact actions, and the displayed
radial operators and their stars preserve the same conditions.
Their equality on that subspace may therefore be checked
against compactly supported equivariant tests.  Such tests
are obtained by compact averaging, and are dense in the
corresponding local space of sections.  Integration by parts
against them shows that two expressions for the same reduced
operator have the same adjoint on this subspace as well.
\end{proof}

\subsection{A four-operator identity}

Assume now \(r\geq1\).  In the radial coordinates set
\[
 p_j=c_j\partial_{c_j},\quad
 \bar p_j=\bar c_j\partial_{\bar c_j},\quad
 T_j(Y)=\begin{pmatrix}Y+p_j&-c_j\\c_j^{-1}&0\end{pmatrix},
 \quad
 \bar T_j(Y)=
       \begin{pmatrix}Y-\bar p_j&-\bar c_j\\\bar c_j^{-1}&0\end{pmatrix}.
\]
Here the derivatives are the Wirtinger derivatives.
Write
\[
 \binom{A(Y)}{B(Y)}=T_{r-1}(Y)\cdots T_0(Y)\binom10,
 \qquad
 \binom{\bar A(Y)}{\bar B(Y)}=
                  \bar T_{r-1}(Y)\cdots\bar T_0(Y)\binom10.
\]
The bar on \(\bar A,\bar B,\bar T\) is part of this
definition; it does not change the order of differential
operators.  Direct integration by parts for \(|c_j|^{-2}dc_j\)
gives
\[
                    p_j^*=-\bar p_j,\qquad c_j^*=\bar c_j.
\]

The actual terminal matrix of Section~\ref{re-section} is
denoted by \(t(Y)\), with
\(\tau_Y=t_{22}(Y)=Gx(Y)\).  Its finite spectral formula is
\(x(Y)=\sigma\sum_\rho p_\rho(Y)P_\rho\), where
\(\sigma=\pm1\), the \(P_\rho\) are orthogonal projections,
and each \(p_\rho\) is a nonzero real polynomial.  Choose
a real \(Y\) for which all these polynomials are nonzero
at \(Y-1,Y,Y+1\).  Such a choice exists, since there are
only finitely many exceptional real numbers.  Put
\[
 \ell=\tau_Y^{-1}t_{21}(Y),\qquad
 \ell^*=\tau_{Y+1}^{-1}t_{12}(Y+1).
\]
The equality in this notation is the proved shifted adjoint
identity of Section~\ref{re-section}.  These terminal
operators commute with all outer coordinate operators.
Define
\begin{equation}\label{jf-four-radial}
 \begin{split}
 \theta_0&=\bar d\,\bar A(Y+1),\qquad
 \theta_3=d\,A(Y),\\
 \theta_1&=d^{-1}\bigl(\ell A(Y)+B(Y)\bigr),\\
 \theta_2&=\bar d^{-1}
               \bigl(\ell^*\bar A(Y+1)+\bar B(Y+1)\bigr).
 \end{split}
\end{equation}
They are differential operators on the radial domain.

\begin{proposition}\label{jf-four-lifts}
There are actual group differential operators
\(\Theta_0,\Theta_1,\Theta_2,\Theta_3\) with radial
operators \(\theta_0,\theta_1,\theta_2,\theta_3\).
They commute with \(\rho(p)\), and they and their Haar
formal adjoints preserve every \(\mathscr S_\varepsilon\).
For their actual Jacquet transforms, both the operator and
adjoint identities hold pointwise on the radial domain.
\end{proposition}

\begin{proof}
Take \(\Theta_0=\overline{\mathcal A}_{Y+1}\) and
\(\Theta_3=\mathcal A_Y\) from
Theorem~\ref{rt-actual-lifts}.  Theorem~\ref{rc-actual-lifts}
supplies the operators with radial parts
\[
 \bar d\,(t_{21}(Y)A(Y)+\tau_Y B(Y)),\qquad
 d\,(t_{12}(Y+1)\bar A(Y+1)+\tau_{Y+1}\bar B(Y+1)).
\]
Compose on the left with the regular group coefficients
of Theorem~\ref{rl-regular-inverse}, whose radial values
are respectively
\(\mathcal G^{-1}x(Y)^{-1}\) and
\(\mathcal G^{-1}x(Y+1)^{-1}\).
Since \(\mathcal G=d\bar dG\), the results are exactly
\(\theta_1,\theta_2\).  In particular, this construction
does not use \(d^{-1}\) as a coefficient on the full group:
the proved full-group coefficient has only powers of
\(\mathcal G\) in its denominator and is smooth even at
vanishing intermediate minors.  Covariance, growth and
support preservation follow from the three cited theorems
and Leibniz's rule.  Lemmas~\ref{jf-pointwise-smooth} and
\ref{jf-radial-adjoint} give the final assertion.
\end{proof}

For \(s\in\mathbb C\) put
\(f_s(c)=\exp(sc_0-\bar s\bar c_0)\), so that
\(|f_s|=1\).  Multiplication by this function is also
denoted by \(f_s\).

\begin{lemma}\label{jf-four-identity}
On the radial domain one has the exact differential identity
\begin{equation}\label{jf-four-equation}
             \theta_0^*f_s\theta_1
                   -\theta_2^*f_s\theta_3=s f_s.
\end{equation}
\end{lemma}

\begin{proof}
Let \(J=\bigl(\begin{smallmatrix}0&1\\-1&0\end{smallmatrix}\bigr)\).
The relation \([p_j,c_j^{-1}]=-c_j^{-1}\) gives
\[
                  \bar T_j(Y+1)^*JT_j(Y)=J.
\]
The only nonconstant entry to check is the upper-left one:
\[
       (Y+1+p_j)c_j^{-1}-c_j^{-1}(Y+p_j)=0.
\]
For \(j=0\), use \([p_0,f_s]=sc_0f_s\) to obtain
\[
             \bar T_0(Y+1)^*f_sJT_0(Y)=f_sJ+s f_sE_{11}.
\]
For every other \(j\), the wave \(f_s\) commutes with
the entries of \(T_j,\bar T_j\).  Successively cancelling
these factors and taking the upper-left entry yields
\begin{equation}\label{jf-telescoping}
          \bar A(Y+1)^*f_s B(Y)
             -\bar B(Y+1)^*f_s A(Y)=s f_s.
\end{equation}
On substituting \eqref{jf-four-radial}, the left side of
\eqref{jf-four-equation} becomes
\[
 \bar A(Y+1)^*f_s(\ell A(Y)+B(Y))
        -(\bar A(Y+1)^*\ell+\bar B(Y+1)^*)f_sA(Y).
\]
The two terms involving \(\ell\) cancel because \(\ell\)
commutes with \(f_s\).  The remainder is
\eqref{jf-telescoping}.  This calculation retains the
order of every pair of noncommuting operators.
\end{proof}

\subsection{Rapid Fourier decay}

Let \(D_1(g)\) be the first leading minor.  It is a linear
polynomial in the entries of \(g\), invariant under the two
integrating groups, and satisfies \(D_1(b(c)m)=c_0\).
Set
\[
       W_sF(g)=e^{sD_1(g)-\bar s\overline{D_1(g)}}F(g).
\]
For each fixed \(s\), both \(W_s\) and multiplication
by any polynomial in \(D_1,\overline{D_1}\) preserve
\(\mathscr S_\varepsilon\).  This follows by differentiating:
each derivative is the same unit-modulus exponential times
a polynomial in the matrix entries.  Invariance gives
\begin{equation}\label{jf-wave}
                        \Phi_{W_sF}=f_s\Phi_F.
\end{equation}

\begin{theorem}\label{jf-fourier-decay}
For the actual place data and any \(F,H\in\mathscr S_+\),
the function
\[
                       s\longmapsto\beta(W_sF,H)
\]
is a Schwartz function on \(\mathbb C=\mathbb R^2\).
In particular, for every integer \(N\geq0\),
\begin{equation}\label{jf-decay-bound}
        |\beta(W_sF,H)|\leq C_N(F,H)(1+|s|)^{-N}.
\end{equation}
All constants are finite and depend only on the indicated
functions, the fixed place data and the chosen operators.
\end{theorem}

\begin{proof}
Apply \eqref{jf-four-equation} to \(\Phi_F\).  By
Proposition~\ref{jf-four-lifts}, Lemma~\ref{jf-pointwise-smooth}
and \eqref{jf-wave}, its left side is the radial Jacquet
transform of
\[
           \Theta_0^\dagger W_s\Theta_1F
                        -\Theta_2^\dagger W_s\Theta_3F.
\]
Every function here belongs to \(\mathscr S_+\).
Pair with \(\Phi_H\), use \eqref{jf-radial-pairing},
and apply the proved actual adjointness of
Corollary~\ref{jp-poly-adjoint}.  All its hypotheses hold:
the operators commute with \(\rho\), all derivatives of
their coefficients and adjoint coefficients have polynomial
growth on the common positive Gram support, and every input
has derivatives of arbitrarily fast decay.  We obtain
\begin{equation}\label{jf-beta-recursion}
 s\beta(W_sF,H)
       =\beta(W_s\Theta_1F,\Theta_0H)
                    -\beta(W_s\Theta_3F,\Theta_2H).
\end{equation}
No integration by parts on a noncompact radial function has
been assumed in this step; it is supplied by the group
adjointness theorem with its proved integrable majorants.

For the initial uniform bound, use \eqref{jf-wave} and
Cauchy--Schwarz:
\[
 |\beta(W_sF,H)|
 \leq C_r^{-1}\|\Phi_F\|_2\|\Phi_H\|_2.
\]
This bound is independent of \(s\).  Define recursively
\[
 K_0(F,H)=C_r^{-1}\|\Phi_F\|_2\|\Phi_H\|_2,\qquad
 K_{n+1}(F,H)=K_n(\Theta_1F,\Theta_0H)
                           +K_n(\Theta_3F,\Theta_2H).
\]
These numbers are finite, by preservation of \(\mathscr S_+\).
Induction in \eqref{jf-beta-recursion} gives
\( |s|^n|\beta(W_sF,H)|\leq K_n(F,H)\).
Together with the bound for \(n=0\) and
\((1+x)^n\leq2^n(1+x^n)\), this proves
\eqref{jf-decay-bound}.

To include derivatives in \(s\), write \(s=u+iv\) and
\(D_1=x+iy\).  The phase is \(2i(uy+vx)\).
Each real derivative therefore replaces \(F\) by a
polynomial in \(D_1,\overline{D_1}\) times \(F\).
These functions belong to \(\mathscr S_+\).
Their radial transforms are square integrable, so
Cauchy--Schwarz and dominated convergence justify all these
derivatives of the pairing.  Applying the just proved
estimate to each differentiated input proves that every
derivative decays to every order.  This is the Schwartz
assertion.
\end{proof}

\subsection{Compact inputs and smooth lower-rank families}

The following argument provides pointwise regularity of
the marginal densities; a Fourier representative alone
would only give almost everywhere information.
Let \(r\geq1\), and decompose
\(N_r=U_r\rtimes N_{r-1}\), where \(U_r\) is the
one-pair nilpotent group of Lemma~\ref{ju-heisenberg}.
For \(F\in C_c^\infty(G_r,E)\cap\mathscr S_\varepsilon\)
and \(t\ne0\), define a coefficient on \(G_{r-1}\) by
\begin{equation}\label{jf-slice-family}
 \begin{split}
 F_t(g)={}&|t|^{2r+k}
   \int_{U_r\times U_r}
       \psi_z(u_L)\overline{\psi_{z'}(u_R)}\\[-2pt]
       &\hspace{34mm}\cdot
          F(u_L^{-1}w(t)\iota(g)u_R)\,du_L\,du_R.
 \end{split}
\end{equation}
The characters are the restrictions of the original
unipotent characters.  Use recursively compatible Haar
measures; other choices only insert fixed positive constants.

\begin{proposition}\label{jf-compact-family}
The function \((t,g)\mapsto F_t(g)\) is smooth on
\(\mathbb C^\times\times G_{r-1}\).  On every compact
\(K\subset\mathbb C^\times\), all \(F_t\) have support
in one compact subset of \(G_{r-1}\), and there is
\(\varepsilon_K>0\) such that \(F_t=0\) on
\(\mathcal G_{r-1}<\varepsilon_K\) for every \(t\in K\).
All their parameter derivatives have the same two properties.
The exact radial factorization is
\begin{equation}\label{jf-factorization}
                  \Phi_F(t,c',m)=\Phi^{(r-1)}_{F_t}(c',m).
\end{equation}
\end{proposition}

\begin{proof}
Apply the inverse coordinate formulas in the proof of
Lemma~\ref{jf-pointwise-smooth} to
\(u_L^{-1}w(t)\iota(g)u_R\) in \(\operatorname{supp}F\).
When \(t\) remains in \(K\), its modulus has a positive
lower bound.  All inverse coordinates \(u_L,u_R,g\) are
then bounded uniformly.  They lie in closed subsets of
finite dimensional matrix spaces, so they lie in fixed
compact sets.  Thus the integrand is supported in one
compact subset of the integration variables and \(g\),
locally uniformly in \(t\).  Differentiation under this
compact integral proves smoothness and the support assertion.
Derivatives do not enlarge this support; taking a slightly
larger compact neighborhood of \(K\) gives the assertion
also at its boundary.

Leading minors in these coordinates satisfy
\(D_{j+1}^{(r)}=tD_j^{(r-1)}\).  In particular,
\[
 \mathcal G_r(u_L^{-1}w(t)\iota(g)u_R)
                                  =|t|^2\mathcal G_{r-1}(g).
\]
Consequently one may take
\(\varepsilon_K=\varepsilon/\max_{t\in K}|t|^2\).
For derivatives, use a compact neighborhood of \(K\) in
this formula.  A smooth function which vanishes on the
corresponding open set has all derivatives zero there.

The semidirect-product Haar formula and the character
identity factor the full Jacquet integral into the outer
integral in \eqref{jf-slice-family} followed by the smaller
Jacquet integral.  To check the ordering, move the two
\(N_{r-1}\) variables and the compact factors across
\(w(t)\), which commutes with the embedded smaller group.
Conjugate \(u_L,u_R\) by those factors.  Their action on
\(U_r\) preserves Haar measure: the unipotent factors
have determinant one, and the compact factors act unitarily.
It preserves the restricted characters because the original
character is a homomorphism and is invariant under the
compact stabilizers.  No Jacobian or phase remains.
All integrals at a regular cell point have compact support
in their nilpotent coordinates, so Fubini is legitimate.
Finally
\(\gamma_r(t,c')=|t|^{2r+k}\gamma_{r-1}(c')\),
which proves \eqref{jf-factorization} including its scalar
normalization.
\end{proof}

\begin{corollary}\label{jf-smooth-slices}
For compactly supported \(F\) as above, the map
\[
       t\longmapsto\Phi_F(t,\cdot,\cdot)
       \quad\text{from }\mathbb C^\times
                  \text{ to }L^2(d\eta_{r-1},E)
\]
is smooth.  Its derivatives are its pointwise derivatives.
The same holds after any finite word in
\[
       t\partial_t,\quad\bar t\partial_{\bar t},\quad
             t,\quad t^{-1},\quad\bar t,\quad\bar t^{-1},
\]
and after multiplication by any fixed wave
\(e^{st-\bar s\bar t}\).  All corresponding Hilbert norms
are bounded locally uniformly for \(t\ne0\).
\end{corollary}

\begin{proof}
On a compact parameter neighborhood, the family in
Proposition~\ref{jf-compact-family} and all its derivatives
have common compact support and a common positive Gram
bound.  Taylor's formula with its integral remainder,
differentiated a fixed number of times on the smaller group,
shows that the family is smooth for every seminorm
\(\|\cdot\|_{A,m}\).  The Jacquet map into
\(L^2(d\eta_{r-1},E)\) is continuous for
\(\|\cdot\|_{A,4}\), by
Theorem~\ref{jp-pairing-equality} and its bound.
It therefore sends this smooth family to a smooth Hilbert
space valued family, with the asserted derivatives.
Pointwise differentiation is also valid by the compact
integral argument.  All six displayed operators and the
wave act on the parameter by smooth coefficients on
\(\mathbb C^\times\), which proves the remaining claims.
\end{proof}

\begin{theorem}\label{jf-marginal-schwartz}
For \(F,H\in\mathscr S_+\), the integrable marginal density
\begin{equation}\label{jf-density}
 D_{F,H}(t)=|t|^{-2}
     \int\langle\Phi_F(t,c',m),\Phi_H(t,c',m)\rangle
                                    \,d\eta_{r-1}(c',m)
\end{equation}
has a Schwartz representative on \(\mathbb C\).
If \(F,H\) are compactly supported, the integral in
\eqref{jf-density} exists for every \(t\ne0\), is smooth
there, and agrees there with that representative.  It
therefore has a unique smooth extension across \(t=0\),
and that extension and all its derivatives decay faster
than every inverse power at infinity.
\end{theorem}

\begin{proof}
Fubini and Cauchy--Schwarz show that the density exists
almost everywhere and belongs to \(L^1(\mathbb C)\).
Its Fourier transform, using the real bilinear phase
\(2\operatorname{Im}(st)\), is
\[
 \int_{\mathbb C}e^{st-\bar s\bar t}D_{F,H}(t)\,dt
                             =C_r\beta(W_sF,H).
\]
Theorem~\ref{jf-fourier-decay} makes this a Schwartz
function.  Ordinary Fourier inversion gives a Schwartz
representative of \(D_{F,H}\).  To recall why the assertion
includes every weight and derivative, differentiation of
the inverse transform inserts a polynomial in the Fourier
variable, and multiplication by a coordinate is integration
by parts in that variable.  All resulting integrals converge
absolutely for a Schwartz Fourier transform.  Approximate
identities identify the inverse transform with the original
\(L^1\) density almost everywhere.

For compact \(F,H\), Corollary~\ref{jf-smooth-slices}
shows that the inner product in \eqref{jf-density} is a
smooth function of \(t\ne0\), with all derivatives given
by the finite Leibniz sums of Hilbert inner products.
It is therefore continuous and agrees almost everywhere
with the smooth Fourier representative.  Two continuous
functions on an open subset of the plane which agree almost
everywhere agree everywhere: otherwise their difference
would stay nonzero on a ball of positive measure.
This proves the claimed pointwise extension and all its
decay properties.
\end{proof}

\subsection{Terminal support and the remaining application}

For compactly supported \(F\), terminal support is also
locally compact in the outer coordinates.  Indeed, on any
compact subset of \((\mathbb C^\times)^r\), the quantity
\(|d(c)|\) is bounded above and away from zero.  The invariant
minor \(D_{r+1}=d(c)\zeta\) is bounded on
\(\operatorname{supp}F\), and
\(\mathcal G=|d(c)|^2G\geq\varepsilon\) on any contributing
orbit.  Thus \(|\zeta|\) is bounded and \(G\) is bounded
away from zero.  By the compact double-coset description
in Section~\ref{jg-section}, the allowed terminal matrices
lie in a compact subset of \(G_0\).  Every terminal
differential operator and every smooth scalar multiplier
preserves this support.  It follows that all finite terminal
operator words applied to these vectors belong to terminal
\(L^2\), and their actual Haar formal adjoints may be used
by ordinary compactly supported integration by parts.
The assertion includes \(r=0\), when there are no outer
variables and no Fourier-decay statement is required.

The source's generated compact class starts with these
vectors, applies finite differential words in the current
outer coordinate when passing to the next level, and
multiplies by unit-modulus waves.  Repeating
Proposition~\ref{jf-compact-family} proves its actual local
slice regularity and terminal compactness at every fixed
nonzero outer parameter.  This also proves the local
Leibniz formulas for its partial pairings.  More explicitly,
after fixing the first \(i\) nonzero coordinates, inductively
represent the vector as the radial Jacquet transform of a
compactly supported smooth coefficient on \(G_{r-i}\).
The coefficients depend smoothly on those fixed parameters,
with common compact support and positive Gram bound on
each compact parameter set.  A wave in the next coordinate
is multiplication of that coefficient by the corresponding
wave in its first leading minor, which preserves these
properties.  Every entry of the source's matrix \(bT\)
is a finite word in the six displayed coordinate operators.
Applying such an entry and then fixing that coordinate is
handled by differentiating the family
\eqref{jf-slice-family}; Proposition~\ref{jf-compact-family}
and Corollary~\ref{jf-smooth-slices} preserve exactly the
inductive properties.  Finite linear combinations do so
as well.  At \(i=r\), a terminal wave is a smooth multiplier
on the already compact terminal support.  This proves the
claim by induction, including the empty sequence of steps.

For the full residue matrices, one further algebraic step
must be kept distinct: all four entries of the initial
matrix \(Q_0(Y)\) must be identified with their full-group
operators, and that identification must intertwine the
successive one-pair restrictions.  The four operators
\(\Theta_j\) proved above do not by themselves establish
that larger matrix identity.  Once it is established, its
polynomial differential operators preserve compact support
and the positive Gram cutoff, so the present theorems apply
to every finite product of those entries.  Until that
identification is supplied, this section proves the actual
Fourier decay, pointwise smooth densities, generated local
slice regularity and terminal integrability stated above;
it does not assert all of the remaining residue-pairing
identities by assumption.


\section{The remaining mathematical construction}
\label{sec:remaining}

The preceding results give a complete analytic implication from
\eqref{eq:remaining-estimate} to Dirichlet GRH, an explicit quadratic
determinant inequality with its consequences for real zeros, and
conditional simplifications that follow from a uniform quasi-Riemann
hypothesis. Section~\ref{cs-section} also proves the abstract slope
comparison directly, including its passage from matrix estimates.
Section~\ref{qs-section} proves the complete finite-family
selected-character estimate under the stated QRH hypothesis.
Section~\ref{sec:ct-intervals} proves the cyclotomic descent,
uniform generator bound, and common-interval estimate, the last
initially conditional on \eqref{ct-coarse-hypothesis}, which
Theorem~\ref{kt-fourteen} now supplies from the stated classical
$K$-theory and duality inputs.
Section~\ref{ht-section} proves the homological deduction from the
stated geometric input and the passage from a strict Selmer estimate
to the coarse $H^2$ estimate. Section~\ref{ac-section} constructs
that complex and proves its exactness, stabilizer bound, and full
orbit estimate, so the arithmetic homology bound is unconditional.
Sections~\ref{sl-section}--\ref{sp-section} prove the class-number
and unit estimates, the comparison with equivariant ideal-class
characters, and the cyclotomic enlargement and descent.
Section~\ref{kt-section} supplies a shorter independent proof of
the coarse $H^2$ and strict Selmer bounds.
Section~\ref{ms-section} also proves the full opposite-parity
Selmer estimate, including primes dividing the field degree,
from the stated classical special-value and duality theorems.
Section~\ref{sz-selmer-section} proves the exceptional trivial-character
case of the complementary parity by a direct residual class-group
injection. Section~\ref{iv-section} proves the fixed-conductor
case by restriction to a fixed cyclotomic field, with an exact
invariant-module bound and an additive common-interval bound.
Section~\ref{sv-section} proves the bounded-order exponent-one case
and its uniform character sum. The other automorphic character
components remain a separate problem. Section~\ref{fcmp-section}
proves the whole-field comparison with exactly one copy of each
normalized character length and a polynomial torsion error at the
fixed primes dividing the original character order.
Section~\ref{av-section} proves the actual adjoint-vanishing implication
from Thorne's classical theorem, including coefficient descent,
polarization and the exact cyclotomic-kernel irreducibility condition.
Automorphy and quantitative integral congruence bounds are not
consequences of rational vanishing and remain separate.
Sections~\ref{sf-section}--\ref{sw-section}
also prove the original finite-field homology route: explicit transfer
functionals, integral arithmetic cycles with prescribed residue values,
the odd-degree rank bound, and the finite-image estimate establish
assertion~\ref{sc-few}. The resulting exponents $86$ and $88$ follow
from the stated classical inputs. The independent $K$-theory argument
retains the stronger exponent $14$.
Sections~\ref{sd-section} and \ref{sh-section} also prove the
degree-one-prime detection theorem and the actual cyclic bar
cycle and its relation used in the original construction.
Sections~\ref{si-section} and \ref{se-section} prove the exact
isomorphism of ideal lattices with prescribed determinant and
the integral interpolation of cyclic characters needed to
assemble those cycles.
Section~\ref{sec:selection-assembly}
checks the order of choices, conductor bounds, field degrees, and
rescaling of the interval.
Section~\ref{coef:section} proves the integral coefficient-lattice
duality, including determinant characters and transported lattices.
Section~\ref{fh-section} computes the exact Fock normalizer as a
product of hook lengths. It gives a clearing integer and proves the
actual equality of Weyl and dual-Weyl lattices at the selected primes
without using a modular simplicity theorem.
Sections~\ref{sm-section}--\ref{ib-section} prove integral spanning
by products of initial minors, the compact branching law, and the
exact normalized integral branching maps. The scalar that makes these
maps integral is retained; an explicit example shows that deleting it
would invalidate integrality.
Section~\ref{div:section} derives the divisor formula, the actual
finite-group dual exponent, and the chain pairing bounds.
Section~\ref{am-section} constructs the coefficient maps and proves
their exact elementary-divisor and singular-value bounds, including
their normalized metrics and determinant-frame cancellation.
Section~\ref{al-section} constructs the standard decorated lattice,
computes the kernel lattice's degree, proves the actual two-prime
stabilizer condition, and supplies short independent Hermitian dual
vectors by an explicit Gaussian transference argument.
Sections~\ref{acl-section} and \ref{acc-section} identify the actual
adelic classes with their decorated lattice coordinates and prove the
additional coset partition, including its stabilizers and Haar weights.
Section~\ref{ip-section} proves the exact integral intersection
under coefficient-field extension and the passage of the period
valuation bound to all test vectors.
Section~\ref{cp-section} proves the initial finite class-period bound
with its automorphism denominator, and Section~\ref{atl-section}
constructs the integral lowering maps for every auxiliary character.
The latter includes the exact conjugate-uniformizer identity and
descent from complex evaluations to the rational integral lattice.
Sections~\ref{ot-section} and \ref{ct-section} prove the ordinary
moment criterion and the three one-dimensional CM theta laws at
the actual large split primes. The proof fixes one unit period,
retains cancellations in arbitrary finite linear combinations,
and checks the exact selected-torsion phase. The integral theory
of the good ordinary CM model is an explicitly cited classical
input. Sections~\ref{lm-section} and \ref{tc-section}
prove the finite measure deductions, the determinant valuation
of the normalizing scalar, and the exact analytic change of
theta lattice. The latter includes the finite quadratic Gauss
sum, its eighth power, and the convergent Gaussian trace.
Sections~\ref{ci-section} and \ref{ts-section} extend the one
ordinary period to nonprincipal ideal bases and prove simultaneous
local standardization of the fourfold lattice. They establish
integral moments of the normalized self-dual lattice vector at
every split prime at least five prime to the ideal norm.
Section~\ref{te-section} proves the actual arithmetic
equivariance, with both lattice moves and the exact determinant
characters. The simultaneous local hyperbolic decomposition,
including at dyadic places, proves the required positive sign
of the finite quadratic sum.
Section~\ref{cv-section} supplies the finite-dimensional
valuation argument and its iteration with the precise quadratic
loss. Sections~\ref{gm-section} and \ref{csp-section} prove the actual
higher-dimensional moment laws and both matrix supports, retaining
the exact denominator and period coordinates. Section~\ref{hec-section}
proves the incidence identity, the exact change of theta normalization,
and the weighted Hecke commutation law. Section~\ref{jcf-section} constructs one
finite coefficient field for all the prescribed jets, with the prime,
level and derivative data fixed before adjoining coefficients.
Section~\ref{tvc-section} identifies the actual polynomial kernel,
including its ordinary mixed derivatives and the harmless target
projection, and combines these results to prove the selected-prime
column bound on the prescribed modified initial lattice. Seed
nonvanishing and the auxiliary-level-prime bounds remain separate.
Section~\ref{frg-section} proves the
scalar rigidity of the relevant Hermitian flags.
Section~\ref{df-section} constructs the common positive complement
for relevant degenerate flags and proves off-diagonal vanishing
by Schur's lemma. Its quotient argument does not require the
compact group to preserve the arithmetic stabilizer lattice.
Absolute convergence remains an explicit hypothesis.
Section~\ref{gv-section} computes exact negative moments of
projected Gram determinants, proves the optimal small-threshold
exponent, and improves both orbital volume estimates.
Section~\ref{ods-section} proves a stronger final dyadic summation.
The coordinate-size reduction and the weighted lattice-pair count
are separate inputs to the diagonal degenerate estimate.
Section~\ref{sec:theta-local-audit} supplies the local Gaussian
and stationary-phase formulas, convolution normalization,
unfolding, and the final combination of orbital estimates.
Section~\ref{sec:theta-orbital-identity} proves finiteness,
uniform local constancy and inverse-parameter support for the
actual local kernel. It also proves the truncated orbital identity
from the explicit germ formula, including cancellation of every
mixed critical term. Proposition~\ref{gu-uniform-germ} reduces the
germ formula, uniformly over arbitrary comparison configurations,
to scalar identities against finitely many fixed marked modules.
Sections~\ref{sec:theta-scalar-germ} and \ref{hr-section} prove
those scalar identities by a Hessenberg determinant calculation
and a finite-module reduction. Section~\ref{tp-section} proves
the single-block Parseval identity, averaged diagonal bound, and
the complete local majorant induction from the now-established germ.
Section~\ref{gt-section} proves a stronger smoothed ideal sum and
its application to the actual period, trace-dual, and lift lattices,
their exact covolumes at the unramified levels, and a Poisson formula
under the regularity and decay conditions used in the application.
Section~\ref{as-section} proves the archimedean polynomial and
width calculations, with a stronger numerical boundary estimate;
the source-space dimension and formal-degree formulas used there
are now proved in Sections~\ref{cd-section} and \ref{fd-section}.
The latter identifies both the source's Haar measure and its
lowest-compact-type condition, using the stated classical
discrete-series theorems. Section~\ref{md-section} supplies the exact mass
recurrence from the stated classical mass theorem, proves all
primitive local densities, and checks the field, discriminant, Euler,
and dyadic conventions.
Section~\ref{amb-section} proves the rational unipotence,
lattice separation, distant-orbit sum and elementary coefficient
decay, identifying the lowest compact-type integral with Sun's
classical formula. Section~\ref{ak-section} now constructs the
finite-rank projection and proves its continuous kernel and exact
pointwise trace, using the stated classical automorphic finiteness
theorem. Section~\ref{rv-section} proves the rare-volume estimate
by counting one primitive isotropic sublattice, with all Haar
constants retained. Section~\ref{nu-section} proves the integrated
near-unipotent estimate by a canonical isotropic flag and block packing.
Section~\ref{dc-section} proves the deep-cusp estimate: Euclidean
transference gives a small radical domain, an explicit local Sobolev
bound controls evaluation, and repeated zero constant terms give
exponential decay. Corollary~\ref{dc-ambient-comparison} completes
the ambient multiplicity comparison with the actual parameters.
Sections~\ref{ca-section} and \ref{cg-section} prove the archimedean
coefficient identities, extraction inequalities, and separated
interpolation sets, with improved numerical constants.
Sections~\ref{rp-section} and \ref{rd-section} prove the terminal
radial operator identity, positivity, and the actual finite orthogonal
decomposition. These are pointwise identities; the remaining
noncompact integral passage is identified separately.
Section~\ref{rf-section} proves that passage under an explicit weighted
Fourier-decay hypothesis, including the extension across zero and the
vanishing of every positive-order derivative integral. It also proves
the actual one-coordinate differential identities and their expansion
over subsets of an admissible word. Section~\ref{jf-section} now proves
the Fourier decay for the actual Jacquet transforms and the local
slice regularity of the compact generated class. The identification
of the full initial matrix and its complete restriction chain remains
separate from the four operators used in that Fourier argument.
Sections~\ref{rl-section}--\ref{rc-section} construct all four
polynomial differential lifts and the regular spectral inverse,
including their global covariance and exact radial actions.
The cofactor construction remains regular where an intermediate
minor vanishes. Sections~\ref{rq-section} and \ref{rr-section}
prove the terminal quantum determinant, every terminal commutation
relation, and both shifted adjugate identities, including
exceptional parameter values. Section~\ref{re-section} constructs
the initial first-order connection and proves both recursive frame
identities, their precise domains, and the resulting positive form
for compactly supported vectors. Section~\ref{jf-section} supplies
actual analytic adjointness and Fourier decay for the indicated
group operators. The complete residue-pairing application still
needs the full initial matrix and its successive restriction identity.
Section~\ref{jk-section} proves the actual coordinate integrals
used for those estimates. The basic integral is evaluated exactly
with necessary and sufficient convergence conditions; the weighted
estimate retains the nonzero character coefficient and all strict
exponent bounds.
Section~\ref{ju-section} proves the complete weighted unipotent
induction, its Gram determinant and Haar coordinates, and the
nonvanishing character coefficient supplied by the negative norm
of the terminal vector.
Section~\ref{jm-section} proves the full-group integral, including
its open-cell Haar density and the compact rank-zero case. It gives
the actual separated coefficient bound and the continuous majorant
with both required compact-stabilizer integrals.
Sections~\ref{jg-section} and \ref{jp-section} prove the geometric
integration formula, regularized pairing, averaged unfolding and
adjointness. The fifth derivative gives the necessary equality at
every positive character. Section~\ref{jf-section} verifies the
operator, regularity and support hypotheses for the four Fourier
lifts. Smooth compact families on successively smaller groups give
pointwise marginal densities and terminal compactness at each fixed
nonzero outer parameter. The remaining full-matrix identification
and integrated residue-form assembly are explicitly distinguished.
These results do not yet give a conventional proof of
\eqref{eq:remaining-estimate}. This section specifies the remaining
work without disguising it as a standard theorem.

For a fixed primitive character $\chi$, the construction has to produce
intervals $[Y,Y^b]$, with fixed $b>1$ and arbitrarily large $Y$, on
which the positive combination \eqref{eq:Hchi} has the stated upper
bound for every prescribed positive error exponent. All constants
that multiply that exponent must remain bounded independently of
the exponent. Constants that merely change the starting point of
an interval may depend on it. The actual functions $G_\psi$ and
their normalizing logarithmic derivatives must be those defined in
\eqref{eq:Gchi}--\eqref{eq:Gone}.

The announcement proposes chains of theta lifts between definite
unitary groups. Here a definite unitary group means the group of
isometries of a Hermitian space over a CM field whose Hermitian form
is definite at every infinite place. A CM field is a totally imaginary
quadratic extension of a totally real number field. A theta lift is
an automorphic integral formed from the Weil representation for a
pair of such groups. The chain length is allowed to grow with $x$.
These definitions describe the objects; they do not assert uniform
estimates for the integrals.

A full reconstruction must provide the following four arguments
with compatible normalizations and uniformity.
\begin{enumerate}
\item Construct the integral modules of theta lifts, specify their
      coefficient rings, pairings, levels, weights, and measures,
      and prove the requisite nonvanishing. Iteration must be
      justified when the chain length is of order $\sqrt{x}$.
\item Prove the upper bounds for the archimedean norms and the lower
      bounds for divisibility of the same integral objects. Compare
      the covolumes, lattice indices, and projections with all
      denominators retained. Separate bounds for unrelated lattices
      do not suffice.
\item Complete the negative-twist Galois cohomology estimates in the
      varying fields. The relevant groups impose vanishing on
      full decomposition groups at the coefficient prime, not only
      an unramified condition. The common-interval mechanism is now
      proved in Section~\ref{sec:ct-intervals}, and its coarse input
      follows from Theorem~\ref{kt-fourteen}. The opposite-parity
      estimate is proved in Theorem~\ref{ms-main}. The remaining
      complementary-parity components need constants uniform in
      the allowed auxiliary primes.
\item Use the character averages and the preceding estimates to
      bound the particular fixed function $H_\chi$. In this step
      the selected-character estimates are supplied, conditional
      on QRH, by Theorem~\ref{qs-selection}, for the entire finite
      family of fixed base characters and their quadratic twists.
      The scale choices and the common interval are compatible
      by Section~\ref{sec:selection-assembly}. One must still verify
      that the losses from
      exceptional subspaces and changes of integral structure
      fit the numerical bounds.
\end{enumerate}

The separate audit contains independent proofs of some algebraic
and cohomological ingredients, including a Bernoulli-number unit
criterion, a restricted cyclotomic Selmer-group vanishing result,
and the existence of regular semisimple elements in the relevant
four-dimensional cosets. Those proofs are useful completed pieces.
They do not settle the uniform constructions just listed.

The most concrete next targets are the remaining actual orbital-sum estimates
\eqref{th-four-estimates} and \eqref{th-diag-input}, and
\eqref{cs-required-divisor} for the integral maps they define.
The local germ, its compact-parameter uniformity, and the averaged
diagonal estimate have now been proved; the positive-majorant
argument is complete in Section~\ref{tp-section}. The next
analytic task is the uniform global passage from these local
majorants to the particular orbital sums and their energy bounds.
The ideal-summation part is now supplied by Section~\ref{gt-section};
the exact trace/Fourier normalization and Poisson summation are
also proved there. The archimedean width calculation is proved in
Section~\ref{as-section}, with the mass, column-dimension and
formal-degree identities supplied in the preceding sections.
The archimedean coefficient identities and the
pointwise radial decomposition are now proved in
Sections~\ref{ca-section}--\ref{rd-section}. The comparison of
ambient multiplicity with its mass/formal-degree leading term is now
proved in Corollary~\ref{dc-ambient-comparison}. The noncompact
integrated radial identity for the actual generated functions and
the remaining orbital estimates still require proof.
The quotient metrics, local dual exponent, and formal divisor
comparison are now proved in Section~\ref{div:section}.
The remaining construction must establish the sharp level-prime
divisibility and the archimedean operator-norm bounds for the same
theta maps. The initial period bound and the all-twist lowering map
are now established. The exact finite valuation iteration and
one-dimensional theta laws combine with the proved incidence relation
for the same columns. The higher-dimensional
theta laws and their supports are now proved directly in
Sections~\ref{gm-section} and \ref{csp-section}.
The polynomial coefficient dualities themselves
are now proved in Section~\ref{coef:section}.
The coefficient-map image ideals and archimedean bounds themselves
are proved in Section~\ref{am-section}; their degree and short-dual
inputs are supplied by Section~\ref{al-section}. The exact integral
intersection used for period tests is supplied by Section~\ref{ip-section}.
After the remaining period and spectral bounds are supplied,
Lemma~\ref{cs-combination} performs the cancellation of the local
factors without a further automorphic argument. In parallel, the
remaining cohomological bounds and the global theta sums must be
proved with constants compatible with the growing number of lifts.

The proposed simplifications have a definite scope. The uniform
quasi-Riemann theorem makes a fixed weaker Dirichlet zero-free
condition automatic and supplies both the ordinary progression-prime
choices and the stronger simultaneous selection in
Section~\ref{qs-section}. The quadratic determinant argument
removes an exceptional-real-character restriction when only the
logarithmic derivative at $1$ is needed. Neither statement covers
arbitrary Hecke characters of the varying CM fields. Neither gives
the norm and divisibility estimates for growing chains of theta
lifts. The exact prime-removal identity
\eqref{eq:qrh-prime-removal} is a possible analytic alternative,
and Proposition~\ref{prop:qrh-exponent-improvement} improves its
deletion error by retaining previously known M\"obius cancellation.
The shorter mean-square estimate required at each improved exponent
has not been proved here. Its validity cannot be inferred from the
existing estimates in their longer ranges.

\begin{example}[A fixed zero-free region and symmetry do not suffice]
Let $\vartheta>1/2$, choose
$0<a<\min(\vartheta-1/2,1/2)$, and put
\[
 P(s)=\bigl((s-\tfrac12-a)^2+1\bigr)
       \bigl((s-\tfrac12+a)^2+1\bigr).
\]
Then $P(1-s)=P(s)$ and
$P(\bar s)=\overline{P(s)}$. It has no real zero and no zero
with $\Re s>\vartheta$, but has four zeros off the critical
line, at $1/2\pm a\pm i$. Every assertion follows by direct
factorization. This polynomial is not a Dirichlet $L$-function
and is not offered as a counterexample to GRH. It shows that
a fixed zero-free region, absence of real zeros, and the usual
reflection symmetries are logically insufficient by themselves.
The arithmetic features of the $L$-functions must supply the
remaining argument.
\end{example}

The next substantial advance in this manuscript should therefore
be a proof of one of the four uniform constructions above, or a
replacement of them by an independently proved estimate strong
enough to yield \eqref{eq:remaining-estimate}. Until then, the
conditional statements retain their stated hypotheses and this
document must not be cited as a completed proof of GRH.

\begin{thebibliography}{99}
\bibitem{alpo-announcement}
L. Alp\"oge, announcement of a proof of RH and Dirichlet GRH,
20 November 2026, 4 pp.
\url{https://alpo.ge/remark.pdf}.
The source archive consists of the fourteen parts linked there.

\bibitem{dlmf-dirichlet}
NIST Digital Library of Mathematical Functions,
\S25.15, Dirichlet $L$-functions, especially equations
25.15.1--25.15.5 and the discussion of zeros.
\url{https://dlmf.nist.gov/25.15}.

\bibitem{dlmf-zeta}
NIST Digital Library of Mathematical Functions,
\S25.2, Definition and Expansions of the Riemann zeta function,
especially equations 25.2.1 and 25.2.4.
\url{https://dlmf.nist.gov/25.2}.

\bibitem{dlmf-zeros}
NIST Digital Library of Mathematical Functions,
\S25.10, Zeros of the Riemann zeta function.
\url{https://dlmf.nist.gov/25.10}.
% Insert these items inside the host document's bibliography.
\bibitem{qrh-september}
OpenAI, \emph{The Quasi-Riemann Hypothesis: A Zero-Free Half-Plane
$\operatorname{Re}s>7/8$}, 30 September 2026, 199 pp.
Version at repository commit
\path{4eaee87548250b2682affda7d05d947c05554bae}.
Theorem 1.1; Proposition 2.1; Sections 17--20.
\url{https://github.com/openai/math/blob/4eaee87548250b2682affda7d05d947c05554bae/preprints/The-Quasi-Riemann-Hypothesis-September-30-2026/paper.pdf}.

\bibitem{qrh-october}
OpenAI, \emph{The Quasi-Riemann Hypothesis}, 5 October 2026,
48 pp. Version at the same repository commit.
Theorem 1.1, Proposition 3.1, Section 3, Propositions 5.1, 5.2,
5.4, and Appendix A.2.
\url{https://github.com/openai/math/blob/4eaee87548250b2682affda7d05d947c05554bae/preprints/The-Quasi-Riemann-Hypothesis-October-5-2026/paper2.pdf}.

\bibitem{qrh-catalogue}
H. Collet and contributors, \emph{QRH Catalogue}, repository
commit \path{5b7217781269248f39fbcec2432394d3d1610344},
consulted 20 November 2026. In particular the README's evidence
qualifications and the mathematical entries G276 and G280.
\url{https://github.com/HaileyStorm/QRH-Catalogue/tree/5b7217781269248f39fbcec2432394d3d1610344}.
This reference identifies a catalogue of source claims and
conditional reductions; it is not a certificate for every
listed theorem.

\bibitem{alpo-prime-supply}
L. Alp\"oge, submitted source archive, the ordinary progression
and quadratic-field prime conditions in
\path{Formalism/SharpRestriction/InertJacobiLevels/PrimeSupply.lean};
the varying quadratic ray-class conditions in
\path{Formalism/SharpRestriction/GlobalSchurBounds/ShortIntervalPrimes.lean}.
Archive SHA-256:
\path{9631b4f020199e580ba86178cb56c6e9b09f690c2369a80dc4334aa93f49c7fc}.

\bibitem{alpo-final-scalar}
L. Alp\"oge, same submitted archive,
\path{Formalism/SizeWindowBrackets/GrowthBoundGivesRiemannHypothesis.lean},
in particular the height hypotheses, the conditions on the
three auxiliary characters, and the final size inequalities.
\bibitem{sz-openai}
OpenAI, \emph{Uniform exclusion of Landau--Siegel zeros},
October 1, 2026, 9 pp.
\url{https://github.com/openai/math/blob/main/preprints/Uniform-exclusion-of-Landau-Siegel-zeros-October-1-2026/paper.pdf}.

\bibitem{sz-guo}
Victor Zhenyu Guo, \emph{An explicit elementary proof of no Landau--Siegel zeros},
arXiv:2610.30476v1, October 21, 2026, 17 pp.
\url{https://arxiv.org/html/2610.30476v1}.

\bibitem{sz-tao}
Terence Tao, \emph{A sketch of the proof of non-existence of Siegel zeroes},
What's New, October 19, 2026.
\url{https://terrytao.wordpress.com/2026/10/19/a-sketch-of-the-proof-of-non-existence-of-siegel-zeroes/}.

\bibitem{sz-pescador}
GPT-6 Pro and Claude Opus 5.5, posted and maintained by Jacob Pescador,
\emph{Landau--Siegel zeros and a Stepanov-type nonvanishing lemma},
version 3, October 22, 2026.
\url{https://jacobpescador.github.io/ai-math/papers/landau-siegel-stepanov/}.

\bibitem{sz-vandijk}
Roelof van Dijk,
\emph{A Frobenius determinant proof of the non-existence of Landau--Siegel zeros},
hexagon:2610.00409v2, 6 pp.
\url{https://hexagonmath.org/pdf/2610.00409v2}.

\bibitem{sz-kerr}
Bryce Kerr, \emph{A simpler proof that there are no Landau--Siegel zeros},
arXiv:2610.21099v1, October 15, 2026, 6 pp.
\url{https://arxiv.org/abs/2610.21099}.
\bibitem{ct-fw}
B. Ferrero and L. C. Washington,
The Iwasawa invariant $\mu_p$ vanishes for abelian number fields,
\emph{Annals of Mathematics} (2) \textbf{109} (1979), 377--395.
\url{https://doi.org/10.2307/1971116}.
Primary bibliographic record:
\url{https://annals.math.princeton.edu/1979/190-2/p05}.

\bibitem{ct-sharifi}
R. Sharifi, \emph{Iwasawa Theory}, author's lecture notes,
Chapter 3 (cyclotomic extensions and limits of class groups) and
Theorem 6.2.1 (Ferrero--Washington).
\url{https://www.math.ucla.edu/~sharifi/iwasawa.pdf}.
Author's reading edition, checked 21 November 2026:
\url{https://www.math.ucla.edu/~sharifi/notes/iwasawa-ch03.html},
\url{https://www.math.ucla.edu/~sharifi/notes/iwasawa-ch06.html}.

\bibitem{milne-ant}
J. S. Milne, \emph{Algebraic Number Theory}, version 3.08,
19 July 2020, Chapter 5 (the unit theorem).
\url{https://www.jmilne.org/math/CourseNotes/ANT.pdf}.

\bibitem{milne-cft}
J. S. Milne, \emph{Class Field Theory}, version 4.03,
6 August 2020, Chapter V, \S3, Theorems 3.5--3.6 and
Example 3.9 (Hilbert class field reciprocity).
\url{https://www.jmilne.org/math/CourseNotes/CFT.pdf}.

\bibitem{milne-adt}
J. S. Milne, \emph{Arithmetic Duality Theorems}, second edition,
BookSurge, 2006. In particular, Chapter II, Proposition 2.9
and Theorem 2.13 give the \'etale--Galois comparison and finite
coefficient cohomology used here.
\url{https://www.jmilne.org/math/Books/ADTnot.pdf}.

\bibitem{weibel-kbook}
C. A. Weibel, \emph{The $K$-book: An Introduction to Algebraic $K$-theory},
Graduate Studies in Mathematics \textbf{145}, American Mathematical
Society, 2013. Author's chapters III, IV and VI, especially III,
Example 1.3.6; IV,
Corollary 1.14.1, Theorems 1.17--1.18 and 6.9, and VI,
equation (8.1.1), Theorem 8.2 and Corollary 8.3.
\url{https://sites.math.rutgers.edu/~weibel/Kbook/Kbook.III.pdf};
\url{https://sites.math.rutgers.edu/~weibel/Kbook/Kbook.IV.pdf};
\url{https://sites.math.rutgers.edu/~weibel/Kbook/Kbook.VI.pdf}.

\bibitem{milnor-moore}
J. W. Milnor and J. C. Moore, On the structure of Hopf algebras,
\emph{Annals of Mathematics} (2) \textbf{81} (1965), 211--264,
especially Corollary 4.18, Theorem 5.18 and the appendix,
pp.\ 262--264.
\url{https://doi.org/10.2307/1970615}.
Author's article in the University of Rochester collection:
\url{https://www.sas.rochester.edu/mth/sites/doug-ravenel/otherpapers/Milnor-Moore.pdf}.

\bibitem{mathew-hurewicz}
A. Mathew, Torsion exponents in stable homotopy and the Hurewicz
homomorphism, \emph{Algebraic \& Geometric Topology}
\textbf{16} (2016), 1025--1041, especially Theorem 6.3.
\url{https://doi.org/10.2140/agt.2016.16.1025}.
Author's article:
\url{https://msp.org/agt/2016/16-2/agt-v16-n2-p12-s.pdf}.

\bibitem{alpo-uniform-source}
L. Alp\"oge, submitted source archive identified in the audit,
especially the directories
\path{Development/AutomorphicArithmetic/AutomorphicVarieties/LayeredIsomorphisms},
\path{Development/AutomorphicArithmetic/AutomorphicVarieties/AdditiveValuationNorms},
\path{Development/AutomorphicArithmetic/AutomorphicVarieties/QuadraticExtensions},
\path{Formalism/LocalToGlobal/Divisibility}, and the theta and
selected-character files listed in the accompanying source manifests.
The source traces specify the exact scope read; these locators do not
assert that every file in the directories was reviewed.

\bibitem{kirschmer-herm}
M. Kirschmer, \emph{Definite quadratic and hermitian forms with
small class number}, habilitation thesis, RWTH Aachen,
\S4.2, especially Theorem 4.2.3 and the definition of the local density,
pp.\ 50--51.
\url{https://www.math.uni-bielefeld.de/~mkirschm/papers/herm.pdf}.
The general mass formula there is attributed to
W. T. Gan and J.-K. Yu, \emph{Group schemes and local densities},
Duke Math.\ J.\ \textbf{105} (2000), 497--524, \S10.

\bibitem{terakado-yu-mass}
Y. Terakado and C.-F. Yu, \emph{Mass formulas and the basic locus
of unitary Shimura varieties}, arXiv:2210.04054v1 (2022),
\S2.3, especially equation (2.6).
\url{https://arxiv.org/pdf/2210.04054}.

\bibitem{harish-chandra-ds-ii}
Harish--Chandra, Discrete series for semisimple Lie groups. II:
Explicit determination of the characters,
\emph{Acta Mathematica} \textbf{116} (1966), 1--111,
especially Theorem 16 and \S41, pp.\ 96--100.
\url{https://doi.org/10.1007/BF02392813}.
Primary article in the Yau Mathematical Sciences Center collection:
\url{https://archive.ymsc.tsinghua.edu.cn/pacm_download/117/6013-11511_2006_Article_BF02392813.pdf}.

\bibitem{duflo-heckman-vergne}
M. Duflo, G. Heckman and M. Vergne,
Projection d'orbites, formule de Kirillov et formule de Blattner,
\emph{M\'emoires de la Soci\'et\'e Math\'ematique de France},
new series, no.\ 15 (1984), 65--128, especially III.1,
theorem (2), p.\ 111.
\url{https://www.numdam.org/article/MSMF_1984_2_15__65_0.pdf}.
This article gives another proof of the formula established by
H. Hecht and W. Schmid, A proof of Blattner's conjecture,
\emph{Inventiones Mathematicae} \textbf{31} (1975), 129--154,
\url{https://doi.org/10.1007/BF01404112}.

\bibitem{baldoni-vergne}
V. Baldoni and M. Vergne,
Discrete series representations and $K$ multiplicities for
$U(p,q)$. User's guide, arXiv:1008.5360v1 (2010),
especially \S2.2, pp.\ 12--14.
\url{https://webusers.imj-prg.fr/~michele.vergne/publications/discreteseries.pdf}.

\bibitem{parthasarathy-minimal}
R. Parthasarathy, Classification of discrete series by minimal $K$-type,
\emph{Representation Theory} \textbf{19} (2015), 167--185,
especially Theorem 1.1, pp.\ 167--168.
\url{https://doi.org/10.1090/ert/467}.
Author-posted full text, used to verify the theorem and its hypotheses:
\url{https://www.researchgate.net/publication/282775275_Classification_of_discrete_series_by_minimal_-type}.
\bibitem{bannai-kobayashi-2007}
K. Bannai and S. Kobayashi,
\emph{Algebraic theta functions and $p$-adic interpolation of
Eisenstein--Kronecker numbers},
arXiv:math/0610163v4, 11 December 2007;
published in Duke Mathematical Journal \textbf{153} (2010), 229--295.
The theorem numbers used here refer to the arXiv version:
Propositions 2.14--2.15, Corollary 2.16(i),
the reduced Mumford-translation formula in Lemma 2.7 and
Proposition 2.8, and \S3.1.
\url{https://arxiv.org/abs/math/0610163v4}.
\url{https://doi.org/10.1215/00127094-2010-024}.
\bibitem{stacks-07S7}
The Stacks Project Authors,
\emph{The Stacks Project}, Lemma 66.14.2, Tag 07S7,
free finite locally free quotients of projective schemes.
\url{https://stacks.math.columbia.edu/tag/07S7}.
\bibitem{sun-eisenstein}
Binyong Sun, \emph{Eisenstein integrals for theta stable parabolic
subalgebras}, arXiv:math/0403440v2, 28 November 2004.
Theorem 6.4 and \S6.4 give the compact-type integral;
\S\S1.1--1.4 fix the group, measure and conjugation conventions.
\url{https://arxiv.org/abs/math/0403440v2}.

\bibitem{vogan-zuckerman}
David A. Vogan, Jr. and Gregg J. Zuckerman,
\emph{Unitary representations with non-zero cohomology},
Compositio Mathematica \textbf{53} (1984), 51--90.
See (5.1), Theorem 5.3 and Proposition 6.1.
\url{https://www.numdam.org/item/CM_1984__53_1_51_0/}.
\bibitem{greither-popescu}
Cornelius Greither and Cristian D. Popescu,
\emph{An equivariant main conjecture in Iwasawa theory and applications},
arXiv:1103.3069v1, 16 March 2011.
Theorem 6.11 and Lemma 6.9, with the conventions in \S6.2.
\url{https://arxiv.org/abs/1103.3069v1}.

\bibitem{getz-hahn-2022}
J. R. Getz and H. Hahn,
\emph{An Introduction to Automorphic Representations},
author manuscript dated 22 April 2022, Theorems 6.2.3
and 6.3.2, with their fixed-level and compact-type conventions.
The theorem numbers here refer to this manuscript.
\url{https://sites.duke.edu/jgetz/files/2022/04/Graduate_Text.pdf}.
\bibitem{weil-siegel}
Andr\'e Weil,
\emph{La formule de Siegel et les groupes classiques},
Acta Mathematica \textbf{113} (1965), 1--87.
The unitary case is \S27, type $(I_2)$; see \S54 for
independence of the normalized quotient volume and
\S57 for the special linear Tamagawa number.
\url{https://archive.ymsc.tsinghua.edu.cn/pacm_download/117/5988-11511_2006_Article_BF02391774.pdf}.

\bibitem{louboutin-residue}
St\'ephane Louboutin,
\emph{Explicit upper bounds for residues of Dedekind zeta functions
and values of $L$-functions at $s=1$, and explicit lower bounds
for relative class numbers of CM-fields},
Canadian Journal of Mathematics \textbf{53} (2001), 1194--1222.
Theorem~1(1).
\url{https://doi.org/10.4153/CJM-2001-045-5}.
\bibitem{banaszczyk-transference}
W. Banaszczyk, New bounds in some transference theorems
in the geometry of numbers,
\emph{Mathematische Annalen} \textbf{296} (1993), 625--635,
especially Theorem~2.2.
\url{https://doi.org/10.1007/BF01445125}.

\bibitem{bost-kunnemann-extensions}
J.-B. Bost and K. K\"unnemann,
Hermitian vector bundles and extension groups on arithmetic
schemes. I. Geometry of numbers,
\emph{Advances in Mathematics} \textbf{223} (2010), 987--1106.
Author preprint, 11 January 2007, especially
Theorem~3.7.2, (3.54), printed p.~46.
\url{https://epub.uni-regensburg.de/13828/1/MP77.pdf}.

\bibitem{khare-larsen-savin}
C. Khare, M. Larsen and G. Savin,
Functoriality and the inverse Galois problem,
\emph{Compositio Mathematica} \textbf{144} (2008), 541--564.
Author version arXiv:math/0610860v3 (29 May 2015),
especially \S4.1, pp.~14--15.
\url{https://arxiv.org/abs/math/0610860}.

\bibitem{zaman-hecke-estimates}
A. Zaman, Explicit estimates for the zeros of Hecke
\(L\)-functions,
\emph{Journal of Number Theory} \textbf{162} (2016), 312--375.
Author version arXiv:1502.05679v4 (29 March 2016),
especially \S2.1 and Lemma~2.1.
\url{https://arxiv.org/abs/1502.05679}.
\bibitem{thorne-adjoint-2022}
J. A. Thorne, \emph{On the vanishing of adjoint Bloch--Kato
Selmer groups of irreducible automorphic Galois representations},
arXiv:2207.04925v1, 11 July 2022, Theorem~6.1.
\url{https://arxiv.org/pdf/2207.04925}.

\bibitem{newton-thorne-adjoint-2023}
J. Newton and J. A. Thorne, \emph{Adjoint Selmer groups of
automorphic Galois representations of unitary type},
J. Eur. Math. Soc. \textbf{25} (2023), 1919--1967.
doi:10.4171/JEMS/1228. Conventions on pp.~1920 and~1923.
\url{https://ems.press/content/serial-article-files/32858}.

\bibitem{alpo-central-comparison}
L. Alp\"oge, submitted archive, the integral family slope comparison and
the divisor inequality, especially
\path{Formalism/RankGapBounds/MeanBoundSecondForm.lean},
\path{Formalism/SizeWindowBrackets/TraceSquareAndDiagBounds.lean},
\path{Formalism/SizeWindowBrackets/ChainSlopeWindows/LevelOfCongruence.lean},
and
\path{Formalism/LocalToGlobal/Divisibility/ConstantGeTwentyEight.lean}.
The independently rewritten comparison in Section~\ref{cs-section}
does not assume the existence of the required arithmetic family.

\end{thebibliography}
\end{document}
